% Options for packages loaded elsewhere
% Options for packages loaded elsewhere
\PassOptionsToPackage{unicode}{hyperref}
\PassOptionsToPackage{hyphens}{url}
\PassOptionsToPackage{dvipsnames,svgnames,x11names}{xcolor}
%
\documentclass[
  letterpaper,
]{scrbook}
\usepackage{xcolor}
\usepackage[top=25mm,bottom=25mm,left=30mm,right=25mm]{geometry}
\usepackage{amsmath,amssymb}
\setcounter{secnumdepth}{5}
\usepackage{iftex}
\ifPDFTeX
  \usepackage[T1]{fontenc}
  \usepackage[utf8]{inputenc}
  \usepackage{textcomp} % provide euro and other symbols
\else % if luatex or xetex
  \usepackage{unicode-math} % this also loads fontspec
  \defaultfontfeatures{Scale=MatchLowercase}
  \defaultfontfeatures[\rmfamily]{Ligatures=TeX,Scale=1}
\fi
\usepackage[]{libertinus}
\ifPDFTeX\else
  % xetex/luatex font selection
\fi
% Use upquote if available, for straight quotes in verbatim environments
\IfFileExists{upquote.sty}{\usepackage{upquote}}{}
\IfFileExists{microtype.sty}{% use microtype if available
  \usepackage[]{microtype}
  \UseMicrotypeSet[protrusion]{basicmath} % disable protrusion for tt fonts
}{}
\makeatletter
\@ifundefined{KOMAClassName}{% if non-KOMA class
  \IfFileExists{parskip.sty}{%
    \usepackage{parskip}
  }{% else
    \setlength{\parindent}{0pt}
    \setlength{\parskip}{6pt plus 2pt minus 1pt}}
}{% if KOMA class
  \KOMAoptions{parskip=half}}
\makeatother
% Make \paragraph and \subparagraph free-standing
\makeatletter
\ifx\paragraph\undefined\else
  \let\oldparagraph\paragraph
  \renewcommand{\paragraph}{
    \@ifstar
      \xxxParagraphStar
      \xxxParagraphNoStar
  }
  \newcommand{\xxxParagraphStar}[1]{\oldparagraph*{#1}\mbox{}}
  \newcommand{\xxxParagraphNoStar}[1]{\oldparagraph{#1}\mbox{}}
\fi
\ifx\subparagraph\undefined\else
  \let\oldsubparagraph\subparagraph
  \renewcommand{\subparagraph}{
    \@ifstar
      \xxxSubParagraphStar
      \xxxSubParagraphNoStar
  }
  \newcommand{\xxxSubParagraphStar}[1]{\oldsubparagraph*{#1}\mbox{}}
  \newcommand{\xxxSubParagraphNoStar}[1]{\oldsubparagraph{#1}\mbox{}}
\fi
\makeatother

\usepackage{color}
\usepackage{fancyvrb}
\newcommand{\VerbBar}{|}
\newcommand{\VERB}{\Verb[commandchars=\\\{\}]}
\DefineVerbatimEnvironment{Highlighting}{Verbatim}{commandchars=\\\{\}}
% Add ',fontsize=\small' for more characters per line
\usepackage{framed}
\definecolor{shadecolor}{RGB}{241,243,245}
\newenvironment{Shaded}{\begin{snugshade}}{\end{snugshade}}
\newcommand{\AlertTok}[1]{\textcolor[rgb]{0.68,0.00,0.00}{#1}}
\newcommand{\AnnotationTok}[1]{\textcolor[rgb]{0.37,0.37,0.37}{#1}}
\newcommand{\AttributeTok}[1]{\textcolor[rgb]{0.40,0.45,0.13}{#1}}
\newcommand{\BaseNTok}[1]{\textcolor[rgb]{0.68,0.00,0.00}{#1}}
\newcommand{\BuiltInTok}[1]{\textcolor[rgb]{0.00,0.23,0.31}{#1}}
\newcommand{\CharTok}[1]{\textcolor[rgb]{0.13,0.47,0.30}{#1}}
\newcommand{\CommentTok}[1]{\textcolor[rgb]{0.37,0.37,0.37}{#1}}
\newcommand{\CommentVarTok}[1]{\textcolor[rgb]{0.37,0.37,0.37}{\textit{#1}}}
\newcommand{\ConstantTok}[1]{\textcolor[rgb]{0.56,0.35,0.01}{#1}}
\newcommand{\ControlFlowTok}[1]{\textcolor[rgb]{0.00,0.23,0.31}{\textbf{#1}}}
\newcommand{\DataTypeTok}[1]{\textcolor[rgb]{0.68,0.00,0.00}{#1}}
\newcommand{\DecValTok}[1]{\textcolor[rgb]{0.68,0.00,0.00}{#1}}
\newcommand{\DocumentationTok}[1]{\textcolor[rgb]{0.37,0.37,0.37}{\textit{#1}}}
\newcommand{\ErrorTok}[1]{\textcolor[rgb]{0.68,0.00,0.00}{#1}}
\newcommand{\ExtensionTok}[1]{\textcolor[rgb]{0.00,0.23,0.31}{#1}}
\newcommand{\FloatTok}[1]{\textcolor[rgb]{0.68,0.00,0.00}{#1}}
\newcommand{\FunctionTok}[1]{\textcolor[rgb]{0.28,0.35,0.67}{#1}}
\newcommand{\ImportTok}[1]{\textcolor[rgb]{0.00,0.46,0.62}{#1}}
\newcommand{\InformationTok}[1]{\textcolor[rgb]{0.37,0.37,0.37}{#1}}
\newcommand{\KeywordTok}[1]{\textcolor[rgb]{0.00,0.23,0.31}{\textbf{#1}}}
\newcommand{\NormalTok}[1]{\textcolor[rgb]{0.00,0.23,0.31}{#1}}
\newcommand{\OperatorTok}[1]{\textcolor[rgb]{0.37,0.37,0.37}{#1}}
\newcommand{\OtherTok}[1]{\textcolor[rgb]{0.00,0.23,0.31}{#1}}
\newcommand{\PreprocessorTok}[1]{\textcolor[rgb]{0.68,0.00,0.00}{#1}}
\newcommand{\RegionMarkerTok}[1]{\textcolor[rgb]{0.00,0.23,0.31}{#1}}
\newcommand{\SpecialCharTok}[1]{\textcolor[rgb]{0.37,0.37,0.37}{#1}}
\newcommand{\SpecialStringTok}[1]{\textcolor[rgb]{0.13,0.47,0.30}{#1}}
\newcommand{\StringTok}[1]{\textcolor[rgb]{0.13,0.47,0.30}{#1}}
\newcommand{\VariableTok}[1]{\textcolor[rgb]{0.07,0.07,0.07}{#1}}
\newcommand{\VerbatimStringTok}[1]{\textcolor[rgb]{0.13,0.47,0.30}{#1}}
\newcommand{\WarningTok}[1]{\textcolor[rgb]{0.37,0.37,0.37}{\textit{#1}}}

\usepackage{longtable,booktabs,array}
\usepackage{calc} % for calculating minipage widths
% Correct order of tables after \paragraph or \subparagraph
\usepackage{etoolbox}
\makeatletter
\patchcmd\longtable{\par}{\if@noskipsec\mbox{}\fi\par}{}{}
\makeatother
% Allow footnotes in longtable head/foot
\IfFileExists{footnotehyper.sty}{\usepackage{footnotehyper}}{\usepackage{footnote}}
\makesavenoteenv{longtable}
\usepackage{graphicx}
\makeatletter
\newsavebox\pandoc@box
\newcommand*\pandocbounded[1]{% scales image to fit in text height/width
  \sbox\pandoc@box{#1}%
  \Gscale@div\@tempa{\textheight}{\dimexpr\ht\pandoc@box+\dp\pandoc@box\relax}%
  \Gscale@div\@tempb{\linewidth}{\wd\pandoc@box}%
  \ifdim\@tempb\p@<\@tempa\p@\let\@tempa\@tempb\fi% select the smaller of both
  \ifdim\@tempa\p@<\p@\scalebox{\@tempa}{\usebox\pandoc@box}%
  \else\usebox{\pandoc@box}%
  \fi%
}
% Set default figure placement to htbp
\def\fps@figure{htbp}
\makeatother





\setlength{\emergencystretch}{3em} % prevent overfull lines

\providecommand{\tightlist}{%
  \setlength{\itemsep}{0pt}\setlength{\parskip}{0pt}}



 


\makeatletter
\@ifpackageloaded{tcolorbox}{}{\usepackage[skins,breakable]{tcolorbox}}
\@ifpackageloaded{fontawesome5}{}{\usepackage{fontawesome5}}
\definecolor{quarto-callout-color}{HTML}{909090}
\definecolor{quarto-callout-note-color}{HTML}{0758E5}
\definecolor{quarto-callout-important-color}{HTML}{CC1914}
\definecolor{quarto-callout-warning-color}{HTML}{EB9113}
\definecolor{quarto-callout-tip-color}{HTML}{00A047}
\definecolor{quarto-callout-caution-color}{HTML}{FC5300}
\definecolor{quarto-callout-color-frame}{HTML}{acacac}
\definecolor{quarto-callout-note-color-frame}{HTML}{4582ec}
\definecolor{quarto-callout-important-color-frame}{HTML}{d9534f}
\definecolor{quarto-callout-warning-color-frame}{HTML}{f0ad4e}
\definecolor{quarto-callout-tip-color-frame}{HTML}{02b875}
\definecolor{quarto-callout-caution-color-frame}{HTML}{fd7e14}
\makeatother
\makeatletter
\@ifpackageloaded{bookmark}{}{\usepackage{bookmark}}
\makeatother
\makeatletter
\@ifpackageloaded{caption}{}{\usepackage{caption}}
\AtBeginDocument{%
\ifdefined\contentsname
  \renewcommand*\contentsname{Table of contents}
\else
  \newcommand\contentsname{Table of contents}
\fi
\ifdefined\listfigurename
  \renewcommand*\listfigurename{List of Figures}
\else
  \newcommand\listfigurename{List of Figures}
\fi
\ifdefined\listtablename
  \renewcommand*\listtablename{List of Tables}
\else
  \newcommand\listtablename{List of Tables}
\fi
\ifdefined\figurename
  \renewcommand*\figurename{Figure}
\else
  \newcommand\figurename{Figure}
\fi
\ifdefined\tablename
  \renewcommand*\tablename{Table}
\else
  \newcommand\tablename{Table}
\fi
}
\@ifpackageloaded{float}{}{\usepackage{float}}
\floatstyle{ruled}
\@ifundefined{c@chapter}{\newfloat{codelisting}{h}{lop}}{\newfloat{codelisting}{h}{lop}[chapter]}
\floatname{codelisting}{Listing}
\newcommand*\listoflistings{\listof{codelisting}{List of Listings}}
\makeatother
\makeatletter
\makeatother
\makeatletter
\@ifpackageloaded{caption}{}{\usepackage{caption}}
\@ifpackageloaded{subcaption}{}{\usepackage{subcaption}}
\makeatother
\usepackage{bookmark}
\IfFileExists{xurl.sty}{\usepackage{xurl}}{} % add URL line breaks if available
\urlstyle{same}
\hypersetup{
  pdftitle={Wayward House Mathematics},
  pdfauthor={Paul Downey},
  colorlinks=true,
  linkcolor={Maroon},
  filecolor={Maroon},
  citecolor={Blue},
  urlcolor={Blue},
  pdfcreator={LaTeX via pandoc}}


\title{Wayward House Mathematics}
\usepackage{etoolbox}
\makeatletter
\providecommand{\subtitle}[1]{% add subtitle to \maketitle
  \apptocmd{\@title}{\par {\large #1 \par}}{}{}
}
\makeatother
\subtitle{A series for people who want to understand, not just pass}
\author{Paul Downey}
\date{2026-04-01}
\begin{document}
\frontmatter
\maketitle

\renewcommand*\contentsname{Table of contents}
{
\hypersetup{linkcolor=}
\setcounter{tocdepth}{2}
\tableofcontents
}

\mainmatter
\bookmarksetup{startatroot}

\chapter{Wayward House Mathematics}\label{wayward-house-mathematics}

A series for people who want to understand, not just pass

\hfill\break

Most mathematics education answers the question \emph{how}. This series
tries to answer \emph{what}, \emph{why}, and \emph{where this goes}.

It starts where the difficulty usually starts --- not at the beginning
of mathematics, but at the point where the notation begins to feel like
it belongs to someone else. Each volume takes a set of tools, shows what
they actually are, and points toward where they're used in the world.

The series runs from the foundations of arithmetic through to upper-year
engineering mathematics. You don't have to start at Volume 1. Start
where you lost the thread.

\section{The map}\label{the-map}

Before you begin, it's worth seeing the whole territory. The
\href{map/wayfinding.qmd}{wayfinding map} shows how the topics connect,
which fields they lead to, and how the mathematics opens up as you move
through the series.

\section{The volumes}\label{the-volumes}

\begin{longtable}[]{@{}
  >{\raggedright\arraybackslash}p{(\linewidth - 4\tabcolsep) * \real{0.2051}}
  >{\raggedright\arraybackslash}p{(\linewidth - 4\tabcolsep) * \real{0.2308}}
  >{\raggedright\arraybackslash}p{(\linewidth - 4\tabcolsep) * \real{0.5641}}@{}}
\toprule\noalign{}
\begin{minipage}[b]{\linewidth}\raggedright
Volume
\end{minipage} & \begin{minipage}[b]{\linewidth}\raggedright
Content
\end{minipage} & \begin{minipage}[b]{\linewidth}\raggedright
Roughly equivalent to
\end{minipage} \\
\midrule\noalign{}
\endhead
\bottomrule\noalign{}
\endlastfoot
1 & Numbers and operations & Alberta Gr 7--8 \\
2 & Patterns and algebra & Alberta Gr 8--9 \\
3 & Shape and measure & Alberta Gr 9--10 \\
4 & Functions and change & Alberta Gr 10--11 \\
5 & Analysis and proof & Alberta Gr 11--12 \\
6 & Probability and data & Alberta Gr 11--12 \\
7 & Engineering mathematics & U of A first and second year engineering
maths \\
8 & Upper-year engineering mathematics & U of A years 3 and 4
mathematical methods \\
\end{longtable}

The grade equivalents are approximate and are there for orientation
only. They are not labels.

\section{How to use this series}\label{how-to-use-this-series}

Each chapter opens with what you need coming in, and closes with what it
opens up. The exercises are puzzles, not drills --- they're designed to
be interesting rather than repetitive.

If a piece of notation stops you, look for the plain-language reading
immediately beside it. Notation is a compression tool, not a gatekeeping
device. This series treats it as such.

\part{The map}

\chapter{The Map}\label{the-map-2}

How the mathematics connects, and where it leads

\hfill\break

Every chapter in this series has a position in the map below. Nodes are
grouped by layer --- foundations at the centre, engineering mathematics
outside that, and discipline-facing upper-year methods on the outer
ring. Colours show which destination fields each topic serves.

Click any node to see what it is, where it leads, and which fields use
it.

\begin{Shaded}
\begin{Highlighting}[]
\NormalTok{//| echo: false}
\NormalTok{jsyaml = require("js{-}yaml@4");}
\NormalTok{d3 = require("d3@7");}
\end{Highlighting}
\end{Shaded}

\begin{Shaded}
\begin{Highlighting}[]
\NormalTok{//| echo: false}
\NormalTok{topo = fetch("/\_data/topology.yml")}
\NormalTok{  .then(r =\textgreater{} r.text())}
\NormalTok{  .then(text =\textgreater{} jsyaml.load(text))}
\NormalTok{  .catch(() =\textgreater{} null);}
\end{Highlighting}
\end{Shaded}

\begin{Shaded}
\begin{Highlighting}[]
\NormalTok{//| echo: false}
\NormalTok{fieldColors = (\{}
\NormalTok{  eng:  "\#185FA5",}
\NormalTok{  sci:  "\#0F6E56",}
\NormalTok{  biz:  "\#854F0B",}
\NormalTok{  comp: "\#534AB7",}
\NormalTok{  geo:  "\#993C1D"}
\NormalTok{\});}

\NormalTok{layerRadius   = (\{ 1: 120, 2: 220, 3: 310, 4: 410, 5: 510 \});}
\NormalTok{layerNodeSize = (\{ 1: 8,   2: 10,  3: 11,  4: 11,  5: 11  \});}
\end{Highlighting}
\end{Shaded}

\begin{Shaded}
\begin{Highlighting}[]
\NormalTok{//| echo: false}
\NormalTok{viewof selectedNode = \{}
\NormalTok{  const el = html\textasciigrave{}\textless{}div style="min{-}height:72px;padding:0.75rem 0 0.25rem;font{-}size:0.88rem;}
\NormalTok{    color:var({-}{-}bs{-}body{-}color,\#333);border{-}top:1px solid \#eee;margin{-}top:0.5rem"\textgreater{}}
\NormalTok{    \textless{}em style="color:\#999"\textgreater{}Click a node to see details.\textless{}/em\textgreater{}}
\NormalTok{  \textless{}/div\textgreater{}\textasciigrave{};}
\NormalTok{  el.value = null;}
\NormalTok{  return el;}
\NormalTok{\}}
\end{Highlighting}
\end{Shaded}

\begin{Shaded}
\begin{Highlighting}[]
\NormalTok{//| echo: false}
\NormalTok{bloom = \{}
\NormalTok{  if (!topo) return html\textasciigrave{}\textless{}p\textgreater{}\textless{}em\textgreater{}Loading…\textless{}/em\textgreater{}\textless{}/p\textgreater{}\textasciigrave{};}

\NormalTok{  const outerRadius = layerRadius[5] + 36;}
\NormalTok{  const W = Math.min(}
\NormalTok{    Math.max((window.innerWidth || 800) {-} 48, outerRadius * 2 + 40),}
\NormalTok{    1120}
\NormalTok{  );}
\NormalTok{  const H = W;}
\NormalTok{  const cx = W / 2, cy = H / 2;}

\NormalTok{  // Index nodes by layer for initial ring placement}
\NormalTok{  const byLayer = \{\};}
\NormalTok{  for (const n of topo.nodes) \{}
\NormalTok{    if (!byLayer[n.layer]) byLayer[n.layer] = [];}
\NormalTok{    byLayer[n.layer].push(n);}
\NormalTok{  \}}

\NormalTok{  const nodes = topo.nodes.map(n =\textgreater{} \{}
\NormalTok{    const peers = byLayer[n.layer];}
\NormalTok{    const i = peers.indexOf(n);}
\NormalTok{    const angle = (i / peers.length) * 2 * Math.PI {-} Math.PI / 2;}
\NormalTok{    const r = layerRadius[n.layer] || 200;}
\NormalTok{    return \{}
\NormalTok{      ...n,}
\NormalTok{      color: fieldColors[n.fields?.[0]] || "\#888",}
\NormalTok{      nodeR: layerNodeSize[n.layer] || 8,}
\NormalTok{      x: cx + r * Math.cos(angle),}
\NormalTok{      y: cy + r * Math.sin(angle)}
\NormalTok{    \};}
\NormalTok{  \});}

\NormalTok{  const nodeById = Object.fromEntries(nodes.map(n =\textgreater{} [n.id, n]));}

\NormalTok{  const links = (topo.edges || [])}
\NormalTok{    .map(([s, t, rel]) =\textgreater{} (\{ source: nodeById[s], target: nodeById[t], rel \}))}
\NormalTok{    .filter(l =\textgreater{} l.source \&\& l.target);}

\NormalTok{  const svg = d3.create("svg")}
\NormalTok{    .attr("viewBox", \textasciigrave{}0 0 $\{W\} $\{H\}\textasciigrave{})}
\NormalTok{    .attr("width", "100\%")}
\NormalTok{    .style("max{-}width", W + "px")}
\NormalTok{    .style("display", "block")}
\NormalTok{    .style("margin", "0 auto");}

\NormalTok{  // Guide rings}
\NormalTok{  svg.append("g").selectAll("circle")}
\NormalTok{    .data([1, 2, 3, 4, 5])}
\NormalTok{    .join("circle")}
\NormalTok{    .attr("cx", cx).attr("cy", cy)}
\NormalTok{    .attr("r", d =\textgreater{} layerRadius[d] + 18)}
\NormalTok{    .attr("fill", "none")}
\NormalTok{    .attr("stroke", "\#ddd")}
\NormalTok{    .attr("stroke{-}width", 0.75)}
\NormalTok{    .attr("stroke{-}dasharray", "4 4");}

\NormalTok{  // Ring labels}
\NormalTok{  svg.append("g").selectAll("text")}
\NormalTok{    .data([}
\NormalTok{      \{ layer: 1, label: "Pre{-}university" \},}
\NormalTok{      \{ layer: 2, label: "Bridge" \},}
\NormalTok{      \{ layer: 3, label: "University" \},}
\NormalTok{      \{ layer: 4, label: "Engineering" \},}
\NormalTok{      \{ layer: 5, label: "Discipline" \}}
\NormalTok{    ])}
\NormalTok{    .join("text")}
\NormalTok{    .attr("x", cx)}
\NormalTok{    .attr("y", d =\textgreater{} cy {-} layerRadius[d.layer] {-} 22)}
\NormalTok{    .attr("text{-}anchor", "middle")}
\NormalTok{    .attr("font{-}size", 10)}
\NormalTok{    .attr("fill", "\#bbb")}
\NormalTok{    .text(d =\textgreater{} d.label);}

\NormalTok{  // Links}
\NormalTok{  const linkSel = svg.append("g").selectAll("line")}
\NormalTok{    .data(links)}
\NormalTok{    .join("line")}
\NormalTok{    .attr("stroke", "\#ccc")}
\NormalTok{    .attr("stroke{-}width", d =\textgreater{} d.rel === "enables" ? 1.5 : 0.75)}
\NormalTok{    .attr("stroke{-}opacity", d =\textgreater{} d.rel === "enables" ? 0.7 : 0.4)}
\NormalTok{    .attr("stroke{-}dasharray", d =\textgreater{} d.rel === "motivates" ? "3 3" : null);}

\NormalTok{  // Node groups}
\NormalTok{  const nodeSel = svg.append("g").selectAll("g")}
\NormalTok{    .data(nodes)}
\NormalTok{    .join("g")}
\NormalTok{    .attr("cursor", "pointer")}
\NormalTok{    .on("click", (event, d) =\textgreater{} \{}
\NormalTok{      event.stopPropagation();}
\NormalTok{      nodeSel.select("circle")}
\NormalTok{        .attr("stroke", n =\textgreater{} n === d ? "\#222" : "\#fff")}
\NormalTok{        .attr("stroke{-}width", n =\textgreater{} n === d ? 2.5 : 1.5);}

\NormalTok{      const fields = (d.fields || []).map(fid =\textgreater{} \{}
\NormalTok{        const fd = (topo.fields || []).find(x =\textgreater{} x.id === fid);}
\NormalTok{        return fd}
\NormalTok{          ? \textasciigrave{}\textless{}span style="display:inline{-}block;margin:0 4px 4px 0;padding:2px 9px;}
\NormalTok{              border{-}radius:99px;font{-}size:0.75rem;background:$\{fieldColors[fid]\};color:\#fff"\textgreater{}$\{fd.label\}\textless{}/span\textgreater{}\textasciigrave{}}
\NormalTok{          : "";}
\NormalTok{      \}).join("");}

\NormalTok{      const link = d.chapter\_ref}
\NormalTok{        ? \textasciigrave{}\textless{}a href="/$\{d.chapter\_ref.replace(".qmd", ".html")\}" style="font{-}size:0.83rem"\textgreater{}}
\NormalTok{            → Go to chapter\textless{}/a\textgreater{}\textasciigrave{}}
\NormalTok{        : "";}

\NormalTok{      selectedNode.innerHTML = \textasciigrave{}}
\NormalTok{        \textless{}strong style="font{-}size:0.95rem"\textgreater{}$\{d.label\}\textless{}/strong\textgreater{}}
\NormalTok{        \textless{}div style="margin:4px 0 6px"\textgreater{}$\{fields\}\textless{}/div\textgreater{}}
\NormalTok{        \textless{}p style="margin:0 0 6px;font{-}size:0.83rem;line{-}height:1.55"\textgreater{}$\{(d.hook || "").trim()\}\textless{}/p\textgreater{}}
\NormalTok{        $\{link\}}
\NormalTok{      \textasciigrave{};}
\NormalTok{      selectedNode.value = d;}
\NormalTok{      selectedNode.dispatchEvent(new Event("input"));}
\NormalTok{    \});}

\NormalTok{  nodeSel.append("circle")}
\NormalTok{    .attr("r", d =\textgreater{} d.nodeR)}
\NormalTok{    .attr("fill", d =\textgreater{} d.color)}
\NormalTok{    .attr("fill{-}opacity", 0.88)}
\NormalTok{    .attr("stroke", "\#fff")}
\NormalTok{    .attr("stroke{-}width", 1.5);}

\NormalTok{  nodeSel.append("title").text(d =\textgreater{} d.label);}

\NormalTok{  // Labels for outer layers}
\NormalTok{  nodeSel.filter(d =\textgreater{} d.layer \textgreater{}= 3)}
\NormalTok{    .append("text")}
\NormalTok{    .attr("dy", d =\textgreater{} d.nodeR + 11)}
\NormalTok{    .attr("text{-}anchor", "middle")}
\NormalTok{    .attr("font{-}size", 8.5)}
\NormalTok{    .attr("fill", "\#666")}
\NormalTok{    .text(d =\textgreater{} d.label.length \textgreater{} 24 ? d.label.slice(0, 22) + "…" : d.label);}

\NormalTok{  // Force simulation}
\NormalTok{  d3.forceSimulation(nodes)}
\NormalTok{    .force("link", d3.forceLink(links)}
\NormalTok{      .distance(d =\textgreater{} Math.abs(d.target.layer {-} d.source.layer) === 0 ? 75 : 95)}
\NormalTok{      .strength(0.25))}
\NormalTok{    .force("charge", d3.forceManyBody().strength({-}100))}
\NormalTok{    .force("radial", d3.forceRadial(d =\textgreater{} layerRadius[d.layer], cx, cy).strength(0.85))}
\NormalTok{    .force("collision", d3.forceCollide(d =\textgreater{} d.nodeR + 9))}
\NormalTok{    .on("tick", () =\textgreater{} \{}
\NormalTok{      linkSel}
\NormalTok{        .attr("x1", d =\textgreater{} d.source.x).attr("y1", d =\textgreater{} d.source.y)}
\NormalTok{        .attr("x2", d =\textgreater{} d.target.x).attr("y2", d =\textgreater{} d.target.y);}
\NormalTok{      nodeSel.attr("transform", d =\textgreater{} \textasciigrave{}translate($\{d.x\},$\{d.y\})\textasciigrave{});}
\NormalTok{    \});}

\NormalTok{  return svg.node();}
\NormalTok{\}}
\end{Highlighting}
\end{Shaded}

\begin{center}\rule{0.5\linewidth}{0.5pt}\end{center}

\textbf{Field key:}

\begin{Shaded}
\begin{Highlighting}[]
\NormalTok{//| echo: false}
\NormalTok{\{}
\NormalTok{  if (!topo) return html\textasciigrave{}\textasciigrave{};}
\NormalTok{  return html\textasciigrave{}\textless{}div style="display:flex;flex{-}wrap:wrap;gap:12px 20px;margin:0.25rem 0 1rem"\textgreater{}}
\NormalTok{    $\{(topo.fields || []).map(f =\textgreater{} html\textasciigrave{}}
\NormalTok{      \textless{}span style="display:flex;align{-}items:center;gap:6px;font{-}size:0.85rem"\textgreater{}}
\NormalTok{        \textless{}span style="width:11px;height:11px;border{-}radius:50\%;}
\NormalTok{          background:$\{fieldColors[f.id]\};display:inline{-}block"\textgreater{}\textless{}/span\textgreater{}}
\NormalTok{        $\{f.label\}}
\NormalTok{      \textless{}/span\textgreater{}}
\NormalTok{    \textasciigrave{})\}}
\NormalTok{  \textless{}/div\textgreater{}\textasciigrave{};}
\NormalTok{\}}
\end{Highlighting}
\end{Shaded}

Solid edges --- direct prerequisite (\texttt{enables}). Lighter solid
--- deepens understanding (\texttt{enriches}). Dashed --- motivational
connection (\texttt{motivates}).

Layers run outward: pre-university → bridge → university → engineering.

\chapter{How to Use This Series}\label{how-to-use-this-series-1}

Where to start, how to navigate, what the notation means

\hfill\break

\section{Where to start}\label{where-to-start}

You do not have to start at Volume 1.

The series runs from arithmetic to engineering mathematics, but it is
not a ladder --- it is a map. Find the place where you lost the thread
and start there. The wayfinding map shows you the prerequisites for any
chapter so you can backfill exactly what you need.

If you are not sure where that is, a rough guide:

\begin{longtable}[]{@{}
  >{\raggedright\arraybackslash}p{(\linewidth - 2\tabcolsep) * \real{0.5000}}
  >{\raggedright\arraybackslash}p{(\linewidth - 2\tabcolsep) * \real{0.5000}}@{}}
\toprule\noalign{}
\begin{minipage}[b]{\linewidth}\raggedright
Where you are
\end{minipage} & \begin{minipage}[b]{\linewidth}\raggedright
Start here
\end{minipage} \\
\midrule\noalign{}
\endhead
\bottomrule\noalign{}
\endlastfoot
Struggled with high school maths & Volume 1 or 2 \\
Passed high school, never felt sure why & Volume 3 or 4 \\
About to start university maths & Volume 5 \\
In first- or second-year engineering, finding the maths hard & Volume
7 \\
In third- or fourth-year engineering, the maths has become embedded in
discipline courses & Volume 8 \\
\end{longtable}

\section{How chapters are structured}\label{how-chapters-are-structured}

Every chapter opens with two or three concrete situations that share the
same underlying mathematical structure. The mathematics comes from the
situation, not the other way round.

After the opening:

\begin{enumerate}
\def\labelenumi{\arabic{enumi}.}
\tightlist
\item
  The notation is introduced and read in plain English beside it
\item
  The method is shown, stepped through, with a verification
\item
  A brief explanation of why the method works
\item
  Two to four worked examples from different fields
\item
  Where the chapter leads --- the specific topics it opens up
\item
  Exercises
\end{enumerate}

\section{The exercises}\label{the-exercises}

Exercises come in four types. The type is stated at the start of each
exercise set.

\textbf{Puzzle} --- convergent, satisfying, one clean answer. The
interesting work is setting up the equation, not solving it.

\textbf{Drill} --- repetitive by design. Some things require fluency,
and fluency requires repetition.

\textbf{Investigation} --- open-ended, more than one valid answer. These
ask you to explore.

\textbf{Project} --- extended, usually involving a real dataset or
context.

\section{What the notation is doing}\label{what-the-notation-is-doing}

Notation is a compression tool. It lets you write in three symbols what
would otherwise take three sentences. This series introduces each symbol
when it is first needed, reads it in plain English immediately beside
it, and notes how heavy a lift it is --- some symbols need one sentence,
some need a whole section.

If a piece of notation stops you, look for the plain-language reading in
the same paragraph. It will be there.

\section{The destination fields}\label{the-destination-fields}

Every chapter is tagged with the fields it serves:

\begin{itemize}
\tightlist
\item
  \textbf{Engineering} --- structural, electrical, mechanical, chemical,
  software
\item
  \textbf{Hard sciences} --- physics, chemistry, biology, earth sciences
\item
  \textbf{Finance and quant} --- financial modelling, actuarial science,
  econometrics
\item
  \textbf{Computing and data} --- CS, machine learning, data science,
  cryptography
\item
  \textbf{Geography and environment} --- GIS, climate science,
  environmental modelling
\end{itemize}

The tags are there to answer the question ``why am I learning this''.
Each chapter's worked examples cover at least three of its tagged
fields.

\section{The prerequisites}\label{the-prerequisites}

Each chapter states what you need coming in. If a prerequisite chapter
is incomplete in your knowledge, the wayfinding map will tell you which
earlier chapter covers it.

The prerequisites are genuine --- not padding. If a chapter says it
requires limits, it means you will need limits to follow the argument.

\part{Volume 1 --- Numbers and operations}

\chapter{Volume 1 --- Numbers and
operations}\label{volume-1-numbers-and-operations-1}

The central move of this volume is placement.

Every number has a location: on a line, relative to zero, relative to
other numbers. This volume is about understanding that structure --- not
as abstract symbol manipulation, but as a description of real things.
Debts and credits. Divided portions. Rates of change. Samples from
populations.

You already know that numbers exist. What this volume builds is fluency
with how they relate --- how to scale them (ratio), express them as
parts of wholes (fractions and decimals), talk about change (percents
and rates), move in both directions from zero (integers), and reason
about collections of them (statistics and probability).

None of this requires anything beyond counting. But it requires counting
carefully, in the right direction, with the right units. That discipline
is what the rest of the series stands on.

\section{Chapters}\label{chapters}

\begin{longtable}[]{@{}
  >{\raggedright\arraybackslash}p{(\linewidth - 4\tabcolsep) * \real{0.3333}}
  >{\raggedright\arraybackslash}p{(\linewidth - 4\tabcolsep) * \real{0.2593}}
  >{\raggedright\arraybackslash}p{(\linewidth - 4\tabcolsep) * \real{0.4074}}@{}}
\toprule\noalign{}
\begin{minipage}[b]{\linewidth}\raggedright
Chapter
\end{minipage} & \begin{minipage}[b]{\linewidth}\raggedright
Title
\end{minipage} & \begin{minipage}[b]{\linewidth}\raggedright
Core move
\end{minipage} \\
\midrule\noalign{}
\endhead
\bottomrule\noalign{}
\endlastfoot
1 & Numbers and the number line & Placing any number exactly \\
2 & Ratio and proportion & Comparing two quantities of the same kind \\
3 & Percent and rates & Expressing change as a fraction of something \\
4 & Integers and signed arithmetic & Extending the number line in both
directions \\
5 & Fractions and decimals & Arithmetic with parts of wholes \\
6 & Statistics and probability & Reasoning honestly from limited data \\
\end{longtable}

\chapter{Numbers and the number line}\label{numbers-and-the-number-line}

Placing any number exactly

\hfill\break

A bank balance of −£40 is not the same as no money --- it is a real
position, 40 units to the left of zero. An altitude of −85 metres is the
Dead Sea shore. A temperature of −18°C is your freezer.

Numbers below zero are not a trick. They are the language of anything
that can be in deficit, in debt, below a reference point. The number
line is the map that shows where every number lives --- and why the
rules of arithmetic have the shape they do.

\section{What the notation is saying}\label{what-the-notation-is-saying}

Draw a horizontal line. Mark a point in the middle and call it
\textbf{0}. Every positive number lives to the right; every negative
number lives to the left. The further from zero, the larger the absolute
value.

\[\cdots \quad -4 \quad -3 \quad -2 \quad -1 \quad 0 \quad 1 \quad 2 \quad 3 \quad 4 \quad \cdots\]

\textbf{Absolute value} \(|x|\) strips the direction: \(|{-3}| = 3\) and
\(|3| = 3\). It answers ``how far from zero?'' --- not ``which side?''

The inequality symbols \(<\) and \(>\) describe position on the line:
\(-5 < -1\) because \(-5\) sits further left. This trips people up ---
\(-5\) feels ``bigger'' in magnitude, but it is smaller as a number.

Fractions and decimals also live on this line. \(\frac{1}{2}\) sits
halfway between 0 and 1. \(-0.7\) sits between \(-1\) and \(0\). The
line is continuous: there is a location for every number, however
precise.

\begin{center}\rule{0.5\linewidth}{0.5pt}\end{center}

\emph{Use the slider below to place any number on the line. The arrow
shows its distance from zero --- that distance is its absolute value.}

\begin{Shaded}
\begin{Highlighting}[]
\NormalTok{//| echo: false}
\NormalTok{nlExplorer = \{}
\NormalTok{  const d3 = await require("d3@7");}

\NormalTok{  const W = 560, H = 160;}
\NormalTok{  const ML = 40, MR = 40, MT = 50, MB = 40;}
\NormalTok{  const innerW = W {-} ML {-} MR;}

\NormalTok{  const xMin = {-}10, xMax = 10;}
\NormalTok{  const xScale = d3.scaleLinear([xMin, xMax], [0, innerW]);}

\NormalTok{  const container = d3.create("div")}
\NormalTok{    .style("font{-}family", "sans{-}serif")}
\NormalTok{    .style("max{-}width", W + "px");}

\NormalTok{  // {-}{-}{-} slider {-}{-}{-}}
\NormalTok{  const sliderLabel = container.append("label")}
\NormalTok{    .style("display", "block")}
\NormalTok{    .style("margin{-}bottom", "6px")}
\NormalTok{    .style("font{-}size", "0.9rem");}
\NormalTok{  sliderLabel.append("span").text("Position: ");}
\NormalTok{  const valueSpan = sliderLabel.append("span")}
\NormalTok{    .style("font{-}weight", "bold");}

\NormalTok{  const slider = container.append("input")}
\NormalTok{    .attr("type", "range")}
\NormalTok{    .attr("min", xMin)}
\NormalTok{    .attr("max", xMax)}
\NormalTok{    .attr("step", 1)}
\NormalTok{    .attr("value", {-}3)}
\NormalTok{    .style("width", "100\%")}
\NormalTok{    .style("margin{-}bottom", "6px");}

\NormalTok{  // {-}{-}{-} SVG {-}{-}{-}}
\NormalTok{  const svg = container.append("svg")}
\NormalTok{    .attr("width", W)}
\NormalTok{    .attr("height", H)}
\NormalTok{    .attr("viewBox", \textasciigrave{}0 0 $\{W\} $\{H\}\textasciigrave{})}
\NormalTok{    .attr("aria{-}label", "Interactive number line");}

\NormalTok{  const g = svg.append("g")}
\NormalTok{    .attr("transform", \textasciigrave{}translate($\{ML\},$\{MT\})\textasciigrave{});}

\NormalTok{  // axis line}
\NormalTok{  g.append("line")}
\NormalTok{    .attr("x1", 0).attr("y1", 0)}
\NormalTok{    .attr("x2", innerW).attr("y2", 0)}
\NormalTok{    .attr("stroke", "\#333").attr("stroke{-}width", 2);}

\NormalTok{  // arrowheads on axis}
\NormalTok{  g.append("polygon")}
\NormalTok{    .attr("points", \textasciigrave{}$\{innerW\},0 $\{innerW{-}8\},{-}4 $\{innerW{-}8\},4\textasciigrave{})}
\NormalTok{    .attr("fill", "\#333");}
\NormalTok{  g.append("polygon")}
\NormalTok{    .attr("points", \textasciigrave{}0,0 8,{-}4 8,4\textasciigrave{})}
\NormalTok{    .attr("fill", "\#333");}

\NormalTok{  // tick marks and labels}
\NormalTok{  for (let v = xMin; v \textless{}= xMax; v++) \{}
\NormalTok{    const x = xScale(v);}
\NormalTok{    g.append("line")}
\NormalTok{      .attr("x1", x).attr("y1", {-}6)}
\NormalTok{      .attr("x2", x).attr("y2", 6)}
\NormalTok{      .attr("stroke", v === 0 ? "\#222" : "\#777")}
\NormalTok{      .attr("stroke{-}width", v === 0 ? 2 : 1);}
\NormalTok{    g.append("text")}
\NormalTok{      .attr("x", x)}
\NormalTok{      .attr("y", 22)}
\NormalTok{      .attr("text{-}anchor", "middle")}
\NormalTok{      .attr("font{-}size", v === 0 ? "13px" : "11px")}
\NormalTok{      .attr("font{-}weight", v === 0 ? "bold" : "normal")}
\NormalTok{      .attr("fill", "\#333")}
\NormalTok{      .text(v);}
\NormalTok{  \}}

\NormalTok{  // zero label}
\NormalTok{  g.append("text")}
\NormalTok{    .attr("x", xScale(0))}
\NormalTok{    .attr("y", {-}14)}
\NormalTok{    .attr("text{-}anchor", "middle")}
\NormalTok{    .attr("font{-}size", "11px")}
\NormalTok{    .attr("fill", "\#555")}
\NormalTok{    .text("zero");}

\NormalTok{  // distance arrow (from zero to point) — drawn under the axis}
\NormalTok{  const arrowGroup = g.append("g").attr("transform", "translate(0,0)");}

\NormalTok{  const arrowLine = arrowGroup.append("line")}
\NormalTok{    .attr("y1", 36).attr("y2", 36)}
\NormalTok{    .attr("stroke", "\#d97706")}
\NormalTok{    .attr("stroke{-}width", 2.5)}
\NormalTok{    .attr("stroke{-}dasharray", "5,3");}

\NormalTok{  const arrowHead = arrowGroup.append("polygon")}
\NormalTok{    .attr("fill", "\#d97706");}

\NormalTok{  const distLabel = arrowGroup.append("text")}
\NormalTok{    .attr("y", 56)}
\NormalTok{    .attr("text{-}anchor", "middle")}
\NormalTok{    .attr("font{-}size", "12px")}
\NormalTok{    .attr("fill", "\#b45309")}
\NormalTok{    .attr("font{-}weight", "bold");}

\NormalTok{  // point marker}
\NormalTok{  const pointGroup = g.append("g");}
\NormalTok{  const pointCircle = pointGroup.append("circle")}
\NormalTok{    .attr("r", 9)}
\NormalTok{    .attr("fill", "\#2563eb")}
\NormalTok{    .attr("stroke", "\#1e3a8a")}
\NormalTok{    .attr("stroke{-}width", 2);}

\NormalTok{  const pointLabel = pointGroup.append("text")}
\NormalTok{    .attr("text{-}anchor", "middle")}
\NormalTok{    .attr("y", {-}16)}
\NormalTok{    .attr("font{-}size", "13px")}
\NormalTok{    .attr("font{-}weight", "bold")}
\NormalTok{    .attr("fill", "\#1e3a8a");}

\NormalTok{  const dirLabel = pointGroup.append("text")}
\NormalTok{    .attr("text{-}anchor", "middle")}
\NormalTok{    .attr("y", {-}30)}
\NormalTok{    .attr("font{-}size", "11px")}
\NormalTok{    .attr("fill", "\#475569");}

\NormalTok{  // absolute value text below SVG}
\NormalTok{  const absDiv = container.append("div")}
\NormalTok{    .style("font{-}size", "1rem")}
\NormalTok{    .style("text{-}align", "center")}
\NormalTok{    .style("margin{-}top", "4px")}
\NormalTok{    .style("color", "\#1e293b");}

\NormalTok{  function update(v) \{}
\NormalTok{    const x = xScale(v);}
\NormalTok{    const x0 = xScale(0);}
\NormalTok{    const absV = Math.abs(v);}

\NormalTok{    valueSpan.text(v);}

\NormalTok{    // point}
\NormalTok{    pointGroup.attr("transform", \textasciigrave{}translate($\{x\},0)\textasciigrave{});}
\NormalTok{    pointLabel.text(v);}
\NormalTok{    dirLabel.text(v \textgreater{} 0 ? "positive \textbackslash{}u2192" : v \textless{} 0 ? "\textbackslash{}u2190 negative" : "zero");}

\NormalTok{    // distance arrow}
\NormalTok{    if (v !== 0) \{}
\NormalTok{      const left = Math.min(x, x0);}
\NormalTok{      const right = Math.max(x, x0);}
\NormalTok{      arrowLine}
\NormalTok{        .attr("x1", left + 4)}
\NormalTok{        .attr("x2", right {-} 4)}
\NormalTok{        .attr("visibility", "visible");}
\NormalTok{      if (v \textgreater{} 0) \{}
\NormalTok{        arrowHead}
\NormalTok{          .attr("points", \textasciigrave{}$\{right\},36 $\{right{-}8\},32 $\{right{-}8\},40\textasciigrave{})}
\NormalTok{          .attr("visibility", "visible");}
\NormalTok{      \} else \{}
\NormalTok{        arrowHead}
\NormalTok{          .attr("points", \textasciigrave{}$\{left\},36 $\{left+8\},32 $\{left+8\},40\textasciigrave{})}
\NormalTok{          .attr("visibility", "visible");}
\NormalTok{      \}}
\NormalTok{      distLabel}
\NormalTok{        .attr("x", (left + right) / 2)}
\NormalTok{        .attr("visibility", "visible")}
\NormalTok{        .text(\textasciigrave{}distance from zero = $\{absV\}\textasciigrave{});}
\NormalTok{    \} else \{}
\NormalTok{      arrowLine.attr("visibility", "hidden");}
\NormalTok{      arrowHead.attr("visibility", "hidden");}
\NormalTok{      distLabel.attr("visibility", "hidden");}
\NormalTok{    \}}

\NormalTok{    absDiv.text(\textasciigrave{}|$\{v\}| = $\{absV\}\textasciigrave{});}
\NormalTok{  \}}

\NormalTok{  update(+slider.property("value"));}

\NormalTok{  slider.on("input", function() \{}
\NormalTok{    update(+this.value);}
\NormalTok{  \});}

\NormalTok{  return container.node();}
\NormalTok{\}}
\end{Highlighting}
\end{Shaded}

\begin{center}\rule{0.5\linewidth}{0.5pt}\end{center}

\section{The method}\label{the-method}

\textbf{Placing a number on the line}

\begin{enumerate}
\def\labelenumi{\arabic{enumi}.}
\tightlist
\item
  Determine the sign: positive (right of zero) or negative (left) ---
  this tells you which direction to go.
\item
  Determine the magnitude: how far from zero --- this tells you how far
  to go.
\item
  Locate it. Every number has exactly one position.
\end{enumerate}

\textbf{Ordering numbers}

To order a list of numbers, locate each one on the line. Left is
smaller, right is larger. For negative numbers:
\(-10 < -2 < 0 < 3 < 7\).

\textbf{Absolute value}

The formula below looks slightly odd. Here is what it is doing: if \(x\)
is already positive, its distance from zero is just \(x\). If \(x\) is
negative, its distance from zero is the positive version of it --- which
is \(-x\). Negating a negative gives a positive.

\(|x| = x\) if \(x \geq 0\); \(|x| = -x\) if \(x < 0\).

So \(|{-7}| = -(-7) = 7\). We will prove in Chapter 4 why \(-(-7) = 7\).

\begin{tcolorbox}[enhanced jigsaw, opacityback=0, colback=white, breakable, colframe=quarto-callout-note-color-frame, arc=.35mm, left=2mm, toprule=.15mm, leftrule=.75mm, bottomrule=.15mm, rightrule=.15mm]
\begin{minipage}[t]{5.5mm}
\textcolor{quarto-callout-note-color}{\faInfo}
\end{minipage}%
\begin{minipage}[t]{\textwidth - 5.5mm}

\textbf{Why this works}

``Negative'' means ``opposite direction from positive'' --- not ``less
real'' or ``doesn't exist.'' Absolute value removes direction because
sometimes you only care about size (distance, error) not which way. When
you multiply two negatives to get a positive, it is because reversing
direction twice returns you to where you started.

\end{minipage}%
\end{tcolorbox}

\section{Worked examples}\label{worked-examples}

\textbf{Example 1 --- Sport.} A football team's goal difference across
four matches: won by 3, lost by 2, won by 1, lost by 4. Express these as
signed numbers and order them from worst to best.

As signed numbers: +3, −2, +1, −4.

Place each on the number line and read left to right:

\[-4 < -2 < 1 < 3\]

Worst to best: lost by 4, lost by 2, won by 1, won by 3.

\begin{center}\rule{0.5\linewidth}{0.5pt}\end{center}

\textbf{Example 2 --- Geography.} A group records these elevations
relative to sea level: mountain peak +2,450 m; valley floor −60 m;
coastal town +8 m; lake bed −340 m. What is the vertical distance from
the lake bed to the peak?

Vertical distance = \(|2450 - (-340)|\).

That subtraction involves two negatives, so pause: subtracting a
negative is the same as adding.

\(|2450 - (-340)| = |2450 + 340| = 2790\) m.

The absolute value ensures we get a distance (positive), not a directed
displacement.

\begin{center}\rule{0.5\linewidth}{0.5pt}\end{center}

\textbf{Example 3 --- Weather.} Five cities report overnight
temperatures: −12°C, +3°C, −7°C, −19°C, +1°C. Which temperature is
furthest from zero? Which is closest?

Find the absolute value of each: 12, 3, 7, 19, 1.

Furthest from zero: −19°C (absolute value 19). Closest to zero: +1°C
(absolute value 1).

\begin{center}\rule{0.5\linewidth}{0.5pt}\end{center}

\textbf{Example 4 --- Money.} Your weekly pocket money is £15. You owe a
friend £8 from last week. You spend £4 on a snack. Express your running
total as signed numbers from the start of the week.

Start: 0. Receive £15: position +15. Repay debt −8: position +7. Spend
£4: position +3.

Each transaction is a movement along the number line. You end up at +3,
meaning £3 remaining.

\begin{center}\rule{0.5\linewidth}{0.5pt}\end{center}

\emph{The interactive below lets you explore addition and subtraction as
movement along the line. Set a starting value A, choose an operation,
and adjust B to see where you land.}

\begin{Shaded}
\begin{Highlighting}[]
\NormalTok{//| echo: false}
\NormalTok{nlArithmetic = \{}
\NormalTok{  const d3 = await require("d3@7");}

\NormalTok{  const W = 560, H = 170;}
\NormalTok{  const ML = 40, MR = 40, MT = 60, MB = 30;}
\NormalTok{  const innerW = W {-} ML {-} MR;}
\NormalTok{  const xMin = {-}10, xMax = 10;}
\NormalTok{  const xScale = d3.scaleLinear([xMin, xMax], [0, innerW]);}

\NormalTok{  const container = d3.create("div")}
\NormalTok{    .style("font{-}family", "sans{-}serif")}
\NormalTok{    .style("max{-}width", W + "px");}

\NormalTok{  // {-}{-}{-} controls row {-}{-}{-}}
\NormalTok{  const controls = container.append("div")}
\NormalTok{    .style("display", "flex")}
\NormalTok{    .style("gap", "16px")}
\NormalTok{    .style("align{-}items", "flex{-}end")}
\NormalTok{    .style("flex{-}wrap", "wrap")}
\NormalTok{    .style("margin{-}bottom", "8px");}

\NormalTok{  function makeNumberInput(parent, labelText, defaultVal) \{}
\NormalTok{    const wrap = parent.append("label")}
\NormalTok{      .style("display", "flex")}
\NormalTok{      .style("flex{-}direction", "column")}
\NormalTok{      .style("font{-}size", "0.85rem")}
\NormalTok{      .style("gap", "3px");}
\NormalTok{    wrap.append("span").text(labelText);}
\NormalTok{    const inp = wrap.append("input")}
\NormalTok{      .attr("type", "range")}
\NormalTok{      .attr("min", {-}10)}
\NormalTok{      .attr("max", 10)}
\NormalTok{      .attr("step", 1)}
\NormalTok{      .attr("value", defaultVal)}
\NormalTok{      .style("width", "120px");}
\NormalTok{    const val = wrap.append("span")}
\NormalTok{      .style("text{-}align", "center")}
\NormalTok{      .style("font{-}weight", "bold");}
\NormalTok{    return \{ inp, val \};}
\NormalTok{  \}}

\NormalTok{  const \{ inp: inpA, val: valA \} = makeNumberInput(controls, "A (start)", 3);}
\NormalTok{  const opWrap = controls.append("label")}
\NormalTok{    .style("display", "flex")}
\NormalTok{    .style("flex{-}direction", "column")}
\NormalTok{    .style("font{-}size", "0.85rem")}
\NormalTok{    .style("gap", "3px");}
\NormalTok{  opWrap.append("span").text("Operation");}
\NormalTok{  const opSel = opWrap.append("select")}
\NormalTok{    .style("font{-}size", "1rem")}
\NormalTok{    .style("padding", "2px 6px");}
\NormalTok{  opSel.append("option").attr("value", "add").text("+ (add)");}
\NormalTok{  opSel.append("option").attr("value", "sub").text("− (subtract)");}

\NormalTok{  const \{ inp: inpB, val: valB \} = makeNumberInput(controls, "B (amount)", {-}5);}

\NormalTok{  // {-}{-}{-} result line {-}{-}{-}}
\NormalTok{  const resultDiv = container.append("div")}
\NormalTok{    .style("font{-}size", "1rem")}
\NormalTok{    .style("text{-}align", "center")}
\NormalTok{    .style("margin{-}bottom", "6px")}
\NormalTok{    .style("font{-}weight", "bold")}
\NormalTok{    .style("color", "\#1e293b");}

\NormalTok{  // {-}{-}{-} SVG {-}{-}{-}}
\NormalTok{  const svg = container.append("svg")}
\NormalTok{    .attr("width", W)}
\NormalTok{    .attr("height", H)}
\NormalTok{    .attr("viewBox", \textasciigrave{}0 0 $\{W\} $\{H\}\textasciigrave{})}
\NormalTok{    .attr("aria{-}label", "Number line showing arithmetic as movement");}

\NormalTok{  const g = svg.append("g")}
\NormalTok{    .attr("transform", \textasciigrave{}translate($\{ML\},$\{MT\})\textasciigrave{});}

\NormalTok{  // axis}
\NormalTok{  g.append("line")}
\NormalTok{    .attr("x1", 0).attr("y1", 0)}
\NormalTok{    .attr("x2", innerW).attr("y2", 0)}
\NormalTok{    .attr("stroke", "\#333").attr("stroke{-}width", 2);}
\NormalTok{  g.append("polygon")}
\NormalTok{    .attr("points", \textasciigrave{}$\{innerW\},0 $\{innerW{-}8\},{-}4 $\{innerW{-}8\},4\textasciigrave{})}
\NormalTok{    .attr("fill", "\#333");}
\NormalTok{  g.append("polygon")}
\NormalTok{    .attr("points", \textasciigrave{}0,0 8,{-}4 8,4\textasciigrave{})}
\NormalTok{    .attr("fill", "\#333");}

\NormalTok{  // ticks}
\NormalTok{  for (let v = xMin; v \textless{}= xMax; v++) \{}
\NormalTok{    const x = xScale(v);}
\NormalTok{    g.append("line")}
\NormalTok{      .attr("x1", x).attr("y1", {-}5).attr("x2", x).attr("y2", 5)}
\NormalTok{      .attr("stroke", v === 0 ? "\#222" : "\#888")}
\NormalTok{      .attr("stroke{-}width", v === 0 ? 2 : 1);}
\NormalTok{    g.append("text")}
\NormalTok{      .attr("x", x).attr("y", 20)}
\NormalTok{      .attr("text{-}anchor", "middle")}
\NormalTok{      .attr("font{-}size", "11px")}
\NormalTok{      .attr("fill", "\#444")}
\NormalTok{      .text(v);}
\NormalTok{  \}}

\NormalTok{  // movement arrow (drawn below axis)}
\NormalTok{  const moveArrow = g.append("g");}
\NormalTok{  const moveLine = moveArrow.append("line")}
\NormalTok{    .attr("y1", 34).attr("y2", 34)}
\NormalTok{    .attr("stroke", "\#7c3aed").attr("stroke{-}width", 3);}
\NormalTok{  const moveHead = moveArrow.append("polygon")}
\NormalTok{    .attr("fill", "\#7c3aed");}
\NormalTok{  const moveLabel = moveArrow.append("text")}
\NormalTok{    .attr("y", 55)}
\NormalTok{    .attr("text{-}anchor", "middle")}
\NormalTok{    .attr("font{-}size", "11px")}
\NormalTok{    .attr("fill", "\#5b21b6");}

\NormalTok{  // point A}
\NormalTok{  const ptA = g.append("g");}
\NormalTok{  ptA.append("circle").attr("r", 8)}
\NormalTok{    .attr("fill", "\#0369a1").attr("stroke", "\#075985").attr("stroke{-}width", 2);}
\NormalTok{  ptA.append("text")}
\NormalTok{    .attr("y", {-}14).attr("text{-}anchor", "middle")}
\NormalTok{    .attr("font{-}size", "12px").attr("fill", "\#075985").attr("font{-}weight", "bold")}
\NormalTok{    .attr("class", "ptA{-}label");}
\NormalTok{  ptA.append("text")}
\NormalTok{    .attr("y", {-}26).attr("text{-}anchor", "middle")}
\NormalTok{    .attr("font{-}size", "10px").attr("fill", "\#0369a1")}
\NormalTok{    .text("A");}

\NormalTok{  // point result}
\NormalTok{  const ptR = g.append("g");}
\NormalTok{  ptR.append("circle").attr("r", 8)}
\NormalTok{    .attr("fill", "\#16a34a").attr("stroke", "\#15803d").attr("stroke{-}width", 2);}
\NormalTok{  ptR.append("text")}
\NormalTok{    .attr("y", {-}14).attr("text{-}anchor", "middle")}
\NormalTok{    .attr("font{-}size", "12px").attr("fill", "\#15803d").attr("font{-}weight", "bold")}
\NormalTok{    .attr("class", "ptR{-}label");}
\NormalTok{  ptR.append("text")}
\NormalTok{    .attr("y", {-}26).attr("text{-}anchor", "middle")}
\NormalTok{    .attr("font{-}size", "10px").attr("fill", "\#16a34a")}
\NormalTok{    .text("Result");}

\NormalTok{  function clamp(v) \{}
\NormalTok{    return Math.max(xMin, Math.min(xMax, v));}
\NormalTok{  \}}

\NormalTok{  function update() \{}
\NormalTok{    const a = +inpA.property("value");}
\NormalTok{    const b = +inpB.property("value");}
\NormalTok{    const op = opSel.property("value");}
\NormalTok{    const result = clamp(op === "add" ? a + b : a {-} b);}
\NormalTok{    const opChar = op === "add" ? "+" : "−";}
\NormalTok{    const bDisplay = op === "sub" ? b : b;}

\NormalTok{    valA.text(a);}
\NormalTok{    valB.text(b);}

\NormalTok{    const xA = xScale(a);}
\NormalTok{    const xR = xScale(result);}

\NormalTok{    ptA.attr("transform", \textasciigrave{}translate($\{xA\},0)\textasciigrave{});}
\NormalTok{    ptA.select(".ptA{-}label").text(a);}

\NormalTok{    ptR.attr("transform", \textasciigrave{}translate($\{xR\},0)\textasciigrave{});}
\NormalTok{    ptR.select(".ptR{-}label").text(result);}

\NormalTok{    // movement arrow}
\NormalTok{    if (a !== result) \{}
\NormalTok{      const leftX = Math.min(xA, xR);}
\NormalTok{      const rightX = Math.max(xA, xR);}
\NormalTok{      moveLine.attr("x1", leftX + 4).attr("x2", rightX {-} 4).attr("visibility", "visible");}
\NormalTok{      if (result \textgreater{} a) \{}
\NormalTok{        moveHead.attr("points", \textasciigrave{}$\{rightX\},34 $\{rightX{-}8\},30 $\{rightX{-}8\},38\textasciigrave{}).attr("visibility", "visible");}
\NormalTok{      \} else \{}
\NormalTok{        moveHead.attr("points", \textasciigrave{}$\{leftX\},34 $\{leftX+8\},30 $\{leftX+8\},38\textasciigrave{}).attr("visibility", "visible");}
\NormalTok{      \}}
\NormalTok{      moveLabel}
\NormalTok{        .attr("x", (leftX + rightX) / 2)}
\NormalTok{        .attr("visibility", "visible")}
\NormalTok{        .text(\textasciigrave{}move $\{Math.abs(result {-} a)\} $\{result \textgreater{} a ? "right" : "left"\}\textasciigrave{});}
\NormalTok{    \} else \{}
\NormalTok{      moveLine.attr("visibility", "hidden");}
\NormalTok{      moveHead.attr("visibility", "hidden");}
\NormalTok{      moveLabel.attr("visibility", "hidden");}
\NormalTok{    \}}

\NormalTok{    resultDiv.text(\textasciigrave{}$\{a\} $\{opChar\} $\{b\} = $\{result\}$\{Math.abs(a + (op === "add" ? b : {-}b)) \textgreater{} 10 ? " (clamped to line range)" : ""\}\textasciigrave{});}
\NormalTok{  \}}

\NormalTok{  update();}
\NormalTok{  inpA.on("input", update);}
\NormalTok{  inpB.on("input", update);}
\NormalTok{  opSel.on("change", update);}

\NormalTok{  // legend}
\NormalTok{  const legend = container.append("div")}
\NormalTok{    .style("display", "flex").style("gap", "16px").style("font{-}size", "0.8rem")}
\NormalTok{    .style("margin{-}top", "4px").style("flex{-}wrap", "wrap");}
\NormalTok{  const items = [}
\NormalTok{    \{ color: "\#0369a1", label: "A — starting value" \},}
\NormalTok{    \{ color: "\#16a34a", label: "Result" \},}
\NormalTok{    \{ color: "\#7c3aed", label: "Movement (the operation)" \}}
\NormalTok{  ];}
\NormalTok{  items.forEach((\{ color, label \}) =\textgreater{} \{}
\NormalTok{    const item = legend.append("div").style("display", "flex").style("align{-}items", "center").style("gap", "5px");}
\NormalTok{    item.append("div")}
\NormalTok{      .style("width", "14px").style("height", "14px")}
\NormalTok{      .style("border{-}radius", "50\%").style("background", color).style("flex{-}shrink", "0");}
\NormalTok{    item.append("span").text(label);}
\NormalTok{  \});}

\NormalTok{  return container.node();}
\NormalTok{\}}
\end{Highlighting}
\end{Shaded}

\begin{center}\rule{0.5\linewidth}{0.5pt}\end{center}

\section{Where this goes}\label{where-this-goes}

The number line is the foundation every later chapter assumes.
\textbf{Integers and signed arithmetic} (Chapter 4) takes addition,
subtraction, multiplication, and division into negative territory ---
all of which are just movements along this line. \textbf{Fractions and
decimals} (Chapter 5) fill in the gaps between the integers.

Further on, the line becomes an axis. In Volume 2, the horizontal axis
of a graph is the number line, and a linear equation traces a path along
it. In Volume 5, limits are defined precisely as behaviour near a point
on the line. The model you are building here does not change --- it
accumulates meaning.

\begin{tcolorbox}[enhanced jigsaw, opacityback=0, colback=white, breakable, colframe=quarto-callout-tip-color-frame, arc=.35mm, left=2mm, toprule=.15mm, leftrule=.75mm, bottomrule=.15mm, rightrule=.15mm]
\begin{minipage}[t]{5.5mm}
\textcolor{quarto-callout-tip-color}{\faLightbulb}
\end{minipage}%
\begin{minipage}[t]{\textwidth - 5.5mm}

\textbf{Where this shows up}

\begin{itemize}
\tightlist
\item
  A weather forecaster reading −18°C is placing a temperature on the
  number line with zero at freezing.
\item
  A diver tracking depth uses a signed number with zero at the surface.
\item
  A games scoreboard showing −5 (penalty points) is using the same idea.
\item
  An app showing your bank balance in red is showing you a negative
  number.
\end{itemize}

Same structure. Different labels.

\end{minipage}%
\end{tcolorbox}

\section{Exercises}\label{exercises}

Each problem has a clean answer. The interesting part is setting up
which numbers go where before you compute.

\begin{enumerate}
\def\labelenumi{\arabic{enumi}.}
\item
  A diver descends from the surface to −28 m, then swims to a coral
  shelf at −15 m, then ascends to −5 m before surfacing. List these
  depths in order from deepest to shallowest.
\item
  Four cities report overnight temperatures: Edmonton −23°C, Vancouver
  +4°C, Winnipeg −31°C, Toronto −8°C. Order from coldest to warmest.
  What is the difference in temperature between the warmest and coldest
  cities?
\item
  A football team's weekly goal differences are: Mon −1, Tue +2, Wed −3,
  Thu +4, Fri −1. On which day was the result furthest from zero?
  \emph{(Use absolute values to compare.)}
\item
  Five friends pool money for a trip. Their balances in the group fund
  are: +£12, −£5, +£3, −£18, +£8. Order from least to most. Who owes the
  most?
\item
  Place these numbers on a number line sketch and order them from least
  to greatest: \(-1.5\), \(2\), \(0\), \(-3\), \(0.7\), \(-0.5\).
\item
  A submarine is at −180 m. A fish swims at −45 m. A buoy floats at +2
  m. What is the distance between the submarine and the buoy? Between
  the fish and the submarine?
\end{enumerate}

\chapter{Ratio and proportion}\label{ratio-and-proportion}

Comparing two quantities of the same kind

\hfill\break

A recipe for 4 people needs 300 g of flour. You're cooking for 10. How
much flour do you need?

A map is drawn at 1:25,000. A road measures 8 cm on the map. How long is
it on the ground?

You and two friends split a streaming subscription. The monthly cost is
£14.99 for three people. One friend joins. What does each person pay
now?

These three questions share the same calculation underneath.

\section{What the notation is
saying}\label{what-the-notation-is-saying-1}

The ratio \(a:b\) means ``for every \(a\) of the first thing, there are
\(b\) of the second.'' It is a comparison by division. Writing \(a:b\)
is the same as writing the fraction \(\frac{a}{b}\).

Two quantities \(y\) and \(x\) are \textbf{proportional} when:

\[y = kx\]

where \(k\) is a constant. Double \(x\), you double \(y\). Halve \(x\),
you halve \(y\). The ratio \(\frac{y}{x} = k\) stays fixed.

Write this as \(y \propto x\) (read: ``\(y\) is proportional to
\(x\)''). The constant \(k\) is called the \textbf{constant of
proportionality} --- it is the fixed multiplier that links the two
quantities.

A proportion equation sets two ratios equal:

\[\frac{a}{b} = \frac{c}{d}\]

Read: ``\(a\) is to \(b\) as \(c\) is to \(d\).'' If you know three of
the four values, you can find the fourth. The next section shows how.

\begin{center}\rule{0.5\linewidth}{0.5pt}\end{center}

\emph{Adjust the two parts of a ratio below. The bars show how A and B
relate to each other and to the total.}

\begin{Shaded}
\begin{Highlighting}[]
\NormalTok{//| echo: false}
\NormalTok{ratioVis = \{}
\NormalTok{  const d3 = await require("d3@7");}

\NormalTok{  const W = 560, H = 200;}
\NormalTok{  const ML = 20, MR = 20, MT = 20, MB = 20;}
\NormalTok{  const innerW = W {-} ML {-} MR;}

\NormalTok{  const container = d3.create("div")}
\NormalTok{    .style("font{-}family", "sans{-}serif")}
\NormalTok{    .style("max{-}width", W + "px");}

\NormalTok{  // controls}
\NormalTok{  const controls = container.append("div")}
\NormalTok{    .style("display", "flex")}
\NormalTok{    .style("gap", "24px")}
\NormalTok{    .style("align{-}items", "flex{-}end")}
\NormalTok{    .style("flex{-}wrap", "wrap")}
\NormalTok{    .style("margin{-}bottom", "10px");}

\NormalTok{  function makeSlider(parent, labelText, defaultVal) \{}
\NormalTok{    const wrap = parent.append("label")}
\NormalTok{      .style("display", "flex")}
\NormalTok{      .style("flex{-}direction", "column")}
\NormalTok{      .style("font{-}size", "0.85rem")}
\NormalTok{      .style("gap", "3px");}
\NormalTok{    const header = wrap.append("div")}
\NormalTok{      .style("display", "flex").style("justify{-}content", "space{-}between");}
\NormalTok{    header.append("span").text(labelText);}
\NormalTok{    const valSpan = header.append("span").style("font{-}weight", "bold");}
\NormalTok{    const inp = wrap.append("input")}
\NormalTok{      .attr("type", "range")}
\NormalTok{      .attr("min", 1).attr("max", 20).attr("step", 1)}
\NormalTok{      .attr("value", defaultVal)}
\NormalTok{      .style("width", "160px");}
\NormalTok{    return \{ inp, valSpan \};}
\NormalTok{  \}}

\NormalTok{  const \{ inp: inpA, valSpan: vsA \} = makeSlider(controls, "Part A", 3);}
\NormalTok{  const \{ inp: inpB, valSpan: vsB \} = makeSlider(controls, "Part B", 5);}

\NormalTok{  // info line}
\NormalTok{  const infoDiv = container.append("div")}
\NormalTok{    .style("font{-}size", "0.95rem")}
\NormalTok{    .style("text{-}align", "center")}
\NormalTok{    .style("margin{-}bottom", "10px")}
\NormalTok{    .style("color", "\#1e293b");}

\NormalTok{  // SVG}
\NormalTok{  const svg = container.append("svg")}
\NormalTok{    .attr("width", W)}
\NormalTok{    .attr("height", H)}
\NormalTok{    .attr("viewBox", \textasciigrave{}0 0 $\{W\} $\{H\}\textasciigrave{})}
\NormalTok{    .attr("aria{-}label", "Ratio bar visualisation");}

\NormalTok{  const g = svg.append("g").attr("transform", \textasciigrave{}translate($\{ML\},$\{MT\})\textasciigrave{});}

\NormalTok{  const BAR\_H = 36;}
\NormalTok{  const GAP = 16;}
\NormalTok{  const rows = [}
\NormalTok{    \{ label: "A", y: 0, fill: "\#2563eb", stroke: "\#1e3a8a" \},}
\NormalTok{    \{ label: "B", y: BAR\_H + GAP, fill: "\#d97706", stroke: "\#92400e" \},}
\NormalTok{    \{ label: "Total (A+B)", y: (BAR\_H + GAP) * 2, fill: "\#64748b", stroke: "\#334155" \}}
\NormalTok{  ];}

\NormalTok{  // background tracks}
\NormalTok{  rows.forEach(row =\textgreater{} \{}
\NormalTok{    g.append("rect")}
\NormalTok{      .attr("x", 0).attr("y", row.y)}
\NormalTok{      .attr("width", innerW).attr("height", BAR\_H)}
\NormalTok{      .attr("fill", "\#f1f5f9").attr("rx", 4);}
\NormalTok{  \});}

\NormalTok{  // bar rects}
\NormalTok{  const barRects = rows.map(row =\textgreater{} \{}
\NormalTok{    return g.append("rect")}
\NormalTok{      .attr("x", 0).attr("y", row.y)}
\NormalTok{      .attr("height", BAR\_H)}
\NormalTok{      .attr("fill", row.fill).attr("rx", 4);}
\NormalTok{  \});}

\NormalTok{  // row labels (left side, outside)}
\NormalTok{  rows.forEach(row =\textgreater{} \{}
\NormalTok{    g.append("text")}
\NormalTok{      .attr("x", {-}8)}
\NormalTok{      .attr("y", row.y + BAR\_H / 2 + 4)}
\NormalTok{      .attr("text{-}anchor", "end")}
\NormalTok{      .attr("font{-}size", "12px")}
\NormalTok{      .attr("fill", "\#334155")}
\NormalTok{      .attr("font{-}weight", "bold")}
\NormalTok{      .text(row.label);}
\NormalTok{  \});}

\NormalTok{  // value labels inside bars}
\NormalTok{  const barLabels = rows.map(row =\textgreater{} \{}
\NormalTok{    return g.append("text")}
\NormalTok{      .attr("y", row.y + BAR\_H / 2 + 5)}
\NormalTok{      .attr("font{-}size", "13px")}
\NormalTok{      .attr("font{-}weight", "bold")}
\NormalTok{      .attr("fill", "\#fff");}
\NormalTok{  \});}

\NormalTok{  function update() \{}
\NormalTok{    const a = +inpA.property("value");}
\NormalTok{    const b = +inpB.property("value");}
\NormalTok{    const total = a + b;}
\NormalTok{    const maxVal = 20; // slider max for A and B; total can reach 40}

\NormalTok{    vsA.text(a);}
\NormalTok{    vsB.text(b);}

\NormalTok{    const values = [a, b, total];}
\NormalTok{    const maxDisplay = 40; // total bar always represents full width when a=b=20}

\NormalTok{    values.forEach((v, i) =\textgreater{} \{}
\NormalTok{      const w = (v / maxDisplay) * innerW;}
\NormalTok{      barRects[i].attr("width", Math.max(w, 4));}
\NormalTok{      barLabels[i]}
\NormalTok{        .attr("x", Math.max(w / 2, 20))}
\NormalTok{        .attr("text{-}anchor", "middle")}
\NormalTok{        .text(v);}
\NormalTok{    \});}

\NormalTok{    const decimal = (a / b).toFixed(3);}
\NormalTok{    infoDiv.html(}
\NormalTok{      \textasciigrave{}A : B = \textless{}strong\textgreater{}$\{a\} : $\{b\}\textless{}/strong\textgreater{} \&nbsp;|\&nbsp; \textasciigrave{} +}
\NormalTok{      \textasciigrave{}A / B = \textless{}strong\textgreater{}$\{decimal\}\textless{}/strong\textgreater{} \&nbsp;|\&nbsp; \textasciigrave{} +}
\NormalTok{      \textasciigrave{}Total parts = \textless{}strong\textgreater{}$\{total\}\textless{}/strong\textgreater{}\textasciigrave{}}
\NormalTok{    );}
\NormalTok{  \}}

\NormalTok{  update();}
\NormalTok{  inpA.on("input", update);}
\NormalTok{  inpB.on("input", update);}

\NormalTok{  // legend}
\NormalTok{  const legend = container.append("div")}
\NormalTok{    .style("display", "flex").style("gap", "16px").style("font{-}size", "0.8rem")}
\NormalTok{    .style("margin{-}top", "6px").style("flex{-}wrap", "wrap");}
\NormalTok{  [}
\NormalTok{    \{ color: "\#2563eb", label: "Part A" \},}
\NormalTok{    \{ color: "\#d97706", label: "Part B" \},}
\NormalTok{    \{ color: "\#64748b", label: "Total (A + B)" \}}
\NormalTok{  ].forEach((\{ color, label \}) =\textgreater{} \{}
\NormalTok{    const item = legend.append("div").style("display", "flex").style("align{-}items", "center").style("gap", "5px");}
\NormalTok{    item.append("div")}
\NormalTok{      .style("width", "14px").style("height", "14px").style("border{-}radius", "3px")}
\NormalTok{      .style("background", color).style("flex{-}shrink", "0");}
\NormalTok{    item.append("span").text(label);}
\NormalTok{  \});}

\NormalTok{  return container.node();}
\NormalTok{\}}
\end{Highlighting}
\end{Shaded}

\begin{center}\rule{0.5\linewidth}{0.5pt}\end{center}

\section{The method}\label{the-method-1}

\textbf{Finding a missing value in a proportion}

Given \(\frac{a}{b} = \frac{c}{d}\), find \(d\).

Multiply both sides by \(bd\) --- this clears both denominators at once,
leaving \(ad = bc\). Rearranging for \(d\):

\[d = \frac{bc}{a}\]

\textbf{Dividing a quantity in a given ratio}

To divide total \(T\) in the ratio \(a:b\):

\begin{enumerate}
\def\labelenumi{\arabic{enumi}.}
\tightlist
\item
  Count the total number of parts: \(a + b\) --- this tells you how many
  equal slices the whole is cut into.
\item
  Value of one part: \(T \div (a + b)\) --- the size of each slice.
\item
  First share: \(a \times\) (one part); second share: \(b \times\) (one
  part) --- multiply each person's slices by the slice size.
\end{enumerate}

\textbf{Cascaded ratios}

Two ratios applied in series multiply. A 3:1 ratio followed by a 4:1
ratio gives a combined ratio of \(3 \times 4 = 12:1\).

\textbf{Gear ratio from tooth counts}

For meshing gears, the ratio equals the number of teeth on the driven
gear divided by the teeth on the driving gear.

\textbf{Finding a constant of proportionality}

Given one known pair \((x_0, y_0)\) where \(y \propto x\):

\[k = \frac{y_0}{x_0}\]

Then for any other \(x\): \(y = kx\).

\begin{tcolorbox}[enhanced jigsaw, opacityback=0, colback=white, breakable, colframe=quarto-callout-note-color-frame, arc=.35mm, left=2mm, toprule=.15mm, leftrule=.75mm, bottomrule=.15mm, rightrule=.15mm]
\begin{minipage}[t]{5.5mm}
\textcolor{quarto-callout-note-color}{\faInfo}
\end{minipage}%
\begin{minipage}[t]{\textwidth - 5.5mm}

\textbf{Why this works}

A ratio is a relative relationship --- it does not depend on the size of
either quantity, only on their comparison. When two quantities are
proportional, their ratio \(k\) is constant by definition.
Cross-multiplication works because multiplying both sides of
\(\frac{a}{b} = \frac{c}{d}\) by \(bd\) gives \(ad = bc\) --- a balanced
equation.

\end{minipage}%
\end{tcolorbox}

\begin{center}\rule{0.5\linewidth}{0.5pt}\end{center}

\emph{Enter three known values in a proportion. The fourth is calculated
for you. Use this to check your worked-example answers.}

\begin{Shaded}
\begin{Highlighting}[]
\NormalTok{//| echo: false}
\NormalTok{proportionSolver = \{}
\NormalTok{  const d3 = await require("d3@7");}

\NormalTok{  const container = d3.create("div")}
\NormalTok{    .style("font{-}family", "sans{-}serif")}
\NormalTok{    .style("max{-}width", "560px");}

\NormalTok{  container.append("p")}
\NormalTok{    .style("font{-}size", "0.85rem")}
\NormalTok{    .style("margin{-}bottom", "10px")}
\NormalTok{    .style("color", "\#475569")}
\NormalTok{    .text("Given a/b = c/d, enter a, b, and c — d is calculated. Enter any three values to find the fourth.");}

\NormalTok{  // controls grid}
\NormalTok{  const grid = container.append("div")}
\NormalTok{    .style("display", "grid")}
\NormalTok{    .style("grid{-}template{-}columns", "repeat(4, 1fr)")}
\NormalTok{    .style("gap", "12px")}
\NormalTok{    .style("align{-}items", "end")}
\NormalTok{    .style("margin{-}bottom", "14px");}

\NormalTok{  function makeField(parent, labelText, defaultVal, isResult) \{}
\NormalTok{    const wrap = parent.append("label")}
\NormalTok{      .style("display", "flex")}
\NormalTok{      .style("flex{-}direction", "column")}
\NormalTok{      .style("font{-}size", "0.85rem")}
\NormalTok{      .style("gap", "4px");}
\NormalTok{    wrap.append("span")}
\NormalTok{      .style("font{-}weight", "bold")}
\NormalTok{      .style("color", isResult ? "\#15803d" : "\#1e293b")}
\NormalTok{      .text(labelText);}
\NormalTok{    const inp = wrap.append("input")}
\NormalTok{      .attr("type", "number")}
\NormalTok{      .attr("step", "any")}
\NormalTok{      .attr("value", defaultVal)}
\NormalTok{      .style("width", "100\%")}
\NormalTok{      .style("font{-}size", "1rem")}
\NormalTok{      .style("padding", "4px 6px")}
\NormalTok{      .style("border", isResult ? "2px solid \#16a34a" : "1px solid \#cbd5e1")}
\NormalTok{      .style("border{-}radius", "4px")}
\NormalTok{      .style("background", isResult ? "\#f0fdf4" : "\#fff")}
\NormalTok{      .property("readOnly", isResult);}
\NormalTok{    return inp;}
\NormalTok{  \}}

\NormalTok{  const inpA2 = makeField(grid, "a", 3, false);}
\NormalTok{  const inpB2 = makeField(grid, "b", 4, false);}
\NormalTok{  const inpC2 = makeField(grid, "c", 6, false);}
\NormalTok{  const inpD2 = makeField(grid, "d = b·c ÷ a", "?", true);}

\NormalTok{  // equation display}
\NormalTok{  const eqDiv = container.append("div")}
\NormalTok{    .style("font{-}size", "1.1rem")}
\NormalTok{    .style("text{-}align", "center")}
\NormalTok{    .style("margin{-}bottom", "10px")}
\NormalTok{    .style("padding", "10px")}
\NormalTok{    .style("background", "\#f8fafc")}
\NormalTok{    .style("border{-}radius", "6px")}
\NormalTok{    .style("border", "1px solid \#e2e8f0");}

\NormalTok{  // explanation line}
\NormalTok{  const explDiv = container.append("div")}
\NormalTok{    .style("font{-}size", "0.85rem")}
\NormalTok{    .style("color", "\#475569")}
\NormalTok{    .style("text{-}align", "center");}

\NormalTok{  function update() \{}
\NormalTok{    const a = parseFloat(inpA2.property("value"));}
\NormalTok{    const b = parseFloat(inpB2.property("value"));}
\NormalTok{    const c = parseFloat(inpC2.property("value"));}

\NormalTok{    if (isNaN(a) || isNaN(b) || isNaN(c) || a === 0) \{}
\NormalTok{      inpD2.property("value", "—");}
\NormalTok{      eqDiv.text("Enter values for a, b, and c above (a must not be zero).");}
\NormalTok{      explDiv.text("");}
\NormalTok{      return;}
\NormalTok{    \}}

\NormalTok{    const d = (b * c) / a;}
\NormalTok{    const dDisplay = Number.isInteger(d) ? d : d.toFixed(4);}
\NormalTok{    inpD2.property("value", dDisplay);}

\NormalTok{    eqDiv.html(}
\NormalTok{      \textasciigrave{}$\{a\} / $\{b\} = $\{c\} / \textless{}strong style="color:\#15803d"\textgreater{}$\{dDisplay\}\textless{}/strong\textgreater{}\textasciigrave{}}
\NormalTok{    );}
\NormalTok{    explDiv.text(\textasciigrave{}d = b × c ÷ a = $\{b\} × $\{c\} ÷ $\{a\} = $\{dDisplay\}\textasciigrave{});}
\NormalTok{  \}}

\NormalTok{  update();}
\NormalTok{  inpA2.on("input", update);}
\NormalTok{  inpB2.on("input", update);}
\NormalTok{  inpC2.on("input", update);}

\NormalTok{  return container.node();}
\NormalTok{\}}
\end{Highlighting}
\end{Shaded}

\begin{center}\rule{0.5\linewidth}{0.5pt}\end{center}

\section{Worked examples}\label{worked-examples-1}

\textbf{Example 1 --- Recipe scaling.} A recipe for 4 people needs 300 g
flour, 150 g butter, and 80 g sugar. How much of each do you need for 10
people?

The quantities scale in proportion with the number of people. Find the
ratio of new to old: \(\frac{10}{4} = 2.5\).

Multiply each ingredient by 2.5 --- this keeps everything in the same
ratio:

\[\text{flour: } 300 \times 2.5 = 750 \text{ g}\]
\[\text{butter: } 150 \times 2.5 = 375 \text{ g}\]
\[\text{sugar: } 80 \times 2.5 = 200 \text{ g}\]

\begin{center}\rule{0.5\linewidth}{0.5pt}\end{center}

\textbf{Example 2 --- Geography: map scale.} A road appears as 7.4 cm on
a 1:50,000 map. What is the actual length of the road?

Map scale \(1:50,000\) means 1 unit on the map = 50,000 units on the
ground. Multiply to convert:

\[\text{actual length} = 7.4 \text{ cm} \times 50,000 = 370,000 \text{ cm} = 3.7 \text{ km}\]

\begin{center}\rule{0.5\linewidth}{0.5pt}\end{center}

\textbf{Example 3 --- Splitting costs.} Three friends split a monthly
phone plan costing £36 in the ratio 2:3:1 (based on how much data each
uses). How much does each pay?

Total parts: \(2 + 3 + 1 = 6\). One part = \(£36 \div 6 = £6\).

Friend A (2 parts): \(2 \times £6 = £12\) Friend B (3 parts):
\(3 \times £6 = £18\) Friend C (1 part): \(1 \times £6 = £6\)

Check: \(£12 + £18 + £6 = £36\). Yes.

\begin{center}\rule{0.5\linewidth}{0.5pt}\end{center}

\textbf{Example 4 --- Finance: dividing profit.} Two partners invest in
a business in the ratio 3:5. The year's profit is £24,000. How much does
each receive?

Total parts: \(3 + 5 = 8\). One part = \(£24,000 \div 8 = £3,000\).

Partner A (3 parts): \(3 \times £3,000 = £9,000\) Partner B (5 parts):
\(5 \times £3,000 = £15,000\)

Check: \(£9,000 + £15,000 = £24,000\). Yes.

\section{Where this goes}\label{where-this-goes-1}

Proportion is the first multiplicative relationship you encounter.
\textbf{Percent and rates} (Chapter 3) is proportion restricted to a
denominator of 100 --- a special case of this chapter.

Later, proportion reappears inside trigonometry (the sine ratio is a
ratio of lengths) and in logarithms (where proportional growth becomes
additive). In engineering mathematics, dimensional analysis --- checking
that equations are consistent in their units --- is a systematic
application of ratio.

The notation \(y \propto x\) is also the simplest case of a functional
relationship. Volume 3's treatment of functions generalises this to any
rule that maps inputs to outputs.

\begin{tcolorbox}[enhanced jigsaw, opacityback=0, colback=white, breakable, colframe=quarto-callout-tip-color-frame, arc=.35mm, left=2mm, toprule=.15mm, leftrule=.75mm, bottomrule=.15mm, rightrule=.15mm]
\begin{minipage}[t]{5.5mm}
\textcolor{quarto-callout-tip-color}{\faLightbulb}
\end{minipage}%
\begin{minipage}[t]{\textwidth - 5.5mm}

\textbf{Where this shows up}

\begin{itemize}
\tightlist
\item
  Scaling a recipe up or down is proportion --- every ingredient changes
  by the same factor.
\item
  Reading a map involves ratio: every measurement on paper is a
  proportional reduction of the real distance.
\item
  Splitting a bill by how much each person ordered is dividing a
  quantity in a ratio.
\item
  A gear ratio on a bike tells you how many times the back wheel turns
  for each turn of the pedals.
\end{itemize}

The arithmetic is identical in every case.

\end{minipage}%
\end{tcolorbox}

\section{Exercises}\label{exercises-1}

\begin{enumerate}
\def\labelenumi{\arabic{enumi}.}
\item
  A recipe for 4 people needs 300 g flour, 150 g butter, and 80 g sugar.
  Scale the recipe to serve 7 people. \emph{(Exact grams to one decimal
  place.)}
\item
  A map has scale 1:25,000. Two towns are 14.6 cm apart on the map. What
  is the actual distance in kilometres?
\item
  A car engine produces 250 Nm of torque. It drives through a gearbox
  with a 3.8:1 ratio, then a differential with a 3.2:1 ratio. What is
  the torque at the rear wheels?
\item
  Silver and copper are mixed in the ratio 7:3 by mass to make an alloy.
  You need 2.5 kg of alloy. How many grams of each metal do you need?
\item
  A 500 mL bottle of concentrated cleaning fluid requires diluting 1:15
  before use. How many litres of diluted product can you make from one
  bottle?
\item
  Three investors share profits in the ratio 2:3:7. Total profit this
  year is £48,000. How much does each investor receive?
\item
  A gear system has a driver gear with 24 teeth and a driven gear with
  60 teeth. What is the gear ratio? If the driver rotates at 900 rpm,
  what is the speed of the driven gear?
\end{enumerate}

\chapter{Percent and rates}\label{percent-and-rates}

Expressing change as a fraction of something

\hfill\break

Your phone battery drains 3\% per hour on standby. A savings account
pays 4.5\% annual interest. A streaming service puts its prices up by
8\%. Your favourite team's win rate this season is 60\%.

All four are the same operation: a fraction expressed out of a hundred,
applied to a quantity. The difference is what the quantity is, and
whether the percent is applied once or repeatedly.

\section{What the notation is
saying}\label{what-the-notation-is-saying-2}

\textbf{Percent} means per hundred. \(p\%\) is the same as
\(\frac{p}{100}\).

To find \(p\%\) of a quantity \(Q\): multiply \(Q\) by
\(\frac{p}{100}\).

\[\text{amount} = Q \times \frac{p}{100}\]

\textbf{Percent change} compares a change to the original:

\[\text{percent change} = \frac{\text{new} - \text{original}}{\text{original}} \times 100\]

\textbf{A rate} is a ratio with different units: speed is km/h, data use
is MB/day, interest is £/year. The structure is the same as percent ---
one quantity compared to another.

\textbf{Compound growth} applies a rate repeatedly. After \(n\) periods
at rate \(r\) (expressed as a decimal), a quantity \(Q_0\) becomes:

\[Q_n = Q_0 \times (1 + r)^n\]

\textbf{Compound decay} uses \(1 - r\) instead:

\[Q_n = Q_0 \times (1 - r)^n\]

The exponent \(n\) is why compound growth and decay are not linear ---
they accelerate in a way that simple percent change does not.

\section{The method}\label{the-method-2}

\textbf{Simple percent of a quantity}

\(p\%\) of \(Q\) = \(Q \times \frac{p}{100}\)

Convert percent to decimal first: \(7\% = 0.07\); \(120\% = 1.20\).

\textbf{Percent increase / decrease}

To increase \(Q\) by \(p\%\): multiply by \((1 + \frac{p}{100})\) ---
one step does the same job as finding the percent and then adding it.

To decrease \(Q\) by \(p\%\): multiply by \((1 - \frac{p}{100})\).

\textbf{Finding the original after a percent change}

If after a \(p\%\) increase the value is \(V\), work backwards by
dividing by the multiplier rather than multiplying:

\[Q_0 = \frac{V}{1 + \frac{p}{100}}\]

\textbf{Compound growth}

\(Q_n = Q_0 \times (1 + r)^n\) where \(r = \frac{p}{100}\).

\begin{tcolorbox}[enhanced jigsaw, opacityback=0, colback=white, breakable, colframe=quarto-callout-note-color-frame, arc=.35mm, left=2mm, toprule=.15mm, leftrule=.75mm, bottomrule=.15mm, rightrule=.15mm]
\begin{minipage}[t]{5.5mm}
\textcolor{quarto-callout-note-color}{\faInfo}
\end{minipage}%
\begin{minipage}[t]{\textwidth - 5.5mm}

\textbf{Why this works}

Each compounding period, the new starting point is the previous result.
Multiplying by \((1+r)\) once applies the rate for one period. Applying
it again gives \((1+r)^2\) --- the rate is now applied to the
already-grown quantity. The exponent is just repeated multiplication.
This is why compound growth eventually overtakes any flat trend: the
base grows, so the next period's growth is larger than the last.

\end{minipage}%
\end{tcolorbox}

\begin{center}\rule{0.5\linewidth}{0.5pt}\end{center}

\textbf{Interactive: Percent change visualiser.} Adjust the original
value and the percent change. The bar shows the result and the
single-step multiplier used.

\begin{Shaded}
\begin{Highlighting}[]
\NormalTok{//| echo: false}
\NormalTok{d3 = require("d3@7")}
\end{Highlighting}
\end{Shaded}

\begin{Shaded}
\begin{Highlighting}[]
\NormalTok{//| echo: false}
\NormalTok{viewof origVal = Inputs.number(}
\NormalTok{  \{ label: "Original value", min: 1, max: 10000, step: 1, value: 100 \}}
\NormalTok{)}
\end{Highlighting}
\end{Shaded}

\begin{Shaded}
\begin{Highlighting}[]
\NormalTok{//| echo: false}
\NormalTok{viewof pctChange = Inputs.range(}
\NormalTok{  [{-}50, 100], \{ label: "Percent change (\%)", step: 1, value: 20 \}}
\NormalTok{)}
\end{Highlighting}
\end{Shaded}

\begin{Shaded}
\begin{Highlighting}[]
\NormalTok{//| echo: false}
\NormalTok{\{}
\NormalTok{  const W = 560;}
\NormalTok{  const H = 160;}
\NormalTok{  const BAR\_Y = 60;}
\NormalTok{  const BAR\_H = 36;}
\NormalTok{  const LABEL\_Y = BAR\_Y + BAR\_H + 20;}
\NormalTok{  const MULT\_Y = LABEL\_Y + 22;}
\NormalTok{  const PAD\_L = 12;}
\NormalTok{  const PAD\_R = 12;}

\NormalTok{  // Derived values — guard against NaN/zero}
\NormalTok{  const orig = Number.isFinite(origVal) \&\& origVal \textgreater{} 0 ? origVal : 100;}
\NormalTok{  const pct  = Number.isFinite(pctChange) ? pctChange : 0;}
\NormalTok{  const multiplier = 1 + pct / 100;}
\NormalTok{  const newVal = orig * multiplier;}

\NormalTok{  // Scale: always show at least max(orig, newVal); min width 1}
\NormalTok{  const maxVal = Math.max(orig, newVal, 1);}
\NormalTok{  const usableW = W {-} PAD\_L {-} PAD\_R;}
\NormalTok{  const scale = usableW / maxVal;}

\NormalTok{  const origW  = Math.max(orig  * scale, 1);}
\NormalTok{  const newW   = Math.max(newVal * scale, 1);}

\NormalTok{  // Colours: increase = teal, decrease = amber, neutral = slate}
\NormalTok{  const isIncrease = pct \textgreater{} 0;}
\NormalTok{  const isDecrease = pct \textless{} 0;}
\NormalTok{  const origColour = "\#4b5563";  // slate{-}600}
\NormalTok{  const newColour  = isIncrease ? "\#0d9488"  // teal{-}600}
\NormalTok{                   : isDecrease ? "\#d97706"  // amber{-}600}
\NormalTok{                   : "\#4b5563";}

\NormalTok{  const sign = pct \textgreater{}= 0 ? "+" : "";}
\NormalTok{  const multSign = multiplier \textgreater{}= 1 ? "×" : "×";}

\NormalTok{  const multStr = multiplier.toFixed(2);}
\NormalTok{  const newValStr = Number.isInteger(newVal)}
\NormalTok{    ? newVal.toString()}
\NormalTok{    : newVal.toFixed(2);}

\NormalTok{  const svg = d3.create("svg")}
\NormalTok{    .attr("viewBox", \textasciigrave{}0 0 $\{W\} $\{H\}\textasciigrave{})}
\NormalTok{    .attr("width", "100\%")}
\NormalTok{    .attr("style", "max{-}width:" + W + "px; font{-}family: inherit;");}

\NormalTok{  // Original bar}
\NormalTok{  svg.append("rect")}
\NormalTok{    .attr("x", PAD\_L)}
\NormalTok{    .attr("y", BAR\_Y)}
\NormalTok{    .attr("width", origW)}
\NormalTok{    .attr("height", BAR\_H)}
\NormalTok{    .attr("fill", origColour)}
\NormalTok{    .attr("rx", 3);}

\NormalTok{  // Original label (above bar, left{-}aligned)}
\NormalTok{  svg.append("text")}
\NormalTok{    .attr("x", PAD\_L)}
\NormalTok{    .attr("y", BAR\_Y {-} 6)}
\NormalTok{    .attr("fill", origColour)}
\NormalTok{    .attr("font{-}size", "12px")}
\NormalTok{    .attr("font{-}weight", "600")}
\NormalTok{    .text(\textasciigrave{}Original: $\{orig\}\textasciigrave{});}

\NormalTok{  // New value bar (drawn below original as a second row)}
\NormalTok{  const BAR2\_Y = BAR\_Y + BAR\_H + 10;}

\NormalTok{  svg.append("rect")}
\NormalTok{    .attr("x", PAD\_L)}
\NormalTok{    .attr("y", BAR2\_Y)}
\NormalTok{    .attr("width", newW)}
\NormalTok{    .attr("height", BAR\_H)}
\NormalTok{    .attr("fill", newColour)}
\NormalTok{    .attr("rx", 3);}

\NormalTok{  // New value label (above second bar, right of bar if space allows)}
\NormalTok{  const newLabelX = PAD\_L + newW + 6;}
\NormalTok{  const newLabelAnchor = newLabelX + 120 \textless{} W ? "start" : "end";}
\NormalTok{  const newLabelXAdj   = newLabelAnchor === "end" ? PAD\_L + newW {-} 4 : newLabelX;}

\NormalTok{  svg.append("text")}
\NormalTok{    .attr("x", newLabelXAdj)}
\NormalTok{    .attr("y", BAR2\_Y + BAR\_H / 2 + 5)}
\NormalTok{    .attr("fill", newColour)}
\NormalTok{    .attr("font{-}size", "13px")}
\NormalTok{    .attr("font{-}weight", "700")}
\NormalTok{    .attr("text{-}anchor", newLabelAnchor)}
\NormalTok{    .text(\textasciigrave{}New value: $\{newValStr\}\textasciigrave{});}

\NormalTok{  // Multiplier label at bottom}
\NormalTok{  const changeLabel = isIncrease}
\NormalTok{    ? \textasciigrave{}$\{sign\}$\{pct\}\% means × $\{multStr\}  (multiply by $\{multStr\} in one step)\textasciigrave{}}
\NormalTok{    : isDecrease}
\NormalTok{    ? \textasciigrave{}$\{sign\}$\{pct\}\% means × $\{multStr\}  (multiply by $\{multStr\} in one step)\textasciigrave{}}
\NormalTok{    : \textasciigrave{}0\% means × 1.00  (no change)\textasciigrave{};}

\NormalTok{  svg.append("text")}
\NormalTok{    .attr("x", PAD\_L)}
\NormalTok{    .attr("y", BAR2\_Y + BAR\_H + 20)}
\NormalTok{    .attr("fill", "\#374151")}
\NormalTok{    .attr("font{-}size", "12px")}
\NormalTok{    .text(changeLabel);}

\NormalTok{  return svg.node();}
\NormalTok{\}}
\end{Highlighting}
\end{Shaded}

\begin{center}\rule{0.5\linewidth}{0.5pt}\end{center}

\section{Worked examples}\label{worked-examples-2}

\textbf{Example 1 --- Savings: compound interest.} You put £2,000 into a
savings account at 3.5\% annual interest, compounded every year. What is
the balance after 5 years?

Set up the compound formula with \(Q_0 = 2000\), \(r = 0.035\),
\(n = 5\):

\[Q_5 = 2000 \times (1.035)^5\]

Work out \((1.035)^5\) --- that is 1.035 multiplied by itself five
times. Using a calculator: \((1.035)^5 = 1.1877\) (to 4 d.p.).

\[Q_5 = 2000 \times 1.1877 = £2,375.40\]

Compare to simple interest: \(5 \times 3.5\% \times 2000 = £350\),
giving £2,350. The compounded version is £25.40 more --- the gap grows
larger the longer you leave it.

\begin{center}\rule{0.5\linewidth}{0.5pt}\end{center}

\textbf{Example 2 --- Phone plan.} A mobile plan costs £25 per month.
The provider raises prices by 12\%. What is the new monthly cost? How
much extra do you pay over a year?

New cost: \(25 \times (1 + 0.12) = 25 \times 1.12 = £28\).

Extra per month: \(£28 - £25 = £3\).

Extra over a year: \(£3 \times 12 = £36\).

\begin{center}\rule{0.5\linewidth}{0.5pt}\end{center}

\textbf{Example 3 --- Geography: population.} A city has a population of
850,000. The annual growth rate is 1.4\%. What is the projected
population in 10 years?

\[P_{10} = 850{,}000 \times (1.014)^{10}\]

\((1.014)^{10} = 1.1495\) (to 4 d.p.)

\[P_{10} = 850{,}000 \times 1.1495 \approx 977{,}000\]

\begin{center}\rule{0.5\linewidth}{0.5pt}\end{center}

\textbf{Example 4 --- Science: radioactive decay.} A sample contains
5,000 atoms of a radioactive isotope that decays at 12\% per hour. How
many atoms remain after 4 hours?

This is compound decay --- each hour, 12\% is lost, so 88\% remains.

\[N_4 = 5000 \times (1 - 0.12)^4 = 5000 \times (0.88)^4\]

\((0.88)^4 = 0.5997\)

\[N_4 = 5000 \times 0.5997 \approx 2,999 \text{ atoms}\]

In 4 hours, roughly 40\% of the atoms have decayed.

\begin{center}\rule{0.5\linewidth}{0.5pt}\end{center}

\textbf{Interactive: Compound growth explorer.} Adjust the annual rate
and number of years. The bars show the compounded value year by year.
The dashed line shows what simple interest would give over the same
period.

\begin{Shaded}
\begin{Highlighting}[]
\NormalTok{//| echo: false}
\NormalTok{viewof annualRate = Inputs.range(}
\NormalTok{  [0, 20], \{ label: "Annual rate (\%)", step: 0.5, value: 5 \}}
\NormalTok{)}
\end{Highlighting}
\end{Shaded}

\begin{Shaded}
\begin{Highlighting}[]
\NormalTok{//| echo: false}
\NormalTok{viewof numYears = Inputs.range(}
\NormalTok{  [1, 20], \{ label: "Number of years", step: 1, value: 10 \}}
\NormalTok{)}
\end{Highlighting}
\end{Shaded}

\begin{Shaded}
\begin{Highlighting}[]
\NormalTok{//| echo: false}
\NormalTok{\{}
\NormalTok{  const W = 560;}
\NormalTok{  const H = 300;}
\NormalTok{  const MARGIN = \{ top: 30, right: 20, bottom: 50, left: 62 \};}
\NormalTok{  const innerW = W {-} MARGIN.left {-} MARGIN.right;}
\NormalTok{  const innerH = H {-} MARGIN.top {-} MARGIN.bottom;}

\NormalTok{  const Q0 = 1000;  // base value — shown as index, labelled clearly}
\NormalTok{  const r  = (Number.isFinite(annualRate) ? annualRate : 5) / 100;}
\NormalTok{  const N  = Number.isFinite(numYears) \&\& numYears \textgreater{}= 1 ? Math.round(numYears) : 10;}

\NormalTok{  // Build data arrays (year 0 to N)}
\NormalTok{  const years = d3.range(0, N + 1);}
\NormalTok{  const compound = years.map(n =\textgreater{} Q0 * Math.pow(1 + r, n));}
\NormalTok{  const simple   = years.map(n =\textgreater{} Q0 * (1 + r * n));}

\NormalTok{  const finalCompound = compound[N];}
\NormalTok{  const finalSimple   = simple[N];}
\NormalTok{  const extraStr = (finalCompound {-} finalSimple).toFixed(2);}

\NormalTok{  const yMax = Math.max(...compound, ...simple) * 1.08;}

\NormalTok{  const xScale = d3.scaleBand()}
\NormalTok{    .domain(years)}
\NormalTok{    .range([0, innerW])}
\NormalTok{    .padding(0.18);}

\NormalTok{  const yScale = d3.scaleLinear()}
\NormalTok{    .domain([0, yMax])}
\NormalTok{    .range([innerH, 0]);}

\NormalTok{  const svg = d3.create("svg")}
\NormalTok{    .attr("viewBox", \textasciigrave{}0 0 $\{W\} $\{H\}\textasciigrave{})}
\NormalTok{    .attr("width", "100\%")}
\NormalTok{    .attr("style", \textasciigrave{}max{-}width:$\{W\}px; font{-}family: inherit;\textasciigrave{});}

\NormalTok{  const g = svg.append("g")}
\NormalTok{    .attr("transform", \textasciigrave{}translate($\{MARGIN.left\},$\{MARGIN.top\})\textasciigrave{});}

\NormalTok{  // Gridlines}
\NormalTok{  g.append("g")}
\NormalTok{    .attr("class", "grid")}
\NormalTok{    .call(}
\NormalTok{      d3.axisLeft(yScale)}
\NormalTok{        .ticks(5)}
\NormalTok{        .tickSize({-}innerW)}
\NormalTok{        .tickFormat("")}
\NormalTok{    )}
\NormalTok{    .call(gg =\textgreater{} gg.select(".domain").remove())}
\NormalTok{    .call(gg =\textgreater{} gg.selectAll("line").attr("stroke", "\#e5e7eb").attr("stroke{-}dasharray", "3,3"));}

\NormalTok{  // Compound bars}
\NormalTok{  g.selectAll(".bar{-}compound")}
\NormalTok{    .data(years)}
\NormalTok{    .enter()}
\NormalTok{    .append("rect")}
\NormalTok{      .attr("class", "bar{-}compound")}
\NormalTok{      .attr("x", d =\textgreater{} xScale(d))}
\NormalTok{      .attr("y", d =\textgreater{} yScale(compound[d]))}
\NormalTok{      .attr("width", xScale.bandwidth())}
\NormalTok{      .attr("height", d =\textgreater{} innerH {-} yScale(compound[d]))}
\NormalTok{      .attr("fill", "\#0d9488")}
\NormalTok{      .attr("rx", 2);}

\NormalTok{  // Simple interest line}
\NormalTok{  const lineGen = d3.line()}
\NormalTok{    .x(d =\textgreater{} xScale(d) + xScale.bandwidth() / 2)}
\NormalTok{    .y(d =\textgreater{} yScale(simple[d]));}

\NormalTok{  g.append("path")}
\NormalTok{    .datum(years)}
\NormalTok{    .attr("fill", "none")}
\NormalTok{    .attr("stroke", "\#d97706")}
\NormalTok{    .attr("stroke{-}width", 2)}
\NormalTok{    .attr("stroke{-}dasharray", "5,4")}
\NormalTok{    .attr("d", lineGen);}

\NormalTok{  // Simple interest dots}
\NormalTok{  g.selectAll(".dot{-}simple")}
\NormalTok{    .data(years)}
\NormalTok{    .enter()}
\NormalTok{    .append("circle")}
\NormalTok{      .attr("cx", d =\textgreater{} xScale(d) + xScale.bandwidth() / 2)}
\NormalTok{      .attr("cy", d =\textgreater{} yScale(simple[d]))}
\NormalTok{      .attr("r", 3)}
\NormalTok{      .attr("fill", "\#d97706");}

\NormalTok{  // X axis}
\NormalTok{  g.append("g")}
\NormalTok{    .attr("transform", \textasciigrave{}translate(0,$\{innerH\})\textasciigrave{})}
\NormalTok{    .call(}
\NormalTok{      d3.axisBottom(xScale)}
\NormalTok{        .tickValues(years.filter((\_, i) =\textgreater{} N \textless{}= 10 || i \% 2 === 0))}
\NormalTok{        .tickFormat(d =\textgreater{} \textasciigrave{}Y$\{d\}\textasciigrave{})}
\NormalTok{    )}
\NormalTok{    .call(gg =\textgreater{} gg.select(".domain").attr("stroke", "\#9ca3af"))}
\NormalTok{    .call(gg =\textgreater{} gg.selectAll("text").attr("font{-}size", "11px"));}

\NormalTok{  // X axis label}
\NormalTok{  g.append("text")}
\NormalTok{    .attr("x", innerW / 2)}
\NormalTok{    .attr("y", innerH + 40)}
\NormalTok{    .attr("text{-}anchor", "middle")}
\NormalTok{    .attr("fill", "\#6b7280")}
\NormalTok{    .attr("font{-}size", "12px")}
\NormalTok{    .text("Year");}

\NormalTok{  // Y axis}
\NormalTok{  g.append("g")}
\NormalTok{    .call(d3.axisLeft(yScale).ticks(5).tickFormat(d =\textgreater{} \textasciigrave{}£$\{d\}\textasciigrave{}))}
\NormalTok{    .call(gg =\textgreater{} gg.select(".domain").attr("stroke", "\#9ca3af"))}
\NormalTok{    .call(gg =\textgreater{} gg.selectAll("text").attr("font{-}size", "11px"));}

\NormalTok{  // Final value label on last bar}
\NormalTok{  g.append("text")}
\NormalTok{    .attr("x", xScale(N) + xScale.bandwidth() / 2)}
\NormalTok{    .attr("y", yScale(finalCompound) {-} 6)}
\NormalTok{    .attr("text{-}anchor", "middle")}
\NormalTok{    .attr("fill", "\#065f46")}
\NormalTok{    .attr("font{-}size", "11px")}
\NormalTok{    .attr("font{-}weight", "700")}
\NormalTok{    .text(\textasciigrave{}£$\{finalCompound.toFixed(0)\}\textasciigrave{});}

\NormalTok{  // Legend}
\NormalTok{  const legendY = {-}18;}
\NormalTok{  // Compound bar swatch}
\NormalTok{  g.append("rect")}
\NormalTok{    .attr("x", 0).attr("y", legendY {-} 9).attr("width", 12).attr("height", 12)}
\NormalTok{    .attr("fill", "\#0d9488").attr("rx", 2);}
\NormalTok{  g.append("text")}
\NormalTok{    .attr("x", 16).attr("y", legendY)}
\NormalTok{    .attr("fill", "\#374151").attr("font{-}size", "11px")}
\NormalTok{    .text("Compound growth");}

\NormalTok{  // Simple line swatch}
\NormalTok{  g.append("line")}
\NormalTok{    .attr("x1", 140).attr("y1", legendY {-} 3)}
\NormalTok{    .attr("x2", 152).attr("y2", legendY {-} 3)}
\NormalTok{    .attr("stroke", "\#d97706").attr("stroke{-}width", 2).attr("stroke{-}dasharray", "4,3");}
\NormalTok{  g.append("text")}
\NormalTok{    .attr("x", 156).attr("y", legendY)}
\NormalTok{    .attr("fill", "\#374151").attr("font{-}size", "11px")}
\NormalTok{    .text("Simple interest");}

\NormalTok{  // Extra label}
\NormalTok{  g.append("text")}
\NormalTok{    .attr("x", innerW)}
\NormalTok{    .attr("y", innerH + 40)}
\NormalTok{    .attr("text{-}anchor", "end")}
\NormalTok{    .attr("fill", "\#0d9488")}
\NormalTok{    .attr("font{-}size", "12px")}
\NormalTok{    .attr("font{-}weight", "600")}
\NormalTok{    .text(\textasciigrave{}Compounding adds £$\{extraStr\} extra over $\{N\} year$\{N === 1 ? "" : "s"\}\textasciigrave{});}

\NormalTok{  return svg.node();}
\NormalTok{\}}
\end{Highlighting}
\end{Shaded}

\begin{center}\rule{0.5\linewidth}{0.5pt}\end{center}

\section{Where this goes}\label{where-this-goes-2}

Percent applied once is straightforward proportion. Applied repeatedly,
it produces exponential behaviour --- and \textbf{exponents and
logarithms} (Volume 4, Chapter 1) is where that behaviour is studied
properly. Logarithms undo the compounding: if you want to know how long
it takes for a quantity to double, you are solving an exponential
equation, which requires a logarithm.

\textbf{Statistics and probability} (Chapter 6) uses rates constantly:
probability is a rate between 0 and 1, and relative frequency is a
percent.

In finance, virtually every calculation --- present value, future value,
yield, annuity --- is a variant of the compound formula. In engineering,
efficiency, signal attenuation, and structural safety factors are all
percent-based ratios. The formula here is simple. The consequences are
not.

\begin{tcolorbox}[enhanced jigsaw, opacityback=0, colback=white, breakable, colframe=quarto-callout-tip-color-frame, arc=.35mm, left=2mm, toprule=.15mm, leftrule=.75mm, bottomrule=.15mm, rightrule=.15mm]
\begin{minipage}[t]{5.5mm}
\textcolor{quarto-callout-tip-color}{\faLightbulb}
\end{minipage}%
\begin{minipage}[t]{\textwidth - 5.5mm}

\textbf{Where this shows up}

\begin{itemize}
\tightlist
\item
  A savings account paying 4\% a year is using the compound formula ---
  the interest each year is calculated on the new, larger balance.
\item
  A streaming service raising prices by 8\% is applying a percent
  increase to a base price.
\item
  A sports commentator quoting a 65\% win rate is expressing a fraction
  as a percent.
\item
  A phone showing ``battery 23\%'' is telling you what fraction of
  capacity remains.
\end{itemize}

The arithmetic is identical. The stakes vary.

\end{minipage}%
\end{tcolorbox}

\section{Exercises}\label{exercises-2}

\begin{enumerate}
\def\labelenumi{\arabic{enumi}.}
\item
  A laptop costs £1,149. It is on sale at 15\% off. What is the sale
  price?
\item
  A city had 320,000 residents in 2015 and 387,000 in 2025. What was the
  total percent change over that decade?
\item
  An investment grows from £4,500 to £5,200 over 3 years. What was the
  approximate annual growth rate? \emph{(Estimate; you do not yet need
  logarithms --- try 3--4 percent and refine.)}
\item
  A water tank contains 12,000 litres. It leaks at a rate of 5\% of its
  current volume per day. How much water remains after 1 week?
\item
  A car's drivetrain transfers engine power to the wheels with 87\%
  efficiency. If the engine produces 110 kW, what power reaches the
  wheels? What percentage is lost?
\item
  A country's annual inflation rate is 2.8\%. A loaf of bread costs
  £1.40 today. What will it cost in 8 years at this rate?
\item
  A business owner offers a 10\% early-payment discount. A customer
  receives an invoice for £3,400. They pay early. How much do they pay?
  What percent of the original price was the discount in pounds?
\end{enumerate}

\chapter{Integers and signed
arithmetic}\label{integers-and-signed-arithmetic}

Moving in both directions

\hfill\break

Your bank balance is £40 and you spend £60. You're now at −£20.

It's −3°C at night. By morning it's dropped another 8 degrees. Where are
you on the thermometer?

You're on the ground floor of a car park. The exit is two floors up; the
lower deck is three floors down. Which direction is floor −3?

All three questions use the same operation: addition on the integer
number line. Negative numbers aren't strange --- they're just the left
side of zero.

\section{What the notation is
saying}\label{what-the-notation-is-saying-3}

An \textbf{integer} is any whole number: \ldots, −3, −2, −1, 0, 1, 2, 3,
\ldots{}

The negative sign in front of a number means ``opposite direction'':
\(-5\) is 5 units to the left of zero.

Adding a negative moves you left on the number line:
\[7 + (-3) = 7 - 3 = 4\]

Subtracting a negative moves you right --- you reverse direction twice:
\[5 - (-2) = 5 + 2 = 7\]

That second one trips people up. Think of it this way: if someone
cancels a debt you owe, you end up better off. Taking away a negative is
the same as adding a positive.

\textbf{The double negative rule:} \(-(-a) = a\). Two negatives cancel
out and you get the original positive back.

\textbf{Sign rules for multiplication and division:}

\begin{longtable}[]{@{}ll@{}}
\toprule\noalign{}
Operation & Sign result \\
\midrule\noalign{}
\endhead
\bottomrule\noalign{}
\endlastfoot
\((+) \times (+)\) & \(+\) \\
\((-) \times (-)\) & \(+\) \\
\((+) \times (-)\) & \(-\) \\
\((-) \times (+)\) & \(-\) \\
\end{longtable}

Same signs give positive. Different signs give negative. The same rules
apply to division.

\begin{center}\rule{0.5\linewidth}{0.5pt}\end{center}

\textbf{Interactive: Signed number line.} Set a starting value and a
change. The arrow shows the movement; the equation updates live.

\begin{Shaded}
\begin{Highlighting}[]
\NormalTok{//| echo: false}
\NormalTok{d3 = require("d3@7")}
\end{Highlighting}
\end{Shaded}

\begin{Shaded}
\begin{Highlighting}[]
\NormalTok{//| echo: false}
\NormalTok{viewof startVal = Inputs.range(}
\NormalTok{  [{-}10, 10], \{ label: "Starting value", step: 1, value: 3 \}}
\NormalTok{)}
\end{Highlighting}
\end{Shaded}

\begin{Shaded}
\begin{Highlighting}[]
\NormalTok{//| echo: false}
\NormalTok{viewof changeVal = Inputs.range(}
\NormalTok{  [{-}10, 10], \{ label: "Change (add this)", step: 1, value: {-}5 \}}
\NormalTok{)}
\end{Highlighting}
\end{Shaded}

\begin{Shaded}
\begin{Highlighting}[]
\NormalTok{//| echo: false}
\NormalTok{\{}
\NormalTok{  const W = 560;}
\NormalTok{  const H = 130;}
\NormalTok{  const LINE\_Y = 58;}
\NormalTok{  const TICK\_H = 10;}
\NormalTok{  const PAD\_L = 28;}
\NormalTok{  const PAD\_R = 28;}
\NormalTok{  const MIN\_VAL = {-}15;}
\NormalTok{  const MAX\_VAL = 15;}

\NormalTok{  const start  = Number.isFinite(startVal)  ? Math.round(startVal)  : 3;}
\NormalTok{  const change = Number.isFinite(changeVal) ? Math.round(changeVal) : {-}5;}
\NormalTok{  const result = start + change;}

\NormalTok{  // Clamp result display — result can go outside visible range; still show equation}
\NormalTok{  const visibleResult = Math.max(MIN\_VAL, Math.min(MAX\_VAL, result));}

\NormalTok{  const usableW = W {-} PAD\_L {-} PAD\_R;}
\NormalTok{  const xScale = d3.scaleLinear()}
\NormalTok{    .domain([MIN\_VAL, MAX\_VAL])}
\NormalTok{    .range([PAD\_L, PAD\_L + usableW]);}

\NormalTok{  const svg = d3.create("svg")}
\NormalTok{    .attr("viewBox", \textasciigrave{}0 0 $\{W\} $\{H\}\textasciigrave{})}
\NormalTok{    .attr("width", "100\%")}
\NormalTok{    .attr("style", \textasciigrave{}max{-}width:$\{W\}px; font{-}family: inherit;\textasciigrave{});}

\NormalTok{  // Number line}
\NormalTok{  svg.append("line")}
\NormalTok{    .attr("x1", PAD\_L {-} 8).attr("y1", LINE\_Y)}
\NormalTok{    .attr("x2", PAD\_L + usableW + 8).attr("y2", LINE\_Y)}
\NormalTok{    .attr("stroke", "\#9ca3af").attr("stroke{-}width", 1.5);}

\NormalTok{  // Arrowheads on line ends}
\NormalTok{  const arrowPath = (x, y, dir) =\textgreater{}}
\NormalTok{    \textasciigrave{}M$\{x\},$\{y\} L$\{x + dir * 7\},$\{y {-} 5\} L$\{x + dir * 7\},$\{y + 5\} Z\textasciigrave{};}
\NormalTok{  svg.append("path").attr("d", arrowPath(PAD\_L {-} 8, LINE\_Y, {-}1)).attr("fill", "\#9ca3af");}
\NormalTok{  svg.append("path").attr("d", arrowPath(PAD\_L + usableW + 8, LINE\_Y, 1)).attr("fill", "\#9ca3af");}

\NormalTok{  // Tick marks and labels}
\NormalTok{  for (let v = MIN\_VAL; v \textless{}= MAX\_VAL; v++) \{}
\NormalTok{    const x = xScale(v);}
\NormalTok{    const isMajor = v \% 5 === 0;}
\NormalTok{    svg.append("line")}
\NormalTok{      .attr("x1", x).attr("y1", LINE\_Y {-} (isMajor ? TICK\_H : TICK\_H * 0.6))}
\NormalTok{      .attr("x2", x).attr("y2", LINE\_Y + (isMajor ? TICK\_H : TICK\_H * 0.6))}
\NormalTok{      .attr("stroke", isMajor ? "\#6b7280" : "\#d1d5db")}
\NormalTok{      .attr("stroke{-}width", isMajor ? 1.5 : 1);}
\NormalTok{    if (isMajor) \{}
\NormalTok{      svg.append("text")}
\NormalTok{        .attr("x", x).attr("y", LINE\_Y + TICK\_H + 13)}
\NormalTok{        .attr("text{-}anchor", "middle")}
\NormalTok{        .attr("fill", "\#6b7280").attr("font{-}size", "11px")}
\NormalTok{        .text(v);}
\NormalTok{    \}}
\NormalTok{  \}}

\NormalTok{  // Arrow showing movement (from start to result, clamped to visible range)}
\NormalTok{  const arrowStartX  = xScale(Math.max(MIN\_VAL, Math.min(MAX\_VAL, start)));}
\NormalTok{  const arrowEndX    = xScale(visibleResult);}
\NormalTok{  const ARROW\_Y      = LINE\_Y {-} 22;}
\NormalTok{  const isOutOfRange = result \textless{} MIN\_VAL || result \textgreater{} MAX\_VAL;}

\NormalTok{  if (change !== 0) \{}
\NormalTok{    const colour = change \textgreater{} 0 ? "\#0d9488" : "\#d97706";}
\NormalTok{    svg.append("line")}
\NormalTok{      .attr("x1", arrowStartX).attr("y1", ARROW\_Y)}
\NormalTok{      .attr("x2", arrowEndX).attr("y2", ARROW\_Y)}
\NormalTok{      .attr("stroke", colour).attr("stroke{-}width", 2.5);}

\NormalTok{    // Arrowhead at end}
\NormalTok{    const dir = change \textgreater{} 0 ? 1 : {-}1;}
\NormalTok{    svg.append("path")}
\NormalTok{      .attr("d", \textasciigrave{}M$\{arrowEndX\},$\{ARROW\_Y\} L$\{arrowEndX {-} dir * 8\},$\{ARROW\_Y {-} 5\} L$\{arrowEndX {-} dir * 8\},$\{ARROW\_Y + 5\} Z\textasciigrave{})}
\NormalTok{      .attr("fill", colour);}
\NormalTok{  \}}

\NormalTok{  // Start position dot}
\NormalTok{  svg.append("circle")}
\NormalTok{    .attr("cx", arrowStartX).attr("cy", LINE\_Y)}
\NormalTok{    .attr("r", 7)}
\NormalTok{    .attr("fill", "\#4b5563")}
\NormalTok{    .attr("stroke", "\#fff").attr("stroke{-}width", 1.5);}

\NormalTok{  svg.append("text")}
\NormalTok{    .attr("x", arrowStartX).attr("cy", LINE\_Y)}
\NormalTok{    .attr("y", LINE\_Y {-} 12)}
\NormalTok{    .attr("text{-}anchor", "middle")}
\NormalTok{    .attr("fill", "\#374151").attr("font{-}size", "11px").attr("font{-}weight", "600")}
\NormalTok{    .text(start);}

\NormalTok{  // Result position dot}
\NormalTok{  const resultColour = result \textgreater{} 0 ? "\#0d9488" : result \textless{} 0 ? "\#d97706" : "\#4b5563";}
\NormalTok{  svg.append("circle")}
\NormalTok{    .attr("cx", xScale(visibleResult)).attr("cy", LINE\_Y)}
\NormalTok{    .attr("r", 7)}
\NormalTok{    .attr("fill", resultColour)}
\NormalTok{    .attr("stroke", "\#fff").attr("stroke{-}width", 1.5);}

\NormalTok{  svg.append("text")}
\NormalTok{    .attr("x", xScale(visibleResult))}
\NormalTok{    .attr("y", LINE\_Y + TICK\_H + 28)}
\NormalTok{    .attr("text{-}anchor", "middle")}
\NormalTok{    .attr("fill", resultColour).attr("font{-}size", "11px").attr("font{-}weight", "700")}
\NormalTok{    .text(isOutOfRange ? \textasciigrave{}$\{result\} (off scale)\textasciigrave{} : result);}

\NormalTok{  // Equation label}
\NormalTok{  const changeSign = change \textgreater{}= 0 ? \textasciigrave{}+ $\{change\}\textasciigrave{} : \textasciigrave{}+ ($\{change\})\textasciigrave{};}
\NormalTok{  const eqLabel    = \textasciigrave{}$\{start\}  $\{changeSign\}  =  $\{result\}\textasciigrave{};}
\NormalTok{  svg.append("text")}
\NormalTok{    .attr("x", W / 2).attr("y", H {-} 6)}
\NormalTok{    .attr("text{-}anchor", "middle")}
\NormalTok{    .attr("fill", "\#1f2937").attr("font{-}size", "13px").attr("font{-}weight", "700")}
\NormalTok{    .text(eqLabel);}

\NormalTok{  return svg.node();}
\NormalTok{\}}
\end{Highlighting}
\end{Shaded}

\begin{center}\rule{0.5\linewidth}{0.5pt}\end{center}

\section{The method}\label{the-method-3}

\textbf{Addition and subtraction}

Rewrite every subtraction as addition of the opposite:
\(a - b = a + (-b)\).

Then use the number line: positive means move right, negative means move
left. Your result is where you land.

\textbf{Multiplication}

\begin{enumerate}
\def\labelenumi{\arabic{enumi}.}
\tightlist
\item
  Multiply the absolute values (ignore signs for now).
\item
  Count the number of negative factors.

  \begin{itemize}
  \tightlist
  \item
    Even number of negatives → positive result.
  \item
    Odd number of negatives → negative result.
  \end{itemize}
\end{enumerate}

\textbf{Division}

Same rule: divide the absolute values, then apply the sign rule.

\textbf{Order of operations with signed numbers}

Brackets first. Then signs attached to numbers. Then multiply/divide
left to right. Then add/subtract left to right.

\begin{tcolorbox}[enhanced jigsaw, opacityback=0, colback=white, breakable, colframe=quarto-callout-note-color-frame, arc=.35mm, left=2mm, toprule=.15mm, leftrule=.75mm, bottomrule=.15mm, rightrule=.15mm]
\begin{minipage}[t]{5.5mm}
\textcolor{quarto-callout-note-color}{\faInfo}
\end{minipage}%
\begin{minipage}[t]{\textwidth - 5.5mm}

\textbf{Why this works}

Think of multiplication as repeated addition. \(3 \times (-4)\) means
``add \(-4\) three times'': \((-4) + (-4) + (-4) = -12\). That gives
\((+) \times (-) = (-)\).

For \((-) \times (-)\): \(-3 \times (-4)\) means ``add \(-4\) negative
three times.'' Adding a negative thing a negative number of times is
removing it --- which is positive. The algebraic proof uses the
distributive law:
\(0 = (-3)(4 + (-4)) = (-3)(4) + (-3)(-4) = -12 + (-3)(-4)\), so
\((-3)(-4) = 12\).

\end{minipage}%
\end{tcolorbox}

\begin{center}\rule{0.5\linewidth}{0.5pt}\end{center}

\textbf{Interactive: Sign rules for multiplication.} Click any cell in
the table to see the calculation and the rule in action. All four
combinations are visible at once.

\begin{Shaded}
\begin{Highlighting}[]
\NormalTok{//| echo: false}
\NormalTok{viewof selectedCell = \{}
\NormalTok{  // Default to top{-}left cell on load}
\NormalTok{  const init = \{ a: 1, b: 1 \};}
\NormalTok{  return Inputs.input(init);}
\NormalTok{\}}
\end{Highlighting}
\end{Shaded}

\begin{Shaded}
\begin{Highlighting}[]
\NormalTok{//| echo: false}
\NormalTok{\{}
\NormalTok{  // Fixed example magnitudes for each sign combination}
\NormalTok{  // Rows: sign of first number (+/{-})}
\NormalTok{  // Cols: sign of second number (+/{-})}
\NormalTok{  const COMBOS = [}
\NormalTok{    \{ aSign: +1, bSign: +1, aAbs: 3, bAbs: 4, row: 0, col: 0 \},}
\NormalTok{    \{ aSign: +1, bSign: {-}1, aAbs: 3, bAbs: 4, row: 0, col: 1 \},}
\NormalTok{    \{ aSign: {-}1, bSign: +1, aAbs: 3, bAbs: 4, row: 1, col: 0 \},}
\NormalTok{    \{ aSign: {-}1, bSign: {-}1, aAbs: 3, bAbs: 4, row: 1, col: 1 \},}
\NormalTok{  ];}

\NormalTok{  const W = 560;}
\NormalTok{  const H = 280;}
\NormalTok{  const CELL\_W = 200;}
\NormalTok{  const CELL\_H = 88;}
\NormalTok{  const GRID\_X = 100;}
\NormalTok{  const GRID\_Y = 50;}
\NormalTok{  const COL\_HEADER\_Y = 28;}
\NormalTok{  const ROW\_HEADER\_X = 52;}

\NormalTok{  const posColour  = "\#0d9488";  // teal{-}600 — positive result}
\NormalTok{  const negColour  = "\#d97706";  // amber{-}600 — negative result}
\NormalTok{  const selStroke  = "\#1f2937";}
\NormalTok{  const dimFill    = "\#f9fafb";}

\NormalTok{  const svg = d3.create("svg")}
\NormalTok{    .attr("viewBox", \textasciigrave{}0 0 $\{W\} $\{H\}\textasciigrave{})}
\NormalTok{    .attr("width", "100\%")}
\NormalTok{    .attr("style", \textasciigrave{}max{-}width:$\{W\}px; font{-}family: inherit; cursor: pointer;\textasciigrave{});}

\NormalTok{  // Column headers: sign of second number}
\NormalTok{  const colLabels = ["Second number: +", "Second number: −"];}
\NormalTok{  colLabels.forEach((label, ci) =\textgreater{} \{}
\NormalTok{    svg.append("text")}
\NormalTok{      .attr("x", GRID\_X + ci * CELL\_W + CELL\_W / 2)}
\NormalTok{      .attr("y", COL\_HEADER\_Y)}
\NormalTok{      .attr("text{-}anchor", "middle")}
\NormalTok{      .attr("fill", "\#374151")}
\NormalTok{      .attr("font{-}size", "12px")}
\NormalTok{      .attr("font{-}weight", "600")}
\NormalTok{      .text(label);}
\NormalTok{  \});}

\NormalTok{  // Row headers: sign of first number}
\NormalTok{  const rowLabels = ["First: +", "First: −"];}
\NormalTok{  rowLabels.forEach((label, ri) =\textgreater{} \{}
\NormalTok{    svg.append("text")}
\NormalTok{      .attr("x", ROW\_HEADER\_X)}
\NormalTok{      .attr("y", GRID\_Y + ri * CELL\_H + CELL\_H / 2 + 4)}
\NormalTok{      .attr("text{-}anchor", "middle")}
\NormalTok{      .attr("fill", "\#374151")}
\NormalTok{      .attr("font{-}size", "12px")}
\NormalTok{      .attr("font{-}weight", "600")}
\NormalTok{      .text(label);}
\NormalTok{  \});}

\NormalTok{  COMBOS.forEach(c =\textgreater{} \{}
\NormalTok{    const aVal   = c.aSign * c.aAbs;}
\NormalTok{    const bVal   = c.bSign * c.bAbs;}
\NormalTok{    const result = aVal * bVal;}
\NormalTok{    const isPos  = result \textgreater{} 0;}
\NormalTok{    const colour = isPos ? posColour : negColour;}

\NormalTok{    const sel = selectedCell;}
\NormalTok{    const isSelected = sel \&\& sel.a === c.aSign \&\& sel.b === c.bSign;}

\NormalTok{    const cellX = GRID\_X + c.col * CELL\_W;}
\NormalTok{    const cellY = GRID\_Y + c.row * CELL\_H;}

\NormalTok{    // Cell background}
\NormalTok{    const g = svg.append("g")}
\NormalTok{      .attr("transform", \textasciigrave{}translate($\{cellX\},$\{cellY\})\textasciigrave{})}
\NormalTok{      .style("cursor", "pointer")}
\NormalTok{      .on("click", () =\textgreater{} \{}
\NormalTok{        selectedCell.value = \{ a: c.aSign, b: c.bSign \};}
\NormalTok{        selectedCell.dispatchEvent(new Event("input"));}
\NormalTok{      \});}

\NormalTok{    g.append("rect")}
\NormalTok{      .attr("width", CELL\_W {-} 4)}
\NormalTok{      .attr("height", CELL\_H {-} 4)}
\NormalTok{      .attr("x", 2).attr("y", 2)}
\NormalTok{      .attr("rx", 6)}
\NormalTok{      .attr("fill", isSelected ? colour : dimFill)}
\NormalTok{      .attr("stroke", isSelected ? selStroke : "\#d1d5db")}
\NormalTok{      .attr("stroke{-}width", isSelected ? 2 : 1);}

\NormalTok{    // Calculation line: e.g. "+3 × −4 = −12"}
\NormalTok{    const aStr = c.aSign \textgreater{} 0 ? \textasciigrave{}+$\{c.aAbs\}\textasciigrave{} : \textasciigrave{}−$\{c.aAbs\}\textasciigrave{};}
\NormalTok{    const bStr = c.bSign \textgreater{} 0 ? \textasciigrave{}+$\{c.bAbs\}\textasciigrave{} : \textasciigrave{}−$\{c.bAbs\}\textasciigrave{};}
\NormalTok{    const rStr = result \textgreater{} 0 ? \textasciigrave{}+$\{Math.abs(result)\}\textasciigrave{} : \textasciigrave{}−$\{Math.abs(result)\}\textasciigrave{};}

\NormalTok{    g.append("text")}
\NormalTok{      .attr("x", (CELL\_W {-} 4) / 2)}
\NormalTok{      .attr("y", 28)}
\NormalTok{      .attr("text{-}anchor", "middle")}
\NormalTok{      .attr("fill", isSelected ? "\#fff" : "\#1f2937")}
\NormalTok{      .attr("font{-}size", "16px")}
\NormalTok{      .attr("font{-}weight", "700")}
\NormalTok{      .text(\textasciigrave{}$\{aStr\} × $\{bStr\} = $\{rStr\}\textasciigrave{});}

\NormalTok{    // Sign rule label}
\NormalTok{    const aSignWord = c.aSign \textgreater{} 0 ? "positive" : "negative";}
\NormalTok{    const bSignWord = c.bSign \textgreater{} 0 ? "positive" : "negative";}
\NormalTok{    const rSignWord = isPos ? "positive" : "negative";}
\NormalTok{    const ruleLine1 = \textasciigrave{}$\{aSignWord\} × $\{bSignWord\}\textasciigrave{};}
\NormalTok{    const ruleLine2 = \textasciigrave{}= $\{rSignWord\}\textasciigrave{};}

\NormalTok{    g.append("text")}
\NormalTok{      .attr("x", (CELL\_W {-} 4) / 2)}
\NormalTok{      .attr("y", 50)}
\NormalTok{      .attr("text{-}anchor", "middle")}
\NormalTok{      .attr("fill", isSelected ? "\#fff" : "\#6b7280")}
\NormalTok{      .attr("font{-}size", "11px")}
\NormalTok{      .text(ruleLine1);}

\NormalTok{    g.append("text")}
\NormalTok{      .attr("x", (CELL\_W {-} 4) / 2)}
\NormalTok{      .attr("y", 64)}
\NormalTok{      .attr("text{-}anchor", "middle")}
\NormalTok{      .attr("fill", isSelected ? "\#fff" : colour)}
\NormalTok{      .attr("font{-}size", "11px")}
\NormalTok{      .attr("font{-}weight", "600")}
\NormalTok{      .text(ruleLine2);}

\NormalTok{    // Shape badge: circle = positive, diamond = negative (redundant encoding)}
\NormalTok{    const BADGE\_X = (CELL\_W {-} 4) / 2;}
\NormalTok{    const BADGE\_Y = 78;}
\NormalTok{    if (isPos) \{}
\NormalTok{      g.append("circle")}
\NormalTok{        .attr("cx", BADGE\_X).attr("cy", BADGE\_Y {-} 4)}
\NormalTok{        .attr("r", 5)}
\NormalTok{        .attr("fill", "none")}
\NormalTok{        .attr("stroke", isSelected ? "\#fff" : colour)}
\NormalTok{        .attr("stroke{-}width", 1.5);}
\NormalTok{    \} else \{}
\NormalTok{      g.append("polygon")}
\NormalTok{        .attr("points", \textasciigrave{}$\{BADGE\_X\},$\{BADGE\_Y {-} 9\} $\{BADGE\_X + 6\},$\{BADGE\_Y {-} 4\} $\{BADGE\_X\},$\{BADGE\_Y + 1\} $\{BADGE\_X {-} 6\},$\{BADGE\_Y {-} 4\}\textasciigrave{})}
\NormalTok{        .attr("fill", "none")}
\NormalTok{        .attr("stroke", isSelected ? "\#fff" : colour)}
\NormalTok{        .attr("stroke{-}width", 1.5);}
\NormalTok{    \}}
\NormalTok{  \});}

\NormalTok{  // Summary rule line at bottom}
\NormalTok{  svg.append("text")}
\NormalTok{    .attr("x", W / 2)}
\NormalTok{    .attr("y", H {-} 10)}
\NormalTok{    .attr("text{-}anchor", "middle")}
\NormalTok{    .attr("fill", "\#374151")}
\NormalTok{    .attr("font{-}size", "12px")}
\NormalTok{    .text("Same signs → positive result   |   Different signs → negative result   |   Click any cell to highlight it");}

\NormalTok{  return svg.node();}
\NormalTok{\}}
\end{Highlighting}
\end{Shaded}

\begin{center}\rule{0.5\linewidth}{0.5pt}\end{center}

\section{Worked examples}\label{worked-examples-3}

\textbf{Example 1 --- Bank balance.} You start the week with £50 in your
account. You spend £30 on food, £25 on a game, and then get paid £60
from a part-time shift. What's your balance at the end of the week?

Add up all the changes --- spends are negative, income is positive:
\[50 + (-30) + (-25) + 60 = 50 - 30 - 25 + 60 = 55\]

End balance: £55.

\begin{center}\rule{0.5\linewidth}{0.5pt}\end{center}

\textbf{Example 2 --- Temperature change.} The temperature at 6 am is
−5°C. By midday it has risen 14°C. By midnight it has fallen 18°C from
the midday high. What is the midnight temperature?

Work through the changes in order:

\[-5 + 14 = 9°\text{C} \quad \text{(midday)}\]
\[9 - 18 = -9°\text{C} \quad \text{(midnight)}\]

The midnight temperature is −9°C.

\begin{center}\rule{0.5\linewidth}{0.5pt}\end{center}

\textbf{Example 3 --- Score in a game.} A quiz game gives +5 for a
correct answer and −3 for a wrong one. A player gets 4 right and 6
wrong. What is their score?

Correct answers: \(4 \times 5 = 20\). Wrong answers:
\(6 \times (-3) = -18\) --- six negatives, so the result is negative.

Total score: \(20 + (-18) = 2\).

\begin{center}\rule{0.5\linewidth}{0.5pt}\end{center}

\textbf{Example 4 --- Running total with negatives.} Evaluate
\((-4) \times (-3) + 2 \times (-5)\) without a calculator.

Handle the multiplications first --- left to right:

\((-4) \times (-3) = 12\) (same signs → positive), giving \(12\).

\(2 \times (-5) = -10\) (different signs → negative), giving \(-10\).

Now add: \(12 + (-10) = 2\).

\section{Where this goes}\label{where-this-goes-3}

Signed arithmetic is the prerequisite for \textbf{linear equations}
(Volume 2, Chapter 1). When you solve \(3x - 7 = -19\), you need to add
7 to both sides and handle the resulting \(-12\) correctly. Every
algebraic manipulation involves signed numbers.

In physics, signed quantities appear everywhere: displacement, velocity,
acceleration, force, work, charge, and potential are all signed. The
sign carries physical meaning --- it encodes direction or polarity.
Getting the sign wrong is not just a numerical error; it is a physical
error.

\begin{tcolorbox}[enhanced jigsaw, opacityback=0, colback=white, breakable, colframe=quarto-callout-tip-color-frame, arc=.35mm, left=2mm, toprule=.15mm, leftrule=.75mm, bottomrule=.15mm, rightrule=.15mm]
\begin{minipage}[t]{5.5mm}
\textcolor{quarto-callout-tip-color}{\faLightbulb}
\end{minipage}%
\begin{minipage}[t]{\textwidth - 5.5mm}

\textbf{Where this shows up}

\begin{itemize}
\tightlist
\item
  A bank account or budget uses signed arithmetic every time money moves
  in or out --- the balance is an integer sum.
\item
  A weather forecast compares signed temperatures to report how much
  warmer or colder one day was than another.
\item
  A physicist writing Newton's second law includes sign: \(F = ma\),
  where \(F < 0\) means the force acts in the negative direction.
\item
  A programmer writing low-level code must track whether arithmetic will
  overflow a signed integer's range.
\end{itemize}

The rules are short. The mistakes, when you forget them, are expensive.

\end{minipage}%
\end{tcolorbox}

\section{Exercises}\label{exercises-3}

\begin{enumerate}
\def\labelenumi{\arabic{enumi}.}
\item
  A lift starts at floor 0 (ground). It goes down 3 floors, then up 7
  floors, then down 2 floors. What floor is it on?
\item
  A player's scores across five rounds of a quiz game: −4, +7, −2, +9,
  −3. What is their total score?
\item
  Evaluate without a calculator:

  \begin{enumerate}
  \def\labelenumii{\alph{enumii})}
  \tightlist
  \item
    \((-8) \times (-6)\)
  \item
    \((-15) \div 3\)
  \item
    \(4 \times (-9) \times (-2)\)
  \item
    \((-3)^3\)
  \end{enumerate}
\item
  The temperature at midnight is −6°C. It rises by 11°C by noon, then
  drops 15°C by the following midnight. What is the temperature at the
  end?
\item
  The lowest temperature ever recorded in Canada is −63°C (Snag, Yukon).
  The highest is +45°C (Lytton, BC). What is the range --- the
  difference between highest and lowest?
\item
  A student owes three friends £8 each. They earn £40 doing odd jobs.
  After paying back all three debts, how much money do they have?
\item
  A share price opens at £14.50. Changes over five days: −£0.80, +£1.20,
  −£2.40, +£0.35, −£0.90. What is the closing price on day 5?
\end{enumerate}

\chapter{Fractions and decimals}\label{fractions-and-decimals}

Arithmetic with parts of wholes

\hfill\break

A recipe needs \(\frac{3}{4}\) of a cup of flour and you want to make
half the recipe. How much flour do you use?

Four friends split a restaurant bill of £37.60 equally. What does each
person pay?

You run 2.5 km on Monday and 1.75 km on Wednesday. How far have you run
in total?

These are all fraction and decimal problems. The numbers sit between the
integers --- and the arithmetic for handling them is a small set of
rules you can learn once and use everywhere.

\section{What the notation is
saying}\label{what-the-notation-is-saying-4}

The fraction \(\frac{a}{b}\) means ``\(a\) divided by \(b\).'' The top
number is the \textbf{numerator} (how many parts you have); the bottom
number is the \textbf{denominator} (how many equal parts make the
whole).

Decimals are fractions whose denominators are powers of ten:
\[0.3 = \frac{3}{10}, \quad 0.47 = \frac{47}{100}, \quad 0.006 = \frac{6}{1000}\]

They are interchangeable. \(\frac{3}{4} = 0.75\).
\(\frac{1}{3} = 0.\overline{3}\) (the bar means the digit repeats
forever).

\textbf{Equivalent fractions:} \(\frac{2}{4} = \frac{1}{2}\) because you
can divide numerator and denominator by 2. A fraction is in
\textbf{lowest terms} when numerator and denominator share no common
factor other than 1.

\begin{center}\rule{0.5\linewidth}{0.5pt}\end{center}

\emph{Use the controls below to explore what a fraction looks like ---
as a divided rectangle and as a position on the number line.}

\begin{Shaded}
\begin{Highlighting}[]
\NormalTok{//| echo: false}
\NormalTok{viewof frac1\_num = Inputs.range([1, 12], \{}
\NormalTok{  value: 3,}
\NormalTok{  step: 1,}
\NormalTok{  label: "Numerator (top)"}
\NormalTok{\})}
\end{Highlighting}
\end{Shaded}

\begin{Shaded}
\begin{Highlighting}[]
\NormalTok{//| echo: false}
\NormalTok{viewof frac1\_den\_raw = Inputs.range([1, 12], \{}
\NormalTok{  value: 4,}
\NormalTok{  step: 1,}
\NormalTok{  label: "Denominator (bottom)"}
\NormalTok{\})}
\end{Highlighting}
\end{Shaded}

\begin{Shaded}
\begin{Highlighting}[]
\NormalTok{//| echo: false}
\NormalTok{\{}
\NormalTok{  // Enforce denominator \textgreater{}= numerator by silently clamping}
\NormalTok{  const num = frac1\_num;}
\NormalTok{  const den = Math.max(frac1\_den\_raw, num);}

\NormalTok{  // {-}{-}{-} GCD helper for display in lowest terms {-}{-}{-}}
\NormalTok{  function gcd(a, b) \{ return b === 0 ? a : gcd(b, a \% b); \}}
\NormalTok{  const g = gcd(num, den);}
\NormalTok{  const rn = num / g;}
\NormalTok{  const rd = den / g;}
\NormalTok{  const decimal = num / den;}
\NormalTok{  const lowestLabel = rn === rd}
\NormalTok{    ? \textasciigrave{}$\{rn\} (= 1)\textasciigrave{}}
\NormalTok{    : rd === 1}
\NormalTok{      ? \textasciigrave{}$\{rn\}\textasciigrave{}}
\NormalTok{      : \textasciigrave{}$\{rn\}/$\{rd\}\textasciigrave{};}

\NormalTok{  const W = 560, H\_rect = 80, H\_line = 70;}
\NormalTok{  const PAD = 40;}

\NormalTok{  const svg = d3.create("svg")}
\NormalTok{    .attr("viewBox", \textasciigrave{}0 0 $\{W\} $\{H\_rect + H\_line + 20\}\textasciigrave{})}
\NormalTok{    .attr("width", "100\%")}
\NormalTok{    .attr("style", "max{-}width:560px; display:block; margin:0 auto; font{-}family:sans{-}serif;");}

\NormalTok{  // ── Area model ──────────────────────────────────────────────}
\NormalTok{  const rectW = W {-} 2 * PAD;}
\NormalTok{  const partW = rectW / den;}

\NormalTok{  const rectG = svg.append("g").attr("transform", \textasciigrave{}translate($\{PAD\}, 10)\textasciigrave{});}

\NormalTok{  for (let i = 0; i \textless{} den; i++) \{}
\NormalTok{    rectG.append("rect")}
\NormalTok{      .attr("x", i * partW)}
\NormalTok{      .attr("y", 0)}
\NormalTok{      .attr("width", partW {-} 1)}
\NormalTok{      .attr("height", H\_rect)}
\NormalTok{      .attr("fill", i \textless{} num ? "\#2563eb" : "\#e5e7eb")}
\NormalTok{      .attr("stroke", "\#94a3b8")}
\NormalTok{      .attr("stroke{-}width", 0.5)}
\NormalTok{      .attr("rx", 2);}
\NormalTok{  \}}

\NormalTok{  // Fraction label centred above the shaded region}
\NormalTok{  rectG.append("text")}
\NormalTok{    .attr("x", (num * partW) / 2)}
\NormalTok{    .attr("y", {-}4)}
\NormalTok{    .attr("text{-}anchor", "middle")}
\NormalTok{    .attr("font{-}size", 13)}
\NormalTok{    .attr("fill", "\#1e3a5f")}
\NormalTok{    .text(\textasciigrave{}$\{num\}/$\{den\}$\{g \textgreater{} 1 ? \textasciigrave{} = $\{lowestLabel\}\textasciigrave{} : ""\} = $\{decimal.toFixed(6).replace(/\textbackslash{}.?0+$/, "") || "0"\}\textasciigrave{});}

\NormalTok{  // Shaded count label inside rectangle}
\NormalTok{  if (num \textgreater{} 0) \{}
\NormalTok{    rectG.append("text")}
\NormalTok{      .attr("x", (num * partW) / 2)}
\NormalTok{      .attr("y", H\_rect / 2 + 5)}
\NormalTok{      .attr("text{-}anchor", "middle")}
\NormalTok{      .attr("font{-}size", Math.min(14, partW * 0.7))}
\NormalTok{      .attr("fill", "\#ffffff")}
\NormalTok{      .attr("font{-}weight", "bold")}
\NormalTok{      .text(num === den ? "whole" : \textasciigrave{}$\{num\} part$\{num \textgreater{} 1 ? "s" : ""\}\textasciigrave{});}
\NormalTok{  \}}

\NormalTok{  // Unshaded count label}
\NormalTok{  if (num \textless{} den) \{}
\NormalTok{    rectG.append("text")}
\NormalTok{      .attr("x", num * partW + ((den {-} num) * partW) / 2)}
\NormalTok{      .attr("y", H\_rect / 2 + 5)}
\NormalTok{      .attr("text{-}anchor", "middle")}
\NormalTok{      .attr("font{-}size", Math.min(13, partW * 0.65))}
\NormalTok{      .attr("fill", "\#64748b")}
\NormalTok{      .text(\textasciigrave{}$\{den {-} num\} part$\{den {-} num \textgreater{} 1 ? "s" : ""\}\textasciigrave{});}
\NormalTok{  \}}

\NormalTok{  // Area model axis label}
\NormalTok{  rectG.append("text")}
\NormalTok{    .attr("x", rectW / 2)}
\NormalTok{    .attr("y", H\_rect + 14)}
\NormalTok{    .attr("text{-}anchor", "middle")}
\NormalTok{    .attr("font{-}size", 11)}
\NormalTok{    .attr("fill", "\#64748b")}
\NormalTok{    .text(\textasciigrave{}Rectangle divided into $\{den\} equal part$\{den \textgreater{} 1 ? "s" : ""\} — $\{num\} shaded\textasciigrave{});}

\NormalTok{  // ── Number line ──────────────────────────────────────────────}
\NormalTok{  const lineY = H\_rect + 20;}
\NormalTok{  const lineG = svg.append("g").attr("transform", \textasciigrave{}translate($\{PAD\}, $\{lineY + 10\})\textasciigrave{});}

\NormalTok{  // Track line}
\NormalTok{  lineG.append("line")}
\NormalTok{    .attr("x1", 0).attr("y1", 20)}
\NormalTok{    .attr("x2", rectW).attr("y2", 20)}
\NormalTok{    .attr("stroke", "\#94a3b8").attr("stroke{-}width", 2);}

\NormalTok{  // Tick at 0}
\NormalTok{  lineG.append("line")}
\NormalTok{    .attr("x1", 0).attr("y1", 13).attr("x2", 0).attr("y2", 27)}
\NormalTok{    .attr("stroke", "\#374151").attr("stroke{-}width", 2);}
\NormalTok{  lineG.append("text")}
\NormalTok{    .attr("x", 0).attr("y", 40)}
\NormalTok{    .attr("text{-}anchor", "middle").attr("font{-}size", 12).attr("fill", "\#374151")}
\NormalTok{    .text("0");}

\NormalTok{  // Tick at 1}
\NormalTok{  lineG.append("line")}
\NormalTok{    .attr("x1", rectW).attr("y1", 13).attr("x2", rectW).attr("y2", 27)}
\NormalTok{    .attr("stroke", "\#374151").attr("stroke{-}width", 2);}
\NormalTok{  lineG.append("text")}
\NormalTok{    .attr("x", rectW).attr("y", 40)}
\NormalTok{    .attr("text{-}anchor", "middle").attr("font{-}size", 12).attr("fill", "\#374151")}
\NormalTok{    .text("1");}

\NormalTok{  // Fraction position marker}
\NormalTok{  const markerX = decimal * rectW;}
\NormalTok{  lineG.append("circle")}
\NormalTok{    .attr("cx", markerX).attr("cy", 20)}
\NormalTok{    .attr("r", 7)}
\NormalTok{    .attr("fill", "\#2563eb")}
\NormalTok{    .attr("stroke", "\#1e3a5f")}
\NormalTok{    .attr("stroke{-}width", 1.5);}

\NormalTok{  // Value label above marker}
\NormalTok{  lineG.append("text")}
\NormalTok{    .attr("x", markerX)}
\NormalTok{    .attr("y", 10)}
\NormalTok{    .attr("text{-}anchor", "middle")}
\NormalTok{    .attr("font{-}size", 12)}
\NormalTok{    .attr("fill", "\#1e3a5f")}
\NormalTok{    .attr("font{-}weight", "bold")}
\NormalTok{    .text(decimal.toFixed(4).replace(/\textbackslash{}.?0+$/, "") || "0");}

\NormalTok{  // Axis label}
\NormalTok{  lineG.append("text")}
\NormalTok{    .attr("x", rectW / 2).attr("y", 55)}
\NormalTok{    .attr("text{-}anchor", "middle").attr("font{-}size", 11).attr("fill", "\#64748b")}
\NormalTok{    .text("Number line from 0 to 1 — dot shows the same value");}

\NormalTok{  return svg.node();}
\NormalTok{\}}
\end{Highlighting}
\end{Shaded}

\begin{tcolorbox}[enhanced jigsaw, opacityback=0, colback=white, breakable, colframe=quarto-callout-note-color-frame, arc=.35mm, left=2mm, toprule=.15mm, leftrule=.75mm, bottomrule=.15mm, rightrule=.15mm]
\begin{minipage}[t]{5.5mm}
\textcolor{quarto-callout-note-color}{\faInfo}
\end{minipage}%
\begin{minipage}[t]{\textwidth - 5.5mm}

If the denominator input goes below the numerator, the interactive
clamps it to keep the fraction valid (≤ 1). Try setting numerator and
denominator to the same value to see what ``1 whole'' looks like in both
views.

\end{minipage}%
\end{tcolorbox}

\begin{center}\rule{0.5\linewidth}{0.5pt}\end{center}

\section{The method}\label{the-method-4}

\textbf{Multiplying fractions}

Multiply numerators together; multiply denominators together --- that's
it: \[\frac{a}{b} \times \frac{c}{d} = \frac{ac}{bd}\]

Simplify before or after. Before is usually easier.

\textbf{Dividing fractions}

Flip the second fraction and multiply --- this is called multiplying by
the reciprocal:
\[\frac{a}{b} \div \frac{c}{d} = \frac{a}{b} \times \frac{d}{c} = \frac{ad}{bc}\]

``Flip and multiply'' works because
\(\frac{c}{d} \times \frac{d}{c} = 1\): the two fractions cancel each
other out, leaving you with a clean multiplication.

\textbf{Adding and subtracting fractions}

Mixed numbers: convert first. Multiply the whole number by the
denominator and add the numerator:
\(1\frac{1}{2} = \frac{(1 \times 2) + 1}{2} = \frac{3}{2}\).

You need a common denominator before you can add. The simplest route:
multiply both denominators to get one that works for both.
\[\frac{a}{b} + \frac{c}{d} = \frac{ad}{bd} + \frac{bc}{bd} = \frac{ad + bc}{bd}\]

Simplify the result.

\textbf{Converting between fractions and decimals}

Fraction → decimal: divide numerator by denominator. Decimal → fraction:
write the decimal over its place-value denominator, then simplify.

\begin{tcolorbox}[enhanced jigsaw, opacityback=0, colback=white, breakable, colframe=quarto-callout-note-color-frame, arc=.35mm, left=2mm, toprule=.15mm, leftrule=.75mm, bottomrule=.15mm, rightrule=.15mm]
\begin{minipage}[t]{5.5mm}
\textcolor{quarto-callout-note-color}{\faInfo}
\end{minipage}%
\begin{minipage}[t]{\textwidth - 5.5mm}

\textbf{Why this works}

A fraction \(\frac{a}{b}\) is a number. The rules for arithmetic with
fractions follow from the same rules that govern whole numbers ---
distribution, commutativity, and the fact that multiplying by 1 changes
nothing. To add \(\frac{a}{b} + \frac{c}{d}\), rewrite the first
fraction as \(\frac{a}{b} \times \frac{d}{d} = \frac{ad}{bd}\) and the
second as \(\frac{c}{d} \times \frac{b}{b} = \frac{bc}{bd}\). Both
multiplications are by 1, so neither fraction's value changes --- only
its representation. The result is \(\frac{ad + bc}{bd}\). Division by a
fraction is multiplication by its reciprocal because
\(\frac{a}{b} \times \frac{b}{a} = 1\) --- they are inverses of each
other.

\end{minipage}%
\end{tcolorbox}

\begin{center}\rule{0.5\linewidth}{0.5pt}\end{center}

\emph{Use the calculator below to check any fraction arithmetic. Choose
two fractions and an operation; the working is shown step by step.}

\begin{Shaded}
\begin{Highlighting}[]
\NormalTok{//| echo: false}
\NormalTok{viewof op\_a\_num = Inputs.range([1, 12], \{ value: 3, step: 1, label: "Fraction A — numerator" \})}
\end{Highlighting}
\end{Shaded}

\begin{Shaded}
\begin{Highlighting}[]
\NormalTok{//| echo: false}
\NormalTok{viewof op\_a\_den = Inputs.range([1, 12], \{ value: 4, step: 1, label: "Fraction A — denominator" \})}
\end{Highlighting}
\end{Shaded}

\begin{Shaded}
\begin{Highlighting}[]
\NormalTok{//| echo: false}
\NormalTok{viewof op\_b\_num = Inputs.range([1, 12], \{ value: 1, step: 1, label: "Fraction B — numerator" \})}
\end{Highlighting}
\end{Shaded}

\begin{Shaded}
\begin{Highlighting}[]
\NormalTok{//| echo: false}
\NormalTok{viewof op\_b\_den = Inputs.range([1, 12], \{ value: 2, step: 1, label: "Fraction B — denominator" \})}
\end{Highlighting}
\end{Shaded}

\begin{Shaded}
\begin{Highlighting}[]
\NormalTok{//| echo: false}
\NormalTok{viewof operation = Inputs.select(["×", "÷", "+", "−"], \{ label: "Operation", value: "×" \})}
\end{Highlighting}
\end{Shaded}

\begin{Shaded}
\begin{Highlighting}[]
\NormalTok{//| echo: false}
\NormalTok{\{}
\NormalTok{  function gcd(a, b) \{ return b === 0 ? a : gcd(b, a \% b); \}}
\NormalTok{  function simplify(n, d) \{}
\NormalTok{    if (d === 0) return [n, d];}
\NormalTok{    const g = gcd(Math.abs(n), Math.abs(d));}
\NormalTok{    return [n / g, d / g];}
\NormalTok{  \}}
\NormalTok{  function fracStr(n, d) \{}
\NormalTok{    if (d === 1) return \textasciigrave{}$\{n\}\textasciigrave{};}
\NormalTok{    return \textasciigrave{}$\{n\}/$\{d\}\textasciigrave{};}
\NormalTok{  \}}

\NormalTok{  const an = op\_a\_num, ad = op\_a\_den;}
\NormalTok{  const bn = op\_b\_num, bd = op\_b\_den;}

\NormalTok{  let steps = [];}
\NormalTok{  let rn, rd;}

\NormalTok{  if (operation === "×") \{}
\NormalTok{    rn = an * bn;}
\NormalTok{    rd = ad * bd;}
\NormalTok{    const [sn, sd] = simplify(rn, rd);}
\NormalTok{    steps = [}
\NormalTok{      \textasciigrave{}Multiply numerators: $\{an\} × $\{bn\} = $\{rn\}\textasciigrave{},}
\NormalTok{      \textasciigrave{}Multiply denominators: $\{ad\} × $\{bd\} = $\{rd\}\textasciigrave{},}
\NormalTok{      \textasciigrave{}Result: $\{rn\}/$\{rd\}\textasciigrave{},}
\NormalTok{      (sn !== rn || sd !== rd) ? \textasciigrave{}Simplified: $\{fracStr(sn, sd)\}\textasciigrave{} : \textasciigrave{}Already in lowest terms\textasciigrave{},}
\NormalTok{      \textasciigrave{}As a decimal: $\{(sn / sd).toFixed(6).replace(/\textbackslash{}.?0+$/, "") || "0"\}\textasciigrave{}}
\NormalTok{    ];}
\NormalTok{    [rn, rd] = [sn, sd];}

\NormalTok{  \} else if (operation === "÷") \{}
\NormalTok{    // Multiply by reciprocal of B}
\NormalTok{    const flipN = bd, flipD = bn;}
\NormalTok{    if (bn === 0) \{}
\NormalTok{      steps = ["Cannot divide by zero — set Fraction B numerator to at least 1."];}
\NormalTok{      rn = NaN; rd = 1;}
\NormalTok{    \} else \{}
\NormalTok{      const rawN = an * flipN;}
\NormalTok{      const rawD = ad * flipD;}
\NormalTok{      const [sn, sd] = simplify(rawN, rawD);}
\NormalTok{      steps = [}
\NormalTok{        \textasciigrave{}Flip Fraction B: $\{bn\}/$\{bd\} becomes $\{flipN\}/$\{flipD\}\textasciigrave{},}
\NormalTok{        \textasciigrave{}Multiply: $\{an\}/$\{ad\} × $\{flipN\}/$\{flipD\}\textasciigrave{},}
\NormalTok{        \textasciigrave{}Numerators: $\{an\} × $\{flipN\} = $\{rawN\}\textasciigrave{},}
\NormalTok{        \textasciigrave{}Denominators: $\{ad\} × $\{flipD\} = $\{rawD\}\textasciigrave{},}
\NormalTok{        \textasciigrave{}Result: $\{rawN\}/$\{rawD\}\textasciigrave{},}
\NormalTok{        (sn !== rawN || sd !== rawD) ? \textasciigrave{}Simplified: $\{fracStr(sn, sd)\}\textasciigrave{} : \textasciigrave{}Already in lowest terms\textasciigrave{},}
\NormalTok{        \textasciigrave{}As a decimal: $\{(sn / sd).toFixed(6).replace(/\textbackslash{}.?0+$/, "") || "0"\}\textasciigrave{}}
\NormalTok{      ];}
\NormalTok{      [rn, rd] = [sn, sd];}
\NormalTok{    \}}

\NormalTok{  \} else \{}
\NormalTok{    // Addition or subtraction — find common denominator}
\NormalTok{    const cd = ad * bd;  // common denominator (cross{-}multiply)}
\NormalTok{    const aN = an * bd;}
\NormalTok{    const bN = bn * ad;}
\NormalTok{    const rawN = operation === "+" ? aN + bN : aN {-} bN;}
\NormalTok{    const [sn, sd] = simplify(rawN, cd);}
\NormalTok{    steps = [}
\NormalTok{      \textasciigrave{}Common denominator: $\{ad\} × $\{bd\} = $\{cd\}\textasciigrave{},}
\NormalTok{      \textasciigrave{}Rewrite A: $\{an\}/$\{ad\} = $\{aN\}/$\{cd\}\textasciigrave{},}
\NormalTok{      \textasciigrave{}Rewrite B: $\{bn\}/$\{bd\} = $\{bN\}/$\{cd\}\textasciigrave{},}
\NormalTok{      operation === "+"}
\NormalTok{        ? \textasciigrave{}Add numerators: $\{aN\} + $\{bN\} = $\{rawN\}\textasciigrave{}}
\NormalTok{        : \textasciigrave{}Subtract numerators: $\{aN\} − $\{bN\} = $\{rawN\}\textasciigrave{},}
\NormalTok{      \textasciigrave{}Result: $\{rawN\}/$\{cd\}\textasciigrave{},}
\NormalTok{      (sn !== rawN || sd !== cd) ? \textasciigrave{}Simplified: $\{fracStr(sn, sd)\}\textasciigrave{} : \textasciigrave{}Already in lowest terms\textasciigrave{},}
\NormalTok{      \textasciigrave{}As a decimal: $\{isFinite(sn / sd) ? (sn / sd).toFixed(6).replace(/\textbackslash{}.?0+$/, "") || "0" : "undefined"\}\textasciigrave{}}
\NormalTok{    ];}
\NormalTok{    [rn, rd] = [sn, sd];}
\NormalTok{  \}}

\NormalTok{  // ── Build SVG output ─────────────────────────────────────────}
\NormalTok{  const W = 560, lineH = 24, topPad = 14, sidePad = 20;}
\NormalTok{  const H = topPad + steps.length * lineH + 60 + 20;}

\NormalTok{  const svg = d3.create("svg")}
\NormalTok{    .attr("viewBox", \textasciigrave{}0 0 $\{W\} $\{H\}\textasciigrave{})}
\NormalTok{    .attr("width", "100\%")}
\NormalTok{    .attr("style", "max{-}width:560px; display:block; margin:0 auto; font{-}family:sans{-}serif; background:\#f8fafc; border{-}radius:6px;");}

\NormalTok{  // Header: A op B =}
\NormalTok{  const headerText = \textasciigrave{}$\{an\}/$\{ad\}  $\{operation\}  $\{bn\}/$\{bd\}\textasciigrave{};}
\NormalTok{  svg.append("text")}
\NormalTok{    .attr("x", W / 2).attr("y", topPad + 16)}
\NormalTok{    .attr("text{-}anchor", "middle")}
\NormalTok{    .attr("font{-}size", 16).attr("font{-}weight", "bold").attr("fill", "\#1e3a5f")}
\NormalTok{    .text(headerText);}

\NormalTok{  // Step{-}by{-}step lines}
\NormalTok{  steps.filter(Boolean).forEach((s, i) =\textgreater{} \{}
\NormalTok{    const isLast = i === steps.filter(Boolean).length {-} 1;}
\NormalTok{    svg.append("text")}
\NormalTok{      .attr("x", sidePad).attr("y", topPad + 40 + i * lineH)}
\NormalTok{      .attr("font{-}size", 13)}
\NormalTok{      .attr("fill", isLast ? "\#166534" : "\#374151")}
\NormalTok{      .attr("font{-}weight", isLast ? "bold" : "normal")}
\NormalTok{      .text(\textasciigrave{}$\{i + 1\}. $\{s\}\textasciigrave{});}
\NormalTok{  \});}

\NormalTok{  // Result box (only when result is valid)}
\NormalTok{  if (isFinite(rn / rd)) \{}
\NormalTok{    const boxY = topPad + 40 + steps.filter(Boolean).length * lineH + 4;}
\NormalTok{    svg.append("rect")}
\NormalTok{      .attr("x", sidePad {-} 6).attr("y", boxY {-} 18)}
\NormalTok{      .attr("width", W {-} 2 * sidePad + 12).attr("height", 36)}
\NormalTok{      .attr("fill", "\#dbeafe").attr("rx", 4).attr("stroke", "\#93c5fd");}

\NormalTok{    // Area model for + and −: two coloured bars side by side in proportion}
\NormalTok{    if (operation === "+" || operation === "−") \{}
\NormalTok{      // Show A and B as proportional bars, stacked / subtracted}
\NormalTok{      const barW = W {-} 2 * sidePad {-} 150;}
\NormalTok{      const barH = 18;}
\NormalTok{      const barX = sidePad + 145;}
\NormalTok{      const barY = boxY {-} 16;}
\NormalTok{      // A bar}
\NormalTok{      const aFrac = Math.min(an / ad, 1);}
\NormalTok{      const bFrac = Math.min(bn / bd, 1);}
\NormalTok{      svg.append("rect")}
\NormalTok{        .attr("x", barX).attr("y", barY)}
\NormalTok{        .attr("width", aFrac * barW).attr("height", barH)}
\NormalTok{        .attr("fill", "\#2563eb").attr("opacity", 0.8).attr("rx", 2);}
\NormalTok{      if (operation === "+") \{}
\NormalTok{        svg.append("rect")}
\NormalTok{          .attr("x", barX + aFrac * barW).attr("y", barY)}
\NormalTok{          .attr("width", bFrac * barW).attr("height", barH)}
\NormalTok{          .attr("fill", "\#16a34a").attr("opacity", 0.8).attr("rx", 2);}
\NormalTok{      \} else \{}
\NormalTok{        // Subtraction: show B removed from A}
\NormalTok{        const removeW = Math.min(bFrac, aFrac) * barW;}
\NormalTok{        svg.append("rect")}
\NormalTok{          .attr("x", barX + aFrac * barW {-} removeW).attr("y", barY)}
\NormalTok{          .attr("width", removeW).attr("height", barH)}
\NormalTok{          .attr("fill", "\#dc2626").attr("opacity", 0.5).attr("rx", 2)}
\NormalTok{          .attr("stroke", "\#dc2626").attr("stroke{-}dasharray", "3,2");}
\NormalTok{      \}}
\NormalTok{      svg.append("rect")}
\NormalTok{        .attr("x", barX).attr("y", barY)}
\NormalTok{        .attr("width", barW).attr("height", barH)}
\NormalTok{        .attr("fill", "none").attr("stroke", "\#94a3b8").attr("rx", 2);}
\NormalTok{      svg.append("text")}
\NormalTok{        .attr("x", barX {-} 4).attr("y", barY + 13)}
\NormalTok{        .attr("text{-}anchor", "end").attr("font{-}size", 10).attr("fill", "\#64748b")}
\NormalTok{        .text("0");}
\NormalTok{      svg.append("text")}
\NormalTok{        .attr("x", barX + barW + 4).attr("y", barY + 13)}
\NormalTok{        .attr("text{-}anchor", "start").attr("font{-}size", 10).attr("fill", "\#64748b")}
\NormalTok{        .text("1");}
\NormalTok{    \}}
\NormalTok{  \}}

\NormalTok{  return svg.node();}
\NormalTok{\}}
\end{Highlighting}
\end{Shaded}

\begin{center}\rule{0.5\linewidth}{0.5pt}\end{center}

\section{Worked examples}\label{worked-examples-4}

\textbf{Example 1 --- Halving a recipe.} A recipe needs \(\frac{3}{4}\)
cup of flour, \(1\frac{1}{2}\) cups of oats, and \(\frac{5}{8}\) cup of
sugar. You want to make half the recipe. What quantities do you need?

Multiply each ingredient by \(\frac{1}{2}\).

Convert the mixed number first: \(1\frac{1}{2} = \frac{3}{2}\).

\[\frac{3}{4} \times \frac{1}{2} = \frac{3}{8} \text{ cup flour}\]

\[\frac{3}{2} \times \frac{1}{2} = \frac{3}{4} \text{ cup oats}\]

\[\frac{5}{8} \times \frac{1}{2} = \frac{5}{16} \text{ cup sugar}\]

\begin{center}\rule{0.5\linewidth}{0.5pt}\end{center}

\textbf{Example 2 --- Splitting costs.} Three friends split the cost of
a camping trip. They agree to share costs in proportion to how many
nights each stays: one friend stays \(\frac{1}{3}\) of the nights,
another \(\frac{5}{12}\), and the third \(\frac{1}{4}\). The total cost
is £120. What does each pay?

Check the fractions sum to 1 first --- to make sure everyone's share is
accounted for:
\[\frac{4}{12} + \frac{5}{12} + \frac{3}{12} = \frac{12}{12} = 1 \checkmark\]

Friend A: \(\frac{1}{3} \times 120 = £40\)

Friend B: \(\frac{5}{12} \times 120 = £50\)

Friend C: \(\frac{1}{4} \times 120 = £30\)

Total: £120. Check.

\begin{center}\rule{0.5\linewidth}{0.5pt}\end{center}

\textbf{Example 3 --- Running distances.} You run 2.5 km on Monday, 1.75
km on Wednesday, and \(\frac{7}{4}\) km on Friday. What is your total
distance for the week?

Convert \(\frac{7}{4}\) to a decimal so all three are in the same form:
\(\frac{7}{4} = 1.75\) km.

\[2.5 + 1.75 + 1.75 = 6.0 \text{ km}\]

\begin{center}\rule{0.5\linewidth}{0.5pt}\end{center}

\textbf{Example 4 --- What fraction is left?} A bag of rice is 0.38 used
up. What fraction remains? Express it as a decimal and as a fraction in
lowest terms.

Remaining: \(1 - 0.38 = 0.62\).

As a fraction: \(0.62 = \frac{62}{100} = \frac{31}{50}\).

\(\frac{31}{50}\) of the bag is left.

\section{Where this goes}\label{where-this-goes-4}

Fraction arithmetic is the foundation of \textbf{algebra}. When you
solve \(\frac{2x}{3} = 8\), you multiply both sides by \(\frac{3}{2}\)
--- a fraction operation. Every algebraic fraction, rational expression,
and limit in calculus is built on the rules in this chapter.

In calculus, the derivative is defined as a limit of a fraction:
\(\frac{\Delta y}{\Delta x}\) as \(\Delta x \to 0\). That fraction needs
to be understood, not just computed. The conceptual step from
``fraction'' to ``rate of change'' is not large; it requires only that
you see the notation clearly.

\begin{tcolorbox}[enhanced jigsaw, opacityback=0, colback=white, breakable, colframe=quarto-callout-tip-color-frame, arc=.35mm, left=2mm, toprule=.15mm, leftrule=.75mm, bottomrule=.15mm, rightrule=.15mm]
\begin{minipage}[t]{5.5mm}
\textcolor{quarto-callout-tip-color}{\faLightbulb}
\end{minipage}%
\begin{minipage}[t]{\textwidth - 5.5mm}

\textbf{Where this shows up}

\begin{itemize}
\tightlist
\item
  Any time you split something equally --- a bill, a recipe, a distance
  --- you're doing fraction arithmetic.
\item
  A structural engineer applies a safety factor of \(\frac{3}{2}\) (or
  1.5) to every load --- multiplying by a fraction.
\item
  A financial model apportions costs, equity, and interest using decimal
  fractions at every step.
\item
  A GIS analyst computing land-cover proportions is adding fractions
  that must sum to 1.
\end{itemize}

The arithmetic is identical. The precision matters.

\end{minipage}%
\end{tcolorbox}

\section{Exercises}\label{exercises-4}

\begin{enumerate}
\def\labelenumi{\arabic{enumi}.}
\item
  A metal bar is \(4\frac{3}{8}\) m long. A \(1\frac{5}{8}\) m section
  is cut off. How much remains?
\item
  A recipe calls for \(\frac{2}{3}\) cup of milk. You want to make
  \(2\frac{1}{2}\) times the recipe. How much milk do you need?
\item
  A sports pitch is 0.75 km long and 0.48 km wide. What is its area in
  km²? Express the answer as a decimal and as a fraction in lowest
  terms.
\item
  Three taps fill a tank. Each supplies \(\frac{3}{8}\) of the tank's
  capacity per hour. A drain removes \(\frac{5}{8}\) of the tank's
  capacity per hour. Is the tank filling or emptying? At what net
  fraction per hour?
\item
  A bag of trail mix is analysed: 0.44 nuts, 0.29 dried fruit, 0.18
  chocolate chips, and the rest is seeds. What fraction is seeds?
  Express as a decimal and a simplified fraction.
\item
  You have a 500 mL bottle of juice concentrate. The instructions say to
  mix it at \(\frac{1}{4}\) strength. What fraction of the final drink
  will be concentrate? How many mL of concentrate goes into a 250 mL
  glass?
\item
  Convert each to a decimal, then order from least to greatest:
  \(\frac{7}{12}\), \(\frac{3}{5}\), \(0.58\), \(\frac{11}{20}\),
  \(0.\overline{6}\).
\end{enumerate}

\chapter{Statistics and probability}\label{statistics-and-probability}

Reasoning honestly from limited data

\hfill\break

A football team scores 1, 3, 0, 2, 1, 4, 1 goals across seven matches.
What is their typical score?

A weather forecast says 70\% chance of rain. What does that actually
mean?

Two students both average 65\% on their tests, but one scores between
60--70\% every time and the other swings from 20\% to 100\%. Are they
the same?

These questions all need the same toolkit: averages, spread, and
probability. The numbers tell you something --- but only if you know
which number to look at.

\section{What the notation is
saying}\label{what-the-notation-is-saying-5}

\textbf{The mean} \(\bar{x}\) is the sum of all values divided by the
number of values. Given data \(x_1, x_2, \ldots, x_n\):

\[\bar{x} = \frac{x_1 + x_2 + \cdots + x_n}{n} = \frac{\sum x_i}{n}\]

The \(\Sigma\) (capital sigma) means ``sum all of the following.'' It is
notation for a loop: add every term \(x_i\) for \(i = 1\) to \(n\).

\textbf{The median} is the middle value when the data are ordered. If
there are two middle values, take their mean. The median is less
sensitive to extreme values than the mean.

\textbf{The range} is maximum minus minimum. It measures total spread
but is heavily influenced by outliers.

\textbf{Probability} of event \(A\) is:

\[P(A) = \frac{\text{number of favourable outcomes}}{\text{total number of equally likely outcomes}}\]

\(P(A) = 0\): impossible. \(P(A) = 1\): certain. All probabilities for a
complete set of outcomes sum to 1.

\textbf{Relative frequency}: when you can't count outcomes
theoretically, estimate probability from data:
\[P(A) \approx \frac{\text{number of times A occurred}}{\text{total trials}}\]

\section{The method}\label{the-method-5}

\textbf{Computing the mean}

\begin{enumerate}
\def\labelenumi{\arabic{enumi}.}
\tightlist
\item
  Add all values --- to get the total.
\item
  Divide by the count --- to find the fair share per value.
\end{enumerate}

\textbf{Computing the median}

\begin{enumerate}
\def\labelenumi{\arabic{enumi}.}
\tightlist
\item
  Order the values from least to greatest.
\item
  If \(n\) is odd, the median is the value at position
  \(\frac{n+1}{2}\).
\item
  If \(n\) is even, average the values at positions \(\frac{n}{2}\) and
  \(\frac{n}{2}+1\).
\end{enumerate}

\textbf{Computing the range and IQR}

Range: \(\max - \min\).

IQR (interquartile range): split the ordered data in half. Lower
quartile Q1 = median of the lower half. Upper quartile Q3 = median of
the upper half. \(\text{IQR} = Q3 - Q1\).

The IQR describes the spread of the middle 50\% of the data --- it
ignores the extremes at both ends.

\textbf{Computing probability from equally likely outcomes}

\begin{enumerate}
\def\labelenumi{\arabic{enumi}.}
\tightlist
\item
  Count total possible outcomes.
\item
  Count outcomes that match event \(A\).
\item
  Divide.
\end{enumerate}

\begin{tcolorbox}[enhanced jigsaw, opacityback=0, colback=white, breakable, colframe=quarto-callout-note-color-frame, arc=.35mm, left=2mm, toprule=.15mm, leftrule=.75mm, bottomrule=.15mm, rightrule=.15mm]
\begin{minipage}[t]{5.5mm}
\textcolor{quarto-callout-note-color}{\faInfo}
\end{minipage}%
\begin{minipage}[t]{\textwidth - 5.5mm}

\textbf{Why this works}

The mean is the ``balance point'' of a data set --- if you placed equal
weights at each value on a number line, the mean is where the line would
balance. The median is the ``middle'' --- it splits the data 50/50. They
give the same answer for symmetric data; they diverge when the data is
skewed. The median is more robust: one extreme outlier shifts the mean
but not the median.

Probability is defined between 0 and 1 because it is a fraction of a
complete set of outcomes. When probabilities sum to more than 1, you
have double-counted. When they sum to less than 1, you have missed some
outcomes.

\end{minipage}%
\end{tcolorbox}

\begin{center}\rule{0.5\linewidth}{0.5pt}\end{center}

\emph{Edit any of the seven scores below. Watch the mean, median, range,
and IQR update instantly --- and see the difference between mean and
median when you push one value to an extreme.}

\begin{Shaded}
\begin{Highlighting}[]
\NormalTok{//| echo: false}
\NormalTok{viewof v1 = Inputs.number(\{ label: "Score 1", value: 62, min: 0, max: 100, step: 1 \})}
\end{Highlighting}
\end{Shaded}

\begin{Shaded}
\begin{Highlighting}[]
\NormalTok{//| echo: false}
\NormalTok{viewof v2 = Inputs.number(\{ label: "Score 2", value: 71, min: 0, max: 100, step: 1 \})}
\end{Highlighting}
\end{Shaded}

\begin{Shaded}
\begin{Highlighting}[]
\NormalTok{//| echo: false}
\NormalTok{viewof v3 = Inputs.number(\{ label: "Score 3", value: 58, min: 0, max: 100, step: 1 \})}
\end{Highlighting}
\end{Shaded}

\begin{Shaded}
\begin{Highlighting}[]
\NormalTok{//| echo: false}
\NormalTok{viewof v4 = Inputs.number(\{ label: "Score 4", value: 74, min: 0, max: 100, step: 1 \})}
\end{Highlighting}
\end{Shaded}

\begin{Shaded}
\begin{Highlighting}[]
\NormalTok{//| echo: false}
\NormalTok{viewof v5 = Inputs.number(\{ label: "Score 5", value: 66, min: 0, max: 100, step: 1 \})}
\end{Highlighting}
\end{Shaded}

\begin{Shaded}
\begin{Highlighting}[]
\NormalTok{//| echo: false}
\NormalTok{viewof v6 = Inputs.number(\{ label: "Score 6", value: 70, min: 0, max: 100, step: 1 \})}
\end{Highlighting}
\end{Shaded}

\begin{Shaded}
\begin{Highlighting}[]
\NormalTok{//| echo: false}
\NormalTok{viewof v7 = Inputs.number(\{ label: "Score 7", value: 63, min: 0, max: 100, step: 1 \})}
\end{Highlighting}
\end{Shaded}

\begin{Shaded}
\begin{Highlighting}[]
\NormalTok{//| echo: false}
\NormalTok{\{}
\NormalTok{  // ── Helpers ──────────────────────────────────────────────────}
\NormalTok{  function medianOf(arr) \{}
\NormalTok{    const s = [...arr].sort((a, b) =\textgreater{} a {-} b);}
\NormalTok{    const mid = Math.floor(s.length / 2);}
\NormalTok{    return s.length \% 2 === 1 ? s[mid] : (s[mid {-} 1] + s[mid]) / 2;}
\NormalTok{  \}}
\NormalTok{  function quartiles(arr) \{}
\NormalTok{    const s = [...arr].sort((a, b) =\textgreater{} a {-} b);}
\NormalTok{    const n = s.length;}
\NormalTok{    const lower = s.slice(0, Math.floor(n / 2));}
\NormalTok{    const upper = n \% 2 === 0 ? s.slice(n / 2) : s.slice(Math.ceil(n / 2));}
\NormalTok{    return \{ q1: medianOf(lower), q3: medianOf(upper) \};}
\NormalTok{  \}}

\NormalTok{  // ── Data ─────────────────────────────────────────────────────}
\NormalTok{  const raw = [v1, v2, v3, v4, v5, v6, v7].map(x =\textgreater{} isFinite(x) ? x : 0);}
\NormalTok{  const sorted = [...raw].sort((a, b) =\textgreater{} a {-} b);}
\NormalTok{  const n = raw.length;}
\NormalTok{  const sum = raw.reduce((a, b) =\textgreater{} a + b, 0);}
\NormalTok{  const mean = sum / n;}
\NormalTok{  const med = medianOf(raw);}
\NormalTok{  const range\_ = sorted[n {-} 1] {-} sorted[0];}
\NormalTok{  const \{ q1, q3 \} = quartiles(raw);}
\NormalTok{  const iqr = q3 {-} q1;}
\NormalTok{  const medIdx = Math.floor(n / 2); // index in sorted array (0{-}based)}

\NormalTok{  // ── Layout ───────────────────────────────────────────────────}
\NormalTok{  const W = 560, PAD = 40, DOT\_AREA\_H = 120, STATS\_H = 110;}
\NormalTok{  const H = DOT\_AREA\_H + STATS\_H + 20;}

\NormalTok{  const svg = d3.create("svg")}
\NormalTok{    .attr("viewBox", \textasciigrave{}0 0 $\{W\} $\{H\}\textasciigrave{})}
\NormalTok{    .attr("width", "100\%")}
\NormalTok{    .attr("style", "max{-}width:560px; display:block; margin:0 auto; font{-}family:sans{-}serif;");}

\NormalTok{  // ── Dot plot ─────────────────────────────────────────────────}
\NormalTok{  const plotW = W {-} 2 * PAD;}
\NormalTok{  const xMin = 0, xMax = 100;}
\NormalTok{  const xScale = v =\textgreater{} PAD + (v {-} xMin) / (xMax {-} xMin) * plotW;}
\NormalTok{  const axisY = 70;}

\NormalTok{  // Axis line}
\NormalTok{  svg.append("line")}
\NormalTok{    .attr("x1", PAD).attr("y1", axisY)}
\NormalTok{    .attr("x2", W {-} PAD).attr("y2", axisY)}
\NormalTok{    .attr("stroke", "\#94a3b8").attr("stroke{-}width", 1.5);}

\NormalTok{  // Tick marks every 10}
\NormalTok{  for (let t = 0; t \textless{}= 100; t += 10) \{}
\NormalTok{    const x = xScale(t);}
\NormalTok{    svg.append("line")}
\NormalTok{      .attr("x1", x).attr("y1", axisY {-} 4)}
\NormalTok{      .attr("x2", x).attr("y2", axisY + 4)}
\NormalTok{      .attr("stroke", "\#94a3b8").attr("stroke{-}width", 1);}
\NormalTok{    svg.append("text")}
\NormalTok{      .attr("x", x).attr("y", axisY + 15)}
\NormalTok{      .attr("text{-}anchor", "middle").attr("font{-}size", 10).attr("fill", "\#64748b")}
\NormalTok{      .text(t);}
\NormalTok{  \}}

\NormalTok{  // Axis label}
\NormalTok{  svg.append("text")}
\NormalTok{    .attr("x", W / 2).attr("y", axisY + 28)}
\NormalTok{    .attr("text{-}anchor", "middle").attr("font{-}size", 11).attr("fill", "\#64748b")}
\NormalTok{    .text("Test score (0–100)");}

\NormalTok{  // Dots — stack duplicates}
\NormalTok{  const stackCount = \{\};}
\NormalTok{  raw.forEach(v =\textgreater{} \{ stackCount[v] = (stackCount[v] || 0) + 1; \});}
\NormalTok{  const placed = \{\};}
\NormalTok{  raw.forEach(v =\textgreater{} \{}
\NormalTok{    placed[v] = (placed[v] || 0);}
\NormalTok{    const cx = xScale(v);}
\NormalTok{    const cy = axisY {-} 12 {-} placed[v] * 16;}
\NormalTok{    svg.append("circle")}
\NormalTok{      .attr("cx", cx).attr("cy", cy).attr("r", 7)}
\NormalTok{      .attr("fill", "\#475569").attr("stroke", "\#1e293b").attr("stroke{-}width", 1);}
\NormalTok{    placed[v]++;}
\NormalTok{  \});}

\NormalTok{  // Mean vertical line}
\NormalTok{  svg.append("line")}
\NormalTok{    .attr("x1", xScale(mean)).attr("y1", 8)}
\NormalTok{    .attr("x2", xScale(mean)).attr("y2", axisY)}
\NormalTok{    .attr("stroke", "\#2563eb").attr("stroke{-}width", 2)}
\NormalTok{    .attr("stroke{-}dasharray", "4,2");}
\NormalTok{  svg.append("text")}
\NormalTok{    .attr("x", xScale(mean)).attr("y", 6)}
\NormalTok{    .attr("text{-}anchor", "middle").attr("font{-}size", 11).attr("fill", "\#2563eb")}
\NormalTok{    .attr("font{-}weight", "bold")}
\NormalTok{    .text(\textasciigrave{}Mean $\{mean.toFixed(1)\}\textasciigrave{});}

\NormalTok{  // Median vertical line}
\NormalTok{  const medOffset = Math.abs(xScale(med) {-} xScale(mean)) \textless{} 20 ? 12 : 0;}
\NormalTok{  svg.append("line")}
\NormalTok{    .attr("x1", xScale(med)).attr("y1", 8)}
\NormalTok{    .attr("x2", xScale(med)).attr("y2", axisY)}
\NormalTok{    .attr("stroke", "\#dc2626").attr("stroke{-}width", 2)}
\NormalTok{    .attr("stroke{-}dasharray", "6,2");}
\NormalTok{  svg.append("text")}
\NormalTok{    .attr("x", xScale(med)).attr("y", 6 + medOffset)}
\NormalTok{    .attr("text{-}anchor", "middle").attr("font{-}size", 11).attr("fill", "\#dc2626")}
\NormalTok{    .attr("font{-}weight", "bold")}
\NormalTok{    .text(\textasciigrave{}Median $\{med\}\textasciigrave{});}

\NormalTok{  // ── Stats table ───────────────────────────────────────────────}
\NormalTok{  const tY = DOT\_AREA\_H + 4;}

\NormalTok{  // Background panel}
\NormalTok{  svg.append("rect")}
\NormalTok{    .attr("x", PAD {-} 10).attr("y", tY)}
\NormalTok{    .attr("width", W {-} 2 * PAD + 20).attr("height", STATS\_H)}
\NormalTok{    .attr("fill", "\#f1f5f9").attr("rx", 4);}

\NormalTok{  const col1 = PAD, col2 = PAD + 180, col3 = PAD + 360;}
\NormalTok{  const rowH = 22;}

\NormalTok{  const stats = [}
\NormalTok{    [\textasciigrave{}Mean  x̄  =  ($\{raw.join(" + ")\}) ÷ $\{n\}\textasciigrave{}, \textasciigrave{}= $\{mean.toFixed(2)\}\textasciigrave{}],}
\NormalTok{    [\textasciigrave{}Median (middle of sorted list)\textasciigrave{}, \textasciigrave{}= $\{med\}\textasciigrave{}],}
\NormalTok{    [\textasciigrave{}Range  =  $\{sorted[n{-}1]\} − $\{sorted[0]\}\textasciigrave{}, \textasciigrave{}= $\{range\_\}\textasciigrave{}],}
\NormalTok{    [\textasciigrave{}Q1 = $\{q1\}   Q3 = $\{q3\}   IQR = Q3 − Q1\textasciigrave{}, \textasciigrave{}= $\{iqr\}\textasciigrave{}],}
\NormalTok{  ];}

\NormalTok{  stats.forEach(([label, val], i) =\textgreater{} \{}
\NormalTok{    svg.append("text")}
\NormalTok{      .attr("x", col1).attr("y", tY + 20 + i * rowH)}
\NormalTok{      .attr("font{-}size", 12).attr("fill", "\#1e293b")}
\NormalTok{      .text(label);}
\NormalTok{    svg.append("text")}
\NormalTok{      .attr("x", W {-} PAD).attr("y", tY + 20 + i * rowH)}
\NormalTok{      .attr("text{-}anchor", "end").attr("font{-}size", 12)}
\NormalTok{      .attr("fill", i === 0 ? "\#2563eb" : i === 1 ? "\#dc2626" : "\#374151")}
\NormalTok{      .attr("font{-}weight", "bold")}
\NormalTok{      .text(val);}
\NormalTok{  \});}

\NormalTok{  // Sorted list display}
\NormalTok{  svg.append("text")}
\NormalTok{    .attr("x", PAD).attr("y", tY + STATS\_H {-} 8)}
\NormalTok{    .attr("font{-}size", 11).attr("fill", "\#64748b")}
\NormalTok{    .text(\textasciigrave{}Sorted: $\{sorted.join(", ")\}\textasciigrave{});}

\NormalTok{  return svg.node();}
\NormalTok{\}}
\end{Highlighting}
\end{Shaded}

\begin{tcolorbox}[enhanced jigsaw, opacityback=0, colback=white, breakable, colframe=quarto-callout-tip-color-frame, arc=.35mm, left=2mm, toprule=.15mm, leftrule=.75mm, bottomrule=.15mm, rightrule=.15mm]
\begin{minipage}[t]{5.5mm}
\textcolor{quarto-callout-tip-color}{\faLightbulb}
\end{minipage}%
\begin{minipage}[t]{\textwidth - 5.5mm}

\textbf{Try this:} change Score 1 from 62 to 2. Watch the mean drop
noticeably while the median barely shifts. That is the difference
between the two measures --- and why it matters which one you use.

\end{minipage}%
\end{tcolorbox}

\begin{center}\rule{0.5\linewidth}{0.5pt}\end{center}

\section{Worked examples}\label{worked-examples-5}

\textbf{Example 1 --- Test scores.} A student scores the following in
seven tests: 62, 71, 58, 74, 66, 70, 63. Find the mean, median, and
range.

Ordered: 58, 62, 63, 66, 70, 71, 74.

Mean --- add all values, then divide by 7:
\[\bar{x} = \frac{62 + 71 + 58 + 74 + 66 + 70 + 63}{7} = \frac{464}{7} = 66.3\]

Range --- highest minus lowest: \[74 - 58 = 16\]

Median --- the 4th value in the ordered list: \[66\]

The mean and median are close, which tells you the scores are fairly
evenly spread --- no extreme result is pulling the average up or down.

\begin{center}\rule{0.5\linewidth}{0.5pt}\end{center}

\textbf{Example 1b --- IQR for the same dataset.} Ordered data (7
values): 58, 62, 63, \textbf{66}, 70, 71, 74.

The median is the 4th value: 66. For an odd-count dataset, exclude the
median when splitting into halves.

Lower half (values 1--3): 58, 62, 63. Q1 = median = 62.

Upper half (values 5--7): 70, 71, 74. Q3 = median = 71.

\[\text{IQR} = Q3 - Q1 = 71 - 62 = 9\]

The middle 50\% of scores sit within a 9-point range.

\begin{center}\rule{0.5\linewidth}{0.5pt}\end{center}

\textbf{Example 2 --- Weather probability.} Over the past 60 days, it
has rained on 15 of them. Estimate the probability of rain on any given
day.

\[P(\text{rain}) \approx \frac{15}{60} = \frac{1}{4} = 0.25 = 25\%\]

This is a relative frequency estimate --- it uses past data rather than
counting equally likely outcomes. The more data you have, the more
reliable the estimate.

\begin{center}\rule{0.5\linewidth}{0.5pt}\end{center}

\textbf{Example 3 --- Picking from a group.} A class has 12 students who
play sport: 5 play football, 4 play basketball, and 3 play tennis. One
student is picked at random to represent the class. What is the
probability they play football? What is the probability they do not play
basketball?

\[P(\text{football}) = \frac{5}{12}\]

\[P(\text{not basketball}) = \frac{12 - 4}{12} = \frac{8}{12} = \frac{2}{3}\]

\begin{center}\rule{0.5\linewidth}{0.5pt}\end{center}

\textbf{Example 4 --- Expected score.} A game show gives two options:
take £300 for certain, or spin a wheel that pays £800 with probability
0.5 and £0 with probability 0.5. What is the expected value of spinning
the wheel?

\[E = (0.5 \times 800) + (0.5 \times 0) = 400 + 0 = £400\]

The expected value of spinning is £400 --- higher than the certain £300.
But ``expected value'' means the average across many spins, not a
guarantee for this one spin. You could walk away with nothing.

\begin{center}\rule{0.5\linewidth}{0.5pt}\end{center}

\emph{The simulator below shows the difference between theoretical
probability and what actually happens in a finite number of draws.
Results vary each time you click --- that is the point.}

\begin{Shaded}
\begin{Highlighting}[]
\NormalTok{//| echo: false}
\NormalTok{viewof red\_count = Inputs.range([1, 10], \{}
\NormalTok{  value: 3,}
\NormalTok{  step: 1,}
\NormalTok{  label: "Red counters in bag"}
\NormalTok{\})}
\end{Highlighting}
\end{Shaded}

\begin{Shaded}
\begin{Highlighting}[]
\NormalTok{//| echo: false}
\NormalTok{viewof blue\_count = Inputs.range([1, 10], \{}
\NormalTok{  value: 5,}
\NormalTok{  step: 1,}
\NormalTok{  label: "Blue counters in bag"}
\NormalTok{\})}
\end{Highlighting}
\end{Shaded}

\begin{Shaded}
\begin{Highlighting}[]
\NormalTok{//| echo: false}
\NormalTok{viewof sim\_button = Inputs.button("Simulate 100 draws")}
\end{Highlighting}
\end{Shaded}

\begin{Shaded}
\begin{Highlighting}[]
\NormalTok{//| echo: false}
\NormalTok{\{}
\NormalTok{  // Trigger on button press (sim\_button changes each click)}
\NormalTok{  sim\_button;}

\NormalTok{  const total = red\_count + blue\_count;}
\NormalTok{  const pRed = red\_count / total;}
\NormalTok{  const pBlue = blue\_count / total;}

\NormalTok{  // Run 100 draws using Math.random()}
\NormalTok{  // Note: results vary on each click — this is intentional and illustrates}
\NormalTok{  // that probability is a long{-}run prediction, not a guarantee for any}
\NormalTok{  // finite sample.}
\NormalTok{  const N = 100;}
\NormalTok{  let simRed = 0, simBlue = 0;}
\NormalTok{  for (let i = 0; i \textless{} N; i++) \{}
\NormalTok{    if (Math.random() \textless{} pRed) simRed++; else simBlue++;}
\NormalTok{  \}}
\NormalTok{  const simPRed = simRed / N;}
\NormalTok{  const simPBlue = simBlue / N;}

\NormalTok{  // ── Layout ───────────────────────────────────────────────────}
\NormalTok{  const W = 560, H = 300, PAD\_L = 120, PAD\_R = 30, PAD\_T = 30, PAD\_B = 60;}
\NormalTok{  const plotW = W {-} PAD\_L {-} PAD\_R;}
\NormalTok{  const plotH = H {-} PAD\_T {-} PAD\_B;}

\NormalTok{  const svg = d3.create("svg")}
\NormalTok{    .attr("viewBox", \textasciigrave{}0 0 $\{W\} $\{H\}\textasciigrave{})}
\NormalTok{    .attr("width", "100\%")}
\NormalTok{    .attr("style", "max{-}width:560px; display:block; margin:0 auto; font{-}family:sans{-}serif;");}

\NormalTok{  // ── Bag visualisation (left panel) ───────────────────────────}
\NormalTok{  // Show counters as small circles in a 4{-}column grid, coloured and labelled}
\NormalTok{  const allCounters = [}
\NormalTok{    ...Array(red\_count).fill("red"),}
\NormalTok{    ...Array(blue\_count).fill("blue")}
\NormalTok{  ];}
\NormalTok{  const COLS = 4;}
\NormalTok{  const CR = 10, CG = 5; // circle radius, gap}
\NormalTok{  const gridX0 = 6, gridY0 = PAD\_T + 10;}
\NormalTok{  allCounters.forEach((colour, i) =\textgreater{} \{}
\NormalTok{    const col = i \% COLS;}
\NormalTok{    const row = Math.floor(i / COLS);}
\NormalTok{    const cx = gridX0 + col * (CR * 2 + CG) + CR;}
\NormalTok{    const cy = gridY0 + row * (CR * 2 + CG) + CR;}
\NormalTok{    svg.append("circle")}
\NormalTok{      .attr("cx", cx).attr("cy", cy).attr("r", CR)}
\NormalTok{      .attr("fill", colour === "red" ? "\#dc2626" : "\#2563eb")}
\NormalTok{      .attr("stroke", colour === "red" ? "\#7f1d1d" : "\#1e3a5f")}
\NormalTok{      .attr("stroke{-}width", 1.5);}
\NormalTok{  \});}

\NormalTok{  // Bag label}
\NormalTok{  svg.append("text")}
\NormalTok{    .attr("x", gridX0 + (COLS * (CR * 2 + CG)) / 2)}
\NormalTok{    .attr("y", H {-} PAD\_B + 15)}
\NormalTok{    .attr("text{-}anchor", "middle").attr("font{-}size", 11).attr("fill", "\#374151")}
\NormalTok{    .text(\textasciigrave{}Bag: $\{red\_count\} red, $\{blue\_count\} blue\textasciigrave{});}

\NormalTok{  svg.append("text")}
\NormalTok{    .attr("x", gridX0 + (COLS * (CR * 2 + CG)) / 2)}
\NormalTok{    .attr("y", H {-} PAD\_B + 27)}
\NormalTok{    .attr("text{-}anchor", "middle").attr("font{-}size", 10).attr("fill", "\#64748b")}
\NormalTok{    .text(\textasciigrave{}Total: $\{total\}\textasciigrave{});}

\NormalTok{  // ── Bar chart ────────────────────────────────────────────────}
\NormalTok{  const barData = [}
\NormalTok{    \{ label: "Red",  theory: pRed,  sim: simPRed,  colour: "\#dc2626", stripe: "\#7f1d1d" \},}
\NormalTok{    \{ label: "Blue", theory: pBlue, sim: simPBlue, colour: "\#2563eb", stripe: "\#1e3a5f" \}}
\NormalTok{  ];}

\NormalTok{  const groupW = plotW / barData.length;}
\NormalTok{  const barW = groupW * 0.35;}
\NormalTok{  const yScale = v =\textgreater{} PAD\_T + plotH * (1 {-} v);}

\NormalTok{  // Y axis}
\NormalTok{  svg.append("line")}
\NormalTok{    .attr("x1", PAD\_L).attr("y1", PAD\_T)}
\NormalTok{    .attr("x2", PAD\_L).attr("y2", PAD\_T + plotH)}
\NormalTok{    .attr("stroke", "\#94a3b8").attr("stroke{-}width", 1.5);}

\NormalTok{  // Y ticks}
\NormalTok{  [0, 0.25, 0.5, 0.75, 1.0].forEach(v =\textgreater{} \{}
\NormalTok{    const y = yScale(v);}
\NormalTok{    svg.append("line")}
\NormalTok{      .attr("x1", PAD\_L {-} 5).attr("y1", y)}
\NormalTok{      .attr("x2", PAD\_L).attr("y2", y)}
\NormalTok{      .attr("stroke", "\#94a3b8").attr("stroke{-}width", 1);}
\NormalTok{    svg.append("text")}
\NormalTok{      .attr("x", PAD\_L {-} 8).attr("y", y + 4)}
\NormalTok{      .attr("text{-}anchor", "end").attr("font{-}size", 10).attr("fill", "\#64748b")}
\NormalTok{      .text(v.toFixed(2));}
\NormalTok{    // Gridline}
\NormalTok{    svg.append("line")}
\NormalTok{      .attr("x1", PAD\_L).attr("y1", y)}
\NormalTok{      .attr("x2", PAD\_L + plotW).attr("y2", y)}
\NormalTok{      .attr("stroke", "\#e2e8f0").attr("stroke{-}width", 1);}
\NormalTok{  \});}

\NormalTok{  // Y axis label}
\NormalTok{  svg.append("text")}
\NormalTok{    .attr("transform", \textasciigrave{}translate($\{PAD\_L {-} 40\}, $\{PAD\_T + plotH / 2\}) rotate({-}90)\textasciigrave{})}
\NormalTok{    .attr("text{-}anchor", "middle").attr("font{-}size", 11).attr("fill", "\#374151")}
\NormalTok{    .text("Probability / relative frequency");}

\NormalTok{  // X axis}
\NormalTok{  svg.append("line")}
\NormalTok{    .attr("x1", PAD\_L).attr("y1", PAD\_T + plotH)}
\NormalTok{    .attr("x2", PAD\_L + plotW).attr("y2", PAD\_T + plotH)}
\NormalTok{    .attr("stroke", "\#94a3b8").attr("stroke{-}width", 1.5);}

\NormalTok{  barData.forEach((d, gi) =\textgreater{} \{}
\NormalTok{    const gx = PAD\_L + gi * groupW + groupW / 2;}

\NormalTok{    // Theory bar (solid fill)}
\NormalTok{    const th = yScale(d.theory);}
\NormalTok{    const thH = (PAD\_T + plotH) {-} th;}
\NormalTok{    svg.append("rect")}
\NormalTok{      .attr("x", gx {-} barW {-} 2).attr("y", th)}
\NormalTok{      .attr("width", barW).attr("height", thH)}
\NormalTok{      .attr("fill", d.colour).attr("opacity", 0.9)}
\NormalTok{      .attr("rx", 2);}

\NormalTok{    // Sim bar (hatched: lighter fill + dashed outline)}
\NormalTok{    const sh = yScale(d.sim);}
\NormalTok{    const shH = (PAD\_T + plotH) {-} sh;}
\NormalTok{    svg.append("rect")}
\NormalTok{      .attr("x", gx + 2).attr("y", sh)}
\NormalTok{      .attr("width", barW).attr("height", shH)}
\NormalTok{      .attr("fill", d.colour).attr("opacity", 0.35)}
\NormalTok{      .attr("stroke", d.colour).attr("stroke{-}width", 1.5)}
\NormalTok{      .attr("stroke{-}dasharray", "4,2")}
\NormalTok{      .attr("rx", 2);}

\NormalTok{    // Value labels on bars}
\NormalTok{    svg.append("text")}
\NormalTok{      .attr("x", gx {-} barW / 2 {-} 2).attr("y", th {-} 3)}
\NormalTok{      .attr("text{-}anchor", "middle").attr("font{-}size", 10)}
\NormalTok{      .attr("fill", d.stripe).attr("font{-}weight", "bold")}
\NormalTok{      .text(\textasciigrave{}$\{(d.theory * 100).toFixed(0)\}\%\textasciigrave{});}

\NormalTok{    svg.append("text")}
\NormalTok{      .attr("x", gx + barW / 2 + 2).attr("y", sh {-} 3)}
\NormalTok{      .attr("text{-}anchor", "middle").attr("font{-}size", 10)}
\NormalTok{      .attr("fill", d.stripe).attr("font{-}weight", "bold")}
\NormalTok{      .text(\textasciigrave{}$\{(d.sim * 100).toFixed(0)\}\%\textasciigrave{});}

\NormalTok{    // Group label (colour name)}
\NormalTok{    svg.append("text")}
\NormalTok{      .attr("x", gx).attr("y", PAD\_T + plotH + 16)}
\NormalTok{      .attr("text{-}anchor", "middle").attr("font{-}size", 12).attr("fill", "\#374151")}
\NormalTok{      .text(d.label);}
\NormalTok{  \});}

\NormalTok{  // ── Legend ───────────────────────────────────────────────────}
\NormalTok{  const legX = PAD\_L + 10, legY = PAD\_T + 6;}
\NormalTok{  // Solid = theory}
\NormalTok{  svg.append("rect")}
\NormalTok{    .attr("x", legX).attr("y", legY).attr("width", 14).attr("height", 10)}
\NormalTok{    .attr("fill", "\#374151").attr("opacity", 0.8).attr("rx", 1);}
\NormalTok{  svg.append("text")}
\NormalTok{    .attr("x", legX + 18).attr("y", legY + 9)}
\NormalTok{    .attr("font{-}size", 10).attr("fill", "\#374151")}
\NormalTok{    .text("Theoretical probability");}

\NormalTok{  svg.append("rect")}
\NormalTok{    .attr("x", legX + 150).attr("y", legY).attr("width", 14).attr("height", 10)}
\NormalTok{    .attr("fill", "\#374151").attr("opacity", 0.35)}
\NormalTok{    .attr("stroke", "\#374151").attr("stroke{-}width", 1.5)}
\NormalTok{    .attr("stroke{-}dasharray", "3,2").attr("rx", 1);}
\NormalTok{  svg.append("text")}
\NormalTok{    .attr("x", legX + 168).attr("y", legY + 9)}
\NormalTok{    .attr("font{-}size", 10).attr("fill", "\#374151")}
\NormalTok{    .text(\textasciigrave{}Simulated ($\{N\} draws)\textasciigrave{});}

\NormalTok{  // ── Summary line ─────────────────────────────────────────────}
\NormalTok{  svg.append("text")}
\NormalTok{    .attr("x", W / 2).attr("y", H {-} 10)}
\NormalTok{    .attr("text{-}anchor", "middle").attr("font{-}size", 11).attr("fill", "\#374151")}
\NormalTok{    .text(\textasciigrave{}Red: theory $\{(pRed*100).toFixed(1)\}\%  vs  simulated $\{(simPRed*100).toFixed(0)\}\%  |  Blue: theory $\{(pBlue*100).toFixed(1)\}\%  vs  simulated $\{(simPBlue*100).toFixed(0)\}\%\textasciigrave{});}

\NormalTok{  return svg.node();}
\NormalTok{\}}
\end{Highlighting}
\end{Shaded}

\begin{tcolorbox}[enhanced jigsaw, opacityback=0, colback=white, breakable, colframe=quarto-callout-note-color-frame, arc=.35mm, left=2mm, toprule=.15mm, leftrule=.75mm, bottomrule=.15mm, rightrule=.15mm]
\begin{minipage}[t]{5.5mm}
\textcolor{quarto-callout-note-color}{\faInfo}
\end{minipage}%
\begin{minipage}[t]{\textwidth - 5.5mm}

\textbf{Simulation note:} each click on ``Simulate 100 draws'' generates
a fresh set of random draws. The simulated bars will be close to the
theoretical bars but rarely identical --- that gap between prediction
and observation is real and expected. With 1 000 draws the bars would be
closer; with 10 draws they might be very far apart. More data means a
more reliable estimate.

\end{minipage}%
\end{tcolorbox}

\section{Where this goes}\label{where-this-goes-5}

This chapter gives you the vocabulary --- mean, median, spread,
probability --- that makes statistical arguments legible. \textbf{Volume
6} develops this substantially: probability distributions, hypothesis
testing, and regression. That is where statistics becomes a tool for
making decisions under uncertainty rather than just summarising data.

The expected value calculation in Example 4 is also the foundation of
decision theory, financial option pricing (Black-Scholes is built on
expected values under a probability distribution), and machine learning
loss functions. The concept is old. The applications are new.

\begin{tcolorbox}[enhanced jigsaw, opacityback=0, colback=white, breakable, colframe=quarto-callout-tip-color-frame, arc=.35mm, left=2mm, toprule=.15mm, leftrule=.75mm, bottomrule=.15mm, rightrule=.15mm]
\begin{minipage}[t]{5.5mm}
\textcolor{quarto-callout-tip-color}{\faLightbulb}
\end{minipage}%
\begin{minipage}[t]{\textwidth - 5.5mm}

\textbf{Where this shows up}

\begin{itemize}
\tightlist
\item
  Sports commentators and coaches use averages and ranges to compare
  players' performance over a season.
\item
  A weather forecast's percentage chance of rain is a probability
  estimated from thousands of past days with similar conditions.
\item
  A lab scientist reports every measurement with ± uncertainty --- that
  is a statement about spread.
\item
  A data scientist evaluates a classifier by its error rate --- a
  probability computed from test data.
\end{itemize}

Statistics is not a course you take. It is a way of reading claims about
the world.

\end{minipage}%
\end{tcolorbox}

\section{Exercises}\label{exercises-5}

\begin{enumerate}
\def\labelenumi{\arabic{enumi}.}
\item
  The daily high temperatures for one week (°C): −4, −7, −2, 1, 3, −1,
  −5. Find the mean, median, and range.
\item
  A student scores 62, 71, 85, 68, 90, 74, and 55 on seven tests. Find
  the mean and median. A new student joins the group with a score of 4.
  Recalculate the mean. Does the median change? What does this show
  about which measure is more stable under an extreme value?
\item
  A standard die has faces 1--6. What is the probability of rolling a
  number greater than 4? Of rolling an even number? Of rolling a 7?
\item
  A bag contains 4 red, 6 blue, and 2 green marbles. A marble is drawn
  at random.

  \begin{enumerate}
  \def\labelenumii{\alph{enumii})}
  \tightlist
  \item
    What is P(red)?
  \item
    What is P(not blue)?
  \item
    What is P(red or green)?
  \end{enumerate}
\item
  A school canteen tracks how many students choose each lunch option
  over 200 days: hot meal 110 times, sandwich 60 times, salad 30 times.
  Estimate the probability a randomly chosen day has more than half the
  students choosing a hot meal. Out of the next 50 days, how many would
  you expect?
\item
  A quiz game has three options:

  \begin{itemize}
  \tightlist
  \item
    Option A: 80\% chance of +£500, 20\% chance of −£200
  \item
    Option B: 50\% chance of +£900, 50\% chance of −£100
  \item
    Option C: certain return of +£280
  \end{itemize}

  Compute the expected value of each option. Which is highest? Which
  would you choose, and why might expected value alone not be the only
  consideration?
\item
  A set of 9 test scores has a median of 14.2 and a mean of 15.1. One
  outlier is identified and removed, leaving 8 values with a mean of
  14.6. What was the outlier value? \emph{(Hint: total of 9 values =
  mean × 9. Total of remaining 8 = new mean × 8. Outlier = first total −
  second total.)}
\end{enumerate}

\part{Volume 2 --- Patterns and algebra}

\chapter{Volume 2 --- Patterns and
algebra}\label{volume-2-patterns-and-algebra-1}

The central move of this volume is naming the gap.

You already know how numbers behave. This volume is about what happens
when one of the numbers is missing --- and about the notation that lets
you work with that absence as if it were a thing you could hold.

By the end of this volume you will be able to write a relationship down
symbolically, solve for unknowns, recognise when a relationship is
linear versus something else, and read a graph as a statement about how
two quantities relate.

These are not abstract skills. They are the toolkit that every
quantitative field --- engineering, science, finance, computing,
geography --- uses before it reaches for anything more sophisticated.

\section{Chapters}\label{chapters-1}

\begin{longtable}[]{@{}lll@{}}
\toprule\noalign{}
Chapter & Title & Core move \\
\midrule\noalign{}
\endhead
\bottomrule\noalign{}
\endlastfoot
1 & Linear equations & Isolating the unknown \\
2 & Inequalities & When the answer is a range \\
3 & Patterns and sequences & When the relationship repeats \\
4 & Graphing relations & When the relationship has a shape \\
\end{longtable}

\chapter{Linear equations}\label{linear-equations}

Finding the unknown

\hfill\break

You and three friends want to split a bill evenly, but one person
already paid for something earlier. How much does each person owe now?

You earn £8.50 an hour at a weekend job. You need £120 for a gig ticket.
How many shifts do you need?

You're buying snacks for a group. Each bag costs the same. You spend
£14.40 in total. How much was each bag?

These are all the same question underneath: you know the relationship,
you're missing one of the values. A linear equation is just that
relationship written down, with a placeholder for the missing part.

\section{What the notation is
saying}\label{what-the-notation-is-saying-6}

\[3x + 7 = 22\]

Read it as: three of something, plus seven, gives twenty-two. The \(x\)
is the placeholder --- a name for the gap in your knowledge. The equals
sign means both sides describe the same amount.

Peel away what's around the unknown until it stands alone. Three of
something plus seven is twenty-two, so three of something is fifteen, so
the something is five. Check: \(3 \times 5 + 7 = 22\). Yes.

That move --- isolating the unknown --- is what this chapter is about.

The concept stays this simple. What gets harder is the notation
accumulating, and equations appearing inside larger problems where it
isn't obvious that's what they are.

\begin{center}\rule{0.5\linewidth}{0.5pt}\end{center}

\textbf{Interactive: Equation balance.} The scale below shows both sides
of an equation. Choose an operation and a value. Apply it to both sides
and watch the equation update --- the balance holds because you're
changing both sides equally.

\begin{Shaded}
\begin{Highlighting}[]
\NormalTok{//| echo: false}
\NormalTok{d3 = require("d3@7")}
\end{Highlighting}
\end{Shaded}

\begin{Shaded}
\begin{Highlighting}[]
\NormalTok{//| echo: false}
\NormalTok{viewof balanceOp = Inputs.select(["+", "−", "×", "÷"], \{}
\NormalTok{  label: "Operation",}
\NormalTok{  value: "+"}
\NormalTok{\})}
\end{Highlighting}
\end{Shaded}

\begin{Shaded}
\begin{Highlighting}[]
\NormalTok{//| echo: false}
\NormalTok{viewof balanceVal = Inputs.range([{-}10, 10], \{}
\NormalTok{  label: "Value to apply to both sides",}
\NormalTok{  step: 1,}
\NormalTok{  value: 3}
\NormalTok{\})}
\end{Highlighting}
\end{Shaded}

\begin{Shaded}
\begin{Highlighting}[]
\NormalTok{//| echo: false}
\NormalTok{\{}
\NormalTok{  // Starting equation: 3x + 7 = 22}
\NormalTok{  // We track the numeric sides as they stand after one operation.}
\NormalTok{  // Left side: 3x + 7  →  represented symbolically}
\NormalTok{  // Right side: 22}
\NormalTok{  // For display we show the symbolic equation text and what each side becomes.}

\NormalTok{  const op  = balanceOp;}
\NormalTok{  const val = Number.isFinite(balanceVal) ? Math.round(balanceVal) : 3;}

\NormalTok{  // Guard against division by zero}
\NormalTok{  const divByZero = op === "÷" \&\& val === 0;}

\NormalTok{  // Represent the original equation symbolically}
\NormalTok{  const origLeft  = "3x + 7";}
\NormalTok{  const origRight = "22";}

\NormalTok{  // Compute what happens to the right{-}hand side (a pure number)}
\NormalTok{  let newRight;}
\NormalTok{  if (op === "+")       newRight = 22 + val;}
\NormalTok{  else if (op === "−")  newRight = 22 {-} val;}
\NormalTok{  else if (op === "×")  newRight = 22 * val;}
\NormalTok{  else if (op === "÷")  newRight = divByZero ? null : 22 / val;}

\NormalTok{  // Build a symbolic description of the new left side}
\NormalTok{  let newLeftDesc;}
\NormalTok{  const valStr = val \textless{} 0 ? \textasciigrave{}($\{val\})\textasciigrave{} : \textasciigrave{}$\{val\}\textasciigrave{};}
\NormalTok{  if (op === "+")       newLeftDesc = \textasciigrave{}3x + 7 + $\{valStr\}\textasciigrave{};}
\NormalTok{  else if (op === "−")  newLeftDesc = \textasciigrave{}3x + 7 − $\{valStr\}\textasciigrave{};}
\NormalTok{  else if (op === "×")  newLeftDesc = \textasciigrave{}$\{valStr\}(3x + 7)\textasciigrave{};}
\NormalTok{  else if (op === "÷")  newLeftDesc = divByZero ? "3x + 7 ÷ 0" : \textasciigrave{}(3x + 7) ÷ $\{valStr\}\textasciigrave{};}

\NormalTok{  // Display}
\NormalTok{  const W = 560;}
\NormalTok{  const H = 260;}
\NormalTok{  const MID\_X = W / 2;}
\NormalTok{  const BEAM\_Y = 120;}
\NormalTok{  const PIVOT\_Y = 80;}
\NormalTok{  const PAN\_Y   = 170;}
\NormalTok{  const PAN\_W   = 150;}
\NormalTok{  const PAN\_H   = 44;}
\NormalTok{  const ROPE\_H  = 50;}

\NormalTok{  const svg = d3.create("svg")}
\NormalTok{    .attr("viewBox", \textasciigrave{}0 0 $\{W\} $\{H\}\textasciigrave{})}
\NormalTok{    .attr("width", "100\%")}
\NormalTok{    .attr("style", \textasciigrave{}max{-}width:$\{W\}px; font{-}family: inherit;\textasciigrave{});}

\NormalTok{  // ── Pivot post ──}
\NormalTok{  svg.append("line")}
\NormalTok{    .attr("x1", MID\_X).attr("y1", PIVOT\_Y {-} 10)}
\NormalTok{    .attr("x2", MID\_X).attr("y2", BEAM\_Y)}
\NormalTok{    .attr("stroke", "\#9ca3af").attr("stroke{-}width", 3);}

\NormalTok{  svg.append("circle")}
\NormalTok{    .attr("cx", MID\_X).attr("cy", PIVOT\_Y {-} 10)}
\NormalTok{    .attr("r", 6)}
\NormalTok{    .attr("fill", "\#6b7280");}

\NormalTok{  // ── Beam ──}
\NormalTok{  svg.append("line")}
\NormalTok{    .attr("x1", 80).attr("y1", BEAM\_Y)}
\NormalTok{    .attr("x2", W {-} 80).attr("y2", BEAM\_Y)}
\NormalTok{    .attr("stroke", "\#374151").attr("stroke{-}width", 4)}
\NormalTok{    .attr("stroke{-}linecap", "round");}

\NormalTok{  // ── Ropes ──}
\NormalTok{  const leftPanCX  = 130;}
\NormalTok{  const rightPanCX = W {-} 130;}

\NormalTok{  svg.append("line")}
\NormalTok{    .attr("x1", leftPanCX).attr("y1", BEAM\_Y)}
\NormalTok{    .attr("x2", leftPanCX).attr("y2", PAN\_Y {-} PAN\_H / 2)}
\NormalTok{    .attr("stroke", "\#9ca3af").attr("stroke{-}width", 2);}

\NormalTok{  svg.append("line")}
\NormalTok{    .attr("x1", rightPanCX).attr("y1", BEAM\_Y)}
\NormalTok{    .attr("x2", rightPanCX).attr("y2", PAN\_Y {-} PAN\_H / 2)}
\NormalTok{    .attr("stroke", "\#9ca3af").attr("stroke{-}width", 2);}

\NormalTok{  // ── Pans ──}
\NormalTok{  const panFill = "\#f3f4f6";}
\NormalTok{  const panStroke = "\#d1d5db";}

\NormalTok{  svg.append("rect")}
\NormalTok{    .attr("x", leftPanCX {-} PAN\_W / 2).attr("y", PAN\_Y {-} PAN\_H / 2)}
\NormalTok{    .attr("width", PAN\_W).attr("height", PAN\_H)}
\NormalTok{    .attr("rx", 6)}
\NormalTok{    .attr("fill", panFill).attr("stroke", panStroke).attr("stroke{-}width", 1.5);}

\NormalTok{  svg.append("rect")}
\NormalTok{    .attr("x", rightPanCX {-} PAN\_W / 2).attr("y", PAN\_Y {-} PAN\_H / 2)}
\NormalTok{    .attr("width", PAN\_W).attr("height", PAN\_H)}
\NormalTok{    .attr("rx", 6)}
\NormalTok{    .attr("fill", panFill).attr("stroke", panStroke).attr("stroke{-}width", 1.5);}

\NormalTok{  // ── Pan labels: new expressions ──}
\NormalTok{  const leftLabel  = divByZero ? "÷ 0 not defined" : newLeftDesc;}
\NormalTok{  const rightLabel = divByZero ? "÷ 0 not defined" :}
\NormalTok{    (Number.isInteger(newRight) ? \textasciigrave{}$\{newRight\}\textasciigrave{} : newRight.toFixed(2));}

\NormalTok{  svg.append("text")}
\NormalTok{    .attr("x", leftPanCX).attr("y", PAN\_Y + 5)}
\NormalTok{    .attr("text{-}anchor", "middle")}
\NormalTok{    .attr("fill", divByZero ? "\#dc2626" : "\#1f2937")}
\NormalTok{    .attr("font{-}size", divByZero ? "10px" : "12px")}
\NormalTok{    .attr("font{-}weight", "700")}
\NormalTok{    .text(leftLabel);}

\NormalTok{  svg.append("text")}
\NormalTok{    .attr("x", rightPanCX).attr("y", PAN\_Y + 5)}
\NormalTok{    .attr("text{-}anchor", "middle")}
\NormalTok{    .attr("fill", divByZero ? "\#dc2626" : "\#1f2937")}
\NormalTok{    .attr("font{-}size", "14px")}
\NormalTok{    .attr("font{-}weight", "700")}
\NormalTok{    .text(rightLabel);}

\NormalTok{  // ── Equation line above balance ──}
\NormalTok{  const origEq = \textasciigrave{}Original: $\{origLeft\} = $\{origRight\}\textasciigrave{};}
\NormalTok{  svg.append("text")}
\NormalTok{    .attr("x", MID\_X).attr("y", 32)}
\NormalTok{    .attr("text{-}anchor", "middle")}
\NormalTok{    .attr("fill", "\#374151").attr("font{-}size", "13px").attr("font{-}weight", "600")}
\NormalTok{    .text(origEq);}

\NormalTok{  // ── Operation applied ──}
\NormalTok{  const opSymbol = op === "÷" ? "÷" : op === "×" ? "×" : op;}
\NormalTok{  const opDesc = divByZero}
\NormalTok{    ? "Cannot divide by zero — the balance breaks"}
\NormalTok{    : \textasciigrave{}Apply  $\{opSymbol\} $\{val\}  to both sides\textasciigrave{};}

\NormalTok{  svg.append("text")}
\NormalTok{    .attr("x", MID\_X).attr("y", 52)}
\NormalTok{    .attr("text{-}anchor", "middle")}
\NormalTok{    .attr("fill", divByZero ? "\#dc2626" : "\#0d9488")}
\NormalTok{    .attr("font{-}size", "12px")}
\NormalTok{    .text(opDesc);}

\NormalTok{  // ── "Both sides equal" label ──}
\NormalTok{  if (!divByZero) \{}
\NormalTok{    svg.append("text")}
\NormalTok{      .attr("x", MID\_X).attr("y", H {-} 12)}
\NormalTok{      .attr("text{-}anchor", "middle")}
\NormalTok{      .attr("fill", "\#6b7280").attr("font{-}size", "11px")}
\NormalTok{      .text("Both sides still describe the same amount — the balance holds.");}
\NormalTok{  \}}

\NormalTok{  // ── "=" centred between pans ──}
\NormalTok{  svg.append("text")}
\NormalTok{    .attr("x", MID\_X).attr("y", PAN\_Y + 6)}
\NormalTok{    .attr("text{-}anchor", "middle")}
\NormalTok{    .attr("fill", "\#6b7280").attr("font{-}size", "18px").attr("font{-}weight", "700")}
\NormalTok{    .text("=");}

\NormalTok{  return svg.node();}
\NormalTok{\}}
\end{Highlighting}
\end{Shaded}

\begin{center}\rule{0.5\linewidth}{0.5pt}\end{center}

\section{The method}\label{the-method-6}

The rule is: whatever you do to one side of the equals sign, do to the
other. The equation is a balance. You can add weight to both sides,
remove weight from both sides, or multiply both sides by the same amount
--- as long as you do the same thing to both, the balance holds.

\[3x + 7 = 22\]

Subtract 7 from both sides --- this removes the 7 from the left, leaving
just the term with \(x\):

\[3x = 15\]

Divide both sides by 3 --- this peels away the coefficient, leaving
\(x\) alone:

\[x = 5\]

That's it. Two operations, one unknown, done.

\begin{tcolorbox}[enhanced jigsaw, opacityback=0, colback=white, breakable, colframe=quarto-callout-note-color-frame, arc=.35mm, left=2mm, toprule=.15mm, leftrule=.75mm, bottomrule=.15mm, rightrule=.15mm]
\begin{minipage}[t]{5.5mm}
\textcolor{quarto-callout-note-color}{\faInfo}
\end{minipage}%
\begin{minipage}[t]{\textwidth - 5.5mm}

\textbf{Why this works}

The equals sign is a statement that two things are the same amount. If
you change one side without changing the other, they're no longer the
same amount --- you've broken the equation. Doing the same thing to both
sides preserves the equality. That's the whole method.

\end{minipage}%
\end{tcolorbox}

\begin{center}\rule{0.5\linewidth}{0.5pt}\end{center}

\textbf{Interactive: Equation solver.} Set the coefficients \(a\),
\(b\), and \(c\) in the equation \(ax + b = c\). The solution steps
update live with your numbers.

\begin{Shaded}
\begin{Highlighting}[]
\NormalTok{//| echo: false}
\NormalTok{viewof solveA = Inputs.range([{-}10, 10], \{}
\NormalTok{  label: "a  (coefficient of x)",}
\NormalTok{  step: 1,}
\NormalTok{  value: 3}
\NormalTok{\})}
\end{Highlighting}
\end{Shaded}

\begin{Shaded}
\begin{Highlighting}[]
\NormalTok{//| echo: false}
\NormalTok{viewof solveB = Inputs.range([{-}20, 20], \{}
\NormalTok{  label: "b  (constant on left)",}
\NormalTok{  step: 1,}
\NormalTok{  value: 7}
\NormalTok{\})}
\end{Highlighting}
\end{Shaded}

\begin{Shaded}
\begin{Highlighting}[]
\NormalTok{//| echo: false}
\NormalTok{viewof solveC = Inputs.range([{-}50, 50], \{}
\NormalTok{  label: "c  (right{-}hand side)",}
\NormalTok{  step: 1,}
\NormalTok{  value: 22}
\NormalTok{\})}
\end{Highlighting}
\end{Shaded}

\begin{Shaded}
\begin{Highlighting}[]
\NormalTok{//| echo: false}
\NormalTok{\{}
\NormalTok{  const a = Number.isFinite(solveA) ? Math.round(solveA) : 3;}
\NormalTok{  const b = Number.isFinite(solveB) ? Math.round(solveB) : 7;}
\NormalTok{  const c = Number.isFinite(solveC) ? Math.round(solveC) : 22;}

\NormalTok{  const noSolution = a === 0;}
\NormalTok{  const cMinusB    = c {-} b;}
\NormalTok{  const x          = noSolution ? null : cMinusB / a;}

\NormalTok{  // Format a coefficient: omit "1x" → "x", handle negatives neatly}
\NormalTok{  const coeffStr = (n) =\textgreater{} \{}
\NormalTok{    if (n === 1)  return "x";}
\NormalTok{    if (n === {-}1) return "−x";}
\NormalTok{    if (n \textless{} 0)    return \textasciigrave{}$\{n\}x\textasciigrave{};}
\NormalTok{    return \textasciigrave{}$\{n\}x\textasciigrave{};}
\NormalTok{  \};}

\NormalTok{  // Format the left side "ax + b" nicely}
\NormalTok{  const leftSide = () =\textgreater{} \{}
\NormalTok{    const aPart = coeffStr(a);}
\NormalTok{    if (b === 0) return \textasciigrave{}$\{aPart\}\textasciigrave{};}
\NormalTok{    if (b \textgreater{} 0)   return \textasciigrave{}$\{aPart\} + $\{b\}\textasciigrave{};}
\NormalTok{    return \textasciigrave{}$\{aPart\} − $\{Math.abs(b)\}\textasciigrave{};}
\NormalTok{  \};}

\NormalTok{  const W = 560;}
\NormalTok{  const H = noSolution ? 130 : 200;}
\NormalTok{  const ROW\_H = 38;}
\NormalTok{  const LABEL\_X = 20;}
\NormalTok{  const EQ\_X = 180;}

\NormalTok{  const svg = d3.create("svg")}
\NormalTok{    .attr("viewBox", \textasciigrave{}0 0 $\{W\} $\{H\}\textasciigrave{})}
\NormalTok{    .attr("width", "100\%")}
\NormalTok{    .attr("style", \textasciigrave{}max{-}width:$\{W\}px; font{-}family: inherit;\textasciigrave{});}

\NormalTok{  // Header: equation as given}
\NormalTok{  svg.append("text")}
\NormalTok{    .attr("x", LABEL\_X).attr("y", 28)}
\NormalTok{    .attr("fill", "\#1f2937").attr("font{-}size", "15px").attr("font{-}weight", "700")}
\NormalTok{    .text(\textasciigrave{}Solve:  $\{leftSide()\} = $\{c\}\textasciigrave{});}

\NormalTok{  if (noSolution) \{}
\NormalTok{    svg.append("text")}
\NormalTok{      .attr("x", LABEL\_X).attr("y", 70)}
\NormalTok{      .attr("fill", "\#dc2626").attr("font{-}size", "13px").attr("font{-}weight", "600")}
\NormalTok{      .text("a = 0: no unique solution (x drops out of the equation).");}

\NormalTok{    svg.append("text")}
\NormalTok{      .attr("x", LABEL\_X).attr("y", 96)}
\NormalTok{      .attr("fill", "\#6b7280").attr("font{-}size", "12px")}
\NormalTok{      .text(b === c}
\NormalTok{        ? \textasciigrave{}With a = 0 the equation reads $\{b\} = $\{c\}, which is always true — any x works.\textasciigrave{}}
\NormalTok{        : \textasciigrave{}With a = 0 the equation reads $\{b\} = $\{c\}, which is never true — no x works.\textasciigrave{}}
\NormalTok{      );}
\NormalTok{  \} else \{}
\NormalTok{    // Step 1}
\NormalTok{    const step1Y = 64;}
\NormalTok{    svg.append("text")}
\NormalTok{      .attr("x", LABEL\_X).attr("y", step1Y)}
\NormalTok{      .attr("fill", "\#6b7280").attr("font{-}size", "12px")}
\NormalTok{      .text(\textasciigrave{}Step 1 — subtract $\{b \textless{} 0 ? \textasciigrave{}($\{b\})\textasciigrave{} : b\} from both sides\textasciigrave{}}
\NormalTok{            + (b \textless{} 0 ? "  [adding a negative]" : "")}
\NormalTok{            + ":");}

\NormalTok{    const bStr = b \textless{} 0 ? \textasciigrave{}+ ($\{b\})\textasciigrave{} : \textasciigrave{}− $\{b\}\textasciigrave{};}
\NormalTok{    svg.append("text")}
\NormalTok{      .attr("x", EQ\_X).attr("y", step1Y + 20)}
\NormalTok{      .attr("fill", "\#1f2937").attr("font{-}size", "14px").attr("font{-}weight", "600")}
\NormalTok{      .text(\textasciigrave{}$\{coeffStr(a)\} = $\{c\} $\{bStr\}  =  $\{cMinusB\}\textasciigrave{});}

\NormalTok{    // Step 2}
\NormalTok{    const step2Y = step1Y + 56;}
\NormalTok{    svg.append("text")}
\NormalTok{      .attr("x", LABEL\_X).attr("y", step2Y)}
\NormalTok{      .attr("fill", "\#6b7280").attr("font{-}size", "12px")}
\NormalTok{      .text(\textasciigrave{}Step 2 — divide both sides by $\{a \textless{} 0 ? \textasciigrave{}($\{a\})\textasciigrave{} : a\}:\textasciigrave{});}

\NormalTok{    const xDisplay = Number.isInteger(x)}
\NormalTok{      ? \textasciigrave{}$\{x\}\textasciigrave{}}
\NormalTok{      : x.toFixed(4).replace(/\textbackslash{}.?0+$/, "");}

\NormalTok{    svg.append("text")}
\NormalTok{      .attr("x", EQ\_X).attr("y", step2Y + 20)}
\NormalTok{      .attr("fill", "\#1f2937").attr("font{-}size", "14px").attr("font{-}weight", "600")}
\NormalTok{      .text(\textasciigrave{}x = $\{cMinusB\} ÷ $\{a\}  =  $\{xDisplay\}\textasciigrave{});}

\NormalTok{    // Answer box}
\NormalTok{    const ansY = step2Y + 46;}
\NormalTok{    svg.append("rect")}
\NormalTok{      .attr("x", EQ\_X {-} 8).attr("y", ansY {-} 20)}
\NormalTok{      .attr("width", 220).attr("height", 32)}
\NormalTok{      .attr("rx", 6)}
\NormalTok{      .attr("fill", "\#0d9488").attr("opacity", 0.12)}
\NormalTok{      .attr("stroke", "\#0d9488").attr("stroke{-}width", 1.5);}

\NormalTok{    svg.append("text")}
\NormalTok{      .attr("x", EQ\_X + 2).attr("y", ansY)}
\NormalTok{      .attr("fill", "\#0d9488").attr("font{-}size", "15px").attr("font{-}weight", "700")}
\NormalTok{      .text(\textasciigrave{}Answer:  x = $\{xDisplay\}\textasciigrave{});}

\NormalTok{    // Check}
\NormalTok{    const checkY = ansY + 28;}
\NormalTok{    const checkLeft  = a * (Number.isInteger(x) ? x : parseFloat(xDisplay)) + b;}
\NormalTok{    const checkRight = c;}
\NormalTok{    const checkOk    = Math.abs(checkLeft {-} checkRight) \textless{} 1e{-}9;}
\NormalTok{    svg.append("text")}
\NormalTok{      .attr("x", LABEL\_X).attr("y", checkY)}
\NormalTok{      .attr("fill", checkOk ? "\#6b7280" : "\#dc2626")}
\NormalTok{      .attr("font{-}size", "11px")}
\NormalTok{      .text(checkOk}
\NormalTok{        ? \textasciigrave{}Check: $\{a\}×$\{xDisplay\} + $\{b\} = $\{checkLeft\} ✓\textasciigrave{}}
\NormalTok{        : \textasciigrave{}Check: $\{a\}×$\{xDisplay\} + $\{b\} = $\{checkLeft\}  (expected $\{c\}) — rounding\textasciigrave{}}
\NormalTok{      );}
\NormalTok{  \}}

\NormalTok{  return svg.node();}
\NormalTok{\}}
\end{Highlighting}
\end{Shaded}

\begin{center}\rule{0.5\linewidth}{0.5pt}\end{center}

\section{Worked examples}\label{worked-examples-6}

\textbf{Example 1.} You're saving up for a festival wristband. The
wristband costs £85. You already have £20 saved and earn £13 per hour
doing odd jobs. How many hours do you need to work?

Let \(h\) = hours worked.

\[20 + 13h = 85\]

Subtract 20 from both sides --- to clear the fixed amount already saved:

\[13h = 65\]

Divide both sides by 13 --- to get from ``13 of something'' to ``one of
something'':

\[h = 5\]

Five hours. Check: \(20 + 13 \times 5 = 20 + 65 = 85\). Yes.

\begin{center}\rule{0.5\linewidth}{0.5pt}\end{center}

\textbf{Example 2.} Two streaming services. Service A costs £8 per month
plus £2 per film you rent. Service B costs £18 per month with unlimited
films. How many film rentals per month makes the two services equal in
cost?

Let \(f\) = films rented per month.

Service A costs: \(8 + 2f\)

Service B costs: \(18\)

They're equal when:

\[8 + 2f = 18\]

\[2f = 10\]

\[f = 5\]

At 5 films per month they cost the same. Above 5 films, Service B is
cheaper.

\begin{center}\rule{0.5\linewidth}{0.5pt}\end{center}

\textbf{Example 3.} You're making a fruit drink by mixing juice
concentrate with water. The concentrate is 60\% real juice. You want the
final drink to be 20\% real juice. You start with 1 litre of
concentrate. How much water do you add?

The amount of real juice stays fixed at \(0.60 \times 1 = 0.6\) litres.
After adding \(w\) litres of water, the total volume is \(1 + w\)
litres, and you want the juice fraction to be 20\%:

\[\frac{0.6}{1 + w} = 0.2\]

Multiply both sides by \((1 + w)\) --- to clear the fraction:

\[0.6 = 0.2(1 + w)\]

Expand the right side:

\[0.6 = 0.2 + 0.2w\]

Subtract 0.2 from both sides:

\[0.4 = 0.2w\]

Divide both sides by 0.2:

\[w = 2\]

Add 2 litres of water.

\section{Where this goes}\label{where-this-goes-6}

This chapter covered one unknown on one side. The next two chapters in
this volume --- inequalities and graphing --- extend the same idea in
different directions: inequalities ask which values \emph{satisfy} a
relationship rather than pinning down one exact answer, and graphing
shows what happens when you let the unknown vary continuously. The
further horizon worth knowing about is \textbf{systems of equations}
(Volume 4): two unknowns, two equations, solved together. The method is
the same balancing act you've already learned, applied twice at once.

\begin{tcolorbox}[enhanced jigsaw, opacityback=0, colback=white, breakable, colframe=quarto-callout-tip-color-frame, arc=.35mm, left=2mm, toprule=.15mm, leftrule=.75mm, bottomrule=.15mm, rightrule=.15mm]
\begin{minipage}[t]{5.5mm}
\textcolor{quarto-callout-tip-color}{\faLightbulb}
\end{minipage}%
\begin{minipage}[t]{\textwidth - 5.5mm}

\textbf{Where this shows up}

\begin{itemize}
\tightlist
\item
  An electrical engineer writing Kirchhoff's voltage law is solving a
  linear equation for current or voltage.
\item
  A nutritionist calculating how much of an ingredient hits a target
  percentage is solving a linear equation.
\item
  A financial analyst finding a break-even point is solving a linear
  equation.
\item
  A data scientist fitting a regression line is solving a system of
  linear equations.
\end{itemize}

The tool is the same. The domain changes.

\end{minipage}%
\end{tcolorbox}

\section{Exercises}\label{exercises-6}

These are puzzles, not drills. Each one has a clean answer, but the
interesting part is setting up the equation --- translating the
situation into the notation before you solve it.

\begin{Shaded}
\begin{Highlighting}[]
\NormalTok{//| echo: false}
\NormalTok{katex = require("katex@0.16")}
\end{Highlighting}
\end{Shaded}

\begin{Shaded}
\begin{Highlighting}[]
\NormalTok{//| echo: false}
\NormalTok{function makeStepperHTML(exerciseNum, steps) \{}
\NormalTok{  let current = 0;}
\NormalTok{  const totalSteps = steps.length;}
\NormalTok{  const container = document.createElement("div");}
\NormalTok{  container.style.cssText = "border:1px solid \#e5e7eb; border{-}radius:8px; padding:1rem 1.25rem; margin:0.75rem 0; font{-}family:inherit;";}
\NormalTok{  const stepsDiv = document.createElement("div");}
\NormalTok{  stepsDiv.style.marginBottom = "0.75rem";}

\NormalTok{  function renderSteps() \{}
\NormalTok{    stepsDiv.innerHTML = "";}
\NormalTok{    for (let i = 0; i \textless{} current; i++) \{}
\NormalTok{      const s = steps[i];}
\NormalTok{      const row = document.createElement("div");}
\NormalTok{      row.style.cssText = "display:grid; grid{-}template{-}columns:200px 1fr; gap:0.35rem 1rem; align{-}items:baseline; padding:0.35rem 0; border{-}top:1px solid \#f3f4f6;";}
\NormalTok{      const opCell = document.createElement("span");}
\NormalTok{      opCell.style.cssText = "font{-}size:0.85em; color:\#6b7280; font{-}style:italic;";}
\NormalTok{      opCell.textContent = s.op;}
\NormalTok{      const eqCell = document.createElement("span");}
\NormalTok{      eqCell.innerHTML = katex.renderToString(s.eq, \{ throwOnError: false, displayMode: false \});}
\NormalTok{      row.appendChild(opCell);}
\NormalTok{      row.appendChild(eqCell);}
\NormalTok{      if (s.note) \{}
\NormalTok{        const noteRow = document.createElement("div");}
\NormalTok{        noteRow.style.cssText = "grid{-}column:2; font{-}size:0.82em; color:\#6b7280; padding{-}bottom:0.2rem;";}
\NormalTok{        noteRow.textContent = s.note;}
\NormalTok{        row.appendChild(noteRow);}
\NormalTok{      \}}
\NormalTok{      stepsDiv.appendChild(row);}
\NormalTok{    \}}
\NormalTok{    if (current === totalSteps) \{}
\NormalTok{      const done = document.createElement("div");}
\NormalTok{      done.style.cssText = "margin{-}top:0.5rem; font{-}size:0.9em; color:\#059669; font{-}weight:500;";}
\NormalTok{      done.textContent = "✓ Solution complete";}
\NormalTok{      stepsDiv.appendChild(done);}
\NormalTok{    \}}
\NormalTok{  \}}

\NormalTok{  const controls = document.createElement("div");}
\NormalTok{  controls.style.cssText = "display:flex; gap:0.5rem; align{-}items:center;";}
\NormalTok{  const nextBtn = document.createElement("button");}
\NormalTok{  nextBtn.textContent = "Next step →";}
\NormalTok{  nextBtn.style.cssText = "padding:0.35rem 0.85rem; border:1px solid \#d1d5db; border{-}radius:4px; background:\#fff; cursor:pointer; font{-}size:0.9em;";}
\NormalTok{  const resetBtn = document.createElement("button");}
\NormalTok{  resetBtn.textContent = "Reset";}
\NormalTok{  resetBtn.style.cssText = "padding:0.35rem 0.75rem; border:1px solid \#e5e7eb; border{-}radius:4px; background:\#f9fafb; cursor:pointer; font{-}size:0.9em; color:\#6b7280;";}
\NormalTok{  const counter = document.createElement("span");}
\NormalTok{  counter.style.cssText = "font{-}size:0.82em; color:\#9ca3af; margin{-}left:0.25rem;";}

\NormalTok{  function updateButtons() \{}
\NormalTok{    nextBtn.disabled = current === totalSteps;}
\NormalTok{    nextBtn.style.opacity = current === totalSteps ? "0.4" : "1";}
\NormalTok{    counter.textContent = current === 0 ? "Click to reveal steps" : \textasciigrave{}Step $\{current\} of $\{totalSteps\}\textasciigrave{};}
\NormalTok{  \}}

\NormalTok{  nextBtn.onclick = () =\textgreater{} \{ if (current \textless{} totalSteps) \{ current++; renderSteps(); updateButtons(); \} \};}
\NormalTok{  resetBtn.onclick = () =\textgreater{} \{ current = 0; renderSteps(); updateButtons(); \};}
\NormalTok{  controls.appendChild(nextBtn);}
\NormalTok{  controls.appendChild(resetBtn);}
\NormalTok{  controls.appendChild(counter);}
\NormalTok{  container.appendChild(stepsDiv);}
\NormalTok{  container.appendChild(controls);}
\NormalTok{  renderSteps();}
\NormalTok{  updateButtons();}
\NormalTok{  return container;}
\NormalTok{\}}
\end{Highlighting}
\end{Shaded}

\begin{enumerate}
\def\labelenumi{\arabic{enumi}.}
\tightlist
\item
  A bus pass costs £1.50 to buy, then £0.90 per journey. You spend £9.60
  in total. How many journeys did you take?
\end{enumerate}

\begin{Shaded}
\begin{Highlighting}[]
\NormalTok{//| echo: false}
\NormalTok{makeStepperHTML(1, [}
\NormalTok{  \{ op: "Write as an equation",           eq: "1.50 + 0.90j = 9.60" \},}
\NormalTok{  \{ op: "Subtract 1.50 from both sides",  eq: "0.90j = 8.10" \},}
\NormalTok{  \{ op: "Divide both sides by 0.90",      eq: "j = 9" \},}
\NormalTok{  \{ op: "Check",                           eq: "1.50 + 0.90 \textbackslash{}\textbackslash{}times 9 = 1.50 + 8.10 = 9.60 \textbackslash{}\textbackslash{}checkmark" \}}
\NormalTok{])}
\end{Highlighting}
\end{Shaded}

\begin{enumerate}
\def\labelenumi{\arabic{enumi}.}
\setcounter{enumi}{1}
\tightlist
\item
  A recipe makes 12 cookies using 180g of flour. You have 480g of flour.
  How many cookies can you make? \emph{(Set up an equation first. Don't
  just divide.)}
\end{enumerate}

\begin{Shaded}
\begin{Highlighting}[]
\NormalTok{//| echo: false}
\NormalTok{makeStepperHTML(2, [}
\NormalTok{  \{ op: "Write as an equation",          eq: "\textbackslash{}\textbackslash{}frac\{180\}\{12\} = \textbackslash{}\textbackslash{}frac\{480\}\{c\}", note: "Flour per cookie is constant, so set up equal ratios." \},}
\NormalTok{  \{ op: "Simplify left side",             eq: "15 = \textbackslash{}\textbackslash{}frac\{480\}\{c\}" \},}
\NormalTok{  \{ op: "Multiply both sides by c",       eq: "15c = 480" \},}
\NormalTok{  \{ op: "Divide both sides by 15",        eq: "c = 32" \},}
\NormalTok{  \{ op: "Check",                          eq: "\textbackslash{}\textbackslash{}frac\{180\}\{12\} = 15 \textbackslash{}\textbackslash{}text\{ g/cookie\},\textbackslash{}\textbackslash{}quad 15 \textbackslash{}\textbackslash{}times 32 = 480 \textbackslash{}\textbackslash{}checkmark" \}}
\NormalTok{])}
\end{Highlighting}
\end{Shaded}

\begin{enumerate}
\def\labelenumi{\arabic{enumi}.}
\setcounter{enumi}{2}
\tightlist
\item
  Two cyclists start 120km apart and ride toward each other. One rides
  at 20km/h, the other at 25km/h. After how many hours do they meet?
\end{enumerate}

\begin{Shaded}
\begin{Highlighting}[]
\NormalTok{//| echo: false}
\NormalTok{makeStepperHTML(3, [}
\NormalTok{  \{ op: "Write as an equation",           eq: "20t + 25t = 120", note: "Together they close 45 km each hour, starting 120 km apart." \},}
\NormalTok{  \{ op: "Collect like terms",             eq: "45t = 120" \},}
\NormalTok{  \{ op: "Divide both sides by 45",        eq: "t = \textbackslash{}\textbackslash{}frac\{120\}\{45\} = \textbackslash{}\textbackslash{}frac\{8\}\{3\} \textbackslash{}\textbackslash{}approx 2.67 \textbackslash{}\textbackslash{}text\{ h\}" \},}
\NormalTok{  \{ op: "Check",                          eq: "20 \textbackslash{}\textbackslash{}times \textbackslash{}\textbackslash{}tfrac\{8\}\{3\} + 25 \textbackslash{}\textbackslash{}times \textbackslash{}\textbackslash{}tfrac\{8\}\{3\} = \textbackslash{}\textbackslash{}tfrac\{160\}\{3\} + \textbackslash{}\textbackslash{}tfrac\{200\}\{3\} = \textbackslash{}\textbackslash{}tfrac\{360\}\{3\} = 120 \textbackslash{}\textbackslash{}checkmark" \}}
\NormalTok{])}
\end{Highlighting}
\end{Shaded}

When two quantities are changing in opposite directions and you want to
know when they're equal, the unknown ends up on both sides of the
equation. The fix is simple: collect all the terms with the unknown on
one side, and all the plain numbers on the other.

For example, if one thing is shrinking and another is growing, you might
get something like \(500 - 35t = 200 + 15t\). Subtract \(15t\) from both
sides to collect the variable terms: \(500 - 50t = 200\). Then subtract
200: \(300 = 50t\). Divide by 50: \(t = 6\). The same balance rule as
before --- it just takes one extra step to gather the unknowns first.

\begin{enumerate}
\def\labelenumi{\arabic{enumi}.}
\setcounter{enumi}{3}
\tightlist
\item
  A water tank holds 800 litres and drains at 35 litres per minute. A
  second tank holds 200 litres and fills at 15 litres per minute. After
  how many minutes do they contain the same amount?
\end{enumerate}

\begin{Shaded}
\begin{Highlighting}[]
\NormalTok{//| echo: false}
\NormalTok{makeStepperHTML(4, [}
\NormalTok{  \{ op: "Write as an equation",           eq: "800 {-} 35t = 200 + 15t", note: "Tank 1 decreases; tank 2 increases. Set them equal." \},}
\NormalTok{  \{ op: "Subtract 15t from both sides",   eq: "800 {-} 50t = 200" \},}
\NormalTok{  \{ op: "Subtract 800 from both sides",   eq: "{-}50t = {-}600" \},}
\NormalTok{  \{ op: "Divide both sides by −50",       eq: "t = 12" \},}
\NormalTok{  \{ op: "Check",                          eq: "800 {-} 35 \textbackslash{}\textbackslash{}times 12 = 800 {-} 420 = 380;\textbackslash{}\textbackslash{}quad 200 + 15 \textbackslash{}\textbackslash{}times 12 = 200 + 180 = 380 \textbackslash{}\textbackslash{}checkmark" \}}
\NormalTok{])}
\end{Highlighting}
\end{Shaded}

\begin{enumerate}
\def\labelenumi{\arabic{enumi}.}
\setcounter{enumi}{4}
\tightlist
\item
  You and two friends are splitting the cost of a group gift. The total
  cost is £96, but one friend already put in £18 as a deposit. How much
  does each of the three of you pay now?
\end{enumerate}

\begin{Shaded}
\begin{Highlighting}[]
\NormalTok{//| echo: false}
\NormalTok{makeStepperHTML(5, [}
\NormalTok{  \{ op: "Write as an equation",           eq: "3p + 18 = 96", note: "Three people each pay p, plus the £18 already paid, equals £96." \},}
\NormalTok{  \{ op: "Subtract 18 from both sides",    eq: "3p = 78" \},}
\NormalTok{  \{ op: "Divide both sides by 3",         eq: "p = 26" \},}
\NormalTok{  \{ op: "Check",                          eq: "3 \textbackslash{}\textbackslash{}times 26 + 18 = 78 + 18 = 96 \textbackslash{}\textbackslash{}checkmark" \}}
\NormalTok{])}
\end{Highlighting}
\end{Shaded}

\chapter{Inequalities}\label{inequalities}

The range of answers that still work

\hfill\break

Your monthly budget for food is £120. There are dozens of ways to spend
it --- some that work, many that don't. The question isn't ``what
exactly do I spend?'' It's ``what combinations are even possible?''

A shelf is rated to hold 50 kg. The boxes you're stacking weigh 8 kg
each. You don't need the exact weight --- you need to know how many
boxes you can put on before it becomes unsafe.

A password field requires between 8 and 20 characters. No single length
is the answer. Any length in that range works.

These are all inequalities. An equation gives you one answer. An
inequality gives you a whole region of answers --- and the boundary that
separates ``works'' from ``doesn't work.''

\section{What the notation is
saying}\label{what-the-notation-is-saying-7}

An equation uses \(=\) to say two things are exactly equal. An
inequality uses a different symbol to say one thing is \emph{bigger} or
\emph{smaller} than another.

The four symbols:

\begin{longtable}[]{@{}lll@{}}
\toprule\noalign{}
Symbol & Reads as & Example \\
\midrule\noalign{}
\endhead
\bottomrule\noalign{}
\endlastfoot
\(<\) & is less than & \(3 < 7\) \\
\(>\) & is greater than & \(12 > 5\) \\
\(\leq\) & is less than or equal to & \(x \leq 10\) \\
\(\geq\) & is greater than or equal to & \(n \geq 4\) \\
\end{longtable}

The pointy end of the symbol always faces the smaller value. So
\(3 < 7\) reads ``3 is less than 7'' --- the small end points at 3, the
open end opens toward 7. If you forget which way round it goes, just
picture the symbol as an arrow squishing the smaller number.

The difference between \(<\) and \(\leq\) matters. If a shelf holds
\emph{at most} 50 kg, exactly 50 kg is fine --- so you'd write
\(w \leq 50\). If something must be \emph{strictly less than} 50,
exactly 50 is not fine --- so you'd write \(w < 50\).

A \textbf{compound inequality} combines two inequalities into one
statement. \(8 \leq n \leq 20\) reads ``n is at least 8 and at most 20''
--- both conditions must hold at once. This is the password length rule:
any whole number from 8 to 20 satisfies it.

\section{The method}\label{the-method-7}

Solving an inequality is almost identical to solving an equation. You do
the same things: add, subtract, multiply, divide --- keeping both sides
balanced until the unknown is isolated.

\textbf{There is one exception. It's the only genuinely surprising thing
in this chapter. It deserves its own section.}

\subsection{The sign-flip rule}\label{the-sign-flip-rule}

When you multiply or divide both sides of an inequality by a
\textbf{negative number}, the direction of the inequality reverses.

Here's why, shown directly. Start with something you know is true:

\[3 < 5\]

Three is less than five. Yes. Now multiply both sides by \(-1\):

\[-3 \quad \text{and} \quad -5\]

Which of these is smaller? \(-5\) is. It's further left on the number
line. So the relationship has flipped:

\[-3 > -5\]

The original said ``3 is less than 5.'' After multiplying by \(-1\), the
true statement is ``\(-3\) is greater than \(-5\).''

This isn't a rule someone invented. It's what actually happens on the
number line when you negate both sides. Multiplying by \(-1\) reflects
every number across zero --- what was to the left of zero ends up to the
right, and the ordering reverses.

\begin{tcolorbox}[enhanced jigsaw, opacityback=0, colback=white, breakable, colframe=quarto-callout-warning-color-frame, arc=.35mm, left=2mm, toprule=.15mm, leftrule=.75mm, bottomrule=.15mm, rightrule=.15mm]
\begin{minipage}[t]{5.5mm}
\textcolor{quarto-callout-warning-color}{\faExclamationTriangle}
\end{minipage}%
\begin{minipage}[t]{\textwidth - 5.5mm}

\textbf{The sign-flip rule}

Multiplying or dividing both sides of an inequality by a \emph{negative}
number reverses the direction of the inequality sign.

\[3 < 5 \implies -3 > -5\]

If you forget to flip, your answer will point the wrong way --- and
every value in your ``solution'' will actually be wrong.

\end{minipage}%
\end{tcolorbox}

\subsection{Solving a linear inequality, step by
step}\label{solving-a-linear-inequality-step-by-step}

\[2x - 3 < 7\]

Add 3 to both sides (adding is safe --- no sign flip):

\[2x < 10\]

Divide both sides by 2 (dividing by a \emph{positive} number --- no flip
needed):

\[x < 5\]

The solution is every number less than 5. Not just one number: all of
them.

Now compare with one that triggers the flip:

\[-3x + 1 \geq 10\]

Subtract 1 from both sides:

\[-3x \geq 9\]

Divide both sides by \(-3\). This is negative --- flip the sign:

\[x \leq -3\]

The solution is every number less than or equal to \(-3\). Check: try
\(x = -4\). Then \(-3(-4) + 1 = 12 + 1 = 13\), and \(13 \geq 10\). Yes.
Try \(x = 0\). Then \(-3(0) + 1 = 1\), and \(1 \geq 10\). No --- and
\(x = 0\) is not in our solution set, so that's correct.

\begin{tcolorbox}[enhanced jigsaw, opacityback=0, colback=white, breakable, colframe=quarto-callout-note-color-frame, arc=.35mm, left=2mm, toprule=.15mm, leftrule=.75mm, bottomrule=.15mm, rightrule=.15mm]
\begin{minipage}[t]{5.5mm}
\textcolor{quarto-callout-note-color}{\faInfo}
\end{minipage}%
\begin{minipage}[t]{\textwidth - 5.5mm}

\textbf{Why solving inequalities works}

An inequality is a statement about the relationship between two
quantities. Doing the same operation to both sides preserves that
relationship --- as long as the operation doesn't reverse the ordering.
Addition and subtraction never reverse ordering. Multiplication and
division by a \emph{positive} number preserve ordering. Multiplication
or division by a \emph{negative} number reverses it. That's the whole
story.

\end{minipage}%
\end{tcolorbox}

\subsection{Reading the solution on a number
line}\label{reading-the-solution-on-a-number-line}

A number line makes the solution visible. You mark the boundary value
and shade the region that satisfies the inequality.

\textbf{Open circle} (\(\circ\)) at the boundary: the boundary value is
\emph{not} included. Use this for strict inequalities (\(<\) and \(>\)).

\textbf{Closed circle} (\(\bullet\)) at the boundary: the boundary value
\emph{is} included. Use this for \(\leq\) and \(\geq\).

For \(x < 5\): draw an open circle at 5, shade to the left.

For \(x \leq 5\): draw a closed circle at 5, shade to the left.

For \(x > -3\): draw an open circle at \(-3\), shade to the right.

For a compound inequality like \(-2 < x \leq 4\): open circle at \(-2\),
closed circle at \(4\), shade the segment between them.

\textbf{Interactive: Number line inequality explorer.} Use the slider to
set the boundary value. Choose the inequality direction and whether the
boundary is included. The shaded region and circle type update live.

\begin{Shaded}
\begin{Highlighting}[]
\NormalTok{//| echo: false}
\NormalTok{viewof nlBoundary = Inputs.range([{-}5, 20], \{}
\NormalTok{  label: "Boundary value",}
\NormalTok{  step: 0.5,}
\NormalTok{  value: 5}
\NormalTok{\})}
\end{Highlighting}
\end{Shaded}

\begin{Shaded}
\begin{Highlighting}[]
\NormalTok{//| echo: false}
\NormalTok{viewof nlSymbol = Inputs.select(["\textless{}", "≤", "\textgreater{}", "≥"], \{}
\NormalTok{  label: "Inequality symbol  (x  \_\_ boundary)",}
\NormalTok{  value: "\textless{}"}
\NormalTok{\})}
\end{Highlighting}
\end{Shaded}

\begin{Shaded}
\begin{Highlighting}[]
\NormalTok{//| echo: false}
\NormalTok{\{}
\NormalTok{  const boundary = nlBoundary;}
\NormalTok{  const sym = nlSymbol;}
\NormalTok{  const inclusive = sym === "≤" || sym === "≥";}
\NormalTok{  const shadedLeft = sym === "\textless{}" || sym === "≤";}

\NormalTok{  const W = 560, H = 110;}
\NormalTok{  const PAD = 40;}
\NormalTok{  const AXIS\_Y = 60;}
\NormalTok{  const TICK\_H = 8;}
\NormalTok{  const domainMin = {-}5, domainMax = 20;}
\NormalTok{  const scale = (v) =\textgreater{} PAD + (v {-} domainMin) / (domainMax {-} domainMin) * (W {-} 2 * PAD);}

\NormalTok{  const svg = d3.create("svg")}
\NormalTok{    .attr("viewBox", \textasciigrave{}0 0 $\{W\} $\{H\}\textasciigrave{})}
\NormalTok{    .attr("width", "100\%")}
\NormalTok{    .attr("style", \textasciigrave{}max{-}width:$\{W\}px; font{-}family:inherit;\textasciigrave{});}

\NormalTok{  // Axis line}
\NormalTok{  svg.append("line")}
\NormalTok{    .attr("x1", PAD).attr("y1", AXIS\_Y)}
\NormalTok{    .attr("x2", W {-} PAD).attr("y2", AXIS\_Y)}
\NormalTok{    .attr("stroke", "\#9ca3af").attr("stroke{-}width", 2);}

\NormalTok{  // Arrow heads}
\NormalTok{  svg.append("polygon")}
\NormalTok{    .attr("points", \textasciigrave{}$\{W {-} PAD\},$\{AXIS\_Y\} $\{W {-} PAD {-} 8\},$\{AXIS\_Y {-} 4\} $\{W {-} PAD {-} 8\},$\{AXIS\_Y + 4\}\textasciigrave{})}
\NormalTok{    .attr("fill", "\#9ca3af");}
\NormalTok{  svg.append("polygon")}
\NormalTok{    .attr("points", \textasciigrave{}$\{PAD\},$\{AXIS\_Y\} $\{PAD + 8\},$\{AXIS\_Y {-} 4\} $\{PAD + 8\},$\{AXIS\_Y + 4\}\textasciigrave{})}
\NormalTok{    .attr("fill", "\#9ca3af");}

\NormalTok{  // Tick marks and labels for integers in domain}
\NormalTok{  for (let v = domainMin; v \textless{}= domainMax; v += 5) \{}
\NormalTok{    const x = scale(v);}
\NormalTok{    svg.append("line")}
\NormalTok{      .attr("x1", x).attr("y1", AXIS\_Y {-} TICK\_H / 2)}
\NormalTok{      .attr("x2", x).attr("y2", AXIS\_Y + TICK\_H / 2)}
\NormalTok{      .attr("stroke", "\#9ca3af").attr("stroke{-}width", 1.5);}
\NormalTok{    svg.append("text")}
\NormalTok{      .attr("x", x).attr("y", AXIS\_Y + 22)}
\NormalTok{      .attr("text{-}anchor", "middle")}
\NormalTok{      .attr("font{-}size", "11px").attr("fill", "\#6b7280")}
\NormalTok{      .text(v);}
\NormalTok{  \}}

\NormalTok{  // Shaded region}
\NormalTok{  const bx = scale(boundary);}
\NormalTok{  const shadeColor = "\#0d9488";}
\NormalTok{  if (shadedLeft) \{}
\NormalTok{    svg.append("rect")}
\NormalTok{      .attr("x", PAD).attr("y", AXIS\_Y {-} 6)}
\NormalTok{      .attr("width", Math.max(0, bx {-} PAD)).attr("height", 12)}
\NormalTok{      .attr("fill", shadeColor).attr("opacity", 0.25);}
\NormalTok{    // Arrow showing infinite extent to the left}
\NormalTok{    svg.append("polygon")}
\NormalTok{      .attr("points", \textasciigrave{}$\{PAD\},$\{AXIS\_Y\} $\{PAD + 8\},$\{AXIS\_Y {-} 4\} $\{PAD + 8\},$\{AXIS\_Y + 4\}\textasciigrave{})}
\NormalTok{      .attr("fill", shadeColor).attr("opacity", 0.6);}
\NormalTok{  \} else \{}
\NormalTok{    svg.append("rect")}
\NormalTok{      .attr("x", bx).attr("y", AXIS\_Y {-} 6)}
\NormalTok{      .attr("width", Math.max(0, W {-} PAD {-} bx)).attr("height", 12)}
\NormalTok{      .attr("fill", shadeColor).attr("opacity", 0.25);}
\NormalTok{    // Arrow showing infinite extent to the right}
\NormalTok{    svg.append("polygon")}
\NormalTok{      .attr("points", \textasciigrave{}$\{W {-} PAD\},$\{AXIS\_Y\} $\{W {-} PAD {-} 8\},$\{AXIS\_Y {-} 4\} $\{W {-} PAD {-} 8\},$\{AXIS\_Y + 4\}\textasciigrave{})}
\NormalTok{      .attr("fill", shadeColor).attr("opacity", 0.6);}
\NormalTok{  \}}

\NormalTok{  // Boundary circle}
\NormalTok{  if (inclusive) \{}
\NormalTok{    svg.append("circle")}
\NormalTok{      .attr("cx", bx).attr("cy", AXIS\_Y)}
\NormalTok{      .attr("r", 7).attr("fill", shadeColor).attr("stroke", shadeColor).attr("stroke{-}width", 2);}
\NormalTok{  \} else \{}
\NormalTok{    svg.append("circle")}
\NormalTok{      .attr("cx", bx).attr("cy", AXIS\_Y)}
\NormalTok{      .attr("r", 7).attr("fill", "\#fff").attr("stroke", shadeColor).attr("stroke{-}width", 2.5);}
\NormalTok{  \}}

\NormalTok{  // Boundary label}
\NormalTok{  svg.append("text")}
\NormalTok{    .attr("x", bx).attr("y", AXIS\_Y {-} 18)}
\NormalTok{    .attr("text{-}anchor", "middle")}
\NormalTok{    .attr("font{-}size", "12px").attr("fill", shadeColor).attr("font{-}weight", "600")}
\NormalTok{    .text(boundary \% 1 === 0 ? boundary : boundary.toFixed(1));}

\NormalTok{  // Inequality statement}
\NormalTok{  const circleDesc = inclusive ? "closed circle (boundary included)" : "open circle (boundary excluded)";}
\NormalTok{  svg.append("text")}
\NormalTok{    .attr("x", W / 2).attr("y", H {-} 6)}
\NormalTok{    .attr("text{-}anchor", "middle")}
\NormalTok{    .attr("font{-}size", "11px").attr("fill", "\#6b7280")}
\NormalTok{    .text(\textasciigrave{}x $\{sym\} $\{boundary \% 1 === 0 ? boundary : boundary.toFixed(1)\}   —   $\{circleDesc\}\textasciigrave{});}

\NormalTok{  return svg.node();}
\NormalTok{\}}
\end{Highlighting}
\end{Shaded}

\section{Worked examples}\label{worked-examples-7}

\textbf{Example 1 --- Budgeting.} You have £120 to spend on food this
month. You've already spent £47. You want to buy some items that cost £8
each. How many can you buy without going over budget?

Let \(n\) = number of items.

\[47 + 8n \leq 120\]

Subtract 47 from both sides:

\[8n \leq 73\]

Divide both sides by 8 (positive --- no flip):

\[n \leq 9.125\]

Since \(n\) must be a whole number, the answer is \(n \leq 9\). You can
buy at most 9 items.

Check: \(47 + 8 \times 9 = 47 + 72 = 119\). Within budget. And
\(47 + 8 \times 10 = 127\). Over. Correct.

\begin{center}\rule{0.5\linewidth}{0.5pt}\end{center}

\textbf{Example 2 --- Shelf load.} A wall-mounted shelf is rated to hold
at most 50 kg. Each box you place on it weighs 8 kg.

How many boxes can the shelf safely hold?

The shelf load rating tells you how much weight the shelf can carry on
top of it. The brackets that mount the shelf to the wall bear their own
weight separately --- they don't count against the shelf's carrying
capacity.

Let \(b\) = number of boxes.

\[8b \leq 50\]

\[b \leq 6.25\]

Since \(b\) must be a whole number: at most 6 boxes.

Check: \(8 \times 6 = 48 \leq 50\). ✓ Seven boxes:
\(8 \times 7 = 56 > 50\). ✗

\begin{center}\rule{0.5\linewidth}{0.5pt}\end{center}

\textbf{Example 3 --- Password validation.} A website requires passwords
to be at least 8 characters long but no longer than 20 characters. A
user types a password. The code checks whether its length \(n\) is
acceptable.

The valid range is:

\[8 \leq n \leq 20\]

Any value of \(n\) satisfying both conditions is acceptable. \(n = 7\)?
Too short. \(n = 8\)? Fine --- the lower boundary is included.
\(n = 20\)? Fine. \(n = 21\)? Too long.

This is a compound inequality. In code, you'd check both conditions at
once: \texttt{8\ \textless{}=\ n\ \&\&\ n\ \textless{}=\ 20} (or the
equivalent in whatever language you're using). The mathematics and the
code are the same structure.

\begin{center}\rule{0.5\linewidth}{0.5pt}\end{center}

\textbf{Example 4 --- An inequality with a negative coefficient.} The
temperature in a storage room must be kept strictly below
\(5°\text{C}\). The thermostat setting \(t\) (in °C) satisfies:

\[-2t + 14 > 4\]

Subtract 14 from both sides:

\[-2t > -10\]

Divide by \(-2\) --- negative, so flip:

\[t < 5\]

The thermostat must be set below 5°C. The sign flip was necessary here:
if you forgot it and wrote \(t > 5\), you'd be saying temperatures
\emph{above} 5°C are acceptable, which is exactly wrong.

\textbf{Interactive: Inequality explorer.} Set the coefficients \(a\)
and \(b\), choose the constant \(c\) and the inequality direction. The
solution is shown algebraically and on a number line. The sign-flip
warning lights up when \(a\) is negative.

\begin{Shaded}
\begin{Highlighting}[]
\NormalTok{//| echo: false}
\NormalTok{viewof ieA = Inputs.range([{-}5, 5], \{ label: "a  (coefficient of x)", step: 1, value: 2 \})}
\end{Highlighting}
\end{Shaded}

\begin{Shaded}
\begin{Highlighting}[]
\NormalTok{//| echo: false}
\NormalTok{viewof ieB = Inputs.range([{-}10, 10], \{ label: "b  (constant, left side)", step: 1, value: {-}3 \})}
\end{Highlighting}
\end{Shaded}

\begin{Shaded}
\begin{Highlighting}[]
\NormalTok{//| echo: false}
\NormalTok{viewof ieC = Inputs.range([{-}20, 20], \{ label: "c  (right{-}hand side)", step: 1, value: 7 \})}
\end{Highlighting}
\end{Shaded}

\begin{Shaded}
\begin{Highlighting}[]
\NormalTok{//| echo: false}
\NormalTok{viewof ieSym = Inputs.select(["\textless{}", "≤", "\textgreater{}", "≥"], \{ label: "Inequality symbol", value: "\textless{}" \})}
\end{Highlighting}
\end{Shaded}

\begin{Shaded}
\begin{Highlighting}[]
\NormalTok{//| echo: false}
\NormalTok{\{}
\NormalTok{  const a = Math.round(ieA), b = Math.round(ieB), c = Math.round(ieC), sym = ieSym;}

\NormalTok{  // Solve:  ax + b  sym  c   →   ax  sym  c {-} b   →   x  sym\textquotesingle{}  (c{-}b)/a}
\NormalTok{  // If a \textless{} 0, sym flips.}
\NormalTok{  const rhs = c {-} b;}
\NormalTok{  const needsFlip = a \textless{} 0;}
\NormalTok{  const flipped = \{ "\textless{}": "\textgreater{}", "≤": "≥", "\textgreater{}": "\textless{}", "≥": "≤" \};}
\NormalTok{  const solvedSym = needsFlip ? flipped[sym] : sym;}

\NormalTok{  const W = 560, H = 200;}
\NormalTok{  const PAD = 40;}
\NormalTok{  const AXIS\_Y = 150;}
\NormalTok{  const TICK\_H = 8;}

\NormalTok{  const leftStr = a === 1 ? \textasciigrave{}x + $\{b\}\textasciigrave{} : a === {-}1 ? \textasciigrave{}{-}x + $\{b\}\textasciigrave{} : \textasciigrave{}$\{a\}x + $\{b\}\textasciigrave{};}
\NormalTok{  const lhsClean = b \textgreater{}= 0}
\NormalTok{    ? (a === 1 ? \textasciigrave{}x + $\{b\}\textasciigrave{} : a === {-}1 ? \textasciigrave{}{-}x + $\{b\}\textasciigrave{} : \textasciigrave{}$\{a\}x + $\{b\}\textasciigrave{})}
\NormalTok{    : (a === 1 ? \textasciigrave{}x {-} $\{Math.abs(b)\}\textasciigrave{} : a === {-}1 ? \textasciigrave{}{-}x {-} $\{Math.abs(b)\}\textasciigrave{} : \textasciigrave{}$\{a\}x {-} $\{Math.abs(b)\}\textasciigrave{});}

\NormalTok{  const svg = d3.create("svg")}
\NormalTok{    .attr("viewBox", \textasciigrave{}0 0 $\{W\} $\{H\}\textasciigrave{})}
\NormalTok{    .attr("width", "100\%")}
\NormalTok{    .attr("style", \textasciigrave{}max{-}width:$\{W\}px; font{-}family:inherit;\textasciigrave{});}

\NormalTok{  // Step labels}
\NormalTok{  const steps = [}
\NormalTok{    \textasciigrave{}Original:  $\{lhsClean\}  $\{sym\}  $\{c\}\textasciigrave{},}
\NormalTok{    \textasciigrave{}Subtract $\{b\} from both sides:  $\{a === 1 ? "x" : a === {-}1 ? "{-}x" : \textasciigrave{}$\{a\}x\textasciigrave{}\}  $\{sym\}  $\{rhs\}\textasciigrave{},}
\NormalTok{    a === 0}
\NormalTok{      ? (rhs \textgreater{} 0 \&\& (sym === "\textless{}" || sym === "≤") ? "No solution (0x = 0, always)" : "All real numbers satisfy this")}
\NormalTok{      : needsFlip}
\NormalTok{        ? \textasciigrave{}Divide by $\{a\}  (negative — flip sign):  x  $\{solvedSym\}  $\{rhs / a\}\textasciigrave{}}
\NormalTok{        : \textasciigrave{}Divide by $\{a\}:  x  $\{solvedSym\}  $\{rhs / a\}\textasciigrave{}}
\NormalTok{  ];}

\NormalTok{  steps.forEach((txt, i) =\textgreater{} \{}
\NormalTok{    const isFlipStep = i === 2 \&\& needsFlip;}
\NormalTok{    svg.append("text")}
\NormalTok{      .attr("x", 16).attr("y", 24 + i * 28)}
\NormalTok{      .attr("font{-}size", "12px")}
\NormalTok{      .attr("fill", isFlipStep ? "\#b45309" : i === steps.length {-} 1 ? "\#0d9488" : "\#374151")}
\NormalTok{      .attr("font{-}weight", isFlipStep || i === steps.length {-} 1 ? "600" : "400")}
\NormalTok{      .text(txt);}
\NormalTok{  \});}

\NormalTok{  if (needsFlip) \{}
\NormalTok{    svg.append("text")}
\NormalTok{      .attr("x", 16).attr("y", 24 + steps.length * 28)}
\NormalTok{      .attr("font{-}size", "11px").attr("fill", "\#b45309")}
\NormalTok{      .text("⚠ Sign flip applied — dividing by a negative number reverses the inequality.");}
\NormalTok{  \}}

\NormalTok{  if (a !== 0) \{}
\NormalTok{    const solution = rhs / a;}
\NormalTok{    const domainMin = Math.floor(solution) {-} 8;}
\NormalTok{    const domainMax = Math.ceil(solution) + 8;}
\NormalTok{    const scale = (v) =\textgreater{} PAD + (v {-} domainMin) / (domainMax {-} domainMin) * (W {-} 2 * PAD);}
\NormalTok{    const bx = scale(solution);}
\NormalTok{    const shadedLeft = solvedSym === "\textless{}" || solvedSym === "≤";}
\NormalTok{    const inclusive = solvedSym === "≤" || solvedSym === "≥";}
\NormalTok{    const shadeColor = "\#0d9488";}

\NormalTok{    // Axis}
\NormalTok{    svg.append("line")}
\NormalTok{      .attr("x1", PAD).attr("y1", AXIS\_Y)}
\NormalTok{      .attr("x2", W {-} PAD).attr("y2", AXIS\_Y)}
\NormalTok{      .attr("stroke", "\#9ca3af").attr("stroke{-}width", 2);}

\NormalTok{    // Arrows}
\NormalTok{    svg.append("polygon")}
\NormalTok{      .attr("points", \textasciigrave{}$\{W {-} PAD\},$\{AXIS\_Y\} $\{W {-} PAD {-} 8\},$\{AXIS\_Y {-} 4\} $\{W {-} PAD {-} 8\},$\{AXIS\_Y + 4\}\textasciigrave{})}
\NormalTok{      .attr("fill", "\#9ca3af");}
\NormalTok{    svg.append("polygon")}
\NormalTok{      .attr("points", \textasciigrave{}$\{PAD\},$\{AXIS\_Y\} $\{PAD + 8\},$\{AXIS\_Y {-} 4\} $\{PAD + 8\},$\{AXIS\_Y + 4\}\textasciigrave{})}
\NormalTok{      .attr("fill", "\#9ca3af");}

\NormalTok{    // Ticks}
\NormalTok{    for (let v = Math.ceil(domainMin); v \textless{}= Math.floor(domainMax); v += Math.ceil((domainMax {-} domainMin) / 8)) \{}
\NormalTok{      const x = scale(v);}
\NormalTok{      svg.append("line")}
\NormalTok{        .attr("x1", x).attr("y1", AXIS\_Y {-} TICK\_H / 2)}
\NormalTok{        .attr("x2", x).attr("y2", AXIS\_Y + TICK\_H / 2)}
\NormalTok{        .attr("stroke", "\#9ca3af").attr("stroke{-}width", 1.5);}
\NormalTok{      svg.append("text")}
\NormalTok{        .attr("x", x).attr("y", AXIS\_Y + 20)}
\NormalTok{        .attr("text{-}anchor", "middle")}
\NormalTok{        .attr("font{-}size", "10px").attr("fill", "\#6b7280")}
\NormalTok{        .text(v);}
\NormalTok{    \}}

\NormalTok{    // Shade}
\NormalTok{    if (shadedLeft) \{}
\NormalTok{      svg.append("rect")}
\NormalTok{        .attr("x", PAD).attr("y", AXIS\_Y {-} 5)}
\NormalTok{        .attr("width", Math.max(0, bx {-} PAD)).attr("height", 10)}
\NormalTok{        .attr("fill", shadeColor).attr("opacity", 0.25);}
\NormalTok{    \} else \{}
\NormalTok{      svg.append("rect")}
\NormalTok{        .attr("x", bx).attr("y", AXIS\_Y {-} 5)}
\NormalTok{        .attr("width", Math.max(0, W {-} PAD {-} bx)).attr("height", 10)}
\NormalTok{        .attr("fill", shadeColor).attr("opacity", 0.25);}
\NormalTok{    \}}

\NormalTok{    // Boundary circle}
\NormalTok{    if (inclusive) \{}
\NormalTok{      svg.append("circle")}
\NormalTok{        .attr("cx", bx).attr("cy", AXIS\_Y)}
\NormalTok{        .attr("r", 6).attr("fill", shadeColor).attr("stroke", shadeColor).attr("stroke{-}width", 2);}
\NormalTok{    \} else \{}
\NormalTok{      svg.append("circle")}
\NormalTok{        .attr("cx", bx).attr("cy", AXIS\_Y)}
\NormalTok{        .attr("r", 6).attr("fill", "\#fff").attr("stroke", shadeColor).attr("stroke{-}width", 2.5);}
\NormalTok{    \}}

\NormalTok{    svg.append("text")}
\NormalTok{      .attr("x", bx).attr("y", AXIS\_Y {-} 14)}
\NormalTok{      .attr("text{-}anchor", "middle")}
\NormalTok{      .attr("font{-}size", "11px").attr("fill", shadeColor).attr("font{-}weight", "600")}
\NormalTok{      .text(Number.isInteger(solution) ? solution : solution.toFixed(2));}
\NormalTok{  \}}

\NormalTok{  return svg.node();}
\NormalTok{\}}
\end{Highlighting}
\end{Shaded}

\section{Where this goes}\label{where-this-goes-7}

Inequalities are where the idea of a \emph{constraint} enters
mathematics. An equation pins down one value. An inequality describes a
boundary --- everything on one side is feasible, everything on the other
isn't.

The most direct continuation is \textbf{linear programming} --- how to
find the best point within a region defined by multiple inequalities at
once. That's in the Optimisation chapter in Volume 7. A manufacturer
choosing how much of two products to make, subject to labour and
material constraints, is solving a linear programme. The feasible region
is defined entirely by inequalities like the ones in this chapter.

The second direction is \textbf{systems of inequalities} and their role
in applied linear algebra. When a machine learning model classifies
data, it's often drawing a boundary in a high-dimensional space --- one
side satisfies the classification condition, the other doesn't. That
boundary is expressed as an inequality. The linear equations you learned
in the previous chapter define the boundary; the inequality says which
side is which. That thread runs from here through linear algebra (Volume
4) all the way to optimisation and machine learning.

\begin{tcolorbox}[enhanced jigsaw, opacityback=0, colback=white, breakable, colframe=quarto-callout-tip-color-frame, arc=.35mm, left=2mm, toprule=.15mm, leftrule=.75mm, bottomrule=.15mm, rightrule=.15mm]
\begin{minipage}[t]{5.5mm}
\textcolor{quarto-callout-tip-color}{\faLightbulb}
\end{minipage}%
\begin{minipage}[t]{\textwidth - 5.5mm}

\textbf{Where this shows up}

\begin{itemize}
\tightlist
\item
  A financial analyst building a budget model sets upper and lower
  bounds on each line of spending --- those bounds are inequalities.
\item
  A structural engineer checks that the load on every beam stays within
  its rated capacity --- an inequality checked against the design at
  every point.
\item
  A software developer writing input validation checks that every
  user-entered value lies in an acceptable range --- a compound
  inequality executed in code.
\item
  A supply chain planner ensures stock levels stay between a minimum (to
  avoid running out) and a maximum (to avoid overflow) --- two
  inequalities, applied simultaneously.
\end{itemize}

\end{minipage}%
\end{tcolorbox}

\section{Exercises}\label{exercises-7}

Each problem has a clean answer. The work is in translating the
situation into the notation before you solve it.

\begin{Shaded}
\begin{Highlighting}[]
\NormalTok{//| echo: false}
\NormalTok{katex = require("katex@0.16")}
\end{Highlighting}
\end{Shaded}

\begin{Shaded}
\begin{Highlighting}[]
\NormalTok{//| echo: false}
\NormalTok{function makeStepperHTML(exerciseNum, steps) \{}
\NormalTok{  let current = 0;}
\NormalTok{  const totalSteps = steps.length;}
\NormalTok{  const container = document.createElement("div");}
\NormalTok{  container.style.cssText = "border:1px solid \#e5e7eb; border{-}radius:8px; padding:1rem 1.25rem; margin:0.75rem 0; font{-}family:inherit;";}
\NormalTok{  const stepsDiv = document.createElement("div");}
\NormalTok{  stepsDiv.style.marginBottom = "0.75rem";}

\NormalTok{  function renderSteps() \{}
\NormalTok{    stepsDiv.innerHTML = "";}
\NormalTok{    for (let i = 0; i \textless{} current; i++) \{}
\NormalTok{      const s = steps[i];}
\NormalTok{      const row = document.createElement("div");}
\NormalTok{      row.style.cssText = "display:grid; grid{-}template{-}columns:200px 1fr; gap:0.35rem 1rem; align{-}items:baseline; padding:0.35rem 0; border{-}top:1px solid \#f3f4f6;";}
\NormalTok{      const opCell = document.createElement("span");}
\NormalTok{      opCell.style.cssText = "font{-}size:0.85em; color:\#6b7280; font{-}style:italic;";}
\NormalTok{      opCell.textContent = s.op;}
\NormalTok{      const eqCell = document.createElement("span");}
\NormalTok{      eqCell.innerHTML = katex.renderToString(s.eq, \{ throwOnError: false, displayMode: false \});}
\NormalTok{      row.appendChild(opCell);}
\NormalTok{      row.appendChild(eqCell);}
\NormalTok{      if (s.note) \{}
\NormalTok{        const noteRow = document.createElement("div");}
\NormalTok{        noteRow.style.cssText = "grid{-}column:2; font{-}size:0.82em; color:\#6b7280; padding{-}bottom:0.2rem;";}
\NormalTok{        noteRow.textContent = s.note;}
\NormalTok{        row.appendChild(noteRow);}
\NormalTok{      \}}
\NormalTok{      stepsDiv.appendChild(row);}
\NormalTok{    \}}
\NormalTok{    if (current === totalSteps) \{}
\NormalTok{      const done = document.createElement("div");}
\NormalTok{      done.style.cssText = "margin{-}top:0.5rem; font{-}size:0.9em; color:\#059669; font{-}weight:500;";}
\NormalTok{      done.textContent = "✓ Solution complete";}
\NormalTok{      stepsDiv.appendChild(done);}
\NormalTok{    \}}
\NormalTok{  \}}

\NormalTok{  const controls = document.createElement("div");}
\NormalTok{  controls.style.cssText = "display:flex; gap:0.5rem; align{-}items:center;";}
\NormalTok{  const nextBtn = document.createElement("button");}
\NormalTok{  nextBtn.textContent = "Next step →";}
\NormalTok{  nextBtn.style.cssText = "padding:0.35rem 0.85rem; border:1px solid \#d1d5db; border{-}radius:4px; background:\#fff; cursor:pointer; font{-}size:0.9em;";}
\NormalTok{  const resetBtn = document.createElement("button");}
\NormalTok{  resetBtn.textContent = "Reset";}
\NormalTok{  resetBtn.style.cssText = "padding:0.35rem 0.75rem; border:1px solid \#e5e7eb; border{-}radius:4px; background:\#f9fafb; cursor:pointer; font{-}size:0.9em; color:\#6b7280;";}
\NormalTok{  const counter = document.createElement("span");}
\NormalTok{  counter.style.cssText = "font{-}size:0.82em; color:\#9ca3af; margin{-}left:0.25rem;";}

\NormalTok{  function updateButtons() \{}
\NormalTok{    nextBtn.disabled = current === totalSteps;}
\NormalTok{    nextBtn.style.opacity = current === totalSteps ? "0.4" : "1";}
\NormalTok{    counter.textContent = current === 0 ? "Click to reveal steps" : \textasciigrave{}Step $\{current\} of $\{totalSteps\}\textasciigrave{};}
\NormalTok{  \}}

\NormalTok{  nextBtn.onclick = () =\textgreater{} \{ if (current \textless{} totalSteps) \{ current++; renderSteps(); updateButtons(); \} \};}
\NormalTok{  resetBtn.onclick = () =\textgreater{} \{ current = 0; renderSteps(); updateButtons(); \};}
\NormalTok{  controls.appendChild(nextBtn);}
\NormalTok{  controls.appendChild(resetBtn);}
\NormalTok{  controls.appendChild(counter);}
\NormalTok{  container.appendChild(stepsDiv);}
\NormalTok{  container.appendChild(controls);}
\NormalTok{  renderSteps();}
\NormalTok{  updateButtons();}
\NormalTok{  return container;}
\NormalTok{\}}
\end{Highlighting}
\end{Shaded}

\begin{enumerate}
\def\labelenumi{\arabic{enumi}.}
\tightlist
\item
  You can spend at most £50 per month on streaming services. You already
  pay £14.97 for two services. How many additional services at £9.99 per
  month can you add and still stay within your budget?
\end{enumerate}

\begin{Shaded}
\begin{Highlighting}[]
\NormalTok{//| echo: false}
\NormalTok{makeStepperHTML(1, [}
\NormalTok{  \{ op: "Write the inequality",            eq: "14.97 + 9.99s \textbackslash{}\textbackslash{}leq 50" \},}
\NormalTok{  \{ op: "Subtract 14.97 from both sides",  eq: "9.99s \textbackslash{}\textbackslash{}leq 35.03" \},}
\NormalTok{  \{ op: "Divide both sides by 9.99",       eq: "s \textbackslash{}\textbackslash{}leq 3.506\textbackslash{}\textbackslash{}ldots" \},}
\NormalTok{  \{ op: "Round down (whole services only)", eq: "s \textbackslash{}\textbackslash{}leq 3" \},}
\NormalTok{  \{ op: "Check",                            eq: "14.97 + 9.99 \textbackslash{}\textbackslash{}times 3 = 14.97 + 29.97 = 44.94 \textbackslash{}\textbackslash{}leq 50 \textbackslash{}\textbackslash{}checkmark" \}}
\NormalTok{])}
\end{Highlighting}
\end{Shaded}

\begin{enumerate}
\def\labelenumi{\arabic{enumi}.}
\setcounter{enumi}{1}
\tightlist
\item
  A lift (elevator) has a maximum capacity of 630 kg. A group of adults
  with an average weight of 78 kg wants to use it. What is the largest
  number of people who can ride together safely?
\end{enumerate}

\begin{Shaded}
\begin{Highlighting}[]
\NormalTok{//| echo: false}
\NormalTok{makeStepperHTML(2, [}
\NormalTok{  \{ op: "Write the inequality",          eq: "78p \textbackslash{}\textbackslash{}leq 630" \},}
\NormalTok{  \{ op: "Divide both sides by 78",       eq: "p \textbackslash{}\textbackslash{}leq 8.076\textbackslash{}\textbackslash{}ldots" \},}
\NormalTok{  \{ op: "Round down (whole people only)", eq: "p \textbackslash{}\textbackslash{}leq 8" \},}
\NormalTok{  \{ op: "Check",                          eq: "78 \textbackslash{}\textbackslash{}times 8 = 624 \textbackslash{}\textbackslash{}leq 630 \textbackslash{}\textbackslash{}checkmark;\textbackslash{}\textbackslash{}quad 78 \textbackslash{}\textbackslash{}times 9 = 702 \textgreater{} 630 \textbackslash{}\textbackslash{}checkmark" \}}
\NormalTok{])}
\end{Highlighting}
\end{Shaded}

\begin{enumerate}
\def\labelenumi{\arabic{enumi}.}
\setcounter{enumi}{2}
\tightlist
\item
  Solve for \(x\) and mark the solution on a number line:
  \[-4x + 5 > 17\] Explain in one sentence what happens to the
  inequality sign and why.
\end{enumerate}

\begin{Shaded}
\begin{Highlighting}[]
\NormalTok{//| echo: false}
\NormalTok{makeStepperHTML(3, [}
\NormalTok{  \{ op: "Write the inequality",           eq: "{-}4x + 5 \textgreater{} 17" \},}
\NormalTok{  \{ op: "Subtract 5 from both sides",     eq: "{-}4x \textgreater{} 12" \},}
\NormalTok{  \{ op: "Divide by −4  (flip sign)",      eq: "x \textless{} {-}3", note: "Dividing by a negative number reverses the inequality — what was greater becomes less once all values are negated." \},}
\NormalTok{  \{ op: "Check: try x = −4",             eq: "{-}4({-}4) + 5 = 16 + 5 = 21 \textgreater{} 17 \textbackslash{}\textbackslash{}checkmark" \},}
\NormalTok{  \{ op: "Check: try x = −2 (not in solution)", eq: "{-}4({-}2) + 5 = 8 + 5 = 13 \textbackslash{}\textbackslash{}not\textgreater{} 17 \textbackslash{}\textbackslash{}checkmark" \}}
\NormalTok{])}
\end{Highlighting}
\end{Shaded}

\begin{enumerate}
\def\labelenumi{\arabic{enumi}.}
\setcounter{enumi}{3}
\tightlist
\item
  A password must be at least 8 characters. It must also be fewer than 3
  times the minimum length. Write this as a compound inequality for the
  password length \(n\), then state the range of valid lengths.
\end{enumerate}

\begin{Shaded}
\begin{Highlighting}[]
\NormalTok{//| echo: false}
\NormalTok{makeStepperHTML(4, [}
\NormalTok{  \{ op: "Write the lower bound",          eq: "n \textbackslash{}\textbackslash{}geq 8" \},}
\NormalTok{  \{ op: "Write the upper bound",          eq: "n \textless{} 3 \textbackslash{}\textbackslash{}times 8 = 24" \},}
\NormalTok{  \{ op: "Combine as compound inequality", eq: "8 \textbackslash{}\textbackslash{}leq n \textless{} 24" \},}
\NormalTok{  \{ op: "State the valid lengths",        eq: "n \textbackslash{}\textbackslash{}in \textbackslash{}\textbackslash{}\{8, 9, 10, \textbackslash{}\textbackslash{}ldots, 23\textbackslash{}\textbackslash{}\}", note: "23 is the largest valid length because n must be strictly less than 24." \},}
\NormalTok{  \{ op: "Check boundaries",              eq: "n=8: \textbackslash{}\textbackslash{}geq 8 \textbackslash{}\textbackslash{}checkmark,\textbackslash{}\textbackslash{}; n=23: \textless{} 24 \textbackslash{}\textbackslash{}checkmark,\textbackslash{}\textbackslash{}; n=24: \textbackslash{}\textbackslash{}text\{ fails\}" \}}
\NormalTok{])}
\end{Highlighting}
\end{Shaded}

\begin{enumerate}
\def\labelenumi{\arabic{enumi}.}
\setcounter{enumi}{4}
\tightlist
\item
  Solve the compound inequality and state which whole numbers satisfy
  it: \[1 < 3x - 2 \leq 13\]
\end{enumerate}

\begin{Shaded}
\begin{Highlighting}[]
\NormalTok{//| echo: false}
\NormalTok{makeStepperHTML(5, [}
\NormalTok{  \{ op: "Write the compound inequality",  eq: "1 \textless{} 3x {-} 2 \textbackslash{}\textbackslash{}leq 13" \},}
\NormalTok{  \{ op: "Add 2 to all three parts",       eq: "3 \textless{} 3x \textbackslash{}\textbackslash{}leq 15", note: "Adding the same amount to all parts preserves both inequalities." \},}
\NormalTok{  \{ op: "Divide all three parts by 3",    eq: "1 \textless{} x \textbackslash{}\textbackslash{}leq 5" \},}
\NormalTok{  \{ op: "State whole{-}number solutions",   eq: "x \textbackslash{}\textbackslash{}in \textbackslash{}\textbackslash{}\{2, 3, 4, 5\textbackslash{}\textbackslash{}\}", note: "x must be strictly greater than 1, so 1 is excluded; x = 5 is included because ≤." \},}
\NormalTok{  \{ op: "Check endpoints",               eq: "x=2:\textbackslash{}\textbackslash{}; 3(2){-}2=4,\textbackslash{}\textbackslash{}; 1 \textless{} 4 \textbackslash{}\textbackslash{}leq 13 \textbackslash{}\textbackslash{}checkmark;\textbackslash{}\textbackslash{}quad x=5:\textbackslash{}\textbackslash{}; 3(5){-}2=13,\textbackslash{}\textbackslash{}; 1 \textless{} 13 \textbackslash{}\textbackslash{}leq 13 \textbackslash{}\textbackslash{}checkmark" \}}
\NormalTok{])}
\end{Highlighting}
\end{Shaded}

\begin{enumerate}
\def\labelenumi{\arabic{enumi}.}
\setcounter{enumi}{5}
\tightlist
\item
  A phone plan charges £0.08 per minute after a 200-minute monthly
  allowance. You want your total bill to be under £25. The base plan
  costs £12. How many additional minutes can you use?
\end{enumerate}

\begin{Shaded}
\begin{Highlighting}[]
\NormalTok{//| echo: false}
\NormalTok{makeStepperHTML(6, [}
\NormalTok{  \{ op: "Write the inequality",           eq: "12 + 0.08m \textless{} 25" \},}
\NormalTok{  \{ op: "Subtract 12 from both sides",    eq: "0.08m \textless{} 13" \},}
\NormalTok{  \{ op: "Divide both sides by 0.08",      eq: "m \textless{} 162.5" \},}
\NormalTok{  \{ op: "State the answer",               eq: "m \textbackslash{}\textbackslash{}leq 162 \textbackslash{}\textbackslash{}text\{ additional minutes\}", note: "Since minutes are whole numbers, the most you can use is 162." \},}
\NormalTok{  \{ op: "Check",                          eq: "12 + 0.08 \textbackslash{}\textbackslash{}times 162 = 12 + 12.96 = 24.96 \textless{} 25 \textbackslash{}\textbackslash{}checkmark" \}}
\NormalTok{])}
\end{Highlighting}
\end{Shaded}

\begin{enumerate}
\def\labelenumi{\arabic{enumi}.}
\setcounter{enumi}{6}
\tightlist
\item
  A company's safety guideline says a storage shelf should never hold
  more than 75\% of its rated maximum. The shelf is rated to 80 kg.
  Boxes weigh 6 kg each. What is the maximum number of boxes you can
  store on this shelf within the guideline?
\end{enumerate}

\begin{Shaded}
\begin{Highlighting}[]
\NormalTok{//| echo: false}
\NormalTok{makeStepperHTML(7, [}
\NormalTok{  \{ op: "Find the safe weight limit",     eq: "0.75 \textbackslash{}\textbackslash{}times 80 = 60 \textbackslash{}\textbackslash{}text\{ kg\}" \},}
\NormalTok{  \{ op: "Write the inequality",           eq: "6b \textbackslash{}\textbackslash{}leq 60" \},}
\NormalTok{  \{ op: "Divide both sides by 6",         eq: "b \textbackslash{}\textbackslash{}leq 10" \},}
\NormalTok{  \{ op: "Check",                          eq: "6 \textbackslash{}\textbackslash{}times 10 = 60 \textbackslash{}\textbackslash{}leq 60 \textbackslash{}\textbackslash{}checkmark;\textbackslash{}\textbackslash{}quad 6 \textbackslash{}\textbackslash{}times 11 = 66 \textgreater{} 60 \textbackslash{}\textbackslash{}checkmark" \}}
\NormalTok{])}
\end{Highlighting}
\end{Shaded}

\chapter{Patterns and sequences}\label{patterns-and-sequences}

When the next step follows a rule

\hfill\break

You decide to save £20 a week. After one week: £20. After two weeks:
£40. After ten weeks: £200. You didn't need a formula for that --- you
just kept adding twenty.

A chain message says: send this to three friends. Each of those three
sends it to three more. After one round: 3 people. After two rounds: 9
people. After three rounds: 27 people. That one grows differently.
Faster. Worryingly faster.

Both of those are sequences. The first adds the same amount each time.
The second multiplies by the same amount each time. Knowing which kind
you're looking at --- and being able to jump straight to any term
without listing all the ones before it --- is what this chapter is
about.

\section{What makes a sequence a
sequence}\label{what-makes-a-sequence-a-sequence}

A sequence is a list of numbers in a definite order, where each position
has exactly one value. Position 1 has a value. Position 2 has a value.
Position \(n\) has a value. The order matters --- you can't shuffle the
terms and call it the same sequence.

What makes a sequence \emph{useful} is when there's a pattern: a rule
that connects each term to the one before it, or better yet, a rule that
connects any term directly to its position number.

That second kind --- where you can calculate any term from its position
alone --- is a function. You put a position in, you get a value out.
That idea is one of the most important in all of mathematics, and
sequences are where it shows up first in a concrete, tangible way.

For now: spot the pattern, name the rule, find any term.

\section{What the notation is
saying}\label{what-the-notation-is-saying-8}

We write the terms of a sequence as:

\[a_1, a_2, a_3, \ldots, a_n\]

Read that as: ``a-one, a-two, a-three, up to a-n.'' The letter \(a\)
names the sequence; the subscript number is the position. So \(a_1\) ---
read ``a-one'' --- is the first term, and \(a_n\) --- read ``a-n'' ---
is the term at position \(n\), wherever that happens to be.

Two sequences, two different rules:

\begin{itemize}
\tightlist
\item
  \textbf{Arithmetic sequences} add the same amount each step. That
  fixed amount is called the \textbf{common difference}, written \(d\).
\item
  \textbf{Geometric sequences} multiply by the same amount each step.
  That fixed multiplier is called the \textbf{common ratio}, written
  \(r\).
\end{itemize}

\section{Arithmetic sequences}\label{arithmetic-sequences}

\subsection{Recognising one}\label{recognising-one}

Look at the gaps between consecutive terms. If the gap is always the
same, it's arithmetic.

\[3, \; 7, \; 11, \; 15, \; 19, \ldots\]

The gap is 4 every time. So \(d = 4\) and this is arithmetic.

\[100, \; 94, \; 88, \; 82, \ldots\]

The gap is \(-6\) every time. So \(d = -6\). A decreasing sequence still
counts --- the difference is just negative.

\subsection{The nth term formula}\label{the-nth-term-formula}

Listing terms one by one works fine if you only want the first few. But
what if you want the 50th term? You don't want to list 50 items.

Notice what's happening with each term:

\begin{longtable}[]{@{}lll@{}}
\toprule\noalign{}
Position & Term & Working \\
\midrule\noalign{}
\endhead
\bottomrule\noalign{}
\endlastfoot
1 & \(a_1\) & just the starting value \\
2 & \(a_1 + d\) & added the difference once \\
3 & \(a_1 + 2d\) & added the difference twice \\
4 & \(a_1 + 3d\) & added the difference three times \\
\(n\) & \(a_1 + (n-1)d\) & added the difference \(n - 1\) times \\
\end{longtable}

The pattern: to get to position \(n\) from position 1, you add \(d\)
exactly \((n - 1)\) times. So:

\[a_n = a_1 + (n - 1)d\]

Read that aloud: ``the nth term equals the first term, plus n minus one
lots of the common difference.''

The \((n - 1)\) looks odd at first. Here's why it isn't: at position 1
you add the difference \emph{zero} times, at position 2 you add it
\emph{once}, at position \(n\) you add it \(n - 1\) times. The count of
additions is always one less than the position.

\textbf{Interactive: Arithmetic sequence builder.} Set the first term
and common difference. The dot plot shows the first 10 terms --- notice
the points fall on a straight line.

\begin{Shaded}
\begin{Highlighting}[]
\NormalTok{//| echo: false}
\NormalTok{viewof arithA1 = Inputs.range([1, 20], \{ label: "First term a₁", step: 1, value: 3 \})}
\end{Highlighting}
\end{Shaded}

\begin{Shaded}
\begin{Highlighting}[]
\NormalTok{//| echo: false}
\NormalTok{viewof arithD = Inputs.range([{-}5, 10], \{ label: "Common difference d", step: 1, value: 4 \})}
\end{Highlighting}
\end{Shaded}

\begin{Shaded}
\begin{Highlighting}[]
\NormalTok{//| echo: false}
\NormalTok{\{}
\NormalTok{  const a1 = Math.round(arithA1);}
\NormalTok{  const d  = Math.round(arithD);}
\NormalTok{  const N  = 10;}
\NormalTok{  const terms = Array.from(\{ length: N \}, (\_, i) =\textgreater{} (\{ n: i + 1, v: a1 + i * d \}));}

\NormalTok{  const W = 560, H = 280;}
\NormalTok{  const PAD\_L = 60, PAD\_R = 30, PAD\_T = 40, PAD\_B = 50;}
\NormalTok{  const plotW = W {-} PAD\_L {-} PAD\_R;}
\NormalTok{  const plotH = H {-} PAD\_T {-} PAD\_B;}

\NormalTok{  const vals = terms.map(t =\textgreater{} t.v);}
\NormalTok{  const vMin = Math.min(...vals);}
\NormalTok{  const vMax = Math.max(...vals);}
\NormalTok{  const vRange = vMax === vMin ? 1 : vMax {-} vMin;}
\NormalTok{  const vLo = vMin {-} vRange * 0.15;}
\NormalTok{  const vHi = vMax + vRange * 0.15;}

\NormalTok{  const xScale = (n) =\textgreater{} PAD\_L + (n {-} 1) / (N {-} 1) * plotW;}
\NormalTok{  const yScale = (v) =\textgreater{} PAD\_T + plotH {-} (v {-} vLo) / (vHi {-} vLo) * plotH;}

\NormalTok{  const d3local = await require("d3@7");}
\NormalTok{  const svg = d3local.create("svg")}
\NormalTok{    .attr("viewBox", \textasciigrave{}0 0 $\{W\} $\{H\}\textasciigrave{})}
\NormalTok{    .attr("width", "100\%")}
\NormalTok{    .attr("style", \textasciigrave{}max{-}width:$\{W\}px; font{-}family:inherit;\textasciigrave{});}

\NormalTok{  // Gridlines}
\NormalTok{  for (let i = 1; i \textless{}= N; i++) \{}
\NormalTok{    svg.append("line")}
\NormalTok{      .attr("x1", xScale(i)).attr("y1", PAD\_T)}
\NormalTok{      .attr("x2", xScale(i)).attr("y2", PAD\_T + plotH)}
\NormalTok{      .attr("stroke", "\#f3f4f6").attr("stroke{-}width", 1);}
\NormalTok{  \}}

\NormalTok{  // Axes}
\NormalTok{  svg.append("line")}
\NormalTok{    .attr("x1", PAD\_L).attr("y1", PAD\_T + plotH)}
\NormalTok{    .attr("x2", PAD\_L + plotW).attr("y2", PAD\_T + plotH)}
\NormalTok{    .attr("stroke", "\#9ca3af").attr("stroke{-}width", 1.5);}
\NormalTok{  svg.append("line")}
\NormalTok{    .attr("x1", PAD\_L).attr("y1", PAD\_T)}
\NormalTok{    .attr("x2", PAD\_L).attr("y2", PAD\_T + plotH)}
\NormalTok{    .attr("stroke", "\#9ca3af").attr("stroke{-}width", 1.5);}

\NormalTok{  // x{-}axis labels}
\NormalTok{  terms.forEach(t =\textgreater{} \{}
\NormalTok{    svg.append("text")}
\NormalTok{      .attr("x", xScale(t.n)).attr("y", PAD\_T + plotH + 18)}
\NormalTok{      .attr("text{-}anchor", "middle").attr("font{-}size", "11px").attr("fill", "\#6b7280")}
\NormalTok{      .text(t.n);}
\NormalTok{  \});}
\NormalTok{  svg.append("text")}
\NormalTok{    .attr("x", PAD\_L + plotW / 2).attr("y", H {-} 6)}
\NormalTok{    .attr("text{-}anchor", "middle").attr("font{-}size", "12px").attr("fill", "\#6b7280")}
\NormalTok{    .text("Position (n)");}

\NormalTok{  // y{-}axis label}
\NormalTok{  svg.append("text")}
\NormalTok{    .attr("transform", \textasciigrave{}translate(14, $\{PAD\_T + plotH / 2\}) rotate({-}90)\textasciigrave{})}
\NormalTok{    .attr("text{-}anchor", "middle").attr("font{-}size", "12px").attr("fill", "\#6b7280")}
\NormalTok{    .text("Value (aₙ)");}

\NormalTok{  // Line connecting dots}
\NormalTok{  const lineGen = d3local.line().x(t =\textgreater{} xScale(t.n)).y(t =\textgreater{} yScale(t.v));}
\NormalTok{  svg.append("path")}
\NormalTok{    .datum(terms)}
\NormalTok{    .attr("d", lineGen)}
\NormalTok{    .attr("fill", "none").attr("stroke", "\#0d9488").attr("stroke{-}width", 1.5)}
\NormalTok{    .attr("stroke{-}dasharray", "4 3").attr("opacity", 0.5);}

\NormalTok{  // Dots + value labels}
\NormalTok{  terms.forEach(t =\textgreater{} \{}
\NormalTok{    svg.append("circle")}
\NormalTok{      .attr("cx", xScale(t.n)).attr("cy", yScale(t.v))}
\NormalTok{      .attr("r", 5).attr("fill", "\#0d9488").attr("stroke", "\#fff").attr("stroke{-}width", 1.5);}
\NormalTok{    svg.append("text")}
\NormalTok{      .attr("x", xScale(t.n)).attr("y", yScale(t.v) {-} 10)}
\NormalTok{      .attr("text{-}anchor", "middle").attr("font{-}size", "10px").attr("fill", "\#374151")}
\NormalTok{      .text(t.v);}
\NormalTok{  \});}

\NormalTok{  // Formula label}
\NormalTok{  svg.append("text")}
\NormalTok{    .attr("x", PAD\_L + plotW {-} 4).attr("y", PAD\_T + 14)}
\NormalTok{    .attr("text{-}anchor", "end").attr("font{-}size", "12px").attr("fill", "\#0d9488").attr("font{-}weight", "600")}
\NormalTok{    .text(\textasciigrave{}a₁ = $\{a1\},  d = $\{d\}\textasciigrave{});}

\NormalTok{  return svg.node();}
\NormalTok{\}}
\end{Highlighting}
\end{Shaded}

\section{Worked examples ---
arithmetic}\label{worked-examples-arithmetic}

\textbf{Example 1 (sci).} Bacterial measurements are taken every 4 hours
starting at hour 0. The recorded counts (in thousands) are: 12, 15, 18,
21, \ldots{}

The question is: what is the count at hour 36?

First, work out which measurement number that is. Hour 0 is measurement
1, hour 4 is measurement 2, and so on --- each measurement is 4 hours
after the previous one. So hour 36 is measurement number
\(\frac{36}{4} + 1 = 10\).

Now check whether the sequence is arithmetic: \(15 - 12 = 3\),
\(18 - 15 = 3\), \(21 - 18 = 3\). Yes, \(d = 3\).

\[a_{10} = 12 + (10 - 1) \times 3 = 12 + 27 = 39\]

The count at hour 36 is 39 thousand. (This is arithmetic growth ---
constant addition. Real bacterial growth is geometric, but the measuring
interval can create arithmetic-looking data over short spans.)

\begin{center}\rule{0.5\linewidth}{0.5pt}\end{center}

\textbf{Example 2 (biz).} A gamer earns 150 points for completing level
1. Each subsequent level awards 25 more points than the previous one.
How many points does level 18 award?

Here \(a_1 = 150\) and \(d = 25\).

\[a_{18} = 150 + (18 - 1) \times 25 = 150 + 425 = 575\]

Level 18 awards 575 points.

\begin{center}\rule{0.5\linewidth}{0.5pt}\end{center}

\textbf{Example 3 (comp).} A programmer is adding rows to a triangle
pattern: row 1 has 1 tile, row 2 has 3 tiles, row 3 has 5 tiles.

The sequence is 1, 3, 5, 7, \ldots{} with \(a_1 = 1\) and \(d = 2\).

How many tiles does row 20 need?

\[a_{20} = 1 + (20 - 1) \times 2 = 1 + 38 = 39\]

Row 20 needs 39 tiles. A programmer building this pattern in a loop
would use exactly this calculation to allocate memory or set loop
bounds.

\begin{center}\rule{0.5\linewidth}{0.5pt}\end{center}

\textbf{Example 4 (eng).} A steel cable loses 3mm of useful length each
time it's re-terminated. It starts at 500mm. After how many
re-terminations is it shorter than 450mm?

The sequence is \(500, 497, 494, \ldots\) with \(a_1 = 500\) and
\(d = -3\).

We want \(a_n < 450\):

\[500 + (n - 1)(-3) < 450\] \[-(n-1) \times 3 < -50\]
\[(n - 1) \times 3 > 50\] \[n - 1 > 16.67\] \[n > 17.67\]

So after the 18th re-termination the cable is below 450mm. Check:
\(a_{18} = 500 + 17 \times (-3) = 500 - 51 = 449\). Yes.

\section{Geometric sequences}\label{geometric-sequences}

\subsection{Recognising one}\label{recognising-one-1}

Look at the \emph{ratios} between consecutive terms instead of the gaps.
If the ratio is always the same, it's geometric.

\[2, \; 6, \; 18, \; 54, \ldots\]

Divide each term by the previous one: \(6 \div 2 = 3\),
\(18 \div 6 = 3\), \(54 \div 18 = 3\). The ratio is always 3, so
\(r = 3\).

\[80, \; 40, \; 20, \; 10, \ldots\]

Each term is half the previous one: \(r = 0.5\). A shrinking sequence
--- the ratio just has to be positive and consistent.

\subsection{The nth term formula}\label{the-nth-term-formula-1}

The same logic as arithmetic, but with multiplication:

\begin{longtable}[]{@{}lll@{}}
\toprule\noalign{}
Position & Term & Working \\
\midrule\noalign{}
\endhead
\bottomrule\noalign{}
\endlastfoot
1 & \(a_1\) & starting value \\
2 & \(a_1 \times r\) & multiplied once \\
3 & \(a_1 \times r^2\) & multiplied twice \\
4 & \(a_1 \times r^3\) & multiplied three times \\
\(n\) & \(a_1 \times r^{n-1}\) & multiplied \(n - 1\) times \\
\end{longtable}

So:

\[a_n = a_1 \times r^{n-1}\]

The \((n - 1)\) is there for the same reason as in arithmetic: at
position 1 you multiply by \(r\) zero times. Any number raised to the
power zero equals 1 --- so \(r^0 = 1\), and the first term is just
\(a_1 \times 1 = a_1\). Each step up adds one more multiplication.

\textbf{Connection to Vol 1 Ch 3:} if this looks like the compound
growth formula, that's because it is. Compound interest is a geometric
sequence. Each year's balance is the first year's amount multiplied by
the growth factor \((1 + \text{rate})\), raised to the number of years.
The structure is identical.

\textbf{Interactive: Sequence explorer.} Toggle between arithmetic and
geometric. Adjust the first term and the common difference or ratio. The
dot plot makes the shape of growth immediately visible --- linear versus
curved.

\begin{Shaded}
\begin{Highlighting}[]
\NormalTok{//| echo: false}
\NormalTok{viewof seqType = Inputs.toggle(\{ label: "Show geometric sequence (off = arithmetic)", value: false \})}
\end{Highlighting}
\end{Shaded}

\begin{Shaded}
\begin{Highlighting}[]
\NormalTok{//| echo: false}
\NormalTok{viewof seqA1 = Inputs.range([1, 100], \{ label: "First term a₁", step: 1, value: 5 \})}
\end{Highlighting}
\end{Shaded}

\begin{Shaded}
\begin{Highlighting}[]
\NormalTok{//| echo: false}
\NormalTok{viewof seqParam = Inputs.range(seqType ? [0.5, 3.0] : [1, 10], \{}
\NormalTok{  label: seqType ? "Common ratio r" : "Common difference d",}
\NormalTok{  step: seqType ? 0.1 : 1,}
\NormalTok{  value: seqType ? 1.8 : 4}
\NormalTok{\})}
\end{Highlighting}
\end{Shaded}

\begin{Shaded}
\begin{Highlighting}[]
\NormalTok{//| echo: false}
\NormalTok{\{}
\NormalTok{  const isGeo = seqType;}
\NormalTok{  const a1 = Math.round(seqA1);}
\NormalTok{  const param = seqParam;}
\NormalTok{  const N = 10;}

\NormalTok{  const terms = Array.from(\{ length: N \}, (\_, i) =\textgreater{} (\{}
\NormalTok{    n: i + 1,}
\NormalTok{    v: isGeo ? a1 * Math.pow(param, i) : a1 + i * param}
\NormalTok{  \}));}

\NormalTok{  const W = 560, H = 280;}
\NormalTok{  const PAD\_L = 65, PAD\_R = 30, PAD\_T = 40, PAD\_B = 50;}
\NormalTok{  const plotW = W {-} PAD\_L {-} PAD\_R;}
\NormalTok{  const plotH = H {-} PAD\_T {-} PAD\_B;}

\NormalTok{  const vals = terms.map(t =\textgreater{} t.v);}
\NormalTok{  const vMin = Math.min(...vals);}
\NormalTok{  const vMax = Math.max(...vals);}
\NormalTok{  const vRange = vMax === vMin ? 1 : vMax {-} vMin;}
\NormalTok{  const vLo = Math.min(0, vMin {-} vRange * 0.1);}
\NormalTok{  const vHi = vMax + vRange * 0.15;}

\NormalTok{  const xScale = (n) =\textgreater{} PAD\_L + (n {-} 1) / (N {-} 1) * plotW;}
\NormalTok{  const yScale = (v) =\textgreater{} PAD\_T + plotH {-} (v {-} vLo) / (vHi {-} vLo) * plotH;}

\NormalTok{  const d3local = await require("d3@7");}
\NormalTok{  const svg = d3local.create("svg")}
\NormalTok{    .attr("viewBox", \textasciigrave{}0 0 $\{W\} $\{H\}\textasciigrave{})}
\NormalTok{    .attr("width", "100\%")}
\NormalTok{    .attr("style", \textasciigrave{}max{-}width:$\{W\}px; font{-}family:inherit;\textasciigrave{});}

\NormalTok{  const colour = isGeo ? "\#7c3aed" : "\#0d9488";}
\NormalTok{  const typeLabel = isGeo}
\NormalTok{    ? \textasciigrave{}Geometric: a₁ = $\{a1\}, r = $\{param.toFixed(1)\}\textasciigrave{}}
\NormalTok{    : \textasciigrave{}Arithmetic: a₁ = $\{a1\}, d = $\{param\}\textasciigrave{};}

\NormalTok{  // Gridlines}
\NormalTok{  for (let i = 1; i \textless{}= N; i++) \{}
\NormalTok{    svg.append("line")}
\NormalTok{      .attr("x1", xScale(i)).attr("y1", PAD\_T)}
\NormalTok{      .attr("x2", xScale(i)).attr("y2", PAD\_T + plotH)}
\NormalTok{      .attr("stroke", "\#f3f4f6").attr("stroke{-}width", 1);}
\NormalTok{  \}}

\NormalTok{  // Axes}
\NormalTok{  svg.append("line")}
\NormalTok{    .attr("x1", PAD\_L).attr("y1", PAD\_T + plotH)}
\NormalTok{    .attr("x2", PAD\_L + plotW).attr("y2", PAD\_T + plotH)}
\NormalTok{    .attr("stroke", "\#9ca3af").attr("stroke{-}width", 1.5);}
\NormalTok{  svg.append("line")}
\NormalTok{    .attr("x1", PAD\_L).attr("y1", PAD\_T)}
\NormalTok{    .attr("x2", PAD\_L).attr("y2", PAD\_T + plotH)}
\NormalTok{    .attr("stroke", "\#9ca3af").attr("stroke{-}width", 1.5);}

\NormalTok{  // x{-}axis labels}
\NormalTok{  terms.forEach(t =\textgreater{} \{}
\NormalTok{    svg.append("text")}
\NormalTok{      .attr("x", xScale(t.n)).attr("y", PAD\_T + plotH + 18)}
\NormalTok{      .attr("text{-}anchor", "middle").attr("font{-}size", "11px").attr("fill", "\#6b7280")}
\NormalTok{      .text(t.n);}
\NormalTok{  \});}
\NormalTok{  svg.append("text")}
\NormalTok{    .attr("x", PAD\_L + plotW / 2).attr("y", H {-} 6)}
\NormalTok{    .attr("text{-}anchor", "middle").attr("font{-}size", "12px").attr("fill", "\#6b7280")}
\NormalTok{    .text("Position (n)");}

\NormalTok{  // y{-}axis label}
\NormalTok{  svg.append("text")}
\NormalTok{    .attr("transform", \textasciigrave{}translate(14, $\{PAD\_T + plotH / 2\}) rotate({-}90)\textasciigrave{})}
\NormalTok{    .attr("text{-}anchor", "middle").attr("font{-}size", "12px").attr("fill", "\#6b7280")}
\NormalTok{    .text("Value (aₙ)");}

\NormalTok{  // Line}
\NormalTok{  const lineGen = d3local.line().x(t =\textgreater{} xScale(t.n)).y(t =\textgreater{} yScale(t.v));}
\NormalTok{  svg.append("path")}
\NormalTok{    .datum(terms)}
\NormalTok{    .attr("d", lineGen)}
\NormalTok{    .attr("fill", "none").attr("stroke", colour).attr("stroke{-}width", 1.5)}
\NormalTok{    .attr("stroke{-}dasharray", isGeo ? "none" : "4 3").attr("opacity", 0.5);}

\NormalTok{  // Dots}
\NormalTok{  terms.forEach(t =\textgreater{} \{}
\NormalTok{    const cy = yScale(t.v);}
\NormalTok{    svg.append("circle")}
\NormalTok{      .attr("cx", xScale(t.n)).attr("cy", cy)}
\NormalTok{      .attr("r", 5).attr("fill", colour).attr("stroke", "\#fff").attr("stroke{-}width", 1.5);}
\NormalTok{    const labelVal = t.v \textgreater{}= 1000 ? \textasciigrave{}$\{(t.v / 1000).toFixed(1)\}k\textasciigrave{} :}
\NormalTok{                     t.v \textgreater{}= 10 ? Math.round(t.v) : t.v.toFixed(1);}
\NormalTok{    svg.append("text")}
\NormalTok{      .attr("x", xScale(t.n)).attr("y", cy {-} 10)}
\NormalTok{      .attr("text{-}anchor", "middle").attr("font{-}size", "9px").attr("fill", "\#374151")}
\NormalTok{      .text(labelVal);}
\NormalTok{  \});}

\NormalTok{  // Type label}
\NormalTok{  svg.append("text")}
\NormalTok{    .attr("x", PAD\_L + 4).attr("y", PAD\_T + 14)}
\NormalTok{    .attr("font{-}size", "12px").attr("fill", colour).attr("font{-}weight", "600")}
\NormalTok{    .text(typeLabel);}

\NormalTok{  return svg.node();}
\NormalTok{\}}
\end{Highlighting}
\end{Shaded}

\section{Worked examples --- geometric}\label{worked-examples-geometric}

\textbf{Example 5 (biz).} A video gets 200 views on day 1. Each day it
gets three times as many views as the day before. How many views on day
6?

Here \(a_1 = 200\) and \(r = 3\).

\[a_6 = 200 \times 3^{6-1} = 200 \times 3^5 = 200 \times 243 = 48{,}600\]

48,600 views on day 6. That's the shape of viral growth --- it looks
manageable for the first few steps, then becomes enormous.

\begin{center}\rule{0.5\linewidth}{0.5pt}\end{center}

\textbf{Example 6 (sci).} A radioactive sample has 640 micrograms at the
start. Each hour, it retains half its mass. How much remains after 5
hours?

Here \(a_1 = 640\) and \(r = 0.5\). Term 1 is at \(t = 0\) (the start),
so after 5 hours we're looking at term 6.

\[a_6 = 640 \times 0.5^{6-1} = 640 \times 0.5^5 = 640 \times \frac{1}{32} = 20\]

20 micrograms remain after 5 hours.

\section{Telling them apart}\label{telling-them-apart}

The question to ask: does the sequence \emph{add} a fixed amount each
step, or \emph{multiply} by a fixed amount?

\begin{longtable}[]{@{}
  >{\raggedright\arraybackslash}p{(\linewidth - 4\tabcolsep) * \real{0.2500}}
  >{\raggedright\arraybackslash}p{(\linewidth - 4\tabcolsep) * \real{0.3929}}
  >{\raggedright\arraybackslash}p{(\linewidth - 4\tabcolsep) * \real{0.3571}}@{}}
\toprule\noalign{}
\begin{minipage}[b]{\linewidth}\raggedright
Check
\end{minipage} & \begin{minipage}[b]{\linewidth}\raggedright
Arithmetic
\end{minipage} & \begin{minipage}[b]{\linewidth}\raggedright
Geometric
\end{minipage} \\
\midrule\noalign{}
\endhead
\bottomrule\noalign{}
\endlastfoot
What's constant? & the difference between terms & the ratio between
terms \\
Each term is\ldots{} & previous term \(+ d\) & previous term
\(\times r\) \\
On a graph it looks like\ldots{} & a straight line & a curve \\
Formula & \(a_n = a_1 + (n-1)d\) & \(a_n = a_1 \times r^{n-1}\) \\
\end{longtable}

If neither test passes --- the differences aren't constant and neither
are the ratios --- the sequence isn't arithmetic or geometric. There are
other kinds, but these two cover most of what you'll see at this level.

\section{Sequences as functions}\label{sequences-as-functions}

Here's the thing to take away from this chapter.

A sequence is a function. You feed in a position number --- 1, 2, 3,
\ldots, \(n\) --- and the formula gives you back a value. One input, one
output, a definite rule. That's the definition of a function.

For an arithmetic sequence with \(a_1 = 5\) and \(d = 3\):

\[a_n = 5 + (n-1) \times 3 = 5 + 3n - 3 = 3n + 2\]

That last form --- \(3n + 2\) --- should look familiar. It's a linear
expression. The position \(n\) plays the same role as \(x\) in a linear
equation. If you plotted the terms on a graph with position on the
horizontal axis and value on the vertical, you'd get points lying on a
straight line.

Geometric sequences, plotted the same way, give points lying on a curve
--- the same kind of exponential curve you saw in compound growth.

The bigger idea --- that a rule can take any input and return a
predictable output --- is what Chapter 4 is about. Sequences are the
clearest introduction to it because the inputs are always whole numbers,
so you can list them and see the pattern directly. When Chapter 4
generalises to real-number inputs, you'll already have the idea.

\section{Where this goes}\label{where-this-goes-8}

\textbf{Chapter 4 --- Graphing relations} takes the connection from this
chapter seriously. If \(a_n = 3n + 2\) is a function of \(n\), then it
has a graph. You'll learn to draw that graph, read information off it,
and understand what the slope and intercept actually mean. Arithmetic
sequences produce straight-line graphs; geometric sequences produce
curves. Seeing them visually makes the difference between the two types
immediately obvious.

\textbf{Volume 4 --- Sequences and series} returns to sequences at a
higher level. A \emph{series} is what you get when you add up the terms
of a sequence: how much have you saved total after \(n\) weeks? How many
people has the chain message reached in total? The formulas become more
involved, but the foundation --- arithmetic or geometric, position
determines value --- is exactly what this chapter builds.

\begin{tcolorbox}[enhanced jigsaw, opacityback=0, colback=white, breakable, colframe=quarto-callout-tip-color-frame, arc=.35mm, left=2mm, toprule=.15mm, leftrule=.75mm, bottomrule=.15mm, rightrule=.15mm]
\begin{minipage}[t]{5.5mm}
\textcolor{quarto-callout-tip-color}{\faLightbulb}
\end{minipage}%
\begin{minipage}[t]{\textwidth - 5.5mm}

\textbf{Where this shows up}

\begin{itemize}
\tightlist
\item
  A financial analyst projecting investment returns is working with a
  geometric sequence --- each period's value is the previous one scaled
  by a growth factor.
\item
  A software engineer estimating how long an algorithm takes as the
  input grows is often asking whether the cost is arithmetic (grows
  linearly) or geometric (grows exponentially).
\item
  A biologist tracking bacterial colony doublings is reading a geometric
  sequence, even if the lab notebook just shows raw counts.
\item
  A structural engineer numbering equally-spaced sensor stations along a
  bridge is building an arithmetic sequence into the measurement system.
\end{itemize}

The difference between ``adds the same amount'' and ``multiplies by the
same amount'' is the difference between manageable and overwhelming.
Recognising which one you're dealing with is a genuinely useful skill.

\end{minipage}%
\end{tcolorbox}

\section{Exercises}\label{exercises-8}

These are puzzles. Each has a clean answer, but the interesting part is
identifying the type of sequence and setting up the formula before you
calculate. Don't just list terms.

\begin{Shaded}
\begin{Highlighting}[]
\NormalTok{//| echo: false}
\NormalTok{katex = require("katex@0.16")}
\end{Highlighting}
\end{Shaded}

\begin{Shaded}
\begin{Highlighting}[]
\NormalTok{//| echo: false}
\NormalTok{function makeStepperHTML(exerciseNum, steps) \{}
\NormalTok{  let current = 0;}
\NormalTok{  const totalSteps = steps.length;}
\NormalTok{  const container = document.createElement("div");}
\NormalTok{  container.style.cssText = "border:1px solid \#e5e7eb; border{-}radius:8px; padding:1rem 1.25rem; margin:0.75rem 0; font{-}family:inherit;";}
\NormalTok{  const stepsDiv = document.createElement("div");}
\NormalTok{  stepsDiv.style.marginBottom = "0.75rem";}

\NormalTok{  function renderSteps() \{}
\NormalTok{    stepsDiv.innerHTML = "";}
\NormalTok{    for (let i = 0; i \textless{} current; i++) \{}
\NormalTok{      const s = steps[i];}
\NormalTok{      const row = document.createElement("div");}
\NormalTok{      row.style.cssText = "display:grid; grid{-}template{-}columns:200px 1fr; gap:0.35rem 1rem; align{-}items:baseline; padding:0.35rem 0; border{-}top:1px solid \#f3f4f6;";}
\NormalTok{      const opCell = document.createElement("span");}
\NormalTok{      opCell.style.cssText = "font{-}size:0.85em; color:\#6b7280; font{-}style:italic;";}
\NormalTok{      opCell.textContent = s.op;}
\NormalTok{      const eqCell = document.createElement("span");}
\NormalTok{      eqCell.innerHTML = katex.renderToString(s.eq, \{ throwOnError: false, displayMode: false \});}
\NormalTok{      row.appendChild(opCell);}
\NormalTok{      row.appendChild(eqCell);}
\NormalTok{      if (s.note) \{}
\NormalTok{        const noteRow = document.createElement("div");}
\NormalTok{        noteRow.style.cssText = "grid{-}column:2; font{-}size:0.82em; color:\#6b7280; padding{-}bottom:0.2rem;";}
\NormalTok{        noteRow.textContent = s.note;}
\NormalTok{        row.appendChild(noteRow);}
\NormalTok{      \}}
\NormalTok{      stepsDiv.appendChild(row);}
\NormalTok{    \}}
\NormalTok{    if (current === totalSteps) \{}
\NormalTok{      const done = document.createElement("div");}
\NormalTok{      done.style.cssText = "margin{-}top:0.5rem; font{-}size:0.9em; color:\#059669; font{-}weight:500;";}
\NormalTok{      done.textContent = "✓ Solution complete";}
\NormalTok{      stepsDiv.appendChild(done);}
\NormalTok{    \}}
\NormalTok{  \}}

\NormalTok{  const controls = document.createElement("div");}
\NormalTok{  controls.style.cssText = "display:flex; gap:0.5rem; align{-}items:center;";}
\NormalTok{  const nextBtn = document.createElement("button");}
\NormalTok{  nextBtn.textContent = "Next step →";}
\NormalTok{  nextBtn.style.cssText = "padding:0.35rem 0.85rem; border:1px solid \#d1d5db; border{-}radius:4px; background:\#fff; cursor:pointer; font{-}size:0.9em;";}
\NormalTok{  const resetBtn = document.createElement("button");}
\NormalTok{  resetBtn.textContent = "Reset";}
\NormalTok{  resetBtn.style.cssText = "padding:0.35rem 0.75rem; border:1px solid \#e5e7eb; border{-}radius:4px; background:\#f9fafb; cursor:pointer; font{-}size:0.9em; color:\#6b7280;";}
\NormalTok{  const counter = document.createElement("span");}
\NormalTok{  counter.style.cssText = "font{-}size:0.82em; color:\#9ca3af; margin{-}left:0.25rem;";}

\NormalTok{  function updateButtons() \{}
\NormalTok{    nextBtn.disabled = current === totalSteps;}
\NormalTok{    nextBtn.style.opacity = current === totalSteps ? "0.4" : "1";}
\NormalTok{    counter.textContent = current === 0 ? "Click to reveal steps" : \textasciigrave{}Step $\{current\} of $\{totalSteps\}\textasciigrave{};}
\NormalTok{  \}}

\NormalTok{  nextBtn.onclick = () =\textgreater{} \{ if (current \textless{} totalSteps) \{ current++; renderSteps(); updateButtons(); \} \};}
\NormalTok{  resetBtn.onclick = () =\textgreater{} \{ current = 0; renderSteps(); updateButtons(); \};}
\NormalTok{  controls.appendChild(nextBtn);}
\NormalTok{  controls.appendChild(resetBtn);}
\NormalTok{  controls.appendChild(counter);}
\NormalTok{  container.appendChild(stepsDiv);}
\NormalTok{  container.appendChild(controls);}
\NormalTok{  renderSteps();}
\NormalTok{  updateButtons();}
\NormalTok{  return container;}
\NormalTok{\}}
\end{Highlighting}
\end{Shaded}

\begin{enumerate}
\def\labelenumi{\arabic{enumi}.}
\tightlist
\item
  A social media account gains 40 followers per day. It starts with 120
  followers. How many followers does it have after 3 weeks (21 days)?
  \emph{(Identify: arithmetic or geometric? Write the formula, then
  calculate.)}
\end{enumerate}

\begin{Shaded}
\begin{Highlighting}[]
\NormalTok{//| echo: false}
\NormalTok{makeStepperHTML(1, [}
\NormalTok{  \{ op: "Identify sequence type",        eq: "\textbackslash{}\textbackslash{}text\{Arithmetic — gains a fixed 40 each day\}", note: "Constant addition = arithmetic." \},}
\NormalTok{  \{ op: "Write the formula",             eq: "a\_n = 120 + (n{-}1) \textbackslash{}\textbackslash{}times 40" \},}
\NormalTok{  \{ op: "Find term at day 21 (n = 21)",  eq: "a\_\{21\} = 120 + 20 \textbackslash{}\textbackslash{}times 40 = 120 + 800 = 920" \},}
\NormalTok{  \{ op: "Check",                         eq: "a\_1 = 120 \textbackslash{}\textbackslash{}checkmark,\textbackslash{}\textbackslash{}quad a\_2 = 160 \textbackslash{}\textbackslash{}checkmark,\textbackslash{}\textbackslash{}quad a\_\{21\} = 920 \textbackslash{}\textbackslash{}checkmark" \}}
\NormalTok{])}
\end{Highlighting}
\end{Shaded}

\begin{enumerate}
\def\labelenumi{\arabic{enumi}.}
\setcounter{enumi}{1}
\tightlist
\item
  The sequence 5, 9, 13, 17, \ldots{} --- what is the 25th term?
\end{enumerate}

\begin{Shaded}
\begin{Highlighting}[]
\NormalTok{//| echo: false}
\NormalTok{makeStepperHTML(2, [}
\NormalTok{  \{ op: "Find the common difference",    eq: "d = 9 {-} 5 = 4" \},}
\NormalTok{  \{ op: "Write the formula",             eq: "a\_n = 5 + (n{-}1) \textbackslash{}\textbackslash{}times 4" \},}
\NormalTok{  \{ op: "Substitute n = 25",             eq: "a\_\{25\} = 5 + 24 \textbackslash{}\textbackslash{}times 4 = 5 + 96 = 101" \},}
\NormalTok{  \{ op: "Check",                         eq: "a\_1 = 5 \textbackslash{}\textbackslash{}checkmark,\textbackslash{}\textbackslash{}quad a\_2 = 9 \textbackslash{}\textbackslash{}checkmark,\textbackslash{}\textbackslash{}quad a\_\{25\} = 101 \textbackslash{}\textbackslash{}checkmark" \}}
\NormalTok{])}
\end{Highlighting}
\end{Shaded}

\begin{enumerate}
\def\labelenumi{\arabic{enumi}.}
\setcounter{enumi}{2}
\tightlist
\item
  A ball is dropped from 160cm. Each bounce reaches exactly 75\% of the
  previous height. How high is the 5th bounce? Round to the nearest
  centimetre.
\end{enumerate}

\begin{Shaded}
\begin{Highlighting}[]
\NormalTok{//| echo: false}
\NormalTok{makeStepperHTML(3, [}
\NormalTok{  \{ op: "Identify sequence type",        eq: "\textbackslash{}\textbackslash{}text\{Geometric — multiplied by 0.75 each time\}" \},}
\NormalTok{  \{ op: "Write the formula",             eq: "a\_n = 160 \textbackslash{}\textbackslash{}times 0.75\^{}\{n{-}1\}", note: "The first bounce is n = 1 (immediately after the drop)." \},}
\NormalTok{  \{ op: "Substitute n = 5",              eq: "a\_5 = 160 \textbackslash{}\textbackslash{}times 0.75\^{}4 = 160 \textbackslash{}\textbackslash{}times 0.3164\textbackslash{}\textbackslash{}ldots \textbackslash{}\textbackslash{}approx 50.6" \},}
\NormalTok{  \{ op: "Round to nearest cm",           eq: "a\_5 \textbackslash{}\textbackslash{}approx 51 \textbackslash{}\textbackslash{}text\{ cm\}" \},}
\NormalTok{  \{ op: "Check",                         eq: "160, 120, 90, 67.5, 50.6 \textbackslash{}\textbackslash{}ldots \textbackslash{}\textbackslash{}checkmark" \}}
\NormalTok{])}
\end{Highlighting}
\end{Shaded}

\begin{enumerate}
\def\labelenumi{\arabic{enumi}.}
\setcounter{enumi}{3}
\tightlist
\item
  A sequence begins 3, 6, 12, 24, \ldots{}

  \begin{enumerate}
  \def\labelenumii{\alph{enumii}.}
  \tightlist
  \item
    Write the formula for the nth term.
  \item
    Which term first exceeds 3,000?
  \end{enumerate}
\end{enumerate}

\begin{Shaded}
\begin{Highlighting}[]
\NormalTok{//| echo: false}
\NormalTok{makeStepperHTML(4, [}
\NormalTok{  \{ op: "Find the common ratio",         eq: "r = 6 \textbackslash{}\textbackslash{}div 3 = 2" \},}
\NormalTok{  \{ op: "Write the formula (part a)",    eq: "a\_n = 3 \textbackslash{}\textbackslash{}times 2\^{}\{n{-}1\}" \},}
\NormalTok{  \{ op: "Try n = 10",                    eq: "a\_\{10\} = 3 \textbackslash{}\textbackslash{}times 2\^{}9 = 3 \textbackslash{}\textbackslash{}times 512 = 1536 \textless{} 3000" \},}
\NormalTok{  \{ op: "Try n = 11",                    eq: "a\_\{11\} = 3 \textbackslash{}\textbackslash{}times 2\^{}\{10\} = 3 \textbackslash{}\textbackslash{}times 1024 = 3072 \textgreater{} 3000 \textbackslash{}\textbackslash{}checkmark" \},}
\NormalTok{  \{ op: "Answer (part b)",               eq: "\textbackslash{}\textbackslash{}text\{Term 11 is the first to exceed 3,000\}" \}}
\NormalTok{])}
\end{Highlighting}
\end{Shaded}

\begin{enumerate}
\def\labelenumi{\arabic{enumi}.}
\setcounter{enumi}{4}
\tightlist
\item
  Exam scores in a practice series improve by 4 marks each attempt: 46,
  50, 54, 58, \ldots{} After how many attempts does a student first
  reach or exceed 80 marks?
\end{enumerate}

\begin{Shaded}
\begin{Highlighting}[]
\NormalTok{//| echo: false}
\NormalTok{makeStepperHTML(5, [}
\NormalTok{  \{ op: "Write the formula",             eq: "a\_n = 46 + (n{-}1) \textbackslash{}\textbackslash{}times 4" \},}
\NormalTok{  \{ op: "Set up inequality",             eq: "46 + (n{-}1) \textbackslash{}\textbackslash{}times 4 \textbackslash{}\textbackslash{}geq 80" \},}
\NormalTok{  \{ op: "Subtract 46",                   eq: "(n{-}1) \textbackslash{}\textbackslash{}times 4 \textbackslash{}\textbackslash{}geq 34" \},}
\NormalTok{  \{ op: "Divide by 4",                   eq: "n {-} 1 \textbackslash{}\textbackslash{}geq 8.5" \},}
\NormalTok{  \{ op: "Add 1, round up",               eq: "n \textbackslash{}\textbackslash{}geq 9.5 \textbackslash{}\textbackslash{}Rightarrow n = 10" \},}
\NormalTok{  \{ op: "Check",                         eq: "a\_\{10\} = 46 + 9 \textbackslash{}\textbackslash{}times 4 = 46 + 36 = 82 \textbackslash{}\textbackslash{}geq 80 \textbackslash{}\textbackslash{}checkmark" \}}
\NormalTok{])}
\end{Highlighting}
\end{Shaded}

\begin{enumerate}
\def\labelenumi{\arabic{enumi}.}
\setcounter{enumi}{5}
\tightlist
\item
  A fundraiser's social posts are shared in a pattern. On day 1 it
  reaches 5 people. Each day it reaches 4 times as many people as the
  day before. On which day does it first reach more than 50,000 people?
\end{enumerate}

\begin{Shaded}
\begin{Highlighting}[]
\NormalTok{//| echo: false}
\NormalTok{makeStepperHTML(6, [}
\NormalTok{  \{ op: "Write the formula",             eq: "a\_n = 5 \textbackslash{}\textbackslash{}times 4\^{}\{n{-}1\}" \},}
\NormalTok{  \{ op: "Try n = 7",                     eq: "a\_7 = 5 \textbackslash{}\textbackslash{}times 4\^{}6 = 5 \textbackslash{}\textbackslash{}times 4096 = 20\{,\}480 \textless{} 50\{,\}000" \},}
\NormalTok{  \{ op: "Try n = 8",                     eq: "a\_8 = 5 \textbackslash{}\textbackslash{}times 4\^{}7 = 5 \textbackslash{}\textbackslash{}times 16\{,\}384 = 81\{,\}920 \textgreater{} 50\{,\}000 \textbackslash{}\textbackslash{}checkmark" \},}
\NormalTok{  \{ op: "Answer",                        eq: "\textbackslash{}\textbackslash{}text\{Day 8 is the first day it exceeds 50,000 people\}" \}}
\NormalTok{])}
\end{Highlighting}
\end{Shaded}

\begin{enumerate}
\def\labelenumi{\arabic{enumi}.}
\setcounter{enumi}{6}
\tightlist
\item
  Two sequences both start at 10.

  \begin{itemize}
  \tightlist
  \item
    Sequence A: add 15 each term.
  \item
    Sequence B: multiply by 1.5 each term.
  \end{itemize}

  By the 10th term, which is larger? By how much?
\end{enumerate}

\begin{Shaded}
\begin{Highlighting}[]
\NormalTok{//| echo: false}
\NormalTok{makeStepperHTML(7, [}
\NormalTok{  \{ op: "Sequence A — formula",          eq: "a\_n = 10 + (n{-}1) \textbackslash{}\textbackslash{}times 15" \},}
\NormalTok{  \{ op: "Sequence A — 10th term",        eq: "a\_\{10\} = 10 + 9 \textbackslash{}\textbackslash{}times 15 = 10 + 135 = 145" \},}
\NormalTok{  \{ op: "Sequence B — formula",          eq: "b\_n = 10 \textbackslash{}\textbackslash{}times 1.5\^{}\{n{-}1\}" \},}
\NormalTok{  \{ op: "Sequence B — 10th term",        eq: "b\_\{10\} = 10 \textbackslash{}\textbackslash{}times 1.5\^{}9 = 10 \textbackslash{}\textbackslash{}times 38.44\textbackslash{}\textbackslash{}ldots \textbackslash{}\textbackslash{}approx 384.4" \},}
\NormalTok{  \{ op: "Compare",                       eq: "384.4 {-} 145 \textbackslash{}\textbackslash{}approx 239.4", note: "Geometric growth overtakes arithmetic — by the 10th term, Sequence B is more than twice as large." \}}
\NormalTok{])}
\end{Highlighting}
\end{Shaded}

\chapter{Graphing relations}\label{graphing-relations}

Seeing equations as pictures

\hfill\break

You've checked a mobile data plan: £10 a month fixed, then £2 per extra
GB. You want to know, at a glance, what you'd pay for 0, 1, 3, 5
gigabytes. You could solve an equation for each value --- or you could
draw a line and read off every answer at once.

That line is the graph of the equation. It doesn't replace the algebra.
It shows you the whole relationship in one picture.

\section{What the notation is
saying}\label{what-the-notation-is-saying-9}

\subsection{The coordinate plane}\label{the-coordinate-plane}

Draw two number lines crossing at right angles. The horizontal one is
the \textbf{x-axis}. The vertical one is the \textbf{y-axis}. The point
where they cross is the \textbf{origin} --- the zero of both axes.

Any point in the plane has an address: \((x, y)\), read ``x across, y
up''. The \(x\) tells you how far to move left or right from the origin.
The \(y\) tells you how far to move up or down.

So \((3, 2)\) means: go 3 to the right, then 2 up. And \((-1, 4)\)
means: go 1 to the left, then 4 up. And \((0, -3)\) means: stay on the
vertical axis, go 3 down.

The two axes divide the plane into four regions, called
\textbf{quadrants}. Top-right is Quadrant I (both \(x\) and \(y\)
positive). Top-left is Quadrant II (\(x\) negative, \(y\) positive).
Bottom-left is Quadrant III (both negative). Bottom-right is Quadrant IV
(\(x\) positive, \(y\) negative). Most of the graphs in this chapter
live in Quadrant I, where both values are positive --- because most
real-world quantities like cost, distance, and time don't go negative.

\textbf{Interactive: Coordinate plane explorer.} Click anywhere on the
grid to place a point. Watch the \((x, y)\) coordinates update and
animated arrows trace the ``go right, then up'' path from the origin.

\begin{Shaded}
\begin{Highlighting}[]
\NormalTok{//| echo: false}
\NormalTok{\{}
\NormalTok{  const W = 500, H = 500;}
\NormalTok{  const ORIGIN\_X = W / 2, ORIGIN\_Y = H / 2;}
\NormalTok{  const GRID\_STEP = 40; // pixels per unit}
\NormalTok{  const UNITS = 5;      // axis runs from {-}5 to +5 in each direction (but we\textquotesingle{}ll clip at W/H)}

\NormalTok{  let pointX = null, pointY = null;}

\NormalTok{  const svg = d3.create("svg")}
\NormalTok{    .attr("viewBox", \textasciigrave{}0 0 $\{W\} $\{H\}\textasciigrave{})}
\NormalTok{    .attr("width", "100\%")}
\NormalTok{    .attr("style", \textasciigrave{}max{-}width:$\{W\}px; font{-}family:inherit; cursor:crosshair;\textasciigrave{});}

\NormalTok{  const toSvg = (mathX, mathY) =\textgreater{} (\{}
\NormalTok{    sx: ORIGIN\_X + mathX * GRID\_STEP,}
\NormalTok{    sy: ORIGIN\_Y {-} mathY * GRID\_STEP}
\NormalTok{  \});}
\NormalTok{  const toMath = (sx, sy) =\textgreater{} (\{}
\NormalTok{    mx: (sx {-} ORIGIN\_X) / GRID\_STEP,}
\NormalTok{    my: (ORIGIN\_Y {-} sy) / GRID\_STEP}
\NormalTok{  \});}

\NormalTok{  // Quadrant labels}
\NormalTok{  const quadrants = [}
\NormalTok{    \{ label: "I", mx: 3, my: 3 \},}
\NormalTok{    \{ label: "II", mx: {-}3, my: 3 \},}
\NormalTok{    \{ label: "III", mx: {-}3, my: {-}3 \},}
\NormalTok{    \{ label: "IV", mx: 3, my: {-}3 \}}
\NormalTok{  ];}
\NormalTok{  quadrants.forEach(q =\textgreater{} \{}
\NormalTok{    const \{ sx, sy \} = toSvg(q.mx, q.my);}
\NormalTok{    svg.append("text")}
\NormalTok{      .attr("x", sx).attr("y", sy)}
\NormalTok{      .attr("text{-}anchor", "middle").attr("dominant{-}baseline", "middle")}
\NormalTok{      .attr("font{-}size", "13px").attr("fill", "\#e5e7eb").attr("font{-}weight", "600")}
\NormalTok{      .text(\textasciigrave{}Q$\{q.label\}\textasciigrave{});}
\NormalTok{  \});}

\NormalTok{  // Grid lines}
\NormalTok{  for (let i = {-}UNITS; i \textless{}= UNITS; i++) \{}
\NormalTok{    const \{ sx \} = toSvg(i, 0);}
\NormalTok{    const \{ sy \} = toSvg(0, i);}
\NormalTok{    svg.append("line")}
\NormalTok{      .attr("x1", sx).attr("y1", 0).attr("x2", sx).attr("y2", H)}
\NormalTok{      .attr("stroke", i === 0 ? "\#6b7280" : "\#f3f4f6").attr("stroke{-}width", i === 0 ? 2 : 1);}
\NormalTok{    svg.append("line")}
\NormalTok{      .attr("x1", 0).attr("y1", sy).attr("x2", W).attr("y2", sy)}
\NormalTok{      .attr("stroke", i === 0 ? "\#6b7280" : "\#f3f4f6").attr("stroke{-}width", i === 0 ? 2 : 1);}
\NormalTok{  \}}

\NormalTok{  // Axis tick labels}
\NormalTok{  for (let i = {-}UNITS; i \textless{}= UNITS; i++) \{}
\NormalTok{    if (i === 0) continue;}
\NormalTok{    const \{ sx \} = toSvg(i, 0);}
\NormalTok{    const \{ sy \} = toSvg(0, i);}
\NormalTok{    svg.append("text")}
\NormalTok{      .attr("x", sx).attr("y", ORIGIN\_Y + 16)}
\NormalTok{      .attr("text{-}anchor", "middle").attr("font{-}size", "10px").attr("fill", "\#9ca3af")}
\NormalTok{      .text(i);}
\NormalTok{    svg.append("text")}
\NormalTok{      .attr("x", ORIGIN\_X {-} 14).attr("y", sy + 4)}
\NormalTok{      .attr("text{-}anchor", "end").attr("font{-}size", "10px").attr("fill", "\#9ca3af")}
\NormalTok{      .text(i);}
\NormalTok{  \}}

\NormalTok{  // Axis labels}
\NormalTok{  svg.append("text").attr("x", W {-} 8).attr("y", ORIGIN\_Y {-} 8)}
\NormalTok{    .attr("font{-}size", "13px").attr("fill", "\#6b7280").attr("font{-}weight", "600").text("x");}
\NormalTok{  svg.append("text").attr("x", ORIGIN\_X + 8).attr("y", 14)}
\NormalTok{    .attr("font{-}size", "13px").attr("fill", "\#6b7280").attr("font{-}weight", "600").text("y");}

\NormalTok{  // Arrow group (updated on click)}
\NormalTok{  const arrowGroup = svg.append("g").attr("class", "arrows");}
\NormalTok{  const pointGroup = svg.append("g").attr("class", "point");}
\NormalTok{  const labelGroup = svg.append("g").attr("class", "label");}

\NormalTok{  // Instruction text}
\NormalTok{  const instrText = svg.append("text")}
\NormalTok{    .attr("x", W / 2).attr("y", H {-} 10)}
\NormalTok{    .attr("text{-}anchor", "middle").attr("font{-}size", "12px").attr("fill", "\#9ca3af")}
\NormalTok{    .text("Click anywhere on the grid to place a point");}

\NormalTok{  svg.on("click", function(event) \{}
\NormalTok{    const [sx, sy] = d3.pointer(event, svg.node());}
\NormalTok{    const \{ mx, my \} = toMath(sx, sy);}
\NormalTok{    const rx = Math.round(mx * 2) / 2;}
\NormalTok{    const ry = Math.round(my * 2) / 2;}

\NormalTok{    const \{ sx: px, sy: py \} = toSvg(rx, ry);}
\NormalTok{    const \{ sy: axisY \} = toSvg(0, 0);}
\NormalTok{    const \{ sx: axisX \} = toSvg(0, 0);}

\NormalTok{    arrowGroup.selectAll("*").remove();}
\NormalTok{    pointGroup.selectAll("*").remove();}
\NormalTok{    labelGroup.selectAll("*").remove();}
\NormalTok{    instrText.text("");}

\NormalTok{    // Horizontal arrow from origin to (rx, axisY)}
\NormalTok{    if (Math.abs(rx) \textgreater{} 0.1) \{}
\NormalTok{      arrowGroup.append("line")}
\NormalTok{        .attr("x1", ORIGIN\_X).attr("y1", ORIGIN\_Y)}
\NormalTok{        .attr("x2", px).attr("y2", ORIGIN\_Y)}
\NormalTok{        .attr("stroke", "\#0d9488").attr("stroke{-}width", 2).attr("stroke{-}dasharray", "5 3");}
\NormalTok{      arrowGroup.append("text")}
\NormalTok{        .attr("x", (ORIGIN\_X + px) / 2).attr("y", ORIGIN\_Y {-} 8)}
\NormalTok{        .attr("text{-}anchor", "middle").attr("font{-}size", "11px").attr("fill", "\#0d9488")}
\NormalTok{        .text(rx \textgreater{} 0 ? \textasciigrave{}→ $\{rx\}\textasciigrave{} : \textasciigrave{}← $\{Math.abs(rx)\}\textasciigrave{});}
\NormalTok{    \}}

\NormalTok{    // Vertical arrow from (rx, 0) to (rx, ry)}
\NormalTok{    if (Math.abs(ry) \textgreater{} 0.1) \{}
\NormalTok{      arrowGroup.append("line")}
\NormalTok{        .attr("x1", px).attr("y1", ORIGIN\_Y)}
\NormalTok{        .attr("x2", px).attr("y2", py)}
\NormalTok{        .attr("stroke", "\#7c3aed").attr("stroke{-}width", 2).attr("stroke{-}dasharray", "5 3");}
\NormalTok{      arrowGroup.append("text")}
\NormalTok{        .attr("x", px + 10).attr("y", (ORIGIN\_Y + py) / 2)}
\NormalTok{        .attr("font{-}size", "11px").attr("fill", "\#7c3aed")}
\NormalTok{        .text(ry \textgreater{} 0 ? \textasciigrave{}↑ $\{ry\}\textasciigrave{} : \textasciigrave{}↓ $\{Math.abs(ry)\}\textasciigrave{});}
\NormalTok{    \}}

\NormalTok{    // Point}
\NormalTok{    pointGroup.append("circle")}
\NormalTok{      .attr("cx", px).attr("cy", py).attr("r", 7)}
\NormalTok{      .attr("fill", "\#f59e0b").attr("stroke", "\#fff").attr("stroke{-}width", 2);}

\NormalTok{    // Coordinate label}
\NormalTok{    const labelX = px + (rx \textgreater{} 3 ? {-}60 : 12);}
\NormalTok{    const labelY = py + (ry \textless{} {-}3 ? 20 : {-}12);}
\NormalTok{    labelGroup.append("rect")}
\NormalTok{      .attr("x", labelX {-} 4).attr("y", labelY {-} 14)}
\NormalTok{      .attr("width", 68).attr("height", 20).attr("rx", 4)}
\NormalTok{      .attr("fill", "\#fff").attr("stroke", "\#e5e7eb").attr("stroke{-}width", 1);}
\NormalTok{    labelGroup.append("text")}
\NormalTok{      .attr("x", labelX).attr("y", labelY)}
\NormalTok{      .attr("font{-}size", "12px").attr("fill", "\#1f2937").attr("font{-}weight", "600")}
\NormalTok{      .text(\textasciigrave{}($\{rx\}, $\{ry\})\textasciigrave{});}
\NormalTok{  \});}

\NormalTok{  return svg.node();}
\NormalTok{\}}
\end{Highlighting}
\end{Shaded}

\subsection{A relation as a table of
pairs}\label{a-relation-as-a-table-of-pairs}

A \textbf{relation} is any set of input-output pairs. You give it an
input value, it gives back an output value.

The phone plan from the opening is a relation. Call the number of extra
GB the input \(x\), and the monthly cost in pounds the output \(y\). The
rule is:

\[y = 2x + 10\]

Read it as: ``y equals 2 times x, plus 10.'' Plug in \(x = 0\): you get
\(y = 10\). Plug in \(x = 3\): you get \(y = 16\). That pair of values
is the point \((3, 16)\).

Here's the table for a few values:

\begin{longtable}[]{@{}ccc@{}}
\toprule\noalign{}
\(x\) (extra GB) & \(y\) (cost, £) & Point \\
\midrule\noalign{}
\endhead
\bottomrule\noalign{}
\endlastfoot
0 & 10 & \((0, 10)\) \\
1 & 12 & \((1, 12)\) \\
2 & 14 & \((2, 14)\) \\
3 & 16 & \((3, 16)\) \\
5 & 20 & \((5, 20)\) \\
\end{longtable}

Plot those five points on the coordinate plane and they fall in a
straight line. Every point on that line satisfies the equation. Every
point off it doesn't.

That's the key idea: the graph and the equation are the same
relationship in two different languages.

\section{The method}\label{the-method-8}

\subsection{Reading slope: rise over
run}\label{reading-slope-rise-over-run}

The equation \(y = 2x + 10\) has the form \(y = mx + b\).

The letter \(m\) is the \textbf{slope} --- how steeply the line rises or
falls. Read it as: for every 1 step to the right, the line goes \(m\)
steps up (or down, if \(m\) is negative).

To measure slope from any two points on the line, count:

\[m = \frac{\text{rise}}{\text{run}} = \frac{\text{change in } y}{\text{change in } x}\]

For the phone plan: pick \((1, 12)\) and \((3, 16)\). Rise =
\(16 - 12 = 4\). Run = \(3 - 1 = 2\). Slope = \(4 \div 2 = 2\). That
matches the \(2\) in front of \(x\).

\textbf{Positive slope} means the line goes up left-to-right --- the
quantity grows as \(x\) increases.

\textbf{Negative slope} means the line goes down left-to-right --- the
quantity shrinks as \(x\) increases. For a negative slope like \(-3\),
the number of units down for each unit across is 3 --- the size of the
slope ignoring the sign. The sign tells you the direction; the number
tells you how steep.

\textbf{Zero slope} means a flat horizontal line --- the output doesn't
change no matter what the input is.

The letter \(b\) is the \textbf{y-intercept} --- where the line crosses
the y-axis. Substitute \(x = 0\) into any linear equation and you get
\(y = b\). For the phone plan, \(b = 10\): the line crosses the y-axis
at \((0, 10)\), meaning you pay £10 even with zero extra data. That's
the standing charge.

\begin{tcolorbox}[enhanced jigsaw, opacityback=0, colback=white, breakable, colframe=quarto-callout-note-color-frame, arc=.35mm, left=2mm, toprule=.15mm, leftrule=.75mm, bottomrule=.15mm, rightrule=.15mm]
\begin{minipage}[t]{5.5mm}
\textcolor{quarto-callout-note-color}{\faInfo}
\end{minipage}%
\begin{minipage}[t]{\textwidth - 5.5mm}

\textbf{Why this works}

Every point on a line satisfies its equation --- that's the definition
of the line. So if you want to draw the line, you only need two points:
the equation holds at both, and a straight line through two points is
determined completely. The y-intercept \((0, b)\) is almost always the
easiest first point. From there, the slope tells you how to find the
next one.

\end{minipage}%
\end{tcolorbox}

\subsection{Method 1: Graphing from a table of
values}\label{method-1-graphing-from-a-table-of-values}

\begin{enumerate}
\def\labelenumi{\arabic{enumi}.}
\tightlist
\item
  Choose three or four values of \(x\) --- spread them out.
\item
  Substitute each into the equation to get the matching \(y\).
\item
  Write down the \((x, y)\) pairs.
\item
  Plot the points on the coordinate plane.
\item
  Draw a straight line through them.
\end{enumerate}

Use at least three points. If the third one doesn't fall on the line
through the first two, you've made an arithmetic error somewhere.

\subsection{Method 2: Graphing from slope-intercept
form}\label{method-2-graphing-from-slope-intercept-form}

When the equation is written as \(y = mx + b\), you can graph it without
a table:

\begin{enumerate}
\def\labelenumi{\arabic{enumi}.}
\tightlist
\item
  Plot the y-intercept \((0, b)\) --- that's your starting point.
\item
  From that point, apply the slope: go 1 unit right, go \(m\) units up
  (or, if \(m\) is negative, go that many units down).
\item
  Mark that second point.
\item
  Draw the line through both points.
\item
  Extend the line in both directions.
\end{enumerate}

This is faster than a table once you're comfortable with it.

\subsection{Method 3: Finding the equation from a point and a
slope}\label{method-3-finding-the-equation-from-a-point-and-a-slope}

Sometimes you know the slope and a point on the line, but not the
y-intercept. To find \(b\), substitute the known point into
\(y = mx + b\) and solve.

For example: slope \(m = 2\), and the line passes through \((3, 11)\).
Substitute into \(y = mx + b\):

\[11 = 2(3) + b\] \[11 = 6 + b\] \[b = 5\]

So the equation is \(y = 2x + 5\). Check: does \((3, 11)\) satisfy it?
\(2(3) + 5 = 11\). Yes.

The same method works when you're given two points instead of a point
and a slope --- first calculate the slope from the two points using rise
over run, then substitute one of the points to find \(b\).

\textbf{Interactive: Linear graph builder.} Drag the sliders to change
slope \(m\) and y-intercept \(b\). The line \(y = mx + b\) redraws live.
The slope triangle shows the rise and run. The equation updates above
the graph.

\begin{Shaded}
\begin{Highlighting}[]
\NormalTok{//| echo: false}
\NormalTok{viewof graphM = Inputs.range([{-}3, 3], \{ label: "Slope m", step: 0.5, value: 1 \})}
\end{Highlighting}
\end{Shaded}

\begin{Shaded}
\begin{Highlighting}[]
\NormalTok{//| echo: false}
\NormalTok{viewof graphB = Inputs.range([{-}5, 5], \{ label: "y{-}intercept b", step: 0.5, value: 2 \})}
\end{Highlighting}
\end{Shaded}

\begin{Shaded}
\begin{Highlighting}[]
\NormalTok{//| echo: false}
\NormalTok{\{}
\NormalTok{  const m = graphM;}
\NormalTok{  const b = graphB;}

\NormalTok{  const W = 500, H = 500;}
\NormalTok{  const X\_MIN = {-}5, X\_MAX = 10, Y\_MIN = {-}5, Y\_MAX = 10;}
\NormalTok{  const PAD\_L = 48, PAD\_R = 20, PAD\_T = 48, PAD\_B = 40;}
\NormalTok{  const plotW = W {-} PAD\_L {-} PAD\_R;}
\NormalTok{  const plotH = H {-} PAD\_T {-} PAD\_B;}

\NormalTok{  const xScale = (v) =\textgreater{} PAD\_L + (v {-} X\_MIN) / (X\_MAX {-} X\_MIN) * plotW;}
\NormalTok{  const yScale = (v) =\textgreater{} PAD\_T + plotH {-} (v {-} Y\_MIN) / (Y\_MAX {-} Y\_MIN) * plotH;}

\NormalTok{  const d3local = await require("d3@7");}
\NormalTok{  const svg = d3local.create("svg")}
\NormalTok{    .attr("viewBox", \textasciigrave{}0 0 $\{W\} $\{H\}\textasciigrave{})}
\NormalTok{    .attr("width", "100\%")}
\NormalTok{    .attr("style", \textasciigrave{}max{-}width:$\{W\}px; font{-}family:inherit;\textasciigrave{});}

\NormalTok{  // Background}
\NormalTok{  svg.append("rect")}
\NormalTok{    .attr("x", PAD\_L).attr("y", PAD\_T)}
\NormalTok{    .attr("width", plotW).attr("height", plotH)}
\NormalTok{    .attr("fill", "\#f9fafb");}

\NormalTok{  // Grid}
\NormalTok{  for (let x = X\_MIN; x \textless{}= X\_MAX; x++) \{}
\NormalTok{    svg.append("line")}
\NormalTok{      .attr("x1", xScale(x)).attr("y1", PAD\_T)}
\NormalTok{      .attr("x2", xScale(x)).attr("y2", PAD\_T + plotH)}
\NormalTok{      .attr("stroke", x === 0 ? "\#9ca3af" : "\#e5e7eb")}
\NormalTok{      .attr("stroke{-}width", x === 0 ? 2 : 1);}
\NormalTok{  \}}
\NormalTok{  for (let y = Y\_MIN; y \textless{}= Y\_MAX; y++) \{}
\NormalTok{    svg.append("line")}
\NormalTok{      .attr("x1", PAD\_L).attr("y1", yScale(y))}
\NormalTok{      .attr("x2", PAD\_L + plotW).attr("y2", yScale(y))}
\NormalTok{      .attr("stroke", y === 0 ? "\#9ca3af" : "\#e5e7eb")}
\NormalTok{      .attr("stroke{-}width", y === 0 ? 2 : 1);}
\NormalTok{  \}}

\NormalTok{  // Axis tick labels}
\NormalTok{  for (let x = X\_MIN; x \textless{}= X\_MAX; x += 2) \{}
\NormalTok{    if (x === 0) continue;}
\NormalTok{    svg.append("text")}
\NormalTok{      .attr("x", xScale(x)).attr("y", yScale(0) + 16)}
\NormalTok{      .attr("text{-}anchor", "middle").attr("font{-}size", "10px").attr("fill", "\#9ca3af")}
\NormalTok{      .text(x);}
\NormalTok{  \}}
\NormalTok{  for (let y = Y\_MIN; y \textless{}= Y\_MAX; y += 2) \{}
\NormalTok{    if (y === 0) continue;}
\NormalTok{    svg.append("text")}
\NormalTok{      .attr("x", xScale(0) {-} 8).attr("y", yScale(y) + 4)}
\NormalTok{      .attr("text{-}anchor", "end").attr("font{-}size", "10px").attr("fill", "\#9ca3af")}
\NormalTok{      .text(y);}
\NormalTok{  \}}
\NormalTok{  // Axis labels}
\NormalTok{  svg.append("text").attr("x", PAD\_L + plotW {-} 4).attr("y", yScale(0) {-} 6)}
\NormalTok{    .attr("font{-}size", "12px").attr("fill", "\#6b7280").attr("font{-}weight", "600").text("x");}
\NormalTok{  svg.append("text").attr("x", xScale(0) + 6).attr("y", PAD\_T + 10)}
\NormalTok{    .attr("font{-}size", "12px").attr("fill", "\#6b7280").attr("font{-}weight", "600").text("y");}

\NormalTok{  // Clip the line to the plot area}
\NormalTok{  const lineY = (x) =\textgreater{} m * x + b;}
\NormalTok{  // Find x values where line enters/exits the plot area}
\NormalTok{  const candidates = [}
\NormalTok{    \{ x: X\_MIN, y: lineY(X\_MIN) \},}
\NormalTok{    \{ x: X\_MAX, y: lineY(X\_MAX) \},}
\NormalTok{    m !== 0 ? \{ x: (Y\_MIN {-} b) / m, y: Y\_MIN \} : null,}
\NormalTok{    m !== 0 ? \{ x: (Y\_MAX {-} b) / m, y: Y\_MAX \} : null}
\NormalTok{  ].filter(p =\textgreater{} p \&\& p.x \textgreater{}= X\_MIN {-} 0.01 \&\& p.x \textless{}= X\_MAX + 0.01 \&\& p.y \textgreater{}= Y\_MIN {-} 0.01 \&\& p.y \textless{}= Y\_MAX + 0.01);}
\NormalTok{  candidates.sort((a, c) =\textgreater{} a.x {-} c.x);}

\NormalTok{  if (candidates.length \textgreater{}= 2) \{}
\NormalTok{    const p1 = candidates[0], p2 = candidates[candidates.length {-} 1];}
\NormalTok{    svg.append("line")}
\NormalTok{      .attr("x1", xScale(p1.x)).attr("y1", yScale(p1.y))}
\NormalTok{      .attr("x2", xScale(p2.x)).attr("y2", yScale(p2.y))}
\NormalTok{      .attr("stroke", "\#0d9488").attr("stroke{-}width", 2.5);}
\NormalTok{  \}}

\NormalTok{  // y{-}intercept point}
\NormalTok{  const bInView = b \textgreater{}= Y\_MIN \&\& b \textless{}= Y\_MAX;}
\NormalTok{  if (bInView) \{}
\NormalTok{    svg.append("circle")}
\NormalTok{      .attr("cx", xScale(0)).attr("cy", yScale(b))}
\NormalTok{      .attr("r", 6).attr("fill", "\#0d9488").attr("stroke", "\#fff").attr("stroke{-}width", 2);}
\NormalTok{    svg.append("text")}
\NormalTok{      .attr("x", xScale(0) + 10).attr("y", yScale(b) {-} 8)}
\NormalTok{      .attr("font{-}size", "11px").attr("fill", "\#0d9488")}
\NormalTok{      .text(\textasciigrave{}(0, $\{b\})\textasciigrave{});}
\NormalTok{  \}}

\NormalTok{  // Slope triangle — drawn from (1, b+m) showing rise and run}
\NormalTok{  const triX0 = 1, triY0 = m * triX0 + b, triX1 = triX0 + 1, triY1 = m * triX1 + b;}
\NormalTok{  const triVisible = triX0 \textgreater{}= X\_MIN \&\& triX1 \textless{}= X\_MAX \&\&}
\NormalTok{    Math.min(triY0, triY1) \textgreater{}= Y\_MIN \&\& Math.max(triY0, triY1) \textless{}= Y\_MAX;}

\NormalTok{  if (triVisible \&\& m !== 0) \{}
\NormalTok{    // Horizontal leg (run = 1)}
\NormalTok{    svg.append("line")}
\NormalTok{      .attr("x1", xScale(triX0)).attr("y1", yScale(triY0))}
\NormalTok{      .attr("x2", xScale(triX1)).attr("y2", yScale(triY0))}
\NormalTok{      .attr("stroke", "\#f59e0b").attr("stroke{-}width", 2).attr("stroke{-}dasharray", "3 2");}
\NormalTok{    svg.append("text")}
\NormalTok{      .attr("x", xScale(triX0 + 0.5)).attr("y", yScale(triY0) + 14)}
\NormalTok{      .attr("text{-}anchor", "middle").attr("font{-}size", "10px").attr("fill", "\#f59e0b")}
\NormalTok{      .text("run = 1");}
\NormalTok{    // Vertical leg (rise = m)}
\NormalTok{    svg.append("line")}
\NormalTok{      .attr("x1", xScale(triX1)).attr("y1", yScale(triY0))}
\NormalTok{      .attr("x2", xScale(triX1)).attr("y2", yScale(triY1))}
\NormalTok{      .attr("stroke", "\#7c3aed").attr("stroke{-}width", 2).attr("stroke{-}dasharray", "3 2");}
\NormalTok{    svg.append("text")}
\NormalTok{      .attr("x", xScale(triX1) + (m \textless{} 0 ? {-}48 : 10)).attr("y", yScale((triY0 + triY1) / 2) + 4)}
\NormalTok{      .attr("font{-}size", "10px").attr("fill", "\#7c3aed")}
\NormalTok{      .text(\textasciigrave{}rise = $\{m\}\textasciigrave{});}
\NormalTok{  \}}

\NormalTok{  // Equation label}
\NormalTok{  const mStr = m === 1 ? "" : m === {-}1 ? "{-}" : \textasciigrave{}$\{m\}\textasciigrave{};}
\NormalTok{  const bStr = b === 0 ? "" : b \textgreater{} 0 ? \textasciigrave{} + $\{b\}\textasciigrave{} : \textasciigrave{} − $\{Math.abs(b)\}\textasciigrave{};}
\NormalTok{  const eqStr = \textasciigrave{}y = $\{mStr\}x$\{bStr\}\textasciigrave{};}
\NormalTok{  svg.append("text")}
\NormalTok{    .attr("x", W / 2).attr("y", 28)}
\NormalTok{    .attr("text{-}anchor", "middle")}
\NormalTok{    .attr("font{-}size", "16px").attr("fill", "\#0d9488").attr("font{-}weight", "700")}
\NormalTok{    .text(eqStr);}

\NormalTok{  return svg.node();}
\NormalTok{\}}
\end{Highlighting}
\end{Shaded}

\section{Why this works}\label{why-this-works}

\begin{tcolorbox}[enhanced jigsaw, opacityback=0, colback=white, breakable, colframe=quarto-callout-note-color-frame, arc=.35mm, left=2mm, toprule=.15mm, leftrule=.75mm, bottomrule=.15mm, rightrule=.15mm]
\begin{minipage}[t]{5.5mm}
\textcolor{quarto-callout-note-color}{\faInfo}
\end{minipage}%
\begin{minipage}[t]{\textwidth - 5.5mm}

\textbf{The graph is the equation}

Every pair \((x, y)\) that satisfies \(y = mx + b\) lies on the line.
Every pair that doesn't satisfy it lies off the line. The line isn't a
picture of the equation --- it \emph{is} the equation, drawn in space.
This means you can read information off the graph the same way you'd get
it from the algebra: find an \(x\) on the horizontal axis, trace up to
the line, and read the matching \(y\) on the vertical axis. Or go the
other way: find a \(y\), trace across to the line, read the \(x\).

This two-way reading becomes powerful when you have two lines on the
same graph. Where they cross, both equations are satisfied at once ---
that intersection point is the solution to a system of two equations.

\end{minipage}%
\end{tcolorbox}

\section{Worked examples}\label{worked-examples-8}

\textbf{Example 1 --- Distance and time (Sciences).}

You cycle to school at a steady 12 km/h. Your school is 6 km away. Write
an equation for your distance from school, and graph it.

Let \(x\) = time in hours since you left home. Let \(y\) = distance from
home, in kilometres.

At a steady speed:

\[y = 12x\]

This is \(y = mx + b\) with \(m = 12\) and \(b = 0\).

\begin{longtable}[]{@{}cc@{}}
\toprule\noalign{}
\(x\) (hours) & \(y\) (km from home) \\
\midrule\noalign{}
\endhead
\bottomrule\noalign{}
\endlastfoot
0 & 0 \\
0.25 & 3 \\
0.5 & 6 \\
\end{longtable}

The slope is 12 --- for every hour, you cover 12 km. The y-intercept is
0 --- at time zero you're still at home. The graph is a line from the
origin, rising steeply. Reading from the graph: after 0.25 hours (15
minutes) you're 3 km along.

You arrive at school when \(y = 6\). Solve: \(12x = 6\), so \(x = 0.5\)
hours. That's 30 minutes. Confirmed by the table.

\textbf{Slope = speed} in a distance-time graph. This is not a
coincidence --- it's exactly what slope means physically.

\begin{center}\rule{0.5\linewidth}{0.5pt}\end{center}

\textbf{Example 2 --- Phone plan (Finance).}

Your data plan costs £8 per month fixed, plus £3 for each extra GB you
use. Write and graph the monthly cost equation.

\[y = 3x + 8\]

where \(x\) = extra GB used, \(y\) = total cost in £.

Method 2 (slope-intercept): Plot \((0, 8)\) --- that's the standing
charge. Slope is 3, so from \((0, 8)\) go 1 right and 3 up to
\((1, 11)\). Draw the line.

\begin{longtable}[]{@{}cc@{}}
\toprule\noalign{}
\(x\) (extra GB) & \(y\) (cost, £) \\
\midrule\noalign{}
\endhead
\bottomrule\noalign{}
\endlastfoot
0 & 8 \\
1 & 11 \\
2 & 14 \\
4 & 20 \\
\end{longtable}

Reading from the graph: 3 extra GB costs £17. When does the bill hit
£26? Set \(y = 26\): \(3x + 8 = 26\), so \(3x = 18\), \(x = 6\) GB. The
algebra and the graph tell the same story.

\begin{center}\rule{0.5\linewidth}{0.5pt}\end{center}

\textbf{Example 3 --- Points scored (Sport).}

A basketball player scores an average of 18 points per game. She has
already played 4 games. Write a formula for her total points, and find
when she'll have 144 total.

\[y = 18x + 72\]

Wait --- where does the 72 come from? She already has
\(18 \times 4 = 72\) points from those 4 games. So \(b = 72\) (her
running total at the start of tracking), and \(m = 18\) (she adds 18
each game).

Plot: start at \((0, 72)\) --- that's game zero from now, 72 points
already. Each game adds 18. From the graph, you can see she hits 144
points after 4 more games.

Check: \(18x + 72 = 144\), so \(18x = 72\), \(x = 4\). Yes.

\textbf{Slope = rate} --- in any graph where the y-axis tracks a total
and the x-axis tracks time or count, slope is the rate of change. How
fast the total grows.

\begin{center}\rule{0.5\linewidth}{0.5pt}\end{center}

\textbf{Example 4 --- Computing: screen coordinates.}

One thing worth knowing before this example: in maths, \(y\) increases
upward. On a screen --- in game engines, browser graphics, and most
image software --- \(y\) increases \emph{downward}, so \((0, 0)\) is the
top-left corner. This example uses mathematical coordinates. If you're
writing actual game code, a character ``moving up'' would have a
\emph{decreasing} \(y\) value in the engine's system.

A character starts at position \((x, y) = (50, 80)\) in mathematical
coordinates and moves right at 5 pixels per frame and down at 3 pixels
per frame.

After \(t\) frames: \[x = 50 + 5t \qquad y = 80 + 3t\]

These are two separate linear equations, one for each axis. After 10
frames: \(x = 100\), \(y = 110\). After 20 frames: \(x = 150\),
\(y = 140\). The path is a straight line, and its slope (in these
coordinates) is \(\frac{3}{5}\) --- 3 units down for every 5 across.

\section{A note on curves}\label{a-note-on-curves}

All the examples above produce straight lines. But not all relationships
are linear.

If you track your savings when you're earning interest on interest, the
graph curves upward --- it gets steeper as time goes on. If you track
how a hot drink cools, it curves downward and flattens out. A ball
thrown in the air traces a curve that dips back down.

A curve on a graph means the rate of change isn't constant. The slope
isn't fixed --- it depends on where you are. That's a very different
situation from a straight line, where the slope is the same everywhere.

You'll meet curved graphs properly when we reach polynomials and
functions in Volume 3. For now, notice the difference: a straight line
means a constant rate. A curve means the rate is changing.

\section{Where this goes}\label{where-this-goes-9}

The slope of a line is a constant --- pick any two points and you get
the same number. In Volume 3, \textbf{Functions and relations} makes
this more precise: you'll define a function formally, think about domain
and range, and meet graphs that aren't lines. From there, in Volume 5,
\textbf{Differential calculus} asks a harder question: if the slope of a
curve isn't constant, what is the slope \emph{at a single point}? The
answer --- the derivative --- is built on exactly the rise-over-run idea
from this chapter. You're already holding the seed of calculus in the
phrase ``change in \(y\) over change in \(x\).''

The second direction is \textbf{systems of equations}. You've already
seen that a linear equation graphs as a line. Two linear equations graph
as two lines. If they cross, the crossing point solves both equations at
once --- you can read the answer off the graph, or find it
algebraically. That's the foundation of linear algebra, which appears in
Volume 4 and underpins nearly everything in engineering and data
science.

\begin{tcolorbox}[enhanced jigsaw, opacityback=0, colback=white, breakable, colframe=quarto-callout-tip-color-frame, arc=.35mm, left=2mm, toprule=.15mm, leftrule=.75mm, bottomrule=.15mm, rightrule=.15mm]
\begin{minipage}[t]{5.5mm}
\textcolor{quarto-callout-tip-color}{\faLightbulb}
\end{minipage}%
\begin{minipage}[t]{\textwidth - 5.5mm}

\textbf{Where this shows up}

\begin{itemize}
\tightlist
\item
  A civil engineer plots a road's vertical profile as a distance-height
  graph. Slope is gradient, measured in percent. Regulations specify
  maximum gradients for different road types.
\item
  A physicist reads the slope of a velocity-time graph to find
  acceleration. The area under the graph gives distance --- a preview of
  integration.
\item
  A business analyst draws cost and revenue lines on the same graph. The
  crossing point is break-even: the number of units where cost equals
  revenue.
\item
  A game developer uses \((x, y)\) coordinates to position every object
  on screen. Movement is addition; rotation involves trig --- which also
  starts with the coordinate plane.
\end{itemize}

\end{minipage}%
\end{tcolorbox}

\section{Exercises}\label{exercises-9}

These are puzzles. Each one has a clean answer. The work is reading the
situation carefully, building the equation, and deciding whether the
graph or the algebra gives you what you need faster.

\begin{Shaded}
\begin{Highlighting}[]
\NormalTok{//| echo: false}
\NormalTok{katex = require("katex@0.16")}
\end{Highlighting}
\end{Shaded}

\begin{Shaded}
\begin{Highlighting}[]
\NormalTok{//| echo: false}
\NormalTok{function makeStepperHTML(exerciseNum, steps) \{}
\NormalTok{  let current = 0;}
\NormalTok{  const totalSteps = steps.length;}
\NormalTok{  const container = document.createElement("div");}
\NormalTok{  container.style.cssText = "border:1px solid \#e5e7eb; border{-}radius:8px; padding:1rem 1.25rem; margin:0.75rem 0; font{-}family:inherit;";}
\NormalTok{  const stepsDiv = document.createElement("div");}
\NormalTok{  stepsDiv.style.marginBottom = "0.75rem";}

\NormalTok{  function renderSteps() \{}
\NormalTok{    stepsDiv.innerHTML = "";}
\NormalTok{    for (let i = 0; i \textless{} current; i++) \{}
\NormalTok{      const s = steps[i];}
\NormalTok{      const row = document.createElement("div");}
\NormalTok{      row.style.cssText = "display:grid; grid{-}template{-}columns:200px 1fr; gap:0.35rem 1rem; align{-}items:baseline; padding:0.35rem 0; border{-}top:1px solid \#f3f4f6;";}
\NormalTok{      const opCell = document.createElement("span");}
\NormalTok{      opCell.style.cssText = "font{-}size:0.85em; color:\#6b7280; font{-}style:italic;";}
\NormalTok{      opCell.textContent = s.op;}
\NormalTok{      const eqCell = document.createElement("span");}
\NormalTok{      eqCell.innerHTML = katex.renderToString(s.eq, \{ throwOnError: false, displayMode: false \});}
\NormalTok{      row.appendChild(opCell);}
\NormalTok{      row.appendChild(eqCell);}
\NormalTok{      if (s.note) \{}
\NormalTok{        const noteRow = document.createElement("div");}
\NormalTok{        noteRow.style.cssText = "grid{-}column:2; font{-}size:0.82em; color:\#6b7280; padding{-}bottom:0.2rem;";}
\NormalTok{        noteRow.textContent = s.note;}
\NormalTok{        row.appendChild(noteRow);}
\NormalTok{      \}}
\NormalTok{      stepsDiv.appendChild(row);}
\NormalTok{    \}}
\NormalTok{    if (current === totalSteps) \{}
\NormalTok{      const done = document.createElement("div");}
\NormalTok{      done.style.cssText = "margin{-}top:0.5rem; font{-}size:0.9em; color:\#059669; font{-}weight:500;";}
\NormalTok{      done.textContent = "✓ Solution complete";}
\NormalTok{      stepsDiv.appendChild(done);}
\NormalTok{    \}}
\NormalTok{  \}}

\NormalTok{  const controls = document.createElement("div");}
\NormalTok{  controls.style.cssText = "display:flex; gap:0.5rem; align{-}items:center;";}
\NormalTok{  const nextBtn = document.createElement("button");}
\NormalTok{  nextBtn.textContent = "Next step →";}
\NormalTok{  nextBtn.style.cssText = "padding:0.35rem 0.85rem; border:1px solid \#d1d5db; border{-}radius:4px; background:\#fff; cursor:pointer; font{-}size:0.9em;";}
\NormalTok{  const resetBtn = document.createElement("button");}
\NormalTok{  resetBtn.textContent = "Reset";}
\NormalTok{  resetBtn.style.cssText = "padding:0.35rem 0.75rem; border:1px solid \#e5e7eb; border{-}radius:4px; background:\#f9fafb; cursor:pointer; font{-}size:0.9em; color:\#6b7280;";}
\NormalTok{  const counter = document.createElement("span");}
\NormalTok{  counter.style.cssText = "font{-}size:0.82em; color:\#9ca3af; margin{-}left:0.25rem;";}

\NormalTok{  function updateButtons() \{}
\NormalTok{    nextBtn.disabled = current === totalSteps;}
\NormalTok{    nextBtn.style.opacity = current === totalSteps ? "0.4" : "1";}
\NormalTok{    counter.textContent = current === 0 ? "Click to reveal steps" : \textasciigrave{}Step $\{current\} of $\{totalSteps\}\textasciigrave{};}
\NormalTok{  \}}

\NormalTok{  nextBtn.onclick = () =\textgreater{} \{ if (current \textless{} totalSteps) \{ current++; renderSteps(); updateButtons(); \} \};}
\NormalTok{  resetBtn.onclick = () =\textgreater{} \{ current = 0; renderSteps(); updateButtons(); \};}
\NormalTok{  controls.appendChild(nextBtn);}
\NormalTok{  controls.appendChild(resetBtn);}
\NormalTok{  controls.appendChild(counter);}
\NormalTok{  container.appendChild(stepsDiv);}
\NormalTok{  container.appendChild(controls);}
\NormalTok{  renderSteps();}
\NormalTok{  updateButtons();}
\NormalTok{  return container;}
\NormalTok{\}}
\end{Highlighting}
\end{Shaded}

\begin{enumerate}
\def\labelenumi{\arabic{enumi}.}
\tightlist
\item
  A streaming service costs £6 per month. A rival service costs £2 per
  month plus £1 per film watched. Write equations for both costs and
  find how many films per month makes them equal. Which service is
  cheaper if you watch 5 films a month?
\end{enumerate}

\begin{Shaded}
\begin{Highlighting}[]
\NormalTok{//| echo: false}
\NormalTok{makeStepperHTML(1, [}
\NormalTok{  \{ op: "Write as an equation: Service A",  eq: "y = 6" \},}
\NormalTok{  \{ op: "Write as an equation: Service B",  eq: "y = x + 2" \},}
\NormalTok{  \{ op: "Set them equal",                   eq: "6 = x + 2" \},}
\NormalTok{  \{ op: "Solve for x",                      eq: "x = 4 \textbackslash{}\textbackslash{}text\{ films\}" \},}
\NormalTok{  \{ op: "Compare at 5 films",               eq: "A: £6,\textbackslash{}\textbackslash{}quad B: 5 + 2 = £7", note: "Service A is cheaper above 4 films per month." \},}
\NormalTok{  \{ op: "Check",                            eq: "x=4:\textbackslash{}\textbackslash{}; 4+2=6 \textbackslash{}\textbackslash{}checkmark" \}}
\NormalTok{])}
\end{Highlighting}
\end{Shaded}

\begin{enumerate}
\def\labelenumi{\arabic{enumi}.}
\setcounter{enumi}{1}
\tightlist
\item
  A delivery driver starts a shift with a full tank holding 60 litres.
  She uses fuel at 0.08 litres per kilometre. Write an equation for fuel
  remaining after \(x\) km. What is the slope, and what does it mean in
  context? How far can she drive before the tank is empty?
\end{enumerate}

\begin{Shaded}
\begin{Highlighting}[]
\NormalTok{//| echo: false}
\NormalTok{makeStepperHTML(2, [}
\NormalTok{  \{ op: "Write as an equation",           eq: "y = 60 {-} 0.08x" \},}
\NormalTok{  \{ op: "Identify the slope",             eq: "m = {-}0.08 \textbackslash{}\textbackslash{}text\{ litres per km\}", note: "Negative slope means the fuel decreases as distance increases." \},}
\NormalTok{  \{ op: "Set y = 0 (tank empty)",         eq: "0 = 60 {-} 0.08x" \},}
\NormalTok{  \{ op: "Solve for x",                    eq: "0.08x = 60 \textbackslash{}\textbackslash{}Rightarrow x = 750 \textbackslash{}\textbackslash{}text\{ km\}" \},}
\NormalTok{  \{ op: "Check",                          eq: "60 {-} 0.08 \textbackslash{}\textbackslash{}times 750 = 60 {-} 60 = 0 \textbackslash{}\textbackslash{}checkmark" \}}
\NormalTok{])}
\end{Highlighting}
\end{Shaded}

\begin{enumerate}
\def\labelenumi{\arabic{enumi}.}
\setcounter{enumi}{2}
\tightlist
\item
  Two friends start at the same place and walk in opposite directions.
  Priya walks at 5 km/h; Jordan walks at 4 km/h. Write an equation for
  the distance between them after \(t\) hours. What is the slope? After
  how long are they 13.5 km apart?
\end{enumerate}

\begin{Shaded}
\begin{Highlighting}[]
\NormalTok{//| echo: false}
\NormalTok{makeStepperHTML(3, [}
\NormalTok{  \{ op: "Write as an equation",           eq: "d = 5t + 4t = 9t", note: "They move away from each other, so total distance grows at 9 km/h." \},}
\NormalTok{  \{ op: "Identify the slope",             eq: "m = 9 \textbackslash{}\textbackslash{}text\{ km/h (the rate at which they separate)\}" \},}
\NormalTok{  \{ op: "Set d = 13.5",                   eq: "9t = 13.5" \},}
\NormalTok{  \{ op: "Solve for t",                    eq: "t = 1.5 \textbackslash{}\textbackslash{}text\{ hours\}" \},}
\NormalTok{  \{ op: "Check",                          eq: "9 \textbackslash{}\textbackslash{}times 1.5 = 13.5 \textbackslash{}\textbackslash{}checkmark" \}}
\NormalTok{])}
\end{Highlighting}
\end{Shaded}

\begin{enumerate}
\def\labelenumi{\arabic{enumi}.}
\setcounter{enumi}{3}
\tightlist
\item
  A runner's training log shows she ran 15 km in week 1 and increases
  her distance by 2 km each week. Write an equation for total distance
  run after \(w\) weeks of training (not counting week 1 as week 0 ---
  let week 1 be \(w = 1\)). Plot the first six weeks. In which week does
  she first exceed 100 km total?
\end{enumerate}

\begin{Shaded}
\begin{Highlighting}[]
\NormalTok{//| echo: false}
\NormalTok{makeStepperHTML(4, [}
\NormalTok{  \{ op: "Write the weekly distance",       eq: "d\_w = 15 + (w{-}1) \textbackslash{}\textbackslash{}times 2 = 2w + 13" \},}
\NormalTok{  \{ op: "Sum the first w weeks",           eq: "T\_w = \textbackslash{}\textbackslash{}sum\_\{k=1\}\^{}\{w\}(2k+13) = w\^{}2 + 14w", note: "Sum of 2k+13 for k=1 to w equals w²+14w using the arithmetic series formula." \},}
\NormalTok{  \{ op: "Set T\_w \textgreater{} 100",                   eq: "w\^{}2 + 14w \textgreater{} 100" \},}
\NormalTok{  \{ op: "Try w = 6",                       eq: "36 + 84 = 120 \textgreater{} 100 \textbackslash{}\textbackslash{}checkmark" \},}
\NormalTok{  \{ op: "Try w = 5",                       eq: "25 + 70 = 95 \textless{} 100" \},}
\NormalTok{  \{ op: "Answer",                          eq: "\textbackslash{}\textbackslash{}text\{She first exceeds 100 km total in week 6\}" \}}
\NormalTok{])}
\end{Highlighting}
\end{Shaded}

\begin{enumerate}
\def\labelenumi{\arabic{enumi}.}
\setcounter{enumi}{4}
\tightlist
\item
  A line passes through the points \((2, 7)\) and \((5, 13)\). Find the
  slope. Then substitute one of the points into \(y = mx + b\) to solve
  for the y-intercept \(b\), and write the full equation. Check your
  equation works for both original points.
\end{enumerate}

\begin{Shaded}
\begin{Highlighting}[]
\NormalTok{//| echo: false}
\NormalTok{makeStepperHTML(5, [}
\NormalTok{  \{ op: "Calculate slope",                eq: "m = \textbackslash{}\textbackslash{}frac\{13{-}7\}\{5{-}2\} = \textbackslash{}\textbackslash{}frac\{6\}\{3\} = 2" \},}
\NormalTok{  \{ op: "Substitute (2, 7) into y=mx+b", eq: "7 = 2(2) + b \textbackslash{}\textbackslash{}Rightarrow 7 = 4 + b \textbackslash{}\textbackslash{}Rightarrow b = 3" \},}
\NormalTok{  \{ op: "Write the equation",             eq: "y = 2x + 3" \},}
\NormalTok{  \{ op: "Check point (2, 7)",             eq: "2(2)+3 = 7 \textbackslash{}\textbackslash{}checkmark" \},}
\NormalTok{  \{ op: "Check point (5, 13)",            eq: "2(5)+3 = 13 \textbackslash{}\textbackslash{}checkmark" \}}
\NormalTok{])}
\end{Highlighting}
\end{Shaded}

\begin{enumerate}
\def\labelenumi{\arabic{enumi}.}
\setcounter{enumi}{5}
\tightlist
\item
  A line has slope \(-2\) and passes through \((3, 4)\). Substitute the
  point into \(y = mx + b\) to find \(b\), then write the full equation.
  Find where the line crosses the x-axis (the point where \(y = 0\)).
\end{enumerate}

\begin{Shaded}
\begin{Highlighting}[]
\NormalTok{//| echo: false}
\NormalTok{makeStepperHTML(6, [}
\NormalTok{  \{ op: "Substitute into y = mx + b",    eq: "4 = {-}2(3) + b \textbackslash{}\textbackslash{}Rightarrow 4 = {-}6 + b \textbackslash{}\textbackslash{}Rightarrow b = 10" \},}
\NormalTok{  \{ op: "Write the equation",             eq: "y = {-}2x + 10" \},}
\NormalTok{  \{ op: "Find x{-}intercept (set y = 0)",  eq: "0 = {-}2x + 10 \textbackslash{}\textbackslash{}Rightarrow 2x = 10 \textbackslash{}\textbackslash{}Rightarrow x = 5" \},}
\NormalTok{  \{ op: "x{-}intercept",                   eq: "(5,\textbackslash{}\textbackslash{}, 0)" \},}
\NormalTok{  \{ op: "Check",                          eq: "{-}2(5)+10 = 0 \textbackslash{}\textbackslash{}checkmark;\textbackslash{}\textbackslash{}quad {-}2(3)+10 = 4 \textbackslash{}\textbackslash{}checkmark" \}}
\NormalTok{])}
\end{Highlighting}
\end{Shaded}

\begin{enumerate}
\def\labelenumi{\arabic{enumi}.}
\setcounter{enumi}{6}
\tightlist
\item
  Two phone plans are graphed below --- but the graph has been removed.
  Plan A: \(y = 4x + 5\). Plan B: \(y = 2x + 15\). Without drawing the
  graph, find the value of \(x\) where the plans cost the same. Which
  plan is cheaper for someone using 4 units? Which for someone using 8
  units?
\end{enumerate}

\begin{Shaded}
\begin{Highlighting}[]
\NormalTok{//| echo: false}
\NormalTok{makeStepperHTML(7, [}
\NormalTok{  \{ op: "Write as an equation",           eq: "4x + 5 = 2x + 15" \},}
\NormalTok{  \{ op: "Subtract 2x from both sides",    eq: "2x + 5 = 15" \},}
\NormalTok{  \{ op: "Subtract 5 from both sides",     eq: "2x = 10 \textbackslash{}\textbackslash{}Rightarrow x = 5" \},}
\NormalTok{  \{ op: "Compare at x = 4",              eq: "A: 4(4)+5=21,\textbackslash{}\textbackslash{}quad B: 2(4)+15=23 \textbackslash{}\textbackslash{}Rightarrow \textbackslash{}\textbackslash{}text\{Plan A cheaper\}" \},}
\NormalTok{  \{ op: "Compare at x = 8",              eq: "A: 4(8)+5=37,\textbackslash{}\textbackslash{}quad B: 2(8)+15=31 \textbackslash{}\textbackslash{}Rightarrow \textbackslash{}\textbackslash{}text\{Plan B cheaper\}" \},}
\NormalTok{  \{ op: "Check crossover at x = 5",      eq: "A: 4(5)+5=25,\textbackslash{}\textbackslash{}quad B: 2(5)+15=25 \textbackslash{}\textbackslash{}checkmark" \}}
\NormalTok{])}
\end{Highlighting}
\end{Shaded}

\part{Volume 3 --- Shape and measure}

\chapter{Volume 3 --- Shape and
measure}\label{volume-3-shape-and-measure-1}

The central move of this volume is giving coordinates to space.

You already know how to calculate and balance equations. This volume
shows you how algebraic tools can be used to describe the physical world
and how the world, in turn, can give meaning to abstract math.

By the end of this volume, you will be able to map distances, relate
angles to side lengths, understand the structures behind non-linear
relationships, and formally describe how different parameters interact.

This is the bridge where algebra stops being just letters on a page and
starts becoming a map of reality.

\section{Chapters}\label{chapters-2}

\begin{longtable}[]{@{}lll@{}}
\toprule\noalign{}
Chapter & Title & Core move \\
\midrule\noalign{}
\endhead
\bottomrule\noalign{}
\endlastfoot
1 & Geometry and measurement & Quantifying physical space \\
2 & Trigonometry & Relating angles to side lengths \\
3 & Polynomials and factoring & Deconstructing complex equations \\
4 & Functions and relations & Mapping inputs to outputs \\
\end{longtable}

\chapter{Geometry and measurement}\label{geometry-and-measurement}

Reading the shape of the world

\hfill\break

Your bedroom floor has a certain area. You know this even without
measuring it --- you can picture whether a rug would fit, whether the
furniture leaves enough space, whether it's bigger or smaller than a
friend's room.

Your phone screen is described by a diagonal measurement in inches. But
the area that matters --- the bit you actually look at --- is the width
times the height, not the diagonal. The diagonal is a shortcut that
hides the real shape.

The skateboard ramp at the park is a triangle in cross-section. If you
built one, you'd need to know the length of the slope, the height it
reaches, and the distance it runs along the ground --- and how those
three numbers relate to each other.

These are all geometry. This chapter gives you the tools to read the
shape of any object you encounter --- and to find the measurement you
need from the ones you have.

\section{What the notation is
saying}\label{what-the-notation-is-saying-10}

Every formula in this chapter is shorthand for a relationship between
measurements. Before using any of them, here is what each symbol means.

\textbf{\(A\)} --- area. The amount of flat space a shape covers.
Measured in square units: square centimetres (cm²), square metres (m²),
square kilometres (km²). If a length is in centimetres, the area will be
in square centimetres.

\textbf{\(P\)} --- perimeter. The total distance around the outside edge
of a flat shape. Measured in the same units as the lengths themselves
--- not squared.

\textbf{\(r\)} --- radius. The distance from the centre of a circle to
any point on its edge. The diameter \(d\) is twice the radius:
\(d = 2r\).

\textbf{\(\pi\)} --- pi. Read it as ``pie''. It's approximately
\(3.14159\), but it never terminates or repeats --- it's an irrational
number. It appears in every circle formula because it captures the
relationship between a circle's diameter and the distance around it.
Your calculator has a \(\pi\) button; use it rather than the
approximation unless you're told otherwise.

\textbf{\(V\)} --- volume. The amount of three-dimensional space a solid
object fills. Measured in cubic units: cm³, m³, litres (1 litre = 1000
cm³).

\textbf{\(h\)} --- height. In a prism or cylinder, this is the
perpendicular distance between the two bases --- how tall the solid is,
measured straight through, at a right angle to the base.

\textbf{\(a^2 + b^2 = c^2\)} --- the Pythagorean theorem. Read it as:
the square of \(a\) plus the square of \(b\) equals the square of \(c\).
Here \(a\) and \(b\) are the two shorter sides of a right-angled
triangle and \(c\) is the longest side --- the one opposite the right
angle, called the \textbf{hypotenuse}. This only works for right
triangles. The right angle is marked with a small square in diagrams.

\section{The method}\label{the-method-9}

\subsection{Rectangles}\label{rectangles}

A rectangle has two pairs of equal sides. Call the longer side the
length \(l\) and the shorter side the width \(w\).

\textbf{Area} --- multiply length by width:

\[A = l \times w\]

A bedroom that is 4 metres long and 3 metres wide has area
\(4 \times 3 = 12\) m².

\textbf{Perimeter} --- add up all four sides. Since opposite sides are
equal, that's two lengths plus two widths:

\[P = 2l + 2w\]

The same bedroom has perimeter \(2(4) + 2(3) = 8 + 6 = 14\) m.

\subsection{Triangles}\label{triangles}

A triangle has three sides and three angles. For area, you need a
\textbf{base} and a \textbf{height}. The base is any side you choose.
The height is the perpendicular distance from that base to the opposite
corner --- straight up, at a right angle to the base.

\textbf{Area:}

\[A = \frac{1}{2} \times b \times h\]

Read it as: half of base times height. The halving is there because a
triangle is half a rectangle --- if you put two identical triangles
together, they make a parallelogram with the same base and height.

A triangle with base 6 cm and height 4 cm has area
\(\frac{1}{2} \times 6 \times 4 = 12\) cm².

\textbf{Perimeter} --- just add all three sides. There's no shortcut
unless the triangle has special properties.

\subsection{Circles}\label{circles}

\textbf{Circumference} --- the distance around the edge of a circle:

\[C = 2\pi r\]

Read it as: two times pi times radius. Since \(d = 2r\), this is the
same as \(\pi d\) --- pi times diameter. That's actually where \(\pi\)
comes from: the ratio of any circle's circumference to its diameter.
Every circle, everywhere, has \(C = \pi d\).

\textbf{Area} --- the amount of flat space a circle covers:

\[A = \pi r^2\]

Read it as: pi times radius squared. A circle with radius 5 cm has area
\(\pi \times 5^2 = 25\pi \approx 78.5\) cm².

Note: leave the answer as \(25\pi\) when exact form is needed. Use the
decimal approximation when you need a number to compare or measure with.

\begin{Shaded}
\begin{Highlighting}[]
\NormalTok{//| echo: false}
\NormalTok{viewof circleRadius = Inputs.range([1, 10], \{}
\NormalTok{  value: 5,}
\NormalTok{  step: 0.5,}
\NormalTok{  label: "Radius r (units)"}
\NormalTok{\})}
\end{Highlighting}
\end{Shaded}

\begin{Shaded}
\begin{Highlighting}[]
\NormalTok{//| echo: false}
\NormalTok{\{}
\NormalTok{  const r = circleRadius;}
\NormalTok{  const scale = 20;}
\NormalTok{  const svgW = 480, svgH = 260;}
\NormalTok{  const cx = 140, cy = svgH / 2;}

\NormalTok{  const area = Math.PI * r * r;}
\NormalTok{  const circ = 2 * Math.PI * r;}

\NormalTok{  const svg = d3.create("svg")}
\NormalTok{    .attr("width", svgW)}
\NormalTok{    .attr("height", svgH)}
\NormalTok{    .style("font{-}family", "inherit");}

\NormalTok{  // Circle}
\NormalTok{  svg.append("circle")}
\NormalTok{    .attr("cx", cx)}
\NormalTok{    .attr("cy", cy)}
\NormalTok{    .attr("r", r * scale)}
\NormalTok{    .attr("fill", "\#dbeafe")}
\NormalTok{    .attr("stroke", "\#3b82f6")}
\NormalTok{    .attr("stroke{-}width", 2);}

\NormalTok{  // Radius line}
\NormalTok{  svg.append("line")}
\NormalTok{    .attr("x1", cx).attr("y1", cy)}
\NormalTok{    .attr("x2", cx + r * scale).attr("y2", cy)}
\NormalTok{    .attr("stroke", "\#1d4ed8")}
\NormalTok{    .attr("stroke{-}width", 1.5)}
\NormalTok{    .attr("stroke{-}dasharray", "4,3");}

\NormalTok{  // Radius label}
\NormalTok{  svg.append("text")}
\NormalTok{    .attr("x", cx + (r * scale) / 2)}
\NormalTok{    .attr("y", cy {-} 6)}
\NormalTok{    .attr("text{-}anchor", "middle")}
\NormalTok{    .attr("font{-}size", "13px")}
\NormalTok{    .attr("fill", "\#1d4ed8")}
\NormalTok{    .text(\textasciigrave{}r = $\{r\}\textasciigrave{});}

\NormalTok{  // Stats panel}
\NormalTok{  const px = 310, py = 60;}
\NormalTok{  svg.append("rect")}
\NormalTok{    .attr("x", px {-} 10).attr("y", py {-} 24)}
\NormalTok{    .attr("width", 170).attr("height", 110)}
\NormalTok{    .attr("fill", "\#f8fafc")}
\NormalTok{    .attr("stroke", "\#e2e8f0")}
\NormalTok{    .attr("rx", 6);}

\NormalTok{  svg.append("text").attr("x", px).attr("y", py)}
\NormalTok{    .attr("font{-}size", "13px").attr("font{-}weight", "600").attr("fill", "\#374151")}
\NormalTok{    .text("Measurements");}

\NormalTok{  svg.append("text").attr("x", px).attr("y", py + 26)}
\NormalTok{    .attr("font{-}size", "13px").attr("fill", "\#374151")}
\NormalTok{    .text(\textasciigrave{}Area = πr²\textasciigrave{});}

\NormalTok{  svg.append("text").attr("x", px).attr("y", py + 44)}
\NormalTok{    .attr("font{-}size", "13px").attr("fill", "\#2563eb")}
\NormalTok{    .text(\textasciigrave{}= $\{area.toFixed(2)\} sq units\textasciigrave{});}

\NormalTok{  svg.append("text").attr("x", px).attr("y", py + 68)}
\NormalTok{    .attr("font{-}size", "13px").attr("fill", "\#374151")}
\NormalTok{    .text(\textasciigrave{}Circumference = 2πr\textasciigrave{});}

\NormalTok{  svg.append("text").attr("x", px).attr("y", py + 86)}
\NormalTok{    .attr("font{-}size", "13px").attr("fill", "\#059669")}
\NormalTok{    .text(\textasciigrave{}= $\{circ.toFixed(2)\} units\textasciigrave{});}

\NormalTok{  return svg.node();}
\NormalTok{\}}
\end{Highlighting}
\end{Shaded}

\subsection{The Pythagorean theorem}\label{the-pythagorean-theorem}

For any right-angled triangle with legs \(a\) and \(b\) and hypotenuse
\(c\):

\[a^2 + b^2 = c^2\]

\textbf{Finding the hypotenuse} when you know the two legs:

\[c = \sqrt{a^2 + b^2}\]

A right triangle with legs 3 cm and 4 cm:

\[c = \sqrt{3^2 + 4^2} = \sqrt{9 + 16} = \sqrt{25} = 5 \text{ cm}\]

\textbf{Finding a leg} when you know the hypotenuse and the other leg:

\[a = \sqrt{c^2 - b^2}\]

A right triangle with hypotenuse 13 cm and one leg 5 cm:

\[a = \sqrt{13^2 - 5^2} = \sqrt{169 - 25} = \sqrt{144} = 12 \text{ cm}\]

The next step looks slightly mechanical. Square both known values, add
(or subtract), then take the square root. In that order, every time.

\begin{Shaded}
\begin{Highlighting}[]
\NormalTok{//| echo: false}
\NormalTok{viewof legA = Inputs.range([1, 8], \{ value: 3, step: 1, label: "Leg a" \})}
\end{Highlighting}
\end{Shaded}

\begin{Shaded}
\begin{Highlighting}[]
\NormalTok{//| echo: false}
\NormalTok{viewof legB = Inputs.range([1, 8], \{ value: 4, step: 1, label: "Leg b" \})}
\end{Highlighting}
\end{Shaded}

\begin{Shaded}
\begin{Highlighting}[]
\NormalTok{//| echo: false}
\NormalTok{\{}
\NormalTok{  const a = legA, b = legB;}
\NormalTok{  const c = Math.sqrt(a * a + b * b);}
\NormalTok{  const scale = 28;}
\NormalTok{  const margin = 60;}
\NormalTok{  const svgW = 520, svgH = 340;}

\NormalTok{  // Triangle vertices: right angle at bottom{-}left}
\NormalTok{  const ox = margin + 20, oy = svgH {-} margin {-} 20;}
\NormalTok{  const ax = ox + b * scale, ay = oy;       // bottom{-}right}
\NormalTok{  const bx = ox, by = oy {-} a * scale;       // top{-}left}

\NormalTok{  const svg = d3.create("svg")}
\NormalTok{    .attr("width", svgW)}
\NormalTok{    .attr("height", svgH)}
\NormalTok{    .style("font{-}family", "inherit");}

\NormalTok{  // Square on side b (bottom, along x{-}axis)}
\NormalTok{  svg.append("rect")}
\NormalTok{    .attr("x", ox).attr("y", oy)}
\NormalTok{    .attr("width", b * scale).attr("height", b * scale)}
\NormalTok{    .attr("fill", "\#fef9c3").attr("stroke", "\#ca8a04").attr("stroke{-}width", 1.5);}
\NormalTok{  svg.append("text")}
\NormalTok{    .attr("x", ox + b * scale / 2).attr("y", oy + b * scale / 2 + 5)}
\NormalTok{    .attr("text{-}anchor", "middle").attr("font{-}size", "13px").attr("fill", "\#92400e")}
\NormalTok{    .text(\textasciigrave{}b² = $\{b*b\}\textasciigrave{});}

\NormalTok{  // Square on side a (left, along y{-}axis) — goes to the left}
\NormalTok{  svg.append("rect")}
\NormalTok{    .attr("x", ox {-} a * scale).attr("y", by)}
\NormalTok{    .attr("width", a * scale).attr("height", a * scale)}
\NormalTok{    .attr("fill", "\#dcfce7").attr("stroke", "\#16a34a").attr("stroke{-}width", 1.5);}
\NormalTok{  svg.append("text")}
\NormalTok{    .attr("x", ox {-} a * scale / 2).attr("y", by + a * scale / 2 + 5)}
\NormalTok{    .attr("text{-}anchor", "middle").attr("font{-}size", "13px").attr("fill", "\#166534")}
\NormalTok{    .text(\textasciigrave{}a² = $\{a*a\}\textasciigrave{});}

\NormalTok{  // Square on hypotenuse c — rotated square drawn as polygon}
\NormalTok{  // hypotenuse goes from (ax,ay) to (bx,by)}
\NormalTok{  // perpendicular direction (outward): normalise then rotate 90°}
\NormalTok{  const dx = bx {-} ax, dy = by {-} ay;}
\NormalTok{  const len = Math.sqrt(dx*dx + dy*dy);}
\NormalTok{  const nx = {-}dy / len, ny = dx / len; // outward normal}
\NormalTok{  const p1 = \{x: ax, y: ay\};}
\NormalTok{  const p2 = \{x: bx, y: by\};}
\NormalTok{  const p3 = \{x: bx + nx * len, y: by + ny * len\};}
\NormalTok{  const p4 = \{x: ax + nx * len, y: ay + ny * len\};}
\NormalTok{  const pts = [p1,p2,p3,p4].map(p =\textgreater{} \textasciigrave{}$\{p.x\},$\{p.y\}\textasciigrave{}).join(" ");}

\NormalTok{  svg.append("polygon")}
\NormalTok{    .attr("points", pts)}
\NormalTok{    .attr("fill", "\#dbeafe").attr("stroke", "\#2563eb").attr("stroke{-}width", 1.5);}

\NormalTok{  const midCx = (p1.x + p2.x + p3.x + p4.x) / 4;}
\NormalTok{  const midCy = (p1.y + p2.y + p3.y + p4.y) / 4;}
\NormalTok{  svg.append("text")}
\NormalTok{    .attr("x", midCx).attr("y", midCy + 5)}
\NormalTok{    .attr("text{-}anchor", "middle").attr("font{-}size", "13px").attr("fill", "\#1e3a8a")}
\NormalTok{    .text(\textasciigrave{}c² = $\{(c*c).toFixed(2)\}\textasciigrave{});}

\NormalTok{  // Triangle fill}
\NormalTok{  svg.append("polygon")}
\NormalTok{    .attr("points", \textasciigrave{}$\{ox\},$\{oy\} $\{ax\},$\{ay\} $\{bx\},$\{by\}\textasciigrave{})}
\NormalTok{    .attr("fill", "\#f1f5f9").attr("stroke", "\#475569").attr("stroke{-}width", 2);}

\NormalTok{  // Right{-}angle marker}
\NormalTok{  const sq = 8;}
\NormalTok{  svg.append("rect")}
\NormalTok{    .attr("x", ox).attr("y", oy {-} sq)}
\NormalTok{    .attr("width", sq).attr("height", sq)}
\NormalTok{    .attr("fill", "none").attr("stroke", "\#475569").attr("stroke{-}width", 1.5);}

\NormalTok{  // Side labels}
\NormalTok{  svg.append("text")}
\NormalTok{    .attr("x", ox + b * scale / 2).attr("y", oy {-} 8)}
\NormalTok{    .attr("text{-}anchor", "middle").attr("font{-}size", "13px").attr("fill", "\#374151")}
\NormalTok{    .text(\textasciigrave{}b = $\{b\}\textasciigrave{});}

\NormalTok{  svg.append("text")}
\NormalTok{    .attr("x", ox {-} 14).attr("y", oy {-} a * scale / 2)}
\NormalTok{    .attr("text{-}anchor", "middle").attr("font{-}size", "13px").attr("fill", "\#374151")}
\NormalTok{    .text(\textasciigrave{}a = $\{a\}\textasciigrave{});}

\NormalTok{  // Equation below}
\NormalTok{  svg.append("text")}
\NormalTok{    .attr("x", svgW / 2).attr("y", svgH {-} 16)}
\NormalTok{    .attr("text{-}anchor", "middle").attr("font{-}size", "14px").attr("fill", "\#111827")}
\NormalTok{    .text(\textasciigrave{}$\{a\}² + $\{b\}² = $\{a*a\} + $\{b*b\} = $\{(a*a+b*b).toFixed(2)\}   →   c = $\{c.toFixed(3)\}\textasciigrave{});}

\NormalTok{  return svg.node();}
\NormalTok{\}}
\end{Highlighting}
\end{Shaded}

\begin{tcolorbox}[enhanced jigsaw, opacityback=0, colback=white, breakable, colframe=quarto-callout-note-color-frame, arc=.35mm, left=2mm, toprule=.15mm, leftrule=.75mm, bottomrule=.15mm, rightrule=.15mm]
\begin{minipage}[t]{5.5mm}
\textcolor{quarto-callout-note-color}{\faInfo}
\end{minipage}%
\begin{minipage}[t]{\textwidth - 5.5mm}

\textbf{Why the Pythagorean theorem works}

Draw a right triangle. Now draw a square on each side --- one square on
leg \(a\), one on leg \(b\), one on the hypotenuse \(c\). The theorem
says the area of the big square exactly equals the total area of the two
smaller squares: \(a^2 + b^2 = c^2\).

You can show this is true by cutting up the two smaller squares and
rearranging the pieces to exactly fill the big one --- no gaps, no
overlaps. That's the core of most visual proofs of this result. The
algebra follows from the geometry.

This only holds when the angle between sides \(a\) and \(b\) is exactly
90°. Change that angle and the relationship changes --- that's what
trigonometry (the next chapter) is about.

\end{minipage}%
\end{tcolorbox}

\subsection{Volume of a rectangular
prism}\label{volume-of-a-rectangular-prism}

A rectangular prism (a box shape) has length \(l\), width \(w\), and
height \(h\).

\[V = l \times w \times h\]

Read it as: length times width times height. A box 10 cm long, 5 cm
wide, and 3 cm tall has volume \(10 \times 5 \times 3 = 150\) cm³.

A useful way to think about it: the base area is \(l \times w\), and
you're stacking that base up through height \(h\). Volume is always base
area times height for any prism.

\subsection{Volume of a cylinder}\label{volume-of-a-cylinder}

A cylinder is a circle stacked up through a height. The base is a circle
of radius \(r\), and the height is \(h\).

\[V = \pi r^2 h\]

Read it as: pi times radius squared times height. The \(\pi r^2\) part
is just the area of the circular base; multiplying by \(h\) stacks it
up.

A can of drink is a cylinder. A standard 330 ml can has a radius of
about 3.3 cm and a height of about 9.6 cm:

\[V = \pi \times 3.3^2 \times 9.6 = \pi \times 10.89 \times 9.6 \approx 328 \text{ cm}^3\]

Close to 330 ml (330 cm³). The small difference is rounding on the
actual dimensions.

\subsection{Units and unit conversion}\label{units-and-unit-conversion}

Measurements only make sense with their units attached. \(15\) is not an
area --- \(15 \text{ m}^2\) is.

When you multiply lengths, the units multiply too:

\[3 \text{ m} \times 4 \text{ m} = 12 \text{ m}^2\]

\[2 \text{ cm} \times 5 \text{ cm} \times 3 \text{ cm} = 30 \text{ cm}^3\]

\textbf{Converting between units:} multiply by a conversion factor equal
to 1. For example, \(1 \text{ m} = 100 \text{ cm}\), so
\(\frac{100 \text{ cm}}{1 \text{ m}} = 1\).

\[5 \text{ m} = 5 \times 100 \text{ cm} = 500 \text{ cm}\]

For areas, the conversion factor gets squared:

\[1 \text{ m}^2 = (100 \text{ cm})^2 = 10{,}000 \text{ cm}^2\]

So \(3 \text{ m}^2 = 3 \times 10{,}000 = 30{,}000 \text{ cm}^2\).

Why squared? A square metre is a square that is 100 cm wide and 100 cm
tall --- that's \(100 \times 100 = 10{,}000\) square centimetres. The
conversion factor gets multiplied by itself once for each dimension
you're measuring.

For volumes, it gets cubed:

\[1 \text{ m}^3 = (100)^3 \text{ cm}^3 = 1{,}000{,}000 \text{ cm}^3\]

A cubic metre is a cube 100 cm × 100 cm × 100 cm, so the conversion
factor multiplies by itself three times. This trips people up. One
square metre is not a hundred square centimetres --- it's ten thousand.
The number gets big fast because you're covering a two-dimensional (or
three-dimensional) space.

\section{Worked examples}\label{worked-examples-9}

\textbf{Example 1 (sci) --- Area of a circular clock face.}

A clock has a circular face with a diameter of 20 cm. What is the area
of the face?

Diameter = 20 cm, so radius = 10 cm.

\[A = \pi r^2 = \pi \times 10^2 = 100\pi \approx 314 \text{ cm}^2\]

The clock face covers approximately 314 cm² --- about the same as a
standard sheet of A4 paper turned on its side. Doubling the diameter
would quadruple the area, because radius is squared in the formula.

\begin{center}\rule{0.5\linewidth}{0.5pt}\end{center}

\textbf{Example 2 (eng) --- Length of a ramp.}

A skateboard ramp rises 1.2 m vertically and runs 1.6 m horizontally
along the ground. The ramp surface is the hypotenuse of a right
triangle. How long is the ramp surface?

\[c = \sqrt{a^2 + b^2} = \sqrt{1.2^2 + 1.6^2} = \sqrt{1.44 + 2.56} = \sqrt{4.00} = 2.0 \text{ m}\]

The ramp surface is exactly 2 m long.

\emph{Note: 1.2, 1.6, 2.0 is a scaled version of the 3-4-5 right
triangle (multiply each by 0.4). When you see a clean answer, check
whether it's a Pythagorean triple.}

\begin{center}\rule{0.5\linewidth}{0.5pt}\end{center}

\textbf{Example 3 (comp) --- Screen area from diagonal and aspect
ratio.}

A phone screen has a 6.1-inch diagonal. The screen is 16:9 aspect ratio
(16 units wide for every 9 units tall). What is the actual screen area?

One approach: call the width \(16k\) and height \(9k\), where \(k\) is a
scaling factor that sets the actual size. The Pythagorean theorem then
gives an equation in \(k\):

\[(16k)^2 + (9k)^2 = 6.1^2\]

\[256k^2 + 81k^2 = 37.21\]

\[337k^2 = 37.21\]

\[k^2 = \frac{37.21}{337} \approx 0.11041\]

\[k \approx 0.3323 \text{ inches}\]

Width \(= 16 \times 0.3323 \approx 5.32\) inches. Height
\(= 9 \times 0.3323 \approx 2.99\) inches.

Area \(= 5.32 \times 2.99 \approx 15.9\) square inches.

The 6.1-inch measurement is the diagonal --- the actual viewing area is
about 15.9 square inches, which is roughly the size of a playing card.

\begin{center}\rule{0.5\linewidth}{0.5pt}\end{center}

\textbf{Example 4 (geo) --- Comparing two screen sizes.}

You're deciding between a phone screen and a tablet screen. The phone
screen measures 6 cm by 12 cm on the poster. The tablet screen is shown
on a second poster at double the linear scale --- so every centimetre on
the first poster represents 2 cm on the second. How does the tablet
screen area compare to the phone screen area, and what is each real area
on the posters?

\textbf{Phone screen area:} \(6 \times 12 = 72\) cm² on the poster.

\textbf{Tablet screen on the second poster:} doubling the linear scale
means each dimension doubles.

Width \(= 2 \times 6 = 12\) cm. Height \(= 2 \times 12 = 24\) cm.

Tablet screen area \(= 12 \times 24 = 288\) cm² on the second poster.

\(288 \div 72 = 4\) --- the tablet poster area is four times the phone
poster area, even though only the scale doubled. That's because area
scales as the square of the linear scale: \(2^2 = 4\).

\section{Where this goes}\label{where-this-goes-10}

The next chapter in this volume is \textbf{trigonometry} (Vol 3, Ch 2).
Trigonometry extends the geometry here by handling triangles where you
don't know the right angle --- or where you need to find angles as well
as lengths. The Pythagorean theorem handles one specific case (right
triangles, lengths only); trigonometry handles the general case using
the sine, cosine, and tangent ratios. Every oscillating system ---
sound, light, signals, tides --- is described using trigonometry.

Further out, the measurement ideas here feed directly into
\textbf{functions and relations} (Vol 3, Ch 4). Area and volume formulas
are functions: plug in the radius, get out the area. Once you're working
with functions, you can ask questions like ``how fast does the area
change as the radius grows?'' --- and that question leads to calculus.

\begin{tcolorbox}[enhanced jigsaw, opacityback=0, colback=white, breakable, colframe=quarto-callout-tip-color-frame, arc=.35mm, left=2mm, toprule=.15mm, leftrule=.75mm, bottomrule=.15mm, rightrule=.15mm]
\begin{minipage}[t]{5.5mm}
\textcolor{quarto-callout-tip-color}{\faLightbulb}
\end{minipage}%
\begin{minipage}[t]{\textwidth - 5.5mm}

\textbf{Where this shows up}

\begin{itemize}
\tightlist
\item
  A structural engineer calculating how much steel plate is needed for a
  bridge deck is computing surface area.
\item
  A physicist calculating fluid pressure on a pipe cross-section is
  computing circular area.
\item
  A GIS analyst measuring a protected forest region from a satellite
  image is computing polygon areas at scale.
\item
  A game developer checking whether two sprites have collided is
  comparing rectangular bounding boxes --- area and distance
  calculations happening thousands of times per second.
\item
  A packaging designer minimising the amount of cardboard in a cereal
  box is minimising surface area for a fixed volume.
\end{itemize}

The formulas are always the same. What changes is what you're measuring
and why it matters.

\end{minipage}%
\end{tcolorbox}

\section{Exercises}\label{exercises-10}

These are puzzles. Each one has a clean answer, but the interesting part
is deciding which formula to use and what counts as the base, height,
radius, or leg.

\begin{enumerate}
\def\labelenumi{\arabic{enumi}.}
\tightlist
\item
  A rectangular phone screen is 13.8 cm wide and 7.0 cm tall. What is
  its area? What is its diagonal length (to one decimal place)?
\end{enumerate}

\begin{Shaded}
\begin{Highlighting}[]
\NormalTok{//| echo: false}
\NormalTok{katex = require("katex@0.16")}
\end{Highlighting}
\end{Shaded}

\begin{Shaded}
\begin{Highlighting}[]
\NormalTok{//| echo: false}
\NormalTok{function makeStepperHTML(exerciseNum, steps) \{}
\NormalTok{  let current = 0;}
\NormalTok{  const totalSteps = steps.length;}
\NormalTok{  const container = document.createElement("div");}
\NormalTok{  container.style.cssText = "border:1px solid \#e5e7eb; border{-}radius:8px; padding:1rem 1.25rem; margin:0.75rem 0; font{-}family:inherit;";}
\NormalTok{  const stepsDiv = document.createElement("div");}
\NormalTok{  stepsDiv.style.marginBottom = "0.75rem";}
\NormalTok{  function renderSteps() \{}
\NormalTok{    stepsDiv.innerHTML = "";}
\NormalTok{    for (let i = 0; i \textless{} current; i++) \{}
\NormalTok{      const s = steps[i];}
\NormalTok{      const row = document.createElement("div");}
\NormalTok{      row.style.cssText = "display:grid; grid{-}template{-}columns:200px 1fr; gap:0.35rem 1rem; align{-}items:baseline; padding:0.35rem 0; border{-}top:1px solid \#f3f4f6;";}
\NormalTok{      const opCell = document.createElement("span");}
\NormalTok{      opCell.style.cssText = "font{-}size:0.85em; color:\#6b7280; font{-}style:italic;";}
\NormalTok{      opCell.textContent = s.op;}
\NormalTok{      const eqCell = document.createElement("span");}
\NormalTok{      eqCell.innerHTML = katex.renderToString(s.eq, \{ throwOnError: false, displayMode: false \});}
\NormalTok{      row.appendChild(opCell);}
\NormalTok{      row.appendChild(eqCell);}
\NormalTok{      if (s.note) \{}
\NormalTok{        const noteRow = document.createElement("div");}
\NormalTok{        noteRow.style.cssText = "grid{-}column:2; font{-}size:0.82em; color:\#6b7280; padding{-}bottom:0.2rem;";}
\NormalTok{        noteRow.textContent = s.note;}
\NormalTok{        row.appendChild(noteRow);}
\NormalTok{      \}}
\NormalTok{      stepsDiv.appendChild(row);}
\NormalTok{    \}}
\NormalTok{    if (current === totalSteps) \{}
\NormalTok{      const done = document.createElement("div");}
\NormalTok{      done.style.cssText = "margin{-}top:0.5rem; font{-}size:0.9em; color:\#059669; font{-}weight:500;";}
\NormalTok{      done.textContent = "✓ Solution complete";}
\NormalTok{      stepsDiv.appendChild(done);}
\NormalTok{    \}}
\NormalTok{  \}}
\NormalTok{  const controls = document.createElement("div");}
\NormalTok{  controls.style.cssText = "display:flex; gap:0.5rem; align{-}items:center;";}
\NormalTok{  const nextBtn = document.createElement("button");}
\NormalTok{  nextBtn.textContent = "Next step →";}
\NormalTok{  nextBtn.style.cssText = "padding:0.35rem 0.85rem; border:1px solid \#d1d5db; border{-}radius:4px; background:\#fff; cursor:pointer; font{-}size:0.9em;";}
\NormalTok{  const resetBtn = document.createElement("button");}
\NormalTok{  resetBtn.textContent = "Reset";}
\NormalTok{  resetBtn.style.cssText = "padding:0.35rem 0.75rem; border:1px solid \#e5e7eb; border{-}radius:4px; background:\#f9fafb; cursor:pointer; font{-}size:0.9em; color:\#6b7280;";}
\NormalTok{  const counter = document.createElement("span");}
\NormalTok{  counter.style.cssText = "font{-}size:0.82em; color:\#9ca3af; margin{-}left:0.25rem;";}
\NormalTok{  function updateButtons() \{}
\NormalTok{    nextBtn.disabled = current === totalSteps;}
\NormalTok{    nextBtn.style.opacity = current === totalSteps ? "0.4" : "1";}
\NormalTok{    counter.textContent = current === 0 ? "Click to reveal steps" : \textasciigrave{}Step $\{current\} of $\{totalSteps\}\textasciigrave{};}
\NormalTok{  \}}
\NormalTok{  nextBtn.onclick = () =\textgreater{} \{ if (current \textless{} totalSteps) \{ current++; renderSteps(); updateButtons(); \} \};}
\NormalTok{  resetBtn.onclick = () =\textgreater{} \{ current = 0; renderSteps(); updateButtons(); \};}
\NormalTok{  controls.appendChild(nextBtn);}
\NormalTok{  controls.appendChild(resetBtn);}
\NormalTok{  controls.appendChild(counter);}
\NormalTok{  container.appendChild(stepsDiv);}
\NormalTok{  container.appendChild(controls);}
\NormalTok{  renderSteps();}
\NormalTok{  updateButtons();}
\NormalTok{  return container;}
\NormalTok{\}}
\end{Highlighting}
\end{Shaded}

\begin{Shaded}
\begin{Highlighting}[]
\NormalTok{//| echo: false}
\NormalTok{makeStepperHTML(1, [}
\NormalTok{  \{ op: "Find the area", eq: "A = l \textbackslash{}\textbackslash{}times w = 13.8 \textbackslash{}\textbackslash{}times 7.0 = 96.6 \textbackslash{}\textbackslash{}text\{ cm\}\^{}2", note: null \},}
\NormalTok{  \{ op: "Write the Pythagorean theorem for the diagonal", eq: "d = \textbackslash{}\textbackslash{}sqrt\{l\^{}2 + w\^{}2\}", note: "The diagonal of a rectangle is the hypotenuse of a right triangle whose legs are the width and height." \},}
\NormalTok{  \{ op: "Substitute values", eq: "d = \textbackslash{}\textbackslash{}sqrt\{13.8\^{}2 + 7.0\^{}2\} = \textbackslash{}\textbackslash{}sqrt\{190.44 + 49.00\} = \textbackslash{}\textbackslash{}sqrt\{239.44\}", note: null \},}
\NormalTok{  \{ op: "Calculate", eq: "d = 15.5 \textbackslash{}\textbackslash{}text\{ cm\}", note: null \},}
\NormalTok{  \{ op: "Check units", eq: "\textbackslash{}\textbackslash{}text\{Area: \} 96.6 \textbackslash{}\textbackslash{}text\{ cm\}\^{}2 \textbackslash{}\textbackslash{}quad \textbackslash{}\textbackslash{}text\{Diagonal: \} 15.5 \textbackslash{}\textbackslash{}text\{ cm\}", note: null \}}
\NormalTok{])}
\end{Highlighting}
\end{Shaded}

\begin{enumerate}
\def\labelenumi{\arabic{enumi}.}
\setcounter{enumi}{1}
\tightlist
\item
  A circular pizza has a diameter of 30 cm. What is its area (leave your
  answer in terms of \(\pi\), then give a decimal approximation to the
  nearest whole number)?
\end{enumerate}

\begin{Shaded}
\begin{Highlighting}[]
\NormalTok{//| echo: false}
\NormalTok{makeStepperHTML(2, [}
\NormalTok{  \{ op: "Find the radius from the diameter", eq: "r = \textbackslash{}\textbackslash{}frac\{d\}\{2\} = \textbackslash{}\textbackslash{}frac\{30\}\{2\} = 15 \textbackslash{}\textbackslash{}text\{ cm\}", note: null \},}
\NormalTok{  \{ op: "Write the area formula", eq: "A = \textbackslash{}\textbackslash{}pi r\^{}2", note: null \},}
\NormalTok{  \{ op: "Substitute", eq: "A = \textbackslash{}\textbackslash{}pi \textbackslash{}\textbackslash{}times 15\^{}2 = 225\textbackslash{}\textbackslash{}pi \textbackslash{}\textbackslash{}text\{ cm\}\^{}2", note: null \},}
\NormalTok{  \{ op: "Decimal approximation", eq: "A \textbackslash{}\textbackslash{}approx 225 \textbackslash{}\textbackslash{}times 3.14159 \textbackslash{}\textbackslash{}approx 707 \textbackslash{}\textbackslash{}text\{ cm\}\^{}2", note: null \},}
\NormalTok{  \{ op: "Check units", eq: "A = 225\textbackslash{}\textbackslash{}pi \textbackslash{}\textbackslash{}approx 707 \textbackslash{}\textbackslash{}text\{ cm\}\^{}2", note: null \}}
\NormalTok{])}
\end{Highlighting}
\end{Shaded}

\begin{enumerate}
\def\labelenumi{\arabic{enumi}.}
\setcounter{enumi}{2}
\tightlist
\item
  A ramp rises 0.9 m vertically over a horizontal distance of 1.2 m. How
  long is the sloped surface? (Give your answer to one decimal place.)
\end{enumerate}

\begin{Shaded}
\begin{Highlighting}[]
\NormalTok{//| echo: false}
\NormalTok{makeStepperHTML(3, [}
\NormalTok{  \{ op: "Identify the right triangle: legs are the vertical rise and horizontal run, hypotenuse is the slope", eq: "c = \textbackslash{}\textbackslash{}sqrt\{a\^{}2 + b\^{}2\}", note: null \},}
\NormalTok{  \{ op: "Substitute values", eq: "c = \textbackslash{}\textbackslash{}sqrt\{0.9\^{}2 + 1.2\^{}2\} = \textbackslash{}\textbackslash{}sqrt\{0.81 + 1.44\} = \textbackslash{}\textbackslash{}sqrt\{2.25\}", note: null \},}
\NormalTok{  \{ op: "Calculate", eq: "c = 1.5 \textbackslash{}\textbackslash{}text\{ m\}", note: "0.9, 1.2, 1.5 is a scaled 3{-}4{-}5 triple (each side multiplied by 0.3)." \},}
\NormalTok{  \{ op: "Check units", eq: "\textbackslash{}\textbackslash{}text\{Slope length\} = 1.5 \textbackslash{}\textbackslash{}text\{ m\}", note: null \}}
\NormalTok{])}
\end{Highlighting}
\end{Shaded}

\begin{enumerate}
\def\labelenumi{\arabic{enumi}.}
\setcounter{enumi}{3}
\tightlist
\item
  A cylindrical water bottle has a radius of 4 cm and a height of 20 cm.
  What is its volume in cm³? What is its capacity in millilitres? (1 cm³
  = 1 ml.)
\end{enumerate}

\begin{Shaded}
\begin{Highlighting}[]
\NormalTok{//| echo: false}
\NormalTok{makeStepperHTML(4, [}
\NormalTok{  \{ op: "Write the volume formula for a cylinder", eq: "V = \textbackslash{}\textbackslash{}pi r\^{}2 h", note: null \},}
\NormalTok{  \{ op: "Substitute values", eq: "V = \textbackslash{}\textbackslash{}pi \textbackslash{}\textbackslash{}times 4\^{}2 \textbackslash{}\textbackslash{}times 20 = \textbackslash{}\textbackslash{}pi \textbackslash{}\textbackslash{}times 16 \textbackslash{}\textbackslash{}times 20 = 320\textbackslash{}\textbackslash{}pi \textbackslash{}\textbackslash{}text\{ cm\}\^{}3", note: null \},}
\NormalTok{  \{ op: "Decimal approximation", eq: "V \textbackslash{}\textbackslash{}approx 320 \textbackslash{}\textbackslash{}times 3.14159 \textbackslash{}\textbackslash{}approx 1005 \textbackslash{}\textbackslash{}text\{ cm\}\^{}3", note: null \},}
\NormalTok{  \{ op: "Convert to millilitres", eq: "1005 \textbackslash{}\textbackslash{}text\{ cm\}\^{}3 = 1005 \textbackslash{}\textbackslash{}text\{ ml\} \textbackslash{}\textbackslash{}approx 1 \textbackslash{}\textbackslash{}text\{ litre\}", note: "1 cm³ = 1 ml, so the numeric value stays the same." \},}
\NormalTok{  \{ op: "Check units", eq: "V = 320\textbackslash{}\textbackslash{}pi \textbackslash{}\textbackslash{}approx 1005 \textbackslash{}\textbackslash{}text\{ ml\}", note: null \}}
\NormalTok{])}
\end{Highlighting}
\end{Shaded}

\begin{enumerate}
\def\labelenumi{\arabic{enumi}.}
\setcounter{enumi}{4}
\tightlist
\item
  A triangular garden bed has a base of 5 m and a perpendicular height
  of 3.6 m. What area of soil does it cover? If soil costs \$12 per m²,
  what is the total cost?
\end{enumerate}

\begin{Shaded}
\begin{Highlighting}[]
\NormalTok{//| echo: false}
\NormalTok{makeStepperHTML(5, [}
\NormalTok{  \{ op: "Write the area formula for a triangle", eq: "A = \textbackslash{}\textbackslash{}frac\{1\}\{2\} \textbackslash{}\textbackslash{}times b \textbackslash{}\textbackslash{}times h", note: null \},}
\NormalTok{  \{ op: "Substitute values", eq: "A = \textbackslash{}\textbackslash{}frac\{1\}\{2\} \textbackslash{}\textbackslash{}times 5 \textbackslash{}\textbackslash{}times 3.6 = \textbackslash{}\textbackslash{}frac\{18\}\{2\} = 9 \textbackslash{}\textbackslash{}text\{ m\}\^{}2", note: null \},}
\NormalTok{  \{ op: "Calculate the cost", eq: "\textbackslash{}\textbackslash{}text\{Cost\} = 9 \textbackslash{}\textbackslash{}times 12 = \textbackslash{}\textbackslash{}$108", note: null \},}
\NormalTok{  \{ op: "Check units", eq: "\textbackslash{}\textbackslash{}text\{Area\} = 9 \textbackslash{}\textbackslash{}text\{ m\}\^{}2, \textbackslash{}\textbackslash{}quad \textbackslash{}\textbackslash{}text\{Cost\} = \textbackslash{}\textbackslash{}$108", note: null \}}
\NormalTok{])}
\end{Highlighting}
\end{Shaded}

\begin{enumerate}
\def\labelenumi{\arabic{enumi}.}
\setcounter{enumi}{5}
\tightlist
\item
  A rectangular swimming pool is 25 m long and 10 m wide. It has a
  uniform depth of 1.8 m. How many cubic metres of water does it hold?
  Give your answer in litres. (1 m³ = 1000 litres.)
\end{enumerate}

\begin{Shaded}
\begin{Highlighting}[]
\NormalTok{//| echo: false}
\NormalTok{makeStepperHTML(6, [}
\NormalTok{  \{ op: "Write the volume formula for a rectangular prism", eq: "V = l \textbackslash{}\textbackslash{}times w \textbackslash{}\textbackslash{}times h", note: null \},}
\NormalTok{  \{ op: "Substitute values", eq: "V = 25 \textbackslash{}\textbackslash{}times 10 \textbackslash{}\textbackslash{}times 1.8 = 450 \textbackslash{}\textbackslash{}text\{ m\}\^{}3", note: null \},}
\NormalTok{  \{ op: "Convert to litres", eq: "450 \textbackslash{}\textbackslash{}text\{ m\}\^{}3 \textbackslash{}\textbackslash{}times 1000 \textbackslash{}\textbackslash{}text\{ litres/m\}\^{}3 = 450\{,\}000 \textbackslash{}\textbackslash{}text\{ litres\}", note: "1 m³ = 1000 litres, so multiply by 1000." \},}
\NormalTok{  \{ op: "Check units", eq: "V = 450 \textbackslash{}\textbackslash{}text\{ m\}\^{}3 = 450\{,\}000 \textbackslash{}\textbackslash{}text\{ litres\}", note: null \}}
\NormalTok{])}
\end{Highlighting}
\end{Shaded}

\chapter{Trigonometry}\label{trigonometry}

Angles and the ratios they fix

\hfill\break

You're building a wheelchair ramp. The code says the angle can't exceed
4.8°. You know how high the step is. How long does the ramp need to be?

You're standing 30 metres from a building, looking up at the roof. You
want to know how tall the building is, but you can't climb it. You can
measure the angle you're looking up at. Is that enough?

Your phone screen is 15.6 cm tall and 7.2 cm wide. What's the diagonal?
Easy --- Pythagoras. But what's the angle in the corner? That's where
trig comes in.

Each of these situations has a right triangle in it somewhere. You know
some sides and angles, and you want to find others. Trigonometry gives
you three ratios that connect them.

\section{What the notation is
saying}\label{what-the-notation-is-saying-11}

Start with a right triangle. One angle is exactly 90°. Pick one of the
other angles and call it \(\theta\) (the Greek letter \emph{theta} ---
just a name for ``the angle you're working with'').

Now name the three sides relative to \(\theta\):

\begin{itemize}
\tightlist
\item
  The \textbf{hypotenuse} is the longest side, always opposite the right
  angle.
\item
  The \textbf{opposite} side is the one directly across the triangle
  from \(\theta\) --- the side \(\theta\) ``faces''.
\item
  The \textbf{adjacent} side is the one next to \(\theta\) that isn't
  the hypotenuse --- the side \(\theta\) ``leans against''.
\end{itemize}

The three trig ratios are:

\[\sin(\theta) = \frac{\text{opposite}}{\text{hypotenuse}}\]

Read it: ``the sine of theta equals opposite over hypotenuse.''

\[\cos(\theta) = \frac{\text{adjacent}}{\text{hypotenuse}}\]

Read it: ``the cosine of theta equals adjacent over hypotenuse.''

\[\tan(\theta) = \frac{\text{opposite}}{\text{adjacent}}\]

Read it: ``the tangent of theta equals opposite over adjacent.''

The memory aid is \textbf{SOH-CAH-TOA}: Sine = Opposite/Hypotenuse,
Cosine = Adjacent/Hypotenuse, Tangent = Opposite/Adjacent. It's a
mnemonic, not mathematics --- use it until you don't need it any more.

\begin{Shaded}
\begin{Highlighting}[]
\NormalTok{//| echo: false}
\NormalTok{viewof triAngle = Inputs.range([5, 85], \{}
\NormalTok{  value: 35,}
\NormalTok{  step: 1,}
\NormalTok{  label: "Angle θ (degrees)"}
\NormalTok{\})}
\end{Highlighting}
\end{Shaded}

\begin{Shaded}
\begin{Highlighting}[]
\NormalTok{//| echo: false}
\NormalTok{viewof triHyp = Inputs.range([1, 10], \{}
\NormalTok{  value: 6,}
\NormalTok{  step: 0.5,}
\NormalTok{  label: "Hypotenuse length (units)"}
\NormalTok{\})}
\end{Highlighting}
\end{Shaded}

\begin{Shaded}
\begin{Highlighting}[]
\NormalTok{//| echo: false}
\NormalTok{\{}
\NormalTok{  const deg = triAngle;}
\NormalTok{  const hyp = triHyp;}
\NormalTok{  const rad = deg * Math.PI / 180;}
\NormalTok{  const opp = hyp * Math.sin(rad);}
\NormalTok{  const adj = hyp * Math.cos(rad);}

\NormalTok{  const scale = 30;}
\NormalTok{  const svgW = 520, svgH = 300;}
\NormalTok{  const margin = 50;}

\NormalTok{  // Triangle: right angle at bottom{-}left}
\NormalTok{  const ox = margin, oy = svgH {-} margin {-} 20;}
\NormalTok{  const ax = ox + adj * scale, ay = oy;      // bottom{-}right}
\NormalTok{  const bx = ox, by = oy {-} opp * scale;      // top{-}left}

\NormalTok{  const svg = d3.create("svg")}
\NormalTok{    .attr("width", svgW)}
\NormalTok{    .attr("height", svgH)}
\NormalTok{    .style("font{-}family", "inherit");}

\NormalTok{  // Triangle fill}
\NormalTok{  svg.append("polygon")}
\NormalTok{    .attr("points", \textasciigrave{}$\{ox\},$\{oy\} $\{ax\},$\{ay\} $\{bx\},$\{by\}\textasciigrave{})}
\NormalTok{    .attr("fill", "\#f0f9ff")}
\NormalTok{    .attr("stroke", "\#0369a1")}
\NormalTok{    .attr("stroke{-}width", 2);}

\NormalTok{  // Right{-}angle marker}
\NormalTok{  const sq = 10;}
\NormalTok{  svg.append("rect")}
\NormalTok{    .attr("x", ox).attr("y", oy {-} sq)}
\NormalTok{    .attr("width", sq).attr("height", sq)}
\NormalTok{    .attr("fill", "none").attr("stroke", "\#374151").attr("stroke{-}width", 1.5);}

\NormalTok{  // Angle arc at bottom{-}right}
\NormalTok{  const arcR = 22;}
\NormalTok{  const arcPath = d3.arc()(\{}
\NormalTok{    innerRadius: 0, outerRadius: arcR,}
\NormalTok{    startAngle: Math.PI,}
\NormalTok{    endAngle: Math.PI + rad}
\NormalTok{  \});}
\NormalTok{  svg.append("path")}
\NormalTok{    .attr("d", arcPath)}
\NormalTok{    .attr("transform", \textasciigrave{}translate($\{ax\},$\{ay\}) rotate({-}90)\textasciigrave{})}
\NormalTok{    .attr("fill", "\#fde68a")}
\NormalTok{    .attr("stroke", "\#d97706")}
\NormalTok{    .attr("stroke{-}width", 1);}

\NormalTok{  svg.append("text")}
\NormalTok{    .attr("x", ax {-} arcR {-} 14).attr("y", ay {-} 8)}
\NormalTok{    .attr("font{-}size", "13px").attr("fill", "\#92400e")}
\NormalTok{    .text(\textasciigrave{}θ=$\{deg\}°\textasciigrave{});}

\NormalTok{  // Side labels}
\NormalTok{  // Hypotenuse midpoint}
\NormalTok{  const hmx = (ax + bx) / 2 + 14, hmy = (ay + by) / 2;}
\NormalTok{  svg.append("text")}
\NormalTok{    .attr("x", hmx).attr("y", hmy)}
\NormalTok{    .attr("font{-}size", "12px").attr("fill", "\#0369a1").attr("font{-}weight", "600")}
\NormalTok{    .text(\textasciigrave{}hyp = $\{hyp.toFixed(1)\}\textasciigrave{});}

\NormalTok{  // Adjacent (bottom)}
\NormalTok{  svg.append("text")}
\NormalTok{    .attr("x", ox + adj * scale / 2).attr("y", oy + 18)}
\NormalTok{    .attr("text{-}anchor", "middle").attr("font{-}size", "12px").attr("fill", "\#059669")}
\NormalTok{    .text(\textasciigrave{}adj = $\{adj.toFixed(2)\}\textasciigrave{});}

\NormalTok{  // Opposite (left side)}
\NormalTok{  svg.append("text")}
\NormalTok{    .attr("x", ox {-} 14).attr("y", oy {-} opp * scale / 2)}
\NormalTok{    .attr("text{-}anchor", "middle").attr("font{-}size", "12px").attr("fill", "\#dc2626")}
\NormalTok{    .text(\textasciigrave{}opp = $\{opp.toFixed(2)\}\textasciigrave{});}

\NormalTok{  // Ratios panel}
\NormalTok{  const px = 310, py = 40;}
\NormalTok{  svg.append("rect")}
\NormalTok{    .attr("x", px {-} 10).attr("y", py {-} 20)}
\NormalTok{    .attr("width", 190).attr("height", 140)}
\NormalTok{    .attr("fill", "\#f8fafc").attr("stroke", "\#e2e8f0").attr("rx", 6);}

\NormalTok{  svg.append("text").attr("x", px).attr("y", py)}
\NormalTok{    .attr("font{-}size", "13px").attr("font{-}weight", "600").attr("fill", "\#374151")}
\NormalTok{    .text("Ratios (fixed by θ alone)");}

\NormalTok{  svg.append("text").attr("x", px).attr("y", py + 26)}
\NormalTok{    .attr("font{-}size", "13px").attr("fill", "\#374151")}
\NormalTok{    .text(\textasciigrave{}sin($\{deg\}°) = opp/hyp\textasciigrave{});}
\NormalTok{  svg.append("text").attr("x", px).attr("y", py + 44)}
\NormalTok{    .attr("font{-}size", "13px").attr("fill", "\#dc2626")}
\NormalTok{    .text(\textasciigrave{}= $\{Math.sin(rad).toFixed(4)\}\textasciigrave{});}

\NormalTok{  svg.append("text").attr("x", px).attr("y", py + 66)}
\NormalTok{    .attr("font{-}size", "13px").attr("fill", "\#374151")}
\NormalTok{    .text(\textasciigrave{}cos($\{deg\}°) = adj/hyp\textasciigrave{});}
\NormalTok{  svg.append("text").attr("x", px).attr("y", py + 84)}
\NormalTok{    .attr("font{-}size", "13px").attr("fill", "\#059669")}
\NormalTok{    .text(\textasciigrave{}= $\{Math.cos(rad).toFixed(4)\}\textasciigrave{});}

\NormalTok{  svg.append("text").attr("x", px).attr("y", py + 106)}
\NormalTok{    .attr("font{-}size", "13px").attr("fill", "\#374151")}
\NormalTok{    .text(\textasciigrave{}tan($\{deg\}°) = opp/adj\textasciigrave{});}
\NormalTok{  svg.append("text").attr("x", px).attr("y", py + 124)}
\NormalTok{    .attr("font{-}size", "13px").attr("fill", "\#7c3aed")}
\NormalTok{    .text(\textasciigrave{}= $\{Math.tan(rad).toFixed(4)\}\textasciigrave{});}

\NormalTok{  return svg.node();}
\NormalTok{\}}
\end{Highlighting}
\end{Shaded}

\section{The method}\label{the-method-10}

\subsection{Finding an unknown side}\label{finding-an-unknown-side}

You know an angle and one side. You want another side. Pick the ratio
that connects what you know to what you want.

\textbf{Example.} A ramp rises 0.5 m. The angle at the base is 4°. How
long is the ramp surface?

The 0.5 m is the side opposite the 4° angle. The ramp surface is the
hypotenuse. The ratio connecting opposite and hypotenuse is sine:

\[\sin(4°) = \frac{0.5}{\text{hyp}}\]

Rearrange step by step:

\[\sin(4°) \times \text{hyp} = 0.5 \quad \text{[multiply both sides by hyp]}\]

\[\text{hyp} = \frac{0.5}{\sin(4°)} \quad \text{[divide both sides by } \sin(4°)\text{]}\]

\[\text{hyp} \approx \frac{0.5}{0.0698} \approx 7.2 \text{ m}\]

The ramp needs to be about 7.2 m long.

\subsection{Finding an unknown angle}\label{finding-an-unknown-angle}

You know two sides. You want the angle between them. Use the inverse
trig function.

\(\sin^{-1}\) reads as ``inverse sine'' or ``arcsine''. It asks: what
angle has this sine value? If \(\sin(\theta) = 0.5\), then
\(\theta = \sin^{-1}(0.5) = 30°\). The \(^{-1}\) notation means ``undo
the sine'' --- it does not mean ``one over sine.''

Similarly, \(\cos^{-1}\) and \(\tan^{-1}\) undo the cosine and tangent
respectively. All three inverse functions work the same way: give them a
ratio, get back an angle.

\textbf{Example.} Your phone screen is 15.6 cm tall and 7.2 cm wide.
What angle does the diagonal make with the long side?

The two known sides are 15.6 cm (adjacent to the angle at the
bottom-left corner) and 7.2 cm (opposite). The ratio connecting opposite
and adjacent is tangent:

\[\tan(\theta) = \frac{7.2}{15.6} \approx 0.4615\]

Apply inverse tangent to both sides:

\[\theta = \tan^{-1}(0.4615) \approx 24.8°\]

The diagonal meets the long side at about 25°.

\begin{tcolorbox}[enhanced jigsaw, opacityback=0, colback=white, breakable, colframe=quarto-callout-note-color-frame, arc=.35mm, left=2mm, toprule=.15mm, leftrule=.75mm, bottomrule=.15mm, rightrule=.15mm]
\begin{minipage}[t]{5.5mm}
\textcolor{quarto-callout-note-color}{\faInfo}
\end{minipage}%
\begin{minipage}[t]{\textwidth - 5.5mm}

\textbf{Why this works}

The ratios are fixed by the angle alone --- not by how big the triangle
is. Any right triangle with a 30° angle has the same ratio of opposite
to hypotenuse: exactly 0.5. Scale the triangle up by ten times and both
sides get ten times longer, but the ratio stays the same.

This is why the ratios are useful. They depend only on the angle, so you
can look up \(\sin(30°) = 0.5\) once and use it for any triangle
containing a 30° angle, whatever its size.

\end{minipage}%
\end{tcolorbox}

\section{Worked examples}\label{worked-examples-10}

\textbf{Example 1 --- Finding the height of a building
(geography/surveying)}

A surveyor stands 40 m from the base of a radio mast. She measures the
angle of elevation to the top as 62°. How tall is the mast?

The 40 m is the side adjacent to the 62° angle. The mast height is the
opposite side. The ratio connecting adjacent and opposite is tangent:

\[\tan(62°) = \frac{\text{height}}{40}\]

Multiply both sides by 40:

\[\text{height} = 40 \times \tan(62°)\]

\(\tan(62°) \approx 1.8807\), so:

\[\text{height} = 40 \times 1.8807 \approx 75.2 \text{ m}\]

The mast is approximately 75 m tall. Check: \(\tan^{-1}(75.2/40)
= \tan^{-1}(1.880) \approx 62°\). Yes.

\begin{center}\rule{0.5\linewidth}{0.5pt}\end{center}

\textbf{Example 2 --- Rope from a fence post (geometry/practical)}

A rope is tied from the top of a fence post to a peg in the ground. The
post is vertical and the rope makes an angle of 35° with the post. The
rope has a tension of 80 N --- the force acting along the rope. What is
the component of that tension pulling horizontally away from the post?

The 80 N tension acts along the rope --- that's the hypotenuse of a
right triangle. The angle at the top of the post is 35°. The horizontal
pull is the side opposite that angle:

\[\sin(35°) = \frac{\text{horizontal component}}{80}\]

\[\text{horizontal component} = 80 \times \sin(35°)\]

\(\sin(35°) \approx 0.5736\), so:

\[\text{horizontal component} = 80 \times 0.5736 \approx 45.9 \text{ N}\]

The peg must hold against about 46 N of horizontal pull.

\begin{center}\rule{0.5\linewidth}{0.5pt}\end{center}

\textbf{Example 3 --- Finding the ramp angle for a target slope
(science/engineering)}

A physiotherapy exercise ramp needs a slope of 1 in 8 --- for every 8 cm
of horizontal run, it rises 1 cm. What angle does this correspond to?

The slope ratio gives you opposite over adjacent:
\(\frac{1}{8} = 0.125\). That's the tangent of the angle:

\[\tan(\theta) = \frac{1}{8} = 0.125\]

\[\theta = \tan^{-1}(0.125) \approx 7.1°\]

The ramp makes an angle of about 7° with the floor.

\begin{center}\rule{0.5\linewidth}{0.5pt}\end{center}

\textbf{Example 4 --- Rotating a point in 2D (computing/game
development)}

A game object is at position \((5, 0)\) --- five units to the right of
the origin, on the horizontal axis. The game rotates everything 40°
anticlockwise. Where does the object end up?

The object starts at distance 5 from the origin (its ``hypotenuse'' from
the rotation centre). After a 40° rotation, the new \(x\)-coordinate is
the adjacent side of the 40° angle, and the new \(y\)-coordinate is the
opposite side:

\[x' = 5 \times \cos(40°) \approx 5 \times 0.7660 \approx 3.83\]

\[y' = 5 \times \sin(40°) \approx 5 \times 0.6428 \approx 3.21\]

The object moves to approximately \((3.83,\ 3.21)\). This is the core of
every 2D rotation in a game engine.

\begin{Shaded}
\begin{Highlighting}[]
\NormalTok{//| echo: false}
\NormalTok{viewof elevDist = Inputs.range([5, 30], \{}
\NormalTok{  value: 15,}
\NormalTok{  step: 1,}
\NormalTok{  label: "Horizontal distance from building (m)"}
\NormalTok{\})}
\end{Highlighting}
\end{Shaded}

\begin{Shaded}
\begin{Highlighting}[]
\NormalTok{//| echo: false}
\NormalTok{\{}
\NormalTok{  const dist = elevDist;}
\NormalTok{  const buildingH = 20; // fixed metres}
\NormalTok{  const angleRad = Math.atan(buildingH / dist);}
\NormalTok{  const angleDeg = angleRad * 180 / Math.PI;}

\NormalTok{  const svgW = 500, svgH = 300;}
\NormalTok{  const margin = 40;}
\NormalTok{  const scale = 7; // px per metre}

\NormalTok{  // Scene coordinates}
\NormalTok{  const groundY = svgH {-} margin {-} 20;}
\NormalTok{  const buildingX = svgW {-} margin {-} 60;}
\NormalTok{  const observerX = buildingX {-} dist * scale;}
\NormalTok{  const buildingTopY = groundY {-} buildingH * scale;}

\NormalTok{  const svg = d3.create("svg")}
\NormalTok{    .attr("width", svgW)}
\NormalTok{    .attr("height", svgH)}
\NormalTok{    .style("font{-}family", "inherit");}

\NormalTok{  // Ground line}
\NormalTok{  svg.append("line")}
\NormalTok{    .attr("x1", margin).attr("y1", groundY)}
\NormalTok{    .attr("x2", svgW {-} margin).attr("y2", groundY)}
\NormalTok{    .attr("stroke", "\#9ca3af").attr("stroke{-}width", 1.5);}

\NormalTok{  // Building (vertical line)}
\NormalTok{  svg.append("line")}
\NormalTok{    .attr("x1", buildingX).attr("y1", groundY)}
\NormalTok{    .attr("x2", buildingX).attr("y2", buildingTopY)}
\NormalTok{    .attr("stroke", "\#374151").attr("stroke{-}width", 4);}

\NormalTok{  // Building label}
\NormalTok{  svg.append("text")}
\NormalTok{    .attr("x", buildingX + 8).attr("y", (groundY + buildingTopY) / 2)}
\NormalTok{    .attr("font{-}size", "12px").attr("fill", "\#374151")}
\NormalTok{    .text("20 m");}

\NormalTok{  // Observer dot}
\NormalTok{  svg.append("circle")}
\NormalTok{    .attr("cx", observerX).attr("cy", groundY)}
\NormalTok{    .attr("r", 6)}
\NormalTok{    .attr("fill", "\#f59e0b").attr("stroke", "\#d97706").attr("stroke{-}width", 1.5);}

\NormalTok{  svg.append("text")}
\NormalTok{    .attr("x", observerX).attr("y", groundY + 18)}
\NormalTok{    .attr("text{-}anchor", "middle").attr("font{-}size", "11px").attr("fill", "\#374151")}
\NormalTok{    .text("observer");}

\NormalTok{  // Hypotenuse line}
\NormalTok{  svg.append("line")}
\NormalTok{    .attr("x1", observerX).attr("y1", groundY)}
\NormalTok{    .attr("x2", buildingX).attr("y2", buildingTopY)}
\NormalTok{    .attr("stroke", "\#2563eb").attr("stroke{-}width", 2)}
\NormalTok{    .attr("stroke{-}dasharray", "5,3");}

\NormalTok{  // Horizontal distance label}
\NormalTok{  svg.append("line")}
\NormalTok{    .attr("x1", observerX).attr("y1", groundY + 28)}
\NormalTok{    .attr("x2", buildingX).attr("y2", groundY + 28)}
\NormalTok{    .attr("stroke", "\#6b7280").attr("stroke{-}width", 1)}
\NormalTok{    .attr("marker{-}end", "url(\#arrow)");}

\NormalTok{  svg.append("text")}
\NormalTok{    .attr("x", (observerX + buildingX) / 2).attr("y", groundY + 42)}
\NormalTok{    .attr("text{-}anchor", "middle").attr("font{-}size", "12px").attr("fill", "\#6b7280")}
\NormalTok{    .text(\textasciigrave{}d = $\{dist\} m\textasciigrave{});}

\NormalTok{  // Angle arc at observer}
\NormalTok{  const arcR = 28;}
\NormalTok{  const arcPath = d3.arc()(\{}
\NormalTok{    innerRadius: 0, outerRadius: arcR,}
\NormalTok{    startAngle: {-}Math.PI / 2,}
\NormalTok{    endAngle: {-}Math.PI / 2 + angleRad}
\NormalTok{  \});}
\NormalTok{  svg.append("path")}
\NormalTok{    .attr("d", arcPath)}
\NormalTok{    .attr("transform", \textasciigrave{}translate($\{observerX\},$\{groundY\}) rotate($\{90 {-} angleDeg\})\textasciigrave{})}
\NormalTok{    .attr("fill", "\#fde68a").attr("stroke", "\#d97706").attr("stroke{-}width", 1);}

\NormalTok{  // Opposite label (building height)}
\NormalTok{  svg.append("text")}
\NormalTok{    .attr("x", buildingX {-} 10).attr("y", groundY {-} buildingH * scale / 2)}
\NormalTok{    .attr("text{-}anchor", "end").attr("font{-}size", "12px").attr("fill", "\#dc2626")}
\NormalTok{    .text("opp = 20 m");}

\NormalTok{  // Adjacent label}
\NormalTok{  svg.append("text")}
\NormalTok{    .attr("x", (observerX + buildingX) / 2).attr("y", groundY {-} 8)}
\NormalTok{    .attr("text{-}anchor", "middle").attr("font{-}size", "12px").attr("fill", "\#059669")}
\NormalTok{    .text(\textasciigrave{}adj = $\{dist\} m\textasciigrave{});}

\NormalTok{  // Results panel}
\NormalTok{  const px = margin, py = margin;}
\NormalTok{  svg.append("rect")}
\NormalTok{    .attr("x", px {-} 4).attr("y", py {-} 18)}
\NormalTok{    .attr("width", 210).attr("height", 80)}
\NormalTok{    .attr("fill", "\#f8fafc").attr("stroke", "\#e2e8f0").attr("rx", 6);}

\NormalTok{  svg.append("text").attr("x", px).attr("y", py)}
\NormalTok{    .attr("font{-}size", "13px").attr("font{-}weight", "600").attr("fill", "\#374151")}
\NormalTok{    .text("Angle of elevation");}

\NormalTok{  svg.append("text").attr("x", px).attr("y", py + 22)}
\NormalTok{    .attr("font{-}size", "13px").attr("fill", "\#374151")}
\NormalTok{    .text(\textasciigrave{}θ = tan⁻¹(20 / $\{dist\})\textasciigrave{});}

\NormalTok{  svg.append("text").attr("x", px).attr("y", py + 42)}
\NormalTok{    .attr("font{-}size", "14px").attr("fill", "\#2563eb").attr("font{-}weight", "600")}
\NormalTok{    .text(\textasciigrave{}θ = $\{angleDeg.toFixed(1)\}°\textasciigrave{});}

\NormalTok{  svg.append("text").attr("x", px).attr("y", py + 60)}
\NormalTok{    .attr("font{-}size", "12px").attr("fill", "\#6b7280")}
\NormalTok{    .text(\textasciigrave{}(tan⁻¹($\{(buildingH/dist).toFixed(3)\}))\textasciigrave{});}

\NormalTok{  return svg.node();}
\NormalTok{\}}
\end{Highlighting}
\end{Shaded}

\section{Where this goes}\label{where-this-goes-11}

The three ratios in this chapter are functions --- you put in an angle
and get out a number. \textbf{Volume 3, Chapter 4} (Functions and
relations) gives that idea its full treatment: a function is a rule that
maps one value to another, and sine, cosine, and tangent are the first
examples most people meet of functions that are not straight lines. Once
you see trig through the function lens you can ask new questions: what
does the graph look like? Where is it increasing? Does it repeat?

It repeats --- and that turns out to matter enormously.
\textbf{Differential calculus} (Volume 5) uses sine and cosine as its
central worked examples because they are the simplest functions whose
rate of change is also a trig function. The derivative of \(\sin\) is
\(\cos\); the derivative of \(\cos\) is \(-\sin\). That self-referential
property is why trig is indispensable in physics and engineering, and
it's the first real glimpse of how calculus and trigonometry are
connected at a deep level.

\begin{tcolorbox}[enhanced jigsaw, opacityback=0, colback=white, breakable, colframe=quarto-callout-tip-color-frame, arc=.35mm, left=2mm, toprule=.15mm, leftrule=.75mm, bottomrule=.15mm, rightrule=.15mm]
\begin{minipage}[t]{5.5mm}
\textcolor{quarto-callout-tip-color}{\faLightbulb}
\end{minipage}%
\begin{minipage}[t]{\textwidth - 5.5mm}

\textbf{Where this shows up}

\begin{itemize}
\tightlist
\item
  A structural engineer resolves a cable force into horizontal and
  vertical components before checking whether a joint will hold.
\item
  A physicist writes the equation of a pendulum's motion using sine ---
  the angle oscillates like a sine wave over time.
\item
  A navigator uses bearing angles and the sine rule to triangulate a
  ship's position from two known landmarks.
\item
  A game developer rotates sprites, aims projectiles, and swings camera
  views using \(\sin\) and \(\cos\) applied to coordinates.
\item
  A sound engineer represents audio waveforms as sums of sine waves ---
  the starting point of Fourier analysis.
\end{itemize}

\end{minipage}%
\end{tcolorbox}

\section{Exercises}\label{exercises-11}

These are puzzles. Each one has a clean answer. The work is in
identifying which sides and which ratio to use before you reach for the
calculator.

\begin{center}\rule{0.5\linewidth}{0.5pt}\end{center}

\textbf{1.} A ski slope is 800 m long and the vertical drop is 240 m.
What angle does the slope make with the horizontal? Give your answer to
one decimal place.

\begin{Shaded}
\begin{Highlighting}[]
\NormalTok{//| echo: false}
\NormalTok{katex = require("katex@0.16")}
\end{Highlighting}
\end{Shaded}

\begin{Shaded}
\begin{Highlighting}[]
\NormalTok{//| echo: false}
\NormalTok{function makeStepperHTML(exerciseNum, steps) \{}
\NormalTok{  let current = 0;}
\NormalTok{  const totalSteps = steps.length;}
\NormalTok{  const container = document.createElement("div");}
\NormalTok{  container.style.cssText = "border:1px solid \#e5e7eb; border{-}radius:8px; padding:1rem 1.25rem; margin:0.75rem 0; font{-}family:inherit;";}
\NormalTok{  const stepsDiv = document.createElement("div");}
\NormalTok{  stepsDiv.style.marginBottom = "0.75rem";}
\NormalTok{  function renderSteps() \{}
\NormalTok{    stepsDiv.innerHTML = "";}
\NormalTok{    for (let i = 0; i \textless{} current; i++) \{}
\NormalTok{      const s = steps[i];}
\NormalTok{      const row = document.createElement("div");}
\NormalTok{      row.style.cssText = "display:grid; grid{-}template{-}columns:200px 1fr; gap:0.35rem 1rem; align{-}items:baseline; padding:0.35rem 0; border{-}top:1px solid \#f3f4f6;";}
\NormalTok{      const opCell = document.createElement("span");}
\NormalTok{      opCell.style.cssText = "font{-}size:0.85em; color:\#6b7280; font{-}style:italic;";}
\NormalTok{      opCell.textContent = s.op;}
\NormalTok{      const eqCell = document.createElement("span");}
\NormalTok{      eqCell.innerHTML = katex.renderToString(s.eq, \{ throwOnError: false, displayMode: false \});}
\NormalTok{      row.appendChild(opCell);}
\NormalTok{      row.appendChild(eqCell);}
\NormalTok{      if (s.note) \{}
\NormalTok{        const noteRow = document.createElement("div");}
\NormalTok{        noteRow.style.cssText = "grid{-}column:2; font{-}size:0.82em; color:\#6b7280; padding{-}bottom:0.2rem;";}
\NormalTok{        noteRow.textContent = s.note;}
\NormalTok{        row.appendChild(noteRow);}
\NormalTok{      \}}
\NormalTok{      stepsDiv.appendChild(row);}
\NormalTok{    \}}
\NormalTok{    if (current === totalSteps) \{}
\NormalTok{      const done = document.createElement("div");}
\NormalTok{      done.style.cssText = "margin{-}top:0.5rem; font{-}size:0.9em; color:\#059669; font{-}weight:500;";}
\NormalTok{      done.textContent = "✓ Solution complete";}
\NormalTok{      stepsDiv.appendChild(done);}
\NormalTok{    \}}
\NormalTok{  \}}
\NormalTok{  const controls = document.createElement("div");}
\NormalTok{  controls.style.cssText = "display:flex; gap:0.5rem; align{-}items:center;";}
\NormalTok{  const nextBtn = document.createElement("button");}
\NormalTok{  nextBtn.textContent = "Next step →";}
\NormalTok{  nextBtn.style.cssText = "padding:0.35rem 0.85rem; border:1px solid \#d1d5db; border{-}radius:4px; background:\#fff; cursor:pointer; font{-}size:0.9em;";}
\NormalTok{  const resetBtn = document.createElement("button");}
\NormalTok{  resetBtn.textContent = "Reset";}
\NormalTok{  resetBtn.style.cssText = "padding:0.35rem 0.75rem; border:1px solid \#e5e7eb; border{-}radius:4px; background:\#f9fafb; cursor:pointer; font{-}size:0.9em; color:\#6b7280;";}
\NormalTok{  const counter = document.createElement("span");}
\NormalTok{  counter.style.cssText = "font{-}size:0.82em; color:\#9ca3af; margin{-}left:0.25rem;";}
\NormalTok{  function updateButtons() \{}
\NormalTok{    nextBtn.disabled = current === totalSteps;}
\NormalTok{    nextBtn.style.opacity = current === totalSteps ? "0.4" : "1";}
\NormalTok{    counter.textContent = current === 0 ? "Click to reveal steps" : \textasciigrave{}Step $\{current\} of $\{totalSteps\}\textasciigrave{};}
\NormalTok{  \}}
\NormalTok{  nextBtn.onclick = () =\textgreater{} \{ if (current \textless{} totalSteps) \{ current++; renderSteps(); updateButtons(); \} \};}
\NormalTok{  resetBtn.onclick = () =\textgreater{} \{ current = 0; renderSteps(); updateButtons(); \};}
\NormalTok{  controls.appendChild(nextBtn);}
\NormalTok{  controls.appendChild(resetBtn);}
\NormalTok{  controls.appendChild(counter);}
\NormalTok{  container.appendChild(stepsDiv);}
\NormalTok{  container.appendChild(controls);}
\NormalTok{  renderSteps();}
\NormalTok{  updateButtons();}
\NormalTok{  return container;}
\NormalTok{\}}
\end{Highlighting}
\end{Shaded}

\begin{Shaded}
\begin{Highlighting}[]
\NormalTok{//| echo: false}
\NormalTok{makeStepperHTML(1, [}
\NormalTok{  \{ op: "Identify opposite, adjacent, hypotenuse", eq: "\textbackslash{}\textbackslash{}text\{hypotenuse\} = 800\textbackslash{}\textbackslash{}text\{ m\},\textbackslash{}\textbackslash{}quad \textbackslash{}\textbackslash{}text\{opposite\} = 240\textbackslash{}\textbackslash{}text\{ m\}", note: "The slope surface is the hypotenuse; the vertical drop is opposite the angle at the base." \},}
\NormalTok{  \{ op: "Choose the ratio", eq: "\textbackslash{}\textbackslash{}sin(\textbackslash{}\textbackslash{}theta) = \textbackslash{}\textbackslash{}frac\{\textbackslash{}\textbackslash{}text\{opposite\}\}\{\textbackslash{}\textbackslash{}text\{hypotenuse\}\} = \textbackslash{}\textbackslash{}frac\{240\}\{800\}", note: null \},}
\NormalTok{  \{ op: "Simplify the fraction", eq: "\textbackslash{}\textbackslash{}sin(\textbackslash{}\textbackslash{}theta) = 0.3", note: null \},}
\NormalTok{  \{ op: "Apply inverse sine", eq: "\textbackslash{}\textbackslash{}theta = \textbackslash{}\textbackslash{}sin\^{}\{{-}1\}(0.3)", note: "Inverse sine asks: what angle has a sine of 0.3?" \},}
\NormalTok{  \{ op: "Calculate", eq: "\textbackslash{}\textbackslash{}theta \textbackslash{}\textbackslash{}approx 17.5°", note: null \},}
\NormalTok{  \{ op: "Check", eq: "\textbackslash{}\textbackslash{}sin(17.5°) \textbackslash{}\textbackslash{}approx 0.3007 \textbackslash{}\textbackslash{}approx 0.3 \textbackslash{}\textbackslash{}checkmark", note: null \}}
\NormalTok{])}
\end{Highlighting}
\end{Shaded}

\begin{center}\rule{0.5\linewidth}{0.5pt}\end{center}

\textbf{2.} A ladder 6 m long leans against a wall. It makes an angle of
72° with the ground. How high up the wall does the top of the ladder
reach? Give your answer to two decimal places.

\begin{Shaded}
\begin{Highlighting}[]
\NormalTok{//| echo: false}
\NormalTok{makeStepperHTML(2, [}
\NormalTok{  \{ op: "Identify opposite, adjacent, hypotenuse", eq: "\textbackslash{}\textbackslash{}text\{hypotenuse\} = 6\textbackslash{}\textbackslash{}text\{ m\},\textbackslash{}\textbackslash{}quad \textbackslash{}\textbackslash{}theta = 72°,\textbackslash{}\textbackslash{}quad \textbackslash{}\textbackslash{}text\{opposite\} = \textbackslash{}\textbackslash{}text\{wall height\}", note: "The ladder is the hypotenuse; the wall height is the side opposite the 72° angle at the base." \},}
\NormalTok{  \{ op: "Choose the ratio", eq: "\textbackslash{}\textbackslash{}sin(72°) = \textbackslash{}\textbackslash{}frac\{\textbackslash{}\textbackslash{}text\{wall height\}\}\{6\}", note: null \},}
\NormalTok{  \{ op: "Rearrange", eq: "\textbackslash{}\textbackslash{}text\{wall height\} = 6 \textbackslash{}\textbackslash{}times \textbackslash{}\textbackslash{}sin(72°)", note: null \},}
\NormalTok{  \{ op: "Calculate", eq: "\textbackslash{}\textbackslash{}text\{wall height\} = 6 \textbackslash{}\textbackslash{}times 0.9511 \textbackslash{}\textbackslash{}approx 5.71\textbackslash{}\textbackslash{}text\{ m\}", note: null \},}
\NormalTok{  \{ op: "Check", eq: "\textbackslash{}\textbackslash{}sin\^{}\{{-}1\}\textbackslash{}\textbackslash{}!\textbackslash{}\textbackslash{}left(\textbackslash{}\textbackslash{}frac\{5.71\}\{6\}\textbackslash{}\textbackslash{}right) = \textbackslash{}\textbackslash{}sin\^{}\{{-}1\}(0.9517) \textbackslash{}\textbackslash{}approx 72.1° \textbackslash{}\textbackslash{}approx 72° \textbackslash{}\textbackslash{}checkmark", note: null \}}
\NormalTok{])}
\end{Highlighting}
\end{Shaded}

\begin{center}\rule{0.5\linewidth}{0.5pt}\end{center}

\textbf{3.} A phone screen has a diagonal of 16.5 cm and the screen is
8.0 cm wide. What angle does the diagonal make with the bottom edge of
the screen?

\begin{Shaded}
\begin{Highlighting}[]
\NormalTok{//| echo: false}
\NormalTok{makeStepperHTML(3, [}
\NormalTok{  \{ op: "Identify opposite, adjacent, hypotenuse", eq: "\textbackslash{}\textbackslash{}text\{hypotenuse\} = 16.5\textbackslash{}\textbackslash{}text\{ cm\},\textbackslash{}\textbackslash{}quad \textbackslash{}\textbackslash{}text\{adjacent\} = 8.0\textbackslash{}\textbackslash{}text\{ cm (width, beside the angle)\}", note: "The angle is at the bottom corner; the width is adjacent, the diagonal is the hypotenuse." \},}
\NormalTok{  \{ op: "Choose the ratio", eq: "\textbackslash{}\textbackslash{}cos(\textbackslash{}\textbackslash{}theta) = \textbackslash{}\textbackslash{}frac\{\textbackslash{}\textbackslash{}text\{adjacent\}\}\{\textbackslash{}\textbackslash{}text\{hypotenuse\}\} = \textbackslash{}\textbackslash{}frac\{8.0\}\{16.5\}", note: null \},}
\NormalTok{  \{ op: "Compute the ratio", eq: "\textbackslash{}\textbackslash{}cos(\textbackslash{}\textbackslash{}theta) \textbackslash{}\textbackslash{}approx 0.4848", note: null \},}
\NormalTok{  \{ op: "Apply inverse cosine", eq: "\textbackslash{}\textbackslash{}theta = \textbackslash{}\textbackslash{}cos\^{}\{{-}1\}(0.4848) \textbackslash{}\textbackslash{}approx 61.0°", note: null \},}
\NormalTok{  \{ op: "Check", eq: "\textbackslash{}\textbackslash{}cos(61°) \textbackslash{}\textbackslash{}approx 0.4848,\textbackslash{}\textbackslash{}quad 16.5 \textbackslash{}\textbackslash{}times 0.4848 \textbackslash{}\textbackslash{}approx 8.0\textbackslash{}\textbackslash{}text\{ cm\} \textbackslash{}\textbackslash{}checkmark", note: null \}}
\NormalTok{])}
\end{Highlighting}
\end{Shaded}

\begin{center}\rule{0.5\linewidth}{0.5pt}\end{center}

\textbf{4.} A solar panel is mounted on a roof. To maximise energy
collection in Alberta, the panel should face the sun at an angle of
elevation of 53°. The vertical rise from the base of the panel to its
top edge is 1.2 m. How long is the panel surface?

\begin{Shaded}
\begin{Highlighting}[]
\NormalTok{//| echo: false}
\NormalTok{makeStepperHTML(4, [}
\NormalTok{  \{ op: "Identify opposite, adjacent, hypotenuse", eq: "\textbackslash{}\textbackslash{}theta = 53°,\textbackslash{}\textbackslash{}quad \textbackslash{}\textbackslash{}text\{opposite\} = 1.2\textbackslash{}\textbackslash{}text\{ m (vertical rise)\},\textbackslash{}\textbackslash{}quad \textbackslash{}\textbackslash{}text\{hypotenuse\} = \textbackslash{}\textbackslash{}text\{panel length\}", note: "The vertical rise is opposite the angle; the panel surface is the hypotenuse." \},}
\NormalTok{  \{ op: "Choose the ratio", eq: "\textbackslash{}\textbackslash{}sin(53°) = \textbackslash{}\textbackslash{}frac\{1.2\}\{\textbackslash{}\textbackslash{}text\{panel length\}\}", note: null \},}
\NormalTok{  \{ op: "Rearrange", eq: "\textbackslash{}\textbackslash{}text\{panel length\} = \textbackslash{}\textbackslash{}frac\{1.2\}\{\textbackslash{}\textbackslash{}sin(53°)\}", note: null \},}
\NormalTok{  \{ op: "Calculate", eq: "\textbackslash{}\textbackslash{}text\{panel length\} = \textbackslash{}\textbackslash{}frac\{1.2\}\{0.7986\} \textbackslash{}\textbackslash{}approx 1.50\textbackslash{}\textbackslash{}text\{ m\}", note: null \},}
\NormalTok{  \{ op: "Check", eq: "\textbackslash{}\textbackslash{}sin(53°) \textbackslash{}\textbackslash{}times 1.50 \textbackslash{}\textbackslash{}approx 0.7986 \textbackslash{}\textbackslash{}times 1.50 \textbackslash{}\textbackslash{}approx 1.20\textbackslash{}\textbackslash{}text\{ m\} \textbackslash{}\textbackslash{}checkmark", note: null \}}
\NormalTok{])}
\end{Highlighting}
\end{Shaded}

\begin{center}\rule{0.5\linewidth}{0.5pt}\end{center}

\textbf{5.} A zip line descends from a platform 18 m high. The cable
makes an angle of 35° with the horizontal. How long is the cable, and
how far from the base of the platform does it end?

\begin{Shaded}
\begin{Highlighting}[]
\NormalTok{//| echo: false}
\NormalTok{makeStepperHTML(5, [}
\NormalTok{  \{ op: "Identify opposite, adjacent, hypotenuse", eq: "\textbackslash{}\textbackslash{}text\{opposite\} = 18\textbackslash{}\textbackslash{}text\{ m (height)\},\textbackslash{}\textbackslash{}quad \textbackslash{}\textbackslash{}theta = 35°,\textbackslash{}\textbackslash{}quad \textbackslash{}\textbackslash{}text\{hypotenuse\} = \textbackslash{}\textbackslash{}text\{cable length\}", note: "The height of the platform is the side opposite the 35° angle at the far end of the cable." \},}
\NormalTok{  \{ op: "Find the cable length using sine", eq: "\textbackslash{}\textbackslash{}sin(35°) = \textbackslash{}\textbackslash{}frac\{18\}\{\textbackslash{}\textbackslash{}text\{cable\}\}", note: null \},}
\NormalTok{  \{ op: "Rearrange", eq: "\textbackslash{}\textbackslash{}text\{cable\} = \textbackslash{}\textbackslash{}frac\{18\}\{\textbackslash{}\textbackslash{}sin(35°)\} = \textbackslash{}\textbackslash{}frac\{18\}\{0.5736\} \textbackslash{}\textbackslash{}approx 31.4\textbackslash{}\textbackslash{}text\{ m\}", note: null \},}
\NormalTok{  \{ op: "Find the horizontal distance using cosine", eq: "\textbackslash{}\textbackslash{}cos(35°) = \textbackslash{}\textbackslash{}frac\{\textbackslash{}\textbackslash{}text\{horizontal\}\}\{\textbackslash{}\textbackslash{}text\{cable\}\} \textbackslash{}\textbackslash{}implies \textbackslash{}\textbackslash{}text\{horizontal\} = 31.4 \textbackslash{}\textbackslash{}times \textbackslash{}\textbackslash{}cos(35°)", note: null \},}
\NormalTok{  \{ op: "Calculate horizontal distance", eq: "\textbackslash{}\textbackslash{}text\{horizontal\} = 31.4 \textbackslash{}\textbackslash{}times 0.8192 \textbackslash{}\textbackslash{}approx 25.7\textbackslash{}\textbackslash{}text\{ m\}", note: null \},}
\NormalTok{  \{ op: "Check using Pythagoras", eq: "\textbackslash{}\textbackslash{}sqrt\{18\^{}2 + 25.7\^{}2\} = \textbackslash{}\textbackslash{}sqrt\{324 + 660.5\} = \textbackslash{}\textbackslash{}sqrt\{984.5\} \textbackslash{}\textbackslash{}approx 31.4\textbackslash{}\textbackslash{}text\{ m\} \textbackslash{}\textbackslash{}checkmark", note: null \}}
\NormalTok{])}
\end{Highlighting}
\end{Shaded}

\begin{center}\rule{0.5\linewidth}{0.5pt}\end{center}

\textbf{6.} A ladder 5 m long leans against a wall. Its base is 1.8 m
from the wall. (a) What angle does the ladder make with the ground? (b)
How high up the wall does it reach?

\begin{Shaded}
\begin{Highlighting}[]
\NormalTok{//| echo: false}
\NormalTok{makeStepperHTML(6, [}
\NormalTok{  \{ op: "Identify the sides", eq: "\textbackslash{}\textbackslash{}text\{hypotenuse\} = 5\textbackslash{}\textbackslash{}text\{ m\},\textbackslash{}\textbackslash{}quad \textbackslash{}\textbackslash{}text\{adjacent\} = 1.8\textbackslash{}\textbackslash{}text\{ m (base distance)\}", note: "The ladder is the hypotenuse; the base distance from the wall is adjacent to the angle at the ground." \},}
\NormalTok{  \{ op: "Find the angle using inverse cosine", eq: "\textbackslash{}\textbackslash{}cos(\textbackslash{}\textbackslash{}theta) = \textbackslash{}\textbackslash{}frac\{1.8\}\{5\} = 0.36", note: null \},}
\NormalTok{  \{ op: "Apply inverse cosine", eq: "\textbackslash{}\textbackslash{}theta = \textbackslash{}\textbackslash{}cos\^{}\{{-}1\}(0.36) \textbackslash{}\textbackslash{}approx 68.9°", note: null \},}
\NormalTok{  \{ op: "Find the height using sine", eq: "\textbackslash{}\textbackslash{}sin(68.9°) = \textbackslash{}\textbackslash{}frac\{\textbackslash{}\textbackslash{}text\{height\}\}\{5\}", note: null \},}
\NormalTok{  \{ op: "Calculate height", eq: "\textbackslash{}\textbackslash{}text\{height\} = 5 \textbackslash{}\textbackslash{}times \textbackslash{}\textbackslash{}sin(68.9°) \textbackslash{}\textbackslash{}approx 5 \textbackslash{}\textbackslash{}times 0.9331 \textbackslash{}\textbackslash{}approx 4.67\textbackslash{}\textbackslash{}text\{ m\}", note: null \},}
\NormalTok{  \{ op: "Check using Pythagoras", eq: "\textbackslash{}\textbackslash{}sqrt\{1.8\^{}2 + 4.67\^{}2\} = \textbackslash{}\textbackslash{}sqrt\{3.24 + 21.81\} = \textbackslash{}\textbackslash{}sqrt\{25.05\} \textbackslash{}\textbackslash{}approx 5.0\textbackslash{}\textbackslash{}text\{ m\} \textbackslash{}\textbackslash{}checkmark", note: null \}}
\NormalTok{])}
\end{Highlighting}
\end{Shaded}

\chapter{Polynomials and factoring}\label{polynomials-and-factoring}

When the curve is the point

\hfill\break

You kick a ball straight up. For a moment it keeps going --- then it
slows, stops, and falls back down. The height isn't a straight line; it
curves up and back. That curve has an equation.

You're fencing off a rectangular garden and you've got 40 metres of
fencing. You want the biggest possible area. If you make it long and
thin, the area shrinks. If you make it square, the area is as large as
it can be. Finding that optimum --- that maximum --- is a polynomial
problem.

Your school sells tickets to a performance. At a lower price more people
buy, but revenue per ticket drops. At a higher price fewer people buy.
Somewhere in the middle is the price that brings in the most money. That
relationship between price and total revenue follows a curve --- a
polynomial.

Straight lines describe constant rates. Polynomials describe the rest.

\section{What the notation is
saying}\label{what-the-notation-is-saying-12}

Here is a polynomial:

\[ax^2 + bx + c\]

Read it aloud: \emph{a times x-squared, plus b times x, plus c.}

The letters \(a\), \(b\), and \(c\) are \textbf{coefficients} --- the
numbers multiplying each power of \(x\). They're just fixed numbers;
you'll swap them for actual values in any particular problem. The letter
\(x\) is the unknown --- the variable you're either solving for, or
letting run freely while you watch what the output does. When you write
\(y = ax^2 + bx + c\), \(y\) is the output --- the value the polynomial
produces for each value of \(x\).

Each piece --- \(ax^2\), \(bx\), \(c\) --- is called a \textbf{term}.
Count the terms and you know you have three. The \(c\) term has no \(x\)
at all; it's a plain number, a constant.

The \textbf{degree} of a polynomial is the highest power of \(x\)
present. In \(ax^2 + bx + c\) the highest power is 2, so the degree is
2. Degree-2 polynomials are called \textbf{quadratics}. Degree-1
polynomials --- like \(3x + 7\) --- are the linear equations from Vol 2.
Degree-3 polynomials are called cubics, and so on.

Why does the \(x^2\) term make the graph curve? Because \(x^2\) grows
faster than \(x\). When \(x = 2\), \(x^2 = 4\). When \(x = 10\),
\(x^2 = 100\). The squared term pulls the output away from the straight
line. Take \(a = 0\) and you're back to a line. Leave \(a\) nonzero and
the whole shape bends.

A specific example: \(y = x^2 - 5x + 6\). Here \(a = 1\), \(b = -5\),
\(c = 6\). The negative \(b\) means the straight-line part is pulling
down, and the positive \(c\) shifts the whole thing up by 6 units. The
coefficient \(a = 1\) means the parabola opens upward --- a positive
\(a\) always curves upward, a negative \(a\) always curves downward.

\section{The method}\label{the-method-11}

\subsection{Part a: Expanding}\label{part-a-expanding}

Before you can factor, it helps to know how to expand --- to go from the
factored form back to the full polynomial. This makes the reverse step
(factoring) easier to recognise.

The distributive law says: multiply everything inside the brackets by
everything outside.

\[3(x + 4) = 3x + 12\]

When you have two brackets multiplied together, you apply the same idea
twice --- every term in the first bracket multiplies every term in the
second. This is sometimes called \textbf{FOIL} (First, Outer, Inner,
Last), though the name matters less than the idea: don't miss any pair.

\[(x + 3)(x + 2)\]

\begin{itemize}
\tightlist
\item
  First terms: \(x \cdot x = x^2\)
\item
  Outer terms: \(x \cdot 2 = 2x\)
\item
  Inner terms: \(3 \cdot x = 3x\)
\item
  Last terms: \(3 \cdot 2 = 6\)
\end{itemize}

Add them up: \(x^2 + 2x + 3x + 6 = x^2 + 5x + 6\).

The pattern is worth noticing. With \((x + 3)(x + 2)\), the number in
the middle (\(5\)) is the sum of 3 and 2. The constant term (\(6\)) is
their product. That observation is what makes factoring work in reverse.

\subsection{Part b: Factoring}\label{part-b-factoring}

Factoring means breaking a polynomial back into the brackets that
produced it. There are three main methods; apply them in this order.

\textbf{Method 1: Common factor.} Check whether every term shares a
factor you can pull out.

\[6x^2 + 9x = 3x(2x + 3)\]

The common factor here is \(3x\). Dividing both terms by \(3x\) gives
the bracket. Always check for a common factor first --- it simplifies
what comes next.

\textbf{Method 2: Difference of squares.} When you have two perfect
squares separated by a minus sign, it factors as:

\[a^2 - b^2 = (a + b)(a - b)\]

For example:

\[x^2 - 25 = (x + 5)(x - 5)\]

Check it by expanding:
\((x + 5)(x - 5) = x^2 - 5x + 5x - 25 = x^2 - 25\). Yes.

This works because the middle terms cancel. It doesn't work if the sign
between the squares is positive --- \(x^2 + 25\) doesn't factor over the
real numbers.

\textbf{Method 3: Factoring a trinomial \(ax^2 + bx + c\).} When
\(a = 1\), look for two numbers that multiply to \(c\) and add to \(b\).

\[x^2 - 5x + 6\]

You need two numbers that multiply to \(6\) and add to \(-5\). Try:
\(-2\) and \(-3\). Product: \((-2)(-3) = 6\). Sum: \((-2) + (-3) = -5\).
Yes.

\[x^2 - 5x + 6 = (x - 2)(x - 3)\]

When \(a \neq 1\), try the \textbf{grouping method} (also called the
\(ac\) method). Multiply \(a\) by \(c\), find two numbers with that
product and sum \(b\), then split the middle term and factor in pairs.

\[2x^2 + 7x + 3\]

Multiply \(a \cdot c = 2 \cdot 3 = 6\). Find two numbers that multiply
to \(6\) and add to \(7\): those are \(1\) and \(6\).

Split the middle term:

\[2x^2 + x + 6x + 3\]

Factor in pairs:

\[x(2x + 1) + 3(2x + 1) = (x + 3)(2x + 1)\]

Both terms contain the factor \((2x + 1)\), so it can be pulled out ---
the same distributive law used in expanding, run backwards.

Check by expanding:
\((x + 3)(2x + 1) = 2x^2 + x + 6x + 3 = 2x^2 + 7x + 3\). Yes.

\subsection{Part c: Solving quadratic
equations}\label{part-c-solving-quadratic-equations}

A \textbf{quadratic equation} is when you set a quadratic equal to zero:

\[ax^2 + bx + c = 0\]

There are two methods.

\textbf{Solving by factoring.} If you can factor the left side, you can
use the fact that if two things multiply to zero, at least one of them
must be zero.

\[x^2 - 5x + 6 = 0\]

\[(x - 2)(x - 3) = 0\]

Either \((x - 2) = 0\) or \((x - 3) = 0\). So \(x = 2\) or \(x = 3\).

Check both: \(4 - 10 + 6 = 0\). Yes. \(9 - 15 + 6 = 0\). Yes.

\textbf{Solving with the quadratic formula.} Not every quadratic factors
neatly. For any quadratic \(ax^2 + bx + c = 0\), the solutions are
always:

\[x = \frac{-b \pm \sqrt{b^2 - 4ac}}{2a}\]

Read it: \emph{x equals negative b, plus or minus the square root of
b-squared minus four-a-c, all over two-a.}

The \(\pm\) sign means there are two candidate answers --- one where you
add the square root term, one where you subtract it (whether there are
two, one, or no real answers depends on the discriminant --- that's the
next section).

The piece under the square root --- \(b^2 - 4ac\) --- is called the
\textbf{discriminant}, often written \(\Delta\) (delta). It tells you
how many solutions exist before you do any arithmetic:

\begin{itemize}
\tightlist
\item
  \(\Delta > 0\): two real solutions (the parabola crosses the x-axis
  twice)
\item
  \(\Delta = 0\): one solution (the parabola just touches the x-axis at
  one point)
\item
  \(\Delta < 0\): no real solutions (the parabola doesn't reach the
  x-axis)
\end{itemize}

\subsection{Part d: Reading the graph}\label{part-d-reading-the-graph}

Every quadratic \(y = ax^2 + bx + c\) draws a \textbf{parabola} --- a
symmetric U-shape (or inverted U if \(a < 0\)).

The \textbf{roots} are where the parabola crosses the x-axis --- the
values of \(x\) where \(y = 0\). These are exactly the solutions to
\(ax^2 + bx + c = 0\). Factoring gives you the roots. The roots are the
x-intercepts of the graph.

The \textbf{vertex} is the turning point --- the very top of a downward
parabola (the maximum) or the very bottom of an upward one (the
minimum). The x-coordinate of the vertex is:

\[x_v = -\frac{b}{2a}\]

Read it: \emph{negative b over two a.} Substitute this back into the
equation to find the y-coordinate --- the maximum or minimum value
itself.

For \(y = x^2 - 5x + 6\): \(x_v = -(-5)/(2 \cdot 1) = 5/2 = 2.5\). Then
\(y = (2.5)^2 - 5(2.5) + 6 = 6.25 - 12.5 + 6 = -0.25\). The vertex is at
\((2.5, -0.25)\) --- the minimum point, halfway between the two roots at
\(x = 2\) and \(x = 3\).

\begin{Shaded}
\begin{Highlighting}[]
\NormalTok{//| echo: false}
\NormalTok{viewof qa = Inputs.range([{-}2, 2], \{ value: 1, step: 0.5, label: "a" \})}
\end{Highlighting}
\end{Shaded}

\begin{Shaded}
\begin{Highlighting}[]
\NormalTok{//| echo: false}
\NormalTok{viewof qb = Inputs.range([{-}6, 6], \{ value: {-}5, step: 1, label: "b" \})}
\end{Highlighting}
\end{Shaded}

\begin{Shaded}
\begin{Highlighting}[]
\NormalTok{//| echo: false}
\NormalTok{viewof qc = Inputs.range([{-}9, 9], \{ value: 6, step: 1, label: "c" \})}
\end{Highlighting}
\end{Shaded}

\begin{Shaded}
\begin{Highlighting}[]
\NormalTok{//| echo: false}
\NormalTok{\{}
\NormalTok{  // Guard: skip if a is too close to zero}
\NormalTok{  const a = Math.abs(qa) \textless{} 0.1 ? 0.5 : qa;}
\NormalTok{  const b = qb, c = qc;}

\NormalTok{  const svgW = 500, svgH = 360;}
\NormalTok{  const margin = \{ top: 50, right: 30, bottom: 40, left: 50 \};}
\NormalTok{  const W = svgW {-} margin.left {-} margin.right;}
\NormalTok{  const H = svgH {-} margin.top {-} margin.bottom;}

\NormalTok{  const xMin = {-}6, xMax = 6;}
\NormalTok{  const xs = d3.range(xMin, xMax + 0.05, 0.05);}
\NormalTok{  const ys = xs.map(x =\textgreater{} a * x * x + b * x + c);}
\NormalTok{  const yMin = Math.min(...ys, {-}10), yMax = Math.max(...ys, 10);}

\NormalTok{  const xScale = d3.scaleLinear([xMin, xMax], [0, W]);}
\NormalTok{  const yScale = d3.scaleLinear([yMin, yMax], [H, 0]);}

\NormalTok{  const disc = b * b {-} 4 * a * c;}

\NormalTok{  // Roots}
\NormalTok{  let roots = [];}
\NormalTok{  if (disc \textgreater{} 0) \{}
\NormalTok{    roots = [}
\NormalTok{      ({-}b + Math.sqrt(disc)) / (2 * a),}
\NormalTok{      ({-}b {-} Math.sqrt(disc)) / (2 * a)}
\NormalTok{    ].filter(r =\textgreater{} r \textgreater{}= xMin \&\& r \textless{}= xMax);}
\NormalTok{  \} else if (Math.abs(disc) \textless{} 1e{-}9) \{}
\NormalTok{    const r = {-}b / (2 * a);}
\NormalTok{    if (r \textgreater{}= xMin \&\& r \textless{}= xMax) roots = [r];}
\NormalTok{  \}}

\NormalTok{  // Vertex}
\NormalTok{  const vx = {-}b / (2 * a);}
\NormalTok{  const vy = a * vx * vx + b * vx + c;}

\NormalTok{  const svg = d3.create("svg")}
\NormalTok{    .attr("width", svgW).attr("height", svgH)}
\NormalTok{    .style("font{-}family", "inherit");}

\NormalTok{  const g = svg.append("g")}
\NormalTok{    .attr("transform", \textasciigrave{}translate($\{margin.left\},$\{margin.top\})\textasciigrave{});}

\NormalTok{  // Axes}
\NormalTok{  if (yScale(0) \textgreater{}= 0 \&\& yScale(0) \textless{}= H) \{}
\NormalTok{    g.append("line")}
\NormalTok{      .attr("x1", 0).attr("y1", yScale(0))}
\NormalTok{      .attr("x2", W).attr("y2", yScale(0))}
\NormalTok{      .attr("stroke", "\#9ca3af").attr("stroke{-}width", 1);}
\NormalTok{  \}}
\NormalTok{  if (xScale(0) \textgreater{}= 0 \&\& xScale(0) \textless{}= W) \{}
\NormalTok{    g.append("line")}
\NormalTok{      .attr("x1", xScale(0)).attr("y1", 0)}
\NormalTok{      .attr("x2", xScale(0)).attr("y2", H)}
\NormalTok{      .attr("stroke", "\#9ca3af").attr("stroke{-}width", 1);}
\NormalTok{  \}}

\NormalTok{  // Grid lines}
\NormalTok{  for (let xi = xMin; xi \textless{}= xMax; xi++) \{}
\NormalTok{    g.append("line")}
\NormalTok{      .attr("x1", xScale(xi)).attr("y1", 0)}
\NormalTok{      .attr("x2", xScale(xi)).attr("y2", H)}
\NormalTok{      .attr("stroke", "\#f3f4f6").attr("stroke{-}width", 1);}
\NormalTok{    g.append("text")}
\NormalTok{      .attr("x", xScale(xi)).attr("y", yScale(0) + 14)}
\NormalTok{      .attr("text{-}anchor", "middle").attr("font{-}size", "11px").attr("fill", "\#9ca3af")}
\NormalTok{      .text(xi === 0 ? "" : xi);}
\NormalTok{  \}}

\NormalTok{  // Curve}
\NormalTok{  const line = d3.line()}
\NormalTok{    .x(x =\textgreater{} xScale(x))}
\NormalTok{    .y(x =\textgreater{} yScale(a * x * x + b * x + c))}
\NormalTok{    .curve(d3.curveCatmullRom);}

\NormalTok{  g.append("path")}
\NormalTok{    .datum(xs)}
\NormalTok{    .attr("d", line)}
\NormalTok{    .attr("fill", "none")}
\NormalTok{    .attr("stroke", "\#2563eb")}
\NormalTok{    .attr("stroke{-}width", 2.5);}

\NormalTok{  // Roots}
\NormalTok{  roots.forEach(r =\textgreater{} \{}
\NormalTok{    g.append("circle")}
\NormalTok{      .attr("cx", xScale(r)).attr("cy", yScale(0))}
\NormalTok{      .attr("r", 5)}
\NormalTok{      .attr("fill", "\#dc2626").attr("stroke", "\#fff").attr("stroke{-}width", 1.5);}
\NormalTok{  \});}

\NormalTok{  // Vertex}
\NormalTok{  if (vx \textgreater{}= xMin \&\& vx \textless{}= xMax) \{}
\NormalTok{    g.append("circle")}
\NormalTok{      .attr("cx", xScale(vx)).attr("cy", yScale(vy))}
\NormalTok{      .attr("r", 6)}
\NormalTok{      .attr("fill", "\#f59e0b").attr("stroke", "\#fff").attr("stroke{-}width", 1.5);}
\NormalTok{  \}}

\NormalTok{  // "No real roots" label}
\NormalTok{  if (disc \textless{} 0) \{}
\NormalTok{    g.append("text")}
\NormalTok{      .attr("x", W / 2).attr("y", H / 2)}
\NormalTok{      .attr("text{-}anchor", "middle").attr("font{-}size", "13px").attr("fill", "\#dc2626")}
\NormalTok{      .text("no real roots (Δ \textless{} 0)");}
\NormalTok{  \}}

\NormalTok{  // Info panel at top}
\NormalTok{  const eqStr = \textasciigrave{}y = $\{a\}x² $\{b \textgreater{}= 0 ? "+" : ""\}$\{b\}x $\{c \textgreater{}= 0 ? "+" : ""\}$\{c\}\textasciigrave{};}
\NormalTok{  svg.append("text")}
\NormalTok{    .attr("x", margin.left).attr("y", 22)}
\NormalTok{    .attr("font{-}size", "13px").attr("font{-}weight", "600").attr("fill", "\#111827")}
\NormalTok{    .text(eqStr);}

\NormalTok{  svg.append("text")}
\NormalTok{    .attr("x", margin.left + 220).attr("y", 22)}
\NormalTok{    .attr("font{-}size", "13px").attr("fill", disc \textgreater{}= 0 ? "\#059669" : "\#dc2626")}
\NormalTok{    .text(\textasciigrave{}Δ = $\{disc.toFixed(2)\}  ($\{disc \textgreater{} 0 ? "2 roots" : disc === 0 ? "1 root" : "no real roots"\})\textasciigrave{});}

\NormalTok{  // Vertex label}
\NormalTok{  if (vx \textgreater{}= xMin \&\& vx \textless{}= xMax) \{}
\NormalTok{    svg.append("text")}
\NormalTok{      .attr("x", margin.left + xScale(vx) + 8)}
\NormalTok{      .attr("y", margin.top + yScale(vy) {-} 4)}
\NormalTok{      .attr("font{-}size", "11px").attr("fill", "\#d97706")}
\NormalTok{      .text(\textasciigrave{}vertex ($\{vx.toFixed(2)\}, $\{vy.toFixed(2)\})\textasciigrave{});}
\NormalTok{  \}}

\NormalTok{  return svg.node();}
\NormalTok{\}}
\end{Highlighting}
\end{Shaded}

\section{Why this works}\label{why-this-works-1}

\begin{tcolorbox}[enhanced jigsaw, opacityback=0, colback=white, breakable, colframe=quarto-callout-note-color-frame, arc=.35mm, left=2mm, toprule=.15mm, leftrule=.75mm, bottomrule=.15mm, rightrule=.15mm]
\begin{minipage}[t]{5.5mm}
\textcolor{quarto-callout-note-color}{\faInfo}
\end{minipage}%
\begin{minipage}[t]{\textwidth - 5.5mm}

\textbf{Where the quadratic formula comes from}

The formula isn't magic --- it's derived by solving
\(ax^2 + bx + c = 0\) directly, using a technique called
\textbf{completing the square}.

The key move: rewrite \(x^2 + \frac{b}{a}x\) as a perfect square minus a
correction term.

\[x^2 + \frac{b}{a}x = \left(x + \frac{b}{2a}\right)^2 - \frac{b^2}{4a^2}\]

Substitute this into \(ax^2 + bx + c = 0\), divide through by \(a\),
then isolate the squared bracket:

\[\left(x + \frac{b}{2a}\right)^2 = \frac{b^2 - 4ac}{4a^2}\]

Take the square root of both sides (giving \(\pm\)), then subtract
\(\frac{b}{2a}\):

\[x = \frac{-b \pm \sqrt{b^2 - 4ac}}{2a}\]

That's it. The formula is just completing the square, done once in
general so you don't have to do it every time.

\end{minipage}%
\end{tcolorbox}

\section{Worked examples}\label{worked-examples-11}

\textbf{Example 1 (sci) --- height of a thrown ball}

A ball is thrown upward from ground level with an initial speed of 14
m/s. Ignoring air resistance, its height in metres after \(t\) seconds
is:

\[h(t) = 14t - 4.9t^2\]

When does it land?

Landing means \(h = 0\):

\[14t - 4.9t^2 = 0\]

Factor out a common term of \(t\):

\[t(14 - 4.9t) = 0\]

Either \(t = 0\) (the moment it's thrown) or \(14 - 4.9t = 0\). Solving
the second:

\[4.9t = 14 \implies t = \frac{14}{4.9} \approx 2.86 \text{ s}\]

The ball lands after about 2.86 seconds.

The vertex --- the highest point --- is at
\(t = 14/(2 \times 4.9) = 14/9.8
\approx 1.43\) s, at height \(h(1.43) = 14(1.43) - 4.9(1.43)^2 \approx
20.02 - 10.02 \approx 10\) m.

\begin{Shaded}
\begin{Highlighting}[]
\NormalTok{//| echo: false}
\NormalTok{viewof ballH0 = Inputs.range([0, 5], \{ value: 0, step: 0.5, label: "Initial height h₀ (m)" \})}
\end{Highlighting}
\end{Shaded}

\begin{Shaded}
\begin{Highlighting}[]
\NormalTok{//| echo: false}
\NormalTok{viewof ballV = Inputs.range([1, 20], \{ value: 14, step: 1, label: "Initial speed v (m/s)" \})}
\end{Highlighting}
\end{Shaded}

\begin{Shaded}
\begin{Highlighting}[]
\NormalTok{//| echo: false}
\NormalTok{\{}
\NormalTok{  const h0 = ballH0, v = ballV;}
\NormalTok{  // h(t) = h0 + v*t {-} 4.9*t\^{}2}
\NormalTok{  // Landing: positive root of {-}4.9t\^{}2 + v*t + h0 = 0}
\NormalTok{  const disc = v * v + 4 * 4.9 * h0;}
\NormalTok{  const tLand = (v + Math.sqrt(disc)) / (2 * 4.9);}
\NormalTok{  const tPeak = v / (2 * 4.9);}
\NormalTok{  const hPeak = h0 + v * tPeak {-} 4.9 * tPeak * tPeak;}

\NormalTok{  const svgW = 500, svgH = 320;}
\NormalTok{  const margin = \{ top: 40, right: 30, bottom: 40, left: 55 \};}
\NormalTok{  const W = svgW {-} margin.left {-} margin.right;}
\NormalTok{  const H = svgH {-} margin.top {-} margin.bottom;}

\NormalTok{  const tMax = tLand * 1.05;}
\NormalTok{  const tArr = d3.range(0, tLand + 0.01, tLand / 200);}
\NormalTok{  const hArr = tArr.map(t =\textgreater{} Math.max(0, h0 + v * t {-} 4.9 * t * t));}

\NormalTok{  const xScale = d3.scaleLinear([0, tMax], [0, W]);}
\NormalTok{  const yMax = hPeak * 1.15 + 1;}
\NormalTok{  const yScale = d3.scaleLinear([0, yMax], [H, 0]);}

\NormalTok{  const svg = d3.create("svg")}
\NormalTok{    .attr("width", svgW).attr("height", svgH)}
\NormalTok{    .style("font{-}family", "inherit");}

\NormalTok{  const g = svg.append("g")}
\NormalTok{    .attr("transform", \textasciigrave{}translate($\{margin.left\},$\{margin.top\})\textasciigrave{});}

\NormalTok{  // Axes}
\NormalTok{  g.append("line").attr("x1", 0).attr("y1", H).attr("x2", W).attr("y2", H)}
\NormalTok{    .attr("stroke", "\#6b7280").attr("stroke{-}width", 1.5);}
\NormalTok{  g.append("line").attr("x1", 0).attr("y1", 0).attr("x2", 0).attr("y2", H)}
\NormalTok{    .attr("stroke", "\#6b7280").attr("stroke{-}width", 1.5);}

\NormalTok{  // Axis labels}
\NormalTok{  g.append("text").attr("x", W / 2).attr("y", H + 32)}
\NormalTok{    .attr("text{-}anchor", "middle").attr("font{-}size", "12px").attr("fill", "\#374151")}
\NormalTok{    .text("time t (seconds)");}
\NormalTok{  g.append("text")}
\NormalTok{    .attr("transform", "rotate({-}90)")}
\NormalTok{    .attr("x", {-}H / 2).attr("y", {-}40)}
\NormalTok{    .attr("text{-}anchor", "middle").attr("font{-}size", "12px").attr("fill", "\#374151")}
\NormalTok{    .text("height h (metres)");}

\NormalTok{  // Curve}
\NormalTok{  const lineGen = d3.line()}
\NormalTok{    .x((\_, i) =\textgreater{} xScale(tArr[i]))}
\NormalTok{    .y(h =\textgreater{} yScale(h));}

\NormalTok{  g.append("path")}
\NormalTok{    .datum(hArr)}
\NormalTok{    .attr("d", lineGen)}
\NormalTok{    .attr("fill", "none")}
\NormalTok{    .attr("stroke", "\#2563eb")}
\NormalTok{    .attr("stroke{-}width", 2.5);}

\NormalTok{  // Peak marker}
\NormalTok{  g.append("circle")}
\NormalTok{    .attr("cx", xScale(tPeak)).attr("cy", yScale(hPeak))}
\NormalTok{    .attr("r", 6).attr("fill", "\#f59e0b").attr("stroke", "\#fff").attr("stroke{-}width", 1.5);}

\NormalTok{  // Landing marker}
\NormalTok{  g.append("circle")}
\NormalTok{    .attr("cx", xScale(tLand)).attr("cy", yScale(0))}
\NormalTok{    .attr("r", 6).attr("fill", "\#dc2626").attr("stroke", "\#fff").attr("stroke{-}width", 1.5);}

\NormalTok{  // Labels}
\NormalTok{  g.append("text")}
\NormalTok{    .attr("x", xScale(tPeak) + 8).attr("y", yScale(hPeak) {-} 8)}
\NormalTok{    .attr("font{-}size", "12px").attr("fill", "\#d97706")}
\NormalTok{    .text(\textasciigrave{}peak: $\{hPeak.toFixed(2)\} m at t=$\{tPeak.toFixed(2)\} s\textasciigrave{});}

\NormalTok{  g.append("text")}
\NormalTok{    .attr("x", xScale(tLand) {-} 8).attr("y", yScale(0) + 18)}
\NormalTok{    .attr("text{-}anchor", "end").attr("font{-}size", "12px").attr("fill", "\#dc2626")}
\NormalTok{    .text(\textasciigrave{}lands at t=$\{tLand.toFixed(2)\} s\textasciigrave{});}

\NormalTok{  // Title}
\NormalTok{  svg.append("text")}
\NormalTok{    .attr("x", margin.left).attr("y", 24)}
\NormalTok{    .attr("font{-}size", "13px").attr("font{-}weight", "600").attr("fill", "\#111827")}
\NormalTok{    .text(\textasciigrave{}h(t) = $\{h0\} + $\{v\}t − 4.9t²\textasciigrave{});}

\NormalTok{  return svg.node();}
\NormalTok{\}}
\end{Highlighting}
\end{Shaded}

\begin{center}\rule{0.5\linewidth}{0.5pt}\end{center}

\textbf{Example 2 (eng) --- maximum area from fixed fencing}

You have 40 m of fencing and want to enclose a rectangular garden with
the largest possible area. What dimensions give the maximum area?

Let the width be \(x\) metres. Since the total perimeter is 40 m, the
two widths use \(2x\) metres, leaving \(40 - 2x\) for the two lengths.
Each length is \((40 - 2x)/2 = 20 - x\) metres.

Area:

\[A = x(20 - x) = 20x - x^2\]

This is a downward parabola (\(a = -1\)). Maximum at:

\[x_v = -\frac{20}{2 \times (-1)} = 10 \text{ m}\]

Length \(= 20 - 10 = 10\) m. Area \(= 10 \times 10 = 100\) m². A square
encloses the most area for a given perimeter.

\begin{center}\rule{0.5\linewidth}{0.5pt}\end{center}

\textbf{Example 3 (biz) --- maximum revenue from a pricing function}

A small business sells handmade candles. Market research shows that at a
price of \(p\) dollars each, the number sold per week is
\(q = 60 - 2p\). Revenue is:

\[R = p \cdot q = p(60 - 2p) = 60p - 2p^2\]

This is a downward parabola (\(a = -2\)). Maximum revenue at:

\[p_v = -\frac{60}{2 \times (-2)} = -\frac{60}{-4} = 15\]

Price of \$15. Revenue: \(R = 15 \times (60 - 30) = 15 \times 30 = 450\)
per week.

To confirm it's a maximum and not a minimum: \(a = -2 < 0\), so the
parabola opens downward. The vertex is the top. Any price above or below
\$15 gives less revenue.

\begin{center}\rule{0.5\linewidth}{0.5pt}\end{center}

\textbf{Example 4 (comp) --- character jump arc}

A game character jumps. The height of the character above the ground in
metres at time \(t\) seconds after the jump is:

\[h(t) = -5t^2 + 10t\]

Find the maximum height and the time the character lands.

The parabola opens downward (\(a = -5\)), so the vertex is the maximum.

Vertex time: \(t_v = -10/(2 \times -5) = 1\) second.

Maximum height: \(h(1) = -5(1)^2 + 10(1) = -5 + 10 = 5\) metres.

Landing: set \(h = 0\).

\[-5t^2 + 10t = 0\] \[t(-5t + 10) = 0\]

\(t = 0\) (take-off) or \(t = 2\) seconds (landing).

Evaluating a polynomial at a specific value --- like checking the height
at \(t = 1\) or \(t = 2\) --- is a calculation every game engine runs
many times per second to determine where characters and objects are.

\section{Where this goes}\label{where-this-goes-12}

The natural next step is \textbf{Vol 3, Chapter 4: Functions and
relations}. A quadratic is a function --- it takes an \(x\) value and
produces a \(y\) value, one output for each input. Once you have the
function concept, the parabola becomes one member of a large family:
linear functions, quadratic functions, exponential functions,
trigonometric functions. The machinery you've built here --- reading
coefficients, finding roots, locating vertices --- applies across all of
them.

The bigger horizon is \textbf{calculus}. The vertex formula
\(x_v = -b/(2a)\) gives you the maximum or minimum, but it only works
for quadratics. What about a cubic, or a more complicated curve? That's
what differential calculus (Vol 5) answers: a general method for finding
the exact peak or trough of any smooth curve. The quadratic vertex is
the simplest case of a much more powerful idea. When you reach it, what
you learned here will already be familiar.

\begin{tcolorbox}[enhanced jigsaw, opacityback=0, colback=white, breakable, colframe=quarto-callout-tip-color-frame, arc=.35mm, left=2mm, toprule=.15mm, leftrule=.75mm, bottomrule=.15mm, rightrule=.15mm]
\begin{minipage}[t]{5.5mm}
\textcolor{quarto-callout-tip-color}{\faLightbulb}
\end{minipage}%
\begin{minipage}[t]{\textwidth - 5.5mm}

\textbf{Where this shows up}

\begin{itemize}
\tightlist
\item
  A structural engineer calculating deflection in a beam under a
  distributed load solves a polynomial equation.
\item
  A physicist modelling projectile trajectory --- range, maximum height,
  time of flight --- works with quadratics throughout.
\item
  A pricing analyst finding the revenue-maximising price point solves
  \(dR/dp = 0\), which for a linear demand function produces a
  quadratic.
\item
  A computer graphics programmer evaluating a Bezier spline is computing
  a polynomial at a parameter value thousands of times per frame.
\item
  A data scientist fitting a parabola to experimental data (polynomial
  regression of degree 2) uses the same \(ax^2 + bx + c\) form.
\end{itemize}

\end{minipage}%
\end{tcolorbox}

\section{Exercises}\label{exercises-12}

These are puzzles. Each has a clean answer. The interesting part is
setting up the equation first, then solving it.

\begin{enumerate}
\def\labelenumi{\arabic{enumi}.}
\tightlist
\item
  A ball is thrown upward from a cliff 20 m above the ground with an
  initial upward speed of 10 m/s. Its height above the ground after
  \(t\) seconds is \(h(t) = 20 + 10t - 4.9t^2\). When does it hit the
  ground?
\end{enumerate}

\begin{Shaded}
\begin{Highlighting}[]
\NormalTok{//| echo: false}
\NormalTok{katex = require("katex@0.16")}
\end{Highlighting}
\end{Shaded}

\begin{Shaded}
\begin{Highlighting}[]
\NormalTok{//| echo: false}
\NormalTok{function makeStepperHTML(exerciseNum, steps) \{}
\NormalTok{  let current = 0;}
\NormalTok{  const totalSteps = steps.length;}
\NormalTok{  const container = document.createElement("div");}
\NormalTok{  container.style.cssText = "border:1px solid \#e5e7eb; border{-}radius:8px; padding:1rem 1.25rem; margin:0.75rem 0; font{-}family:inherit;";}
\NormalTok{  const stepsDiv = document.createElement("div");}
\NormalTok{  stepsDiv.style.marginBottom = "0.75rem";}
\NormalTok{  function renderSteps() \{}
\NormalTok{    stepsDiv.innerHTML = "";}
\NormalTok{    for (let i = 0; i \textless{} current; i++) \{}
\NormalTok{      const s = steps[i];}
\NormalTok{      const row = document.createElement("div");}
\NormalTok{      row.style.cssText = "display:grid; grid{-}template{-}columns:200px 1fr; gap:0.35rem 1rem; align{-}items:baseline; padding:0.35rem 0; border{-}top:1px solid \#f3f4f6;";}
\NormalTok{      const opCell = document.createElement("span");}
\NormalTok{      opCell.style.cssText = "font{-}size:0.85em; color:\#6b7280; font{-}style:italic;";}
\NormalTok{      opCell.textContent = s.op;}
\NormalTok{      const eqCell = document.createElement("span");}
\NormalTok{      eqCell.innerHTML = katex.renderToString(s.eq, \{ throwOnError: false, displayMode: false \});}
\NormalTok{      row.appendChild(opCell);}
\NormalTok{      row.appendChild(eqCell);}
\NormalTok{      if (s.note) \{}
\NormalTok{        const noteRow = document.createElement("div");}
\NormalTok{        noteRow.style.cssText = "grid{-}column:2; font{-}size:0.82em; color:\#6b7280; padding{-}bottom:0.2rem;";}
\NormalTok{        noteRow.textContent = s.note;}
\NormalTok{        row.appendChild(noteRow);}
\NormalTok{      \}}
\NormalTok{      stepsDiv.appendChild(row);}
\NormalTok{    \}}
\NormalTok{    if (current === totalSteps) \{}
\NormalTok{      const done = document.createElement("div");}
\NormalTok{      done.style.cssText = "margin{-}top:0.5rem; font{-}size:0.9em; color:\#059669; font{-}weight:500;";}
\NormalTok{      done.textContent = "✓ Solution complete";}
\NormalTok{      stepsDiv.appendChild(done);}
\NormalTok{    \}}
\NormalTok{  \}}
\NormalTok{  const controls = document.createElement("div");}
\NormalTok{  controls.style.cssText = "display:flex; gap:0.5rem; align{-}items:center;";}
\NormalTok{  const nextBtn = document.createElement("button");}
\NormalTok{  nextBtn.textContent = "Next step →";}
\NormalTok{  nextBtn.style.cssText = "padding:0.35rem 0.85rem; border:1px solid \#d1d5db; border{-}radius:4px; background:\#fff; cursor:pointer; font{-}size:0.9em;";}
\NormalTok{  const resetBtn = document.createElement("button");}
\NormalTok{  resetBtn.textContent = "Reset";}
\NormalTok{  resetBtn.style.cssText = "padding:0.35rem 0.75rem; border:1px solid \#e5e7eb; border{-}radius:4px; background:\#f9fafb; cursor:pointer; font{-}size:0.9em; color:\#6b7280;";}
\NormalTok{  const counter = document.createElement("span");}
\NormalTok{  counter.style.cssText = "font{-}size:0.82em; color:\#9ca3af; margin{-}left:0.25rem;";}
\NormalTok{  function updateButtons() \{}
\NormalTok{    nextBtn.disabled = current === totalSteps;}
\NormalTok{    nextBtn.style.opacity = current === totalSteps ? "0.4" : "1";}
\NormalTok{    counter.textContent = current === 0 ? "Click to reveal steps" : \textasciigrave{}Step $\{current\} of $\{totalSteps\}\textasciigrave{};}
\NormalTok{  \}}
\NormalTok{  nextBtn.onclick = () =\textgreater{} \{ if (current \textless{} totalSteps) \{ current++; renderSteps(); updateButtons(); \} \};}
\NormalTok{  resetBtn.onclick = () =\textgreater{} \{ current = 0; renderSteps(); updateButtons(); \};}
\NormalTok{  controls.appendChild(nextBtn);}
\NormalTok{  controls.appendChild(resetBtn);}
\NormalTok{  controls.appendChild(counter);}
\NormalTok{  container.appendChild(stepsDiv);}
\NormalTok{  container.appendChild(controls);}
\NormalTok{  renderSteps();}
\NormalTok{  updateButtons();}
\NormalTok{  return container;}
\NormalTok{\}}
\end{Highlighting}
\end{Shaded}

\begin{Shaded}
\begin{Highlighting}[]
\NormalTok{//| echo: false}
\NormalTok{makeStepperHTML(1, [}
\NormalTok{  \{ op: "Set height equal to zero", eq: "20 + 10t {-} 4.9t\^{}2 = 0", note: "The ball hits the ground when h = 0." \},}
\NormalTok{  \{ op: "Rearrange into standard form", eq: "{-}4.9t\^{}2 + 10t + 20 = 0", note: null \},}
\NormalTok{  \{ op: "Identify coefficients (the quadratic formula handles any sign of a)", eq: "a = {-}4.9,\textbackslash{}\textbackslash{}quad b = 10,\textbackslash{}\textbackslash{}quad c = 20", note: "No need to rearrange further — the formula works with a negative leading coefficient." \},}
\NormalTok{  \{ op: "Compute the discriminant", eq: "\textbackslash{}\textbackslash{}Delta = 10\^{}2 {-} 4({-}4.9)(20) = 100 + 392 = 492", note: "Positive discriminant means two real solutions; we take the one giving positive t." \},}
\NormalTok{  \{ op: "Apply the quadratic formula", eq: "t = \textbackslash{}\textbackslash{}frac\{{-}10 \textbackslash{}\textbackslash{}pm \textbackslash{}\textbackslash{}sqrt\{492\}\}\{2 \textbackslash{}\textbackslash{}times ({-}4.9)\} = \textbackslash{}\textbackslash{}frac\{{-}10 \textbackslash{}\textbackslash{}pm 22.18\}\{{-}9.8\}", note: null \},}
\NormalTok{  \{ op: "Take the root that gives positive t", eq: "t = \textbackslash{}\textbackslash{}frac\{{-}10 {-} 22.18\}\{{-}9.8\} = \textbackslash{}\textbackslash{}frac\{{-}32.18\}\{{-}9.8\} \textbackslash{}\textbackslash{}approx 3.28 \textbackslash{}\textbackslash{}text\{ s\}", note: "The other root gives t ≈ −1.24 s, which is before the ball was thrown." \},}
\NormalTok{  \{ op: "Check", eq: "h(3.28) = 20 + 10(3.28) {-} 4.9(3.28)\^{}2 \textbackslash{}\textbackslash{}approx 20 + 32.8 {-} 52.7 \textbackslash{}\textbackslash{}approx 0.1 \textbackslash{}\textbackslash{}approx 0 \textbackslash{}\textbackslash{}checkmark", note: null \}}
\NormalTok{])}
\end{Highlighting}
\end{Shaded}

\begin{enumerate}
\def\labelenumi{\arabic{enumi}.}
\setcounter{enumi}{1}
\tightlist
\item
  Factor \(x^2 - 7x + 12\) completely, then write down the two roots of
  \(x^2 - 7x + 12 = 0\).
\end{enumerate}

\begin{Shaded}
\begin{Highlighting}[]
\NormalTok{//| echo: false}
\NormalTok{makeStepperHTML(2, [}
\NormalTok{  \{ op: "Write the trinomial", eq: "x\^{}2 {-} 7x + 12", note: null \},}
\NormalTok{  \{ op: "Find two numbers that multiply to 12 and add to {-}7", eq: "({-}3) \textbackslash{}\textbackslash{}times ({-}4) = 12 \textbackslash{}\textbackslash{}quad\textbackslash{}\textbackslash{}text\{and\}\textbackslash{}\textbackslash{}quad ({-}3) + ({-}4) = {-}7", note: null \},}
\NormalTok{  \{ op: "Write the factored form", eq: "x\^{}2 {-} 7x + 12 = (x {-} 3)(x {-} 4)", note: null \},}
\NormalTok{  \{ op: "Set each factor to zero", eq: "x {-} 3 = 0 \textbackslash{}\textbackslash{}implies x = 3 \textbackslash{}\textbackslash{}qquad x {-} 4 = 0 \textbackslash{}\textbackslash{}implies x = 4", note: null \},}
\NormalTok{  \{ op: "Check", eq: "(3)\^{}2 {-} 7(3) + 12 = 9 {-} 21 + 12 = 0 \textbackslash{}\textbackslash{}checkmark \textbackslash{}\textbackslash{}qquad (4)\^{}2 {-} 7(4) + 12 = 16 {-} 28 + 12 = 0 \textbackslash{}\textbackslash{}checkmark", note: null \}}
\NormalTok{])}
\end{Highlighting}
\end{Shaded}

\begin{enumerate}
\def\labelenumi{\arabic{enumi}.}
\setcounter{enumi}{2}
\tightlist
\item
  A rectangular garden has an area of 48 m². Its length is 2 m more than
  its width. Find the dimensions.
\end{enumerate}

\begin{Shaded}
\begin{Highlighting}[]
\NormalTok{//| echo: false}
\NormalTok{makeStepperHTML(3, [}
\NormalTok{  \{ op: "Let width = x, then length = x + 2", eq: "x(x + 2) = 48", note: null \},}
\NormalTok{  \{ op: "Expand", eq: "x\^{}2 + 2x = 48", note: null \},}
\NormalTok{  \{ op: "Rearrange into standard form", eq: "x\^{}2 + 2x {-} 48 = 0", note: null \},}
\NormalTok{  \{ op: "Find two numbers that multiply to {-}48 and add to 2", eq: "8 \textbackslash{}\textbackslash{}times ({-}6) = {-}48 \textbackslash{}\textbackslash{}quad\textbackslash{}\textbackslash{}text\{and\}\textbackslash{}\textbackslash{}quad 8 + ({-}6) = 2", note: null \},}
\NormalTok{  \{ op: "Factor", eq: "(x + 8)(x {-} 6) = 0", note: null \},}
\NormalTok{  \{ op: "Solve — take the positive root", eq: "x = 6 \textbackslash{}\textbackslash{}text\{ (width)\}, \textbackslash{}\textbackslash{}quad x + 2 = 8 \textbackslash{}\textbackslash{}text\{ (length)\}", note: "x = {-}8 is rejected since a width cannot be negative." \},}
\NormalTok{  \{ op: "Check", eq: "6 \textbackslash{}\textbackslash{}times 8 = 48 \textbackslash{}\textbackslash{}checkmark", note: null \}}
\NormalTok{])}
\end{Highlighting}
\end{Shaded}

\begin{enumerate}
\def\labelenumi{\arabic{enumi}.}
\setcounter{enumi}{3}
\tightlist
\item
  Use the quadratic formula to solve \(2x^2 - 3x - 2 = 0\).
\end{enumerate}

\begin{Shaded}
\begin{Highlighting}[]
\NormalTok{//| echo: false}
\NormalTok{makeStepperHTML(4, [}
\NormalTok{  \{ op: "Identify coefficients", eq: "a = 2,\textbackslash{}\textbackslash{}quad b = {-}3,\textbackslash{}\textbackslash{}quad c = {-}2", note: null \},}
\NormalTok{  \{ op: "Compute the discriminant", eq: "\textbackslash{}\textbackslash{}Delta = ({-}3)\^{}2 {-} 4(2)({-}2) = 9 + 16 = 25", note: null \},}
\NormalTok{  \{ op: "Apply the quadratic formula", eq: "x = \textbackslash{}\textbackslash{}frac\{3 \textbackslash{}\textbackslash{}pm \textbackslash{}\textbackslash{}sqrt\{25\}\}\{4\} = \textbackslash{}\textbackslash{}frac\{3 \textbackslash{}\textbackslash{}pm 5\}\{4\}", note: null \},}
\NormalTok{  \{ op: "Find both solutions", eq: "x = \textbackslash{}\textbackslash{}frac\{3 + 5\}\{4\} = 2 \textbackslash{}\textbackslash{}qquad\textbackslash{}\textbackslash{}text\{or\}\textbackslash{}\textbackslash{}qquad x = \textbackslash{}\textbackslash{}frac\{3 {-} 5\}\{4\} = {-}\textbackslash{}\textbackslash{}frac\{1\}\{2\}", note: null \},}
\NormalTok{  \{ op: "Check x = 2", eq: "2(4) {-} 3(2) {-} 2 = 8 {-} 6 {-} 2 = 0 \textbackslash{}\textbackslash{}checkmark", note: null \},}
\NormalTok{  \{ op: "Check x = {-}1/2", eq: "2\textbackslash{}\textbackslash{}left(\textbackslash{}\textbackslash{}frac\{1\}\{4\}\textbackslash{}\textbackslash{}right) {-} 3\textbackslash{}\textbackslash{}left({-}\textbackslash{}\textbackslash{}frac\{1\}\{2\}\textbackslash{}\textbackslash{}right) {-} 2 = \textbackslash{}\textbackslash{}frac\{1\}\{2\} + \textbackslash{}\textbackslash{}frac\{3\}\{2\} {-} 2 = 0 \textbackslash{}\textbackslash{}checkmark", note: null \}}
\NormalTok{])}
\end{Highlighting}
\end{Shaded}

\begin{enumerate}
\def\labelenumi{\arabic{enumi}.}
\setcounter{enumi}{4}
\tightlist
\item
  A candle company finds its weekly profit in dollars is
  \(P(p) = -3p^2 + 48p - 165\), where \(p\) is the price per candle. At
  what price is profit maximised, and what is that profit?
\end{enumerate}

\begin{Shaded}
\begin{Highlighting}[]
\NormalTok{//| echo: false}
\NormalTok{makeStepperHTML(5, [}
\NormalTok{  \{ op: "Identify coefficients", eq: "a = {-}3,\textbackslash{}\textbackslash{}quad b = 48,\textbackslash{}\textbackslash{}quad c = {-}165", note: "The negative a confirms this is a downward parabola — it has a maximum." \},}
\NormalTok{  \{ op: "Find the price at the vertex", eq: "p\_v = {-}\textbackslash{}\textbackslash{}frac\{b\}\{2a\} = {-}\textbackslash{}\textbackslash{}frac\{48\}\{2({-}3)\} = {-}\textbackslash{}\textbackslash{}frac\{48\}\{{-}6\} = 8", note: null \},}
\NormalTok{  \{ op: "Substitute to find maximum profit", eq: "P(8) = {-}3(64) + 48(8) {-} 165 = {-}192 + 384 {-} 165 = 27", note: null \},}
\NormalTok{  \{ op: "State the answer", eq: "\textbackslash{}\textbackslash{}text\{Maximum profit: \} \textbackslash{}\textbackslash{}$27 \textbackslash{}\textbackslash{}text\{ per week at a price of \} \textbackslash{}\textbackslash{}$8", note: null \},}
\NormalTok{  \{ op: "Check", eq: "P(7) = {-}3(49) + 48(7) {-} 165 = {-}147 + 336 {-} 165 = 24 \textless{} 27 \textbackslash{}\textbackslash{}checkmark", note: "Checking a neighbouring value confirms 8 is the peak." \}}
\NormalTok{])}
\end{Highlighting}
\end{Shaded}

\begin{enumerate}
\def\labelenumi{\arabic{enumi}.}
\setcounter{enumi}{5}
\tightlist
\item
  Factor \(4x^2 - 25\) completely.
\end{enumerate}

\begin{Shaded}
\begin{Highlighting}[]
\NormalTok{//| echo: false}
\NormalTok{makeStepperHTML(6, [}
\NormalTok{  \{ op: "Recognise the structure", eq: "4x\^{}2 {-} 25 = (2x)\^{}2 {-} 5\^{}2", note: "Both terms are perfect squares separated by a minus sign — use the difference of squares identity." \},}
\NormalTok{  \{ op: "Apply a² {-} b² = (a + b)(a {-} b)", eq: "4x\^{}2 {-} 25 = (2x + 5)(2x {-} 5)", note: null \},}
\NormalTok{  \{ op: "Check by expanding", eq: "(2x + 5)(2x {-} 5) = 4x\^{}2 {-} 10x + 10x {-} 25 = 4x\^{}2 {-} 25 \textbackslash{}\textbackslash{}checkmark", note: null \}}
\NormalTok{])}
\end{Highlighting}
\end{Shaded}

\chapter{Functions and relations}\label{functions-and-relations}

Given this, produce that

\hfill\break

Press B4 on a vending machine. One specific drink comes out --- not a
random choice, not two drinks. One input, one output, every time.

Type 100 into a Celsius-to-Fahrenheit converter. You get 212. Type 100
again tomorrow and you still get 212. The relationship doesn't change.
The output is determined entirely by the input.

A tip calculator takes one number --- the bill total --- and produces
one number: the suggested tip. It doesn't need to know anything else.
One input, one output, every time.

These are all the same structure: a rule that takes an input and
produces exactly one output. Mathematicians call that a function. This
chapter is about what that means precisely --- and why the restriction
matters.

\section{What the notation is
saying}\label{what-the-notation-is-saying-13}

Here is a function:

\[f(x) = 2x + 3\]

Read it aloud: ``f of x equals 2x plus 3.'' The letter \(f\) is the name
of the function --- you could call it \(g\) or \(T\) or anything, but
\(f\) is the default. The \(x\) inside the parentheses is the input. The
right-hand side, \(2x + 3\), is the rule: whatever you put in, double it
and add 3.

The notation \(f(x)\) does not mean \(f\) multiplied by \(x\). It means
``the output of \(f\) when the input is \(x\).'' That's an easy thing to
misread the first time you see it. Say it in words every time until it's
automatic.

You may also see function notation written as \(f: A \to B\), which says
``\(f\) maps from set \(A\) to set \(B\)'' --- every element of \(A\)
gets sent to exactly one element of \(B\). For example, the temperature
converter maps the set of Celsius values to the set of Fahrenheit
values.

One thing the definition requires: \textbf{every input produces exactly
one output.} That's what separates a function from a general relation. A
relation might pair one input with two or more outputs --- but then it's
not a function. The vending machine analogy holds: if pressing B4
sometimes gave you a cola and sometimes gave you nothing, or randomly
gave you two things, you wouldn't call it a reliable machine. A function
is a reliable machine.

\begin{Shaded}
\begin{Highlighting}[]
\NormalTok{//| echo: false}
\NormalTok{viewof fmInput = Inputs.range([{-}5, 5], \{}
\NormalTok{  value: 2,}
\NormalTok{  step: 0.5,}
\NormalTok{  label: "Input x"}
\NormalTok{\})}
\end{Highlighting}
\end{Shaded}

\begin{Shaded}
\begin{Highlighting}[]
\NormalTok{//| echo: false}
\NormalTok{viewof fmChoice = Inputs.select(["2x + 1", "x²", "|x|", "√(x+5)"], \{}
\NormalTok{  label: "Function f(x)"}
\NormalTok{\})}
\end{Highlighting}
\end{Shaded}

\begin{Shaded}
\begin{Highlighting}[]
\NormalTok{//| echo: false}
\NormalTok{\{}
\NormalTok{  const x = fmInput;}
\NormalTok{  const choice = fmChoice;}

\NormalTok{  let formula, result, resultStr;}

\NormalTok{  if (choice === "2x + 1") \{}
\NormalTok{    formula = "f(x) = 2x + 1";}
\NormalTok{    result = 2 * x + 1;}
\NormalTok{    resultStr = result.toFixed(3);}
\NormalTok{  \} else if (choice === "x²") \{}
\NormalTok{    formula = "f(x) = x²";}
\NormalTok{    result = x * x;}
\NormalTok{    resultStr = result.toFixed(3);}
\NormalTok{  \} else if (choice === "|x|") \{}
\NormalTok{    formula = "f(x) = |x|";}
\NormalTok{    result = Math.abs(x);}
\NormalTok{    resultStr = result.toFixed(3);}
\NormalTok{  \} else \{}
\NormalTok{    formula = "f(x) = √(x + 5)";}
\NormalTok{    if (x + 5 \textless{} 0) \{}
\NormalTok{      resultStr = "undefined (√ of negative)";}
\NormalTok{      result = null;}
\NormalTok{    \} else \{}
\NormalTok{      result = Math.sqrt(x + 5);}
\NormalTok{      resultStr = result.toFixed(3);}
\NormalTok{    \}}
\NormalTok{  \}}

\NormalTok{  const svgW = 460, svgH = 160;}
\NormalTok{  const svg = d3.create("svg")}
\NormalTok{    .attr("width", svgW).attr("height", svgH)}
\NormalTok{    .style("font{-}family", "inherit");}

\NormalTok{  // Machine box}
\NormalTok{  svg.append("rect")}
\NormalTok{    .attr("x", 160).attr("y", 40)}
\NormalTok{    .attr("width", 140).attr("height", 70)}
\NormalTok{    .attr("fill", "\#eff6ff").attr("stroke", "\#3b82f6").attr("stroke{-}width", 2)}
\NormalTok{    .attr("rx", 8);}

\NormalTok{  svg.append("text")}
\NormalTok{    .attr("x", 230).attr("y", 72)}
\NormalTok{    .attr("text{-}anchor", "middle").attr("font{-}size", "13px").attr("font{-}weight", "600").attr("fill", "\#1d4ed8")}
\NormalTok{    .text(formula);}

\NormalTok{  svg.append("text")}
\NormalTok{    .attr("x", 230).attr("y", 94)}
\NormalTok{    .attr("text{-}anchor", "middle").attr("font{-}size", "11px").attr("fill", "\#6b7280")}
\NormalTok{    .text("function machine");}

\NormalTok{  // Input arrow}
\NormalTok{  svg.append("line")}
\NormalTok{    .attr("x1", 60).attr("y1", 75)}
\NormalTok{    .attr("x2", 155).attr("y2", 75)}
\NormalTok{    .attr("stroke", "\#374151").attr("stroke{-}width", 2)}
\NormalTok{    .attr("marker{-}end", "url(\#arrowhead)");}

\NormalTok{  svg.append("text")}
\NormalTok{    .attr("x", 90).attr("y", 65)}
\NormalTok{    .attr("text{-}anchor", "middle").attr("font{-}size", "13px").attr("fill", "\#374151")}
\NormalTok{    .text(\textasciigrave{}x = $\{x\}\textasciigrave{});}

\NormalTok{  // Output arrow}
\NormalTok{  svg.append("line")}
\NormalTok{    .attr("x1", 305).attr("y1", 75)}
\NormalTok{    .attr("x2", 390).attr("y2", 75)}
\NormalTok{    .attr("stroke", "\#374151").attr("stroke{-}width", 2)}
\NormalTok{    .attr("marker{-}end", "url(\#arrowhead)");}

\NormalTok{  svg.append("text")}
\NormalTok{    .attr("x", 350).attr("y", 65)}
\NormalTok{    .attr("text{-}anchor", "middle")}
\NormalTok{    .attr("font{-}size", "13px")}
\NormalTok{    .attr("fill", result === null ? "\#dc2626" : "\#059669")}
\NormalTok{    .text(\textasciigrave{}f(x) = $\{resultStr\}\textasciigrave{});}

\NormalTok{  // Arrowhead marker}
\NormalTok{  const defs = svg.append("defs");}
\NormalTok{  defs.append("marker")}
\NormalTok{    .attr("id", "arrowhead")}
\NormalTok{    .attr("markerWidth", 8).attr("markerHeight", 6)}
\NormalTok{    .attr("refX", 8).attr("refY", 3)}
\NormalTok{    .attr("orient", "auto")}
\NormalTok{    .append("polygon")}
\NormalTok{    .attr("points", "0 0, 8 3, 0 6")}
\NormalTok{    .attr("fill", "\#374151");}

\NormalTok{  // Pair display}
\NormalTok{  svg.append("text")}
\NormalTok{    .attr("x", svgW / 2).attr("y", 140)}
\NormalTok{    .attr("text{-}anchor", "middle").attr("font{-}size", "13px").attr("fill", "\#374151")}
\NormalTok{    .text(result === null}
\NormalTok{      ? \textasciigrave{}($\{x\}, undefined)\textasciigrave{}}
\NormalTok{      : \textasciigrave{}Input: $\{x\}   →   Output: $\{resultStr\}   ⟹   ($\{x\}, $\{resultStr\})\textasciigrave{});}

\NormalTok{  return svg.node();}
\NormalTok{\}}
\end{Highlighting}
\end{Shaded}

\section{The method}\label{the-method-12}

\subsection{Part 1: Evaluating a
function}\label{part-1-evaluating-a-function}

To evaluate \(f(x)\) at a specific input, replace every \(x\) in the
rule with that value and simplify.

If \(f(x) = 2x + 3\), then \(f(5)\) means: put 5 in place of \(x\).

\[f(5) = 2(5) + 3 = 10 + 3 = 13\]

The input was 5, the output is 13.

If the input is negative or a fraction, the same substitution applies:

\[f(-4) = 2(-4) + 3 = -8 + 3 = -5\]

The parentheses around \(-4\) are there to keep the sign attached to the
number. Leave them out and \(2-4+3\) is ambiguous. Keep them in and it's
clear.

\subsection{Part 2: Domain and range}\label{part-2-domain-and-range}

The \textbf{domain} of a function is the set of all valid inputs --- the
values of \(x\) you're allowed to use.

The \textbf{range} is the set of all possible outputs --- every value
\(f(x)\) can produce when you run through the domain.

For \(f(x) = 2x + 3\), any real number is a valid input and any real
number is a possible output. The domain and range are both all real
numbers, written \((-\infty, \infty)\) in interval notation.

Domain restrictions come in two situations:

\textbf{Division by zero is not allowed.} For \(g(x) = \dfrac{1}{x-2}\),
the input \(x = 2\) makes the denominator zero, which is undefined. The
domain is all real numbers except 2, written
\((-\infty, 2) \cup (2, \infty)\). The \(\cup\) symbol means ``or'' ---
the domain includes numbers from the first interval or the second.

\textbf{Square roots of negative numbers don't produce real outputs.}
For \(h(x) = \sqrt{x - 1}\), the expression under the root must be
non-negative: \(x - 1 \geq 0\), so \(x \geq 1\). The domain is
\([1, \infty)\). Why does the square root require this? No real number
multiplied by itself gives a negative result --- so the square root of a
negative number doesn't land on the number line.

Interval notation in brief: - \([a, b]\) --- includes both endpoints
(\(a \leq x \leq b\)) - \((a, b)\) --- excludes both endpoints
(\(a < x < b\)) - \([a, \infty)\) --- from \(a\) upward, including \(a\)
- \((-\infty, b)\) --- everything below \(b\), not including \(b\)

The domain can also be restricted by context --- a function describing
something you can only sell in whole numbers, or a physical measurement
that can't be negative. That kind of restriction comes from the
situation, not the algebra.

\subsection{Part 3: Composition}\label{part-3-composition}

You can chain functions together. If \(g\) converts one thing and \(f\)
converts another, you can run an input through \(g\) first, then feed
\(g\)'s output into \(f\).

The notation for this is \(f(g(x))\), read as ``f of g of x.'' You may
also see it written \(f \circ g\), read as ``f composed with g.'' Both
mean the same thing: apply \(g\) first, then apply \(f\) to the result.

If \(g(x) = x + 3\) and \(f(x) = 2x + 1\):

\[f(g(x)) = f(x + 3) = 2(x + 3) + 1 = 2x + 6 + 1 = 2x + 7\]

The order matters. \(f(g(x))\) and \(g(f(x))\) are generally different.

\(g(f(x)) = g(2x + 1) = (2x + 1) + 3 = 2x + 4\)

That is a different function. Get the order right by reading from the
inside out: in \(f(g(x))\), \(g\) acts first.

\begin{Shaded}
\begin{Highlighting}[]
\NormalTok{//| echo: false}
\NormalTok{viewof compF = Inputs.select(["x + 3", "2x", "x²"], \{ label: "f(x)" \})}
\end{Highlighting}
\end{Shaded}

\begin{Shaded}
\begin{Highlighting}[]
\NormalTok{//| echo: false}
\NormalTok{viewof compG = Inputs.select(["x − 1", "3x", "x + 2"], \{ label: "g(x)" \})}
\end{Highlighting}
\end{Shaded}

\begin{Shaded}
\begin{Highlighting}[]
\NormalTok{//| echo: false}
\NormalTok{viewof compX = Inputs.range([{-}4, 4], \{ value: 1, step: 0.5, label: "Input x" \})}
\end{Highlighting}
\end{Shaded}

\begin{Shaded}
\begin{Highlighting}[]
\NormalTok{//| echo: false}
\NormalTok{\{}
\NormalTok{  const x = compX;}
\NormalTok{  const fName = compF;}
\NormalTok{  const gName = compG;}

\NormalTok{  // Evaluate f and g numerically}
\NormalTok{  function evalF(v) \{}
\NormalTok{    if (fName === "x + 3") return v + 3;}
\NormalTok{    if (fName === "2x") return 2 * v;}
\NormalTok{    if (fName === "x²") return v * v;}
\NormalTok{  \}}
\NormalTok{  function evalG(v) \{}
\NormalTok{    if (gName === "x − 1") return v {-} 1;}
\NormalTok{    if (gName === "3x") return 3 * v;}
\NormalTok{    if (gName === "x + 2") return v + 2;}
\NormalTok{  \}}

\NormalTok{  const gx = evalG(x);}
\NormalTok{  const fgx = evalF(gx);}
\NormalTok{  const fx = evalF(x);}
\NormalTok{  const gfx = evalG(fx);}

\NormalTok{  // Symbolic composed forms for the 9 combinations}
\NormalTok{  const symbolic = \{}
\NormalTok{    "x + 3|x − 1": \{ fg: "f(g(x)) = (x−1)+3 = x+2", gf: "g(f(x)) = (x+3)−1 = x+2" \},}
\NormalTok{    "x + 3|3x":    \{ fg: "f(g(x)) = 3x+3", gf: "g(f(x)) = 3(x+3) = 3x+9" \},}
\NormalTok{    "x + 3|x + 2": \{ fg: "f(g(x)) = (x+2)+3 = x+5", gf: "g(f(x)) = (x+3)+2 = x+5" \},}
\NormalTok{    "2x|x − 1":    \{ fg: "f(g(x)) = 2(x−1) = 2x−2", gf: "g(f(x)) = 2x−1" \},}
\NormalTok{    "2x|3x":       \{ fg: "f(g(x)) = 2(3x) = 6x", gf: "g(f(x)) = 3(2x) = 6x" \},}
\NormalTok{    "2x|x + 2":    \{ fg: "f(g(x)) = 2(x+2) = 2x+4", gf: "g(f(x)) = 2x+2" \},}
\NormalTok{    "x²|x − 1":   \{ fg: "f(g(x)) = (x−1)²", gf: "g(f(x)) = x²−1" \},}
\NormalTok{    "x²|3x":       \{ fg: "f(g(x)) = (3x)² = 9x²", gf: "g(f(x)) = 3x²" \},}
\NormalTok{    "x²|x + 2":    \{ fg: "f(g(x)) = (x+2)²", gf: "g(f(x)) = x²+2" \},}
\NormalTok{  \};}
\NormalTok{  const key = \textasciigrave{}$\{fName\}|$\{gName\}\textasciigrave{};}
\NormalTok{  const sym = symbolic[key] || \{ fg: "f(g(x))", gf: "g(f(x))" \};}

\NormalTok{  const svgW = 480, svgH = 210;}
\NormalTok{  const svg = d3.create("svg")}
\NormalTok{    .attr("width", svgW).attr("height", svgH)}
\NormalTok{    .style("font{-}family", "inherit");}

\NormalTok{  // Panel background}
\NormalTok{  svg.append("rect").attr("x", 10).attr("y", 10)}
\NormalTok{    .attr("width", svgW {-} 20).attr("height", svgH {-} 20)}
\NormalTok{    .attr("fill", "\#f8fafc").attr("stroke", "\#e2e8f0").attr("rx", 8);}

\NormalTok{  // Heading}
\NormalTok{  svg.append("text").attr("x", 24).attr("y", 36)}
\NormalTok{    .attr("font{-}size", "13px").attr("font{-}weight", "600").attr("fill", "\#111827")}
\NormalTok{    .text(\textasciigrave{}f(x) = $\{fName\}   g(x) = $\{gName\}   x = $\{x\}\textasciigrave{});}

\NormalTok{  // g(x)}
\NormalTok{  svg.append("text").attr("x", 24).attr("y", 64)}
\NormalTok{    .attr("font{-}size", "13px").attr("fill", "\#374151")}
\NormalTok{    .text(\textasciigrave{}g(x) = $\{gx.toFixed(3)\}\textasciigrave{});}

\NormalTok{  // f(g(x))}
\NormalTok{  svg.append("text").attr("x", 24).attr("y", 88)}
\NormalTok{    .attr("font{-}size", "13px").attr("fill", "\#2563eb")}
\NormalTok{    .text(\textasciigrave{}f(g(x)) = f($\{gx.toFixed(2)\}) = $\{fgx.toFixed(3)\}\textasciigrave{});}

\NormalTok{  svg.append("text").attr("x", 24).attr("y", 108)}
\NormalTok{    .attr("font{-}size", "12px").attr("fill", "\#6b7280")}
\NormalTok{    .text(sym.fg);}

\NormalTok{  // Divider}
\NormalTok{  svg.append("line").attr("x1", 24).attr("y1", 120).attr("x2", svgW {-} 24).attr("y2", 120)}
\NormalTok{    .attr("stroke", "\#e5e7eb").attr("stroke{-}width", 1);}

\NormalTok{  // f(x)}
\NormalTok{  svg.append("text").attr("x", 24).attr("y", 142)}
\NormalTok{    .attr("font{-}size", "13px").attr("fill", "\#374151")}
\NormalTok{    .text(\textasciigrave{}f(x) = $\{fx.toFixed(3)\}\textasciigrave{});}

\NormalTok{  // g(f(x))}
\NormalTok{  svg.append("text").attr("x", 24).attr("y", 166)}
\NormalTok{    .attr("font{-}size", "13px").attr("fill", "\#7c3aed")}
\NormalTok{    .text(\textasciigrave{}g(f(x)) = g($\{fx.toFixed(2)\}) = $\{gfx.toFixed(3)\}\textasciigrave{});}

\NormalTok{  svg.append("text").attr("x", 24).attr("y", 186)}
\NormalTok{    .attr("font{-}size", "12px").attr("fill", "\#6b7280")}
\NormalTok{    .text(sym.gf);}

\NormalTok{  return svg.node();}
\NormalTok{\}}
\end{Highlighting}
\end{Shaded}

\subsection{Part 4: Inverses}\label{part-4-inverses}

The inverse of a function undoes it. If \(f\) takes 5 to 13, then
\(f^{-1}\) --- read as ``f inverse'' --- takes 13 back to 5.

The notation \(f^{-1}(x)\) does not mean \(\dfrac{1}{f(x)}\). That's
another thing that looks like it should mean one thing but means
another. The superscript \(-1\) on a function name means ``the inverse
function,'' not ``one over.''

Not every function has an inverse. For an inverse to exist, each output
must have come from exactly one input --- otherwise, when you try to go
backwards, you don't know which input to return to. A function with this
property is called \textbf{one-to-one}: different inputs produce
different outputs.

To find the inverse algebraically: write \(y = f(x)\), swap \(x\) and
\(y\), then solve for \(y\).

For \(f(x) = 2x + 3\):

Write as: \(y = 2x + 3\)

Swap \(x\) and \(y\): \(x = 2y + 3\)

Solve for \(y\):

\[x - 3 = 2y\] \[y = \frac{x - 3}{2}\]

So \(f^{-1}(x) = \dfrac{x - 3}{2}\).

Check: \(f(5) = 13\), and
\(f^{-1}(13) = \dfrac{13 - 3}{2} = \dfrac{10}{2} = 5\). Yes.

\begin{tcolorbox}[enhanced jigsaw, opacityback=0, colback=white, breakable, colframe=quarto-callout-note-color-frame, arc=.35mm, left=2mm, toprule=.15mm, leftrule=.75mm, bottomrule=.15mm, rightrule=.15mm]
\begin{minipage}[t]{5.5mm}
\textcolor{quarto-callout-note-color}{\faInfo}
\end{minipage}%
\begin{minipage}[t]{\textwidth - 5.5mm}

\textbf{Why this works: the vertical line test}

Draw the graph of a relation. Draw a vertical line anywhere on the
graph.

If any vertical line crosses the graph more than once, the relation is
not a function --- because that input \(x\) (the horizontal position of
your line) maps to two or more outputs (the two crossing points).

If every vertical line crosses the graph at most once, the relation is a
function. The test is visual and immediate: one crossing means one
output for that input. That is the requirement.

\end{minipage}%
\end{tcolorbox}

\section{Worked examples}\label{worked-examples-12}

\begin{center}\rule{0.5\linewidth}{0.5pt}\end{center}

\textbf{Example 1 (Science) --- Temperature conversion and its inverse.}

The formula for converting Celsius to Fahrenheit is:

\[F(c) = \frac{9}{5}c + 32\]

Here \(c\) is the input (temperature in °C) and \(F(c)\) is the output
(temperature in °F).

Evaluate at the boiling point of water, \(c = 100\):

\[F(100) = \frac{9}{5}(100) + 32 = 180 + 32 = 212\]

Water boils at 212°F. Check: widely known to be correct.

Now find the inverse --- the converter run backwards, converting °F to
°C.

Write \(y = \dfrac{9}{5}c + 32\), swap variables:
\(c = \dfrac{9}{5}y + 32\).

Solve for \(y\):

\[c - 32 = \frac{9}{5}y\] \[y = \frac{5}{9}(c - 32)\]

So \(F^{-1}(T) = \dfrac{5}{9}(T - 32)\), where \(T\) is a temperature in
°F.

Verify:
\(F^{-1}(212) = \dfrac{5}{9}(212 - 32) = \dfrac{5}{9}(180) = 100\). Back
to 100°C. Correct.

\begin{center}\rule{0.5\linewidth}{0.5pt}\end{center}

\textbf{Example 2 (Computing) --- Songs downloaded and storage used.}

Every song downloaded takes up space on a device. For a particular music
format, the storage used \(S\) in megabytes as a function of songs
downloaded \(n\) is:

\[S(n) = 4.5n + 20\]

The 20 MB is the app's base footprint; each song adds 4.5 MB.

At \(n = 10\) songs: \[S(10) = 4.5(10) + 20 = 45 + 20 = 65 \text{ MB}\]

At \(n = 50\) songs:
\[S(50) = 4.5(50) + 20 = 225 + 20 = 245 \text{ MB}\]

At \(n = 100\) songs:
\[S(100) = 4.5(100) + 20 = 450 + 20 = 470 \text{ MB}\]

The domain has a physical restriction: \(n \geq 0\) and \(n\) must be a
whole number. The maximum \(n\) depends on the device's storage limit
--- the model stops being useful beyond that point.

\begin{center}\rule{0.5\linewidth}{0.5pt}\end{center}

\textbf{Example 3 (Computing) --- Function composition in a data
pipeline.}

A data pipeline applies two transformations to a raw temperature reading
from a sensor.

Step 1: the sensor outputs a raw voltage \(v\). The first function
converts voltage to Celsius:

\[g(v) = 25v - 10\]

Step 2: the second function converts Celsius to Fahrenheit (from Example
1):

\[f(c) = \frac{9}{5}c + 32\]

The pipeline composes them: given a raw voltage, produce a Fahrenheit
reading.

\[f(g(v)) = f(25v - 10) = \frac{9}{5}(25v - 10) + 32\]
\[= 45v - 18 + 32 = 45v + 14\]

So the combined function is \(h(v) = 45v + 14\).

Test it: at \(v = 1.2\) volts:

\[g(1.2) = 25(1.2) - 10 = 30 - 10 = 20 \text{°C}\]
\[f(20) = \frac{9}{5}(20) + 32 = 36 + 32 = 68 \text{°F}\]
\[h(1.2) = 45(1.2) + 14 = 54 + 14 = 68 \text{°F}\]

Both paths give the same answer. The composed function works.

\begin{center}\rule{0.5\linewidth}{0.5pt}\end{center}

\textbf{Example 4 (Business) --- Revenue function with restricted
domain.}

A market stall sells handmade notebooks. The revenue \(R\) in dollars as
a function of notebooks sold \(n\) is:

\[R(n) = 12n\]

But the domain is restricted: \(n\) must be a non-negative integer. You
can't sell −3 notebooks or 4.7 notebooks.

Revenue at \(n = 0\): \(R(0) = 0\). Nothing sold, nothing earned.

Revenue at \(n = 25\): \(R(25) = 12(25) = 300\).

Revenue at \(n = 50\): \(R(50) = 12(50) = 600\).

The range is \(\{0, 12, 24, 36, \ldots\}\) --- multiples of 12, starting
at 0. This is not a continuous range because the domain is discrete.

If the stall owner wants to know how many notebooks to sell to reach a
revenue target of \$180, that's the inverse:

\[R^{-1}(180) = \frac{180}{12} = 15 \text{ notebooks}\]

\section{Where this goes}\label{where-this-goes-13}

Exponents and logarithms (Volume 4, Chapter 1) are specific function
families with their own shapes and properties. Once you know what a
function is and how to evaluate one, invert one, and compose two, you
can work with any family --- exponential, logarithmic, trigonometric ---
using the same framework from this chapter. The function concept doesn't
change; only the rule inside changes.

Calculus (Volume 5) asks a new question about functions: not ``what is
the output?'' but ``how fast is the output changing?'' That question ---
the derivative --- only makes sense if you have a function to begin
with. Limits (Volume 5, Chapter 1) ask what a function's output
approaches as the input gets close to a particular value, which requires
the same precise notion of function you now have. Everything in analysis
builds on this chapter.

\begin{tcolorbox}[enhanced jigsaw, opacityback=0, colback=white, breakable, colframe=quarto-callout-tip-color-frame, arc=.35mm, left=2mm, toprule=.15mm, leftrule=.75mm, bottomrule=.15mm, rightrule=.15mm]
\begin{minipage}[t]{5.5mm}
\textcolor{quarto-callout-tip-color}{\faLightbulb}
\end{minipage}%
\begin{minipage}[t]{\textwidth - 5.5mm}

\textbf{Where this shows up}

\begin{itemize}
\tightlist
\item
  A structural engineer's load-deflection formula is a function: given
  load, compute displacement.
\item
  Every physics formula (\(F = ma\), \(E = mc^2\), \(v = u + at\)) is a
  function expressing one quantity in terms of others.
\item
  A pricing algorithm that takes customer behaviour and outputs a price
  is a function --- ideally a deterministic one.
\item
  A compiler transforms source code (input) into machine code (output):
  a function from programs to programs.
\item
  A chemist using Beer--Lambert law computes absorbance as a function of
  concentration. The inverse gives concentration from a measured
  absorbance.
\end{itemize}

\end{minipage}%
\end{tcolorbox}

\section{Exercises}\label{exercises-13}

These are puzzles. Each one has a clean answer. The interesting part is
reading the situation carefully before reaching for algebra.

\begin{center}\rule{0.5\linewidth}{0.5pt}\end{center}

\textbf{1.} Let \(f(x) = 3x^2 - 2x + 1\).

Evaluate: (a) \(f(0)\), (b) \(f(2)\), (c) \(f(-3)\).

\begin{Shaded}
\begin{Highlighting}[]
\NormalTok{//| echo: false}
\NormalTok{katex = require("katex@0.16")}
\end{Highlighting}
\end{Shaded}

\begin{Shaded}
\begin{Highlighting}[]
\NormalTok{//| echo: false}
\NormalTok{function makeStepperHTML(exerciseNum, steps) \{}
\NormalTok{  let current = 0;}
\NormalTok{  const totalSteps = steps.length;}
\NormalTok{  const container = document.createElement("div");}
\NormalTok{  container.style.cssText = "border:1px solid \#e5e7eb; border{-}radius:8px; padding:1rem 1.25rem; margin:0.75rem 0; font{-}family:inherit;";}
\NormalTok{  const stepsDiv = document.createElement("div");}
\NormalTok{  stepsDiv.style.marginBottom = "0.75rem";}
\NormalTok{  function renderSteps() \{}
\NormalTok{    stepsDiv.innerHTML = "";}
\NormalTok{    for (let i = 0; i \textless{} current; i++) \{}
\NormalTok{      const s = steps[i];}
\NormalTok{      const row = document.createElement("div");}
\NormalTok{      row.style.cssText = "display:grid; grid{-}template{-}columns:200px 1fr; gap:0.35rem 1rem; align{-}items:baseline; padding:0.35rem 0; border{-}top:1px solid \#f3f4f6;";}
\NormalTok{      const opCell = document.createElement("span");}
\NormalTok{      opCell.style.cssText = "font{-}size:0.85em; color:\#6b7280; font{-}style:italic;";}
\NormalTok{      opCell.textContent = s.op;}
\NormalTok{      const eqCell = document.createElement("span");}
\NormalTok{      eqCell.innerHTML = katex.renderToString(s.eq, \{ throwOnError: false, displayMode: false \});}
\NormalTok{      row.appendChild(opCell);}
\NormalTok{      row.appendChild(eqCell);}
\NormalTok{      if (s.note) \{}
\NormalTok{        const noteRow = document.createElement("div");}
\NormalTok{        noteRow.style.cssText = "grid{-}column:2; font{-}size:0.82em; color:\#6b7280; padding{-}bottom:0.2rem;";}
\NormalTok{        noteRow.textContent = s.note;}
\NormalTok{        row.appendChild(noteRow);}
\NormalTok{      \}}
\NormalTok{      stepsDiv.appendChild(row);}
\NormalTok{    \}}
\NormalTok{    if (current === totalSteps) \{}
\NormalTok{      const done = document.createElement("div");}
\NormalTok{      done.style.cssText = "margin{-}top:0.5rem; font{-}size:0.9em; color:\#059669; font{-}weight:500;";}
\NormalTok{      done.textContent = "✓ Solution complete";}
\NormalTok{      stepsDiv.appendChild(done);}
\NormalTok{    \}}
\NormalTok{  \}}
\NormalTok{  const controls = document.createElement("div");}
\NormalTok{  controls.style.cssText = "display:flex; gap:0.5rem; align{-}items:center;";}
\NormalTok{  const nextBtn = document.createElement("button");}
\NormalTok{  nextBtn.textContent = "Next step →";}
\NormalTok{  nextBtn.style.cssText = "padding:0.35rem 0.85rem; border:1px solid \#d1d5db; border{-}radius:4px; background:\#fff; cursor:pointer; font{-}size:0.9em;";}
\NormalTok{  const resetBtn = document.createElement("button");}
\NormalTok{  resetBtn.textContent = "Reset";}
\NormalTok{  resetBtn.style.cssText = "padding:0.35rem 0.75rem; border:1px solid \#e5e7eb; border{-}radius:4px; background:\#f9fafb; cursor:pointer; font{-}size:0.9em; color:\#6b7280;";}
\NormalTok{  const counter = document.createElement("span");}
\NormalTok{  counter.style.cssText = "font{-}size:0.82em; color:\#9ca3af; margin{-}left:0.25rem;";}
\NormalTok{  function updateButtons() \{}
\NormalTok{    nextBtn.disabled = current === totalSteps;}
\NormalTok{    nextBtn.style.opacity = current === totalSteps ? "0.4" : "1";}
\NormalTok{    counter.textContent = current === 0 ? "Click to reveal steps" : \textasciigrave{}Step $\{current\} of $\{totalSteps\}\textasciigrave{};}
\NormalTok{  \}}
\NormalTok{  nextBtn.onclick = () =\textgreater{} \{ if (current \textless{} totalSteps) \{ current++; renderSteps(); updateButtons(); \} \};}
\NormalTok{  resetBtn.onclick = () =\textgreater{} \{ current = 0; renderSteps(); updateButtons(); \};}
\NormalTok{  controls.appendChild(nextBtn);}
\NormalTok{  controls.appendChild(resetBtn);}
\NormalTok{  controls.appendChild(counter);}
\NormalTok{  container.appendChild(stepsDiv);}
\NormalTok{  container.appendChild(controls);}
\NormalTok{  renderSteps();}
\NormalTok{  updateButtons();}
\NormalTok{  return container;}
\NormalTok{\}}
\end{Highlighting}
\end{Shaded}

\begin{Shaded}
\begin{Highlighting}[]
\NormalTok{//| echo: false}
\NormalTok{makeStepperHTML(1, [}
\NormalTok{  \{ op: "Write the function", eq: "f(x) = 3x\^{}2 {-} 2x + 1", note: null \},}
\NormalTok{  \{ op: "Evaluate f(0): substitute x = 0", eq: "f(0) = 3(0)\^{}2 {-} 2(0) + 1 = 0 {-} 0 + 1 = 1", note: null \},}
\NormalTok{  \{ op: "Evaluate f(2): substitute x = 2", eq: "f(2) = 3(2)\^{}2 {-} 2(2) + 1 = 12 {-} 4 + 1 = 9", note: null \},}
\NormalTok{  \{ op: "Evaluate f({-}3): substitute x = {-}3", eq: "f({-}3) = 3({-}3)\^{}2 {-} 2({-}3) + 1 = 27 + 6 + 1 = 34", note: "Bracket the negative number before squaring: ({-}3)\^{}2 = 9, not {-}9." \},}
\NormalTok{  \{ op: "Check by re{-}substituting f(2)", eq: "3(4) {-} 2(2) + 1 = 12 {-} 4 + 1 = 9 \textbackslash{}\textbackslash{}checkmark", note: null \}}
\NormalTok{])}
\end{Highlighting}
\end{Shaded}

\begin{center}\rule{0.5\linewidth}{0.5pt}\end{center}

\textbf{2.} State the domain of each function, using interval notation.

\begin{enumerate}
\def\labelenumi{(\alph{enumi})}
\item
  \(g(x) = \sqrt{x + 4}\)
\item
  \(h(x) = \dfrac{x}{x^2 - 9}\)
\item
  \(k(x) = x^2 + 5\)
\end{enumerate}

\begin{Shaded}
\begin{Highlighting}[]
\NormalTok{//| echo: false}
\NormalTok{makeStepperHTML(2, [}
\NormalTok{  \{ op: "For g(x): the expression under the root must be non{-}negative", eq: "x + 4 \textbackslash{}\textbackslash{}geq 0", note: null \},}
\NormalTok{  \{ op: "Solve the inequality", eq: "x \textbackslash{}\textbackslash{}geq {-}4", note: null \},}
\NormalTok{  \{ op: "Write domain of g in interval notation", eq: "[{-}4, \textbackslash{}\textbackslash{}infty)", note: null \},}
\NormalTok{  \{ op: "For h(x): the denominator cannot be zero", eq: "x\^{}2 {-} 9 \textbackslash{}\textbackslash{}neq 0 \textbackslash{}\textbackslash{}implies x \textbackslash{}\textbackslash{}neq 3 \textbackslash{}\textbackslash{}text\{ and \} x \textbackslash{}\textbackslash{}neq {-}3", note: "Factor as (x{-}3)(x+3) to find the excluded values." \},}
\NormalTok{  \{ op: "Write domain of h in interval notation", eq: "({-}\textbackslash{}\textbackslash{}infty, {-}3) \textbackslash{}\textbackslash{}cup ({-}3, 3) \textbackslash{}\textbackslash{}cup (3, \textbackslash{}\textbackslash{}infty)", note: null \},}
\NormalTok{  \{ op: "For k(x): any real input is valid (squaring never causes division by zero or square roots of negatives)", eq: "({-}\textbackslash{}\textbackslash{}infty, \textbackslash{}\textbackslash{}infty)", note: null \},}
\NormalTok{  \{ op: "Check g at x = {-}4 (boundary)", eq: "g({-}4) = \textbackslash{}\textbackslash{}sqrt\{{-}4 + 4\} = \textbackslash{}\textbackslash{}sqrt\{0\} = 0 \textbackslash{}\textbackslash{}checkmark", note: null \}}
\NormalTok{])}
\end{Highlighting}
\end{Shaded}

\begin{center}\rule{0.5\linewidth}{0.5pt}\end{center}

\textbf{3.} Let \(f(x) = x^2\) and \(g(x) = x - 5\).

\begin{enumerate}
\def\labelenumi{(\alph{enumi})}
\item
  Find \(f(g(x))\) and simplify.
\item
  Find \(g(f(x))\) and simplify.
\item
  Evaluate both \(f(g(x))\) and \(g(f(x))\) at \(x = 3\). Are they equal
  or different? Then try \(x = 2\) --- are they still the same?
\end{enumerate}

\begin{Shaded}
\begin{Highlighting}[]
\NormalTok{//| echo: false}
\NormalTok{makeStepperHTML(3, [}
\NormalTok{  \{ op: "Write both functions", eq: "f(x) = x\^{}2, \textbackslash{}\textbackslash{}quad g(x) = x {-} 5", note: null \},}
\NormalTok{  \{ op: "Find f(g(x)): substitute g(x) into f", eq: "f(g(x)) = f(x {-} 5) = (x {-} 5)\^{}2", note: "Replace every x in f(x) = x\^{}2 with the expression (x {-} 5)." \},}
\NormalTok{  \{ op: "Expand f(g(x))", eq: "(x {-} 5)\^{}2 = x\^{}2 {-} 10x + 25", note: null \},}
\NormalTok{  \{ op: "Find g(f(x)): substitute f(x) into g", eq: "g(f(x)) = g(x\^{}2) = x\^{}2 {-} 5", note: "Replace every x in g(x) = x {-} 5 with x\^{}2." \},}
\NormalTok{  \{ op: "Evaluate both at x = 3", eq: "f(g(3)) = (3{-}5)\^{}2 = 4, \textbackslash{}\textbackslash{}quad g(f(3)) = 9 {-} 5 = 4", note: "They are equal at x = 3 — but that is a coincidence for this specific value." \},}
\NormalTok{  \{ op: "Evaluate both at x = 2", eq: "f(g(2)) = (2{-}5)\^{}2 = 9, \textbackslash{}\textbackslash{}quad g(f(2)) = 4 {-} 5 = {-}1", note: "Different at x = 2. The functions are genuinely distinct — x = 3 was a coincidence." \}}
\NormalTok{])}
\end{Highlighting}
\end{Shaded}

\begin{center}\rule{0.5\linewidth}{0.5pt}\end{center}

\textbf{4.} A company's profit \(P\) in dollars from selling \(n\) units
is:

\[P(n) = 8n - 150\]

The factory can produce between 0 and 200 units per day.

\begin{enumerate}
\def\labelenumi{(\alph{enumi})}
\item
  State the domain given the physical context.
\item
  At what number of units does the company break even (profit = 0)?
\item
  What is the range of \(P\) over this domain?
\end{enumerate}

\begin{Shaded}
\begin{Highlighting}[]
\NormalTok{//| echo: false}
\NormalTok{makeStepperHTML(4, [}
\NormalTok{  \{ op: "Write the function and note the context", eq: "P(n) = 8n {-} 150, \textbackslash{}\textbackslash{}quad n \textbackslash{}\textbackslash{}in \textbackslash{}\textbackslash{}\{0, 1, 2, \textbackslash{}\textbackslash{}ldots, 200\textbackslash{}\textbackslash{}\}", note: "n must be a non{-}negative integer since you can\textquotesingle{}t sell a fraction of a unit." \},}
\NormalTok{  \{ op: "Find break{-}even: set P(n) = 0 and solve for n", eq: "8n {-} 150 = 0", note: null \},}
\NormalTok{  \{ op: "Add 150 to both sides", eq: "8n = 150", note: null \},}
\NormalTok{  \{ op: "Divide both sides by 8", eq: "n = 18.75", note: "Since n must be a whole number, the company needs at least 19 units to be in profit." \},}
\NormalTok{  \{ op: "Find the range: evaluate P at n = 0 and n = 200", eq: "P(0) = {-}150, \textbackslash{}\textbackslash{}quad P(200) = 8(200) {-} 150 = 1450", note: null \},}
\NormalTok{  \{ op: "State the range", eq: "P \textbackslash{}\textbackslash{}in \textbackslash{}\textbackslash{}\{{-}150, {-}142, {-}134, \textbackslash{}\textbackslash{}ldots, 1450\textbackslash{}\textbackslash{}\}", note: "Multiples of 8 offset by {-}150, from {-}150 to 1450." \},}
\NormalTok{  \{ op: "Check break{-}even at n = 19", eq: "P(19) = 8(19) {-} 150 = 152 {-} 150 = 2 \textbackslash{}\textbackslash{}checkmark", note: "Positive profit at n = 19, confirming 19 units is past break{-}even." \}}
\NormalTok{])}
\end{Highlighting}
\end{Shaded}

\begin{center}\rule{0.5\linewidth}{0.5pt}\end{center}

\textbf{5.} Find the inverse of \(h(x) = \dfrac{x + 4}{3}\).

Then verify that \(h(h^{-1}(x)) = x\).

\begin{Shaded}
\begin{Highlighting}[]
\NormalTok{//| echo: false}
\NormalTok{makeStepperHTML(5, [}
\NormalTok{  \{ op: "Write as y = h(x)", eq: "y = \textbackslash{}\textbackslash{}frac\{x + 4\}\{3\}", note: null \},}
\NormalTok{  \{ op: "Swap x and y to begin finding the inverse", eq: "x = \textbackslash{}\textbackslash{}frac\{y + 4\}\{3\}", note: "Swapping the roles of input and output is what finding an inverse means." \},}
\NormalTok{  \{ op: "Multiply both sides by 3", eq: "3x = y + 4", note: null \},}
\NormalTok{  \{ op: "Subtract 4 from both sides", eq: "y = 3x {-} 4", note: null \},}
\NormalTok{  \{ op: "Write the inverse function", eq: "h\^{}\{{-}1\}(x) = 3x {-} 4", note: null \},}
\NormalTok{  \{ op: "Verify: compute h(h⁻¹(x))", eq: "h(h\^{}\{{-}1\}(x)) = h(3x {-} 4) = \textbackslash{}\textbackslash{}frac\{(3x {-} 4) + 4\}\{3\} = \textbackslash{}\textbackslash{}frac\{3x\}\{3\} = x \textbackslash{}\textbackslash{}checkmark", note: null \}}
\NormalTok{])}
\end{Highlighting}
\end{Shaded}

\begin{center}\rule{0.5\linewidth}{0.5pt}\end{center}

\textbf{6.} The temperature inside a greenhouse over the course of a
morning is modelled by:

\[T(t) = 0.5t^2 + 2t + 15\]

where \(T\) is in °C and \(t\) is time in hours after 6 am, with domain
\(0 \leq t \leq 6\).

\begin{enumerate}
\def\labelenumi{(\alph{enumi})}
\item
  What is the temperature at 6 am (\(t = 0\))?
\item
  What is the temperature at 9 am (\(t = 3\))?
\item
  What is the temperature at noon (\(t = 6\))?
\item
  Is this function one-to-one over this domain? If so, briefly explain
  why.
\end{enumerate}

\begin{Shaded}
\begin{Highlighting}[]
\NormalTok{//| echo: false}
\NormalTok{makeStepperHTML(6, [}
\NormalTok{  \{ op: "Write the function", eq: "T(t) = 0.5t\^{}2 + 2t + 15", note: null \},}
\NormalTok{  \{ op: "Evaluate at t = 0 (6 am)", eq: "T(0) = 0.5(0)\^{}2 + 2(0) + 15 = 15 \textbackslash{}\textbackslash{}text\{ °C\}", note: null \},}
\NormalTok{  \{ op: "Evaluate at t = 3 (9 am)", eq: "T(3) = 0.5(9) + 2(3) + 15 = 4.5 + 6 + 15 = 25.5 \textbackslash{}\textbackslash{}text\{ °C\}", note: null \},}
\NormalTok{  \{ op: "Evaluate at t = 6 (noon)", eq: "T(6) = 0.5(36) + 2(6) + 15 = 18 + 12 + 15 = 45 \textbackslash{}\textbackslash{}text\{ °C\}", note: null \},}
\NormalTok{  \{ op: "Assess whether T is one{-}to{-}one on [0, 6]", eq: "x\_v = {-}\textbackslash{}\textbackslash{}frac\{b\}\{2a\} = {-}\textbackslash{}\textbackslash{}frac\{2\}\{2 \textbackslash{}\textbackslash{}times 0.5\} = {-}2", note: "The vertex of T(t) = 0.5t² + 2t + 15 is at t = −2, which is outside the domain [0, 6]. The parabola opens upward (a = 0.5 \textgreater{} 0) and its vertex is to the left of the domain, so T is strictly increasing on [0, 6] — each input gives a different output." \},}
\NormalTok{  \{ op: "Check: T is increasing throughout", eq: "T(0) = 15 \textless{} T(3) = 25.5 \textless{} T(6) = 45 \textbackslash{}\textbackslash{}checkmark", note: null \}}
\NormalTok{])}
\end{Highlighting}
\end{Shaded}

\part{Volume 4 --- Functions and change}

\chapter{Volume 4 --- Functions and
change}\label{volume-4-functions-and-change-1}

The central move of this volume is tracking scale and structure.

In previous volumes, you studied relationships where scaling happens in
simple or periodic ways. This volume is about what happens when growth
accelerates, when dimensions expand, and when you need to handle many
different, changing variables all at once.

By the end of this volume, you will be comfortable navigating
exponential growth and decay, working within the complex plane,
organizing large sets of linear constraints using matrices, and making
sense of infinite step-by-step accumulations.

\section{Chapters}\label{chapters-3}

\begin{longtable}[]{@{}lll@{}}
\toprule\noalign{}
Chapter & Title & Core move \\
\midrule\noalign{}
\endhead
\bottomrule\noalign{}
\endlastfoot
1 & Exponents and logarithms & Managing multiplicative growth \\
2 & Complex numbers & Expanding the number plane \\
3 & Systems and matrices & Handling multiple relationships at once \\
4 & Sequences and series & Identifying and summing discrete patterns \\
\end{longtable}

\chapter{Exponents and logarithms}\label{exponents-and-logarithms}

Growth, decay, and the scale of things

\hfill\break

You put £500 in a savings account. Each year it grows by 6\%. After one
year you have £530. After two years, £561.80. After twenty years,
something like £1604. That's not 20 × £30 = £600 added --- it's the
interest compounding on itself, quietly doubling your money while you
sleep.

A phone battery starts fully charged. Every hour it drains by half. An
hour in: 50\%. Two hours: 25\%. Three: 12.5\%. That's the same structure
as the savings account, just running backward.

An earthquake measuring 6 on the Richter scale releases ten times more
energy than one measuring 5, and a hundred times more than one measuring
4. The scale compresses a range from ``barely felt'' to ``levels
cities'' into ten manageable numbers.

These are not three different ideas. They're the same mathematical
structure --- exponential growth, exponential decay, and the logarithm
that reads off what power was used.

\section{What the notation is
saying}\label{what-the-notation-is-saying-14}

\subsection{The exponent laws}\label{the-exponent-laws}

Start with what you already know: \(2^3\) means
\(2 \times 2 \times 2 = 8\). The \(3\) is the \textbf{exponent} --- it
counts how many copies of 2 you multiply. The \(2\) is the
\textbf{base}.

The notation \(a^x\) --- read ``a to the power x'' --- means the base
\(a\) multiplied by itself \(x\) times. When you multiply two of these
together with the same base, the exponents add:

\[a^m \times a^n = a^{m+n}\]

because you're just running \(m\) copies followed by \(n\) copies. Four
laws cover everything you need:

\begin{longtable}[]{@{}
  >{\raggedright\arraybackslash}p{(\linewidth - 4\tabcolsep) * \real{0.1852}}
  >{\raggedright\arraybackslash}p{(\linewidth - 4\tabcolsep) * \real{0.4815}}
  >{\raggedright\arraybackslash}p{(\linewidth - 4\tabcolsep) * \real{0.3333}}@{}}
\toprule\noalign{}
\begin{minipage}[b]{\linewidth}\raggedright
Law
\end{minipage} & \begin{minipage}[b]{\linewidth}\raggedright
What it says
\end{minipage} & \begin{minipage}[b]{\linewidth}\raggedright
Example
\end{minipage} \\
\midrule\noalign{}
\endhead
\bottomrule\noalign{}
\endlastfoot
\(a^m \times a^n = a^{m+n}\) & Multiplying: add exponents &
\(2^3 \times 2^4 = 2^7 = 128\) \\
\(a^m \div a^n = a^{m-n}\) & Dividing: subtract exponents &
\(2^5 \div 2^2 = 2^3 = 8\) \\
\((a^m)^n = a^{mn}\) & Power of a power: multiply exponents &
\((2^3)^2 = 2^6 = 64\) \\
\(a^0 = 1\) & Zero power: always one & \(17^0 = 1\) \\
\end{longtable}

Two more for completeness:

\[a^{-n} = \frac{1}{a^n} \qquad a^{1/n} = \sqrt[n]{a}\]

A negative exponent means ``divide by'' instead of ``multiply by'' ---
\(2^{-3} = \frac{1}{8}\). A fractional exponent means a root ---
\(8^{1/3} = \sqrt[3]{8} = 2\). These look odd until you see that they
follow directly from the first law above:
\(a^3 \times a^{-3} = a^0 = 1\), so \(a^{-3}\) must equal
\(\frac{1}{a^3}\).

\subsection{The exponential function}\label{the-exponential-function}

When the exponent is a variable, you have a function:

\[f(x) = a^x\]

Read it as: ``the output is the base \(a\) raised to the power \(x\).''

\begin{itemize}
\tightlist
\item
  If \(a > 1\): the function grows as \(x\) increases. Larger \(a\)
  means faster growth.
\item
  If \(0 < a < 1\): the function decays --- same shape but reflected.
\item
  The y-intercept is always 1: \(f(0) = a^0 = 1\) for any base.
\end{itemize}

The battery example: \(f(t) = (0.5)^t\) gives the fraction of charge
remaining after \(t\) hours. The savings example: \(f(t) = 1.06^t\)
gives the growth factor after \(t\) years (multiply by the initial
amount to get the actual balance).

\subsection{\texorpdfstring{The natural base
\(e\)}{The natural base e}}\label{the-natural-base-e}

You can build an exponential function with any base. But there is one
base that is special.

Consider this question: if a quantity grows at a rate exactly equal to
its current size --- so the bigger it gets, the faster it grows, and the
growth rate is always exactly the current amount --- what is the base?

The answer is \(e \approx 2.718\). This number arises whenever a process
is described by its own rate of change. You can't deduce its exact value
without calculus, but you can use it freely right now. The function
\(f(x) = e^x\) is called the \textbf{natural exponential}. It obeys all
the same exponent laws, with \(e\) as the base.

\(e\) is not arbitrary. It's not a convenient approximation. When you
reach differential calculus (Volume 5), you'll see that \(e^x\) is the
only function that is its own derivative --- the only curve whose slope
at every point equals its own height. That property is what makes it
turn up everywhere in science and engineering.

\begin{Shaded}
\begin{Highlighting}[]
\NormalTok{//| echo: false}
\NormalTok{\{}
\NormalTok{  const div = document.createElement("div");}
\NormalTok{  div.style.cssText = "font{-}family:inherit; padding:0.5rem 0;";}

\NormalTok{  const controls = document.createElement("div");}
\NormalTok{  controls.style.cssText = "display:flex; gap:1.5rem; flex{-}wrap:wrap; margin{-}bottom:0.75rem;";}

\NormalTok{  function makeSlider(label, min, max, step, val) \{}
\NormalTok{    const wrap = document.createElement("div");}
\NormalTok{    const lbl = document.createElement("label");}
\NormalTok{    lbl.style.cssText = "font{-}size:0.85rem; color:\#374151; display:block; margin{-}bottom:0.25rem;";}
\NormalTok{    const span = document.createElement("span");}
\NormalTok{    span.textContent = val;}
\NormalTok{    lbl.appendChild(document.createTextNode(label + " = "));}
\NormalTok{    lbl.appendChild(span);}
\NormalTok{    const input = document.createElement("input");}
\NormalTok{    input.type = "range"; input.min = min; input.max = max; input.step = step; input.value = val;}
\NormalTok{    input.style.cssText = "width:180px; display:block;";}
\NormalTok{    wrap.appendChild(lbl); wrap.appendChild(input);}
\NormalTok{    return \{wrap, input, span\};}
\NormalTok{  \}}

\NormalTok{  const aS = makeSlider("base a", 0.5, 3.0, 0.1, 2.0);}
\NormalTok{  const AS = makeSlider("initial value A", 1, 100, 1, 1);}
\NormalTok{  controls.appendChild(aS.wrap); controls.appendChild(AS.wrap);}

\NormalTok{  const labelEl = document.createElement("div");}
\NormalTok{  labelEl.style.cssText = "font{-}size:0.9rem; font{-}weight:600; margin{-}bottom:0.5rem; color:\#1d4ed8;";}

\NormalTok{  const canvas = document.createElement("canvas");}
\NormalTok{  canvas.width = 500; canvas.height = 300;}
\NormalTok{  canvas.style.cssText = "border:1px solid \#e5e7eb; border{-}radius:6px; width:100\%; max{-}width:500px; display:block;";}

\NormalTok{  div.appendChild(controls); div.appendChild(labelEl); div.appendChild(canvas);}

\NormalTok{  function draw() \{}
\NormalTok{    const a = parseFloat(aS.input.value);}
\NormalTok{    const A = parseFloat(AS.input.value);}
\NormalTok{    aS.span.textContent = a.toFixed(1);}
\NormalTok{    AS.span.textContent = A;}
\NormalTok{    const trend = a \textgreater{} 1 ? "growth" : a \textless{} 1 ? "decay" : "constant";}
\NormalTok{    labelEl.textContent = "f(x) = " + A + " · " + a.toFixed(1) + "ˣ  — exponential " + trend;}

\NormalTok{    const ctx = canvas.getContext("2d");}
\NormalTok{    const W = canvas.width, H = canvas.height;}
\NormalTok{    ctx.clearRect(0, 0, W, H);}

\NormalTok{    const xMin = {-}5, xMax = 5;}
\NormalTok{    const yMax = Math.max(A * Math.pow(a, 5), A * Math.pow(a, {-}5), 10) * 1.2;}
\NormalTok{    const yMin = 0;}

\NormalTok{    function px(x) \{ return (x {-} xMin) / (xMax {-} xMin) * W; \}}
\NormalTok{    function py(y) \{ return H {-} (y {-} yMin) / (yMax {-} yMin) * H; \}}

\NormalTok{    // Axes}
\NormalTok{    ctx.strokeStyle = "\#9ca3af"; ctx.lineWidth = 1;}
\NormalTok{    ctx.beginPath(); ctx.moveTo(px(xMin), py(0)); ctx.lineTo(px(xMax), py(0)); ctx.stroke();}
\NormalTok{    ctx.beginPath(); ctx.moveTo(px(0), 0); ctx.lineTo(px(0), H); ctx.stroke();}

\NormalTok{    // Curve}
\NormalTok{    const isSpecial = Math.abs(a {-} Math.E) \textless{} 0.05;}
\NormalTok{    ctx.strokeStyle = isSpecial ? "\#7c3aed" : "\#2563eb";}
\NormalTok{    ctx.lineWidth = 2.5;}
\NormalTok{    ctx.beginPath();}
\NormalTok{    let first = true;}
\NormalTok{    for (let px\_ = 0; px\_ \textless{}= W; px\_++) \{}
\NormalTok{      const x = xMin + (px\_ / W) * (xMax {-} xMin);}
\NormalTok{      const y = A * Math.pow(a, x);}
\NormalTok{      if (y \textless{} 0 || y \textgreater{} yMax * 1.1) \{ first = true; continue; \}}
\NormalTok{      if (first) \{ ctx.moveTo(px\_, py(y)); first = false; \}}
\NormalTok{      else ctx.lineTo(px\_, py(y));}
\NormalTok{    \}}
\NormalTok{    ctx.stroke();}

\NormalTok{    // y{-}intercept}
\NormalTok{    ctx.fillStyle = "\#ef4444";}
\NormalTok{    ctx.beginPath(); ctx.arc(px(0), py(A), 5, 0, 2*Math.PI); ctx.fill();}
\NormalTok{    ctx.fillStyle = "\#374151"; ctx.font = "12px sans{-}serif";}
\NormalTok{    ctx.fillText("(0, " + A + ")", px(0) + 7, py(A) {-} 5);}

\NormalTok{    // Point at x=1}
\NormalTok{    const y1 = A * a;}
\NormalTok{    if (y1 \textless{}= yMax) \{}
\NormalTok{      ctx.fillStyle = "\#059669";}
\NormalTok{      ctx.beginPath(); ctx.arc(px(1), py(y1), 5, 0, 2*Math.PI); ctx.fill();}
\NormalTok{      ctx.fillStyle = "\#374151";}
\NormalTok{      ctx.fillText("(1, " + y1.toFixed(1) + ")", px(1) + 7, py(y1) {-} 5);}
\NormalTok{    \}}

\NormalTok{    if (isSpecial) \{}
\NormalTok{      ctx.fillStyle = "\#7c3aed"; ctx.font = "bold 12px sans{-}serif";}
\NormalTok{      ctx.fillText("a = e ≈ 2.718", px(1.5), py(A * Math.E) + 15);}
\NormalTok{    \}}
\NormalTok{  \}}

\NormalTok{  aS.input.addEventListener("input", draw);}
\NormalTok{  AS.input.addEventListener("input", draw);}
\NormalTok{  draw();}
\NormalTok{  return div;}
\NormalTok{\}}
\end{Highlighting}
\end{Shaded}

\section{The method}\label{the-method-13}

\subsection{Part 1: Working with exponent
laws}\label{part-1-working-with-exponent-laws}

The laws in the table above are the full toolkit. A typical
simplification chains two or three of them together.

\textbf{Simplify \(\dfrac{x^5 \cdot x^{-2}}{x^3}\).}

Top: \(x^5 \cdot x^{-2} = x^{5 + (-2)} = x^3\).

Then: \(x^3 \div x^3 = x^{3-3} = x^0 = 1\).

\textbf{Simplify \((2a^3)^4\).}

The power applies to everything inside the bracket:
\((2)^4 \cdot (a^3)^4 = 16 \cdot a^{12} = 16a^{12}\).

\subsection{Part 2: Solving exponential equations by
inspection}\label{part-2-solving-exponential-equations-by-inspection}

Some equations can be solved by rewriting both sides as powers of the
same base.

\textbf{Solve \(2^x = 16\).}

Rewrite 16 as a power of 2: \(16 = 2^4\).

So \(2^x = 2^4\), which means \(x = 4\).

\textbf{Solve \(9^x = 27\).}

Both 9 and 27 are powers of 3: \(9 = 3^2\) and \(27 = 3^3\).

So \((3^2)^x = 3^3\), giving \(3^{2x} = 3^3\), so \(2x = 3\), and
\(x = \frac{3}{2}\).

This only works when you can express both sides with the same base. For
\(2^x = 10\), you can't. That's where logarithms come in.

\subsection{Part 3: Logarithms}\label{part-3-logarithms}

A logarithm is the inverse of an exponential. It answers the question:
``what power do I need?''

\[\log_a(b) = c \quad \text{means} \quad a^c = b\]

Read it aloud: ``log base \(a\) of \(b\) equals \(c\)'' means ``\(a\)
raised to the power \(c\) gives \(b\).''

Three examples to fix this:

\[\log_2(8) = 3 \quad \text{because} \quad 2^3 = 8\]

\[\log_{10}(1000) = 3 \quad \text{because} \quad 10^3 = 1000\]

\[\log_3(9) = 2 \quad \text{because} \quad 3^2 = 9\]

Two special bases appear so often they have short names:

\begin{itemize}
\tightlist
\item
  \(\log_{10}(x)\) is often written \(\log(x)\) --- the \textbf{common
  logarithm}
\item
  \(\log_e(x)\) is written \(\ln(x)\) --- the \textbf{natural
  logarithm}, read ``ell-en of x''
\end{itemize}

\begin{tcolorbox}[enhanced jigsaw, opacityback=0, colback=white, breakable, colframe=quarto-callout-note-color-frame, arc=.35mm, left=2mm, toprule=.15mm, leftrule=.75mm, bottomrule=.15mm, rightrule=.15mm]
\begin{minipage}[t]{5.5mm}
\textcolor{quarto-callout-note-color}{\faInfo}
\end{minipage}%
\begin{minipage}[t]{\textwidth - 5.5mm}

\textbf{Why this works}

Logarithms are the inverse of exponentials the same way division undoes
multiplication. If \(2^3 = 8\), then \(\log_2(8) = 3\). One statement,
two ways of writing it. The exponential says ``start with the power, get
the result.'' The logarithm says ``start with the result, get the
power.'' They are the same relationship read in opposite directions.

\end{minipage}%
\end{tcolorbox}

\begin{Shaded}
\begin{Highlighting}[]
\NormalTok{//| echo: false}
\NormalTok{\{}
\NormalTok{  const div = document.createElement("div");}
\NormalTok{  div.style.cssText = "font{-}family:inherit; padding:0.5rem 0;";}

\NormalTok{  const controls = document.createElement("div");}
\NormalTok{  controls.style.cssText = "display:flex; align{-}items:center; gap:1rem; margin{-}bottom:0.75rem; flex{-}wrap:wrap;";}
\NormalTok{  const lbl = document.createElement("label");}
\NormalTok{  lbl.style.cssText = "font{-}size:0.85rem; color:\#374151;";}
\NormalTok{  const spanVal = document.createElement("span"); spanVal.textContent = "2";}
\NormalTok{  lbl.appendChild(document.createTextNode("base a = ")); lbl.appendChild(spanVal);}
\NormalTok{  const sl = document.createElement("input");}
\NormalTok{  sl.type = "range"; sl.min = 2; sl.max = 10; sl.step = 1; sl.value = 2;}
\NormalTok{  sl.style.width = "160px";}
\NormalTok{  controls.appendChild(lbl); controls.appendChild(sl);}

\NormalTok{  const infoEl = document.createElement("div");}
\NormalTok{  infoEl.id = "viz2\_info";}
\NormalTok{  infoEl.style.cssText = "font{-}size:0.88rem; color:\#1d4ed8; margin{-}bottom:0.5rem; min{-}height:1.4em;";}

\NormalTok{  const container = document.createElement("div");}
\NormalTok{  container.style.cssText = "display:flex; gap:1rem; flex{-}wrap:wrap;";}
\NormalTok{  const c1 = document.createElement("canvas"); c1.width = 220; c1.height = 220;}
\NormalTok{  c1.style.cssText = "border:1px solid \#e5e7eb; border{-}radius:6px;";}
\NormalTok{  const c2 = document.createElement("canvas"); c2.width = 220; c2.height = 220;}
\NormalTok{  c2.style.cssText = "border:1px solid \#e5e7eb; border{-}radius:6px;";}
\NormalTok{  container.appendChild(c1); container.appendChild(c2);}

\NormalTok{  div.appendChild(controls); div.appendChild(infoEl); div.appendChild(container);}

\NormalTok{  function drawPanel(canvas, xMin, xMax, yMin, yMax, fn, color, ptX) \{}
\NormalTok{    const ctx = canvas.getContext("2d");}
\NormalTok{    const W = canvas.width, H = canvas.height;}
\NormalTok{    ctx.clearRect(0, 0, W, H);}
\NormalTok{    function px(x) \{ return (x {-} xMin) / (xMax {-} xMin) * W; \}}
\NormalTok{    function py(y) \{ return H {-} (y {-} yMin) / (yMax {-} yMin) * H; \}}

\NormalTok{    // Reflection line y=x}
\NormalTok{    ctx.strokeStyle = "\#d1d5db"; ctx.lineWidth = 1; ctx.setLineDash([4,4]);}
\NormalTok{    ctx.beginPath(); ctx.moveTo(px(xMin), py(yMin)); ctx.lineTo(px(xMax), py(yMax)); ctx.stroke();}
\NormalTok{    ctx.setLineDash([]);}

\NormalTok{    // Axes}
\NormalTok{    ctx.strokeStyle = "\#9ca3af"; ctx.lineWidth = 1;}
\NormalTok{    ctx.beginPath(); ctx.moveTo(px(xMin), py(0)); ctx.lineTo(px(xMax), py(0)); ctx.stroke();}
\NormalTok{    ctx.beginPath(); ctx.moveTo(px(0), 0); ctx.lineTo(px(0), H); ctx.stroke();}

\NormalTok{    // Curve}
\NormalTok{    ctx.strokeStyle = color; ctx.lineWidth = 2;}
\NormalTok{    ctx.beginPath();}
\NormalTok{    let first = true;}
\NormalTok{    for (let xi = 0; xi \textless{}= W; xi++) \{}
\NormalTok{      const x = xMin + (xi / W) * (xMax {-} xMin);}
\NormalTok{      const y = fn(x);}
\NormalTok{      if (isNaN(y) || !isFinite(y) || y \textless{} yMin {-} 0.5 || y \textgreater{} yMax + 0.5) \{ first = true; continue; \}}
\NormalTok{      if (first) \{ ctx.moveTo(xi, py(y)); first = false; \}}
\NormalTok{      else ctx.lineTo(xi, py(y));}
\NormalTok{    \}}
\NormalTok{    ctx.stroke();}

\NormalTok{    // Highlight point}
\NormalTok{    if (ptX !== null) \{}
\NormalTok{      const ptY = fn(ptX);}
\NormalTok{      if (!isNaN(ptY) \&\& isFinite(ptY) \&\& ptY \textgreater{}= yMin \&\& ptY \textless{}= yMax) \{}
\NormalTok{        ctx.fillStyle = "\#ef4444";}
\NormalTok{        ctx.beginPath(); ctx.arc(px(ptX), py(ptY), 5, 0, 2*Math.PI); ctx.fill();}
\NormalTok{        ctx.fillStyle = "\#374151"; ctx.font = "11px sans{-}serif";}
\NormalTok{        ctx.fillText("(" + ptX.toFixed(0) + ", " + ptY.toFixed(1) + ")", px(ptX)+7, py(ptY){-}5);}
\NormalTok{      \}}
\NormalTok{    \}}
\NormalTok{  \}}

\NormalTok{  function update() \{}
\NormalTok{    const b = parseInt(sl.value);}
\NormalTok{    spanVal.textContent = b;}
\NormalTok{    const ptX = b; // highlight x=b on the exp curve}
\NormalTok{    drawPanel(c1, {-}2, 4, {-}1, Math.pow(b, 4)+1, x =\textgreater{} Math.pow(b, x), "\#2563eb", ptX);}
\NormalTok{    const logPtX = Math.pow(b, ptX);}
\NormalTok{    drawPanel(c2, 0.01, Math.pow(b,4)+1, {-}2, 4, x =\textgreater{} Math.log(x)/Math.log(b), "\#059669", logPtX \textgreater{} 0 ? logPtX : null);}
\NormalTok{    infoEl.textContent = b + "\^{}" + ptX + " = " + Math.pow(b, ptX).toFixed(0) + "  ↔  log\_" + b + "(" + Math.pow(b, ptX).toFixed(0) + ") = " + ptX;}
\NormalTok{  \}}

\NormalTok{  sl.addEventListener("input", update);}
\NormalTok{  update();}
\NormalTok{  return div;}
\NormalTok{\}}
\end{Highlighting}
\end{Shaded}

\subsection{Part 4: Log laws and solving logarithmic
equations}\label{part-4-log-laws-and-solving-logarithmic-equations}

Logarithms inherit structure from the exponent laws. The three log laws:

\begin{longtable}[]{@{}
  >{\raggedright\arraybackslash}p{(\linewidth - 4\tabcolsep) * \real{0.2903}}
  >{\raggedright\arraybackslash}p{(\linewidth - 4\tabcolsep) * \real{0.4194}}
  >{\raggedright\arraybackslash}p{(\linewidth - 4\tabcolsep) * \real{0.2903}}@{}}
\toprule\noalign{}
\begin{minipage}[b]{\linewidth}\raggedright
Log law
\end{minipage} & \begin{minipage}[b]{\linewidth}\raggedright
What it says
\end{minipage} & \begin{minipage}[b]{\linewidth}\raggedright
Example
\end{minipage} \\
\midrule\noalign{}
\endhead
\bottomrule\noalign{}
\endlastfoot
\(\log_a(mn) = \log_a(m) + \log_a(n)\) & Log of a product: add logs &
\(\log_2(4 \times 8) = \log_2(4) + \log_2(8) = 2 + 3 = 5\) \\
\(\log_a\!\left(\dfrac{m}{n}\right) = \log_a(m) - \log_a(n)\) & Log of a
quotient: subtract logs & \(\log_{10}(1000/10) = 3 - 1 = 2\) \\
\(\log_a(m^n) = n \cdot \log_a(m)\) & Log of a power: bring down the
exponent & \(\log_2(8^4) = 4 \cdot \log_2(8) = 4 \times 3 = 12\) \\
\end{longtable}

One more, needed for calculators --- the \textbf{change of base
formula}:

\[\log_a(x) = \frac{\ln(x)}{\ln(a)} = \frac{\log(x)}{\log(a)}\]

Your calculator has \(\ln\) and \(\log_{10}\) buttons. This formula lets
you compute any log from them.

\textbf{Solving \(2^x = 10\).}

Take \(\log_{10}\) (or \(\ln\)) of both sides:

\[\log(2^x) = \log(10)\]

Use the power law to bring \(x\) down:

\[x \cdot \log(2) = 1\]

Divide:

\[x = \frac{1}{\log(2)} = \frac{1}{0.3010\ldots} \approx 3.32\]

Check: \(2^{3.32} \approx 9.97 \approx 10\). Yes.

\textbf{Solving \(\log_3(x) + \log_3(x - 2) = 1\).}

Use the product law to combine the left side:

\[\log_3\bigl(x(x-2)\bigr) = 1\]

Convert back to exponential form (\(3^1 = x(x-2)\)):

\[x^2 - 2x = 3\]

\[x^2 - 2x - 3 = 0\]

\[(x-3)(x+1) = 0\]

So \(x = 3\) or \(x = -1\). But \(\log_3(-1)\) is not defined --- you
can't take the log of a negative number. Discard \(x = -1\).

Answer: \(x = 3\). Check: \(\log_3(3) + \log_3(1) = 1 + 0 = 1\). Yes.

\section{Worked examples}\label{worked-examples-13}

\textbf{Example 1 --- Science: radioactive half-life.}

Carbon-14 has a half-life of 5730 years --- meaning every 5730 years,
half the remaining carbon-14 decays. An organic sample contains 1.0 g of
carbon-14 today. How much will remain after 17 190 years?

17 190 years = 3 × 5730, so three half-lives have passed.

\[\text{Remaining} = 1.0 \times \left(\frac{1}{2}\right)^3 = \frac{1}{8} = 0.125 \text{ g}\]

If instead you need to find \emph{when} the sample drops to 0.10 g, set
up the equation:

\[0.10 = 1.0 \times \left(\frac{1}{2}\right)^{t/5730}\]

Divide both sides by 1.0 to isolate the exponential:

\[0.10 = (0.5)^{t/5730}\]

Take \(\ln\) of both sides:

\[\ln(0.10) = \ln\!\left((0.5)^{t/5730}\right)\]

Apply the power law --- bring the exponent down:

\[\ln(0.10) = \frac{t}{5730} \cdot \ln(0.5)\]

\[t = 5730 \times \frac{\ln(0.10)}{\ln(0.5)} = 5730 \times \frac{-2.3026}{-0.6931} \approx 5730 \times 3.322 \approx 19\,030 \text{ years}\]

\begin{center}\rule{0.5\linewidth}{0.5pt}\end{center}

\textbf{Example 2 --- Finance: compound interest doubling time.}

An investment earns 7\% per year, compounded annually. How long until it
doubles?

Let the initial amount be \(P\). You want the time \(t\) when
\(P \times 1.07^t = 2P\).

Divide both sides by \(P\):

\[1.07^t = 2\]

Take \(\ln\) of both sides and bring \(t\) down:

\[t \cdot \ln(1.07) = \ln(2)\]

\[t = \frac{\ln(2)}{\ln(1.07)} = \frac{0.6931}{0.06766} \approx 10.2 \text{ years}\]

The Rule of 72 is a quick mental estimate: divide 72 by the interest
rate. \(72 / 7 \approx 10.3\) --- very close. That rule is a log
approximation.

\begin{center}\rule{0.5\linewidth}{0.5pt}\end{center}

\textbf{Example 3 --- Engineering: the decibel.}

Sound intensity is measured on a logarithmic scale. The decibel (dB)
level is:

\[L = 10 \cdot \log_{10}\!\left(\frac{I}{I_0}\right)\]

where \(I\) is the intensity of the sound and \(I_0 = 10^{-12}\) W/m² is
the reference intensity. The reference intensity \(I_0 = 10^{-12}\) W/m²
is the quietest sound a human ear can detect --- it is the standard
baseline from which decibels are measured.

Normal conversation is about \(I = 10^{-6}\) W/m². What is that in dB?

\[L = 10 \cdot \log_{10}\!\left(\frac{10^{-6}}{10^{-12}}\right) = 10 \cdot \log_{10}(10^6) = 10 \times 6 = 60 \text{ dB}\]

A sound at 70 dB has what intensity ratio compared to 60 dB?

The difference is 10 dB, which corresponds to a factor of
\(10^{10/10} = 10^1 = 10\). A 70 dB sound is 10 times as intense as a 60
dB sound. The logarithm compresses very large ratios into small,
manageable numbers.

\begin{center}\rule{0.5\linewidth}{0.5pt}\end{center}

\textbf{Example 4 --- Computing: binary search step count.}

Binary search finds a target in a sorted list by repeatedly halving the
search space. How many steps to search a list of \(n = 1{,}000{,}000\)
items?

After \(k\) steps, the search space has been halved \(k\) times, leaving
at most \(n / 2^k\) items. The search ends when \(n / 2^k = 1\):

\[2^k = n \implies k = \log_2(n)\]

For one million items: \(k = \log_2(1{,}000{,}000) \approx 19.93\), so
at most 20 steps. That's why \(O(\log n)\) algorithms are fast ---
doubling the input adds only one more step.

\section{Where this goes}\label{where-this-goes-14}

The derivative of \(e^x\) is \(e^x\) --- the only function that is its
own slope. That property, which you can only properly state and prove in
differential calculus (Volume 5), is why \(e\) appears in every model of
continuous growth and decay. You already have the function; calculus
will explain what makes it uniquely suited to describe rates of change.

From there, ordinary differential equations (Volume 7) make the
connection concrete. The simplest ODE --- ``a quantity grows
proportionally to itself'' --- has \(e^{kt}\) as its solution.
Radioactive decay, population growth, cooling, charging capacitors: all
of them are that one equation in different physical clothes. This
chapter is the front end of that pipeline.

\begin{tcolorbox}[enhanced jigsaw, opacityback=0, colback=white, breakable, colframe=quarto-callout-tip-color-frame, arc=.35mm, left=2mm, toprule=.15mm, leftrule=.75mm, bottomrule=.15mm, rightrule=.15mm]
\begin{minipage}[t]{5.5mm}
\textcolor{quarto-callout-tip-color}{\faLightbulb}
\end{minipage}%
\begin{minipage}[t]{\textwidth - 5.5mm}

\textbf{Where this shows up}

\begin{itemize}
\tightlist
\item
  A physicist uses exponential decay to date ancient organic material
  via carbon-14 ratios.
\item
  A financial analyst uses logarithms to find how long an investment
  takes to reach a target value.
\item
  An audio engineer specifies amplifier gain and speaker sensitivity in
  decibels --- a logarithmic unit.
\item
  A software engineer analyses why binary search handles a billion
  records in about 30 comparisons, while linear search needs a billion.
\end{itemize}

The base changes. The structure does not.

\end{minipage}%
\end{tcolorbox}

\section{Exercises}\label{exercises-14}

These are puzzles. Each has a clean answer. The interesting part is
identifying whether you're working with an exponential, a log, or both
--- and setting up the equation before you solve it.

\begin{Shaded}
\begin{Highlighting}[]
\NormalTok{//| echo: false}
\NormalTok{katex = require("katex@0.16")}
\end{Highlighting}
\end{Shaded}

\begin{Shaded}
\begin{Highlighting}[]
\NormalTok{//| echo: false}
\NormalTok{function makeStepperHTML(exerciseNum, steps) \{}
\NormalTok{  let current = 0;}
\NormalTok{  const totalSteps = steps.length;}
\NormalTok{  const container = document.createElement("div");}
\NormalTok{  container.style.cssText = "border:1px solid \#e5e7eb; border{-}radius:8px; padding:1rem 1.25rem; margin:0.75rem 0; font{-}family:inherit;";}
\NormalTok{  const stepsDiv = document.createElement("div");}
\NormalTok{  stepsDiv.style.marginBottom = "0.75rem";}
\NormalTok{  function renderSteps() \{}
\NormalTok{    stepsDiv.innerHTML = "";}
\NormalTok{    for (let i = 0; i \textless{} current; i++) \{}
\NormalTok{      const s = steps[i];}
\NormalTok{      const row = document.createElement("div");}
\NormalTok{      row.style.cssText = "display:grid; grid{-}template{-}columns:200px 1fr; gap:0.35rem 1rem; align{-}items:baseline; padding:0.35rem 0; border{-}top:1px solid \#f3f4f6;";}
\NormalTok{      const opEl = document.createElement("span");}
\NormalTok{      opEl.style.cssText = "font{-}size:0.8rem; color:\#6b7280; font{-}style:italic;";}
\NormalTok{      opEl.textContent = s.op;}
\NormalTok{      const eqEl = document.createElement("span");}
\NormalTok{      eqEl.innerHTML = katex.renderToString(s.eq, \{throwOnError: false, displayMode: false\});}
\NormalTok{      row.appendChild(opEl);}
\NormalTok{      row.appendChild(eqEl);}
\NormalTok{      container.appendChild(row);}
\NormalTok{      if (s.note) \{}
\NormalTok{        const noteEl = document.createElement("p");}
\NormalTok{        noteEl.style.cssText = "margin:0.2rem 0 0.5rem 200px; font{-}size:0.82rem; color:\#374151;";}
\NormalTok{        noteEl.textContent = s.note;}
\NormalTok{        container.appendChild(noteEl);}
\NormalTok{      \}}
\NormalTok{    \}}
\NormalTok{    if (current === totalSteps) \{}
\NormalTok{      const done = document.createElement("p");}
\NormalTok{      done.style.cssText = "margin{-}top:0.5rem; font{-}size:0.85rem; color:\#059669; font{-}weight:600;";}
\NormalTok{      done.textContent = "✓ Solution complete";}
\NormalTok{      stepsDiv.appendChild(done);}
\NormalTok{    \}}
\NormalTok{  \}}
\NormalTok{  const controls = document.createElement("div");}
\NormalTok{  controls.style.cssText = "display:flex; align{-}items:center; gap:0.75rem; margin{-}top:0.5rem;";}
\NormalTok{  const nextBtn = document.createElement("button");}
\NormalTok{  nextBtn.textContent = "Next step →";}
\NormalTok{  nextBtn.style.cssText = "padding:0.35rem 0.85rem; background:\#2563eb; color:\#fff; border:none; border{-}radius:6px; cursor:pointer; font{-}size:0.85rem;";}
\NormalTok{  const resetBtn = document.createElement("button");}
\NormalTok{  resetBtn.textContent = "Reset";}
\NormalTok{  resetBtn.style.cssText = "padding:0.35rem 0.85rem; background:\#f3f4f6; color:\#374151; border:1px solid \#d1d5db; border{-}radius:6px; cursor:pointer; font{-}size:0.85rem;";}
\NormalTok{  const counter = document.createElement("span");}
\NormalTok{  counter.style.cssText = "font{-}size:0.8rem; color:\#6b7280;";}
\NormalTok{  nextBtn.onclick = () =\textgreater{} \{ if (current \textless{} totalSteps) \{ current++; renderSteps(); counter.textContent = current + " / " + totalSteps; if (current === totalSteps) nextBtn.disabled = true; \} \};}
\NormalTok{  resetBtn.onclick = () =\textgreater{} \{ current = 0; nextBtn.disabled = false; renderSteps(); counter.textContent = "0 / " + totalSteps; \};}
\NormalTok{  counter.textContent = "0 / " + totalSteps;}
\NormalTok{  controls.appendChild(nextBtn);}
\NormalTok{  controls.appendChild(resetBtn);}
\NormalTok{  controls.appendChild(counter);}
\NormalTok{  container.appendChild(stepsDiv);}
\NormalTok{  container.appendChild(controls);}
\NormalTok{  return container;}
\NormalTok{\}}
\end{Highlighting}
\end{Shaded}

\textbf{1.} Simplify, writing your answer as a single power:

\[\frac{3^8 \times 3^{-3}}{3^2}\]

\begin{Shaded}
\begin{Highlighting}[]
\NormalTok{//| echo: false}
\NormalTok{\{}
\NormalTok{  const el = makeStepperHTML(1, [}
\NormalTok{    \{ op: "Write as an equation", eq: "\textbackslash{}\textbackslash{}frac\{3\^{}8 \textbackslash{}\textbackslash{}times 3\^{}\{{-}3\}\}\{3\^{}2\}", note: null \},}
\NormalTok{    \{ op: "Multiply in the numerator — add exponents", eq: "\textbackslash{}\textbackslash{}frac\{3\^{}\{8 + ({-}3)\}\}\{3\^{}2\} = \textbackslash{}\textbackslash{}frac\{3\^{}5\}\{3\^{}2\}", note: "Product law: same base, so add the exponents." \},}
\NormalTok{    \{ op: "Divide — subtract exponents", eq: "3\^{}\{5 {-} 2\} = 3\^{}3", note: "Quotient law: same base, so subtract the exponents." \},}
\NormalTok{    \{ op: "Evaluate", eq: "3\^{}3 = 27", note: null \},}
\NormalTok{    \{ op: "Check", eq: "\textbackslash{}\textbackslash{}frac\{3\^{}8 \textbackslash{}\textbackslash{}times 3\^{}\{{-}3\}\}\{3\^{}2\} = \textbackslash{}\textbackslash{}frac\{6561 \textbackslash{}\textbackslash{}times (1/27)\}\{9\} = \textbackslash{}\textbackslash{}frac\{243\}\{9\} = 27 \textbackslash{}\textbackslash{}checkmark", note: null \},}
\NormalTok{  ]);}
\NormalTok{  return el;}
\NormalTok{\}}
\end{Highlighting}
\end{Shaded}

\begin{center}\rule{0.5\linewidth}{0.5pt}\end{center}

\textbf{2.} Solve for \(x\):

\[4^x = 128\]

\begin{Shaded}
\begin{Highlighting}[]
\NormalTok{//| echo: false}
\NormalTok{\{}
\NormalTok{  const el = makeStepperHTML(2, [}
\NormalTok{    \{ op: "Write as an equation", eq: "4\^{}x = 128", note: null \},}
\NormalTok{    \{ op: "Express both sides as powers of 2", eq: "(2\^{}2)\^{}x = 2\^{}7", note: "4 = 2² and 128 = 2⁷, so rewrite with a common base." \},}
\NormalTok{    \{ op: "Use the power{-}of{-}a{-}power law", eq: "2\^{}\{2x\} = 2\^{}7", note: null \},}
\NormalTok{    \{ op: "Match exponents — same base, so the powers must be equal", eq: "2x = 7", note: null \},}
\NormalTok{    \{ op: "Divide both sides by 2", eq: "x = \textbackslash{}\textbackslash{}frac\{7\}\{2\} = 3.5", note: null \},}
\NormalTok{    \{ op: "Check", eq: "4\^{}\{3.5\} = 4\^{}3 \textbackslash{}\textbackslash{}times 4\^{}\{0.5\} = 64 \textbackslash{}\textbackslash{}times 2 = 128 \textbackslash{}\textbackslash{}checkmark", note: null \},}
\NormalTok{  ]);}
\NormalTok{  return el;}
\NormalTok{\}}
\end{Highlighting}
\end{Shaded}

\begin{center}\rule{0.5\linewidth}{0.5pt}\end{center}

\textbf{3.} A population of bacteria doubles every 3 hours. You start
with 400 bacteria. How many hours until there are more than 25 600?

\begin{Shaded}
\begin{Highlighting}[]
\NormalTok{//| echo: false}
\NormalTok{\{}
\NormalTok{  const el = makeStepperHTML(3, [}
\NormalTok{    \{ op: "Write as an equation", eq: "400 \textbackslash{}\textbackslash{}times 2\^{}\{t/3\} \textgreater{} 25600", note: "Every 3 hours is one doubling cycle, so the exponent is t/3." \},}
\NormalTok{    \{ op: "Divide both sides by 400", eq: "2\^{}\{t/3\} \textgreater{} 64", note: null \},}
\NormalTok{    \{ op: "Express 64 as a power of 2", eq: "2\^{}\{t/3\} \textgreater{} 2\^{}6", note: null \},}
\NormalTok{    \{ op: "Match exponents", eq: "\textbackslash{}\textbackslash{}frac\{t\}\{3\} \textgreater{} 6", note: null \},}
\NormalTok{    \{ op: "Multiply both sides by 3", eq: "t \textgreater{} 18 \textbackslash{}\textbackslash{}text\{ hours\}", note: null \},}
\NormalTok{    \{ op: "Check", eq: "400 \textbackslash{}\textbackslash{}times 2\^{}\{18/3\} = 400 \textbackslash{}\textbackslash{}times 2\^{}6 = 400 \textbackslash{}\textbackslash{}times 64 = 25600 \textbackslash{}\textbackslash{}checkmark \textbackslash{}\textbackslash{}text\{ (exactly at 18 h, so \textgreater{}25600 means after 18 h)\}", note: null \},}
\NormalTok{  ]);}
\NormalTok{  return el;}
\NormalTok{\}}
\end{Highlighting}
\end{Shaded}

\begin{center}\rule{0.5\linewidth}{0.5pt}\end{center}

\textbf{4.} Use logarithms to solve for \(t\):

\[5^t = 200\]

Give your answer to two decimal places.

\begin{Shaded}
\begin{Highlighting}[]
\NormalTok{//| echo: false}
\NormalTok{\{}
\NormalTok{  const el = makeStepperHTML(4, [}
\NormalTok{    \{ op: "Write as an equation", eq: "5\^{}t = 200", note: null \},}
\NormalTok{    \{ op: "Take ln of both sides", eq: "\textbackslash{}\textbackslash{}ln(5\^{}t) = \textbackslash{}\textbackslash{}ln(200)", note: "You can use log base 10 equally well — the ratio will be the same." \},}
\NormalTok{    \{ op: "Use the power law — bring t down", eq: "t \textbackslash{}\textbackslash{}cdot \textbackslash{}\textbackslash{}ln(5) = \textbackslash{}\textbackslash{}ln(200)", note: "\textbackslash{}\textbackslash{}ln(m\^{}n) = n \textbackslash{}\textbackslash{}cdot \textbackslash{}\textbackslash{}ln(m)" \},}
\NormalTok{    \{ op: "Divide both sides by ln(5)", eq: "t = \textbackslash{}\textbackslash{}frac\{\textbackslash{}\textbackslash{}ln(200)\}\{\textbackslash{}\textbackslash{}ln(5)\} = \textbackslash{}\textbackslash{}frac\{5.2983\}\{1.6094\}", note: null \},}
\NormalTok{    \{ op: "Evaluate", eq: "t \textbackslash{}\textbackslash{}approx 3.29", note: null \},}
\NormalTok{    \{ op: "Check", eq: "5\^{}\{3.29\} \textbackslash{}\textbackslash{}approx 5\^{}3 \textbackslash{}\textbackslash{}times 5\^{}\{0.29\} \textbackslash{}\textbackslash{}approx 125 \textbackslash{}\textbackslash{}times 1.60 \textbackslash{}\textbackslash{}approx 200 \textbackslash{}\textbackslash{}checkmark", note: null \},}
\NormalTok{  ]);}
\NormalTok{  return el;}
\NormalTok{\}}
\end{Highlighting}
\end{Shaded}

\begin{center}\rule{0.5\linewidth}{0.5pt}\end{center}

\textbf{5.} Solve the logarithmic equation:

\[\log_2(x+4) + \log_2(x-4) = 5\]

\begin{Shaded}
\begin{Highlighting}[]
\NormalTok{//| echo: false}
\NormalTok{\{}
\NormalTok{  const el = makeStepperHTML(5, [}
\NormalTok{    \{ op: "Write as an equation", eq: "\textbackslash{}\textbackslash{}log\_2(x+4) + \textbackslash{}\textbackslash{}log\_2(x{-}4) = 5", note: null \},}
\NormalTok{    \{ op: "Use the product law — combine the two logs", eq: "\textbackslash{}\textbackslash{}log\_2\textbackslash{}\textbackslash{}bigl((x+4)(x{-}4)\textbackslash{}\textbackslash{}bigr) = 5", note: "\textbackslash{}\textbackslash{}log\_a(m) + \textbackslash{}\textbackslash{}log\_a(n) = \textbackslash{}\textbackslash{}log\_a(mn)" \},}
\NormalTok{    \{ op: "Expand the bracket", eq: "\textbackslash{}\textbackslash{}log\_2(x\^{}2 {-} 16) = 5", note: "(x+4)(x{-}4) = x\^{}2 {-} 16. This uses the difference of two squares shortcut: (a + b)(a − b) = a² − b². You may have seen this in Vol 3 Ch 3." \},}
\NormalTok{    \{ op: "Convert to exponential form", eq: "x\^{}2 {-} 16 = 2\^{}5 = 32", note: "\textbackslash{}\textbackslash{}log\_2(k) = 5 \textbackslash{}\textbackslash{}text\{ means \} 2\^{}5 = k" \},}
\NormalTok{    \{ op: "Solve the quadratic", eq: "x\^{}2 = 48 \textbackslash{}\textbackslash{}implies x = \textbackslash{}\textbackslash{}pm\textbackslash{}\textbackslash{}sqrt\{48\} = \textbackslash{}\textbackslash{}pm 4\textbackslash{}\textbackslash{}sqrt\{3\} \textbackslash{}\textbackslash{}approx \textbackslash{}\textbackslash{}pm 6.93", note: null \},}
\NormalTok{    \{ op: "Reject the negative solution", eq: "x = 4\textbackslash{}\textbackslash{}sqrt\{3\} \textbackslash{}\textbackslash{}approx 6.93", note: "x = {-}4\textbackslash{}\textbackslash{}sqrt\{3\} gives \textbackslash{}\textbackslash{}log\_2({-}4\textbackslash{}\textbackslash{}sqrt\{3\}+4), which is log of a negative — undefined." \},}
\NormalTok{    \{ op: "Check", eq: "\textbackslash{}\textbackslash{}log\_2(4\textbackslash{}\textbackslash{}sqrt\{3\}+4) + \textbackslash{}\textbackslash{}log\_2(4\textbackslash{}\textbackslash{}sqrt\{3\}{-}4) = \textbackslash{}\textbackslash{}log\_2\textbackslash{}\textbackslash{}bigl((4\textbackslash{}\textbackslash{}sqrt\{3\})\^{}2 {-} 16\textbackslash{}\textbackslash{}bigr) = \textbackslash{}\textbackslash{}log\_2(48{-}16) = \textbackslash{}\textbackslash{}log\_2(32) = 5 \textbackslash{}\textbackslash{}checkmark", note: null \},}
\NormalTok{  ]);}
\NormalTok{  return el;}
\NormalTok{\}}
\end{Highlighting}
\end{Shaded}

\begin{center}\rule{0.5\linewidth}{0.5pt}\end{center}

\textbf{6.} The pH of a solution is defined as
\(\text{pH} = -\log_{10}[\text{H}^+]\), where \([\text{H}^+]\) denotes
the concentration of hydrogen ions in moles per litre --- this is
standard chemistry notation. Pure water has \([\text{H}^+] = 10^{-7}\)
mol/L, giving a pH of 7. A lemon juice sample has
\([\text{H}^+] = 10^{-2.4}\) mol/L. What is its pH, and how many times
more acidic is it than pure water?

\begin{Shaded}
\begin{Highlighting}[]
\NormalTok{//| echo: false}
\NormalTok{\{}
\NormalTok{  const el = makeStepperHTML(6, [}
\NormalTok{    \{ op: "Write as an equation", eq: "\textbackslash{}\textbackslash{}text\{pH\} = {-}\textbackslash{}\textbackslash{}log\_\{10\}(10\^{}\{{-}2.4\})", note: null \},}
\NormalTok{    \{ op: "Evaluate the log", eq: "\textbackslash{}\textbackslash{}text\{pH\} = {-}({-}2.4) = 2.4", note: "\textbackslash{}\textbackslash{}log\_\{10\}(10\^{}k) = k" \},}
\NormalTok{    \{ op: "Find the ratio of hydrogen ion concentrations", eq: "\textbackslash{}\textbackslash{}frac\{[\textbackslash{}\textbackslash{}text\{H\}\^{}+]\_\{\textbackslash{}\textbackslash{}text\{lemon\}\}\}\{[\textbackslash{}\textbackslash{}text\{H\}\^{}+]\_\{\textbackslash{}\textbackslash{}text\{water\}\}\} = \textbackslash{}\textbackslash{}frac\{10\^{}\{{-}2.4\}\}\{10\^{}\{{-}7\}\}", note: null \},}
\NormalTok{    \{ op: "Subtract exponents — dividing with the same base", eq: "= 10\^{}\{{-}2.4 {-} ({-}7)\} = 10\^{}\{4.6\}", note: null \},}
\NormalTok{    \{ op: "Evaluate", eq: "10\^{}\{4.6\} \textbackslash{}\textbackslash{}approx 39\textbackslash{}\textbackslash{},811", note: null \},}
\NormalTok{    \{ op: "Check", eq: "\textbackslash{}\textbackslash{}text\{pH of lemon\} = 2.4;\textbackslash{}\textbackslash{} 10\^{}\{7 {-} 2.4\} = 10\^{}\{4.6\} \textbackslash{}\textbackslash{}approx 39\textbackslash{}\textbackslash{},811 \textbackslash{}\textbackslash{}text\{ times more acidic than water\} \textbackslash{}\textbackslash{}checkmark", note: null \},}
\NormalTok{  ]);}
\NormalTok{  return el;}
\NormalTok{\}}
\end{Highlighting}
\end{Shaded}

\chapter{Complex numbers}\label{complex-numbers}

What lives where the number line doesn't go

\hfill\break

You've been solving quadratic equations. Most of them have real
solutions. But every now and then the discriminant --- the bit under the
square root --- comes out negative.

Try \(x^2 + 1 = 0\). Rearrange: \(x^2 = -1\). Take the square root:
\(x = \sqrt{-1}\).

That's impossible on the number line. There is no real number whose
square is negative. Positive times positive is positive. Negative times
negative is also positive. Nothing multiplied by itself gives \(-1\).

So for a long time, equations like this were said to have ``no
solution.''

That answer was unsatisfying, and mathematicians knew it. The fix was
the same fix that introduced negative numbers, fractions, and irrational
numbers: \emph{give the new thing a name, and work out its arithmetic}.

The name for \(\sqrt{-1}\) is \(i\).

Once you do that, something unexpected happens. The same number system
that lets you give \(x^2 + 1 = 0\) an answer also turns out to describe
oscillating voltages in electrical circuits, the quantum states of
particles, and the algorithm that your phone uses to process audio. None
of that was the motivation. It fell out of the arithmetic.

There's also this: \(e^{i\pi} + 1 = 0\). Five of the most important
numbers in mathematics --- \(e\), \(i\), \(\pi\), \(1\), \(0\) --- in
one equation. We'll see why by the end of this chapter.

\section{What the notation is
saying}\label{what-the-notation-is-saying-15}

\subsection{The imaginary unit}\label{the-imaginary-unit}

Define \(i\) by one rule:

\[i^2 = -1\]

That's the whole definition. \(i\) is the number whose square is \(-1\).
It's not a number on the real number line --- it's something new.

From this one rule, powers of \(i\) cycle in a pattern:

\[i^1 = i \qquad i^2 = -1 \qquad i^3 = -i \qquad i^4 = 1 \qquad i^5 = i \qquad \ldots\]

Every four powers, the cycle repeats. That's a hint that \(i\) has
something to do with rotation --- four quarter-turns bring you back to
where you started.

\subsection{Complex numbers}\label{complex-numbers-1}

A \textbf{complex number} is a number of the form:

\[z = a + bi\]

Read it as: ``the real part \(a\), plus \(b\) times \(i\).'' Here \(a\)
and \(b\) are ordinary real numbers. \(a\) is called the \textbf{real
part} of \(z\), written \(\operatorname{Re}(z)\). \(b\) is called the
\textbf{imaginary part}, written \(\operatorname{Im}(z)\).

Examples: \(3 + 4i\), \(-2 + i\), \(5\) (which is \(5 + 0i\)), \(3i\)
(which is \(0 + 3i\)).

Real numbers are a special case: they're complex numbers where
\(b = 0\).

\subsection{The complex plane}\label{the-complex-plane}

Every complex number \(a + bi\) corresponds to a point in a plane. Put
the real axis horizontal and the imaginary axis vertical. The number
\(3 + 4i\) is the point three units right and four units up.

This is the \textbf{Argand diagram}, or the \textbf{complex plane}. The
horizontal axis is the real line you already know. The vertical axis is
entirely imaginary numbers. Every complex number lives at exactly one
point in this plane.

\begin{Shaded}
\begin{Highlighting}[]
\NormalTok{//| echo: false}
\NormalTok{\{}
\NormalTok{  const div = document.createElement("div");}
\NormalTok{  div.style.cssText = "font{-}family:inherit; padding:0.5rem 0;";}

\NormalTok{  const controls = document.createElement("div");}
\NormalTok{  controls.style.cssText = "display:flex; gap:1.5rem; flex{-}wrap:wrap; margin{-}bottom:0.75rem;";}

\NormalTok{  function makeSlider(label, min, max, step, val) \{}
\NormalTok{    const wrap = document.createElement("div");}
\NormalTok{    const lbl = document.createElement("label");}
\NormalTok{    lbl.style.cssText = "font{-}size:0.85rem; color:\#374151; display:block; margin{-}bottom:0.2rem;";}
\NormalTok{    const span = document.createElement("span"); span.textContent = val;}
\NormalTok{    lbl.appendChild(document.createTextNode(label + " = ")); lbl.appendChild(span);}
\NormalTok{    const input = document.createElement("input");}
\NormalTok{    input.type = "range"; input.min = min; input.max = max; input.step = step; input.value = val;}
\NormalTok{    input.style.width = "150px";}
\NormalTok{    wrap.appendChild(lbl); wrap.appendChild(input);}
\NormalTok{    return \{wrap, input, span\};}
\NormalTok{  \}}

\NormalTok{  const aS = makeSlider("a (real part)", {-}5, 5, 0.5, 3);}
\NormalTok{  const bS = makeSlider("b (imaginary part)", {-}5, 5, 0.5, 2);}
\NormalTok{  controls.appendChild(aS.wrap);}
\NormalTok{  controls.appendChild(bS.wrap);}

\NormalTok{  const infoEl = document.createElement("div");}
\NormalTok{  infoEl.style.cssText = "font{-}size:0.9rem; color:\#1d4ed8; margin{-}bottom:0.5rem; min{-}height:1.4em;";}

\NormalTok{  const canvas = document.createElement("canvas");}
\NormalTok{  canvas.width = 340; canvas.height = 340;}
\NormalTok{  canvas.style.cssText = "border:1px solid \#e5e7eb; border{-}radius:6px; display:block; max{-}width:340px;";}

\NormalTok{  div.appendChild(controls); div.appendChild(infoEl); div.appendChild(canvas);}

\NormalTok{  function draw() \{}
\NormalTok{    const a = parseFloat(aS.input.value);}
\NormalTok{    const b = parseFloat(bS.input.value);}
\NormalTok{    aS.span.textContent = a.toFixed(1);}
\NormalTok{    bS.span.textContent = b.toFixed(1);}

\NormalTok{    const r = Math.sqrt(a*a + b*b);}
\NormalTok{    const theta = Math.atan2(b, a);}
\NormalTok{    const thetaDeg = (theta * 180 / Math.PI).toFixed(1);}
\NormalTok{    infoEl.textContent = "z = " + a.toFixed(1) + (b\textgreater{}=0?"+":"") + b.toFixed(1) + "i   |z| = " + r.toFixed(2) + "   arg(z) = " + thetaDeg + "°";}

\NormalTok{    const ctx = canvas.getContext("2d");}
\NormalTok{    const W = canvas.width, H = canvas.height;}
\NormalTok{    ctx.clearRect(0, 0, W, H);}

\NormalTok{    const range = 6;}
\NormalTok{    function px(x) \{ return (x + range) / (2*range) * W; \}}
\NormalTok{    function py(y) \{ return H {-} (y + range) / (2*range) * H; \}}

\NormalTok{    // Grid}
\NormalTok{    ctx.strokeStyle = "\#f3f4f6"; ctx.lineWidth = 1;}
\NormalTok{    for (let i = {-}5; i \textless{}= 5; i++) \{}
\NormalTok{      ctx.beginPath(); ctx.moveTo(px(i), 0); ctx.lineTo(px(i), H); ctx.stroke();}
\NormalTok{      ctx.beginPath(); ctx.moveTo(0, py(i)); ctx.lineTo(W, py(i)); ctx.stroke();}
\NormalTok{    \}}

\NormalTok{    // Axes}
\NormalTok{    ctx.strokeStyle = "\#9ca3af"; ctx.lineWidth = 1.5;}
\NormalTok{    ctx.beginPath(); ctx.moveTo(0, py(0)); ctx.lineTo(W, py(0)); ctx.stroke();}
\NormalTok{    ctx.beginPath(); ctx.moveTo(px(0), 0); ctx.lineTo(px(0), H); ctx.stroke();}
\NormalTok{    ctx.fillStyle = "\#6b7280"; ctx.font = "11px sans{-}serif";}
\NormalTok{    ctx.fillText("Re", W{-}20, py(0){-}5); ctx.fillText("Im", px(0)+5, 12);}

\NormalTok{    // Modulus line}
\NormalTok{    ctx.strokeStyle = "\#2563eb"; ctx.lineWidth = 2;}
\NormalTok{    ctx.beginPath(); ctx.moveTo(px(0), py(0)); ctx.lineTo(px(a), py(b)); ctx.stroke();}

\NormalTok{    // Angle arc}
\NormalTok{    if (r \textgreater{} 0.2) \{}
\NormalTok{      ctx.strokeStyle = "\#f59e0b"; ctx.lineWidth = 1.5;}
\NormalTok{      ctx.beginPath();}
\NormalTok{      ctx.arc(px(0), py(0), 30, {-}Math.atan2(b, a), 0, b \textgreater{} 0);}
\NormalTok{      ctx.stroke();}
\NormalTok{      ctx.fillStyle = "\#d97706"; ctx.font = "11px sans{-}serif";}
\NormalTok{      ctx.fillText("θ", px(0)+34, py(0){-}8);}
\NormalTok{    \}}

\NormalTok{    // Dashed projection lines}
\NormalTok{    ctx.strokeStyle = "\#cbd5e1"; ctx.lineWidth = 1; ctx.setLineDash([4,4]);}
\NormalTok{    ctx.beginPath(); ctx.moveTo(px(a), py(0)); ctx.lineTo(px(a), py(b)); ctx.stroke();}
\NormalTok{    ctx.beginPath(); ctx.moveTo(px(0), py(b)); ctx.lineTo(px(a), py(b)); ctx.stroke();}
\NormalTok{    ctx.setLineDash([]);}

\NormalTok{    // Point}
\NormalTok{    ctx.fillStyle = "\#ef4444";}
\NormalTok{    ctx.beginPath(); ctx.arc(px(a), py(b), 6, 0, 2*Math.PI); ctx.fill();}

\NormalTok{    // Labels}
\NormalTok{    ctx.fillStyle = "\#374151"; ctx.font = "bold 12px sans{-}serif";}
\NormalTok{    ctx.fillText("z = " + a.toFixed(1) + (b\textgreater{}=0?"+":"") + b.toFixed(1) + "i", px(a)+8, py(b){-}8);}
\NormalTok{    ctx.fillStyle = "\#2563eb"; ctx.font = "11px sans{-}serif";}
\NormalTok{    ctx.fillText("|z| = " + r.toFixed(2), px(a/2){-}5, py(b/2){-}10);}
\NormalTok{  \}}

\NormalTok{  aS.input.addEventListener("input", draw);}
\NormalTok{  bS.input.addEventListener("input", draw);}
\NormalTok{  draw();}
\NormalTok{  return div;}
\NormalTok{\}}
\end{Highlighting}
\end{Shaded}

\section{The method}\label{the-method-14}

\begin{tcolorbox}[enhanced jigsaw, opacityback=0, colback=white, breakable, colframe=quarto-callout-note-color-frame, arc=.35mm, left=2mm, toprule=.15mm, leftrule=.75mm, bottomrule=.15mm, rightrule=.15mm]
\begin{minipage}[t]{5.5mm}
\textcolor{quarto-callout-note-color}{\faInfo}
\end{minipage}%
\begin{minipage}[t]{\textwidth - 5.5mm}

\textbf{Complex numbers in 2D graphics}

Complex numbers appear in physics and engineering, but their geometry is
already useful in 2D graphics. Rotating a sprite on a screen by angle
\(\theta\) is the same as multiplying its position (written as a complex
number \(x + yi\)) by \(\cos\theta + i\sin\theta\). Game engines doing
this in bulk use the same arithmetic you are about to practise.

\end{minipage}%
\end{tcolorbox}

\subsection{Part 1: Arithmetic}\label{part-1-arithmetic}

\textbf{Adding and subtracting} complex numbers works the way you'd
expect: add the real parts together, add the imaginary parts together.

\[(3 + 4i) + (1 - 2i) = (3 + 1) + (4 - 2)i = 4 + 2i\]

\[(3 + 4i) - (1 - 2i) = (3 - 1) + (4 - (-2))i = 2 + 6i\]

\textbf{Multiplying} uses FOIL (expand like brackets), then replace
every \(i^2\) with \(-1\).

\[(3 + 4i)(1 - 2i) = 3 \cdot 1 + 3 \cdot (-2i) + 4i \cdot 1 + 4i \cdot (-2i)\]

\[= 3 - 6i + 4i - 8i^2\]

Now replace \(i^2 = -1\):

\[= 3 - 6i + 4i - 8(-1) = 3 - 6i + 4i + 8 = 11 - 2i\]

The key move: \(i^2\) always becomes \(-1\). That's the entire rule for
complex multiplication.

\textbf{Dividing} requires a trick. A complex denominator means the
fraction can't be split into real and imaginary parts yet --- you need a
real number on the bottom. Multiplying by the conjugate achieves this,
because \((a + bi)(a - bi) = a^2 + b^2\), which is always real. The tool
is the \textbf{complex conjugate}: if \(z = a + bi\), its conjugate is
\(\bar{z} = a - bi\) --- same real part, imaginary part negated.

Multiplying a complex number by its conjugate always gives a real
number:

\[(a + bi)(a - bi) = a^2 - (bi)^2 = a^2 - b^2 i^2 = a^2 + b^2\]

To divide, multiply numerator and denominator by the conjugate of the
denominator.

\[\frac{3 + 4i}{1 - 2i} \cdot \frac{1 + 2i}{1 + 2i} = \frac{(3 + 4i)(1 + 2i)}{(1 - 2i)(1 + 2i)}\]

Denominator: \((1)^2 + (2)^2 = 5\)

Numerator: \(3 + 6i + 4i + 8i^2 = 3 + 10i - 8 = -5 + 10i\)

Result: \(\dfrac{-5 + 10i}{5} = -1 + 2i\)

\subsection{Part 2: Modulus and
argument}\label{part-2-modulus-and-argument}

The \textbf{modulus} of \(z = a + bi\), written \(|z|\), is its distance
from the origin in the complex plane. By Pythagoras:

\[|z| = \sqrt{a^2 + b^2}\]

For \(z = 3 + 4i\): \(|z| = \sqrt{9 + 16} = \sqrt{25} = 5\).

The \textbf{argument} of \(z\), written \(\arg(z)\), is the angle the
line from the origin to \(z\) makes with the positive real axis.
Measured in radians (or degrees):

\[\arg(z) = \arctan\!\left(\frac{b}{a}\right)\]

Be careful with the quadrant. The \(\arctan\) function returns values in
the range \((-\pi/2, \pi/2)\), so you need to check whether \(z\) is in
the left half of the plane and adjust by \(\pm \pi\) if so.

For \(z = 3 + 4i\): \(\arg(z) = \arctan(4/3) \approx 0.927\) radians
(\(\approx 53.1°\)).

Together, modulus and argument give two ways to describe the same point:
Cartesian form \(a + bi\) and polar form \(r \angle \theta\).

\subsection{Part 3: Polar form and Euler's
formula}\label{part-3-polar-form-and-eulers-formula}

On the unit circle (where \(r = 1\)), the real part of a complex number
at angle \(\theta\) is \(\cos\theta\) and the imaginary part is
\(\sin\theta\) --- by the same definition of cosine and sine used in
trigonometry. So any point at distance \(r\) from the origin at angle
\(\theta\) is \(r\) times that unit-circle point.

A complex number with modulus \(r\) and argument \(\theta\) can be
written as:

\[z = r(\cos\theta + i\sin\theta)\]

This is \textbf{polar form}. Read it as: go distance \(r\) from the
origin, at angle \(\theta\) from the positive real axis.

Now for the extraordinary fact. \textbf{Euler's formula} says:

\[e^{i\theta} = \cos\theta + i\sin\theta\]

This connects the exponential function --- which you know from growth
and decay --- to the trigonometric functions. The proof belongs to
calculus (it comes from comparing the Taylor series of \(e^x\),
\(\cos x\), and \(\sin x\)), but you can use the result here.

With Euler's formula, polar form becomes simply:

\[z = re^{i\theta}\]

This is often the cleanest way to write a complex number when you're
multiplying, dividing, or raising to a power.

The famous identity \(e^{i\pi} + 1 = 0\) is just the special case
\(\theta = \pi\): \(e^{i\pi} = \cos\pi + i\sin\pi = -1 + 0i = -1\).

In particular, \(i = e^{i\pi/2}\) (since \(\cos(\pi/2) = 0\) and
\(\sin(\pi/2) = 1\)). Because multiplication adds arguments, multiplying
any complex number by \(i\) adds \(\pi/2\) to its angle --- a 90°
anticlockwise rotation in the plane.

\begin{Shaded}
\begin{Highlighting}[]
\NormalTok{//| echo: false}
\NormalTok{\{}
\NormalTok{  const div = document.createElement("div");}
\NormalTok{  div.style.cssText = "font{-}family:inherit; padding:0.5rem 0;";}

\NormalTok{  const controls = document.createElement("div");}
\NormalTok{  controls.style.cssText = "display:flex; gap:1.5rem; flex{-}wrap:wrap; margin{-}bottom:0.75rem;";}

\NormalTok{  function makeSlider(label, min, max, step, val) \{}
\NormalTok{    const wrap = document.createElement("div");}
\NormalTok{    const lbl = document.createElement("label");}
\NormalTok{    lbl.style.cssText = "font{-}size:0.85rem; color:\#374151; display:block; margin{-}bottom:0.2rem;";}
\NormalTok{    const span = document.createElement("span"); span.textContent = val;}
\NormalTok{    lbl.appendChild(document.createTextNode(label + " = ")); lbl.appendChild(span);}
\NormalTok{    const input = document.createElement("input");}
\NormalTok{    input.type = "range"; input.min = min; input.max = max; input.step = step; input.value = val;}
\NormalTok{    input.style.width = "140px";}
\NormalTok{    wrap.appendChild(lbl); wrap.appendChild(input);}
\NormalTok{    return \{wrap, input, span\};}
\NormalTok{  \}}

\NormalTok{  const r1S = makeSlider("r₁ (|z₁|)", 0.5, 3, 0.5, 1.5);}
\NormalTok{  const t1S = makeSlider("θ₁ (arg z₁, °)", 0, 360, 15, 45);}
\NormalTok{  const r2S = makeSlider("r₂ (|z₂|)", 0.5, 3, 0.5, 1);}
\NormalTok{  const t2S = makeSlider("θ₂ (arg z₂, °)", 0, 360, 15, 30);}
\NormalTok{  controls.appendChild(r1S.wrap);}
\NormalTok{  controls.appendChild(t1S.wrap);}
\NormalTok{  controls.appendChild(r2S.wrap);}
\NormalTok{  controls.appendChild(t2S.wrap);}

\NormalTok{  const infoEl = document.createElement("div");}
\NormalTok{  infoEl.style.cssText = "font{-}size:0.85rem; color:\#374151; margin{-}bottom:0.5rem; min{-}height:2em; background:\#f8fafc; border{-}radius:6px; padding:0.4rem 0.75rem;";}

\NormalTok{  const canvas = document.createElement("canvas");}
\NormalTok{  canvas.width = 360; canvas.height = 360;}
\NormalTok{  canvas.style.cssText = "border:1px solid \#e5e7eb; border{-}radius:6px; display:block; max{-}width:360px;";}

\NormalTok{  div.appendChild(controls); div.appendChild(infoEl); div.appendChild(canvas);}

\NormalTok{  function draw() \{}
\NormalTok{    const r1 = parseFloat(r1S.input.value);}
\NormalTok{    const t1d = parseFloat(t1S.input.value);}
\NormalTok{    const r2 = parseFloat(r2S.input.value);}
\NormalTok{    const t2d = parseFloat(t2S.input.value);}
\NormalTok{    r1S.span.textContent = r1S.input.value;}
\NormalTok{    t1S.span.textContent = t1S.input.value + "°";}
\NormalTok{    r2S.span.textContent = r2S.input.value;}
\NormalTok{    t2S.span.textContent = t2S.input.value + "°";}

\NormalTok{    const t1 = t1d * Math.PI / 180;}
\NormalTok{    const t2 = t2d * Math.PI / 180;}
\NormalTok{    const z1 = \{r: r1, a: r1*Math.cos(t1), b: r1*Math.sin(t1)\};}
\NormalTok{    const z2 = \{r: r2, a: r2*Math.cos(t2), b: r2*Math.sin(t2)\};}
\NormalTok{    const rp = r1*r2;}
\NormalTok{    const tp = t1+t2;}
\NormalTok{    const prod = \{a: rp*Math.cos(tp), b: rp*Math.sin(tp)\};}

\NormalTok{    infoEl.innerHTML =}
\NormalTok{      "z₁ = " + z1.a.toFixed(2) + (z1.b\textgreater{}=0?"+":"") + z1.b.toFixed(2) + "i  |z₁|=" + r1 + "  arg=" + t1d + "°\textless{}br\textgreater{}" +}
\NormalTok{      "z₂ = " + z2.a.toFixed(2) + (z2.b\textgreater{}=0?"+":"") + z2.b.toFixed(2) + "i  |z₂|=" + r2 + "  arg=" + t2d + "°\textless{}br\textgreater{}" +}
\NormalTok{      "z₁·z₂ = " + prod.a.toFixed(2) + (prod.b\textgreater{}=0?"+":"") + prod.b.toFixed(2) + "i  |z₁·z₂|=" + rp.toFixed(2) + "  arg=" + (t1d+t2d) + "°";}

\NormalTok{    const ctx = canvas.getContext("2d");}
\NormalTok{    const W = canvas.width, H = canvas.height;}
\NormalTok{    ctx.clearRect(0, 0, W, H);}
\NormalTok{    const range = 5;}
\NormalTok{    function px(x) \{ return (x+range)/(2*range)*W; \}}
\NormalTok{    function py(y) \{ return H{-}(y+range)/(2*range)*H; \}}

\NormalTok{    ctx.strokeStyle = "\#9ca3af"; ctx.lineWidth = 1;}
\NormalTok{    ctx.beginPath(); ctx.moveTo(0, py(0)); ctx.lineTo(W, py(0)); ctx.stroke();}
\NormalTok{    ctx.beginPath(); ctx.moveTo(px(0), 0); ctx.lineTo(px(0), H); ctx.stroke();}
\NormalTok{    ctx.fillStyle = "\#9ca3af"; ctx.font = "11px sans{-}serif";}
\NormalTok{    ctx.fillText("Re", W{-}20, py(0){-}4); ctx.fillText("Im", px(0)+4, 12);}

\NormalTok{    const pts = [}
\NormalTok{      \{x: z1.a, y: z1.b, color: "\#2563eb", label: "z₁"\},}
\NormalTok{      \{x: z2.a, y: z2.b, color: "\#059669", label: "z₂"\},}
\NormalTok{      \{x: prod.a, y: prod.b, color: "\#ef4444", label: "z₁·z₂"\},}
\NormalTok{    ];}
\NormalTok{    pts.forEach(p =\textgreater{} \{}
\NormalTok{      ctx.strokeStyle = p.color; ctx.lineWidth = 2;}
\NormalTok{      ctx.beginPath(); ctx.moveTo(px(0), py(0)); ctx.lineTo(px(p.x), py(p.y)); ctx.stroke();}
\NormalTok{      ctx.fillStyle = p.color;}
\NormalTok{      ctx.beginPath(); ctx.arc(px(p.x), py(p.y), 5, 0, 2*Math.PI); ctx.fill();}
\NormalTok{      ctx.fillStyle = "\#374151"; ctx.font = "bold 11px sans{-}serif";}
\NormalTok{      ctx.fillText(p.label, px(p.x)+7, py(p.y){-}5);}
\NormalTok{    \});}

\NormalTok{    // Angle arcs}
\NormalTok{    [t1, t2, tp].forEach((t, i) =\textgreater{} \{}
\NormalTok{      const colors = ["\#2563eb","\#059669","\#ef4444"];}
\NormalTok{      ctx.strokeStyle = colors[i]; ctx.lineWidth = 1.5;}
\NormalTok{      const arcR = 25 + i*10;}
\NormalTok{      ctx.beginPath();}
\NormalTok{      ctx.arc(px(0), py(0), arcR, {-}t, 0, t \textless{} 0);}
\NormalTok{      ctx.stroke();}
\NormalTok{    \});}
\NormalTok{  \}}

\NormalTok{  r1S.input.addEventListener("input", draw);}
\NormalTok{  t1S.input.addEventListener("input", draw);}
\NormalTok{  r2S.input.addEventListener("input", draw);}
\NormalTok{  t2S.input.addEventListener("input", draw);}
\NormalTok{  draw();}
\NormalTok{  return div;}
\NormalTok{\}}
\end{Highlighting}
\end{Shaded}

\subsection{Part 4: De Moivre's
theorem}\label{part-4-de-moivres-theorem}

Because \(z = re^{i\theta}\), raising \(z\) to the power \(n\) is clean:

\[z^n = (re^{i\theta})^n = r^n e^{in\theta} = r^n(\cos n\theta + i\sin n\theta)\]

This is \textbf{de Moivre's theorem}. To find the \(n\)th power of a
complex number: raise the modulus to the \(n\)th power, multiply the
argument by \(n\).

Example: find \((1 + i)^4\).

First, convert to polar form. \(|1 + i| = \sqrt{2}\). For the argument:
\(\arg(1 + i) = \arctan(1/1) = \arctan(1) = \pi/4\).

So \(1 + i = \sqrt{2}\,e^{i\pi/4}\).

Then:

\[(1 + i)^4 = (\sqrt{2})^4 \cdot e^{i \cdot 4 \cdot \pi/4} = 4 \cdot e^{i\pi} = 4 \cdot (-1) = -4\]

Check by direct multiplication: \((1+i)^2 = 1 + 2i - 1 = 2i\). Then
\((2i)^2 = 4i^2 = -4\). Matches.

\begin{tcolorbox}[enhanced jigsaw, opacityback=0, colback=white, breakable, colframe=quarto-callout-note-color-frame, arc=.35mm, left=2mm, toprule=.15mm, leftrule=.75mm, bottomrule=.15mm, rightrule=.15mm]
\begin{minipage}[t]{5.5mm}
\textcolor{quarto-callout-note-color}{\faInfo}
\end{minipage}%
\begin{minipage}[t]{\textwidth - 5.5mm}

\textbf{Why this works}

Multiplying any complex number by \(i\) adds \(\pi/2\) to its angle ---
a 90° anticlockwise rotation in the complex plane. That's why
\(i^4 = 1\): four quarter-turns bring you back.

Complex multiplication in general is a \emph{rotation and scaling}
operation. When you multiply \(z_1 \cdot z_2\), the moduli multiply
(scaling) and the arguments add (rotation). This is exactly why complex
numbers are the natural language for anything involving oscillation,
phase, and rotation --- from AC circuits to quantum wavefunctions to
signal processing. The arithmetic of complex numbers is the arithmetic
of rotations.

\end{minipage}%
\end{tcolorbox}

\section{Worked examples}\label{worked-examples-14}

\subsection{Example 1 --- Engineering: AC circuit
impedance}\label{example-1-engineering-ac-circuit-impedance}

An AC circuit has a resistor and a capacitor in series. The resistor has
impedance \(Z_R = 50\,\Omega\) (purely real). The capacitor has
impedance \(Z_C = -30i\,\Omega\) (purely imaginary, negative because
capacitors shift the phase of current).

The total series impedance is:

\[Z = Z_R + Z_C = 50 + (-30i) = 50 - 30i\,\Omega\]

The magnitude of the impedance --- how much the circuit resists current
overall --- is the modulus:

\[|Z| = \sqrt{50^2 + 30^2} = \sqrt{2500 + 900} = \sqrt{3400} \approx 58.3\,\Omega\]

The phase angle --- how much voltage leads current --- is the argument:

\[\arg(Z) = \arctan\!\left(\frac{-30}{50}\right) \approx -30.96°\]

The negative angle means the voltage lags behind the current by about
31°, which is the expected behaviour for a capacitive circuit.

\begin{center}\rule{0.5\linewidth}{0.5pt}\end{center}

\subsection{Example 2 --- Physics: Quantum wavefunction
amplitude}\label{example-2-physics-quantum-wavefunction-amplitude}

In quantum mechanics, a particle in a particular state has a probability
amplitude that is a complex number. The probability of measuring that
state is the square of the modulus of the amplitude.

Suppose the amplitude is \(\psi = \frac{3}{5} + \frac{4}{5}i\).

The probability is:

\[|\psi|^2 = \left(\frac{3}{5}\right)^2 + \left(\frac{4}{5}\right)^2 = \frac{9}{25} + \frac{16}{25} = \frac{25}{25} = 1\]

A probability of 1 means the particle is certain to be found in this
state. The amplitude is normalised --- its modulus is 1. This is a
constraint quantum mechanics imposes on all valid probability
amplitudes: \(|\psi| = 1\).

\begin{center}\rule{0.5\linewidth}{0.5pt}\end{center}

\subsection{Example 3 --- Computing: FFT
butterfly}\label{example-3-computing-fft-butterfly}

The Fast Fourier Transform (FFT) is an algorithm that decomposes a
signal into its frequency components. Its key step multiplies a complex
number by a `twiddle factor' \(W\) --- which is just a complex number on
the unit circle --- to rotate it by a precise angle.

The FFT processes a signal by repeatedly multiplying values by twiddle
factors of the form \(W = e^{-2\pi i k/N}\) for various integers \(k\)
and \(N\).

For an 8-point FFT with \(k = 1\), \(N = 8\):

\[W = e^{-2\pi i / 8} = e^{-i\pi/4} = \cos\!\left(-\frac{\pi}{4}\right) + i\sin\!\left(-\frac{\pi}{4}\right) = \frac{1}{\sqrt{2}} - \frac{i}{\sqrt{2}}\]

One butterfly step multiplies a signal value \(X = 3 + 2i\) by this
twiddle factor:

\[X \cdot W = (3 + 2i)\!\left(\frac{1}{\sqrt{2}} - \frac{i}{\sqrt{2}}\right)\]

\[= \frac{1}{\sqrt{2}}\bigl(3 + 2i\bigr)(1 - i) = \frac{1}{\sqrt{2}}(3 - 3i + 2i - 2i^2)\]

\[= \frac{1}{\sqrt{2}}(3 - i + 2) = \frac{1}{\sqrt{2}}(5 - i) = \frac{5}{\sqrt{2}} - \frac{i}{\sqrt{2}} \approx 3.54 - 0.71i\]

In polar terms: multiplication by \(W\) kept the modulus the same (since
\(|W| = 1\)) and rotated the point by \(-45°\). The FFT is, at its core,
a very efficient way to apply many such rotations simultaneously.

\section{Where this goes}\label{where-this-goes-15}

\textbf{Polynomials (Vol 3 Ch 3)} --- every polynomial with real
coefficients that has no real roots has complex roots. The quadratic
formula \(x = \frac{-b \pm \sqrt{b^2 - 4ac}}{2a}\) gives complex answers
when \(b^2 - 4ac < 0\). You can now say exactly what those answers are
and verify them by substituting back in. The discriminant is no longer a
dead end.

\textbf{Complex analysis (Vol 7)} extends calculus to complex-valued
functions. It turns out that calculus in the complex plane is
extraordinarily powerful --- more constrained than real calculus in ways
that produce surprising results, including a method for evaluating real
integrals that seem impossible by ordinary means. The Cauchy integral
theorem, residues, and conformal mapping all live here. Euler's formula
is the entry point.

\textbf{Differential calculus (Vol 5)} --- Euler's formula
\(e^{i\theta} = \cos\theta + i\sin\theta\) is proved by comparing the
Taylor series expansions of each side. That proof belongs to calculus,
but once you have it, it explains why \(e\), \(i\), \(\pi\), \(1\), and
\(0\) are connected: they're all special values of the same exponential
function operating on the complex plane.

\begin{tcolorbox}[enhanced jigsaw, opacityback=0, colback=white, breakable, colframe=quarto-callout-tip-color-frame, arc=.35mm, left=2mm, toprule=.15mm, leftrule=.75mm, bottomrule=.15mm, rightrule=.15mm]
\begin{minipage}[t]{5.5mm}
\textcolor{quarto-callout-tip-color}{\faLightbulb}
\end{minipage}%
\begin{minipage}[t]{\textwidth - 5.5mm}

\textbf{Where this shows up}

\begin{itemize}
\item
  An electrical engineer analysing an AC circuit writes impedance as a
  complex number, computes its modulus to find the magnitude, and its
  argument to find the phase shift. The letter \(j\) is used instead of
  \(i\) in electrical engineering (because \(i\) is already taken for
  current), but the mathematics is identical.
\item
  A physicist writing down the Schrödinger equation is working with a
  complex-valued wavefunction. The probability of measurement outcomes
  is the squared modulus of that function.
\item
  A software engineer implementing a digital audio filter uses the
  Discrete Fourier Transform --- which is an algorithm entirely built on
  complex multiplication. The FFT makes this computation fast enough to
  run in real time.
\item
  A control engineer analysing whether a feedback system is stable
  places poles in the complex plane. Poles in the left half-plane mean
  the system is stable; poles in the right half-plane mean it oscillates
  out of control.
\end{itemize}

\end{minipage}%
\end{tcolorbox}

\section{Exercises}\label{exercises-15}

These are puzzles. Each one has a clean answer. The interesting part is
knowing which form of a complex number to use for which operation.

\begin{Shaded}
\begin{Highlighting}[]
\NormalTok{//| echo: false}
\NormalTok{katex = require("katex@0.16")}
\end{Highlighting}
\end{Shaded}

\begin{Shaded}
\begin{Highlighting}[]
\NormalTok{//| echo: false}
\NormalTok{function makeStepperHTML(exerciseNum, steps) \{}
\NormalTok{  let current = 0;}
\NormalTok{  const totalSteps = steps.length;}
\NormalTok{  const container = document.createElement("div");}
\NormalTok{  container.style.cssText = "border:1px solid \#e5e7eb; border{-}radius:8px; padding:1rem 1.25rem; margin:0.75rem 0; font{-}family:inherit;";}
\NormalTok{  const stepsDiv = document.createElement("div");}
\NormalTok{  stepsDiv.style.marginBottom = "0.75rem";}
\NormalTok{  function renderSteps() \{}
\NormalTok{    stepsDiv.innerHTML = "";}
\NormalTok{    for (let i = 0; i \textless{} current; i++) \{}
\NormalTok{      const s = steps[i];}
\NormalTok{      const row = document.createElement("div");}
\NormalTok{      row.style.cssText = "display:grid; grid{-}template{-}columns:200px 1fr; gap:0.35rem 1rem; align{-}items:baseline; padding:0.35rem 0; border{-}top:1px solid \#f3f4f6;";}
\NormalTok{      const opEl = document.createElement("span");}
\NormalTok{      opEl.style.cssText = "font{-}size:0.8rem; color:\#6b7280; font{-}style:italic;";}
\NormalTok{      opEl.textContent = s.op;}
\NormalTok{      const eqEl = document.createElement("span");}
\NormalTok{      eqEl.innerHTML = katex.renderToString(s.eq, \{throwOnError: false, displayMode: false\});}
\NormalTok{      row.appendChild(opEl);}
\NormalTok{      row.appendChild(eqEl);}
\NormalTok{      container.appendChild(row);}
\NormalTok{      if (s.note) \{}
\NormalTok{        const noteEl = document.createElement("p");}
\NormalTok{        noteEl.style.cssText = "margin:0.2rem 0 0.5rem 200px; font{-}size:0.82rem; color:\#374151;";}
\NormalTok{        noteEl.textContent = s.note;}
\NormalTok{        container.appendChild(noteEl);}
\NormalTok{      \}}
\NormalTok{    \}}
\NormalTok{    if (current === totalSteps) \{}
\NormalTok{      const done = document.createElement("p");}
\NormalTok{      done.style.cssText = "margin{-}top:0.5rem; font{-}size:0.85rem; color:\#059669; font{-}weight:600;";}
\NormalTok{      done.textContent = "✓ Solution complete";}
\NormalTok{      stepsDiv.appendChild(done);}
\NormalTok{    \}}
\NormalTok{  \}}
\NormalTok{  const controls = document.createElement("div");}
\NormalTok{  controls.style.cssText = "display:flex; align{-}items:center; gap:0.75rem; margin{-}top:0.5rem;";}
\NormalTok{  const nextBtn = document.createElement("button");}
\NormalTok{  nextBtn.textContent = "Next step →";}
\NormalTok{  nextBtn.style.cssText = "padding:0.35rem 0.85rem; background:\#2563eb; color:\#fff; border:none; border{-}radius:6px; cursor:pointer; font{-}size:0.85rem;";}
\NormalTok{  const resetBtn = document.createElement("button");}
\NormalTok{  resetBtn.textContent = "Reset";}
\NormalTok{  resetBtn.style.cssText = "padding:0.35rem 0.85rem; background:\#f3f4f6; color:\#374151; border:1px solid \#d1d5db; border{-}radius:6px; cursor:pointer; font{-}size:0.85rem;";}
\NormalTok{  const counter = document.createElement("span");}
\NormalTok{  counter.style.cssText = "font{-}size:0.8rem; color:\#6b7280;";}
\NormalTok{  nextBtn.onclick = () =\textgreater{} \{ if (current \textless{} totalSteps) \{ current++; renderSteps(); counter.textContent = current + " / " + totalSteps; if (current === totalSteps) nextBtn.disabled = true; \} \};}
\NormalTok{  resetBtn.onclick = () =\textgreater{} \{ current = 0; nextBtn.disabled = false; renderSteps(); counter.textContent = "0 / " + totalSteps; \};}
\NormalTok{  counter.textContent = "0 / " + totalSteps;}
\NormalTok{  controls.appendChild(nextBtn);}
\NormalTok{  controls.appendChild(resetBtn);}
\NormalTok{  controls.appendChild(counter);}
\NormalTok{  container.appendChild(stepsDiv);}
\NormalTok{  container.appendChild(controls);}
\NormalTok{  return container;}
\NormalTok{\}}
\end{Highlighting}
\end{Shaded}

\textbf{1.} Compute \((2 + 3i)(4 - i)\).

\begin{Shaded}
\begin{Highlighting}[]
\NormalTok{//| echo: false}
\NormalTok{\{}
\NormalTok{  const el = makeStepperHTML(1, [}
\NormalTok{    \{ op: "Expand using FOIL", eq: "(2 + 3i)(4 {-} i) = 2 \textbackslash{}\textbackslash{}cdot 4 + 2 \textbackslash{}\textbackslash{}cdot ({-}i) + 3i \textbackslash{}\textbackslash{}cdot 4 + 3i \textbackslash{}\textbackslash{}cdot ({-}i)", note: null \},}
\NormalTok{    \{ op: "Collect terms", eq: "= 8 {-} 2i + 12i {-} 3i\^{}2", note: null \},}
\NormalTok{    \{ op: "Replace i² = −1", eq: "= 8 {-} 2i + 12i {-} 3({-}1) = 8 + 3 + ({-}2 + 12)i", note: "Every i² becomes −1; this is the only rule that makes complex arithmetic close." \},}
\NormalTok{    \{ op: "Write the answer", eq: "= 11 + 10i", note: null \},}
\NormalTok{    \{ op: "Check", eq: "(2 + 3i)(4 {-} i) = 11 + 10i \textbackslash{}\textbackslash{}checkmark \textbackslash{}\textbackslash{}text\{ — verify by expanding again\}", note: null \},}
\NormalTok{  ]);}
\NormalTok{  return el;}
\NormalTok{\}}
\end{Highlighting}
\end{Shaded}

\textbf{2.} Find the modulus and argument of \(z = -1 + \sqrt{3}\,i\).

\begin{Shaded}
\begin{Highlighting}[]
\NormalTok{//| echo: false}
\NormalTok{\{}
\NormalTok{  const el = makeStepperHTML(2, [}
\NormalTok{    \{ op: "Identify real and imaginary parts", eq: "a = {-}1,\textbackslash{}\textbackslash{}quad b = \textbackslash{}\textbackslash{}sqrt\{3\}", note: null \},}
\NormalTok{    \{ op: "Compute the modulus", eq: "|z| = \textbackslash{}\textbackslash{}sqrt\{({-}1)\^{}2 + (\textbackslash{}\textbackslash{}sqrt\{3\})\^{}2\} = \textbackslash{}\textbackslash{}sqrt\{1 + 3\} = \textbackslash{}\textbackslash{}sqrt\{4\} = 2", note: null \},}
\NormalTok{    \{ op: "Compute arctan(b/a)", eq: "\textbackslash{}\textbackslash{}arctan\textbackslash{}\textbackslash{}!\textbackslash{}\textbackslash{}left(\textbackslash{}\textbackslash{}frac\{\textbackslash{}\textbackslash{}sqrt\{3\}\}\{{-}1\}\textbackslash{}\textbackslash{}right) = \textbackslash{}\textbackslash{}arctan({-}\textbackslash{}\textbackslash{}sqrt\{3\}) = {-}\textbackslash{}\textbackslash{}frac\{\textbackslash{}\textbackslash{}pi\}\{3\}", note: "arctan gives a value in (−π/2, π/2). We must check the quadrant." \},}
\NormalTok{    \{ op: "Adjust for correct quadrant", eq: "z \textbackslash{}\textbackslash{}text\{ is in the second quadrant (a \textless{} 0, b \textgreater{} 0), so \} \textbackslash{}\textbackslash{}arg(z) = {-}\textbackslash{}\textbackslash{}frac\{\textbackslash{}\textbackslash{}pi\}\{3\} + \textbackslash{}\textbackslash{}pi = \textbackslash{}\textbackslash{}frac\{2\textbackslash{}\textbackslash{}pi\}\{3\}", note: "Add π when the real part is negative." \},}
\NormalTok{    \{ op: "State the answer", eq: "|z| = 2,\textbackslash{}\textbackslash{}quad \textbackslash{}\textbackslash{}arg(z) = \textbackslash{}\textbackslash{}frac\{2\textbackslash{}\textbackslash{}pi\}\{3\} \textbackslash{}\textbackslash{}approx 120°", note: null \},}
\NormalTok{    \{ op: "Check", eq: "2\textbackslash{}\textbackslash{}!\textbackslash{}\textbackslash{}left(\textbackslash{}\textbackslash{}cos\textbackslash{}\textbackslash{}frac\{2\textbackslash{}\textbackslash{}pi\}\{3\} + i\textbackslash{}\textbackslash{}sin\textbackslash{}\textbackslash{}frac\{2\textbackslash{}\textbackslash{}pi\}\{3\}\textbackslash{}\textbackslash{}right) = 2\textbackslash{}\textbackslash{}!\textbackslash{}\textbackslash{}left({-}\textbackslash{}\textbackslash{}frac\{1\}\{2\} + \textbackslash{}\textbackslash{}frac\{\textbackslash{}\textbackslash{}sqrt\{3\}\}\{2\}i\textbackslash{}\textbackslash{}right) = {-}1 + \textbackslash{}\textbackslash{}sqrt\{3\}\textbackslash{}\textbackslash{},i \textbackslash{}\textbackslash{}checkmark", note: null \},}
\NormalTok{  ]);}
\NormalTok{  return el;}
\NormalTok{\}}
\end{Highlighting}
\end{Shaded}

\textbf{3.} Convert \(z = 4e^{i\pi/6}\) to Cartesian form \(a + bi\).

\begin{Shaded}
\begin{Highlighting}[]
\NormalTok{//| echo: false}
\NormalTok{\{}
\NormalTok{  const el = makeStepperHTML(3, [}
\NormalTok{    \{ op: "Apply Euler\textquotesingle{}s formula", eq: "z = 4\textbackslash{}\textbackslash{}!\textbackslash{}\textbackslash{}left(\textbackslash{}\textbackslash{}cos\textbackslash{}\textbackslash{}frac\{\textbackslash{}\textbackslash{}pi\}\{6\} + i\textbackslash{}\textbackslash{}sin\textbackslash{}\textbackslash{}frac\{\textbackslash{}\textbackslash{}pi\}\{6\}\textbackslash{}\textbackslash{}right)", note: "e\^{}\{iθ\} = cos θ + i sin θ — this is the definition of the exponential on complex numbers." \},}
\NormalTok{    \{ op: "Evaluate the trig values", eq: "\textbackslash{}\textbackslash{}cos\textbackslash{}\textbackslash{}frac\{\textbackslash{}\textbackslash{}pi\}\{6\} = \textbackslash{}\textbackslash{}frac\{\textbackslash{}\textbackslash{}sqrt\{3\}\}\{2\},\textbackslash{}\textbackslash{}quad \textbackslash{}\textbackslash{}sin\textbackslash{}\textbackslash{}frac\{\textbackslash{}\textbackslash{}pi\}\{6\} = \textbackslash{}\textbackslash{}frac\{1\}\{2\}", note: null \},}
\NormalTok{    \{ op: "Multiply through by 4", eq: "z = 4 \textbackslash{}\textbackslash{}cdot \textbackslash{}\textbackslash{}frac\{\textbackslash{}\textbackslash{}sqrt\{3\}\}\{2\} + 4 \textbackslash{}\textbackslash{}cdot \textbackslash{}\textbackslash{}frac\{1\}\{2\}\textbackslash{}\textbackslash{},i = 2\textbackslash{}\textbackslash{}sqrt\{3\} + 2i", note: null \},}
\NormalTok{    \{ op: "Check", eq: "|2\textbackslash{}\textbackslash{}sqrt\{3\} + 2i| = \textbackslash{}\textbackslash{}sqrt\{(2\textbackslash{}\textbackslash{}sqrt\{3\})\^{}2 + 2\^{}2\} = \textbackslash{}\textbackslash{}sqrt\{12 + 4\} = \textbackslash{}\textbackslash{}sqrt\{16\} = 4 \textbackslash{}\textbackslash{}checkmark", note: null \},}
\NormalTok{  ]);}
\NormalTok{  return el;}
\NormalTok{\}}
\end{Highlighting}
\end{Shaded}

\textbf{4.} Compute \((-1 + i)^6\) using de Moivre's theorem.

\begin{Shaded}
\begin{Highlighting}[]
\NormalTok{//| echo: false}
\NormalTok{\{}
\NormalTok{  const el = makeStepperHTML(4, [}
\NormalTok{    \{ op: "Convert to polar form", eq: "|{-}1 + i| = \textbackslash{}\textbackslash{}sqrt\{1 + 1\} = \textbackslash{}\textbackslash{}sqrt\{2\},\textbackslash{}\textbackslash{}quad \textbackslash{}\textbackslash{}arg({-}1 + i) = \textbackslash{}\textbackslash{}frac\{3\textbackslash{}\textbackslash{}pi\}\{4\}", note: "Real part negative, imaginary part positive: second quadrant." \},}
\NormalTok{    \{ op: "Write in exponential form", eq: "{-}1 + i = \textbackslash{}\textbackslash{}sqrt\{2\}\textbackslash{}\textbackslash{},e\^{}\{i \textbackslash{}\textbackslash{}cdot 3\textbackslash{}\textbackslash{}pi/4\}", note: null \},}
\NormalTok{    \{ op: "Apply de Moivre: raise modulus to n, multiply argument by n", eq: "({-}1 + i)\^{}6 = (\textbackslash{}\textbackslash{}sqrt\{2\})\^{}6 \textbackslash{}\textbackslash{}cdot e\^{}\{i \textbackslash{}\textbackslash{}cdot 6 \textbackslash{}\textbackslash{}cdot 3\textbackslash{}\textbackslash{}pi/4\} = 8\textbackslash{}\textbackslash{},e\^{}\{i \textbackslash{}\textbackslash{}cdot 9\textbackslash{}\textbackslash{}pi/2\}", note: "(\textbackslash{}\textbackslash{}sqrt\{2\})\^{}6 = 2\^{}\{6/2\} = 2\^{}3 = 8." \},}
\NormalTok{    \{ op: "Reduce the angle mod 2π", eq: "\textbackslash{}\textbackslash{}frac\{9\textbackslash{}\textbackslash{}pi\}\{2\} = 4\textbackslash{}\textbackslash{}pi + \textbackslash{}\textbackslash{}frac\{\textbackslash{}\textbackslash{}pi\}\{2\},\textbackslash{}\textbackslash{}quad \textbackslash{}\textbackslash{}text\{so \} e\^{}\{i \textbackslash{}\textbackslash{}cdot 9\textbackslash{}\textbackslash{}pi/2\} = e\^{}\{i\textbackslash{}\textbackslash{}pi/2\} = i", note: "Angles differing by 2π are identical in the complex plane." \},}
\NormalTok{    \{ op: "Write the answer", eq: "({-}1 + i)\^{}6 = 8i", note: null \},}
\NormalTok{    \{ op: "Check", eq: "({-}1+i)\^{}2 = 1 {-} 2i + i\^{}2 = {-}2i;\textbackslash{}\textbackslash{}quad ({-}2i)\^{}3 = {-}8i\^{}3 = {-}8({-}i) = 8i \textbackslash{}\textbackslash{}checkmark", note: null \},}
\NormalTok{  ]);}
\NormalTok{  return el;}
\NormalTok{\}}
\end{Highlighting}
\end{Shaded}

\textbf{5.} Divide \(\dfrac{2 + i}{3 - 4i}\).

\begin{Shaded}
\begin{Highlighting}[]
\NormalTok{//| echo: false}
\NormalTok{\{}
\NormalTok{  const el = makeStepperHTML(5, [}
\NormalTok{    \{ op: "Multiply numerator and denominator by the conjugate of the denominator", eq: "\textbackslash{}\textbackslash{}frac\{2 + i\}\{3 {-} 4i\} \textbackslash{}\textbackslash{}cdot \textbackslash{}\textbackslash{}frac\{3 + 4i\}\{3 + 4i\}", note: "The conjugate of 3 − 4i is 3 + 4i. This clears i from the denominator." \},}
\NormalTok{    \{ op: "Expand the denominator", eq: "(3 {-} 4i)(3 + 4i) = 9 + 16 = 25", note: "(a − bi)(a + bi) = a² + b² — always real." \},}
\NormalTok{    \{ op: "Expand the numerator", eq: "(2 + i)(3 + 4i) = 6 + 8i + 3i + 4i\^{}2 = 6 + 11i {-} 4 = 2 + 11i", note: null \},}
\NormalTok{    \{ op: "Write the answer", eq: "\textbackslash{}\textbackslash{}frac\{2 + 11i\}\{25\} = \textbackslash{}\textbackslash{}frac\{2\}\{25\} + \textbackslash{}\textbackslash{}frac\{11\}\{25\}i", note: null \},}
\NormalTok{    \{ op: "Check", eq: "\textbackslash{}\textbackslash{}left(\textbackslash{}\textbackslash{}frac\{2\}\{25\} + \textbackslash{}\textbackslash{}frac\{11\}\{25\}i\textbackslash{}\textbackslash{}right)(3 {-} 4i) = \textbackslash{}\textbackslash{}frac\{1\}\{25\}(2 + 11i)(3 {-} 4i) = \textbackslash{}\textbackslash{}frac\{1\}\{25\}(6 {-} 8i + 33i {-} 44i\^{}2) = \textbackslash{}\textbackslash{}frac\{1\}\{25\}(50 + 25i) = 2 + i \textbackslash{}\textbackslash{}checkmark", note: null \},}
\NormalTok{  ]);}
\NormalTok{  return el;}
\NormalTok{\}}
\end{Highlighting}
\end{Shaded}

\chapter{Systems of equations}\label{systems-of-equations}

Two unknowns, two constraints

\hfill\break

Two streaming plans. Plan A charges \$6 per month plus \$3 per film.
Plan B charges \$15 per month with unlimited films. At some number of
films per month they cost exactly the same. Where is that crossover?

Two walkers start from opposite ends of a 9 km path. One walks at 4
km/h, the other at 5 km/h. When do they meet, and where?

Both questions have the same shape: two things you don't know, two
separate facts about them. A single equation can pin down one unknown.
To pin down two, you need two.

\section{What the notation is
saying}\label{what-the-notation-is-saying-16}

Here are the two streaming plans written as equations. Let \(m\) =
months (fixed at 1, so we can drop it) and \(f\) = number of films
rented:

\[\begin{cases} \text{Cost}_A = 6 + 3f \\ \text{Cost}_B = 15 \end{cases}\]

When the costs are equal: \(6 + 3f = 15\). One unknown, one equation ---
you already know how to solve this.

Now for the walkers. Let \(t\) = time in hours, \(d\) = distance walked
by the first walker. Two facts:

\[\begin{cases} d + (9 - d) = 9 \\ d = 4t \quad \text{and} \quad 9 - d = 5t \end{cases}\]

Rewriting cleanly:

\[\begin{cases} d + 5t = 9 \\ d - 4t = 0 \end{cases}\]

Read the curly brace as ``both of these must be true at the same time.''
That is a system. The solution is the pair of values \((d,\, t)\) that
satisfies \emph{both} equations simultaneously.

\subsection{The general form}\label{the-general-form}

The standard 2×2 system looks like this:

\[\begin{cases} a_1 x + b_1 y = c_1 \\ a_2 x + b_2 y = c_2 \end{cases}\]

Here \(x\) and \(y\) are the unknowns. The numbers \(a_1, b_1, c_1\) and
\(a_2, b_2, c_2\) are given --- they are the \textbf{coefficients} (the
multipliers of \(x\) and \(y\)) and the \textbf{constant terms} (the
values on the right).

\subsection{Compact matrix notation}\label{compact-matrix-notation}

The same system can be written using a \textbf{matrix} --- just the
coefficients arranged in a rectangular array:

\[\begin{pmatrix} a_1 & b_1 \\ a_2 & b_2 \end{pmatrix} \begin{pmatrix} x \\ y \end{pmatrix} = \begin{pmatrix} c_1 \\ c_2 \end{pmatrix}\]

Say it as: the 2×2 matrix of coefficients, multiplied by the column of
unknowns, equals the column of constants. This notation is compact ---
it stores the same information without repeating \(x\) and \(y\) on
every line. For a 2×2 system it's a convenience. For a 100×100 system
(as in circuit analysis or machine learning), it's essential.

The matrix itself --- the rectangular array of coefficients --- is
written with parentheses around it:

\[A = \begin{pmatrix} a_1 & b_1 \\ a_2 & b_2 \end{pmatrix}\]

Read ``\(A\)'' as ``the coefficient matrix.'' The entries are: top-left
\(a_1\) (coefficient of \(x\) in the first equation), top-right \(b_1\)
(coefficient of \(y\) in the first equation), and so on.

\section{The method}\label{the-method-15}

Three methods solve a 2×2 system. Each is worth knowing: the graphical
method shows you what you're doing, substitution is efficient for small
systems, and elimination is the one that scales.

\subsection{Method 1 --- Graphical}\label{method-1-graphical}

Each equation is a straight line. The solution is where the two lines
cross.

\[\begin{cases} x + y = 5 \\ x - y = 1 \end{cases}\]

Rearrange each equation into \(y = mx + b\) form:

\begin{itemize}
\tightlist
\item
  Line 1: \(y = -x + 5\) (slope \(-1\), intercept \(5\))
\item
  Line 2: \(y = x - 1\) (slope \(1\), intercept \(-1\))
\end{itemize}

Plot both. They cross at \((3, 2)\). Check: \(3 + 2 = 5\) ✓ and
\(3 - 2 = 1\) ✓.

\textbf{Three possible outcomes:}

\begin{longtable}[]{@{}
  >{\raggedright\arraybackslash}p{(\linewidth - 2\tabcolsep) * \real{0.5000}}
  >{\raggedright\arraybackslash}p{(\linewidth - 2\tabcolsep) * \real{0.5000}}@{}}
\toprule\noalign{}
\begin{minipage}[b]{\linewidth}\raggedright
What the lines do
\end{minipage} & \begin{minipage}[b]{\linewidth}\raggedright
What it means
\end{minipage} \\
\midrule\noalign{}
\endhead
\bottomrule\noalign{}
\endlastfoot
Cross at one point & One unique solution \\
Are parallel (never cross) & No solution --- the equations contradict
each other \\
Are the same line & Infinitely many solutions --- one equation gives no
new information \\
\end{longtable}

\begin{Shaded}
\begin{Highlighting}[]
\NormalTok{//| echo: false}
\NormalTok{\{}
\NormalTok{  const div = document.createElement("div");}
\NormalTok{  div.style.cssText = "font{-}family:inherit; padding:0.5rem 0;";}

\NormalTok{  const controls = document.createElement("div");}
\NormalTok{  controls.style.cssText = "display:grid; grid{-}template{-}columns:repeat(3,1fr); gap:0.5rem 1.5rem; margin{-}bottom:0.75rem; max{-}width:600px;";}

\NormalTok{  const params = [}
\NormalTok{    \{id:"sa1", label:"a₁", min:{-}5, max:5, step:0.5, val:1\},}
\NormalTok{    \{id:"sb1", label:"b₁", min:{-}5, max:5, step:0.5, val:1\},}
\NormalTok{    \{id:"sc1", label:"c₁", min:{-}10, max:10, step:1, val:5\},}
\NormalTok{    \{id:"sa2", label:"a₂", min:{-}5, max:5, step:0.5, val:2\},}
\NormalTok{    \{id:"sb2", label:"b₂", min:{-}5, max:5, step:0.5, val:{-}1\},}
\NormalTok{    \{id:"sc2", label:"c₂", min:{-}10, max:10, step:1, val:4\},}
\NormalTok{  ];}

\NormalTok{  params.forEach(p =\textgreater{} \{}
\NormalTok{    const wrap = document.createElement("div");}
\NormalTok{    const lbl = document.createElement("label");}
\NormalTok{    lbl.style.cssText = "font{-}size:0.8rem; color:\#374151; display:block;";}
\NormalTok{    const span = document.createElement("span"); span.textContent = p.val;}
\NormalTok{    lbl.appendChild(document.createTextNode(p.label + " = ")); lbl.appendChild(span);}
\NormalTok{    const input = document.createElement("input");}
\NormalTok{    input.type = "range"; input.min = p.min; input.max = p.max; input.step = p.step; input.value = p.val;}
\NormalTok{    input.style.width = "100\%";}
\NormalTok{    wrap.appendChild(lbl); wrap.appendChild(input);}
\NormalTok{    controls.appendChild(wrap);}
\NormalTok{    p.input = input;}
\NormalTok{    p.span = span;}
\NormalTok{  \});}

\NormalTok{  const infoEl = document.createElement("div");}
\NormalTok{  infoEl.style.cssText = "font{-}size:0.88rem; margin{-}bottom:0.5rem; min{-}height:1.4em; font{-}weight:600;";}

\NormalTok{  const canvas = document.createElement("canvas");}
\NormalTok{  canvas.width = 400; canvas.height = 400;}
\NormalTok{  canvas.style.cssText = "border:1px solid \#e5e7eb; border{-}radius:6px; display:block; max{-}width:400px;";}

\NormalTok{  div.appendChild(controls); div.appendChild(infoEl); div.appendChild(canvas);}

\NormalTok{  function draw() \{}
\NormalTok{    const vals = \{\};}
\NormalTok{    params.forEach(p =\textgreater{} \{}
\NormalTok{      const v = parseFloat(p.input.value);}
\NormalTok{      p.span.textContent = v;}
\NormalTok{      vals[p.id] = v;}
\NormalTok{    \});}

\NormalTok{    const \{sa1:a1, sb1:b1, sc1:c1, sa2:a2, sb2:b2, sc2:c2\} = vals;}
\NormalTok{    const det = a1*b2 {-} a2*b1;}

\NormalTok{    const ctx = canvas.getContext("2d");}
\NormalTok{    const W = canvas.width, H = canvas.height;}
\NormalTok{    ctx.clearRect(0, 0, W, H);}
\NormalTok{    const range = 12;}
\NormalTok{    function px(x) \{ return (x+range)/(2*range)*W; \}}
\NormalTok{    function py(y) \{ return H{-}(y+range)/(2*range)*H; \}}

\NormalTok{    // Grid}
\NormalTok{    ctx.strokeStyle = "\#f3f4f6"; ctx.lineWidth = 1;}
\NormalTok{    for (let i = {-}10; i \textless{}= 10; i += 2) \{}
\NormalTok{      ctx.beginPath(); ctx.moveTo(px(i), 0); ctx.lineTo(px(i), H); ctx.stroke();}
\NormalTok{      ctx.beginPath(); ctx.moveTo(0, py(i)); ctx.lineTo(W, py(i)); ctx.stroke();}
\NormalTok{    \}}
\NormalTok{    ctx.strokeStyle = "\#9ca3af"; ctx.lineWidth = 1.5;}
\NormalTok{    ctx.beginPath(); ctx.moveTo(0, py(0)); ctx.lineTo(W, py(0)); ctx.stroke();}
\NormalTok{    ctx.beginPath(); ctx.moveTo(px(0), 0); ctx.lineTo(px(0), H); ctx.stroke();}

\NormalTok{    // Line 1: a1*x + b1*y = c1  =\textgreater{} y = (c1 {-} a1*x)/b1}
\NormalTok{    function drawLine(a, b, c, color) \{}
\NormalTok{      ctx.strokeStyle = color; ctx.lineWidth = 2.5;}
\NormalTok{      ctx.beginPath();}
\NormalTok{      let first = true;}
\NormalTok{      for (let xi = 0; xi \textless{}= W; xi++) \{}
\NormalTok{        const x = {-}range + (xi / W) * 2 * range;}
\NormalTok{        let y;}
\NormalTok{        if (Math.abs(b) \textgreater{} 0.001) y = (c {-} a*x) / b;}
\NormalTok{        else \{ if (xi === 0) \{ ctx.moveTo(px(c/a), 0); ctx.lineTo(px(c/a), H); ctx.stroke(); \} return; \}}
\NormalTok{        if (first) \{ ctx.moveTo(xi, py(y)); first = false; \}}
\NormalTok{        else ctx.lineTo(xi, py(y));}
\NormalTok{      \}}
\NormalTok{      ctx.stroke();}
\NormalTok{    \}}

\NormalTok{    drawLine(a1, b1, c1, "\#2563eb");}
\NormalTok{    drawLine(a2, b2, c2, "\#ef4444");}

\NormalTok{    if (Math.abs(det) \textless{} 0.001) \{}
\NormalTok{      const parallel = Math.abs(a1*c2 {-} a2*c1) \textgreater{} 0.001 || Math.abs(b1*c2 {-} b2*c1) \textgreater{} 0.001;}
\NormalTok{      infoEl.textContent = parallel ? "No unique solution — lines are parallel" : "Infinitely many solutions — same line";}
\NormalTok{      infoEl.style.color = "\#dc2626";}
\NormalTok{    \} else \{}
\NormalTok{      const x = (c1*b2 {-} c2*b1) / det;}
\NormalTok{      const y = (a1*c2 {-} a2*c1) / det;}
\NormalTok{      ctx.fillStyle = "\#7c3aed";}
\NormalTok{      ctx.beginPath(); ctx.arc(px(x), py(y), 6, 0, 2*Math.PI); ctx.fill();}
\NormalTok{      ctx.fillStyle = "\#374151"; ctx.font = "bold 12px sans{-}serif";}
\NormalTok{      ctx.fillText("(" + x.toFixed(2) + ", " + y.toFixed(2) + ")", px(x)+8, py(y){-}8);}
\NormalTok{      infoEl.textContent = "Solution: x = " + x.toFixed(2) + ", y = " + y.toFixed(2);}
\NormalTok{      infoEl.style.color = "\#7c3aed";}
\NormalTok{    \}}

\NormalTok{    // Equation labels}
\NormalTok{    ctx.font = "12px sans{-}serif";}
\NormalTok{    ctx.fillStyle = "\#2563eb";}
\NormalTok{    ctx.fillText(a1+"x + "+b1+"y = "+c1, 8, 16);}
\NormalTok{    ctx.fillStyle = "\#ef4444";}
\NormalTok{    ctx.fillText(a2+"x + "+b2+"y = "+c2, 8, 32);}
\NormalTok{  \}}

\NormalTok{  params.forEach(p =\textgreater{} p.input.addEventListener("input", draw));}
\NormalTok{  draw();}
\NormalTok{  return div;}
\NormalTok{\}}
\end{Highlighting}
\end{Shaded}

\subsection{Method 2 --- Substitution}\label{method-2-substitution}

Solve one equation for one variable, then substitute into the other.
This eliminates one unknown and leaves a single equation you can solve
directly.

\[\begin{cases} x + y = 5 \qquad (1) \\ x - y = 1 \qquad (2) \end{cases}\]

From equation (1): \(x = 5 - y\).

Substitute into equation (2):

\[(5 - y) - y = 1\]

\[5 - 2y = 1\]

\[-2y = -4 \implies y = 2\]

Now substitute \(y = 2\) back into \(x = 5 - y\):

\[x = 5 - 2 = 3\]

Solution: \(x = 3\), \(y = 2\).

\textbf{When to use it:} substitution works cleanly when one equation
already has a variable with coefficient 1 (so there's no fraction to
deal with when you isolate it).

\subsection{Method 3 --- Elimination}\label{method-3-elimination}

Multiply one or both equations by constants to make the coefficient of
one variable the same (or opposite) in both equations. Then add or
subtract to eliminate that variable.

\[\begin{cases} 2x + 3y = 12 \qquad (1) \\ 4x - y = \phantom{1}5 \qquad\;\; (2) \end{cases}\]

The coefficient of \(y\) is \(3\) in equation (1) and \(-1\) in equation
(2). Multiply equation (2) by \(3\) to make the \(y\)-coefficients match
in magnitude:

\[3 \times (2): \quad 12x - 3y = 15 \qquad (2')\]

Now add equations (1) and (2'):

\[(2x + 3y) + (12x - 3y) = 12 + 15\]

\[14x = 27 \implies x = \frac{27}{14}\]

Substitute into equation (1):

\[2 \cdot \frac{27}{14} + 3y = 12 \implies \frac{27}{7} + 3y = 12 \implies 3y = 12 - \frac{27}{7} = \frac{57}{7} \implies y = \frac{19}{7}\]

Check in equation (2):
\(4 \cdot \frac{27}{14} - \frac{19}{7} = \frac{54}{7} - \frac{19}{7} = \frac{35}{7} = 5\)
✓

\textbf{When to use it:} elimination is best when no variable has a
coefficient of 1, or when the numbers make substitution awkward. It's
also the method that generalises directly to 3×3 and larger systems
(Gaussian elimination).

\section{Why this works}\label{why-this-works-2}

\begin{tcolorbox}[enhanced jigsaw, opacityback=0, colback=white, breakable, colframe=quarto-callout-note-color-frame, arc=.35mm, left=2mm, toprule=.15mm, leftrule=.75mm, bottomrule=.15mm, rightrule=.15mm]
\begin{minipage}[t]{5.5mm}
\textcolor{quarto-callout-note-color}{\faInfo}
\end{minipage}%
\begin{minipage}[t]{\textwidth - 5.5mm}

\textbf{The determinant tells you whether a unique solution exists}

For the system
\(\begin{pmatrix} a & b \\ c & d \end{pmatrix} \begin{pmatrix} x \\ y \end{pmatrix} = \begin{pmatrix} e \\ f \end{pmatrix}\),
define the \textbf{determinant}:

\[\det(A) = ad - bc\]

Read it as: top-left times bottom-right, minus top-right times
bottom-left.

\begin{itemize}
\tightlist
\item
  If \(\det(A) \neq 0\): one unique solution exists. The two lines
  cross.
\item
  If \(\det(A) = 0\): no unique solution. The lines are parallel (no
  solution) or identical (infinitely many).
\end{itemize}

For the system \(\begin{cases} x + y = 5 \\ x - y = 1 \end{cases}\), the
matrix is \(\begin{pmatrix} 1 & 1 \\ 1 & -1 \end{pmatrix}\) and
\(\det = (1)(-1) - (1)(1) = -2 \neq 0\). One solution. ✓

The determinant is a quick test before you do any work. If it's zero,
stop --- no amount of algebra will produce a unique answer.

\end{minipage}%
\end{tcolorbox}

\section{Worked examples}\label{worked-examples-15}

\begin{center}\rule{0.5\linewidth}{0.5pt}\end{center}

\textbf{Example 1 --- Business: supply and demand equilibrium.}

A market for concert tickets has a supply curve and a demand curve:

\[\begin{cases} p = 2q + 10 \qquad \text{(supply: price rises with quantity)} \\ p = -3q + 60 \qquad \text{(demand: price falls as quantity rises)} \end{cases}\]

Here \(p\) = price (dollars), \(q\) = quantity (thousands of tickets).
The equilibrium is where supply price equals demand price.

Set the right-hand sides equal (substitution):

\[2q + 10 = -3q + 60\]

\[5q = 50 \implies q = 10\]

\[p = 2(10) + 10 = 30\]

\textbf{Equilibrium:} 10,000 tickets at \$30 each.

Check in the demand equation: \(-3(10) + 60 = -30 + 60 = 30\) ✓

The determinant check: rewrite as
\(\begin{cases} 2q - p = -10 \\ -3q - p = -60 \end{cases}\). Matrix
\(\begin{pmatrix} 2 & -1 \\ -3 & -1 \end{pmatrix}\),
\(\det = (2)(-1) - (-1)(-3) = -2 - 3 = -5 \neq 0\). Unique solution
confirmed.

\begin{Shaded}
\begin{Highlighting}[]
\NormalTok{//| echo: false}
\NormalTok{\{}
\NormalTok{  const canvas = document.createElement("canvas");}
\NormalTok{  canvas.width = 400; canvas.height = 380;}
\NormalTok{  canvas.style.cssText = "border:1px solid \#e5e7eb; border{-}radius:6px; display:block; max{-}width:400px;";}

\NormalTok{  const ctx = canvas.getContext("2d");}
\NormalTok{  const W = canvas.width, H = canvas.height;}
\NormalTok{  const qMax = 20, pMax = 70;}
\NormalTok{  const marg = \{l:50, r:20, t:20, b:40\};}

\NormalTok{  function px(q) \{ return marg.l + (q / qMax) * (W {-} marg.l {-} marg.r); \}}
\NormalTok{  function py(p) \{ return marg.t + (1 {-} p / pMax) * (H {-} marg.t {-} marg.b); \}}

\NormalTok{  // Background shade: demand \textgreater{} supply (q \textless{} 10, left of intersection)}
\NormalTok{  ctx.fillStyle = "rgba(219,234,254,0.4)";}
\NormalTok{  ctx.fillRect(px(0), marg.t, px(10){-}px(0), H{-}marg.t{-}marg.b);}
\NormalTok{  ctx.fillStyle = "rgba(220,252,231,0.4)";}
\NormalTok{  ctx.fillRect(px(10), marg.t, px(qMax){-}px(10), H{-}marg.t{-}marg.b);}

\NormalTok{  // Axes}
\NormalTok{  ctx.strokeStyle = "\#9ca3af"; ctx.lineWidth = 1.5;}
\NormalTok{  ctx.beginPath(); ctx.moveTo(px(0), py(0)); ctx.lineTo(px(qMax), py(0)); ctx.stroke();}
\NormalTok{  ctx.beginPath(); ctx.moveTo(px(0), py(0)); ctx.lineTo(px(0), py(pMax)); ctx.stroke();}
\NormalTok{  ctx.fillStyle = "\#6b7280"; ctx.font = "12px sans{-}serif";}
\NormalTok{  ctx.fillText("q", px(qMax){-}10, py(0)+20);}
\NormalTok{  ctx.fillText("p", px(0){-}20, py(pMax)+5);}

\NormalTok{  // Tick marks}
\NormalTok{  for (let q = 0; q \textless{}= 20; q += 5) \{}
\NormalTok{    ctx.beginPath(); ctx.moveTo(px(q), py(0)); ctx.lineTo(px(q), py(0)+5); ctx.stroke();}
\NormalTok{    ctx.fillText(q, px(q){-}5, py(0)+18);}
\NormalTok{  \}}
\NormalTok{  for (let p = 0; p \textless{}= 70; p += 10) \{}
\NormalTok{    ctx.beginPath(); ctx.moveTo(px(0), py(p)); ctx.lineTo(px(0){-}5, py(p)); ctx.stroke();}
\NormalTok{    ctx.fillText(p, px(0){-}28, py(p)+4);}
\NormalTok{  \}}

\NormalTok{  // Supply: p = 2q + 10}
\NormalTok{  ctx.strokeStyle = "\#ef4444"; ctx.lineWidth = 2.5;}
\NormalTok{  ctx.beginPath(); ctx.moveTo(px(0), py(10)); ctx.lineTo(px(qMax), py(2*qMax+10)); ctx.stroke();}
\NormalTok{  ctx.fillStyle = "\#ef4444"; ctx.font = "bold 11px sans{-}serif";}
\NormalTok{  ctx.fillText("Supply: p = 2q + 10", px(12), py(2*12+10){-}8);}

\NormalTok{  // Demand: p = {-}3q + 60}
\NormalTok{  ctx.strokeStyle = "\#2563eb"; ctx.lineWidth = 2.5;}
\NormalTok{  ctx.beginPath(); ctx.moveTo(px(0), py(60)); ctx.lineTo(px(20), py({-}3*20+60)); ctx.stroke();}
\NormalTok{  ctx.fillStyle = "\#2563eb";}
\NormalTok{  ctx.fillText("Demand: p = −3q + 60", px(1), py({-}3*1+60){-}8);}

\NormalTok{  // Intersection at (10, 30)}
\NormalTok{  ctx.fillStyle = "\#7c3aed";}
\NormalTok{  ctx.beginPath(); ctx.arc(px(10), py(30), 6, 0, 2*Math.PI); ctx.fill();}
\NormalTok{  ctx.fillStyle = "\#374151"; ctx.font = "bold 12px sans{-}serif";}
\NormalTok{  ctx.fillText("Equilibrium (10, 30)", px(10)+8, py(30){-}8);}

\NormalTok{  // Shade labels}
\NormalTok{  ctx.font = "11px sans{-}serif";}
\NormalTok{  ctx.fillStyle = "\#1d4ed8"; ctx.fillText("demand \textgreater{} supply", px(1), py(50));}
\NormalTok{  ctx.fillStyle = "\#166534"; ctx.fillText("supply \textgreater{} demand", px(12), py(50));}

\NormalTok{  return canvas;}
\NormalTok{\}}
\end{Highlighting}
\end{Shaded}

\begin{center}\rule{0.5\linewidth}{0.5pt}\end{center}

\textbf{Example 2 --- Engineering: two-loop circuit (Kirchhoff's laws).}

A simple circuit has two loops. Kirchhoff's voltage law says the voltage
drops around each loop must sum to zero. After applying the law to each
loop, the unknowns are the two loop currents \(I_1\) and \(I_2\) (in
amps):

\[\begin{cases} 3I_1 + 2I_2 = 12 \qquad \text{(loop 1: total 12 V)} \\ 2I_1 + 5I_2 = 14 \qquad \text{(loop 2: total 14 V)} \end{cases}\]

Use elimination. Multiply equation (1) by 5 and equation (2) by 2:

\[\begin{cases} 15I_1 + 10I_2 = 60 \\ 4I_1 + 10I_2 = 28 \end{cases}\]

Subtract the second from the first:

\[11I_1 = 32 \implies I_1 = \frac{32}{11} \approx 2.91 \text{ A}\]

Substitute into equation (1):

\[3 \cdot \frac{32}{11} + 2I_2 = 12 \implies \frac{96}{11} + 2I_2 = 12 \implies 2I_2 = 12 - \frac{96}{11} = \frac{36}{11} \implies I_2 = \frac{18}{11} \approx 1.64 \text{ A}\]

Check in equation (2):
\(2 \cdot \frac{32}{11} + 5 \cdot \frac{18}{11} = \frac{64 + 90}{11} = \frac{154}{11} = 14\)
✓

\begin{center}\rule{0.5\linewidth}{0.5pt}\end{center}

\textbf{Example 3 --- Science: mixture problem.}

A lab needs 80 mL of a 35\% saline solution. The lab has a 20\% solution
and a 50\% solution. How many mL of each should be mixed?

Let \(x\) = mL of the 20\% solution, \(y\) = mL of the 50\% solution.

Two constraints:

\[\begin{cases} x + y = 80 \qquad \text{(total volume)} \\ 0.20x + 0.50y = 0.35 \times 80 = 28 \qquad \text{(total salt)} \end{cases}\]

From equation (1): \(x = 80 - y\). Substitute into equation (2):

\[0.20(80 - y) + 0.50y = 28\]

\[16 - 0.20y + 0.50y = 28\]

\[0.30y = 12 \implies y = 40\]

\[x = 80 - 40 = 40\]

Mix 40 mL of the 20\% solution with 40 mL of the 50\% solution.

Check: salt = \(0.20(40) + 0.50(40) = 8 + 20 = 28\) and \(28/80 = 0.35\)
✓

\begin{center}\rule{0.5\linewidth}{0.5pt}\end{center}

\textbf{Example 4 --- Computing: game resource constraints.}

A game character can perform two actions: attacking costs 3 energy and 1
second, defending costs 1 energy and 2 seconds. In one round, the
character uses exactly 14 energy and 8 seconds. How many attacks and
defends did they perform?

Let \(a\) = attacks, \(d\) = defends.

\[\begin{cases} 3a + d = 14 \qquad \text{(energy)} \\ a + 2d = 8 \qquad\;\, \text{(time)} \end{cases}\]

From equation (1): \(d = 14 - 3a\). Substitute into equation (2):

\[a + 2(14 - 3a) = 8\]

\[a + 28 - 6a = 8\]

\[-5a = -20 \implies a = 4\]

\[d = 14 - 3(4) = 2\]

Four attacks and two defends.

Check energy: \(3(4) + 2 = 14\) ✓. Check time: \(4 + 2(2) = 8\) ✓.

Determinant check: matrix
\(\begin{pmatrix} 3 & 1 \\ 1 & 2 \end{pmatrix}\),
\(\det = (3)(2) - (1)(1) = 5 \neq 0\). Unique solution. ✓

\section{Where this goes}\label{where-this-goes-16}

This chapter handled two unknowns and two equations. Real systems are
often much larger --- 10 unknowns, 100, or more. The method that scales
is \textbf{Gaussian elimination}: the same idea as the elimination you
learned here, applied row by row to a matrix. That is the core of
\textbf{linear algebra} (Vol 7). Matrix notation, which you saw as a
shorthand here, becomes indispensable there --- you'll work with
eigenvalues, matrix inverses, and linear transformations that are
impossible to write without it.

When systems get large enough, solving them exactly becomes
computationally expensive. \textbf{Numerical methods} (Vol 7) replace
exact algebra with iterative algorithms --- methods that get close
enough, fast enough, for engineering and data science work. Systems of
equations are the motivating problem that drives a large part of that
volume. The determinant you met here generalises too: for an
\(n \times n\) matrix, the same zero/non-zero test predicts whether a
unique solution exists.

\begin{tcolorbox}[enhanced jigsaw, opacityback=0, colback=white, breakable, colframe=quarto-callout-tip-color-frame, arc=.35mm, left=2mm, toprule=.15mm, leftrule=.75mm, bottomrule=.15mm, rightrule=.15mm]
\begin{minipage}[t]{5.5mm}
\textcolor{quarto-callout-tip-color}{\faLightbulb}
\end{minipage}%
\begin{minipage}[t]{\textwidth - 5.5mm}

\textbf{Where this shows up}

\begin{itemize}
\tightlist
\item
  \textbf{Engineering:} Kirchhoff's voltage and current laws applied to
  a circuit produce one equation per loop or node. A circuit with two
  loops gives a 2×2 system; a real circuit might give a 50×50 system
  solved by computer using the same method.
\item
  \textbf{Science:} Chemical stoichiometry --- balancing an equation
  with two unknown coefficients --- is a 2×2 system. Equilibrium
  calculations in physical chemistry are similar.
\item
  \textbf{Finance:} Supply and demand equilibrium is the intersection of
  two linear functions. Every introductory economics model of
  price-setting is a system of equations.
\item
  \textbf{Computing:} Linear regression with two parameters solves a 2×2
  system. Graphics engines apply matrix transformations (systems of
  equations) to every vertex on screen, thousands of times per second.
\end{itemize}

\end{minipage}%
\end{tcolorbox}

\section{Exercises}\label{exercises-16}

These are puzzles. The interesting part is setting up the system ---
translating the situation into two equations before you solve.

\begin{center}\rule{0.5\linewidth}{0.5pt}\end{center}

\begin{Shaded}
\begin{Highlighting}[]
\NormalTok{//| echo: false}
\NormalTok{katex = require("katex@0.16")}
\end{Highlighting}
\end{Shaded}

\begin{Shaded}
\begin{Highlighting}[]
\NormalTok{//| echo: false}
\NormalTok{function makeStepperHTML(exerciseNum, steps) \{}
\NormalTok{  let current = 0;}
\NormalTok{  const totalSteps = steps.length;}
\NormalTok{  const container = document.createElement("div");}
\NormalTok{  container.style.cssText = "border:1px solid \#e5e7eb; border{-}radius:8px; padding:1rem 1.25rem; margin:0.75rem 0; font{-}family:inherit;";}
\NormalTok{  const stepsDiv = document.createElement("div");}
\NormalTok{  stepsDiv.style.marginBottom = "0.75rem";}
\NormalTok{  function renderSteps() \{}
\NormalTok{    stepsDiv.innerHTML = "";}
\NormalTok{    for (let i = 0; i \textless{} current; i++) \{}
\NormalTok{      const s = steps[i];}
\NormalTok{      const row = document.createElement("div");}
\NormalTok{      row.style.cssText = "display:grid; grid{-}template{-}columns:200px 1fr; gap:0.35rem 1rem; align{-}items:baseline; padding:0.35rem 0; border{-}top:1px solid \#f3f4f6;";}
\NormalTok{      const opEl = document.createElement("span");}
\NormalTok{      opEl.style.cssText = "font{-}size:0.8rem; color:\#6b7280; font{-}style:italic;";}
\NormalTok{      opEl.textContent = s.op;}
\NormalTok{      const eqEl = document.createElement("span");}
\NormalTok{      eqEl.innerHTML = katex.renderToString(s.eq, \{throwOnError: false, displayMode: false\});}
\NormalTok{      row.appendChild(opEl);}
\NormalTok{      row.appendChild(eqEl);}
\NormalTok{      container.appendChild(row);}
\NormalTok{      if (s.note) \{}
\NormalTok{        const noteEl = document.createElement("p");}
\NormalTok{        noteEl.style.cssText = "margin:0.2rem 0 0.5rem 200px; font{-}size:0.82rem; color:\#374151;";}
\NormalTok{        noteEl.textContent = s.note;}
\NormalTok{        container.appendChild(noteEl);}
\NormalTok{      \}}
\NormalTok{    \}}
\NormalTok{    if (current === totalSteps) \{}
\NormalTok{      const done = document.createElement("p");}
\NormalTok{      done.style.cssText = "margin{-}top:0.5rem; font{-}size:0.85rem; color:\#059669; font{-}weight:600;";}
\NormalTok{      done.textContent = "✓ Solution complete";}
\NormalTok{      stepsDiv.appendChild(done);}
\NormalTok{    \}}
\NormalTok{  \}}
\NormalTok{  const controls = document.createElement("div");}
\NormalTok{  controls.style.cssText = "display:flex; align{-}items:center; gap:0.75rem; margin{-}top:0.5rem;";}
\NormalTok{  const nextBtn = document.createElement("button");}
\NormalTok{  nextBtn.textContent = "Next step →";}
\NormalTok{  nextBtn.style.cssText = "padding:0.35rem 0.85rem; background:\#2563eb; color:\#fff; border:none; border{-}radius:6px; cursor:pointer; font{-}size:0.85rem;";}
\NormalTok{  const resetBtn = document.createElement("button");}
\NormalTok{  resetBtn.textContent = "Reset";}
\NormalTok{  resetBtn.style.cssText = "padding:0.35rem 0.85rem; background:\#f3f4f6; color:\#374151; border:1px solid \#d1d5db; border{-}radius:6px; cursor:pointer; font{-}size:0.85rem;";}
\NormalTok{  const counter = document.createElement("span");}
\NormalTok{  counter.style.cssText = "font{-}size:0.8rem; color:\#6b7280;";}
\NormalTok{  nextBtn.onclick = () =\textgreater{} \{ if (current \textless{} totalSteps) \{ current++; renderSteps(); counter.textContent = current + " / " + totalSteps; if (current === totalSteps) nextBtn.disabled = true; \} \};}
\NormalTok{  resetBtn.onclick = () =\textgreater{} \{ current = 0; nextBtn.disabled = false; renderSteps(); counter.textContent = "0 / " + totalSteps; \};}
\NormalTok{  counter.textContent = "0 / " + totalSteps;}
\NormalTok{  controls.appendChild(nextBtn);}
\NormalTok{  controls.appendChild(resetBtn);}
\NormalTok{  controls.appendChild(counter);}
\NormalTok{  container.appendChild(stepsDiv);}
\NormalTok{  container.appendChild(controls);}
\NormalTok{  return container;}
\NormalTok{\}}
\end{Highlighting}
\end{Shaded}

\textbf{1.} Two friends are buying supplies for a camping trip. Rashida
buys 3 bags of trail mix and 2 water bottles for \$19. Tomasz buys 1 bag
of trail mix and 4 water bottles for \$17. What does one bag of trail
mix cost, and what does one water bottle cost?

\begin{Shaded}
\begin{Highlighting}[]
\NormalTok{//| echo: false}
\NormalTok{\{}
\NormalTok{  const el = makeStepperHTML(1, [}
\NormalTok{    \{ op: "Write as a system", eq: "\textbackslash{}\textbackslash{}begin\{cases\} 3t + 2w = 19 \textbackslash{}\textbackslash{}\textbackslash{}\textbackslash{} t + 4w = 17 \textbackslash{}\textbackslash{}end\{cases\}", note: "t = price of trail mix, w = price of water bottle" \},}
\NormalTok{    \{ op: "Solve equation 2 for t", eq: "t = 17 {-} 4w", note: null \},}
\NormalTok{    \{ op: "Substitute into equation 1", eq: "3(17 {-} 4w) + 2w = 19", note: null \},}
\NormalTok{    \{ op: "Expand and simplify", eq: "51 {-} 12w + 2w = 19 \textbackslash{}\textbackslash{}implies {-}10w = {-}32 \textbackslash{}\textbackslash{}implies w = 3.20", note: null \},}
\NormalTok{    \{ op: "Back{-}substitute to find t", eq: "t = 17 {-} 4(3.20) = 17 {-} 12.80 = 4.20", note: null \},}
\NormalTok{    \{ op: "Check", eq: "3(4.20) + 2(3.20) = 12.60 + 6.40 = 19 \textbackslash{}\textbackslash{}checkmark \textbackslash{}\textbackslash{}quad 4.20 + 4(3.20) = 4.20 + 12.80 = 17 \textbackslash{}\textbackslash{}checkmark", note: null \},}
\NormalTok{  ]);}
\NormalTok{  return el;}
\NormalTok{\}}
\end{Highlighting}
\end{Shaded}

\begin{center}\rule{0.5\linewidth}{0.5pt}\end{center}

\textbf{2.} A cinema charges one price for adults and another for
children. Two families bought tickets last Saturday:

\begin{itemize}
\tightlist
\item
  The Nguyens: 3 adults and 2 children, total \$32.
\item
  The Garcias: 1 adult and 4 children, total \$24.
\end{itemize}

What is the adult ticket price? What is the child ticket price?

\begin{Shaded}
\begin{Highlighting}[]
\NormalTok{//| echo: false}
\NormalTok{\{}
\NormalTok{  const el = makeStepperHTML(2, [}
\NormalTok{    \{ op: "Write as a system", eq: "\textbackslash{}\textbackslash{}begin\{cases\} 3a + 2c = 32 \textbackslash{}\textbackslash{}\textbackslash{}\textbackslash{} a + 4c = 24 \textbackslash{}\textbackslash{}end\{cases\}", note: "a = adult price, c = child price" \},}
\NormalTok{    \{ op: "Solve equation 2 for a", eq: "a = 24 {-} 4c", note: null \},}
\NormalTok{    \{ op: "Substitute into equation 1", eq: "3(24 {-} 4c) + 2c = 32 \textbackslash{}\textbackslash{}implies 72 {-} 12c + 2c = 32 \textbackslash{}\textbackslash{}implies {-}10c = {-}40 \textbackslash{}\textbackslash{}implies c = 4", note: null \},}
\NormalTok{    \{ op: "Back{-}substitute to find a", eq: "a = 24 {-} 4(4) = 24 {-} 16 = 8", note: null \},}
\NormalTok{    \{ op: "Check", eq: "3(8) + 2(4) = 24 + 8 = 32 \textbackslash{}\textbackslash{}checkmark \textbackslash{}\textbackslash{}quad 1(8) + 4(4) = 8 + 16 = 24 \textbackslash{}\textbackslash{}checkmark", note: "Child ticket = £4, adult ticket = £8." \},}
\NormalTok{  ]);}
\NormalTok{  return el;}
\NormalTok{\}}
\end{Highlighting}
\end{Shaded}

\begin{center}\rule{0.5\linewidth}{0.5pt}\end{center}

\textbf{3.} Two cyclists start from the same point at the same time and
ride in opposite directions. After 3 hours they are 132 km apart. One
rides 6 km/h faster than the other. What is each cyclist's speed?

\begin{Shaded}
\begin{Highlighting}[]
\NormalTok{//| echo: false}
\NormalTok{\{}
\NormalTok{  const el = makeStepperHTML(3, [}
\NormalTok{    \{ op: "Write as a system", eq: "\textbackslash{}\textbackslash{}begin\{cases\} 3r + 3s = 132 \textbackslash{}\textbackslash{}\textbackslash{}\textbackslash{} r {-} s = 6 \textbackslash{}\textbackslash{}end\{cases\}", note: "r and s are the two speeds; together they cover 132 km in 3 hours; one is 6 km/h faster" \},}
\NormalTok{    \{ op: "Simplify equation 1", eq: "r + s = 44", note: null \},}
\NormalTok{    \{ op: "Add the two equations", eq: "(r + s) + (r {-} s) = 44 + 6 \textbackslash{}\textbackslash{}implies 2r = 50 \textbackslash{}\textbackslash{}implies r = 25 \textbackslash{}\textbackslash{}text\{ km/h\}", note: null \},}
\NormalTok{    \{ op: "Back{-}substitute to find s", eq: "s = 44 {-} 25 = 19 \textbackslash{}\textbackslash{}text\{ km/h\}", note: null \},}
\NormalTok{    \{ op: "Check", eq: "3(25) + 3(19) = 75 + 57 = 132 \textbackslash{}\textbackslash{}checkmark \textbackslash{}\textbackslash{}quad 25 {-} 19 = 6 \textbackslash{}\textbackslash{}checkmark", note: null \},}
\NormalTok{  ]);}
\NormalTok{  return el;}
\NormalTok{\}}
\end{Highlighting}
\end{Shaded}

\begin{center}\rule{0.5\linewidth}{0.5pt}\end{center}

\textbf{4.} Use the determinant to decide, without solving, whether each
system has a unique solution, no solution, or infinitely many solutions.

\textbf{(a)} \[\begin{cases} 2x + 4y = 8 \\ x + 2y = 4 \end{cases}\]

\textbf{(b)} \[\begin{cases} 3x - y = 5 \\ 6x - 2y = 7 \end{cases}\]

\textbf{(c)} \[\begin{cases} 5x + 3y = 11 \\ 2x - y = 3 \end{cases}\]

\begin{Shaded}
\begin{Highlighting}[]
\NormalTok{//| echo: false}
\NormalTok{\{}
\NormalTok{  const el = makeStepperHTML(4, [}
\NormalTok{    \{ op: "Write the coefficient matrix for (a)", eq: "A = \textbackslash{}\textbackslash{}begin\{pmatrix\} 2 \& 4 \textbackslash{}\textbackslash{}\textbackslash{}\textbackslash{} 1 \& 2 \textbackslash{}\textbackslash{}end\{pmatrix\}", note: null \},}
\NormalTok{    \{ op: "Compute det(A) for (a)", eq: "\textbackslash{}\textbackslash{}det(A) = (2)(2) {-} (4)(1) = 4 {-} 4 = 0", note: "det = 0 means no unique solution; check if the equations are the same line or parallel" \},}
\NormalTok{    \{ op: "Inspect (a): are the equations the same?", eq: "\textbackslash{}\textbackslash{}text\{Row 1\} = 2 \textbackslash{}\textbackslash{}times \textbackslash{}\textbackslash{}text\{Row 2\}: \textbackslash{}\textbackslash{}quad 2x + 4y = 8 \textbackslash{}\textbackslash{}text\{ is exactly \} 2(x + 2y = 4). \textbackslash{}\textbackslash{}text\{ Same line.\}", note: "Infinitely many solutions." \},}
\NormalTok{    \{ op: "Compute det for (b)", eq: "\textbackslash{}\textbackslash{}det\textbackslash{}\textbackslash{}begin\{pmatrix\} 3 \& {-}1 \textbackslash{}\textbackslash{}\textbackslash{}\textbackslash{} 6 \& {-}2 \textbackslash{}\textbackslash{}end\{pmatrix\} = (3)({-}2) {-} ({-}1)(6) = {-}6 + 6 = 0", note: null \},}
\NormalTok{    \{ op: "Inspect (b): same line or parallel?", eq: "\textbackslash{}\textbackslash{}text\{Row 1\} = \textbackslash{}\textbackslash{}tfrac\{1\}\{2\} \textbackslash{}\textbackslash{}times \textbackslash{}\textbackslash{}text\{Row 2 LHS, but \} \textbackslash{}\textbackslash{}tfrac\{7\}\{2\} \textbackslash{}\textbackslash{}neq 5. \textbackslash{}\textbackslash{}text\{ Parallel lines.\}", note: "No solution." \},}
\NormalTok{    \{ op: "Compute det for (c)", eq: "\textbackslash{}\textbackslash{}det\textbackslash{}\textbackslash{}begin\{pmatrix\} 5 \& 3 \textbackslash{}\textbackslash{}\textbackslash{}\textbackslash{} 2 \& {-}1 \textbackslash{}\textbackslash{}end\{pmatrix\} = (5)({-}1) {-} (3)(2) = {-}5 {-} 6 = {-}11 \textbackslash{}\textbackslash{}neq 0", note: "Unique solution exists." \},}
\NormalTok{    \{ op: "Solve by substitution", eq: "2x {-} y = 3 \textbackslash{}\textbackslash{}Rightarrow y = 2x {-} 3", note: null \},}
\NormalTok{    \{ op: "Substitute into first equation", eq: "5x + 3(2x{-}3) = 11 \textbackslash{}\textbackslash{}Rightarrow 11x = 20 \textbackslash{}\textbackslash{}Rightarrow x = \textbackslash{}\textbackslash{}frac\{20\}\{11\}", note: null \},}
\NormalTok{    \{ op: "Find y", eq: "y = 2 \textbackslash{}\textbackslash{}cdot \textbackslash{}\textbackslash{}frac\{20\}\{11\} {-} 3 = \textbackslash{}\textbackslash{}frac\{40\}\{11\} {-} \textbackslash{}\textbackslash{}frac\{33\}\{11\} = \textbackslash{}\textbackslash{}frac\{7\}\{11\}", note: null \},}
\NormalTok{    \{ op: "Check", eq: "5 \textbackslash{}\textbackslash{}cdot \textbackslash{}\textbackslash{}frac\{20\}\{11\} + 3 \textbackslash{}\textbackslash{}cdot \textbackslash{}\textbackslash{}frac\{7\}\{11\} = \textbackslash{}\textbackslash{}frac\{100\}\{11\} + \textbackslash{}\textbackslash{}frac\{21\}\{11\} = \textbackslash{}\textbackslash{}frac\{121\}\{11\} = 11 \textbackslash{}\textbackslash{}checkmark", note: null \},}
\NormalTok{  ]);}
\NormalTok{  return el;}
\NormalTok{\}}
\end{Highlighting}
\end{Shaded}

\begin{center}\rule{0.5\linewidth}{0.5pt}\end{center}

\textbf{5.} Two phone plans:

\begin{itemize}
\tightlist
\item
  Plan X: \$20 per month plus \$0.10 per text.
\item
  Plan Y: \$35 per month, unlimited texts.
\end{itemize}

\emph{(a)} Set up a system and find the number of texts per month where
both plans cost the same.

\emph{(b)} If you send 200 texts per month, which plan is cheaper and by
how much?

\begin{Shaded}
\begin{Highlighting}[]
\NormalTok{//| echo: false}
\NormalTok{\{}
\NormalTok{  const el = makeStepperHTML(5, [}
\NormalTok{    \{ op: "Write the cost equations", eq: "\textbackslash{}\textbackslash{}begin\{cases\} C\_X = 20 + 0.10n \textbackslash{}\textbackslash{}\textbackslash{}\textbackslash{} C\_Y = 35 \textbackslash{}\textbackslash{}end\{cases\}", note: "n = number of texts per month" \},}
\NormalTok{    \{ op: "Set costs equal and solve", eq: "20 + 0.10n = 35 \textbackslash{}\textbackslash{}implies 0.10n = 15 \textbackslash{}\textbackslash{}implies n = 150", note: null \},}
\NormalTok{    \{ op: "State the crossover", eq: "\textbackslash{}\textbackslash{}text\{At \} n = 150 \textbackslash{}\textbackslash{}text\{ texts, both plans cost \} \textbackslash{}\textbackslash{}$35.", note: null \},}
\NormalTok{    \{ op: "Evaluate at n = 200 for part (b)", eq: "C\_X = 20 + 0.10(200) = 20 + 20 = \textbackslash{}\textbackslash{}$40; \textbackslash{}\textbackslash{}quad C\_Y = \textbackslash{}\textbackslash{}$35", note: null \},}
\NormalTok{    \{ op: "Check", eq: "C\_X {-} C\_Y = 40 {-} 35 = \textbackslash{}\textbackslash{}$5. \textbackslash{}\textbackslash{}text\{ Plan Y is cheaper by \} \textbackslash{}\textbackslash{}$5 \textbackslash{}\textbackslash{}text\{ at 200 texts.\} \textbackslash{}\textbackslash{}checkmark", note: null \},}
\NormalTok{  ]);}
\NormalTok{  return el;}
\NormalTok{\}}
\end{Highlighting}
\end{Shaded}

\begin{center}\rule{0.5\linewidth}{0.5pt}\end{center}

\textbf{6.} A chemist has two alcohol solutions: one is 30\% alcohol,
the other is 70\% alcohol. She needs 200 mL of a 45\% solution.

How many mL of each solution should she mix?

\emph{(Hint: write one equation for total volume and one for total
amount of alcohol.)}

\begin{Shaded}
\begin{Highlighting}[]
\NormalTok{//| echo: false}
\NormalTok{\{}
\NormalTok{  const el = makeStepperHTML(6, [}
\NormalTok{    \{ op: "Write as a system", eq: "\textbackslash{}\textbackslash{}begin\{cases\} x + y = 200 \textbackslash{}\textbackslash{}\textbackslash{}\textbackslash{} 0.30x + 0.70y = 0.45 \textbackslash{}\textbackslash{}times 200 = 90 \textbackslash{}\textbackslash{}end\{cases\}", note: "x = mL of 30\% solution, y = mL of 70\% solution" \},}
\NormalTok{    \{ op: "Solve equation 1 for x", eq: "x = 200 {-} y", note: null \},}
\NormalTok{    \{ op: "Substitute into equation 2", eq: "0.30(200 {-} y) + 0.70y = 90 \textbackslash{}\textbackslash{}implies 60 {-} 0.30y + 0.70y = 90 \textbackslash{}\textbackslash{}implies 0.40y = 30 \textbackslash{}\textbackslash{}implies y = 75", note: null \},}
\NormalTok{    \{ op: "Back{-}substitute to find x", eq: "x = 200 {-} 75 = 125", note: null \},}
\NormalTok{    \{ op: "Check", eq: "0.30(125) + 0.70(75) = 37.50 + 52.50 = 90 \textbackslash{}\textbackslash{}quad \textbackslash{}\textbackslash{}text\{and\} \textbackslash{}\textbackslash{}quad 90/200 = 0.45 \textbackslash{}\textbackslash{}checkmark", note: null \},}
\NormalTok{  ]);}
\NormalTok{  return el;}
\NormalTok{\}}
\end{Highlighting}
\end{Shaded}

\chapter{Sequences and series}\label{sequences-and-series}

When adding up becomes its own question

\hfill\break

You start a sponsored read-a-thon. On day 1 you read 10 pages. Each day
you read 3 more pages than the day before. After 30 days, how many pages
have you read in total? Adding them up one by one would take a while ---
but there's a pattern.

To walk from here to a wall, you must first walk halfway. Then half of
the remaining distance. Then half again. You take infinitely many steps
--- so how do you ever arrive? The answer is that the distances form a
series whose infinite sum is exactly the distance to the wall.

Both problems need the same tool. This chapter builds it.

\section{What the notation is
saying}\label{what-the-notation-is-saying-17}

A \textbf{sequence} is a list of numbers following a rule. A
\textbf{series} is the sum of those numbers, in order.

Write the sequence \(3, 7, 11, 15, 19\). The corresponding series is:

\[3 + 7 + 11 + 15 + 19 = 45\]

That works when you have a short list. For 300 terms, writing out every
addition is unrealistic. Mathematicians use a shorthand --- the
\textbf{sigma notation} --- named after the Greek letter \(\Sigma\)
(capital sigma).

\[\sum_{i=1}^{n} a_i\]

Read it as: \emph{the sum of \(a_i\), as \(i\) runs from 1 to \(n\).}

The \(i\) below the sigma is the \textbf{index} --- it's a counter that
starts at 1 and steps up by 1 until it reaches \(n\). At each step, you
substitute that value of \(i\) into the expression \(a_i\) and add the
result to your running total.

A concrete example before using sigma in formulas. Write the sum
\(1 + 4 + 9 + 16 + 25\) --- the first five perfect squares --- using
sigma:

\[\sum_{i=1}^{5} i^2 = 1^2 + 2^2 + 3^2 + 4^2 + 5^2 = 55\]

The expression after \(\Sigma\) is the rule. The limits below and above
tell you where to start and stop. That is the whole notation.

\begin{Shaded}
\begin{Highlighting}[]
\NormalTok{//| echo: false}
\NormalTok{\{}
\NormalTok{  const div = document.createElement("div");}
\NormalTok{  div.style.cssText = "font{-}family:inherit; padding:0.5rem 0;";}

\NormalTok{  const controls = document.createElement("div");}
\NormalTok{  controls.style.cssText = "display:flex; gap:1.5rem; flex{-}wrap:wrap; align{-}items:flex{-}start; margin{-}bottom:0.75rem;";}

\NormalTok{  // Type selector}
\NormalTok{  const typeWrap = document.createElement("div");}
\NormalTok{  const typeLbl = document.createElement("label");}
\NormalTok{  typeLbl.style.cssText = "font{-}size:0.85rem; color:\#374151; display:block; margin{-}bottom:0.2rem;";}
\NormalTok{  typeLbl.textContent = "Type";}
\NormalTok{  const typeSelect = document.createElement("select");}
\NormalTok{  typeSelect.style.cssText = "padding:0.2rem 0.4rem; border:1px solid \#d1d5db; border{-}radius:4px;";}
\NormalTok{  ["arithmetic", "geometric"].forEach(t =\textgreater{} \{}
\NormalTok{    const o = document.createElement("option"); o.value = t; o.textContent = t;}
\NormalTok{    typeSelect.appendChild(o);}
\NormalTok{  \});}
\NormalTok{  typeWrap.appendChild(typeLbl); typeWrap.appendChild(typeSelect);}
\NormalTok{  controls.appendChild(typeWrap);}

\NormalTok{  // makeSlider returns \{wrap, input, span\} so callers keep direct refs}
\NormalTok{  function makeSlider(label, min, max, step, val) \{}
\NormalTok{    const wrap = document.createElement("div");}
\NormalTok{    const lbl = document.createElement("label");}
\NormalTok{    lbl.style.cssText = "font{-}size:0.85rem; color:\#374151; display:block; margin{-}bottom:0.2rem;";}
\NormalTok{    const span = document.createElement("span"); span.textContent = val;}
\NormalTok{    lbl.appendChild(document.createTextNode(label + " = ")); lbl.appendChild(span);}
\NormalTok{    const input = document.createElement("input");}
\NormalTok{    input.type = "range"; input.min = min; input.max = max; input.step = step; input.value = val;}
\NormalTok{    input.style.width = "140px";}
\NormalTok{    wrap.appendChild(lbl); wrap.appendChild(input);}
\NormalTok{    return \{wrap, input, span\};}
\NormalTok{  \}}

\NormalTok{  const a1S = makeSlider("a₁ (first term)", 1, 20, 1, 2);}
\NormalTok{  controls.appendChild(a1S.wrap);}
\NormalTok{  const dS = makeSlider("d (common diff)", {-}5, 10, 1, 3);}
\NormalTok{  dS.wrap.style.display = "";}
\NormalTok{  controls.appendChild(dS.wrap);}
\NormalTok{  const rS = makeSlider("r (ratio)", {-}3, 3, 0.5, 2);}
\NormalTok{  rS.wrap.style.display = "none";}
\NormalTok{  controls.appendChild(rS.wrap);}
\NormalTok{  const nS = makeSlider("n (terms)", 3, 12, 1, 6);}
\NormalTok{  controls.appendChild(nS.wrap);}

\NormalTok{  const displayEl = document.createElement("div");}
\NormalTok{  displayEl.style.cssText = "background:\#f8fafc; border:1px solid \#e5e7eb; border{-}radius:6px; padding:0.75rem 1rem; font{-}size:0.9rem; line{-}height:1.8;";}

\NormalTok{  div.appendChild(controls); div.appendChild(displayEl);}

\NormalTok{  function update() \{}
\NormalTok{    const type = typeSelect.value;}
\NormalTok{    const a1 = parseFloat(a1S.input.value);}
\NormalTok{    const d = parseFloat(dS.input.value);}
\NormalTok{    const r = parseFloat(rS.input.value);}
\NormalTok{    const n = parseInt(nS.input.value);}
\NormalTok{    a1S.span.textContent = a1;}
\NormalTok{    dS.span.textContent = d;}
\NormalTok{    rS.span.textContent = r;}
\NormalTok{    nS.span.textContent = n;}
\NormalTok{    dS.wrap.style.display = type === "arithmetic" ? "" : "none";}
\NormalTok{    rS.wrap.style.display = type === "geometric" ? "" : "none";}

\NormalTok{    const terms = [];}
\NormalTok{    let sum = 0;}
\NormalTok{    for (let k = 1; k \textless{}= n; k++) \{}
\NormalTok{      const t = type === "arithmetic" ? a1 + (k{-}1)*d : a1 * Math.pow(r, k{-}1);}
\NormalTok{      terms.push(t);}
\NormalTok{      sum += t;}
\NormalTok{    \}}

\NormalTok{    const sigmaLabel = type === "arithmetic"}
\NormalTok{      ? "Σ(k=1 to "+n+") ["+a1+(d\textgreater{}=0?"+":"")+d+"(k−1)]"}
\NormalTok{      : "Σ(k=1 to "+n+") "+a1+"·"+r+"\^{}(k−1)";}

\NormalTok{    let expanded;}
\NormalTok{    if (n \textless{}= 6) expanded = terms.map(t =\textgreater{} t.toFixed(1)).join(" + ");}
\NormalTok{    else expanded = terms.slice(0,3).map(t=\textgreater{}t.toFixed(1)).join(" + ") + " + … + " + terms[n{-}1].toFixed(1);}

\NormalTok{    displayEl.innerHTML =}
\NormalTok{      "\textless{}b\textgreater{}Sigma notation:\textless{}/b\textgreater{} " + sigmaLabel + "\textless{}br\textgreater{}" +}
\NormalTok{      "\textless{}b\textgreater{}Expanded:\textless{}/b\textgreater{} " + expanded + "\textless{}br\textgreater{}" +}
\NormalTok{      "\textless{}b\textgreater{}Sum:\textless{}/b\textgreater{} S\textless{}sub\textgreater{}" + n + "\textless{}/sub\textgreater{} = " + sum.toFixed(2);}
\NormalTok{  \}}

\NormalTok{  typeSelect.addEventListener("input", update);}
\NormalTok{  [a1S, dS, rS, nS].forEach(s =\textgreater{} s.input.addEventListener("input", update));}
\NormalTok{  update();}
\NormalTok{  return div;}
\NormalTok{\}}
\end{Highlighting}
\end{Shaded}

\section{The method}\label{the-method-16}

\subsection{Arithmetic series}\label{arithmetic-series}

An arithmetic sequence has a first term \(a_1\) and a common difference
\(d\). The terms are:

\[a_1,\quad a_1 + d,\quad a_1 + 2d,\quad \ldots,\quad a_1 + (n-1)d\]

Adding all \(n\) of them gives the \textbf{arithmetic series} \(S_n\),
read as \emph{the sum of the first \(n\) terms}.

To see why the formula works, try a small example first. Take the
arithmetic sequence \(3, 6, 9, 12\) (four terms, \(d = 3\)). Write the
sum forwards and backwards, then add column by column:

\[\begin{array}{r} S_4 = \phantom{1}3 + \phantom{1}6 + \phantom{1}9 + 12 \\ S_4 = 12 + \phantom{1}9 + \phantom{1}6 + \phantom{1}3 \\ \hline 2S_4 = 15 + 15 + 15 + 15 = 4 \times 15 = 60 \end{array}\]

\[S_4 = 30\]

The same pairing trick works for any arithmetic series. Label the sum
forwards and backwards:

\[S_n = a_1 + (a_1 + d) + (a_1 + 2d) + \cdots + a_n\]

\[S_n = a_n + (a_n - d) + (a_n - 2d) + \cdots + a_1\]

Add the two lines together. Each column pairs a term from the top with a
term from the bottom. Every pair sums to \(a_1 + a_n\), and there are
\(n\) such pairs:

\[2S_n = n(a_1 + a_n)\]

\[\boxed{S_n = \frac{n}{2}(a_1 + a_n)}\]

Since \(a_n = a_1 + (n-1)d\), you can write this entirely in terms of
the first term and the common difference:

\[S_n = \frac{n}{2}\bigl(2a_1 + (n-1)d\bigr)\]

Use whichever form is more convenient. If you know the first and last
term, use the first. If you only know \(a_1\) and \(d\), use the second.

\begin{tcolorbox}[enhanced jigsaw, opacityback=0, colback=white, breakable, colframe=quarto-callout-note-color-frame, arc=.35mm, left=2mm, toprule=.15mm, leftrule=.75mm, bottomrule=.15mm, rightrule=.15mm]
\begin{minipage}[t]{5.5mm}
\textcolor{quarto-callout-note-color}{\faInfo}
\end{minipage}%
\begin{minipage}[t]{\textwidth - 5.5mm}

\textbf{The Gauss pairing trick}

Carl Friedrich Gauss was reportedly seven years old when his teacher set
the class the task of summing all integers from 1 to 100, hoping for an
hour of quiet. Gauss had the answer in seconds: 5050.

His insight: pair the first term with the last, the second with the
second-to-last, and so on. Each pair sums to 101. There are 50 pairs. So
\(50 \times 101 = 5050\).

This is not a shortcut --- it is the formula.
\(S_{100} = \frac{100}{2}(1 + 100) = 50 \times 101 = 5050\).

The formula encodes the same pairing, for any arithmetic series. That is
what elegant mathematical thinking looks like: notice a structure, use
it, and get there in seconds instead of hours.

\end{minipage}%
\end{tcolorbox}

\subsection{Geometric series}\label{geometric-series}

A geometric sequence has a first term \(a_1\) and a common ratio \(r\).
The terms are:

\[a_1,\quad a_1 r,\quad a_1 r^2,\quad \ldots,\quad a_1 r^{n-1}\]

Adding all \(n\) terms:

\[S_n = a_1 + a_1 r + a_1 r^2 + \cdots + a_1 r^{n-1}\]

If you multiply the sum by \(r\), each term shifts up by one power ---
almost matching the original sum. Subtracting cancels all but the first
and last terms.

The trick here is to multiply both sides by \(r\):

\[r S_n = a_1 r + a_1 r^2 + a_1 r^3 + \cdots + a_1 r^n\]

Now subtract the first line from the second. Nearly everything cancels:

\[r S_n - S_n = a_1 r^n - a_1\]

\[S_n(r - 1) = a_1(r^n - 1)\]

\[\boxed{S_n = \frac{a_1(1 - r^n)}{1 - r}} \qquad (r \neq 1)\]

The form with \((1 - r^n)\) in the numerator is the convention --- both
forms are algebraically identical, but this one is easier to read when
\(|r| < 1\).

When \(r = 1\) all terms are equal to \(a_1\), so \(S_n = n \cdot a_1\).
That case doesn't need the formula.

\subsection{Infinite geometric series}\label{infinite-geometric-series}

Here is something that sounds impossible: add infinitely many numbers
and get a finite result.

Look at the partial sums of the series
\(1 + \frac{1}{2} + \frac{1}{4} + \frac{1}{8} + \cdots\)

\[S_1 = 1\] \[S_2 = 1.5\] \[S_3 = 1.75\] \[S_4 = 1.875\]
\[S_5 = 1.9375\]

The sums are getting closer to 2 --- but never exceeding it, and never
quite reaching it. The terms are shrinking. Each new term adds less than
the last. The sum is converging toward a fixed value.

Look at the formula \(S_n = \frac{a_1(1 - r^n)}{1 - r}\) when
\(|r| < 1\). As \(n\) grows, \(r^n\) gets closer and closer to 0 --- and
when \(|r| < 1\), eventually it is so small we treat it as 0:

\[S_\infty = \frac{a_1(1 - 0)}{1 - r} = \frac{a_1}{1 - r}\]

\[\boxed{S_\infty = \frac{a_1}{1 - r}} \qquad (|r| < 1)\]

Read \(S_\infty\) as \emph{the sum of infinitely many terms}.

For the series above: \(a_1 = 1\), \(r = \frac{1}{2}\). So
\(S_\infty = \frac{1}{1 - \frac{1}{2}} = \frac{1}{\frac{1}{2}} = 2\).
Exactly 2.

\textbf{What if \(|r| \geq 1\)?} Then \(r^n\) does not approach 0 --- it
either stays at 1 (when \(r = 1\)) or grows without bound (when
\(|r| > 1\)) or oscillates growing (when \(r \leq -1\)). The terms never
shrink to nothing, so the sum never settles. It \textbf{diverges} --- it
has no finite value. You can add as many terms as you like and the sum
will keep growing (or oscillating further and further from zero).

This is worth sitting with for a moment. It's genuinely strange that you
can add infinitely many things and get a finite answer. The reason it
works is not that infinity is somehow small --- it's that the terms
shrink fast enough. Halving at every step is fast enough. Shrinking by
one percent at every step is fast enough. Staying the same size is not.

\begin{Shaded}
\begin{Highlighting}[]
\NormalTok{//| echo: false}
\NormalTok{\{}
\NormalTok{  const div = document.createElement("div");}
\NormalTok{  div.style.cssText = "font{-}family:inherit; padding:0.5rem 0;";}

\NormalTok{  const controls = document.createElement("div");}
\NormalTok{  controls.style.cssText = "display:flex; gap:1.5rem; flex{-}wrap:wrap; margin{-}bottom:0.75rem;";}

\NormalTok{  function makeSlider2(label, min, max, step, val) \{}
\NormalTok{    const wrap = document.createElement("div");}
\NormalTok{    const lbl = document.createElement("label");}
\NormalTok{    lbl.style.cssText = "font{-}size:0.85rem; color:\#374151; display:block; margin{-}bottom:0.2rem;";}
\NormalTok{    const span = document.createElement("span"); span.textContent = val;}
\NormalTok{    lbl.appendChild(document.createTextNode(label + " = ")); lbl.appendChild(span);}
\NormalTok{    const input = document.createElement("input");}
\NormalTok{    input.type = "range"; input.min = min; input.max = max; input.step = step; input.value = val;}
\NormalTok{    input.style.width = "160px";}
\NormalTok{    wrap.appendChild(lbl); wrap.appendChild(input);}
\NormalTok{    return \{wrap, input, span\};}
\NormalTok{  \}}

\NormalTok{  const rSlider = makeSlider2("r (ratio)", {-}99, 99, 1, 50);}
\NormalTok{  const a1Slider = makeSlider2("a₁ (first term)", 1, 10, 1, 1);}
\NormalTok{  controls.appendChild(rSlider.wrap);}
\NormalTok{  controls.appendChild(a1Slider.wrap);}

\NormalTok{  const infoEl = document.createElement("div");}
\NormalTok{  infoEl.style.cssText = "font{-}size:0.88rem; color:\#374151; margin{-}bottom:0.5rem;";}

\NormalTok{  const canvas = document.createElement("canvas");}
\NormalTok{  canvas.width = 500; canvas.height = 260;}
\NormalTok{  canvas.style.cssText = "border:1px solid \#e5e7eb; border{-}radius:6px; display:block; max{-}width:500px;";}

\NormalTok{  div.appendChild(controls); div.appendChild(infoEl); div.appendChild(canvas);}

\NormalTok{  function draw() \{}
\NormalTok{    const rRaw = parseFloat(rSlider.input.value) / 100;}
\NormalTok{    const a1 = parseFloat(a1Slider.input.value);}
\NormalTok{    rSlider.span.textContent = rRaw.toFixed(2);}
\NormalTok{    a1Slider.span.textContent = a1;}

\NormalTok{    const N = 20;}
\NormalTok{    const partials = [];}
\NormalTok{    let sum = 0;}
\NormalTok{    for (let k = 1; k \textless{}= N; k++) \{}
\NormalTok{      sum += a1 * Math.pow(rRaw, k{-}1);}
\NormalTok{      partials.push(sum);}
\NormalTok{    \}}

\NormalTok{    const sInf = Math.abs(rRaw) \textless{} 1 ? a1 / (1 {-} rRaw) : null;}

\NormalTok{    if (sInf !== null) \{}
\NormalTok{      infoEl.textContent = "S∞ = " + a1 + " / (1 − " + rRaw.toFixed(2) + ") = " + sInf.toFixed(2) + "  (|r| = " + Math.abs(rRaw).toFixed(2) + " \textless{} 1, converges)";}
\NormalTok{    \} else \{}
\NormalTok{      infoEl.textContent = "|r| = " + Math.abs(rRaw).toFixed(2) + " ≥ 1 — series diverges";}
\NormalTok{    \}}

\NormalTok{    const ctx = canvas.getContext("2d");}
\NormalTok{    const W = canvas.width, H = canvas.height;}
\NormalTok{    ctx.clearRect(0, 0, W, H);}
\NormalTok{    const marg = \{l:55, r:20, t:15, b:35\};}
\NormalTok{    const pHeight = H {-} marg.t {-} marg.b;}
\NormalTok{    const pWidth = W {-} marg.l {-} marg.r;}

\NormalTok{    const allVals = sInf !== null ? [...partials, sInf] : partials;}
\NormalTok{    const yMin = Math.min(...allVals) * (Math.min(...allVals) \textless{} 0 ? 1.2 : 0.8);}
\NormalTok{    const yMax = Math.max(...allVals) * (Math.max(...allVals) \textgreater{} 0 ? 1.2 : 0.8);}
\NormalTok{    const yRange = yMax {-} yMin || 1;}

\NormalTok{    function px(k) \{ return marg.l + ((k{-}1) / (N{-}1)) * pWidth; \}}
\NormalTok{    function py(v) \{ return marg.t + (1 {-} (v {-} yMin)/yRange) * pHeight; \}}

\NormalTok{    // y{-}axis}
\NormalTok{    ctx.strokeStyle = "\#9ca3af"; ctx.lineWidth = 1;}
\NormalTok{    ctx.beginPath(); ctx.moveTo(marg.l, marg.t); ctx.lineTo(marg.l, marg.t+pHeight); ctx.stroke();}
\NormalTok{    ctx.beginPath(); ctx.moveTo(marg.l, marg.t+pHeight); ctx.lineTo(marg.l+pWidth, marg.t+pHeight); ctx.stroke();}

\NormalTok{    // Dashed S∞ line}
\NormalTok{    if (sInf !== null \&\& sInf \textgreater{} yMin \&\& sInf \textless{} yMax) \{}
\NormalTok{      ctx.strokeStyle = "\#f59e0b"; ctx.lineWidth = 1.5; ctx.setLineDash([6,4]);}
\NormalTok{      ctx.beginPath(); ctx.moveTo(marg.l, py(sInf)); ctx.lineTo(marg.l+pWidth, py(sInf)); ctx.stroke();}
\NormalTok{      ctx.setLineDash([]);}
\NormalTok{      ctx.fillStyle = "\#d97706"; ctx.font = "11px sans{-}serif";}
\NormalTok{      ctx.fillText("S∞ = " + sInf.toFixed(2), marg.l+pWidth{-}60, py(sInf){-}5);}
\NormalTok{    \}}

\NormalTok{    // Partial sums dots + line}
\NormalTok{    ctx.strokeStyle = "\#2563eb"; ctx.lineWidth = 1.5;}
\NormalTok{    ctx.beginPath();}
\NormalTok{    partials.forEach((v, i) =\textgreater{} \{}
\NormalTok{      if (i === 0) ctx.moveTo(px(i+1), py(v)); else ctx.lineTo(px(i+1), py(v));}
\NormalTok{    \});}
\NormalTok{    ctx.stroke();}

\NormalTok{    ctx.fillStyle = "\#2563eb";}
\NormalTok{    partials.forEach((v, i) =\textgreater{} \{}
\NormalTok{      ctx.beginPath(); ctx.arc(px(i+1), py(v), 3.5, 0, 2*Math.PI); ctx.fill();}
\NormalTok{    \});}

\NormalTok{    // x{-}axis labels}
\NormalTok{    ctx.fillStyle = "\#6b7280"; ctx.font = "10px sans{-}serif";}
\NormalTok{    [1,5,10,15,20].forEach(k =\textgreater{} \{}
\NormalTok{      ctx.fillText(k, px(k){-}4, marg.t+pHeight+14);}
\NormalTok{    \});}
\NormalTok{    ctx.fillText("n", marg.l+pWidth+5, marg.t+pHeight+14);}
\NormalTok{    ctx.fillText("Sₙ", marg.l{-}30, marg.t+pHeight/2);}
\NormalTok{  \}}

\NormalTok{  [rSlider, a1Slider].forEach(s =\textgreater{} s.input.addEventListener("input", draw));}
\NormalTok{  draw();}
\NormalTok{  return div;}
\NormalTok{\}}
\end{Highlighting}
\end{Shaded}

\subsection{Sigma notation: working with
sums}\label{sigma-notation-working-with-sums}

A few practical moves.

\textbf{Evaluating a sum.} To evaluate
\(\displaystyle\sum_{i=1}^{4}(3i - 1)\), substitute \(i = 1, 2, 3, 4\)
and add:

\[(3 \cdot 1 - 1) + (3 \cdot 2 - 1) + (3 \cdot 3 - 1) + (3 \cdot 4 - 1)\]
\[= 2 + 5 + 8 + 11 = 26\]

You can check this with the arithmetic series formula: \(a_1 = 2\),
\(d = 3\), \(n = 4\). \(S_4 = \frac{4}{2}(2 + 11) = 2 \times 13 = 26\).
Same answer.

\textbf{Changing the index.} The index variable is a dummy --- it
doesn't matter what letter you use. \(\displaystyle\sum_{i=1}^{n} a_i\)
and \(\displaystyle\sum_{k=1}^{n} a_k\) are the same sum.

\textbf{Splitting a sum.} You can split a sum that has two parts:

\[\sum_{i=1}^{n}(a_i + b_i) = \sum_{i=1}^{n} a_i + \sum_{i=1}^{n} b_i\]

And pull a constant out:

\[\sum_{i=1}^{n} c \cdot a_i = c \sum_{i=1}^{n} a_i\]

\section{Worked examples}\label{worked-examples-16}

\textbf{Example 1 (Finance).} A car loan of \$12 000 charges 0.5\%
interest per month. You plan to repay in a lump sum after exactly 24
months --- no intermediate payments, interest compounds monthly. What is
the total you owe at the end?

Each month the balance is multiplied by \(1.005\). After 24 months:

\[A = 12000 \times (1.005)^{24}\]

\((1.005)^{24} \approx 1.1272\), so
\(A \approx 12000 \times 1.1272 = \$13\,526.40\).

Now ask: what is the total interest paid? That is
\(A - 12000 = \$1\,526.40\).

The balance itself is a geometric sequence with first term \(12000\) and
ratio \(1.005\). The total interest is what the geometric growth added
over 24 steps.

\begin{center}\rule{0.5\linewidth}{0.5pt}\end{center}

\textbf{Example 2 (Computing).} Merge sort splits an array of \(n\)
elements in half repeatedly, sorts each half, and merges. The merging
step at each level does work proportional to \(n\). The number of levels
is \(\log_2 n\). The total work can be written as:

\[T = n + n + n + \cdots \quad (\log_2 n \text{ terms})\]

This is an arithmetic series with all terms equal to \(n\):

\[T = n \cdot \log_2 n\]

This is the source of the famous \(O(n \log n)\) complexity of merge
sort. Each level of the recursion tree contributes \(n\) operations;
summing \(\log_2 n\) such levels gives the total.

\begin{center}\rule{0.5\linewidth}{0.5pt}\end{center}

\textbf{Example 3 (Science).} Zeno's paradox resolved. A runner is 8
metres from a wall. She runs to the halfway point (4 m), then halfway
again (2 m), then halfway again (1 m), and so on. The distances form a
geometric sequence:

\[4,\quad 2,\quad 1,\quad \frac{1}{2},\quad \frac{1}{4},\quad \ldots\]

Here \(a_1 = 4\), \(r = \frac{1}{2}\), and \(|r| < 1\) so the infinite
sum converges:

\[S_\infty = \frac{4}{1 - \frac{1}{2}} = \frac{4}{\frac{1}{2}} = 8\]

The infinite sum of all the halfway-steps is exactly 8 metres --- the
distance to the wall. The paradox dissolves: infinitely many steps do
not require infinite time or distance. They require exactly the right
amount of both.

\begin{center}\rule{0.5\linewidth}{0.5pt}\end{center}

\textbf{Example 4 (Engineering).} A simple periodic signal can be
approximated by a partial sum of sine waves. Consider the first five odd
harmonics of a square wave at frequency \(f\):

\[s(t) \approx \sum_{k=0}^{3} \frac{1}{2k+1}\sin\bigl((2k+1) \cdot 2\pi f t\bigr)\]

Expanding the first four terms (\(k = 0, 1, 2, 3\)):

\[s(t) \approx \sin(2\pi f t) + \frac{1}{3}\sin(6\pi f t) + \frac{1}{5}\sin(10\pi f t) + \frac{1}{7}\sin(14\pi f t)\]

The coefficients \(1, \frac{1}{3}, \frac{1}{5}, \frac{1}{7}\) form a
series --- the terms shrink, so the infinite sum converges. More terms
give a closer approximation to the square wave. This is Fourier series:
decomposing a complex signal as an infinite sum of simple ones. The tool
you need to make sense of it is exactly the infinite series convergence
you just learned.

You might notice the coefficients
\(1, \frac{1}{3}, \frac{1}{5}, \ldots\) shrink slowly --- the series
\(1 + \frac{1}{3} + \frac{1}{5} + \cdots\) does diverge on its own (like
the harmonic series). But in the Fourier series each coefficient
multiplies an oscillating sine function. Those sines partially cancel
each other out, which is what allows the series to converge. The full
explanation requires calculus.

\section{Where this goes}\label{where-this-goes-17}

The phrase ``approaches a finite value as \(n \to \infty\)'' was used
informally here. \textbf{Limits and continuity} (Vol 5, Ch 1) makes that
precise: what does it mean, formally, for a sequence to converge to a
number? The epsilon-delta definition is the rigorous version of the
intuition you've been building. Every infinite series you have seen in
this chapter is, underneath, a limit.

\textbf{Integral calculus} (Vol 5, Ch 3) connects series to continuous
accumulation. A Riemann sum --- the standard way to define the integral
--- is a series: you divide an area into strips, compute the area of
each strip, and sum. As the number of strips goes to infinity, the
series converges to the exact area. The passage from ``finite sum'' to
``integral'' is the same passage from ``series'' to ``infinite series''
you made in this chapter. \textbf{Fourier series} (Vol 7) is the direct
destination of example 4: a systematic way to represent any periodic
function as an infinite sum of sines and cosines. Every result in
Fourier analysis stands on the geometric series convergence you proved
here.

\begin{tcolorbox}[enhanced jigsaw, opacityback=0, colback=white, breakable, colframe=quarto-callout-tip-color-frame, arc=.35mm, left=2mm, toprule=.15mm, leftrule=.75mm, bottomrule=.15mm, rightrule=.15mm]
\begin{minipage}[t]{5.5mm}
\textcolor{quarto-callout-tip-color}{\faLightbulb}
\end{minipage}%
\begin{minipage}[t]{\textwidth - 5.5mm}

\textbf{Where this shows up}

\begin{itemize}
\tightlist
\item
  A mortgage calculator computes the total repayment using the geometric
  series formula for present value of an annuity --- every bank does
  this calculation thousands of times a day.
\item
  Merge sort, binary search, and most divide-and-conquer algorithms have
  costs that are sums of a geometric series --- this is why their Big-O
  complexity involves logarithms.
\item
  Numerical methods for differential equations (Runge-Kutta, finite
  differences) rely on Taylor series to approximate function values ---
  an infinite polynomial series truncated at a practical number of
  terms.
\item
  Audio compression (MP3, AAC) works by computing a Fourier series
  decomposition, discarding components the ear cannot hear, and
  reconstructing the signal from the remainder.
\end{itemize}

\end{minipage}%
\end{tcolorbox}

\section{Exercises}\label{exercises-17}

These are puzzles. Each has a clean numerical answer, but the
interesting part is identifying the series type and setting up the
formula before you calculate.

\begin{Shaded}
\begin{Highlighting}[]
\NormalTok{//| echo: false}
\NormalTok{katex = require("katex@0.16")}
\end{Highlighting}
\end{Shaded}

\begin{Shaded}
\begin{Highlighting}[]
\NormalTok{//| echo: false}
\NormalTok{function makeStepperHTML(exerciseNum, steps) \{}
\NormalTok{  let current = 0;}
\NormalTok{  const totalSteps = steps.length;}
\NormalTok{  const container = document.createElement("div");}
\NormalTok{  container.style.cssText = "border:1px solid \#e5e7eb; border{-}radius:8px; padding:1rem 1.25rem; margin:0.75rem 0; font{-}family:inherit;";}
\NormalTok{  const stepsDiv = document.createElement("div");}
\NormalTok{  stepsDiv.style.marginBottom = "0.75rem";}
\NormalTok{  function renderSteps() \{}
\NormalTok{    stepsDiv.innerHTML = "";}
\NormalTok{    for (let i = 0; i \textless{} current; i++) \{}
\NormalTok{      const s = steps[i];}
\NormalTok{      const row = document.createElement("div");}
\NormalTok{      row.style.cssText = "display:grid; grid{-}template{-}columns:200px 1fr; gap:0.35rem 1rem; align{-}items:baseline; padding:0.35rem 0; border{-}top:1px solid \#f3f4f6;";}
\NormalTok{      const opEl = document.createElement("span");}
\NormalTok{      opEl.style.cssText = "font{-}size:0.8rem; color:\#6b7280; font{-}style:italic;";}
\NormalTok{      opEl.textContent = s.op;}
\NormalTok{      const eqEl = document.createElement("span");}
\NormalTok{      eqEl.innerHTML = katex.renderToString(s.eq, \{throwOnError: false, displayMode: false\});}
\NormalTok{      row.appendChild(opEl);}
\NormalTok{      row.appendChild(eqEl);}
\NormalTok{      container.appendChild(row);}
\NormalTok{      if (s.note) \{}
\NormalTok{        const noteEl = document.createElement("p");}
\NormalTok{        noteEl.style.cssText = "margin:0.2rem 0 0.5rem 200px; font{-}size:0.82rem; color:\#374151;";}
\NormalTok{        noteEl.textContent = s.note;}
\NormalTok{        container.appendChild(noteEl);}
\NormalTok{      \}}
\NormalTok{    \}}
\NormalTok{    if (current === totalSteps) \{}
\NormalTok{      const done = document.createElement("p");}
\NormalTok{      done.style.cssText = "margin{-}top:0.5rem; font{-}size:0.85rem; color:\#059669; font{-}weight:600;";}
\NormalTok{      done.textContent = "✓ Solution complete";}
\NormalTok{      stepsDiv.appendChild(done);}
\NormalTok{    \}}
\NormalTok{  \}}
\NormalTok{  const controls = document.createElement("div");}
\NormalTok{  controls.style.cssText = "display:flex; align{-}items:center; gap:0.75rem; margin{-}top:0.5rem;";}
\NormalTok{  const nextBtn = document.createElement("button");}
\NormalTok{  nextBtn.textContent = "Next step →";}
\NormalTok{  nextBtn.style.cssText = "padding:0.35rem 0.85rem; background:\#2563eb; color:\#fff; border:none; border{-}radius:6px; cursor:pointer; font{-}size:0.85rem;";}
\NormalTok{  const resetBtn = document.createElement("button");}
\NormalTok{  resetBtn.textContent = "Reset";}
\NormalTok{  resetBtn.style.cssText = "padding:0.35rem 0.85rem; background:\#f3f4f6; color:\#374151; border:1px solid \#d1d5db; border{-}radius:6px; cursor:pointer; font{-}size:0.85rem;";}
\NormalTok{  const counter = document.createElement("span");}
\NormalTok{  counter.style.cssText = "font{-}size:0.8rem; color:\#6b7280;";}
\NormalTok{  nextBtn.onclick = () =\textgreater{} \{ if (current \textless{} totalSteps) \{ current++; renderSteps(); counter.textContent = current + " / " + totalSteps; if (current === totalSteps) nextBtn.disabled = true; \} \};}
\NormalTok{  resetBtn.onclick = () =\textgreater{} \{ current = 0; nextBtn.disabled = false; renderSteps(); counter.textContent = "0 / " + totalSteps; \};}
\NormalTok{  counter.textContent = "0 / " + totalSteps;}
\NormalTok{  controls.appendChild(nextBtn);}
\NormalTok{  controls.appendChild(resetBtn);}
\NormalTok{  controls.appendChild(counter);}
\NormalTok{  container.appendChild(stepsDiv);}
\NormalTok{  container.appendChild(controls);}
\NormalTok{  return container;}
\NormalTok{\}}
\end{Highlighting}
\end{Shaded}

\textbf{Exercise 1.} The first term of an arithmetic series is 5, the
common difference is 4, and there are 20 terms. Find the sum.

\begin{Shaded}
\begin{Highlighting}[]
\NormalTok{//| echo: false}
\NormalTok{\{}
\NormalTok{  const el = makeStepperHTML(1, [}
\NormalTok{    \{ op: "Identify the values", eq: "a\_1 = 5,\textbackslash{}\textbackslash{}quad d = 4,\textbackslash{}\textbackslash{}quad n = 20", note: null \},}
\NormalTok{    \{ op: "Write the sum formula", eq: "S\_n = \textbackslash{}\textbackslash{}frac\{n\}\{2\}\textbackslash{}\textbackslash{}bigl(2a\_1 + (n{-}1)d\textbackslash{}\textbackslash{}bigr)", note: "Use this form when you know a₁ and d but not the last term." \},}
\NormalTok{    \{ op: "Substitute", eq: "S\_\{20\} = \textbackslash{}\textbackslash{}frac\{20\}\{2\}\textbackslash{}\textbackslash{}bigl(2 \textbackslash{}\textbackslash{}times 5 + 19 \textbackslash{}\textbackslash{}times 4\textbackslash{}\textbackslash{}bigr) = 10\textbackslash{}\textbackslash{}bigl(10 + 76\textbackslash{}\textbackslash{}bigr)", note: null \},}
\NormalTok{    \{ op: "Simplify", eq: "S\_\{20\} = 10 \textbackslash{}\textbackslash{}times 86 = 860", note: null \},}
\NormalTok{    \{ op: "Check", eq: "a\_\{20\} = 5 + 19 \textbackslash{}\textbackslash{}times 4 = 81;\textbackslash{}\textbackslash{}quad S\_\{20\} = \textbackslash{}\textbackslash{}frac\{20\}\{2\}(5 + 81) = 10 \textbackslash{}\textbackslash{}times 86 = 860 \textbackslash{}\textbackslash{}checkmark", note: null \},}
\NormalTok{  ]);}
\NormalTok{  return el;}
\NormalTok{\}}
\end{Highlighting}
\end{Shaded}

\begin{center}\rule{0.5\linewidth}{0.5pt}\end{center}

\textbf{Exercise 2.} A geometric series has first term 3 and common
ratio 2. Find the sum of the first 8 terms.

\begin{Shaded}
\begin{Highlighting}[]
\NormalTok{//| echo: false}
\NormalTok{\{}
\NormalTok{  const el = makeStepperHTML(2, [}
\NormalTok{    \{ op: "Identify the values", eq: "a\_1 = 3,\textbackslash{}\textbackslash{}quad r = 2,\textbackslash{}\textbackslash{}quad n = 8", note: null \},}
\NormalTok{    \{ op: "Write the sum formula", eq: "S\_n = \textbackslash{}\textbackslash{}frac\{a\_1(1 {-} r\^{}n)\}\{1 {-} r\}", note: null \},}
\NormalTok{    \{ op: "Substitute", eq: "S\_8 = \textbackslash{}\textbackslash{}frac\{3(1 {-} 2\^{}8)\}\{1 {-} 2\} = \textbackslash{}\textbackslash{}frac\{3(1 {-} 256)\}\{{-}1\}", note: null \},}
\NormalTok{    \{ op: "Simplify numerator", eq: "S\_8 = \textbackslash{}\textbackslash{}frac\{3 \textbackslash{}\textbackslash{}times ({-}255)\}\{{-}1\} = \textbackslash{}\textbackslash{}frac\{{-}765\}\{{-}1\}", note: "Dividing a negative by a negative gives a positive." \},}
\NormalTok{    \{ op: "Compute", eq: "S\_8 = 765", note: null \},}
\NormalTok{    \{ op: "Check", eq: "3 + 6 + 12 + 24 + 48 + 96 + 192 + 384 = 765 \textbackslash{}\textbackslash{}checkmark", note: null \},}
\NormalTok{  ]);}
\NormalTok{  return el;}
\NormalTok{\}}
\end{Highlighting}
\end{Shaded}

\begin{center}\rule{0.5\linewidth}{0.5pt}\end{center}

\textbf{Exercise 3.} A bouncing ball is dropped from 10 m. Each bounce
reaches 60\% of the previous height. Find the total distance travelled
by the ball before it comes to rest.

\emph{(Hint: the ball travels down, up, down, up, \ldots{} Be careful to
count both the downward and upward legs of each bounce.)}

\begin{Shaded}
\begin{Highlighting}[]
\NormalTok{//| echo: false}
\NormalTok{\{}
\NormalTok{  const el = makeStepperHTML(3, [}
\NormalTok{    \{ op: "Set up the distances", eq: "\textbackslash{}\textbackslash{}text\{Down: \} 10 + 6 + 3.6 + \textbackslash{}\textbackslash{}cdots;\textbackslash{}\textbackslash{}quad \textbackslash{}\textbackslash{}text\{Up: \} 6 + 3.6 + 2.16 + \textbackslash{}\textbackslash{}cdots", note: "The first drop (10 m) has no matching upward leg, so split the sum." \},}
\NormalTok{    \{ op: "Sum the downward legs", eq: "S\_\{\textbackslash{}\textbackslash{}text\{down\}\} = \textbackslash{}\textbackslash{}frac\{10\}\{1 {-} 0.6\} = \textbackslash{}\textbackslash{}frac\{10\}\{0.4\} = 25 \textbackslash{}\textbackslash{}text\{ m\}", note: "Infinite geometric series: a₁ = 10, r = 0.6." \},}
\NormalTok{    \{ op: "Sum the upward legs", eq: "S\_\{\textbackslash{}\textbackslash{}text\{up\}\} = \textbackslash{}\textbackslash{}frac\{6\}\{1 {-} 0.6\} = \textbackslash{}\textbackslash{}frac\{6\}\{0.4\} = 15 \textbackslash{}\textbackslash{}text\{ m\}", note: "The ball first rises to 6 m after the first bounce." \},}
\NormalTok{    \{ op: "Total distance", eq: "S\_\{\textbackslash{}\textbackslash{}text\{total\}\} = 25 + 15 = 40 \textbackslash{}\textbackslash{}text\{ m\}", note: null \},}
\NormalTok{    \{ op: "Check (alternative formula)", eq: "S\_\{\textbackslash{}\textbackslash{}text\{total\}\} = 10 + 2 \textbackslash{}\textbackslash{}times \textbackslash{}\textbackslash{}frac\{6\}\{1{-}0.6\} = 10 + 30 = 40 \textbackslash{}\textbackslash{}text\{ m\} \textbackslash{}\textbackslash{}checkmark", note: "Adding the initial drop once and all bounce arcs (up + down) twice." \},}
\NormalTok{  ]);}
\NormalTok{  return el;}
\NormalTok{\}}
\end{Highlighting}
\end{Shaded}

\begin{center}\rule{0.5\linewidth}{0.5pt}\end{center}

\textbf{Exercise 4.} Evaluate the sum
\(\displaystyle\sum_{i=1}^{6}(2i + 3)\).

\emph{(First expand and add term by term to check your answer, then
verify using the arithmetic series formula.)}

\begin{Shaded}
\begin{Highlighting}[]
\NormalTok{//| echo: false}
\NormalTok{\{}
\NormalTok{  const el = makeStepperHTML(4, [}
\NormalTok{    \{ op: "Expand the sum", eq: "(2 \textbackslash{}\textbackslash{}cdot 1 + 3) + (2 \textbackslash{}\textbackslash{}cdot 2 + 3) + (2 \textbackslash{}\textbackslash{}cdot 3 + 3) + (2 \textbackslash{}\textbackslash{}cdot 4 + 3) + (2 \textbackslash{}\textbackslash{}cdot 5 + 3) + (2 \textbackslash{}\textbackslash{}cdot 6 + 3)", note: null \},}
\NormalTok{    \{ op: "Compute each term", eq: "5 + 7 + 9 + 11 + 13 + 15", note: null \},}
\NormalTok{    \{ op: "Add", eq: "\textbackslash{}\textbackslash{}sum\_\{i=1\}\^{}\{6\}(2i+3) = 60", note: null \},}
\NormalTok{    \{ op: "Verify using the formula", eq: "a\_1 = 5,\textbackslash{}\textbackslash{}quad a\_6 = 15,\textbackslash{}\textbackslash{}quad n = 6 \textbackslash{}\textbackslash{}implies S\_6 = \textbackslash{}\textbackslash{}frac\{6\}\{2\}(5 + 15) = 3 \textbackslash{}\textbackslash{}times 20 = 60 \textbackslash{}\textbackslash{}checkmark", note: "The terms 5, 7, 9, 11, 13, 15 form an arithmetic sequence with d = 2." \},}
\NormalTok{  ]);}
\NormalTok{  return el;}
\NormalTok{\}}
\end{Highlighting}
\end{Shaded}

\begin{center}\rule{0.5\linewidth}{0.5pt}\end{center}

\textbf{Exercise 5.} An apprentice electrician earns \$18 per hour in
their first year and gets a \$1.50 per hour raise every year. They plan
to work in this trade for 12 years. Their annual hours are 1820 each
year. What are their total earnings over the 12 years?

\emph{(Note: compute total annual earnings for each year first, then sum
the 12 annual totals.)}

\begin{Shaded}
\begin{Highlighting}[]
\NormalTok{//| echo: false}
\NormalTok{\{}
\NormalTok{  const el = makeStepperHTML(5, [}
\NormalTok{    \{ op: "Compute year 1 earnings", eq: "E\_1 = 18 \textbackslash{}\textbackslash{}times 1820 = \textbackslash{}\textbackslash{}$32\{,\}760", note: null \},}
\NormalTok{    \{ op: "Identify the series", eq: "a\_1 = 32760,\textbackslash{}\textbackslash{}quad d = 1.50 \textbackslash{}\textbackslash{}times 1820 = 2730,\textbackslash{}\textbackslash{}quad n = 12", note: "Each year\textquotesingle{}s earnings increase by $1.50 × 1820 = $2730 over the previous year." \},}
\NormalTok{    \{ op: "Apply the arithmetic series formula", eq: "S\_\{12\} = \textbackslash{}\textbackslash{}frac\{12\}\{2\}\textbackslash{}\textbackslash{}bigl(2 \textbackslash{}\textbackslash{}times 32760 + 11 \textbackslash{}\textbackslash{}times 2730\textbackslash{}\textbackslash{}bigr)", note: null \},}
\NormalTok{    \{ op: "Simplify", eq: "S\_\{12\} = 6\textbackslash{}\textbackslash{}bigl(65520 + 30030\textbackslash{}\textbackslash{}bigr) = 6 \textbackslash{}\textbackslash{}times 95550", note: null \},}
\NormalTok{    \{ op: "Compute", eq: "S\_\{12\} = \textbackslash{}\textbackslash{}$573\{,\}300", note: null \},}
\NormalTok{    \{ op: "Check", eq: "E\_\{12\} = 18 + 11 \textbackslash{}\textbackslash{}times 1.50 = \textbackslash{}\textbackslash{}$34.50 \textbackslash{}\textbackslash{}text\{ per hour\};\textbackslash{}\textbackslash{}quad S\_\{12\} = \textbackslash{}\textbackslash{}frac\{12\}\{2\}(32760 + 34.50 \textbackslash{}\textbackslash{}times 1820) = 6(32760 + 62790) = 6 \textbackslash{}\textbackslash{}times 95550 = \textbackslash{}\textbackslash{}$573\{,\}300 \textbackslash{}\textbackslash{}checkmark", note: null \},}
\NormalTok{  ]);}
\NormalTok{  return el;}
\NormalTok{\}}
\end{Highlighting}
\end{Shaded}

\begin{center}\rule{0.5\linewidth}{0.5pt}\end{center}

\textbf{Exercise 6.} Show that the infinite series
\(\displaystyle\sum_{k=1}^{\infty} \frac{1}{4^k}\) converges, and find
its sum.

Then explain in one sentence why the harmonic series
\(1 + \frac{1}{2} + \frac{1}{3} + \frac{1}{4} + \cdots\) does \emph{not}
have a finite sum. (Hint: think about how quickly the terms are
shrinking compared to a geometric series.)

\begin{Shaded}
\begin{Highlighting}[]
\NormalTok{//| echo: false}
\NormalTok{\{}
\NormalTok{  const el = makeStepperHTML(6, [}
\NormalTok{    \{ op: "Write the first few terms", eq: "\textbackslash{}\textbackslash{}frac\{1\}\{4\} + \textbackslash{}\textbackslash{}frac\{1\}\{16\} + \textbackslash{}\textbackslash{}frac\{1\}\{64\} + \textbackslash{}\textbackslash{}cdots", note: null \},}
\NormalTok{    \{ op: "Identify a₁ and r", eq: "a\_1 = \textbackslash{}\textbackslash{}frac\{1\}\{4\},\textbackslash{}\textbackslash{}quad r = \textbackslash{}\textbackslash{}frac\{1\}\{4\}", note: "Each term is one{-}quarter of the previous term." \},}
\NormalTok{    \{ op: "Check convergence condition", eq: "|r| = \textbackslash{}\textbackslash{}frac\{1\}\{4\} \textless{} 1 \textbackslash{}\textbackslash{}implies \textbackslash{}\textbackslash{}text\{converges\}", note: null \},}
\NormalTok{    \{ op: "Apply the infinite sum formula", eq: "S\_\{\textbackslash{}\textbackslash{}infty\} = \textbackslash{}\textbackslash{}frac\{a\_1\}\{1 {-} r\} = \textbackslash{}\textbackslash{}frac\{\textbackslash{}\textbackslash{}frac\{1\}\{4\}\}\{1 {-} \textbackslash{}\textbackslash{}frac\{1\}\{4\}\} = \textbackslash{}\textbackslash{}frac\{\textbackslash{}\textbackslash{}frac\{1\}\{4\}\}\{\textbackslash{}\textbackslash{}frac\{3\}\{4\}\} = \textbackslash{}\textbackslash{}frac\{1\}\{3\}", note: null \},}
\NormalTok{    \{ op: "Reasoning about the harmonic series", eq: "\textbackslash{}\textbackslash{}frac\{a\_\{k+1\}\}\{a\_k\} = \textbackslash{}\textbackslash{}frac\{1/(k+1)\}\{1/k\} = \textbackslash{}\textbackslash{}frac\{k\}\{k+1\} \textbackslash{}\textbackslash{}to 1 \textbackslash{}\textbackslash{}text\{ as \} k \textbackslash{}\textbackslash{}to \textbackslash{}\textbackslash{}infty", note: "The harmonic series is not geometric — there is no fixed ratio r. The ratio between consecutive terms is k/(k+1), which approaches 1 as k grows. The terms shrink, but more and more slowly — not fast enough to produce a finite sum, unlike a geometric series where |r| is fixed below 1." \},}
\NormalTok{    \{ op: "Check", eq: "\textbackslash{}\textbackslash{}frac\{1\}\{4\} + \textbackslash{}\textbackslash{}frac\{1\}\{16\} + \textbackslash{}\textbackslash{}frac\{1\}\{64\} + \textbackslash{}\textbackslash{}frac\{1\}\{256\} \textbackslash{}\textbackslash{}approx 0.332 \textbackslash{}\textbackslash{}approx \textbackslash{}\textbackslash{}frac\{1\}\{3\} \textbackslash{}\textbackslash{}checkmark", note: null \},}
\NormalTok{  ]);}
\NormalTok{  return el;}
\NormalTok{\}}
\end{Highlighting}
\end{Shaded}

\part{Volume 5 --- Analysis and proof}

\chapter{Volume 5 --- Analysis and
proof}\label{volume-5-analysis-and-proof-1}

The central move of this volume is taming the infinite.

Up until now, you've relied on assumptions about how numbers behave
between points. This volume strips away those assumptions and rebuilds
mathematics on a rigorous foundation of limits, instantaneous rates of
change, and continuous accumulation.

By the end of this volume, you will understand how to freeze a changing
system precisely at a moment in time, how to add up an infinite number
of infinitely small pieces, and how to construct a rock-solid logical
argument that proves why something must guaranteeably be true.

\section{Chapters}\label{chapters-4}

\begin{longtable}[]{@{}lll@{}}
\toprule\noalign{}
Chapter & Title & Core move \\
\midrule\noalign{}
\endhead
\bottomrule\noalign{}
\endlastfoot
1 & Limits and continuity & Approaching the boundary \\
2 & Differential calculus & Measuring instantaneous change \\
3 & Integral calculus & Accumulating continuous amounts \\
4 & Proof and reasoning & Establishing absolute certainty \\
\end{longtable}

\chapter{Limits and continuity}\label{limits-and-continuity}

What a function approaches, not what it equals

\hfill\break

Your car's speedometer reads 80 km/h at this exact instant. But speed is
distance divided by time --- and at a single instant, no time has
passed, no distance covered. So what is the speedometer actually
measuring?

The answer is a limit. It's the value that the ratio distance/time
approaches as the time interval shrinks toward zero --- not the ratio at
zero, which is 0/0 and meaningless, but the value that ratio gets closer
and closer to.

You've already seen this pattern in another form. In the last chapter
you watched sequences converge to numbers they never quite reached. A
limit is that same idea applied to functions.

\section{The idea before the
notation}\label{the-idea-before-the-notation}

Three situations will make this concrete. Notice what they have in
common.

\textbf{The speedometer.} A car's position at time \(t\) is some
function \(s(t)\). To estimate speed, you record position over a short
interval: speed \(\approx
(s(t + \Delta t) - s(t)) / \Delta t\). The shorter the interval, the
better the estimate. The true instantaneous speed is what this ratio
\emph{approaches} as \(\Delta t\) shrinks toward zero. The ratio itself
is never evaluated at \(\Delta t = 0\) --- that would be \(0/0\) --- but
it heads steadily toward a definite value.

\textbf{The temperature probe.} A weather sensor is mounted on a
boundary between two air masses --- one at 18°C and one at 12°C. On the
warm side, readings approach 18°C. On the cold side, readings approach
12°C. The sensor itself sits exactly at the boundary and might read
either, neither, or something else entirely depending on how it was
calibrated. The function has a well-defined \emph{approach} from each
side. Whether the value at the boundary matches either approach is a
separate question.

\textbf{Sequences converging to a number.} In Vol 4 ch 4 you watched the
sequence \(1, \frac{1}{2}, \frac{1}{4}, \frac{1}{8}, \ldots\) converge
to zero. The terms never equal zero, but they get arbitrarily close.
Functions can do exactly the same thing: a function can get arbitrarily
close to a value \(L\) as \(x\) approaches some point \(a\), without
ever equalling \(L\) at \(x = a\).

All three situations share the same structure: a value is
\emph{approached} without necessarily being \emph{reached}. That is the
concept of a limit.

\section{What the notation is
saying}\label{what-the-notation-is-saying-18}

The symbol

\[\lim_{x \to a} f(x) = L\]

is read: \emph{the limit of \(f(x)\) as \(x\) approaches \(a\) equals
\(L\)}.

Every piece matters:

\begin{itemize}
\tightlist
\item
  \(x \to a\) means \(x\) gets closer and closer to \(a\), from either
  side, but does not equal \(a\).
\item
  \(f(x)\) is the output of the function at each nearby \(x\).
\item
  \(L\) is the value that \(f(x)\) approaches.
\end{itemize}

The limit says nothing about \(f(a)\) itself. The function might be
defined at \(a\), undefined at \(a\), or defined but with the wrong
value. The limit only asks about the \emph{approach}.

\subsection{One-sided limits}\label{one-sided-limits}

Sometimes the approach from the left and the approach from the right
give different values. Two separate notations handle this.

\[\lim_{x \to a^-} f(x)\]

is the \textbf{left-hand limit} --- the value \(f(x)\) approaches as
\(x\) comes in from values \emph{less than} \(a\). Read it: \emph{the
limit from the left}.

\[\lim_{x \to a^+} f(x)\]

is the \textbf{right-hand limit} --- the value \(f(x)\) approaches as
\(x\) comes in from values \emph{greater than} \(a\). Read it: \emph{the
limit from the right}.

\begin{tcolorbox}[enhanced jigsaw, title=\textcolor{quarto-callout-note-color}{\faInfo}\hspace{0.5em}{Note}, opacityback=0, colback=white, breakable, leftrule=.75mm, arc=.35mm, toptitle=1mm, left=2mm, colframe=quarto-callout-note-color-frame, coltitle=black, titlerule=0mm, colbacktitle=quarto-callout-note-color!10!white, opacitybacktitle=0.6, toprule=.15mm, bottomtitle=1mm, bottomrule=.15mm, rightrule=.15mm]

\textbf{The two-sided limit exists if and only if both one-sided limits
exist and are equal.} If the left-hand limit is \(L_1\) and the
right-hand limit is \(L_2\), then:

\[\lim_{x \to a} f(x) = L \iff \lim_{x \to a^-} f(x) = \lim_{x \to a^+} f(x) = L\]

If \(L_1 \neq L_2\), the two-sided limit does not exist.

\end{tcolorbox}

\subsection{The limit versus the value --- three
cases}\label{the-limit-versus-the-value-three-cases}

There is no rule that says \(\lim_{x \to a} f(x)\) must equal \(f(a)\).
In fact, that equality is the definition of continuity, and it can fail
in three genuinely distinct ways.

\textbf{Case 1 --- Continuous point.} The function is defined at \(a\),
the limit exists, and they agree:

\[\lim_{x \to a} f(x) = f(a)\]

Example: \(f(x) = x^2\) at \(a = 3\). The limit is 9, the value is 9.

\textbf{Case 2 --- Removable discontinuity (value mismatch).} The limit
exists, \(f(a)\) is defined, but they are unequal. The function has a
value at \(a\) --- it is simply the \emph{wrong} value relative to what
the surrounding behaviour predicts.

Example: define \(g(x) = x + 2\) for \(x \neq 2\), but \(g(2) = 7\). The
limit as \(x \to 2\) is 4. The function value at \(x = 2\) is 7. There
is a mismatch. The discontinuity can be removed by redefining
\(g(2) = 4\).

\textbf{Case 3 --- Hole (value absent).} The limit exists, but \(f(a)\)
is simply not defined. There is no value to mismatch --- the function
has a gap at \(x = a\).

Example: \(h(x) = \dfrac{x^2 - 4}{x - 2}\). At \(x = 2\) the denominator
is zero, so \(h(2)\) is undefined. But for every \(x \neq 2\):

\[h(x) = \frac{(x-2)(x+2)}{x-2} = x + 2\]

The limit is 4. The function has a hole at \((2, 4)\) --- no value is
there at all. Patching the hole by defining \(h(2) = 4\) would make
\(h\) continuous.

Cases 2 and 3 are genuinely different situations: in Case 2 you have a
value but it conflicts with the limit; in Case 3 you simply have
nothing. Both are called \emph{removable discontinuities} because the
fix --- redefining or defining \(f(a) = L\) --- is the same in both
cases.

\begin{Shaded}
\begin{Highlighting}[]
\NormalTok{//| echo: false}
\NormalTok{d3\_limits = require("d3@7")}
\end{Highlighting}
\end{Shaded}

\begin{Shaded}
\begin{Highlighting}[]
\NormalTok{//| echo: false}
\NormalTok{viewof xSlider = Inputs.range([0.5, 3.5], \{}
\NormalTok{  step: 0.01,}
\NormalTok{  value: 0.5,}
\NormalTok{  label: "x approaching 2"}
\NormalTok{\})}
\end{Highlighting}
\end{Shaded}

\begin{Shaded}
\begin{Highlighting}[]
\NormalTok{//| echo: false}
\NormalTok{\{}
\NormalTok{  const width = 560, height = 300;}
\NormalTok{  const margin = \{ top: 20, right: 20, bottom: 40, left: 45 \};}
\NormalTok{  const innerW = width {-} margin.left {-} margin.right;}
\NormalTok{  const innerH = height {-} margin.top {-} margin.bottom;}

\NormalTok{  const xDom = [{-}1, 5];}
\NormalTok{  const yDom = [0, 6];}

\NormalTok{  const xScale = d3\_limits.scaleLinear().domain(xDom).range([0, innerW]);}
\NormalTok{  const yScale = d3\_limits.scaleLinear().domain(yDom).range([innerH, 0]);}

\NormalTok{  const svg = d3\_limits.create("svg")}
\NormalTok{    .attr("width", width)}
\NormalTok{    .attr("height", height)}
\NormalTok{    .style("font{-}family", "inherit");}

\NormalTok{  const g = svg.append("g")}
\NormalTok{    .attr("transform", \textasciigrave{}translate($\{margin.left\},$\{margin.top\})\textasciigrave{});}

\NormalTok{  // Axes}
\NormalTok{  g.append("g")}
\NormalTok{    .attr("transform", \textasciigrave{}translate(0,$\{innerH\})\textasciigrave{})}
\NormalTok{    .call(d3\_limits.axisBottom(xScale).ticks(6))}
\NormalTok{    .append("text")}
\NormalTok{    .attr("x", innerW / 2).attr("y", 35)}
\NormalTok{    .attr("fill", "\#374151").attr("text{-}anchor", "middle")}
\NormalTok{    .style("font{-}size", "13px").text("x");}

\NormalTok{  g.append("g")}
\NormalTok{    .call(d3\_limits.axisLeft(yScale).ticks(6))}
\NormalTok{    .append("text")}
\NormalTok{    .attr("transform", "rotate({-}90)")}
\NormalTok{    .attr("x", {-}innerH / 2).attr("y", {-}35)}
\NormalTok{    .attr("fill", "\#374151").attr("text{-}anchor", "middle")}
\NormalTok{    .style("font{-}size", "13px").text("f(x)");}

\NormalTok{  // Grid}
\NormalTok{  g.append("g")}
\NormalTok{    .selectAll("line.hgrid")}
\NormalTok{    .data(yScale.ticks(6))}
\NormalTok{    .join("line")}
\NormalTok{    .attr("x1", 0).attr("x2", innerW)}
\NormalTok{    .attr("y1", d =\textgreater{} yScale(d)).attr("y2", d =\textgreater{} yScale(d))}
\NormalTok{    .attr("stroke", "\#f3f4f6").attr("stroke{-}width", 1);}

\NormalTok{  // Line generator}
\NormalTok{  const line = d3\_limits.line()}
\NormalTok{    .x(d =\textgreater{} xScale(d[0]))}
\NormalTok{    .y(d =\textgreater{} yScale(d[1]));}

\NormalTok{  // Left segment: x in [{-}1, 2) — f(x) = x+2}
\NormalTok{  const leftPts = d3\_limits.range({-}1, 1.98, 0.01).map(x =\textgreater{} [x, x + 2]);}
\NormalTok{  g.append("path")}
\NormalTok{    .datum(leftPts)}
\NormalTok{    .attr("d", line)}
\NormalTok{    .attr("fill", "none")}
\NormalTok{    .attr("stroke", "\#2563eb")}
\NormalTok{    .attr("stroke{-}width", 2.5);}

\NormalTok{  // Right segment: x in (2, 5] — f(x) = x+2}
\NormalTok{  const rightPts = d3\_limits.range(2.02, 5.01, 0.01).map(x =\textgreater{} [x, x + 2]);}
\NormalTok{  g.append("path")}
\NormalTok{    .datum(rightPts)}
\NormalTok{    .attr("d", line)}
\NormalTok{    .attr("fill", "none")}
\NormalTok{    .attr("stroke", "\#2563eb")}
\NormalTok{    .attr("stroke{-}width", 2.5);}

\NormalTok{  // Open circle at (2, 4)}
\NormalTok{  g.append("circle")}
\NormalTok{    .attr("cx", xScale(2)).attr("cy", yScale(4))}
\NormalTok{    .attr("r", 6)}
\NormalTok{    .attr("fill", "white")}
\NormalTok{    .attr("stroke", "\#2563eb")}
\NormalTok{    .attr("stroke{-}width", 2.5);}

\NormalTok{  // Moving dot — hide when within 0.02 of x=2}
\NormalTok{  const x = xSlider;}
\NormalTok{  const fx = x + 2;}
\NormalTok{  const nearHole = Math.abs(x {-} 2) \textless{} 0.02;}

\NormalTok{  if (!nearHole) \{}
\NormalTok{    g.append("circle")}
\NormalTok{      .attr("cx", xScale(x)).attr("cy", yScale(fx))}
\NormalTok{      .attr("r", 6)}
\NormalTok{      .attr("fill", "\#dc2626");}
\NormalTok{  \}}

\NormalTok{  // Static label}
\NormalTok{  g.append("text")}
\NormalTok{    .attr("x", xScale(3.6)).attr("y", yScale(5.5))}
\NormalTok{    .attr("fill", "\#6b7280")}
\NormalTok{    .style("font{-}size", "12px")}
\NormalTok{    .text("The dot never lands at x = 2");}

\NormalTok{  // Dynamic readout}
\NormalTok{  const readX = x.toFixed(3);}
\NormalTok{  const readFx = nearHole ? "undefined" : fx.toFixed(3);}

\NormalTok{  g.append("text")}
\NormalTok{    .attr("x", xScale({-}0.8)).attr("y", yScale(5.6))}
\NormalTok{    .attr("fill", "\#374151")}
\NormalTok{    .style("font{-}size", "13px").style("font{-}weight", "500")}
\NormalTok{    .text(\textasciigrave{}x = $\{readX\}\textasciigrave{});}

\NormalTok{  g.append("text")}
\NormalTok{    .attr("x", xScale({-}0.8)).attr("y", yScale(5.1))}
\NormalTok{    .attr("fill", nearHole ? "\#9ca3af" : "\#374151")}
\NormalTok{    .style("font{-}size", "13px").style("font{-}weight", "500")}
\NormalTok{    .text(\textasciigrave{}f(x) ≈ $\{readFx\}\textasciigrave{});}

\NormalTok{  return svg.node();}
\NormalTok{\}}
\end{Highlighting}
\end{Shaded}

\section{The method --- evaluating
limits}\label{the-method-evaluating-limits}

Four techniques cover the vast majority of problems you will encounter.

\subsection{Technique 1: Direct
substitution}\label{technique-1-direct-substitution}

If \(f\) is a polynomial, a rational function with a non-zero
denominator at \(a\), a trigonometric function, or any other function
that is \emph{continuous} at \(a\) --- we'll define exactly what this
means shortly, but the idea is that the function has no break at that
point --- then:

\[\lim_{x \to a} f(x) = f(a)\]

Just substitute. This always works when no algebraic obstruction exists.

\textbf{When to use it:} Try this first. If the substitution gives a
real number, you're done.

\textbf{Worked example:} Find \(\lim_{x \to 3} (x^2 + 2x - 1)\).

Substitute \(x = 3\) directly:

\[\lim_{x \to 3} (x^2 + 2x - 1) = 3^2 + 2(3) - 1 = 9 + 6 - 1 = 14\]

The limit is 14. The function is a polynomial --- continuous everywhere
--- so direct substitution is valid at every real number.

\textbf{When it fails:} If substitution gives \(\frac{0}{0}\),
\(\frac{\infty}{\infty}\), or another indeterminate form, you need one
of the other techniques.

\begin{center}\rule{0.5\linewidth}{0.5pt}\end{center}

\subsection{Technique 2: Factor and
cancel}\label{technique-2-factor-and-cancel}

The \(\frac{0}{0}\) form is a signal: both numerator and denominator
share a common factor that vanishes at \(x = a\). Factor both, cancel
the common factor, and then substitute.

\textbf{When to use it:} When direct substitution gives \(\frac{0}{0}\)
and the expression is a ratio of polynomials.

\textbf{Worked example:} Find \(\lim_{x \to 2} \dfrac{x^2 - 4}{x - 2}\).

\textbf{Step 1.} Attempt direct substitution:
\(\dfrac{2^2 - 4}{2 - 2} = \dfrac{0}{0}\). Indeterminate --- proceed to
step 2.

\textbf{Step 2.} Factor the numerator:

\[\frac{x^2 - 4}{x - 2} = \frac{(x-2)(x+2)}{x-2}\]

\textbf{Step 3.} Cancel the common factor \((x - 2)\). This is valid
because we are evaluating a limit as \(x \to 2\), meaning \(x \neq 2\),
so \((x-2) \neq 0\):

\[\frac{(x-2)(x+2)}{x-2} = x + 2 \quad (x \neq 2)\]

\textbf{Step 4.} Now substitute:

\[\lim_{x \to 2} \frac{x^2 - 4}{x - 2} = \lim_{x \to 2}(x + 2) = 2 + 2 = 4\]

The cancellation is legal only because we never set \(x = 2\) --- the
limit asks what happens \emph{as} \(x\) approaches 2, not \emph{at}
\(x = 2\).

\begin{center}\rule{0.5\linewidth}{0.5pt}\end{center}

\subsection{Technique 3: Rationalise}\label{technique-3-rationalise}

When a square root appears and direct substitution gives
\(\frac{0}{0}\), the standard move is to multiply numerator and
denominator by the \textbf{conjugate} of the expression containing the
square root. The conjugate of \((\sqrt{u} - v)\) is \((\sqrt{u} + v)\).
Multiplying removes the radical from one factor using the
difference-of-squares identity.

\textbf{When to use it:} When the \(\frac{0}{0}\) form arises from a
square root.

\textbf{Worked example:} Find
\(\lim_{x \to 0} \dfrac{\sqrt{x+4} - 2}{x}\).

\textbf{Step 1.} Direct substitution:
\(\dfrac{\sqrt{0+4} - 2}{0} = \dfrac{2 - 2}{0} = \dfrac{0}{0}\).
Indeterminate --- proceed.

\textbf{Step 2.} Multiply numerator and denominator by the conjugate
\((\sqrt{x+4} + 2)\):

\[\frac{\sqrt{x+4} - 2}{x} \cdot \frac{\sqrt{x+4} + 2}{\sqrt{x+4} + 2}\]

\textbf{Step 3.} Expand the numerator using \((a-b)(a+b) = a^2 - b^2\):

\[\frac{(\sqrt{x+4})^2 - 2^2}{x(\sqrt{x+4} + 2)} = \frac{(x+4) - 4}{x(\sqrt{x+4} + 2)} = \frac{x}{x(\sqrt{x+4} + 2)}\]

\textbf{Step 4.} Cancel the common factor \(x\) (valid since \(x \to 0\)
means \(x \neq 0\)):

\[\frac{x}{x(\sqrt{x+4} + 2)} = \frac{1}{\sqrt{x+4} + 2}\]

\textbf{Step 5.} Now substitute \(x = 0\):

\[\lim_{x \to 0} \frac{\sqrt{x+4} - 2}{x} = \frac{1}{\sqrt{0+4} + 2} = \frac{1}{2 + 2} = \frac{1}{4}\]

The limit is \(\frac{1}{4}\).

\begin{center}\rule{0.5\linewidth}{0.5pt}\end{center}

\subsection{Technique 4: One-sided limits from
graphs}\label{technique-4-one-sided-limits-from-graphs}

When you have a graph (or can sketch one), you can read limits directly.
Look at the curve's \emph{approach} as \(x\) moves toward the target
from each side. Do not look at the dot at the target point --- look at
where the curve is \emph{heading}.

\textbf{When to use it:} Piecewise functions, graphs given without
formulas, or any function with a potential jump.

\textbf{Worked example:} Let

\[f(x) = \begin{cases} x + 1 & x < 3 \\ 5 & x = 3 \\ 2x - 4 & x > 3 \end{cases}\]

Find \(\lim_{x \to 3^-} f(x)\), \(\lim_{x \to 3^+} f(x)\), and
\(\lim_{x \to 3} f(x)\).

\textbf{Left-hand limit.} For \(x < 3\), the rule is \(f(x) = x + 1\).
As \(x \to 3^-\):

\[\lim_{x \to 3^-} f(x) = 3 + 1 = 4\]

\textbf{Right-hand limit.} For \(x > 3\), the rule is \(f(x) = 2x - 4\).
As \(x \to 3^+\):

\[\lim_{x \to 3^+} f(x) = 2(3) - 4 = 2\]

\textbf{Two-sided limit.} The left and right limits are \(4 \neq 2\),
so:

\[\lim_{x \to 3} f(x) \text{ does not exist}\]

Note that \(f(3) = 5\) is defined --- but the two-sided limit still
fails to exist because the one-sided limits disagree. The value at the
point is irrelevant to whether the limit exists.

\begin{center}\rule{0.5\linewidth}{0.5pt}\end{center}

\section{Continuity}\label{continuity}

With limits in hand, the definition of continuity is short.

\begin{tcolorbox}[enhanced jigsaw, title=\textcolor{quarto-callout-note-color}{\faInfo}\hspace{0.5em}{Note}, opacityback=0, colback=white, breakable, leftrule=.75mm, arc=.35mm, toptitle=1mm, left=2mm, colframe=quarto-callout-note-color-frame, coltitle=black, titlerule=0mm, colbacktitle=quarto-callout-note-color!10!white, opacitybacktitle=0.6, toprule=.15mm, bottomtitle=1mm, bottomrule=.15mm, rightrule=.15mm]

\textbf{Definition.} A function \(f\) is \textbf{continuous at \(a\)} if
all three conditions hold:

\begin{enumerate}
\def\labelenumi{\arabic{enumi}.}
\tightlist
\item
  \(f(a)\) is defined.
\item
  \(\lim_{x \to a} f(x)\) exists.
\item
  \(\lim_{x \to a} f(x) = f(a)\).
\end{enumerate}

If any condition fails, \(f\) is \textbf{discontinuous at \(a\)}.

\end{tcolorbox}

Informally: a function is continuous at a point if you can draw through
that point without lifting your pen. The limit and the value agree.

\subsection{The three types of
discontinuity}\label{the-three-types-of-discontinuity}

\textbf{Removable discontinuity.} Condition 1 or 3 fails, but condition
2 holds. The limit exists, but either the function is not defined there,
or the function is defined but gives a different value. The graph has a
single isolated hole. The discontinuity can be removed by redefining
\(f(a) = L\).

Example: \(g(x) = \dfrac{x^2 - 4}{x - 2}\) is not defined at \(x = 2\),
but \(\lim_{x \to 2} g(x) = 4\). Removable.

\textbf{Jump discontinuity.} The left-hand and right-hand limits both
exist, but they are not equal. The function ``jumps'' from one value to
another. No amount of redefining a single point can fix this --- the gap
is structural.

Example: The piecewise function from Technique 4 above. At \(x = 3\),
the left limit is 4 and the right limit is 2.

\textbf{Infinite discontinuity.} One or both one-sided limits are
\(\pm\infty\). The graph has a vertical asymptote at \(x = a\).

Example: \(h(x) = \dfrac{1}{x}\) at \(x = 0\). As \(x \to 0^+\),
\(h(x) \to +\infty\). As \(x \to 0^-\), \(h(x) \to -\infty\).

\begin{Shaded}
\begin{Highlighting}[]
\NormalTok{//| echo: false}
\NormalTok{d3\_cont = require("d3@7")}
\end{Highlighting}
\end{Shaded}

\begin{Shaded}
\begin{Highlighting}[]
\NormalTok{//| echo: false}
\NormalTok{viewof selectedPanel = Inputs.radio(["Removable", "Jump", "Infinite"], \{}
\NormalTok{  label: "Select discontinuity type",}
\NormalTok{  value: "Removable"}
\NormalTok{\})}
\end{Highlighting}
\end{Shaded}

\begin{Shaded}
\begin{Highlighting}[]
\NormalTok{//| echo: false}
\NormalTok{\{}
\NormalTok{  const panelW = 170, panelH = 180;}
\NormalTok{  const m = \{ top: 28, right: 12, bottom: 28, left: 30 \};}
\NormalTok{  const iW = panelW {-} m.left {-} m.right;}
\NormalTok{  const iH = panelH {-} m.top {-} m.bottom;}
\NormalTok{  const totalW = 3 * panelW + 2 * 16;}

\NormalTok{  const panels = [}
\NormalTok{    \{ id: "Removable", title: "Removable (hole)" \},}
\NormalTok{    \{ id: "Jump",      title: "Jump" \},}
\NormalTok{    \{ id: "Infinite",  title: "Infinite" \}}
\NormalTok{  ];}

\NormalTok{  const conditions = \{}
\NormalTok{    Removable: [}
\NormalTok{      \{ label: "Condition 1 — f(a) defined", pass: false, note: "f(2) is undefined (0/0)" \},}
\NormalTok{      \{ label: "Condition 2 — limit exists",  pass: true,  note: "lim = 4 ✓" \},}
\NormalTok{      \{ label: "Condition 3 — limit = f(a)",  pass: false, note: "f(a) missing, so cannot agree" \}}
\NormalTok{    ],}
\NormalTok{    Jump: [}
\NormalTok{      \{ label: "Condition 1 — f(a) defined", pass: true,  note: "f(3) is defined ✓" \},}
\NormalTok{      \{ label: "Condition 2 — limit exists",  pass: false, note: "left ≠ right, limit DNE" \},}
\NormalTok{      \{ label: "Condition 3 — limit = f(a)",  pass: false, note: "limit doesn\textquotesingle{}t exist, fails" \}}
\NormalTok{    ],}
\NormalTok{    Infinite: [}
\NormalTok{      \{ label: "Condition 1 — f(a) defined", pass: false, note: "f(0) = 1/0, undefined" \},}
\NormalTok{      \{ label: "Condition 2 — limit exists",  pass: false, note: "limits are ±∞" \},}
\NormalTok{      \{ label: "Condition 3 — limit = f(a)",  pass: false, note: "neither condition holds" \}}
\NormalTok{    ]}
\NormalTok{  \};}

\NormalTok{  const container = document.createElement("div");}
\NormalTok{  container.style.cssText = "font{-}family:inherit;";}

\NormalTok{  // Panels row}
\NormalTok{  const svgRow = document.createElement("div");}
\NormalTok{  svgRow.style.cssText = "display:flex; gap:16px; flex{-}wrap:wrap;";}

\NormalTok{  panels.forEach(p =\textgreater{} \{}
\NormalTok{    const isSelected = p.id === selectedPanel;}
\NormalTok{    const wrap = document.createElement("div");}
\NormalTok{    wrap.style.cssText = \textasciigrave{}border:2px solid $\{isSelected ? "\#2563eb" : "\#e5e7eb"\}; border{-}radius:8px; background:$\{isSelected ? "\#eff6ff" : "\#fafafa"\}; cursor:pointer; transition:border{-}color 0.15s;\textasciigrave{};}

\NormalTok{    const svg = d3\_cont.create("svg")}
\NormalTok{      .attr("width", panelW)}
\NormalTok{      .attr("height", panelH);}

\NormalTok{    const g = svg.append("g").attr("transform", \textasciigrave{}translate($\{m.left\},$\{m.top\})\textasciigrave{});}

\NormalTok{    // Title}
\NormalTok{    svg.append("text")}
\NormalTok{      .attr("x", panelW / 2).attr("y", 16)}
\NormalTok{      .attr("text{-}anchor", "middle")}
\NormalTok{      .attr("fill", isSelected ? "\#1d4ed8" : "\#374151")}
\NormalTok{      .style("font{-}size", "11px").style("font{-}weight", "600")}
\NormalTok{      .text(p.title);}

\NormalTok{    if (p.id === "Removable") \{}
\NormalTok{      const xS = d3\_cont.scaleLinear().domain([0, 4]).range([0, iW]);}
\NormalTok{      const yS = d3\_cont.scaleLinear().domain([0, 6]).range([iH, 0]);}
\NormalTok{      g.append("g").attr("transform", \textasciigrave{}translate(0,$\{iH\})\textasciigrave{}).call(d3\_cont.axisBottom(xS).ticks(4).tickSize(3));}
\NormalTok{      g.append("g").call(d3\_cont.axisLeft(yS).ticks(4).tickSize(3));}
\NormalTok{      const pts = d3\_cont.range(0, 1.97, 0.05).concat(d3\_cont.range(2.03, 4.01, 0.05));}
\NormalTok{      const leftPts = d3\_cont.range(0, 1.97, 0.05).map(x =\textgreater{} [x, x + 2]);}
\NormalTok{      const rightPts = d3\_cont.range(2.03, 4.01, 0.05).map(x =\textgreater{} [x, x + 2]);}
\NormalTok{      const line = d3\_cont.line().x(d =\textgreater{} xS(d[0])).y(d =\textgreater{} yS(d[1]));}
\NormalTok{      g.append("path").datum(leftPts).attr("d", line).attr("fill","none").attr("stroke","\#2563eb").attr("stroke{-}width",2);}
\NormalTok{      g.append("path").datum(rightPts).attr("d", line).attr("fill","none").attr("stroke","\#2563eb").attr("stroke{-}width",2);}
\NormalTok{      g.append("circle").attr("cx", xS(2)).attr("cy", yS(4)).attr("r", 5).attr("fill","white").attr("stroke","\#2563eb").attr("stroke{-}width",2);}
\NormalTok{    \}}

\NormalTok{    if (p.id === "Jump") \{}
\NormalTok{      const xS = d3\_cont.scaleLinear().domain([0, 5]).range([0, iW]);}
\NormalTok{      const yS = d3\_cont.scaleLinear().domain([0, 7]).range([iH, 0]);}
\NormalTok{      g.append("g").attr("transform", \textasciigrave{}translate(0,$\{iH\})\textasciigrave{}).call(d3\_cont.axisBottom(xS).ticks(5).tickSize(3));}
\NormalTok{      g.append("g").call(d3\_cont.axisLeft(yS).ticks(5).tickSize(3));}
\NormalTok{      const line = d3\_cont.line().x(d =\textgreater{} xS(d[0])).y(d =\textgreater{} yS(d[1]));}
\NormalTok{      // y=x for x\textless{}3}
\NormalTok{      g.append("path").datum([[0,0],[2.97,2.97]]).attr("d", line).attr("fill","none").attr("stroke","\#2563eb").attr("stroke{-}width",2);}
\NormalTok{      // y=2x{-}4 for x\textgreater{}3}
\NormalTok{      g.append("path").datum([[3.03,2.06],[5,6]]).attr("d", line).attr("fill","none").attr("stroke","\#2563eb").attr("stroke{-}width",2);}
\NormalTok{      // open circles at gap ends}
\NormalTok{      g.append("circle").attr("cx", xS(3)).attr("cy", yS(3)).attr("r", 4).attr("fill","white").attr("stroke","\#2563eb").attr("stroke{-}width",2);}
\NormalTok{      g.append("circle").attr("cx", xS(3)).attr("cy", yS(2)).attr("r", 4).attr("fill","white").attr("stroke","\#2563eb").attr("stroke{-}width",2);}
\NormalTok{    \}}

\NormalTok{    if (p.id === "Infinite") \{}
\NormalTok{      const xS = d3\_cont.scaleLinear().domain([{-}2, 2]).range([0, iW]);}
\NormalTok{      const yS = d3\_cont.scaleLinear().domain([{-}4, 4]).range([iH, 0]);}
\NormalTok{      g.append("g").attr("transform", \textasciigrave{}translate(0,$\{iH/2\})\textasciigrave{}).call(d3\_cont.axisBottom(xS).ticks(4).tickSize(3));}
\NormalTok{      g.append("g").call(d3\_cont.axisLeft(yS).ticks(4).tickSize(3));}
\NormalTok{      // 1/x left branch}
\NormalTok{      const leftB = d3\_cont.range({-}2, {-}0.28, 0.05).map(x =\textgreater{} [x, 1/x]);}
\NormalTok{      const rightB = d3\_cont.range(0.28, 2.01, 0.05).map(x =\textgreater{} [x, 1/x]);}
\NormalTok{      const line = d3\_cont.line().x(d =\textgreater{} xS(d[0])).y(d =\textgreater{} yS(d[1])).defined(d =\textgreater{} Math.abs(d[1]) \textless{}= 4);}
\NormalTok{      g.append("path").datum(leftB).attr("d", line).attr("fill","none").attr("stroke","\#2563eb").attr("stroke{-}width",2);}
\NormalTok{      g.append("path").datum(rightB).attr("d", line).attr("fill","none").attr("stroke","\#2563eb").attr("stroke{-}width",2);}
\NormalTok{      // vertical dashed asymptote}
\NormalTok{      g.append("line")}
\NormalTok{        .attr("x1", xS(0)).attr("x2", xS(0))}
\NormalTok{        .attr("y1", 0).attr("y2", iH)}
\NormalTok{        .attr("stroke", "\#dc2626").attr("stroke{-}width", 1.5)}
\NormalTok{        .attr("stroke{-}dasharray", "4,3");}
\NormalTok{    \}}

\NormalTok{    wrap.appendChild(svg.node());}
\NormalTok{    svgRow.appendChild(wrap);}
\NormalTok{  \});}

\NormalTok{  container.appendChild(svgRow);}

\NormalTok{  // Condition analysis}
\NormalTok{  const info = document.createElement("div");}
\NormalTok{  info.style.cssText = "margin{-}top:1rem; border:1px solid \#e5e7eb; border{-}radius:8px; padding:0.75rem 1rem; background:\#fff;";}

\NormalTok{  const title = document.createElement("div");}
\NormalTok{  title.style.cssText = "font{-}weight:600; font{-}size:0.95em; color:\#374151; margin{-}bottom:0.5rem;";}
\NormalTok{  title.textContent = \textasciigrave{}$\{selectedPanel\} discontinuity — condition analysis\textasciigrave{};}
\NormalTok{  info.appendChild(title);}

\NormalTok{  conditions[selectedPanel].forEach(c =\textgreater{} \{}
\NormalTok{    const row = document.createElement("div");}
\NormalTok{    row.style.cssText = "display:flex; align{-}items:baseline; gap:0.5rem; padding:0.2rem 0; font{-}size:0.9em;";}
\NormalTok{    const icon = document.createElement("span");}
\NormalTok{    icon.textContent = c.pass ? "✓" : "✗";}
\NormalTok{    icon.style.cssText = \textasciigrave{}font{-}weight:700; color:$\{c.pass ? "\#059669" : "\#dc2626"\}; min{-}width:1.2em;\textasciigrave{};}
\NormalTok{    const text = document.createElement("span");}
\NormalTok{    text.style.color = "\#374151";}
\NormalTok{    text.textContent = c.label + " — " + c.note;}
\NormalTok{    row.appendChild(icon);}
\NormalTok{    row.appendChild(text);}
\NormalTok{    info.appendChild(row);}
\NormalTok{  \});}

\NormalTok{  container.appendChild(info);}
\NormalTok{  return container;}
\NormalTok{\}}
\end{Highlighting}
\end{Shaded}

\section{The epsilon-delta
definition}\label{the-epsilon-delta-definition}

The techniques above let you \emph{compute} limits. The epsilon-delta
definition lets you \emph{prove} they are correct. These are different
activities, and most students --- even in university calculus --- use
the techniques far more often than the proof machinery. But
understanding the definition is what separates a practitioner from a
mathematician.

Return to the speedometer. Suppose you want to estimate speed to within
\(\varepsilon = 0.01\) km/h. You need to take a time interval short
enough --- call it \(\delta\) seconds --- that the average speed over
that interval lands within \(0.01\) km/h of the true instantaneous
speed. Different tolerance \(\varepsilon\), different required
\(\delta\). The claim ``the limit is \(L\)'' means: \emph{for every
tolerance you name, I can find a window small enough to guarantee that
accuracy.}

\begin{tcolorbox}[enhanced jigsaw, title=\textcolor{quarto-callout-important-color}{\faExclamation}\hspace{0.5em}{Important}, opacityback=0, colback=white, breakable, leftrule=.75mm, arc=.35mm, toptitle=1mm, left=2mm, colframe=quarto-callout-important-color-frame, coltitle=black, titlerule=0mm, colbacktitle=quarto-callout-important-color!10!white, opacitybacktitle=0.6, toprule=.15mm, bottomtitle=1mm, bottomrule=.15mm, rightrule=.15mm]

\textbf{Epsilon-delta definition of a limit.}

\[\lim_{x \to a} f(x) = L\]

means: for every \(\varepsilon > 0\), there exists \(\delta > 0\) such
that

\[0 < |x - a| < \delta \implies |f(x) - L| < \varepsilon\]

In plain language: however small a target band you set around \(L\) (the
\(\varepsilon\)-band), I can find a window around \(a\) (the
\(\delta\)-window) such that every \(x\) inside the window --- but not
equal to \(a\) itself --- sends \(f(x)\) inside the target band.

\end{tcolorbox}

\textbf{One worked example.} Prove that \(\lim_{x \to 2} 3x = 6\).

We need to show that for any \(\varepsilon > 0\) we can produce a
\(\delta > 0\) satisfying the definition.

We want \(|f(x) - L| < \varepsilon\), i.e., \(|3x - 6| < \varepsilon\).

Factor: \(|3x - 6| = 3|x - 2|\).

So the condition \(|3x - 6| < \varepsilon\) becomes
\(3|x - 2| < \varepsilon\), i.e., \(|x - 2| < \varepsilon/3\).

This tells us to choose \(\delta = \varepsilon/3\).

\textbf{Verification.} Suppose \(0 < |x - 2| < \delta = \varepsilon/3\).
Then:

\[|3x - 6| = 3|x - 2| < 3 \cdot \frac{\varepsilon}{3} = \varepsilon\]

The condition is satisfied for every \(\varepsilon > 0\). The limit is
6. \(\square\)

\begin{tcolorbox}[enhanced jigsaw, title=\textcolor{quarto-callout-note-color}{\faInfo}\hspace{0.5em}{Note}, opacityback=0, colback=white, breakable, leftrule=.75mm, arc=.35mm, toptitle=1mm, left=2mm, colframe=quarto-callout-note-color-frame, coltitle=black, titlerule=0mm, colbacktitle=quarto-callout-note-color!10!white, opacitybacktitle=0.6, toprule=.15mm, bottomtitle=1mm, bottomrule=.15mm, rightrule=.15mm]

\textbf{Why this works --- the logical structure.}

The phrase ``for every \(\varepsilon > 0\)'' covers every possible
tolerance, however tiny. The phrase ``there exists \(\delta > 0\)'' says
you can always respond to any challenge.

Together they convert the informal idea of ``approaches without
reaching'' --- which involves an infinite process --- into a finite,
checkable algebraic statement. You never have to evaluate an infinite
process; you just have to exhibit the \(\delta\) as a function of
\(\varepsilon\) and verify the implication.

This is the logical structure of most rigorous analysis. The techniques
in the previous section are shortcuts valid because the functions
involved are ``nice enough'' that the \(\varepsilon\)-\(\delta\)
machinery can be applied once and the results packaged into rules.

\end{tcolorbox}

\section{Worked examples}\label{worked-examples-17}

\subsection{Example 1 --- Direct substitution
(polynomial)}\label{example-1-direct-substitution-polynomial}

\textbf{Problem.} Find \(\lim_{x \to 3} (x^2 + 2x - 1)\).

\textbf{Context.} In structural engineering, polynomial functions
describe deflection curves along a beam. Evaluating a deflection at a
specific point is direct substitution.

\textbf{Solution.}

The function \(f(x) = x^2 + 2x - 1\) is a polynomial. Polynomials are
continuous everywhere. Direct substitution is valid at \(x = 3\):

\[\lim_{x \to 3}(x^2 + 2x - 1) = 3^2 + 2(3) - 1 = 9 + 6 - 1 = \boxed{14}\]

\textbf{Technique used:} Direct substitution. Applied whenever the
function is continuous at the target point --- no denominator that
vanishes, no radical causing trouble.

\begin{center}\rule{0.5\linewidth}{0.5pt}\end{center}

\subsection{Example 2 --- Factoring (removable
discontinuity)}\label{example-2-factoring-removable-discontinuity}

\textbf{Problem.} Find \(\lim_{x \to 2} \dfrac{x^2 - 4}{x - 2}\).

\textbf{Context.} In software engineering, a function representing
computational cost might be undefined at a specific input size but still
have a meaningful limiting value. Analysing the limit is the correct
tool.

\textbf{Solution.}

\textbf{Step 1.} Direct substitution gives \(\frac{0}{0}\) ---
indeterminate.

\textbf{Step 2.} Factor:

\[\frac{x^2 - 4}{x - 2} = \frac{(x - 2)(x + 2)}{x - 2}\]

\textbf{Step 3.} Cancel \((x-2)\), valid since \(x \neq 2\) in a limit:

\[= x + 2 \quad (x \neq 2)\]

\textbf{Step 4.} Substitute:

\[\lim_{x \to 2} \frac{x^2 - 4}{x - 2} = \lim_{x \to 2}(x + 2) = 2 + 2 = \boxed{4}\]

The function has a hole at \((2, 4)\) --- a removable discontinuity. The
limit exists and equals 4 even though the function is undefined at
\(x = 2\).

\begin{center}\rule{0.5\linewidth}{0.5pt}\end{center}

\subsection{Example 3 --- Rationalisation (resolving an indeterminate
form)}\label{example-3-rationalisation-resolving-an-indeterminate-form}

\textbf{Problem.} Find \(\lim_{x \to 0} \dfrac{\sqrt{x+4} - 2}{x}\).

\textbf{Context.} In chemistry, reaction rate expressions often involve
differences of concentrations under a square root. Finding the limiting
rate at a reference condition uses exactly this technique.

\textbf{Solution.}

\textbf{Step 1.} Direct substitution:
\(\frac{\sqrt{4} - 2}{0} = \frac{0}{0}\). Indeterminate.

\textbf{Step 2.} Multiply by the conjugate
\(\frac{\sqrt{x+4} + 2}{\sqrt{x+4} + 2}\):

\[\frac{\sqrt{x+4} - 2}{x} \cdot \frac{\sqrt{x+4} + 2}{\sqrt{x+4} + 2} = \frac{(x+4) - 4}{x(\sqrt{x+4} + 2)} = \frac{x}{x(\sqrt{x+4} + 2)}\]

\textbf{Step 3.} Cancel \(x\) (\(x \neq 0\) in the limit):

\[= \frac{1}{\sqrt{x+4} + 2}\]

\textbf{Step 4.} Substitute \(x = 0\):

\[\lim_{x \to 0} \frac{\sqrt{x+4} - 2}{x} = \frac{1}{\sqrt{4} + 2} = \frac{1}{4} = \boxed{0.25}\]

\begin{center}\rule{0.5\linewidth}{0.5pt}\end{center}

\subsection{Example 4 --- One-sided limits (regime
change)}\label{example-4-one-sided-limits-regime-change}

\textbf{Problem.} A continuous-time financial model uses the interest
rate function

\[r(t) = \begin{cases} 0.04 & t < 5 \\ 0.04 & t = 5 \\ 0.065 & t > 5 \end{cases}\]

where \(t\) is years from the start of a contract. Does
\(\lim_{t \to 5} r(t)\) exist?

\textbf{Context.} In finance, many contracts include rate resets at
specified times. Whether the limit exists at the reset point determines
whether the model is mathematically continuous or whether a jump
discontinuity must be handled separately.

\textbf{Solution.}

\textbf{Left-hand limit:} As \(t \to 5^-\), the rule gives
\(r(t) = 0.04\):

\[\lim_{t \to 5^-} r(t) = 0.04\]

\textbf{Right-hand limit:} As \(t \to 5^+\), the rule gives
\(r(t) = 0.065\):

\[\lim_{t \to 5^+} r(t) = 0.065\]

\textbf{Two-sided limit:} The one-sided limits are unequal
(\(0.04 \neq 0.065\)), so:

\[\lim_{t \to 5} r(t) \text{ does not exist}\]

This is a jump discontinuity. Note that \(r(5) = 0.04\) is defined ---
the value exists at the point --- but the two-sided limit does not. The
model has a structural break at \(t = 5\) that must be accounted for,
not patched over.

\section{Where this goes}\label{where-this-goes-18}

\textbf{Toward differential calculus (Chapter 2).} The derivative of
\(f\) at a point \(a\) is defined as:

\[f'(a) = \lim_{h \to 0} \frac{f(a+h) - f(a)}{h}\]

This is precisely a limit of the \(\frac{0}{0}\) form --- the same kind
you evaluated in Examples 2 and 3 above. The technique of factoring,
cancelling, and substituting that you practised here is exactly the
algebra you will apply to evaluate derivatives in the next chapter.
Limits are not preliminary material to get past --- they are the
mechanism the derivative is made of.

\textbf{Back to sequences and series (Vol 4 Ch 4).} The convergence you
studied there was the same idea in a different setting. When you showed
that \(\sum_{k=0}^{\infty} \left(\frac{1}{2}\right)^k = 2\), you were
evaluating a limit of partial sums: \(\lim_{n \to \infty} S_n = 2\). The
same idea --- for any tolerance, there's a point beyond which the
partial sums stay within that tolerance --- reappears in the formal
definition of series convergence. The concepts are one family, not two
separate topics.

\begin{tcolorbox}[enhanced jigsaw, title=\textcolor{quarto-callout-note-color}{\faInfo}\hspace{0.5em}{Note}, opacityback=0, colback=white, breakable, leftrule=.75mm, arc=.35mm, toptitle=1mm, left=2mm, colframe=quarto-callout-note-color-frame, coltitle=black, titlerule=0mm, colbacktitle=quarto-callout-note-color!10!white, opacitybacktitle=0.6, toprule=.15mm, bottomtitle=1mm, bottomrule=.15mm, rightrule=.15mm]

\textbf{Where limits appear in practice}

\begin{itemize}
\item
  \textbf{Speedometers and velocity.} Instantaneous velocity is
  \(\lim_{\Delta t \to 0}
  \frac{\Delta s}{\Delta t}\) --- the starting point for all of
  mechanics.
\item
  \textbf{Continuously compounded interest.}
  \(\lim_{n \to \infty}\left(1 + \frac{r}{n}\right)^n = e^r\). Every
  continuous-time financial model uses this limit.
\item
  \textbf{Numerical root-finding.} The bisection method generates a
  sequence of intervals that converge to a root. Convergence is a limit:
  each iteration halves the error, and the sequence of midpoints limits
  to the root.
\item
  \textbf{Signal processing.} A low-pass filter's attenuation as
  frequency \(\to \infty\) is a limit question. Whether the filter
  ``works'' depends on whether that limit is zero.
\end{itemize}

\end{tcolorbox}

\section{Exercises}\label{exercises-18}

\begin{Shaded}
\begin{Highlighting}[]
\NormalTok{//| echo: false}
\NormalTok{katex = require("katex@0.16")}
\end{Highlighting}
\end{Shaded}

\begin{Shaded}
\begin{Highlighting}[]
\NormalTok{//| echo: false}
\NormalTok{function makeStepperHTML(exerciseNum, steps) \{}
\NormalTok{  let current = 0;}
\NormalTok{  const totalSteps = steps.length;}
\NormalTok{  const container = document.createElement("div");}
\NormalTok{  container.style.cssText = "border:1px solid \#e5e7eb; border{-}radius:8px; padding:1rem 1.25rem; margin:0.75rem 0; font{-}family:inherit;";}
\NormalTok{  const stepsDiv = document.createElement("div");}
\NormalTok{  stepsDiv.style.marginBottom = "0.75rem";}

\NormalTok{  function renderSteps() \{}
\NormalTok{    stepsDiv.innerHTML = "";}
\NormalTok{    for (let i = 0; i \textless{} current; i++) \{}
\NormalTok{      const s = steps[i];}
\NormalTok{      const row = document.createElement("div");}
\NormalTok{      row.style.cssText = "display:grid; grid{-}template{-}columns:200px 1fr; gap:0.35rem 1rem; align{-}items:baseline; padding:0.35rem 0; border{-}top:1px solid \#f3f4f6;";}
\NormalTok{      const opCell = document.createElement("span");}
\NormalTok{      opCell.style.cssText = "font{-}size:0.85em; color:\#6b7280; font{-}style:italic;";}
\NormalTok{      opCell.textContent = s.op;}
\NormalTok{      const eqCell = document.createElement("span");}
\NormalTok{      eqCell.innerHTML = katex.renderToString(s.eq, \{ throwOnError: false, displayMode: false \});}
\NormalTok{      row.appendChild(opCell);}
\NormalTok{      row.appendChild(eqCell);}
\NormalTok{      if (s.note) \{}
\NormalTok{        const noteRow = document.createElement("div");}
\NormalTok{        noteRow.style.cssText = "grid{-}column:2; font{-}size:0.82em; color:\#6b7280; padding{-}bottom:0.2rem;";}
\NormalTok{        noteRow.textContent = s.note;}
\NormalTok{        row.appendChild(noteRow);}
\NormalTok{      \}}
\NormalTok{      stepsDiv.appendChild(row);}
\NormalTok{    \}}
\NormalTok{    if (current === totalSteps) \{}
\NormalTok{      const done = document.createElement("div");}
\NormalTok{      done.style.cssText = "margin{-}top:0.5rem; font{-}size:0.9em; color:\#059669; font{-}weight:500;";}
\NormalTok{      done.textContent = "✓ Solution complete";}
\NormalTok{      stepsDiv.appendChild(done);}
\NormalTok{    \}}
\NormalTok{  \}}

\NormalTok{  const controls = document.createElement("div");}
\NormalTok{  controls.style.cssText = "display:flex; gap:0.5rem; align{-}items:center;";}
\NormalTok{  const nextBtn = document.createElement("button");}
\NormalTok{  nextBtn.textContent = "Next step →";}
\NormalTok{  nextBtn.style.cssText = "padding:0.35rem 0.85rem; border:1px solid \#d1d5db; border{-}radius:4px; background:\#fff; cursor:pointer; font{-}size:0.9em;";}
\NormalTok{  const resetBtn = document.createElement("button");}
\NormalTok{  resetBtn.textContent = "Reset";}
\NormalTok{  resetBtn.style.cssText = "padding:0.35rem 0.75rem; border:1px solid \#e5e7eb; border{-}radius:4px; background:\#f9fafb; cursor:pointer; font{-}size:0.9em; color:\#6b7280;";}
\NormalTok{  const counter = document.createElement("span");}
\NormalTok{  counter.style.cssText = "font{-}size:0.82em; color:\#9ca3af; margin{-}left:0.25rem;";}

\NormalTok{  function updateButtons() \{}
\NormalTok{    nextBtn.disabled = current === totalSteps;}
\NormalTok{    nextBtn.style.opacity = current === totalSteps ? "0.4" : "1";}
\NormalTok{    counter.textContent = current === 0 ? "Click to reveal steps" : \textasciigrave{}Step $\{current\} of $\{totalSteps\}\textasciigrave{};}
\NormalTok{  \}}

\NormalTok{  nextBtn.onclick = () =\textgreater{} \{ if (current \textless{} totalSteps) \{ current++; renderSteps(); updateButtons(); \} \};}
\NormalTok{  resetBtn.onclick = () =\textgreater{} \{ current = 0; renderSteps(); updateButtons(); \};}
\NormalTok{  controls.appendChild(nextBtn);}
\NormalTok{  controls.appendChild(resetBtn);}
\NormalTok{  controls.appendChild(counter);}
\NormalTok{  container.appendChild(stepsDiv);}
\NormalTok{  container.appendChild(controls);}
\NormalTok{  renderSteps();}
\NormalTok{  updateButtons();}
\NormalTok{  return container;}
\NormalTok{\}}
\end{Highlighting}
\end{Shaded}

\textbf{Exercise 1.} A structural engineer models the deflection of a
beam using \(d(x) = 2x^3 - 3x^2 + x + 1\), where \(x\) is position in
metres.

Find \(\lim_{x \to 2} d(x)\).

\begin{Shaded}
\begin{Highlighting}[]
\NormalTok{//| echo: false}
\NormalTok{makeStepperHTML(1, [}
\NormalTok{  \{ op: "Identify the function type",  eq: "\textbackslash{}\textbackslash{}text\{d(x) = 2x\^{}3 {-} 3x\^{}2 + x + 1 is a polynomial\}" \},}
\NormalTok{  \{ op: "Apply direct substitution",   eq: "\textbackslash{}\textbackslash{}text\{Polynomials are continuous everywhere, so \} \textbackslash{}\textbackslash{}lim\_\{x \textbackslash{}\textbackslash{}to 2\} d(x) = d(2)" \},}
\NormalTok{  \{ op: "Evaluate",                    eq: "d(2) = 2(2)\^{}3 {-} 3(2)\^{}2 + 2 + 1 = 16 {-} 12 + 2 + 1 = 7" \},}
\NormalTok{  \{ op: "State the answer",            eq: "\textbackslash{}\textbackslash{}lim\_\{x \textbackslash{}\textbackslash{}to 2\} d(x) = 7 \textbackslash{}\textbackslash{}text\{ metres\}" \}}
\NormalTok{])}
\end{Highlighting}
\end{Shaded}

\begin{center}\rule{0.5\linewidth}{0.5pt}\end{center}

\textbf{Exercise 2.} A computing platform charges users based on data
processed. The cost function has the form

\[C(x) = \frac{x^2 - 9}{x - 3}\]

where \(x\) is gigabytes processed. The pricing model is undefined at
\(x = 3\) GB due to a contract edge case. Find \(\lim_{x \to 3} C(x)\).

\begin{Shaded}
\begin{Highlighting}[]
\NormalTok{//| echo: false}
\NormalTok{makeStepperHTML(2, [}
\NormalTok{  \{ op: "Attempt direct substitution", eq: "C(3) = \textbackslash{}\textbackslash{}frac\{9 {-} 9\}\{3 {-} 3\} = \textbackslash{}\textbackslash{}frac\{0\}\{0\}", note: "Indeterminate — factor instead." \},}
\NormalTok{  \{ op: "Factor the numerator",        eq: "x\^{}2 {-} 9 = (x{-}3)(x+3), \textbackslash{}\textbackslash{}text\{ so \} C(x) = \textbackslash{}\textbackslash{}frac\{(x{-}3)(x+3)\}\{x{-}3\}" \},}
\NormalTok{  \{ op: "Cancel the common factor",    eq: "C(x) = x + 3 \textbackslash{}\textbackslash{}quad (x \textbackslash{}\textbackslash{}neq 3)" \},}
\NormalTok{  \{ op: "Substitute",                  eq: "\textbackslash{}\textbackslash{}lim\_\{x \textbackslash{}\textbackslash{}to 3\} C(x) = 3 + 3 = 6" \},}
\NormalTok{  \{ op: "Interpret",                   eq: "\textbackslash{}\textbackslash{}text\{Cost approaches \textbackslash{}\textbackslash{}$6 per GB as usage approaches 3 GB threshold\}" \}}
\NormalTok{])}
\end{Highlighting}
\end{Shaded}

\begin{center}\rule{0.5\linewidth}{0.5pt}\end{center}

\textbf{Exercise 3.} In biochemistry, a reaction rate expression takes
the form

\[R(c) = \frac{\sqrt{c + 9} - 3}{c}\]

where \(c\) is the change in concentration from a reference value. Find
\(\lim_{c \to 0} R(c)\).

\begin{Shaded}
\begin{Highlighting}[]
\NormalTok{//| echo: false}
\NormalTok{makeStepperHTML(3, [}
\NormalTok{  \{ op: "Attempt direct substitution", eq: "R(0) = \textbackslash{}\textbackslash{}frac\{\textbackslash{}\textbackslash{}sqrt\{9\} {-} 3\}\{0\} = \textbackslash{}\textbackslash{}frac\{0\}\{0\}", note: "Indeterminate — rationalise." \},}
\NormalTok{  \{ op: "Multiply by the conjugate",   eq: "\textbackslash{}\textbackslash{}frac\{\textbackslash{}\textbackslash{}sqrt\{c+9\} {-} 3\}\{c\} \textbackslash{}\textbackslash{}cdot \textbackslash{}\textbackslash{}frac\{\textbackslash{}\textbackslash{}sqrt\{c+9\} + 3\}\{\textbackslash{}\textbackslash{}sqrt\{c+9\} + 3\} = \textbackslash{}\textbackslash{}frac\{(c+9) {-} 9\}\{c(\textbackslash{}\textbackslash{}sqrt\{c+9\} + 3)\} = \textbackslash{}\textbackslash{}frac\{c\}\{c(\textbackslash{}\textbackslash{}sqrt\{c+9\}+3)\}" \},}
\NormalTok{  \{ op: "Cancel c",                    eq: "\textbackslash{}\textbackslash{}frac\{1\}\{\textbackslash{}\textbackslash{}sqrt\{c+9\}+3\}", note: "Valid since c ≠ 0 in the limit." \},}
\NormalTok{  \{ op: "Substitute c = 0",            eq: "\textbackslash{}\textbackslash{}frac\{1\}\{\textbackslash{}\textbackslash{}sqrt\{9\}+3\} = \textbackslash{}\textbackslash{}frac\{1\}\{3+3\} = \textbackslash{}\textbackslash{}frac\{1\}\{6\}" \},}
\NormalTok{  \{ op: "Conclude",                    eq: "\textbackslash{}\textbackslash{}lim\_\{c \textbackslash{}\textbackslash{}to 0\} R(c) = \textbackslash{}\textbackslash{}dfrac\{1\}\{6\}" \}}
\NormalTok{])}
\end{Highlighting}
\end{Shaded}

\begin{center}\rule{0.5\linewidth}{0.5pt}\end{center}

\textbf{Exercise 4.} A piecewise model for a climate variable switches
regime at time \(t = 10\) years:

\[T(t) = \begin{cases} 2t + 1 & t < 10 \\ 30 & t = 10 \\ 3t - 8 & t > 10 \end{cases}\]

\begin{enumerate}
\def\labelenumi{(\alph{enumi})}
\tightlist
\item
  Find \(\lim_{t \to 10^-} T(t)\).
\item
  Find \(\lim_{t \to 10^+} T(t)\).
\item
  Does \(\lim_{t \to 10} T(t)\) exist? Explain.
\item
  Is \(T\) continuous at \(t = 10\)?
\end{enumerate}

\begin{Shaded}
\begin{Highlighting}[]
\NormalTok{//| echo: false}
\NormalTok{makeStepperHTML(4, [}
\NormalTok{  \{ op: "Left{-}hand limit",   eq: "\textbackslash{}\textbackslash{}text\{For \} t \textless{} 10,\textbackslash{}\textbackslash{}; T(t) = 2t+1.\textbackslash{}\textbackslash{}; \textbackslash{}\textbackslash{}lim\_\{t \textbackslash{}\textbackslash{}to 10\^{}{-}\} T(t) = 2(10)+1 = 21" \},}
\NormalTok{  \{ op: "Right{-}hand limit",  eq: "\textbackslash{}\textbackslash{}text\{For \} t \textgreater{} 10,\textbackslash{}\textbackslash{}; T(t) = 3t{-}8.\textbackslash{}\textbackslash{}; \textbackslash{}\textbackslash{}lim\_\{t \textbackslash{}\textbackslash{}to 10\^{}+\} T(t) = 3(10){-}8 = 22" \},}
\NormalTok{  \{ op: "Two{-}sided limit",   eq: "21 \textbackslash{}\textbackslash{}neq 22 \textbackslash{}\textbackslash{}implies \textbackslash{}\textbackslash{}lim\_\{t \textbackslash{}\textbackslash{}to 10\} T(t) \textbackslash{}\textbackslash{}text\{ does not exist\}" \},}
\NormalTok{  \{ op: "Check continuity",  eq: "T(10) = 30 \textbackslash{}\textbackslash{}text\{ (defined ✓), but two{-}sided limit DNE ✗\}", note: "Condition 2 fails — T is not continuous at t = 10." \},}
\NormalTok{  \{ op: "Classify",          eq: "\textbackslash{}\textbackslash{}text\{Jump discontinuity: both one{-}sided limits exist but are unequal\}" \}}
\NormalTok{])}
\end{Highlighting}
\end{Shaded}

\begin{center}\rule{0.5\linewidth}{0.5pt}\end{center}

\textbf{Exercise 5.} Use the three conditions to determine whether

\[f(x) = \begin{cases} x^2 - 1 & x \neq 2 \\ 3 & x = 2 \end{cases}\]

is continuous at \(x = 2\).

\begin{Shaded}
\begin{Highlighting}[]
\NormalTok{//| echo: false}
\NormalTok{makeStepperHTML(5, [}
\NormalTok{  \{ op: "Condition 1: f(2) defined?",           eq: "f(2) = 3 \textbackslash{}\textbackslash{}text\{ — defined ✓\}" \},}
\NormalTok{  \{ op: "Condition 2: limit exists?",            eq: "\textbackslash{}\textbackslash{}lim\_\{x \textbackslash{}\textbackslash{}to 2\}(x\^{}2 {-} 1) = 4 {-} 1 = 3 \textbackslash{}\textbackslash{}text\{ — exists ✓\}", note: "The limit uses the polynomial branch; the override at x=2 does not affect the approach." \},}
\NormalTok{  \{ op: "Condition 3: limit = f(2)?",            eq: "\textbackslash{}\textbackslash{}lim = 3 = f(2) = 3 \textbackslash{}\textbackslash{}text\{ — agree ✓\}" \},}
\NormalTok{  \{ op: "Conclude",                              eq: "\textbackslash{}\textbackslash{}text\{All three conditions hold: \} f \textbackslash{}\textbackslash{}text\{ is continuous at \} x = 2" \},}
\NormalTok{  \{ op: "Note",                                  eq: "\textbackslash{}\textbackslash{}text\{The piecewise value at \} x=2 \textbackslash{}\textbackslash{}text\{ happens to match the limit perfectly\}", note: "No discontinuity despite the piecewise definition." \}}
\NormalTok{])}
\end{Highlighting}
\end{Shaded}

\begin{center}\rule{0.5\linewidth}{0.5pt}\end{center}

\textbf{Exercise 6 (epsilon-delta).} Use the \(\varepsilon\)-\(\delta\)
definition to prove that \(\lim_{x \to 1} (5x - 2) = 3\).

\begin{Shaded}
\begin{Highlighting}[]
\NormalTok{//| echo: false}
\NormalTok{makeStepperHTML(6, [}
\NormalTok{  \{ op: "Write down what to show",     eq: "\textbackslash{}\textbackslash{}forall\textbackslash{}\textbackslash{}, \textbackslash{}\textbackslash{}varepsilon \textgreater{} 0,\textbackslash{}\textbackslash{}; \textbackslash{}\textbackslash{}exists\textbackslash{}\textbackslash{}, \textbackslash{}\textbackslash{}delta \textgreater{} 0 :\textbackslash{}\textbackslash{}; 0 \textless{} |x{-}1| \textless{} \textbackslash{}\textbackslash{}delta \textbackslash{}\textbackslash{}implies |(5x{-}2){-}3| \textless{} \textbackslash{}\textbackslash{}varepsilon" \},}
\NormalTok{  \{ op: "Simplify the target",         eq: "|(5x{-}2){-}3| = |5x{-}5| = 5|x{-}1|" \},}
\NormalTok{  \{ op: "Work backwards to find δ",    eq: "5|x{-}1| \textless{} \textbackslash{}\textbackslash{}varepsilon \textbackslash{}\textbackslash{}implies |x{-}1| \textless{} \textbackslash{}\textbackslash{}varepsilon/5 \textbackslash{}\textbackslash{}implies \textbackslash{}\textbackslash{}text\{choose \} \textbackslash{}\textbackslash{}delta = \textbackslash{}\textbackslash{}varepsilon/5" \},}
\NormalTok{  \{ op: "Write the proof",             eq: "0 \textless{} |x{-}1| \textless{} \textbackslash{}\textbackslash{}delta = \textbackslash{}\textbackslash{}varepsilon/5 \textbackslash{}\textbackslash{}implies |(5x{-}2){-}3| = 5|x{-}1| \textless{} 5 \textbackslash{}\textbackslash{}cdot \textbackslash{}\textbackslash{}tfrac\{\textbackslash{}\textbackslash{}varepsilon\}\{5\} = \textbackslash{}\textbackslash{}varepsilon \textbackslash{}\textbackslash{}quad \textbackslash{}\textbackslash{}square" \},}
\NormalTok{  \{ op: "Reflect",                     eq: "\textbackslash{}\textbackslash{}text\{Key move: simplify \} |f(x){-}L| \textbackslash{}\textbackslash{}text\{ until it contains \} |x{-}a|, \textbackslash{}\textbackslash{}text\{ then read off \} \textbackslash{}\textbackslash{}delta", note: "The same pattern applies to every linear limit proof." \}}
\NormalTok{])}
\end{Highlighting}
\end{Shaded}

\chapter{Differential calculus}\label{differential-calculus}

The instantaneous rate of change

\hfill\break

Your car's GPS records your average speed over the entire trip: 62 km/h.
But your speedometer, right now, reads 90 km/h. These are different
quantities. The GPS average looks back over the whole journey; the
speedometer is telling you something about \emph{this instant}. The
question ``how fast am I going right now?'' is not the same question as
``how fast did I go on average?'' --- and answering it requires a
fundamentally different tool.

Here is a sharper version of the same problem. A drug is injected into
the bloodstream. A doctor monitors its concentration --- measured in
milligrams per litre --- every hour. She can compute the average rate of
clearance over a two-hour window. But what determines the dosing
schedule is the \emph{instantaneous} rate of clearance right now: how
fast is the concentration dropping at this moment? If that rate is too
slow, the drug accumulates to toxic levels. Too fast, and the dose is
ineffective before the next one can be given. The decision is made on
the basis of a rate at a point in time, not an average across an
interval.

Both situations share the same mathematical structure. In chapter 1 you
met limits: the value that a function approaches as the input gets
arbitrarily close to some point. The speedometer reading \emph{is} a
limit --- the limit of average speeds over shorter and shorter time
intervals. This chapter builds that limit into a systematic tool.

\section{The limit definition of the
derivative}\label{the-limit-definition-of-the-derivative}

The derivative has a precise definition as a limit. Start with the
\textbf{average rate of change} of \(f\) over the interval from \(x\) to
\(x + h\):

\[\frac{f(x+h) - f(x)}{h}\]

The numerator \(f(x+h) - f(x)\) is the change in output. The denominator
\(h\) is the change in input. Their ratio is the slope of the
\textbf{secant line} connecting the two points \((x,\, f(x))\) and
\((x+h,\, f(x+h))\) on the graph.

Now let \(h \to 0\). The second point slides toward the first. The
secant line tilts and, in the limit, becomes the \textbf{tangent line}
at \(x\). The slope of that tangent line is the derivative:

\[\boxed{f'(x) = \lim_{h \to 0} \frac{f(x+h) - f(x)}{h}}\]

This is the formal definition. Every differentiation rule derived later
is a consequence of it. The definition itself is the limit you saw in
chapter 1, applied to a specific quotient.

\section{What the notation is
saying}\label{what-the-notation-is-saying-19}

Two notations for the derivative are standard. You will need both.

\textbf{Prime notation} (Newton): \(f'(x)\), read \emph{f prime of x}.
Clean and fast to write. Used when there is no ambiguity about which
variable you're differentiating with respect to.

\textbf{Leibniz notation}: \(\dfrac{dy}{dx}\), read \emph{dy by dx}.
More verbose but carries structural information --- it names both the
output variable (\(y\)) and the input variable (\(x\)). When a problem
involves several variables, or when you need to keep track of units,
Leibniz notation makes the algebra clearer. Most science and engineering
texts default to it.

Both mean the same thing. If \(y = f(x)\), then:

\[f'(x) = \frac{dy}{dx} = \frac{d}{dx}\bigl[f(x)\bigr]\]

The expression \(\dfrac{d}{dx}[\,\cdot\,]\) is an \textbf{operator} ---
it takes a function as input and produces its derivative as output.
Think of \(d/dx\) as an instruction: \emph{differentiate with respect to
\(x\)}.

\begin{Shaded}
\begin{Highlighting}[]
\NormalTok{//| echo: false}
\NormalTok{viewof x2\_secant = Inputs.range([1.01, 3.0], \{}
\NormalTok{  step: 0.01,}
\NormalTok{  value: 3.0,}
\NormalTok{  label: "x₂ (drag toward 1)"}
\NormalTok{\})}
\end{Highlighting}
\end{Shaded}

\begin{Shaded}
\begin{Highlighting}[]
\NormalTok{//| echo: false}
\NormalTok{\{}
\NormalTok{  const d3 = await require("d3@7");}

\NormalTok{  const width = 560, height = 320;}
\NormalTok{  const margin = \{ top: 24, right: 32, bottom: 44, left: 52 \};}
\NormalTok{  const innerW = width {-} margin.left {-} margin.right;}
\NormalTok{  const innerH = height {-} margin.top {-} margin.bottom;}

\NormalTok{  const f = x =\textgreater{} x * x + 1;}
\NormalTok{  const x1 = 1, y1 = f(x1);}
\NormalTok{  const x2 = x2\_secant, y2 = f(x2);}
\NormalTok{  const slope = (y2 {-} y1) / (x2 {-} x1);}
\NormalTok{  const nearTangent = x2 \textless{} 1.05;}

\NormalTok{  const xScale = d3.scaleLinear().domain([{-}0.5, 3]).range([0, innerW]);}
\NormalTok{  const yScale = d3.scaleLinear().domain([0, 11]).range([innerH, 0]);}

\NormalTok{  const svg = d3.create("svg")}
\NormalTok{    .attr("width", width)}
\NormalTok{    .attr("height", height)}
\NormalTok{    .attr("viewBox", \textasciigrave{}0 0 $\{width\} $\{height\}\textasciigrave{})}
\NormalTok{    .style("font{-}family", "inherit");}

\NormalTok{  const g = svg.append("g")}
\NormalTok{    .attr("transform", \textasciigrave{}translate($\{margin.left\},$\{margin.top\})\textasciigrave{});}

\NormalTok{  // Axes}
\NormalTok{  g.append("g")}
\NormalTok{    .attr("transform", \textasciigrave{}translate(0,$\{innerH\})\textasciigrave{})}
\NormalTok{    .call(d3.axisBottom(xScale).ticks(7))}
\NormalTok{    .call(ax =\textgreater{} ax.select(".domain").attr("stroke", "\#9ca3af"))}
\NormalTok{    .call(ax =\textgreater{} ax.selectAll("line").attr("stroke", "\#9ca3af"))}
\NormalTok{    .call(ax =\textgreater{} ax.selectAll("text").attr("fill", "\#6b7280").attr("font{-}size", "11"));}

\NormalTok{  g.append("g")}
\NormalTok{    .call(d3.axisLeft(yScale).ticks(6))}
\NormalTok{    .call(ax =\textgreater{} ax.select(".domain").attr("stroke", "\#9ca3af"))}
\NormalTok{    .call(ax =\textgreater{} ax.selectAll("line").attr("stroke", "\#9ca3af"))}
\NormalTok{    .call(ax =\textgreater{} ax.selectAll("text").attr("fill", "\#6b7280").attr("font{-}size", "11"));}

\NormalTok{  // Axis labels}
\NormalTok{  g.append("text").attr("x", innerW / 2).attr("y", innerH + 38)}
\NormalTok{    .attr("text{-}anchor", "middle").attr("fill", "\#6b7280").attr("font{-}size", "12").text("x");}
\NormalTok{  g.append("text").attr("transform", "rotate({-}90)").attr("x", {-}innerH / 2).attr("y", {-}38)}
\NormalTok{    .attr("text{-}anchor", "middle").attr("fill", "\#6b7280").attr("font{-}size", "12").text("y");}

\NormalTok{  // Grid lines}
\NormalTok{  g.append("g").attr("class", "grid")}
\NormalTok{    .call(d3.axisLeft(yScale).ticks(6).tickSize({-}innerW).tickFormat(""))}
\NormalTok{    .call(ax =\textgreater{} ax.select(".domain").remove())}
\NormalTok{    .call(ax =\textgreater{} ax.selectAll("line").attr("stroke", "\#f3f4f6").attr("stroke{-}dasharray", "3,3"));}

\NormalTok{  // Curve f(x) = x² + 1}
\NormalTok{  const lineGen = d3.line()}
\NormalTok{    .x(d =\textgreater{} xScale(d))}
\NormalTok{    .y(d =\textgreater{} yScale(f(d)))}
\NormalTok{    .curve(d3.curveBasis);}

\NormalTok{  const xs = d3.range({-}0.5, 3.02, 0.02);}
\NormalTok{  g.append("path")}
\NormalTok{    .datum(xs)}
\NormalTok{    .attr("fill", "none")}
\NormalTok{    .attr("stroke", "\#2563eb")}
\NormalTok{    .attr("stroke{-}width", 2.5)}
\NormalTok{    .attr("d", lineGen);}

\NormalTok{  // Curve label}
\NormalTok{  g.append("text")}
\NormalTok{    .attr("x", xScale(2.7))}
\NormalTok{    .attr("y", yScale(f(2.7)) {-} 10)}
\NormalTok{    .attr("fill", "\#2563eb")}
\NormalTok{    .attr("font{-}size", "12")}
\NormalTok{    .text("f(x) = x² + 1");}

\NormalTok{  // Draw line (secant or tangent) extended across the visible range}
\NormalTok{  const drawLine = (slope, passX, passY, color, dash) =\textgreater{} \{}
\NormalTok{    const xLeft = {-}0.5, xRight = 3;}
\NormalTok{    const yLeft = passY + slope * (xLeft {-} passX);}
\NormalTok{    const yRight = passY + slope * (xRight {-} passX);}
\NormalTok{    g.append("line")}
\NormalTok{      .attr("x1", xScale(xLeft)).attr("y1", yScale(yLeft))}
\NormalTok{      .attr("x2", xScale(xRight)).attr("y2", yScale(yRight))}
\NormalTok{      .attr("stroke", color)}
\NormalTok{      .attr("stroke{-}width", nearTangent ? 2 : 1.5)}
\NormalTok{      .attr("stroke{-}dasharray", dash || "none")}
\NormalTok{      .attr("opacity", 0.85);}
\NormalTok{  \};}

\NormalTok{  if (nearTangent) \{}
\NormalTok{    // Show tangent line}
\NormalTok{    drawLine(2, x1, y1, "\#10b981", null);}
\NormalTok{    g.append("text")}
\NormalTok{      .attr("x", xScale(2.2))}
\NormalTok{      .attr("y", yScale(2 * (2.2 {-} 1) + y1) {-} 10)}
\NormalTok{      .attr("fill", "\#10b981")}
\NormalTok{      .attr("font{-}size", "12")}
\NormalTok{      .attr("font{-}weight", "600")}
\NormalTok{      .text("tangent slope → 2");}
\NormalTok{  \} else \{}
\NormalTok{    // Show secant line}
\NormalTok{    drawLine(slope, x1, y1, "\#ef4444", null);}
\NormalTok{  \}}

\NormalTok{  // P₁ dot (fixed)}
\NormalTok{  g.append("circle")}
\NormalTok{    .attr("cx", xScale(x1)).attr("cy", yScale(y1))}
\NormalTok{    .attr("r", 5).attr("fill", "\#2563eb").attr("stroke", "\#fff").attr("stroke{-}width", 1.5);}
\NormalTok{  g.append("text")}
\NormalTok{    .attr("x", xScale(x1) {-} 14).attr("y", yScale(y1) {-} 9)}
\NormalTok{    .attr("fill", "\#2563eb").attr("font{-}size", "12").attr("font{-}weight", "600")}
\NormalTok{    .text("P₁");}

\NormalTok{  // P₂ dot (moving)}
\NormalTok{  if (!nearTangent) \{}
\NormalTok{    g.append("circle")}
\NormalTok{      .attr("cx", xScale(x2)).attr("cy", yScale(y2))}
\NormalTok{      .attr("r", 5).attr("fill", "\#ef4444").attr("stroke", "\#fff").attr("stroke{-}width", 1.5);}
\NormalTok{    g.append("text")}
\NormalTok{      .attr("x", xScale(x2) + 7).attr("y", yScale(y2) {-} 6)}
\NormalTok{      .attr("fill", "\#ef4444").attr("font{-}size", "12").attr("font{-}weight", "600")}
\NormalTok{      .text("P₂");}
\NormalTok{  \}}

\NormalTok{  // Slope readout}
\NormalTok{  const slopeText = nearTangent}
\NormalTok{    ? "secant slope → 2.000  (tangent)"}
\NormalTok{    : \textasciigrave{}secant slope = $\{slope.toFixed(4)\}\textasciigrave{};}
\NormalTok{  const slopeColor = nearTangent ? "\#10b981" : "\#ef4444";}

\NormalTok{  g.append("text")}
\NormalTok{    .attr("x", innerW {-} 4).attr("y", 14)}
\NormalTok{    .attr("text{-}anchor", "end")}
\NormalTok{    .attr("fill", slopeColor)}
\NormalTok{    .attr("font{-}size", "13")}
\NormalTok{    .attr("font{-}weight", "600")}
\NormalTok{    .text(slopeText);}

\NormalTok{  return svg.node();}
\NormalTok{\}}
\end{Highlighting}
\end{Shaded}

\begin{tcolorbox}[enhanced jigsaw, opacityback=0, colback=white, breakable, colframe=quarto-callout-note-color-frame, arc=.35mm, left=2mm, toprule=.15mm, leftrule=.75mm, bottomrule=.15mm, rightrule=.15mm]
\begin{minipage}[t]{5.5mm}
\textcolor{quarto-callout-note-color}{\faInfo}
\end{minipage}%
\begin{minipage}[t]{\textwidth - 5.5mm}

\textbf{The tangent line}

The derivative \(f'(a)\) is the slope of the \textbf{tangent line} to
the graph of \(f\) at the point \((a, f(a))\) --- the straight line
through \((a, f(a))\) with slope \(f'(a)\), which is the best linear
approximation to \(f\) near \(x = a\).

Unlike the informal picture, a tangent line can cross the curve ---
\(f(x) = x^3\) at the origin has a tangent line that crosses. What
defines a tangent is its slope, not whether it crosses.

This is why a function fails to be differentiable at a corner or cusp:
no unique tangent line exists there.

\end{minipage}%
\end{tcolorbox}

\section{The method --- differentiation
rules}\label{the-method-differentiation-rules}

Computing every derivative from the limit definition would be
exhausting. Fortunately, the limit has been worked out once for each
class of function, packaged into rules, and those rules can be applied
mechanically. The skill in differentiation is knowing which rule to
reach for.

\subsection{Power rule}\label{power-rule}

For any function \(f(x) = x^n\), where \(n\) is a real number:

\[\frac{d}{dx}\bigl[x^n\bigr] = n x^{n-1}\]

This is the most-used rule in calculus. It applies to positive integers,
negative integers, fractions --- any real exponent.

\textbf{Derivation for positive integer \(n\).} Start from the limit
definition:

\[f'(x) = \lim_{h \to 0} \frac{(x+h)^n - x^n}{h}\]

To make progress, we need to expand \((x+h)^n\). For small \(n\) you can
do this by multiplying out, but the pattern for general \(n\) comes from
the binomial theorem --- covered in Vol 3 polynomials. For \(n = 2\):
\((x+h)^2 = x^2 + 2xh + h^2\). For \(n = 3\):
\((x+h)^3 = x^3 + 3x^2h + 3xh^2 + h^3\). In general:

\[(x+h)^n = x^n + n x^{n-1} h + \frac{n(n-1)}{2} x^{n-2} h^2 + \cdots + h^n\]

Subtract \(x^n\) and divide by \(h\):

\[\frac{(x+h)^n - x^n}{h} = n x^{n-1} + \frac{n(n-1)}{2} x^{n-2} h + \cdots + h^{n-1}\]

Every term except the first contains at least one factor of \(h\). As
\(h \to 0\), all those terms vanish:

\[f'(x) = \lim_{h \to 0} \left( n x^{n-1} + \text{terms with } h \right) = n x^{n-1}\]

The rule is proved. Notice what happened: the limit killed every term
except the leading one. This is why the power rule is so clean.

\textbf{Examples:}

\begin{longtable}[]{@{}ll@{}}
\toprule\noalign{}
\(f(x)\) & \(f'(x)\) \\
\midrule\noalign{}
\endhead
\bottomrule\noalign{}
\endlastfoot
\(x^5\) & \(5x^4\) \\
\(x^{-2}\) & \(-2x^{-3}\) \\
\(\sqrt{x} = x^{1/2}\) & \(\tfrac{1}{2} x^{-1/2}\) \\
\(1 = x^0\) & \(0\) \\
\end{longtable}

The last row says: the derivative of any constant is zero. A constant
function has zero slope everywhere --- that makes sense.

\subsection{Sum/difference and constant multiple
rules}\label{sumdifference-and-constant-multiple-rules}

These follow directly from the limit definition: because the limit of a
sum is the sum of the limits, and a constant factor passes through a
limit.

\textbf{Constant multiple rule:}
\[\frac{d}{dx}\bigl[c \cdot f(x)\bigr] = c \cdot f'(x)\]

\textbf{Sum/difference rule:}
\[\frac{d}{dx}\bigl[f(x) \pm g(x)\bigr] = f'(x) \pm g'(x)\]

Together these mean you can differentiate a polynomial term by term. For
\(p(x) = 3x^4 - 2x^2 + 5x - 1\):

\[p'(x) = 3 \cdot 4x^3 - 2 \cdot 2x + 5 \cdot 1 - 0 = 12x^3 - 4x + 5\]

\subsection{Product rule}\label{product-rule}

The derivative of a product is \emph{not} the product of the
derivatives. That failure is easy to verify:
\(\tfrac{d}{dx}[x \cdot x] = \tfrac{d}{dx}[x^2] = 2x\), but
\(\tfrac{d}{dx}[x] \times \tfrac{d}{dx}[x] = 1 \times 1 = 1\). Different
answers.

The correct rule: if \(f\) and \(g\) are both differentiable, then:

\[\frac{d}{dx}\bigl[f(x)\cdot g(x)\bigr] = f'(x)\cdot g(x) + f(x)\cdot g'(x)\]

A useful mnemonic: \emph{derivative of first times second, plus first
times derivative of second.} The rule accounts for how both factors are
changing simultaneously.

\textbf{Check:}
\(\tfrac{d}{dx}[x \cdot x] = 1 \cdot x + x \cdot 1 = 2x\). Correct.

\subsection{Chain rule}\label{chain-rule}

The chain rule handles \textbf{composite functions} --- functions of
functions. If \(y = f(u)\) and \(u = g(x)\), so that \(y = f(g(x))\),
then:

\[\frac{dy}{dx} = \frac{dy}{du} \cdot \frac{du}{dx}\]

In prime notation:
\(\bigl(f \circ g\bigr)'(x) = f'\bigl(g(x)\bigr) \cdot g'(x)\).

This is the most important and most commonly misapplied rule in
calculus. The key skill is \textbf{identifying the outer and inner
functions} before differentiating.

\textbf{Example.} Differentiate \(y = \sin(x^2)\).

\begin{itemize}
\tightlist
\item
  Outer function: \(f(u) = \sin(u)\), so \(\dfrac{dy}{du} = \cos(u)\).
\item
  Inner function: \(g(x) = x^2\), so \(\dfrac{du}{dx} = 2x\).
\item
  Chain rule: \(\dfrac{dy}{dx} = \cos(x^2) \cdot 2x = 2x\cos(x^2)\).
\end{itemize}

The notation \(\dfrac{dy}{dx} = \dfrac{dy}{du} \cdot \dfrac{du}{dx}\)
looks like the \(du\)'s cancel --- and you can use it as if they do. The
reason it works is multiplicative structure in the limit, not literal
fraction cancellation, but the practical result is the same: set up the
product and compute.

\begin{Shaded}
\begin{Highlighting}[]
\NormalTok{//| echo: false}
\NormalTok{viewof chainFunc = Inputs.select(}
\NormalTok{  ["sin(x²)", "e\^{}(2x)", "(3x+1)⁴"],}
\NormalTok{  \{ label: "Composite function", value: "sin(x²)" \}}
\NormalTok{)}
\end{Highlighting}
\end{Shaded}

\begin{Shaded}
\begin{Highlighting}[]
\NormalTok{//| echo: false}
\NormalTok{viewof chainX = Inputs.range([{-}2, 2], \{}
\NormalTok{  step: 0.01,}
\NormalTok{  value: 0,}
\NormalTok{  label: "x"}
\NormalTok{\})}
\end{Highlighting}
\end{Shaded}

\begin{Shaded}
\begin{Highlighting}[]
\NormalTok{//| echo: false}
\NormalTok{\{}
\NormalTok{  const katexLib = await require("katex@0.16");}

\NormalTok{  const defs = \{}
\NormalTok{    "sin(x²)": \{}
\NormalTok{      fgx:    "\textbackslash{}\textbackslash{}sin(x\^{}2)",}
\NormalTok{      gx:     "x\^{}2",}
\NormalTok{      fu:     "\textbackslash{}\textbackslash{}sin(u)",}
\NormalTok{      fprimeGx: "\textbackslash{}\textbackslash{}cos(x\^{}2)",}
\NormalTok{      gprime: "2x",}
\NormalTok{      product: "\textbackslash{}\textbackslash{}cos(x\^{}2) \textbackslash{}\textbackslash{}cdot 2x",}
\NormalTok{      numericFn: x =\textgreater{} Math.cos(x * x) * 2 * x,}
\NormalTok{    \},}
\NormalTok{    "e\^{}(2x)": \{}
\NormalTok{      fgx:    "e\^{}\{2x\}",}
\NormalTok{      gx:     "2x",}
\NormalTok{      fu:     "e\^{}\{u\}",}
\NormalTok{      fprimeGx: "e\^{}\{2x\}",}
\NormalTok{      gprime: "2",}
\NormalTok{      product: "e\^{}\{2x\} \textbackslash{}\textbackslash{}cdot 2",}
\NormalTok{      numericFn: x =\textgreater{} Math.exp(2 * x) * 2,}
\NormalTok{    \},}
\NormalTok{    "(3x+1)⁴": \{}
\NormalTok{      fgx:    "(3x+1)\^{}4",}
\NormalTok{      gx:     "3x+1",}
\NormalTok{      fu:     "u\^{}4",}
\NormalTok{      fprimeGx: "4(3x+1)\^{}3",}
\NormalTok{      gprime: "3",}
\NormalTok{      product: "12(3x+1)\^{}3",}
\NormalTok{      numericFn: x =\textgreater{} 12 * Math.pow(3 * x + 1, 3),}
\NormalTok{    \},}
\NormalTok{  \};}

\NormalTok{  const d = defs[chainFunc];}
\NormalTok{  const numVal = d.numericFn(chainX);}
\NormalTok{  const numDisplay = isFinite(numVal) ? numVal.toFixed(4) : "undefined";}

\NormalTok{  const render = (tex, display = false) =\textgreater{}}
\NormalTok{    katexLib.renderToString(tex, \{ throwOnError: false, displayMode: display \});}

\NormalTok{  const box = document.createElement("div");}
\NormalTok{  box.style.cssText = "border:1px solid \#e5e7eb; border{-}radius:8px; padding:1rem 1.25rem; margin:0.75rem 0; font{-}family:inherit; max{-}width:560px;";}

\NormalTok{  const row = (label, tex) =\textgreater{} \{}
\NormalTok{    const r = document.createElement("div");}
\NormalTok{    r.style.cssText = "display:grid; grid{-}template{-}columns:200px 1fr; gap:0.3rem 0.75rem; align{-}items:baseline; padding:0.35rem 0; border{-}top:1px solid \#f3f4f6;";}
\NormalTok{    const lEl = document.createElement("span");}
\NormalTok{    lEl.style.cssText = "font{-}size:0.8rem; color:\#6b7280; font{-}style:italic;";}
\NormalTok{    lEl.textContent = label;}
\NormalTok{    const rEl = document.createElement("span");}
\NormalTok{    rEl.innerHTML = render(tex);}
\NormalTok{    r.appendChild(lEl);}
\NormalTok{    r.appendChild(rEl);}
\NormalTok{    return r;}
\NormalTok{  \};}

\NormalTok{  const heading = document.createElement("p");}
\NormalTok{  heading.style.cssText = "margin:0 0 0.6rem 0; font{-}size:0.9rem; color:\#374151; font{-}weight:600;";}
\NormalTok{  heading.innerHTML = "Chain rule: " + render(\textasciigrave{}\textbackslash{}\textbackslash{}dfrac\{d\}\{dx\}[f(g(x))] = f\textquotesingle{}(g(x)) \textbackslash{}\textbackslash{}cdot g\textquotesingle{}(x)\textasciigrave{});}
\NormalTok{  box.appendChild(heading);}

\NormalTok{  box.appendChild(row("function f(g(x))", d.fgx));}
\NormalTok{  box.appendChild(row("inner: g(x)", d.gx));}
\NormalTok{  box.appendChild(row("outer: f(u)", d.fu));}
\NormalTok{  box.appendChild(row("f′(g(x))", d.fprimeGx));}
\NormalTok{  box.appendChild(row("g′(x)", d.gprime));}
\NormalTok{  box.appendChild(row("product f′(g(x))·g′(x)", d.product));}

\NormalTok{  const divider = document.createElement("div");}
\NormalTok{  divider.style.cssText = "margin:0.6rem 0 0.4rem 0; border{-}top:1px solid \#e5e7eb;";}
\NormalTok{  box.appendChild(divider);}

\NormalTok{  const valRow = document.createElement("div");}
\NormalTok{  valRow.style.cssText = "display:grid; grid{-}template{-}columns:200px 1fr; gap:0.3rem 0.75rem; align{-}items:baseline; padding:0.25rem 0;";}
\NormalTok{  const valLabel = document.createElement("span");}
\NormalTok{  valLabel.style.cssText = "font{-}size:0.8rem; color:\#6b7280; font{-}style:italic;";}
\NormalTok{  valLabel.textContent = \textasciigrave{}numeric value at x = $\{chainX.toFixed(2)\}\textasciigrave{};}
\NormalTok{  const valNum = document.createElement("span");}
\NormalTok{  valNum.style.cssText = "font{-}weight:600; color:\#2563eb; font{-}size:0.95rem;";}
\NormalTok{  valNum.textContent = numDisplay;}
\NormalTok{  valRow.appendChild(valLabel);}
\NormalTok{  valRow.appendChild(valNum);}
\NormalTok{  box.appendChild(valRow);}

\NormalTok{  return box;}
\NormalTok{\}}
\end{Highlighting}
\end{Shaded}

\subsection{Derivatives of key
functions}\label{derivatives-of-key-functions}

These are stated here without proof. Each can be derived from the limit
definition, but the derivations involve limits of trigonometric
quotients and the definition of \(e\) --- both of which are established
results.

\begin{longtable}[]{@{}
  >{\raggedright\arraybackslash}p{(\linewidth - 4\tabcolsep) * \real{0.3333}}
  >{\raggedright\arraybackslash}p{(\linewidth - 4\tabcolsep) * \real{0.3333}}
  >{\raggedright\arraybackslash}p{(\linewidth - 4\tabcolsep) * \real{0.3333}}@{}}
\toprule\noalign{}
\begin{minipage}[b]{\linewidth}\raggedright
Function
\end{minipage} & \begin{minipage}[b]{\linewidth}\raggedright
Derivative
\end{minipage} & \begin{minipage}[b]{\linewidth}\raggedright
Read the derivative as
\end{minipage} \\
\midrule\noalign{}
\endhead
\bottomrule\noalign{}
\endlastfoot
\(\sin x\) & \(\cos x\) & Cosine is the rate of change of sine ---
maximum slope at \(x=0\), zero slope at the peaks \\
\(\cos x\) & \(-\sin x\) & Sine with a sign flip --- sine is decreasing
where cosine reaches its peak \\
\(e^x\) & \(e^x\) & The exponential is its own derivative --- the rate
of growth equals the current value \\
\(\ln x\) & \(\dfrac{1}{x}\) & Valid for \(x > 0\) --- the domain of
\(\ln x\). For \(x < 0\), use \(\ln\lvert x\rvert\), whose derivative is
also \(\frac{1}{x}\) \\
\end{longtable}

The self-referential property of \(e^x\) --- that it equals its own
derivative --- is why \(e\) appears everywhere in physics and
engineering. Any system whose rate of change is proportional to its
current value is governed by an exponential, and the differential
equation describing it is \(\dfrac{dy}{dx} = y\).

\section{Higher derivatives}\label{higher-derivatives}

The derivative of \(f\) is itself a function. You can differentiate it
again. The \textbf{second derivative} \(f''(x)\) (or
\(\dfrac{d^2y}{dx^2}\)) is the rate of change of the rate of change.

The clearest physical example: if \(s(t)\) is position at time \(t\),
then \(s'(t)\) is velocity and \(s''(t)\) is acceleration. Acceleration
is how quickly velocity is changing.

The second derivative also carries geometric information:

\begin{itemize}
\tightlist
\item
  \(f''(x) > 0\): the curve is \textbf{concave up} at \(x\) --- the
  slope is increasing, the curve bends upward like a bowl.
\item
  \(f''(x) < 0\): the curve is \textbf{concave down} at \(x\) --- the
  slope is decreasing, the curve bends downward like an arch.
\end{itemize}

This is what the second derivative test uses.

\section{Applications: optimisation}\label{applications-optimisation}

One of the most immediate uses of calculus is finding where a function
reaches its maximum or minimum value.

\textbf{Critical points.} If \(f\) is differentiable and has a local
maximum or minimum at \(x = c\), then the tangent line at \(c\) must be
horizontal --- a tilted tangent would mean the function is still rising
or falling, so \(c\) couldn't be an extreme point. Therefore:

\[f'(c) = 0 \implies c \text{ is a critical point}\]

Not every critical point is a maximum or minimum --- it could be a
\textbf{saddle point} (like \(x = 0\) for \(f(x) = x^3\), where the
function is neither increasing nor decreasing at that instant but
continues in the same direction on both sides). The second derivative
distinguishes the cases.

\textbf{Second derivative test.} If \(f'(c) = 0\):

\begin{itemize}
\tightlist
\item
  If \(f''(c) > 0\): concave up at \(c\), so the tangent is at the
  bottom of a bowl --- \textbf{local minimum}.
\item
  If \(f''(c) < 0\): concave down at \(c\), so the tangent is at the top
  of an arch --- \textbf{local maximum}.
\item
  If \(f''(c) = 0\): the test is inconclusive. What to do when
  \(f''(c) = 0\) is deferred to multivariable calculus --- the tools for
  that case are richer there. In single-variable calculus, if the second
  derivative test is inconclusive, you check the sign of \(f'\) on
  either side of \(c\) directly.
\end{itemize}

\textbf{Worked procedure.} Find the local extrema of
\(f(x) = x^3 - 6x^2 + 9x + 1\).

Step 1. Differentiate: \(f'(x) = 3x^2 - 12x + 9\).

Step 2. Set \(f'(x) = 0\):
\(3x^2 - 12x + 9 = 0 \implies x^2 - 4x + 3 = 0 \implies (x-1)(x-3) = 0\).

Critical points: \(x = 1\) and \(x = 3\).

Step 3. Compute \(f''(x) = 6x - 12\).

Step 4. Test each critical point:

\begin{itemize}
\tightlist
\item
  \(f''(1) = 6(1) - 12 = -6 < 0\): local maximum at \(x = 1\). Value:
  \(f(1) = 1 - 6 + 9 + 1 = 5\).
\item
  \(f''(3) = 6(3) - 12 = 6 > 0\): local minimum at \(x = 3\). Value:
  \(f(3) = 27 - 54 + 27 + 1 = 1\).
\end{itemize}

The function rises to a local maximum of 5 at \(x = 1\), falls to a
local minimum of 1 at \(x = 3\), then rises again.

\section{Worked examples}\label{worked-examples-18}

\begin{center}\rule{0.5\linewidth}{0.5pt}\end{center}

\textbf{Example 1 (Computing/data).} Differentiate
\(f(x) = 3x^4 - 2x^2 + 5x - 1\).

Apply the power rule term by term, with the constant multiple and sum
rules:

\[f'(x) = 3 \cdot 4x^3 - 2 \cdot 2x^1 + 5 \cdot 1 - 0\]

\[f'(x) = 12x^3 - 4x + 5\]

This is a polynomial --- differentiating a polynomial always produces a
polynomial of degree one lower. The constant term \(-1\) vanishes: a
constant has zero slope everywhere.

\begin{center}\rule{0.5\linewidth}{0.5pt}\end{center}

\textbf{Example 2 (Science).} Differentiate \(f(x) = x^2 \sin(x)\).

This is a product. Identify \(f(x) = x^2\) (the first factor) and
\(g(x) = \sin(x)\) (the second). Apply the product rule:

\[\frac{d}{dx}\bigl[x^2 \sin x\bigr] = 2x \cdot \sin x + x^2 \cdot \cos x\]

\[= 2x\sin x + x^2 \cos x\]

This kind of function appears in wave mechanics: \(x^2 \sin x\) models
an oscillation whose amplitude grows quadratically with distance from
the source.

\begin{center}\rule{0.5\linewidth}{0.5pt}\end{center}

\textbf{Example 3 (Engineering).} Differentiate \(f(x) = e^{3x^2}\).

This is a composite function. Identify the outer and inner functions:

\begin{itemize}
\tightlist
\item
  Outer: \(f(u) = e^u\), which has derivative \(e^u\).
\item
  Inner: \(g(x) = 3x^2\), which has derivative \(6x\).
\end{itemize}

Chain rule:

\[\frac{d}{dx}\bigl[e^{3x^2}\bigr] = e^{3x^2} \cdot 6x = 6x\,e^{3x^2}\]

Functions of this form appear in Gaussian distributions, heat diffusion
kernels, and signal-processing windows. The chain rule is the only way
to differentiate them correctly.

\begin{center}\rule{0.5\linewidth}{0.5pt}\end{center}

\textbf{Example 4 (Real-world optimisation).} A farmer has 80 m of
fencing and wants to enclose a rectangular paddock. One side is a
riverbank, which needs no fence. What dimensions maximise the enclosed
area?

\textbf{Set up.} Let the side perpendicular to the river have length
\(x\) (metres). Two of these sides are needed. The side parallel to the
river, call it \(y\), uses the remaining fencing:

\[2x + y = 80 \implies y = 80 - 2x\]

The area is:

\[A(x) = x \cdot y = x(80 - 2x) = 80x - 2x^2\]

\textbf{Differentiate and find critical points:}

\[A'(x) = 80 - 4x\]

Setting \(A'(x) = 0\): \(80 - 4x = 0 \implies x = 20\).

\textbf{Classify:} \(A''(x) = -4 < 0\) everywhere, so \(x = 20\) is a
local (and global) maximum.

\textbf{Dimensions:} \(x = 20\) m, \(y = 80 - 2(20) = 40\) m.

\textbf{Maximum area:} \(A = 20 \times 40 = 800\text{ m}^2\).

The shape that maximises the area is twice as wide as it is deep --- a
result that appears in architectural and agricultural optimisation
problems wherever one boundary is free.

\section{Where this goes}\label{where-this-goes-19}

\textbf{Integral calculus (ch03):} Integration inverts differentiation.
Where this chapter asked ``given a function, find its rate of change,''
the next chapter asks ``given the rate of change, recover the original
function.'' The Fundamental Theorem of Calculus --- the deepest result
in the subject --- shows these two operations are inverses of each
other. Every calculation of area, volume, work, and accumulated change
runs through it.

\textbf{Ordinary differential equations (Vol 7):} A differential
equation is a relationship between a function and its own derivative ---
an equation that says something like ``the rate of change of this
quantity is proportional to the quantity itself.'' Everything in this
chapter is prerequisite: you need to know what a derivative is, how to
compute one, and how the chain rule works, before a differential
equation makes sense. The entire Vol 7 ODE sequence begins here.

\begin{tcolorbox}[enhanced jigsaw, opacityback=0, colback=white, breakable, colframe=quarto-callout-tip-color-frame, arc=.35mm, left=2mm, toprule=.15mm, leftrule=.75mm, bottomrule=.15mm, rightrule=.15mm]
\begin{minipage}[t]{5.5mm}
\textcolor{quarto-callout-tip-color}{\faLightbulb}
\end{minipage}%
\begin{minipage}[t]{\textwidth - 5.5mm}

\textbf{Where this shows up}

\begin{itemize}
\tightlist
\item
  \textbf{Velocity and acceleration.} If \(s(t)\) is the position of an
  object at time \(t\), then \(s'(t)\) is its velocity and \(s''(t)\) is
  its acceleration. Every problem in Newtonian mechanics is stated in
  these terms.
\item
  \textbf{Marginal analysis.} In economics, marginal cost is the
  derivative of total cost with respect to quantity. Profit is maximised
  where marginal revenue equals marginal cost --- a critical point of
  the profit function.
\item
  \textbf{Newton's law of cooling.} The rate of heat loss of an object
  is proportional to the difference between the object's temperature and
  the ambient temperature. That ``rate of heat loss'' is a derivative.
  The law is \(\tfrac{dT}{dt} = -k(T - T_\text{ambient})\) --- a
  differential equation that lives in Vol 7.
\item
  \textbf{Gradient descent.} The optimisation algorithm that trains
  neural networks works by computing the derivative of the loss function
  with respect to each model parameter, then updating the parameters in
  the direction of steepest descent. The chain rule --- applied
  thousands of times through a deep network --- is called
  backpropagation.
\end{itemize}

\end{minipage}%
\end{tcolorbox}

\section{Exercises}\label{exercises-19}

These are puzzles. Each has a clean answer, but the interesting part is
choosing the right rule and executing the steps carefully before you
reach for the result.

\begin{Shaded}
\begin{Highlighting}[]
\NormalTok{//| echo: false}
\NormalTok{katex = require("katex@0.16")}
\end{Highlighting}
\end{Shaded}

\begin{Shaded}
\begin{Highlighting}[]
\NormalTok{//| echo: false}
\NormalTok{function makeStepperHTML(exerciseNum, steps) \{}
\NormalTok{  let current = 0;}
\NormalTok{  const totalSteps = steps.length;}
\NormalTok{  const container = document.createElement("div");}
\NormalTok{  container.style.cssText = "border:1px solid \#e5e7eb; border{-}radius:8px; padding:1rem 1.25rem; margin:0.75rem 0; font{-}family:inherit;";}
\NormalTok{  const stepsDiv = document.createElement("div");}
\NormalTok{  stepsDiv.style.marginBottom = "0.75rem";}
\NormalTok{  function renderSteps() \{}
\NormalTok{    stepsDiv.innerHTML = "";}
\NormalTok{    for (let i = 0; i \textless{} current; i++) \{}
\NormalTok{      const s = steps[i];}
\NormalTok{      const row = document.createElement("div");}
\NormalTok{      row.style.cssText = "display:grid; grid{-}template{-}columns:220px 1fr; gap:0.35rem 1rem; align{-}items:baseline; padding:0.35rem 0; border{-}top:1px solid \#f3f4f6;";}
\NormalTok{      const opEl = document.createElement("span");}
\NormalTok{      opEl.style.cssText = "font{-}size:0.8rem; color:\#6b7280; font{-}style:italic;";}
\NormalTok{      opEl.textContent = s.op;}
\NormalTok{      const eqEl = document.createElement("span");}
\NormalTok{      eqEl.innerHTML = katex.renderToString(s.eq, \{throwOnError: false, displayMode: false\});}
\NormalTok{      row.appendChild(opEl);}
\NormalTok{      row.appendChild(eqEl);}
\NormalTok{      stepsDiv.appendChild(row);}
\NormalTok{      if (s.note) \{}
\NormalTok{        const noteEl = document.createElement("p");}
\NormalTok{        noteEl.style.cssText = "margin:0.2rem 0 0.5rem 220px; font{-}size:0.82rem; color:\#374151;";}
\NormalTok{        noteEl.textContent = s.note;}
\NormalTok{        stepsDiv.appendChild(noteEl);}
\NormalTok{      \}}
\NormalTok{    \}}
\NormalTok{    if (current === totalSteps) \{}
\NormalTok{      const done = document.createElement("p");}
\NormalTok{      done.style.cssText = "margin{-}top:0.5rem; font{-}size:0.85rem; color:\#059669; font{-}weight:600;";}
\NormalTok{      done.textContent = "✓ Solution complete";}
\NormalTok{      stepsDiv.appendChild(done);}
\NormalTok{    \}}
\NormalTok{  \}}
\NormalTok{  const controls = document.createElement("div");}
\NormalTok{  controls.style.cssText = "display:flex; align{-}items:center; gap:0.75rem; margin{-}top:0.5rem;";}
\NormalTok{  const nextBtn = document.createElement("button");}
\NormalTok{  nextBtn.textContent = "Next step →";}
\NormalTok{  nextBtn.style.cssText = "padding:0.35rem 0.85rem; background:\#2563eb; color:\#fff; border:none; border{-}radius:6px; cursor:pointer; font{-}size:0.85rem;";}
\NormalTok{  const resetBtn = document.createElement("button");}
\NormalTok{  resetBtn.textContent = "Reset";}
\NormalTok{  resetBtn.style.cssText = "padding:0.35rem 0.85rem; background:\#f3f4f6; color:\#374151; border:1px solid \#d1d5db; border{-}radius:6px; cursor:pointer; font{-}size:0.85rem;";}
\NormalTok{  const counter = document.createElement("span");}
\NormalTok{  counter.style.cssText = "font{-}size:0.8rem; color:\#6b7280;";}
\NormalTok{  nextBtn.onclick = () =\textgreater{} \{}
\NormalTok{    if (current \textless{} totalSteps) \{}
\NormalTok{      current++;}
\NormalTok{      renderSteps();}
\NormalTok{      counter.textContent = current + " / " + totalSteps;}
\NormalTok{      if (current === totalSteps) nextBtn.disabled = true;}
\NormalTok{    \}}
\NormalTok{  \};}
\NormalTok{  resetBtn.onclick = () =\textgreater{} \{}
\NormalTok{    current = 0;}
\NormalTok{    nextBtn.disabled = false;}
\NormalTok{    renderSteps();}
\NormalTok{    counter.textContent = "0 / " + totalSteps;}
\NormalTok{  \};}
\NormalTok{  counter.textContent = "0 / " + totalSteps;}
\NormalTok{  controls.appendChild(nextBtn);}
\NormalTok{  controls.appendChild(resetBtn);}
\NormalTok{  controls.appendChild(counter);}
\NormalTok{  container.appendChild(stepsDiv);}
\NormalTok{  container.appendChild(controls);}
\NormalTok{  return container;}
\NormalTok{\}}
\end{Highlighting}
\end{Shaded}

\textbf{Exercise 1.} Differentiate \(f(x) = 4x^5 - 3x^3 + 7x - 2\) using
the power rule.

\begin{Shaded}
\begin{Highlighting}[]
\NormalTok{//| echo: false}
\NormalTok{\{}
\NormalTok{  const el = makeStepperHTML(1, [}
\NormalTok{    \{ op: "Identify each term", eq: "4x\^{}5,\textbackslash{}\textbackslash{}quad {-}3x\^{}3,\textbackslash{}\textbackslash{}quad 7x,\textbackslash{}\textbackslash{}quad {-}2", note: "Differentiate term by term using the sum and constant multiple rules." \},}
\NormalTok{    \{ op: "Power rule on 4x⁵", eq: "\textbackslash{}\textbackslash{}frac\{d\}\{dx\}[4x\^{}5] = 4 \textbackslash{}\textbackslash{}cdot 5x\^{}4 = 20x\^{}4", note: null \},}
\NormalTok{    \{ op: "Power rule on −3x³", eq: "\textbackslash{}\textbackslash{}frac\{d\}\{dx\}[{-}3x\^{}3] = {-}3 \textbackslash{}\textbackslash{}cdot 3x\^{}2 = {-}9x\^{}2", note: null \},}
\NormalTok{    \{ op: "Power rule on 7x", eq: "\textbackslash{}\textbackslash{}frac\{d\}\{dx\}[7x] = 7 \textbackslash{}\textbackslash{}cdot 1 = 7", note: "x = x¹, so the power is 1 and the result is the constant coefficient." \},}
\NormalTok{    \{ op: "Derivative of constant", eq: "\textbackslash{}\textbackslash{}frac\{d\}\{dx\}[{-}2] = 0", note: "A constant has zero slope everywhere." \},}
\NormalTok{    \{ op: "Combine", eq: "f\textquotesingle{}(x) = 20x\^{}4 {-} 9x\^{}2 + 7", note: null \},}
\NormalTok{  ]);}
\NormalTok{  return el;}
\NormalTok{\}}
\end{Highlighting}
\end{Shaded}

\begin{center}\rule{0.5\linewidth}{0.5pt}\end{center}

\textbf{Exercise 2.} Differentiate \(f(x) = x^3 \cos(x)\) using the
product rule.

\begin{Shaded}
\begin{Highlighting}[]
\NormalTok{//| echo: false}
\NormalTok{\{}
\NormalTok{  const el = makeStepperHTML(2, [}
\NormalTok{    \{ op: "Identify the two factors", eq: "u = x\^{}3,\textbackslash{}\textbackslash{}quad v = \textbackslash{}\textbackslash{}cos x", note: null \},}
\NormalTok{    \{ op: "Differentiate each factor", eq: "u\textquotesingle{} = 3x\^{}2,\textbackslash{}\textbackslash{}quad v\textquotesingle{} = {-}\textbackslash{}\textbackslash{}sin x", note: null \},}
\NormalTok{    \{ op: "Apply product rule", eq: "\textbackslash{}\textbackslash{}frac\{d\}\{dx\}[x\^{}3 \textbackslash{}\textbackslash{}cos x] = u\textquotesingle{}v + uv\textquotesingle{}", note: "Derivative of first × second, plus first × derivative of second." \},}
\NormalTok{    \{ op: "Substitute", eq: "= 3x\^{}2 \textbackslash{}\textbackslash{}cos x + x\^{}3 \textbackslash{}\textbackslash{}cdot ({-}\textbackslash{}\textbackslash{}sin x)", note: null \},}
\NormalTok{    \{ op: "Write the result", eq: "f\textquotesingle{}(x) = 3x\^{}2 \textbackslash{}\textbackslash{}cos x {-} x\^{}3 \textbackslash{}\textbackslash{}sin x", note: null \},}
\NormalTok{  ]);}
\NormalTok{  return el;}
\NormalTok{\}}
\end{Highlighting}
\end{Shaded}

\begin{center}\rule{0.5\linewidth}{0.5pt}\end{center}

\textbf{Exercise 3.} Differentiate \(f(x) = (2x^3 + 1)^5\) using the
chain rule.

\begin{Shaded}
\begin{Highlighting}[]
\NormalTok{//| echo: false}
\NormalTok{\{}
\NormalTok{  const el = makeStepperHTML(3, [}
\NormalTok{    \{ op: "Identify outer and inner functions", eq: "f(u) = u\^{}5,\textbackslash{}\textbackslash{}quad g(x) = 2x\^{}3 + 1", note: "The outer function acts on whatever is inside the bracket." \},}
\NormalTok{    \{ op: "Differentiate the outer", eq: "\textbackslash{}\textbackslash{}frac\{d\}\{du\}[u\^{}5] = 5u\^{}4", note: null \},}
\NormalTok{    \{ op: "Differentiate the inner", eq: "\textbackslash{}\textbackslash{}frac\{d\}\{dx\}[2x\^{}3 + 1] = 6x\^{}2", note: null \},}
\NormalTok{    \{ op: "Apply chain rule", eq: "f\textquotesingle{}(x) = 5u\^{}4 \textbackslash{}\textbackslash{}cdot 6x\^{}2", note: "Multiply: derivative of outer (with inner substituted back) × derivative of inner." \},}
\NormalTok{    \{ op: "Substitute u = 2x³ + 1", eq: "f\textquotesingle{}(x) = 5(2x\^{}3+1)\^{}4 \textbackslash{}\textbackslash{}cdot 6x\^{}2", note: null \},}
\NormalTok{    \{ op: "Simplify", eq: "f\textquotesingle{}(x) = 30x\^{}2(2x\^{}3+1)\^{}4", note: null \},}
\NormalTok{  ]);}
\NormalTok{  return el;}
\NormalTok{\}}
\end{Highlighting}
\end{Shaded}

\begin{center}\rule{0.5\linewidth}{0.5pt}\end{center}

\textbf{Exercise 4.} Find all critical points of
\(f(x) = x^3 - 3x^2 - 9x + 2\) and classify each as a local maximum,
local minimum, or neither.

\begin{Shaded}
\begin{Highlighting}[]
\NormalTok{//| echo: false}
\NormalTok{\{}
\NormalTok{  const el = makeStepperHTML(4, [}
\NormalTok{    \{ op: "Differentiate f", eq: "f\textquotesingle{}(x) = 3x\^{}2 {-} 6x {-} 9", note: null \},}
\NormalTok{    \{ op: "Set f\textquotesingle{}(x) = 0", eq: "3x\^{}2 {-} 6x {-} 9 = 0 \textbackslash{}\textbackslash{}implies x\^{}2 {-} 2x {-} 3 = 0", note: "Divide through by 3 to simplify." \},}
\NormalTok{    \{ op: "Factorise", eq: "(x{-}3)(x+1) = 0 \textbackslash{}\textbackslash{}implies x = 3 \textbackslash{}\textbackslash{}text\{ or \} x = {-}1", note: null \},}
\NormalTok{    \{ op: "Compute second derivative", eq: "f\textquotesingle{}\textquotesingle{}(x) = 6x {-} 6", note: null \},}
\NormalTok{    \{ op: "Test x = −1", eq: "f\textquotesingle{}\textquotesingle{}({-}1) = 6({-}1) {-} 6 = {-}12 \textless{} 0", note: "Concave down → local maximum." \},}
\NormalTok{    \{ op: "Value at x = −1", eq: "f({-}1) = ({-}1)\^{}3 {-} 3({-}1)\^{}2 {-} 9({-}1) + 2 = {-}1 {-} 3 + 9 + 2 = 7", note: "Local maximum of 7 at x = −1." \},}
\NormalTok{    \{ op: "Test x = 3", eq: "f\textquotesingle{}\textquotesingle{}(3) = 6(3) {-} 6 = 12 \textgreater{} 0", note: "Concave up → local minimum." \},}
\NormalTok{    \{ op: "Value at x = 3", eq: "f(3) = 27 {-} 27 {-} 27 + 2 = {-}25", note: "Local minimum of −25 at x = 3." \},}
\NormalTok{  ]);}
\NormalTok{  return el;}
\NormalTok{\}}
\end{Highlighting}
\end{Shaded}

\begin{center}\rule{0.5\linewidth}{0.5pt}\end{center}

\textbf{Exercise 5.} A manufacturer's total cost in dollars for
producing \(q\) units per day is \(C(q) = q^3 - 6q^2 + 15q + 50\). At
what production level is the marginal cost minimised? (Marginal cost is
\(C'(q)\).)

\begin{Shaded}
\begin{Highlighting}[]
\NormalTok{//| echo: false}
\NormalTok{\{}
\NormalTok{  const el = makeStepperHTML(5, [}
\NormalTok{    \{ op: "Find marginal cost C\textquotesingle{}(q)", eq: "C\textquotesingle{}(q) = 3q\^{}2 {-} 12q + 15", note: "Marginal cost is the rate of change of total cost." \},}
\NormalTok{    \{ op: "Minimise C\textquotesingle{}(q): differentiate again", eq: "C\textquotesingle{}\textquotesingle{}(q) = 6q {-} 12", note: "To find the minimum of C\textquotesingle{}(q), treat it as a function and find its critical points." \},}
\NormalTok{    \{ op: "Set C\textquotesingle{}\textquotesingle{}(q) = 0", eq: "6q {-} 12 = 0 \textbackslash{}\textbackslash{}implies q = 2", note: null \},}
\NormalTok{    \{ op: "Check sign of C\textquotesingle{}(q) on each side", eq: "C\textquotesingle{}(1) = 3 {-} 12 + 15 = 6,\textbackslash{}\textbackslash{}quad C\textquotesingle{}(3) = 27 {-} 36 + 15 = 6", note: "Evaluate C\textquotesingle{}(q) at points just below and just above q = 2 to confirm this is a minimum." \},}
\NormalTok{    \{ op: "Confirm minimum: check C\textquotesingle{}\textquotesingle{}(q) sign", eq: "C\textquotesingle{}\textquotesingle{}(q) = 6q {-} 12 \textless{} 0 \textbackslash{}\textbackslash{}text\{ for \} q \textless{} 2,\textbackslash{}\textbackslash{}quad C\textquotesingle{}\textquotesingle{}(q) \textgreater{} 0 \textbackslash{}\textbackslash{}text\{ for \} q \textgreater{} 2", note: "C\textquotesingle{}(q) is decreasing then increasing around q = 2, confirming q = 2 is a local minimum of C\textquotesingle{}(q)." \},}
\NormalTok{    \{ op: "Minimum marginal cost value", eq: "C\textquotesingle{}(2) = 3(4) {-} 12(2) + 15 = 12 {-} 24 + 15 = 3", note: null \},}
\NormalTok{    \{ op: "Interpret", eq: "q = 2 \textbackslash{}\textbackslash{}text\{ units/day gives minimum marginal cost of \} \textbackslash{}\textbackslash{}$3 \textbackslash{}\textbackslash{}text\{ per unit\}", note: "Below this level, each extra unit is cheaper than the next one will be; above it, costs are rising faster." \},}
\NormalTok{  ]);}
\NormalTok{  return el;}
\NormalTok{\}}
\end{Highlighting}
\end{Shaded}

\begin{center}\rule{0.5\linewidth}{0.5pt}\end{center}

\textbf{Exercise 6.} A stone is dropped from a bridge. Its height above
the water at time \(t\) seconds is \(h(t) = 45 - 4.9t^2\) metres. At
what speed (in m/s) is it falling at the instant it hits the water?

\begin{Shaded}
\begin{Highlighting}[]
\NormalTok{//| echo: false}
\NormalTok{\{}
\NormalTok{  const el = makeStepperHTML(6, [}
\NormalTok{    \{ op: "Find the time of impact", eq: "h(t) = 0 \textbackslash{}\textbackslash{}implies 45 {-} 4.9t\^{}2 = 0", note: "The stone hits the water when h(t) = 0." \},}
\NormalTok{    \{ op: "Solve for t", eq: "t\^{}2 = \textbackslash{}\textbackslash{}frac\{45\}\{4.9\} \textbackslash{}\textbackslash{}approx 9.184 \textbackslash{}\textbackslash{}implies t \textbackslash{}\textbackslash{}approx 3.031 \textbackslash{}\textbackslash{}text\{ s\}", note: null \},}
\NormalTok{    \{ op: "Find the velocity function", eq: "h\textquotesingle{}(t) = {-}9.8t", note: "Velocity is the derivative of height. The negative sign means the stone is descending." \},}
\NormalTok{    \{ op: "Speed at impact", eq: "h\textquotesingle{}(3.031) = {-}9.8 \textbackslash{}\textbackslash{}times 3.031 \textbackslash{}\textbackslash{}approx {-}29.7 \textbackslash{}\textbackslash{}text\{ m/s\}", note: null \},}
\NormalTok{    \{ op: "Interpret", eq: "\textbackslash{}\textbackslash{}text\{Speed\} = |h\textquotesingle{}(3.031)| \textbackslash{}\textbackslash{}approx 29.7 \textbackslash{}\textbackslash{}text\{ m/s\}", note: "Speed is the magnitude of velocity. The stone is falling at approximately 29.7 m/s when it reaches the water — about 107 km/h." \},}
\NormalTok{  ]);}
\NormalTok{  return el;}
\NormalTok{\}}
\end{Highlighting}
\end{Shaded}

\chapter{Integral calculus}\label{integral-calculus}

Accumulation

\hfill\break

Your car's speedometer shows 80 km/h. Two minutes later it shows 95
km/h. How far have you travelled? You cannot just multiply speed by time
--- the speed was changing the whole way. What you need is a way to add
up all those instantaneous velocities over the interval. That is
integration.

A water tank fills at a rate that varies with time --- fast at first,
slower as back-pressure builds. You want to know how much water is in
the tank after an hour. You know the rate at each moment. You need the
total.

These are the same problem. This chapter builds the tool.

\section{What the notation is
saying}\label{what-the-notation-is-saying-20}

The integral sign \(\int\) is an elongated S --- it stands for
\emph{sum}. The \(dx\) at the end names the variable we are integrating
with respect to. The expression

\[\int_a^b f(x)\,dx\]

reads: \emph{``the integral of \(f(x)\) with respect to \(x\), from
\(a\) to \(b\).''} Picture many thin rectangles, each with height
\(f(x)\) and width \(\Delta x\) --- we'll shrink \(\Delta x\) to zero in
the limit, so the strips get thinner and thinner until they fit exactly
under the curve. The integral adds up the areas of all of them. The
result is the total area between the curve \(y = f(x)\) and the
\(x\)-axis, between \(x = a\) and \(x = b\).

That is the geometric picture. The computational meaning follows from
Riemann sums.

\subsection{Riemann sums}\label{riemann-sums}

Partition the interval \([a, b]\) into \(n\) equal subintervals, each of
width

\[\Delta x = \frac{b - a}{n}\]

Pick a sample point \(x_i\) in each subinterval --- say the left
endpoint, the right endpoint, or the midpoint. At each sample point, the
rectangle has height \(f(x_i)\) and width \(\Delta x\). Sum all \(n\)
rectangles:

\[\text{Area} \approx \sum_{i=1}^{n} f(x_i)\,\Delta x\]

This is a Riemann sum --- you saw the idea of summing up many small
terms in Volume 4, and the same question of what happens as the number
of terms grows without bound applies here too. The definite integral is
the limit of this sum as the rectangles get infinitely thin:

\[\int_a^b f(x)\,dx = \lim_{n \to \infty} \sum_{i=1}^{n} f(x_i)\,\Delta x\]

For well-behaved functions, this limit exists and is a single, precise
number.

\textbf{Interactive: Riemann sum builder.} Adjust the number of
rectangles and choose a sum type to see how the approximation converges
to the exact value.

\begin{Shaded}
\begin{Highlighting}[]
\NormalTok{//| echo: false}
\NormalTok{d3 = require("d3@7")}
\end{Highlighting}
\end{Shaded}

\begin{Shaded}
\begin{Highlighting}[]
\NormalTok{//| echo: false}
\NormalTok{viewof riemannN = Inputs.range([1, 50], \{}
\NormalTok{  label: "n (number of rectangles)",}
\NormalTok{  step: 1,}
\NormalTok{  value: 5}
\NormalTok{\})}
\end{Highlighting}
\end{Shaded}

\begin{Shaded}
\begin{Highlighting}[]
\NormalTok{//| echo: false}
\NormalTok{viewof riemannType = Inputs.select(["Left", "Right", "Midpoint"], \{}
\NormalTok{  label: "Sum type",}
\NormalTok{  value: "Left"}
\NormalTok{\})}
\end{Highlighting}
\end{Shaded}

\begin{Shaded}
\begin{Highlighting}[]
\NormalTok{//| echo: false}
\NormalTok{\{}
\NormalTok{  const W = 560, H = 320;}
\NormalTok{  const margin = \{ top: 20, right: 20, bottom: 50, left: 50 \};}
\NormalTok{  const innerW = W {-} margin.left {-} margin.right;}
\NormalTok{  const innerH = H {-} margin.top {-} margin.bottom;}

\NormalTok{  const a = 0, b = 3;}
\NormalTok{  const f = x =\textgreater{} x * x + 1;}
\NormalTok{  const exactArea = 12; // ∫₀³ (x²+1) dx = [x³/3 + x]₀³ = 9 + 3 = 12}

\NormalTok{  const n = Math.max(1, Math.round(riemannN));}
\NormalTok{  const type = riemannType;}
\NormalTok{  const dx = (b {-} a) / n;}

\NormalTok{  // Compute approximate area}
\NormalTok{  let approxArea = 0;}
\NormalTok{  for (let i = 0; i \textless{} n; i++) \{}
\NormalTok{    let xSample;}
\NormalTok{    if (type === "Left")     xSample = a + i * dx;}
\NormalTok{    else if (type === "Right") xSample = a + (i + 1) * dx;}
\NormalTok{    else                     xSample = a + (i + 0.5) * dx; // Midpoint}
\NormalTok{    approxArea += f(xSample) * dx;}
\NormalTok{  \}}

\NormalTok{  const xScale = d3.scaleLinear().domain([a, b]).range([0, innerW]);}
\NormalTok{  const yScale = d3.scaleLinear().domain([0, 11]).range([innerH, 0]);}

\NormalTok{  const svg = d3.create("svg")}
\NormalTok{    .attr("viewBox", \textasciigrave{}0 0 $\{W\} $\{H\}\textasciigrave{})}
\NormalTok{    .attr("width", "100\%")}
\NormalTok{    .attr("style", \textasciigrave{}max{-}width:$\{W\}px; font{-}family:inherit;\textasciigrave{});}

\NormalTok{  const g = svg.append("g")}
\NormalTok{    .attr("transform", \textasciigrave{}translate($\{margin.left\},$\{margin.top\})\textasciigrave{});}

\NormalTok{  // Draw rectangles}
\NormalTok{  for (let i = 0; i \textless{} n; i++) \{}
\NormalTok{    let xSample;}
\NormalTok{    if (type === "Left")       xSample = a + i * dx;}
\NormalTok{    else if (type === "Right") xSample = a + (i + 1) * dx;}
\NormalTok{    else                       xSample = a + (i + 0.5) * dx;}

\NormalTok{    const rectX = xScale(a + i * dx);}
\NormalTok{    const rectW = xScale(a + (i + 1) * dx) {-} rectX;}
\NormalTok{    const rectH = innerH {-} yScale(f(xSample));}

\NormalTok{    g.append("rect")}
\NormalTok{      .attr("x", rectX)}
\NormalTok{      .attr("y", yScale(f(xSample)))}
\NormalTok{      .attr("width", rectW)}
\NormalTok{      .attr("height", rectH)}
\NormalTok{      .attr("fill", "\#bfdbfe")}
\NormalTok{      .attr("stroke", "\#2563eb")}
\NormalTok{      .attr("stroke{-}width", 0.8)}
\NormalTok{      .attr("opacity", 0.7);}
\NormalTok{  \}}

\NormalTok{  // Draw x{-}axis}
\NormalTok{  g.append("line")}
\NormalTok{    .attr("x1", 0).attr("y1", innerH)}
\NormalTok{    .attr("x2", innerW).attr("y2", innerH)}
\NormalTok{    .attr("stroke", "\#374151").attr("stroke{-}width", 1.5);}

\NormalTok{  // Draw y{-}axis}
\NormalTok{  g.append("line")}
\NormalTok{    .attr("x1", 0).attr("y1", 0)}
\NormalTok{    .attr("x2", 0).attr("y2", innerH)}
\NormalTok{    .attr("stroke", "\#374151").attr("stroke{-}width", 1.5);}

\NormalTok{  // X{-}axis ticks}
\NormalTok{  [0, 1, 2, 3].forEach(v =\textgreater{} \{}
\NormalTok{    g.append("line")}
\NormalTok{      .attr("x1", xScale(v)).attr("y1", innerH)}
\NormalTok{      .attr("x2", xScale(v)).attr("y2", innerH + 5)}
\NormalTok{      .attr("stroke", "\#374151");}
\NormalTok{    g.append("text")}
\NormalTok{      .attr("x", xScale(v)).attr("y", innerH + 18)}
\NormalTok{      .attr("text{-}anchor", "middle")}
\NormalTok{      .attr("fill", "\#374151").attr("font{-}size", "12px")}
\NormalTok{      .text(v);}
\NormalTok{  \});}

\NormalTok{  // Y{-}axis ticks}
\NormalTok{  [0, 2, 4, 6, 8, 10].forEach(v =\textgreater{} \{}
\NormalTok{    g.append("line")}
\NormalTok{      .attr("x1", {-}5).attr("y1", yScale(v))}
\NormalTok{      .attr("x2", 0).attr("y2", yScale(v))}
\NormalTok{      .attr("stroke", "\#374151");}
\NormalTok{    g.append("text")}
\NormalTok{      .attr("x", {-}10).attr("y", yScale(v) + 4)}
\NormalTok{      .attr("text{-}anchor", "end")}
\NormalTok{      .attr("fill", "\#374151").attr("font{-}size", "11px")}
\NormalTok{      .text(v);}
\NormalTok{  \});}

\NormalTok{  // Draw curve on top}
\NormalTok{  const curvePoints = d3.range(0, b + 0.01, 0.02).map(x =\textgreater{} [xScale(x), yScale(f(x))]);}
\NormalTok{  const linePath = d3.line()(curvePoints);}
\NormalTok{  g.append("path")}
\NormalTok{    .attr("d", linePath)}
\NormalTok{    .attr("fill", "none")}
\NormalTok{    .attr("stroke", "\#1e40af")}
\NormalTok{    .attr("stroke{-}width", 2.5);}

\NormalTok{  // Axis labels}
\NormalTok{  g.append("text")}
\NormalTok{    .attr("x", innerW / 2).attr("y", innerH + 40)}
\NormalTok{    .attr("text{-}anchor", "middle")}
\NormalTok{    .attr("fill", "\#374151").attr("font{-}size", "13px")}
\NormalTok{    .text("x");}

\NormalTok{  g.append("text")}
\NormalTok{    .attr("transform", "rotate({-}90)")}
\NormalTok{    .attr("x", {-}innerH / 2).attr("y", {-}36)}
\NormalTok{    .attr("text{-}anchor", "middle")}
\NormalTok{    .attr("fill", "\#374151").attr("font{-}size", "13px")}
\NormalTok{    .text("f(x) = x² + 1");}

\NormalTok{  // Stats display below chart}
\NormalTok{  const statsY = innerH + 52;}
\NormalTok{  const error = Math.abs(approxArea {-} exactArea);}

\NormalTok{  svg.append("text")}
\NormalTok{    .attr("x", margin.left).attr("y", margin.top + statsY)}
\NormalTok{    .attr("fill", "\#1d4ed8").attr("font{-}size", "13px").attr("font{-}weight", "600")}
\NormalTok{    .text(\textasciigrave{}Approximate area ≈ $\{approxArea.toFixed(4)\}\textasciigrave{});}

\NormalTok{  svg.append("text")}
\NormalTok{    .attr("x", margin.left + 220).attr("y", margin.top + statsY)}
\NormalTok{    .attr("fill", "\#374151").attr("font{-}size", "13px")}
\NormalTok{    .text(\textasciigrave{}Exact area = $\{exactArea.toFixed(3)\}\textasciigrave{});}

\NormalTok{  svg.append("text")}
\NormalTok{    .attr("x", margin.left + 400).attr("y", margin.top + statsY)}
\NormalTok{    .attr("fill", error \textless{} 0.01 ? "\#059669" : "\#b45309").attr("font{-}size", "13px")}
\NormalTok{    .text(\textasciigrave{}|Error| = $\{error.toFixed(4)\}\textasciigrave{});}

\NormalTok{  return svg.node();}
\NormalTok{\}}
\end{Highlighting}
\end{Shaded}

\section{The Fundamental Theorem of
Calculus}\label{the-fundamental-theorem-of-calculus}

The Riemann sum definition is conceptually important but computationally
brutal --- you would need to take a limit every time you evaluate an
integral. The Fundamental Theorem of Calculus (FTC) is the result that
makes integration tractable. It has two parts.

\subsection{Part 1 --- Differentiation undoes
integration}\label{part-1-differentiation-undoes-integration}

Define a function \(F\) by accumulating area from a fixed point \(a\):

\[F(x) = \int_a^x f(t)\,dt\]

Part 1 says: \emph{differentiate \(F\) and you get \(f\) back.}

\[F'(x) = f(x)\]

In words: the rate at which accumulated area grows at \(x\) is exactly
\(f(x)\). If \(f(x)\) is large and positive, area accumulates quickly;
if \(f(x)\) is zero, the accumulated area is momentarily flat; if
\(f(x)\) is negative, the running total is decreasing. The integral of a
rate is an accumulation function, and its derivative returns the
original rate.

This is the theoretical bridge between differentiation and integration.
They are inverse operations --- like squaring and taking a square root
--- but with a subtlety: integration produces a \emph{family} of
antiderivatives differing by a constant. We'll see why immediately.

\subsection{Part 2 --- Evaluating definite
integrals}\label{part-2-evaluating-definite-integrals}

Part 2 is the computational workhorse. If \(F\) is any antiderivative of
\(f\) (meaning \(F'(x) = f(x)\)), then:

\[\int_a^b f(x)\,dx = F(b) - F(a)\]

The shorthand \(\bigl[F(x)\bigr]_a^b\) means \emph{evaluate \(F\) at the
top limit, subtract \(F\) at the bottom limit}:

\[\bigl[F(x)\bigr]_a^b = F(b) - F(a)\]

In words: to evaluate a definite integral, find any antiderivative of
the integrand, plug in the upper and lower limits, and subtract. You
never need to sum up rectangles explicitly.

\textbf{A quick example.} Evaluate \(\displaystyle\int_0^3 2x\,dx\).

An antiderivative of \(2x\) is \(x^2\) (because the derivative of
\(x^2\) is \(2x\)). Apply Part 2:

\[\int_0^3 2x\,dx = \bigl[x^2\bigr]_0^3 = 3^2 - 0^2 = 9\]

The area under the line \(y = 2x\) between \(x = 0\) and \(x = 3\) is
exactly \(9\).

\section{Integration rules}\label{integration-rules}

Building antiderivatives is the reverse of differentiation. Where
differentiation had rules for breaking functions apart, integration has
rules for putting them back together. These five cover most of
first-year calculus.

\textbf{1. Power rule (reversed)}

\[\int x^n\,dx = \frac{x^{n+1}}{n+1} + C \qquad (n \neq -1)\]

Check it: differentiate \(\dfrac{x^{n+1}}{n+1}\) using the power rule
and you recover \(x^n\).

\textbf{2. Exponential}

\[\int e^x\,dx = e^x + C\]

The exponential function is its own derivative, so it is also its own
antiderivative.

\textbf{3. Reciprocal}

\[\int \frac{1}{x}\,dx = \ln|x| + C\]

The absolute value handles negative \(x\). The derivative of \(\ln|x|\)
is \(1/x\) for all \(x \neq 0\).

\textbf{4. Sine}

\[\int \sin(x)\,dx = -\cos(x) + C\]

Check: the derivative of \(-\cos(x)\) is \(\sin(x)\).

\textbf{5. Cosine}

\[\int \cos(x)\,dx = \sin(x) + C\]

Check: the derivative of \(\sin(x)\) is \(\cos(x)\).

\begin{tcolorbox}[enhanced jigsaw, title=\textcolor{quarto-callout-note-color}{\faInfo}\hspace{0.5em}{Why the constant of integration?}, opacityback=0, colback=white, breakable, leftrule=.75mm, arc=.35mm, toptitle=1mm, left=2mm, colframe=quarto-callout-note-color-frame, coltitle=black, titlerule=0mm, colbacktitle=quarto-callout-note-color!10!white, opacitybacktitle=0.6, toprule=.15mm, bottomtitle=1mm, bottomrule=.15mm, rightrule=.15mm]

Every antiderivative of a function \(f\) differs from every other by a
constant. If \(F'(x) = f(x)\), then \((F(x) + 7)' = f(x)\) too. And
\((F(x) - 3)' = f(x)\). They all work.

Writing \(+C\) after an indefinite integral captures the entire
\emph{family} of valid antiderivatives at once. The \(+C\) is not
optional decoration --- omitting it loses real solutions. For example,
if you are solving a differential equation and need to apply an initial
condition, the value of \(C\) is determined by that condition. Dropping
\(C\) before you apply the condition makes the problem unsolvable.

For \emph{definite} integrals the constant cancels. When you compute
\(F(b) - F(a)\), the \(+C\) appears in both terms and vanishes:
\((F(b) + C) - (F(a) + C) = F(b) - F(a)\). The area between two limits
is uniquely defined.

\end{tcolorbox}

\section{Integration by substitution}\label{integration-by-substitution}

Substitution is the chain rule running in reverse. Recall that the chain
rule for differentiation says:

\[\frac{d}{dx}\bigl[f(g(x))\bigr] = f'(g(x)) \cdot g'(x)\]

If you see a product where one factor looks like the derivative of what
is inside a function, you can reverse this process.

\textbf{The method, step by step:}

\begin{enumerate}
\def\labelenumi{\arabic{enumi}.}
\tightlist
\item
  Identify a substitution \(u = g(x)\) --- usually the ``inside''
  function.
\item
  Differentiate to find \(du = g'(x)\,dx\), then solve for \(dx\).
\item
  Rewrite the entire integrand in terms of \(u\) and \(du\).
\item
  Integrate with respect to \(u\).
\item
  Substitute back to express the answer in terms of \(x\).
\end{enumerate}

\textbf{Worked example:} Evaluate \(\displaystyle\int 2x\cos(x^2)\,dx\).

\textbf{Step 1.} The inside function is \(x^2\). Let \(u = x^2\).

\textbf{Step 2.} Differentiate: \(\dfrac{du}{dx} = 2x\), so
\(du = 2x\,dx\).

\textbf{Step 3.} Notice that \(2x\,dx\) appears in the integrand.
Substitute:

\[\int 2x\cos(x^2)\,dx = \int \cos(u)\,du\]

\textbf{Step 4.} Integrate using the cosine rule:

\[\int \cos(u)\,du = \sin(u) + C\]

\textbf{Step 5.} Substitute back \(u = x^2\):

\[\int 2x\cos(x^2)\,dx = \sin(x^2) + C\]

\textbf{Check:} Differentiate \(\sin(x^2)\) using the chain rule:
\(\dfrac{d}{dx}\bigl[\sin(x^2)\bigr] = \cos(x^2) \cdot 2x = 2x\cos(x^2)\).
Correct.

The substitution worked because the factor \(2x\) outside the cosine was
exactly the derivative of \(x^2\) inside it. When that match exists,
substitution unravels the integral cleanly.

\textbf{Interactive: Area under a curve with sign.} Move the sliders to
set the limits \(a\) and \(b\). Positive area (above the axis) is shaded
blue; negative area (below the axis) is shaded red. The integral value
updates live.

\begin{Shaded}
\begin{Highlighting}[]
\NormalTok{//| echo: false}
\NormalTok{viewof sinA = Inputs.range([{-}Math.PI, 2 * Math.PI], \{}
\NormalTok{  label: "a (lower limit)",}
\NormalTok{  step: 0.05,}
\NormalTok{  value: 0}
\NormalTok{\})}
\end{Highlighting}
\end{Shaded}

\begin{Shaded}
\begin{Highlighting}[]
\NormalTok{//| echo: false}
\NormalTok{viewof sinB = Inputs.range([{-}Math.PI, 2 * Math.PI], \{}
\NormalTok{  label: "b (upper limit)",}
\NormalTok{  step: 0.05,}
\NormalTok{  value: Math.PI}
\NormalTok{\})}
\end{Highlighting}
\end{Shaded}

\begin{Shaded}
\begin{Highlighting}[]
\NormalTok{//| echo: false}
\NormalTok{\{}
\NormalTok{  const W = 560, H = 300;}
\NormalTok{  const margin = \{ top: 20, right: 20, bottom: 55, left: 50 \};}
\NormalTok{  const innerW = W {-} margin.left {-} margin.right;}
\NormalTok{  const innerH = H {-} margin.top {-} margin.bottom;}

\NormalTok{  const xMin = {-}Math.PI, xMax = 2 * Math.PI;}
\NormalTok{  const f = x =\textgreater{} Math.sin(x);}

\NormalTok{  // Clamp and order a, b}
\NormalTok{  const aRaw = Math.max(xMin, Math.min(xMax, sinA));}
\NormalTok{  const bRaw = Math.max(xMin, Math.min(xMax, sinB));}
\NormalTok{  const lo = Math.min(aRaw, bRaw);}
\NormalTok{  const hi = Math.max(aRaw, bRaw);}

\NormalTok{  // Exact integral: ∫sin(x)dx = {-}cos(x), so F(b) {-} F(a) = {-}cos(b) + cos(a)}
\NormalTok{  const integralValue = {-}Math.cos(bRaw) + Math.cos(aRaw);}

\NormalTok{  const xScale = d3.scaleLinear().domain([xMin, xMax]).range([0, innerW]);}
\NormalTok{  const yScale = d3.scaleLinear().domain([{-}1.4, 1.4]).range([innerH, 0]);}

\NormalTok{  const svg = d3.create("svg")}
\NormalTok{    .attr("viewBox", \textasciigrave{}0 0 $\{W\} $\{H\}\textasciigrave{})}
\NormalTok{    .attr("width", "100\%")}
\NormalTok{    .attr("style", \textasciigrave{}max{-}width:$\{W\}px; font{-}family:inherit;\textasciigrave{});}

\NormalTok{  const g = svg.append("g")}
\NormalTok{    .attr("transform", \textasciigrave{}translate($\{margin.left\},$\{margin.top\})\textasciigrave{});}

\NormalTok{  // Shaded area — split into positive and negative segments}
\NormalTok{  const steps = 600;}
\NormalTok{  const dx = (hi {-} lo) / steps;}

\NormalTok{  if (hi \textgreater{} lo) \{}
\NormalTok{    // Build positive{-}area path}
\NormalTok{    let posPath = "";}
\NormalTok{    let inPos = false;}
\NormalTok{    let posStart = lo;}

\NormalTok{    for (let i = 0; i \textless{}= steps; i++) \{}
\NormalTok{      const x = lo + i * dx;}
\NormalTok{      const y = f(x);}
\NormalTok{      if (y \textgreater{}= 0) \{}
\NormalTok{        if (!inPos) \{ posStart = x; inPos = true; \}}
\NormalTok{        if (i === steps || f(lo + (i+1)*dx) \textless{} 0) \{}
\NormalTok{          // close this positive segment}
\NormalTok{          const segPoints = d3.range(0, 1.001, 1/60).map(t =\textgreater{} \{}
\NormalTok{            const sx = posStart + t * (x {-} posStart);}
\NormalTok{            return [xScale(sx), yScale(f(sx))];}
\NormalTok{          \});}
\NormalTok{          const baseline = [[xScale(x), yScale(0)], [xScale(posStart), yScale(0)]];}
\NormalTok{          g.append("path")}
\NormalTok{            .attr("d", "M" + segPoints.map(p =\textgreater{} p.join(",")).join("L") + "L" + baseline.map(p =\textgreater{} p.join(",")).join("L") + "Z")}
\NormalTok{            .attr("fill", "\#3b82f6").attr("opacity", 0.35);}
\NormalTok{          inPos = false;}
\NormalTok{        \}}
\NormalTok{      \}}
\NormalTok{    \}}

\NormalTok{    // Build negative{-}area path segments}
\NormalTok{    inPos = false;}
\NormalTok{    let negStart = lo;}
\NormalTok{    for (let i = 0; i \textless{}= steps; i++) \{}
\NormalTok{      const x = lo + i * dx;}
\NormalTok{      const y = f(x);}
\NormalTok{      if (y \textless{} 0) \{}
\NormalTok{        if (!inPos) \{ negStart = x; inPos = true; \}}
\NormalTok{        if (i === steps || f(lo + (i+1)*dx) \textgreater{}= 0) \{}
\NormalTok{          const segPoints = d3.range(0, 1.001, 1/60).map(t =\textgreater{} \{}
\NormalTok{            const sx = negStart + t * (x {-} negStart);}
\NormalTok{            return [xScale(sx), yScale(f(sx))];}
\NormalTok{          \});}
\NormalTok{          const baseline = [[xScale(x), yScale(0)], [xScale(negStart), yScale(0)]];}
\NormalTok{          g.append("path")}
\NormalTok{            .attr("d", "M" + segPoints.map(p =\textgreater{} p.join(",")).join("L") + "L" + baseline.map(p =\textgreater{} p.join(",")).join("L") + "Z")}
\NormalTok{            .attr("fill", "\#ef4444").attr("opacity", 0.35);}
\NormalTok{          inPos = false;}
\NormalTok{        \}}
\NormalTok{      \}}
\NormalTok{    \}}
\NormalTok{  \}}

\NormalTok{  // X{-}axis line}
\NormalTok{  g.append("line")}
\NormalTok{    .attr("x1", 0).attr("y1", yScale(0))}
\NormalTok{    .attr("x2", innerW).attr("y2", yScale(0))}
\NormalTok{    .attr("stroke", "\#374151").attr("stroke{-}width", 1.5);}

\NormalTok{  // Y{-}axis line}
\NormalTok{  g.append("line")}
\NormalTok{    .attr("x1", xScale(0)).attr("y1", 0)}
\NormalTok{    .attr("x2", xScale(0)).attr("y2", innerH)}
\NormalTok{    .attr("stroke", "\#9ca3af").attr("stroke{-}width", 1);}

\NormalTok{  // Draw curve}
\NormalTok{  const curvePoints = d3.range(xMin, xMax + 0.01, 0.02).map(x =\textgreater{} [xScale(x), yScale(f(x))]);}
\NormalTok{  g.append("path")}
\NormalTok{    .attr("d", d3.line()(curvePoints))}
\NormalTok{    .attr("fill", "none")}
\NormalTok{    .attr("stroke", "\#1e40af")}
\NormalTok{    .attr("stroke{-}width", 2.5);}

\NormalTok{  // Limit markers}
\NormalTok{  [aRaw, bRaw].forEach((v, idx) =\textgreater{} \{}
\NormalTok{    g.append("line")}
\NormalTok{      .attr("x1", xScale(v)).attr("y1", 0)}
\NormalTok{      .attr("x2", xScale(v)).attr("y2", innerH)}
\NormalTok{      .attr("stroke", "\#6b7280").attr("stroke{-}width", 1)}
\NormalTok{      .attr("stroke{-}dasharray", "4,3");}
\NormalTok{    g.append("text")}
\NormalTok{      .attr("x", xScale(v)).attr("y", {-}6)}
\NormalTok{      .attr("text{-}anchor", "middle")}
\NormalTok{      .attr("fill", "\#374151").attr("font{-}size", "11px").attr("font{-}weight", "600")}
\NormalTok{      .text(idx === 0 ? "a" : "b");}
\NormalTok{  \});}

\NormalTok{  // X{-}axis tick labels}
\NormalTok{  const ticks = [}
\NormalTok{    \{ v: {-}Math.PI, label: "−π" \},}
\NormalTok{    \{ v: 0,        label: "0" \},}
\NormalTok{    \{ v: Math.PI,  label: "π" \},}
\NormalTok{    \{ v: 2 * Math.PI, label: "2π" \}}
\NormalTok{  ];}
\NormalTok{  ticks.forEach((\{ v, label \}) =\textgreater{} \{}
\NormalTok{    g.append("line")}
\NormalTok{      .attr("x1", xScale(v)).attr("y1", yScale(0))}
\NormalTok{      .attr("x2", xScale(v)).attr("y2", yScale(0) + 5)}
\NormalTok{      .attr("stroke", "\#374151");}
\NormalTok{    g.append("text")}
\NormalTok{      .attr("x", xScale(v)).attr("y", yScale(0) + 18)}
\NormalTok{      .attr("text{-}anchor", "middle")}
\NormalTok{      .attr("fill", "\#374151").attr("font{-}size", "12px")}
\NormalTok{      .text(label);}
\NormalTok{  \});}

\NormalTok{  // Integral value display}
\NormalTok{  const absVal = Math.abs(integralValue);}
\NormalTok{  const sign = integralValue \textgreater{} 0.0005 ? "positive" : integralValue \textless{} {-}0.0005 ? "negative" : "zero";}
\NormalTok{  const signColour = sign === "positive" ? "\#2563eb" : sign === "negative" ? "\#dc2626" : "\#059669";}

\NormalTok{  svg.append("text")}
\NormalTok{    .attr("x", margin.left).attr("y", margin.top + innerH + 46)}
\NormalTok{    .attr("fill", signColour).attr("font{-}size", "13px").attr("font{-}weight", "600")}
\NormalTok{    .text(\textasciigrave{}∫ sin(x) dx from a to b  =  $\{integralValue.toFixed(4)\}   ($\{sign\})\textasciigrave{});}

\NormalTok{  return svg.node();}
\NormalTok{\}}
\end{Highlighting}
\end{Shaded}

\section{Worked examples}\label{worked-examples-19}

\subsection{Example 1 --- Power rule and the FTC (computing /
probability)}\label{example-1-power-rule-and-the-ftc-computing-probability}

\textbf{Problem.} Find the indefinite integral
\(\displaystyle\int (3x^2 - 4x + 1)\,dx\), then evaluate the definite
integral from \(x = 0\) to \(x = 2\).

A probability density function for a random variable on \([0, 2]\) might
have this polynomial form; the definite integral would give the total
probability mass over that interval.

\textbf{Indefinite integral.} Apply the power rule term by term:

\[\int (3x^2 - 4x + 1)\,dx = \frac{3x^3}{3} - \frac{4x^2}{2} + x + C = x^3 - 2x^2 + x + C\]

\textbf{Definite integral.} Apply FTC Part 2 with
\(F(x) = x^3 - 2x^2 + x\):

\[\int_0^2 (3x^2 - 4x + 1)\,dx = \bigl[x^3 - 2x^2 + x\bigr]_0^2\]

Evaluate at the upper limit:

\[F(2) = 2^3 - 2(2^2) + 2 = 8 - 8 + 2 = 2\]

Evaluate at the lower limit:

\[F(0) = 0 - 0 + 0 = 0\]

The definite integral equals \(F(2) - F(0) = 2 - 0 = 2\).

\begin{center}\rule{0.5\linewidth}{0.5pt}\end{center}

\subsection{Example 2 --- FTC Part 2 (science /
measurement)}\label{example-2-ftc-part-2-science-measurement}

\textbf{Problem.} A radioactive tracer is administered to a patient and
its concentration in the bloodstream decays in a way modelled by
\(c(x) = x^{-1/2}\) (in appropriate units, for \(x > 0\)). Compute the
total exposure --- the area under the concentration curve --- between
\(x = 1\) and \(x = 4\).

\[\int_1^4 \frac{1}{\sqrt{x}}\,dx = \int_1^4 x^{-1/2}\,dx\]

Apply the power rule with \(n = -\tfrac{1}{2}\):

\[\int x^{-1/2}\,dx = \frac{x^{-1/2 + 1}}{-1/2 + 1} = \frac{x^{1/2}}{1/2} = 2\sqrt{x} + C\]

Apply the limits:

\[\bigl[2\sqrt{x}\bigr]_1^4 = 2\sqrt{4} - 2\sqrt{1} = 2(2) - 2(1) = 4 - 2 = 2\]

Total exposure over the interval is \(2\) (in the model's units).

\begin{center}\rule{0.5\linewidth}{0.5pt}\end{center}

\subsection{Example 3 --- Substitution (engineering /
signals)}\label{example-3-substitution-engineering-signals}

\textbf{Problem.} In signal processing, an engineer needs to evaluate
\(\displaystyle\int x\,e^{x^2}\,dx\) as part of computing the energy in
a modulated waveform.

Let \(u = x^2\), so \(du = 2x\,dx\), which gives
\(x\,dx = \dfrac{du}{2}\).

\[\int x\,e^{x^2}\,dx = \int e^u \cdot \frac{du}{2} = \frac{1}{2}\int e^u\,du = \frac{1}{2}e^u + C\]

Substitute back \(u = x^2\):

\[\int x\,e^{x^2}\,dx = \frac{1}{2}e^{x^2} + C\]

\textbf{Check:} Differentiate \(\tfrac{1}{2}e^{x^2}\) with the chain
rule:
\(\dfrac{d}{dx}\bigl[\tfrac{1}{2}e^{x^2}\bigr] = \tfrac{1}{2} \cdot e^{x^2} \cdot 2x = x\,e^{x^2}\).
Correct.

\begin{center}\rule{0.5\linewidth}{0.5pt}\end{center}

\subsection{Example 3b --- Area between two
curves}\label{example-3b-area-between-two-curves}

\textbf{Problem.} Find the area enclosed between \(f(x) = x + 3\) and
\(g(x) = x^2 + 1\).

\textbf{Step 1 --- Find intersections:}

\[x + 3 = x^2 + 1 \implies x^2 - x - 2 = 0 \implies (x-2)(x+1) = 0\]

Intersections at \(x = -1\) and \(x = 2\).

\textbf{Step 2 --- Identify which is on top:} At \(x = 0\), \(f(0) = 3\)
and \(g(0) = 1\), so \(f > g\) on \([-1, 2]\).

\textbf{Step 3 --- Integrate:}

\[\int_{-1}^{2} \bigl[(x+3) - (x^2+1)\bigr]\,dx = \int_{-1}^{2} (2 + x - x^2)\,dx\]

\[= \left[2x + \frac{x^2}{2} - \frac{x^3}{3}\right]_{-1}^{2}\]

\[= \left(4 + 2 - \frac{8}{3}\right) - \left(-2 + \frac{1}{2} + \frac{1}{3}\right)\]

\[= \frac{10}{3} - \left(-\frac{7}{6}\right) = \frac{20}{6} + \frac{7}{6} = \frac{27}{6} = \frac{9}{2}\]

The enclosed area is \(\dfrac{9}{2}\).

\begin{center}\rule{0.5\linewidth}{0.5pt}\end{center}

\subsection{Example 4 --- Area between two curves (business / economic
surplus)}\label{example-4-area-between-two-curves-business-economic-surplus}

\textbf{Problem.} In a competitive market, the demand curve is
\(p_D(q) = 8 - q^2\) and the supply curve is \(p_S(q) = 2q\). Find the
consumer surplus --- the area between the demand curve and the
equilibrium price, from zero to the equilibrium quantity.

\textbf{Step 1. Find the equilibrium.} Set \(p_D = p_S\):

\[8 - q^2 = 2q \implies q^2 + 2q - 8 = 0 \implies (q+4)(q-2) = 0\]

Equilibrium is at \(q^* = 2\) (taking the positive root), equilibrium
price \(p^* = 4\).

\textbf{Step 2. Consumer surplus} is the area between the demand curve
and the equilibrium price from \(q = 0\) to \(q = 2\):

\[\text{CS} = \int_0^2 \bigl[(8 - q^2) - 4\bigr]\,dq = \int_0^2 (4 - q^2)\,dq\]

\textbf{Step 3. Evaluate:}

\[\int_0^2 (4 - q^2)\,dq = \bigl[4q - \tfrac{q^3}{3}\bigr]_0^2 = \left(8 - \frac{8}{3}\right) - 0 = \frac{16}{3} \approx 5.33\]

Consumer surplus is \(\tfrac{16}{3}\) (in the model's currency units per
unit quantity). This is the aggregate benefit consumers receive by
paying the equilibrium price rather than their maximum willingness to
pay.

\section{Where this goes}\label{where-this-goes-20}

\textbf{Multivariable calculus (Vol 7).} A single integral accumulates a
function over an interval --- a one-dimensional slice. In two and three
dimensions you can integrate over areas, volumes, and surfaces --- the
same idea of accumulation, applied to more variables. Those tools (Vol
7) are the language of electromagnetism, fluid dynamics, and heat flow.

\textbf{Ordinary differential equations (Vol 7).} A differential
equation relates a function to its own derivatives. The simplest case
--- \emph{separable} equations --- is solved by moving all the
\(y\)-terms to one side and all the \(x\)-terms to the other, then
integrating both sides. Every technique in this chapter applies
directly. The ability to integrate is the ability to solve differential
equations.

\section{Applications}\label{applications}

\begin{itemize}
\item
  \textbf{Distance from velocity (physics).} Integrate \(v(t)\) over a
  time interval to recover displacement. If \(v(t) = 9.8t\) (free fall
  from rest), then \(\int_0^T 9.8t\,dt = 4.9T^2\) --- the familiar
  kinematic formula, derived by integration.
\item
  \textbf{Work by a variable force (mechanics).} Work is
  \(W = \int_a^b F(x)\,dx\) when force \(F\) varies with position \(x\).
  Compressing a spring, pumping water from a tank --- both are definite
  integrals.
\item
  \textbf{Consumer and producer surplus (economics).} The area between a
  demand (or supply) curve and the equilibrium price is a definite
  integral. It measures welfare.
\item
  \textbf{Probability from a density function (statistics).} If \(f(x)\)
  is a probability density function, then
  \(P(a \leq X \leq b) = \int_a^b f(x)\,dx\). The cumulative
  distribution function is \(F(x) = \int_{-\infty}^x f(t)\,dt\) ---
  literally the FTC in action.
\item
  \textbf{Average value of a function.} The average value of \(f\) over
  \([a,b]\) is \(\dfrac{1}{b-a}\int_a^b f(x)\,dx\) --- the integral
  normalised by the length of the interval.
\end{itemize}

\begin{center}\rule{0.5\linewidth}{0.5pt}\end{center}

\section{Exercises}\label{exercises-20}

\begin{Shaded}
\begin{Highlighting}[]
\NormalTok{//| echo: false}
\NormalTok{katex = require("katex@0.16")}
\end{Highlighting}
\end{Shaded}

\begin{Shaded}
\begin{Highlighting}[]
\NormalTok{//| echo: false}
\NormalTok{function makeStepperHTML(exerciseNum, steps) \{}
\NormalTok{  let current = 0;}
\NormalTok{  const totalSteps = steps.length;}
\NormalTok{  const container = document.createElement("div");}
\NormalTok{  container.style.cssText = "border:1px solid \#e5e7eb; border{-}radius:8px; padding:1rem 1.25rem; margin:0.75rem 0; font{-}family:inherit;";}
\NormalTok{  const stepsDiv = document.createElement("div");}
\NormalTok{  stepsDiv.style.marginBottom = "0.75rem";}

\NormalTok{  function renderSteps() \{}
\NormalTok{    stepsDiv.innerHTML = "";}
\NormalTok{    for (let i = 0; i \textless{} current; i++) \{}
\NormalTok{      const s = steps[i];}
\NormalTok{      const row = document.createElement("div");}
\NormalTok{      row.style.cssText = "display:grid; grid{-}template{-}columns:200px 1fr; gap:0.35rem 1rem; align{-}items:baseline; padding:0.35rem 0; border{-}top:1px solid \#f3f4f6;";}
\NormalTok{      const opCell = document.createElement("span");}
\NormalTok{      opCell.style.cssText = "font{-}size:0.85em; color:\#6b7280; font{-}style:italic;";}
\NormalTok{      opCell.textContent = s.op;}
\NormalTok{      const eqCell = document.createElement("span");}
\NormalTok{      eqCell.innerHTML = katex.renderToString(s.eq, \{ throwOnError: false, displayMode: false \});}
\NormalTok{      row.appendChild(opCell);}
\NormalTok{      row.appendChild(eqCell);}
\NormalTok{      if (s.note) \{}
\NormalTok{        const noteRow = document.createElement("div");}
\NormalTok{        noteRow.style.cssText = "grid{-}column:2; font{-}size:0.82em; color:\#6b7280; padding{-}bottom:0.2rem;";}
\NormalTok{        noteRow.textContent = s.note;}
\NormalTok{        row.appendChild(noteRow);}
\NormalTok{      \}}
\NormalTok{      stepsDiv.appendChild(row);}
\NormalTok{    \}}
\NormalTok{    if (current === totalSteps) \{}
\NormalTok{      const done = document.createElement("div");}
\NormalTok{      done.style.cssText = "margin{-}top:0.5rem; font{-}size:0.9em; color:\#059669; font{-}weight:500;";}
\NormalTok{      done.textContent = "✓ Solution complete";}
\NormalTok{      stepsDiv.appendChild(done);}
\NormalTok{    \}}
\NormalTok{  \}}

\NormalTok{  const controls = document.createElement("div");}
\NormalTok{  controls.style.cssText = "display:flex; gap:0.5rem; align{-}items:center;";}
\NormalTok{  const nextBtn = document.createElement("button");}
\NormalTok{  nextBtn.textContent = "Next step →";}
\NormalTok{  nextBtn.style.cssText = "padding:0.35rem 0.85rem; border:1px solid \#d1d5db; border{-}radius:4px; background:\#fff; cursor:pointer; font{-}size:0.9em;";}
\NormalTok{  const resetBtn = document.createElement("button");}
\NormalTok{  resetBtn.textContent = "Reset";}
\NormalTok{  resetBtn.style.cssText = "padding:0.35rem 0.75rem; border:1px solid \#e5e7eb; border{-}radius:4px; background:\#f9fafb; cursor:pointer; font{-}size:0.9em; color:\#6b7280;";}
\NormalTok{  const counter = document.createElement("span");}
\NormalTok{  counter.style.cssText = "font{-}size:0.82em; color:\#9ca3af; margin{-}left:0.25rem;";}

\NormalTok{  function updateButtons() \{}
\NormalTok{    nextBtn.disabled = current === totalSteps;}
\NormalTok{    nextBtn.style.opacity = current === totalSteps ? "0.4" : "1";}
\NormalTok{    counter.textContent = current === 0 ? "Click to reveal steps" : \textasciigrave{}Step $\{current\} of $\{totalSteps\}\textasciigrave{};}
\NormalTok{  \}}

\NormalTok{  nextBtn.onclick = () =\textgreater{} \{ if (current \textless{} totalSteps) \{ current++; renderSteps(); updateButtons(); \} \};}
\NormalTok{  resetBtn.onclick = () =\textgreater{} \{ current = 0; renderSteps(); updateButtons(); \};}
\NormalTok{  controls.appendChild(nextBtn);}
\NormalTok{  controls.appendChild(resetBtn);}
\NormalTok{  controls.appendChild(counter);}
\NormalTok{  container.appendChild(stepsDiv);}
\NormalTok{  container.appendChild(controls);}
\NormalTok{  renderSteps();}
\NormalTok{  updateButtons();}
\NormalTok{  return container;}
\NormalTok{\}}
\end{Highlighting}
\end{Shaded}

\textbf{Exercise 1.} Find the indefinite integral:

\[\int \left(5x^4 - 3x^2 + 2\right)dx\]

Include the constant of integration in your answer.

\begin{Shaded}
\begin{Highlighting}[]
\NormalTok{//| echo: false}
\NormalTok{makeStepperHTML(1, [}
\NormalTok{  \{ op: "Apply power rule term by term",       eq: "\textbackslash{}\textbackslash{}int 5x\^{}4\textbackslash{}\textbackslash{},dx {-} \textbackslash{}\textbackslash{}int 3x\^{}2\textbackslash{}\textbackslash{},dx + \textbackslash{}\textbackslash{}int 2\textbackslash{}\textbackslash{},dx", note: "Integration is linear: split across addition and pull out constants." \},}
\NormalTok{  \{ op: "Integrate each term",                 eq: "\textbackslash{}\textbackslash{}frac\{5x\^{}5\}\{5\} {-} \textbackslash{}\textbackslash{}frac\{3x\^{}3\}\{3\} + 2x + C" \},}
\NormalTok{  \{ op: "Simplify coefficients",               eq: "x\^{}5 {-} x\^{}3 + 2x + C" \},}
\NormalTok{  \{ op: "Check by differentiating",            eq: "\textbackslash{}\textbackslash{}frac\{d\}\{dx\}\textbackslash{}\textbackslash{}left(x\^{}5 {-} x\^{}3 + 2x + C\textbackslash{}\textbackslash{}right) = 5x\^{}4 {-} 3x\^{}2 + 2 \textbackslash{}\textbackslash{}checkmark" \}}
\NormalTok{])}
\end{Highlighting}
\end{Shaded}

\begin{center}\rule{0.5\linewidth}{0.5pt}\end{center}

\textbf{Exercise 2.} Evaluate the definite integral using the
Fundamental Theorem of Calculus Part 2:

\[\int_1^9 \frac{3}{\sqrt{x}}\,dx\]

\begin{Shaded}
\begin{Highlighting}[]
\NormalTok{//| echo: false}
\NormalTok{makeStepperHTML(2, [}
\NormalTok{  \{ op: "Rewrite as a power",                  eq: "\textbackslash{}\textbackslash{}int\_1\^{}9 3x\^{}\{{-}1/2\}\textbackslash{}\textbackslash{},dx", note: "1/√x = x\^{}\{{-}1/2\} — apply the power rule." \},}
\NormalTok{  \{ op: "Find antiderivative",                 eq: "3 \textbackslash{}\textbackslash{}cdot \textbackslash{}\textbackslash{}frac\{x\^{}\{1/2\}\}\{1/2\} = 6x\^{}\{1/2\} = 6\textbackslash{}\textbackslash{}sqrt\{x\}" \},}
\NormalTok{  \{ op: "Apply FTC Part 2",                    eq: "\textbackslash{}\textbackslash{}left[6\textbackslash{}\textbackslash{}sqrt\{x\}\textbackslash{}\textbackslash{}right]\_1\^{}9 = 6\textbackslash{}\textbackslash{}sqrt\{9\} {-} 6\textbackslash{}\textbackslash{}sqrt\{1\}" \},}
\NormalTok{  \{ op: "Evaluate at limits",                  eq: "6(3) {-} 6(1) = 18 {-} 6 = 12" \},}
\NormalTok{  \{ op: "Answer",                              eq: "\textbackslash{}\textbackslash{}int\_1\^{}9 \textbackslash{}\textbackslash{}frac\{3\}\{\textbackslash{}\textbackslash{}sqrt\{x\}\}\textbackslash{}\textbackslash{},dx = 12" \}}
\NormalTok{])}
\end{Highlighting}
\end{Shaded}

\begin{center}\rule{0.5\linewidth}{0.5pt}\end{center}

\textbf{Exercise 3.} Use \(u\)-substitution to evaluate:

\[\int 3x^2\sin(x^3)\,dx\]

\begin{Shaded}
\begin{Highlighting}[]
\NormalTok{//| echo: false}
\NormalTok{makeStepperHTML(3, [}
\NormalTok{  \{ op: "Choose substitution",                 eq: "u = x\^{}3", note: "The inside function of sin is x³. Its derivative 3x² appears outside." \},}
\NormalTok{  \{ op: "Differentiate",                       eq: "\textbackslash{}\textbackslash{}frac\{du\}\{dx\} = 3x\^{}2 \textbackslash{}\textbackslash{}implies du = 3x\^{}2\textbackslash{}\textbackslash{},dx" \},}
\NormalTok{  \{ op: "Substitute into integral",            eq: "\textbackslash{}\textbackslash{}int \textbackslash{}\textbackslash{}sin(u)\textbackslash{}\textbackslash{},du" \},}
\NormalTok{  \{ op: "Integrate",                           eq: "{-}\textbackslash{}\textbackslash{}cos(u) + C" \},}
\NormalTok{  \{ op: "Substitute back",                     eq: "{-}\textbackslash{}\textbackslash{}cos(x\^{}3) + C" \},}
\NormalTok{  \{ op: "Check by differentiating",            eq: "\textbackslash{}\textbackslash{}frac\{d\}\{dx\}\textbackslash{}\textbackslash{}left[{-}\textbackslash{}\textbackslash{}cos(x\^{}3)\textbackslash{}\textbackslash{}right] = \textbackslash{}\textbackslash{}sin(x\^{}3) \textbackslash{}\textbackslash{}cdot 3x\^{}2 = 3x\^{}2\textbackslash{}\textbackslash{}sin(x\^{}3) \textbackslash{}\textbackslash{}checkmark" \}}
\NormalTok{])}
\end{Highlighting}
\end{Shaded}

\begin{center}\rule{0.5\linewidth}{0.5pt}\end{center}

\textbf{Exercise 4.} Find the area under the curve \(f(x) = x^2 + 2x\)
between \(x = 0\) and \(x = 3\). Sketch the region (or describe it in
words) before computing.

\begin{Shaded}
\begin{Highlighting}[]
\NormalTok{//| echo: false}
\NormalTok{makeStepperHTML(4, [}
\NormalTok{  \{ op: "Set up the definite integral",        eq: "\textbackslash{}\textbackslash{}int\_0\^{}3 (x\^{}2 + 2x)\textbackslash{}\textbackslash{},dx", note: "f(x) ≥ 0 on [0,3], so the integral gives the geometric area." \},}
\NormalTok{  \{ op: "Find the antiderivative",             eq: "F(x) = \textbackslash{}\textbackslash{}frac\{x\^{}3\}\{3\} + x\^{}2" \},}
\NormalTok{  \{ op: "Apply FTC Part 2",                    eq: "\textbackslash{}\textbackslash{}left[\textbackslash{}\textbackslash{}frac\{x\^{}3\}\{3\} + x\^{}2\textbackslash{}\textbackslash{}right]\_0\^{}3 = F(3) {-} F(0)" \},}
\NormalTok{  \{ op: "Evaluate at upper limit",             eq: "F(3) = \textbackslash{}\textbackslash{}frac\{27\}\{3\} + 9 = 9 + 9 = 18" \},}
\NormalTok{  \{ op: "Evaluate at lower limit",             eq: "F(0) = 0 + 0 = 0" \},}
\NormalTok{  \{ op: "Compute the area",                    eq: "18 {-} 0 = 18" \}}
\NormalTok{])}
\end{Highlighting}
\end{Shaded}

\begin{center}\rule{0.5\linewidth}{0.5pt}\end{center}

\textbf{Exercise 5.} Find the area enclosed between \(f(x) = 4 - x^2\)
and \(g(x) = x + 2\).

\emph{Hint: find the intersection points first. The area between two
curves is \(\int_a^b \bigl[f(x) - g(x)\bigr]dx\) when \(f(x) \geq g(x)\)
on \([a, b]\).}

\begin{Shaded}
\begin{Highlighting}[]
\NormalTok{//| echo: false}
\NormalTok{makeStepperHTML(5, [}
\NormalTok{  \{ op: "Find intersections",                  eq: "4 {-} x\^{}2 = x + 2 \textbackslash{}\textbackslash{}implies x\^{}2 + x {-} 2 = 0 \textbackslash{}\textbackslash{}implies (x+2)(x{-}1) = 0", note: "Intersections at x = −2 and x = 1." \},}
\NormalTok{  \{ op: "Check which is on top",               eq: "\textbackslash{}\textbackslash{}text\{At \} x=0:\textbackslash{}\textbackslash{}; f(0)=4,\textbackslash{}\textbackslash{}; g(0)=2 \textbackslash{}\textbackslash{}implies f \textgreater{} g \textbackslash{}\textbackslash{}text\{ on \} [{-}2,1]" \},}
\NormalTok{  \{ op: "Set up integral of difference",       eq: "\textbackslash{}\textbackslash{}int\_\{{-}2\}\^{}\{1\} \textbackslash{}\textbackslash{}bigl[(4{-}x\^{}2) {-} (x+2)\textbackslash{}\textbackslash{}bigr]\textbackslash{}\textbackslash{},dx = \textbackslash{}\textbackslash{}int\_\{{-}2\}\^{}\{1\} (2 {-} x {-} x\^{}2)\textbackslash{}\textbackslash{},dx" \},}
\NormalTok{  \{ op: "Find antiderivative",                 eq: "F(x) = 2x {-} \textbackslash{}\textbackslash{}frac\{x\^{}2\}\{2\} {-} \textbackslash{}\textbackslash{}frac\{x\^{}3\}\{3\}" \},}
\NormalTok{  \{ op: "Evaluate at upper limit",             eq: "F(1) = 2 {-} \textbackslash{}\textbackslash{}tfrac\{1\}\{2\} {-} \textbackslash{}\textbackslash{}tfrac\{1\}\{3\} = \textbackslash{}\textbackslash{}tfrac\{12{-}3{-}2\}\{6\} = \textbackslash{}\textbackslash{}tfrac\{7\}\{6\}" \},}
\NormalTok{  \{ op: "Evaluate at lower limit",             eq: "F({-}2) = {-}4 {-} 2 + \textbackslash{}\textbackslash{}tfrac\{8\}\{3\} = \textbackslash{}\textbackslash{}tfrac\{{-}18+8\}\{3\} = {-}\textbackslash{}\textbackslash{}tfrac\{10\}\{3\}" \},}
\NormalTok{  \{ op: "Compute the area",                    eq: "\textbackslash{}\textbackslash{}tfrac\{7\}\{6\} {-} \textbackslash{}\textbackslash{}left({-}\textbackslash{}\textbackslash{}tfrac\{10\}\{3\}\textbackslash{}\textbackslash{}right) = \textbackslash{}\textbackslash{}tfrac\{7\}\{6\} + \textbackslash{}\textbackslash{}tfrac\{20\}\{6\} = \textbackslash{}\textbackslash{}tfrac\{27\}\{6\} = \textbackslash{}\textbackslash{}tfrac\{9\}\{2\}" \}}
\NormalTok{])}
\end{Highlighting}
\end{Shaded}

\begin{center}\rule{0.5\linewidth}{0.5pt}\end{center}

\textbf{Exercise 6 --- Accumulation in context.} A solar panel generates
electricity at a rate of
\(P(t) = 4\sin\!\left(\dfrac{\pi t}{12}\right)\) kilowatts, where \(t\)
is hours after sunrise (\(0 \leq t \leq 12\)).

\begin{enumerate}
\def\labelenumi{(\alph{enumi})}
\item
  Find the total energy (in kilowatt-hours) generated over the full
  12-hour day: \(\displaystyle\int_0^{12} P(t)\,dt\).
\item
  Interpret the result: what does the definite integral represent
  physically?
\end{enumerate}

\emph{(Your answer will be in terms of \(\pi\) --- that is expected.)}

\begin{Shaded}
\begin{Highlighting}[]
\NormalTok{//| echo: false}
\NormalTok{makeStepperHTML(6, [}
\NormalTok{  \{ op: "Set up the integral",                 eq: "\textbackslash{}\textbackslash{}int\_0\^{}\{12\} 4\textbackslash{}\textbackslash{}sin\textbackslash{}\textbackslash{}!\textbackslash{}\textbackslash{}left(\textbackslash{}\textbackslash{}frac\{\textbackslash{}\textbackslash{}pi t\}\{12\}\textbackslash{}\textbackslash{}right)dt" \},}
\NormalTok{  \{ op: "Choose substitution",                 eq: "u = \textbackslash{}\textbackslash{}frac\{\textbackslash{}\textbackslash{}pi t\}\{12\} \textbackslash{}\textbackslash{}implies du = \textbackslash{}\textbackslash{}frac\{\textbackslash{}\textbackslash{}pi\}\{12\}\textbackslash{}\textbackslash{},dt \textbackslash{}\textbackslash{}implies dt = \textbackslash{}\textbackslash{}frac\{12\}\{\textbackslash{}\textbackslash{}pi\}\textbackslash{}\textbackslash{},du", note: "u{-}substitution clears the coefficient inside sin." \},}
\NormalTok{  \{ op: "Change limits",                       eq: "t=0 \textbackslash{}\textbackslash{}Rightarrow u=0;\textbackslash{}\textbackslash{}quad t=12 \textbackslash{}\textbackslash{}Rightarrow u=\textbackslash{}\textbackslash{}pi" \},}
\NormalTok{  \{ op: "Substitute and simplify",             eq: "4 \textbackslash{}\textbackslash{}cdot \textbackslash{}\textbackslash{}frac\{12\}\{\textbackslash{}\textbackslash{}pi\} \textbackslash{}\textbackslash{}int\_0\^{}\{\textbackslash{}\textbackslash{}pi\} \textbackslash{}\textbackslash{}sin(u)\textbackslash{}\textbackslash{},du = \textbackslash{}\textbackslash{}frac\{48\}\{\textbackslash{}\textbackslash{}pi\}\textbackslash{}\textbackslash{}int\_0\^{}\{\textbackslash{}\textbackslash{}pi\} \textbackslash{}\textbackslash{}sin(u)\textbackslash{}\textbackslash{},du" \},}
\NormalTok{  \{ op: "Integrate sin",                       eq: "\textbackslash{}\textbackslash{}frac\{48\}\{\textbackslash{}\textbackslash{}pi\}\textbackslash{}\textbackslash{}left[{-}\textbackslash{}\textbackslash{}cos(u)\textbackslash{}\textbackslash{}right]\_0\^{}\{\textbackslash{}\textbackslash{}pi\}" \},}
\NormalTok{  \{ op: "Evaluate at limits",                  eq: "\textbackslash{}\textbackslash{}frac\{48\}\{\textbackslash{}\textbackslash{}pi\}\textbackslash{}\textbackslash{}left({-}\textbackslash{}\textbackslash{}cos(\textbackslash{}\textbackslash{}pi) + \textbackslash{}\textbackslash{}cos(0)\textbackslash{}\textbackslash{}right) = \textbackslash{}\textbackslash{}frac\{48\}\{\textbackslash{}\textbackslash{}pi\}(1+1) = \textbackslash{}\textbackslash{}frac\{96\}\{\textbackslash{}\textbackslash{}pi\}" \},}
\NormalTok{  \{ op: "Interpret the result",                eq: "\textbackslash{}\textbackslash{}text\{Total energy\} = \textbackslash{}\textbackslash{}frac\{96\}\{\textbackslash{}\textbackslash{}pi\} \textbackslash{}\textbackslash{}approx 30.6 \textbackslash{}\textbackslash{}text\{ kWh\}", note: "The definite integral sums the instantaneous power output over the full day — the accumulated energy." \}}
\NormalTok{])}
\end{Highlighting}
\end{Shaded}

\chapter{Proof and mathematical
reasoning}\label{proof-and-mathematical-reasoning}

Establishing that something must be true

\hfill\break

You have used the formula for the sum of an arithmetic series. You have
used the quadratic formula. You have been told that \(\sqrt{2}\) is
irrational. Each time, you accepted the result because it worked on the
examples you tried, or because someone authoritative said so.

That is about to change.

The shift from computing to proving is the sharpest turn in mathematics.
It is the difference between ``this has always worked so far'' and
``this cannot fail.'' Most of the mathematics you have seen up to now
has been the first kind. This chapter introduces the second.

\section{A pattern that lies to you}\label{a-pattern-that-lies-to-you}

Consider the expression \(n^2 - n + 41\). Try a few values:

\begin{longtable}[]{@{}lll@{}}
\toprule\noalign{}
\(n\) & \(n^2 - n + 41\) & Prime? \\
\midrule\noalign{}
\endhead
\bottomrule\noalign{}
\endlastfoot
1 & 41 & Yes \\
2 & 43 & Yes \\
3 & 47 & Yes \\
4 & 53 & Yes \\
10 & 131 & Yes \\
20 & 421 & Yes \\
\end{longtable}

In fact, every value from \(n = 1\) to \(n = 40\) gives a prime number.

Forty consecutive confirmations. Most people would bet on ``always
prime.''

But try \(n = 41\):

\[41^2 - 41 + 41 = 41^2 = 1681 = 41 \times 41\]

Clearly composite. The pattern fails the moment you step past the range
where you checked it.

This is the central lesson before we write a single proof: \textbf{no
number of confirmed cases constitutes a proof.} Not ten, not a hundred,
not a trillion. A proof must cover every case, including those that
haven't been tried.

\begin{Shaded}
\begin{Highlighting}[]
\NormalTok{//| echo: false}

\NormalTok{// ── VIZ 1: Primality checker for n²{-}n+41 ──────────────────────────────────}

\NormalTok{function isPrime(n) \{}
\NormalTok{  if (n \textless{} 2) return false;}
\NormalTok{  for (let i = 2; i \textless{}= Math.sqrt(n); i++) \{}
\NormalTok{    if (n \% i === 0) return false;}
\NormalTok{  \}}
\NormalTok{  return true;}
\NormalTok{\}}

\NormalTok{function factorise(n) \{}
\NormalTok{  const factors = [];}
\NormalTok{  let d = 2;}
\NormalTok{  let m = n;}
\NormalTok{  while (d * d \textless{}= m) \{}
\NormalTok{    while (m \% d === 0) \{}
\NormalTok{      factors.push(d);}
\NormalTok{      m = Math.floor(m / d);}
\NormalTok{    \}}
\NormalTok{    d++;}
\NormalTok{  \}}
\NormalTok{  if (m \textgreater{} 1) factors.push(m);}
\NormalTok{  return factors;}
\NormalTok{\}}

\NormalTok{function formatFactors(n) \{}
\NormalTok{  const f = factorise(n);}
\NormalTok{  if (f.length \textless{}= 1) return String(n);}
\NormalTok{  // group into exponent notation for display}
\NormalTok{  const counts = \{\};}
\NormalTok{  for (const p of f) counts[p] = (counts[p] || 0) + 1;}
\NormalTok{  return Object.entries(counts)}
\NormalTok{    .map(([p, e]) =\textgreater{} e === 1 ? p : \textasciigrave{}$\{p\}\^{}\{$\{e\}\}\textasciigrave{})}
\NormalTok{    .join(" \textbackslash{}\textbackslash{}times ");}
\NormalTok{\}}

\NormalTok{viewof nSlider = Inputs.range([1, 50], \{}
\NormalTok{  step: 1,}
\NormalTok{  value: 1,}
\NormalTok{  label: "n"}
\NormalTok{\})}
\end{Highlighting}
\end{Shaded}

\begin{Shaded}
\begin{Highlighting}[]
\NormalTok{//| echo: false}

\NormalTok{\{}
\NormalTok{  const n = nSlider;}
\NormalTok{  const p = n * n {-} n + 41;}
\NormalTok{  const prime = isPrime(p);}

\NormalTok{  const dotColour = prime ? "\#22c55e" : "\#ef4444";}
\NormalTok{  const statusText = prime ? "prime" : "composite";}
\NormalTok{  const statusColour = prime ? "\#16a34a" : "\#dc2626";}

\NormalTok{  // Factorisation string for composite values}
\NormalTok{  let factorHTML = "";}
\NormalTok{  if (!prime) \{}
\NormalTok{    const counts = \{\};}
\NormalTok{    for (const f of factorise(p)) counts[f] = (counts[f] || 0) + 1;}
\NormalTok{    const fStr = Object.entries(counts)}
\NormalTok{      .map(([f, e]) =\textgreater{} e === 1 ? f : \textasciigrave{}$\{f\}\textless{}sup\textgreater{}$\{e\}\textless{}/sup\textgreater{}\textasciigrave{})}
\NormalTok{      .join(" × ");}
\NormalTok{    factorHTML = \textasciigrave{}\textless{}p style="margin:0.4rem 0 0;font{-}size:0.9rem;color:\#374151;"\textgreater{}}
\NormalTok{      Factorisation: $\{p\} = $\{fStr\}}
\NormalTok{    \textless{}/p\textgreater{}\textasciigrave{};}
\NormalTok{  \}}

\NormalTok{  // Show a snapshot table of n = 1..current n (up to 10 most recent from 1)}
\NormalTok{  const tableStart = Math.max(1, n {-} 9);}
\NormalTok{  const tableRows = Array.from(\{ length: n {-} tableStart + 1 \}, (\_, i) =\textgreater{} \{}
\NormalTok{    const k = tableStart + i;}
\NormalTok{    const pk = k * k {-} k + 41;}
\NormalTok{    const primeK = isPrime(pk);}
\NormalTok{    const c = primeK ? "\#16a34a" : "\#dc2626";}
\NormalTok{    const lbl = primeK ? "prime" : "composite";}
\NormalTok{    const bold = k === n ? "font{-}weight:700;background:\#f5f3ff;" : "";}
\NormalTok{    return \textasciigrave{}\textless{}tr style="$\{bold\}"\textgreater{}}
\NormalTok{      \textless{}td style="padding:3px 10px;text{-}align:center;"\textgreater{}$\{k\}\textless{}/td\textgreater{}}
\NormalTok{      \textless{}td style="padding:3px 10px;text{-}align:center;"\textgreater{}$\{pk\}\textless{}/td\textgreater{}}
\NormalTok{      \textless{}td style="padding:3px 10px;text{-}align:center;color:$\{c\};"\textgreater{}$\{lbl\}\textless{}/td\textgreater{}}
\NormalTok{    \textless{}/tr\textgreater{}\textasciigrave{};}
\NormalTok{  \}).reverse().join("");}

\NormalTok{  return htl.html\textasciigrave{}\textless{}div style="font{-}family:sans{-}serif;max{-}width:480px;"\textgreater{}}
\NormalTok{    \textless{}div style="padding:1rem 1.25rem;border:1px solid \#e5e7eb;border{-}radius:8px;margin{-}bottom:0.75rem;"\textgreater{}}
\NormalTok{      \textless{}p style="margin:0 0 0.35rem;font{-}size:1rem;"\textgreater{}\textless{}em\textgreater{}n\textless{}/em\textgreater{} = \textless{}strong\textgreater{}$\{n\}\textless{}/strong\textgreater{}\textless{}/p\textgreater{}}
\NormalTok{      \textless{}p style="margin:0 0 0.5rem;font{-}size:1rem;"\textgreater{}\textless{}em\textgreater{}n\textless{}/em\textgreater{}² − \textless{}em\textgreater{}n\textless{}/em\textgreater{} + 41 = \textless{}strong\textgreater{}$\{p\}\textless{}/strong\textgreater{}\textless{}/p\textgreater{}}
\NormalTok{      \textless{}p style="margin:0;display:flex;align{-}items:center;gap:6px;"\textgreater{}}
\NormalTok{        \textless{}span style="display:inline{-}block;width:13px;height:13px;border{-}radius:50\%;background:$\{dotColour\};flex{-}shrink:0;"\textgreater{}\textless{}/span\textgreater{}}
\NormalTok{        \textless{}strong style="color:$\{statusColour\};"\textgreater{}$\{p\} is $\{statusText\}\textless{}/strong\textgreater{}}
\NormalTok{      \textless{}/p\textgreater{}}
\NormalTok{      $\{htl.html([factorHTML])\}}
\NormalTok{    \textless{}/div\textgreater{}}
\NormalTok{    \textless{}table style="border{-}collapse:collapse;width:100\%;font{-}size:0.88rem;"\textgreater{}}
\NormalTok{      \textless{}thead\textgreater{}}
\NormalTok{        \textless{}tr style="border{-}bottom:2px solid \#d1d5db;"\textgreater{}}
\NormalTok{          \textless{}th style="padding:3px 10px;font{-}weight:600;"\textgreater{}n\textless{}/th\textgreater{}}
\NormalTok{          \textless{}th style="padding:3px 10px;font{-}weight:600;"\textgreater{}n²−n+41\textless{}/th\textgreater{}}
\NormalTok{          \textless{}th style="padding:3px 10px;font{-}weight:600;"\textgreater{}Status\textless{}/th\textgreater{}}
\NormalTok{        \textless{}/tr\textgreater{}}
\NormalTok{      \textless{}/thead\textgreater{}}
\NormalTok{      \textless{}tbody\textgreater{}$\{htl.html([tableRows])\}\textless{}/tbody\textgreater{}}
\NormalTok{    \textless{}/table\textgreater{}}
\NormalTok{  \textless{}/div\textgreater{}\textasciigrave{};}
\NormalTok{\}}
\end{Highlighting}
\end{Shaded}

\section{What a proof is}\label{what-a-proof-is}

A \textbf{proof} is a sequence of logical deductions from agreed
starting points --- called \textbf{axioms} or previously proved
\textbf{theorems} --- leading to a conclusion. Each step must follow
from the previous ones by a valid logical rule. If every step is valid,
the conclusion is guaranteed. Not likely. Not probably. Guaranteed.

Every proof has the same skeleton:

\begin{itemize}
\tightlist
\item
  \textbf{Statement:} what you are claiming is true
\item
  \textbf{Hypothesis:} what you are allowed to assume
\item
  \textbf{Conclusion:} what you need to establish
\item
  \textbf{Proof:} the chain of deductions from hypothesis to conclusion
\end{itemize}

The hypothesis is given to you. The conclusion is what you must reach.
Everything in between is your work.

\subsection{The notation of logical
structure}\label{the-notation-of-logical-structure}

Two symbols carry most of the logical load:

\(P \Rightarrow Q\) means \textbf{\(P\) implies \(Q\)}: if \(P\) is
true, then \(Q\) must be true. Read it as ``if \(P\), then \(Q\).'' The
arrow shows direction --- \(P \Rightarrow Q\) does not automatically
mean \(Q \Rightarrow P\).

\(P \Leftrightarrow Q\) means \textbf{\(P\) if and only if \(Q\)}: the
implication runs both ways. \(P\) is true exactly when \(Q\) is true,
and vice versa.

Two quantifiers say \emph{which} objects you are making claims about:

\(\forall x\) means \textbf{for all \(x\)} --- the statement holds for
every element in the domain under discussion.

\(\exists x\) means \textbf{there exists an \(x\)} --- the statement
holds for at least one element.

The symbol \(x \in S\) means \textbf{\(x\) is a member of the set
\(S\)}. You will see it most often with the standard number sets:

\begin{itemize}
\tightlist
\item
  \(\mathbb{N}\) --- the natural numbers: \(1, 2, 3, \ldots\)
\item
  \(\mathbb{Z}\) --- the integers: \(\ldots, -2, -1, 0, 1, 2, \ldots\)
\item
  \(\mathbb{Q}\) --- the rationals: all fractions \(p/q\) with
  \(p, q \in \mathbb{Z}\), \(q \neq 0\)
\item
  \(\mathbb{R}\) --- the real numbers: everything on the continuous
  number line
\end{itemize}

So ``\(\forall n \in \mathbb{Z}\)'' reads: ``for every integer \(n\).''
A typical theorem opens with a phrase like this to tell you exactly what
kind of object is being discussed.

\section{Four proof methods}\label{four-proof-methods}

There is no single template. What follows are four approaches that cover
the vast majority of what you will encounter in the first two years of
university mathematics.

\subsection{1. Direct proof}\label{direct-proof}

The simplest strategy: assume the hypothesis, apply definitions and
algebra, arrive at the conclusion.

\begin{tcolorbox}[enhanced jigsaw, opacityback=0, colback=white, breakable, colframe=quarto-callout-note-color-frame, arc=.35mm, left=2mm, toprule=.15mm, leftrule=.75mm, bottomrule=.15mm, rightrule=.15mm]

\textbf{Theorem.} The sum of two even integers is even.

\textbf{Proof.} Let \(m\) and \(n\) be even integers. By definition of
``even,'' there exist integers \(a\) and \(b\) such that \(m = 2a\) and
\(n = 2b\). Then:

\[m + n = 2a + 2b = 2(a + b)\]

Since \(a + b\) is an integer, \(m + n\) is divisible by 2, hence even.
\(\square\)

\end{tcolorbox}

Notice the structure: start with the definitions of the objects
involved, do algebra, end at the conclusion. The key move was unpacking
``even'' into its definition (\(m = 2a\)) so that the algebra could
work.

\subsection{2. Proof by contrapositive}\label{proof-by-contrapositive}

Every statement ``if \(P\) then \(Q\)'' (\(P \Rightarrow Q\)) is
\textbf{logically equivalent} (meaning: any proof of one automatically
proves the other --- they make exactly the same claim in different
words) to its \textbf{contrapositive}: ``if \(\neg Q\) then \(\neg P\)''
(\(\neg Q \Rightarrow \neg P\)), where \(\neg Q\) means ``not \(Q\)''
--- the opposite truth value.

Sometimes the contrapositive is much easier to work with directly.

\begin{tcolorbox}[enhanced jigsaw, opacityback=0, colback=white, breakable, colframe=quarto-callout-note-color-frame, arc=.35mm, left=2mm, toprule=.15mm, leftrule=.75mm, bottomrule=.15mm, rightrule=.15mm]

\textbf{Theorem.} If \(n^2\) is even, then \(n\) is even.

\textbf{Proof.} We prove the contrapositive: if \(n\) is odd, then
\(n^2\) is odd.

Suppose \(n\) is odd. Then \(n = 2k + 1\) for some integer \(k\).
Squaring:

\[n^2 = (2k+1)^2 = 4k^2 + 4k + 1 = 2(2k^2 + 2k) + 1\]

Since \(2k^2 + 2k\) is an integer, \(n^2\) has the form
\(2(\text{integer}) + 1\), so \(n^2\) is odd.

The contrapositive is proved. Therefore the original statement holds.
\(\square\)

\end{tcolorbox}

The reason to choose contrapositive over direct proof here: starting
from ``\(n^2\) is even'' and trying to work towards ``\(n\) is even'' is
awkward --- you'd have to take a square root, and it isn't obvious that
\(\sqrt{2k}\) is an integer. The contrapositive avoids that entirely.

\subsection{3. Proof by contradiction}\label{proof-by-contradiction}

Assume the conclusion is \emph{false}. Show that this assumption,
together with the hypothesis, leads to a logical impossibility --- a
statement that contradicts something you know to be true. Since the
assumption of falseness leads to a contradiction, the conclusion must be
true.

The classic example is one you may have seen stated but never seen
proved:

\begin{tcolorbox}[enhanced jigsaw, opacityback=0, colback=white, breakable, colframe=quarto-callout-note-color-frame, arc=.35mm, left=2mm, toprule=.15mm, leftrule=.75mm, bottomrule=.15mm, rightrule=.15mm]

\textbf{Theorem.} \(\sqrt{2}\) is irrational.

\textbf{Proof.} Suppose for contradiction that \(\sqrt{2}\) is rational.
Then we can write \(\sqrt{2} = p/q\) where \(p, q \in \mathbb{Z}\),
\(q \neq 0\), and \(p\) and \(q\) share no common factor (the fraction
is in \textbf{lowest terms}).

Squaring both sides: \(2 = p^2 / q^2\), so \(p^2 = 2q^2\).

This means \(p^2\) is even. By the theorem proved above
(contrapositive), \(p\) must be even. So write \(p = 2k\) for some
integer \(k\).

Substituting: \((2k)^2 = 2q^2\), which gives \(4k^2 = 2q^2\), so
\(q^2 = 2k^2\).

This means \(q^2\) is even, so \(q\) is even (by the same reasoning we
used for \(p\): if a number's square is even, the number itself must be
even).

But now both \(p\) and \(q\) are even --- they share a factor of 2. This
contradicts our assumption that \(p/q\) was in lowest terms.

The assumption that \(\sqrt{2}\) is rational leads to a contradiction.
Therefore \(\sqrt{2}\) is irrational. \(\square\)

\end{tcolorbox}

Note how the proof assembled earlier results --- ``if \(n^2\) is even
then \(n\) is even'' was used twice. Proofs are built on top of other
proofs. The earlier theorem was not an isolated exercise; it was a tool.

\subsection{4. Mathematical induction}\label{mathematical-induction}

Induction proves statements of the form ``this holds for every natural
number \(n \geq 1\)'' (or \(n \geq 0\), or from any fixed starting
point).

The structure has two parts:

\begin{enumerate}
\def\labelenumi{\arabic{enumi}.}
\tightlist
\item
  \textbf{Base case.} Show the statement holds for the starting value
  (usually \(n = 1\)).
\item
  \textbf{Inductive step.} Show that \emph{if} the statement holds for
  \(n = k\) (the \textbf{inductive hypothesis}), then it holds for
  \(n = k + 1\).
\end{enumerate}

If both hold, the statement holds for all \(n\) from the base case
onward. The logic is a chain: the base case starts the chain, the
inductive step extends it one link at a time, and there is no upper
limit on the chain length.

\begin{tcolorbox}[enhanced jigsaw, opacityback=0, colback=white, breakable, colframe=quarto-callout-note-color-frame, arc=.35mm, left=2mm, toprule=.15mm, leftrule=.75mm, bottomrule=.15mm, rightrule=.15mm]

\textbf{Theorem.} For all \(n \geq 1\):

\[1 + 2 + 3 + \cdots + n = \frac{n(n+1)}{2}\]

\textbf{Proof.}

\textbf{Base case.} When \(n = 1\): the left side is \(1\), and
\(\frac{1 \cdot 2}{2} = 1\). The statement holds.

\textbf{Inductive step.} Assume the statement holds for \(n = k\), i.e.,
assume:

\[1 + 2 + \cdots + k = \frac{k(k+1)}{2}\]

We need to show it holds for \(n = k + 1\), i.e., that:

\[1 + 2 + \cdots + k + (k+1) = \frac{(k+1)(k+2)}{2}\]

Starting from the left side and using the inductive hypothesis:

\[1 + 2 + \cdots + k + (k+1) = \frac{k(k+1)}{2} + (k+1)\]

Factor out \((k+1)\):

\[= (k+1)\left(\frac{k}{2} + 1\right) = (k+1) \cdot \frac{k+2}{2} = \frac{(k+1)(k+2)}{2}\]

This is exactly what we needed to show.

By induction, the formula holds for all \(n \geq 1\). \(\square\)

\end{tcolorbox}

The formula \(n(n+1)/2\) is not new to you --- it appeared when we
studied arithmetic series. But now you have not just a formula and some
confirming examples. You have a proof. It holds for \(n = 10^{100}\) as
surely as for \(n = 5\).

\begin{Shaded}
\begin{Highlighting}[]
\NormalTok{//| echo: false}

\NormalTok{// ── VIZ 2: Induction visualiser for Σk = n(n+1)/2 ────────────────────────}

\NormalTok{viewof nInduct = Inputs.range([1, 20], \{}
\NormalTok{  step: 1,}
\NormalTok{  value: 5,}
\NormalTok{  label: "n"}
\NormalTok{\})}
\end{Highlighting}
\end{Shaded}

\begin{Shaded}
\begin{Highlighting}[]
\NormalTok{//| echo: false}

\NormalTok{\{}
\NormalTok{  const n = nInduct;}

\NormalTok{  // Compute partial sums S\_k = k(k+1)/2 for k = 1..n}
\NormalTok{  const data = Array.from(\{ length: n \}, (\_, i) =\textgreater{} \{}
\NormalTok{    const k = i + 1;}
\NormalTok{    const Sk = k * (k + 1) / 2;}
\NormalTok{    return \{ k, Sk \};}
\NormalTok{  \});}

\NormalTok{  const Sn = n * (n + 1) / 2;}
\NormalTok{  const maxS = data[data.length {-} 1].Sk;}

\NormalTok{  // SVG dimensions}
\NormalTok{  const W = 480, H = 260;}
\NormalTok{  const padL = 48, padR = 20, padT = 20, padB = 40;}
\NormalTok{  const chartW = W {-} padL {-} padR;}
\NormalTok{  const chartH = H {-} padT {-} padB;}

\NormalTok{  const xScale = k =\textgreater{} padL + ((k {-} 1) / Math.max(n {-} 1, 1)) * chartW;}
\NormalTok{  const yScale = v =\textgreater{} padT + chartH {-} (v / maxS) * chartH;}
\NormalTok{  const barW = Math.max(4, Math.min(32, (chartW / n) * 0.6));}

\NormalTok{  // Build bar rects}
\NormalTok{  const bars = data.map((\{ k, Sk \}) =\textgreater{} \{}
\NormalTok{    const x = xScale(k) {-} barW / 2;}
\NormalTok{    const y = yScale(Sk);}
\NormalTok{    const h = padT + chartH {-} y;}
\NormalTok{    return \textasciigrave{}\textless{}rect x="$\{x.toFixed(1)\}" y="$\{y.toFixed(1)\}" width="$\{barW\}" height="$\{h.toFixed(1)\}"}
\NormalTok{              fill="\#6366f1" opacity="0.75" rx="2"/\textgreater{}\textasciigrave{};}
\NormalTok{  \}).join("\textbackslash{}n");}

\NormalTok{  // Build formula dots + connecting line}
\NormalTok{  const pts = data.map((\{ k, Sk \}) =\textgreater{} \textasciigrave{}$\{xScale(k).toFixed(1)\},$\{yScale(Sk).toFixed(1)\}\textasciigrave{}).join(" ");}
\NormalTok{  const dots = data.map((\{ k, Sk \}) =\textgreater{}}
\NormalTok{    \textasciigrave{}\textless{}circle cx="$\{xScale(k).toFixed(1)\}" cy="$\{yScale(Sk).toFixed(1)\}" r="5" fill="\#f97316" stroke="white" stroke{-}width="1.5"/\textgreater{}\textasciigrave{}}
\NormalTok{  ).join("\textbackslash{}n");}

\NormalTok{  // x{-}axis tick labels}
\NormalTok{  const xTicks = data.map((\{ k \}) =\textgreater{}}
\NormalTok{    \textasciigrave{}\textless{}text x="$\{xScale(k).toFixed(1)\}" y="$\{(padT + chartH + 18).toFixed(1)\}" text{-}anchor="middle" font{-}size="11" fill="\#6b7280"\textgreater{}$\{k\}\textless{}/text\textgreater{}\textasciigrave{}}
\NormalTok{  ).join("\textbackslash{}n");}

\NormalTok{  // y{-}axis ticks (3–4 ticks)}
\NormalTok{  const yTickCount = 4;}
\NormalTok{  const yTicks = Array.from(\{ length: yTickCount + 1 \}, (\_, i) =\textgreater{} \{}
\NormalTok{    const v = Math.round((maxS / yTickCount) * i);}
\NormalTok{    const y = yScale(v);}
\NormalTok{    return \textasciigrave{}\textless{}line x1="$\{padL {-} 4\}" y1="$\{y.toFixed(1)\}" x2="$\{padL\}" y2="$\{y.toFixed(1)\}" stroke="\#9ca3af" stroke{-}width="1"/\textgreater{}}
\NormalTok{\textless{}text x="$\{(padL {-} 6).toFixed(1)\}" y="$\{(y + 4).toFixed(1)\}" text{-}anchor="end" font{-}size="10" fill="\#9ca3af"\textgreater{}$\{v\}\textless{}/text\textgreater{}\textasciigrave{};}
\NormalTok{  \}).join("\textbackslash{}n");}

\NormalTok{  const svg = \textasciigrave{}\textless{}svg viewBox="0 0 $\{W\} $\{H\}" style="width:100\%;max{-}width:$\{W\}px;display:block;font{-}family:sans{-}serif;"\textgreater{}}
\NormalTok{  \textless{}!{-}{-} axes {-}{-}\textgreater{}}
\NormalTok{  \textless{}line x1="$\{padL\}" y1="$\{padT\}" x2="$\{padL\}" y2="$\{padT + chartH\}" stroke="\#d1d5db" stroke{-}width="1"/\textgreater{}}
\NormalTok{  \textless{}line x1="$\{padL\}" y1="$\{padT + chartH\}" x2="$\{padL + chartW\}" y2="$\{padT + chartH\}" stroke="\#d1d5db" stroke{-}width="1"/\textgreater{}}
\NormalTok{  $\{yTicks\}}
\NormalTok{  $\{xTicks\}}
\NormalTok{  \textless{}!{-}{-} bars {-}{-}\textgreater{}}
\NormalTok{  $\{bars\}}
\NormalTok{  \textless{}!{-}{-} formula line {-}{-}\textgreater{}}
\NormalTok{  \textless{}polyline points="$\{pts\}" fill="none" stroke="\#f97316" stroke{-}width="2" opacity="0.6"/\textgreater{}}
\NormalTok{  \textless{}!{-}{-} formula dots {-}{-}\textgreater{}}
\NormalTok{  $\{dots\}}
\NormalTok{  \textless{}!{-}{-} legend {-}{-}\textgreater{}}
\NormalTok{  \textless{}rect x="$\{padL + 4\}" y="$\{padT + 4\}" width="12" height="12" fill="\#6366f1" opacity="0.75" rx="2"/\textgreater{}}
\NormalTok{  \textless{}text x="$\{padL + 20\}" y="$\{padT + 14\}" font{-}size="11" fill="\#374151"\textgreater{}partial sum\textless{}/text\textgreater{}}
\NormalTok{  \textless{}circle cx="$\{padL + 130\}" cy="$\{padT + 10\}" r="5" fill="\#f97316"/\textgreater{}}
\NormalTok{  \textless{}text x="$\{padL + 140\}" y="$\{padT + 14\}" font{-}size="11" fill="\#374151"\textgreater{}k(k+1)/2\textless{}/text\textgreater{}}
\NormalTok{  \textless{}!{-}{-} x{-}axis label {-}{-}\textgreater{}}
\NormalTok{  \textless{}text x="$\{(padL + chartW / 2).toFixed(1)\}" y="$\{H {-} 4\}" text{-}anchor="middle" font{-}size="11" fill="\#6b7280"\textgreater{}k\textless{}/text\textgreater{}}
\NormalTok{\textless{}/svg\textgreater{}\textasciigrave{};}

\NormalTok{  const formulaValue = n * (n + 1) / 2;}

\NormalTok{  return htl.html\textasciigrave{}\textless{}div style="font{-}family:sans{-}serif;max{-}width:500px;"\textgreater{}}
\NormalTok{    $\{htl.html([svg])\}}
\NormalTok{    \textless{}p style="margin:0.75rem 0 0.25rem;font{-}size:0.95rem;"\textgreater{}}
\NormalTok{      Sum: \textless{}strong\textgreater{}1 + 2 + ⋯ + $\{n\} = $\{Sn\}\textless{}/strong\textgreater{}}
\NormalTok{    \textless{}/p\textgreater{}}
\NormalTok{    \textless{}p style="margin:0;font{-}size:0.95rem;"\textgreater{}}
\NormalTok{      Formula: \textless{}strong\textgreater{}$\{n\}($\{n\}+1)/2 = $\{n\} × $\{n+1\} / 2 = $\{formulaValue\}\textless{}/strong\textgreater{}}
\NormalTok{    \textless{}/p\textgreater{}}
\NormalTok{  \textless{}/div\textgreater{}\textasciigrave{};}
\NormalTok{\}}
\end{Highlighting}
\end{Shaded}

\section{What makes a proof fail}\label{what-makes-a-proof-fail}

Not every argument that looks like a proof is one. Three failure modes
appear constantly, and learning to recognise them is half the work.

\subsection{1. Assuming the conclusion}\label{assuming-the-conclusion}

Also called \textbf{circular reasoning}: a step in the proof assumes, in
some form, the very thing you are trying to prove.

Here is a compact example. Suppose someone ``proves'' that \(1 = 2\)
like this:

\begin{quote}
Let \(a = b\). Then \(a^2 = ab\), so \(a^2 - b^2 = ab - b^2\), so
\((a+b)(a-b) = b(a-b)\), so \(a + b = b\). Since \(a = b\), this gives
\(2b = b\), so \(2 = 1\).
\end{quote}

The flaw: the step ``so \(a + b = b\)'' divides both sides by
\((a - b)\). But \(a = b\), so \(a - b = 0\). Division by zero is not a
valid algebraic step. It is not that the algebra went wrong in a subtle
way --- the step is simply not allowed.

Circular reasoning is harder to spot than hidden division by zero. Watch
for any step that \emph{uses} the thing you're trying to
\emph{establish}.

\subsection{2. Missing a case}\label{missing-a-case}

A proof that only covers some instances --- even most instances --- is
not a proof. The primality of \(n^2 - n + 41\) was ``verified'' for 40
cases and failed on the 41st. Missing cases can be more subtle: a proof
about integers that silently only works for positive ones, a proof about
triangles that only handles the acute case.

Before finishing a proof, ask: are there configurations I haven't
considered? Edge cases, boundary values, negative inputs, degenerate
shapes?

\subsection{3. An invalid step}\label{an-invalid-step}

An algebraic manipulation that is not universally valid. The most
common:

\begin{itemize}
\tightlist
\item
  \textbf{Dividing by a variable} without confirming it's nonzero (as in
  the \(a = b\) example above)
\item
  \textbf{Taking a square root} and losing the negative case:
  \(x^2 = 4\) implies \(x = 2\) or \(x = -2\), not just \(x = 2\)
\item
  \textbf{Cancelling factors} that might be zero
\item
  \textbf{Swapping the order} of limits, integrals, or sums without
  justification (this matters more in later courses)
\end{itemize}

Every algebraic step has conditions under which it is valid. When those
conditions aren't met, the step is wrong --- even if it gives the right
answer in specific cases.

\begin{tcolorbox}[enhanced jigsaw, title=\textcolor{quarto-callout-tip-color}{\faLightbulb}\hspace{0.5em}{What proofs are actually for}, opacityback=0, colback=white, breakable, leftrule=.75mm, arc=.35mm, toptitle=1mm, left=2mm, colframe=quarto-callout-tip-color-frame, coltitle=black, titlerule=0mm, colbacktitle=quarto-callout-tip-color!10!white, opacitybacktitle=0.6, toprule=.15mm, bottomtitle=1mm, bottomrule=.15mm, rightrule=.15mm]

Proofs are not rituals performed for examiners. They are the reason
mathematicians can trust results for trillions of cases they've never
checked.

When you prove that an algorithm runs in \(O(n \log n)\) time, you're
not estimating --- you're guaranteeing. When a cryptographic security
proof says the encryption is unbreakable under certain assumptions, the
logic is auditable. A regulator, a competitor, a court can read every
step. The proof is the audit trail.

Physical laws are not just patterns observed in experiments. They are
structures whose consequences can be derived --- predicted before they
are measured. The derivation is a proof. When a prediction matches an
observation, it confirms the law; when it doesn't, the law is wrong. The
reliability of science rests on this structure.

Mathematics is the only domain where ``I've checked a huge number of
cases'' is explicitly not sufficient, and where a short argument can
settle a question for every case simultaneously. That is what makes it
powerful.

\end{tcolorbox}

\section{Worked examples}\label{worked-examples-20}

\subsection{Example 1: Direct proof --- the product of two odd
numbers}\label{example-1-direct-proof-the-product-of-two-odd-numbers}

\textbf{Context.} Two players take turns on a board game. Each player's
turn takes an odd number of seconds. Prove that no matter how many turns
have been played, the product of any two turn lengths is always odd.

\begin{tcolorbox}[enhanced jigsaw, opacityback=0, colback=white, breakable, colframe=quarto-callout-note-color-frame, arc=.35mm, left=2mm, toprule=.15mm, leftrule=.75mm, bottomrule=.15mm, rightrule=.15mm]

\textbf{Theorem.} The product of two odd integers is odd.

\textbf{Proof.} Let \(m\) and \(n\) be odd integers. By definition,
there exist integers \(a\) and \(b\) such that \(m = 2a + 1\) and
\(n = 2b + 1\).

Then:
\[m \cdot n = (2a+1)(2b+1) = 4ab + 2a + 2b + 1 = 2(2ab + a + b) + 1\]

Since \(2ab + a + b\) is an integer, \(mn\) has the form
\(2(\text{integer}) + 1\), so \(mn\) is odd. \(\square\)

\end{tcolorbox}

\textbf{Commentary.} The key move was unpacking the definition of
``odd'' into algebra. Once you write \(m = 2a + 1\) and \(n = 2b + 1\),
the result of multiplying them is computable. The only thing left is to
recognise that the answer has the right form --- it is one more than an
even number.

The proof doesn't appeal to intuition, examples, or the fact that you've
tested it. It works because the algebra always produces
\(2(\ldots) + 1\), regardless of what \(a\) and \(b\) are.

\subsection{Example 2: Proof by contradiction --- infinitely many
primes}\label{example-2-proof-by-contradiction-infinitely-many-primes}

\textbf{Context.} Public-key cryptography (RSA, used in every HTTPS
connection) relies on the existence of very large prime numbers. The
security argument assumes there are infinitely many primes to draw from.
Euclid proved this over 2000 years ago. The argument has not changed.

\begin{tcolorbox}[enhanced jigsaw, opacityback=0, colback=white, breakable, colframe=quarto-callout-note-color-frame, arc=.35mm, left=2mm, toprule=.15mm, leftrule=.75mm, bottomrule=.15mm, rightrule=.15mm]

\textbf{Theorem.} There are infinitely many prime numbers.

\textbf{Proof.} Suppose for contradiction that there are only finitely
many primes. List them all: \(p_1, p_2, \ldots, p_k\).

Consider the number: \[N = p_1 \cdot p_2 \cdots p_k + 1\]

\(N\) is greater than 1, so it has at least one prime factor. Call it
\(p\). Since our list contains all primes, \(p\) must be one of
\(p_1, p_2, \ldots, p_k\).

But \(p\) divides \(p_1 \cdot p_2 \cdots p_k\) (since \(p\) is one of
those factors). And \(p\) divides \(N\). So \(p\) must divide the
difference: \[N - p_1 \cdot p_2 \cdots p_k = 1\]

But no prime divides 1. Contradiction.

The assumption that there are finitely many primes leads to a
contradiction. Therefore there are infinitely many primes. \(\square\)

\end{tcolorbox}

\textbf{Commentary.} The number \(N\) is not claimed to be prime. It is
constructed deliberately so that dividing it by any prime on the list
leaves a remainder of 1. The contradiction is that the prime factor of
\(N\) must be on the list but cannot be on the list.

This is one of the most elegant proofs in mathematics. It uses no
advanced machinery --- only the definition of prime, the fact that every
integer greater than 1 has a prime factor, and arithmetic. The idea
(constructing \(N\)) is the creative step; the logic is then
straightforward.

\subsection{\texorpdfstring{Example 3: Proof by induction ---
\(2^n > n\) for all
\(n \geq 1\)}{Example 3: Proof by induction --- 2\^{}n \textgreater{} n for all n \textbackslash geq 1}}\label{example-3-proof-by-induction-2n-n-for-all-n-geq-1}

\textbf{Context.} In digital circuit design, the number of possible
states in a circuit with \(n\) binary switches is \(2^n\). Engineers say
this ``grows exponentially.'' Proving \(2^n > n\) is a basic instance of
establishing that exponential growth outpaces linear growth --- relevant
whenever you're deciding whether a brute-force approach (checking every
state) is feasible.

\begin{tcolorbox}[enhanced jigsaw, opacityback=0, colback=white, breakable, colframe=quarto-callout-note-color-frame, arc=.35mm, left=2mm, toprule=.15mm, leftrule=.75mm, bottomrule=.15mm, rightrule=.15mm]

\textbf{Theorem.} For all integers \(n \geq 1\), \(2^n > n\).

\textbf{Proof.}

\textbf{Base case.} When \(n = 1\): \(2^1 = 2 > 1\). The statement
holds.

\textbf{Inductive step.} Assume the statement holds for \(n = k\), i.e.,
assume \(2^k > k\). We need to show \(2^{k+1} > k + 1\).

Starting from the inductive hypothesis \(2^k > k\), multiply both sides
by 2: \[2^{k+1} = 2 \cdot 2^k > 2k\]

It suffices to show \(2k \geq k + 1\), i.e., \(k \geq 1\). Since
\(k \geq 1\) by our starting point, this holds.

Therefore \(2^{k+1} > 2k \geq k + 1\).

By induction, \(2^n > n\) for all \(n \geq 1\). \(\square\)

\end{tcolorbox}

\textbf{Commentary.} The inductive step needed two pieces: first, using
the inductive hypothesis to get a bound on \(2^{k+1}\); second,
connecting that bound to \(k + 1\). The step ``\(2k \geq k+1\) because
\(k \geq 1\)'' looks almost too obvious to write down. Write it down
anyway. In a proof, ``obvious'' is not a valid reason --- the valid
reason is ``\(k \geq 1\).''

\section{Where this goes}\label{where-this-goes-21}

\textbf{Throughout this series.} Every other volume has been making
claims without full justification --- the sum formula for arithmetic
series, the laws of exponents, properties of functions. You now have the
tools to ask ``how do we know that?'' and to read the answer. The proofs
are in the textbooks; you can now follow them.

\textbf{Complex analysis (Volume 7).} The central results there ---
Cauchy's integral theorem, the residue theorem --- are not computational
tricks. They are theorems with proofs that depend on precise definitions
of limits, continuity, and differentiability. Those proofs are
accessible once you are comfortable with the logical structure
introduced here. Complex analysis is where proof and computation meet in
the most surprising way: a function's behaviour on a closed curve can
determine its behaviour at every point inside.

\begin{tcolorbox}[enhanced jigsaw, title=\textcolor{quarto-callout-tip-color}{\faLightbulb}\hspace{0.5em}{Applications}, opacityback=0, colback=white, breakable, leftrule=.75mm, arc=.35mm, toptitle=1mm, left=2mm, colframe=quarto-callout-tip-color-frame, coltitle=black, titlerule=0mm, colbacktitle=quarto-callout-tip-color!10!white, opacitybacktitle=0.6, toprule=.15mm, bottomtitle=1mm, bottomrule=.15mm, rightrule=.15mm]

\textbf{Algorithm correctness.} Proving that an algorithm produces the
right output for \emph{every} valid input, not just the test cases
you've tried. Induction is the standard tool for array algorithms and
recursive functions.

\textbf{Cryptographic security proofs.} RSA's security is reduced to the
difficulty of factoring large numbers --- the reduction is a proof. If
you could break RSA, you could factor; but factoring appears hard. The
chain of implication is a formal proof.

\textbf{Derivation from physical laws.} The equations governing heat
flow, electromagnetic fields, and quantum mechanics are not empirical
curve fits. They are derived --- proved --- from foundational
principles. The derivations are proofs.

\textbf{Formal verification.} Aviation, medical devices, and nuclear
control systems are subject to formal verification: software is checked
against a mathematical model, and the check is a proof. DO-178C
(aviation software standard) and IEC 62304 (medical software) both
recognise formal methods. The mathematical core is the logic introduced
here.

\end{tcolorbox}

\section{Exercises}\label{exercises-21}

\begin{Shaded}
\begin{Highlighting}[]
\NormalTok{//| echo: false}

\NormalTok{// ── Shared stepper infrastructure ──────────────────────────────────────────}

\NormalTok{katex = require("katex@0.16")}

\NormalTok{function makeStepperHTML(problem, steps) \{}
\NormalTok{  let current = {-}1;}

\NormalTok{  const problemDiv = document.createElement("div");}
\NormalTok{  problemDiv.style.cssText = "margin{-}bottom:0.75rem;font{-}weight:600;font{-}size:0.95rem;";}
\NormalTok{  problemDiv.textContent = problem;}

\NormalTok{  const stepsContainer = document.createElement("div");}

\NormalTok{  const stepEls = steps.map((s, i) =\textgreater{} \{}
\NormalTok{    const el = document.createElement("div");}
\NormalTok{    el.style.cssText = \textasciigrave{}}
\NormalTok{      display:none;}
\NormalTok{      margin{-}top:0.6rem;}
\NormalTok{      padding:0.6rem 0.85rem;}
\NormalTok{      border{-}left:3px solid \#6366f1;}
\NormalTok{      background:\#f5f3ff;}
\NormalTok{      border{-}radius:0 6px 6px 0;}
\NormalTok{      font{-}size:0.92rem;}
\NormalTok{    \textasciigrave{};}
\NormalTok{    const opSpan = document.createElement("div");}
\NormalTok{    opSpan.style.cssText = "color:\#6b7280;font{-}size:0.82rem;margin{-}bottom:0.3rem;font{-}style:italic;";}
\NormalTok{    opSpan.textContent = s.op;}

\NormalTok{    const eqDiv = document.createElement("div");}
\NormalTok{    try \{}
\NormalTok{      eqDiv.innerHTML = katex.renderToString(s.eq, \{ throwOnError: false, displayMode: false \});}
\NormalTok{    \} catch(e) \{}
\NormalTok{      eqDiv.textContent = s.eq;}
\NormalTok{    \}}

\NormalTok{    el.appendChild(opSpan);}
\NormalTok{    el.appendChild(eqDiv);}
\NormalTok{    stepsContainer.appendChild(el);}
\NormalTok{    return el;}
\NormalTok{  \});}

\NormalTok{  const btnRow = document.createElement("div");}
\NormalTok{  btnRow.style.cssText = "display:flex;gap:0.5rem;margin{-}top:0.75rem;align{-}items:center;";}

\NormalTok{  const nextBtn = document.createElement("button");}
\NormalTok{  nextBtn.textContent = "Next step →";}
\NormalTok{  nextBtn.style.cssText = \textasciigrave{}}
\NormalTok{    padding:0.35rem 0.85rem;}
\NormalTok{    background:\#6366f1;}
\NormalTok{    color:white;}
\NormalTok{    border:none;}
\NormalTok{    border{-}radius:5px;}
\NormalTok{    cursor:pointer;}
\NormalTok{    font{-}size:0.88rem;}
\NormalTok{  \textasciigrave{};}

\NormalTok{  const resetBtn = document.createElement("button");}
\NormalTok{  resetBtn.textContent = "Reset";}
\NormalTok{  resetBtn.style.cssText = \textasciigrave{}}
\NormalTok{    padding:0.35rem 0.75rem;}
\NormalTok{    background:transparent;}
\NormalTok{    color:\#6b7280;}
\NormalTok{    border:1px solid \#d1d5db;}
\NormalTok{    border{-}radius:5px;}
\NormalTok{    cursor:pointer;}
\NormalTok{    font{-}size:0.88rem;}
\NormalTok{  \textasciigrave{};}

\NormalTok{  const counter = document.createElement("span");}
\NormalTok{  counter.style.cssText = "font{-}size:0.82rem;color:\#9ca3af;margin{-}left:0.25rem;";}
\NormalTok{  counter.textContent = \textasciigrave{}0 / $\{steps.length\}\textasciigrave{};}

\NormalTok{  function update() \{}
\NormalTok{    stepEls.forEach((el, i) =\textgreater{} \{}
\NormalTok{      el.style.display = i \textless{}= current ? "block" : "none";}
\NormalTok{    \});}
\NormalTok{    counter.textContent = \textasciigrave{}$\{current + 1\} / $\{steps.length\}\textasciigrave{};}
\NormalTok{    nextBtn.disabled = current \textgreater{}= steps.length {-} 1;}
\NormalTok{    nextBtn.style.opacity = nextBtn.disabled ? "0.45" : "1";}
\NormalTok{  \}}

\NormalTok{  nextBtn.addEventListener("click", () =\textgreater{} \{}
\NormalTok{    if (current \textless{} steps.length {-} 1) \{ current++; update(); \}}
\NormalTok{  \});}
\NormalTok{  resetBtn.addEventListener("click", () =\textgreater{} \{}
\NormalTok{    current = {-}1; update();}
\NormalTok{  \});}

\NormalTok{  btnRow.appendChild(nextBtn);}
\NormalTok{  btnRow.appendChild(resetBtn);}
\NormalTok{  btnRow.appendChild(counter);}

\NormalTok{  const wrapper = document.createElement("div");}
\NormalTok{  wrapper.style.cssText = \textasciigrave{}}
\NormalTok{    border:1px solid \#e5e7eb;}
\NormalTok{    border{-}radius:8px;}
\NormalTok{    padding:1rem 1.1rem;}
\NormalTok{    max{-}width:560px;}
\NormalTok{    font{-}family:sans{-}serif;}
\NormalTok{    margin{-}top:0.75rem;}
\NormalTok{  \textasciigrave{};}
\NormalTok{  wrapper.appendChild(problemDiv);}
\NormalTok{  wrapper.appendChild(stepsContainer);}
\NormalTok{  wrapper.appendChild(btnRow);}

\NormalTok{  update();}
\NormalTok{  return wrapper;}
\NormalTok{\}}
\end{Highlighting}
\end{Shaded}

\textbf{Exercise 1.} Prove that if \(n\) is divisible by 3, then \(n^2\)
is divisible by 9.

\emph{Hint.} Write \(n = 3k\) for some integer \(k\), then compute
\(n^2\).

\begin{Shaded}
\begin{Highlighting}[]
\NormalTok{//| echo: false}

\NormalTok{makeStepperHTML(}
\NormalTok{  "If 3 divides n, prove that 9 divides n².",}
\NormalTok{  [}
\NormalTok{    \{}
\NormalTok{      op: "Write out the hypothesis in algebraic form",}
\NormalTok{      eq: "\textbackslash{}\textbackslash{}text\{Since \} 3 \textbackslash{}\textbackslash{}mid n, \textbackslash{}\textbackslash{}text\{ there exists an integer \} k \textbackslash{}\textbackslash{}text\{ such that \} n = 3k."}
\NormalTok{    \},}
\NormalTok{    \{}
\NormalTok{      op: "Compute n²",}
\NormalTok{      eq: "n\^{}2 = (3k)\^{}2 = 9k\^{}2"}
\NormalTok{    \},}
\NormalTok{    \{}
\NormalTok{      op: "Identify the structure of the result",}
\NormalTok{      eq: "9k\^{}2 = 9 \textbackslash{}\textbackslash{}cdot k\^{}2, \textbackslash{}\textbackslash{}text\{ and \} k\^{}2 \textbackslash{}\textbackslash{}in \textbackslash{}\textbackslash{}mathbb\{Z\}, \textbackslash{}\textbackslash{}text\{ so \} 9 \textbackslash{}\textbackslash{}mid n\^{}2."}
\NormalTok{    \},}
\NormalTok{    \{}
\NormalTok{      op: "Conclude",}
\NormalTok{      eq: "\textbackslash{}\textbackslash{}text\{Therefore \} n\^{}2 \textbackslash{}\textbackslash{}text\{ is divisible by \} 9. \textbackslash{}\textbackslash{}quad \textbackslash{}\textbackslash{}square"}
\NormalTok{    \}}
\NormalTok{  ]}
\NormalTok{)}
\end{Highlighting}
\end{Shaded}

\begin{center}\rule{0.5\linewidth}{0.5pt}\end{center}

\textbf{Exercise 2.} Prove by contradiction that \(\log_2 3\) is
irrational.

\emph{Hint.} Assume \(\log_2 3 = p/q\) in lowest terms and use the
definition of logarithm: \(\log_2 3 = p/q\) means \(2^{p/q} = 3\).

\begin{Shaded}
\begin{Highlighting}[]
\NormalTok{//| echo: false}

\NormalTok{makeStepperHTML(}
\NormalTok{  "Prove that log₂(3) is irrational.",}
\NormalTok{  [}
\NormalTok{    \{}
\NormalTok{      op: "Set up the contradiction",}
\NormalTok{      eq: "\textbackslash{}\textbackslash{}text\{Assume \} \textbackslash{}\textbackslash{}log\_2 3 = \textbackslash{}\textbackslash{}tfrac\{p\}\{q\} \textbackslash{}\textbackslash{}text\{ where \} p, q \textbackslash{}\textbackslash{}in \textbackslash{}\textbackslash{}mathbb\{N\} \textbackslash{}\textbackslash{}text\{ share no common factor.\}"}
\NormalTok{    \},}
\NormalTok{    \{}
\NormalTok{      op: "Use the definition of logarithm",}
\NormalTok{      eq: "\textbackslash{}\textbackslash{}log\_2 3 = \textbackslash{}\textbackslash{}tfrac\{p\}\{q\} \textbackslash{}\textbackslash{}implies 2\^{}\{p/q\} = 3 \textbackslash{}\textbackslash{}implies 2\^{}p = 3\^{}q"}
\NormalTok{    \},}
\NormalTok{    \{}
\NormalTok{      op: "Apply the Fundamental Theorem of Arithmetic",}
\NormalTok{      eq: "\textbackslash{}\textbackslash{}text\{Left side: only prime factor is \} 2. \textbackslash{}\textbackslash{}quad \textbackslash{}\textbackslash{}text\{Right side: only prime factor is \} 3."}
\NormalTok{    \},}
\NormalTok{    \{}
\NormalTok{      op: "Derive the contradiction",}
\NormalTok{      eq: "2\^{}p = 3\^{}q \textbackslash{}\textbackslash{}text\{ is impossible for positive integers \} p, q \textbackslash{}\textbackslash{}text\{ (distinct prime bases)\}."}
\NormalTok{    \},}
\NormalTok{    \{}
\NormalTok{      op: "Conclude",}
\NormalTok{      eq: "\textbackslash{}\textbackslash{}text\{The assumption that \} \textbackslash{}\textbackslash{}log\_2 3 \textbackslash{}\textbackslash{}in \textbackslash{}\textbackslash{}mathbb\{Q\} \textbackslash{}\textbackslash{}text\{ leads to a contradiction.\} \textbackslash{}\textbackslash{}quad \textbackslash{}\textbackslash{}square"}
\NormalTok{    \}}
\NormalTok{  ]}
\NormalTok{)}
\end{Highlighting}
\end{Shaded}

\begin{center}\rule{0.5\linewidth}{0.5pt}\end{center}

\textbf{Exercise 3.} Prove by induction that \(n! > 2^n\) for all
\(n \geq 4\).

\emph{Instructions.} (a) Check the base case \(n = 4\) by computing both
sides. (b) Write out the inductive step: assuming \(k! > 2^k\), show
\((k+1)! > 2^{k+1}\).

\begin{Shaded}
\begin{Highlighting}[]
\NormalTok{//| echo: false}

\NormalTok{makeStepperHTML(}
\NormalTok{  "Prove n! \textgreater{} 2ⁿ for all n ≥ 4.",}
\NormalTok{  [}
\NormalTok{    \{}
\NormalTok{      op: "Base case: n = 4",}
\NormalTok{      eq: "4! = 24, \textbackslash{}\textbackslash{}quad 2\^{}4 = 16, \textbackslash{}\textbackslash{}quad 24 \textgreater{} 16. \textbackslash{}\textbackslash{}checkmark"}
\NormalTok{    \},}
\NormalTok{    \{}
\NormalTok{      op: "State the inductive hypothesis",}
\NormalTok{      eq: "\textbackslash{}\textbackslash{}text\{Assume \} k! \textgreater{} 2\^{}k \textbackslash{}\textbackslash{}text\{ for some fixed \} k \textbackslash{}\textbackslash{}geq 4."}
\NormalTok{    \},}
\NormalTok{    \{}
\NormalTok{      op: "Expand the left side of the target",}
\NormalTok{      eq: "(k+1)! = (k+1) \textbackslash{}\textbackslash{}cdot k!"}
\NormalTok{    \},}
\NormalTok{    \{}
\NormalTok{      op: "Apply the inductive hypothesis",}
\NormalTok{      eq: "(k+1)! = (k+1) \textbackslash{}\textbackslash{}cdot k! \textgreater{} (k+1) \textbackslash{}\textbackslash{}cdot 2\^{}k"}
\NormalTok{    \},}
\NormalTok{    \{}
\NormalTok{      op: "Compare with the target 2\^{}(k+1)",}
\NormalTok{      eq: "\textbackslash{}\textbackslash{}text\{Need: \} (k+1) \textbackslash{}\textbackslash{}cdot 2\^{}k \textgreater{} 2\^{}\{k+1\} = 2 \textbackslash{}\textbackslash{}cdot 2\^{}k, \textbackslash{}\textbackslash{}text\{ i.e., \} k+1 \textgreater{} 2. \textbackslash{}\textbackslash{}text\{ True since \} k \textbackslash{}\textbackslash{}geq 4."}
\NormalTok{    \},}
\NormalTok{    \{}
\NormalTok{      op: "Conclude by induction",}
\NormalTok{      eq: "(k+1)! \textgreater{} (k+1) \textbackslash{}\textbackslash{}cdot 2\^{}k \textgreater{} 2 \textbackslash{}\textbackslash{}cdot 2\^{}k = 2\^{}\{k+1\}. \textbackslash{}\textbackslash{}quad \textbackslash{}\textbackslash{}therefore\textbackslash{}\textbackslash{} n! \textgreater{} 2\^{}n \textbackslash{}\textbackslash{} \textbackslash{}\textbackslash{}forall n \textbackslash{}\textbackslash{}geq 4. \textbackslash{}\textbackslash{}quad \textbackslash{}\textbackslash{}square"}
\NormalTok{    \}}
\NormalTok{  ]}
\NormalTok{)}
\end{Highlighting}
\end{Shaded}

\begin{center}\rule{0.5\linewidth}{0.5pt}\end{center}

\textbf{Exercise 4.} Find the flaw in the following ``proof'' that all
horses are the same colour.

\begin{quote}
\textbf{Claim.} In any group of \(n\) horses, all horses are the same
colour.

\textbf{Proof.}

\emph{Base case.} When \(n = 1\): a group of one horse trivially has all
horses the same colour (there is only one). ✓

\emph{Inductive step.} Assume any group of \(k\) horses is the same
colour. Consider a group of \(k+1\) horses:
\(H_1, H_2, \ldots, H_k, H_{k+1}\).

Look at the first \(k\) horses \(\{H_1, \ldots, H_k\}\). By the
inductive hypothesis, they are all the same colour.

Now look at the last \(k\) horses \(\{H_2, \ldots, H_{k+1}\}\). By the
inductive hypothesis, they are all the same colour.

These two groups overlap (they share \(H_2, \ldots, H_k\)), so all
\(k+1\) horses must be the same colour.

By induction, all horses in any group are the same colour. \(\square\)
\end{quote}

The base case is correct. The logic in the inductive step sounds
plausible. But the conclusion is obviously false. Where does the proof
break?

\emph{Hint.} What happens to the ``overlap'' argument when \(k = 1\)?

\begin{Shaded}
\begin{Highlighting}[]
\NormalTok{//| echo: false}

\NormalTok{makeStepperHTML(}
\NormalTok{  "Find the flaw in the horse{-}colour induction.",}
\NormalTok{  [}
\NormalTok{    \{}
\NormalTok{      op: "Accept the base case",}
\NormalTok{      eq: "\textbackslash{}\textbackslash{}text\{For \} n = 1 \textbackslash{}\textbackslash{}text\{: trivially true. No flaw here.\}"}
\NormalTok{    \},}
\NormalTok{    \{}
\NormalTok{      op: "Examine the inductive step at k = 1 (the step from n=1 to n=2)",}
\NormalTok{      eq: "\textbackslash{}\textbackslash{}text\{Groups: \} \textbackslash{}\textbackslash{}\{H\_1\textbackslash{}\textbackslash{}\} \textbackslash{}\textbackslash{}text\{ and \} \textbackslash{}\textbackslash{}\{H\_2\textbackslash{}\textbackslash{}\}. \textbackslash{}\textbackslash{}text\{ The claimed overlap is \} \textbackslash{}\textbackslash{}\{H\_2, \textbackslash{}\textbackslash{}ldots, H\_k\textbackslash{}\textbackslash{}\} = \textbackslash{}\textbackslash{}\{H\_2, \textbackslash{}\textbackslash{}ldots, H\_1\textbackslash{}\textbackslash{}\}."}
\NormalTok{    \},}
\NormalTok{    \{}
\NormalTok{      op: "Identify the flaw: the overlap is empty",}
\NormalTok{      eq: "\textbackslash{}\textbackslash{}text\{Indices from \} 2 \textbackslash{}\textbackslash{}text\{ to \} k = 1 \textbackslash{}\textbackslash{}text\{ form an empty range.\} \textbackslash{}\textbackslash{}quad |\textbackslash{}\textbackslash{}\{H\_2,\textbackslash{}\textbackslash{}ldots,H\_1\textbackslash{}\textbackslash{}\}| = 0."}
\NormalTok{    \},}
\NormalTok{    \{}
\NormalTok{      op: "State why the argument fails",}
\NormalTok{      eq: "\textbackslash{}\textbackslash{}text\{An empty overlap tells us nothing: \} H\_1 \textbackslash{}\textbackslash{}text\{ and \} H\_2 \textbackslash{}\textbackslash{}text\{ need not share a colour.\}"}
\NormalTok{    \},}
\NormalTok{    \{}
\NormalTok{      op: "Conclude",}
\NormalTok{      eq: "\textbackslash{}\textbackslash{}text\{The step \} n=1 \textbackslash{}\textbackslash{}to n=2 \textbackslash{}\textbackslash{}text\{ is not covered. The induction never leaves \} n=1. \textbackslash{}\textbackslash{}quad \textbackslash{}\textbackslash{}square"}
\NormalTok{    \}}
\NormalTok{  ]}
\NormalTok{)}
\end{Highlighting}
\end{Shaded}

\begin{center}\rule{0.5\linewidth}{0.5pt}\end{center}

\textbf{Exercise 5.} \emph{(Harder.)} Prove that between any two
distinct real numbers, there exists a rational number.

More precisely: prove that for any \(a, b \in \mathbb{R}\) with
\(a < b\), there exists \(q \in \mathbb{Q}\) with \(a < q < b\).

\emph{Hint.} You will need the \textbf{Archimedean property}: for any
real number \(x > 0\), there exists a natural number \(n\) such that
\(n > x\). Use it to find an integer \(n\) large enough that
\(1/n < b - a\). Then find an integer \(m\) such that \(m/n\) falls in
the interval \((a, b)\).

\begin{Shaded}
\begin{Highlighting}[]
\NormalTok{//| echo: false}

\NormalTok{makeStepperHTML(}
\NormalTok{  "Prove that between any two real numbers a \textless{} b, there is a rational.",}
\NormalTok{  [}
\NormalTok{    \{}
\NormalTok{      op: "Find a denominator n such that 1/n is smaller than the gap",}
\NormalTok{      eq: "b {-} a \textgreater{} 0. \textbackslash{}\textbackslash{}text\{ By the Archimedean property, \} \textbackslash{}\textbackslash{}exists\textbackslash{}\textbackslash{}, n \textbackslash{}\textbackslash{}in \textbackslash{}\textbackslash{}mathbb\{N\} \textbackslash{}\textbackslash{}text\{ with \} n \textgreater{} \textbackslash{}\textbackslash{}tfrac\{1\}\{b{-}a\}, \textbackslash{}\textbackslash{}text\{ so \} \textbackslash{}\textbackslash{}tfrac\{1\}\{n\} \textless{} b {-} a."}
\NormalTok{    \},}
\NormalTok{    \{}
\NormalTok{      op: "Choose numerator m: smallest integer with m \textgreater{} na",}
\NormalTok{      eq: "\textbackslash{}\textbackslash{}text\{Let \} m = \textbackslash{}\textbackslash{}lfloor na \textbackslash{}\textbackslash{}rfloor + 1, \textbackslash{}\textbackslash{}text\{ so \} m \textgreater{} na, \textbackslash{}\textbackslash{}text\{ i.e., \} \textbackslash{}\textbackslash{}tfrac\{m\}\{n\} \textgreater{} a."}
\NormalTok{    \},}
\NormalTok{    \{}
\NormalTok{      op: "Show m/n \textless{} b",}
\NormalTok{      eq: "m \textbackslash{}\textbackslash{}leq na + 1 \textbackslash{}\textbackslash{}implies \textbackslash{}\textbackslash{}tfrac\{m\}\{n\} \textbackslash{}\textbackslash{}leq a + \textbackslash{}\textbackslash{}tfrac\{1\}\{n\} \textless{} a + (b {-} a) = b."}
\NormalTok{    \},}
\NormalTok{    \{}
\NormalTok{      op: "Conclude",}
\NormalTok{      eq: "a \textless{} \textbackslash{}\textbackslash{}tfrac\{m\}\{n\} \textless{} b, \textbackslash{}\textbackslash{}quad \textbackslash{}\textbackslash{}tfrac\{m\}\{n\} \textbackslash{}\textbackslash{}in \textbackslash{}\textbackslash{}mathbb\{Q\}. \textbackslash{}\textbackslash{}quad \textbackslash{}\textbackslash{}therefore \textbackslash{}\textbackslash{}text\{ a rational exists between \} a \textbackslash{}\textbackslash{}text\{ and \} b. \textbackslash{}\textbackslash{}quad \textbackslash{}\textbackslash{}square"}
\NormalTok{    \}}
\NormalTok{  ]}
\NormalTok{)}
\end{Highlighting}
\end{Shaded}

\part{Volume 6 --- Probability and data}

\chapter{Volume 6 --- Probability and
data}\label{volume-6-probability-and-data-1}

The central move of this volume is quantifying uncertainty.

Often, you won't have the exact formula or the complete set of facts.
This volume provides the tools for operating in those gray areas where
outcomes are determined by chance or obscured by noise.

By the end of this volume, you will learn how to formally measure the
likelihood of future events, recognize the natural shapes that variation
takes, draw robust conclusions from limited samples, and separate true
signals from random fluctuations.

\section{Chapters}\label{chapters-5}

\begin{longtable}[]{@{}lll@{}}
\toprule\noalign{}
Chapter & Title & Core move \\
\midrule\noalign{}
\endhead
\bottomrule\noalign{}
\endlastfoot
1 & Probability theory & Modeling unpredictable events \\
2 & Distributions & Recognizing shapes of variation \\
3 & Statistical inference & Drawing conclusions from valid samples \\
4 & Data analysis & Finding signals in noise \\
\end{longtable}

\chapter{Probability theory}\label{probability-theory}

Reasoning precisely about uncertainty

\hfill\break

A medical test is 99\% accurate. You test positive. How worried should
you be?

Most people say: very --- 99\% accurate means there's only a 1\% chance
this is wrong. But the actual answer depends almost entirely on
something the test result doesn't tell you: how common the disease is in
the first place. If only one person in a thousand has the disease, a
positive result is more likely to be a false alarm than a true positive,
even with a 99\% accurate test. The maths will show you exactly why ---
and it's not complicated once you have the right framework.

Here's a second trap. You toss a fair coin five times and get five heads
in a row. What's the probability the next toss is tails --- higher than
usual, to ``balance things out''? No.~It's still 50\%. The coin has no
memory. Each toss is independent of every toss before it. The belief
that a run of heads makes tails ``overdue'' is called the gambler's
fallacy, and it has emptied a great many wallets.

Third trap. Two events each have probability 0.5. What's the probability
they both happen? Surely 0.25 --- half of half? Only if the events have
nothing to do with each other. If knowing that one happened changes the
probability of the other, the calculation is different.

Probability has real structure. Intuition is not enough to navigate it.
This chapter gives you the rules that make the structure precise --- and
with those rules, the three problems above all have clean, computable
answers.

\section{What the notation is
saying}\label{what-the-notation-is-saying-21}

\subsection{Sample spaces and events}\label{sample-spaces-and-events}

Before you can write \(P(\text{anything})\), you need to know what the
full set of possibilities is.

The \textbf{sample space} \(\Omega\) (the capital Greek letter omega) is
the set of all possible outcomes of an experiment. Roll a standard die:
\(\Omega = \{1,
2, 3, 4, 5, 6\}\). Toss a coin: \(\Omega = \{\text{H}, \text{T}\}\).
These are the complete lists --- every outcome that could happen.

An \textbf{event} is any subset of the sample space --- any collection
of outcomes you're interested in. ``Roll an even number'' is the event
\(\{2, 4, 6\}\). ``Roll a number greater than 4'' is \(\{5, 6\}\).
Events are usually labelled with capital letters: \(A\), \(B\), \(C\).

Now the five symbols:

\textbf{\(P(A)\)} --- read ``the probability of \(A\).'' This is a
number between 0 and 1 that measures how likely event \(A\) is.
\(P(A) = 0\) means impossible; \(P(A) = 1\) means certain.

\emph{Example:} Let \(A\) = ``roll a 3.'' On a fair die,
\(P(A) = \frac{1}{6}\).

\textbf{\(P(A^c)\)} --- read ``the probability of the complement of
\(A\).'' The complement \(A^c\) (also written \(\bar{A}\) or \(A'\)) is
the event ``A does not happen'' --- every outcome in \(\Omega\) that is
not in \(A\).

\emph{Example:} Let \(A\) = ``roll a 3.'' Then \(A^c\) = ``roll anything
but 3'' = \(\{1, 2, 4, 5, 6\}\). Since one of the two must happen:
\(P(A^c) = \frac{5}{6}\).

\textbf{\(P(A \cup B)\)} --- read ``the probability of \(A\) union
\(B\).'' The union \(A \cup B\) is the event ``\(A\) or \(B\) or both''
--- any outcome in at least one of \(A\) and \(B\).

\emph{Example:} Let \(A\) = ``roll an even number'' = \(\{2, 4, 6\}\)
and \(B\) = ``roll a number greater than 4'' = \(\{5, 6\}\). Then
\(A \cup B\) = \(\{2, 4, 5, 6\}\), so
\(P(A \cup B) = \frac{4}{6} = \frac{2}{3}\).

\textbf{\(P(A \cap B)\)} --- read ``the probability of \(A\)
intersection \(B\).'' The intersection \(A \cap B\) is the event ``both
\(A\) and \(B\)'' --- only the outcomes that are in \(A\) and also in
\(B\).

\emph{Example:} With \(A = \{2, 4, 6\}\) and \(B = \{5, 6\}\):
\(A \cap B = \{6\}\), so \(P(A \cap B) = \frac{1}{6}\).

\textbf{\(P(A \mid B)\)} --- read ``the probability of \(A\) given
\(B\).'' This is the conditional probability --- the probability that
\(A\) happens, given that you already know \(B\) happened. The vertical
bar is read ``given.''

\emph{Example:} You roll a die and someone tells you the result was
greater than 4 (so \(B\) happened, and you know the outcome is in
\(\{5, 6\}\)). What's the probability it was a 6? You're now working in
a smaller world --- only two outcomes are possible, and one of them is
6. So \(P(\text{roll 6} \mid B)
= \frac{1}{2}\).

Those five symbols are the entire vocabulary you need. Everything else
in this chapter is built from combinations of them.

\section{The rules}\label{the-rules}

\subsection{Rule 1: Complement}\label{rule-1-complement}

\[P(A^c) = 1 - P(A)\]

Read it as: the probability that \(A\) does not happen equals one minus
the probability that it does.

This is simple, but it earns its place in every calculation where it's
easier to ask ``what are the chances this doesn't happen?'' than ``what
are the chances it does?'' --- which turns out to be surprisingly often.

\emph{Example:} What's the probability of rolling at least one six when
you roll a die four times? Computing directly means listing every
combination that includes at least one six. Computing via the complement
means one step:
\(P(\text{at least one six}) = 1 - P(\text{no sixes}) = 1 - (5/6)^4
\approx 1 - 0.482 = 0.518\) (we're getting ahead slightly --- since each
roll is independent, P(no six on all four rolls) =
(5/6)×(5/6)×(5/6)×(5/6) = (5/6)⁴; we'll see this rule formally in Rule 4
below).

\subsection{Rule 2: Addition rule}\label{rule-2-addition-rule}

\[P(A \cup B) = P(A) + P(B) - P(A \cap B)\]

Read it as: the probability of \(A\) or \(B\) equals the probability of
\(A\) plus the probability of \(B\), minus the probability of both.

Why subtract the intersection? Because if you add \(P(A)\) and \(P(B)\)
separately, any outcome that falls in both \(A\) and \(B\) gets counted
twice --- once for \(A\) and once for \(B\). Subtracting \(P(A \cap B)\)
corrects that double-count.

\emph{Example:} Back to the die. \(P(\text{even}) = \frac{3}{6}\),
\(P(\text{greater than 4}) = \frac{2}{6}\),
\(P(\text{even and greater than 4}) = P(\{6\}) = \frac{1}{6}\).

\[P(\text{even or greater than 4}) = \frac{3}{6} + \frac{2}{6} - \frac{1}{6} = \frac{4}{6} = \frac{2}{3}\]

\textbf{Special case --- mutually exclusive events:} If \(A\) and \(B\)
cannot both happen at the same time --- if \(A \cap B\) is empty ---
then \(P(A \cap B) = 0\), and the rule simplifies to:

\[P(A \cup B) = P(A) + P(B) \qquad (\text{mutually exclusive only})\]

Rolling a 2 and rolling a 5 on the same single die toss are mutually
exclusive. Rolling an even and rolling a 6 are not --- 6 satisfies both.

\subsection{Rule 3: Conditional
probability}\label{rule-3-conditional-probability}

\[P(A \mid B) = \frac{P(A \cap B)}{P(B)}\]

Read it as: the probability of \(A\) given \(B\) equals the probability
of both divided by the probability of \(B\).

The intuition: when you learn that \(B\) happened, you've narrowed down
the sample space. You're no longer considering all of \(\Omega\) ---
only the part of \(\Omega\) where \(B\) is true. Within that restricted
space, how often does \(A\) also happen? That fraction is
\(P(A \mid B)\).

A concrete picture makes this clearer than any formula. Imagine a group
of 1000 people:

\begin{longtable}[]{@{}llll@{}}
\toprule\noalign{}
& Disease positive & Disease negative & Total \\
\midrule\noalign{}
\endhead
\bottomrule\noalign{}
\endlastfoot
\textbf{Test positive} & 95 & 49 & 144 \\
\textbf{Test negative} & 5 & 851 & 856 \\
\textbf{Total} & 100 & 900 & 1000 \\
\end{longtable}

Here disease prevalence is 10\% (100 of 1000 have the disease), test
sensitivity is 95\% (catches 95 of the 100 cases), and specificity (the
probability of a negative result given no disease) is roughly 94.6\%
(misses 49 of 900 healthy people).

If you test positive, you're in the top row: 144 people. Of those 144,
only 95 actually have the disease.
\(P(\text{disease} \mid \text{positive}) = \frac{95}{144} \approx 0.66\).

But what happens when the disease is rarer? That's what the interactive
below explores --- and the numbers become surprising.

\begin{Shaded}
\begin{Highlighting}[]
\NormalTok{//| echo: false}
\NormalTok{\{}
\NormalTok{  // Sliders}
\NormalTok{  const popSlider = Inputs.range([100, 10000], \{value: 1000, step: 100, label: "Population size"\});}
\NormalTok{  const prevSlider = Inputs.range([0.1, 20], \{value: 10, step: 0.1, label: "Disease prevalence (\%)"\});}
\NormalTok{  const sensSlider = Inputs.range([50, 99.9], \{value: 95, step: 0.1, label: "Test sensitivity (\%)"\});}
\NormalTok{  const specSlider = Inputs.range([50, 99.9], \{value: 94.6, step: 0.1, label: "Test specificity (\%)"\});}

\NormalTok{  const container = document.createElement("div");}
\NormalTok{  container.style.cssText = "font{-}family:inherit; max{-}width:640px;";}

\NormalTok{  const controls = document.createElement("div");}
\NormalTok{  controls.style.cssText = "display:grid; gap:0.5rem; margin{-}bottom:1.25rem;";}
\NormalTok{  controls.appendChild(popSlider);}
\NormalTok{  controls.appendChild(prevSlider);}
\NormalTok{  controls.appendChild(sensSlider);}
\NormalTok{  controls.appendChild(specSlider);}

\NormalTok{  const tableWrap = document.createElement("div");}
\NormalTok{  tableWrap.style.cssText = "margin{-}bottom:1rem;";}

\NormalTok{  const ppvDisplay = document.createElement("div");}
\NormalTok{  ppvDisplay.style.cssText = "font{-}size:1.35rem; font{-}weight:600; margin{-}bottom:0.5rem; padding:0.75rem 1rem; background:\#f0fdf4; border:1px solid \#bbf7d0; border{-}radius:6px; color:\#15803d;";}

\NormalTok{  const noteEl = document.createElement("div");}
\NormalTok{  noteEl.style.cssText = "font{-}size:0.88em; color:\#b45309; background:\#fffbeb; border:1px solid \#fde68a; border{-}radius:4px; padding:0.5rem 0.75rem; margin{-}top:0.5rem; display:none;";}
\NormalTok{  noteEl.textContent = "Low prevalence warning: when the disease is rare, even a good test produces more false positives than true positives. A positive result is not a diagnosis.";}

\NormalTok{  function update() \{}
\NormalTok{    const N   = popSlider.value;}
\NormalTok{    const prev = prevSlider.value / 100;}
\NormalTok{    const sens = sensSlider.value / 100;}
\NormalTok{    const spec = specSlider.value / 100;}

\NormalTok{    const diseased   = Math.round(N * prev);}
\NormalTok{    const healthy    = N {-} diseased;}
\NormalTok{    const tp = Math.round(diseased * sens);}
\NormalTok{    const fn = diseased {-} tp;}
\NormalTok{    const tn = Math.round(healthy * spec);}
\NormalTok{    const fp = healthy {-} tn;}

\NormalTok{    const testPos = tp + fp;}
\NormalTok{    const testNeg = fn + tn;}
\NormalTok{    const ppv = testPos \textgreater{} 0 ? tp / testPos : 0;}

\NormalTok{    ppvDisplay.textContent = \textasciigrave{}P(disease | positive test) = $\{(ppv * 100).toFixed(1)\}\%  ($\{tp\} true positives out of $\{testPos\} positive tests)\textasciigrave{};}
\NormalTok{    noteEl.style.display = prev \textless{} 0.05 ? "block" : "none";}

\NormalTok{    tableWrap.innerHTML = "";}
\NormalTok{    const tbl = document.createElement("table");}
\NormalTok{    tbl.style.cssText = "border{-}collapse:collapse; width:100\%; font{-}size:0.92em;";}

\NormalTok{    const headerRow = tbl.insertRow();}
\NormalTok{    ["", "Test positive", "Test negative", "Total"].forEach((h, i) =\textgreater{} \{}
\NormalTok{      const th = document.createElement("th");}
\NormalTok{      th.textContent = h;}
\NormalTok{      th.style.cssText = "padding:0.4rem 0.75rem; border:1px solid \#d1d5db; background:\#f9fafb; text{-}align:" + (i === 0 ? "left" : "right") + ";";}
\NormalTok{      headerRow.appendChild(th);}
\NormalTok{    \});}

\NormalTok{    const rows = [}
\NormalTok{      ["Disease positive", tp, fn, diseased],}
\NormalTok{      ["Disease negative", fp, tn, healthy],}
\NormalTok{      ["Total", testPos, testNeg, N],}
\NormalTok{    ];}
\NormalTok{    rows.forEach((r, ri) =\textgreater{} \{}
\NormalTok{      const tr = tbl.insertRow();}
\NormalTok{      r.forEach((cell, ci) =\textgreater{} \{}
\NormalTok{        const td = tr.insertCell();}
\NormalTok{        td.textContent = cell;}
\NormalTok{        const isTp = ri === 0 \&\& ci === 1;}
\NormalTok{        td.style.cssText = "padding:0.4rem 0.75rem; border:1px solid \#d1d5db; text{-}align:" + (ci === 0 ? "left" : "right") + ";" + (isTp ? " background:\#dcfce7; font{-}weight:600;" : "");}
\NormalTok{      \});}
\NormalTok{    \});}

\NormalTok{    tableWrap.appendChild(tbl);}
\NormalTok{  \}}

\NormalTok{  popSlider.addEventListener("input", update);}
\NormalTok{  prevSlider.addEventListener("input", update);}
\NormalTok{  sensSlider.addEventListener("input", update);}
\NormalTok{  specSlider.addEventListener("input", update);}

\NormalTok{  update();}

\NormalTok{  container.appendChild(controls);}
\NormalTok{  container.appendChild(ppvDisplay);}
\NormalTok{  container.appendChild(tableWrap);}
\NormalTok{  container.appendChild(noteEl);}
\NormalTok{  return container;}
\NormalTok{\}}
\end{Highlighting}
\end{Shaded}

\subsection{Rule 4: Multiplication
rule}\label{rule-4-multiplication-rule}

\[P(A \cap B) = P(A \mid B) \cdot P(B)\]

This is just the conditional probability definition rearranged. Solve
for \(P(A \cap B)\) and you get: to find the probability that both \(A\)
and \(B\) happen, multiply the probability that \(B\) happens by the
probability that \(A\) happens given \(B\) already happened.

\emph{Example:} A bag has 5 red and 3 blue marbles. You draw one without
looking, then draw a second without replacing the first. What's the
probability both are red?

\begin{itemize}
\tightlist
\item
  \(P(\text{first is red}) = \frac{5}{8}\)
\item
  \(P(\text{second is red} \mid \text{first was red}) = \frac{4}{7}\)
  (only 4 red left, only 7 total left)
\item
  \(P(\text{both red}) = \frac{5}{8} \times \frac{4}{7} = \frac{20}{56} = \frac{5}{14}\)
\end{itemize}

\subsection{Rule 5: Independence}\label{rule-5-independence}

Events \(A\) and \(B\) are \textbf{independent} if:

\[P(A \mid B) = P(A)\]

Read it as: knowing that \(B\) happened does not change the probability
of \(A\) at all. \(B\) carries no information about \(A\).

An equivalent (and often more useful) form comes directly from
substituting into the multiplication rule:

\[P(A \cap B) = P(A) \cdot P(B) \qquad (\text{independent events only})\]

This is the definition of independence --- not something you prove from
intuition, but a precise property you can check. Two coin tosses are
independent. Drawing a card from a shuffled deck and then rolling a die
are independent. Drawing two cards without replacement are not --- what
you get first changes what's available for the second.

The gambler's fallacy is exactly the mistake of believing that
independent events are somehow dependent --- that after five heads, the
coin ``owes'' you a tail. The coin's next outcome has probability
\(\frac{1}{2}\) regardless of any previous results, because each toss is
independent of every other.

\begin{Shaded}
\begin{Highlighting}[]
\NormalTok{//| echo: false}
\NormalTok{\{}
\NormalTok{  const pASlider = Inputs.range([0.01, 0.99], \{value: 0.5, step: 0.01, label: "P(A)"\});}
\NormalTok{  const pBSlider = Inputs.range([0.01, 0.99], \{value: 0.5, step: 0.01, label: "P(B)"\});}
\NormalTok{  const corrSelect = Inputs.radio(["Independent", "Positively correlated", "Negatively correlated"], \{value: "Independent", label: "Relationship"\});}

\NormalTok{  const container = document.createElement("div");}
\NormalTok{  container.style.cssText = "font{-}family:inherit; max{-}width:640px;";}

\NormalTok{  const controls = document.createElement("div");}
\NormalTok{  controls.style.cssText = "display:grid; gap:0.5rem; margin{-}bottom:1.25rem;";}
\NormalTok{  controls.appendChild(pASlider);}
\NormalTok{  controls.appendChild(pBSlider);}
\NormalTok{  controls.appendChild(corrSelect);}

\NormalTok{  const summary = document.createElement("div");}
\NormalTok{  summary.style.cssText = "margin{-}bottom:1rem;";}

\NormalTok{  const gridWrap = document.createElement("div");}
\NormalTok{  gridWrap.style.cssText = "margin{-}bottom:0.5rem;";}

\NormalTok{  function clamp(v, lo, hi) \{ return Math.max(lo, Math.min(hi, v)); \}}

\NormalTok{  function jointProbs(pA, pB, mode) \{}
\NormalTok{    let pAB;}
\NormalTok{    if (mode === "Independent") \{}
\NormalTok{      pAB = pA * pB;}
\NormalTok{    \} else if (mode === "Positively correlated") \{}
\NormalTok{      const maxAB = Math.min(pA, pB);}
\NormalTok{      pAB = clamp(pA * pB + 0.15 * Math.min(pA, pB, 1 {-} pA, 1 {-} pB), 0, maxAB);}
\NormalTok{    \} else \{}
\NormalTok{      const minAB = Math.max(0, pA + pB {-} 1);}
\NormalTok{      pAB = clamp(pA * pB {-} 0.15 * Math.min(pA, pB, 1 {-} pA, 1 {-} pB), minAB, Math.min(pA, pB));}
\NormalTok{    \}}
\NormalTok{    return \{}
\NormalTok{      pAB:      pAB,}
\NormalTok{      pAnotB:   pA {-} pAB,}
\NormalTok{      pBnotA:   pB {-} pAB,}
\NormalTok{      pNeither: 1 {-} pA {-} pB + pAB,}
\NormalTok{    \};}
\NormalTok{  \}}

\NormalTok{  function fmtPct(v) \{ return (v * 100).toFixed(1) + "\%"; \}}

\NormalTok{  function renderGrid(pA, pB, mode) \{}
\NormalTok{    const j = jointProbs(pA, pB, mode);}
\NormalTok{    const indepAB = pA * pB;}

\NormalTok{    summary.innerHTML = "";}
\NormalTok{    const cases = [}
\NormalTok{      \{ label: "Independent",            pab: indepAB \},}
\NormalTok{      \{ label: "Positively correlated",  pab: jointProbs(pA, pB, "Positively correlated").pAB \},}
\NormalTok{      \{ label: "Negatively correlated",  pab: jointProbs(pA, pB, "Negatively correlated").pAB \},}
\NormalTok{    ];}
\NormalTok{    const sumDiv = document.createElement("div");}
\NormalTok{    sumDiv.style.cssText = "display:grid; grid{-}template{-}columns:repeat(3,1fr); gap:0.5rem; margin{-}bottom:1rem;";}
\NormalTok{    cases.forEach(c =\textgreater{} \{}
\NormalTok{      const card = document.createElement("div");}
\NormalTok{      const isActive = c.label === mode;}
\NormalTok{      card.style.cssText = "border:1px solid " + (isActive ? "\#3b82f6" : "\#e5e7eb") + "; border{-}radius:6px; padding:0.5rem 0.75rem; background:" + (isActive ? "\#eff6ff" : "\#f9fafb") + ";";}
\NormalTok{      card.innerHTML = \textasciigrave{}\textless{}div style="font{-}size:0.8em;color:\#6b7280;"\textgreater{}$\{c.label\}\textless{}/div\textgreater{}\textless{}div style="font{-}weight:600;font{-}size:1.05em;"\textgreater{}P(A∩B) = $\{fmtPct(c.pab)\}\textless{}/div\textgreater{}\textasciigrave{};}
\NormalTok{      sumDiv.appendChild(card);}
\NormalTok{    \});}
\NormalTok{    summary.appendChild(sumDiv);}

\NormalTok{    gridWrap.innerHTML = "";}
\NormalTok{    const tbl = document.createElement("table");}
\NormalTok{    tbl.style.cssText = "border{-}collapse:collapse; width:100\%; font{-}size:0.9em;";}

\NormalTok{    const hdr = tbl.insertRow();}
\NormalTok{    ["", "B occurs", "B does not occur", "Total"].forEach((h, i) =\textgreater{} \{}
\NormalTok{      const th = document.createElement("th");}
\NormalTok{      th.textContent = h;}
\NormalTok{      th.style.cssText = "padding:0.4rem 0.7rem; border:1px solid \#d1d5db; background:\#f9fafb; text{-}align:" + (i === 0 ? "left" : "right") + ";";}
\NormalTok{      hdr.appendChild(th);}
\NormalTok{    \});}

\NormalTok{    const dataRows = [}
\NormalTok{      ["A occurs",         j.pAB,      j.pAnotB,   pA],}
\NormalTok{      ["A does not occur", j.pBnotA,   j.pNeither, 1 {-} pA],}
\NormalTok{      ["Total",            pB,         1 {-} pB,     1.0],}
\NormalTok{    ];}
\NormalTok{    dataRows.forEach(r =\textgreater{} \{}
\NormalTok{      const tr = tbl.insertRow();}
\NormalTok{      r.forEach((cell, ci) =\textgreater{} \{}
\NormalTok{        const td = tr.insertCell();}
\NormalTok{        td.textContent = ci === 0 ? cell : fmtPct(cell);}
\NormalTok{        td.style.cssText = "padding:0.4rem 0.7rem; border:1px solid \#d1d5db; text{-}align:" + (ci === 0 ? "left" : "right") + ";";}
\NormalTok{      \});}
\NormalTok{    \});}

\NormalTok{    gridWrap.appendChild(tbl);}

\NormalTok{    const note = document.createElement("div");}
\NormalTok{    note.style.cssText = "font{-}size:0.85em; color:\#6b7280; margin{-}top:0.5rem;";}
\NormalTok{    const diff = j.pAB {-} indepAB;}
\NormalTok{    if (Math.abs(diff) \textless{} 0.001) \{}
\NormalTok{      note.textContent = "The joint probability equals P(A) × P(B) = " + fmtPct(indepAB) + " — the events are independent.";}
\NormalTok{    \} else if (diff \textgreater{} 0) \{}
\NormalTok{      note.textContent = "P(A∩B) = " + fmtPct(j.pAB) + " \textgreater{} P(A)×P(B) = " + fmtPct(indepAB) + " — knowing A occurred makes B more likely.";}
\NormalTok{    \} else \{}
\NormalTok{      note.textContent = "P(A∩B) = " + fmtPct(j.pAB) + " \textless{} P(A)×P(B) = " + fmtPct(indepAB) + " — knowing A occurred makes B less likely.";}
\NormalTok{    \}}
\NormalTok{    gridWrap.appendChild(note);}
\NormalTok{  \}}

\NormalTok{  function update() \{}
\NormalTok{    renderGrid(pASlider.value, pBSlider.value, corrSelect.value);}
\NormalTok{  \}}

\NormalTok{  pASlider.addEventListener("input", update);}
\NormalTok{  pBSlider.addEventListener("input", update);}
\NormalTok{  corrSelect.addEventListener("change", update);}

\NormalTok{  update();}

\NormalTok{  container.appendChild(controls);}
\NormalTok{  container.appendChild(summary);}
\NormalTok{  container.appendChild(gridWrap);}
\NormalTok{  return container;}
\NormalTok{\}}
\end{Highlighting}
\end{Shaded}

\subsection{Bayes' theorem}\label{bayes-theorem}

Everything so far has been building to this. Start with the
multiplication rule written two ways:

\[P(A \cap B) = P(A \mid B) \cdot P(B)\]
\[P(A \cap B) = P(B \mid A) \cdot P(A)\]

The left-hand sides are equal, so the right-hand sides are too. Divide
both sides by \(P(B)\):

\[P(A \mid B) = \frac{P(B \mid A) \cdot P(A)}{P(B)}\]

This is \textbf{Bayes' theorem}. Read each part:

\begin{itemize}
\tightlist
\item
  \(P(A)\) --- the \textbf{prior}: what you believed about \(A\) before
  any new information.
\item
  \(P(B \mid A)\) --- the \textbf{likelihood}: how probable was this new
  evidence \(B\), assuming \(A\) is true?
\item
  \(P(A \mid B)\) --- the \textbf{posterior}: what you now believe about
  \(A\), having seen \(B\).
\item
  \(P(B)\) --- the \textbf{normalising constant}: the total probability
  of seeing \(B\), regardless of whether \(A\) is true or not.
\end{itemize}

The theorem tells you how to update a belief in the light of new
evidence. That is, precisely, what reasoning under uncertainty requires.

Now the medical test problem from the opening has a clean solution.

\section{Worked examples}\label{worked-examples-21}

\subsection{Example 1 --- Medical testing
(sciences)}\label{example-1-medical-testing-sciences}

A disease affects 1\% of the population. A test for it has sensitivity
95\% (correctly identifies 95\% of people who have the disease) and
specificity 95\% (correctly identifies 95\% of people who don't have the
disease).

You test positive. What is the probability you actually have the
disease?

\textbf{Set up the events:}

\begin{itemize}
\tightlist
\item
  \(A\) = ``you have the disease''
\item
  \(B\) = ``you test positive''
\end{itemize}

\textbf{What you know:}

\begin{itemize}
\tightlist
\item
  \(P(A) = 0.01\) --- disease prevalence (the prior)
\item
  \(P(B \mid A) = 0.95\) --- sensitivity (probability of a positive test
  given disease)
\item
  \(P(B \mid A^c) = 0.05\) --- false positive rate (probability of a
  positive test given no disease; this is \(1 - \text{specificity}\))
\end{itemize}

\textbf{Find \(P(B)\)} --- to find P(positive test) we use the
\textbf{total probability law}: since every person either has the
disease or doesn't, we can compute P(positive) by considering both
groups separately and adding. Split by whether you have the disease or
not:

\[P(B) = P(B \mid A) \cdot P(A) + P(B \mid A^c) \cdot P(A^c)\]
\[P(B) = (0.95)(0.01) + (0.05)(0.99) = 0.0095 + 0.0495 = 0.059\]

\textbf{Apply Bayes' theorem:}

\[P(A \mid B) = \frac{P(B \mid A) \cdot P(A)}{P(B)} = \frac{0.95 \times 0.01}{0.059} = \frac{0.0095}{0.059} \approx 0.161\]

\textbf{The result:} If you test positive for a disease that affects 1\%
of the population, even with a 95\% accurate test, the probability you
actually have it is about 16\%.

Why so low? Because the disease is rare. In a population of 10 000
people: 100 have the disease and 9900 do not. The test correctly flags
about 95 of the 100 sick people. But it also incorrectly flags about 495
of the 9900 healthy people. Of the 590 positive results in total, only
95 are true positives --- about 16\%.

A positive result is not a diagnosis. It's information that should
prompt further testing. This is why medical screening programmes track
\emph{positive predictive value} --- the probability a positive test
reflects a true positive --- and not just test accuracy.

\begin{center}\rule{0.5\linewidth}{0.5pt}\end{center}

\subsection{Example 2 --- Cards without replacement
(finance)}\label{example-2-cards-without-replacement-finance}

A standard deck has 52 cards, including 4 aces. You draw two cards
without replacing the first. What is the probability that both are aces?

\textbf{Using the multiplication rule for dependent events:}

\[P(\text{both aces}) = P(\text{first is ace}) \times P(\text{second is ace} \mid \text{first was ace})\]

\[= \frac{4}{52} \times \frac{3}{51} = \frac{12}{2652} = \frac{1}{221} \approx 0.0045\]

About 0.45\% --- less than half a percent.

Note the dependence: the second draw depends on the first because the
first card is not returned to the deck. If you replaced the card between
draws (sampling with replacement), the draws would be independent:
\(P(\text{both aces}) = \frac{4}{52} \times \frac{4}{52} = \frac{16}{2704}
\approx 0.0059\) --- a slightly higher probability because you haven't
reduced the pool of aces.

This distinction between sampling with and without replacement appears
constantly in actuarial and financial modelling, where the question is
whether one event (e.g.~one bond defaulting) affects the probability of
another.

\begin{center}\rule{0.5\linewidth}{0.5pt}\end{center}

\subsection{Example 3 --- Spam filter
(computing)}\label{example-3-spam-filter-computing}

A word appears in 80\% of spam emails and in 10\% of legitimate emails.
Overall, 30\% of incoming email is spam.

A new message contains this word. What is the probability it is spam?

\textbf{Set up the events:}

\begin{itemize}
\tightlist
\item
  \(A\) = ``email is spam''
\item
  \(B\) = ``email contains the word''
\end{itemize}

\textbf{What you know:}

\begin{itemize}
\tightlist
\item
  \(P(A) = 0.30\) --- prior probability that email is spam
\item
  \(P(B \mid A) = 0.80\) --- probability the word appears in spam
\item
  \(P(B \mid A^c) = 0.10\) --- probability the word appears in
  legitimate email
\end{itemize}

\textbf{Find \(P(B)\):}

\[P(B) = P(B \mid A) \cdot P(A) + P(B \mid A^c) \cdot P(A^c)\]
\[= (0.80)(0.30) + (0.10)(0.70) = 0.24 + 0.07 = 0.31\]

\textbf{Apply Bayes' theorem:}

\[P(A \mid B) = \frac{P(B \mid A) \cdot P(A)}{P(B)} = \frac{0.80 \times 0.30}{0.31} = \frac{0.24}{0.31} \approx 0.774\]

\textbf{The result:} Given that the email contains this word, the
probability it is spam is about 77\%.

A real spam filter does this for dozens or hundreds of words
simultaneously, multiplying the contributions of each word (under an
independence assumption that is not strictly true but works well in
practice). That extension is called a Naive Bayes classifier --- and it
was one of the most effective spam filters ever built.

\begin{center}\rule{0.5\linewidth}{0.5pt}\end{center}

\subsection{Example 4 --- Weather (geography and
environment)}\label{example-4-weather-geography-and-environment}

A weather station records the following over a year:

\begin{itemize}
\tightlist
\item
  20\% of days are rainy.
\item
  40\% of days have cloudy mornings.
\item
  Of all rainy days, 70\% were preceded by a cloudy morning.
\end{itemize}

Given that today's morning is cloudy, what is the probability of rain?

\textbf{Set up the events:}

\begin{itemize}
\tightlist
\item
  \(A\) = ``today is rainy''
\item
  \(B\) = ``this morning is cloudy''
\end{itemize}

\textbf{What you know:}

\begin{itemize}
\tightlist
\item
  \(P(A) = 0.20\)
\item
  \(P(B) = 0.40\)
\item
  \(P(B \mid A) = 0.70\)
\end{itemize}

\textbf{Apply Bayes' theorem directly:}

\[P(A \mid B) = \frac{P(B \mid A) \cdot P(A)}{P(B)} = \frac{0.70 \times 0.20}{0.40} = \frac{0.14}{0.40} = 0.35\]

\textbf{The result:} A cloudy morning raises the probability of rain
from the baseline 20\% to 35\%.

Environmental scientists and meteorologists work with probability
statements like this constantly --- flood return periods, drought
likelihoods, and extreme precipitation events are all expressed as
conditional probabilities, updated as new observations come in.

\section{Where this goes}\label{where-this-goes-22}

\subsection{Probability distributions}\label{probability-distributions}

You now have a framework for reasoning about individual events and pairs
of events. The next natural question is: what if a random process can
produce many different numerical outcomes --- not just ``heads/tails''
or ``disease/no disease'' but any integer from 0 to 20, or any real
number in an interval? How do probabilities distribute across all those
possible values?

That question leads directly to \textbf{probability distributions}
(chapter 2). A distribution is a complete description of the
probabilities attached to every possible outcome of a random quantity.
The binomial distribution (counting successes in repeated trials), the
normal distribution (the bell curve that appears everywhere in nature
and statistics), and others all build on the conditional probability and
independence machinery you've just learned.

\subsection{Bayesian reasoning in computing and machine
learning}\label{bayesian-reasoning-in-computing-and-machine-learning}

The spam filter example above is a stripped-down version of a
\textbf{Naive Bayes classifier} --- one of the fundamental algorithms in
machine learning. The same structure extends to any classification
problem: given a collection of observed features, compute the
probability that the observation belongs to each possible class, and
predict the most probable one.

Medical diagnosis, document classification, image labelling, and fraud
detection all use variants of Bayesian reasoning. The conditional
probability machinery from this chapter is the mathematical core of each
of them. What changes when you move to a full ML course is scale
(thousands of features, not one word) and the methods for estimating the
prior and likelihood from data.

\section{Applications}\label{applications-2}

\begin{tcolorbox}[enhanced jigsaw, opacityback=0, colback=white, breakable, colframe=quarto-callout-note-color-frame, arc=.35mm, left=2mm, toprule=.15mm, leftrule=.75mm, bottomrule=.15mm, rightrule=.15mm]

\textbf{Where this chapter turns up in practice}

\begin{itemize}
\item
  \textbf{Disease screening programmes} --- a test's positive predictive
  value depends on prevalence, not just accuracy. Screening for rare
  diseases requires high specificity to avoid swamping clinicians with
  false positives.
\item
  \textbf{Forensic evidence} --- DNA evidence in court is often misread.
  The probability that a match occurs by chance (1 in a billion) is not
  the same as the probability that the accused is innocent given the
  match. Getting this wrong is called the prosecutor's fallacy, and it
  has contributed to wrongful convictions.
\item
  \textbf{Spam and content filters} --- Naive Bayes classifiers were the
  first practical spam filters and remain in use today. The same
  structure underlies content moderation systems.
\item
  \textbf{Flood return periods} --- a ``1-in-100-year flood'' means a
  flood of that size has a 1\% probability of occurring in any given
  year --- not that you get one every hundred years and can relax
  afterward. It can occur two years in a row.
\end{itemize}

\end{tcolorbox}

\section{Exercises}\label{exercises-22}

\begin{Shaded}
\begin{Highlighting}[]
\NormalTok{//| echo: false}
\NormalTok{katex = require("katex@0.16")}
\end{Highlighting}
\end{Shaded}

\begin{Shaded}
\begin{Highlighting}[]
\NormalTok{//| echo: false}
\NormalTok{function makeStepperHTML(exerciseNum, steps) \{}
\NormalTok{  let current = 0;}
\NormalTok{  const totalSteps = steps.length;}
\NormalTok{  const container = document.createElement("div");}
\NormalTok{  container.style.cssText = "border:1px solid \#e5e7eb; border{-}radius:8px; padding:1rem 1.25rem; margin:0.75rem 0; font{-}family:inherit;";}
\NormalTok{  const stepsDiv = document.createElement("div");}
\NormalTok{  stepsDiv.style.marginBottom = "0.75rem";}
\NormalTok{  function renderSteps() \{}
\NormalTok{    stepsDiv.innerHTML = "";}
\NormalTok{    for (let i = 0; i \textless{} current; i++) \{}
\NormalTok{      const s = steps[i];}
\NormalTok{      const row = document.createElement("div");}
\NormalTok{      row.style.cssText = "display:grid; grid{-}template{-}columns:200px 1fr; gap:0.5rem 1rem; align{-}items:baseline; padding:0.35rem 0; border{-}top:1px solid \#f3f4f6;";}
\NormalTok{      const opCell = document.createElement("span");}
\NormalTok{      opCell.style.cssText = "font{-}size:0.85em; color:\#6b7280; font{-}style:italic;";}
\NormalTok{      opCell.textContent = s.op;}
\NormalTok{      const eqCell = document.createElement("span");}
\NormalTok{      eqCell.innerHTML = katex.renderToString(s.eq, \{ throwOnError: false, displayMode: false \});}
\NormalTok{      row.appendChild(opCell);}
\NormalTok{      row.appendChild(eqCell);}
\NormalTok{      if (s.note) \{}
\NormalTok{        const noteRow = document.createElement("div");}
\NormalTok{        noteRow.style.cssText = "grid{-}column:2; font{-}size:0.82em; color:\#6b7280; padding{-}bottom:0.2rem;";}
\NormalTok{        noteRow.textContent = s.note;}
\NormalTok{        row.appendChild(noteRow);}
\NormalTok{      \}}
\NormalTok{      stepsDiv.appendChild(row);}
\NormalTok{    \}}
\NormalTok{    if (current === totalSteps) \{}
\NormalTok{      const done = document.createElement("div");}
\NormalTok{      done.style.cssText = "margin{-}top:0.5rem; font{-}size:0.9em; color:\#059669; font{-}weight:500;";}
\NormalTok{      done.textContent = "✓ Solution complete";}
\NormalTok{      stepsDiv.appendChild(done);}
\NormalTok{    \}}
\NormalTok{  \}}
\NormalTok{  const controls = document.createElement("div");}
\NormalTok{  controls.style.cssText = "display:flex; gap:0.5rem; align{-}items:center;";}
\NormalTok{  const nextBtn = document.createElement("button");}
\NormalTok{  nextBtn.textContent = "Next step →";}
\NormalTok{  nextBtn.style.cssText = "padding:0.35rem 0.85rem; border:1px solid \#d1d5db; border{-}radius:4px; background:\#fff; cursor:pointer; font{-}size:0.9em;";}
\NormalTok{  const resetBtn = document.createElement("button");}
\NormalTok{  resetBtn.textContent = "Reset";}
\NormalTok{  resetBtn.style.cssText = "padding:0.35rem 0.75rem; border:1px solid \#e5e7eb; border{-}radius:4px; background:\#f9fafb; cursor:pointer; font{-}size:0.9em; color:\#6b7280;";}
\NormalTok{  const counter = document.createElement("span");}
\NormalTok{  counter.style.cssText = "font{-}size:0.82em; color:\#9ca3af; margin{-}left:0.25rem;";}
\NormalTok{  function updateButtons() \{}
\NormalTok{    nextBtn.disabled = current === totalSteps;}
\NormalTok{    nextBtn.style.opacity = current === totalSteps ? "0.4" : "1";}
\NormalTok{    counter.textContent = current === 0 ? "Click to reveal steps" : \textasciigrave{}Step $\{current\} of $\{totalSteps\}\textasciigrave{};}
\NormalTok{  \}}
\NormalTok{  nextBtn.onclick = () =\textgreater{} \{ if (current \textless{} totalSteps) \{ current++; renderSteps(); updateButtons(); \} \};}
\NormalTok{  resetBtn.onclick = () =\textgreater{} \{ current = 0; renderSteps(); updateButtons(); \};}
\NormalTok{  controls.appendChild(nextBtn);}
\NormalTok{  controls.appendChild(resetBtn);}
\NormalTok{  controls.appendChild(counter);}
\NormalTok{  container.appendChild(stepsDiv);}
\NormalTok{  container.appendChild(controls);}
\NormalTok{  renderSteps();}
\NormalTok{  updateButtons();}
\NormalTok{  return container;}
\NormalTok{\}}
\end{Highlighting}
\end{Shaded}

\subsection{Exercise 1}\label{exercise-1}

The probability that a flight departs on time is 0.78. What is the
probability it does not depart on time?

\begin{Shaded}
\begin{Highlighting}[]
\NormalTok{//| echo: false}
\NormalTok{makeStepperHTML(1, [}
\NormalTok{  \{ op: "Identify the complement rule", eq: "P(\textbackslash{}\textbackslash{}text\{on time\}\^{}c) = 1 {-} P(\textbackslash{}\textbackslash{}text\{on time\})" \},}
\NormalTok{  \{ op: "Substitute", eq: "P(\textbackslash{}\textbackslash{}text\{not on time\}) = 1 {-} 0.78 = 0.22", note: "The probability the flight is delayed is 0.22." \}}
\NormalTok{])}
\end{Highlighting}
\end{Shaded}

\begin{center}\rule{0.5\linewidth}{0.5pt}\end{center}

\subsection{Exercise 2}\label{exercise-2}

In a group of students: 45\% play a sport, 30\% play a musical
instrument, and 15\% do both. What is the probability that a randomly
chosen student plays a sport or a musical instrument (or both)?

\begin{Shaded}
\begin{Highlighting}[]
\NormalTok{//| echo: false}
\NormalTok{makeStepperHTML(2, [}
\NormalTok{  \{ op: "Identify the addition rule", eq: "P(A \textbackslash{}\textbackslash{}cup B) = P(A) + P(B) {-} P(A \textbackslash{}\textbackslash{}cap B)" \},}
\NormalTok{  \{ op: "Label the events", eq: "P(A) = 0.45,\textbackslash{}\textbackslash{}; P(B) = 0.30,\textbackslash{}\textbackslash{}; P(A \textbackslash{}\textbackslash{}cap B) = 0.15", note: "A = plays sport, B = plays instrument." \},}
\NormalTok{  \{ op: "Substitute", eq: "P(A \textbackslash{}\textbackslash{}cup B) = 0.45 + 0.30 {-} 0.15 = 0.60", note: "The probability that a student plays a sport or an instrument is 0.60." \}}
\NormalTok{])}
\end{Highlighting}
\end{Shaded}

\begin{center}\rule{0.5\linewidth}{0.5pt}\end{center}

\subsection{Exercise 3}\label{exercise-3}

\emph{This exercise uses the full Bayes machinery --- prior, likelihood,
total probability, and posterior --- as in the medical test example
above. Work through each piece in turn.}

In a factory, 3\% of parts are defective. A quality inspector catches
90\% of defective parts and incorrectly flags 4\% of good parts as
defective.

A part is flagged by the inspector. What is the probability it is
actually defective?

\begin{Shaded}
\begin{Highlighting}[]
\NormalTok{//| echo: false}
\NormalTok{makeStepperHTML(3, [}
\NormalTok{  \{ op: "Set up events and knowns", eq: "P(A) = 0.03,\textbackslash{}\textbackslash{}; P(A\^{}c) = 0.97,\textbackslash{}\textbackslash{}; P(B \textbackslash{}\textbackslash{}mid A) = 0.90,\textbackslash{}\textbackslash{}; P(B \textbackslash{}\textbackslash{}mid A\^{}c) = 0.04", note: "A = part is defective, B = part is flagged." \},}
\NormalTok{  \{ op: "Total probability for P(B)", eq: "P(B) = (0.90)(0.03) + (0.04)(0.97) = 0.027 + 0.0388 = 0.0658" \},}
\NormalTok{  \{ op: "Apply Bayes\textquotesingle{} theorem", eq: "P(A \textbackslash{}\textbackslash{}mid B) = \textbackslash{}\textbackslash{}frac\{0.90 \textbackslash{}\textbackslash{}times 0.03\}\{0.0658\} = \textbackslash{}\textbackslash{}frac\{0.027\}\{0.0658\} \textbackslash{}\textbackslash{}approx 0.410", note: "A flagged part has about a 41\% chance of being genuinely defective. The other 59\% of flags are false alarms on good parts." \}}
\NormalTok{])}
\end{Highlighting}
\end{Shaded}

\begin{center}\rule{0.5\linewidth}{0.5pt}\end{center}

\subsection{Exercise 4}\label{exercise-4}

A bag contains 6 red and 4 blue counters. Two counters are drawn one at
a time without replacement. What is the probability of drawing a red
counter followed by a blue counter?

\begin{Shaded}
\begin{Highlighting}[]
\NormalTok{//| echo: false}
\NormalTok{makeStepperHTML(4, [}
\NormalTok{  \{ op: "Identify dependent draws", eq: "P(\textbackslash{}\textbackslash{}text\{red then blue\}) = P(\textbackslash{}\textbackslash{}text\{first red\}) \textbackslash{}\textbackslash{}times P(\textbackslash{}\textbackslash{}text\{second blue\} \textbackslash{}\textbackslash{}mid \textbackslash{}\textbackslash{}text\{first red\})", note: "The first counter is not replaced, so the draws are dependent." \},}
\NormalTok{  \{ op: "First draw probability", eq: "P(\textbackslash{}\textbackslash{}text\{first is red\}) = \textbackslash{}\textbackslash{}frac\{6\}\{10\} = \textbackslash{}\textbackslash{}frac\{3\}\{5\}" \},}
\NormalTok{  \{ op: "Conditional probability of second draw", eq: "P(\textbackslash{}\textbackslash{}text\{second is blue\} \textbackslash{}\textbackslash{}mid \textbackslash{}\textbackslash{}text\{first was red\}) = \textbackslash{}\textbackslash{}frac\{4\}\{9\}", note: "After drawing red: 9 counters remain, 4 of which are blue." \},}
\NormalTok{  \{ op: "Multiply", eq: "P(\textbackslash{}\textbackslash{}text\{red then blue\}) = \textbackslash{}\textbackslash{}frac\{3\}\{5\} \textbackslash{}\textbackslash{}times \textbackslash{}\textbackslash{}frac\{4\}\{9\} = \textbackslash{}\textbackslash{}frac\{12\}\{45\} = \textbackslash{}\textbackslash{}frac\{4\}\{15\} \textbackslash{}\textbackslash{}approx 0.267" \}}
\NormalTok{])}
\end{Highlighting}
\end{Shaded}

\begin{center}\rule{0.5\linewidth}{0.5pt}\end{center}

\subsection{Exercise 5}\label{exercise-5}

A card is drawn at random from a standard 52-card deck.

Event \(A\): the card is a heart. \(P(A) = \frac{13}{52} = \frac{1}{4}\)

Event \(B\): the card is a face card (jack, queen, or king).
\(P(B) = \frac{12}{52} = \frac{3}{13}\)

There are 3 face cards that are also hearts (jack, queen, and king of
hearts).

Are events \(A\) and \(B\) independent? Show your reasoning.

\begin{Shaded}
\begin{Highlighting}[]
\NormalTok{//| echo: false}
\NormalTok{makeStepperHTML(5, [}
\NormalTok{  \{ op: "State the independence criterion", eq: "\textbackslash{}\textbackslash{}text\{Independent iff \} P(A \textbackslash{}\textbackslash{}cap B) = P(A) \textbackslash{}\textbackslash{}cdot P(B)" \},}
\NormalTok{  \{ op: "Find P(A ∩ B)", eq: "P(A \textbackslash{}\textbackslash{}cap B) = \textbackslash{}\textbackslash{}frac\{3\}\{52\}", note: "A ∩ B = card is a heart and a face card = 3 cards (J, Q, K of hearts)." \},}
\NormalTok{  \{ op: "Compute P(A) · P(B)", eq: "P(A) \textbackslash{}\textbackslash{}cdot P(B) = \textbackslash{}\textbackslash{}frac\{1\}\{4\} \textbackslash{}\textbackslash{}times \textbackslash{}\textbackslash{}frac\{3\}\{13\} = \textbackslash{}\textbackslash{}frac\{3\}\{52\}" \},}
\NormalTok{  \{ op: "Compare", eq: "P(A \textbackslash{}\textbackslash{}cap B) = \textbackslash{}\textbackslash{}frac\{3\}\{52\} = P(A) \textbackslash{}\textbackslash{}cdot P(B) \textbackslash{}\textbackslash{}implies \textbackslash{}\textbackslash{}text\{independent\}", note: "Knowing a card is a heart tells you nothing about whether it is a face card, and vice versa." \}}
\NormalTok{])}
\end{Highlighting}
\end{Shaded}

\begin{center}\rule{0.5\linewidth}{0.5pt}\end{center}

\subsection{Exercise 6 --- Multi-step
problem}\label{exercise-6-multi-step-problem}

A university entrance exam has two papers. Based on historical data:

\begin{itemize}
\tightlist
\item
  60\% of students pass Paper 1.
\item
  Of students who pass Paper 1, 80\% also pass Paper 2.
\item
  Of students who fail Paper 1, 25\% still pass Paper 2.
\end{itemize}

To be admitted, a student must pass at least one paper.

\begin{enumerate}
\def\labelenumi{(\alph{enumi})}
\item
  What is the probability that a student passes both papers?
\item
  What is the probability that a student passes at least one paper?
\item
  Given that a student passes Paper 2, what is the probability they also
  passed Paper 1?
\end{enumerate}

\begin{Shaded}
\begin{Highlighting}[]
\NormalTok{//| echo: false}
\NormalTok{makeStepperHTML(6, [}
\NormalTok{  \{ op: "Part (a) — both papers", eq: "P(A \textbackslash{}\textbackslash{}cap B) = P(B \textbackslash{}\textbackslash{}mid A) \textbackslash{}\textbackslash{}cdot P(A) = 0.80 \textbackslash{}\textbackslash{}times 0.60 = 0.48", note: "A = passes Paper 1, B = passes Paper 2." \},}
\NormalTok{  \{ op: "Part (b) — find P(fails both)", eq: "P(B\^{}c \textbackslash{}\textbackslash{}cap A\^{}c) = P(B\^{}c \textbackslash{}\textbackslash{}mid A\^{}c) \textbackslash{}\textbackslash{}cdot P(A\^{}c) = 0.75 \textbackslash{}\textbackslash{}times 0.40 = 0.30" \},}
\NormalTok{  \{ op: "Part (b) — at least one paper", eq: "P(\textbackslash{}\textbackslash{}text\{at least one\}) = 1 {-} 0.30 = 0.70", note: "Using the complement rule on the \textquotesingle{}fails both\textquotesingle{} probability." \},}
\NormalTok{  \{ op: "Part (c) — find P(B) by total probability", eq: "P(B) = (0.80)(0.60) + (0.25)(0.40) = 0.48 + 0.10 = 0.58" \},}
\NormalTok{  \{ op: "Part (c) — apply Bayes\textquotesingle{} theorem", eq: "P(A \textbackslash{}\textbackslash{}mid B) = \textbackslash{}\textbackslash{}frac\{0.80 \textbackslash{}\textbackslash{}times 0.60\}\{0.58\} = \textbackslash{}\textbackslash{}frac\{0.48\}\{0.58\} \textbackslash{}\textbackslash{}approx 0.828", note: "Given passing Paper 2, there is about an 83\% chance the student also passed Paper 1." \}}
\NormalTok{])}
\end{Highlighting}
\end{Shaded}

\chapter{Probability distributions}\label{probability-distributions-1}

How probability spreads across possible values

\hfill\break

You flip a fair coin 10 times. Before you do it, you know 5 heads seems
most likely --- but what exactly is the probability of getting 5? What
about 8? And should 8 heads make you suspicious, or does it happen often
enough that you'd expect it occasionally? Without knowing the full shape
of the probability --- who gets what share --- you can't answer any of
those questions.

Here is a different situation. You're standing in a school gymnasium
with every 17-year-old in the district. Their heights spread out around
some average. A few people are very tall, a few very short, most cluster
near the middle. That familiar bell shape is not a coincidence --- it
emerges every time many small independent factors (genetics, nutrition,
sleep, dozens of other things) pile on top of each other.

One more. A website gets an average of 100 visitors per hour. One night
at 2am, the server logs show 140. Is that suspicious --- a bot attack,
maybe --- or just the kind of natural swing you'd expect? You need to
know how much variation is normal before you can say whether 140 is
alarming.

All three situations share the same underlying need: a complete picture
of how probability is spread across every possible value. That picture
is a \textbf{probability distribution}.

\section{What the notation is
saying}\label{what-the-notation-is-saying-22}

\subsection{Random variables}\label{random-variables}

A \textbf{random variable} is a quantity whose value is determined by a
random process. We write it with a capital letter --- usually \(X\).

\begin{itemize}
\tightlist
\item
  \(X\) = the number of heads in 10 coin flips. \(X\) could be 0, 1, 2,
  \ldots, 10.
\item
  \(X\) = the height of a randomly chosen 17-year-old. \(X\) could be
  any value in some continuous range.
\end{itemize}

The notation \(P(X = k)\) means: the probability that the random
variable \(X\) takes the specific value \(k\).

For 10 coin flips, \(P(X = 5)\) is the probability of getting exactly 5
heads. You'll calculate this shortly.

One constraint: all the probabilities must add up to 1. If you list
every value \(X\) could possibly take, the probabilities must account
for the whole picture.

We'll write \(\sum\) (capital sigma, the Greek letter S) to mean ``add
up all the terms that follow --- here, over every possible value
\(k\)'':

\[\sum_{\text{all } k} P(X = k) = 1\]

This is just saying that \emph{something} must happen.

\subsection{Mean, variance, and standard
deviation}\label{mean-variance-and-standard-deviation}

Once you have a distribution, you can summarise it with two numbers.

\textbf{The mean} (also called the \textbf{expected value}) tells you
the long-run average --- what value you'd get if you ran the random
process over and over and averaged the results. Written \(\mu\) (the
Greek letter \emph{mu}):

\[\mu = \sum_k k \cdot P(X = k)\]

Read this as: multiply each possible value by its probability, then add
everything up. It is a probability-weighted average.

\textbf{The variance} \(\sigma^2\) (sigma-squared) measures how spread
out the distribution is --- how far values tend to stray from the mean
on average. \textbf{The standard deviation} \(\sigma\) is the square
root of the variance, putting it back in the same units as \(X\) itself.

The key intuition: \(\sigma\) measures spread, not location. Two
distributions can have the same mean but look completely different.

\begin{longtable}[]{@{}llll@{}}
\toprule\noalign{}
Distribution & Mean \(\mu\) & Std dev \(\sigma\) & What it looks like \\
\midrule\noalign{}
\endhead
\bottomrule\noalign{}
\endlastfoot
A & 50 & 2 & Narrow spike near 50 \\
B & 50 & 15 & Wide, flat spread \\
\end{longtable}

If you're measuring exam scores and \(\sigma = 2\), almost everyone
scored within 4 marks of the average. If \(\sigma = 15\), scores are all
over the place. Same average, completely different distributions.

\section{The binomial distribution}\label{the-binomial-distribution}

The binomial distribution answers one specific question:

\begin{quote}
If I repeat an experiment \(n\) times, and each time the probability of
success is \(p\), what is the probability of getting exactly \(k\)
successes?
\end{quote}

\textbf{The conditions:} 1. Fixed number of trials: \(n\) 2. Each trial
has the same two outcomes: success or failure 3. The probability of
success is the same on every trial: \(p\) 4. The trials are independent
--- the result of one doesn't affect another

If your situation fits all four, the count of successes \(X\) follows a
\textbf{binomial distribution}, written \(X \sim B(n, p)\).

\subsection{Deriving the formula}\label{deriving-the-formula}

Suppose \(n = 3\) and you want exactly \(k = 2\) successes (call them S)
and 1 failure (F). One specific sequence of outcomes that gives this is:
S, S, F.

The probability of that exact sequence is
\(p \cdot p \cdot (1-p) = p^2(1-p)\).

But that's only one arrangement. You could also have: S, F, S or F, S,
S. There are 3 arrangements that give exactly 2 successes in 3 trials.

So the total probability is \(3 \times p^2(1-p)\).

The number 3 came from counting arrangements. In general, the number of
ways to arrange \(k\) successes in \(n\) trials is written
\(\binom{n}{k}\) (read ``n choose k'') and calculated as:

\[\binom{n}{k} = \frac{n!}{k!(n-k)!}\]

where \(n! = n \times (n-1) \times \cdots \times 2 \times 1\) is ``n
factorial'' --- the number of ways to order \(n\) distinct items.

\begin{tcolorbox}[enhanced jigsaw, title=\textcolor{quarto-callout-tip-color}{\faLightbulb}\hspace{0.5em}{Tip}, opacityback=0, colback=white, breakable, leftrule=.75mm, arc=.35mm, toptitle=1mm, left=2mm, colframe=quarto-callout-tip-color-frame, coltitle=black, titlerule=0mm, colbacktitle=quarto-callout-tip-color!10!white, opacitybacktitle=0.6, toprule=.15mm, bottomtitle=1mm, bottomrule=.15mm, rightrule=.15mm]

\textbf{Calculating \(\binom{n}{k}\)}

\(\binom{5}{2} = \frac{5!}{2! \cdot 3!} = \frac{5 \times 4 \times 3 \times 2 \times 1}{(2 \times 1)(3 \times 2 \times 1)} = \frac{120}{2 \times 6} = \frac{120}{12} = 10\)

A shortcut:
\(\binom{n}{k} = \frac{n \times (n-1) \times \cdots \times (n-k+1)}{k!}\)
--- multiply \(k\) descending numbers from \(n\), divide by \(k!\).

For \(\binom{5}{2}\):
\(\frac{5 \times 4}{2 \times 1} = \frac{20}{2} = 10\). Same answer, less
work.

\end{tcolorbox}

Putting it together:

\[\boxed{P(X = k) = \binom{n}{k} p^k (1-p)^{n-k}}\]

Three parts: - \(\binom{n}{k}\): the number of arrangements of \(k\)
successes in \(n\) trials - \(p^k\): the probability that each of those
\(k\) trials is a success - \((1-p)^{n-k}\): the probability that each
of the remaining trials is a failure

Multiply the parts: count of arrangements × probability of each
arrangement.

\subsection{Mean and variance of the
binomial}\label{mean-and-variance-of-the-binomial}

You could compute \(\mu = \sum k \cdot P(X = k)\) directly, but the
algebra is tedious. The results come out cleanly:

\[\mu = np \qquad \sigma^2 = np(1-p) \qquad \sigma = \sqrt{np(1-p)}\]

These make intuitive sense. If you flip a fair coin 100 times
(\(n=100\), \(p=0.5\)), you'd expect 50 heads on average:
\(\mu = 100 \times 0.5 = 50\). The standard deviation is
\(\sqrt{100 \times 0.5 \times 0.5} = \sqrt{25} = 5\), meaning a
``typical'' result is within about 5 heads of 50.

\begin{Shaded}
\begin{Highlighting}[]
\NormalTok{//| echo: false}
\NormalTok{\{}
\NormalTok{  const d3 = await require("d3@7");}

\NormalTok{  // {-}{-}{-} layout constants {-}{-}{-}}
\NormalTok{  const W = 660, H = 340;}
\NormalTok{  const margin = \{ top: 30, right: 20, bottom: 50, left: 55 \};}
\NormalTok{  const innerW = W {-} margin.left {-} margin.right;}
\NormalTok{  const innerH = H {-} margin.top {-} margin.bottom;}

\NormalTok{  // {-}{-}{-} helper: log{-}factorial (for large n) {-}{-}{-}}
\NormalTok{  function logFactorial(n) \{}
\NormalTok{    let v = 0;}
\NormalTok{    for (let i = 2; i \textless{}= n; i++) v += Math.log(i);}
\NormalTok{    return v;}
\NormalTok{  \}}
\NormalTok{  function logBinom(n, k) \{}
\NormalTok{    return logFactorial(n) {-} logFactorial(k) {-} logFactorial(n {-} k);}
\NormalTok{  \}}
\NormalTok{  function binomPMF(n, p, k) \{}
\NormalTok{    if (k \textless{} 0 || k \textgreater{} n) return 0;}
\NormalTok{    return Math.exp(logBinom(n, k) + k * Math.log(p) + (n {-} k) * Math.log(1 {-} p));}
\NormalTok{  \}}
\NormalTok{  function normalPDF(x, mu, sigma) \{}
\NormalTok{    const z = (x {-} mu) / sigma;}
\NormalTok{    return Math.exp({-}0.5 * z * z) / (sigma * Math.sqrt(2 * Math.PI));}
\NormalTok{  \}}

\NormalTok{  // {-}{-}{-} controls {-}{-}{-}}
\NormalTok{  const nSlider = Inputs.range([1, 30], \{ step: 1, value: 10, label: "n (trials)" \});}
\NormalTok{  const pSlider = Inputs.range([0.01, 0.99], \{ step: 0.01, value: 0.5, label: "p (success prob)" \});}
\NormalTok{  const normToggle = Inputs.toggle(\{ label: "Show normal approximation", value: false \});}

\NormalTok{  // {-}{-}{-} SVG {-}{-}{-}}
\NormalTok{  const svg = d3.create("svg")}
\NormalTok{    .attr("viewBox", \textasciigrave{}0 0 $\{W\} $\{H\}\textasciigrave{})}
\NormalTok{    .attr("width", "100\%")}
\NormalTok{    .style("font{-}family", "inherit");}

\NormalTok{  const g = svg.append("g")}
\NormalTok{    .attr("transform", \textasciigrave{}translate($\{margin.left\},$\{margin.top\})\textasciigrave{});}

\NormalTok{  // axes groups}
\NormalTok{  const xAxisG = g.append("g").attr("transform", \textasciigrave{}translate(0,$\{innerH\})\textasciigrave{});}
\NormalTok{  const yAxisG = g.append("g");}

\NormalTok{  // axis labels}
\NormalTok{  g.append("text")}
\NormalTok{    .attr("x", innerW / 2).attr("y", innerH + 42)}
\NormalTok{    .attr("text{-}anchor", "middle").attr("font{-}size", 13)}
\NormalTok{    .text("k (number of successes)");}

\NormalTok{  g.append("text")}
\NormalTok{    .attr("transform", "rotate({-}90)")}
\NormalTok{    .attr("x", {-}innerH / 2).attr("y", {-}44)}
\NormalTok{    .attr("text{-}anchor", "middle").attr("font{-}size", 13)}
\NormalTok{    .text("P(X = k)");}

\NormalTok{  // bars group}
\NormalTok{  const barsG = g.append("g");}
\NormalTok{  // normal curve}
\NormalTok{  const curvePath = g.append("path")}
\NormalTok{    .attr("fill", "none")}
\NormalTok{    .attr("stroke", "\#e55")}
\NormalTok{    .attr("stroke{-}width", 2.5)}
\NormalTok{    .attr("stroke{-}dasharray", "6,3")}
\NormalTok{    .attr("opacity", 0);}

\NormalTok{  // stats text}
\NormalTok{  const statsText = svg.append("text")}
\NormalTok{    .attr("x", margin.left + innerW {-} 4)}
\NormalTok{    .attr("y", margin.top + 16)}
\NormalTok{    .attr("text{-}anchor", "end")}
\NormalTok{    .attr("font{-}size", 13)}
\NormalTok{    .attr("fill", "\#374151");}

\NormalTok{  function update() \{}
\NormalTok{    const n = nSlider.value;}
\NormalTok{    const p = pSlider.value;}
\NormalTok{    const showNorm = normToggle.value;}

\NormalTok{    const mu = n * p;}
\NormalTok{    const sigma = Math.sqrt(n * p * (1 {-} p));}

\NormalTok{    // pmf data}
\NormalTok{    const data = d3.range(0, n + 1).map(k =\textgreater{} (\{ k, prob: binomPMF(n, p, k) \}));}

\NormalTok{    const maxProb = d3.max(data, d =\textgreater{} d.prob);}
\NormalTok{    const yMax = Math.max(maxProb * 1.15, 0.05);}

\NormalTok{    // scales}
\NormalTok{    const xScale = d3.scaleBand()}
\NormalTok{      .domain(d3.range(0, n + 1))}
\NormalTok{      .range([0, innerW])}
\NormalTok{      .padding(0.1);}

\NormalTok{    const yScale = d3.scaleLinear()}
\NormalTok{      .domain([0, yMax])}
\NormalTok{      .range([innerH, 0]);}

\NormalTok{    // axes}
\NormalTok{    const xTick = n \textless{}= 15 ? 1 : n \textless{}= 22 ? 2 : 3;}
\NormalTok{    xAxisG.call(}
\NormalTok{      d3.axisBottom(xScale)}
\NormalTok{        .tickValues(d3.range(0, n + 1).filter(k =\textgreater{} k \% xTick === 0))}
\NormalTok{        .tickFormat(d3.format("d"))}
\NormalTok{    );}
\NormalTok{    yAxisG.call(d3.axisLeft(yScale).ticks(5).tickFormat(d3.format(".3f")));}

\NormalTok{    // bars}
\NormalTok{    const bars = barsG.selectAll("rect").data(data, d =\textgreater{} d.k);}
\NormalTok{    bars.enter().append("rect")}
\NormalTok{      .attr("fill", "\#3b82f6")}
\NormalTok{      .attr("opacity", 0.8)}
\NormalTok{      .merge(bars)}
\NormalTok{      .attr("x", d =\textgreater{} xScale(d.k))}
\NormalTok{      .attr("width", xScale.bandwidth())}
\NormalTok{      .attr("y", d =\textgreater{} yScale(d.prob))}
\NormalTok{      .attr("height", d =\textgreater{} innerH {-} yScale(d.prob));}
\NormalTok{    bars.exit().remove();}

\NormalTok{    // normal overlay}
\NormalTok{    if (showNorm \&\& sigma \textgreater{} 0.3) \{}
\NormalTok{      const xMin = xScale(0) + xScale.bandwidth() / 2;}
\NormalTok{      const xMax = xScale(n) + xScale.bandwidth() / 2;}
\NormalTok{      const pts = d3.range(200).map(i =\textgreater{} \{}
\NormalTok{        const xVal = mu {-} 4 * sigma + (i / 199) * 8 * sigma;}
\NormalTok{        const screenX = (xVal / n) * innerW; // linear position across bar centres}
\NormalTok{        // map xVal to screen: bar centre for integer k is xScale(k) + bw/2}
\NormalTok{        // approximate as linear interpolation}
\NormalTok{        const bw = xScale.bandwidth();}
\NormalTok{        const screenXv = xScale(0) + bw / 2 + (xVal / 1) * (xScale(1) {-} xScale(0));}
\NormalTok{        return [xScale(0) + bw / 2 + xVal * (xScale.step()), yScale(normalPDF(xVal, mu, sigma))];}
\NormalTok{      \});}
\NormalTok{      const line = d3.line().x(d =\textgreater{} d[0]).y(d =\textgreater{} d[1]).curve(d3.curveBasis);}
\NormalTok{      curvePath.attr("d", line(pts)).attr("opacity", 1);}
\NormalTok{    \} else \{}
\NormalTok{      curvePath.attr("opacity", 0);}
\NormalTok{    \}}

\NormalTok{    statsText.text(\textasciigrave{}μ = $\{mu.toFixed(2)\}   σ = $\{sigma.toFixed(3)\}\textasciigrave{});}
\NormalTok{  \}}

\NormalTok{  nSlider.addEventListener("input", update);}
\NormalTok{  pSlider.addEventListener("input", update);}
\NormalTok{  normToggle.addEventListener("input", update);}

\NormalTok{  update();}

\NormalTok{  const container = html\textasciigrave{}\textless{}div style="border:1px solid \#e5e7eb;border{-}radius:8px;padding:1rem 1.25rem;margin:0.75rem 0;"\textgreater{}}
\NormalTok{    \textless{}div style="font{-}weight:600;margin{-}bottom:0.5rem;font{-}size:0.95em;"\textgreater{}Binomial distribution explorer\textless{}/div\textgreater{}}
\NormalTok{    \textless{}div style="display:grid;grid{-}template{-}columns:1fr 1fr 1fr;gap:0.5rem 1.5rem;margin{-}bottom:0.75rem;"\textgreater{}}
\NormalTok{      $\{nSlider\}}
\NormalTok{      $\{pSlider\}}
\NormalTok{      $\{normToggle\}}
\NormalTok{    \textless{}/div\textgreater{}}
\NormalTok{    $\{svg.node()\}}
\NormalTok{    \textless{}p style="font{-}size:0.82em;color:\#6b7280;margin{-}top:0.4rem;"\textgreater{}}
\NormalTok{      Bar heights show P(X = k). The dashed red curve (when toggled) is the normal approximation N(np, np(1−p)).}
\NormalTok{      For large n and moderate p the two shapes converge — this is the central limit theorem in action.}
\NormalTok{    \textless{}/p\textgreater{}}
\NormalTok{  \textless{}/div\textgreater{}\textasciigrave{};}
\NormalTok{  return container;}
\NormalTok{\}}
\end{Highlighting}
\end{Shaded}

\section{The normal distribution}\label{the-normal-distribution}

The binomial counts discrete things --- 0 heads, 1 head, 2 heads. But
heights, weights, temperatures, and exam scores can take any value in a
continuous range. For these, you need a different model.

The \textbf{normal distribution} \(N(\mu, \sigma^2)\) is a continuous
bell curve centred at \(\mu\) and spreading out by \(\sigma\) in each
direction (the second parameter is the variance \(\sigma^2\), not the
standard deviation --- to find the spread in the original units, take
the square root: \(\sigma = \sqrt{\sigma^2}\)). You specify it by two
numbers only: the mean and the variance.

\subsection{Why the bell curve appears
everywhere}\label{why-the-bell-curve-appears-everywhere}

This is one of the most remarkable results in all of probability: when
you add up many independent random influences, the total tends toward
the normal distribution --- regardless of what the individual influences
look like. A person's height is the sum of contributions from hundreds
of genes, years of nutrition, sleep patterns, and dozens of other
factors. None of those individually looks like a bell curve. Their sum
does.

This is the \textbf{central limit theorem}. You won't prove it at this
stage, but you'll use it, and you'll see it confirmed every time you
look at data from a natural process.

\subsection{The 68-95-99.7 rule}\label{the-68-95-99.7-rule}

For any normal distribution \(N(\mu, \sigma^2)\):

\begin{longtable}[]{@{}ll@{}}
\toprule\noalign{}
Range & Probability \\
\midrule\noalign{}
\endhead
\bottomrule\noalign{}
\endlastfoot
\(\mu - \sigma\) to \(\mu + \sigma\) (within 1 std dev) &
\textbf{68\%} \\
\(\mu - 2\sigma\) to \(\mu + 2\sigma\) (within 2 std dev) &
\textbf{95\%} \\
\(\mu - 3\sigma\) to \(\mu + 3\sigma\) (within 3 std dev) &
\textbf{99.7\%} \\
\end{longtable}

This is sometimes called the \textbf{empirical rule}. It gives you a
quick sense of what is ``normal'' variation and what is surprising.

If exam scores are \(N(65, 144)\) --- so \(\mu = 65\) and
\(\sigma = 12\) --- then about 95\% of students score between
\(65 - 24 = 41\) and \(65 + 24 = 89\). A score of 20 would be genuinely
extraordinary (more than 3 standard deviations below the mean),
occurring less than 0.15\% of the time.

\subsection{Standardising: the z-score}\label{standardising-the-z-score}

A practical problem: there is a different normal distribution for every
combination of \(\mu\) and \(\sigma\). You can't carry a different table
for each one.

The solution is to convert any normal distribution to a single standard
one. The \textbf{standard normal} distribution is \(N(0, 1)\) --- mean
0, variance 1.

To convert a value \(x\) from any normal distribution to its standard
normal equivalent, compute:

\[z = \frac{x - \mu}{\sigma}\]

The value \(z\) is called a \textbf{z-score}. It tells you how many
standard deviations above (positive) or below (negative) the mean the
value \(x\) sits.

Once you have \(z\), you look up \(P(Z \leq z)\) in a standard normal
table or compute it on a calculator. This single table handles all
normal distributions.

\textbf{Example.} Scores are \(N(65, 144)\), so \(\mu = 65\),
\(\sigma = 12\). For a score of \(x = 83\):

\[z = \frac{83 - 65}{12} = \frac{18}{12} = 1.5\]

A score of 83 sits 1.5 standard deviations above the mean. A table gives
\(P(Z \leq 1.5) \approx 0.933\), so about 93.3\% of students scored 83
or below --- and about 6.7\% scored above 83.

\begin{Shaded}
\begin{Highlighting}[]
\NormalTok{//| echo: false}
\NormalTok{\{}
\NormalTok{  const d3 = await require("d3@7");}

\NormalTok{  // {-}{-}{-} layout {-}{-}{-}}
\NormalTok{  const W = 660, H = 320;}
\NormalTok{  const margin = \{ top: 30, right: 20, bottom: 50, left: 55 \};}
\NormalTok{  const innerW = W {-} margin.left {-} margin.right;}
\NormalTok{  const innerH = H {-} margin.top {-} margin.bottom;}

\NormalTok{  // {-}{-}{-} normal CDF (erf approximation) {-}{-}{-}}
\NormalTok{  function erf(x) \{}
\NormalTok{    const a1 =  0.254829592, a2 = {-}0.284496736, a3 =  1.421413741;}
\NormalTok{    const a4 = {-}1.453152027, a5 =  1.061405429, p  =  0.3275911;}
\NormalTok{    const sign = x \textless{} 0 ? {-}1 : 1;}
\NormalTok{    const t = 1 / (1 + p * Math.abs(x));}
\NormalTok{    const y = 1 {-} (((((a5 * t + a4) * t) + a3) * t + a2) * t + a1) * t * Math.exp({-}x * x);}
\NormalTok{    return sign * y;}
\NormalTok{  \}}
\NormalTok{  function normalCDF(x, mu, sigma) \{}
\NormalTok{    return 0.5 * (1 + erf((x {-} mu) / (sigma * Math.sqrt(2))));}
\NormalTok{  \}}
\NormalTok{  function normalPDF(x, mu, sigma) \{}
\NormalTok{    const z = (x {-} mu) / sigma;}
\NormalTok{    return Math.exp({-}0.5 * z * z) / (sigma * Math.sqrt(2 * Math.PI));}
\NormalTok{  \}}

\NormalTok{  // {-}{-}{-} mode state {-}{-}{-}}
\NormalTok{  let mode = "sigma"; // "sigma" | "interval"}

\NormalTok{  // {-}{-}{-} controls {-}{-}{-}}
\NormalTok{  const muSlider  = Inputs.range([{-}10, 10],  \{ step: 0.5,  value: 0,   label: "μ (mean)" \});}
\NormalTok{  const sigSlider = Inputs.range([0.5, 5],   \{ step: 0.1,  value: 1,   label: "σ (std dev)" \});}
\NormalTok{  const bandSel   = Inputs.select(["1σ", "2σ", "3σ"], \{ label: "Shade within", value: "1σ" \});}
\NormalTok{  const aInput    = Inputs.number(\{ label: "a", value: {-}1, step: 0.1 \});}
\NormalTok{  const bInput    = Inputs.number(\{ label: "b", value:  1, step: 0.1 \});}
\NormalTok{  const modeBtn   = document.createElement("button");}
\NormalTok{  modeBtn.textContent = "Switch to P(a \textless{} X \textless{} b) mode";}
\NormalTok{  modeBtn.style.cssText = "padding:0.3rem 0.8rem;border:1px solid \#d1d5db;border{-}radius:4px;background:\#fff;cursor:pointer;font{-}size:0.87em;margin{-}top:0.3rem;";}

\NormalTok{  // sigma{-}mode controls wrapper}
\NormalTok{  const sigmaControls = html\textasciigrave{}\textless{}div style="display:grid;grid{-}template{-}columns:1fr 1fr 1fr;gap:0.5rem 1.5rem;"\textgreater{}$\{muSlider\}$\{sigSlider\}$\{bandSel\}\textless{}/div\textgreater{}\textasciigrave{};}
\NormalTok{  const intervalControls = html\textasciigrave{}\textless{}div style="display:grid;grid{-}template{-}columns:1fr 1fr 1fr;gap:0.5rem 1.5rem;"\textgreater{}$\{muSlider\}$\{sigSlider\}\textless{}div style="display:flex;gap:0.75rem;"\textgreater{}$\{aInput\}$\{bInput\}\textless{}/div\textgreater{}\textless{}/div\textgreater{}\textasciigrave{};}
\NormalTok{  intervalControls.style.display = "none";}

\NormalTok{  // {-}{-}{-} SVG {-}{-}{-}}
\NormalTok{  const svg = d3.create("svg")}
\NormalTok{    .attr("viewBox", \textasciigrave{}0 0 $\{W\} $\{H\}\textasciigrave{})}
\NormalTok{    .attr("width", "100\%")}
\NormalTok{    .style("font{-}family", "inherit");}

\NormalTok{  const g = svg.append("g")}
\NormalTok{    .attr("transform", \textasciigrave{}translate($\{margin.left\},$\{margin.top\})\textasciigrave{});}

\NormalTok{  const xAxisG = g.append("g").attr("transform", \textasciigrave{}translate(0,$\{innerH\})\textasciigrave{});}
\NormalTok{  const yAxisG = g.append("g");}

\NormalTok{  g.append("text")}
\NormalTok{    .attr("x", innerW / 2).attr("y", innerH + 42)}
\NormalTok{    .attr("text{-}anchor", "middle").attr("font{-}size", 13).text("x");}
\NormalTok{  g.append("text")}
\NormalTok{    .attr("transform", "rotate({-}90)")}
\NormalTok{    .attr("x", {-}innerH / 2).attr("y", {-}44)}
\NormalTok{    .attr("text{-}anchor", "middle").attr("font{-}size", 13).text("f(x)");}

\NormalTok{  const shadeArea = g.append("path").attr("fill", "\#3b82f6").attr("opacity", 0.25);}
\NormalTok{  const curvePath = g.append("path").attr("fill", "none").attr("stroke", "\#1d4ed8").attr("stroke{-}width", 2.5);}
\NormalTok{  const meanLine  = g.append("line").attr("stroke", "\#374151").attr("stroke{-}width", 1.5).attr("stroke{-}dasharray", "5,3");}

\NormalTok{  // sigma band lines}
\NormalTok{  const bandLines = [{-}3,{-}2,{-}1,1,2,3].map(m =\textgreater{}}
\NormalTok{    g.append("line").attr("stroke", "\#9ca3af").attr("stroke{-}width", 1).attr("stroke{-}dasharray", "3,3").attr("opacity", 0)}
\NormalTok{  );}

\NormalTok{  const probLabel = svg.append("text")}
\NormalTok{    .attr("x", margin.left + innerW {-} 4)}
\NormalTok{    .attr("y", margin.top + 18)}
\NormalTok{    .attr("text{-}anchor", "end")}
\NormalTok{    .attr("font{-}size", 13)}
\NormalTok{    .attr("fill", "\#1d4ed8")}
\NormalTok{    .attr("font{-}weight", "600");}

\NormalTok{  const statsLabel = svg.append("text")}
\NormalTok{    .attr("x", margin.left + 6)}
\NormalTok{    .attr("y", margin.top + 18)}
\NormalTok{    .attr("font{-}size", 12)}
\NormalTok{    .attr("fill", "\#6b7280");}

\NormalTok{  const N\_CURVE = 400;}

\NormalTok{  function getRange(mu, sigma) \{}
\NormalTok{    return [mu {-} 4.5 * sigma, mu + 4.5 * sigma];}
\NormalTok{  \}}

\NormalTok{  function update() \{}
\NormalTok{    const mu  = muSlider.value;}
\NormalTok{    const sig = sigSlider.value;}
\NormalTok{    const [xMin, xMax] = getRange(mu, sig);}
\NormalTok{    const yMax = normalPDF(mu, mu, sig) * 1.15;}

\NormalTok{    const xScale = d3.scaleLinear().domain([xMin, xMax]).range([0, innerW]);}
\NormalTok{    const yScale = d3.scaleLinear().domain([0, yMax]).range([innerH, 0]);}

\NormalTok{    xAxisG.call(d3.axisBottom(xScale).ticks(8).tickFormat(d3.format(".1f")));}
\NormalTok{    yAxisG.call(d3.axisLeft(yScale).ticks(5).tickFormat(d3.format(".3f")));}

\NormalTok{    // curve}
\NormalTok{    const curveData = d3.range(N\_CURVE).map(i =\textgreater{} \{}
\NormalTok{      const x = xMin + (i / (N\_CURVE {-} 1)) * (xMax {-} xMin);}
\NormalTok{      return [xScale(x), yScale(normalPDF(x, mu, sig))];}
\NormalTok{    \});}
\NormalTok{    const line = d3.line().x(d =\textgreater{} d[0]).y(d =\textgreater{} d[1]).curve(d3.curveBasis);}
\NormalTok{    curvePath.attr("d", line(curveData));}

\NormalTok{    // mean line}
\NormalTok{    meanLine}
\NormalTok{      .attr("x1", xScale(mu)).attr("x2", xScale(mu))}
\NormalTok{      .attr("y1", 0).attr("y2", innerH);}

\NormalTok{    statsLabel.text(\textasciigrave{}μ = $\{mu.toFixed(1)\}   σ = $\{sig.toFixed(1)\}\textasciigrave{});}

\NormalTok{    if (mode === "sigma") \{}
\NormalTok{      const bandMap = \{ "1σ": 1, "2σ": 2, "3σ": 3 \};}
\NormalTok{      const k = bandMap[bandSel.value];}
\NormalTok{      const lo = mu {-} k * sig;}
\NormalTok{      const hi = mu + k * sig;}

\NormalTok{      // shade}
\NormalTok{      const shadeData = d3.range(200).map(i =\textgreater{} \{}
\NormalTok{        const x = lo + (i / 199) * (hi {-} lo);}
\NormalTok{        return [xScale(x), yScale(normalPDF(x, mu, sig))];}
\NormalTok{      \});}
\NormalTok{      const area = d3.area()}
\NormalTok{        .x(d =\textgreater{} d[0]).y0(innerH).y1(d =\textgreater{} d[1]).curve(d3.curveBasis);}
\NormalTok{      shadeArea.attr("d", area(shadeData));}

\NormalTok{      const prob = normalCDF(hi, mu, sig) {-} normalCDF(lo, mu, sig);}
\NormalTok{      probLabel.text(\textasciigrave{}P(within $\{bandSel.value\}) = $\{(prob * 100).toFixed(2)\}\%\textasciigrave{});}

\NormalTok{      // band dashed lines}
\NormalTok{      [{-}k, k].forEach((m, i) =\textgreater{} \{}
\NormalTok{        bandLines[i].attr("x1", xScale(mu + m * sig)).attr("x2", xScale(mu + m * sig))}
\NormalTok{          .attr("y1", 0).attr("y2", innerH).attr("opacity", 0.7);}
\NormalTok{      \});}
\NormalTok{      bandLines.slice(2).forEach(l =\textgreater{} l.attr("opacity", 0));}

\NormalTok{    \} else \{}
\NormalTok{      // interval mode}
\NormalTok{      const a = Math.min(aInput.value, bInput.value);}
\NormalTok{      const b = Math.max(aInput.value, bInput.value);}
\NormalTok{      const lo = Math.max(xMin, a);}
\NormalTok{      const hi = Math.min(xMax, b);}
\NormalTok{      if (lo \textless{} hi) \{}
\NormalTok{        const shadeData = d3.range(200).map(i =\textgreater{} \{}
\NormalTok{          const x = lo + (i / 199) * (hi {-} lo);}
\NormalTok{          return [xScale(x), yScale(normalPDF(x, mu, sig))];}
\NormalTok{        \});}
\NormalTok{        const area = d3.area()}
\NormalTok{          .x(d =\textgreater{} d[0]).y0(innerH).y1(d =\textgreater{} d[1]).curve(d3.curveBasis);}
\NormalTok{        shadeArea.attr("d", area(shadeData));}
\NormalTok{        const prob = normalCDF(b, mu, sig) {-} normalCDF(a, mu, sig);}
\NormalTok{        probLabel.text(\textasciigrave{}P($\{a.toFixed(1)\} \textless{} X \textless{} $\{b.toFixed(1)\}) = $\{(prob * 100).toFixed(2)\}\%\textasciigrave{});}
\NormalTok{      \} else \{}
\NormalTok{        shadeArea.attr("d", null);}
\NormalTok{        probLabel.text("—");}
\NormalTok{      \}}
\NormalTok{      bandLines.forEach(l =\textgreater{} l.attr("opacity", 0));}
\NormalTok{    \}}
\NormalTok{  \}}

\NormalTok{  muSlider.addEventListener("input",  update);}
\NormalTok{  sigSlider.addEventListener("input", update);}
\NormalTok{  bandSel.addEventListener("input",   update);}
\NormalTok{  aInput.addEventListener("input",    update);}
\NormalTok{  bInput.addEventListener("input",    update);}

\NormalTok{  modeBtn.addEventListener("click", () =\textgreater{} \{}
\NormalTok{    mode = mode === "sigma" ? "interval" : "sigma";}
\NormalTok{    modeBtn.textContent = mode === "sigma"}
\NormalTok{      ? "Switch to P(a \textless{} X \textless{} b) mode"}
\NormalTok{      : "Switch to σ{-}band mode";}
\NormalTok{    sigmaControls.style.display   = mode === "sigma"    ? "" : "none";}
\NormalTok{    intervalControls.style.display = mode === "interval" ? "" : "none";}
\NormalTok{    update();}
\NormalTok{  \});}

\NormalTok{  update();}

\NormalTok{  const container = html\textasciigrave{}\textless{}div style="border:1px solid \#e5e7eb;border{-}radius:8px;padding:1rem 1.25rem;margin:0.75rem 0;"\textgreater{}}
\NormalTok{    \textless{}div style="font{-}weight:600;margin{-}bottom:0.5rem;font{-}size:0.95em;"\textgreater{}Normal distribution explorer\textless{}/div\textgreater{}}
\NormalTok{    $\{sigmaControls\}}
\NormalTok{    $\{intervalControls\}}
\NormalTok{    \textless{}div style="margin{-}top:0.5rem;"\textgreater{}$\{modeBtn\}\textless{}/div\textgreater{}}
\NormalTok{    $\{svg.node()\}}
\NormalTok{    \textless{}p style="font{-}size:0.82em;color:\#6b7280;margin{-}top:0.4rem;"\textgreater{}}
\NormalTok{      Shaded region shows the selected probability. The dashed vertical line marks μ.}
\NormalTok{      Adjust μ to shift the curve left/right; adjust σ to widen or narrow it.}
\NormalTok{    \textless{}/p\textgreater{}}
\NormalTok{  \textless{}/div\textgreater{}\textasciigrave{};}
\NormalTok{  return container;}
\NormalTok{\}}
\end{Highlighting}
\end{Shaded}

\section{Worked examples}\label{worked-examples-22}

\subsection{Example 1 --- Binomial exact (quality
control)}\label{example-1-binomial-exact-quality-control}

A manufacturing line produces components. Historical data shows that
15\% of components are defective. A quality inspector takes a batch of 8
components. What is the probability that exactly 2 are defective?

Here \(X\) = number of defectives, \(n = 8\), \(p = 0.15\), and we want
\(P(X = 2)\).

\[P(X = 2) = \binom{8}{2}(0.15)^2(0.85)^6\]

Calculate each part:

\[\binom{8}{2} = \frac{8 \times 7}{2 \times 1} = 28\]

\[(0.15)^2 = 0.0225\]

\[(0.85)^6 = 0.85 \times 0.85 \times 0.85 \times 0.85 \times 0.85 \times 0.85 \approx 0.3771\]

\[P(X = 2) = 28 \times 0.0225 \times 0.3771 \approx 28 \times 0.008485 \approx 0.2376\]

About 24\% of batches of 8 will contain exactly 2 defectives.

Check: the mean number of defectives is
\(\mu = np = 8 \times 0.15 = 1.2\), so 2 is slightly above average --- a
24\% probability is plausible.

\begin{Shaded}
\begin{Highlighting}[]
\NormalTok{//| echo: false}
\NormalTok{katex = require("katex@0.16")}
\end{Highlighting}
\end{Shaded}

\begin{Shaded}
\begin{Highlighting}[]
\NormalTok{//| echo: false}
\NormalTok{function makeStepperHTML(exerciseNum, steps) \{}
\NormalTok{  let current = 0;}
\NormalTok{  const totalSteps = steps.length;}
\NormalTok{  const container = document.createElement("div");}
\NormalTok{  container.style.cssText = "border:1px solid \#e5e7eb; border{-}radius:8px; padding:1rem 1.25rem; margin:0.75rem 0; font{-}family:inherit;";}
\NormalTok{  const stepsDiv = document.createElement("div");}
\NormalTok{  stepsDiv.style.marginBottom = "0.75rem";}
\NormalTok{  function renderSteps() \{}
\NormalTok{    stepsDiv.innerHTML = "";}
\NormalTok{    for (let i = 0; i \textless{} current; i++) \{}
\NormalTok{      const s = steps[i];}
\NormalTok{      const row = document.createElement("div");}
\NormalTok{      row.style.cssText = "display:grid; grid{-}template{-}columns:200px 1fr; gap:0.5rem 1rem; align{-}items:baseline; padding:0.35rem 0; border{-}top:1px solid \#f3f4f6;";}
\NormalTok{      const opCell = document.createElement("span");}
\NormalTok{      opCell.style.cssText = "font{-}size:0.85em; color:\#6b7280; font{-}style:italic;";}
\NormalTok{      opCell.textContent = s.op;}
\NormalTok{      const eqCell = document.createElement("span");}
\NormalTok{      eqCell.innerHTML = katex.renderToString(s.eq, \{ throwOnError: false, displayMode: false \});}
\NormalTok{      row.appendChild(opCell);}
\NormalTok{      row.appendChild(eqCell);}
\NormalTok{      if (s.note) \{}
\NormalTok{        const noteRow = document.createElement("div");}
\NormalTok{        noteRow.style.cssText = "grid{-}column:2; font{-}size:0.82em; color:\#6b7280; padding{-}bottom:0.2rem;";}
\NormalTok{        noteRow.textContent = s.note;}
\NormalTok{        row.appendChild(noteRow);}
\NormalTok{      \}}
\NormalTok{      stepsDiv.appendChild(row);}
\NormalTok{    \}}
\NormalTok{    if (current === totalSteps) \{}
\NormalTok{      const done = document.createElement("div");}
\NormalTok{      done.style.cssText = "margin{-}top:0.5rem; font{-}size:0.9em; color:\#059669; font{-}weight:500;";}
\NormalTok{      done.textContent = "✓ Solution complete";}
\NormalTok{      stepsDiv.appendChild(done);}
\NormalTok{    \}}
\NormalTok{  \}}
\NormalTok{  const controls = document.createElement("div");}
\NormalTok{  controls.style.cssText = "display:flex; gap:0.5rem; align{-}items:center;";}
\NormalTok{  const nextBtn = document.createElement("button");}
\NormalTok{  nextBtn.textContent = "Next step →";}
\NormalTok{  nextBtn.style.cssText = "padding:0.35rem 0.85rem; border:1px solid \#d1d5db; border{-}radius:4px; background:\#fff; cursor:pointer; font{-}size:0.9em;";}
\NormalTok{  const resetBtn = document.createElement("button");}
\NormalTok{  resetBtn.textContent = "Reset";}
\NormalTok{  resetBtn.style.cssText = "padding:0.35rem 0.75rem; border:1px solid \#e5e7eb; border{-}radius:4px; background:\#f9fafb; cursor:pointer; font{-}size:0.9em; color:\#6b7280;";}
\NormalTok{  const counter = document.createElement("span");}
\NormalTok{  counter.style.cssText = "font{-}size:0.82em; color:\#9ca3af; margin{-}left:0.25rem;";}
\NormalTok{  function updateButtons() \{}
\NormalTok{    nextBtn.disabled = current === totalSteps;}
\NormalTok{    nextBtn.style.opacity = current === totalSteps ? "0.4" : "1";}
\NormalTok{    counter.textContent = current === 0 ? "Click to reveal steps" : \textasciigrave{}Step $\{current\} of $\{totalSteps\}\textasciigrave{};}
\NormalTok{  \}}
\NormalTok{  nextBtn.onclick = () =\textgreater{} \{ if (current \textless{} totalSteps) \{ current++; renderSteps(); updateButtons(); \} \};}
\NormalTok{  resetBtn.onclick = () =\textgreater{} \{ current = 0; renderSteps(); updateButtons(); \};}
\NormalTok{  controls.appendChild(nextBtn);}
\NormalTok{  controls.appendChild(resetBtn);}
\NormalTok{  controls.appendChild(counter);}
\NormalTok{  container.appendChild(stepsDiv);}
\NormalTok{  container.appendChild(controls);}
\NormalTok{  renderSteps();}
\NormalTok{  updateButtons();}
\NormalTok{  return container;}
\NormalTok{\}}
\end{Highlighting}
\end{Shaded}

\begin{Shaded}
\begin{Highlighting}[]
\NormalTok{//| echo: false}
\NormalTok{makeStepperHTML(1, [}
\NormalTok{  \{ op: "Identify the distribution",    eq: "X \textbackslash{}\textbackslash{}sim B(8,\textbackslash{}\textbackslash{}, 0.15),\textbackslash{}\textbackslash{}quad n=8,\textbackslash{}\textbackslash{}; p=0.15,\textbackslash{}\textbackslash{}; k=2" \},}
\NormalTok{  \{ op: "Write the formula",            eq: "P(X=2) = \textbackslash{}\textbackslash{}binom\{8\}\{2\}(0.15)\^{}2(0.85)\^{}6" \},}
\NormalTok{  \{ op: "Compute the combination",      eq: "\textbackslash{}\textbackslash{}binom\{8\}\{2\} = \textbackslash{}\textbackslash{}frac\{8 \textbackslash{}\textbackslash{}times 7\}\{2 \textbackslash{}\textbackslash{}times 1\} = 28" \},}
\NormalTok{  \{ op: "Compute p\^{}k",                  eq: "(0.15)\^{}2 = 0.0225" \},}
\NormalTok{  \{ op: "Compute (1{-}p)\^{}\{n{-}k\}",          eq: "(0.85)\^{}6 \textbackslash{}\textbackslash{}approx 0.3771" \},}
\NormalTok{  \{ op: "Multiply all three parts",     eq: "P(X=2) = 28 \textbackslash{}\textbackslash{}times 0.0225 \textbackslash{}\textbackslash{}times 0.3771 \textbackslash{}\textbackslash{}approx 0.2376", note: "About 24\% of batches will contain exactly 2 defectives." \},}
\NormalTok{  \{ op: "Sanity check via mean",        eq: "\textbackslash{}\textbackslash{}mu = np = 8 \textbackslash{}\textbackslash{}times 0.15 = 1.2", note: "k = 2 is just above average, so \textasciitilde{}24\% is a plausible probability." \}}
\NormalTok{])}
\end{Highlighting}
\end{Shaded}

\subsection{Example 2 --- Binomial complement
(healthcare)}\label{example-2-binomial-complement-healthcare}

A new drug cures 70\% of patients. In a clinical trial, 10 patients
receive the drug. What is the probability that at least 8 are cured?

\(X \sim B(10, 0.7)\). We want
\(P(X \geq 8) = P(X=8) + P(X=9) + P(X=10)\).

With only 3 terms, direct calculation is manageable:

\[P(X=8) = \binom{10}{8}(0.7)^8(0.3)^2 = 45 \times 0.05765 \times 0.09 \approx 0.2335\]

\[P(X=9) = \binom{10}{9}(0.7)^9(0.3)^1 = 10 \times 0.04035 \times 0.3 \approx 0.1211\]

\[P(X=10) = \binom{10}{10}(0.7)^{10}(0.3)^0 = 1 \times 0.02825 \times 1 \approx 0.0282\]

\[P(X \geq 8) \approx 0.2335 + 0.1211 + 0.0282 = 0.3828\]

There is about a 38\% probability that at least 8 of 10 patients are
cured. Note: you would use the complement approach when the ``at least''
threshold is low --- e.g., \(P(X \geq 2)\) is much easier to compute as
\(1 - P(X=0) - P(X=1)\) than to sum 9 terms directly.

\begin{Shaded}
\begin{Highlighting}[]
\NormalTok{//| echo: false}
\NormalTok{makeStepperHTML(2, [}
\NormalTok{  \{ op: "Identify the distribution",    eq: "X \textbackslash{}\textbackslash{}sim B(10,\textbackslash{}\textbackslash{}, 0.7)" \},}
\NormalTok{  \{ op: "Write as a sum",               eq: "P(X \textbackslash{}\textbackslash{}geq 8) = P(X=8) + P(X=9) + P(X=10)" \},}
\NormalTok{  \{ op: "Compute P(X = 8)",             eq: "\textbackslash{}\textbackslash{}binom\{10\}\{8\}(0.7)\^{}8(0.3)\^{}2 = 45 \textbackslash{}\textbackslash{}times 0.05765 \textbackslash{}\textbackslash{}times 0.09 \textbackslash{}\textbackslash{}approx 0.2335" \},}
\NormalTok{  \{ op: "Compute P(X = 9)",             eq: "\textbackslash{}\textbackslash{}binom\{10\}\{9\}(0.7)\^{}9(0.3)\^{}1 = 10 \textbackslash{}\textbackslash{}times 0.04035 \textbackslash{}\textbackslash{}times 0.3 \textbackslash{}\textbackslash{}approx 0.1211" \},}
\NormalTok{  \{ op: "Compute P(X = 10)",            eq: "\textbackslash{}\textbackslash{}binom\{10\}\{10\}(0.7)\^{}\{10\}(0.3)\^{}0 = 1 \textbackslash{}\textbackslash{}times 0.02825 \textbackslash{}\textbackslash{}times 1 \textbackslash{}\textbackslash{}approx 0.0282" \},}
\NormalTok{  \{ op: "Sum the three terms",          eq: "P(X \textbackslash{}\textbackslash{}geq 8) \textbackslash{}\textbackslash{}approx 0.2335 + 0.1211 + 0.0282 = 0.3828", note: "About 38\% probability at least 8 patients are cured." \}}
\NormalTok{])}
\end{Highlighting}
\end{Shaded}

\subsection{Example 3 --- Normal z-score
(education)}\label{example-3-normal-z-score-education}

End-of-year exam scores in a large district are approximately normally
distributed with mean \(\mu = 65\) and standard deviation
\(\sigma = 12\). What fraction of students scored above 80?

We want \(P(X > 80)\).

\textbf{Step 1: Standardise.}

\[z = \frac{80 - 65}{12} = \frac{15}{12} = 1.25\]

\textbf{Step 2: Look up \(P(Z \leq 1.25)\).}

From a standard normal table (or calculator):
\(P(Z \leq 1.25) \approx 0.8944\).

\textbf{Step 3: Convert to the required probability.}

\[P(X > 80) = 1 - P(X \leq 80) = 1 - 0.8944 = 0.1056\]

About 10.6\% of students scored above 80.

A quick check with the 68-95-99.7 rule: 80 is \(\frac{15}{12} = 1.25\)
standard deviations above the mean, which is between 1\(\sigma\) and
2\(\sigma\). The rule says about 16\% of values lie beyond 1\(\sigma\),
so 10.6\% beyond 1.25\(\sigma\) is consistent.

\begin{Shaded}
\begin{Highlighting}[]
\NormalTok{//| echo: false}
\NormalTok{makeStepperHTML(3, [}
\NormalTok{  \{ op: "State the distribution",       eq: "X \textbackslash{}\textbackslash{}sim N(65,\textbackslash{}\textbackslash{}, 144),\textbackslash{}\textbackslash{}quad \textbackslash{}\textbackslash{}mu=65,\textbackslash{}\textbackslash{}; \textbackslash{}\textbackslash{}sigma=12" \},}
\NormalTok{  \{ op: "Standardise x = 80",           eq: "z = \textbackslash{}\textbackslash{}frac\{80 {-} 65\}\{12\} = \textbackslash{}\textbackslash{}frac\{15\}\{12\} = 1.25" \},}
\NormalTok{  \{ op: "Look up cumulative prob",      eq: "P(Z \textbackslash{}\textbackslash{}leq 1.25) \textbackslash{}\textbackslash{}approx 0.8944", note: "From a standard normal table or calculator." \},}
\NormalTok{  \{ op: "Use the complement",           eq: "P(X \textgreater{} 80) = 1 {-} 0.8944 = 0.1056", note: "About 10.6\% of students scored above 80." \},}
\NormalTok{  \{ op: "Sanity check",                 eq: "1.25\textbackslash{}\textbackslash{}sigma \textgreater{} \textbackslash{}\textbackslash{}mu \textbackslash{}\textbackslash{}Rightarrow \textbackslash{}\textbackslash{}text\{between 16\textbackslash{}\textbackslash{}\% and 2.5\textbackslash{}\textbackslash{}\% tail\} \textbackslash{}\textbackslash{}checkmark", note: "10.6\% sits between those bounds, so the answer is consistent." \}}
\NormalTok{])}
\end{Highlighting}
\end{Shaded}

\subsection{Example 4 --- Normal reverse lookup (pass with
distinction)}\label{example-4-normal-reverse-lookup-pass-with-distinction}

Using the same exam distribution (\(\mu = 65\), \(\sigma = 12\)), the
school awards ``distinction'' to the top 10\% of students. What score is
the distinction threshold?

We want the value \(x\) such that \(P(X > x) = 0.10\), i.e.,
\(P(X \leq x) = 0.90\).

\textbf{Step 1: Find the z-score for the 90th percentile.}

From a standard normal table, find \(z\) such that
\(P(Z \leq z) = 0.90\).

Looking up 0.90: \(z \approx 1.282\).

\textbf{Step 2: Unstandardise --- convert back to the original scale.}

Since \(z = \frac{x - \mu}{\sigma}\), rearranging gives:

\[x = \mu + z\sigma = 65 + 1.282 \times 12 = 65 + 15.38 \approx 80.4\]

A score of about 80 is the distinction threshold.

Note the structure: for a forward problem you go \(x \to z \to\)
probability. For a reverse problem you go probability \(\to z \to x\).
The formula \(x = \mu + z\sigma\) is just the standardisation formula
solved for \(x\).

\begin{Shaded}
\begin{Highlighting}[]
\NormalTok{//| echo: false}
\NormalTok{makeStepperHTML(4, [}
\NormalTok{  \{ op: "State the goal",               eq: "P(X \textgreater{} x) = 0.10 \textbackslash{}\textbackslash{}Leftrightarrow P(X \textbackslash{}\textbackslash{}leq x) = 0.90" \},}
\NormalTok{  \{ op: "Find the 90th{-}percentile z",   eq: "z\_\{0.90\} \textbackslash{}\textbackslash{}approx 1.282", note: "Look up 0.90 in the body of a standard normal table." \},}
\NormalTok{  \{ op: "Unstandardise",               eq: "x = \textbackslash{}\textbackslash{}mu + z\textbackslash{}\textbackslash{}sigma = 65 + 1.282 \textbackslash{}\textbackslash{}times 12", note: "Rearrange z = (x − μ)/σ to get x = μ + zσ." \},}
\NormalTok{  \{ op: "Compute the threshold",        eq: "x = 65 + 15.38 \textbackslash{}\textbackslash{}approx 80.4", note: "A score of about 80 is the distinction threshold." \}}
\NormalTok{])}
\end{Highlighting}
\end{Shaded}

\section{Where this goes}\label{where-this-goes-23}

\textbf{Statistical inference (Chapter 3)} is the direct next step. A
distribution gives you the expected shape of data from a random process.
Inference asks the reverse question: your actual data has a shape --- is
that consistent with what you'd expect from the assumed distribution, or
is something else going on? Every hypothesis test you'll ever run is
answering that question.

\textbf{Volume 7 probability} builds the mathematical foundations
underneath what you've been using here. Why does the normal distribution
have exactly the bell-curve formula it does? What does ``continuous
probability'' mean rigorously? The answers require calculus ---
integration gives you the area under the bell curve --- but you can use
the results now.

\section{Applications}\label{applications-3}

\begin{tcolorbox}[enhanced jigsaw, title=\textcolor{quarto-callout-note-color}{\faInfo}\hspace{0.5em}{Note}, opacityback=0, colback=white, breakable, leftrule=.75mm, arc=.35mm, toptitle=1mm, left=2mm, colframe=quarto-callout-note-color-frame, coltitle=black, titlerule=0mm, colbacktitle=quarto-callout-note-color!10!white, opacitybacktitle=0.6, toprule=.15mm, bottomtitle=1mm, bottomrule=.15mm, rightrule=.15mm]

\textbf{Where these distributions appear in practice}

\textbf{A/B testing.} A website tests two versions of a button. 200
users see version A, 200 see version B. Is the difference in click rates
real or just chance? Binomial distributions tell you what click-rate
differences to expect from chance alone.

\textbf{Quality control charts.} Manufacturing uses Shewhart control
charts: plot each batch measurement, mark lines at \(\mu \pm 2\sigma\)
and \(\mu \pm 3\sigma\). A point outside the \(3\sigma\) line has less
than a 0.3\% chance of occurring naturally --- strong evidence something
has gone wrong in the process.

\textbf{Weather forecasting.} A ``70\% chance of rain'' is a probability
from a distribution over possible outcomes. Forecast \emph{skill} ---
whether forecasts are actually calibrated --- is assessed using
statistics built on the normal distribution.

\textbf{Clinical trial power.} Before running a drug trial,
statisticians calculate how many patients they need to have a reasonable
chance of detecting a real effect if one exists. That calculation
requires knowing the distribution of outcomes under both the null
hypothesis and the alternative. You're one chapter away from doing this.

\end{tcolorbox}

\section{Exercises}\label{exercises-23}

\textbf{Exercise 1.} A student guesses randomly on a 12-question
true/false quiz. Let \(X\) be the number of correct answers.

\begin{enumerate}
\def\labelenumi{(\alph{enumi})}
\tightlist
\item
  State the distribution of \(X\), giving \(n\) and \(p\).
\item
  Calculate \(P(X = 6)\).
\item
  Calculate the mean and standard deviation of \(X\).
\end{enumerate}

\begin{Shaded}
\begin{Highlighting}[]
\NormalTok{//| echo: false}
\NormalTok{makeStepperHTML(5, [}
\NormalTok{  \{ op: "Identify the distribution",    eq: "X \textbackslash{}\textbackslash{}sim B(12,\textbackslash{}\textbackslash{}, 0.5)" \},}
\NormalTok{  \{ op: "Write the formula for P(X=6)", eq: "P(X=6) = \textbackslash{}\textbackslash{}binom\{12\}\{6\}(0.5)\^{}6(0.5)\^{}6 = \textbackslash{}\textbackslash{}binom\{12\}\{6\}(0.5)\^{}\{12\}" \},}
\NormalTok{  \{ op: "Compute the combination",      eq: "\textbackslash{}\textbackslash{}binom\{12\}\{6\} = \textbackslash{}\textbackslash{}frac\{12!\}\{6!\textbackslash{}\textbackslash{},6!\} = 924" \},}
\NormalTok{  \{ op: "Compute (0.5)\^{}\{12\}",           eq: "(0.5)\^{}\{12\} = \textbackslash{}\textbackslash{}frac\{1\}\{4096\} \textbackslash{}\textbackslash{}approx 0.000244" \},}
\NormalTok{  \{ op: "Multiply",                     eq: "P(X=6) = 924 \textbackslash{}\textbackslash{}times 0.000244 \textbackslash{}\textbackslash{}approx 0.2256", note: "About 22.6\% chance of getting exactly 6 correct by pure guessing." \},}
\NormalTok{  \{ op: "Mean and standard deviation",  eq: "\textbackslash{}\textbackslash{}mu = np = 12 \textbackslash{}\textbackslash{}times 0.5 = 6,\textbackslash{}\textbackslash{}quad \textbackslash{}\textbackslash{}sigma = \textbackslash{}\textbackslash{}sqrt\{12 \textbackslash{}\textbackslash{}times 0.5 \textbackslash{}\textbackslash{}times 0.5\} = \textbackslash{}\textbackslash{}sqrt\{3\} \textbackslash{}\textbackslash{}approx 1.73" \}}
\NormalTok{])}
\end{Highlighting}
\end{Shaded}

\begin{center}\rule{0.5\linewidth}{0.5pt}\end{center}

\textbf{Exercise 2.} In a city, 20\% of households own an electric
vehicle. A researcher surveys 15 randomly chosen households.

\begin{enumerate}
\def\labelenumi{(\alph{enumi})}
\tightlist
\item
  Let \(X\) be the number of households with an electric vehicle. State
  the distribution of \(X\).
\item
  Find \(P(X = 3)\).
\item
  Find \(P(X \leq 1)\) using the complement or direct calculation.
\item
  Find the mean and variance of \(X\).
\end{enumerate}

\begin{Shaded}
\begin{Highlighting}[]
\NormalTok{//| echo: false}
\NormalTok{makeStepperHTML(6, [}
\NormalTok{  \{ op: "Identify the distribution",    eq: "X \textbackslash{}\textbackslash{}sim B(15,\textbackslash{}\textbackslash{}, 0.2)" \},}
\NormalTok{  \{ op: "Formula for P(X = 3)",         eq: "P(X=3) = \textbackslash{}\textbackslash{}binom\{15\}\{3\}(0.2)\^{}3(0.8)\^{}\{12\}" \},}
\NormalTok{  \{ op: "Compute combination",          eq: "\textbackslash{}\textbackslash{}binom\{15\}\{3\} = \textbackslash{}\textbackslash{}frac\{15 \textbackslash{}\textbackslash{}times 14 \textbackslash{}\textbackslash{}times 13\}\{3 \textbackslash{}\textbackslash{}times 2 \textbackslash{}\textbackslash{}times 1\} = 455" \},}
\NormalTok{  \{ op: "Compute each power",           eq: "(0.2)\^{}3 = 0.008,\textbackslash{}\textbackslash{}quad (0.8)\^{}\{12\} \textbackslash{}\textbackslash{}approx 0.06872" \},}
\NormalTok{  \{ op: "Multiply",                     eq: "P(X=3) = 455 \textbackslash{}\textbackslash{}times 0.008 \textbackslash{}\textbackslash{}times 0.06872 \textbackslash{}\textbackslash{}approx 0.2501" \},}
\NormalTok{  \{ op: "P(X ≤ 1) directly",            eq: "P(X=0) = (0.8)\^{}\{15\} \textbackslash{}\textbackslash{}approx 0.03518,\textbackslash{}\textbackslash{}quad P(X=1) = 15(0.2)(0.8)\^{}\{14\} \textbackslash{}\textbackslash{}approx 0.13194" \},}
\NormalTok{  \{ op: "Sum for P(X ≤ 1)",             eq: "P(X \textbackslash{}\textbackslash{}leq 1) \textbackslash{}\textbackslash{}approx 0.03518 + 0.13194 = 0.1671", note: "About 16.7\% chance of at most 1 EV household." \},}
\NormalTok{  \{ op: "Mean and variance",            eq: "\textbackslash{}\textbackslash{}mu = np = 15 \textbackslash{}\textbackslash{}times 0.2 = 3,\textbackslash{}\textbackslash{}quad \textbackslash{}\textbackslash{}sigma\^{}2 = np(1{-}p) = 15 \textbackslash{}\textbackslash{}times 0.2 \textbackslash{}\textbackslash{}times 0.8 = 2.4" \}}
\NormalTok{])}
\end{Highlighting}
\end{Shaded}

\begin{center}\rule{0.5\linewidth}{0.5pt}\end{center}

\textbf{Exercise 3.} A coin is known to be biased: it shows heads with
probability \(p = 0.4\). The coin is flipped 20 times.

\begin{enumerate}
\def\labelenumi{(\alph{enumi})}
\tightlist
\item
  Write down the mean and variance of the number of heads.
\item
  Without calculating individual probabilities, explain whether you
  would be surprised to get 14 heads. (Use \(\sigma\) to guide your
  reasoning.)
\end{enumerate}

\begin{Shaded}
\begin{Highlighting}[]
\NormalTok{//| echo: false}
\NormalTok{makeStepperHTML(7, [}
\NormalTok{  \{ op: "Identify the distribution",    eq: "X \textbackslash{}\textbackslash{}sim B(20,\textbackslash{}\textbackslash{}, 0.4)" \},}
\NormalTok{  \{ op: "Mean",                         eq: "\textbackslash{}\textbackslash{}mu = np = 20 \textbackslash{}\textbackslash{}times 0.4 = 8" \},}
\NormalTok{  \{ op: "Variance",                     eq: "\textbackslash{}\textbackslash{}sigma\^{}2 = np(1{-}p) = 20 \textbackslash{}\textbackslash{}times 0.4 \textbackslash{}\textbackslash{}times 0.6 = 4.8" \},}
\NormalTok{  \{ op: "Standard deviation",          eq: "\textbackslash{}\textbackslash{}sigma = \textbackslash{}\textbackslash{}sqrt\{4.8\} \textbackslash{}\textbackslash{}approx 2.19" \},}
\NormalTok{  \{ op: "How far is 14 from the mean?", eq: "14 {-} 8 = 6 = \textbackslash{}\textbackslash{}frac\{6\}\{2.19\} \textbackslash{}\textbackslash{}approx 2.74\textbackslash{}\textbackslash{},\textbackslash{}\textbackslash{}sigma \textbackslash{}\textbackslash{}text\{ above the mean\}" \},}
\NormalTok{  \{ op: "Assess surprise",              eq: "P(|X {-} \textbackslash{}\textbackslash{}mu| \textbackslash{}\textbackslash{}geq 2.74\textbackslash{}\textbackslash{}sigma) \textless{} 0.006", note: "A value almost 3σ from the mean is very unusual — you would be surprised to see 14 heads with a p = 0.4 coin." \}}
\NormalTok{])}
\end{Highlighting}
\end{Shaded}

\begin{center}\rule{0.5\linewidth}{0.5pt}\end{center}

\textbf{Exercise 4.} The resting heart rate of adults in a large study
is approximately normally distributed with mean \(\mu = 72\) beats per
minute and standard deviation \(\sigma = 10\) bpm.

\begin{enumerate}
\def\labelenumi{(\alph{enumi})}
\tightlist
\item
  Find the probability that a randomly chosen adult has a resting heart
  rate above 90 bpm.
\item
  Find the probability that a randomly chosen adult has a resting heart
  rate between 62 and 82 bpm. (You may use the 68-95-99.7 rule for this
  part.)
\item
  What percentage of adults have resting heart rates below 55 bpm?
\end{enumerate}

\begin{Shaded}
\begin{Highlighting}[]
\NormalTok{//| echo: false}
\NormalTok{makeStepperHTML(8, [}
\NormalTok{  \{ op: "State the distribution",       eq: "X \textbackslash{}\textbackslash{}sim N(72,\textbackslash{}\textbackslash{}, 100),\textbackslash{}\textbackslash{}quad \textbackslash{}\textbackslash{}mu=72,\textbackslash{}\textbackslash{}; \textbackslash{}\textbackslash{}sigma=10" \},}
\NormalTok{  \{ op: "(a) Standardise x = 90",       eq: "z = \textbackslash{}\textbackslash{}frac\{90 {-} 72\}\{10\} = 1.8" \},}
\NormalTok{  \{ op: "(a) Look up and complement",   eq: "P(X \textgreater{} 90) = 1 {-} P(Z \textbackslash{}\textbackslash{}leq 1.8) = 1 {-} 0.9641 = 0.0359", note: "About 3.6\% of adults have a resting heart rate above 90 bpm." \},}
\NormalTok{  \{ op: "(b) Recognise the interval",   eq: "[62,\textbackslash{}\textbackslash{},82] = [\textbackslash{}\textbackslash{}mu {-} \textbackslash{}\textbackslash{}sigma,\textbackslash{}\textbackslash{}; \textbackslash{}\textbackslash{}mu + \textbackslash{}\textbackslash{}sigma]", note: "62 = 72 − 10 and 82 = 72 + 10, so this is exactly 1σ either side." \},}
\NormalTok{  \{ op: "(b) Apply empirical rule",     eq: "P(62 \textless{} X \textless{} 82) \textbackslash{}\textbackslash{}approx 0.68", note: "68\% of adults have resting heart rates in this range." \},}
\NormalTok{  \{ op: "(c) Standardise x = 55",       eq: "z = \textbackslash{}\textbackslash{}frac\{55 {-} 72\}\{10\} = {-}1.7" \},}
\NormalTok{  \{ op: "(c) Look up",                  eq: "P(X \textless{} 55) = P(Z \textless{} {-}1.7) = 1 {-} P(Z \textbackslash{}\textbackslash{}leq 1.7) \textbackslash{}\textbackslash{}approx 1 {-} 0.9554 = 0.0446", note: "About 4.5\% of adults have resting heart rates below 55 bpm." \}}
\NormalTok{])}
\end{Highlighting}
\end{Shaded}

\begin{center}\rule{0.5\linewidth}{0.5pt}\end{center}

\textbf{Exercise 5.} Graduate entry to a university programme requires a
score in the top 5\% of a standardised test. Scores are normally
distributed with \(\mu = 500\) and \(\sigma = 100\).

\begin{enumerate}
\def\labelenumi{(\alph{enumi})}
\tightlist
\item
  Find the z-score corresponding to the 95th percentile. (Use
  \(z_{0.95} \approx 1.645\).)
\item
  What is the minimum score needed for entry?
\item
  A student scores 620. What percentile does this correspond to? (Find
  \(z\), then use a table or calculator.)
\end{enumerate}

\begin{Shaded}
\begin{Highlighting}[]
\NormalTok{//| echo: false}
\NormalTok{makeStepperHTML(9, [}
\NormalTok{  \{ op: "State the distribution",       eq: "X \textbackslash{}\textbackslash{}sim N(500,\textbackslash{}\textbackslash{}, 100\^{}2),\textbackslash{}\textbackslash{}quad \textbackslash{}\textbackslash{}mu=500,\textbackslash{}\textbackslash{}; \textbackslash{}\textbackslash{}sigma=100" \},}
\NormalTok{  \{ op: "(a) 95th{-}percentile z{-}score",  eq: "z\_\{0.95\} \textbackslash{}\textbackslash{}approx 1.645", note: "Given in the problem; top 5\% means P(Z ≤ z) = 0.95." \},}
\NormalTok{  \{ op: "(b) Unstandardise",           eq: "x = \textbackslash{}\textbackslash{}mu + z\textbackslash{}\textbackslash{}sigma = 500 + 1.645 \textbackslash{}\textbackslash{}times 100 = 664.5" \},}
\NormalTok{  \{ op: "(b) Minimum entry score",      eq: "x \textbackslash{}\textbackslash{}approx 665", note: "A score of at least 665 is needed to be in the top 5\%." \},}
\NormalTok{  \{ op: "(c) Standardise score 620",    eq: "z = \textbackslash{}\textbackslash{}frac\{620 {-} 500\}\{100\} = 1.2" \},}
\NormalTok{  \{ op: "(c) Look up CDF",             eq: "P(Z \textbackslash{}\textbackslash{}leq 1.2) \textbackslash{}\textbackslash{}approx 0.8849", note: "A score of 620 sits at roughly the 88th percentile." \}}
\NormalTok{])}
\end{Highlighting}
\end{Shaded}

\begin{center}\rule{0.5\linewidth}{0.5pt}\end{center}

\textbf{Exercise 6.} A marine biologist is studying kelp-eating urchins
on a reef. She records the number of urchins found in each of 50
one-metre quadrats. Historical data for this reef type suggests urchin
counts per quadrat follow a distribution with \(\mu = 8\) and
\(\sigma = 2.8\).

\begin{enumerate}
\def\labelenumi{(\alph{enumi})}
\tightlist
\item
  A count of 15 urchins is recorded in one quadrat. How many standard
  deviations above the mean is this?
\item
  Using the 68-95-99.7 rule, roughly what fraction of quadrats would you
  expect to contain more than 13.6 urchins?
\item
  The biologist suspects a disturbance has shifted urchin density. She
  finds 12 quadrats (out of 50) with counts above 13.6. Based on your
  answer to (b), does this seem consistent with the historical
  distribution? Explain briefly.
\end{enumerate}

\begin{Shaded}
\begin{Highlighting}[]
\NormalTok{//| echo: false}
\NormalTok{makeStepperHTML(10, [}
\NormalTok{  \{ op: "State the parameters",         eq: "\textbackslash{}\textbackslash{}mu = 8,\textbackslash{}\textbackslash{}quad \textbackslash{}\textbackslash{}sigma = 2.8" \},}
\NormalTok{  \{ op: "(a) z{-}score for count 15",     eq: "z = \textbackslash{}\textbackslash{}frac\{15 {-} 8\}\{2.8\} = \textbackslash{}\textbackslash{}frac\{7\}\{2.8\} = 2.5", note: "A count of 15 is 2.5 standard deviations above the mean." \},}
\NormalTok{  \{ op: "(b) Identify 13.6 in σ units", eq: "13.6 = \textbackslash{}\textbackslash{}mu + 2\textbackslash{}\textbackslash{}sigma = 8 + 2 \textbackslash{}\textbackslash{}times 2.8 = 13.6" \},}
\NormalTok{  \{ op: "(b) Apply 95\% rule",           eq: "P(\textbackslash{}\textbackslash{}mu {-} 2\textbackslash{}\textbackslash{}sigma \textless{} X \textless{} \textbackslash{}\textbackslash{}mu + 2\textbackslash{}\textbackslash{}sigma) \textbackslash{}\textbackslash{}approx 0.95" \},}
\NormalTok{  \{ op: "(b) Tail fraction",            eq: "P(X \textgreater{} \textbackslash{}\textbackslash{}mu + 2\textbackslash{}\textbackslash{}sigma) \textbackslash{}\textbackslash{}approx \textbackslash{}\textbackslash{}frac\{1 {-} 0.95\}\{2\} = 0.025", note: "About 2.5\% of quadrats expected above 13.6 under historical conditions." \},}
\NormalTok{  \{ op: "(c) Expected vs observed",     eq: "0.025 \textbackslash{}\textbackslash{}times 50 = 1.25 \textbackslash{}\textbackslash{}text\{ quadrats expected\};\textbackslash{}\textbackslash{}quad 12 \textbackslash{}\textbackslash{}text\{ observed\}" \},}
\NormalTok{  \{ op: "(c) Conclusion",               eq: "\textbackslash{}\textbackslash{}frac\{12\}\{50\} = 24\textbackslash{}\textbackslash{}\% \textbackslash{}\textbackslash{}gg 2.5\textbackslash{}\textbackslash{}\%", note: "12 quadrats is far more than the 1–2 expected by chance. This is strong evidence the distribution has shifted — a disturbance is a plausible explanation." \}}
\NormalTok{])}
\end{Highlighting}
\end{Shaded}

\chapter{Statistical Inference}\label{statistical-inference}

From sample to conclusion

\hfill\break

A study makes the news: a new energy drink improves reaction time by 12
milliseconds on average in a sample of 40 university students. The
company behind it claims the drink genuinely works. A sceptic says the
12 ms gap could just be luck.

Both people are looking at the same number. How can they reach opposite
conclusions?

The answer is that a single number from a single sample cannot settle
the question on its own. What matters is what numbers like that look
like when the drink \emph{does nothing} --- when the only thing
producing variation is chance. If 12 ms is the kind of gap you'd often
see from chance alone, the result is unremarkable. If 12 ms is extremely
unusual under chance, something else is probably going on.

This chapter is about how to make that reasoning precise. The tools are
called \textbf{statistical inference} --- the machinery for drawing
conclusions about populations from samples, while being honest about
uncertainty.

\section{Populations and samples}\label{populations-and-samples}

A \textbf{population} is the complete set of individuals or measurements
you care about. A \textbf{sample} is the subset you can actually
observe.

In most real situations, measuring the entire population is impossible.
You cannot test every batch of tablets that will ever leave a factory.
You cannot ask every voter in an election what they think today. You
cannot measure the resting heart rate of every adult in the country. So
you take a sample --- as many observations as you can afford --- and
reason from there.

The fundamental challenge is that the sample will not perfectly reflect
the population. Take a different sample and you'll get a different mean.
Take enough different samples and the sample means form their own
distribution --- one that carries useful information about the
population.

A few terms used throughout:

\begin{itemize}
\tightlist
\item
  \(N\) (or sometimes the full description): the \textbf{population
  size} --- total number of individuals
\item
  \(n\): the \textbf{sample size} --- how many you actually measured
\item
  \(\mu\) (mu): the \textbf{population mean} --- what you want to know
  but usually cannot compute directly
\item
  \(\bar{x}\) (x-bar): the \textbf{sample mean} --- what you computed
  from your \(n\) observations
\item
  \(\sigma\): the \textbf{population standard deviation} --- spread of
  the whole population
\item
  \(s\): the \textbf{sample standard deviation} --- spread estimated
  from your sample
\end{itemize}

The goal of inference is to say something credible about \(\mu\) using
\(\bar{x}\), knowing that \(\bar{x}\) is a good guess but not an exact
answer.

\subsection{What makes a sample
useful}\label{what-makes-a-sample-useful}

Not every sample tells you much. A sample of convenience --- asking only
your friends, measuring only the easiest-to-reach parts of a forest,
testing only patients who volunteer --- may be biased in ways that
invalidate any conclusion about the broader population.

A \textbf{representative sample} is one where each member of the
population has a fair chance of being included. The cleanest version is
a \textbf{simple random sample}: every member of the population is
equally likely to be chosen, and choices are independent. In practice,
this is an ideal rather than a guarantee, but it is the assumption that
makes everything in this chapter valid. Whenever you see a claim built
on inference, ask: was the sample actually representative?

\section{The sampling distribution}\label{the-sampling-distribution}

Suppose heights in a population are normally distributed with mean
\(\mu = 170\) cm and standard deviation \(\sigma = 8\) cm. You take a
random sample of \(n = 25\) people and compute their mean height
\(\bar{x}\).

Now imagine doing this many times: take 25 people, compute \(\bar{x}\),
record it. Take another 25, compute \(\bar{x}\), record it. Repeat this
thousands of times.

What would you see if you plotted all those sample means?

Three things, and all three are predictable:

\begin{enumerate}
\def\labelenumi{\arabic{enumi}.}
\item
  The sample means cluster around \(\mu = 170\). No surprise --- if your
  sampling is unbiased, the average of the averages should track the
  true population mean.
\item
  The sample means are \emph{less} spread out than individual heights.
  Averaging smooths out extremes. A single person might be 155 cm or 185
  cm --- that's plausible. But a \emph{mean} of 25 people being 155 cm
  would require nearly everyone in the sample to be very short, which is
  far less likely.
\item
  The distribution of sample means is approximately normal --- even if
  the population itself is not perfectly normal.
\end{enumerate}

This distribution --- the distribution of \(\bar{x}\) across all
possible samples of size \(n\) --- is called the \textbf{sampling
distribution of the mean}.

\begin{Shaded}
\begin{Highlighting}[]
\NormalTok{//| echo: false}
\NormalTok{// Sampling Distribution Simulator}
\NormalTok{// Reader sets μ, σ, n; draws 1000 samples and watches x̄ pile up into}
\NormalTok{// a histogram with the theoretical N(μ, σ²/n) curve overlaid.}

\NormalTok{viewof sim\_mu = Inputs.range([50, 150], \{value: 100, step: 1, label: "Population mean μ"\})}
\end{Highlighting}
\end{Shaded}

\begin{Shaded}
\begin{Highlighting}[]
\NormalTok{//| echo: false}
\NormalTok{viewof sim\_sigma = Inputs.range([5, 30], \{value: 12, step: 1, label: "Population SD σ"\})}
\end{Highlighting}
\end{Shaded}

\begin{Shaded}
\begin{Highlighting}[]
\NormalTok{//| echo: false}
\NormalTok{viewof sim\_n = Inputs.range([1, 100], \{value: 25, step: 1, label: "Sample size n"\})}
\end{Highlighting}
\end{Shaded}

\begin{Shaded}
\begin{Highlighting}[]
\NormalTok{//| echo: false}
\NormalTok{viewof sim\_drawBtn = Inputs.button("Draw 1000 samples")}
\end{Highlighting}
\end{Shaded}

\begin{Shaded}
\begin{Highlighting}[]
\NormalTok{//| echo: false}
\NormalTok{\{}
\NormalTok{  // Accumulate sample means across button presses}
\NormalTok{  const stored = \{ means: [] \};}

\NormalTok{  // Re{-}run whenever drawBtn is clicked or parameters change}
\NormalTok{  sim\_drawBtn;}
\NormalTok{  const mu = sim\_mu, sigma = sim\_sigma, n = sim\_n;}

\NormalTok{  // Box{-}Muller normal random generator}
\NormalTok{  function randNorm(mu, sigma) \{}
\NormalTok{    let u, v;}
\NormalTok{    do \{ u = Math.random(); \} while (u === 0);}
\NormalTok{    v = Math.random();}
\NormalTok{    return mu + sigma * Math.sqrt({-}2 * Math.log(u)) * Math.cos(2 * Math.PI * v);}
\NormalTok{  \}}

\NormalTok{  // Draw 1000 new sample means and append}
\NormalTok{  for (let i = 0; i \textless{} 1000; i++) \{}
\NormalTok{    let sum = 0;}
\NormalTok{    for (let j = 0; j \textless{} n; j++) sum += randNorm(mu, sigma);}
\NormalTok{    stored.means.push(sum / n);}
\NormalTok{  \}}

\NormalTok{  const means = stored.means;}
\NormalTok{  const SE = sigma / Math.sqrt(n);}
\NormalTok{  const totalDrawn = means.length;}

\NormalTok{  // Layout}
\NormalTok{  const W = 600, H = 320;}
\NormalTok{  const marginL = 55, marginR = 20, marginT = 30, marginB = 50;}
\NormalTok{  const plotW = W {-} marginL {-} marginR;}
\NormalTok{  const plotH = H {-} marginT {-} marginB;}

\NormalTok{  // Axis range: μ ± 4*SE but at least μ ± 3*sigma/sqrt(1)}
\NormalTok{  const xMin = mu {-} 4.5 * Math.max(SE, 1);}
\NormalTok{  const xMax = mu + 4.5 * Math.max(SE, 1);}

\NormalTok{  // Histogram bins}
\NormalTok{  const nBins = 50;}
\NormalTok{  const binW = (xMax {-} xMin) / nBins;}
\NormalTok{  const counts = new Array(nBins).fill(0);}
\NormalTok{  for (const m of means) \{}
\NormalTok{    const b = Math.floor((m {-} xMin) / binW);}
\NormalTok{    if (b \textgreater{}= 0 \&\& b \textless{} nBins) counts[b]++;}
\NormalTok{  \}}
\NormalTok{  const maxCount = Math.max(...counts, 1);}

\NormalTok{  // Scale helpers}
\NormalTok{  const xScale = x =\textgreater{} marginL + ((x {-} xMin) / (xMax {-} xMin)) * plotW;}
\NormalTok{  const yScaleCount = c =\textgreater{} marginT + plotH {-} (c / maxCount) * plotH;}

\NormalTok{  // Theoretical normal PDF, scaled to match histogram}
\NormalTok{  function normalPDF(x, mean, sd) \{}
\NormalTok{    return Math.exp({-}0.5 * ((x {-} mean) / sd) ** 2) / (sd * Math.sqrt(2 * Math.PI));}
\NormalTok{  \}}
\NormalTok{  const pdfScale = totalDrawn * binW; // converts density → expected count}

\NormalTok{  // Build SVG}
\NormalTok{  const svg = d3.create("svg").attr("width", W).attr("height", H)}
\NormalTok{    .attr("viewBox", \textasciigrave{}0 0 $\{W\} $\{H\}\textasciigrave{})}
\NormalTok{    .style("font{-}family", "inherit")}
\NormalTok{    .style("overflow", "visible");}

\NormalTok{  // Background}
\NormalTok{  svg.append("rect").attr("x", marginL).attr("y", marginT)}
\NormalTok{    .attr("width", plotW).attr("height", plotH)}
\NormalTok{    .attr("fill", "\#f9fafb").attr("rx", 3);}

\NormalTok{  // Histogram bars}
\NormalTok{  for (let i = 0; i \textless{} nBins; i++) \{}
\NormalTok{    const bx = xMin + i * binW;}
\NormalTok{    svg.append("rect")}
\NormalTok{      .attr("x", xScale(bx) + 0.5)}
\NormalTok{      .attr("y", yScaleCount(counts[i]))}
\NormalTok{      .attr("width", Math.max(1, xScale(bx + binW) {-} xScale(bx) {-} 1))}
\NormalTok{      .attr("height", marginT + plotH {-} yScaleCount(counts[i]))}
\NormalTok{      .attr("fill", "\#93c5fd")}
\NormalTok{      .attr("stroke", "\#60a5fa")}
\NormalTok{      .attr("stroke{-}width", 0.4);}
\NormalTok{  \}}

\NormalTok{  // Theoretical curve N(μ, σ²/n)}
\NormalTok{  const curvePoints = [];}
\NormalTok{  const steps = 200;}
\NormalTok{  for (let i = 0; i \textless{}= steps; i++) \{}
\NormalTok{    const cx = xMin + (i / steps) * (xMax {-} xMin);}
\NormalTok{    const cy = normalPDF(cx, mu, SE) * pdfScale;}
\NormalTok{    const scaledY = marginT + plotH {-} (cy / maxCount) * plotH;}
\NormalTok{    curvePoints.push([xScale(cx), Math.max(marginT, scaledY)]);}
\NormalTok{  \}}
\NormalTok{  const lineGen = d3.line().x(d =\textgreater{} d[0]).y(d =\textgreater{} d[1]).curve(d3.curveBasis);}
\NormalTok{  svg.append("path").attr("d", lineGen(curvePoints))}
\NormalTok{    .attr("fill", "none").attr("stroke", "\#1d4ed8").attr("stroke{-}width", 2);}

\NormalTok{  // Mean line}
\NormalTok{  svg.append("line")}
\NormalTok{    .attr("x1", xScale(mu)).attr("x2", xScale(mu))}
\NormalTok{    .attr("y1", marginT).attr("y2", marginT + plotH)}
\NormalTok{    .attr("stroke", "\#6b7280").attr("stroke{-}width", 1).attr("stroke{-}dasharray", "4,3");}

\NormalTok{  // x{-}axis}
\NormalTok{  const xTicks = 6;}
\NormalTok{  const xStep = (xMax {-} xMin) / xTicks;}
\NormalTok{  svg.append("line").attr("x1", marginL).attr("x2", marginL + plotW)}
\NormalTok{    .attr("y1", marginT + plotH).attr("y2", marginT + plotH)}
\NormalTok{    .attr("stroke", "\#374151").attr("stroke{-}width", 1);}
\NormalTok{  for (let i = 0; i \textless{}= xTicks; i++) \{}
\NormalTok{    const val = xMin + i * xStep;}
\NormalTok{    const px = xScale(val);}
\NormalTok{    svg.append("line").attr("x1", px).attr("x2", px)}
\NormalTok{      .attr("y1", marginT + plotH).attr("y2", marginT + plotH + 4)}
\NormalTok{      .attr("stroke", "\#374151");}
\NormalTok{    svg.append("text").attr("x", px).attr("y", marginT + plotH + 16)}
\NormalTok{      .attr("text{-}anchor", "middle").attr("font{-}size", 11).attr("fill", "\#374151")}
\NormalTok{      .text(val.toFixed(1));}
\NormalTok{  \}}
\NormalTok{  svg.append("text").attr("x", marginL + plotW / 2).attr("y", H {-} 6)}
\NormalTok{    .attr("text{-}anchor", "middle").attr("font{-}size", 12).attr("fill", "\#374151")}
\NormalTok{    .text("Sample mean x̄");}

\NormalTok{  // y{-}axis label}
\NormalTok{  svg.append("text")}
\NormalTok{    .attr("transform", \textasciigrave{}translate(14, $\{marginT + plotH / 2\}) rotate({-}90)\textasciigrave{})}
\NormalTok{    .attr("text{-}anchor", "middle").attr("font{-}size", 12).attr("fill", "\#374151")}
\NormalTok{    .text("Count");}

\NormalTok{  // Info panel}
\NormalTok{  const infoX = marginL + plotW {-} 160, infoY = marginT + 8;}
\NormalTok{  svg.append("rect").attr("x", infoX {-} 6).attr("y", infoY {-} 4)}
\NormalTok{    .attr("width", 164).attr("height", 62).attr("rx", 4)}
\NormalTok{    .attr("fill", "white").attr("stroke", "\#e5e7eb");}
\NormalTok{  svg.append("text").attr("x", infoX).attr("y", infoY + 10)}
\NormalTok{    .attr("font{-}size", 12).attr("fill", "\#1d4ed8").attr("font{-}weight", "600")}
\NormalTok{    .text(\textasciigrave{}SE = σ/√n = $\{SE.toFixed(3)\}\textasciigrave{});}
\NormalTok{  svg.append("text").attr("x", infoX).attr("y", infoY + 28)}
\NormalTok{    .attr("font{-}size", 11).attr("fill", "\#374151")}
\NormalTok{    .text(\textasciigrave{}μ = $\{mu\},  σ = $\{sigma\},  n = $\{n\}\textasciigrave{});}
\NormalTok{  svg.append("text").attr("x", infoX).attr("y", infoY + 46)}
\NormalTok{    .attr("font{-}size", 11).attr("fill", "\#6b7280")}
\NormalTok{    .text(\textasciigrave{}$\{totalDrawn.toLocaleString()\} sample means drawn\textasciigrave{});}

\NormalTok{  // Legend}
\NormalTok{  svg.append("rect").attr("x", marginL + 6).attr("y", marginT + 8)}
\NormalTok{    .attr("width", 12).attr("height", 12).attr("fill", "\#93c5fd").attr("stroke", "\#60a5fa");}
\NormalTok{  svg.append("text").attr("x", marginL + 22).attr("y", marginT + 19)}
\NormalTok{    .attr("font{-}size", 11).attr("fill", "\#374151").text("Simulated x̄ values");}
\NormalTok{  svg.append("line").attr("x1", marginL + 6).attr("x2", marginL + 18)}
\NormalTok{    .attr("y1", marginT + 33).attr("y2", marginT + 33)}
\NormalTok{    .attr("stroke", "\#1d4ed8").attr("stroke{-}width", 2);}
\NormalTok{  svg.append("text").attr("x", marginL + 22).attr("y", marginT + 37)}
\NormalTok{    .attr("font{-}size", 11).attr("fill", "\#374151").text("Theoretical N(μ, σ²/n)");}

\NormalTok{  // Title}
\NormalTok{  svg.append("text").attr("x", W / 2).attr("y", marginT {-} 10)}
\NormalTok{    .attr("text{-}anchor", "middle").attr("font{-}size", 13).attr("fill", "\#111827").attr("font{-}weight", "600")}
\NormalTok{    .text("Sampling distribution of x̄");}

\NormalTok{  return svg.node();}
\NormalTok{\}}
\end{Highlighting}
\end{Shaded}

\section{The Central Limit Theorem}\label{the-central-limit-theorem}

The three observations above are not separate facts. They are
consequences of a single result called the \textbf{Central Limit
Theorem} (CLT).

\begin{tcolorbox}[enhanced jigsaw, title=\textcolor{quarto-callout-note-color}{\faInfo}\hspace{0.5em}{Note}, opacityback=0, colback=white, breakable, leftrule=.75mm, arc=.35mm, toptitle=1mm, left=2mm, colframe=quarto-callout-note-color-frame, coltitle=black, titlerule=0mm, colbacktitle=quarto-callout-note-color!10!white, opacitybacktitle=0.6, toprule=.15mm, bottomtitle=1mm, bottomrule=.15mm, rightrule=.15mm]

\textbf{The Central Limit Theorem}

If you take random samples of size \(n\) from a population with mean
\(\mu\) and standard deviation \(\sigma\), then for sufficiently large
\(n\), the sampling distribution of \(\bar{x}\) is approximately normal:

\[\bar{x} \sim N\!\left(\mu,\, \frac{\sigma^2}{n}\right)\]

The mean of the sampling distribution equals \(\mu\). The standard
deviation of the sampling distribution is \(\dfrac{\sigma}{\sqrt{n}}\).

\end{tcolorbox}

The standard deviation of the sampling distribution has its own name:
the \textbf{standard error} (SE).

\[SE = \frac{\sigma}{\sqrt{n}}\]

The standard error is not the spread of individual measurements --- it
is the spread of \emph{sample means}. It tells you how much \(\bar{x}\)
typically varies from one sample to the next.

\subsection{Why this is surprising}\label{why-this-is-surprising}

The CLT makes no assumption about the shape of the population
distribution. Whether the individual measurements are skewed, uniform,
bimodal, or anything else, the sample means pile up into a bell curve
once \(n\) is large enough.

This is the reason the normal distribution appears so often in
statistical work. It is not because every quantity in nature is normally
distributed. It is because \emph{sample means} tend to be normally
distributed, and most of the quantities we measure and compare are
averages of some kind.

As a rough guide, \(n \geq 30\) is often enough for the CLT to give a
good approximation. For symmetric populations, even smaller samples work
well. For heavily skewed populations, larger \(n\) may be needed.

\subsection{Working with the standard error --- a numerical
example}\label{working-with-the-standard-error-a-numerical-example}

Return to the height example: \(\mu = 170\) cm, \(\sigma = 8\) cm.

\textbf{For samples of size \(n = 25\):}

\[SE = \frac{8}{\sqrt{25}} = \frac{8}{5} = 1.6 \text{ cm}\]

The sampling distribution is \(\bar{x} \sim N(170, 1.6^2)\). A sample
mean of 173 cm is:

\[z = \frac{173 - 170}{1.6} = \frac{3}{1.6} = 1.875 \text{ standard errors above the mean}\]

\textbf{For samples of size \(n = 100\):}

\[SE = \frac{8}{\sqrt{100}} = \frac{8}{10} = 0.8 \text{ cm}\]

The same sample mean of 173 cm is now:

\[z = \frac{173 - 170}{0.8} = \frac{3}{0.8} = 3.75 \text{ standard errors above the mean}\]

With \(n = 25\), a sample mean of 173 cm is not very surprising --- it
is less than 2 standard errors out. With \(n = 100\), the same 173 cm is
very surprising --- nearly 4 standard errors out. Larger samples give
you more information, so unusual results become more informative.

This is the key insight: \textbf{more data makes your estimate more
precise}, and the improvement scales as \(1/\sqrt{n}\), not \(1/n\). To
halve the standard error, you need four times as much data.

\section{Confidence intervals}\label{confidence-intervals}

You take a sample, compute \(\bar{x}\), and want to say something about
\(\mu\). A single number --- ``\(\mu\) is approximately 170 cm'' --- is
called a \textbf{point estimate}. It is useful but incomplete, because
it gives no indication of how uncertain you are.

An \textbf{interval estimate} gives a range of plausible values for
\(\mu\). A \textbf{confidence interval} attaches a stated level of
confidence to that range --- typically 90\%, 95\%, or 99\%.

\subsection{Building a 95\% confidence
interval}\label{building-a-95-confidence-interval}

From Chapter 2, you know that for a normal distribution, 95\% of values
fall within approximately 1.96 standard deviations of the mean.
Precisely:

\[P(-1.96 \leq Z \leq 1.96) = 0.95\]

where \(Z \sim N(0,1)\).

Since \(\bar{x} \sim N(\mu, \sigma^2/n)\), standardising gives:

\[P\!\left(-1.96 \leq \frac{\bar{x} - \mu}{SE} \leq 1.96\right) = 0.95\]

Rearranging the inequality to isolate \(\mu\): multiply through by SE to
get \(\bar{x} - 1.96 \cdot SE \leq \mu \leq \bar{x} + 1.96 \cdot SE\).
This is the same algebra as rearranging \(-1.96 \leq (x-5)/2 \leq 1.96\)
for \(x\) --- multiply through by 2, then add 5 throughout.

\[P\!\left(\bar{x} - 1.96 \cdot SE \leq \mu \leq \bar{x} + 1.96 \cdot SE\right) = 0.95\]

This is the \textbf{95\% confidence interval for \(\mu\)} (when
\(\sigma\) is known):

\[\boxed{\bar{x} \pm 1.96 \cdot \frac{\sigma}{\sqrt{n}}}\]

The value 1.96 is called the \textbf{critical value} for 95\%
confidence. For other confidence levels, the critical value changes:

\begin{longtable}[]{@{}ll@{}}
\toprule\noalign{}
Confidence level & Critical value \(z^*\) \\
\midrule\noalign{}
\endhead
\bottomrule\noalign{}
\endlastfoot
90\% & 1.645 \\
95\% & 1.960 \\
99\% & 2.576 \\
\end{longtable}

\subsection{Worked example --- confidence
interval}\label{worked-example-confidence-interval}

A nutritionist measures the daily sodium intake (mg) of a random sample
of \(n = 36\) adults from a region. The sample mean is
\(\bar{x} = 2{,}450\) mg. Population studies suggest \(\sigma = 360\) mg
for daily sodium intake.

Construct a 95\% confidence interval for the true mean daily sodium
intake \(\mu\) in this region.

\textbf{Step 1: Calculate the standard error.}

\[SE = \frac{\sigma}{\sqrt{n}} = \frac{360}{\sqrt{36}} = \frac{360}{6} = 60 \text{ mg}\]

\textbf{Step 2: Find the margin of error.}

\[\text{Margin of error} = 1.96 \times SE = 1.96 \times 60 = 117.6 \text{ mg}\]

\textbf{Step 3: Construct the interval.}

\[\bar{x} \pm \text{margin of error} = 2{,}450 \pm 117.6\]

\[\text{CI: } (2{,}332.4 \text{ mg},\ 2{,}567.6 \text{ mg})\]

\textbf{Step 4: Interpret.}

\begin{quote}
We are 95\% confident that the true mean daily sodium intake in this
region is between 2,332 mg and 2,568 mg.
\end{quote}

\subsection{What ``95\% confident'' actually
means}\label{what-95-confident-actually-means}

This is one of the most commonly misunderstood ideas in statistics.

\textbf{The incorrect interpretation:} ``There is a 95\% probability
that \(\mu\) lies in this interval.''

This sounds reasonable but is wrong. The population mean \(\mu\) is a
fixed (if unknown) number --- it does not have a probability
distribution. It either is or is not in the interval. The interval
itself is random, because \(\bar{x}\) varies from sample to sample.

\textbf{The correct interpretation:} If you were to repeat this sampling
procedure many times --- each time collecting 36 adults, computing
\(\bar{x}\), and constructing the interval --- then 95\% of the
resulting intervals would contain the true \(\mu\).

In other words, confidence describes the procedure, not any single
interval. The 95\% refers to the long-run success rate of the method.

In practice, you have one interval from one sample, and you cannot know
whether yours is one of the 95\% that work or one of the 5\% that don't.
The confidence level tells you the method is reliable, not that any
particular result is guaranteed.

\begin{tcolorbox}[enhanced jigsaw, title=\textcolor{quarto-callout-tip-color}{\faLightbulb}\hspace{0.5em}{Tip}, opacityback=0, colback=white, breakable, leftrule=.75mm, arc=.35mm, toptitle=1mm, left=2mm, colframe=quarto-callout-tip-color-frame, coltitle=black, titlerule=0mm, colbacktitle=quarto-callout-tip-color!10!white, opacitybacktitle=0.6, toprule=.15mm, bottomtitle=1mm, bottomrule=.15mm, rightrule=.15mm]

\textbf{A useful sentence template for interpreting confidence
intervals}

``We are {[}level{]}\% confident that the true population mean
{[}quantity{]} is between {[}lower{]} and {[}upper{]} {[}units{]}.''

Avoid: ``There is a 95\% chance the mean is in this interval.'' Use:
``We are 95\% confident the mean is in this interval.''

\end{tcolorbox}

\section{Hypothesis testing}\label{hypothesis-testing}

A confidence interval tells you what values of \(\mu\) are plausible
given your data. A \textbf{hypothesis test} answers a more direct
question: is the evidence strong enough to reject a specific claim about
\(\mu\)?

\subsection{The logic of the test}\label{the-logic-of-the-test}

The framework borrows from legal reasoning. A court presumes innocence
until the evidence makes guilt sufficiently clear. In statistics:

\begin{itemize}
\tightlist
\item
  You start by assuming a specific null state --- the \textbf{null
  hypothesis} \(H_0\).
\item
  You ask: if the null hypothesis were true, how likely is it that you'd
  see data at least as extreme as what you observed?
\item
  If the answer is ``very unlikely,'' the data is evidence against
  \(H_0\).
\end{itemize}

What counts as ``very unlikely'' is a threshold you set in advance,
called the \textbf{significance level} \(\alpha\). The most common
choice is \(\alpha = 0.05\).

\subsection{Setting up the hypotheses}\label{setting-up-the-hypotheses}

Every test involves two competing hypotheses:

\textbf{The null hypothesis \(H_0\)} is the claim being tested. It
usually represents ``no effect,'' ``no difference,'' or a specific
benchmark value. You write it as an equality: \(H_0: \mu = \mu_0\),
where \(\mu_0\) is the claimed value.

\textbf{The alternative hypothesis \(H_1\)} (sometimes written \(H_a\))
is what you believe might be true instead. It can be:

\begin{itemize}
\tightlist
\item
  \textbf{Two-tailed:} \(H_1: \mu \neq \mu_0\) --- the mean could be
  either higher or lower than claimed
\item
  \textbf{One-tailed (upper):} \(H_1: \mu > \mu_0\) --- you suspect the
  mean is higher
\item
  \textbf{One-tailed (lower):} \(H_1: \mu < \mu_0\) --- you suspect the
  mean is lower
\end{itemize}

The choice of one- or two-tailed test should reflect your research
question \emph{before} you look at the data. Choosing the direction
after seeing the data inflates the apparent significance and undermines
the logic of the test.

\subsection{The test statistic}\label{the-test-statistic}

If \(H_0: \mu = \mu_0\) is true, then the sampling distribution of
\(\bar{x}\) is:

\[\bar{x} \sim N\!\left(\mu_0,\, \frac{\sigma^2}{n}\right)\]

Standardising gives the \textbf{z-test statistic}:

\[z = \frac{\bar{x} - \mu_0}{\sigma / \sqrt{n}}\]

This \(z\) value measures how many standard errors your observed sample
mean is from the null hypothesis value. If \(|z|\) is small, your data
is consistent with \(H_0\). If \(|z|\) is large, your data is unusual
under \(H_0\).

\subsection{The p-value}\label{the-p-value}

The \textbf{p-value} is the probability of observing a test statistic at
least as extreme as the one you computed, \emph{if \(H_0\) is true}.

\begin{itemize}
\tightlist
\item
  For a two-tailed test (\(H_1: \mu \neq \mu_0\)):
  \(p\text{-value} = 2 \times P(Z \geq |z|)\)
\item
  For an upper-tailed test (\(H_1: \mu > \mu_0\)):
  \(p\text{-value} = P(Z \geq z)\)
\item
  For a lower-tailed test (\(H_1: \mu < \mu_0\)):
  \(p\text{-value} = P(Z \leq z)\)
\end{itemize}

A small p-value means your observed data would be rare under \(H_0\) ---
the data is hard to explain by chance alone.

\subsection{The decision rule}\label{the-decision-rule}

Compare the p-value to the significance level \(\alpha\):

\begin{itemize}
\tightlist
\item
  If \(p\text{-value} \leq \alpha\): \textbf{reject \(H_0\)}. The data
  is sufficiently unusual under \(H_0\) to conclude there is evidence
  for \(H_1\).
\item
  If \(p\text{-value} > \alpha\): \textbf{fail to reject \(H_0\)}. The
  data is consistent with \(H_0\) --- not proof that \(H_0\) is true,
  only that there is insufficient evidence to reject it.
\end{itemize}

Note the careful language: you \emph{fail to reject}, not \emph{accept}
\(H_0\). Absence of evidence is not evidence of absence.

\begin{Shaded}
\begin{Highlighting}[]
\NormalTok{//| echo: false}
\NormalTok{// Hypothesis Test Explorer}
\NormalTok{// Reader sets μ₀, x̄, σ, n; chooses tail type and α.}
\NormalTok{// The viz draws N(μ₀, σ²/n), marks x̄, shades rejection region,}
\NormalTok{// and shows z, p{-}value, and the reject/fail{-}to{-}reject decision.}

\NormalTok{viewof ht\_mu0 = Inputs.range([50, 250], \{value: 650, step: 1, label: "Null mean μ₀"\})}
\end{Highlighting}
\end{Shaded}

\begin{Shaded}
\begin{Highlighting}[]
\NormalTok{//| echo: false}
\NormalTok{viewof ht\_xbar = Inputs.range([50, 250], \{value: 672, step: 0.1, label: "Observed mean x̄"\})}
\end{Highlighting}
\end{Shaded}

\begin{Shaded}
\begin{Highlighting}[]
\NormalTok{//| echo: false}
\NormalTok{viewof ht\_sigma = Inputs.range([1, 100], \{value: 80, step: 1, label: "Population SD σ"\})}
\end{Highlighting}
\end{Shaded}

\begin{Shaded}
\begin{Highlighting}[]
\NormalTok{//| echo: false}
\NormalTok{viewof ht\_n = Inputs.range([2, 200], \{value: 40, step: 1, label: "Sample size n"\})}
\end{Highlighting}
\end{Shaded}

\begin{Shaded}
\begin{Highlighting}[]
\NormalTok{//| echo: false}
\NormalTok{viewof ht\_tail = Inputs.select(["two{-}tailed", "upper (H₁: μ \textgreater{} μ₀)", "lower (H₁: μ \textless{} μ₀)"], \{}
\NormalTok{  value: "upper (H₁: μ \textgreater{} μ₀)",}
\NormalTok{  label: "Test type"}
\NormalTok{\})}
\end{Highlighting}
\end{Shaded}

\begin{Shaded}
\begin{Highlighting}[]
\NormalTok{//| echo: false}
\NormalTok{viewof ht\_alpha = Inputs.select([0.01, 0.05, 0.10], \{value: 0.05, label: "Significance level α"\})}
\end{Highlighting}
\end{Shaded}

\begin{Shaded}
\begin{Highlighting}[]
\NormalTok{//| echo: false}
\NormalTok{\{}
\NormalTok{  const mu0 = ht\_mu0, xbar = ht\_xbar, sigma = ht\_sigma, n = ht\_n;}
\NormalTok{  const alpha = ht\_alpha;}
\NormalTok{  const tail = ht\_tail;}

\NormalTok{  const SE = sigma / Math.sqrt(n);}
\NormalTok{  const z = (xbar {-} mu0) / SE;}

\NormalTok{  // Standard normal CDF via error function approximation}
\NormalTok{  function erf(x) \{}
\NormalTok{    const a1=0.254829592, a2={-}0.284496736, a3=1.421413741, a4={-}1.453152027, a5=1.061405429, p=0.3275911;}
\NormalTok{    const sign = x \textless{} 0 ? {-}1 : 1;}
\NormalTok{    const t = 1 / (1 + p * Math.abs(x));}
\NormalTok{    const y = 1 {-} (((((a5*t + a4)*t + a3)*t + a2)*t + a1)*t) * Math.exp({-}x*x);}
\NormalTok{    return sign * y;}
\NormalTok{  \}}
\NormalTok{  function normalCDF(x) \{ return 0.5 * (1 + erf(x / Math.sqrt(2))); \}}
\NormalTok{  function normalPDF(x, mu, sd) \{}
\NormalTok{    return Math.exp({-}0.5 * ((x {-} mu) / sd) ** 2) / (sd * Math.sqrt(2 * Math.PI));}
\NormalTok{  \}}

\NormalTok{  // p{-}value}
\NormalTok{  let pval;}
\NormalTok{  if (tail === "two{-}tailed") pval = 2 * (1 {-} normalCDF(Math.abs(z)));}
\NormalTok{  else if (tail.startsWith("upper")) pval = 1 {-} normalCDF(z);}
\NormalTok{  else pval = normalCDF(z);}
\NormalTok{  pval = Math.max(0, Math.min(1, pval));}

\NormalTok{  // Critical z values}
\NormalTok{  let zCritUpper = Infinity, zCritLower = {-}Infinity;}
\NormalTok{  if (tail === "two{-}tailed") \{}
\NormalTok{    zCritUpper = {-}normalCDFInv(alpha / 2);}
\NormalTok{    zCritLower = normalCDFInv(alpha / 2);}
\NormalTok{  \} else if (tail.startsWith("upper")) \{}
\NormalTok{    zCritUpper = {-}normalCDFInv(alpha);}
\NormalTok{  \} else \{}
\NormalTok{    zCritLower = normalCDFInv(alpha);}
\NormalTok{  \}}

\NormalTok{  // Inverse normal CDF (bisection) — good enough for critical value display}
\NormalTok{  function normalCDFInv(p) \{}
\NormalTok{    let lo = {-}10, hi = 10;}
\NormalTok{    for (let i = 0; i \textless{} 60; i++) \{}
\NormalTok{      const mid = (lo + hi) / 2;}
\NormalTok{      if (normalCDF(mid) \textless{} p) lo = mid; else hi = mid;}
\NormalTok{    \}}
\NormalTok{    return (lo + hi) / 2;}
\NormalTok{  \}}

\NormalTok{  // Fix: compute critical values after defining normalCDFInv}
\NormalTok{  let critLo = {-}Infinity, critHi = Infinity;}
\NormalTok{  if (tail === "two{-}tailed") \{}
\NormalTok{    critLo = normalCDFInv(alpha / 2);}
\NormalTok{    critHi = normalCDFInv(1 {-} alpha / 2);}
\NormalTok{  \} else if (tail.startsWith("upper")) \{}
\NormalTok{    critHi = normalCDFInv(1 {-} alpha);}
\NormalTok{  \} else \{}
\NormalTok{    critLo = normalCDFInv(alpha);}
\NormalTok{  \}}

\NormalTok{  const decision = pval \textless{}= alpha ? "Reject H₀" : "Fail to reject H₀";}
\NormalTok{  const decisionColor = pval \textless{}= alpha ? "\#dc2626" : "\#059669";}

\NormalTok{  // Layout}
\NormalTok{  const W = 620, H = 320;}
\NormalTok{  const mL = 55, mR = 20, mT = 40, mB = 55;}
\NormalTok{  const pW = W {-} mL {-} mR, pH = H {-} mT {-} mB;}

\NormalTok{  // x{-}axis range: μ₀ ± 4 SE, ensure xbar is visible}
\NormalTok{  const spread = Math.max(4 * SE, Math.abs(xbar {-} mu0) + 2 * SE);}
\NormalTok{  const xMin = mu0 {-} spread;}
\NormalTok{  const xMax = mu0 + spread;}

\NormalTok{  const xScale = x =\textgreater{} mL + ((x {-} xMin) / (xMax {-} xMin)) * pW;}
\NormalTok{  const maxPDF = normalPDF(mu0, mu0, SE);}
\NormalTok{  const yScale = y =\textgreater{} mT + pH {-} (y / maxPDF) * pH * 0.92;}

\NormalTok{  const svg = d3.create("svg").attr("width", W).attr("height", H)}
\NormalTok{    .attr("viewBox", \textasciigrave{}0 0 $\{W\} $\{H\}\textasciigrave{})}
\NormalTok{    .style("font{-}family", "inherit").style("overflow", "visible");}

\NormalTok{  svg.append("rect").attr("x", mL).attr("y", mT).attr("width", pW).attr("height", pH)}
\NormalTok{    .attr("fill", "\#f9fafb").attr("rx", 3);}

\NormalTok{  // Build curve path points}
\NormalTok{  const steps = 400;}
\NormalTok{  function curvePoints(xFrom, xTo) \{}
\NormalTok{    const pts = [];}
\NormalTok{    for (let i = 0; i \textless{}= steps; i++) \{}
\NormalTok{      const cx = xFrom + (i / steps) * (xTo {-} xFrom);}
\NormalTok{      pts.push([xScale(cx), yScale(normalPDF(cx, mu0, SE))]);}
\NormalTok{    \}}
\NormalTok{    return pts;}
\NormalTok{  \}}

\NormalTok{  const lineGen = d3.line().x(d =\textgreater{} d[0]).y(d =\textgreater{} d[1]).curve(d3.curveBasis);}

\NormalTok{  // Shade rejection region(s)}
\NormalTok{  function shadedArea(xFrom, xTo, color) \{}
\NormalTok{    const pts = curvePoints(xFrom, xTo);}
\NormalTok{    const area = [...pts, [xScale(xTo), mT + pH], [xScale(xFrom), mT + pH]];}
\NormalTok{    svg.append("polygon")}
\NormalTok{      .attr("points", area.map(p =\textgreater{} p.join(",")).join(" "))}
\NormalTok{      .attr("fill", color).attr("opacity", 0.35);}
\NormalTok{  \}}

\NormalTok{  const rejColor = "\#fca5a5";}
\NormalTok{  if (tail === "two{-}tailed") \{}
\NormalTok{    shadedArea(xMin, mu0 + critLo * SE, rejColor);}
\NormalTok{    shadedArea(mu0 + critHi * SE, xMax, rejColor);}
\NormalTok{  \} else if (tail.startsWith("upper")) \{}
\NormalTok{    shadedArea(mu0 + critHi * SE, xMax, rejColor);}
\NormalTok{  \} else \{}
\NormalTok{    shadedArea(xMin, mu0 + critLo * SE, rejColor);}
\NormalTok{  \}}

\NormalTok{  // Main curve}
\NormalTok{  svg.append("path").attr("d", lineGen(curvePoints(xMin, xMax)))}
\NormalTok{    .attr("fill", "none").attr("stroke", "\#1d4ed8").attr("stroke{-}width", 2.5);}

\NormalTok{  // Critical value lines}
\NormalTok{  function critLine(xVal, label) \{}
\NormalTok{    if (!isFinite(xVal)) return;}
\NormalTok{    const px = xScale(xVal);}
\NormalTok{    svg.append("line").attr("x1", px).attr("x2", px).attr("y1", mT).attr("y2", mT + pH)}
\NormalTok{      .attr("stroke", "\#dc2626").attr("stroke{-}width", 1.5).attr("stroke{-}dasharray", "5,3");}
\NormalTok{    svg.append("text").attr("x", px).attr("y", mT {-} 6)}
\NormalTok{      .attr("text{-}anchor", "middle").attr("font{-}size", 10).attr("fill", "\#dc2626")}
\NormalTok{      .text(label);}
\NormalTok{  \}}
\NormalTok{  if (tail === "two{-}tailed") \{}
\NormalTok{    critLine(mu0 + critLo * SE, \textasciigrave{}z* = $\{critLo.toFixed(2)\}\textasciigrave{});}
\NormalTok{    critLine(mu0 + critHi * SE, \textasciigrave{}z* = $\{critHi.toFixed(2)\}\textasciigrave{});}
\NormalTok{  \} else if (tail.startsWith("upper")) \{}
\NormalTok{    critLine(mu0 + critHi * SE, \textasciigrave{}z* = $\{critHi.toFixed(2)\}\textasciigrave{});}
\NormalTok{  \} else \{}
\NormalTok{    critLine(mu0 + critLo * SE, \textasciigrave{}z* = $\{critLo.toFixed(2)\}\textasciigrave{});}
\NormalTok{  \}}

\NormalTok{  // Observed x̄ line}
\NormalTok{  const xbarPx = xScale(xbar);}
\NormalTok{  svg.append("line").attr("x1", xbarPx).attr("x2", xbarPx).attr("y1", mT).attr("y2", mT + pH)}
\NormalTok{    .attr("stroke", "\#7c3aed").attr("stroke{-}width", 2);}
\NormalTok{  svg.append("text").attr("x", xbarPx).attr("y", mT + pH + 16)}
\NormalTok{    .attr("text{-}anchor", "middle").attr("font{-}size", 11).attr("fill", "\#7c3aed")}
\NormalTok{    .text(\textasciigrave{}x̄ = $\{xbar\}\textasciigrave{});}

\NormalTok{  // μ₀ line}
\NormalTok{  const mu0Px = xScale(mu0);}
\NormalTok{  svg.append("line").attr("x1", mu0Px).attr("x2", mu0Px).attr("y1", mT).attr("y2", mT + pH)}
\NormalTok{    .attr("stroke", "\#374151").attr("stroke{-}width", 1).attr("stroke{-}dasharray", "3,3");}
\NormalTok{  svg.append("text").attr("x", mu0Px).attr("y", mT + pH + 28)}
\NormalTok{    .attr("text{-}anchor", "middle").attr("font{-}size", 10).attr("fill", "\#374151")}
\NormalTok{    .text(\textasciigrave{}μ₀ = $\{mu0\}\textasciigrave{});}

\NormalTok{  // x{-}axis}
\NormalTok{  svg.append("line").attr("x1", mL).attr("x2", mL + pW)}
\NormalTok{    .attr("y1", mT + pH).attr("y2", mT + pH).attr("stroke", "\#374151");}
\NormalTok{  const nTicks = 7;}
\NormalTok{  for (let i = 0; i \textless{}= nTicks; i++) \{}
\NormalTok{    const val = xMin + (i / nTicks) * (xMax {-} xMin);}
\NormalTok{    const px = xScale(val);}
\NormalTok{    svg.append("line").attr("x1", px).attr("x2", px)}
\NormalTok{      .attr("y1", mT + pH).attr("y2", mT + pH + 4).attr("stroke", "\#374151");}
\NormalTok{    svg.append("text").attr("x", px).attr("y", mT + pH + 14)}
\NormalTok{      .attr("text{-}anchor", "middle").attr("font{-}size", 10).attr("fill", "\#374151")}
\NormalTok{      .text(val.toFixed(1));}
\NormalTok{  \}}

\NormalTok{  // Info box}
\NormalTok{  const iX = mL + 8, iY = mT + 6;}
\NormalTok{  svg.append("rect").attr("x", iX {-} 4).attr("y", iY {-} 4)}
\NormalTok{    .attr("width", 210).attr("height", 80).attr("rx", 4)}
\NormalTok{    .attr("fill", "white").attr("stroke", "\#e5e7eb");}
\NormalTok{  svg.append("text").attr("x", iX).attr("y", iY + 12).attr("font{-}size", 12).attr("fill", "\#374151")}
\NormalTok{    .text(\textasciigrave{}SE = $\{SE.toFixed(3)\}\textasciigrave{});}
\NormalTok{  svg.append("text").attr("x", iX).attr("y", iY + 28).attr("font{-}size", 12).attr("fill", "\#374151")}
\NormalTok{    .text(\textasciigrave{}z = $\{z.toFixed(3)\}\textasciigrave{});}
\NormalTok{  svg.append("text").attr("x", iX).attr("y", iY + 44).attr("font{-}size", 12).attr("fill", "\#374151")}
\NormalTok{    .text(\textasciigrave{}p{-}value = $\{pval \textless{} 0.001 ? "\textless{} 0.001" : pval.toFixed(4)\}\textasciigrave{});}
\NormalTok{  svg.append("text").attr("x", iX).attr("y", iY + 62).attr("font{-}size", 13)}
\NormalTok{    .attr("fill", decisionColor).attr("font{-}weight", "700").text(decision);}

\NormalTok{  // Title}
\NormalTok{  svg.append("text").attr("x", W / 2).attr("y", mT {-} 18)}
\NormalTok{    .attr("text{-}anchor", "middle").attr("font{-}size", 13).attr("fill", "\#111827").attr("font{-}weight", "600")}
\NormalTok{    .text("Hypothesis test: N(μ₀, σ²/n) sampling distribution");}

\NormalTok{  return svg.node();}
\NormalTok{\}}
\end{Highlighting}
\end{Shaded}

\section{Worked example --- full hypothesis
test}\label{worked-example-full-hypothesis-test}

The school cafeteria claims that a standard portion of their pasta
contains 650 mg of sodium. A nutritionist suspects the actual amount is
higher. She measures 40 randomly chosen portions and finds a sample mean
of \(\bar{x} = 672\) mg. From manufacturer data, \(\sigma = 80\) mg.

Test the nutritionist's claim at the 5\% significance level.

\textbf{Step 1: State the hypotheses.}

The cafeteria claims \(\mu = 650\) mg. The nutritionist suspects it is
\emph{higher}, so this is a one-tailed (upper) test.

\[H_0: \mu = 650 \quad H_1: \mu > 650\]

\textbf{Step 2: Choose the significance level.}

\[\alpha = 0.05\]

\textbf{Step 3: Compute the test statistic.}

\[z = \frac{\bar{x} - \mu_0}{\sigma / \sqrt{n}} = \frac{672 - 650}{80 / \sqrt{40}} = \frac{22}{80/6.325} = \frac{22}{12.65} \approx 1.74\]

\textbf{Step 4: Find the p-value.}

For an upper-tailed test:

\[p\text{-value} = P(Z \geq 1.74)\]

From a standard normal table: \(P(Z \leq 1.74) \approx 0.9591\).

\[p\text{-value} = 1 - 0.9591 = 0.0409\]

\textbf{Step 5: Make a decision.}

\(p\text{-value} = 0.041 < \alpha = 0.05\), so we \textbf{reject
\(H_0\)}.

\textbf{Step 6: State the conclusion in plain language.}

\begin{quote}
At the 5\% significance level, there is sufficient evidence to conclude
that the mean sodium content of cafeteria pasta portions exceeds the
claimed 650 mg.
\end{quote}

Notice what this conclusion does \emph{not} say. It does not say the
true mean is 672 mg. It does not say the cafeteria is lying. It says the
evidence is strong enough to act on the suspicion that something is off.
The strength of that conclusion depends on the assumptions: random
sampling, known \(\sigma\), and that the significance level was chosen
before looking at the data.

\section{Type I and Type II errors}\label{type-i-and-type-ii-errors}

Every hypothesis test makes a binary decision under uncertainty. Two
kinds of mistakes are possible:

\begin{longtable}[]{@{}
  >{\raggedright\arraybackslash}p{(\linewidth - 4\tabcolsep) * \real{0.3333}}
  >{\raggedright\arraybackslash}p{(\linewidth - 4\tabcolsep) * \real{0.3333}}
  >{\raggedright\arraybackslash}p{(\linewidth - 4\tabcolsep) * \real{0.3333}}@{}}
\toprule\noalign{}
\begin{minipage}[b]{\linewidth}\raggedright
\end{minipage} & \begin{minipage}[b]{\linewidth}\raggedright
\(H_0\) is actually true
\end{minipage} & \begin{minipage}[b]{\linewidth}\raggedright
\(H_0\) is actually false
\end{minipage} \\
\midrule\noalign{}
\endhead
\bottomrule\noalign{}
\endlastfoot
\textbf{Reject \(H_0\)} & \textbf{Type I error} (false positive) &
Correct decision \\
\textbf{Fail to reject \(H_0\)} & Correct decision & \textbf{Type II
error} (false negative) \\
\end{longtable}

A \textbf{Type I error} is rejecting \(H_0\) when it is true ---
concluding there is an effect when there is none. The probability of a
Type I error is exactly \(\alpha\), the significance level you chose.
Setting \(\alpha = 0.05\) means you accept a 5\% chance of a false
positive if \(H_0\) is true.

A \textbf{Type II error} is failing to reject \(H_0\) when it is false
--- missing a real effect. Its probability is denoted \(\beta\). The
\textbf{power} of a test is \(1 - \beta\): the probability of correctly
detecting a true effect.

\subsection{Why this tradeoff matters}\label{why-this-tradeoff-matters}

Reducing \(\alpha\) (demanding stronger evidence before rejecting
\(H_0\)) decreases the chance of a false positive but increases the
chance of a false negative --- you become more conservative and more
likely to miss real effects.

The stakes differ by context. A Type I error in a drug trial means
approving an ineffective drug --- real harm. A Type II error means
failing to approve an effective one --- also real harm, but different in
character. Before running a study, statisticians specify both \(\alpha\)
and the desired power, then calculate the sample size \(n\) required to
achieve both. Larger samples reduce both types of error simultaneously,
which is why clinical trials are expensive.

\begin{tcolorbox}[enhanced jigsaw, title=\textcolor{quarto-callout-tip-color}{\faLightbulb}\hspace{0.5em}{Tip}, opacityback=0, colback=white, breakable, leftrule=.75mm, arc=.35mm, toptitle=1mm, left=2mm, colframe=quarto-callout-tip-color-frame, coltitle=black, titlerule=0mm, colbacktitle=quarto-callout-tip-color!10!white, opacitybacktitle=0.6, toprule=.15mm, bottomtitle=1mm, bottomrule=.15mm, rightrule=.15mm]

\textbf{A concrete framing}

A fire alarm that goes off every time someone burns toast (Type I error:
false positive) is annoying and teaches people to ignore it. A fire
alarm that never sounds during an actual fire (Type II error: false
negative) is catastrophic. Good alarm design --- and good statistical
design --- requires thinking carefully about which error is more costly.

\end{tcolorbox}

\section{Connecting to the bigger
picture}\label{connecting-to-the-bigger-picture}

Every time you see a published result with a p-value, you are looking at
the output of this machinery. The researcher had a null hypothesis,
collected data, computed a test statistic, and compared the p-value to a
threshold.

When \(\sigma\) is unknown and estimated from the sample (using \(s\)
instead), the test statistic no longer follows the standard normal
distribution. It follows a \textbf{t-distribution}, which is slightly
wider and depends on the sample size. For large \(n\), the difference is
small; for small samples, it matters. The t-distribution is the natural
next step beyond this chapter.

Two further extensions are essential for most real work:

\textbf{Comparing two groups.} Does treatment A work better than
treatment B? Here you test whether the \emph{difference} in means
\(\mu_1 - \mu_2 = 0\) rather than testing a single mean. The two-sample
z-test or t-test extends the framework from this chapter directly.

\textbf{Regression.} Once you move to modelling how one variable
predicts another, you are still doing inference --- testing whether the
slope of a line is zero, constructing confidence intervals around
predictions. The language of \(H_0\), test statistics, and p-values is
identical.

The machinery in this chapter is the foundation for all of it.

\section{Where this goes}\label{where-this-goes-24}

\begin{tcolorbox}[enhanced jigsaw, opacityback=0, colback=white, breakable, colframe=quarto-callout-note-color-frame, arc=.35mm, left=2mm, toprule=.15mm, leftrule=.75mm, bottomrule=.15mm, rightrule=.15mm]
\begin{minipage}[t]{5.5mm}
\textcolor{quarto-callout-note-color}{\faInfo}
\end{minipage}%
\begin{minipage}[t]{\textwidth - 5.5mm}

\textbf{What this chapter enables}

\begin{itemize}
\tightlist
\item
  \textbf{Computing (comp):} Machine learning validation --- train/test
  splits, significance of benchmark improvements --- is applied
  inference. Every time a researcher reports that model A beats model B,
  the question is whether the difference is real or sampling noise.
\item
  \textbf{Hard sciences (sci):} Every published experiment reports a
  p-value. Clinical trial analysis, drug approval decisions, and
  environmental monitoring all run on the framework you have just
  learned.
\item
  \textbf{Finance and business (biz):} A/B testing at scale, quality
  control charts, and actuarial inference are direct applications. The
  decision to change a product is a hypothesis test with real stakes.
\item
  \textbf{Geography and environment (geo):} Climate trend detection ---
  distinguishing a genuine warming signal from year-to-year variation
  --- requires formal inference. So does environmental monitoring for
  pollution events and species population changes.
\end{itemize}

\end{minipage}%
\end{tcolorbox}

\section{Exercises}\label{exercises-24}

\begin{tcolorbox}[enhanced jigsaw, opacityback=0, colback=white, breakable, colframe=quarto-callout-note-color-frame, arc=.35mm, left=2mm, toprule=.15mm, leftrule=.75mm, bottomrule=.15mm, rightrule=.15mm]
\begin{minipage}[t]{5.5mm}
\textcolor{quarto-callout-note-color}{\faInfo}
\end{minipage}%
\begin{minipage}[t]{\textwidth - 5.5mm}

\textbf{A note on these exercises}

All exercises in this chapter give you \(\sigma\), the population
standard deviation. In practice, \(\sigma\) is almost never known ---
you estimate it from the sample using \(s\). That extension (the
t-distribution) is the natural next step beyond this chapter. For now,
you will always be given \(\sigma\).

\end{minipage}%
\end{tcolorbox}

\begin{Shaded}
\begin{Highlighting}[]
\NormalTok{//| echo: false}
\NormalTok{katex = require("katex@0.16")}
\end{Highlighting}
\end{Shaded}

\begin{Shaded}
\begin{Highlighting}[]
\NormalTok{//| echo: false}
\NormalTok{function makeStepperHTML(exerciseNum, steps) \{}
\NormalTok{  let current = 0;}
\NormalTok{  const totalSteps = steps.length;}
\NormalTok{  const container = document.createElement("div");}
\NormalTok{  container.style.cssText = "border:1px solid \#e5e7eb; border{-}radius:8px; padding:1rem 1.25rem; margin:0.75rem 0; font{-}family:inherit;";}
\NormalTok{  const stepsDiv = document.createElement("div");}
\NormalTok{  stepsDiv.style.marginBottom = "0.75rem";}
\NormalTok{  function renderSteps() \{}
\NormalTok{    stepsDiv.innerHTML = "";}
\NormalTok{    for (let i = 0; i \textless{} current; i++) \{}
\NormalTok{      const s = steps[i];}
\NormalTok{      const row = document.createElement("div");}
\NormalTok{      row.style.cssText = "display:grid; grid{-}template{-}columns:200px 1fr; gap:0.5rem 1rem; align{-}items:baseline; padding:0.35rem 0; border{-}top:1px solid \#f3f4f6;";}
\NormalTok{      const opCell = document.createElement("span");}
\NormalTok{      opCell.style.cssText = "font{-}size:0.85em; color:\#6b7280; font{-}style:italic;";}
\NormalTok{      opCell.textContent = s.op;}
\NormalTok{      const eqCell = document.createElement("span");}
\NormalTok{      eqCell.innerHTML = katex.renderToString(s.eq, \{ throwOnError: false, displayMode: false \});}
\NormalTok{      row.appendChild(opCell);}
\NormalTok{      row.appendChild(eqCell);}
\NormalTok{      if (s.note) \{}
\NormalTok{        const noteRow = document.createElement("div");}
\NormalTok{        noteRow.style.cssText = "grid{-}column:2; font{-}size:0.82em; color:\#6b7280; padding{-}bottom:0.2rem;";}
\NormalTok{        noteRow.textContent = s.note;}
\NormalTok{        row.appendChild(noteRow);}
\NormalTok{      \}}
\NormalTok{      stepsDiv.appendChild(row);}
\NormalTok{    \}}
\NormalTok{    if (current === totalSteps) \{}
\NormalTok{      const done = document.createElement("div");}
\NormalTok{      done.style.cssText = "margin{-}top:0.5rem; font{-}size:0.9em; color:\#059669; font{-}weight:500;";}
\NormalTok{      done.textContent = "✓ Solution complete";}
\NormalTok{      stepsDiv.appendChild(done);}
\NormalTok{    \}}
\NormalTok{  \}}
\NormalTok{  const controls = document.createElement("div");}
\NormalTok{  controls.style.cssText = "display:flex; gap:0.5rem; align{-}items:center;";}
\NormalTok{  const nextBtn = document.createElement("button");}
\NormalTok{  nextBtn.textContent = "Next step →";}
\NormalTok{  nextBtn.style.cssText = "padding:0.35rem 0.85rem; border:1px solid \#d1d5db; border{-}radius:4px; background:\#fff; cursor:pointer; font{-}size:0.9em;";}
\NormalTok{  const resetBtn = document.createElement("button");}
\NormalTok{  resetBtn.textContent = "Reset";}
\NormalTok{  resetBtn.style.cssText = "padding:0.35rem 0.75rem; border:1px solid \#e5e7eb; border{-}radius:4px; background:\#f9fafb; cursor:pointer; font{-}size:0.9em; color:\#6b7280;";}
\NormalTok{  const counter = document.createElement("span");}
\NormalTok{  counter.style.cssText = "font{-}size:0.82em; color:\#9ca3af; margin{-}left:0.25rem;";}
\NormalTok{  function updateButtons() \{}
\NormalTok{    nextBtn.disabled = current === totalSteps;}
\NormalTok{    nextBtn.style.opacity = current === totalSteps ? "0.4" : "1";}
\NormalTok{    counter.textContent = current === 0 ? "Click to reveal steps" : \textasciigrave{}Step $\{current\} of $\{totalSteps\}\textasciigrave{};}
\NormalTok{  \}}
\NormalTok{  nextBtn.onclick = () =\textgreater{} \{ if (current \textless{} totalSteps) \{ current++; renderSteps(); updateButtons(); \} \};}
\NormalTok{  resetBtn.onclick = () =\textgreater{} \{ current = 0; renderSteps(); updateButtons(); \};}
\NormalTok{  controls.appendChild(nextBtn);}
\NormalTok{  controls.appendChild(resetBtn);}
\NormalTok{  controls.appendChild(counter);}
\NormalTok{  container.appendChild(stepsDiv);}
\NormalTok{  container.appendChild(controls);}
\NormalTok{  renderSteps();}
\NormalTok{  updateButtons();}
\NormalTok{  return container;}
\NormalTok{\}}
\end{Highlighting}
\end{Shaded}

\textbf{Exercise 1.} A population of salmon in a river system has a
known standard deviation of \(\sigma = 42\) mm in body length. A
researcher samples \(n\) fish from this population.

\begin{enumerate}
\def\labelenumi{(\alph{enumi})}
\tightlist
\item
  Calculate the standard error when \(n = 9\).
\item
  Calculate the standard error when \(n = 36\).
\item
  Calculate the standard error when \(n = 144\).
\item
  By what factor does the standard error change each time \(n\)
  increases by a factor of 4? Explain why this makes sense using the
  formula.
\end{enumerate}

\begin{Shaded}
\begin{Highlighting}[]
\NormalTok{//| echo: false}
\NormalTok{makeStepperHTML(1, [}
\NormalTok{  \{ op: "Formula", eq: "SE = \textbackslash{}\textbackslash{}dfrac\{\textbackslash{}\textbackslash{}sigma\}\{\textbackslash{}\textbackslash{}sqrt\{n\}\}" \},}
\NormalTok{  \{ op: "(a) n = 9", eq: "SE = \textbackslash{}\textbackslash{}dfrac\{42\}\{\textbackslash{}\textbackslash{}sqrt\{9\}\} = \textbackslash{}\textbackslash{}dfrac\{42\}\{3\} = 14 \textbackslash{}\textbackslash{}text\{ mm\}" \},}
\NormalTok{  \{ op: "(b) n = 36", eq: "SE = \textbackslash{}\textbackslash{}dfrac\{42\}\{\textbackslash{}\textbackslash{}sqrt\{36\}\} = \textbackslash{}\textbackslash{}dfrac\{42\}\{6\} = 7 \textbackslash{}\textbackslash{}text\{ mm\}" \},}
\NormalTok{  \{ op: "(c) n = 144", eq: "SE = \textbackslash{}\textbackslash{}dfrac\{42\}\{\textbackslash{}\textbackslash{}sqrt\{144\}\} = \textbackslash{}\textbackslash{}dfrac\{42\}\{12\} = 3.5 \textbackslash{}\textbackslash{}text\{ mm\}" \},}
\NormalTok{  \{ op: "(d) Factor each time", eq: "14 \textbackslash{}\textbackslash{}to 7 \textbackslash{}\textbackslash{}to 3.5: \textbackslash{}\textbackslash{}text\{ each time \} SE \textbackslash{}\textbackslash{}text\{ halves\}", note: "n increases by ×4 each time; √4 = 2, so SE divides by 2." \},}
\NormalTok{  \{ op: "General rule", eq: "SE \textbackslash{}\textbackslash{}propto \textbackslash{}\textbackslash{}dfrac\{1\}\{\textbackslash{}\textbackslash{}sqrt\{n\}\}: \textbackslash{}\textbackslash{}text\{ quadruple \} n \textbackslash{}\textbackslash{}Rightarrow \textbackslash{}\textbackslash{}text\{halve \} SE" \}}
\NormalTok{])}
\end{Highlighting}
\end{Shaded}

\begin{center}\rule{0.5\linewidth}{0.5pt}\end{center}

\textbf{Exercise 2.} A climatologist records the daily high temperature
at a monitoring station over 49 randomly selected days. The sample mean
is \(\bar{x} = 18.4°\)C. From long-term records, the population standard
deviation is \(\sigma = 5.6°\)C.

\begin{enumerate}
\def\labelenumi{(\alph{enumi})}
\tightlist
\item
  Calculate the standard error for this sample.
\item
  Construct a 95\% confidence interval for the true mean daily high
  temperature at this location.
\item
  Construct a 99\% confidence interval. (Use \(z^* = 2.576\).)
\item
  Write one sentence interpreting your 95\% CI correctly.
\end{enumerate}

\begin{Shaded}
\begin{Highlighting}[]
\NormalTok{//| echo: false}
\NormalTok{makeStepperHTML(2, [}
\NormalTok{  \{ op: "(a) Standard error", eq: "SE = \textbackslash{}\textbackslash{}dfrac\{\textbackslash{}\textbackslash{}sigma\}\{\textbackslash{}\textbackslash{}sqrt\{n\}\} = \textbackslash{}\textbackslash{}dfrac\{5.6\}\{\textbackslash{}\textbackslash{}sqrt\{49\}\} = \textbackslash{}\textbackslash{}dfrac\{5.6\}\{7\} = 0.8\^{}\{\textbackslash{}\textbackslash{}circ\}\textbackslash{}\textbackslash{}text\{C\}" \},}
\NormalTok{  \{ op: "(b) 95\% margin of error", eq: "1.96 \textbackslash{}\textbackslash{}times 0.8 = 1.568\^{}\{\textbackslash{}\textbackslash{}circ\}\textbackslash{}\textbackslash{}text\{C\}", note: "Critical value z* = 1.96 for 95\% confidence." \},}
\NormalTok{  \{ op: "(b) 95\% CI", eq: "18.4 \textbackslash{}\textbackslash{}pm 1.568 \textbackslash{}\textbackslash{}implies (16.83\^{}\{\textbackslash{}\textbackslash{}circ\}\textbackslash{}\textbackslash{}text\{C\},\textbackslash{}\textbackslash{} 19.97\^{}\{\textbackslash{}\textbackslash{}circ\}\textbackslash{}\textbackslash{}text\{C\})" \},}
\NormalTok{  \{ op: "(c) 99\% margin of error", eq: "2.576 \textbackslash{}\textbackslash{}times 0.8 = 2.061\^{}\{\textbackslash{}\textbackslash{}circ\}\textbackslash{}\textbackslash{}text\{C\}", note: "Critical value z* = 2.576 for 99\% confidence." \},}
\NormalTok{  \{ op: "(c) 99\% CI", eq: "18.4 \textbackslash{}\textbackslash{}pm 2.061 \textbackslash{}\textbackslash{}implies (16.34\^{}\{\textbackslash{}\textbackslash{}circ\}\textbackslash{}\textbackslash{}text\{C\},\textbackslash{}\textbackslash{} 20.46\^{}\{\textbackslash{}\textbackslash{}circ\}\textbackslash{}\textbackslash{}text\{C\})" \},}
\NormalTok{  \{ op: "(d) Interpretation", eq: "\textbackslash{}\textbackslash{}text\{We are 95\textbackslash{}\textbackslash{}\% confident the true mean daily high is between \} 16.83\^{}\{\textbackslash{}\textbackslash{}circ\}\textbackslash{}\textbackslash{}text\{C and \} 19.97\^{}\{\textbackslash{}\textbackslash{}circ\}\textbackslash{}\textbackslash{}text\{C.\}", note: "Confidence describes the procedure, not the probability that this particular interval contains μ." \}}
\NormalTok{])}
\end{Highlighting}
\end{Shaded}

\begin{center}\rule{0.5\linewidth}{0.5pt}\end{center}

\textbf{Exercise 3.} For each of the following scenarios, state the null
and alternative hypotheses. Identify whether the test is one-tailed or
two-tailed. If one-tailed, state the direction.

\begin{enumerate}
\def\labelenumi{(\alph{enumi})}
\item
  A consumer watchdog tests whether a cereal box labelled ``500 g net
  weight'' actually contains 500 g on average. The watchdog is concerned
  the boxes may contain \emph{less} than advertised.
\item
  A researcher believes a new fertiliser changes mean crop yield. She
  does not know whether the effect is positive or negative.
\item
  A traffic engineer claims that average vehicle speed on a stretch of
  highway is 95 km/h. A safety audit suspects drivers are exceeding this
  average.
\end{enumerate}

\begin{Shaded}
\begin{Highlighting}[]
\NormalTok{//| echo: false}
\NormalTok{makeStepperHTML(3, [}
\NormalTok{  \{ op: "(a) Null hypothesis", eq: "H\_0: \textbackslash{}\textbackslash{}mu = 500 \textbackslash{}\textbackslash{}text\{ g\}", note: "The label claim is the null state." \},}
\NormalTok{  \{ op: "(a) Alternative", eq: "H\_1: \textbackslash{}\textbackslash{}mu \textless{} 500 \textbackslash{}\textbackslash{}text\{ g\} \textbackslash{}\textbackslash{}quad \textbackslash{}\textbackslash{}text\{(one{-}tailed, lower)\}", note: "The watchdog suspects less than advertised — lower tail only." \},}
\NormalTok{  \{ op: "(b) Null hypothesis", eq: "H\_0: \textbackslash{}\textbackslash{}mu = \textbackslash{}\textbackslash{}mu\_\{\textbackslash{}\textbackslash{}text\{control\}\}", note: "No effect of the fertiliser on mean yield." \},}
\NormalTok{  \{ op: "(b) Alternative", eq: "H\_1: \textbackslash{}\textbackslash{}mu \textbackslash{}\textbackslash{}neq \textbackslash{}\textbackslash{}mu\_\{\textbackslash{}\textbackslash{}text\{control\}\} \textbackslash{}\textbackslash{}quad \textbackslash{}\textbackslash{}text\{(two{-}tailed)\}", note: "Direction unknown in advance — test both tails." \},}
\NormalTok{  \{ op: "(c) Null hypothesis", eq: "H\_0: \textbackslash{}\textbackslash{}mu = 95 \textbackslash{}\textbackslash{}text\{ km/h\}" \},}
\NormalTok{  \{ op: "(c) Alternative", eq: "H\_1: \textbackslash{}\textbackslash{}mu \textgreater{} 95 \textbackslash{}\textbackslash{}text\{ km/h\} \textbackslash{}\textbackslash{}quad \textbackslash{}\textbackslash{}text\{(one{-}tailed, upper)\}", note: "The safety audit suspects speeds are higher — upper tail only." \}}
\NormalTok{])}
\end{Highlighting}
\end{Shaded}

\begin{center}\rule{0.5\linewidth}{0.5pt}\end{center}

\textbf{Exercise 4.} An athletics coach claims that the training
programme he coaches produces marathon runners with a mean finish time
of 210 minutes. A sports scientist believes the true mean is different
(in either direction). She times 36 runners through the programme. Their
sample mean finish time is \(\bar{x} = 205.2\) minutes. The population
standard deviation is known to be \(\sigma = 18\) minutes.

\begin{enumerate}
\def\labelenumi{(\alph{enumi})}
\tightlist
\item
  State \(H_0\) and \(H_1\).
\item
  Calculate the test statistic \(z\).
\item
  Find the p-value. (This is a two-tailed test; use
  \(P(Z \leq -1.60) \approx 0.055\) as a reference value.)
\item
  At \(\alpha = 0.05\), state your decision and conclusion in plain
  language.
\end{enumerate}

\begin{Shaded}
\begin{Highlighting}[]
\NormalTok{//| echo: false}
\NormalTok{makeStepperHTML(4, [}
\NormalTok{  \{ op: "(a) Hypotheses", eq: "H\_0: \textbackslash{}\textbackslash{}mu = 210 \textbackslash{}\textbackslash{}text\{ min\}, \textbackslash{}\textbackslash{}quad H\_1: \textbackslash{}\textbackslash{}mu \textbackslash{}\textbackslash{}neq 210 \textbackslash{}\textbackslash{}text\{ min\}", note: "Two{-}tailed: the scientist suspects a difference in either direction." \},}
\NormalTok{  \{ op: "(b) Standard error", eq: "SE = \textbackslash{}\textbackslash{}dfrac\{\textbackslash{}\textbackslash{}sigma\}\{\textbackslash{}\textbackslash{}sqrt\{n\}\} = \textbackslash{}\textbackslash{}dfrac\{18\}\{\textbackslash{}\textbackslash{}sqrt\{36\}\} = \textbackslash{}\textbackslash{}dfrac\{18\}\{6\} = 3 \textbackslash{}\textbackslash{}text\{ min\}" \},}
\NormalTok{  \{ op: "(b) Test statistic", eq: "z = \textbackslash{}\textbackslash{}dfrac\{\textbackslash{}\textbackslash{}bar\{x\} {-} \textbackslash{}\textbackslash{}mu\_0\}\{SE\} = \textbackslash{}\textbackslash{}dfrac\{205.2 {-} 210\}\{3\} = \textbackslash{}\textbackslash{}dfrac\{{-}4.8\}\{3\} = {-}1.60" \},}
\NormalTok{  \{ op: "(c) p{-}value (two{-}tailed)", eq: "p = 2 \textbackslash{}\textbackslash{}times P(Z \textbackslash{}\textbackslash{}leq {-}1.60) \textbackslash{}\textbackslash{}approx 2 \textbackslash{}\textbackslash{}times 0.055 = 0.110", note: "Both tails are counted because H₁ is two{-}sided." \},}
\NormalTok{  \{ op: "(d) Decision", eq: "p = 0.110 \textgreater{} \textbackslash{}\textbackslash{}alpha = 0.05 \textbackslash{}\textbackslash{}implies \textbackslash{}\textbackslash{}text\{fail to reject \} H\_0", note: "The evidence is not strong enough at the 5\% level to conclude the mean differs from 210 minutes." \},}
\NormalTok{  \{ op: "(d) Conclusion", eq: "\textbackslash{}\textbackslash{}text\{Insufficient evidence that the true mean finish time differs from 210 min\}", note: "Failing to reject H₀ does not prove the mean is exactly 210 — only that the data are consistent with it." \}}
\NormalTok{])}
\end{Highlighting}
\end{Shaded}

\begin{center}\rule{0.5\linewidth}{0.5pt}\end{center}

\textbf{Exercise 5.} A pharmaceutical company is testing a new
painkiller. The clinical team must choose a significance level before
running the trial.

\begin{enumerate}
\def\labelenumi{(\alph{enumi})}
\tightlist
\item
  Explain what a Type I error would mean in this context.
\item
  Explain what a Type II error would mean in this context.
\item
  The team is debating between \(\alpha = 0.05\) and \(\alpha = 0.01\).
  Describe the tradeoff. Which error becomes more likely if they choose
  the stricter threshold? Which becomes less likely?
\item
  The drug is for a condition where the standard treatment has serious
  side effects. How might this context influence the choice of
  significance level? There is no single right answer --- explain your
  reasoning.
\end{enumerate}

\begin{Shaded}
\begin{Highlighting}[]
\NormalTok{//| echo: false}
\NormalTok{makeStepperHTML(5, [}
\NormalTok{  \{ op: "(a) Type I error here", eq: "\textbackslash{}\textbackslash{}text\{Reject \} H\_0 \textbackslash{}\textbackslash{}text\{ when drug has no real effect\}", note: "The drug is approved and prescribed, but it does not work — patients are exposed to side effects for no benefit." \},}
\NormalTok{  \{ op: "(b) Type II error here", eq: "\textbackslash{}\textbackslash{}text\{Fail to reject \} H\_0 \textbackslash{}\textbackslash{}text\{ when drug does work\}", note: "An effective painkiller is not approved — patients who would have benefited go without it." \},}
\NormalTok{  \{ op: "(c) Stricter threshold (α = 0.01)", eq: "P(\textbackslash{}\textbackslash{}text\{Type I\}) = 0.01 \textbackslash{}\textbackslash{}downarrow, \textbackslash{}\textbackslash{}quad P(\textbackslash{}\textbackslash{}text\{Type II\}) = \textbackslash{}\textbackslash{}beta \textbackslash{}\textbackslash{}uparrow", note: "Demanding stronger evidence reduces false approvals but increases the chance of missing a genuinely effective drug." \},}
\NormalTok{  \{ op: "(c) Looser threshold (α = 0.05)", eq: "P(\textbackslash{}\textbackslash{}text\{Type I\}) = 0.05 \textbackslash{}\textbackslash{}uparrow, \textbackslash{}\textbackslash{}quad P(\textbackslash{}\textbackslash{}text\{Type II\}) = \textbackslash{}\textbackslash{}beta \textbackslash{}\textbackslash{}downarrow", note: "More approvals overall — including more false ones." \},}
\NormalTok{  \{ op: "(d) Context argument", eq: "\textbackslash{}\textbackslash{}text\{Standard treatment has serious side effects\} \textbackslash{}\textbackslash{}Rightarrow \textbackslash{}\textbackslash{}text\{Type II error is costly\}", note: "If patients suffer under current treatment, missing an effective alternative is a serious harm. This argues for a less strict α to reduce β — though the team must weigh both risks explicitly." \}}
\NormalTok{])}
\end{Highlighting}
\end{Shaded}

\begin{center}\rule{0.5\linewidth}{0.5pt}\end{center}

\emph{Exercise 6 goes beyond what this chapter formally covers --- you
don't need the two-sample formula, but you do need to reason carefully
about what would be required.}

\textbf{Exercise 6.} A data scientist is running an A/B test for an
e-commerce site. Version A (the existing checkout flow) has a known mean
order value of \(\mu_A = \$87\) with \(\sigma_A = \$22\), estimated from
thousands of previous transactions. Version B (a redesigned flow) is
tested on a new sample.

\begin{enumerate}
\def\labelenumi{(\alph{enumi})}
\tightlist
\item
  A sample of 100 users on Version B produces \(\bar{x}_B = \$93\). If
  you were to test whether this is significantly different from the
  Version A mean of \$87, what would your null and alternative
  hypotheses be?
\item
  What additional information would you need to conduct a formal
  two-sample test comparing Version A and Version B means?
\item
  The data scientist argues: ``Version B has a higher sample mean, so we
  should switch.'' A statistician responds: ``Not so fast.'' Explain the
  statistician's concern in terms of sampling variability and the
  standard error. What would a proper inference procedure provide that
  the raw comparison of means cannot?
\end{enumerate}

\begin{Shaded}
\begin{Highlighting}[]
\NormalTok{//| echo: false}
\NormalTok{makeStepperHTML(6, [}
\NormalTok{  \{ op: "(a) Null hypothesis", eq: "H\_0: \textbackslash{}\textbackslash{}mu\_B = 87", note: "The new design produces the same mean order value as Version A." \},}
\NormalTok{  \{ op: "(a) Alternative", eq: "H\_1: \textbackslash{}\textbackslash{}mu\_B \textbackslash{}\textbackslash{}neq 87 \textbackslash{}\textbackslash{}quad \textbackslash{}\textbackslash{}text\{(two{-}tailed)\}", note: "The design could increase or decrease mean order value." \},}
\NormalTok{  \{ op: "(b) Additional information needed", eq: "\textbackslash{}\textbackslash{}sigma\_B \textbackslash{}\textbackslash{}text\{ (or \} s\_B\textbackslash{}\textbackslash{}text\{) for Version B\}", note: "You need a measure of spread for the B sample to compute SE for the two{-}sample test. You also need the B sample to be a random, representative sample." \},}
\NormalTok{  \{ op: "(b) SE for Version B alone", eq: "SE\_B = \textbackslash{}\textbackslash{}dfrac\{\textbackslash{}\textbackslash{}sigma\_B\}\{\textbackslash{}\textbackslash{}sqrt\{n\_B\}\} = \textbackslash{}\textbackslash{}dfrac\{\textbackslash{}\textbackslash{}sigma\_B\}\{\textbackslash{}\textbackslash{}sqrt\{100\}\}", note: "With σ\_B known (or estimated by s\_B), the test statistic can be computed." \},}
\NormalTok{  \{ op: "(c) The statistician\textquotesingle{}s concern", eq: "SE\_B = \textbackslash{}\textbackslash{}dfrac\{22\}\{10\} = 2.20 \textbackslash{}\textbackslash{}text\{ (using \} \textbackslash{}\textbackslash{}sigma\_A \textbackslash{}\textbackslash{}text\{ as proxy)\}", note: "Even if σ\_B ≈ σ\_A = 22, the SE of the B sample mean is $2.20. A $6 difference is about 2.7 SEs — notable, but the sampling distribution must be checked formally." \},}
\NormalTok{  \{ op: "(c) What inference adds", eq: "z = \textbackslash{}\textbackslash{}dfrac\{93 {-} 87\}\{2.20\} \textbackslash{}\textbackslash{}approx 2.73, \textbackslash{}\textbackslash{}quad p \textbackslash{}\textbackslash{}approx 0.006", note: "Formal inference gives a p{-}value. The raw comparison of $93 vs $87 tells you nothing about whether such a gap is plausible from chance; the test statistic and p{-}value do." \}}
\NormalTok{])}
\end{Highlighting}
\end{Shaded}

\chapter{Data Analysis}\label{data-analysis}

Patterns in numbers, conclusions about the world

\hfill\break

Spotify has about 100 million songs. When you finish a track and it
auto-plays the next one, you didn't choose that song --- an algorithm
did. The algorithm looked at what you've played, when you skipped, what
similar listeners chose, and built a model of your taste from the
pattern in those numbers.

That is data analysis: you have a pile of numbers that records something
about the world, and you want to find the pattern, measure it, and use
it to predict something new.

The methods in this chapter are the foundation of that entire
enterprise. The streaming recommendation is built on techniques that
grew from the same roots you're about to learn: summarising a
distribution, measuring how two variables move together, fitting a line
that lets you predict one thing from another. Every time a platform says
\emph{we think you'll like this}, the chain of reasoning traces back
here.

The chapter has four moves. First: organise and describe data on a
single variable --- centre, spread, shape. Second: describe the
relationship between two variables visually, before any formula. Third:
measure that relationship with a single number, the correlation
coefficient \(r\). Fourth: model it with a line --- regression --- so
you can make predictions. At each step the goal is not just to produce a
number, but to know what the number is actually saying.

\begin{center}\rule{0.5\linewidth}{0.5pt}\end{center}

\section{1 · Types of data}\label{types-of-data}

Before you summarise data, you need to know what kind of data you have.
The choice of summary and visualisation depends on it.

\textbf{Categorical data} names groups: favourite genre, country of
birth, blood type, whether a patient recovered. The values are labels,
not amounts --- you cannot meaningfully average them.

\textbf{Quantitative data} measures amounts: temperature, height, number
of streams, exam score. These values are numbers where arithmetic makes
sense.

Within quantitative data there is a further distinction.

\emph{Discrete} data takes separate, countable values: number of goals
scored, number of siblings, number of days absent. There is no
meaningful value between 2 and 3.

\emph{Continuous} data can take any value in a range: height, weight,
time, temperature. Between any two values there is always another one.

Statisticians sometimes also describe four \textbf{levels of
measurement}, which capture a slightly different property --- how much
structure the numbers carry.

\begin{longtable}[]{@{}
  >{\raggedright\arraybackslash}p{(\linewidth - 4\tabcolsep) * \real{0.2121}}
  >{\raggedright\arraybackslash}p{(\linewidth - 4\tabcolsep) * \real{0.2727}}
  >{\raggedright\arraybackslash}p{(\linewidth - 4\tabcolsep) * \real{0.5152}}@{}}
\toprule\noalign{}
\begin{minipage}[b]{\linewidth}\raggedright
Level
\end{minipage} & \begin{minipage}[b]{\linewidth}\raggedright
Example
\end{minipage} & \begin{minipage}[b]{\linewidth}\raggedright
What you can do
\end{minipage} \\
\midrule\noalign{}
\endhead
\bottomrule\noalign{}
\endlastfoot
Nominal & Blood type: A, B, AB, O & Count only; order is meaningless \\
Ordinal & Survey: agree / neutral / disagree & Order is meaningful; gaps
between values are not \\
Interval & Temperature in °C & Differences are meaningful; zero is
arbitrary (0°C is not ``no temperature'') \\
Ratio & Distance, mass, income & Differences and ratios both meaningful;
zero means none \\
\end{longtable}

This matters in practice. You can say that 40°C is 20° hotter than 20°C
(interval: differences work). You cannot say it is twice as hot, because
0°C is not ``no heat'' --- it is just the freezing point of water. But
you can say that 40km is twice as far as 20km (ratio: zero means nothing
there).

For this chapter the key split is \textbf{categorical vs.~quantitative},
because it determines every choice that follows: which summary
statistics to compute, which plots to draw, and whether correlation and
regression make sense at all.

\begin{center}\rule{0.5\linewidth}{0.5pt}\end{center}

\section{2 · Describing one variable}\label{describing-one-variable}

\subsection{Centre}\label{centre}

Three measures compete for ``the typical value'' of a dataset:

\textbf{Mean} --- add all values, divide by how many.

\[\bar{x} = \frac{\sum_{i=1}^{n} x_i}{n}\]

\textbf{Median} --- the middle value when the data is sorted. For an
even number of values, average the two middle ones.

\textbf{Mode} --- the most frequently occurring value.

Each one answers a slightly different question, and each has weaknesses.

Consider seven students' weekly hours of paid work:

\[4, \; 6, \; 6, \; 8, \; 9, \; 11, \; 42\]

The mean is \((4 + 6 + 6 + 8 + 9 + 11 + 42) / 7 = 86/7 \approx 12.3\).
The median is the fourth value when sorted: \(8\). The mode is \(6\).

Which is ``right''? That depends on your question. The 42 is clearly
unusual --- one student has a second job. The mean got pulled to 12.3 by
that one outlier, past six of the seven actual values. The median
(\(8\)) is not sensitive to the extreme value. The mode tells you the
most common single value.

\begin{tcolorbox}[enhanced jigsaw, opacityback=0, colback=white, breakable, colframe=quarto-callout-note-color-frame, arc=.35mm, left=2mm, toprule=.15mm, leftrule=.75mm, bottomrule=.15mm, rightrule=.15mm]
\begin{minipage}[t]{5.5mm}
\textcolor{quarto-callout-note-color}{\faInfo}
\end{minipage}%
\begin{minipage}[t]{\textwidth - 5.5mm}

\textbf{The outlier rule of thumb}

When a distribution has outliers or is heavily skewed, report the
\textbf{median} as the centre. The mean is a better summary when the
data is roughly symmetric and free of extreme values. Both together tell
you more than either alone.

\end{minipage}%
\end{tcolorbox}

\subsection{Spread}\label{spread}

A measure of centre on its own is incomplete. Knowing the average
temperature in a city tells you nothing about whether it's stable or
wildly variable. You need a measure of spread.

\textbf{Range} = maximum − minimum. Simple, but sensitive to extremes.
In the work-hours example: \(42 - 4 = 38\).

\textbf{Interquartile range (IQR)} = \(Q_3 - Q_1\). The middle 50\% of
the data, untouched by the extremes.

To find \(Q_1\) and \(Q_3\): split the sorted data at the median, then
find the median of each half.

For our seven values \(4, 6, 6, 8, 9, 11, 42\):

\begin{itemize}
\tightlist
\item
  Median (Q2) = 8
\item
  Lower half: \(4, 6, 6\) → \(Q_1 = 6\)
\item
  Upper half: \(9, 11, 42\) → \(Q_3 = 11\)
\item
  \(\text{IQR} = 11 - 6 = 5\)
\end{itemize}

The middle 50\% of students work between 6 and 11 hours per week. The 42
did not affect this at all.

\textbf{Standard deviation} measures the average distance from the mean.
You will see this written as \(s\) for a sample:

\[s = \sqrt{\frac{\sum_{i=1}^{n}(x_i - \bar{x})^2}{n-1}}\]

The formula squares the deviations (so positive and negative distances
don't cancel), takes their average (roughly --- the \(n-1\) is a
technical correction for samples), then takes the square root to get
back to the original units.

For the work-hours example, with mean \(\approx 12.3\):

\begin{longtable}[]{@{}lll@{}}
\toprule\noalign{}
\(x_i\) & \(x_i - \bar{x}\) & \((x_i - \bar{x})^2\) \\
\midrule\noalign{}
\endhead
\bottomrule\noalign{}
\endlastfoot
4 & −8.3 & 68.89 \\
6 & −6.3 & 39.69 \\
6 & −6.3 & 39.69 \\
8 & −4.3 & 18.49 \\
9 & −3.3 & 10.89 \\
11 & −1.3 & 1.69 \\
42 & +29.7 & 882.09 \\
\end{longtable}

Sum of squared deviations: \(1061.43\). Divide by \(n - 1 = 6\):
\(176.9\). Take the square root: \(s \approx 13.3\) hours.

That large standard deviation reflects the 42-hour outlier distorting
the picture. Again: when a distribution has strong outliers, IQR is the
more honest spread measure.

\subsection{Shape: the five-number summary and box
plot}\label{shape-the-five-number-summary-and-box-plot}

The \textbf{five-number summary} is min, \(Q_1\), median, \(Q_3\), max.
It describes the distribution's shape without computing a mean at all:

\[4, \quad 6, \quad 8, \quad 11, \quad 42\]

A \textbf{box plot} draws this. The box spans from \(Q_1\) to \(Q_3\). A
line inside the box marks the median. ``Whiskers'' extend to the minimum
and maximum (or to the boundary of non-outlier data, with individual
dots for outliers beyond \(1.5 \times \text{IQR}\) from the box edges).

Box plots are excellent for comparing distributions across groups ---
for example, work hours broken down by year of study --- because they
pack a lot of shape information into a small space.

\begin{tcolorbox}[enhanced jigsaw, opacityback=0, colback=white, breakable, colframe=quarto-callout-note-color-frame, arc=.35mm, left=2mm, toprule=.15mm, leftrule=.75mm, bottomrule=.15mm, rightrule=.15mm]
\begin{minipage}[t]{5.5mm}
\textcolor{quarto-callout-note-color}{\faInfo}
\end{minipage}%
\begin{minipage}[t]{\textwidth - 5.5mm}

\textbf{Outlier detection rule}

A common convention: flag a value as an outlier if it falls more than
\(1.5 \times \text{IQR}\) below \(Q_1\) or above \(Q_3\).

Lower fence: \(Q_1 - 1.5 \times \text{IQR} = 6 - 7.5 = -1.5\)

Upper fence: \(Q_3 + 1.5 \times \text{IQR} = 11 + 7.5 = 18.5\)

The value 42 is above 18.5, so it is flagged as an outlier. The other
values all fall within the fences.

\end{minipage}%
\end{tcolorbox}

\begin{center}\rule{0.5\linewidth}{0.5pt}\end{center}

\section{3 · Describing relationships}\label{describing-relationships}

When you have \textbf{two quantitative variables} measured on the same
subjects --- hours studied and exam score, temperature and ice cream
sales, daily steps and resting heart rate --- the question changes from
``what is the distribution of one variable?'' to ``how do these two
variables move together?''

The right starting point is always a \textbf{scatter plot}. Put one
variable on the horizontal axis (\(x\)) and the other on the vertical
axis (\(y\)). Each observation becomes a dot.

Before reaching for any formula, look at the plot and describe what you
see using three properties.

\textbf{Direction.} As \(x\) increases, does \(y\) tend to increase
(positive association) or decrease (negative association), or is there
no consistent tendency?

\textbf{Form.} Does the pattern look like a straight line (linear), a
curve, or something else entirely?

\textbf{Strength.} How tightly do the points cluster around the pattern?
A tight cluster suggests a strong relationship. A diffuse cloud suggests
a weak one.

Only if the form is linear does it make sense to ask ``how strong is
this linear relationship?'' --- which is what the correlation
coefficient measures. If the form is curved, the correlation coefficient
can give a misleading answer (we will see this shortly).

\begin{center}\rule{0.5\linewidth}{0.5pt}\end{center}

\section{4 · Pearson correlation}\label{pearson-correlation}

The \textbf{Pearson correlation coefficient} \(r\) measures the strength
and direction of a \emph{linear} relationship between two quantitative
variables. It is always a number between \(-1\) and \(+1\).

\[r = \frac{1}{n-1} \sum_{i=1}^{n} \left(\frac{x_i - \bar{x}}{s_x}\right)\left(\frac{y_i - \bar{y}}{s_y}\right)\]

This is the average of the products of the standardised \(x\) and \(y\)
values. Standardising means subtracting the mean and dividing by the
standard deviation, so each variable's scale disappears --- \(r\) is
unit-free.

\textbf{The intuition:} when a standardised \(x\) is positive (above
average) and the corresponding standardised \(y\) is also positive,
their product is positive. When both are below average, both are
negative, and the product is still positive. When one is above and one
is below, the product is negative. If there is a consistent positive
trend, positive products dominate and \(r\) is close to \(+1\). If there
is a consistent negative trend, negative products dominate and \(r\) is
close to \(-1\). If there is no consistent linear pattern, positive and
negative products cancel and \(r\) is near \(0\).

\textbf{A worked example.} Five students, with hours of sleep the night
before an exam (\(x\)) and exam score out of 40 (\(y\)):

\begin{longtable}[]{@{}lll@{}}
\toprule\noalign{}
Student & Sleep \(x\) & Score \(y\) \\
\midrule\noalign{}
\endhead
\bottomrule\noalign{}
\endlastfoot
A & 5 & 22 \\
B & 6 & 25 \\
C & 7 & 30 \\
D & 8 & 33 \\
E & 9 & 35 \\
\end{longtable}

\(\bar{x} = 7\), \(\bar{y} = 29\)

\(s_x = \sqrt{\frac{(5-7)^2+(6-7)^2+(7-7)^2+(8-7)^2+(9-7)^2}{4}} = \sqrt{\frac{4+1+0+1+4}{4}} = \sqrt{2.5} \approx 1.581\)

\(s_y = \sqrt{\frac{(22-29)^2+(25-29)^2+(30-29)^2+(33-29)^2+(35-29)^2}{4}} = \sqrt{\frac{49+16+1+16+36}{4}} = \sqrt{29.5} \approx 5.431\)

\begin{longtable}[]{@{}
  >{\raggedright\arraybackslash}p{(\linewidth - 6\tabcolsep) * \real{0.1250}}
  >{\raggedright\arraybackslash}p{(\linewidth - 6\tabcolsep) * \real{0.3750}}
  >{\raggedright\arraybackslash}p{(\linewidth - 6\tabcolsep) * \real{0.3750}}
  >{\raggedright\arraybackslash}p{(\linewidth - 6\tabcolsep) * \real{0.1250}}@{}}
\toprule\noalign{}
\begin{minipage}[b]{\linewidth}\raggedright
Student
\end{minipage} & \begin{minipage}[b]{\linewidth}\raggedright
\(\frac{x_i-\bar{x}}{s_x}\)
\end{minipage} & \begin{minipage}[b]{\linewidth}\raggedright
\(\frac{y_i-\bar{y}}{s_y}\)
\end{minipage} & \begin{minipage}[b]{\linewidth}\raggedright
product
\end{minipage} \\
\midrule\noalign{}
\endhead
\bottomrule\noalign{}
\endlastfoot
A & −1.265 & −1.288 & 1.629 \\
B & −0.632 & −0.737 & 0.466 \\
C & 0 & 0.184 & 0 \\
D & 0.632 & 0.737 & 0.466 \\
E & 1.265 & 1.104 & 1.397 \\
\end{longtable}

Sum of products: \(3.958\). Divide by \(n - 1 = 4\): \(r \approx 0.99\).

That is a very strong positive linear relationship. More sleep, higher
score --- and it is nearly perfectly linear across this small sample.

\textbf{Interpreting \(r\):} there is no single rule, but typical
conventions are:

\begin{longtable}[]{@{}ll@{}}
\toprule\noalign{}
\(|r|\) & Interpretation \\
\midrule\noalign{}
\endhead
\bottomrule\noalign{}
\endlastfoot
0.9 -- 1.0 & Very strong \\
0.7 -- 0.9 & Strong \\
0.5 -- 0.7 & Moderate \\
0.3 -- 0.5 & Weak \\
0.0 -- 0.3 & Very weak or none \\
\end{longtable}

\begin{tcolorbox}[enhanced jigsaw, opacityback=0, colback=white, breakable, colframe=quarto-callout-warning-color-frame, arc=.35mm, left=2mm, toprule=.15mm, leftrule=.75mm, bottomrule=.15mm, rightrule=.15mm]
\begin{minipage}[t]{5.5mm}
\textcolor{quarto-callout-warning-color}{\faExclamationTriangle}
\end{minipage}%
\begin{minipage}[t]{\textwidth - 5.5mm}

\textbf{\(r\) measures only linear association}

If the relationship is curved, \(r\) can be close to zero even when
there is a strong and obvious pattern. A dataset where \(y = x^2\)
centred at zero will give \(r \approx 0\) even though \(y\) is
completely determined by \(x\). This is why you always look at the
scatter plot first.

\end{minipage}%
\end{tcolorbox}

\begin{Shaded}
\begin{Highlighting}[]
\NormalTok{//| echo: false}
\NormalTok{\{}
\NormalTok{  const W = 500, H = 340, PAD = 40;}

\NormalTok{  // Controls}
\NormalTok{  const rSlider = Inputs.range([{-}1, 1], \{}
\NormalTok{    label: "Correlation strength",}
\NormalTok{    step: 0.1,}
\NormalTok{    value: 0.5}
\NormalTok{  \});}
\NormalTok{  const outliersToggle = Inputs.toggle(\{ label: "Add outliers", value: false \});}

\NormalTok{  const presetBar = document.createElement("div");}
\NormalTok{  presetBar.style.cssText = "display:flex; gap:0.5rem; flex{-}wrap:wrap; margin:0.5rem 0 0.25rem;";}
\NormalTok{  const presets = [}
\NormalTok{    \{ label: "r = 0.9",  v: 0.9  \},}
\NormalTok{    \{ label: "r = 0.5",  v: 0.5  \},}
\NormalTok{    \{ label: "r = 0.1",  v: 0.1  \},}
\NormalTok{    \{ label: "r = −0.7", v: {-}0.7 \}}
\NormalTok{  ];}
\NormalTok{  presets.forEach(p =\textgreater{} \{}
\NormalTok{    const btn = document.createElement("button");}
\NormalTok{    btn.textContent = p.label;}
\NormalTok{    btn.style.cssText = "padding:0.25rem 0.7rem; border:1px solid \#d1d5db; border{-}radius:4px; background:\#f9fafb; cursor:pointer; font{-}size:0.85em;";}
\NormalTok{    btn.onclick = () =\textgreater{} \{ rSlider.value = p.v; rSlider.dispatchEvent(new Event("input")); \};}
\NormalTok{    presetBar.appendChild(btn);}
\NormalTok{  \});}

\NormalTok{  // Box{-}Muller bivariate normal with correlation rho}
\NormalTok{  function bivariateNormal(n, rho, seed) \{}
\NormalTok{    // Simple LCG for reproducible{-}ish noise}
\NormalTok{    let s = seed || 42;}
\NormalTok{    function rand() \{ s = (s * 1664525 + 1013904223) \& 0xffffffff; return (s \textgreater{}\textgreater{}\textgreater{} 0) / 0x100000000; \}}
\NormalTok{    function stdNormal() \{}
\NormalTok{      const u1 = rand() || 1e{-}10, u2 = rand();}
\NormalTok{      return Math.sqrt({-}2 * Math.log(u1)) * Math.cos(2 * Math.PI * u2);}
\NormalTok{    \}}
\NormalTok{    const pts = [];}
\NormalTok{    for (let i = 0; i \textless{} n; i++) \{}
\NormalTok{      const z1 = stdNormal(), z2 = stdNormal();}
\NormalTok{      pts.push(\{ x: z1, y: rho * z1 + Math.sqrt(1 {-} rho * rho) * z2 \});}
\NormalTok{    \}}
\NormalTok{    return pts;}
\NormalTok{  \}}

\NormalTok{  function computeR(pts) \{}
\NormalTok{    const n = pts.length;}
\NormalTok{    if (n \textless{} 2) return NaN;}
\NormalTok{    const mx = pts.reduce((a, p) =\textgreater{} a + p.x, 0) / n;}
\NormalTok{    const my = pts.reduce((a, p) =\textgreater{} a + p.y, 0) / n;}
\NormalTok{    let num = 0, dx2 = 0, dy2 = 0;}
\NormalTok{    pts.forEach(p =\textgreater{} \{}
\NormalTok{      const dx = p.x {-} mx, dy = p.y {-} my;}
\NormalTok{      num += dx * dy; dx2 += dx * dx; dy2 += dy * dy;}
\NormalTok{    \});}
\NormalTok{    const denom = Math.sqrt(dx2 * dy2);}
\NormalTok{    return denom \textless{} 1e{-}12 ? 0 : num / denom;}
\NormalTok{  \}}

\NormalTok{  const svg = d3.create("svg").attr("width", W).attr("height", H)}
\NormalTok{    .style("border", "1px solid \#e5e7eb").style("border{-}radius", "6px").style("background", "\#fafafa");}

\NormalTok{  const xScale = d3.scaleLinear().domain([{-}3.5, 3.5]).range([PAD, W {-} PAD]);}
\NormalTok{  const yScale = d3.scaleLinear().domain([{-}3.5, 3.5]).range([H {-} PAD, PAD]);}

\NormalTok{  // Axes}
\NormalTok{  svg.append("g").attr("transform", \textasciigrave{}translate(0,$\{yScale(0)\})\textasciigrave{}).call(d3.axisBottom(xScale).ticks(6)).attr("color","\#9ca3af");}
\NormalTok{  svg.append("g").attr("transform", \textasciigrave{}translate($\{xScale(0)\},0)\textasciigrave{}).call(d3.axisLeft(yScale).ticks(6)).attr("color","\#9ca3af");}
\NormalTok{  svg.append("text").attr("x", W/2).attr("y", H {-} 4).attr("text{-}anchor","middle").attr("fill","\#6b7280").attr("font{-}size","12px").text("x (standardised)");}
\NormalTok{  svg.append("text").attr("transform","rotate({-}90)").attr("x", {-}H/2).attr("y", 12).attr("text{-}anchor","middle").attr("fill","\#6b7280").attr("font{-}size","12px").text("y (standardised)");}

\NormalTok{  const dotsGroup = svg.append("g");}

\NormalTok{  const rDisplay = document.createElement("div");}
\NormalTok{  rDisplay.style.cssText = "text{-}align:center; margin:0.5rem 0; font{-}size:1.6rem; font{-}weight:700; color:\#1e40af; letter{-}spacing:0.02em;";}

\NormalTok{  const noteDisplay = document.createElement("div");}
\NormalTok{  noteDisplay.style.cssText = "text{-}align:center; font{-}size:0.85em; color:\#6b7280; margin{-}bottom:0.5rem;";}

\NormalTok{  function update() \{}
\NormalTok{    const rho = +rSlider.value;}
\NormalTok{    const addOut = outliersToggle.value;}
\NormalTok{    let pts = bivariateNormal(60, rho, 42);}
\NormalTok{    if (addOut) \{}
\NormalTok{      pts = pts.concat([}
\NormalTok{        \{ x: 2.8,  y: {-}2.5 \},}
\NormalTok{        \{ x: {-}2.6, y:  2.7 \},}
\NormalTok{        \{ x:  3.0, y: {-}2.9 \}}
\NormalTok{      ]);}
\NormalTok{    \}}
\NormalTok{    const r = computeR(pts);}

\NormalTok{    dotsGroup.selectAll("circle").data(pts).join("circle")}
\NormalTok{      .attr("cx", d =\textgreater{} xScale(d.x))}
\NormalTok{      .attr("cy", d =\textgreater{} yScale(d.y))}
\NormalTok{      .attr("r", 4)}
\NormalTok{      .attr("fill", (d, i) =\textgreater{} (addOut \&\& i \textgreater{}= 60) ? "\#dc2626" : "\#3b82f6")}
\NormalTok{      .attr("opacity", (d, i) =\textgreater{} (addOut \&\& i \textgreater{}= 60) ? 0.85 : 0.55);}

\NormalTok{    rDisplay.textContent = \textasciigrave{}Computed r = $\{r.toFixed(3)\}\textasciigrave{};}
\NormalTok{    const diff = Math.abs(r {-} rho);}
\NormalTok{    noteDisplay.textContent = addOut}
\NormalTok{      ? \textasciigrave{}Target ρ = $\{rho.toFixed(1)\} · $\{diff \textgreater{} 0.05 ? \textasciigrave{}Outliers shifted r by $\{(r {-} rho).toFixed(3)\}\textasciigrave{} : "Outliers had minimal effect here"\}\textasciigrave{}}
\NormalTok{      : \textasciigrave{}Target ρ = $\{rho.toFixed(1)\}\textasciigrave{};}
\NormalTok{  \}}

\NormalTok{  rSlider.addEventListener("input", update);}
\NormalTok{  outliersToggle.addEventListener("input", update);}
\NormalTok{  update();}

\NormalTok{  const container = document.createElement("div");}
\NormalTok{  container.style.cssText = "margin:1rem 0; font{-}family:inherit;";}
\NormalTok{  const label = document.createElement("p");}
\NormalTok{  label.style.cssText = "font{-}size:0.9em; color:\#374151; margin{-}bottom:0.5rem;";}
\NormalTok{  label.textContent = "Adjust the slider to set a target correlation, then explore how the cloud of points changes. Toggle outliers to see how a few extreme points can distort the computed r.";}
\NormalTok{  container.append(label, rSlider, presetBar, outliersToggle, svg.node(), rDisplay, noteDisplay);}
\NormalTok{  return container;}
\NormalTok{\}}
\end{Highlighting}
\end{Shaded}

\begin{center}\rule{0.5\linewidth}{0.5pt}\end{center}

\section{5 · Correlation does not imply
causation}\label{correlation-does-not-imply-causation}

This section deserves real time. It is not a disclaimer. It is one of
the most practically important ideas in this entire volume.

Finding that \(r = 0.78\) between two variables does not tell you that
one causes the other. It tells you they are associated --- they tend to
move together. Association and causation are different things, and
confusing them is one of the most common errors in public reasoning
about data.

There are three mechanisms that produce correlation without causation.

\begin{center}\rule{0.5\linewidth}{0.5pt}\end{center}

\textbf{Mechanism 1: Common cause (confounding)}

Ice cream sales and drowning deaths are positively correlated. In months
when ice cream sales are high, drowning deaths are also high. Does
eating ice cream cause drowning?

No.~Both are caused by a third variable: hot weather. When it is hot,
more people buy ice cream \emph{and} more people swim. The ice cream and
the drownings are both responding to temperature. Temperature is the
\emph{confounding variable} --- a common cause that creates an apparent
relationship between two things that do not directly affect each other.

This is extremely common in health and social science research. Two
variables that seem related often share an underlying driver.

\begin{center}\rule{0.5\linewidth}{0.5pt}\end{center}

\textbf{Mechanism 2: Reverse causation}

Suppose a survey finds that people who sleep fewer hours are more likely
to be depressed. The tempting interpretation: poor sleep causes
depression.

But the causal arrow might point the other way. Depression is a common
cause of disturbed sleep. The correlation is real; the causal direction
is not established by the correlation alone.

Or both things cause each other in a loop --- poor sleep worsens mood,
worsened mood disrupts sleep. Correlation cannot untangle this.

\begin{center}\rule{0.5\linewidth}{0.5pt}\end{center}

\textbf{Mechanism 3: Coincidence}

Nicolas Cage releases more films in years when more people drown in
swimming pools. The correlation over a twelve-year period was reported
at \(r \approx 0.67\).

There is obviously no causal connection. With a sufficiently large
collection of variables measured over time, some pairs will be
correlated by chance. The internet is full of these ``spurious
correlations'' --- per capita cheese consumption and deaths by bedsheet
tangling, US spending on science and suicides by hanging. These are real
correlations in the data. They mean nothing.

\begin{center}\rule{0.5\linewidth}{0.5pt}\end{center}

\textbf{What does establish causation?}

Observational data (just measuring things in the world) can suggest
causal hypotheses. Establishing causation requires:

\begin{itemize}
\tightlist
\item
  \textbf{Randomised controlled experiments} --- randomly assign
  subjects to conditions so confounders are distributed equally across
  groups.
\item
  \textbf{Strong temporal evidence} --- the cause must precede the
  effect.
\item
  \textbf{Plausible mechanism} --- a credible explanation of \emph{how}
  one thing could cause the other.
\item
  \textbf{Replication} --- the relationship holds across different
  populations and study designs.
\end{itemize}

When you read a headline claiming ``X causes Y,'' your first question
should be: was this a randomised experiment, or an observational study?
If the latter, the correct language is ``X is associated with Y'' ---
not ``X causes Y.''

\begin{tcolorbox}[enhanced jigsaw, opacityback=0, colback=white, breakable, colframe=quarto-callout-note-color-frame, arc=.35mm, left=2mm, toprule=.15mm, leftrule=.75mm, bottomrule=.15mm, rightrule=.15mm]
\begin{minipage}[t]{5.5mm}
\textcolor{quarto-callout-note-color}{\faInfo}
\end{minipage}%
\begin{minipage}[t]{\textwidth - 5.5mm}

\textbf{The student's toolkit for evaluating causal claims}

Ask three questions:

\begin{enumerate}
\def\labelenumi{\arabic{enumi}.}
\tightlist
\item
  Could a third variable (confounder) explain both X and Y?
\item
  Could the direction be reversed --- does Y cause X instead?
\item
  Could this be coincidence in a large dataset?
\end{enumerate}

If any of these is plausible, the causal claim is not established.

\end{minipage}%
\end{tcolorbox}

\begin{center}\rule{0.5\linewidth}{0.5pt}\end{center}

\section{6 · Linear regression}\label{linear-regression}

Once you have established (from a scatter plot) that two variables have
a linear form, you can model that relationship with a line. That line
lets you do two things the correlation coefficient cannot: make
predictions for new values, and quantify how much \(y\) changes per unit
of \(x\).

\textbf{The problem.} Suppose you have a scatter of points that
approximately follow a line. Which line? You could draw many lines that
pass ``through'' the cloud. The \textbf{least-squares line} is the one
that minimises the sum of squared residuals --- where a \emph{residual}
\(e_i = y_i - \hat{y}_i\) is the vertical distance from each point to
the line.

Why squared? Squaring makes all residuals positive (so positive and
negative errors don't cancel), and it penalises large errors more
heavily than small ones. The least-squares line is the unique line that
minimises the total squared error --- it is ``as close as possible'' to
all the points simultaneously.

\textbf{The formulas.} For the least-squares line
\(\hat{y} = b_0 + b_1 x\):

\[b_1 = r \cdot \frac{s_y}{s_x}\]

\[b_0 = \bar{y} - b_1 \bar{x}\]

where \(r\) is the Pearson correlation, \(s_x\) and \(s_y\) are the
standard deviations of \(x\) and \(y\), and \(\bar{x}\), \(\bar{y}\) are
the means.

\textbf{Interpretation.}

\(b_1\) (slope): for every one-unit increase in \(x\), the predicted
value of \(y\) changes by \(b_1\) units. Always interpret this in
context --- the units matter.

\(b_0\) (intercept): the predicted value of \(y\) when \(x = 0\). This
is sometimes meaningful (it might represent a baseline), but often
extrapolating to \(x = 0\) is outside the range of the data and the
intercept is just an algebraic anchor with no practical interpretation.

\textbf{The regression line always passes through
\((\bar{x}, \bar{y})\).} This is a useful sanity check.

\begin{Shaded}
\begin{Highlighting}[]
\NormalTok{//| echo: false}
\NormalTok{\{}
\NormalTok{  const W = 500, H = 340, PAD = 45;}
\NormalTok{  const MAX\_PTS = 20;}

\NormalTok{  let points = [];}
\NormalTok{  let showResiduals = false;}

\NormalTok{  const svg = d3.create("svg")}
\NormalTok{    .attr("width", W).attr("height", H)}
\NormalTok{    .style("border", "1px solid \#e5e7eb").style("border{-}radius", "6px")}
\NormalTok{    .style("background", "\#fafafa").style("cursor", "crosshair");}

\NormalTok{  const xScale = d3.scaleLinear().domain([0, 10]).range([PAD, W {-} PAD]);}
\NormalTok{  const yScale = d3.scaleLinear().domain([0, 10]).range([H {-} PAD, PAD]);}

\NormalTok{  // Grid lines}
\NormalTok{  const grid = svg.append("g").attr("opacity", 0.25);}
\NormalTok{  xScale.ticks(5).forEach(t =\textgreater{} grid.append("line").attr("x1", xScale(t)).attr("x2", xScale(t)).attr("y1", PAD).attr("y2", H {-} PAD).attr("stroke","\#9ca3af").attr("stroke{-}dasharray","3,3"));}
\NormalTok{  yScale.ticks(5).forEach(t =\textgreater{} grid.append("line").attr("x1", PAD).attr("x2", W {-} PAD).attr("y1", yScale(t)).attr("y2", yScale(t)).attr("stroke","\#9ca3af").attr("stroke{-}dasharray","3,3"));}

\NormalTok{  svg.append("g").attr("transform", \textasciigrave{}translate(0,$\{H {-} PAD\})\textasciigrave{}).call(d3.axisBottom(xScale).ticks(5)).attr("color","\#9ca3af");}
\NormalTok{  svg.append("g").attr("transform", \textasciigrave{}translate($\{PAD\},0)\textasciigrave{}).call(d3.axisLeft(yScale).ticks(5)).attr("color","\#9ca3af");}
\NormalTok{  svg.append("text").attr("x", W/2).attr("y", H {-} 5).attr("text{-}anchor","middle").attr("fill","\#6b7280").attr("font{-}size","12px").text("x");}
\NormalTok{  svg.append("text").attr("transform","rotate({-}90)").attr("x", {-}H/2).attr("y", 14).attr("text{-}anchor","middle").attr("fill","\#6b7280").attr("font{-}size","12px").text("y");}

\NormalTok{  const hintText = svg.append("text").attr("x", W/2).attr("y", H/2)}
\NormalTok{    .attr("text{-}anchor","middle").attr("fill","\#9ca3af").attr("font{-}size","13px")}
\NormalTok{    .text("Click anywhere to add a data point");}

\NormalTok{  const residGroup = svg.append("g");}
\NormalTok{  const lineGroup = svg.append("g");}
\NormalTok{  const dotsGroup = svg.append("g");}

\NormalTok{  function leastSquares(pts) \{}
\NormalTok{    const n = pts.length;}
\NormalTok{    if (n \textless{} 2) return null;}
\NormalTok{    const mx = pts.reduce((a, p) =\textgreater{} a + p.x, 0) / n;}
\NormalTok{    const my = pts.reduce((a, p) =\textgreater{} a + p.y, 0) / n;}
\NormalTok{    let num = 0, den = 0;}
\NormalTok{    pts.forEach(p =\textgreater{} \{ const dx = p.x {-} mx; num += dx * (p.y {-} my); den += dx * dx; \});}
\NormalTok{    if (Math.abs(den) \textless{} 1e{-}12) return null;}
\NormalTok{    const b1 = num / den;}
\NormalTok{    const b0 = my {-} b1 * mx;}
\NormalTok{    const yhat = pts.map(p =\textgreater{} b0 + b1 * p.x);}
\NormalTok{    const ssTot = pts.reduce((a, p) =\textgreater{} a + (p.y {-} my) ** 2, 0);}
\NormalTok{    const ssRes = pts.reduce((a, p, i) =\textgreater{} a + (p.y {-} yhat[i]) ** 2, 0);}
\NormalTok{    const r2 = ssTot \textless{} 1e{-}12 ? 1 : 1 {-} ssRes / ssTot;}
\NormalTok{    return \{ b0, b1, r2, yhat \};}
\NormalTok{  \}}

\NormalTok{  function render() \{}
\NormalTok{    hintText.style("display", points.length === 0 ? null : "none");}

\NormalTok{    const fit = leastSquares(points);}

\NormalTok{    // Residual lines}
\NormalTok{    residGroup.selectAll("line").data(showResiduals \&\& fit ? points : [])}
\NormalTok{      .join("line")}
\NormalTok{      .attr("x1", d =\textgreater{} xScale(d.x)).attr("x2", d =\textgreater{} xScale(d.x))}
\NormalTok{      .attr("y1", d =\textgreater{} yScale(d.y))}
\NormalTok{      .attr("y2", (d, i) =\textgreater{} yScale(fit.yhat[i]))}
\NormalTok{      .attr("stroke", "\#ef4444").attr("stroke{-}width", 1.5).attr("opacity", 0.7);}

\NormalTok{    // Regression line}
\NormalTok{    if (fit) \{}
\NormalTok{      const x0 = 0, x1 = 10;}
\NormalTok{      lineGroup.selectAll("line").data([1]).join("line")}
\NormalTok{        .attr("x1", xScale(x0)).attr("x2", xScale(x1))}
\NormalTok{        .attr("y1", yScale(fit.b0 + fit.b1 * x0))}
\NormalTok{        .attr("y2", yScale(fit.b0 + fit.b1 * x1))}
\NormalTok{        .attr("stroke", "\#1d4ed8").attr("stroke{-}width", 2).attr("opacity", 0.85);}
\NormalTok{    \} else \{}
\NormalTok{      lineGroup.selectAll("line").remove();}
\NormalTok{    \}}

\NormalTok{    // Dots}
\NormalTok{    dotsGroup.selectAll("circle").data(points).join("circle")}
\NormalTok{      .attr("cx", d =\textgreater{} xScale(d.x)).attr("cy", d =\textgreater{} yScale(d.y))}
\NormalTok{      .attr("r", 5).attr("fill", "\#3b82f6").attr("stroke", "\#1d4ed8").attr("stroke{-}width", 1).attr("opacity", 0.8);}

\NormalTok{    // Stats display}
\NormalTok{    if (fit) \{}
\NormalTok{      statsDiv.innerHTML = \textasciigrave{}}
\NormalTok{        \textless{}span style="margin{-}right:1.5rem;"\textgreater{}\textless{}em\textgreater{}b\textless{}/em\textgreater{}\textless{}sub\textgreater{}0\textless{}/sub\textgreater{} = \textless{}strong\textgreater{}$\{fit.b0.toFixed(3)\}\textless{}/strong\textgreater{}\textless{}/span\textgreater{}}
\NormalTok{        \textless{}span style="margin{-}right:1.5rem;"\textgreater{}\textless{}em\textgreater{}b\textless{}/em\textgreater{}\textless{}sub\textgreater{}1\textless{}/sub\textgreater{} = \textless{}strong\textgreater{}$\{fit.b1.toFixed(3)\}\textless{}/strong\textgreater{}\textless{}/span\textgreater{}}
\NormalTok{        \textless{}span\textgreater{}\textless{}em\textgreater{}R\textless{}/em\textgreater{}\textless{}sup\textgreater{}2\textless{}/sup\textgreater{} = \textless{}strong\textgreater{}$\{fit.r2.toFixed(3)\}\textless{}/strong\textgreater{}\textless{}/span\textgreater{}}
\NormalTok{        \textless{}span style="margin{-}left:1rem; color:\#6b7280; font{-}size:0.85em;"\textgreater{}($\{points.length\} point$\{points.length===1?"":"s"\})\textless{}/span\textgreater{}\textasciigrave{};}
\NormalTok{      statsDiv.style.color = "\#1e3a8a";}
\NormalTok{    \} else \{}
\NormalTok{      statsDiv.textContent = points.length \textless{} 2 ? "Add at least 2 points to fit a line." : "";}
\NormalTok{      statsDiv.style.color = "\#6b7280";}
\NormalTok{    \}}
\NormalTok{  \}}

\NormalTok{  svg.on("click", function(event) \{}
\NormalTok{    if (points.length \textgreater{}= MAX\_PTS) return;}
\NormalTok{    const [px, py] = d3.pointer(event);}
\NormalTok{    const x = xScale.invert(px), y = yScale.invert(py);}
\NormalTok{    if (x \textless{} 0 || x \textgreater{} 10 || y \textless{} 0 || y \textgreater{} 10) return;}
\NormalTok{    points.push(\{ x, y \});}
\NormalTok{    render();}
\NormalTok{  \});}

\NormalTok{  // Controls}
\NormalTok{  const residBtn = document.createElement("button");}
\NormalTok{  residBtn.textContent = "Show residuals";}
\NormalTok{  residBtn.style.cssText = "padding:0.3rem 0.8rem; border:1px solid \#d1d5db; border{-}radius:4px; background:\#fff; cursor:pointer; font{-}size:0.875em;";}
\NormalTok{  residBtn.onclick = () =\textgreater{} \{}
\NormalTok{    showResiduals = !showResiduals;}
\NormalTok{    residBtn.textContent = showResiduals ? "Hide residuals" : "Show residuals";}
\NormalTok{    residBtn.style.background = showResiduals ? "\#fee2e2" : "\#fff";}
\NormalTok{    render();}
\NormalTok{  \};}

\NormalTok{  const clearBtn = document.createElement("button");}
\NormalTok{  clearBtn.textContent = "Clear";}
\NormalTok{  clearBtn.style.cssText = "padding:0.3rem 0.8rem; border:1px solid \#e5e7eb; border{-}radius:4px; background:\#f9fafb; cursor:pointer; font{-}size:0.875em; color:\#6b7280;";}
\NormalTok{  clearBtn.onclick = () =\textgreater{} \{ points = []; showResiduals = false; residBtn.textContent = "Show residuals"; residBtn.style.background = "\#fff"; render(); \};}

\NormalTok{  const statsDiv = document.createElement("div");}
\NormalTok{  statsDiv.style.cssText = "min{-}height:1.6em; font{-}size:0.95em; text{-}align:center; margin{-}top:0.4rem; padding:0.4rem; background:\#eff6ff; border{-}radius:4px;";}

\NormalTok{  const ctrlRow = document.createElement("div");}
\NormalTok{  ctrlRow.style.cssText = "display:flex; gap:0.5rem; margin:0.5rem 0;";}
\NormalTok{  ctrlRow.append(residBtn, clearBtn);}

\NormalTok{  const hint = document.createElement("p");}
\NormalTok{  hint.style.cssText = "font{-}size:0.85em; color:\#374151; margin{-}bottom:0.4rem;";}
\NormalTok{  hint.textContent = \textasciigrave{}Click on the plot to add points (up to $\{MAX\_PTS\}). The least{-}squares line updates in real time.\textasciigrave{};}

\NormalTok{  render();}

\NormalTok{  const container = document.createElement("div");}
\NormalTok{  container.style.cssText = "margin:1rem 0; font{-}family:inherit;";}
\NormalTok{  container.append(hint, svg.node(), ctrlRow, statsDiv);}
\NormalTok{  return container;}
\NormalTok{\}}
\end{Highlighting}
\end{Shaded}

\begin{center}\rule{0.5\linewidth}{0.5pt}\end{center}

\section{\texorpdfstring{7 · \(R^2\) and residual
plots}{7 · R\^{}2 and residual plots}}\label{r2-and-residual-plots}

\subsection{The coefficient of
determination}\label{the-coefficient-of-determination}

\(R^2\) (R-squared) measures how well the regression line fits the data.
For simple linear regression, \(R^2 = r^2\) --- the square of the
Pearson correlation.

\[R^2 = r^2\]

\(R^2\) is the \textbf{proportion of the variance in \(y\) that is
explained by the linear relationship with \(x\)}. It ranges from 0 to 1.

\(R^2 = 0.72\) means: 72\% of the variation in \(y\) is accounted for by
the linear relationship with \(x\). The remaining 28\% is due to other
factors not in the model, natural variability, or measurement error.

\(R^2 = 0.25\) means the line captures only 25\% of the variation. A lot
is left unexplained --- you might be missing important variables, or the
relationship might not be well-described by a line.

\(R^2 = 0.95\) means 95\% of the variation is explained --- a strong
fit.

\begin{tcolorbox}[enhanced jigsaw, opacityback=0, colback=white, breakable, colframe=quarto-callout-note-color-frame, arc=.35mm, left=2mm, toprule=.15mm, leftrule=.75mm, bottomrule=.15mm, rightrule=.15mm]
\begin{minipage}[t]{5.5mm}
\textcolor{quarto-callout-note-color}{\faInfo}
\end{minipage}%
\begin{minipage}[t]{\textwidth - 5.5mm}

\textbf{\(R^2\) in plain language}

Imagine the simplest possible model for \(y\): just predict \(\bar{y}\)
for everyone. That model ignores \(x\) completely. The total spread
around that flat prediction is the total variance in \(y\).

The regression line does better --- it accounts for \(x\). \(R^2\)
measures how much better: what fraction of the total variance did the
line explain away?

\end{minipage}%
\end{tcolorbox}

\subsection{Residual plots}\label{residual-plots}

Fitting a line and computing \(R^2\) is not the end of the analysis. You
need to check whether the line is actually the right model.

A \textbf{residual plot} shows the residuals \(e_i = y_i - \hat{y}_i\)
on the vertical axis against the fitted values \(\hat{y}_i\) (or against
\(x\)) on the horizontal axis. A good model produces a scatter with no
pattern --- residuals roughly centred at zero, with no curve, fan shape,
or clustering.

If the linear model is appropriate, the residuals should look like
random scatter around zero --- no pattern.

\textbf{What patterns signal problems:}

\begin{itemize}
\tightlist
\item
  A curved pattern (residuals positive at the extremes, negative in the
  middle) means the true relationship is nonlinear. A line is the wrong
  model.
\item
  A funnel shape (residuals spreading out as \(x\) increases) means the
  variability in \(y\) changes with \(x\). The model's uncertainty is
  not constant --- a standard regression assumption is violated.
\item
  A few extreme residuals are outliers in the regression sense. They may
  be data errors or genuinely unusual cases.
\end{itemize}

The residual plot is the diagnostic. A clean residual plot validates
your model; a patterned one tells you to reconsider.

\begin{center}\rule{0.5\linewidth}{0.5pt}\end{center}

\section{8 · End-to-end worked example}\label{end-to-end-worked-example}

A café tracks average daily temperature (°C) and the number of iced
coffee drinks sold. Here are eight days of data:

\begin{longtable}[]{@{}lll@{}}
\toprule\noalign{}
Day & Temp \(x\) & Drinks \(y\) \\
\midrule\noalign{}
\endhead
\bottomrule\noalign{}
\endlastfoot
1 & 18 & 34 \\
2 & 21 & 40 \\
3 & 24 & 48 \\
4 & 27 & 55 \\
5 & 29 & 61 \\
6 & 31 & 65 \\
7 & 33 & 72 \\
8 & 36 & 78 \\
\end{longtable}

\textbf{Step 1: Compute the summary statistics.}

\(n = 8\)

\(\bar{x} = (18 + 21 + 24 + 27 + 29 + 31 + 33 + 36)/8 = 219/8 = 27.375\)

\(\bar{y} = (34 + 40 + 48 + 55 + 61 + 65 + 72 + 78)/8 = 453/8 = 56.625\)

\(s_x \approx 6.117\), \(s_y \approx 15.35\)

(These are computed using the standard deviation formula from Section
2.)

\textbf{Step 2: Compute the Pearson correlation.}

Running through all 8 standardised products (same procedure as the 5-row
table in Section 4) gives:

\[r \approx 0.999\]

A near-perfect positive linear relationship. This is not surprising ---
the scatter plot would show the points nearly on a straight line.

\textbf{Step 3: Find the least-squares line.}

\[b_1 = r \cdot \frac{s_y}{s_x} = 0.999 \times \frac{15.35}{6.117} \approx 0.999 \times 2.509 \approx 2.506\]

\[b_0 = \bar{y} - b_1 \bar{x} = 56.625 - 2.506 \times 27.375 \approx 56.625 - 68.602 \approx -11.97\]

The regression equation is:

\[\hat{y} = -11.97 + 2.506x\]

\textbf{Step 4: Interpret slope and intercept.}

\emph{Slope:} For every 1°C increase in temperature, the model predicts
about 2.5 additional iced coffees sold. This makes intuitive sense.

\emph{Intercept:} The model predicts \(-11.97\) drinks at 0°C. This is
nonsensical --- no café sells negative drinks --- but we should not
expect the intercept to be meaningful here. The data ranges from 18°C to
36°C. Extrapolating to 0°C is far outside that range.

\textbf{Step 5: Predict a new value.}

On a 30°C day, the model predicts:

\[\hat{y} = -11.97 + 2.506 \times 30 = -11.97 + 75.18 = 63.21\]

Approximately 63 iced coffees.

\textbf{Step 6: Compute \(R^2\).}

\[R^2 = r^2 = (0.999)^2 \approx 0.998\]

The linear model explains about 99.8\% of the variation in drinks sold.
Temperature is almost the entire story here (within this range and this
simple model).

\textbf{Step 7: Check a residual.}

For Day 5, the actual value is \(y_5 = 61\). The fitted value is:

\[\hat{y}_5 = -11.97 + 2.506 \times 29 = -11.97 + 72.674 = 60.70\]

Residual: \(e_5 = 61 - 60.70 = 0.30\)

The model is off by less than one drink on this day. The residuals for
all eight days would be small and show no obvious pattern --- this is
exactly what a well-fitting linear model looks like.

\begin{center}\rule{0.5\linewidth}{0.5pt}\end{center}

\begin{tcolorbox}[enhanced jigsaw, opacityback=0, colback=white, breakable, colframe=quarto-callout-tip-color-frame, arc=.35mm, left=2mm, toprule=.15mm, leftrule=.75mm, bottomrule=.15mm, rightrule=.15mm]
\begin{minipage}[t]{5.5mm}
\textcolor{quarto-callout-tip-color}{\faLightbulb}
\end{minipage}%
\begin{minipage}[t]{\textwidth - 5.5mm}

\textbf{Computing form of the regression slope}

An algebraically equivalent form useful for hand calculation:

\[b_1 = \frac{\sum x_i y_i - n\bar{x}\bar{y}}{\sum x_i^2 - n\bar{x}^2}\]

This is the same formula in computing form --- useful when you have the
raw sums but not the deviations. You will use this form in Exercise 4.

\end{minipage}%
\end{tcolorbox}

\section{Exercises}\label{exercises-25}

\begin{Shaded}
\begin{Highlighting}[]
\NormalTok{//| echo: false}
\NormalTok{katex = require("katex@0.16")}
\end{Highlighting}
\end{Shaded}

\begin{Shaded}
\begin{Highlighting}[]
\NormalTok{//| echo: false}
\NormalTok{function makeStepperHTML(exerciseNum, steps) \{}
\NormalTok{  let current = 0;}
\NormalTok{  const totalSteps = steps.length;}
\NormalTok{  const container = document.createElement("div");}
\NormalTok{  container.style.cssText = "border:1px solid \#e5e7eb; border{-}radius:8px; padding:1rem 1.25rem; margin:0.75rem 0; font{-}family:inherit;";}
\NormalTok{  const stepsDiv = document.createElement("div");}
\NormalTok{  stepsDiv.style.marginBottom = "0.75rem";}

\NormalTok{  function renderSteps() \{}
\NormalTok{    stepsDiv.innerHTML = "";}
\NormalTok{    for (let i = 0; i \textless{} current; i++) \{}
\NormalTok{      const s = steps[i];}
\NormalTok{      const row = document.createElement("div");}
\NormalTok{      row.style.cssText = "display:grid; grid{-}template{-}columns:220px 1fr; gap:0.35rem 1rem; align{-}items:baseline; padding:0.35rem 0; border{-}top:1px solid \#f3f4f6;";}
\NormalTok{      const opCell = document.createElement("span");}
\NormalTok{      opCell.style.cssText = "font{-}size:0.85em; color:\#6b7280; font{-}style:italic;";}
\NormalTok{      opCell.textContent = s.op;}
\NormalTok{      const eqCell = document.createElement("span");}
\NormalTok{      eqCell.innerHTML = katex.renderToString(s.eq, \{ throwOnError: false, displayMode: false \});}
\NormalTok{      row.appendChild(opCell);}
\NormalTok{      row.appendChild(eqCell);}
\NormalTok{      if (s.note) \{}
\NormalTok{        const noteRow = document.createElement("div");}
\NormalTok{        noteRow.style.cssText = "grid{-}column:2; font{-}size:0.82em; color:\#6b7280; padding{-}bottom:0.2rem;";}
\NormalTok{        noteRow.textContent = s.note;}
\NormalTok{        row.appendChild(noteRow);}
\NormalTok{      \}}
\NormalTok{      stepsDiv.appendChild(row);}
\NormalTok{    \}}
\NormalTok{    if (current === totalSteps) \{}
\NormalTok{      const done = document.createElement("div");}
\NormalTok{      done.style.cssText = "margin{-}top:0.5rem; font{-}size:0.9em; color:\#059669; font{-}weight:500;";}
\NormalTok{      done.textContent = "✓ Solution complete";}
\NormalTok{      stepsDiv.appendChild(done);}
\NormalTok{    \}}
\NormalTok{  \}}

\NormalTok{  const controls = document.createElement("div");}
\NormalTok{  controls.style.cssText = "display:flex; gap:0.5rem; align{-}items:center;";}
\NormalTok{  const nextBtn = document.createElement("button");}
\NormalTok{  nextBtn.textContent = "Next step →";}
\NormalTok{  nextBtn.style.cssText = "padding:0.35rem 0.85rem; border:1px solid \#d1d5db; border{-}radius:4px; background:\#fff; cursor:pointer; font{-}size:0.9em;";}
\NormalTok{  const resetBtn = document.createElement("button");}
\NormalTok{  resetBtn.textContent = "Reset";}
\NormalTok{  resetBtn.style.cssText = "padding:0.35rem 0.75rem; border:1px solid \#e5e7eb; border{-}radius:4px; background:\#f9fafb; cursor:pointer; font{-}size:0.9em; color:\#6b7280;";}
\NormalTok{  const counter = document.createElement("span");}
\NormalTok{  counter.style.cssText = "font{-}size:0.82em; color:\#9ca3af; margin{-}left:0.25rem;";}

\NormalTok{  function updateButtons() \{}
\NormalTok{    nextBtn.disabled = current === totalSteps;}
\NormalTok{    nextBtn.style.opacity = current === totalSteps ? "0.4" : "1";}
\NormalTok{    counter.textContent = current === 0 ? "Click to reveal steps" : \textasciigrave{}Step $\{current\} of $\{totalSteps\}\textasciigrave{};}
\NormalTok{  \}}

\NormalTok{  nextBtn.onclick = () =\textgreater{} \{ if (current \textless{} totalSteps) \{ current++; renderSteps(); updateButtons(); \} \};}
\NormalTok{  resetBtn.onclick = () =\textgreater{} \{ current = 0; renderSteps(); updateButtons(); \};}
\NormalTok{  controls.appendChild(nextBtn);}
\NormalTok{  controls.appendChild(resetBtn);}
\NormalTok{  controls.appendChild(counter);}
\NormalTok{  container.appendChild(stepsDiv);}
\NormalTok{  container.appendChild(controls);}
\NormalTok{  renderSteps();}
\NormalTok{  updateButtons();}
\NormalTok{  return container;}
\NormalTok{\}}
\end{Highlighting}
\end{Shaded}

\textbf{Exercise 1.} The following dataset shows the number of hours
seven people spent on social media last week:

\[3, \; 7, \; 7, \; 9, \; 12, \; 14, \; 38\]

Compute the mean, median, IQR, and identify any outliers using the
\(1.5 \times \text{IQR}\) rule.

\begin{Shaded}
\begin{Highlighting}[]
\NormalTok{//| echo: false}
\NormalTok{makeStepperHTML(1, [}
\NormalTok{  \{ op: "Sort the data",}
\NormalTok{    eq: "3,\textbackslash{}\textbackslash{} 7,\textbackslash{}\textbackslash{} 7,\textbackslash{}\textbackslash{} 9,\textbackslash{}\textbackslash{} 12,\textbackslash{}\textbackslash{} 14,\textbackslash{}\textbackslash{} 38",}
\NormalTok{    note: "Already sorted. n = 7." \},}
\NormalTok{  \{ op: "Compute the mean",}
\NormalTok{    eq: "\textbackslash{}\textbackslash{}bar\{x\} = \textbackslash{}\textbackslash{}frac\{3+7+7+9+12+14+38\}\{7\} = \textbackslash{}\textbackslash{}frac\{90\}\{7\} \textbackslash{}\textbackslash{}approx 12.9" \},}
\NormalTok{  \{ op: "Find the median",}
\NormalTok{    eq: "\textbackslash{}\textbackslash{}text\{median\} = 9",}
\NormalTok{    note: "The 4th value in 7 sorted values." \},}
\NormalTok{  \{ op: "Find Q₁",}
\NormalTok{    eq: "Q\_1 = 7",}
\NormalTok{    note: "Median of the lower half: 3, 7, 7 → middle value is 7." \},}
\NormalTok{  \{ op: "Find Q₃",}
\NormalTok{    eq: "Q\_3 = 14",}
\NormalTok{    note: "Median of the upper half: 12, 14, 38 → middle value is 14." \},}
\NormalTok{  \{ op: "Compute the IQR",}
\NormalTok{    eq: "\textbackslash{}\textbackslash{}text\{IQR\} = Q\_3 {-} Q\_1 = 14 {-} 7 = 7" \},}
\NormalTok{  \{ op: "Find the fences",}
\NormalTok{    eq: "\textbackslash{}\textbackslash{}text\{Lower: \} 7 {-} 1.5(7) = {-}3.5 \textbackslash{}\textbackslash{}qquad \textbackslash{}\textbackslash{}text\{Upper: \} 14 + 1.5(7) = 24.5" \},}
\NormalTok{  \{ op: "Identify outliers",}
\NormalTok{    eq: "38 \textgreater{} 24.5 \textbackslash{}\textbackslash{}Rightarrow 38 \textbackslash{}\textbackslash{}text\{ is an outlier\}",}
\NormalTok{    note: "No values fall below −3.5. The value 38 is flagged." \},}
\NormalTok{  \{ op: "Summary",}
\NormalTok{    eq: "\textbackslash{}\textbackslash{}bar\{x\} \textbackslash{}\textbackslash{}approx 12.9,\textbackslash{}\textbackslash{} \textbackslash{}\textbackslash{}text\{median\} = 9,\textbackslash{}\textbackslash{} \textbackslash{}\textbackslash{}text\{IQR\} = 7,\textbackslash{}\textbackslash{} \textbackslash{}\textbackslash{}text\{outlier: \} 38",}
\NormalTok{    note: "The mean is pulled far above the median by the outlier. Report median as centre." \}}
\NormalTok{])}
\end{Highlighting}
\end{Shaded}

\begin{center}\rule{0.5\linewidth}{0.5pt}\end{center}

\textbf{Exercise 2.} A scatter plot of students' hours of revision
(\(x\)) and exam percentage (\(y\)) shows:

\begin{itemize}
\tightlist
\item
  Points rise from bottom-left to top-right
\item
  The points cluster fairly tightly around an imaginary straight line
\item
  There are no pronounced curves or bends in the pattern
\end{itemize}

Describe the direction, form, and strength of the association, and
estimate whether \(r\) is closer to 0.1, 0.7, or 0.95. Justify your
estimate.

\begin{Shaded}
\begin{Highlighting}[]
\NormalTok{//| echo: false}
\NormalTok{makeStepperHTML(2, [}
\NormalTok{  \{ op: "Direction",}
\NormalTok{    eq: "\textbackslash{}\textbackslash{}text\{Positive\}",}
\NormalTok{    note: "Points rise from bottom{-}left to top{-}right: as revision hours increase, exam score increases." \},}
\NormalTok{  \{ op: "Form",}
\NormalTok{    eq: "\textbackslash{}\textbackslash{}text\{Linear\}",}
\NormalTok{    note: "The pattern follows a straight line, not a curve." \},}
\NormalTok{  \{ op: "Strength",}
\NormalTok{    eq: "\textbackslash{}\textbackslash{}text\{Strong\}",}
\NormalTok{    note: "The points cluster tightly around the line — little scatter." \},}
\NormalTok{  \{ op: "Estimate r",}
\NormalTok{    eq: "r \textbackslash{}\textbackslash{}approx 0.85\textbackslash{}\textbackslash{}text\{–\}0.95",}
\NormalTok{    note: "Strong positive linear → r well above 0.7. \textquotesingle{}Fairly tight\textquotesingle{} suggests not quite perfect. 0.1 is nearly no relationship; 0.7 is moderate. 0.95 fits best." \}}
\NormalTok{])}
\end{Highlighting}
\end{Shaded}

\begin{center}\rule{0.5\linewidth}{0.5pt}\end{center}

\textbf{Exercise 3.} A study reports \(r = 0.82\) between the number of
books students own and their reading comprehension test scores. For each
of the following three statements, say whether it is correct or
incorrect and explain why.

\emph{(a)} ``Students who own more books tend to score higher on reading
tests.''

\emph{(b)} ``Owning more books causes higher reading comprehension.''

\emph{(c)} ``82\% of students who own books score above average on
reading tests.''

\begin{Shaded}
\begin{Highlighting}[]
\NormalTok{//| echo: false}
\NormalTok{makeStepperHTML(3, [}
\NormalTok{  \{ op: "Statement (a)",}
\NormalTok{    eq: "\textbackslash{}\textbackslash{}text\{Correct\}",}
\NormalTok{    note: "r = 0.82 is a strong positive association. As x (books) increases, y (score) tends to increase. This is exactly what positive correlation means." \},}
\NormalTok{  \{ op: "Statement (b)",}
\NormalTok{    eq: "\textbackslash{}\textbackslash{}text\{Incorrect — confounding likely\}",}
\NormalTok{    note: "This is a causal claim from observational data. A common cause (family income, reading culture at home, parental education) could explain both owning more books and higher scores. The correlation alone does not establish causation." \},}
\NormalTok{  \{ op: "Statement (c)",}
\NormalTok{    eq: "\textbackslash{}\textbackslash{}text\{Incorrect — misreads R\}",}
\NormalTok{    note: "r = 0.82 is not a percentage of people. It is a dimensionless measure of linear association strength. The statement confuses r with R², and R² with a frequency." \}}
\NormalTok{])}
\end{Highlighting}
\end{Shaded}

\begin{center}\rule{0.5\linewidth}{0.5pt}\end{center}

\textbf{Exercise 4.} A dataset of \(n = 6\) observations has the
following summary totals:

\[n = 6, \quad \sum x = 42, \quad \sum y = 78, \quad \sum x^2 = 322, \quad \sum xy = 591\]

Use the computing formulas below to find \(b_1\) and \(b_0\), then write
the regression equation.

\[b_1 = \frac{\sum xy - \frac{(\sum x)(\sum y)}{n}}{\sum x^2 - \frac{(\sum x)^2}{n}} \qquad b_0 = \bar{y} - b_1 \bar{x}\]

\begin{Shaded}
\begin{Highlighting}[]
\NormalTok{//| echo: false}
\NormalTok{makeStepperHTML(4, [}
\NormalTok{  \{ op: "Compute the means",}
\NormalTok{    eq: "\textbackslash{}\textbackslash{}bar\{x\} = \textbackslash{}\textbackslash{}frac\{42\}\{6\} = 7 \textbackslash{}\textbackslash{}qquad \textbackslash{}\textbackslash{}bar\{y\} = \textbackslash{}\textbackslash{}frac\{78\}\{6\} = 13" \},}
\NormalTok{  \{ op: "Compute the numerator of b₁",}
\NormalTok{    eq: "\textbackslash{}\textbackslash{}sum xy {-} \textbackslash{}\textbackslash{}frac\{(\textbackslash{}\textbackslash{}sum x)(\textbackslash{}\textbackslash{}sum y)\}\{n\} = 591 {-} \textbackslash{}\textbackslash{}frac\{42 \textbackslash{}\textbackslash{}times 78\}\{6\} = 591 {-} 546 = 45" \},}
\NormalTok{  \{ op: "Compute the denominator of b₁",}
\NormalTok{    eq: "\textbackslash{}\textbackslash{}sum x\^{}2 {-} \textbackslash{}\textbackslash{}frac\{(\textbackslash{}\textbackslash{}sum x)\^{}2\}\{n\} = 322 {-} \textbackslash{}\textbackslash{}frac\{42\^{}2\}\{6\} = 322 {-} \textbackslash{}\textbackslash{}frac\{1764\}\{6\} = 322 {-} 294 = 28" \},}
\NormalTok{  \{ op: "Compute b₁",}
\NormalTok{    eq: "b\_1 = \textbackslash{}\textbackslash{}frac\{45\}\{28\} \textbackslash{}\textbackslash{}approx 1.607" \},}
\NormalTok{  \{ op: "Compute b₀",}
\NormalTok{    eq: "b\_0 = \textbackslash{}\textbackslash{}bar\{y\} {-} b\_1 \textbackslash{}\textbackslash{}bar\{x\} = 13 {-} 1.607 \textbackslash{}\textbackslash{}times 7 = 13 {-} 11.25 = 1.75" \},}
\NormalTok{  \{ op: "Write the regression equation",}
\NormalTok{    eq: "\textbackslash{}\textbackslash{}hat\{y\} = 1.75 + 1.607x",}
\NormalTok{    note: "Check: at x = x̄ = 7, ŷ = 1.75 + 1.607(7) = 1.75 + 11.25 = 13 = ȳ. ✓" \}}
\NormalTok{])}
\end{Highlighting}
\end{Shaded}

\begin{center}\rule{0.5\linewidth}{0.5pt}\end{center}

\textbf{Exercise 5.} A researcher models the relationship between daily
temperature (\(x\), in °C) and the number of parkrun participants
(\(y\)). The regression equation is:

\[\hat{y} = 120 + 3.8x\]

\begin{enumerate}
\def\labelenumi{(\alph{enumi})}
\item
  Interpret the slope in context.
\item
  Interpret the intercept in context. Is it meaningful here?
\item
  Predict the number of participants on a 15°C day.
\item
  On one 15°C day, the actual count was 182. Compute the residual and
  interpret it.
\end{enumerate}

\begin{Shaded}
\begin{Highlighting}[]
\NormalTok{//| echo: false}
\NormalTok{makeStepperHTML(5, [}
\NormalTok{  \{ op: "(a) Interpret the slope",}
\NormalTok{    eq: "b\_1 = 3.8",}
\NormalTok{    note: "For each additional 1°C increase in temperature, the model predicts approximately 3.8 more parkrun participants." \},}
\NormalTok{  \{ op: "(b) Interpret the intercept",}
\NormalTok{    eq: "b\_0 = 120",}
\NormalTok{    note: "The model predicts 120 participants at 0°C. This may not be meaningful — parkruns in freezing temperatures are unusual, and 0°C may be outside the range of the original data. Treat as an algebraic anchor, not a real prediction." \},}
\NormalTok{  \{ op: "(c) Predict at x = 15",}
\NormalTok{    eq: "\textbackslash{}\textbackslash{}hat\{y\} = 120 + 3.8 \textbackslash{}\textbackslash{}times 15 = 120 + 57 = 177" \},}
\NormalTok{  \{ op: "(d) Compute the residual",}
\NormalTok{    eq: "e = y {-} \textbackslash{}\textbackslash{}hat\{y\} = 182 {-} 177 = 5",}
\NormalTok{    note: "The actual count was 5 higher than the model predicted. On this particular day, something (a local event? good weather beyond temperature alone?) drew extra participants not captured by temperature alone." \}}
\NormalTok{])}
\end{Highlighting}
\end{Shaded}

\begin{center}\rule{0.5\linewidth}{0.5pt}\end{center}

\textbf{Exercise 6 --- Critical thinking.} A newspaper runs the
headline:

\begin{quote}
\emph{Students who eat breakfast get better grades --- schools should
mandate breakfast to raise academic performance.}
\end{quote}

The headline is based on a large survey showing a positive correlation
between breakfast-eating habits and GPA. Using the concepts from this
chapter, evaluate the causal claim. What questions would you ask before
accepting the conclusion? What alternative explanations exist?

\begin{Shaded}
\begin{Highlighting}[]
\NormalTok{//| echo: false}
\NormalTok{makeStepperHTML(6, [}
\NormalTok{  \{ op: "Identify the type of evidence",}
\NormalTok{    eq: "\textbackslash{}\textbackslash{}text\{Observational study (survey data)\}",}
\NormalTok{    note: "Students were not randomly assigned to eat or skip breakfast. The researcher observed existing habits." \},}
\NormalTok{  \{ op: "Ask: could a confounder explain both?",}
\NormalTok{    eq: "\textbackslash{}\textbackslash{}text\{Possible confounder: household stability / socioeconomic status\}",}
\NormalTok{    note: "Students from stable, higher{-}income households are more likely to eat breakfast regularly AND to have resources supporting academic success (tutoring, quiet study space, less stress). Both the breakfast habit and the grades could be caused by the same underlying conditions." \},}
\NormalTok{  \{ op: "Ask: could causation be reversed?",}
\NormalTok{    eq: "\textbackslash{}\textbackslash{}text\{Less likely here, but consider: motivated students may also structure mornings better\}",}
\NormalTok{    note: "Academically motivated students tend to have more structured routines. The structure (including breakfast) could be a symptom of motivation, not a cause of performance." \},}
\NormalTok{  \{ op: "Ask: what would establish causation?",}
\NormalTok{    eq: "\textbackslash{}\textbackslash{}text\{A randomised controlled trial\}",}
\NormalTok{    note: "Randomly assign students to a provided breakfast programme vs control. Measure grades. If the breakfast group improves more, that is causal evidence. The survey cannot do this." \},}
\NormalTok{  \{ op: "Evaluate the policy conclusion",}
\NormalTok{    eq: "\textbackslash{}\textbackslash{}text\{The jump from correlation to mandating breakfast is premature\}",}
\NormalTok{    note: "Even if breakfast does help (plausible mechanistically — blood glucose affects concentration), the survey alone does not establish this. The school intervention might not replicate the effect if the real driver is household stability. Policy decisions deserve stronger evidence." \},}
\NormalTok{  \{ op: "Summary verdict",}
\NormalTok{    eq: "\textbackslash{}\textbackslash{}text\{Association\} \textbackslash{}\textbackslash{}neq \textbackslash{}\textbackslash{}text\{causation\}",}
\NormalTok{    note: "The correlation is real. The causal claim and the policy conclusion are not supported by this evidence alone. The correct language: \textquotesingle{}breakfast{-}eating is associated with higher GPA.\textquotesingle{} The recommended action: investigate with a controlled study before mandating a policy." \}}
\NormalTok{])}
\end{Highlighting}
\end{Shaded}

\begin{center}\rule{0.5\linewidth}{0.5pt}\end{center}

\section{Where this chapter leaves
you}\label{where-this-chapter-leaves-you}

Volume 6 is complete.

You can now describe a distribution in full --- its centre, spread, and
shape --- and choose the right summary for the kind of data you have and
the presence or absence of outliers. You can describe the relationship
between two variables from a scatter plot before any formula touches the
data. You can measure that relationship with \(r\), interpret what the
number actually means, and recognise the one thing it emphatically does
not mean: causation. And you can build the simplest predictive model,
the least-squares line, interpret its slope and intercept in context,
assess its fit with \(R^2\), and diagnose problems with a residual plot.

That is a complete analytical toolkit for real data at this level.

The road ahead is wide. Volume 7 opens with the natural extension of
everything here: \textbf{multiple regression}, where more than one
predictor variable (\(x_1, x_2, \ldots\)) is used to model \(y\)
simultaneously. The least-squares criterion is the same; the algebra
requires linear algebra. Multiple regression is also the foundation of
the simplest supervised machine learning model --- gradient descent and
neural networks are, at their roots, sophisticated generalisations of
fitting a line to data. Everything you have learned in this chapter is
load-bearing for what follows.

\begin{tcolorbox}[enhanced jigsaw, opacityback=0, colback=white, breakable, colframe=quarto-callout-tip-color-frame, arc=.35mm, left=2mm, toprule=.15mm, leftrule=.75mm, bottomrule=.15mm, rightrule=.15mm]
\begin{minipage}[t]{5.5mm}
\textcolor{quarto-callout-tip-color}{\faLightbulb}
\end{minipage}%
\begin{minipage}[t]{\textwidth - 5.5mm}

\textbf{Where this shows up}

\begin{itemize}
\tightlist
\item
  A climate scientist analysing temperature trends over decades is
  fitting a regression line to time-series data and checking the
  residuals.
\item
  A sports analyst building a player rating model is computing
  correlations between performance metrics and outcomes, then building a
  regression to predict results.
\item
  A public health researcher reporting that ``smoking is associated with
  lung cancer'' is using the careful causal language earned by decades
  of controlled studies --- because correlation alone was not enough to
  establish it.
\item
  A data scientist building a recommendation system starts with
  correlation matrices. The step from there to collaborative filtering
  and neural networks is real, but this is where the intuition lives.
\end{itemize}

The ideas are simple. Their reach is extraordinary.

\end{minipage}%
\end{tcolorbox}

\part{Volume 7 --- Engineering mathematics: ODEs}

\chapter{Volume 7 --- Engineering
mathematics}\label{volume-7-engineering-mathematics}

This volume is the core mathematics sequence that first- and second-year
engineering students usually experience as separate courses: ordinary
differential equations, linear algebra, multivariable calculus, vector
calculus, Fourier methods, PDEs, complex analysis, numerical methods,
optimisation, and probability.

The build is intentionally broad before it becomes specialised. The aim
is not to imitate a single textbook chapter order, but to give readers
the mathematical objects that upper-year engineering courses assume they
can already use: state equations, eigenmodes, fields, transforms,
discretised systems, feasible regions, and distributions.

Volume 8 extends this into the mathematics that appears inside third-
and fourth-year discipline courses: control and feedback, discrete-time
systems, continuum models, computational simulation, estimation,
reliability, and nonlinear design optimisation.

\section{Sections}\label{sections}

\begin{longtable}[]{@{}
  >{\raggedright\arraybackslash}p{(\linewidth - 2\tabcolsep) * \real{0.4500}}
  >{\raggedright\arraybackslash}p{(\linewidth - 2\tabcolsep) * \real{0.5500}}@{}}
\toprule\noalign{}
\begin{minipage}[b]{\linewidth}\raggedright
Section
\end{minipage} & \begin{minipage}[b]{\linewidth}\raggedright
Core move
\end{minipage} \\
\midrule\noalign{}
\endhead
\bottomrule\noalign{}
\endlastfoot
ODEs & Model change by relating a system to its own rates \\
Linear algebra and vector calculus & Move from single equations to
spaces, fields, and operators \\
Fourier analysis and PDEs & Decompose signals and solve distributed
physical systems \\
Complex analysis & Use the complex plane as a practical computational
tool \\
Numerical methods & Replace exact formulas with controlled
approximation \\
Optimisation and probability & Decide and infer under constraints and
uncertainty \\
\end{longtable}

\chapter{First-order ODEs}\label{first-order-odes}

When the rate of change depends on the state

\hfill\break

An RC circuit: a resistor and capacitor in series, connected to a
voltage source. The capacitor charges up. At \(t = 0\) you disconnect
the source. The voltage across the capacitor starts dropping.

The question is not ``what is the voltage after 1 second?'' You could
measure that with a multimeter. The question is: what \emph{function}
describes the voltage at every instant? How does it fall --- linearly,
exponentially, as a step? Does it depend on \(R\) and \(C\) individually
or only on their product? What happens if you double the capacitance?

To answer those questions you need an equation that the voltage function
must satisfy. That equation involves the derivative of the voltage ---
it is an ordinary differential equation. Solving it means finding the
function, not a number.

This chapter develops two techniques --- separation of variables and the
integrating factor --- that between them handle a large fraction of
first-order ODEs encountered in engineering. The chapter closes with the
existence and uniqueness theorem, which tells you something crucial: you
should check that a solution exists before spending time looking for
one.

\begin{center}\rule{0.5\linewidth}{0.5pt}\end{center}

\section{The ODE problem}\label{the-ode-problem}

A function \(y(x)\) satisfies an \textbf{ordinary differential equation}
when substituting \(y\) and its derivatives into the equation produces
an identity. The simplest possible example:

\[y' = y\]

This asks: for which functions does the derivative equal the function
itself? The answer is \(y = Ce^x\) for any constant \(C\). There is not
one solution --- there is a family of them, one curve for each value of
\(C\).

This is the fundamental shift from algebra. In algebra you solve for a
number. In an ODE you solve for a function.

\subsection{Order and linearity}\label{order-and-linearity}

The \textbf{order} of an ODE is the order of the highest derivative that
appears. This chapter covers first-order ODEs: equations involving \(y\)
and \(y'\), but not \(y''\) or higher.

An ODE is \textbf{linear} if it can be written in the form

\[a_n(x)\,y^{(n)} + a_{n-1}(x)\,y^{(n-1)} + \cdots + a_1(x)\,y' + a_0(x)\,y = f(x)\]

where all coefficients depend only on \(x\), not on \(y\). For
first-order:

\[a_1(x)\,y' + a_0(x)\,y = f(x)\]

If \(y\) or \(y'\) appear raised to a power, inside a trig function,
multiplied together, or otherwise nonlinearly, the equation is
\textbf{nonlinear}. Nonlinear first-order ODEs are harder; two important
special cases (separable equations and the logistic equation) are
tractable analytically, but most require numerical methods.

\subsection{General and particular
solutions}\label{general-and-particular-solutions}

The \textbf{general solution} of a first-order ODE contains one
arbitrary constant \(C\). Geometrically it is a one-parameter family of
curves --- every value of \(C\) gives one solution curve.

A \textbf{particular solution} is obtained by applying an
\textbf{initial condition}: a specified value \(y(x_0) = y_0\). This
pins one curve from the family.

For example, \(y' = 2x\) has general solution \(y = x^2 + C\). The
initial condition \(y(0) = 3\) gives \(C = 3\), so the particular
solution is \(y = x^2 + 3\).

\subsection{Why existence and uniqueness
matters}\label{why-existence-and-uniqueness-matters}

Before solving an ODE, it is worth knowing whether a solution exists
and, if it does, whether it is unique. If uniqueness fails, a numerical
solver that finds \emph{a} solution may not find \emph{the} solution
relevant to your physical problem. The Picard-Lindelöf theorem (stated
precisely at the end of this chapter) gives sufficient conditions for
both. In practice: check the theorem first, solve second.

\begin{center}\rule{0.5\linewidth}{0.5pt}\end{center}

\section{Separable ODEs}\label{separable-odes}

A first-order ODE is \textbf{separable} if it can be written in the form

\[\frac{dy}{dx} = g(x)\,h(y)\]

The method is to separate all \(y\)-dependence to the left and all
\(x\)-dependence to the right, then integrate both sides.

\textbf{Method.} Divide both sides by \(h(y)\) (assuming
\(h(y) \neq 0\)):

\[\frac{1}{h(y)}\,\frac{dy}{dx} = g(x)\]

Interpret \(\frac{dy}{dx}\,dx = dy\) and integrate. This is shorthand
for a substitution --- treating \(dy/dx\) as a ratio is a notational
convenience that can be made rigorous via the substitution rule for
integrals; the result is always valid for separable ODEs.

\[\int \frac{dy}{h(y)} = \int g(x)\,dx + C\]

The constant of integration appears once --- it is conventional to write
it on the right. The initial condition then determines \(C\).

\subsection{Worked example 1: exponential growth and
decay}\label{worked-example-1-exponential-growth-and-decay}

Solve \(\dfrac{dy}{dt} = ky\), with \(y(0) = y_0\).

This is separable. Divide by \(y\) (assuming \(y \neq 0\)):

\[\frac{dy}{y} = k\,dt\]

Integrate both sides:

\[\ln|y| = kt + C_1\]

Exponentiate:

\[|y| = e^{C_1} e^{kt}\]

Three sub-steps are compressed here: (1) remove the absolute value signs
by introducing a \(\pm\) --- so \(|y| = e^{C_1}e^{kt}\) becomes
\(y = \pm e^{C_1}e^{kt}\); (2) rename \(\pm e^{C_1}\) as \(A\), which is
still a nonzero constant since \(e^{C_1} > 0\) always; (3) check whether
\(y = 0\) is also a solution --- substituting gives \(0 = k \cdot 0\),
which is true, so \(y = 0\) is valid and we extend \(A\) to allow zero
as well. Collecting these observations:

\[y(t) = A e^{kt}\]

Apply the initial condition \(y(0) = y_0\):

\[y_0 = A e^0 = A\]

Particular solution:

\[\boxed{y(t) = y_0\,e^{kt}}\]

The constant \(k\) controls the behaviour. If \(k > 0\) the function
grows without bound --- radioactive build-up, compound interest,
early-stage epidemic spread. If \(k < 0\) (write \(k = -\lambda\),
\(\lambda > 0\)) the function decays to zero --- radioactive decay, drug
clearance, capacitor discharge. The quantity \(\tau = 1/|\lambda|\) is
the \textbf{time constant}: the time for the function to fall to
\(e^{-1} \approx 36.8\%\) of its initial value, or equivalently, the
time at which the tangent at \(t = 0\) would reach zero.

\subsection{Worked example 2: Newton's law of
cooling}\label{worked-example-2-newtons-law-of-cooling}

A body at temperature \(T\) loses heat to an environment at constant
temperature \(T_{\text{env}}\). Newton's law of cooling states that the
rate of heat loss is proportional to the temperature difference:

\[\frac{dT}{dt} = -k(T - T_{\text{env}}), \qquad k > 0\]

The negative sign ensures cooling: if \(T > T_{\text{env}}\), the right
side is negative, so \(T\) decreases. The equation is separable, but a
substitution makes it cleaner. Let \(u = T - T_{\text{env}}\). Then
\(\frac{du}{dt} = \frac{dT}{dt}\) (since \(T_{\text{env}}\) is
constant), and the equation becomes:

\[\frac{du}{dt} = -ku\]

This is \(y' = ky\) with \(k \to -k\). By the previous example:

\[u(t) = u_0\,e^{-kt}\]

where \(u_0 = T(0) - T_{\text{env}}\). Substituting back:

\[\boxed{T(t) = T_{\text{env}} + (T_0 - T_{\text{env}})\,e^{-kt}}\]

The temperature difference decays exponentially. As \(t \to \infty\),
\(T \to T_{\text{env}}\): the body equilibrates with its surroundings,
as expected. The constant \(k\) depends on the thermal properties of the
body and the medium --- a copper block cools faster than a ceramic mug
of the same mass (larger \(k\)).

\begin{center}\rule{0.5\linewidth}{0.5pt}\end{center}

\begin{Shaded}
\begin{Highlighting}[]
\NormalTok{//| echo: false}
\NormalTok{viewof vizEq1 = Inputs.select(["y\textquotesingle{} = ky", "y\textquotesingle{} = y(1−y)"], \{label: "Equation"\})}
\NormalTok{viewof vizK   = Inputs.range([{-}2, 2], \{step: 0.1, value: {-}0.5, label: "k (for y′ = ky)"\})}
\NormalTok{viewof vizY0  = Inputs.range([0.1, 3], \{step: 0.05, value: 1, label: "Initial condition y(0)"\})}
\end{Highlighting}
\end{Shaded}

\begin{Shaded}
\begin{Highlighting}[]
\NormalTok{//| echo: false}
\NormalTok{\{}
\NormalTok{  const d3 = await require("d3@7");}

\NormalTok{  const W = 560, H = 320;}
\NormalTok{  const margin = \{top: 20, right: 120, bottom: 40, left: 50\};}
\NormalTok{  const w = W {-} margin.left {-} margin.right;}
\NormalTok{  const h = H {-} margin.top  {-} margin.bottom;}

\NormalTok{  const tMin = 0, tMax = 5;}
\NormalTok{  const yMin = {-}0.5, yMax = 3;}

\NormalTok{  const xScale = d3.scaleLinear([tMin, tMax], [0, w]);}
\NormalTok{  const yScale = d3.scaleLinear([yMin, yMax], [h, 0]);}

\NormalTok{  const svg = d3.create("svg")}
\NormalTok{      .attr("width",  W)}
\NormalTok{      .attr("height", H)}
\NormalTok{      .style("font{-}family", "sans{-}serif")}
\NormalTok{      .style("font{-}size", "12px");}

\NormalTok{  const g = svg.append("g")}
\NormalTok{      .attr("transform", \textasciigrave{}translate($\{margin.left\},$\{margin.top\})\textasciigrave{});}

\NormalTok{  // axes}
\NormalTok{  g.append("g").attr("transform", \textasciigrave{}translate(0,$\{h\})\textasciigrave{}).call(d3.axisBottom(xScale).ticks(5));}
\NormalTok{  g.append("g").call(d3.axisLeft(yScale).ticks(6));}

\NormalTok{  // axis labels}
\NormalTok{  g.append("text")}
\NormalTok{      .attr("x", w / 2).attr("y", h + 34)}
\NormalTok{      .attr("text{-}anchor", "middle").text("t");}
\NormalTok{  g.append("text")}
\NormalTok{      .attr("transform", "rotate({-}90)")}
\NormalTok{      .attr("x", {-}h / 2).attr("y", {-}38)}
\NormalTok{      .attr("text{-}anchor", "middle").text("y");}

\NormalTok{  // slope function}
\NormalTok{  const isLogistic = vizEq1 === "y\textquotesingle{} = y(1−y)";}
\NormalTok{  const slope = (t, y) =\textgreater{} isLogistic ? y * (1 {-} y) : vizK * y;}

\NormalTok{  // direction field}
\NormalTok{  const nT = 20, nY = 18;}
\NormalTok{  const arrowLen = 0.18;}
\NormalTok{  for (let i = 0; i \textless{}= nT; i++) \{}
\NormalTok{    for (let j = 0; j \textless{}= nY; j++) \{}
\NormalTok{      const t = tMin + (i / nT) * (tMax {-} tMin);}
\NormalTok{      const y = yMin + (j / nY) * (yMax {-} yMin);}
\NormalTok{      const s = slope(t, y);}
\NormalTok{      const mag = Math.sqrt(1 + s * s);}
\NormalTok{      const dt = arrowLen / mag;}
\NormalTok{      const dy = s * arrowLen / mag;}
\NormalTok{      const x1 = xScale(t {-} dt / 2);}
\NormalTok{      const y1 = yScale(y {-} dy / 2);}
\NormalTok{      const x2 = xScale(t + dt / 2);}
\NormalTok{      const y2 = yScale(y + dy / 2);}
\NormalTok{      g.append("line")}
\NormalTok{          .attr("x1", x1).attr("y1", y1)}
\NormalTok{          .attr("x2", x2).attr("y2", y2)}
\NormalTok{          .attr("stroke", "\#ccc").attr("stroke{-}width", 1.2);}
\NormalTok{    \}}
\NormalTok{  \}}

\NormalTok{  // equilibrium line for logistic}
\NormalTok{  if (isLogistic) \{}
\NormalTok{    g.append("line")}
\NormalTok{        .attr("x1", xScale(tMin)).attr("y1", yScale(1))}
\NormalTok{        .attr("x2", xScale(tMax)).attr("y2", yScale(1))}
\NormalTok{        .attr("stroke", "\#e67e22").attr("stroke{-}width", 1.5)}
\NormalTok{        .attr("stroke{-}dasharray", "5,4");}
\NormalTok{    g.append("text")}
\NormalTok{        .attr("x", xScale(tMax) + 4).attr("y", yScale(1) + 4)}
\NormalTok{        .attr("fill", "\#e67e22").text("y=1 (equilibrium)");}
\NormalTok{  \}}

\NormalTok{  // Euler trajectory}
\NormalTok{  const dt = 0.05;}
\NormalTok{  const steps = 100;}
\NormalTok{  const pts = [];}
\NormalTok{  let ty = vizY0;}
\NormalTok{  for (let i = 0; i \textless{}= steps; i++) \{}
\NormalTok{    const t = i * dt;}
\NormalTok{    pts.push([t, ty]);}
\NormalTok{    ty = ty + dt * slope(t, ty);}
\NormalTok{    if (ty \textgreater{} yMax + 1 || ty \textless{} yMin {-} 1) break;}
\NormalTok{  \}}

\NormalTok{  const line = d3.line()}
\NormalTok{      .x(d =\textgreater{} xScale(d[0]))}
\NormalTok{      .y(d =\textgreater{} yScale(d[1]))}
\NormalTok{      .defined(d =\textgreater{} d[1] \textgreater{}= yMin {-} 0.5 \&\& d[1] \textless{}= yMax + 0.5);}

\NormalTok{  g.append("path")}
\NormalTok{      .datum(pts)}
\NormalTok{      .attr("fill", "none")}
\NormalTok{      .attr("stroke", "\#2980b9")}
\NormalTok{      .attr("stroke{-}width", 2)}
\NormalTok{      .attr("d", line);}

\NormalTok{  // formula label top{-}right}
\NormalTok{  const formulaLabel = isLogistic ? "y′ = y(1−y)" : \textasciigrave{}y′ = $\{vizK.toFixed(1)\}y\textasciigrave{};}
\NormalTok{  g.append("text")}
\NormalTok{      .attr("x", w {-} 4).attr("y", 12)}
\NormalTok{      .attr("text{-}anchor", "end")}
\NormalTok{      .attr("fill", "\#333")}
\NormalTok{      .style("font{-}size", "13px")}
\NormalTok{      .text(formulaLabel);}

\NormalTok{  return svg.node();}
\NormalTok{\}}
\end{Highlighting}
\end{Shaded}

\begin{center}\rule{0.5\linewidth}{0.5pt}\end{center}

\section{Linear first-order ODEs}\label{linear-first-order-odes}

Not every first-order ODE is separable. Consider:

\[y' + 2xy = e^{-x^2}\]

The right side prevents simple separation. But the equation is linear
--- it has the form

\[y' + P(x)\,y = Q(x)\]

This is the \textbf{standard form} of a first-order linear ODE.

\subsection{The integrating factor}\label{the-integrating-factor}

The key idea: multiply both sides by a function \(\mu(x)\), chosen so
that the left side becomes the derivative of a product. If

\[\mu(x)\,y' + \mu(x)\,P(x)\,y = \frac{d}{dx}[\mu(x)\,y]\]

then by the product rule the right side expands as
\(\mu\,y' + \mu'\,y\). Matching:

\[\mu(x)\,P(x) = \mu'(x)\]

This is a separable ODE for \(\mu\):

\[\frac{d\mu}{\mu} = P(x)\,dx \implies \ln\mu = \int P(x)\,dx \implies \mu(x) = e^{\int P(x)\,dx}\]

No constant of integration is needed --- any \(\mu\) that works is
sufficient, and the simplest choice sets the constant to zero.

\textbf{Why this is the key step.} After multiplying by \(\mu\), the ODE
becomes:

\[\frac{d}{dx}[\mu(x)\,y] = \mu(x)\,Q(x)\]

The left side is a perfect derivative. Integrate both sides:

\[\mu(x)\,y = \int \mu(x)\,Q(x)\,dx + C\]

\[y = \frac{1}{\mu(x)}\left[\int \mu(x)\,Q(x)\,dx + C\right]\]

The integrating factor converts the ODE into a straightforward
integration problem.

\textbf{Summary of the method:}

\begin{enumerate}
\def\labelenumi{\arabic{enumi}.}
\tightlist
\item
  Write in standard form: \(y' + P(x)\,y = Q(x)\)
\item
  Compute \(\mu(x) = e^{\int P(x)\,dx}\)
\item
  Multiply both sides by \(\mu\): left side becomes
  \(\frac{d}{dx}[\mu y]\)
\item
  Integrate both sides
\item
  Solve for \(y\); apply initial condition to find \(C\)
\end{enumerate}

\subsection{Worked example: RC circuit}\label{worked-example-rc-circuit}

A series RC circuit connected to a constant voltage source \(V_s\)
satisfies Kirchhoff's voltage law:

\[V_R + V_C = V_s\]

The resistor voltage is \(V_R = Ri = RC\,\frac{dV_C}{dt}\) (using
\(i = C\,\frac{dV_C}{dt}\)). Writing \(V = V_C\):

\[RC\,\frac{dV}{dt} + V = V_s\]

Divide by \(RC\) to reach standard form:

\[\frac{dV}{dt} + \frac{1}{RC}\,V = \frac{V_s}{RC}\]

Here \(P(t) = \frac{1}{RC}\) and \(Q(t) = \frac{V_s}{RC}\).

\textbf{Step 2.} Integrating factor:

\[\mu(t) = e^{\int \frac{1}{RC}\,dt} = e^{t/RC}\]

\textbf{Step 3.} Multiply through:

\[\frac{d}{dt}\left[e^{t/RC}\,V\right] = \frac{V_s}{RC}\,e^{t/RC}\]

\textbf{Step 4.} Integrate both sides:

\[e^{t/RC}\,V = V_s\,e^{t/RC} + C\]

\textbf{Step 5.} Divide by \(e^{t/RC}\):

\[V(t) = V_s + C\,e^{-t/RC}\]

Apply \(V(0) = V_0\):

\[V_0 = V_s + C \implies C = V_0 - V_s\]

Particular solution:

\[\boxed{V(t) = V_s + (V_0 - V_s)\,e^{-t/RC}}\]

Reading the solution physically: the voltage starts at \(V_0\) and
approaches \(V_s\) exponentially with time constant \(\tau = RC\). If
\(V_0 = 0\) (capacitor initially uncharged), the voltage rises from zero
to \(V_s\). If \(V_0 > V_s\) (capacitor initially overcharged), it
decays down. The product \(\tau = RC\) --- measured in seconds when
\(R\) is in ohms and \(C\) in farads --- governs how fast the transient
dies. After \(5\tau\) the transient is below 1\% of its initial
magnitude; engineers call this ``essentially complete.''

\begin{center}\rule{0.5\linewidth}{0.5pt}\end{center}

\begin{Shaded}
\begin{Highlighting}[]
\NormalTok{//| echo: false}
\NormalTok{viewof vizR  = Inputs.range([1, 100],   \{step: 1,    value: 10,   label: "R (Ω)"\})}
\NormalTok{viewof vizC  = Inputs.range([100, 10000], \{step: 100, value: 1000, label: "C (μF)"\})}
\NormalTok{viewof vizV0 = Inputs.range([0, 20],    \{step: 0.5,  value: 10,   label: "V₀ (initial voltage, V)"\})}
\NormalTok{viewof vizVs = Inputs.range([0, 20],    \{step: 0.5,  value: 0,    label: "Vₛ (source voltage, V)"\})}
\end{Highlighting}
\end{Shaded}

\begin{Shaded}
\begin{Highlighting}[]
\NormalTok{//| echo: false}
\NormalTok{\{}
\NormalTok{  const d3 = await require("d3@7");}

\NormalTok{  const W = 560, H = 300;}
\NormalTok{  const margin = \{top: 20, right: 140, bottom: 40, left: 55\};}
\NormalTok{  const w = W {-} margin.left {-} margin.right;}
\NormalTok{  const h = H {-} margin.top  {-} margin.bottom;}

\NormalTok{  // τ in seconds  (C is in μF → convert to F)}
\NormalTok{  const tau = vizR * vizC * 1e{-}6;}
\NormalTok{  const tMax = 5 * tau;}

\NormalTok{  const vMin = Math.min(vizVs, vizV0) {-} 0.5;}
\NormalTok{  const vMax = Math.max(vizVs, vizV0) + 0.5;}

\NormalTok{  const xScale = d3.scaleLinear([0, tMax], [0, w]);}
\NormalTok{  const yScale = d3.scaleLinear([vMin, vMax], [h, 0]);}

\NormalTok{  const svg = d3.create("svg")}
\NormalTok{      .attr("width",  W)}
\NormalTok{      .attr("height", H)}
\NormalTok{      .style("font{-}family", "sans{-}serif")}
\NormalTok{      .style("font{-}size", "12px");}

\NormalTok{  const g = svg.append("g")}
\NormalTok{      .attr("transform", \textasciigrave{}translate($\{margin.left\},$\{margin.top\})\textasciigrave{});}

\NormalTok{  // axes}
\NormalTok{  g.append("g").attr("transform", \textasciigrave{}translate(0,$\{h\})\textasciigrave{})}
\NormalTok{      .call(d3.axisBottom(xScale).ticks(5).tickFormat(v =\textgreater{} d3.format(".3g")(v)));}
\NormalTok{  g.append("g").call(d3.axisLeft(yScale).ticks(6));}

\NormalTok{  // axis labels}
\NormalTok{  g.append("text")}
\NormalTok{      .attr("x", w / 2).attr("y", h + 34)}
\NormalTok{      .attr("text{-}anchor", "middle").text("t (s)");}
\NormalTok{  g.append("text")}
\NormalTok{      .attr("transform", "rotate({-}90)")}
\NormalTok{      .attr("x", {-}h / 2).attr("y", {-}44)}
\NormalTok{      .attr("text{-}anchor", "middle").text("V(t) (V)");}

\NormalTok{  // V∞ horizontal dashed line}
\NormalTok{  g.append("line")}
\NormalTok{      .attr("x1", 0).attr("y1", yScale(vizVs))}
\NormalTok{      .attr("x2", w).attr("y2", yScale(vizVs))}
\NormalTok{      .attr("stroke", "\#27ae60").attr("stroke{-}width", 1.4)}
\NormalTok{      .attr("stroke{-}dasharray", "5,4");}
\NormalTok{  g.append("text")}
\NormalTok{      .attr("x", w + 4).attr("y", yScale(vizVs) + 4)}
\NormalTok{      .attr("fill", "\#27ae60")}
\NormalTok{      .style("font{-}size", "11px")}
\NormalTok{      .text(\textasciigrave{}V∞ = Vₛ = $\{vizVs.toFixed(1)\} V\textasciigrave{});}

\NormalTok{  // τ vertical dashed line}
\NormalTok{  g.append("line")}
\NormalTok{      .attr("x1", xScale(tau)).attr("y1", 0)}
\NormalTok{      .attr("x2", xScale(tau)).attr("y2", h)}
\NormalTok{      .attr("stroke", "\#8e44ad").attr("stroke{-}width", 1.4)}
\NormalTok{      .attr("stroke{-}dasharray", "5,4");}
\NormalTok{  g.append("text")}
\NormalTok{      .attr("x", xScale(tau) + 4).attr("y", 14)}
\NormalTok{      .attr("fill", "\#8e44ad")}
\NormalTok{      .style("font{-}size", "11px")}
\NormalTok{      .text("τ = RC");}

\NormalTok{  // V(τ) annotation dot + label}
\NormalTok{  const vAtTau = vizVs + (vizV0 {-} vizVs) * Math.exp({-}1);}
\NormalTok{  g.append("circle")}
\NormalTok{      .attr("cx", xScale(tau)).attr("cy", yScale(vAtTau))}
\NormalTok{      .attr("r", 4).attr("fill", "\#8e44ad");}
\NormalTok{  g.append("text")}
\NormalTok{      .attr("x", xScale(tau) + 6).attr("y", yScale(vAtTau) {-} 6)}
\NormalTok{      .attr("fill", "\#8e44ad")}
\NormalTok{      .style("font{-}size", "11px")}
\NormalTok{      .text(\textasciigrave{}V(τ) ≈ $\{vAtTau.toFixed(2)\} V\textasciigrave{});}

\NormalTok{  // V(t) curve}
\NormalTok{  const nPts = 300;}
\NormalTok{  const pts = d3.range(nPts + 1).map(i =\textgreater{} \{}
\NormalTok{    const t = (i / nPts) * tMax;}
\NormalTok{    const v = vizVs + (vizV0 {-} vizVs) * Math.exp({-}t / tau);}
\NormalTok{    return [t, v];}
\NormalTok{  \});}

\NormalTok{  const line = d3.line()}
\NormalTok{      .x(d =\textgreater{} xScale(d[0]))}
\NormalTok{      .y(d =\textgreater{} yScale(d[1]));}

\NormalTok{  g.append("path")}
\NormalTok{      .datum(pts)}
\NormalTok{      .attr("fill", "none")}
\NormalTok{      .attr("stroke", "\#2980b9")}
\NormalTok{      .attr("stroke{-}width", 2.2)}
\NormalTok{      .attr("d", line);}

\NormalTok{  // τ value label top{-}right}
\NormalTok{  const tauMs = (tau * 1000).toFixed(1);}
\NormalTok{  g.append("text")}
\NormalTok{      .attr("x", w {-} 4).attr("y", 14)}
\NormalTok{      .attr("text{-}anchor", "end")}
\NormalTok{      .attr("fill", "\#333")}
\NormalTok{      .style("font{-}size", "12px")}
\NormalTok{      .text(\textasciigrave{}τ = $\{tauMs\} ms\textasciigrave{});}

\NormalTok{  return svg.node();}
\NormalTok{\}}
\end{Highlighting}
\end{Shaded}

\begin{center}\rule{0.5\linewidth}{0.5pt}\end{center}

\section{Applications}\label{applications-4}

\subsection{Mixing problem}\label{mixing-problem}

A 100 L tank initially contains fresh water. Brine at concentration 2
g/L flows in at 4 L/min. The well-mixed solution flows out at 4 L/min,
so the volume stays constant at 100 L.

Let \(S(t)\) be the salt content in grams at time \(t\). The rate of
change is:

\[\frac{dS}{dt} = \underbrace{2 \times 4}_{\text{in}} - \underbrace{\frac{S}{100} \times 4}_{\text{out}} = 8 - \frac{S}{25}\]

Rearranging to standard form:

\[\frac{dS}{dt} + \frac{1}{25}\,S = 8\]

Integrating factor: \(\mu = e^{t/25}\).

\[\frac{d}{dt}\left[e^{t/25}\,S\right] = 8\,e^{t/25}\]

\[e^{t/25}\,S = 200\,e^{t/25} + C\]

\[S(t) = 200 + C\,e^{-t/25}\]

Initial condition \(S(0) = 0\) gives \(C = -200\):

\[\boxed{S(t) = 200\!\left(1 - e^{-t/25}\right)}\]

As \(t \to \infty\), \(S \to 200\) g --- this is the equilibrium: 2 g/L
in 100 L. The approach is exponential with time constant \(\tau = 25\)
min. The tank is at 99\% of equilibrium after \(5\tau = 125\) min. The
form of this solution is identical to the RC circuit: a system driven
toward equilibrium by a constant input, with exponential relaxation.

\subsection{Drug pharmacokinetics: one-compartment
model}\label{drug-pharmacokinetics-one-compartment-model}

A drug is injected as a bolus (instantaneous dose) of \(D_0\) mg into
the bloodstream. The drug is cleared by the kidneys at a rate
proportional to the blood concentration \(C(t)\) (mg/L), with
elimination rate constant \(k_e\) (min⁻¹). The volume of distribution is
\(V_d\) (L).

The ODE for the amount \(A(t) = C(t)\,V_d\) in the body is:

\[\frac{dA}{dt} = -k_e\,A, \qquad A(0) = D_0\]

This is \(y' = ky\) with \(k = -k_e\). Solution:

\[A(t) = D_0\,e^{-k_e t}\]

The blood concentration is \(C(t) = A(t)/V_d = (D_0/V_d)\,e^{-k_e t}\).

The \textbf{half-life} of the drug is the time for the concentration to
halve:

\[C(t_{1/2}) = \frac{C_0}{2} \implies e^{-k_e t_{1/2}} = \frac{1}{2} \implies t_{1/2} = \frac{\ln 2}{k_e}\]

This is the number quoted on drug package inserts. A drug with
\(t_{1/2} = 4\) hours requires dosing every 4--8 hours to maintain
therapeutic levels. The simplicity of this model makes it foundational
in pharmacology --- more realistic models add compartments (tissue, fat,
CSF) but the same ODE structure applies to each.

\begin{center}\rule{0.5\linewidth}{0.5pt}\end{center}

\section{Existence and uniqueness}\label{existence-and-uniqueness}

Before investing effort in solving an ODE, it is worth asking whether a
solution exists and whether it is unique. The \textbf{Picard-Lindelöf
theorem} gives a sufficient condition.

\textbf{Theorem (Picard-Lindelöf).} Consider the initial value problem

\[y' = f(x, y), \qquad y(x_0) = y_0\]

If \(f\) and \(\partial f / \partial y\) are both continuous on some
open rectangle \(R = \{(x,y) : |x - x_0| < a,\; |y - y_0| < b\}\)
containing \((x_0, y_0)\), then there exists an interval
\(|x - x_0| < h\) (for some \(h > 0\)) on which the IVP has a
\textbf{unique} solution.

The theorem does not tell you the size of \(h\) --- the solution may
only exist locally, not for all \(x\). But it does guarantee that on
some interval, exactly one solution passes through the given initial
point.

\textbf{What failure looks like.} Consider \(y' = y^{2/3}\) with
\(y(0) = 0\).

Here \(f(x,y) = y^{2/3}\), which is continuous everywhere. But
\(\partial f / \partial y = \frac{2}{3} y^{-1/3}\), which is \emph{not}
continuous at \(y = 0\) --- the theorem's hypothesis fails. And indeed,
uniqueness fails: both

\[y_1(x) = 0 \qquad \text{and} \qquad y_2(x) = \left(\frac{x}{3}\right)^3\]

satisfy \(y' = y^{2/3}\) and \(y(0) = 0\). There is not a unique
solution through the origin.

\textbf{Engineering consequence.} When the theorem applies, you know
that any reasonable numerical method converges to the same solution ---
there is no ambiguity about which solution you are approximating. When
it fails, a numerical solver may track different solution branches
depending on the step size or starting perturbation, and the results can
be misleading. Check the theorem first.

\begin{center}\rule{0.5\linewidth}{0.5pt}\end{center}

\section{Exercises}\label{exercises-26}

These exercises are puzzles with physical content. The interesting part
is identifying the method (separable or integrating factor) and setting
up the solution correctly.

\begin{Shaded}
\begin{Highlighting}[]
\NormalTok{//| echo: false}
\NormalTok{katex = require("katex@0.16")}
\end{Highlighting}
\end{Shaded}

\begin{Shaded}
\begin{Highlighting}[]
\NormalTok{//| echo: false}
\NormalTok{function makeStepperHTML(steps) \{}
\NormalTok{  let current = 0;}
\NormalTok{  const totalSteps = steps.length;}
\NormalTok{  const container = document.createElement("div");}
\NormalTok{  container.style.cssText = "border:1px solid \#e5e7eb; border{-}radius:8px; padding:1rem 1.25rem; margin:0.75rem 0; font{-}family:inherit;";}
\NormalTok{  const stepsDiv = document.createElement("div");}
\NormalTok{  stepsDiv.style.marginBottom = "0.75rem";}

\NormalTok{  function renderSteps() \{}
\NormalTok{    stepsDiv.innerHTML = "";}
\NormalTok{    for (let i = 0; i \textless{} current; i++) \{}
\NormalTok{      const s = steps[i];}
\NormalTok{      const row = document.createElement("div");}
\NormalTok{      row.style.cssText = "display:grid; grid{-}template{-}columns:220px 1fr; gap:0.35rem 1rem; align{-}items:baseline; padding:0.35rem 0; border{-}top:1px solid \#f3f4f6;";}
\NormalTok{      const opCell = document.createElement("span");}
\NormalTok{      opCell.style.cssText = "font{-}size:0.85em; color:\#6b7280; font{-}style:italic;";}
\NormalTok{      opCell.textContent = s.op;}
\NormalTok{      const eqCell = document.createElement("span");}
\NormalTok{      eqCell.innerHTML = katex.renderToString(s.eq, \{ throwOnError: false, displayMode: false \});}
\NormalTok{      row.appendChild(opCell);}
\NormalTok{      row.appendChild(eqCell);}
\NormalTok{      if (s.note) \{}
\NormalTok{        const noteRow = document.createElement("div");}
\NormalTok{        noteRow.style.cssText = "grid{-}column:2; font{-}size:0.82em; color:\#6b7280; padding{-}bottom:0.2rem;";}
\NormalTok{        noteRow.textContent = s.note;}
\NormalTok{        row.appendChild(noteRow);}
\NormalTok{      \}}
\NormalTok{      stepsDiv.appendChild(row);}
\NormalTok{    \}}
\NormalTok{    if (current === totalSteps) \{}
\NormalTok{      const done = document.createElement("div");}
\NormalTok{      done.style.cssText = "margin{-}top:0.5rem; font{-}size:0.9em; color:\#059669; font{-}weight:500;";}
\NormalTok{      done.textContent = "Solution complete.";}
\NormalTok{      stepsDiv.appendChild(done);}
\NormalTok{    \}}
\NormalTok{  \}}

\NormalTok{  const controls = document.createElement("div");}
\NormalTok{  controls.style.cssText = "display:flex; gap:0.5rem; align{-}items:center;";}
\NormalTok{  const nextBtn = document.createElement("button");}
\NormalTok{  nextBtn.textContent = "Next step →";}
\NormalTok{  nextBtn.style.cssText = "padding:0.35rem 0.85rem; border:1px solid \#d1d5db; border{-}radius:4px; background:\#fff; cursor:pointer; font{-}size:0.9em;";}
\NormalTok{  const resetBtn = document.createElement("button");}
\NormalTok{  resetBtn.textContent = "Reset";}
\NormalTok{  resetBtn.style.cssText = "padding:0.35rem 0.75rem; border:1px solid \#e5e7eb; border{-}radius:4px; background:\#f9fafb; cursor:pointer; font{-}size:0.9em; color:\#6b7280;";}
\NormalTok{  const counter = document.createElement("span");}
\NormalTok{  counter.style.cssText = "font{-}size:0.82em; color:\#9ca3af; margin{-}left:0.25rem;";}

\NormalTok{  function updateButtons() \{}
\NormalTok{    nextBtn.disabled = current === totalSteps;}
\NormalTok{    nextBtn.style.opacity = current === totalSteps ? "0.4" : "1";}
\NormalTok{    counter.textContent = current === 0 ? "Click to reveal steps" : \textasciigrave{}Step $\{current\} of $\{totalSteps\}\textasciigrave{};}
\NormalTok{  \}}

\NormalTok{  nextBtn.onclick = () =\textgreater{} \{ if (current \textless{} totalSteps) \{ current++; renderSteps(); updateButtons(); \} \};}
\NormalTok{  resetBtn.onclick = () =\textgreater{} \{ current = 0; renderSteps(); updateButtons(); \};}
\NormalTok{  controls.appendChild(nextBtn);}
\NormalTok{  controls.appendChild(resetBtn);}
\NormalTok{  controls.appendChild(counter);}
\NormalTok{  container.appendChild(stepsDiv);}
\NormalTok{  container.appendChild(controls);}
\NormalTok{  renderSteps();}
\NormalTok{  updateButtons();}
\NormalTok{  return container;}
\NormalTok{\}}
\end{Highlighting}
\end{Shaded}

\begin{center}\rule{0.5\linewidth}{0.5pt}\end{center}

\textbf{Exercise 1.} Solve \(\dfrac{dy}{dx} = \dfrac{x}{y}\). Find the
general solution and the particular solution satisfying \(y(0) = 2\).

\begin{Shaded}
\begin{Highlighting}[]
\NormalTok{//| echo: false}
\NormalTok{makeStepperHTML([}
\NormalTok{  \{}
\NormalTok{    op: "Separate variables",}
\NormalTok{    eq: "y\textbackslash{}\textbackslash{},dy = x\textbackslash{}\textbackslash{},dx",}
\NormalTok{    note: "Multiply both sides by y and by dx."}
\NormalTok{  \},}
\NormalTok{  \{}
\NormalTok{    op: "Integrate both sides",}
\NormalTok{    eq: "\textbackslash{}\textbackslash{}int y\textbackslash{}\textbackslash{},dy = \textbackslash{}\textbackslash{}int x\textbackslash{}\textbackslash{},dx \textbackslash{}\textbackslash{}implies \textbackslash{}\textbackslash{}frac\{y\^{}2\}\{2\} = \textbackslash{}\textbackslash{}frac\{x\^{}2\}\{2\} + C\_1"}
\NormalTok{  \},}
\NormalTok{  \{}
\NormalTok{    op: "Multiply through by 2, absorb constant",}
\NormalTok{    eq: "y\^{}2 = x\^{}2 + C \textbackslash{}\textbackslash{}quad \textbackslash{}\textbackslash{}text\{(general solution, family of hyperbolas)\}"}
\NormalTok{  \},}
\NormalTok{  \{}
\NormalTok{    op: "Apply initial condition y(0) = 2",}
\NormalTok{    eq: "4 = 0 + C \textbackslash{}\textbackslash{}implies C = 4"}
\NormalTok{  \},}
\NormalTok{  \{}
\NormalTok{    op: "Particular solution",}
\NormalTok{    eq: "y\^{}2 = x\^{}2 + 4 \textbackslash{}\textbackslash{}implies y = \textbackslash{}\textbackslash{}sqrt\{x\^{}2 + 4\}",}
\NormalTok{    note: "Take the positive root because y(0) = 2 \textgreater{} 0."}
\NormalTok{  \}}
\NormalTok{])}
\end{Highlighting}
\end{Shaded}

\begin{center}\rule{0.5\linewidth}{0.5pt}\end{center}

\textbf{Exercise 2.} Solve \(\dfrac{dy}{dt} = -0.1y\) with
\(y(0) = 500\). Find \(y(t)\), then find the time at which \(y = 100\).

\begin{Shaded}
\begin{Highlighting}[]
\NormalTok{//| echo: false}
\NormalTok{makeStepperHTML([}
\NormalTok{  \{}
\NormalTok{    op: "Recognise form y\textquotesingle{} = ky",}
\NormalTok{    eq: "k = {-}0.1, \textbackslash{}\textbackslash{}quad y\_0 = 500"}
\NormalTok{  \},}
\NormalTok{  \{}
\NormalTok{    op: "Write general solution",}
\NormalTok{    eq: "y(t) = 500\textbackslash{}\textbackslash{},e\^{}\{{-}0.1t\}"}
\NormalTok{  \},}
\NormalTok{  \{}
\NormalTok{    op: "Set y = 100, solve for t",}
\NormalTok{    eq: "500\textbackslash{}\textbackslash{},e\^{}\{{-}0.1t\} = 100 \textbackslash{}\textbackslash{}implies e\^{}\{{-}0.1t\} = 0.2"}
\NormalTok{  \},}
\NormalTok{  \{}
\NormalTok{    op: "Take natural log",}
\NormalTok{    eq: "{-}0.1t = \textbackslash{}\textbackslash{}ln(0.2) = {-}\textbackslash{}\textbackslash{}ln 5"}
\NormalTok{  \},}
\NormalTok{  \{}
\NormalTok{    op: "Solve for t",}
\NormalTok{    eq: "t = \textbackslash{}\textbackslash{}frac\{\textbackslash{}\textbackslash{}ln 5\}\{0.1\} = 10\textbackslash{}\textbackslash{}ln 5 \textbackslash{}\textbackslash{}approx 16.09",}
\NormalTok{    note: "ln 5 ≈ 1.609."}
\NormalTok{  \}}
\NormalTok{])}
\end{Highlighting}
\end{Shaded}

\begin{center}\rule{0.5\linewidth}{0.5pt}\end{center}

\textbf{Exercise 3.} Coffee at 90°C sits in a 20°C room. The cooling
constant is \(k = 0.05\) min⁻¹. Find \(T(t)\). How long until the coffee
reaches 60°C?

\begin{Shaded}
\begin{Highlighting}[]
\NormalTok{//| echo: false}
\NormalTok{makeStepperHTML([}
\NormalTok{  \{}
\NormalTok{    op: "Write Newton\textquotesingle{}s law of cooling",}
\NormalTok{    eq: "\textbackslash{}\textbackslash{}frac\{dT\}\{dt\} = {-}0.05(T {-} 20), \textbackslash{}\textbackslash{}quad T(0) = 90"}
\NormalTok{  \},}
\NormalTok{  \{}
\NormalTok{    op: "Apply earlier result",}
\NormalTok{    eq: "T(t) = 20 + (90 {-} 20)\textbackslash{}\textbackslash{},e\^{}\{{-}0.05t\} = 20 + 70\textbackslash{}\textbackslash{},e\^{}\{{-}0.05t\}"}
\NormalTok{  \},}
\NormalTok{  \{}
\NormalTok{    op: "Set T = 60",}
\NormalTok{    eq: "20 + 70\textbackslash{}\textbackslash{},e\^{}\{{-}0.05t\} = 60 \textbackslash{}\textbackslash{}implies 70\textbackslash{}\textbackslash{},e\^{}\{{-}0.05t\} = 40"}
\NormalTok{  \},}
\NormalTok{  \{}
\NormalTok{    op: "Divide and take ln",}
\NormalTok{    eq: "e\^{}\{{-}0.05t\} = \textbackslash{}\textbackslash{}tfrac\{4\}\{7\} \textbackslash{}\textbackslash{}implies {-}0.05t = \textbackslash{}\textbackslash{}ln\textbackslash{}\textbackslash{}tfrac\{4\}\{7\}"}
\NormalTok{  \},}
\NormalTok{  \{}
\NormalTok{    op: "Solve for t",}
\NormalTok{    eq: "t = \textbackslash{}\textbackslash{}frac\{{-}\textbackslash{}\textbackslash{}ln(4/7)\}\{0.05\} = 20\textbackslash{}\textbackslash{}ln\textbackslash{}\textbackslash{}tfrac\{7\}\{4\} \textbackslash{}\textbackslash{}approx 11.2 \textbackslash{}\textbackslash{}text\{ min\}",}
\NormalTok{    note: "ln(7/4) ≈ 0.5596."}
\NormalTok{  \}}
\NormalTok{])}
\end{Highlighting}
\end{Shaded}

\begin{center}\rule{0.5\linewidth}{0.5pt}\end{center}

\textbf{Exercise 4.} Solve \(y' + 2y = 4\) using the integrating factor
method. Find the general solution, then the particular solution with
\(y(0) = 1\).

\begin{Shaded}
\begin{Highlighting}[]
\NormalTok{//| echo: false}
\NormalTok{makeStepperHTML([}
\NormalTok{  \{}
\NormalTok{    op: "Already in standard form",}
\NormalTok{    eq: "P(x) = 2, \textbackslash{}\textbackslash{}quad Q(x) = 4"}
\NormalTok{  \},}
\NormalTok{  \{}
\NormalTok{    op: "Compute integrating factor",}
\NormalTok{    eq: "\textbackslash{}\textbackslash{}mu = e\^{}\{\textbackslash{}\textbackslash{}int 2\textbackslash{}\textbackslash{},dx\} = e\^{}\{2x\}"}
\NormalTok{  \},}
\NormalTok{  \{}
\NormalTok{    op: "Multiply through; left side is a perfect derivative",}
\NormalTok{    eq: "\textbackslash{}\textbackslash{}frac\{d\}\{dx\}\textbackslash{}\textbackslash{}!\textbackslash{}\textbackslash{}left[e\^{}\{2x\}y\textbackslash{}\textbackslash{}right] = 4e\^{}\{2x\}"}
\NormalTok{  \},}
\NormalTok{  \{}
\NormalTok{    op: "Integrate both sides",}
\NormalTok{    eq: "e\^{}\{2x\}y = 2e\^{}\{2x\} + C"}
\NormalTok{  \},}
\NormalTok{  \{}
\NormalTok{    op: "Divide by integrating factor",}
\NormalTok{    eq: "y = 2 + Ce\^{}\{{-}2x\} \textbackslash{}\textbackslash{}quad \textbackslash{}\textbackslash{}text\{(general solution)\}"}
\NormalTok{  \},}
\NormalTok{  \{}
\NormalTok{    op: "Apply y(0) = 1",}
\NormalTok{    eq: "1 = 2 + C \textbackslash{}\textbackslash{}implies C = {-}1",}
\NormalTok{    note: "As x → ∞, y → 2: the steady state driven by Q = 4."}
\NormalTok{  \},}
\NormalTok{  \{}
\NormalTok{    op: "Particular solution",}
\NormalTok{    eq: "y = 2 {-} e\^{}\{{-}2x\}"}
\NormalTok{  \}}
\NormalTok{])}
\end{Highlighting}
\end{Shaded}

\begin{center}\rule{0.5\linewidth}{0.5pt}\end{center}

\textbf{Exercise 5.} An RC circuit satisfies \(V' + V = 12\) (taking
\(RC = 1\) s). The capacitor is initially uncharged: \(V(0) = 0\). Find
\(V(t)\). What does \(V(t)\) approach as \(t \to \infty\), and why?

\begin{Shaded}
\begin{Highlighting}[]
\NormalTok{//| echo: false}
\NormalTok{makeStepperHTML([}
\NormalTok{  \{}
\NormalTok{    op: "Standard form: P = 1, Q = 12",}
\NormalTok{    eq: "V\textquotesingle{} + V = 12"}
\NormalTok{  \},}
\NormalTok{  \{}
\NormalTok{    op: "Integrating factor",}
\NormalTok{    eq: "\textbackslash{}\textbackslash{}mu = e\^{}\{\textbackslash{}\textbackslash{}int 1\textbackslash{}\textbackslash{},dt\} = e\^{}t"}
\NormalTok{  \},}
\NormalTok{  \{}
\NormalTok{    op: "Multiply through",}
\NormalTok{    eq: "\textbackslash{}\textbackslash{}frac\{d\}\{dt\}\textbackslash{}\textbackslash{}!\textbackslash{}\textbackslash{}left[e\^{}t V\textbackslash{}\textbackslash{}right] = 12\textbackslash{}\textbackslash{},e\^{}t"}
\NormalTok{  \},}
\NormalTok{  \{}
\NormalTok{    op: "Integrate",}
\NormalTok{    eq: "e\^{}t V = 12\textbackslash{}\textbackslash{},e\^{}t + C"}
\NormalTok{  \},}
\NormalTok{  \{}
\NormalTok{    op: "General solution",}
\NormalTok{    eq: "V(t) = 12 + Ce\^{}\{{-}t\}"}
\NormalTok{  \},}
\NormalTok{  \{}
\NormalTok{    op: "Apply V(0) = 0",}
\NormalTok{    eq: "0 = 12 + C \textbackslash{}\textbackslash{}implies C = {-}12"}
\NormalTok{  \},}
\NormalTok{  \{}
\NormalTok{    op: "Particular solution",}
\NormalTok{    eq: "V(t) = 12\textbackslash{}\textbackslash{}!\textbackslash{}\textbackslash{}left(1 {-} e\^{}\{{-}t\}\textbackslash{}\textbackslash{}right)",}
\NormalTok{    note: "As t → ∞, V → 12 V. The capacitor charges to the source voltage — the exponential transient dies out and only the steady state driven by Q = 12 remains."}
\NormalTok{  \}}
\NormalTok{])}
\end{Highlighting}
\end{Shaded}

\begin{center}\rule{0.5\linewidth}{0.5pt}\end{center}

\textbf{Exercise 6.} A 100 L tank initially contains 10 g of salt. Brine
at 0.5 g/L flows in at 2 L/min; the well-mixed solution flows out at 2
L/min. Set up and solve the ODE for total salt \(S(t)\) (grams). Find
\(S(t)\) and the long-run salt content.

\begin{Shaded}
\begin{Highlighting}[]
\NormalTok{//| echo: false}
\NormalTok{makeStepperHTML([}
\NormalTok{  \{}
\NormalTok{    op: "Rate of salt in",}
\NormalTok{    eq: "\textbackslash{}\textbackslash{}text\{rate in\} = 0.5 \textbackslash{}\textbackslash{}times 2 = 1 \textbackslash{}\textbackslash{}text\{ g/min\}"}
\NormalTok{  \},}
\NormalTok{  \{}
\NormalTok{    op: "Rate of salt out (concentration S/100 times flow 2)",}
\NormalTok{    eq: "\textbackslash{}\textbackslash{}text\{rate out\} = \textbackslash{}\textbackslash{}frac\{S\}\{100\} \textbackslash{}\textbackslash{}times 2 = \textbackslash{}\textbackslash{}frac\{S\}\{50\} \textbackslash{}\textbackslash{}text\{ g/min\}"}
\NormalTok{  \},}
\NormalTok{  \{}
\NormalTok{    op: "ODE",}
\NormalTok{    eq: "\textbackslash{}\textbackslash{}frac\{dS\}\{dt\} = 1 {-} \textbackslash{}\textbackslash{}frac\{S\}\{50\}"}
\NormalTok{  \},}
\NormalTok{  \{}
\NormalTok{    op: "Standard form",}
\NormalTok{    eq: "\textbackslash{}\textbackslash{}frac\{dS\}\{dt\} + \textbackslash{}\textbackslash{}frac\{1\}\{50\}S = 1"}
\NormalTok{  \},}
\NormalTok{  \{}
\NormalTok{    op: "Integrating factor",}
\NormalTok{    eq: "\textbackslash{}\textbackslash{}mu = e\^{}\{t/50\}"}
\NormalTok{  \},}
\NormalTok{  \{}
\NormalTok{    op: "Multiply; integrate",}
\NormalTok{    eq: "e\^{}\{t/50\}S = 50\textbackslash{}\textbackslash{},e\^{}\{t/50\} + C"}
\NormalTok{  \},}
\NormalTok{  \{}
\NormalTok{    op: "General solution",}
\NormalTok{    eq: "S(t) = 50 + Ce\^{}\{{-}t/50\}"}
\NormalTok{  \},}
\NormalTok{  \{}
\NormalTok{    op: "Apply S(0) = 10",}
\NormalTok{    eq: "10 = 50 + C \textbackslash{}\textbackslash{}implies C = {-}40"}
\NormalTok{  \},}
\NormalTok{  \{}
\NormalTok{    op: "Particular solution",}
\NormalTok{    eq: "S(t) = 50 {-} 40\textbackslash{}\textbackslash{},e\^{}\{{-}t/50\}",}
\NormalTok{    note: "As t → ∞, S → 50 g = 0.5 g/L × 100 L. The tank equilibrates to the incoming concentration."}
\NormalTok{  \}}
\NormalTok{])}
\end{Highlighting}
\end{Shaded}

\begin{center}\rule{0.5\linewidth}{0.5pt}\end{center}

\section{Where this goes}\label{where-this-goes-25}

This chapter covered first-order ODEs. The two methods here ---
separation of variables and the integrating factor --- handle a large
class of single first-order equations. They are not the whole story.

\textbf{Second-order ODEs} (Chapter 2) introduce new phenomena:
oscillation, resonance, and damping. A mass on a spring, an LC circuit,
a building subject to seismic loading --- these all require second-order
equations and the concept of complementary and particular solutions.

\textbf{Systems of ODEs} (Chapter 3) handle multiple interdependent
quantities: two species in competition, currents in a multi-loop
circuit, the position and velocity of a mechanical system as coupled
first-order equations. The phase plane --- a picture in state space
rather than time --- provides geometric insight that algebraic solutions
alone cannot.

\textbf{Laplace transforms} (Chapter 4) convert ODEs to algebraic
equations in a new variable, solve them there, and invert the transform.
The technique handles discontinuous forcing (switches, impulses) cleanly
and connects naturally to control system design.

The techniques of this chapter will reappear in all three. The RC
circuit solution you derived here is the building block for
understanding filter transients, step responses, and feedback systems.

\begin{tcolorbox}[enhanced jigsaw, opacityback=0, colback=white, breakable, colframe=quarto-callout-tip-color-frame, arc=.35mm, left=2mm, toprule=.15mm, leftrule=.75mm, bottomrule=.15mm, rightrule=.15mm]
\begin{minipage}[t]{5.5mm}
\textcolor{quarto-callout-tip-color}{\faLightbulb}
\end{minipage}%
\begin{minipage}[t]{\textwidth - 5.5mm}

\textbf{Where this shows up}

\begin{itemize}
\tightlist
\item
  An electrical engineer designing an RC filter needs the time constant
  \(\tau = RC\) to set the cutoff frequency --- it comes directly from
  the ODE solution.
\item
  A pharmacokineticist calculating a dosing interval needs the drug
  half-life \(t_{1/2} = (\ln 2)/k_e\) --- from the same exponential
  decay formula.
\item
  A chemical engineer sizing a mixing tank needs to know how long before
  the outlet concentration reaches 99\% of the inlet --- a mixing
  problem of exactly the form in Exercise 6.
\item
  A control engineer specifying a first-order lag in a block diagram is
  working with \(\tau\,\dot{y} + y = u\) --- the RC circuit equation in
  disguise.
\end{itemize}

The ODE changes labels. The mathematics does not.

\end{minipage}%
\end{tcolorbox}

\chapter{Second-order ODEs}\label{second-order-odes}

Oscillation, damping, and resonance

\hfill\break

In November 1940, four months after it opened, the Tacoma Narrows Bridge
began twisting in a 65 km/h wind. The oscillations grew. Within hours
the deck was flexing through metres of arc. Then it collapsed.

The engineers had missed a resonance condition. The wind's forcing
frequency matched the bridge's natural frequency, and the amplitude grew
without bound. A second-order ODE --- mass, damping, stiffness, forcing
--- contained the answer. The question is whether you know how to read
it.

This chapter develops every tool needed to answer that question. By the
end you will be able to look at any spring, circuit, or structure, write
down its governing equation, and classify its response before solving
it.

\section{Structure of the general
equation}\label{structure-of-the-general-equation}

A second-order linear ODE with constant coefficients has the form

\[ay'' + by' + cy = f(x)\]

where \(a \neq 0\). The coefficients \(a\), \(b\), \(c\) are constants;
\(f(x)\) is the \textbf{forcing function} (or input, or source term).

\textbf{Why second order?} Newton's second law says \(F = ma\), where
\(a = \ddot{x}\) is the second derivative of position. Any system
governed by Newton's second law --- masses on springs, charges in
circuits, beams under load --- produces a second-order ODE. The second
derivative is not an accident of notation; it encodes the physics.

\textbf{Linearity and superposition.} The equation is linear because
\(y\), \(y'\), \(y''\) each appear to the first power and are not
multiplied together. Linearity buys you the superposition principle: if
\(y_1\) and \(y_2\) are both solutions to the homogeneous equation
(\(f = 0\)), then so is \(c_1 y_1 + c_2 y_2\) for any constants \(c_1\),
\(c_2\). This principle is the engine of everything that follows.

\textbf{The decomposition: complementary + particular.} The general
solution is

\[y = y_c + y_p\]

where \(y_c\) (the \textbf{complementary solution}) solves the
homogeneous equation \(ay'' + by' + cy = 0\), and \(y_p\) (the
\textbf{particular solution}) is any one solution to the full equation
with \(f(x)\) present. The two free constants in \(y_c\) absorb the
initial conditions.

\textbf{Why two initial conditions?} The ODE is second order, so its
general solution contains two arbitrary constants. To pick a unique
solution from the family, you need two pieces of information ---
typically \(y(x_0)\) and \(y'(x_0)\). Specify position and velocity at
time zero, and the subsequent motion is determined.

\begin{tcolorbox}[enhanced jigsaw, opacityback=0, colback=white, breakable, colframe=quarto-callout-note-color-frame, arc=.35mm, left=2mm, toprule=.15mm, leftrule=.75mm, bottomrule=.15mm, rightrule=.15mm]
\begin{minipage}[t]{5.5mm}
\textcolor{quarto-callout-note-color}{\faInfo}
\end{minipage}%
\begin{minipage}[t]{\textwidth - 5.5mm}

\textbf{The structure in one line}

\[\underbrace{ay'' + by' + cy = 0}_{\text{homogeneous: find } y_c} \quad + \quad \underbrace{ay'' + by' + cy = f(x)}_{\text{find any } y_p} \quad \Longrightarrow \quad y = y_c + y_p\]

\end{minipage}%
\end{tcolorbox}

\section{Homogeneous equations: the characteristic
equation}\label{homogeneous-equations-the-characteristic-equation}

Take the homogeneous equation

\[ay'' + by' + cy = 0\]

and try a solution of the form \(y = e^{rx}\). Substituting:

\[y = e^{rx}, \quad y' = re^{rx}, \quad y'' = r^2 e^{rx}\]

\[a r^2 e^{rx} + b r e^{rx} + c e^{rx} = 0\]

Factor out \(e^{rx}\), which is never zero:

\[e^{rx}(ar^2 + br + c) = 0 \implies ar^2 + br + c = 0\]

This is the \textbf{characteristic equation}. An ODE has become an
algebraic problem: find the roots of a quadratic. The nature of those
roots determines the form of \(y_c\).

The quadratic formula gives

\[r = \frac{-b \pm \sqrt{b^2 - 4ac}}{2a}\]

The discriminant \(\Delta = b^2 - 4ac\) controls everything.

\subsection{\texorpdfstring{Case 1: Two distinct real roots
(\(\Delta > 0\))}{Case 1: Two distinct real roots (\textbackslash Delta \textgreater{} 0)}}\label{case-1-two-distinct-real-roots-delta-0}

Two distinct roots \(r_1 \neq r_2\), both real. The general solution is

\[y_c = C_1 e^{r_1 x} + C_2 e^{r_2 x}\]

These are two linearly independent solutions; their linear combination
spans the solution space.

\textbf{Physical behaviour:} If both roots are negative, both
exponentials decay. The system returns to equilibrium without
oscillating. In a spring-mass context this is the \textbf{overdamped}
regime --- the mass creeps back slowly, never crossing the equilibrium.

\textbf{Example.} Solve \(y'' - 5y' + 6y = 0\).

Characteristic equation: \(r^2 - 5r + 6 = 0\), which factors as
\((r-2)(r-3) = 0\).

Roots: \(r_1 = 2\), \(r_2 = 3\).

General solution: \(y = C_1 e^{2x} + C_2 e^{3x}\).

\subsection{\texorpdfstring{Case 2: Repeated real root
(\(\Delta = 0\))}{Case 2: Repeated real root (\textbackslash Delta = 0)}}\label{case-2-repeated-real-root-delta-0}

One root \(r = -b/(2a)\) repeated. The naive approach gives only
\(e^{rx}\), which is one solution. A second independent solution is
\(xe^{rx}\) (derivable by reduction of order --- see §6 below, or verify
by substitution).

General solution:

\[y_c = (C_1 + C_2 x)e^{rx}\]

\textbf{Physical behaviour:} The system returns to equilibrium as fast
as possible without oscillating. This is \textbf{critical damping} ---
the boundary between decaying and oscillating behaviour. Car suspension
systems are often tuned close to critical damping.

\textbf{Example.} Solve \(y'' + 4y' + 4y = 0\).

Characteristic equation: \(r^2 + 4r + 4 = (r+2)^2 = 0\).

Root: \(r = -2\) (repeated).

General solution: \(y = (C_1 + C_2 x)e^{-2x}\).

\subsection{\texorpdfstring{Case 3: Complex conjugate roots
(\(\Delta < 0\))}{Case 3: Complex conjugate roots (\textbackslash Delta \textless{} 0)}}\label{case-3-complex-conjugate-roots-delta-0}

The roots are \(r = \alpha \pm \beta i\) where \(\alpha = -b/(2a)\) and
\(\beta = \sqrt{4ac - b^2}/(2a) > 0\).

Using Euler's formula \(e^{i\theta} = \cos\theta + i\sin\theta\), the
complex exponentials can be rewritten as real-valued functions:

\[y_c = e^{\alpha x}(C_1 \cos(\beta x) + C_2 \sin(\beta x))\]

\textbf{Physical behaviour:} The \(e^{\alpha x}\) factor is the
envelope. If \(\alpha < 0\), the oscillation decays --- this is
\textbf{underdamped}. If \(\alpha = 0\), oscillation continues forever
at constant amplitude --- \textbf{undamped simple harmonic motion}. If
\(\alpha > 0\), the oscillation grows (unstable system).

\textbf{Example.} Solve \(y'' + 2y' + 5y = 0\).

Characteristic equation: \(r^2 + 2r + 5 = 0\).

Roots:
\(r = (-2 \pm \sqrt{4 - 20})/2 = (-2 \pm \sqrt{-16})/2 = -1 \pm 2i\).

So \(\alpha = -1\), \(\beta = 2\).

General solution: \(y = e^{-x}(C_1 \cos 2x + C_2 \sin 2x)\).

\begin{tcolorbox}[enhanced jigsaw, opacityback=0, colback=white, breakable, colframe=quarto-callout-note-color-frame, arc=.35mm, left=2mm, toprule=.15mm, leftrule=.75mm, bottomrule=.15mm, rightrule=.15mm]
\begin{minipage}[t]{5.5mm}
\textcolor{quarto-callout-note-color}{\faInfo}
\end{minipage}%
\begin{minipage}[t]{\textwidth - 5.5mm}

\textbf{Summary: three cases}

\begin{longtable}[]{@{}
  >{\raggedright\arraybackslash}p{(\linewidth - 6\tabcolsep) * \real{0.2500}}
  >{\raggedright\arraybackslash}p{(\linewidth - 6\tabcolsep) * \real{0.2500}}
  >{\raggedright\arraybackslash}p{(\linewidth - 6\tabcolsep) * \real{0.2500}}
  >{\raggedright\arraybackslash}p{(\linewidth - 6\tabcolsep) * \real{0.2500}}@{}}
\toprule\noalign{}
\begin{minipage}[b]{\linewidth}\raggedright
Discriminant
\end{minipage} & \begin{minipage}[b]{\linewidth}\raggedright
Roots
\end{minipage} & \begin{minipage}[b]{\linewidth}\raggedright
Solution
\end{minipage} & \begin{minipage}[b]{\linewidth}\raggedright
Behaviour
\end{minipage} \\
\midrule\noalign{}
\endhead
\bottomrule\noalign{}
\endlastfoot
\(\Delta > 0\) & Two distinct real & \(C_1 e^{r_1 x} + C_2 e^{r_2 x}\) &
Overdamped / exponential \\
\(\Delta = 0\) & Repeated real & \((C_1 + C_2 x)e^{rx}\) & Critically
damped \\
\(\Delta < 0\) & Complex conjugate &
\(e^{\alpha x}(C_1\cos\beta x + C_2\sin\beta x)\) & Underdamped /
oscillatory \\
\end{longtable}

\end{minipage}%
\end{tcolorbox}

\begin{Shaded}
\begin{Highlighting}[]
\NormalTok{//| echo: false}
\NormalTok{viewof smdM = Inputs.range([0.1, 5], \{step: 0.1, value: 1, label: "m (mass, kg)"\})}
\NormalTok{viewof smdC = Inputs.range([0, 10], \{step: 0.1, value: 1, label: "c (damping, N·s/m)"\})}
\NormalTok{viewof smdK = Inputs.range([0.1, 20], \{step: 0.1, value: 4, label: "k (stiffness, N/m)"\})}
\end{Highlighting}
\end{Shaded}

\begin{Shaded}
\begin{Highlighting}[]
\NormalTok{//| echo: false}
\NormalTok{\{}
\NormalTok{  const d3 = await require("d3@7");}

\NormalTok{  const width = 560, height = 300;}
\NormalTok{  const margin = \{ top: 20, right: 20, bottom: 45, left: 55 \};}
\NormalTok{  const innerW = width {-} margin.left {-} margin.right;}
\NormalTok{  const innerH = height {-} margin.top {-} margin.bottom;}

\NormalTok{  const m = smdM, c = smdC, k = smdK;}
\NormalTok{  const disc = c * c {-} 4 * m * k;}

\NormalTok{  // Classify regime}
\NormalTok{  const eps = 1e{-}6;}
\NormalTok{  let regime;}
\NormalTok{  if (disc \textgreater{} eps) regime = "Overdamped (Δ \textgreater{} 0)";}
\NormalTok{  else if (disc \textgreater{} {-}eps) regime = "Critically damped (Δ = 0)";}
\NormalTok{  else regime = "Underdamped (Δ \textless{} 0)";}

\NormalTok{  // Euler integration: m x\textquotesingle{}\textquotesingle{} + c x\textquotesingle{} + k x = 0, x(0)=1, x\textquotesingle{}(0)=0}
\NormalTok{  const dt = 0.01, steps = 600;}
\NormalTok{  const ts = [], xs = [];}
\NormalTok{  let x = 1, v = 0;}
\NormalTok{  for (let i = 0; i \textless{}= steps; i++) \{}
\NormalTok{    ts.push(i * dt);}
\NormalTok{    xs.push(x);}
\NormalTok{    const a = ({-}c * v {-} k * x) / m;}
\NormalTok{    v += a * dt;}
\NormalTok{    x += v * dt;}
\NormalTok{  \}}

\NormalTok{  const xExtent = d3.extent(xs);}
\NormalTok{  const yMin = Math.min(xExtent[0], {-}1.05);}
\NormalTok{  const yMax = Math.max(xExtent[1], 1.05);}

\NormalTok{  const xScale = d3.scaleLinear().domain([0, steps * dt]).range([0, innerW]);}
\NormalTok{  const yScale = d3.scaleLinear().domain([yMin, yMax]).range([innerH, 0]);}

\NormalTok{  const svg = d3.create("svg")}
\NormalTok{    .attr("viewBox", \textasciigrave{}0 0 $\{width\} $\{height\}\textasciigrave{})}
\NormalTok{    .attr("width", width)}
\NormalTok{    .attr("height", height)}
\NormalTok{    .style("font{-}family", "sans{-}serif");}

\NormalTok{  const g = svg.append("g").attr("transform", \textasciigrave{}translate($\{margin.left\},$\{margin.top\})\textasciigrave{});}

\NormalTok{  // Zero line}
\NormalTok{  g.append("line")}
\NormalTok{    .attr("x1", 0).attr("x2", innerW)}
\NormalTok{    .attr("y1", yScale(0)).attr("y2", yScale(0))}
\NormalTok{    .attr("stroke", "\#999").attr("stroke{-}dasharray", "4,3").attr("stroke{-}width", 1);}

\NormalTok{  // Exponential envelope for underdamped}
\NormalTok{  if (disc \textless{} {-}eps) \{}
\NormalTok{    const alpha = c / (2 * m);}
\NormalTok{    const envLine = d3.line()}
\NormalTok{      .x((\_, i) =\textgreater{} xScale(ts[i]))}
\NormalTok{      .y((v) =\textgreater{} yScale(v));}

\NormalTok{    const envPos = ts.map(t =\textgreater{} Math.exp({-}alpha * t));}
\NormalTok{    const envNeg = ts.map(t =\textgreater{} {-}Math.exp({-}alpha * t));}

\NormalTok{    g.append("path")}
\NormalTok{      .datum(envPos)}
\NormalTok{      .attr("fill", "none")}
\NormalTok{      .attr("stroke", "\#999")}
\NormalTok{      .attr("stroke{-}dasharray", "5,4")}
\NormalTok{      .attr("stroke{-}width", 1.2)}
\NormalTok{      .attr("d", envLine);}

\NormalTok{    g.append("path")}
\NormalTok{      .datum(envNeg)}
\NormalTok{      .attr("fill", "none")}
\NormalTok{      .attr("stroke", "\#999")}
\NormalTok{      .attr("stroke{-}dasharray", "5,4")}
\NormalTok{      .attr("stroke{-}width", 1.2)}
\NormalTok{      .attr("d", envLine);}
\NormalTok{  \}}

\NormalTok{  // Response curve}
\NormalTok{  const line = d3.line()}
\NormalTok{    .x((\_, i) =\textgreater{} xScale(ts[i]))}
\NormalTok{    .y((v) =\textgreater{} yScale(v));}

\NormalTok{  g.append("path")}
\NormalTok{    .datum(xs)}
\NormalTok{    .attr("fill", "none")}
\NormalTok{    .attr("stroke", "\#3b82f6")}
\NormalTok{    .attr("stroke{-}width", 2)}
\NormalTok{    .attr("d", line);}

\NormalTok{  // Axes}
\NormalTok{  g.append("g").attr("transform", \textasciigrave{}translate(0,$\{innerH\})\textasciigrave{}).call(d3.axisBottom(xScale).ticks(6));}
\NormalTok{  g.append("g").call(d3.axisLeft(yScale).ticks(5));}

\NormalTok{  // Axis labels}
\NormalTok{  g.append("text")}
\NormalTok{    .attr("x", innerW / 2).attr("y", innerH + 38)}
\NormalTok{    .attr("text{-}anchor", "middle").attr("font{-}size", 13)}
\NormalTok{    .text("t (s)");}

\NormalTok{  g.append("text")}
\NormalTok{    .attr("transform", "rotate({-}90)")}
\NormalTok{    .attr("x", {-}innerH / 2).attr("y", {-}42)}
\NormalTok{    .attr("text{-}anchor", "middle").attr("font{-}size", 13)}
\NormalTok{    .text("x(t) (m)");}

\NormalTok{  // Regime label (top{-}right)}
\NormalTok{  const regimeColour = disc \textgreater{} eps ? "\#b45309" : disc \textgreater{} {-}eps ? "\#7c3aed" : "\#059669";}
\NormalTok{  g.append("text")}
\NormalTok{    .attr("x", innerW {-} 5).attr("y", 15)}
\NormalTok{    .attr("text{-}anchor", "end").attr("font{-}size", 12).attr("font{-}weight", "bold")}
\NormalTok{    .attr("fill", regimeColour)}
\NormalTok{    .text(regime);}

\NormalTok{  // ω₀ and ζ}
\NormalTok{  const omega0 = Math.sqrt(k / m);}
\NormalTok{  const zeta = c / (2 * Math.sqrt(k * m));}
\NormalTok{  g.append("text")}
\NormalTok{    .attr("x", innerW {-} 5).attr("y", 32)}
\NormalTok{    .attr("text{-}anchor", "end").attr("font{-}size", 11).attr("fill", "\#555")}
\NormalTok{    .text(\textasciigrave{}ω₀ = $\{omega0.toFixed(3)\}  ζ = $\{zeta.toFixed(3)\}\textasciigrave{});}

\NormalTok{  return svg.node();}
\NormalTok{\}}
\end{Highlighting}
\end{Shaded}

\section{The spring-mass-damper
system}\label{the-spring-mass-damper-system}

The physical archetype for every second-order linear ODE is a mass \(m\)
on a spring with stiffness \(k\), connected to a dashpot (viscous
damper) with damping coefficient \(c\):

\[m\ddot{x} + c\dot{x} + kx = 0\]

Here \(x(t)\) is displacement from equilibrium, dots denote time
derivatives. This maps directly onto the standard form
\(ay'' + by' + cy = 0\) with \(a = m\), \(b = c\),
\(c_{\text{coeff}} = k\).

The characteristic equation is \(mr^2 + cr + k = 0\), with roots

\[r = \frac{-c \pm \sqrt{c^2 - 4mk}}{2m}\]

Two natural parameters emerge from dividing the equation by \(m\):

\[\ddot{x} + \frac{c}{m}\dot{x} + \frac{k}{m}x = 0\]

\textbf{Natural frequency} (undamped):

\[\omega_0 = \sqrt{\frac{k}{m}}\]

\textbf{Damping ratio:}

\[\zeta = \frac{c}{2\sqrt{km}} = \frac{c}{2m\omega_0}\]

The equation becomes
\(\ddot{x} + 2\zeta\omega_0 \dot{x} + \omega_0^2 x = 0\), and the
discriminant condition translates cleanly:

\begin{longtable}[]{@{}
  >{\raggedright\arraybackslash}p{(\linewidth - 4\tabcolsep) * \real{0.3333}}
  >{\raggedright\arraybackslash}p{(\linewidth - 4\tabcolsep) * \real{0.3333}}
  >{\raggedright\arraybackslash}p{(\linewidth - 4\tabcolsep) * \real{0.3333}}@{}}
\toprule\noalign{}
\begin{minipage}[b]{\linewidth}\raggedright
\(\zeta\)
\end{minipage} & \begin{minipage}[b]{\linewidth}\raggedright
Regime
\end{minipage} & \begin{minipage}[b]{\linewidth}\raggedright
Behaviour
\end{minipage} \\
\midrule\noalign{}
\endhead
\bottomrule\noalign{}
\endlastfoot
\(\zeta > 1\) & Overdamped & Two negative real roots; slow return to
equilibrium \\
\(\zeta = 1\) & Critically damped & Repeated negative real root; fastest
return without oscillation \\
\(\zeta < 1\) & Underdamped & Complex roots; decaying oscillation at
frequency \(\omega_d = \omega_0\sqrt{1-\zeta^2}\) \\
\(\zeta = 0\) & Undamped & Pure imaginary roots; oscillation at
\(\omega_0\) forever \\
\end{longtable}

\subsection{Worked numerical example}\label{worked-numerical-example}

A spring-mass-damper system has \(m = 1\) kg, \(c = 5\) N·s/m, \(k = 6\)
N/m. Find \(x(t)\) given \(x(0) = 1\), \(\dot{x}(0) = 0\).

\textbf{Step 1: characteristic equation.}

\[r^2 + 5r + 6 = 0 \implies (r+2)(r+3) = 0 \implies r_1 = -2,\ r_2 = -3\]

\textbf{Step 2: classify.} Both roots real and negative. Discriminant
\(= 25 - 24 = 1 > 0\).

Check via \(\zeta\): \(\omega_0 = \sqrt{6} \approx 2.449\),
\(\zeta = 5/(2\sqrt{6}) \approx 1.02 > 1\). Overdamped.

\textbf{Step 3: general solution.}

\[x(t) = C_1 e^{-2t} + C_2 e^{-3t}\]

\textbf{Step 4: apply initial conditions.}

\(x(0) = 1\): \(C_1 + C_2 = 1\)

\(\dot{x}(0) = 0\): \(-2C_1 - 3C_2 = 0 \implies C_1 = -3C_2/2\)

Substituting:
\(-3C_2/2 + C_2 = 1 \implies -C_2/2 = 1 \implies C_2 = -2\),
\(C_1 = 3\).

\textbf{Solution:}

\[x(t) = 3e^{-2t} - 2e^{-3t}\]

The mass starts at \(x = 1\), initially moves slightly before the two
exponentials pull it back to zero without oscillating. Physical
interpretation: heavy damping prevents overshoot.

Now change the damping: \(m = 1\), \(c = 2\), \(k = 1\). Characteristic
equation: \(r^2 + 2r + 1 = (r+1)^2 = 0\). Root \(r = -1\) repeated.
Critically damped. With \(x(0) = 1\), \(\dot{x}(0) = 0\):

\(x = (C_1 + C_2 t)e^{-t}\); \(x(0) = C_1 = 1\);
\(\dot{x}(0) = C_2 - C_1 = 0 \implies C_2 = 1\).

\[x(t) = (1 + t)e^{-t}\]

Now change to \(m = 1\), \(c = 1\), \(k = 4\). Characteristic equation:
\(r^2 + r + 4 = 0\).

Roots: \(r = (-1 \pm \sqrt{1-16})/2 = -1/2 \pm i\sqrt{15}/2\).
Underdamped. General solution:
\(x(t) = e^{-t/2}(C_1\cos\frac{\sqrt{15}}{2}t + C_2\sin\frac{\sqrt{15}}{2}t)\).

With \(x(0) = 1\), \(\dot{x}(0) = 0\):

\(x(0) = C_1 = 1\).

Differentiate using the product rule:
\(\dot{x}(t) = -\tfrac{1}{2}e^{-t/2}(C_1\cos\frac{\sqrt{15}}{2}t + C_2\sin\frac{\sqrt{15}}{2}t) + e^{-t/2}(-C_1\tfrac{\sqrt{15}}{2}\sin\frac{\sqrt{15}}{2}t + C_2\tfrac{\sqrt{15}}{2}\cos\frac{\sqrt{15}}{2}t)\).

At \(t = 0\):
\(\dot{x}(0) = -\tfrac{1}{2}C_1 + C_2\tfrac{\sqrt{15}}{2} = 0\), so
\(C_2 = \tfrac{C_1}{\sqrt{15}} = \tfrac{1}{\sqrt{15}}\).

\[x(t) = e^{-t/2}\!\left(\cos\frac{\sqrt{15}}{2}t + \frac{1}{\sqrt{15}}\sin\frac{\sqrt{15}}{2}t\right)\]

The oscillation is visible; the exponential envelope \(e^{-t/2}\) decays
it.

\begin{tcolorbox}[enhanced jigsaw, opacityback=0, colback=white, breakable, colframe=quarto-callout-tip-color-frame, arc=.35mm, left=2mm, toprule=.15mm, leftrule=.75mm, bottomrule=.15mm, rightrule=.15mm]
\begin{minipage}[t]{5.5mm}
\textcolor{quarto-callout-tip-color}{\faLightbulb}
\end{minipage}%
\begin{minipage}[t]{\textwidth - 5.5mm}

\textbf{Engineering rule of thumb}

\(\zeta > 1\): overdamped --- sluggish, no oscillation. \(\zeta = 1\):
critically damped --- fastest non-oscillatory return. \(\zeta < 1\):
underdamped --- fast but overshoots; for \(\zeta \ll 1\) the oscillation
persists for many cycles.

Automotive suspension targets \(\zeta \approx 0.3\)--\(0.4\): some
overshoot is acceptable; too much damping makes the ride stiff and slow
to recover.

\end{minipage}%
\end{tcolorbox}

\section{Method of undetermined
coefficients}\label{method-of-undetermined-coefficients}

When \(f(x) \neq 0\), we need a particular solution \(y_p\). The method
of undetermined coefficients works when \(f(x)\) belongs to a standard
family: polynomials, exponentials, sines and cosines, and products of
these.

The strategy: guess \(y_p\) has the same form as \(f(x)\), with unknown
coefficients. Substitute into the ODE; match coefficients on both sides.

\subsection{Standard guess table}\label{standard-guess-table}

\begin{longtable}[]{@{}
  >{\raggedright\arraybackslash}p{(\linewidth - 2\tabcolsep) * \real{0.5000}}
  >{\raggedright\arraybackslash}p{(\linewidth - 2\tabcolsep) * \real{0.5000}}@{}}
\toprule\noalign{}
\begin{minipage}[b]{\linewidth}\raggedright
\(f(x)\)
\end{minipage} & \begin{minipage}[b]{\linewidth}\raggedright
Guess for \(y_p\)
\end{minipage} \\
\midrule\noalign{}
\endhead
\bottomrule\noalign{}
\endlastfoot
\(P_n(x)\) (polynomial, degree \(n\)) &
\(A_n x^n + \cdots + A_1 x + A_0\) \\
\(e^{\alpha x}\) & \(Ae^{\alpha x}\) \\
\(\cos(\beta x)\) or \(\sin(\beta x)\) &
\(A\cos(\beta x) + B\sin(\beta x)\) \\
\(e^{\alpha x}\cos(\beta x)\) or similar &
\(e^{\alpha x}(A\cos(\beta x) + B\sin(\beta x))\) \\
\(P_n(x)e^{\alpha x}\) & \(e^{\alpha x}(A_n x^n + \cdots + A_0)\) \\
\end{longtable}

\textbf{The modification rule.} If the guess for \(y_p\) duplicates a
term already present in \(y_c\), multiply the guess by \(x\) (or \(x^2\)
if \(y_c\) contains both \(e^{rx}\) and \(xe^{rx}\), as in the
repeated-root case). Failure to apply this rule produces a contradiction
when you substitute.

\subsection{Worked example 1: exponential
forcing}\label{worked-example-1-exponential-forcing}

Find a particular solution of \(y'' + 3y' + 2y = e^{-3x}\).

Homogeneous equation: \(r^2 + 3r + 2 = (r+1)(r+2) = 0\), roots
\(r = -1, -2\).

\(y_c = C_1 e^{-x} + C_2 e^{-2x}\).

\(f(x) = e^{-3x}\): the exponent \(-3\) does not match any root, so no
modification is needed. Guess \(y_p = Ae^{-3x}\).

Substitute:

\[y_p'' + 3y_p' + 2y_p = 9Ae^{-3x} - 9Ae^{-3x} + 2Ae^{-3x} = 2Ae^{-3x}\]

Set equal to \(e^{-3x}\): \(2A = 1 \implies A = 1/2\).

\[y_p = \tfrac{1}{2}e^{-3x}\]

General solution: \(y = C_1 e^{-x} + C_2 e^{-2x} + \frac{1}{2}e^{-3x}\).

\subsection{Worked example 2: sinusoidal
forcing}\label{worked-example-2-sinusoidal-forcing}

Find a particular solution of \(y'' + 4y' + 5y = \sin(2x)\).

Homogeneous: \(r^2 + 4r + 5 = 0 \implies r = -2 \pm i\). So
\(y_c = e^{-2x}(C_1\cos x + C_2\sin x)\).

\(f(x) = \sin(2x)\): guess \(y_p = A\cos(2x) + B\sin(2x)\).

Compute derivatives:

\[y_p' = -2A\sin(2x) + 2B\cos(2x)\] \[y_p'' = -4A\cos(2x) - 4B\sin(2x)\]

Substitute:

\[(-4A\cos 2x - 4B\sin 2x) + 4(-2A\sin 2x + 2B\cos 2x) + 5(A\cos 2x + B\sin 2x) = \sin 2x\]

Collect \(\cos(2x)\): \(-4A + 8B + 5A = (A + 8B)\cos 2x\)

Collect \(\sin(2x)\): \(-4B - 8A + 5B = (B - 8A)\sin 2x\)

Match to \(0\cdot\cos 2x + 1\cdot\sin 2x\):

\[A + 8B = 0 \implies A = -8B\]
\[B - 8A = 1 \implies B + 64B = 1 \implies B = \tfrac{1}{65},\quad A = -\tfrac{8}{65}\]

\[y_p = -\frac{8}{65}\cos(2x) + \frac{1}{65}\sin(2x)\]

\subsection{Modification rule in
action}\label{modification-rule-in-action}

Find a particular solution of \(y'' + 2y' + y = e^{-x}\).

Characteristic equation: \((r+1)^2 = 0\), root \(r = -1\) repeated.

\(y_c = (C_1 + C_2 x)e^{-x}\).

Naive guess \(Ae^{-x}\) is in \(y_c\); modified guess \(Axe^{-x}\) is
also in \(y_c\) (the \(C_2 x e^{-x}\) term). So multiply again: guess
\(y_p = Ax^2 e^{-x}\).

\[y_p' = (2Ax - Ax^2)e^{-x}, \quad y_p'' = (2A - 4Ax + Ax^2)e^{-x}\]

Substitute (the \(x^2 e^{-x}\) and \(xe^{-x}\) terms cancel):

\[y_p'' + 2y_p' + y_p = 2Ae^{-x}\]

Match: \(2A = 1 \implies A = 1/2\).

\[y_p = \tfrac{1}{2}x^2 e^{-x}\]

\section{Forced oscillation and
resonance}\label{forced-oscillation-and-resonance}

Now force the undamped spring-mass system:

\[m\ddot{x} + kx = F_0 \cos(\omega t), \quad x(0) = 0,\ \dot{x}(0) = 0\]

The homogeneous solution is
\(x_c = C_1\cos(\omega_0 t) + C_2\sin(\omega_0 t)\) where
\(\omega_0 = \sqrt{k/m}\).

\textbf{Case 1: \(\omega \neq \omega_0\) (off resonance).} Guess
\(x_p = A\cos(\omega t)\).

\[-m\omega^2 A\cos(\omega t) + kA\cos(\omega t) = F_0\cos(\omega t)\]

\[A(k - m\omega^2) = F_0 \implies A = \frac{F_0}{k - m\omega^2} = \frac{F_0/m}{\omega_0^2 - \omega^2}\]

Applying initial conditions \(x(0) = 0\), \(\dot{x}(0) = 0\): the
general solution is
\(x = C_1\cos\omega_0 t + C_2\sin\omega_0 t + \frac{F_0/m}{\omega_0^2-\omega^2}\cos\omega t\).
From \(x(0) = 0\): \(C_1 + \frac{F_0/m}{\omega_0^2-\omega^2} = 0\), so
\(C_1 = -\frac{F_0/m}{\omega_0^2-\omega^2}\). From \(\dot{x}(0) = 0\):
\(C_2\omega_0 - \frac{F_0/m}{\omega_0^2-\omega^2}\omega\sin 0 = 0\),
giving \(C_2 = 0\). Substituting:

\[x(t) = \frac{F_0/m}{\omega_0^2 - \omega^2}\bigl(\cos(\omega t) - \cos(\omega_0 t)\bigr)\]

\textbf{Beating.} When \(\omega\) is close to but not equal to
\(\omega_0\), write \(\omega_0 = \omega + \epsilon\) for small
\(\epsilon\). The difference \(\cos(\omega t) - \cos(\omega_0 t)\) can
be rewritten using a product-to-sum identity (from
\(\cos A - \cos B = 2\sin\frac{A+B}{2}\sin\frac{B-A}{2}\)):

\[\cos(\omega t) - \cos(\omega_0 t) = 2\sin\!\left(\frac{\omega_0 + \omega}{2}t\right)\sin\!\left(\frac{\omega_0 - \omega}{2}t\right)\]

This is a rapid oscillation at frequency
\((\omega_0 + \omega)/2 \approx \omega_0\) modulated by a slow envelope
at frequency \((\omega_0 - \omega)/2 \approx \epsilon/2\). The amplitude
swells and collapses periodically --- \textbf{beating}. It is audible
between two tuning forks of slightly different pitch.

\textbf{Case 2: \(\omega = \omega_0\) (resonance).} The guess
\(A\cos(\omega_0 t)\) duplicates a homogeneous solution; apply the
modification rule and guess
\(x_p = t(A\cos(\omega_0 t) + B\sin(\omega_0 t))\).

Substituting and simplifying (the cosine terms cancel; only the sine
survives in the second derivative's cross-terms):

\[x_p = \frac{F_0}{2m\omega_0} t\sin(\omega_0 t)\]

The particular solution at resonance is

\[\boxed{x_p = \frac{F_0}{2m\omega_0}\, t\sin(\omega_0 t)}\]

The factor of \(t\) in front of the sine is critical. The amplitude
\(\frac{F_0}{2m\omega_0} t\) grows \textbf{linearly without bound} as
\(t \to \infty\).

\begin{tcolorbox}[enhanced jigsaw, opacityback=0, colback=white, breakable, colframe=quarto-callout-warning-color-frame, arc=.35mm, left=2mm, toprule=.15mm, leftrule=.75mm, bottomrule=.15mm, rightrule=.15mm]
\begin{minipage}[t]{5.5mm}
\textcolor{quarto-callout-warning-color}{\faExclamationTriangle}
\end{minipage}%
\begin{minipage}[t]{\textwidth - 5.5mm}

\textbf{Resonance is an unbounded response}

At \(\omega = \omega_0\), the forcing and the natural oscillation are
perfectly in phase. Each cycle, the force adds energy to the system.
With no damping to dissipate it, amplitude grows without limit. In a
real system, something fails first --- yielding, fracture, or collapse.

The Tacoma Narrows Bridge, Broughton Suspension Bridge (1831), and
numerous aircraft structures have failed under resonance conditions. In
engineering design, knowing \(\omega_0\) and ensuring forcing
frequencies avoid it --- or adding sufficient damping --- is not
optional.

\end{minipage}%
\end{tcolorbox}

\textbf{With damping present.} Add a damping term \(c\dot{x}\) to the
forced equation. The steady-state amplitude response at driving
frequency \(\omega\) becomes

\[X(\omega) = \frac{F_0/k}{\sqrt{(1 - r^2)^2 + (2\zeta r)^2}}, \quad r = \frac{\omega}{\omega_0}\]

(This formula is derivable by applying undetermined coefficients to
\(m\ddot{x}+c\dot{x}+kx=F_0\cos\omega t\) --- the calculation is
straightforward but lengthy, and we state the result here.)

As \(\zeta \to 0\), the peak amplitude at \(r = 1\) becomes unbounded.
For \(\zeta > 0\), the peak is finite but can be very large for small
\(\zeta\). The peak occurs near \(r = 1\) and grows sharper and taller
as damping decreases --- the amplification factor \(Q = 1/(2\zeta)\) is
called the \textbf{quality factor} in electrical engineering.

\begin{Shaded}
\begin{Highlighting}[]
\NormalTok{//| echo: false}
\NormalTok{viewof forceC = Inputs.range([0, 5], \{step: 0.05, value: 1, label: "Damping c"\})}
\NormalTok{viewof forceW0 = Inputs.range([0.1, 5], \{step: 0.05, value: 2, label: "Natural frequency ω₀"\})}
\end{Highlighting}
\end{Shaded}

\begin{Shaded}
\begin{Highlighting}[]
\NormalTok{//| echo: false}
\NormalTok{\{}
\NormalTok{  const d3 = await require("d3@7");}

\NormalTok{  const width = 560, height = 280;}
\NormalTok{  const margin = \{ top: 30, right: 20, bottom: 48, left: 60 \};}
\NormalTok{  const innerW = width {-} margin.left {-} margin.right;}
\NormalTok{  const innerH = height {-} margin.top {-} margin.bottom;}

\NormalTok{  const c = forceC;}
\NormalTok{  const w0 = forceW0;}
\NormalTok{  // m=1, k=w0\^{}2}
\NormalTok{  const m = 1, k = w0 * w0;}
\NormalTok{  const zeta = c / (2 * Math.sqrt(m * k));}

\NormalTok{  // Amplitude response: X / (F0/k) = 1 / sqrt((1{-}r\^{}2)\^{}2 + (2*zeta*r)\^{}2)}
\NormalTok{  const nPoints = 400;}
\NormalTok{  const wMax = 2.5 * w0;}
\NormalTok{  const data = [];}
\NormalTok{  for (let i = 0; i \textless{}= nPoints; i++) \{}
\NormalTok{    const w = (i / nPoints) * wMax;}
\NormalTok{    const r = w / w0;}
\NormalTok{    const denom = Math.sqrt((1 {-} r * r) * (1 {-} r * r) + (2 * zeta * r) * (2 * zeta * r));}
\NormalTok{    const amp = denom \textless{} 1e{-}9 ? 1e6 : 1 / denom;}
\NormalTok{    data.push(\{ w, amp \});}
\NormalTok{  \}}

\NormalTok{  const ampValues = data.map(d =\textgreater{} d.amp).filter(isFinite);}
\NormalTok{  const yMax = d3.max(ampValues) * 1.10;}

\NormalTok{  const xScale = d3.scaleLinear().domain([0, wMax]).range([0, innerW]);}
\NormalTok{  const yScale = d3.scaleLinear().domain([0, yMax]).range([innerH, 0]);}

\NormalTok{  // Colour: high c → blue, low c → red (c in [0,5])}
\NormalTok{  const colourT = Math.min(1, c / 5);}
\NormalTok{  const curveColour = d3.interpolateRdYlBu(colourT);}

\NormalTok{  const svg = d3.create("svg")}
\NormalTok{    .attr("viewBox", \textasciigrave{}0 0 $\{width\} $\{height\}\textasciigrave{})}
\NormalTok{    .attr("width", width)}
\NormalTok{    .attr("height", height)}
\NormalTok{    .style("font{-}family", "sans{-}serif");}

\NormalTok{  const g = svg.append("g").attr("transform", \textasciigrave{}translate($\{margin.left\},$\{margin.top\})\textasciigrave{});}

\NormalTok{  // ω₀ vertical dashed line}
\NormalTok{  g.append("line")}
\NormalTok{    .attr("x1", xScale(w0)).attr("x2", xScale(w0))}
\NormalTok{    .attr("y1", 0).attr("y2", innerH)}
\NormalTok{    .attr("stroke", "\#777").attr("stroke{-}dasharray", "5,4").attr("stroke{-}width", 1.2);}

\NormalTok{  g.append("text")}
\NormalTok{    .attr("x", xScale(w0) + 4).attr("y", 12)}
\NormalTok{    .attr("font{-}size", 12).attr("fill", "\#555")}
\NormalTok{    .text("ω₀");}

\NormalTok{  // Amplitude curve}
\NormalTok{  const line = d3.line()}
\NormalTok{    .x(d =\textgreater{} xScale(d.w))}
\NormalTok{    .y(d =\textgreater{} Math.max(0, yScale(Math.min(d.amp, yMax * 1.05))))}
\NormalTok{    .defined(d =\textgreater{} isFinite(d.amp));}

\NormalTok{  g.append("path")}
\NormalTok{    .datum(data)}
\NormalTok{    .attr("fill", "none")}
\NormalTok{    .attr("stroke", curveColour)}
\NormalTok{    .attr("stroke{-}width", 2.2)}
\NormalTok{    .attr("d", line);}

\NormalTok{  // Resonance peak}
\NormalTok{  // Peak near w = w0 * sqrt(1 {-} 2*zeta\^{}2), only exists if zeta \textless{} 1/sqrt(2)}
\NormalTok{  let peakW = null, peakAmp = null;}
\NormalTok{  if (zeta \textless{} 1 / Math.SQRT2) \{}
\NormalTok{    peakW = w0 * Math.sqrt(Math.max(0, 1 {-} 2 * zeta * zeta));}
\NormalTok{    const r = peakW / w0;}
\NormalTok{    const denom = Math.sqrt((1 {-} r * r) * (1 {-} r * r) + (2 * zeta * r) * (2 * zeta * r));}
\NormalTok{    peakAmp = denom \textless{} 1e{-}9 ? null : 1 / denom;}
\NormalTok{  \} else \{}
\NormalTok{    // No resonance peak — show value at w=0 (amp=1)}
\NormalTok{    peakW = 0;}
\NormalTok{    peakAmp = 1;}
\NormalTok{  \}}

\NormalTok{  if (peakAmp !== null \&\& isFinite(peakAmp) \&\& peakAmp \textless{}= yMax * 1.05) \{}
\NormalTok{    g.append("circle")}
\NormalTok{      .attr("cx", xScale(peakW)).attr("cy", yScale(Math.min(peakAmp, yMax)))}
\NormalTok{      .attr("r", 4)}
\NormalTok{      .attr("fill", curveColour).attr("stroke", "\#222").attr("stroke{-}width", 1);}
\NormalTok{  \}}

\NormalTok{  // Resonance label}
\NormalTok{  const peakWLabel = zeta \textless{} 1 / Math.SQRT2}
\NormalTok{    ? \textasciigrave{}ω ≈ ω₀√(1−2ζ²) = $\{(w0 * Math.sqrt(Math.max(0, 1 {-} 2 * zeta * zeta))).toFixed(3)\}\textasciigrave{}}
\NormalTok{    : "No resonance peak (ζ ≥ 1/√2)";}

\NormalTok{  g.append("text")}
\NormalTok{    .attr("x", innerW {-} 5).attr("y", 14)}
\NormalTok{    .attr("text{-}anchor", "end").attr("font{-}size", 11).attr("fill", "\#555")}
\NormalTok{    .text(\textasciigrave{}Resonance peak at $\{peakWLabel\}\textasciigrave{});}

\NormalTok{  // Axes}
\NormalTok{  g.append("g").attr("transform", \textasciigrave{}translate(0,$\{innerH\})\textasciigrave{}).call(d3.axisBottom(xScale).ticks(6));}
\NormalTok{  g.append("g").call(d3.axisLeft(yScale).ticks(5));}

\NormalTok{  // Axis labels}
\NormalTok{  g.append("text")}
\NormalTok{    .attr("x", innerW / 2).attr("y", innerH + 40)}
\NormalTok{    .attr("text{-}anchor", "middle").attr("font{-}size", 13)}
\NormalTok{    .text("Driving frequency ω");}

\NormalTok{  g.append("text")}
\NormalTok{    .attr("transform", "rotate({-}90)")}
\NormalTok{    .attr("x", {-}innerH / 2).attr("y", {-}48)}
\NormalTok{    .attr("text{-}anchor", "middle").attr("font{-}size", 13)}
\NormalTok{    .text("X(ω) / (F₀/k)");}

\NormalTok{  // ζ annotation}
\NormalTok{  g.append("text")}
\NormalTok{    .attr("x", innerW {-} 5).attr("y", 28)}
\NormalTok{    .attr("text{-}anchor", "end").attr("font{-}size", 11).attr("fill", "\#555")}
\NormalTok{    .text(\textasciigrave{}ζ = $\{zeta.toFixed(3)\}\textasciigrave{});}

\NormalTok{  return svg.node();}
\NormalTok{\}}
\end{Highlighting}
\end{Shaded}

\section{Reduction of order}\label{reduction-of-order}

Suppose you know one solution \(y_1\) to \(ay'' + by' + cy = 0\) ---
perhaps from inspection or from the characteristic equation in a
non-constant coefficient problem --- but cannot find a second
independent one directly. The method of \textbf{reduction of order}
produces \(y_2\).

Set \(y = v(x)\,y_1(x)\) where \(v(x)\) is an unknown function. Compute
\(y'\) and \(y''\), substitute into the ODE. Terms involving \(v\)
(without derivatives) cancel --- this is guaranteed because \(y_1\)
itself satisfies the equation. The result is a first-order ODE in
\(w = v'\), which can be solved by the integrating factor method from
Chapter 1.

For the constant-coefficient case, this is how the repeated-root
solution \(xe^{rx}\) is derived. Set \(y_1 = e^{rx}\), \(y = ve^{rx}\):

\[y' = (v' + rv)e^{rx}, \quad y'' = (v'' + 2rv' + r^2v)e^{rx}\]

Substituting into \(ay'' + by' + cy = 0\) and using
\(ar^2 + br + c = 0\):

\[a(v'' + 2rv') + bv' = 0 \implies av'' + (2ar + b)v' = 0\]

At the repeated root \(r = -b/(2a)\), we have \(2ar + b = 0\), so
\(av'' = 0\), giving \(v'' = 0\) and thus \(v = C_1 + C_2 x\). Therefore
\(y = (C_1 + C_2 x)e^{rx}\) --- the result we stated earlier, now
derived.

Reduction of order becomes more important when dealing with
variable-coefficient equations (Euler--Cauchy type) or when one solution
is known from a series method. It will reappear in context.

\section{Where this leads}\label{where-this-leads}

You now have the complete toolkit for second-order linear ODEs with
constant coefficients: classify by roots of the characteristic equation;
add a particular solution via undetermined coefficients; understand why
resonance is structurally different from off-resonance forcing.

Three threads open from here:

\textbf{Laplace transforms} (Chapter 4) handle initial-value problems
algorithmically, absorb discontinuous forcing naturally, and are the
standard industrial tool for control systems analysis. The Laplace
transform of \(y''\) introduces the transfer function, which encodes the
characteristic equation in a different algebraic form.

\textbf{Systems and phase plane} (Chapter 3) extend the ideas to pairs
of first-order equations --- which is what a second-order ODE becomes
when you introduce \(v = y'\) as a second unknown. Phase-plane portraits
make the overdamped/underdamped/critically damped classification visual.

\textbf{Fourier series} (Volume 7, Fourier--PDEs section) answer the
question: what if \(f(x)\) is not a simple sinusoid but a periodic
function of arbitrary shape? Decompose it into harmonics, solve for
each, and superpose. The resonance analysis extends directly.

\begin{tcolorbox}[enhanced jigsaw, opacityback=0, colback=white, breakable, colframe=quarto-callout-tip-color-frame, arc=.35mm, left=2mm, toprule=.15mm, leftrule=.75mm, bottomrule=.15mm, rightrule=.15mm]
\begin{minipage}[t]{5.5mm}
\textcolor{quarto-callout-tip-color}{\faLightbulb}
\end{minipage}%
\begin{minipage}[t]{\textwidth - 5.5mm}

\textbf{Where this shows up}

\begin{itemize}
\tightlist
\item
  A structural engineer computing the natural frequencies of a building
  must keep them away from the frequencies of earthquakes and wind ---
  both sources of periodic forcing.
\item
  An electrical engineer designing an RLC filter tunes \(\omega_0\) and
  \(\zeta\) to pass or reject particular frequency bands. The quality
  factor \(Q = 1/(2\zeta)\) measures selectivity.
\item
  A control systems engineer designing a PID controller must ensure the
  closed-loop characteristic equation has roots with negative real parts
  --- the same stability condition as \(\alpha < 0\) for the underdamped
  spring.
\item
  A physicist modelling a damped quantum harmonic oscillator writes down
  the same equation. The mathematics is unchanged; only the
  interpretation of \(x\) differs.
\end{itemize}

\end{minipage}%
\end{tcolorbox}

\section{Exercises}\label{exercises-27}

\begin{Shaded}
\begin{Highlighting}[]
\NormalTok{//| echo: false}
\NormalTok{katex = require("katex@0.16")}
\end{Highlighting}
\end{Shaded}

\begin{Shaded}
\begin{Highlighting}[]
\NormalTok{//| echo: false}
\NormalTok{function makeStepperHTML(exerciseNum, steps) \{}
\NormalTok{  let current = 0;}
\NormalTok{  const totalSteps = steps.length;}
\NormalTok{  const container = document.createElement("div");}
\NormalTok{  container.style.cssText = "border:1px solid \#e5e7eb; border{-}radius:8px; padding:1rem 1.25rem; margin:0.75rem 0; font{-}family:inherit;";}
\NormalTok{  const stepsDiv = document.createElement("div");}
\NormalTok{  stepsDiv.style.marginBottom = "0.75rem";}

\NormalTok{  function renderSteps() \{}
\NormalTok{    stepsDiv.innerHTML = "";}
\NormalTok{    for (let i = 0; i \textless{} current; i++) \{}
\NormalTok{      const s = steps[i];}
\NormalTok{      const row = document.createElement("div");}
\NormalTok{      row.style.cssText = "display:grid; grid{-}template{-}columns:220px 1fr; gap:0.35rem 1rem; align{-}items:baseline; padding:0.35rem 0; border{-}top:1px solid \#f3f4f6;";}
\NormalTok{      const opCell = document.createElement("span");}
\NormalTok{      opCell.style.cssText = "font{-}size:0.85em; color:\#6b7280; font{-}style:italic;";}
\NormalTok{      opCell.textContent = s.op;}
\NormalTok{      const eqCell = document.createElement("span");}
\NormalTok{      eqCell.innerHTML = katex.renderToString(s.eq, \{ throwOnError: false, displayMode: false \});}
\NormalTok{      row.appendChild(opCell);}
\NormalTok{      row.appendChild(eqCell);}
\NormalTok{      if (s.note) \{}
\NormalTok{        const noteRow = document.createElement("div");}
\NormalTok{        noteRow.style.cssText = "grid{-}column:1 / {-}1; font{-}size:0.82em; color:\#6b7280; padding:0.1rem 0 0.25rem 0; border{-}top: none;";}
\NormalTok{        noteRow.textContent = s.note;}
\NormalTok{        row.appendChild(noteRow);}
\NormalTok{      \}}
\NormalTok{      stepsDiv.appendChild(row);}
\NormalTok{    \}}
\NormalTok{    if (current === totalSteps) \{}
\NormalTok{      const done = document.createElement("div");}
\NormalTok{      done.style.cssText = "margin{-}top:0.5rem; font{-}size:0.9em; color:\#059669; font{-}weight:500;";}
\NormalTok{      done.textContent = "✓ Solution complete";}
\NormalTok{      stepsDiv.appendChild(done);}
\NormalTok{    \}}
\NormalTok{  \}}

\NormalTok{  const controls = document.createElement("div");}
\NormalTok{  controls.style.cssText = "display:flex; gap:0.5rem; align{-}items:center;";}
\NormalTok{  const nextBtn = document.createElement("button");}
\NormalTok{  nextBtn.textContent = "Next step →";}
\NormalTok{  nextBtn.style.cssText = "padding:0.35rem 0.85rem; border:1px solid \#d1d5db; border{-}radius:4px; background:\#fff; cursor:pointer; font{-}size:0.9em;";}
\NormalTok{  const resetBtn = document.createElement("button");}
\NormalTok{  resetBtn.textContent = "Reset";}
\NormalTok{  resetBtn.style.cssText = "padding:0.35rem 0.75rem; border:1px solid \#e5e7eb; border{-}radius:4px; background:\#f9fafb; cursor:pointer; font{-}size:0.9em; color:\#6b7280;";}
\NormalTok{  const counter = document.createElement("span");}
\NormalTok{  counter.style.cssText = "font{-}size:0.82em; color:\#9ca3af; margin{-}left:0.25rem;";}

\NormalTok{  function updateButtons() \{}
\NormalTok{    nextBtn.disabled = current === totalSteps;}
\NormalTok{    nextBtn.style.opacity = current === totalSteps ? "0.4" : "1";}
\NormalTok{    counter.textContent = current === 0 ? "Click to reveal steps" : \textasciigrave{}Step $\{current\} of $\{totalSteps\}\textasciigrave{};}
\NormalTok{  \}}

\NormalTok{  nextBtn.onclick = () =\textgreater{} \{ if (current \textless{} totalSteps) \{ current++; renderSteps(); updateButtons(); \} \};}
\NormalTok{  resetBtn.onclick = () =\textgreater{} \{ current = 0; renderSteps(); updateButtons(); \};}
\NormalTok{  controls.appendChild(nextBtn);}
\NormalTok{  controls.appendChild(resetBtn);}
\NormalTok{  controls.appendChild(counter);}
\NormalTok{  container.appendChild(stepsDiv);}
\NormalTok{  container.appendChild(controls);}
\NormalTok{  renderSteps();}
\NormalTok{  updateButtons();}
\NormalTok{  return container;}
\NormalTok{\}}
\end{Highlighting}
\end{Shaded}

\textbf{1.} Find the general solution of \(y'' - 5y' + 6y = 0\).

\begin{Shaded}
\begin{Highlighting}[]
\NormalTok{//| echo: false}
\NormalTok{makeStepperHTML(1, [}
\NormalTok{  \{ op: "Characteristic equation",    eq: "r\^{}2 {-} 5r + 6 = 0" \},}
\NormalTok{  \{ op: "Factor",                      eq: "(r{-}2)(r{-}3) = 0" \},}
\NormalTok{  \{ op: "Roots",                       eq: "r\_1 = 2,\textbackslash{}\textbackslash{}quad r\_2 = 3", note: "Discriminant Δ = 25 − 24 = 1 \textgreater{} 0: two distinct real roots." \},}
\NormalTok{  \{ op: "General solution",            eq: "y = C\_1 e\^{}\{2x\} + C\_2 e\^{}\{3x\}" \}}
\NormalTok{])}
\end{Highlighting}
\end{Shaded}

\begin{center}\rule{0.5\linewidth}{0.5pt}\end{center}

\textbf{2.} Solve \(y'' + 4y' + 4y = 0\), \(y(0) = 1\), \(y'(0) = 0\).

\begin{Shaded}
\begin{Highlighting}[]
\NormalTok{//| echo: false}
\NormalTok{makeStepperHTML(2, [}
\NormalTok{  \{ op: "Characteristic equation",    eq: "r\^{}2 + 4r + 4 = 0" \},}
\NormalTok{  \{ op: "Factor",                      eq: "(r+2)\^{}2 = 0" \},}
\NormalTok{  \{ op: "Repeated root",               eq: "r = {-}2", note: "Discriminant Δ = 16 − 16 = 0: critically damped." \},}
\NormalTok{  \{ op: "General solution",            eq: "y = (C\_1 + C\_2 x)e\^{}\{{-}2x\}" \},}
\NormalTok{  \{ op: "Apply y(0) = 1",              eq: "C\_1 = 1" \},}
\NormalTok{  \{ op: "Differentiate; apply y\textquotesingle{}(0)=0", eq: "y\textquotesingle{}(0) = C\_2 {-} 2C\_1 = 0 \textbackslash{}\textbackslash{}Rightarrow C\_2 = 2" \},}
\NormalTok{  \{ op: "Final solution",              eq: "y = (1 + 2x)e\^{}\{{-}2x\}" \}}
\NormalTok{])}
\end{Highlighting}
\end{Shaded}

\begin{center}\rule{0.5\linewidth}{0.5pt}\end{center}

\textbf{3.} Solve \(y'' + 9y = 0\), \(y(0) = 0\), \(y'(0) = 3\).

\begin{Shaded}
\begin{Highlighting}[]
\NormalTok{//| echo: false}
\NormalTok{makeStepperHTML(3, [}
\NormalTok{  \{ op: "Characteristic equation",    eq: "r\^{}2 + 9 = 0" \},}
\NormalTok{  \{ op: "Roots",                       eq: "r = \textbackslash{}\textbackslash{}pm 3i", note: "Pure imaginary roots: undamped simple harmonic motion at ω₀ = 3." \},}
\NormalTok{  \{ op: "General solution",            eq: "y = C\_1 \textbackslash{}\textbackslash{}cos(3x) + C\_2 \textbackslash{}\textbackslash{}sin(3x)" \},}
\NormalTok{  \{ op: "Apply y(0) = 0",              eq: "C\_1 = 0" \},}
\NormalTok{  \{ op: "Apply y\textquotesingle{}(0) = 3",             eq: "3C\_2 = 3 \textbackslash{}\textbackslash{}Rightarrow C\_2 = 1" \},}
\NormalTok{  \{ op: "Final solution",              eq: "y = \textbackslash{}\textbackslash{}sin(3x)" \}}
\NormalTok{])}
\end{Highlighting}
\end{Shaded}

\begin{center}\rule{0.5\linewidth}{0.5pt}\end{center}

\textbf{4.} Find a particular solution of \(y'' + 3y' + 2y = e^{-3x}\).

\begin{Shaded}
\begin{Highlighting}[]
\NormalTok{//| echo: false}
\NormalTok{makeStepperHTML(4, [}
\NormalTok{  \{ op: "Homogeneous roots",           eq: "r\^{}2+3r+2=(r+1)(r+2)=0 \textbackslash{}\textbackslash{}Rightarrow r={-}1,{-}2" \},}
\NormalTok{  \{ op: "Modification check",          eq: "e\^{}\{{-}3x\} \textbackslash{}\textbackslash{}notin y\_c", note: "Exponent −3 ≠ −1 or −2. Standard guess applies." \},}
\NormalTok{  \{ op: "Guess",                        eq: "y\_p = Ae\^{}\{{-}3x\}" \},}
\NormalTok{  \{ op: "Substitute",                   eq: "9Ae\^{}\{{-}3x\} {-} 9Ae\^{}\{{-}3x\} + 2Ae\^{}\{{-}3x\} = 2Ae\^{}\{{-}3x\}" \},}
\NormalTok{  \{ op: "Match coefficients",           eq: "2A = 1 \textbackslash{}\textbackslash{}Rightarrow A = \textbackslash{}\textbackslash{}tfrac\{1\}\{2\}" \},}
\NormalTok{  \{ op: "Particular solution",          eq: "y\_p = \textbackslash{}\textbackslash{}tfrac\{1\}\{2\}e\^{}\{{-}3x\}" \}}
\NormalTok{])}
\end{Highlighting}
\end{Shaded}

\begin{center}\rule{0.5\linewidth}{0.5pt}\end{center}

\textbf{5.} Solve \(y'' + 4y = \cos(2x)\).

\emph{Note: the driving frequency \(\omega = 2\) equals
\(\omega_0 = \sqrt{4} = 2\) --- this is resonance. The standard guess
fails; the modification rule applies.}

\begin{Shaded}
\begin{Highlighting}[]
\NormalTok{//| echo: false}
\NormalTok{makeStepperHTML(5, [}
\NormalTok{  \{ op: "Characteristic equation",    eq: "r\^{}2 + 4 = 0 \textbackslash{}\textbackslash{}Rightarrow r = \textbackslash{}\textbackslash{}pm 2i" \},}
\NormalTok{  \{ op: "Homogeneous solution",        eq: "y\_c = C\_1\textbackslash{}\textbackslash{}cos(2x) + C\_2\textbackslash{}\textbackslash{}sin(2x)" \},}
\NormalTok{  \{ op: "Standard guess fails — resonance", eq: "\textbackslash{}\textbackslash{}cos(2x) \textbackslash{}\textbackslash{}in y\_c \textbackslash{}\textbackslash{}Rightarrow \textbackslash{}\textbackslash{}text\{multiply by \} x", note: "This is the resonance case: ω = ω₀ = 2." \},}
\NormalTok{  \{ op: "Modified guess",              eq: "y\_p = x(A\textbackslash{}\textbackslash{}cos 2x + B\textbackslash{}\textbackslash{}sin 2x)" \},}
\NormalTok{  \{ op: "Substitute; x{-}terms cancel",  eq: "{-}4A\textbackslash{}\textbackslash{}sin(2x) + 4B\textbackslash{}\textbackslash{}cos(2x) = \textbackslash{}\textbackslash{}cos(2x)" \},}
\NormalTok{  \{ op: "Match coefficients",          eq: "A = 0,\textbackslash{}\textbackslash{}quad B = \textbackslash{}\textbackslash{}tfrac\{1\}\{4\}" \},}
\NormalTok{  \{ op: "Particular solution",         eq: "y\_p = \textbackslash{}\textbackslash{}tfrac\{1\}\{4\}x\textbackslash{}\textbackslash{}sin(2x)", note: "The factor of x causes amplitude to grow linearly with x — the resonance signature." \},}
\NormalTok{  \{ op: "General solution",            eq: "y = C\_1\textbackslash{}\textbackslash{}cos(2x) + C\_2\textbackslash{}\textbackslash{}sin(2x) + \textbackslash{}\textbackslash{}tfrac\{1\}\{4\}x\textbackslash{}\textbackslash{}sin(2x)" \}}
\NormalTok{])}
\end{Highlighting}
\end{Shaded}

\begin{center}\rule{0.5\linewidth}{0.5pt}\end{center}

\textbf{6.} An RLC circuit has \(L = 1\) H, \(R = 2\) Ω,
\(C = \tfrac{1}{3}\) F with no forcing voltage. The governing equation
for charge \(q(t)\) is \(L\ddot{q} + R\dot{q} + q/C = 0\). Find \(q(t)\)
given \(q(0) = 1\) C, \(\dot{q}(0) = 0\). Classify the response.

\begin{Shaded}
\begin{Highlighting}[]
\NormalTok{//| echo: false}
\NormalTok{makeStepperHTML(6, [}
\NormalTok{  \{ op: "ODE with given values",       eq: "\textbackslash{}\textbackslash{}ddot\{q\} + 2\textbackslash{}\textbackslash{}dot\{q\} + 3q = 0", note: "L=1, R=2, 1/C=3." \},}
\NormalTok{  \{ op: "Characteristic equation",    eq: "r\^{}2 + 2r + 3 = 0" \},}
\NormalTok{  \{ op: "Discriminant",                eq: "\textbackslash{}\textbackslash{}Delta = 4 {-} 12 = {-}8 \textless{} 0", note: "Negative discriminant: underdamped oscillation." \},}
\NormalTok{  \{ op: "Roots",                        eq: "r = {-}1 \textbackslash{}\textbackslash{}pm i\textbackslash{}\textbackslash{}sqrt\{2\}" \},}
\NormalTok{  \{ op: "General solution",            eq: "q = e\^{}\{{-}t\}\textbackslash{}\textbackslash{}bigl(C\_1\textbackslash{}\textbackslash{}cos(\textbackslash{}\textbackslash{}sqrt\{2\}\textbackslash{}\textbackslash{},t) + C\_2\textbackslash{}\textbackslash{}sin(\textbackslash{}\textbackslash{}sqrt\{2\}\textbackslash{}\textbackslash{},t)\textbackslash{}\textbackslash{}bigr)" \},}
\NormalTok{  \{ op: "Apply q(0) = 1",              eq: "C\_1 = 1" \},}
\NormalTok{  \{ op: "Apply \textbackslash{}\textbackslash{}dot\{q\}(0) = 0",       eq: "{-}C\_1 + \textbackslash{}\textbackslash{}sqrt\{2\}\textbackslash{}\textbackslash{},C\_2 = 0 \textbackslash{}\textbackslash{}Rightarrow C\_2 = \textbackslash{}\textbackslash{}tfrac\{1\}\{\textbackslash{}\textbackslash{}sqrt\{2\}\}" \},}
\NormalTok{  \{ op: "Final solution",              eq: "q(t) = e\^{}\{{-}t\}\textbackslash{}\textbackslash{}!\textbackslash{}\textbackslash{}left(\textbackslash{}\textbackslash{}cos(\textbackslash{}\textbackslash{}sqrt\{2\}\textbackslash{}\textbackslash{},t) + \textbackslash{}\textbackslash{}tfrac\{1\}\{\textbackslash{}\textbackslash{}sqrt\{2\}\}\textbackslash{}\textbackslash{}sin(\textbackslash{}\textbackslash{}sqrt\{2\}\textbackslash{}\textbackslash{},t)\textbackslash{}\textbackslash{}right)", note: "Underdamped: decaying oscillation. The circuit rings at ωd = √2 rad/s and settles to q = 0." \}}
\NormalTok{])}
\end{Highlighting}
\end{Shaded}

\chapter{Systems and phase plane}\label{systems-and-phase-plane}

The geometry of dynamical systems

\hfill\break

In a wildlife reserve, park wardens track two populations: rabbits and
foxes. Neither changes independently. When rabbits are plentiful, foxes
thrive and multiply. More foxes means more predation, so rabbit numbers
fall. With fewer rabbits, foxes starve and decline. With fewer foxes,
rabbit numbers recover. The cycle repeats.

Write down a rate equation for rabbits alone and you immediately run
into trouble: the rate of change of the rabbit population depends on the
fox population. You cannot separate them. The two populations are a
single coupled system, and they need to be modelled together.

This chapter develops the language and tools for that. The central
object is the \textbf{phase plane} --- a picture where every point
represents a state of the system (both populations at once), and every
solution traces a curve through that space. The geometry of those curves
--- do they spiral inward, fly outward, or loop forever? --- is
controlled by the eigenvalues of a matrix. That connection between
linear algebra and the long-run behaviour of dynamical systems is one of
the most useful in applied mathematics.

\begin{center}\rule{0.5\linewidth}{0.5pt}\end{center}

\section{Systems of first-order ODEs}\label{systems-of-first-order-odes}

\subsection{From a single ODE to a
system}\label{from-a-single-ode-to-a-system}

A second-order ODE contains a second derivative. It can always be
written as two coupled first-order equations by introducing an auxiliary
variable.

Take the damped oscillator:

\[y'' + py' + qy = 0\]

Set \(x_1 = y\) (the displacement) and \(x_2 = y'\) (the velocity).
Then:

\[x_1' = y' = x_2\] \[x_2' = y'' = -qx_1 - px_2\]

This is a \textbf{system of two first-order ODEs}. The original
second-order equation is equivalent to it. Nothing has been lost or
gained --- this is just a change of representation, chosen because
first-order systems are easier to analyse systematically.

\subsection{Matrix form}\label{matrix-form}

Collect the equations above:

\[\begin{pmatrix} x_1' \\ x_2' \end{pmatrix} = \begin{pmatrix} 0 & 1 \\ -q & -p \end{pmatrix} \begin{pmatrix} x_1 \\ x_2 \end{pmatrix}\]

In compact notation, writing \(\mathbf{x} = (x_1, x_2)^T\) for the
\textbf{state vector}:

\[\mathbf{x}' = A\mathbf{x}, \qquad A = \begin{pmatrix} 0 & 1 \\ -q & -p \end{pmatrix}\]

The state vector \(\mathbf{x}(t)\) records everything about the system
at time \(t\): both the displacement and the velocity. Knowing
\(\mathbf{x}(t_0)\) at any one time, together with the equation
\(\mathbf{x}' = A\mathbf{x}\), determines the entire future trajectory.
This is the central payoff of the state-vector formulation: the full
dynamics live in a single matrix equation.

\subsection{More general systems}\label{more-general-systems}

Not every system comes from a single second-order ODE. A pair of coupled
first-order equations can arise directly from modelling:

\[\begin{aligned}
x' &= f(x, y) \\
y' &= g(x, y)
\end{aligned}\]

When \(f\) and \(g\) are linear functions of \(x\) and \(y\), this is a
\textbf{linear system} \(\mathbf{x}' = A\mathbf{x}\). When they are not
linear, the analysis is more involved --- but the linear case gives the
essential tools, and nonlinear systems are treated by linearising around
equilibria (Section 5).

\begin{tcolorbox}[enhanced jigsaw, title=\textcolor{quarto-callout-note-color}{\faInfo}\hspace{0.5em}{Why systems?}, opacityback=0, colback=white, breakable, leftrule=.75mm, arc=.35mm, toptitle=1mm, left=2mm, colframe=quarto-callout-note-color-frame, coltitle=black, titlerule=0mm, colbacktitle=quarto-callout-note-color!10!white, opacitybacktitle=0.6, toprule=.15mm, bottomtitle=1mm, bottomrule=.15mm, rightrule=.15mm]

Real engineering systems are rarely one-dimensional. A coupled spring
has two masses, two positions, two velocities. An RLC circuit has both
charge and current as state variables. A rotating shaft has angular
displacement and angular momentum. The system
\(\mathbf{x}' = A\mathbf{x}\) handles all of these with the same method.

\end{tcolorbox}

\begin{center}\rule{0.5\linewidth}{0.5pt}\end{center}

\section{The phase plane}\label{the-phase-plane}

\subsection{Trajectories as curves}\label{trajectories-as-curves}

Consider any two-dimensional system:

\[x' = f(x, y), \qquad y' = g(x, y)\]

A \textbf{solution} is a pair of functions \(x(t)\), \(y(t)\) satisfying
both equations simultaneously. We can plot this solution as a parametric
curve in the \((x, y)\)-plane --- called the \textbf{phase plane} ---
with \(t\) as the parameter. As \(t\) advances, the point
\((x(t), y(t))\) moves through the plane, tracing a \textbf{trajectory}
or \textbf{orbit}.

The phase plane collapses time: instead of watching \(x\) and \(y\) vary
separately against \(t\), we watch them vary against each other. This
makes the geometry visible. A spiral trajectory in the phase plane means
oscillation with decay. A closed loop means sustained oscillation. A
trajectory heading straight for the origin means exponential decay.

\subsection{Direction fields}\label{direction-fields}

At every point \((x, y)\), the system specifies a velocity vector:

\[\mathbf{v} = \begin{pmatrix} f(x,y) \\ g(x,y) \end{pmatrix}\]

Drawing a short arrow at each point --- pointing in the direction
\((f, g)\) --- gives a \textbf{direction field} (or \textbf{vector
field}) for the system. A trajectory is a curve that is everywhere
tangent to these arrows.

Before computing any solutions, the direction field reveals the topology
of the phase plane: where trajectories converge, where they diverge,
what the long-run behaviour looks like. For engineering applications,
this geometric picture is often more informative than an explicit
formula.

\subsection{Equilibrium points}\label{equilibrium-points}

An \textbf{equilibrium} (also called a \textbf{critical point} or
\textbf{fixed point}) is a point \(\mathbf{x}^*\) where
\(\mathbf{x}' = \mathbf{0}\), i.e., \(f(x^*, y^*) = 0\) and
\(g(x^*, y^*) = 0\). At an equilibrium, the system is stationary --- if
you start there, you stay there.

For a linear system \(\mathbf{x}' = A\mathbf{x}\), the origin
\(\mathbf{x}^* = \mathbf{0}\) is always an equilibrium. If \(A\) is
invertible, it is the only one.

The central question of phase plane analysis is: \textbf{what do
trajectories do near an equilibrium?} Do they approach it, move away
from it, or orbit around it? The answer comes from the eigenvalues of
\(A\).

\begin{center}\rule{0.5\linewidth}{0.5pt}\end{center}

\section{Linear systems and the eigenvalue
method}\label{linear-systems-and-the-eigenvalue-method}

\subsection{The trial solution}\label{the-trial-solution}

For the scalar equation \(x' = ax\), the solution is \(x(t) = Ce^{at}\).
For the system \(\mathbf{x}' = A\mathbf{x}\), try the analogous form:

\[\mathbf{x}(t) = e^{\lambda t} \mathbf{v}\]

where \(\lambda\) is a scalar and \(\mathbf{v}\) is a (constant) vector.
Substituting into \(\mathbf{x}' = A\mathbf{x}\):

\[\lambda e^{\lambda t} \mathbf{v} = A e^{\lambda t} \mathbf{v}\]
\[A\mathbf{v} = \lambda \mathbf{v}\]

This is precisely the \textbf{eigenvalue equation}. The trial solution
\(\mathbf{x}(t) = e^{\lambda t}\mathbf{v}\) solves the system if and
only if \(\lambda\) is an eigenvalue of \(A\) with eigenvector
\(\mathbf{v}\).

\subsection{General solution}\label{general-solution}

A \(2 \times 2\) matrix generically has two eigenvalues
\(\lambda_1, \lambda_2\) with eigenvectors
\(\mathbf{v}_1, \mathbf{v}_2\). By superposition (the equation is
linear), the general solution is:

\[\mathbf{x}(t) = c_1 e^{\lambda_1 t} \mathbf{v}_1 + c_2 e^{\lambda_2 t} \mathbf{v}_2\]

The constants \(c_1, c_2\) are fixed by initial conditions. Each term is
a trajectory along the eigenvector direction, growing or decaying at its
own rate.

\subsection{Classification of
equilibria}\label{classification-of-equilibria}

The behaviour near the origin depends entirely on the eigenvalues. There
are four principal cases.

\begin{center}\rule{0.5\linewidth}{0.5pt}\end{center}

\textbf{Case 1: Real eigenvalues, same sign --- Node}

\textbf{Example:} \(A = \begin{pmatrix} -1 & 0 \\ 0 & -3 \end{pmatrix}\)

Eigenvalues: \(\lambda_1 = -1\), \(\lambda_2 = -3\) (both negative).
Eigenvectors: \(\mathbf{v}_1 = (1,0)^T\), \(\mathbf{v}_2 = (0,1)^T\).

General solution:

\[\mathbf{x}(t) = c_1 e^{-t} \begin{pmatrix}1\\0\end{pmatrix} + c_2 e^{-3t} \begin{pmatrix}0\\1\end{pmatrix}\]

Both terms decay to zero. Every trajectory approaches the origin as
\(t \to \infty\). This is a \textbf{stable node} (also called a
\textbf{stable improper node}). Eigenvalues of the same sign but
positive give an \textbf{unstable node} --- trajectories fly away from
the origin.

The phase portrait: all trajectories converge to the origin, tangent to
the slow eigendirection (\(\lambda_1 = -1\)) as \(t \to \infty\).

\begin{center}\rule{0.5\linewidth}{0.5pt}\end{center}

\textbf{Case 2: Real eigenvalues, opposite signs --- Saddle}

\textbf{Example:} \(A = \begin{pmatrix} 1 & 0 \\ 0 & -2 \end{pmatrix}\)

Eigenvalues: \(\lambda_1 = 1\), \(\lambda_2 = -2\). General solution:

\[\mathbf{x}(t) = c_1 e^{t} \begin{pmatrix}1\\0\end{pmatrix} + c_2 e^{-2t} \begin{pmatrix}0\\1\end{pmatrix}\]

The first term grows without bound; the second decays. Unless
\(c_1 = 0\) exactly, every trajectory eventually flies away from the
origin along the unstable eigendirection. The origin is a \textbf{saddle
point} --- unstable.

The \textbf{stable manifold} (the set of trajectories that do approach
the origin) consists of the single line along
\(\mathbf{v}_2 = (0,1)^T\). The \textbf{unstable manifold} consists of
the line along \(\mathbf{v}_1 = (1,0)^T\). These two lines are the
separatrices of the saddle.

\begin{center}\rule{0.5\linewidth}{0.5pt}\end{center}

\textbf{Case 3: Complex eigenvalues --- Spiral or Centre}

When \(A\) has a pair of complex conjugate eigenvalues
\(\lambda = \alpha \pm i\beta\) (with \(\beta \neq 0\)), the solution
involves \(e^{\alpha t}(\cos\beta t, \sin\beta t)\) --- oscillation
modulated by \(e^{\alpha t}\).

\begin{itemize}
\tightlist
\item
  \(\alpha < 0\): oscillations decay to zero. \textbf{Stable spiral}
  (also: stable focus).
\item
  \(\alpha > 0\): oscillations grow. \textbf{Unstable spiral}.
\item
  \(\alpha = 0\): pure oscillation, no growth or decay. \textbf{Centre}.
\end{itemize}

\textbf{Extracting real solutions from a complex eigenvector.} When the
eigenvalue is \(\alpha + i\beta\) with eigenvector \(\mathbf{v}\), the
complex exponential \(e^{(\alpha+i\beta)t}\mathbf{v}\) is a
(complex-valued) solution. Take its real and imaginary parts to get two
linearly independent real solutions. For example, with eigenvalue
\(-1+2i\) and eigenvector \(\mathbf{v} = (1, i)^T\):

\[e^{(-1+2i)t}\begin{pmatrix}1\\i\end{pmatrix} = e^{-t}\begin{pmatrix}\cos 2t + i\sin 2t\\ i\cos 2t - \sin 2t\end{pmatrix}\]

Real part: \(e^{-t}(\cos 2t,\; -\sin 2t)^T\). Imaginary part:
\(e^{-t}(\sin 2t,\; \cos 2t)^T\). These two real vectors are the two
linearly independent solutions --- the template used in Exercise 3
below.

\textbf{Example (stable spiral):}
\(A = \begin{pmatrix} -1 & 2 \\ -2 & -1 \end{pmatrix}\)

Characteristic polynomial:
\(\det(A - \lambda I) = (\lambda+1)^2 + 4 = 0\), giving
\(\lambda = -1 \pm 2i\).

Since \(\alpha = -1 < 0\), all trajectories spiral inward toward the
origin.

\textbf{Example (centre):}
\(A = \begin{pmatrix} 0 & 1 \\ -4 & 0 \end{pmatrix}\)

Characteristic polynomial: \(\lambda^2 + 4 = 0\), giving
\(\lambda = \pm 2i\). Pure imaginary eigenvalues. All trajectories are
closed ellipses surrounding the origin --- periodic orbits. This is a
\textbf{centre}.

\begin{tcolorbox}[enhanced jigsaw, title=\textcolor{quarto-callout-note-color}{\faInfo}\hspace{0.5em}{Summary: eigenvalue classification}, opacityback=0, colback=white, breakable, leftrule=.75mm, arc=.35mm, toptitle=1mm, left=2mm, colframe=quarto-callout-note-color-frame, coltitle=black, titlerule=0mm, colbacktitle=quarto-callout-note-color!10!white, opacitybacktitle=0.6, toprule=.15mm, bottomtitle=1mm, bottomrule=.15mm, rightrule=.15mm]

\begin{longtable}[]{@{}lll@{}}
\toprule\noalign{}
Eigenvalues & Equilibrium type & Stable? \\
\midrule\noalign{}
\endhead
\bottomrule\noalign{}
\endlastfoot
Real, both negative & Stable node & Yes \\
Real, both positive & Unstable node & No \\
Real, opposite signs & Saddle & No \\
Complex, \(\alpha < 0\) & Stable spiral & Yes \\
Complex, \(\alpha > 0\) & Unstable spiral & No \\
Purely imaginary & Centre & Neutral \\
\end{longtable}

The trace \(\tau = \lambda_1 + \lambda_2 = \text{tr}(A)\) and
determinant \(\delta = \lambda_1\lambda_2 = \det(A)\) give a shortcut:
if \(\delta < 0\) it is a saddle; if \(\delta > 0\) and \(\tau < 0\) it
is a stable node or stable spiral; if \(\delta > 0\) and \(\tau > 0\) it
is unstable.

\end{tcolorbox}

\begin{Shaded}
\begin{Highlighting}[]
\NormalTok{//| echo: false}
\NormalTok{viewof ppParams = Inputs.form(\{}
\NormalTok{  ppA11: Inputs.range([{-}3, 3], \{step: 0.1, value: {-}1, label: "a₁₁"\}),}
\NormalTok{  ppA12: Inputs.range([{-}3, 3], \{step: 0.1, value: 1,  label: "a₁₂"\}),}
\NormalTok{  ppA21: Inputs.range([{-}3, 3], \{step: 0.1, value: {-}2, label: "a₂₁"\}),}
\NormalTok{  ppA22: Inputs.range([{-}3, 3], \{step: 0.1, value: {-}1, label: "a₂₂"\})}
\NormalTok{\})}
\end{Highlighting}
\end{Shaded}

\begin{Shaded}
\begin{Highlighting}[]
\NormalTok{//| echo: false}
\NormalTok{\{}
\NormalTok{  const d3 = await require("d3@7");}

\NormalTok{  const a11 = ppParams.ppA11;}
\NormalTok{  const a12 = ppParams.ppA12;}
\NormalTok{  const a21 = ppParams.ppA21;}
\NormalTok{  const a22 = ppParams.ppA22;}

\NormalTok{  // Eigenvalue computation via quadratic formula}
\NormalTok{  const tr  = a11 + a22;}
\NormalTok{  const det = a11 * a22 {-} a12 * a21;}
\NormalTok{  const disc = tr * tr {-} 4 * det;}

\NormalTok{  let lam1, lam2, lamStr1, lamStr2, classification, stableColor;}

\NormalTok{  if (disc \textgreater{}= 0) \{}
\NormalTok{    // Real eigenvalues}
\NormalTok{    lam1 = (tr + Math.sqrt(disc)) / 2;}
\NormalTok{    lam2 = (tr {-} Math.sqrt(disc)) / 2;}
\NormalTok{    lamStr1 = lam1.toFixed(3);}
\NormalTok{    lamStr2 = lam2.toFixed(3);}
\NormalTok{    if (lam1 * lam2 \textless{} 0) \{}
\NormalTok{      classification = "Saddle (unstable)";}
\NormalTok{      stableColor = "saddle";}
\NormalTok{    \} else if (lam1 \textless{} 0 \&\& lam2 \textless{} 0) \{}
\NormalTok{      classification = "Stable node";}
\NormalTok{      stableColor = "stable";}
\NormalTok{    \} else if (lam1 \textgreater{} 0 \&\& lam2 \textgreater{} 0) \{}
\NormalTok{      classification = "Unstable node";}
\NormalTok{      stableColor = "unstable";}
\NormalTok{    \} else \{}
\NormalTok{      // one or both zero}
\NormalTok{      classification = "Non{-}isolated equilibrium";}
\NormalTok{      stableColor = "neutral";}
\NormalTok{    \}}
\NormalTok{  \} else \{}
\NormalTok{    // Complex eigenvalues: α ± βi}
\NormalTok{    const alpha = tr / 2;}
\NormalTok{    const beta  = Math.sqrt({-}disc) / 2;}
\NormalTok{    lam1 = \{ re: alpha, im:  beta \};}
\NormalTok{    lam2 = \{ re: alpha, im: {-}beta \};}
\NormalTok{    const sign = beta \textgreater{}= 0 ? "+" : "{-}";}
\NormalTok{    lamStr1 = \textasciigrave{}$\{alpha.toFixed(3)\} + $\{beta.toFixed(3)\}i\textasciigrave{};}
\NormalTok{    lamStr2 = \textasciigrave{}$\{alpha.toFixed(3)\} − $\{beta.toFixed(3)\}i\textasciigrave{};}
\NormalTok{    if (Math.abs(alpha) \textless{} 1e{-}9) \{}
\NormalTok{      classification = "Centre (neutrally stable)";}
\NormalTok{      stableColor = "neutral";}
\NormalTok{    \} else if (alpha \textless{} 0) \{}
\NormalTok{      classification = "Stable spiral";}
\NormalTok{      stableColor = "stable";}
\NormalTok{    \} else \{}
\NormalTok{      classification = "Unstable spiral";}
\NormalTok{      stableColor = "unstable";}
\NormalTok{    \}}
\NormalTok{  \}}

\NormalTok{  // SVG dimensions}
\NormalTok{  const W = 560, H = 400;}
\NormalTok{  const margin = \{ top: 30, right: 20, bottom: 40, left: 40 \};}
\NormalTok{  const iW = W {-} margin.left {-} margin.right;}
\NormalTok{  const iH = H {-} margin.top  {-} margin.bottom;}

\NormalTok{  const xDom = [{-}3, 3], yDom = [{-}3, 3];}
\NormalTok{  const xSc = d3.scaleLinear().domain(xDom).range([0, iW]);}
\NormalTok{  const ySc = d3.scaleLinear().domain(yDom).range([iH, 0]);}

\NormalTok{  const svg = d3.create("svg")}
\NormalTok{    .attr("width", W).attr("height", H)}
\NormalTok{    .style("font{-}family", "sans{-}serif");}

\NormalTok{  const g = svg.append("g")}
\NormalTok{    .attr("transform", \textasciigrave{}translate($\{margin.left\},$\{margin.top\})\textasciigrave{});}

\NormalTok{  // Background}
\NormalTok{  g.append("rect")}
\NormalTok{    .attr("width", iW).attr("height", iH)}
\NormalTok{    .attr("fill", "\#fafafa").attr("stroke", "\#ccc");}

\NormalTok{  // Axes}
\NormalTok{  g.append("line")}
\NormalTok{    .attr("x1", xSc(0)).attr("x2", xSc(0))}
\NormalTok{    .attr("y1", 0).attr("y2", iH)}
\NormalTok{    .attr("stroke", "\#aaa").attr("stroke{-}width", 1);}
\NormalTok{  g.append("line")}
\NormalTok{    .attr("x1", 0).attr("x2", iW)}
\NormalTok{    .attr("y1", ySc(0)).attr("y2", ySc(0))}
\NormalTok{    .attr("stroke", "\#aaa").attr("stroke{-}width", 1);}

\NormalTok{  // Axis labels}
\NormalTok{  g.append("text").attr("x", iW + 5).attr("y", ySc(0) + 4).text("x₁").attr("font{-}size", 11).attr("fill", "\#555");}
\NormalTok{  g.append("text").attr("x", xSc(0) + 4).attr("y", {-}8).text("x₂").attr("font{-}size", 11).attr("fill", "\#555");}

\NormalTok{  // Direction field: 15×15 grid}
\NormalTok{  const gridN = 15;}
\NormalTok{  const arrowLen = (iW / gridN) * 0.38;}
\NormalTok{  for (let i = 0; i \textless{}= gridN; i++) \{}
\NormalTok{    for (let j = 0; j \textless{}= gridN; j++) \{}
\NormalTok{      const x = xDom[0] + (i / gridN) * (xDom[1] {-} xDom[0]);}
\NormalTok{      const y = yDom[0] + (j / gridN) * (yDom[1] {-} yDom[0]);}
\NormalTok{      const dx = a11 * x + a12 * y;}
\NormalTok{      const dy = a21 * x + a22 * y;}
\NormalTok{      const mag = Math.sqrt(dx * dx + dy * dy) || 1;}
\NormalTok{      const nx = dx / mag, ny = dy / mag;}
\NormalTok{      const cx = xSc(x), cy = ySc(y);}
\NormalTok{      const ex = cx + nx * arrowLen;}
\NormalTok{      const ey = cy {-} ny * arrowLen;}
\NormalTok{      g.append("line")}
\NormalTok{        .attr("x1", cx {-} nx * arrowLen * 0.5)}
\NormalTok{        .attr("y1", cy + ny * arrowLen * 0.5)}
\NormalTok{        .attr("x2", ex)}
\NormalTok{        .attr("y2", ey)}
\NormalTok{        .attr("stroke", "\#bbb").attr("stroke{-}width", 0.8)}
\NormalTok{        .attr("marker{-}end", "url(\#arrowhead)");}
\NormalTok{    \}}
\NormalTok{  \}}

\NormalTok{  // Arrowhead marker}
\NormalTok{  svg.append("defs").append("marker")}
\NormalTok{    .attr("id", "arrowhead")}
\NormalTok{    .attr("viewBox", "0 {-}3 6 6")}
\NormalTok{    .attr("refX", 5).attr("refY", 0)}
\NormalTok{    .attr("markerWidth", 4).attr("markerHeight", 4)}
\NormalTok{    .attr("orient", "auto")}
\NormalTok{    .append("path").attr("d", "M0,{-}3L6,0L0,3").attr("fill", "\#bbb");}

\NormalTok{  // Trajectory colours}
\NormalTok{  const trajColor = stableColor === "stable"   ? "\#1e6fcc" :}
\NormalTok{                    stableColor === "unstable"  ? "\#cc3333" :}
\NormalTok{                    stableColor === "saddle"    ? null       : "\#888";}

\NormalTok{  // 8 trajectories from radius 2, evenly spaced angles}
\NormalTok{  const nTraj = 8;}
\NormalTok{  const r0 = 2;}
\NormalTok{  const dt = 0.02, nSteps = 300;}

\NormalTok{  function eulerSteps(x0, y0, forward) \{}
\NormalTok{    const pts = [[x0, y0]];}
\NormalTok{    let x = x0, y = y0;}
\NormalTok{    const sign = forward ? 1 : {-}1;}
\NormalTok{    for (let s = 0; s \textless{} nSteps; s++) \{}
\NormalTok{      const dx = a11 * x + a12 * y;}
\NormalTok{      const dy = a21 * x + a22 * y;}
\NormalTok{      x += sign * dt * dx;}
\NormalTok{      y += sign * dt * dy;}
\NormalTok{      if (Math.abs(x) \textgreater{} 6 || Math.abs(y) \textgreater{} 6) break;}
\NormalTok{      pts.push([x, y]);}
\NormalTok{    \}}
\NormalTok{    return pts;}
\NormalTok{  \}}

\NormalTok{  const lineGen = d3.line()}
\NormalTok{    .x(p =\textgreater{} xSc(p[0]))}
\NormalTok{    .y(p =\textgreater{} ySc(p[1]))}
\NormalTok{    .defined(p =\textgreater{} Math.abs(p[0]) \textless{}= 3.2 \&\& Math.abs(p[1]) \textless{}= 3.2);}

\NormalTok{  for (let k = 0; k \textless{} nTraj; k++) \{}
\NormalTok{    const angle = (2 * Math.PI * k) / nTraj;}
\NormalTok{    const x0 = r0 * Math.cos(angle);}
\NormalTok{    const y0 = r0 * Math.sin(angle);}

\NormalTok{    const fwd = eulerSteps(x0, y0, true);}
\NormalTok{    const bwd = eulerSteps(x0, y0, false).slice(1).reverse();}
\NormalTok{    const pts = [...bwd, ...fwd];}

\NormalTok{    // Colour saddle trajectories by quadrant (proxy)}
\NormalTok{    let color;}
\NormalTok{    if (stableColor === "saddle") \{}
\NormalTok{      color = (k \% 2 === 0) ? "\#1e6fcc" : "\#cc3333";}
\NormalTok{    \} else \{}
\NormalTok{      color = trajColor;}
\NormalTok{    \}}

\NormalTok{    g.append("path")}
\NormalTok{      .datum(pts)}
\NormalTok{      .attr("d", lineGen)}
\NormalTok{      .attr("fill", "none")}
\NormalTok{      .attr("stroke", color)}
\NormalTok{      .attr("stroke{-}width", 1.4)}
\NormalTok{      .attr("opacity", 0.75);}
\NormalTok{  \}}

\NormalTok{  // Equilibrium dot at origin}
\NormalTok{  g.append("circle")}
\NormalTok{    .attr("cx", xSc(0)).attr("cy", ySc(0))}
\NormalTok{    .attr("r", 5)}
\NormalTok{    .attr("fill", stableColor === "stable" ? "\#1e6fcc" : stableColor === "unstable" ? "\#cc3333" : "\#888")}
\NormalTok{    .attr("stroke", "\#333").attr("stroke{-}width", 1.5);}

\NormalTok{  // Classification label}
\NormalTok{  const labelColor = stableColor === "stable" ? "\#1e6fcc" : stableColor === "unstable" ? "\#cc3333" : "\#555";}
\NormalTok{  g.append("text")}
\NormalTok{    .attr("x", iW / 2).attr("y", iH + 30)}
\NormalTok{    .attr("text{-}anchor", "middle")}
\NormalTok{    .attr("font{-}size", 13).attr("font{-}weight", "bold")}
\NormalTok{    .attr("fill", labelColor)}
\NormalTok{    .text(\textasciigrave{}$\{classification\}  |  λ₁ = $\{lamStr1\},  λ₂ = $\{lamStr2\}\textasciigrave{});}

\NormalTok{  return svg.node();}
\NormalTok{\}}
\end{Highlighting}
\end{Shaded}

\begin{center}\rule{0.5\linewidth}{0.5pt}\end{center}

\section{Stability}\label{stability}

\subsection{Lyapunov's definition}\label{lyapunovs-definition}

An equilibrium \(\mathbf{x}^*\) is \textbf{stable} (in the sense of
Lyapunov) if trajectories that start close to \(\mathbf{x}^*\) stay
close --- formally: for every \(\varepsilon > 0\) there exists
\(\delta > 0\) such that \(\|\mathbf{x}(0) - \mathbf{x}^*\| < \delta\)
implies \(\|\mathbf{x}(t) - \mathbf{x}^*\| < \varepsilon\) for all
\(t > 0\).

An equilibrium is \textbf{asymptotically stable} if it is stable and,
additionally, trajectories starting near it converge to it:
\(\mathbf{x}(t) \to \mathbf{x}^*\) as \(t \to \infty\).

An equilibrium is \textbf{unstable} if it fails the stability condition
--- there exist trajectories starting arbitrarily close that eventually
leave a fixed neighbourhood.

For linear systems, the connection to eigenvalues is exact:

\begin{itemize}
\tightlist
\item
  All eigenvalues have negative real part \(\Rightarrow\) asymptotically
  stable.
\item
  Any eigenvalue with positive real part \(\Rightarrow\) unstable.
\item
  Purely imaginary eigenvalues \(\Rightarrow\) stable (but not
  asymptotically stable) for the linear system; the nonlinear case
  requires more care.
\end{itemize}

\subsection{Engineering implications}\label{engineering-implications}

The stability of an equilibrium is the primary question in control
engineering. A feedback controller is designed to make the closed-loop
eigenvalues lie in the left half of the complex plane
(\(\text{Re}(\lambda) < 0\)). How far left they sit determines how
quickly the system returns to equilibrium after a disturbance --- the
damping rate is \(|\text{Re}(\lambda)|\), and the oscillation frequency
is \(|\text{Im}(\lambda)|\).

A saddle equilibrium in a control loop is dangerous: the system will
diverge for almost any initial condition, with only the stable manifold
--- a set of measure zero --- as the exception.

\begin{center}\rule{0.5\linewidth}{0.5pt}\end{center}

\section{Nonlinear systems:
linearisation}\label{nonlinear-systems-linearisation}

\subsection{Why linearise?}\label{why-linearise}

For a nonlinear system \(\mathbf{x}' = \mathbf{f}(\mathbf{x})\), we
generally cannot write down explicit solutions. But the behaviour near
an equilibrium \(\mathbf{x}^*\) can still be understood by approximating
\(\mathbf{f}\) with its best linear approximation there --- a Taylor
expansion.

\subsection{The Jacobian}\label{the-jacobian}

Let \(\mathbf{x}^*\) satisfy \(\mathbf{f}(\mathbf{x}^*) = \mathbf{0}\).
Write \(\mathbf{x}(t) = \mathbf{x}^* + \mathbf{u}(t)\) where
\(\mathbf{u}\) is a small perturbation. Expanding \(\mathbf{f}\) to
first order:

\[\mathbf{x}' = \mathbf{f}(\mathbf{x}^* + \mathbf{u}) \approx \mathbf{f}(\mathbf{x}^*) + J(\mathbf{x}^*)\,\mathbf{u} = J(\mathbf{x}^*)\,\mathbf{u}\]

since \(\mathbf{f}(\mathbf{x}^*) = \mathbf{0}\). The matrix \(J\) is the
\textbf{Jacobian} of \(\mathbf{f}\):

\[J = \begin{pmatrix} \partial f_1/\partial x_1 & \partial f_1/\partial x_2 \\ \partial f_2/\partial x_1 & \partial f_2/\partial x_2 \end{pmatrix}\]

evaluated at \(\mathbf{x}^*\). So the perturbation \(\mathbf{u}\)
satisfies the linear system \(\mathbf{u}' = J\mathbf{u}\). The
eigenvalues of \(J\) classify the equilibrium, exactly as before.

This works reliably when the eigenvalues of \(J\) are not purely
imaginary. If \(J\) has purely imaginary eigenvalues, the linearisation
is inconclusive about the nonlinear behaviour --- the equilibrium could
be a centre, a spiral, or even something more exotic.

\subsection{The pendulum}\label{the-pendulum}

The undamped pendulum satisfies:

\[\theta'' + \sin\theta = 0\]

(with appropriate normalisation of length and gravity). Converting to a
system: \(x_1 = \theta\), \(x_2 = \theta'\):

\[x_1' = x_2, \qquad x_2' = -\sin x_1\]

\textbf{Equilibria:} Set \(x_2 = 0\) and \(-\sin x_1 = 0\). Equilibria
are at \((n\pi, 0)\) for integer \(n\).

\begin{itemize}
\tightlist
\item
  \((0, 0)\): the pendulum hanging straight down.
\item
  \((\pi, 0)\): the pendulum balanced straight up.
\end{itemize}

\textbf{Jacobian:}

\[J = \begin{pmatrix} 0 & 1 \\ -\cos x_1 & 0 \end{pmatrix}\]

\textbf{At \((0,0)\):} \(J = \begin{pmatrix}0&1\\-1&0\end{pmatrix}\).
Eigenvalues: \(\lambda = \pm i\). The linearisation gives a centre ---
consistent with the familiar small-angle oscillation. (The full
nonlinear pendulum also has closed orbits near this point, which
requires a more careful argument to confirm.)

\textbf{At \((\pi,0)\):} \(J = \begin{pmatrix}0&1\\1&0\end{pmatrix}\).
Eigenvalues: \(\lambda = \pm 1\). A saddle --- the upright position is
unstable, as expected. The tiniest perturbation sends the pendulum
swinging away.

The phase portrait of the pendulum has closed orbits (libration) near
\((0,0)\) and open orbits (rotation) beyond a separatrix that passes
through the saddle point \((\pm\pi, 0)\). The separatrix is the boundary
between ``swings back'' and ``goes over the top''.

\begin{center}\rule{0.5\linewidth}{0.5pt}\end{center}

\section{Lotka-Volterra
predator-prey}\label{lotka-volterra-predator-prey}

\subsection{The equations}\label{the-equations}

The Lotka-Volterra model for a prey population \(x\) (rabbits) and
predator population \(y\) (foxes) is:

\[\begin{aligned}
x' &= ax - bxy \\
y' &= -dy + cxy
\end{aligned}\]

All parameters \(a, b, c, d > 0\).

\begin{itemize}
\tightlist
\item
  \(ax\): prey grow exponentially in the absence of predators.
\item
  \(-bxy\): prey are removed at a rate proportional to encounters with
  predators.
\item
  \(-dy\): predators decline exponentially in the absence of prey.
\item
  \(cxy\): predators grow from encounters with prey.
\end{itemize}

\subsection{Equilibria}\label{equilibria}

Set \(x' = 0\) and \(y' = 0\):

\[x(a - by) = 0 \quad \text{and} \quad y(-d + cx) = 0\]

This gives two equilibria:

\begin{enumerate}
\def\labelenumi{\arabic{enumi}.}
\tightlist
\item
  \((0, 0)\): both populations extinct.
\item
  \((d/c,\; a/b)\): coexistence equilibrium --- both populations
  positive.
\end{enumerate}

\subsection{Classifying via the
Jacobian}\label{classifying-via-the-jacobian}

\[\mathbf{f}(x,y) = \begin{pmatrix} ax - bxy \\ -dy + cxy \end{pmatrix}, \qquad J = \begin{pmatrix} a - by & -bx \\ cy & -d + cx \end{pmatrix}\]

\textbf{At \((0,0)\):}

\[J(0,0) = \begin{pmatrix} a & 0 \\ 0 & -d \end{pmatrix}\]

Eigenvalues: \(\lambda_1 = a > 0\), \(\lambda_2 = -d < 0\). A saddle.
The extinction equilibrium is unstable --- any positive initial
population will move away from it.

\textbf{At \((d/c,\; a/b)\):}

Substituting \(x^* = d/c\), \(y^* = a/b\):

\[J(x^*, y^*) = \begin{pmatrix} 0 & -bd/c \\ ca/b & 0 \end{pmatrix}\]

Characteristic polynomial: \(\lambda^2 + ad = 0\), giving
\(\lambda = \pm i\sqrt{ad}\).

Purely imaginary eigenvalues. The linearisation gives a centre. (For the
full Lotka-Volterra system it can be shown, using a conserved quantity,
that the nonlinear equilibrium is also a true centre --- the populations
cycle indefinitely around the coexistence point.)

\subsection{Physical interpretation}\label{physical-interpretation}

The coexistence equilibrium \((d/c, a/b)\) is a neutrally stable centre.
The populations do not converge to these values; instead, they cycle
around them. The amplitude of the cycle depends on the initial
conditions. Start with many rabbits and few foxes: the fox population
grows, rabbit numbers fall, fox numbers eventually fall too, then
rabbits recover. The cycle is closed.

This explains the wildlife warden's observation: the populations
oscillate, neither dying out nor converging to a constant value, with
the phase plane orbit a closed curve around the coexistence point.

\begin{Shaded}
\begin{Highlighting}[]
\NormalTok{//| echo: false}
\NormalTok{viewof lvParams = Inputs.form(\{}
\NormalTok{  lvA:  Inputs.range([0.1, 2],    \{step: 0.05, value: 0.8, label: "a (prey growth rate)"\}),}
\NormalTok{  lvB:  Inputs.range([0.01, 0.5], \{step: 0.01, value: 0.1, label: "b (predation rate)"\}),}
\NormalTok{  lvC:  Inputs.range([0.01, 0.5], \{step: 0.01, value: 0.1, label: "c (predator efficiency)"\}),}
\NormalTok{  lvD:  Inputs.range([0.1, 2],    \{step: 0.05, value: 0.5, label: "d (predator death rate)"\}),}
\NormalTok{  lvX0: Inputs.range([1, 20],     \{step: 0.5,  value: 8,   label: "x₀ (initial prey)"\}),}
\NormalTok{  lvY0: Inputs.range([1, 20],     \{step: 0.5,  value: 4,   label: "y₀ (initial predators)"\})}
\NormalTok{\})}
\end{Highlighting}
\end{Shaded}

\begin{Shaded}
\begin{Highlighting}[]
\NormalTok{//| echo: false}
\NormalTok{\{}
\NormalTok{  const d3 = await require("d3@7");}

\NormalTok{  const a  = lvParams.lvA;}
\NormalTok{  const b  = lvParams.lvB;}
\NormalTok{  const c  = lvParams.lvC;}
\NormalTok{  const d  = lvParams.lvD;}
\NormalTok{  const x0 = lvParams.lvX0;}
\NormalTok{  const y0 = lvParams.lvY0;}

\NormalTok{  // RK4 integration of Lotka{-}Volterra}
\NormalTok{  const dt = 0.05, nSteps = 2000;}
\NormalTok{  const tMax = 100;}
\NormalTok{  const storeEvery = 1;  // store every step}

\NormalTok{  const xs = [x0], ys = [y0], ts = [0];}
\NormalTok{  let cx = x0, cy = y0, ct = 0;}

\NormalTok{  for (let i = 0; i \textless{} nSteps; i++) \{}
\NormalTok{    // RK4}
\NormalTok{    const k1x = a * cx       {-} b * cx * cy;}
\NormalTok{    const k1y = c * cx * cy  {-} d * cy;}

\NormalTok{    const mx2 = cx + 0.5 * dt * k1x;}
\NormalTok{    const my2 = cy + 0.5 * dt * k1y;}
\NormalTok{    const k2x = a * mx2      {-} b * mx2 * my2;}
\NormalTok{    const k2y = c * mx2*my2  {-} d * my2;}

\NormalTok{    const mx3 = cx + 0.5 * dt * k2x;}
\NormalTok{    const my3 = cy + 0.5 * dt * k2y;}
\NormalTok{    const k3x = a * mx3      {-} b * mx3 * my3;}
\NormalTok{    const k3y = c * mx3*my3  {-} d * my3;}

\NormalTok{    const mx4 = cx + dt * k3x;}
\NormalTok{    const my4 = cy + dt * k3y;}
\NormalTok{    const k4x = a * mx4      {-} b * mx4 * my4;}
\NormalTok{    const k4y = c * mx4*my4  {-} d * my4;}

\NormalTok{    cx = cx + (dt / 6) * (k1x + 2*k2x + 2*k3x + k4x);}
\NormalTok{    cy = cy + (dt / 6) * (k1y + 2*k2y + 2*k3y + k4y);}
\NormalTok{    ct += dt;}

\NormalTok{    // clamp to avoid blow{-}up}
\NormalTok{    if (cx \textless{} 0) cx = 0;}
\NormalTok{    if (cy \textless{} 0) cy = 0;}

\NormalTok{    xs.push(cx);}
\NormalTok{    ys.push(cy);}
\NormalTok{    ts.push(ct);}

\NormalTok{    if (ct \textgreater{}= tMax) break;}
\NormalTok{  \}}

\NormalTok{  // Equilibria}
\NormalTok{  const eqX = d / c;}
\NormalTok{  const eqY = a / b;}

\NormalTok{  // Layout: two panels side by side, total 560px wide}
\NormalTok{  const totalW = 560, panelW = 270, H = 300;}
\NormalTok{  const margin = \{ top: 24, right: 18, bottom: 40, left: 44 \};}
\NormalTok{  const iW = panelW {-} margin.left {-} margin.right;}
\NormalTok{  const iH = H      {-} margin.top  {-} margin.bottom;}

\NormalTok{  const svg = d3.create("svg")}
\NormalTok{    .attr("width", totalW).attr("height", H)}
\NormalTok{    .style("font{-}family", "sans{-}serif");}

\NormalTok{  // ── Left panel: Phase portrait ──────────────────────────────────────────}
\NormalTok{  const gL = svg.append("g")}
\NormalTok{    .attr("transform", \textasciigrave{}translate($\{margin.left\}, $\{margin.top\})\textasciigrave{});}

\NormalTok{  const xMax = Math.max(...xs) * 1.15 || 20;}
\NormalTok{  const yMax = Math.max(...ys) * 1.15 || 20;}

\NormalTok{  const xSc = d3.scaleLinear().domain([0, xMax]).range([0, iW]);}
\NormalTok{  const ySc = d3.scaleLinear().domain([0, yMax]).range([iH, 0]);}

\NormalTok{  // Background}
\NormalTok{  gL.append("rect").attr("width", iW).attr("height", iH)}
\NormalTok{    .attr("fill", "\#fafafa").attr("stroke", "\#ccc");}

\NormalTok{  // Trajectory}
\NormalTok{  const phaseLine = d3.line()}
\NormalTok{    .x((\_, i) =\textgreater{} xSc(xs[i]))}
\NormalTok{    .y((\_, i) =\textgreater{} ySc(ys[i]));}

\NormalTok{  gL.append("path")}
\NormalTok{    .datum(xs)}
\NormalTok{    .attr("d", phaseLine)}
\NormalTok{    .attr("fill", "none")}
\NormalTok{    .attr("stroke", "\#5b8dd9")}
\NormalTok{    .attr("stroke{-}width", 1.6)}
\NormalTok{    .attr("opacity", 0.85);}

\NormalTok{  // Equilibrium (0, 0) — open circle}
\NormalTok{  gL.append("circle")}
\NormalTok{    .attr("cx", xSc(0)).attr("cy", ySc(0))}
\NormalTok{    .attr("r", 5).attr("fill", "white")}
\NormalTok{    .attr("stroke", "\#555").attr("stroke{-}width", 1.5);}

\NormalTok{  // Equilibrium (d/c, a/b) — filled circle}
\NormalTok{  if (eqX \textless{}= xMax \&\& eqY \textless{}= yMax) \{}
\NormalTok{    gL.append("circle")}
\NormalTok{      .attr("cx", xSc(eqX)).attr("cy", ySc(eqY))}
\NormalTok{      .attr("r", 5).attr("fill", "\#e07b2a")}
\NormalTok{      .attr("stroke", "\#333").attr("stroke{-}width", 1.5);}
\NormalTok{  \}}

\NormalTok{  // Starting point}
\NormalTok{  gL.append("circle")}
\NormalTok{    .attr("cx", xSc(x0)).attr("cy", ySc(y0))}
\NormalTok{    .attr("r", 4).attr("fill", "\#333");}

\NormalTok{  // Axes}
\NormalTok{  const xAxisL = d3.axisBottom(xSc).ticks(5);}
\NormalTok{  const yAxisL = d3.axisLeft(ySc).ticks(5);}
\NormalTok{  gL.append("g").attr("transform", \textasciigrave{}translate(0,$\{iH\})\textasciigrave{}).call(xAxisL);}
\NormalTok{  gL.append("g").call(yAxisL);}

\NormalTok{  // Axis labels}
\NormalTok{  gL.append("text").attr("x", iW / 2).attr("y", iH + 32)}
\NormalTok{    .attr("text{-}anchor", "middle").attr("font{-}size", 11).attr("fill", "\#444")}
\NormalTok{    .text("x (prey)");}
\NormalTok{  gL.append("text")}
\NormalTok{    .attr("transform", \textasciigrave{}translate({-}32,$\{iH/2\}) rotate({-}90)\textasciigrave{})}
\NormalTok{    .attr("text{-}anchor", "middle").attr("font{-}size", 11).attr("fill", "\#444")}
\NormalTok{    .text("y (predators)");}
\NormalTok{  gL.append("text").attr("x", iW / 2).attr("y", {-}8)}
\NormalTok{    .attr("text{-}anchor", "middle").attr("font{-}size", 12).attr("font{-}weight", "bold").attr("fill", "\#333")}
\NormalTok{    .text("Phase portrait");}

\NormalTok{  // ── Right panel: Time series ────────────────────────────────────────────}
\NormalTok{  const gR = svg.append("g")}
\NormalTok{    .attr("transform", \textasciigrave{}translate($\{panelW + margin.left\}, $\{margin.top\})\textasciigrave{});}

\NormalTok{  const tEnd = ts[ts.length {-} 1];}
\NormalTok{  const popMax = Math.max(...xs, ...ys) * 1.15 || 20;}

\NormalTok{  const tSc   = d3.scaleLinear().domain([0, tEnd]).range([0, iW]);}
\NormalTok{  const popSc = d3.scaleLinear().domain([0, popMax]).range([iH, 0]);}

\NormalTok{  gR.append("rect").attr("width", iW).attr("height", iH)}
\NormalTok{    .attr("fill", "\#fafafa").attr("stroke", "\#ccc");}

\NormalTok{  // Prey x(t) in green}
\NormalTok{  const xLine = d3.line()}
\NormalTok{    .x((\_, i) =\textgreater{} tSc(ts[i]))}
\NormalTok{    .y((\_, i) =\textgreater{} popSc(xs[i]));}

\NormalTok{  gR.append("path")}
\NormalTok{    .datum(xs)}
\NormalTok{    .attr("d", xLine)}
\NormalTok{    .attr("fill", "none")}
\NormalTok{    .attr("stroke", "\#2e8b57")}
\NormalTok{    .attr("stroke{-}width", 1.6);}

\NormalTok{  // Predator y(t) in orange}
\NormalTok{  const yLine = d3.line()}
\NormalTok{    .x((\_, i) =\textgreater{} tSc(ts[i]))}
\NormalTok{    .y((\_, i) =\textgreater{} popSc(ys[i]));}

\NormalTok{  gR.append("path")}
\NormalTok{    .datum(ys)}
\NormalTok{    .attr("d", yLine)}
\NormalTok{    .attr("fill", "none")}
\NormalTok{    .attr("stroke", "\#e07b2a")}
\NormalTok{    .attr("stroke{-}width", 1.6);}

\NormalTok{  // Axes}
\NormalTok{  const xAxisR = d3.axisBottom(tSc).ticks(5);}
\NormalTok{  const yAxisR = d3.axisLeft(popSc).ticks(5);}
\NormalTok{  gR.append("g").attr("transform", \textasciigrave{}translate(0,$\{iH\})\textasciigrave{}).call(xAxisR);}
\NormalTok{  gR.append("g").call(yAxisR);}

\NormalTok{  // Axis labels}
\NormalTok{  gR.append("text").attr("x", iW / 2).attr("y", iH + 32)}
\NormalTok{    .attr("text{-}anchor", "middle").attr("font{-}size", 11).attr("fill", "\#444")}
\NormalTok{    .text("t");}
\NormalTok{  gR.append("text")}
\NormalTok{    .attr("transform", \textasciigrave{}translate({-}32,$\{iH/2\}) rotate({-}90)\textasciigrave{})}
\NormalTok{    .attr("text{-}anchor", "middle").attr("font{-}size", 11).attr("fill", "\#444")}
\NormalTok{    .text("population");}
\NormalTok{  gR.append("text").attr("x", iW / 2).attr("y", {-}8)}
\NormalTok{    .attr("text{-}anchor", "middle").attr("font{-}size", 12).attr("font{-}weight", "bold").attr("fill", "\#333")}
\NormalTok{    .text("Time series");}

\NormalTok{  // Legend}
\NormalTok{  const legX = iW {-} 100, legY = 8;}
\NormalTok{  gR.append("line").attr("x1", legX).attr("x2", legX+18).attr("y1", legY).attr("y2", legY)}
\NormalTok{    .attr("stroke", "\#2e8b57").attr("stroke{-}width", 2);}
\NormalTok{  gR.append("text").attr("x", legX+22).attr("y", legY+4)}
\NormalTok{    .attr("font{-}size", 10).attr("fill", "\#2e8b57").text("prey (x)");}
\NormalTok{  gR.append("line").attr("x1", legX).attr("x2", legX+18).attr("y1", legY+16).attr("y2", legY+16)}
\NormalTok{    .attr("stroke", "\#e07b2a").attr("stroke{-}width", 2);}
\NormalTok{  gR.append("text").attr("x", legX+22).attr("y", legY+20)}
\NormalTok{    .attr("font{-}size", 10).attr("fill", "\#e07b2a").text("predators (y)");}

\NormalTok{  return svg.node();}
\NormalTok{\}}
\end{Highlighting}
\end{Shaded}

\begin{center}\rule{0.5\linewidth}{0.5pt}\end{center}

\section{Application: Coupled spring
system}\label{application-coupled-spring-system}

Two masses \(m_1\) and \(m_2\) are connected by three springs with
constants \(k_1\), \(k\), \(k_2\), attached to fixed walls at both ends.
The middle spring couples the two masses. Let \(x_1\) and \(x_2\) be
displacements from equilibrium.

Newton's second law for each mass:

\[m_1 x_1'' = -k_1 x_1 - k(x_1 - x_2) = -(k_1+k)x_1 + kx_2\]
\[m_2 x_2'' = -k_2 x_2 + k(x_1 - x_2) = kx_1 - (k_2+k)x_2\]

For equal masses \(m_1 = m_2 = 1\) and equal springs
\(k_1 = k_2 = k = 1\):

\[x_1'' = -2x_1 + x_2, \qquad x_2'' = x_1 - 2x_2\]

Convert to a first-order system with state vector
\((x_1, x_2, x_1', x_2')^T\):

\[A = \begin{pmatrix} 0 & 0 & 1 & 0 \\ 0 & 0 & 0 & 1 \\ -2 & 1 & 0 & 0 \\ 1 & -2 & 0 & 0 \end{pmatrix}\]

The eigenvalues of \(A\) are \(\pm i\) and \(\pm i\sqrt{3}\) --- two
pairs of purely imaginary eigenvalues. The system oscillates at two
\textbf{natural frequencies}: \(\omega_1 = 1\) (the two masses moving
together, in phase) and \(\omega_2 = \sqrt{3}\) (the two masses moving
in opposition, out of phase). These are the \textbf{normal modes} of the
system.

Every motion of the coupled system is a superposition of these two
modes. If the initial conditions excite only one mode, that mode
persists unchanged. If both are excited, the motion is a superposition
--- a beating pattern in which energy transfers back and forth between
the two masses.

\begin{center}\rule{0.5\linewidth}{0.5pt}\end{center}

\section{Exercises}\label{exercises-28}

\begin{Shaded}
\begin{Highlighting}[]
\NormalTok{//| echo: false}
\NormalTok{katex = require("katex@0.16")}
\end{Highlighting}
\end{Shaded}

\begin{Shaded}
\begin{Highlighting}[]
\NormalTok{//| echo: false}
\NormalTok{\{}
\NormalTok{  const css = document.createElement("link");}
\NormalTok{  css.rel = "stylesheet";}
\NormalTok{  css.href = "https://cdn.jsdelivr.net/npm/katex@0.16/dist/katex.min.css";}
\NormalTok{  document.head.appendChild(css);}
\NormalTok{\}}
\end{Highlighting}
\end{Shaded}

\begin{Shaded}
\begin{Highlighting}[]
\NormalTok{//| echo: false}
\NormalTok{function makeStepperHTML(steps) \{}
\NormalTok{  let current = 0;}
\NormalTok{  const totalSteps = steps.length;}
\NormalTok{  const container = document.createElement("div");}
\NormalTok{  container.style.cssText = "border:1px solid \#e5e7eb; border{-}radius:8px; padding:1rem 1.25rem; margin:0.75rem 0; font{-}family:inherit;";}
\NormalTok{  const stepsDiv = document.createElement("div");}
\NormalTok{  stepsDiv.style.marginBottom = "0.75rem";}

\NormalTok{  function renderSteps() \{}
\NormalTok{    stepsDiv.innerHTML = "";}
\NormalTok{    for (let i = 0; i \textless{} current; i++) \{}
\NormalTok{      const s = steps[i];}
\NormalTok{      const row = document.createElement("div");}
\NormalTok{      row.style.cssText = "display:grid; grid{-}template{-}columns:220px 1fr; gap:0.35rem 1rem; align{-}items:baseline; padding:0.35rem 0; border{-}top:1px solid \#f3f4f6;";}
\NormalTok{      const opCell = document.createElement("span");}
\NormalTok{      opCell.style.cssText = "font{-}size:0.85em; color:\#6b7280; font{-}style:italic;";}
\NormalTok{      opCell.textContent = s.op;}
\NormalTok{      const eqCell = document.createElement("span");}
\NormalTok{      eqCell.innerHTML = katex.renderToString(s.eq, \{ throwOnError: false, displayMode: false \});}
\NormalTok{      row.appendChild(opCell);}
\NormalTok{      row.appendChild(eqCell);}
\NormalTok{      if (s.note) \{}
\NormalTok{        const noteRow = document.createElement("div");}
\NormalTok{        noteRow.style.cssText = "grid{-}column:2; font{-}size:0.82em; color:\#6b7280; padding{-}bottom:0.2rem;";}
\NormalTok{        noteRow.textContent = s.note;}
\NormalTok{        row.appendChild(noteRow);}
\NormalTok{      \}}
\NormalTok{      stepsDiv.appendChild(row);}
\NormalTok{    \}}
\NormalTok{    if (current === totalSteps) \{}
\NormalTok{      const done = document.createElement("div");}
\NormalTok{      done.style.cssText = "margin{-}top:0.5rem; font{-}size:0.9em; color:\#059669; font{-}weight:500;";}
\NormalTok{      done.textContent = "Solution complete.";}
\NormalTok{      stepsDiv.appendChild(done);}
\NormalTok{    \}}
\NormalTok{  \}}

\NormalTok{  const controls = document.createElement("div");}
\NormalTok{  controls.style.cssText = "display:flex; gap:0.5rem; align{-}items:center;";}
\NormalTok{  const nextBtn = document.createElement("button");}
\NormalTok{  nextBtn.textContent = "Next step →";}
\NormalTok{  nextBtn.style.cssText = "padding:0.35rem 0.85rem; border:1px solid \#d1d5db; border{-}radius:4px; background:\#fff; cursor:pointer; font{-}size:0.9em;";}
\NormalTok{  const resetBtn = document.createElement("button");}
\NormalTok{  resetBtn.textContent = "Reset";}
\NormalTok{  resetBtn.style.cssText = "padding:0.35rem 0.75rem; border:1px solid \#e5e7eb; border{-}radius:4px; background:\#f9fafb; cursor:pointer; font{-}size:0.9em; color:\#6b7280;";}
\NormalTok{  const counter = document.createElement("span");}
\NormalTok{  counter.style.cssText = "font{-}size:0.82em; color:\#9ca3af; margin{-}left:0.25rem;";}

\NormalTok{  function updateButtons() \{}
\NormalTok{    nextBtn.disabled = current === totalSteps;}
\NormalTok{    nextBtn.style.opacity = current === totalSteps ? "0.4" : "1";}
\NormalTok{    counter.textContent = current === 0 ? "Click to reveal steps" : \textasciigrave{}Step $\{current\} of $\{totalSteps\}\textasciigrave{};}
\NormalTok{  \}}

\NormalTok{  nextBtn.onclick = () =\textgreater{} \{ if (current \textless{} totalSteps) \{ current++; renderSteps(); updateButtons(); \} \};}
\NormalTok{  resetBtn.onclick = () =\textgreater{} \{ current = 0; renderSteps(); updateButtons(); \};}
\NormalTok{  controls.appendChild(nextBtn);}
\NormalTok{  controls.appendChild(resetBtn);}
\NormalTok{  controls.appendChild(counter);}
\NormalTok{  container.appendChild(stepsDiv);}
\NormalTok{  container.appendChild(controls);}
\NormalTok{  updateButtons();}
\NormalTok{  return container;}
\NormalTok{\}}
\end{Highlighting}
\end{Shaded}

\textbf{Exercise 1.} Convert \(y'' + 3y' + 2y = 0\) to a first-order
system; write in matrix form; find the eigenvalues.

\begin{Shaded}
\begin{Highlighting}[]
\NormalTok{//| echo: false}
\NormalTok{makeStepperHTML([}
\NormalTok{  \{ op: "Introduce state variables", eq: "x\_1 = y,\textbackslash{}\textbackslash{}quad x\_2 = y\textquotesingle{} \textbackslash{}\textbackslash{}implies x\_1\textquotesingle{} = x\_2,\textbackslash{}\textbackslash{}quad x\_2\textquotesingle{} = {-}2x\_1 {-} 3x\_2" \},}
\NormalTok{  \{ op: "Write in matrix form", eq: "\textbackslash{}\textbackslash{}begin\{pmatrix\}x\_1\textquotesingle{}\textbackslash{}\textbackslash{}\textbackslash{}\textbackslash{}x\_2\textquotesingle{}\textbackslash{}\textbackslash{}end\{pmatrix\} = \textbackslash{}\textbackslash{}begin\{pmatrix\}0\&1\textbackslash{}\textbackslash{}\textbackslash{}\textbackslash{}{-}2\&{-}3\textbackslash{}\textbackslash{}end\{pmatrix\}\textbackslash{}\textbackslash{}begin\{pmatrix\}x\_1\textbackslash{}\textbackslash{}\textbackslash{}\textbackslash{}x\_2\textbackslash{}\textbackslash{}end\{pmatrix\}" \},}
\NormalTok{  \{ op: "Characteristic polynomial", eq: "\textbackslash{}\textbackslash{}det(A {-} \textbackslash{}\textbackslash{}lambda I) = \textbackslash{}\textbackslash{}lambda\^{}2 + 3\textbackslash{}\textbackslash{}lambda + 2" \},}
\NormalTok{  \{ op: "Factor and solve", eq: "(\textbackslash{}\textbackslash{}lambda+1)(\textbackslash{}\textbackslash{}lambda+2) = 0 \textbackslash{}\textbackslash{}implies \textbackslash{}\textbackslash{}lambda\_1 = {-}1,\textbackslash{}\textbackslash{}quad \textbackslash{}\textbackslash{}lambda\_2 = {-}2" \},}
\NormalTok{  \{ op: "Interpret", eq: "\textbackslash{}\textbackslash{}lambda\_1, \textbackslash{}\textbackslash{}lambda\_2 \textless{} 0 \textbackslash{}\textbackslash{}implies \textbackslash{}\textbackslash{}text\{stable node — solutions decay to zero\}", note: "Both eigenvalues real and negative: consistent with the damped oscillator." \}}
\NormalTok{])}
\end{Highlighting}
\end{Shaded}

\begin{center}\rule{0.5\linewidth}{0.5pt}\end{center}

\textbf{Exercise 2.} Given
\(A = \begin{pmatrix}2&1\\0&-1\end{pmatrix}\), find the eigenvalues and
classify the equilibrium at the origin.

\begin{Shaded}
\begin{Highlighting}[]
\NormalTok{//| echo: false}
\NormalTok{makeStepperHTML([}
\NormalTok{  \{ op: "Read eigenvalues (triangular)", eq: "\textbackslash{}\textbackslash{}lambda\_1 = 2,\textbackslash{}\textbackslash{}quad \textbackslash{}\textbackslash{}lambda\_2 = {-}1", note: "Upper triangular: eigenvalues are the diagonal entries." \},}
\NormalTok{  \{ op: "Check signs", eq: "\textbackslash{}\textbackslash{}lambda\_1 \textgreater{} 0,\textbackslash{}\textbackslash{}quad \textbackslash{}\textbackslash{}lambda\_2 \textless{} 0 \textbackslash{}\textbackslash{}implies \textbackslash{}\textbackslash{}text\{saddle point — unstable\}" \},}
\NormalTok{  \{ op: "Eigenvector for λ₁ = 2", eq: "(A{-}2I)\textbackslash{}\textbackslash{}mathbf\{v\}=0 \textbackslash{}\textbackslash{}implies \textbackslash{}\textbackslash{}mathbf\{v\}\_1 = (1,0)\^{}T" \},}
\NormalTok{  \{ op: "Eigenvector for λ₂ = −1", eq: "(A+I)\textbackslash{}\textbackslash{}mathbf\{v\}=0 \textbackslash{}\textbackslash{}implies \textbackslash{}\textbackslash{}mathbf\{v\}\_2 = (1,{-}3)\^{}T" \},}
\NormalTok{  \{ op: "Phase portrait", eq: "\textbackslash{}\textbackslash{}text\{Unstable manifold along \}(1,0)\^{}T\textbackslash{}\textbackslash{}text\{; stable manifold along \}(1,{-}3)\^{}T", note: "All trajectories approach along the stable manifold and depart along the unstable manifold." \}}
\NormalTok{])}
\end{Highlighting}
\end{Shaded}

\begin{center}\rule{0.5\linewidth}{0.5pt}\end{center}

\textbf{Exercise 3.} Given
\(A = \begin{pmatrix}-1&2\\-2&-1\end{pmatrix}\), find the eigenvalues,
classify the equilibrium, and write the general solution.

Step 1 \textbar{} Compute the characteristic polynomial \textbar{}
\(\det(A - \lambda I) = (-1-\lambda)^2 + 4 = \lambda^2 + 2\lambda + 5\)
Step 2 \textbar{} Solve \textbar{}
\(\lambda = \dfrac{-2 \pm \sqrt{4-20}}{2} = -1 \pm 2i\) Step 3
\textbar{} Classify \textbar{} Complex eigenvalues with
\(\alpha = -1 < 0\): \textbf{stable spiral}. All trajectories spiral
into the origin. Step 4 \textbar{} Find a complex eigenvector \textbar{}
For \(\lambda = -1+2i\): \((A - \lambda I)\mathbf{v} = 0\) gives
\(\begin{pmatrix}-2i&2\\-2&-2i\end{pmatrix}\mathbf{v}=0\), so
\(\mathbf{v} = \begin{pmatrix}1\\i\end{pmatrix}\). Step 5 \textbar{}
Write the real general solution \textbar{} Separate real and imaginary
parts of \(e^{(-1+2i)t}\begin{pmatrix}1\\i\end{pmatrix}\): the two real
solutions are \(e^{-t}\begin{pmatrix}\cos 2t\\-\sin 2t\end{pmatrix}\)
and \(e^{-t}\begin{pmatrix}\sin 2t\\\cos 2t\end{pmatrix}\). General
solution:
\(\mathbf{x}(t) = e^{-t}\left[c_1\begin{pmatrix}\cos 2t\\-\sin 2t\end{pmatrix} + c_2\begin{pmatrix}\sin 2t\\\cos 2t\end{pmatrix}\right]\)

\begin{center}\rule{0.5\linewidth}{0.5pt}\end{center}

\textbf{Exercise 4.} Given
\(A = \begin{pmatrix}0&1\\-4&0\end{pmatrix}\), find the eigenvalues,
classify as a centre, write the general solution, and describe the phase
portrait.

Step 1 \textbar{} Compute the characteristic polynomial \textbar{}
\(\det(A - \lambda I) = -\lambda(-\lambda) + 4 = \lambda^2 + 4\) Step 2
\textbar{} Solve \textbar{} \(\lambda = \pm 2i\): purely imaginary
eigenvalues. Step 3 \textbar{} Classify \textbar{} \textbf{Centre.} All
trajectories are closed curves --- periodic orbits. Step 4 \textbar{}
Find a complex eigenvector \textbar{} For \(\lambda = 2i\):
\((A - 2iI)\mathbf{v} = 0\) gives
\(\begin{pmatrix}-2i&1\\-4&-2i\end{pmatrix}\mathbf{v}=0\), so
\(\mathbf{v} = \begin{pmatrix}1\\2i\end{pmatrix}\). Step 5 \textbar{}
Write the general solution \textbar{}
\(e^{2it}\begin{pmatrix}1\\2i\end{pmatrix}\) has real part
\(\begin{pmatrix}\cos 2t\\-2\sin 2t\end{pmatrix}\) and imaginary part
\(\begin{pmatrix}\sin 2t\\2\cos 2t\end{pmatrix}\). General solution:
\(\mathbf{x}(t) = c_1\begin{pmatrix}\cos 2t\\-2\sin 2t\end{pmatrix} + c_2\begin{pmatrix}\sin 2t\\2\cos 2t\end{pmatrix}\)
Step 6 \textbar{} Describe the phase portrait \textbar{} Eliminating
\(t\): if \(x = c_1\cos 2t + c_2\sin 2t\) and
\(y = -2c_1\sin 2t + 2c_2\cos 2t\), then
\(x^2 + (y/2)^2 = c_1^2 + c_2^2\). The trajectories are ellipses, with
semi-axes in ratio \(1:2\) (elongated vertically), traversed in the
clockwise direction.

\begin{center}\rule{0.5\linewidth}{0.5pt}\end{center}

\textbf{Exercise 5.} Linearise \(\{x' = x(1-y),\; y' = y(x-1)\}\) at the
equilibrium \((1,1)\) and classify it.

Step 1 \textbar{} Verify equilibrium \textbar{} At \((1,1)\):
\(x' = 1\cdot(1-1) = 0\). \(y' = 1\cdot(1-1) = 0\). Confirmed. Step 2
\textbar{} Compute the Jacobian \textbar{} \(f(x,y) = x - xy\),
\(g(x,y) = xy - y\).
\(J = \begin{pmatrix}\partial f/\partial x & \partial f/\partial y \\ \partial g/\partial x & \partial g/\partial y\end{pmatrix} = \begin{pmatrix}1-y & -x \\ y & x-1\end{pmatrix}\)
Step 3 \textbar{} Evaluate at \((1,1)\) \textbar{}
\(J(1,1) = \begin{pmatrix}0 & -1 \\ 1 & 0\end{pmatrix}\) Step 4
\textbar{} Find eigenvalues \textbar{}
\(\det(J - \lambda I) = \lambda^2 + 1 = 0 \Rightarrow \lambda = \pm i\)
Step 5 \textbar{} Classify \textbar{} Purely imaginary eigenvalues: the
linearisation is a \textbf{centre}. The system is equivalent to
Lotka-Volterra with \(a = b = c = d = 1\); the true nonlinear
equilibrium is also a centre (orbits are closed). The linearisation is
inconclusive about the nonlinear stability --- a Lyapunov function
argument is needed to confirm the orbits are closed.

\begin{center}\rule{0.5\linewidth}{0.5pt}\end{center}

\emph{Exercise 6 extends the methods to a system with a continuum of
equilibria --- one of which has a zero eigenvalue. The approach is the
same: Jacobian, eigenvalues, interpret. The \(R_0\) interpretation is
the key result.}

\textbf{Exercise 6.} An epidemic model: \(S' = -0.002SI\),
\(I' = 0.002SI - 0.5I\) (SIR without recovery back to \(S\)). Find all
equilibria. Classify the disease-free equilibrium \((S_0, 0)\).
Interpret in terms of outbreak dynamics.

Step 1 \textbar{} Find equilibria \textbar{} Set \(S' = 0\):
\(-0.002SI = 0 \Rightarrow S = 0\) or \(I = 0\). Set \(I' = 0\):
\(I(0.002S - 0.5) = 0 \Rightarrow I = 0\) or \(S = 250\). Equilibria:
\((S, 0)\) for any \(S\) (disease-free line) and \((250, 0)\), which is
on the disease-free line. There is no interior equilibrium with
\(I > 0\). Step 2 \textbar{} Compute the Jacobian \textbar{}
\(f = -0.002SI\), \(g = 0.002SI - 0.5I\).
\(J = \begin{pmatrix}-0.002I & -0.002S \\ 0.002I & 0.002S - 0.5\end{pmatrix}\)
Step 3 \textbar{} Evaluate at \((S_0, 0)\) \textbar{}
\(J(S_0, 0) = \begin{pmatrix}0 & -0.002S_0 \\ 0 & 0.002S_0 - 0.5\end{pmatrix}\)
Step 4 \textbar{} Read eigenvalues (triangular matrix) \textbar{}
\(\lambda_1 = 0\), \(\lambda_2 = 0.002S_0 - 0.5\) Step 5 \textbar{}
Classify \textbar{} If \(S_0 < 250\): \(\lambda_2 < 0\), the
disease-free equilibrium is stable (no outbreak). If \(S_0 > 250\):
\(\lambda_2 > 0\), the equilibrium is unstable (outbreak occurs). The
threshold is \(S_0 = 250 = 0.5/0.002\). Step 6 \textbar{} Interpret
\textbar{} The ratio \(R_0 = 0.002S_0 / 0.5 = S_0/250\) is the basic
reproduction number. If \(R_0 > 1\) (more than 250 susceptibles), an
infection introduced into the population will spread. If \(R_0 < 1\),
any outbreak dies out. This is the fundamental threshold theorem of
mathematical epidemiology, here derived as a stability condition on the
disease-free equilibrium.

\begin{center}\rule{0.5\linewidth}{0.5pt}\end{center}

\begin{tcolorbox}[enhanced jigsaw, title=\textcolor{quarto-callout-tip-color}{\faLightbulb}\hspace{0.5em}{Key results}, opacityback=0, colback=white, breakable, leftrule=.75mm, arc=.35mm, toptitle=1mm, left=2mm, colframe=quarto-callout-tip-color-frame, coltitle=black, titlerule=0mm, colbacktitle=quarto-callout-tip-color!10!white, opacitybacktitle=0.6, toprule=.15mm, bottomtitle=1mm, bottomrule=.15mm, rightrule=.15mm]

\begin{enumerate}
\def\labelenumi{\arabic{enumi}.}
\item
  A second-order ODE \(y'' + py' + qy = 0\) is equivalent to the linear
  system \(\mathbf{x}' = A\mathbf{x}\) with
  \(A = \begin{pmatrix}0&1\\-q&-p\end{pmatrix}\).
\item
  The general solution is
  \(\mathbf{x}(t) = c_1 e^{\lambda_1 t}\mathbf{v}_1 + c_2 e^{\lambda_2 t}\mathbf{v}_2\).
\item
  Eigenvalue classification: real same-sign \(\to\) node; real
  opposite-sign \(\to\) saddle; complex \(\to\) spiral; purely imaginary
  \(\to\) centre.
\item
  Stability: asymptotically stable if and only if all eigenvalues have
  negative real part.
\item
  For nonlinear systems, linearise via the Jacobian \(J\) at each
  equilibrium and classify from the eigenvalues of \(J\).
\end{enumerate}

\end{tcolorbox}

\section{Where this goes}\label{where-this-goes-26}

The eigenvalue geometry developed here connects directly to three next
steps.

\textbf{Laplace transforms} (Chapter 4 of this volume) re-express the
same eigenvalue information as poles of a transfer function in the
complex \(s\)-plane. A stable spiral in the phase plane corresponds to a
pair of complex conjugate poles with negative real part --- same
mathematics, different language.

\textbf{Linear algebra} (Vol 7 linear algebra section) develops the
eigenvalue theory further: diagonalisation, the spectral theorem, and
Jordan form handle cases where eigenvectors are not a full basis. The
phase plane intuition built here makes those abstract results concrete.

\textbf{Partial differential equations} (Vol 7 Fourier--PDEs section)
extend the phase plane idea to spatial systems: the stability analysis
of a PDE linearised around a steady state uses the same eigenvalue
logic, now applied to differential operators rather than matrices.

\chapter{Laplace transforms}\label{laplace-transforms}

Converting differential equations to algebra

\hfill\break

A solenoid valve in a hydraulic system is closed. At exactly \(t = 2\)
seconds it opens, sending a pressure pulse through the system. The ODE
governing the pressure response has discontinuous forcing --- zero for
\(t < 2\), nonzero for \(t \geq 2\). The characteristic equation method
does not handle this cleanly; you would need to split the problem at
\(t = 2\), solve each piece separately, then match conditions at the
join. For two segments that is tedious. For a system that switches at
\(t = 2\), \(t = 5\), and \(t = 11\), it becomes impractical.

The Laplace transform handles it in one pass. The discontinuity is
encoded as an exponential factor \(e^{-2s}\) inside an algebraic
expression in \(s\). Solve the algebra. Invert. Done.

This chapter develops that mechanism from the integral definition
through to impulse response and the convolution theorem.

\section{Definition and motivation}\label{definition-and-motivation}

The Laplace transform of a function \(f(t)\) defined for \(t \geq 0\) is

\[\mathcal{L}\{f(t)\} = F(s) = \int_0^\infty e^{-st} f(t)\,dt\]

The variable \(s\) is a complex number, though for most engineering
applications we treat it as a real parameter and work in the right
half-plane where the integral converges.

\textbf{Why does this turn a differential equation into algebra?} Two
reasons.

First, the transform converts differentiation in \(t\) into
multiplication by \(s\) in the frequency domain (plus boundary terms
carrying the initial conditions). We derive this precisely in Section 3.

Second, the transform is linear:

\[\mathcal{L}\{af(t) + bg(t)\} = aF(s) + bG(s)\]

This follows directly from linearity of the integral. So a linear ODE in
\(t\) becomes a linear algebraic equation in \(s\). Solve for \(Y(s)\),
then invert.

The four-step method, which we will use throughout this chapter:

\begin{enumerate}
\def\labelenumi{\arabic{enumi}.}
\tightlist
\item
  \textbf{Take the Laplace transform} of both sides of the ODE.
\item
  \textbf{Apply the initial conditions} --- they appear automatically in
  the boundary terms.
\item
  \textbf{Solve algebraically} for \(Y(s)\).
\item
  \textbf{Invert} \(Y(s)\) to find \(y(t)\).
\end{enumerate}

\section{Table of transforms}\label{table-of-transforms}

We build the table entry by entry from the integral definition.

\textbf{Transform of 1}

\[\mathcal{L}\{1\} = \int_0^\infty e^{-st}\,dt = \left[-\frac{e^{-st}}{s}\right]_0^\infty = \frac{1}{s}, \quad s > 0\]

The boundary term at \(t \to \infty\) vanishes because \(e^{-st} \to 0\)
for \(s > 0\).

\textbf{Transform of \(e^{at}\)}

\[\mathcal{L}\{e^{at}\} = \int_0^\infty e^{-st} e^{at}\,dt = \int_0^\infty e^{-(s-a)t}\,dt = \frac{1}{s-a}, \quad s > a\]

The same calculation as above but with \(s\) replaced by \(s - a\).

\textbf{Transform of \(\sin(\omega t)\)}

Use the identity
\(\sin(\omega t) = \frac{e^{i\omega t} - e^{-i\omega t}}{2i}\), or
integrate directly using integration by parts twice. The direct route:

\[\mathcal{L}\{\sin(\omega t)\} = \int_0^\infty e^{-st}\sin(\omega t)\,dt\]

Let \(I\) denote the integral. For the first integration by parts, set
\(u = e^{-st}\) and \(dv = \sin(\omega t)\,dt\), so
\(du = -se^{-st}\,dt\) and \(v = -\cos(\omega t)/\omega\):

\[I = \left[-\frac{e^{-st}\cos(\omega t)}{\omega}\right]_0^\infty - \int_0^\infty \frac{s}{\omega}e^{-st}\cos(\omega t)\,dt = \frac{1}{\omega} - \frac{s}{\omega}\int_0^\infty e^{-st}\cos(\omega t)\,dt\]

(The boundary term at \(t = 0\) gives \(1/\omega\); at \(t \to \infty\)
the exponential kills the cosine.)

Applying IBP again to \(\int e^{-st}\cos(\omega t)\,dt\) by symmetry
(set \(u = e^{-st}\), \(dv = \cos(\omega t)\,dt\),
\(v = \sin(\omega t)/\omega\)) returns the original integral \(I\) with
a coefficient. After two applications the integral \(I\) reappears on
the right:

\[I = \frac{\omega}{s^2} - \frac{\omega^2}{s^2}I\]

Solving for \(I\):

\[I\left(1 + \frac{\omega^2}{s^2}\right) = \frac{\omega}{s^2} \implies I = \frac{\omega}{s^2 + \omega^2}\]

Similarly, integrating \(e^{-st}\cos(\omega t)\) by parts (or using the
derivative relation) gives
\(\mathcal{L}\{\cos(\omega t)\} = \dfrac{s}{s^2 + \omega^2}\).

\textbf{Transform of \(t^n\)}

Integrate by parts, reducing the power of \(t\) by one at each step. The
result is

\[\mathcal{L}\{t^n\} = \frac{n!}{s^{n+1}}, \quad s > 0\]

For \(n = 1\): \(\mathcal{L}\{t\} = 1/s^2\). For \(n = 2\):
\(\mathcal{L}\{t^2\} = 2/s^3\).

\textbf{Reference table}

\begin{longtable}[]{@{}
  >{\raggedright\arraybackslash}p{(\linewidth - 4\tabcolsep) * \real{0.3333}}
  >{\raggedright\arraybackslash}p{(\linewidth - 4\tabcolsep) * \real{0.3333}}
  >{\raggedright\arraybackslash}p{(\linewidth - 4\tabcolsep) * \real{0.3333}}@{}}
\toprule\noalign{}
\begin{minipage}[b]{\linewidth}\raggedright
\(f(t)\)
\end{minipage} & \begin{minipage}[b]{\linewidth}\raggedright
\(F(s) = \mathcal{L}\{f(t)\}\)
\end{minipage} & \begin{minipage}[b]{\linewidth}\raggedright
Condition
\end{minipage} \\
\midrule\noalign{}
\endhead
\bottomrule\noalign{}
\endlastfoot
\(1\) & \(\dfrac{1}{s}\) & \(s > 0\) \\
\(t\) & \(\dfrac{1}{s^2}\) & \(s > 0\) \\
\(t^n\) & \(\dfrac{n!}{s^{n+1}}\) & \(s > 0\) \\
\(e^{at}\) & \(\dfrac{1}{s-a}\) & \(s > a\) \\
\(te^{at}\) & \(\dfrac{1}{(s-a)^2}\) & \(s > a\) \\
\(\sin(\omega t)\) & \(\dfrac{\omega}{s^2+\omega^2}\) & \(s > 0\) \\
\(\cos(\omega t)\) & \(\dfrac{s}{s^2+\omega^2}\) & \(s > 0\) \\
\(e^{at}\sin(\omega t)\) & \(\dfrac{\omega}{(s-a)^2+\omega^2}\) &
\(s > a\) \\
\(e^{at}\cos(\omega t)\) & \(\dfrac{s-a}{(s-a)^2+\omega^2}\) &
\(s > a\) \\
\(u(t-a)\) & \(\dfrac{e^{-as}}{s}\) & \(s > 0\) \\
\(\delta(t-a)\) & \(e^{-as}\) & all \(s\) \\
\end{longtable}

\begin{Shaded}
\begin{Highlighting}[]
\NormalTok{//| echo: false}
\NormalTok{viewof ltFunc = Inputs.select(}
\NormalTok{  ["1 (unit step)", "e\^{}\{at\} (a={-}1)", "sin(ωt) (ω=2)", "cos(ωt) (ω=2)", "t"],}
\NormalTok{  \{label: "f(t)"\}}
\NormalTok{)}
\end{Highlighting}
\end{Shaded}

\begin{Shaded}
\begin{Highlighting}[]
\NormalTok{//| echo: false}
\NormalTok{viewof ltS = Inputs.range([0.1, 5], \{step: 0.1, value: 1, label: "s"\})}
\end{Highlighting}
\end{Shaded}

\begin{Shaded}
\begin{Highlighting}[]
\NormalTok{//| echo: false}
\NormalTok{\{}
\NormalTok{  const d3 = await require("d3@7");}

\NormalTok{  // Function definitions}
\NormalTok{  const funcs = \{}
\NormalTok{    "1 (unit step)":    \{ ft: t =\textgreater{} 1,          Fs: s =\textgreater{} 1/s,              yMax: 1.4 \},}
\NormalTok{    "e\^{}\{at\} (a={-}1)":   \{ ft: t =\textgreater{} Math.exp({-}t), Fs: s =\textgreater{} 1/(s+1),         yMax: 1.4 \},}
\NormalTok{    "sin(ωt) (ω=2)":   \{ ft: t =\textgreater{} Math.sin(2*t), Fs: s =\textgreater{} 2/(s*s+4),      yMax: 1.4 \},}
\NormalTok{    "cos(ωt) (ω=2)":   \{ ft: t =\textgreater{} Math.cos(2*t), Fs: s =\textgreater{} s/(s*s+4),      yMax: 1.4 \},}
\NormalTok{    "t":               \{ ft: t =\textgreater{} t,            Fs: s =\textgreater{} 1/(s*s),          yMax: 6.5 \}}
\NormalTok{  \};}

\NormalTok{  const def = funcs[ltFunc];}
\NormalTok{  const ft = def.ft;}
\NormalTok{  const Fs = def.Fs;}
\NormalTok{  const s  = ltS;}

\NormalTok{  // Layout}
\NormalTok{  const panelW = 260, panelH = 240;}
\NormalTok{  const margin = \{top: 20, right: 15, bottom: 35, left: 42\};}
\NormalTok{  const innerW = panelW {-} margin.left {-} margin.right;}
\NormalTok{  const innerH = panelH {-} margin.top {-} margin.bottom;}
\NormalTok{  const totalW = panelW * 2 + 20;}

\NormalTok{  const svg = d3.create("svg")}
\NormalTok{    .attr("viewBox", \textasciigrave{}0 0 $\{totalW\} $\{panelH\}\textasciigrave{})}
\NormalTok{    .attr("width", totalW)}
\NormalTok{    .attr("style", "font{-}family: sans{-}serif; font{-}size: 11px;");}

\NormalTok{  // ── LEFT PANEL: f(t) and e\^{}\{{-}st\}·f(t) ────────────────────────────────}
\NormalTok{  const tDomain = [0, 6];}
\NormalTok{  const tValues = d3.range(tDomain[0], tDomain[1], 0.02);}

\NormalTok{  const yMax = def.yMax;}
\NormalTok{  const yMin = ltFunc === "cos(ωt) (ω=2)" || ltFunc === "sin(ωt) (ω=2)" ? {-}yMax * 0.6 : {-}0.15;}

\NormalTok{  const xL = d3.scaleLinear().domain(tDomain).range([0, innerW]);}
\NormalTok{  const yL = d3.scaleLinear().domain([yMin, yMax]).range([innerH, 0]);}

\NormalTok{  const gL = svg.append("g")}
\NormalTok{    .attr("transform", \textasciigrave{}translate($\{margin.left\},$\{margin.top\})\textasciigrave{});}

\NormalTok{  // Shaded area under e\^{}\{{-}st\}·f(t)}
\NormalTok{  const areaGen = d3.area()}
\NormalTok{    .x(t =\textgreater{} xL(t))}
\NormalTok{    .y0(yL(0))}
\NormalTok{    .y1(t =\textgreater{} yL(Math.exp({-}s * t) * ft(t)));}

\NormalTok{  gL.append("path")}
\NormalTok{    .datum(tValues)}
\NormalTok{    .attr("d", areaGen)}
\NormalTok{    .attr("fill", "orange")}
\NormalTok{    .attr("opacity", 0.18);}

\NormalTok{  // Axes}
\NormalTok{  gL.append("g").attr("transform", \textasciigrave{}translate(0,$\{yL(0)\})\textasciigrave{}).call(d3.axisBottom(xL).ticks(5));}
\NormalTok{  gL.append("g").call(d3.axisLeft(yL).ticks(5));}

\NormalTok{  // Axis labels}
\NormalTok{  gL.append("text").attr("x", innerW / 2).attr("y", innerH + 28)}
\NormalTok{    .attr("text{-}anchor", "middle").text("t");}
\NormalTok{  gL.append("text")}
\NormalTok{    .attr("transform", "rotate({-}90)")}
\NormalTok{    .attr("x", {-}innerH / 2).attr("y", {-}32)}
\NormalTok{    .attr("text{-}anchor", "middle").text("value");}

\NormalTok{  // f(t) in blue}
\NormalTok{  const lineL = d3.line().x(t =\textgreater{} xL(t)).y(t =\textgreater{} yL(ft(t)));}
\NormalTok{  gL.append("path")}
\NormalTok{    .datum(tValues)}
\NormalTok{    .attr("d", lineL)}
\NormalTok{    .attr("fill", "none")}
\NormalTok{    .attr("stroke", "steelblue")}
\NormalTok{    .attr("stroke{-}width", 2);}

\NormalTok{  // e\^{}\{{-}st\}·f(t) in orange}
\NormalTok{  const lineI = d3.line().x(t =\textgreater{} xL(t)).y(t =\textgreater{} yL(Math.exp({-}s * t) * ft(t)));}
\NormalTok{  gL.append("path")}
\NormalTok{    .datum(tValues)}
\NormalTok{    .attr("d", lineI)}
\NormalTok{    .attr("fill", "none")}
\NormalTok{    .attr("stroke", "orange")}
\NormalTok{    .attr("stroke{-}width", 2);}

\NormalTok{  // Legend}
\NormalTok{  gL.append("line").attr("x1", 4).attr("y1", 8).attr("x2", 20).attr("y2", 8)}
\NormalTok{    .attr("stroke", "steelblue").attr("stroke{-}width", 2);}
\NormalTok{  gL.append("text").attr("x", 23).attr("y", 12).text("f(t)").attr("fill", "steelblue");}
\NormalTok{  gL.append("line").attr("x1", 55).attr("y1", 8).attr("x2", 71).attr("y2", 8)}
\NormalTok{    .attr("stroke", "orange").attr("stroke{-}width", 2);}
\NormalTok{  gL.append("text").attr("x", 74).attr("y", 12).text("e⁻ˢᵗf(t)").attr("fill", "orange");}

\NormalTok{  // Panel title}
\NormalTok{  gL.append("text").attr("x", innerW / 2).attr("y", {-}6)}
\NormalTok{    .attr("text{-}anchor", "middle").attr("font{-}size", "12px").attr("font{-}weight", "bold")}
\NormalTok{    .text("Time domain");}

\NormalTok{  // ── RIGHT PANEL: F(s) vs s ────────────────────────────────────────────}
\NormalTok{  const gapX = panelW + 20;}
\NormalTok{  const sValues = d3.range(0.1, 5.05, 0.05);}

\NormalTok{  const fsValues = sValues.map(sv =\textgreater{} Fs(sv));}
\NormalTok{  const fsFinite = fsValues.filter(v =\textgreater{} isFinite(v) \&\& !isNaN(v));}
\NormalTok{  const fsMax = Math.min(d3.max(fsFinite) * 1.1, 10);}
\NormalTok{  const fsMin = Math.max(d3.min(fsFinite) * 1.1, {-}2);}

\NormalTok{  const xR = d3.scaleLinear().domain([0.1, 5]).range([0, innerW]);}
\NormalTok{  const yR = d3.scaleLinear().domain([fsMin, fsMax]).range([innerH, 0]);}

\NormalTok{  const gR = svg.append("g")}
\NormalTok{    .attr("transform", \textasciigrave{}translate($\{gapX + margin.left\},$\{margin.top\})\textasciigrave{});}

\NormalTok{  // Axes}
\NormalTok{  gR.append("g").attr("transform", \textasciigrave{}translate(0,$\{yR(0)\})\textasciigrave{}).call(d3.axisBottom(xR).ticks(5));}
\NormalTok{  gR.append("g").call(d3.axisLeft(yR).ticks(5));}

\NormalTok{  // Axis labels}
\NormalTok{  gR.append("text").attr("x", innerW / 2).attr("y", innerH + 28)}
\NormalTok{    .attr("text{-}anchor", "middle").text("s");}
\NormalTok{  gR.append("text")}
\NormalTok{    .attr("transform", "rotate({-}90)")}
\NormalTok{    .attr("x", {-}innerH / 2).attr("y", {-}32)}
\NormalTok{    .attr("text{-}anchor", "middle").text("F(s)");}

\NormalTok{  // F(s) curve}
\NormalTok{  const lineR = d3.line()}
\NormalTok{    .defined((\_, i) =\textgreater{} isFinite(fsValues[i]) \&\& fsValues[i] \textless{}= fsMax + 1 \&\& fsValues[i] \textgreater{}= fsMin {-} 1)}
\NormalTok{    .x((\_, i) =\textgreater{} xR(sValues[i]))}
\NormalTok{    .y((\_, i) =\textgreater{} yR(fsValues[i]));}

\NormalTok{  gR.append("path")}
\NormalTok{    .datum(sValues)}
\NormalTok{    .attr("d", lineR)}
\NormalTok{    .attr("fill", "none")}
\NormalTok{    .attr("stroke", "steelblue")}
\NormalTok{    .attr("stroke{-}width", 2);}

\NormalTok{  // Vertical dashed line at current s}
\NormalTok{  const currentFs = Fs(s);}
\NormalTok{  if (isFinite(currentFs) \&\& currentFs \textless{}= fsMax + 0.5) \{}
\NormalTok{    gR.append("line")}
\NormalTok{      .attr("x1", xR(s)).attr("y1", 0)}
\NormalTok{      .attr("x2", xR(s)).attr("y2", innerH)}
\NormalTok{      .attr("stroke", "\#999")}
\NormalTok{      .attr("stroke{-}width", 1.5)}
\NormalTok{      .attr("stroke{-}dasharray", "4,3");}

\NormalTok{    // Dot on curve}
\NormalTok{    const dotY = Math.max(0, Math.min(innerH, yR(currentFs)));}
\NormalTok{    gR.append("circle")}
\NormalTok{      .attr("cx", xR(s)).attr("cy", dotY)}
\NormalTok{      .attr("r", 5)}
\NormalTok{      .attr("fill", "steelblue");}

\NormalTok{    // Label}
\NormalTok{    const labelVal = currentFs.toFixed(3);}
\NormalTok{    gR.append("text")}
\NormalTok{      .attr("x", Math.min(xR(s) + 6, innerW {-} 80))}
\NormalTok{      .attr("y", Math.max(dotY {-} 8, 14))}
\NormalTok{      .attr("fill", "steelblue")}
\NormalTok{      .attr("font{-}size", "11px")}
\NormalTok{      .text(\textasciigrave{}F($\{s.toFixed(1)\}) = $\{labelVal\}\textasciigrave{});}
\NormalTok{  \}}

\NormalTok{  // Panel title}
\NormalTok{  gR.append("text").attr("x", innerW / 2).attr("y", {-}6)}
\NormalTok{    .attr("text{-}anchor", "middle").attr("font{-}size", "12px").attr("font{-}weight", "bold")}
\NormalTok{    .text("s{-}domain: F(s)");}

\NormalTok{  return svg.node();}
\NormalTok{\}}
\end{Highlighting}
\end{Shaded}

\section{The transform of
derivatives}\label{the-transform-of-derivatives}

This is the mechanism that converts ODEs into algebra.

\textbf{First derivative.} Integrate by parts with \(u = f(t)\),
\(dv = e^{-st}\,dt\):

\[\mathcal{L}\{f'(t)\} = \int_0^\infty e^{-st} f'(t)\,dt = \left[e^{-st}f(t)\right]_0^\infty + s\int_0^\infty e^{-st}f(t)\,dt\]

The boundary term: at \(t \to \infty\) the factor \(e^{-st}\) kills
\(f(t)\) (for functions of exponential order --- meaning
\(|f(t)| \leq Me^{ct}\) for some constants \(M > 0\) and \(c\);
essentially, \(f(t)\) does not grow faster than some exponential, and
all functions you will encounter in practice satisfy this); at \(t = 0\)
we get \(-f(0)\). Therefore:

\[\boxed{\mathcal{L}\{f'(t)\} = sF(s) - f(0)}\]

The initial condition \(f(0)\) appears automatically as part of the
algebra.

\textbf{Second derivative.} Apply the first-derivative rule to \(f'\),
treating \(f'\) as a function in its own right:

\[\mathcal{L}\{f''(t)\} = s\mathcal{L}\{f'(t)\} - f'(0) = s[sF(s) - f(0)] - f'(0)\]

\[\boxed{\mathcal{L}\{f''(t)\} = s^2 F(s) - s f(0) - f'(0)}\]

Both initial conditions appear. This is why the Laplace transform method
is particularly clean for IVPs: the boundary data enters the algebra
rather than being appended at the end.

The pattern continues:
\(\mathcal{L}\{f^{(n)}\} = s^n F(s) - s^{n-1}f(0) - \cdots - f^{(n-1)}(0)\).

\section{Solving initial value
problems}\label{solving-initial-value-problems}

We now apply the four-step method.

\subsection{First-order example}\label{first-order-example}

\textbf{Solve} \(y' + 3y = 6\), \(y(0) = 1\).

\textbf{Step 1: Take the Laplace transform of both sides.}

\[\mathcal{L}\{y'\} + 3\mathcal{L}\{y\} = \mathcal{L}\{6\}\]

\[[sY(s) - y(0)] + 3Y(s) = \frac{6}{s}\]

\textbf{Step 2: Apply the initial condition \(y(0) = 1\).}

\[sY(s) - 1 + 3Y(s) = \frac{6}{s}\]

\textbf{Step 3: Solve algebraically for \(Y(s)\).}

\[Y(s)(s + 3) = \frac{6}{s} + 1 = \frac{6 + s}{s}\]

\[Y(s) = \frac{s + 6}{s(s+3)}\]

\textbf{Step 4: Invert.} We need partial fractions:

\[\frac{s+6}{s(s+3)} = \frac{A}{s} + \frac{B}{s+3}\]

Multiply through: \(s + 6 = A(s+3) + Bs\).

\begin{itemize}
\tightlist
\item
  Set \(s = 0\): \(6 = 3A\), so \(A = 2\).
\item
  Set \(s = -3\): \(3 = -3B\), so \(B = -1\).
\end{itemize}

Therefore:

\[Y(s) = \frac{2}{s} - \frac{1}{s+3}\]

Inverting term by term using the table (\(\mathcal{L}^{-1}\{1/s\} = 1\),
\(\mathcal{L}^{-1}\{1/(s+3)\} = e^{-3t}\)):

\[y(t) = 2 - e^{-3t}\]

Check: \(y(0) = 2 - 1 = 1\) ✓. As \(t \to \infty\), \(y \to 2\) (the
steady state, consistent with \(3 \times 2 = 6\)) ✓.

\subsection{Second-order example}\label{second-order-example}

\textbf{Solve} \(y'' + 5y' + 6y = e^{-t}\), \(y(0) = 1\), \(y'(0) = 0\).

\textbf{Step 1: Transform.}

\[[s^2 Y - sy(0) - y'(0)] + 5[sY - y(0)] + 6Y = \frac{1}{s+1}\]

\textbf{Step 2: Apply initial conditions.}

\[s^2 Y - s - 0 + 5sY - 5 + 6Y = \frac{1}{s+1}\]

\[(s^2 + 5s + 6)Y = \frac{1}{s+1} + s + 5\]

\textbf{Step 3: Solve for \(Y(s)\).}

\[Y(s) = \frac{1}{(s+1)(s+2)(s+3)} + \frac{s+5}{(s+2)(s+3)}\]

Combine over the common denominator \((s+1)(s+2)(s+3)\):

\[Y(s) = \frac{1 + (s+5)(s+1)}{(s+1)(s+2)(s+3)} = \frac{s^2 + 6s + 6}{(s+1)(s+2)(s+3)}\]

\textbf{Step 4: Partial fractions.}

\[\frac{s^2+6s+6}{(s+1)(s+2)(s+3)} = \frac{A}{s+1} + \frac{B}{s+2} + \frac{C}{s+3}\]

\begin{itemize}
\tightlist
\item
  \(s = -1\): \(1 - 6 + 6 = A(1)(2) \Rightarrow A = \tfrac{1}{2}\)
\item
  \(s = -2\): \(4 - 12 + 6 = B(-1)(1) \Rightarrow B = 2\)
\item
  \(s = -3\): \(9 - 18 + 6 = C(-2)(-1) \Rightarrow C = -\tfrac{3}{2}\)
\end{itemize}

\[y(t) = \frac{1}{2}e^{-t} + 2e^{-2t} - \frac{3}{2}e^{-3t}\]

\section{Partial fractions}\label{partial-fractions}

Inverting \(Y(s)\) almost always requires decomposing a rational
function into terms matching the table. Two cases arise repeatedly.

\subsection{Distinct real poles}\label{distinct-real-poles}

Given \(\dfrac{N(s)}{(s - r_1)(s - r_2)\cdots(s - r_n)}\) with distinct
roots, write

\[\frac{N(s)}{(s-r_1)\cdots(s-r_n)} = \frac{A_1}{s-r_1} + \cdots + \frac{A_n}{s-r_n}\]

Each coefficient: \(A_k = \dfrac{N(r_k)}{\prod_{j \neq k}(r_k - r_j)}\),
or equivalently multiply both sides by \((s - r_k)\) and set
\(s = r_k\). We did this in both worked examples above.

\subsection{Complex conjugate poles}\label{complex-conjugate-poles}

When the denominator has a quadratic factor \(s^2 + bs + c\) with no
real roots (complex conjugate poles), keep the quadratic together and
write a linear numerator:

\[\frac{N(s)}{(s-r)(s^2+bs+c)} = \frac{A}{s-r} + \frac{Bs + C}{s^2+bs+c}\]

After finding \(A\), \(B\), \(C\) by equating coefficients or
substituting convenient values of \(s\), complete the square in the
denominator and match to the \(\sin\) and \(\cos\) entries in the table.

\textbf{Example.} Invert \(\dfrac{s}{(s+1)(s^2+4)}\).

\[\frac{s}{(s+1)(s^2+4)} = \frac{A}{s+1} + \frac{Bs+C}{s^2+4}\]

Multiply through: \(s = A(s^2+4) + (Bs+C)(s+1)\).

\begin{itemize}
\tightlist
\item
  \(s = -1\): \(-1 = 5A\), so \(A = -\tfrac{1}{5}\).
\item
  \(s = 0\): \(0 = 4A + C\), so \(C = \tfrac{4}{5}\).
\item
  \(s^2\) coefficient: here we equate coefficients of each power of
  \(s\) rather than substituting a specific value --- both methods work;
  coefficient matching is more systematic when substitution would give
  \(0 = 0\). Equating \(s^2\) terms: \(0 = A + B\), so
  \(B = \tfrac{1}{5}\).
\end{itemize}

\[Y(s) = \frac{-1/5}{s+1} + \frac{(1/5)s + 4/5}{s^2+4} = -\frac{1}{5}\cdot\frac{1}{s+1} + \frac{1}{5}\cdot\frac{s}{s^2+4} + \frac{2}{5}\cdot\frac{2}{s^2+4}\]

Inverting:

\[y(t) = -\frac{1}{5}e^{-t} + \frac{1}{5}\cos(2t) + \frac{2}{5}\sin(2t)\]

\section{The Heaviside step function and piecewise
forcing}\label{the-heaviside-step-function-and-piecewise-forcing}

Engineering systems are often forced by signals that switch on or off at
a specific time. The Heaviside unit step function captures this:

\[u(t-a) = \begin{cases} 0 & t < a \\ 1 & t \geq a \end{cases}\]

Its Laplace transform is

\[\mathcal{L}\{u(t-a)\} = \int_a^\infty e^{-st}\,dt = \frac{e^{-as}}{s}\]

\textbf{Writing piecewise functions using \(u(t-a)\).} Any function that
switches on at \(t = a\) can be written using the Heaviside function.
The piecewise function

\[f(t) = \begin{cases} g(t) & 0 \leq t < a \\ h(t) & t \geq a \end{cases}\]

equals \(g(t) + [h(t) - g(t)]u(t-a)\).

\textbf{Second shifting theorem.} If \(\mathcal{L}\{f(t)\} = F(s)\),
then

\[\mathcal{L}\{f(t-a)\,u(t-a)\} = e^{-as}F(s)\]

A time delay by \(a\) in the \(t\)-domain multiplies the transform by
\(e^{-as}\). This is the mechanism that keeps piecewise forcing
contained in a single algebraic expression.

\subsection{Worked example: RC circuit with a step
input}\label{worked-example-rc-circuit-with-a-step-input}

\textbf{Problem.} An RC circuit with \(R = 1\,\Omega\) and
\(C = 1\,\text{F}\) has the voltage equation

\[\frac{dv}{dt} + v = V_\text{in}(t), \quad v(0) = 0\]

where \(V_\text{in}(t) = u(t - 2)\) --- a unit voltage step applied at
\(t = 2\).

\textbf{Step 1: Transform.}

\[sV(s) - v(0) + V(s) = \frac{e^{-2s}}{s}\]

\textbf{Step 2: Apply \(v(0) = 0\).}

\[V(s)(s+1) = \frac{e^{-2s}}{s}\]

\textbf{Step 3: Solve for \(V(s)\).}

\[V(s) = \frac{e^{-2s}}{s(s+1)}\]

\textbf{Step 4: Invert.} Factor out \(e^{-2s}\), invert the rest, then
apply the shifting theorem.

The function without the exponential factor is
\(G(s) = \dfrac{1}{s(s+1)}\). Partial fractions:
\(\dfrac{1}{s(s+1)} = \dfrac{1}{s} - \dfrac{1}{s+1}\), so
\(g(t) = 1 - e^{-t}\).

By the second shifting theorem:

\[v(t) = g(t-2)\,u(t-2) = \left(1 - e^{-(t-2)}\right)u(t-2)\]

Explicitly:

\[v(t) = \begin{cases} 0 & t < 2 \\ 1 - e^{-(t-2)} & t \geq 2 \end{cases}\]

The circuit is at rest until \(t = 2\). After the switch closes, the
voltage rises exponentially toward 1 with time constant
\(\tau = RC = 1\) second.

\begin{Shaded}
\begin{Highlighting}[]
\NormalTok{//| echo: false}
\NormalTok{viewof srW0 = Inputs.range([0.5, 5], \{step: 0.1, value: 2, label: "ω₀ (natural frequency)"\})}
\end{Highlighting}
\end{Shaded}

\begin{Shaded}
\begin{Highlighting}[]
\NormalTok{//| echo: false}
\NormalTok{viewof srZeta = Inputs.range([0.05, 2], \{step: 0.05, value: 0.3, label: "ζ (damping ratio)"\})}
\end{Highlighting}
\end{Shaded}

\begin{Shaded}
\begin{Highlighting}[]
\NormalTok{//| echo: false}
\NormalTok{viewof srA = Inputs.range([0, 8], \{step: 0.5, value: 3, label: "Switch time a (s)"\})}
\end{Highlighting}
\end{Shaded}

\begin{Shaded}
\begin{Highlighting}[]
\NormalTok{//| echo: false}
\NormalTok{\{}
\NormalTok{  const d3 = await require("d3@7");}

\NormalTok{  const w0   = srW0;}
\NormalTok{  const zeta = srZeta;}
\NormalTok{  const a    = srA;}

\NormalTok{  // RK4 numerical integration of y\textquotesingle{}\textquotesingle{} + 2ζω₀y\textquotesingle{} + ω₀²y = ω₀²·u(t{-}a)}
\NormalTok{  // State: [y, dy/dt]}
\NormalTok{  const dt   = 0.01;}
\NormalTok{  const tEnd = 20;}
\NormalTok{  const N    = Math.ceil(tEnd / dt);}

\NormalTok{  const ts = new Float64Array(N + 1);}
\NormalTok{  const ys = new Float64Array(N + 1);}

\NormalTok{  const deriv = (t, y, v) =\textgreater{} \{}
\NormalTok{    const u = t \textgreater{}= a ? 1.0 : 0.0;}
\NormalTok{    const dvdt = w0 * w0 * u {-} 2 * zeta * w0 * v {-} w0 * w0 * y;}
\NormalTok{    return [v, dvdt];}
\NormalTok{  \};}

\NormalTok{  let y = 0, v = 0;}
\NormalTok{  ts[0] = 0; ys[0] = 0;}
\NormalTok{  for (let i = 0; i \textless{} N; i++) \{}
\NormalTok{    const t = ts[i];}
\NormalTok{    const [k1y, k1v] = deriv(t,          y,                    v);}
\NormalTok{    const [k2y, k2v] = deriv(t + dt/2,   y + dt/2 * k1y,      v + dt/2 * k1v);}
\NormalTok{    const [k3y, k3v] = deriv(t + dt/2,   y + dt/2 * k2y,      v + dt/2 * k2v);}
\NormalTok{    const [k4y, k4v] = deriv(t + dt,     y + dt * k3y,         v + dt * k3v);}
\NormalTok{    y += (dt / 6) * (k1y + 2*k2y + 2*k3y + k4y);}
\NormalTok{    v += (dt / 6) * (k1v + 2*k2v + 2*k3v + k4v);}
\NormalTok{    ts[i+1] = t + dt;}
\NormalTok{    ys[i+1] = y;}
\NormalTok{  \}}

\NormalTok{  // Layout}
\NormalTok{  const width = 560, height = 300;}
\NormalTok{  const margin = \{top: 24, right: 24, bottom: 40, left: 50\};}
\NormalTok{  const innerW = width  {-} margin.left {-} margin.right;}
\NormalTok{  const innerH = height {-} margin.top  {-} margin.bottom;}

\NormalTok{  const yAll = Array.from(ys);}
\NormalTok{  const yMax = Math.max(d3.max(yAll) * 1.15, 1.2);}
\NormalTok{  const yMin = Math.min(d3.min(yAll) * 1.1, {-}0.15);}

\NormalTok{  const xScale = d3.scaleLinear().domain([0, tEnd]).range([0, innerW]);}
\NormalTok{  const yScale = d3.scaleLinear().domain([yMin, yMax]).range([innerH, 0]);}

\NormalTok{  const svg = d3.create("svg")}
\NormalTok{    .attr("viewBox", \textasciigrave{}0 0 $\{width\} $\{height\}\textasciigrave{})}
\NormalTok{    .attr("width", width)}
\NormalTok{    .attr("style", "font{-}family: sans{-}serif; font{-}size: 11px;");}

\NormalTok{  const g = svg.append("g")}
\NormalTok{    .attr("transform", \textasciigrave{}translate($\{margin.left\},$\{margin.top\})\textasciigrave{});}

\NormalTok{  // Steady{-}state dashed line}
\NormalTok{  g.append("line")}
\NormalTok{    .attr("x1", 0).attr("y1", yScale(1))}
\NormalTok{    .attr("x2", innerW).attr("y2", yScale(1))}
\NormalTok{    .attr("stroke", "\#aaa")}
\NormalTok{    .attr("stroke{-}width", 1.5)}
\NormalTok{    .attr("stroke{-}dasharray", "6,4");}

\NormalTok{  g.append("text")}
\NormalTok{    .attr("x", innerW {-} 4).attr("y", yScale(1) {-} 5)}
\NormalTok{    .attr("text{-}anchor", "end")}
\NormalTok{    .attr("fill", "\#888")}
\NormalTok{    .text("Steady state (y = 1)");}

\NormalTok{  // Switch time dashed line}
\NormalTok{  if (a \textless{}= tEnd) \{}
\NormalTok{    g.append("line")}
\NormalTok{      .attr("x1", xScale(a)).attr("y1", 0)}
\NormalTok{      .attr("x2", xScale(a)).attr("y2", innerH)}
\NormalTok{      .attr("stroke", "\#bbb")}
\NormalTok{      .attr("stroke{-}width", 1.5)}
\NormalTok{      .attr("stroke{-}dasharray", "4,3");}

\NormalTok{    g.append("text")}
\NormalTok{      .attr("x", xScale(a) + 4).attr("y", 14)}
\NormalTok{      .attr("fill", "\#999")}
\NormalTok{      .text(\textasciigrave{}Switch at t = $\{a\}\textasciigrave{});}
\NormalTok{  \}}

\NormalTok{  // Axes}
\NormalTok{  g.append("g").attr("transform", \textasciigrave{}translate(0,$\{yScale(0)\})\textasciigrave{}).call(d3.axisBottom(xScale).ticks(10));}
\NormalTok{  g.append("g").call(d3.axisLeft(yScale).ticks(6));}

\NormalTok{  // Axis labels}
\NormalTok{  g.append("text").attr("x", innerW / 2).attr("y", innerH + 32)}
\NormalTok{    .attr("text{-}anchor", "middle").text("t (s)");}
\NormalTok{  g.append("text")}
\NormalTok{    .attr("transform", "rotate({-}90)")}
\NormalTok{    .attr("x", {-}innerH / 2).attr("y", {-}38)}
\NormalTok{    .attr("text{-}anchor", "middle").text("y(t)");}

\NormalTok{  // Response curve}
\NormalTok{  const lineGen = d3.line()}
\NormalTok{    .x((\_, i) =\textgreater{} xScale(ts[i]))}
\NormalTok{    .y((v) =\textgreater{} yScale(v));}

\NormalTok{  g.append("path")}
\NormalTok{    .datum(yAll)}
\NormalTok{    .attr("d", lineGen)}
\NormalTok{    .attr("fill", "none")}
\NormalTok{    .attr("stroke", "steelblue")}
\NormalTok{    .attr("stroke{-}width", 2.5);}

\NormalTok{  // Damping label}
\NormalTok{  const dampLabel = zeta \textless{} 1 ? \textasciigrave{}Underdamped (ζ = $\{zeta.toFixed(2)\})\textasciigrave{} :}
\NormalTok{                    zeta === 1 ? "Critically damped" :}
\NormalTok{                    \textasciigrave{}Overdamped (ζ = $\{zeta.toFixed(2)\})\textasciigrave{};}
\NormalTok{  g.append("text")}
\NormalTok{    .attr("x", 6).attr("y", 14)}
\NormalTok{    .attr("fill", "steelblue")}
\NormalTok{    .attr("font{-}size", "12px")}
\NormalTok{    .text(dampLabel);}

\NormalTok{  return svg.node();}
\NormalTok{\}}
\end{Highlighting}
\end{Shaded}

\section{The Dirac delta function}\label{the-dirac-delta-function}

A solenoid strikes a piston at \(t = a\). The contact lasts a fraction
of a millisecond but delivers a definite impulse --- a large force for a
very short time. The Dirac delta function \(\delta(t - a)\) is the
mathematical idealisation: an infinitely concentrated spike at \(t = a\)
with unit area.

Formally, \(\delta(t-a)\) is defined through its sifting property: for
any continuous function \(f\),

\[\int_{-\infty}^\infty f(t)\,\delta(t-a)\,dt = f(a)\]

It is not a function in the classical sense --- it is a distribution.
But for engineering purposes, the sifting property is all we need.

\textbf{Laplace transform of \(\delta(t-a)\)} for \(a > 0\):

\[\mathcal{L}\{\delta(t-a)\} = \int_0^\infty e^{-st}\delta(t-a)\,dt = e^{-as}\]

The sifting property pulls out the factor \(e^{-st}\) evaluated at
\(t = a\). Notably, the transform of an impulse is just an exponential
--- a constant for any given \(a\).

\subsection{Impulse response of a second-order
system}\label{impulse-response-of-a-second-order-system}

\textbf{Problem.} A mechanical oscillator receives a unit impulse at
\(t = 1\):

\[\ddot{x} + 2\dot{x} + 5x = \delta(t-1), \quad x(0) = 0,\ \dot{x}(0) = 0\]

\textbf{Step 1: Transform.}

\[s^2 X + 2sX + 5X = e^{-s}\]

\textbf{Step 2: Initial conditions already zero.}

\[(s^2 + 2s + 5)X(s) = e^{-s}\]

\textbf{Step 3: Solve for \(X(s)\).}

\[X(s) = \frac{e^{-s}}{s^2 + 2s + 5}\]

\textbf{Step 4: Invert using the second shifting theorem.} Complete the
square in the denominator: \(s^2 + 2s + 5 = (s+1)^2 + 4\).

The function without \(e^{-s}\) is \(G(s) = \dfrac{1}{(s+1)^2+4}\).

From the table:
\(\mathcal{L}^{-1}\!\left\{\dfrac{\omega}{(s-a)^2+\omega^2}\right\} = e^{at}\sin(\omega t)\),
so with \(a = -1\) and \(\omega = 2\):

\[g(t) = \frac{1}{2}e^{-t}\sin(2t)\]

By the second shifting theorem:

\[x(t) = g(t-1)\,u(t-1) = \frac{1}{2}e^{-(t-1)}\sin(2(t-1))\,u(t-1)\]

The system is undisturbed until the impulse at \(t = 1\). After the
impulse, it responds with a damped sinusoidal oscillation --- the
system's natural behaviour, triggered instantaneously.

\begin{tcolorbox}[enhanced jigsaw, opacityback=0, colback=white, breakable, colframe=quarto-callout-note-color-frame, arc=.35mm, left=2mm, toprule=.15mm, leftrule=.75mm, bottomrule=.15mm, rightrule=.15mm]
\begin{minipage}[t]{5.5mm}
\textcolor{quarto-callout-note-color}{\faInfo}
\end{minipage}%
\begin{minipage}[t]{\textwidth - 5.5mm}

\textbf{Impulse response and system characterisation}

The impulse response \(h(t)\) --- the output when the input is
\(\delta(t)\) with zero initial conditions --- completely characterises
a linear time-invariant system. Any input \(u(t)\) produces the output
\(\int_0^t h(t-\tau)u(\tau)\,d\tau\). This is the convolution integral,
and it connects directly to the convolution theorem below.

\end{minipage}%
\end{tcolorbox}

\section{Convolution theorem}\label{convolution-theorem}

For two functions \(f\) and \(g\), the convolution is

\[(f * g)(t) = \int_0^t f(t-\tau)\,g(\tau)\,d\tau\]

The convolution theorem states:

\[\mathcal{L}\{f * g\} = F(s)\cdot G(s)\]

Multiplication in the \(s\)-domain corresponds to convolution in the
\(t\)-domain.

\textbf{Transfer function interpretation.} If a system has impulse
response \(h(t)\) with transform \(H(s)\), then for input \(U(s)\) the
output is \(Y(s) = H(s)\,U(s)\). Here \(H(s)\) is the \textbf{transfer
function} of the system. Control engineers characterise systems entirely
through their transfer functions --- poles of \(H(s)\) determine
stability, zeros determine how the system shapes the input signal. The
step from a differential equation to a transfer function is the step
from the \(t\)-domain to the \(s\)-domain via the Laplace transform.

\section{Where this goes}\label{where-this-goes-27}

The Laplace transform is the bridge from time-domain dynamics to the
frequency domain. The next natural step is \textbf{Fourier series} and
the \textbf{Fourier transform}, which decompose periodic or
square-integrable signals into frequency components directly --- without
the convergence factor \(e^{-st}\). The Fourier transform is obtained by
restricting the Laplace transform to the imaginary axis \(s = i\omega\).
That connection, explored in the Fourier chapter, motivates the full
theory of \textbf{complex analysis} (Vol 7): the inverse Laplace
transform is a contour integral in the complex \(s\)-plane, the Bromwich
integral, and understanding it properly requires residue calculus.

For PDEs, the Laplace transform in \(t\) reduces certain parabolic and
hyperbolic PDEs to ODEs in the spatial variable --- an approach that
reuses everything developed here.

\begin{tcolorbox}[enhanced jigsaw, opacityback=0, colback=white, breakable, colframe=quarto-callout-tip-color-frame, arc=.35mm, left=2mm, toprule=.15mm, leftrule=.75mm, bottomrule=.15mm, rightrule=.15mm]
\begin{minipage}[t]{5.5mm}
\textcolor{quarto-callout-tip-color}{\faLightbulb}
\end{minipage}%
\begin{minipage}[t]{\textwidth - 5.5mm}

\textbf{Where this shows up}

\begin{itemize}
\tightlist
\item
  A control engineer designing a PID controller writes the plant as a
  transfer function \(G(s)\) and shapes the closed-loop response in the
  \(s\)-domain --- poles, zeros, Bode plots, all Laplace-transform
  concepts.
\item
  A nuclear engineer modelling a decay chain with time-varying neutron
  flux uses Laplace transforms to handle the discontinuous source term.
\item
  A signal processing engineer designing a digital filter uses the
  Z-transform, the discrete analogue of the Laplace transform:
  \(\mathcal{Z}\{x[n]\} = \sum_{n=0}^\infty x[n]z^{-n}\).
\item
  A quant analyst pricing exotic options uses the characteristic
  function of the asset price process --- effectively the Laplace
  transform of the risk-neutral density --- with fast Fourier transform
  inversion.
\end{itemize}

\end{minipage}%
\end{tcolorbox}

\section{Exercises}\label{exercises-29}

\begin{Shaded}
\begin{Highlighting}[]
\NormalTok{//| echo: false}
\NormalTok{katex = require("katex@0.16")}
\end{Highlighting}
\end{Shaded}

\begin{Shaded}
\begin{Highlighting}[]
\NormalTok{//| echo: false}
\NormalTok{function makeStepperHTML(exerciseNum, steps) \{}
\NormalTok{  let current = 0;}
\NormalTok{  const totalSteps = steps.length;}
\NormalTok{  const container = document.createElement("div");}
\NormalTok{  container.style.cssText = "border:1px solid \#e5e7eb; border{-}radius:8px; padding:1rem 1.25rem; margin:0.75rem 0; font{-}family:inherit;";}
\NormalTok{  const stepsDiv = document.createElement("div");}
\NormalTok{  stepsDiv.style.marginBottom = "0.75rem";}

\NormalTok{  function renderSteps() \{}
\NormalTok{    stepsDiv.innerHTML = "";}
\NormalTok{    for (let i = 0; i \textless{} current; i++) \{}
\NormalTok{      const s = steps[i];}
\NormalTok{      const row = document.createElement("div");}
\NormalTok{      row.style.cssText = "display:grid; grid{-}template{-}columns:220px 1fr; gap:0.35rem 1rem; align{-}items:baseline; padding:0.35rem 0; border{-}top:1px solid \#f3f4f6;";}
\NormalTok{      const opCell = document.createElement("span");}
\NormalTok{      opCell.style.cssText = "font{-}size:0.85em; color:\#6b7280; font{-}style:italic;";}
\NormalTok{      opCell.textContent = s.op;}
\NormalTok{      const eqCell = document.createElement("span");}
\NormalTok{      eqCell.innerHTML = katex.renderToString(s.eq, \{ throwOnError: false, displayMode: false \});}
\NormalTok{      row.appendChild(opCell);}
\NormalTok{      row.appendChild(eqCell);}
\NormalTok{      if (s.note) \{}
\NormalTok{        const noteRow = document.createElement("div");}
\NormalTok{        noteRow.style.cssText = "grid{-}column:2; font{-}size:0.82em; color:\#6b7280; padding{-}bottom:0.2rem;";}
\NormalTok{        noteRow.textContent = s.note;}
\NormalTok{        row.appendChild(noteRow);}
\NormalTok{      \}}
\NormalTok{      stepsDiv.appendChild(row);}
\NormalTok{    \}}
\NormalTok{    if (current === totalSteps) \{}
\NormalTok{      const done = document.createElement("div");}
\NormalTok{      done.style.cssText = "margin{-}top:0.5rem; font{-}size:0.9em; color:\#059669; font{-}weight:500;";}
\NormalTok{      done.textContent = "✓ Solution complete";}
\NormalTok{      stepsDiv.appendChild(done);}
\NormalTok{    \}}
\NormalTok{  \}}

\NormalTok{  const controls = document.createElement("div");}
\NormalTok{  controls.style.cssText = "display:flex; gap:0.5rem; align{-}items:center;";}
\NormalTok{  const nextBtn = document.createElement("button");}
\NormalTok{  nextBtn.textContent = "Next step →";}
\NormalTok{  nextBtn.style.cssText = "padding:0.35rem 0.85rem; border:1px solid \#d1d5db; border{-}radius:4px; background:\#fff; cursor:pointer; font{-}size:0.9em;";}
\NormalTok{  const resetBtn = document.createElement("button");}
\NormalTok{  resetBtn.textContent = "Reset";}
\NormalTok{  resetBtn.style.cssText = "padding:0.35rem 0.75rem; border:1px solid \#e5e7eb; border{-}radius:4px; background:\#f9fafb; cursor:pointer; font{-}size:0.9em; color:\#6b7280;";}
\NormalTok{  const counter = document.createElement("span");}
\NormalTok{  counter.style.cssText = "font{-}size:0.82em; color:\#9ca3af; margin{-}left:0.25rem;";}

\NormalTok{  function updateButtons() \{}
\NormalTok{    nextBtn.disabled = current === totalSteps;}
\NormalTok{    nextBtn.style.opacity = current === totalSteps ? "0.4" : "1";}
\NormalTok{    counter.textContent = current === 0 ? "Click to reveal steps" : \textasciigrave{}Step $\{current\} of $\{totalSteps\}\textasciigrave{};}
\NormalTok{  \}}

\NormalTok{  nextBtn.onclick = () =\textgreater{} \{ if (current \textless{} totalSteps) \{ current++; renderSteps(); updateButtons(); \} \};}
\NormalTok{  resetBtn.onclick = () =\textgreater{} \{ current = 0; renderSteps(); updateButtons(); \};}
\NormalTok{  controls.appendChild(nextBtn);}
\NormalTok{  controls.appendChild(resetBtn);}
\NormalTok{  controls.appendChild(counter);}
\NormalTok{  container.appendChild(stepsDiv);}
\NormalTok{  container.appendChild(controls);}
\NormalTok{  renderSteps();}
\NormalTok{  updateButtons();}
\NormalTok{  return container;}
\NormalTok{\}}
\end{Highlighting}
\end{Shaded}

\textbf{1.} Compute \(\mathcal{L}\{3e^{2t} - 4\sin(3t)\}\) using the
table and linearity.

\begin{Shaded}
\begin{Highlighting}[]
\NormalTok{//| echo: false}
\NormalTok{makeStepperHTML(1, [}
\NormalTok{  \{ op: "Apply linearity",}
\NormalTok{    eq: "\textbackslash{}\textbackslash{}mathcal\{L\}\textbackslash{}\textbackslash{}\{3e\^{}\{2t\} {-} 4\textbackslash{}\textbackslash{}sin(3t)\textbackslash{}\textbackslash{}\} = 3\textbackslash{}\textbackslash{}mathcal\{L\}\textbackslash{}\textbackslash{}\{e\^{}\{2t\}\textbackslash{}\textbackslash{}\} {-} 4\textbackslash{}\textbackslash{}mathcal\{L\}\textbackslash{}\textbackslash{}\{\textbackslash{}\textbackslash{}sin(3t)\textbackslash{}\textbackslash{}\}" \},}
\NormalTok{  \{ op: "Read e\^{}\{2t\} from table",}
\NormalTok{    eq: "\textbackslash{}\textbackslash{}mathcal\{L\}\textbackslash{}\textbackslash{}\{e\^{}\{2t\}\textbackslash{}\textbackslash{}\} = \textbackslash{}\textbackslash{}frac\{1\}\{s{-}2\}",}
\NormalTok{    note: "Standard entry: ℒ\{e\^{}\{at\}\} = 1/(s−a), here a = 2." \},}
\NormalTok{  \{ op: "Read sin(3t) from table",}
\NormalTok{    eq: "\textbackslash{}\textbackslash{}mathcal\{L\}\textbackslash{}\textbackslash{}\{\textbackslash{}\textbackslash{}sin(3t)\textbackslash{}\textbackslash{}\} = \textbackslash{}\textbackslash{}frac\{3\}\{s\^{}2+9\}",}
\NormalTok{    note: "Standard entry: ℒ\{sin(ωt)\} = ω/(s²+ω²), here ω = 3." \},}
\NormalTok{  \{ op: "Combine",}
\NormalTok{    eq: "\textbackslash{}\textbackslash{}frac\{3\}\{s{-}2\} {-} \textbackslash{}\textbackslash{}frac\{12\}\{s\^{}2+9\}" \}}
\NormalTok{])}
\end{Highlighting}
\end{Shaded}

\begin{center}\rule{0.5\linewidth}{0.5pt}\end{center}

\textbf{2.} Solve \(y' + 3y = 6\), \(y(0) = 1\).

\begin{Shaded}
\begin{Highlighting}[]
\NormalTok{//| echo: false}
\NormalTok{makeStepperHTML(2, [}
\NormalTok{  \{ op: "Transform both sides",}
\NormalTok{    eq: "sY {-} 1 + 3Y = \textbackslash{}\textbackslash{}frac\{6\}\{s\}",}
\NormalTok{    note: "ℒ\{y\textquotesingle{}\} = sY − y(0) = sY − 1. Right side: ℒ\{6\} = 6/s." \},}
\NormalTok{  \{ op: "Collect Y terms",}
\NormalTok{    eq: "(s+3)Y = \textbackslash{}\textbackslash{}frac\{6\}\{s\} + 1 = \textbackslash{}\textbackslash{}frac\{s+6\}\{s\}" \},}
\NormalTok{  \{ op: "Solve for Y(s)",}
\NormalTok{    eq: "Y(s) = \textbackslash{}\textbackslash{}frac\{s+6\}\{s(s+3)\}" \},}
\NormalTok{  \{ op: "Partial fractions: A/s + B/(s+3)",}
\NormalTok{    eq: "A = 2,\textbackslash{}\textbackslash{}quad B = {-}1 \textbackslash{}\textbackslash{}implies Y(s) = \textbackslash{}\textbackslash{}frac\{2\}\{s\} {-} \textbackslash{}\textbackslash{}frac\{1\}\{s+3\}",}
\NormalTok{    note: "Set s=0: 6=3A; set s=−3: 3=−3B." \},}
\NormalTok{  \{ op: "Invert term by term",}
\NormalTok{    eq: "y(t) = 2 {-} e\^{}\{{-}3t\}" \}}
\NormalTok{])}
\end{Highlighting}
\end{Shaded}

\begin{center}\rule{0.5\linewidth}{0.5pt}\end{center}

\textbf{3.} Solve \(y'' + 4y = \sin(2t)\), \(y(0) = 0\), \(y'(0) = 1\).

\begin{Shaded}
\begin{Highlighting}[]
\NormalTok{//| echo: false}
\NormalTok{makeStepperHTML(3, [}
\NormalTok{  \{ op: "Transform both sides",}
\NormalTok{    eq: "s\^{}2 Y {-} s\textbackslash{}\textbackslash{}cdot 0 {-} 1 + 4Y = \textbackslash{}\textbackslash{}frac\{2\}\{s\^{}2+4\}",}
\NormalTok{    note: "ℒ\{y\textquotesingle{}\textquotesingle{}\} = s²Y − sy(0) − y\textquotesingle{}(0) = s²Y − 1." \},}
\NormalTok{  \{ op: "Collect and solve for Y",}
\NormalTok{    eq: "(s\^{}2+4)Y = \textbackslash{}\textbackslash{}frac\{2\}\{s\^{}2+4\} + 1 \textbackslash{}\textbackslash{}implies Y(s) = \textbackslash{}\textbackslash{}frac\{1\}\{s\^{}2+4\} + \textbackslash{}\textbackslash{}frac\{2\}\{(s\^{}2+4)\^{}2\}" \},}
\NormalTok{  \{ op: "Invert first term",}
\NormalTok{    eq: "\textbackslash{}\textbackslash{}mathcal\{L\}\^{}\{{-}1\}\textbackslash{}\textbackslash{}!\textbackslash{}\textbackslash{}left\textbackslash{}\textbackslash{}\{\textbackslash{}\textbackslash{}frac\{1\}\{s\^{}2+4\}\textbackslash{}\textbackslash{}right\textbackslash{}\textbackslash{}\} = \textbackslash{}\textbackslash{}frac\{1\}\{2\}\textbackslash{}\textbackslash{}sin(2t)",}
\NormalTok{    note: "ℒ\{sin(ωt)\} = ω/(s²+ω²) with ω=2." \},}
\NormalTok{  \{ op: "Invert second term",}
\NormalTok{    eq: "\textbackslash{}\textbackslash{}mathcal\{L\}\^{}\{{-}1\}\textbackslash{}\textbackslash{}!\textbackslash{}\textbackslash{}left\textbackslash{}\textbackslash{}\{\textbackslash{}\textbackslash{}frac\{2\}\{(s\^{}2+4)\^{}2\}\textbackslash{}\textbackslash{}right\textbackslash{}\textbackslash{}\} = \textbackslash{}\textbackslash{}frac\{1\}\{8\}(\textbackslash{}\textbackslash{}sin 2t {-} 2t\textbackslash{}\textbackslash{}cos 2t)",}
\NormalTok{    note: "Using ℒ\{sin(ωt) − ωt cos(ωt)\} = 2ω³/(s²+ω²)² with ω=2; scale: 2/(s²+4)² = (1/16)·(2·8)/(s²+4)², giving factor 1/8." \},}
\NormalTok{  \{ op: "Combine",}
\NormalTok{    eq: "y(t) = \textbackslash{}\textbackslash{}tfrac\{5\}\{8\}\textbackslash{}\textbackslash{}sin(2t) {-} \textbackslash{}\textbackslash{}tfrac\{1\}\{4\}t\textbackslash{}\textbackslash{}cos(2t)" \}}
\NormalTok{])}
\end{Highlighting}
\end{Shaded}

\begin{center}\rule{0.5\linewidth}{0.5pt}\end{center}

\textbf{4.} Write
\(f(t) = \begin{cases}0 & t<2\\t-2 & t\geq 2\end{cases}\) in terms of
\(u(t-2)\) and find \(\mathcal{L}\{f(t)\}\).

\begin{Shaded}
\begin{Highlighting}[]
\NormalTok{//| echo: false}
\NormalTok{makeStepperHTML(4, [}
\NormalTok{  \{ op: "Express using Heaviside",}
\NormalTok{    eq: "f(t) = (t{-}2)\textbackslash{}\textbackslash{},u(t{-}2)",}
\NormalTok{    note: "The factor u(t−2) switches the function on at t=2." \},}
\NormalTok{  \{ op: "Identify g(t)",}
\NormalTok{    eq: "f(t) = g(t{-}2)\textbackslash{}\textbackslash{},u(t{-}2) \textbackslash{}\textbackslash{}text\{ where \} g(t) = t" \},}
\NormalTok{  \{ op: "Apply second shifting theorem",}
\NormalTok{    eq: "\textbackslash{}\textbackslash{}mathcal\{L\}\textbackslash{}\textbackslash{}\{g(t{-}2)\textbackslash{}\textbackslash{},u(t{-}2)\textbackslash{}\textbackslash{}\} = e\^{}\{{-}2s\}\textbackslash{}\textbackslash{}mathcal\{L\}\textbackslash{}\textbackslash{}\{g(t)\textbackslash{}\textbackslash{}\} = e\^{}\{{-}2s\}\textbackslash{}\textbackslash{}mathcal\{L\}\textbackslash{}\textbackslash{}\{t\textbackslash{}\textbackslash{}\}",}
\NormalTok{    note: "ℒ\{f(t−a)u(t−a)\} = e\^{}\{−as\}F(s)." \},}
\NormalTok{  \{ op: "Read ℒ\{t\} from table",}
\NormalTok{    eq: "\textbackslash{}\textbackslash{}mathcal\{L\}\textbackslash{}\textbackslash{}\{t\textbackslash{}\textbackslash{}\} = \textbackslash{}\textbackslash{}frac\{1\}\{s\^{}2\}" \},}
\NormalTok{  \{ op: "Final answer",}
\NormalTok{    eq: "\textbackslash{}\textbackslash{}mathcal\{L\}\textbackslash{}\textbackslash{}\{f(t)\textbackslash{}\textbackslash{}\} = \textbackslash{}\textbackslash{}frac\{e\^{}\{{-}2s\}\}\{s\^{}2\}" \}}
\NormalTok{])}
\end{Highlighting}
\end{Shaded}

\begin{center}\rule{0.5\linewidth}{0.5pt}\end{center}

\textbf{5.} Solve \(y'' + y = u(t - \pi)\), \(y(0) = 0\), \(y'(0) = 0\).
The system is at rest; a step input arrives at \(t = \pi\).

\begin{Shaded}
\begin{Highlighting}[]
\NormalTok{//| echo: false}
\NormalTok{makeStepperHTML(5, [}
\NormalTok{  \{ op: "Transform (zero ICs)",}
\NormalTok{    eq: "s\^{}2 Y + Y = \textbackslash{}\textbackslash{}frac\{e\^{}\{{-}\textbackslash{}\textbackslash{}pi s\}\}\{s\}",}
\NormalTok{    note: "Both initial conditions are zero, so no boundary terms." \},}
\NormalTok{  \{ op: "Solve for Y(s)",}
\NormalTok{    eq: "Y(s) = \textbackslash{}\textbackslash{}frac\{e\^{}\{{-}\textbackslash{}\textbackslash{}pi s\}\}\{s(s\^{}2+1)\}" \},}
\NormalTok{  \{ op: "Partial fractions on G(s) = 1/[s(s²+1)]",}
\NormalTok{    eq: "\textbackslash{}\textbackslash{}frac\{1\}\{s(s\^{}2+1)\} = \textbackslash{}\textbackslash{}frac\{1\}\{s\} {-} \textbackslash{}\textbackslash{}frac\{s\}\{s\^{}2+1\}",}
\NormalTok{    note: "Check: 1/(s·…) at s=0 gives A=1; equate s² gives B=−1." \},}
\NormalTok{  \{ op: "Invert G(s)",}
\NormalTok{    eq: "g(t) = 1 {-} \textbackslash{}\textbackslash{}cos t" \},}
\NormalTok{  \{ op: "Apply second shifting theorem",}
\NormalTok{    eq: "y(t) = [1 {-} \textbackslash{}\textbackslash{}cos(t{-}\textbackslash{}\textbackslash{}pi)]\textbackslash{}\textbackslash{},u(t{-}\textbackslash{}\textbackslash{}pi)" \},}
\NormalTok{  \{ op: "Simplify: cos(t−π) = −cos t",}
\NormalTok{    eq: "y(t) = [1 + \textbackslash{}\textbackslash{}cos t]\textbackslash{}\textbackslash{},u(t{-}\textbackslash{}\textbackslash{}pi)",}
\NormalTok{    note: "Using cos(t−π) = −cos t, from the cosine addition formula. For t \textgreater{} π: y oscillates about 1. Before t=π: y=0." \}}
\NormalTok{])}
\end{Highlighting}
\end{Shaded}

\begin{center}\rule{0.5\linewidth}{0.5pt}\end{center}

\textbf{6.} A mechanical system receives a unit impulse at \(t = 1\):
\(\ddot{x} + 2\dot{x} + 5x = \delta(t-1)\), \(x(0) = 0\),
\(\dot{x}(0) = 0\). Find \(x(t)\) for \(t > 1\).

\begin{Shaded}
\begin{Highlighting}[]
\NormalTok{//| echo: false}
\NormalTok{makeStepperHTML(6, [}
\NormalTok{  \{ op: "Transform: ℒ\{δ(t−1)\} = e\^{}\{−s\}",}
\NormalTok{    eq: "(s\^{}2 + 2s + 5)X = e\^{}\{{-}s\}",}
\NormalTok{    note: "Zero initial conditions; impulse transform is e\^{}\{−as\} with a=1." \},}
\NormalTok{  \{ op: "Solve for X(s)",}
\NormalTok{    eq: "X(s) = \textbackslash{}\textbackslash{}frac\{e\^{}\{{-}s\}\}\{s\^{}2+2s+5\}" \},}
\NormalTok{  \{ op: "Complete the square",}
\NormalTok{    eq: "s\^{}2+2s+5 = (s+1)\^{}2+4" \},}
\NormalTok{  \{ op: "Invert G(s) = 1/[(s+1)²+4]",}
\NormalTok{    eq: "g(t) = \textbackslash{}\textbackslash{}frac\{1\}\{2\}e\^{}\{{-}t\}\textbackslash{}\textbackslash{}sin(2t)",}
\NormalTok{    note: "Match to ℒ\{e\^{}\{at\}sin(ωt)\} = ω/[(s−a)²+ω²] with a=−1, ω=2." \},}
\NormalTok{  \{ op: "Apply second shifting theorem",}
\NormalTok{    eq: "x(t) = \textbackslash{}\textbackslash{}tfrac\{1\}\{2\}e\^{}\{{-}(t{-}1)\}\textbackslash{}\textbackslash{}sin\textbackslash{}\textbackslash{}bigl(2(t{-}1)\textbackslash{}\textbackslash{}bigr)\textbackslash{}\textbackslash{},u(t{-}1)",}
\NormalTok{    note: "For t \textgreater{} 1: damped sinusoid with natural frequency 2 rad/s, decay rate 1." \}}
\NormalTok{])}
\end{Highlighting}
\end{Shaded}

\part{Volume 7 --- Linear algebra and vector calculus}

\chapter{Matrices and linear systems}\label{matrices-and-linear-systems}

Organising equations so machines --- and humans --- can solve them

\hfill\break

The structural analysis of the Tacoma Narrows Bridge involved linear
systems of equations --- the same fundamental form as any system of
resistors, any truss, any set of chemical mass-balance equations.
Thousands of numbers, hundreds of equations, but the same underlying
question every time: does this system have a solution, is it unique, and
what is it?

For small systems --- two or three equations --- you can solve by
substitution, the way you learned in school. For larger systems, that
approach bogs down in bookkeeping errors. The method in this chapter,
Gaussian elimination, is systematic substitution: a sequence of
operations that any careful person (or machine) can apply without
creativity, and that always terminates with a definitive answer.

Before any of that, we need the language: matrices and the algebra of
matrix operations. The notation is not decoration --- it compresses a
system of twenty equations into a single line, and makes the structure
of the solution visible.

\begin{center}\rule{0.5\linewidth}{0.5pt}\end{center}

\section{Matrices and operations}\label{matrices-and-operations}

A \textbf{matrix} is a rectangular array of numbers. The matrix
\(\mathbf{A}\) with \(m\) rows and \(n\) columns is called an
\(m \times n\) matrix:

\[\mathbf{A} = \begin{pmatrix} a_{11} & a_{12} & \cdots & a_{1n} \\ a_{21} & a_{22} & \cdots & a_{2n} \\ \vdots & \vdots & \ddots & \vdots \\ a_{m1} & a_{m2} & \cdots & a_{mn} \end{pmatrix}\]

The element in row \(i\) and column \(j\) is \(a_{ij}\). The first
subscript is always the row, the second is always the column. A matrix
with one column (\(n = 1\)) is a \textbf{column vector}; one with one
row (\(m = 1\)) is a \textbf{row vector}.

\subsection{Special types}\label{special-types}

\textbf{Square matrix.} \(m = n\). The elements
\(a_{11}, a_{22}, \ldots, a_{nn}\) form the \textbf{main diagonal}.

\textbf{Symmetric matrix.} \(a_{ij} = a_{ji}\) for all \(i, j\).
Equivalently, \(\mathbf{A} = \mathbf{A}^T\). Symmetric matrices arise
naturally: the stiffness matrix of a structure is symmetric, as is any
correlation matrix in statistics.

\textbf{Diagonal matrix.} All off-diagonal entries are zero:
\(a_{ij} = 0\) whenever \(i \neq j\). Multiplication by a diagonal
matrix scales each component independently.

\textbf{Identity matrix \(\mathbf{I}_n\).} Diagonal with all diagonal
entries equal to 1. It plays the role of 1 in matrix arithmetic:
\(\mathbf{A}\mathbf{I} = \mathbf{I}\mathbf{A} = \mathbf{A}\).

\textbf{Triangular matrices.} A matrix is \textbf{upper triangular} if
\(a_{ij} = 0\) whenever \(i > j\) --- all entries below the main
diagonal are zero. It is \textbf{lower triangular} if \(a_{ij} = 0\)
whenever \(i < j\). Gaussian elimination produces an upper triangular
matrix; this is why triangular forms matter.

\textbf{Transpose.} The transpose \(\mathbf{A}^T\) swaps rows and
columns: \((\mathbf{A}^T)_{ij} = a_{ji}\). A \(3 \times 2\) matrix
transposes to a \(2 \times 3\) matrix. Key identities:
\((\mathbf{A}^T)^T = \mathbf{A}\);
\((\mathbf{A}\mathbf{B})^T = \mathbf{B}^T \mathbf{A}^T\) (order
reverses).

\subsection{Matrix addition and scalar
multiplication}\label{matrix-addition-and-scalar-multiplication}

Two matrices of the same size can be added element by element:

\[(\mathbf{A} + \mathbf{B})_{ij} = a_{ij} + b_{ij}\]

Scalar multiplication scales every entry:
\((\lambda \mathbf{A})_{ij} = \lambda a_{ij}\).

These operations satisfy the expected algebraic properties ---
commutativity, associativity, distributivity. Nothing surprising here.

\subsection{Matrix multiplication}\label{matrix-multiplication}

Multiplying two matrices is less obvious. The product
\(\mathbf{C} = \mathbf{A}\mathbf{B}\) requires that the number of
columns of \(\mathbf{A}\) equals the number of rows of \(\mathbf{B}\).
If \(\mathbf{A}\) is \(m \times n\) and \(\mathbf{B}\) is
\(n \times p\), then \(\mathbf{C}\) is \(m \times p\).

The \((i, j)\) entry of \(\mathbf{C}\) is the dot product of row \(i\)
of \(\mathbf{A}\) with column \(j\) of \(\mathbf{B}\):

\[c_{ij} = \sum_{k=1}^{n} a_{ik}\,b_{kj}\]

\textbf{Worked example.} Let

\[\mathbf{A} = \begin{pmatrix} 2 & 1 \\ 0 & 3 \end{pmatrix}, \qquad \mathbf{B} = \begin{pmatrix} 1 & 4 \\ 2 & 1 \end{pmatrix}\]

Then \(c_{11} = 2 \cdot 1 + 1 \cdot 2 = 4\);
\(c_{12} = 2 \cdot 4 + 1 \cdot 1 = 9\);
\(c_{21} = 0 \cdot 1 + 3 \cdot 2 = 6\);
\(c_{22} = 0 \cdot 4 + 3 \cdot 1 = 3\).

\[\mathbf{A}\mathbf{B} = \begin{pmatrix} 4 & 9 \\ 6 & 3 \end{pmatrix}\]

\textbf{Important:} matrix multiplication is not commutative.
\(\mathbf{A}\mathbf{B} \neq \mathbf{B}\mathbf{A}\) in general --- and
sometimes \(\mathbf{B}\mathbf{A}\) is not even defined when
\(\mathbf{A}\mathbf{B}\) is. This is a genuine departure from ordinary
arithmetic, and it matters whenever you are rearranging matrix
equations.

\textbf{What multiplication means geometrically.} Each column of
\(\mathbf{A}\mathbf{B}\) is a linear combination of the columns of
\(\mathbf{A}\), with coefficients from the corresponding column of
\(\mathbf{B}\). This column-space interpretation is the key to
understanding why the structure of \(\mathbf{A}\) determines the
solvability of \(\mathbf{A}\mathbf{x} = \mathbf{b}\).

\begin{center}\rule{0.5\linewidth}{0.5pt}\end{center}

\section{Linear systems: Ax = b}\label{linear-systems-ax-b}

A \textbf{system of \(m\) linear equations in \(n\) unknowns} has the
form

\[\begin{aligned} a_{11}x_1 + a_{12}x_2 + \cdots + a_{1n}x_n &= b_1 \\ a_{21}x_1 + a_{22}x_2 + \cdots + a_{2n}x_n &= b_2 \\ &\vdots \\ a_{m1}x_1 + a_{m2}x_2 + \cdots + a_{mn}x_n &= b_m \end{aligned}\]

In matrix notation this is \(\mathbf{A}\mathbf{x} = \mathbf{b}\), where
\(\mathbf{A}\) is the \(m \times n\) \textbf{coefficient matrix},
\(\mathbf{x}\) is the \(n \times 1\) vector of unknowns, and
\(\mathbf{b}\) is the \(m \times 1\) right-hand side vector.

The compactness of this notation is not cosmetic. Once you write a
system as \(\mathbf{A}\mathbf{x} = \mathbf{b}\), all the structure of
the solution is encoded in \(\mathbf{A}\) --- its rank, its column
space, its null space. You can reason about solvability without touching
the specific values in \(\mathbf{b}\).

\subsection{The augmented matrix}\label{the-augmented-matrix}

The working object for solving a linear system is the \textbf{augmented
matrix} \([\mathbf{A}\,|\,\mathbf{b}]\): the coefficient matrix with
\(\mathbf{b}\) appended as an extra column, separated by a vertical bar.

\[[\mathbf{A}\,|\,\mathbf{b}] = \left(\begin{array}{ccc|c} a_{11} & \cdots & a_{1n} & b_1 \\ \vdots & \ddots & \vdots & \vdots \\ a_{m1} & \cdots & a_{mn} & b_m \end{array}\right)\]

Working on the augmented matrix instead of the equations directly keeps
things compact and avoids rewriting variable names at every step.

\subsection{What a solution means
geometrically}\label{what-a-solution-means-geometrically}

Each equation \(a_{i1}x_1 + a_{i2}x_2 + \cdots + a_{in}x_n = b_i\) is a
hyperplane in \(\mathbb{R}^n\). Solving the system means finding the
intersection of all \(m\) hyperplanes.

For \(n = 2\) (two unknowns): each equation is a line in the plane. Two
lines either intersect at one point (unique solution), are parallel (no
solution), or are the same line (infinitely many solutions).

For \(n = 3\): each equation is a plane in three-dimensional space.
Three planes can intersect at one point, at a line, at a plane, or not
at all. Gaussian elimination finds which case applies and produces the
solution in each case.

\begin{center}\rule{0.5\linewidth}{0.5pt}\end{center}

\section{Gaussian elimination}\label{gaussian-elimination}

\textbf{Row operations.} The three operations that can be applied to
rows of an augmented matrix without changing the solution set are:

\begin{enumerate}
\def\labelenumi{\arabic{enumi}.}
\tightlist
\item
  \textbf{Swap} two rows: \(R_i \leftrightarrow R_j\)
\item
  \textbf{Scale} a row by a nonzero constant:
  \(R_i \leftarrow \lambda R_i\) (\(\lambda \neq 0\))
\item
  \textbf{Add a multiple of one row to another}:
  \(R_i \leftarrow R_i + \mu R_j\)
\end{enumerate}

These are called \textbf{elementary row operations}. They are
reversible, so they transform the system into an equivalent one --- same
solution set, different representation.

\textbf{The algorithm.} Gaussian elimination applies row operations to
reduce \([\mathbf{A}\,|\,\mathbf{b}]\) to \textbf{upper triangular}
(also called \textbf{row echelon}) form, where all entries below the
main diagonal are zero. Back-substitution then extracts the unknowns
from bottom to top.

\subsection{Worked example: 3×3 system}\label{worked-example-33-system}

Solve:

\[\begin{aligned} 2x_1 + 4x_2 - 2x_3 &= 2 \\ x_1 + 3x_2 + x_3 &= 10 \\ -x_1 + 2x_2 + 3x_3 &= 8 \end{aligned}\]

\textbf{Step 1.} Write the augmented matrix:

\[\left(\begin{array}{rrr|r} 2 & 4 & -2 & 2 \\ 1 & 3 & 1 & 10 \\ -1 & 2 & 3 & 8 \end{array}\right)\]

\textbf{Step 2.} Use row 1 as the pivot row to eliminate \(x_1\) from
rows 2 and 3.

\(R_2 \leftarrow R_2 - \tfrac{1}{2}R_1\):

\[\left(\begin{array}{rrr|r} 2 & 4 & -2 & 2 \\ 0 & 1 & 2 & 9 \\ -1 & 2 & 3 & 8 \end{array}\right)\]

\(R_3 \leftarrow R_3 + \tfrac{1}{2}R_1\):

\[\left(\begin{array}{rrr|r} 2 & 4 & -2 & 2 \\ 0 & 1 & 2 & 9 \\ 0 & 4 & 2 & 9 \end{array}\right)\]

\textbf{Step 3.} Use row 2 as the pivot row to eliminate \(x_2\) from
row 3.

\(R_3 \leftarrow R_3 - 4R_2\):

\[\left(\begin{array}{rrr|r} 2 & 4 & -2 & 2 \\ 0 & 1 & 2 & 9 \\ 0 & 0 & -6 & -27 \end{array}\right)\]

This is upper triangular form. The system now reads:

\[\begin{aligned} 2x_1 + 4x_2 - 2x_3 &= 2 \\ x_2 + 2x_3 &= 9 \\ -6x_3 &= -27 \end{aligned}\]

\textbf{Step 4. Back-substitution.}

From row 3: \(x_3 = \tfrac{-27}{-6} = \tfrac{9}{2}\).

From row 2: \(x_2 = 9 - 2 \cdot \tfrac{9}{2} = 9 - 9 = 0\).

From row 1: \(2x_1 = 2 - 4(0) + 2(\tfrac{9}{2}) = 2 + 9 = 11\), so
\(x_1 = \tfrac{11}{2}\).

Solution: \(x_1 = \tfrac{11}{2}\), \(x_2 = 0\), \(x_3 = \tfrac{9}{2}\).

\subsection{The pivot}\label{the-pivot}

At each elimination step, the diagonal element used to eliminate the
column below it is called the \textbf{pivot}. The pivot must be nonzero.
If it is zero, swap the current row with any row below it that has a
nonzero entry in that column (\textbf{partial pivoting}). If no such row
exists, the matrix is singular --- the system either has no solution or
infinitely many.

In practice, numerical software always uses partial pivoting even when
the pivot is technically nonzero, to avoid large round-off errors from
dividing by very small numbers.

\emph{Step through the elimination --- the amber cell is the active
pivot.}

\begin{Shaded}
\begin{Highlighting}[]
\NormalTok{//| echo: false}
\NormalTok{\{}
\NormalTok{  const states = [}
\NormalTok{    \{}
\NormalTok{      label: "Initial matrix — pivot x₁ = 2 (amber) will zero the column below it.",}
\NormalTok{      op: null,}
\NormalTok{      M: [[2,4,{-}2,2],[1,3,1,10],[{-}1,2,3,8]],}
\NormalTok{      pr: 0, pc: 0}
\NormalTok{    \},}
\NormalTok{    \{}
\NormalTok{      label: "x₁ cleared from row 2. Pivot still active — row 3 still needs clearing.",}
\NormalTok{      op: "R₂  ←  R₂ − ½R₁    (m₂₁ = ½)",}
\NormalTok{      M: [[2,4,{-}2,2],[0,1,2,9],[{-}1,2,3,8]],}
\NormalTok{      pr: 0, pc: 0}
\NormalTok{    \},}
\NormalTok{    \{}
\NormalTok{      label: "Column 1 fully cleared. Pivot shifts to x₂ = 1 in row 2.",}
\NormalTok{      op: "R₃  ←  R₃ + ½R₁    (m₃₁ = −½)",}
\NormalTok{      M: [[2,4,{-}2,2],[0,1,2,9],[0,4,2,9]],}
\NormalTok{      pr: 1, pc: 1}
\NormalTok{    \},}
\NormalTok{    \{}
\NormalTok{      label: "Upper triangular form reached. Back{-}substitution proceeds from the bottom.",}
\NormalTok{      op: "R₃  ←  R₃ − 4R₂    (m₃₂ = 4)",}
\NormalTok{      M: [[2,4,{-}2,2],[0,1,2,9],[0,0,{-}6,{-}27]],}
\NormalTok{      pr: 2, pc: 2}
\NormalTok{    \}}
\NormalTok{  ];}

\NormalTok{  let step = 0;}

\NormalTok{  const wrap = document.createElement("div");}
\NormalTok{  wrap.style.cssText = "border:1px solid \#e5e7eb;border{-}radius:8px;padding:1rem 1.25rem;margin:0.5rem 0;max{-}width:480px;font{-}family:inherit;";}

\NormalTok{  const opEl = document.createElement("div");}
\NormalTok{  opEl.style.cssText = "font{-}size:0.9em;font{-}family:monospace;color:\#374151;min{-}height:1.5em;margin{-}bottom:0.6rem;white{-}space:nowrap;";}

\NormalTok{  const matWrap = document.createElement("div");}
\NormalTok{  matWrap.style.cssText = "margin{-}bottom:0.5rem;";}

\NormalTok{  const lblEl = document.createElement("div");}
\NormalTok{  lblEl.style.cssText = "font{-}size:0.82em;color:\#6b7280;min{-}height:1.4em;margin{-}bottom:0.75rem;";}

\NormalTok{  function renderMatrix(M, pr, pc) \{}
\NormalTok{    const nc = M[0].length;}
\NormalTok{    let html = \textquotesingle{}\textless{}div style="display:inline{-}flex;align{-}items:stretch;"\textgreater{}\textquotesingle{};}
\NormalTok{    html += \textquotesingle{}\textless{}div style="border{-}left:2px solid \#374151;border{-}top:2px solid \#374151;border{-}bottom:2px solid \#374151;width:5px;border{-}radius:3px 0 0 3px;margin{-}right:2px;"\textgreater{}\textless{}/div\textgreater{}\textquotesingle{};}
\NormalTok{    html += \textquotesingle{}\textless{}table style="border{-}collapse:collapse;font{-}size:1em;line{-}height:1.8;"\textgreater{}\textquotesingle{};}
\NormalTok{    for (let r = 0; r \textless{} M.length; r++) \{}
\NormalTok{      html += "\textless{}tr\textgreater{}";}
\NormalTok{      for (let c = 0; c \textless{} nc; c++) \{}
\NormalTok{        const v = M[r][c];}
\NormalTok{        const disp = v \textless{} 0 ? "−" + Math.abs(v) : String(v);}
\NormalTok{        const augBorder = c === nc {-} 1 ? "border{-}left:2px solid \#9ca3af;" : "";}
\NormalTok{        const pivBg = r === pr \&\& c === pc ? "background:\#fef3c7;" : "";}
\NormalTok{        const pivFw = r === pr \&\& c === pc ? "font{-}weight:700;color:\#92400e;" : "";}
\NormalTok{        html += \textasciigrave{}\textless{}td style="padding:0.15rem 0.65rem;text{-}align:right;min{-}width:1.8em;$\{augBorder\}$\{pivBg\}$\{pivFw\}"\textgreater{}$\{disp\}\textless{}/td\textgreater{}\textasciigrave{};}
\NormalTok{      \}}
\NormalTok{      html += "\textless{}/tr\textgreater{}";}
\NormalTok{    \}}
\NormalTok{    html += "\textless{}/table\textgreater{}";}
\NormalTok{    html += \textquotesingle{}\textless{}div style="border{-}right:2px solid \#374151;border{-}top:2px solid \#374151;border{-}bottom:2px solid \#374151;width:5px;border{-}radius:0 3px 3px 0;margin{-}left:2px;"\textgreater{}\textless{}/div\textgreater{}\textquotesingle{};}
\NormalTok{    html += \textquotesingle{}\textless{}/div\textgreater{}\textquotesingle{};}
\NormalTok{    return html;}
\NormalTok{  \}}

\NormalTok{  function draw() \{}
\NormalTok{    const s = states[step];}
\NormalTok{    opEl.textContent = s.op || "";}
\NormalTok{    matWrap.innerHTML = renderMatrix(s.M, s.pr, s.pc);}
\NormalTok{    lblEl.textContent = s.label;}
\NormalTok{    nextBtn.disabled = step \textgreater{}= states.length {-} 1;}
\NormalTok{    nextBtn.style.opacity = step \textgreater{}= states.length {-} 1 ? "0.4" : "1";}
\NormalTok{    ctrEl.textContent = step === 0 ? "Click to begin" : \textasciigrave{}Step $\{step\} of $\{states.length {-} 1\}\textasciigrave{};}
\NormalTok{    doneEl.style.display = step === states.length {-} 1 ? "block" : "none";}
\NormalTok{  \}}

\NormalTok{  const ctrlRow = document.createElement("div");}
\NormalTok{  ctrlRow.style.cssText = "display:flex;gap:0.5rem;align{-}items:center;";}

\NormalTok{  const nextBtn = document.createElement("button");}
\NormalTok{  nextBtn.textContent = "Next operation →";}
\NormalTok{  nextBtn.style.cssText = "padding:0.35rem 0.85rem;border:1px solid \#d1d5db;border{-}radius:4px;background:\#fff;cursor:pointer;font{-}size:0.9em;";}

\NormalTok{  const rstBtn = document.createElement("button");}
\NormalTok{  rstBtn.textContent = "Reset";}
\NormalTok{  rstBtn.style.cssText = "padding:0.35rem 0.75rem;border:1px solid \#e5e7eb;border{-}radius:4px;background:\#f9fafb;cursor:pointer;font{-}size:0.9em;color:\#6b7280;";}

\NormalTok{  const ctrEl = document.createElement("span");}
\NormalTok{  ctrEl.style.cssText = "font{-}size:0.82em;color:\#9ca3af;margin{-}left:0.25rem;";}

\NormalTok{  const doneEl = document.createElement("div");}
\NormalTok{  doneEl.style.cssText = "margin{-}top:0.5rem;font{-}size:0.9em;color:\#059669;font{-}weight:500;display:none;";}
\NormalTok{  doneEl.textContent = "Upper triangular form reached. Back{-}substitution proceeds from row 3 upward.";}

\NormalTok{  nextBtn.onclick = () =\textgreater{} \{ if (step \textless{} states.length {-} 1) \{ step++; draw(); \} \};}
\NormalTok{  rstBtn.onclick = () =\textgreater{} \{ step = 0; draw(); \};}

\NormalTok{  ctrlRow.append(nextBtn, rstBtn, ctrEl);}
\NormalTok{  wrap.append(opEl, matWrap, lblEl, ctrlRow, doneEl);}
\NormalTok{  draw();}
\NormalTok{  return wrap;}
\NormalTok{\}}
\end{Highlighting}
\end{Shaded}

\subsection{Gauss-Jordan elimination and
RREF}\label{gauss-jordan-elimination-and-rref}

\textbf{Gauss-Jordan elimination} continues beyond upper triangular
form: it also eliminates entries \emph{above} each pivot, reducing the
matrix to \textbf{reduced row echelon form (RREF)}. In RREF every pivot
is 1, and every other entry in the pivot column is 0. The unknowns can
be read off directly --- no back-substitution needed.

For large systems Gaussian elimination with back-substitution is more
efficient than full Gauss-Jordan. For small systems, and especially for
finding matrix inverses, RREF is convenient.

\textbf{Continuing the worked example.} Starting from the upper
triangular form and scaling:

\(R_3 \leftarrow -\tfrac{1}{6}R_3\), \(R_2 \leftarrow R_2\),
\(R_1 \leftarrow \tfrac{1}{2}R_1\):

\[\left(\begin{array}{rrr|r} 1 & 2 & -1 & 1 \\ 0 & 1 & 2 & 9 \\ 0 & 0 & 1 & 9/2 \end{array}\right)\]

\(R_2 \leftarrow R_2 - 2R_3\):

\[\left(\begin{array}{rrr|r} 1 & 2 & -1 & 1 \\ 0 & 1 & 0 & 0 \\ 0 & 0 & 1 & 9/2 \end{array}\right)\]

\(R_1 \leftarrow R_1 + R_3\), then \(R_1 \leftarrow R_1 - 2R_2\):

\[\left(\begin{array}{rrr|r} 1 & 0 & 0 & 11/2 \\ 0 & 1 & 0 & 0 \\ 0 & 0 & 1 & 9/2 \end{array}\right)\]

The solution \(x_1 = \tfrac{11}{2}\), \(x_2 = 0\),
\(x_3 = \tfrac{9}{2}\) is visible directly in the right-hand column.
Same answer, different path.

\begin{center}\rule{0.5\linewidth}{0.5pt}\end{center}

\section{Solution types and rank}\label{solution-types-and-rank}

Not every system \(\mathbf{A}\mathbf{x} = \mathbf{b}\) has a unique
solution. Three outcomes are possible, and the row-reduced form reveals
which one applies.

\subsection{The three cases}\label{the-three-cases}

\textbf{Case 1: unique solution.} Every column of \(\mathbf{A}\) has a
pivot. After row reduction there is one leading 1 per row and per
column. The solution is determined completely.

\textbf{Case 2: no solution (inconsistent).} After row reduction, at
least one row has the form \([0 \; 0 \; \cdots \; 0 \;|\; c]\) with
\(c \neq 0\). This represents the equation \(0 = c\), which is
impossible. Geometrically: the hyperplanes do not share a common point.

\textbf{Case 3: infinitely many solutions (underdetermined).} After row
reduction, at least one column has no pivot --- the corresponding
variable is \textbf{free}, meaning it can take any value. Each free
variable introduces one dimension of solutions. Geometrically: the
hyperplanes intersect along a line, a plane, or a higher-dimensional
subspace.

\emph{Move the sliders to see how the geometry changes. Try slope −0.50
for the parallel and coincident cases.}

\begin{Shaded}
\begin{Highlighting}[]
\NormalTok{//| echo: false}
\NormalTok{viewof geomSlope = Inputs.range([{-}1.5, 0.5], \{}
\NormalTok{  value: 0.3,}
\NormalTok{  step: 0.05,}
\NormalTok{  label: "Slope of line 2 (m₂)"}
\NormalTok{\})}
\end{Highlighting}
\end{Shaded}

\begin{Shaded}
\begin{Highlighting}[]
\NormalTok{//| echo: false}
\NormalTok{viewof geomIntercept = Inputs.range([{-}1, 5], \{}
\NormalTok{  value: 3.5,}
\NormalTok{  step: 0.1,}
\NormalTok{  label: "y{-}intercept of line 2 (b₂)"}
\NormalTok{\})}
\end{Highlighting}
\end{Shaded}

\begin{Shaded}
\begin{Highlighting}[]
\NormalTok{//| echo: false}
\NormalTok{\{}
\NormalTok{  // Line 1 (fixed): x + 2y = 4  →  y = −0.5x + 2}
\NormalTok{  const m1 = {-}0.5, b1 = 2;}
\NormalTok{  const m2 = geomSlope, b2 = geomIntercept;}

\NormalTok{  const EPS\_S = 0.03;   // slope tolerance (slider step is 0.05)}
\NormalTok{  const EPS\_B = 0.06;   // intercept tolerance (slider step is 0.1)}
\NormalTok{  const sameSlope  = Math.abs(m2 {-} m1) \textless{} EPS\_S;}
\NormalTok{  const coincident = sameSlope \&\& Math.abs(b2 {-} b1) \textless{} EPS\_B;}
\NormalTok{  const parallel   = sameSlope \&\& !coincident;}

\NormalTok{  let caseLabel, caseColor;}
\NormalTok{  if (coincident) \{}
\NormalTok{    caseLabel = "Infinitely many solutions — same line";}
\NormalTok{    caseColor = "\#059669";}
\NormalTok{  \} else if (parallel) \{}
\NormalTok{    caseLabel = "No solution — parallel lines (0 = c, impossible)";}
\NormalTok{    caseColor = "\#dc2626";}
\NormalTok{  \} else \{}
\NormalTok{    const xi = (b2 {-} b1) / (m1 {-} m2);}
\NormalTok{    const yi = m1 * xi + b1;}
\NormalTok{    caseLabel = \textasciigrave{}Unique solution — intersection at ($\{xi.toFixed(2)\}, $\{yi.toFixed(2)\})\textasciigrave{};}
\NormalTok{    caseColor = "\#2563eb";}
\NormalTok{  \}}

\NormalTok{  const W = 340, H = 280, PAD = 44;}
\NormalTok{  const xMin = {-}5, xMax = 5, yMin = {-}3, yMax = 6;}

\NormalTok{  const svg = document.createElementNS("http://www.w3.org/2000/svg", "svg");}
\NormalTok{  svg.setAttribute("width", W);}
\NormalTok{  svg.setAttribute("height", H);}
\NormalTok{  svg.setAttribute("viewBox", \textasciigrave{}0 0 $\{W\} $\{H\}\textasciigrave{});}
\NormalTok{  svg.style.cssText = "display:block;font{-}family:inherit;border:1px solid \#e5e7eb;border{-}radius:6px;";}

\NormalTok{  const ns = "http://www.w3.org/2000/svg";}
\NormalTok{  function mk(tag, attrs) \{}
\NormalTok{    const e = document.createElementNS(ns, tag);}
\NormalTok{    if (attrs) for (const [k, v] of Object.entries(attrs)) e.setAttribute(k, v);}
\NormalTok{    return e;}
\NormalTok{  \}}
\NormalTok{  function mkT(tag, attrs, text) \{ const e = mk(tag, attrs); e.textContent = text; return e; \}}

\NormalTok{  function spx(x) \{ return PAD + (x {-} xMin) / (xMax {-} xMin) * (W {-} 2 * PAD); \}}
\NormalTok{  function spy(y) \{ return H {-} PAD {-} (y {-} yMin) / (yMax {-} yMin) * (H {-} 2 * PAD); \}}

\NormalTok{  const clipId = "geom{-}clip";}
\NormalTok{  const defs = mk("defs");}
\NormalTok{  const cp = mk("clipPath", \{ id: clipId \});}
\NormalTok{  cp.appendChild(mk("rect", \{ x: PAD, y: PAD, width: W {-} 2*PAD, height: H {-} 2*PAD \}));}
\NormalTok{  defs.appendChild(cp);}
\NormalTok{  svg.appendChild(defs);}

\NormalTok{  svg.appendChild(mk("rect", \{ x: 0, y: 0, width: W, height: H, fill: "\#f9fafb", rx: 6 \}));}
\NormalTok{  svg.appendChild(mk("rect", \{ x: PAD, y: PAD, width: W{-}2*PAD, height: H{-}2*PAD, fill: "\#fff", stroke: "\#e5e7eb", "stroke{-}width": 1 \}));}

\NormalTok{  // Grid}
\NormalTok{  const gridG = mk("g", \{ "clip{-}path": \textasciigrave{}url(\#$\{clipId\})\textasciigrave{} \});}
\NormalTok{  for (let x = Math.ceil(xMin); x \textless{}= Math.floor(xMax); x++)}
\NormalTok{    gridG.appendChild(mk("line", \{ x1: spx(x), y1: PAD, x2: spx(x), y2: H{-}PAD, stroke: "\#f3f4f6", "stroke{-}width": 1 \}));}
\NormalTok{  for (let y = Math.ceil(yMin); y \textless{}= Math.floor(yMax); y++)}
\NormalTok{    gridG.appendChild(mk("line", \{ x1: PAD, y1: spy(y), x2: W{-}PAD, y2: spy(y), stroke: "\#f3f4f6", "stroke{-}width": 1 \}));}
\NormalTok{  svg.appendChild(gridG);}

\NormalTok{  // Axes}
\NormalTok{  const clip = \{ "clip{-}path": \textasciigrave{}url(\#$\{clipId\})\textasciigrave{} \};}
\NormalTok{  svg.appendChild(mk("line", \{ x1: PAD, y1: spy(0), x2: W{-}PAD, y2: spy(0), stroke: "\#d1d5db", "stroke{-}width": 1.5, ...clip \}));}
\NormalTok{  svg.appendChild(mk("line", \{ x1: spx(0), y1: PAD, x2: spx(0), y2: H{-}PAD, stroke: "\#d1d5db", "stroke{-}width": 1.5, ...clip \}));}
\NormalTok{  svg.appendChild(mkT("text", \{ x: W{-}PAD+6, y: spy(0)+4, fill: "\#9ca3af", "font{-}size": 12 \}, "x"));}
\NormalTok{  svg.appendChild(mkT("text", \{ x: spx(0)+5, y: PAD{-}4, fill: "\#9ca3af", "font{-}size": 12 \}, "y"));}

\NormalTok{  function addLine(m, b, color, dash, width, opacity) \{}
\NormalTok{    svg.appendChild(mk("line", \{}
\NormalTok{      x1: spx(xMin), y1: spy(m*xMin+b),}
\NormalTok{      x2: spx(xMax), y2: spy(m*xMax+b),}
\NormalTok{      stroke: color, "stroke{-}width": width,}
\NormalTok{      "stroke{-}dasharray": dash || "",}
\NormalTok{      opacity: opacity || 1,}
\NormalTok{      "clip{-}path": \textasciigrave{}url(\#$\{clipId\})\textasciigrave{}}
\NormalTok{    \}));}
\NormalTok{  \}}

\NormalTok{  // Line 1 — fixed, solid blue}
\NormalTok{  addLine(m1, b1, "\#3b82f6", "", 2.5, 1);}
\NormalTok{  const lx1 = 2.5;}
\NormalTok{  svg.appendChild(mkT("text", \{ x: spx(lx1)+4, y: spy(m1*lx1+b1){-}5, fill: "\#3b82f6", "font{-}size": 11, "font{-}weight": "600" \}, "x + 2y = 4"));}

\NormalTok{  // Line 2 — controlled}
\NormalTok{  if (coincident) \{}
\NormalTok{    addLine(m2, b2, "\#059669", "8 4", 5, 0.55);}
\NormalTok{  \} else \{}
\NormalTok{    addLine(m2, b2, "\#f59e0b", parallel ? "" : "7 4", 2.5, 1);}
\NormalTok{    // Label line 2 near its left end}
\NormalTok{    const lx2 = {-}3.5, ly2 = m2*lx2+b2;}
\NormalTok{    if (ly2 \textgreater{}= yMin \&\& ly2 \textless{}= yMax)}
\NormalTok{      svg.appendChild(mkT("text", \{ x: spx(lx2), y: spy(ly2){-}6, fill: "\#f59e0b", "font{-}size": 11, "font{-}weight": "600" \}, "line 2"));}
\NormalTok{  \}}

\NormalTok{  // Intersection dot}
\NormalTok{  if (!parallel \&\& !coincident) \{}
\NormalTok{    const xi = (b2 {-} b1) / (m1 {-} m2);}
\NormalTok{    const yi = m1 * xi + b1;}
\NormalTok{    if (xi \textgreater{}= xMin \&\& xi \textless{}= xMax \&\& yi \textgreater{}= yMin \&\& yi \textless{}= yMax)}
\NormalTok{      svg.appendChild(mk("circle", \{ cx: spx(xi), cy: spy(yi), r: 5, fill: "\#2563eb", stroke: "\#fff", "stroke{-}width": 2 \}));}
\NormalTok{  \}}

\NormalTok{  const footer = document.createElement("div");}
\NormalTok{  footer.style.cssText = \textasciigrave{}font{-}size:0.88em;font{-}weight:600;color:$\{caseColor\};margin{-}top:0.4rem;\textasciigrave{};}
\NormalTok{  footer.textContent = caseLabel;}

\NormalTok{  const container = document.createElement("div");}
\NormalTok{  container.style.marginTop = "0.5rem";}
\NormalTok{  container.append(svg, footer);}
\NormalTok{  return container;}
\NormalTok{\}}
\end{Highlighting}
\end{Shaded}

\subsection{Rank}\label{rank}

The \textbf{rank} of a matrix \(\mathbf{A}\), written
\(\text{rank}(\mathbf{A})\), is the number of pivots in its row-reduced
form --- equivalently, the number of linearly independent rows (or
columns).

For an \(m \times n\) matrix:
\(\text{rank}(\mathbf{A}) \leq \min(m, n)\).

\textbf{Solvability criterion.} The system
\(\mathbf{A}\mathbf{x} = \mathbf{b}\) has at least one solution if and
only if
\(\text{rank}(\mathbf{A}) = \text{rank}([\mathbf{A}\,|\,\mathbf{b}])\).
When this holds:

\begin{itemize}
\tightlist
\item
  If \(\text{rank}(\mathbf{A}) = n\) (all unknowns are pivot variables),
  the solution is unique.
\item
  If \(\text{rank}(\mathbf{A}) < n\), there are
  \(n - \text{rank}(\mathbf{A})\) free variables and infinitely many
  solutions.
\end{itemize}

\subsection{Rank-nullity theorem}\label{rank-nullity-theorem}

For an \(n\)-column matrix \(\mathbf{A}\):

\[\text{rank}(\mathbf{A}) + \text{nullity}(\mathbf{A}) = n\]

The \textbf{nullity} is the number of free variables --- the dimension
of the \textbf{null space} (the set of all solutions to
\(\mathbf{A}\mathbf{x} = \mathbf{0}\)). This theorem says that the
number of ``determined'' directions (rank) plus the number of
``undetermined'' directions (nullity) always adds up to the total number
of unknowns. It is stated here for completeness; it becomes central in
the chapter on eigenvalues.

\subsection{What rank tells you before you
compute}\label{what-rank-tells-you-before-you-compute}

If you have a \(3 \times 3\) system and you can see by inspection that
one equation is a linear combination of the other two (for example,
equation 3 = equation 1 + equation 2), then
\(\text{rank}(\mathbf{A}) \leq 2 < 3\) regardless of \(\mathbf{b}\). The
system is either inconsistent or underdetermined. Checking rank first
can save significant computation.

\begin{center}\rule{0.5\linewidth}{0.5pt}\end{center}

\section{LU factorisation}\label{lu-factorisation}

Gaussian elimination solves one system
\(\mathbf{A}\mathbf{x} = \mathbf{b}\) at a time. In engineering, the
same matrix \(\mathbf{A}\) often appears with many different right-hand
sides \(\mathbf{b}_1, \mathbf{b}_2, \ldots, \mathbf{b}_k\) --- for
example, a structural stiffness matrix under different load patterns.
Re-running Gaussian elimination from scratch each time wastes the work
already done.

\textbf{LU factorisation} stores the elimination steps. It factors
\(\mathbf{A}\) into:

\[\mathbf{A} = \mathbf{L}\mathbf{U}\]

where \(\mathbf{L}\) is a lower triangular matrix (with 1s on the
diagonal) and \(\mathbf{U}\) is the upper triangular matrix produced by
elimination. The elimination multipliers --- the values
\(m_{ij} = a_{ij}/a_{jj}\) used to zero out entries below each pivot ---
become the off-diagonal entries of \(\mathbf{L}\).

\subsection{Solving with LU}\label{solving-with-lu}

Given \(\mathbf{A} = \mathbf{L}\mathbf{U}\), solving
\(\mathbf{A}\mathbf{x} = \mathbf{b}\) becomes two triangular solves:

\begin{enumerate}
\def\labelenumi{\arabic{enumi}.}
\tightlist
\item
  Solve \(\mathbf{L}\mathbf{y} = \mathbf{b}\) for \(\mathbf{y}\) by
  \textbf{forward substitution} (top to bottom, since \(\mathbf{L}\) is
  lower triangular).
\item
  Solve \(\mathbf{U}\mathbf{x} = \mathbf{y}\) for \(\mathbf{x}\) by
  \textbf{back-substitution} (bottom to top, since \(\mathbf{U}\) is
  upper triangular).
\end{enumerate}

Each triangular solve costs \(O(n^2)\) operations. The factorisation
itself costs \(O(n^3)\) --- the same as Gaussian elimination --- but it
only needs to happen once per matrix. Subsequent solves for different
\(\mathbf{b}\) cost only \(O(n^2)\) each.

\subsection{Worked example: 3×3 LU
factorisation}\label{worked-example-33-lu-factorisation}

Factorise:

\[\mathbf{A} = \begin{pmatrix} 2 & 4 & -2 \\ 1 & 3 & 1 \\ -1 & 2 & 3 \end{pmatrix}\]

\textbf{Step 1.} Column 1. Use \(a_{11} = 2\) as pivot.

Multipliers: \(m_{21} = \tfrac{1}{2}\), \(m_{31} = \tfrac{-1}{2}\).

Apply \(R_2 \leftarrow R_2 - \tfrac{1}{2}R_1\) and
\(R_3 \leftarrow R_3 - (-\tfrac{1}{2})R_1 = R_3 + \tfrac{1}{2}R_1\):

\[\begin{pmatrix} 2 & 4 & -2 \\ 0 & 1 & 2 \\ 0 & 4 & 2 \end{pmatrix}\]

\textbf{Step 2.} Column 2. Use \(a_{22}^{(1)} = 1\) as pivot.

Multiplier: \(m_{32} = \tfrac{4}{1} = 4\).

Apply \(R_3 \leftarrow R_3 - 4R_2\):

\[\mathbf{U} = \begin{pmatrix} 2 & 4 & -2 \\ 0 & 1 & 2 \\ 0 & 0 & -6 \end{pmatrix}\]

\textbf{Assembling L.} The multipliers fill the lower triangle of
\(\mathbf{L}\), with 1s on the diagonal:

\[\mathbf{L} = \begin{pmatrix} 1 & 0 & 0 \\ 1/2 & 1 & 0 \\ -1/2 & 4 & 1 \end{pmatrix}\]

\textbf{Verification:} \(\mathbf{L}\mathbf{U}\) should recover
\(\mathbf{A}\). Check the \((3,2)\) entry:
\((-\tfrac{1}{2})(4) + (4)(1) + (1)(0) = -2 + 4 = 2\) ✓.

\subsection{When LU fails}\label{when-lu-fails}

LU factorisation in the form above requires that every pivot encountered
during elimination is nonzero. When a zero pivot arises --- even if the
matrix is not singular --- rows must be swapped before continuing. This
leads to \textbf{PLU factorisation}
(\(\mathbf{P}\mathbf{A} = \mathbf{L}\mathbf{U}\) where \(\mathbf{P}\) is
a permutation matrix recording the row swaps), which is what production
numerical libraries always compute. MATLAB's \texttt{lu()} and NumPy's
\texttt{scipy.linalg.lu()} both return the PLU form.

\subsection{Numerical perspective}\label{numerical-perspective}

For \(n > 4\) or 5, Gaussian elimination and LU factorisation are almost
always done by software. Understanding the algorithm tells you:

\begin{itemize}
\tightlist
\item
  \textbf{When it fails}: a zero pivot signals a singular or
  near-singular matrix --- the system is ill-conditioned and the
  solution, if it exists, may be numerically unreliable.
\item
  \textbf{Why pivoting matters}: without partial pivoting, round-off
  errors compound in a predictable direction and can make the computed
  solution wildly wrong even when a true solution exists.
\item
  \textbf{What to trust}: if NumPy's \texttt{linalg.solve()} raises a
  \texttt{LinAlgError:\ Singular\ matrix}, the coefficient matrix is
  rank-deficient. The fix is not a different solver --- it is a
  reconsideration of the physical model.
\end{itemize}

\begin{center}\rule{0.5\linewidth}{0.5pt}\end{center}

\section{Applications}\label{applications-5}

\subsection{Circuit mesh analysis: Kirchhoff's voltage
law}\label{circuit-mesh-analysis-kirchhoffs-voltage-law}

Consider a three-mesh resistive circuit. Applying Kirchhoff's voltage
law (KVL) around each mesh --- the sum of voltage drops around any
closed loop is zero --- yields one linear equation per mesh in the mesh
currents \(I_1, I_2, I_3\).

For a specific circuit with resistances \(R_1 = 2\,\Omega\),
\(R_2 = 4\,\Omega\), \(R_3 = 3\,\Omega\), \(R_4 = 6\,\Omega\),
\(R_5 = 1\,\Omega\) and voltage sources \(V_1 = 12\,\text{V}\),
\(V_2 = 6\,\text{V}\), the KVL equations might take the form:

\[\begin{aligned} (R_1 + R_2)I_1 - R_2 I_2 &= V_1 \\ -R_2 I_1 + (R_2 + R_3 + R_4)I_2 - R_4 I_3 &= 0 \\ -R_4 I_2 + (R_4 + R_5)I_3 &= -V_2 \end{aligned}\]

Substituting values:

\[\begin{aligned} 6I_1 - 4I_2 &= 12 \\ -4I_1 + 13I_2 - 6I_3 &= 0 \\ -6I_2 + 7I_3 &= -6 \end{aligned}\]

In matrix form \(\mathbf{A}\mathbf{x} = \mathbf{b}\):

\[\begin{pmatrix} 6 & -4 & 0 \\ -4 & 13 & -6 \\ 0 & -6 & 7 \end{pmatrix} \begin{pmatrix} I_1 \\ I_2 \\ I_3 \end{pmatrix} = \begin{pmatrix} 12 \\ 0 \\ -6 \end{pmatrix}\]

The coefficient matrix is symmetric --- a consequence of reciprocity in
passive resistive networks. This is always the case when you apply KVL
systematically; it is not a coincidence to exploit, it is a theorem to
rely on.

Applying Gaussian elimination yields (after row reduction):
\(I_1 = 3.2\,\text{A}\), \(I_2 = 1.8\,\text{A}\),
\(I_3 = 0.686\,\text{A}\) (to three significant figures). Once the mesh
currents are known, every voltage and power in the circuit is computable
--- the linear system is the complete solution to the circuit analysis
problem.

\textbf{Scaling remark.} A practical circuit board may have 50 to 500
mesh loops. The matrix is still symmetric, but \(50 \times 50\) by hand
is not feasible. SPICE --- the standard circuit simulator used by every
electronics engineer --- solves such systems using LU factorisation with
partial pivoting, exactly the method in this chapter.

\subsection{Structural force
distribution}\label{structural-force-distribution}

A statically determinate truss consists of members (bars) and joints
(nodes). At each joint, equilibrium requires that the sum of forces in
the horizontal direction is zero, and the sum in the vertical direction
is zero. Each equilibrium equation is linear in the member forces. For a
truss with \(n\) joints and \(m\) members plus three reaction forces,
static determinacy gives \(m + 3 = 2n\), and the full set of equilibrium
equations is a square linear system.

For a simple three-member triangular truss with a vertical load \(P\) at
the apex:

\[\begin{pmatrix} \cos\theta_1 & \cos\theta_2 & 0 \\ \sin\theta_1 & \sin\theta_2 & 0 \\ 0 & -1 & 1 \end{pmatrix} \begin{pmatrix} F_1 \\ F_2 \\ R \end{pmatrix} = \begin{pmatrix} 0 \\ P \\ 0 \end{pmatrix}\]

where \(F_1, F_2\) are member forces, \(R\) is a reaction force, and
\(\theta_1, \theta_2\) are member angles. Solving this gives the force
in each member --- positive values indicate tension, negative indicate
compression. A member under too much compression will buckle; too much
tension and it will yield. The linear system does not just find numbers:
it tells the structural engineer which members to upsize.

\begin{center}\rule{0.5\linewidth}{0.5pt}\end{center}

\section{Exercises}\label{exercises-30}

These exercises use real engineering and science contexts. The
interesting part is setting up the augmented matrix correctly and
identifying which solution case applies.

\begin{Shaded}
\begin{Highlighting}[]
\NormalTok{//| echo: false}
\NormalTok{katex = require("katex@0.16")}
\end{Highlighting}
\end{Shaded}

\begin{Shaded}
\begin{Highlighting}[]
\NormalTok{//| echo: false}
\NormalTok{function makeStepperHTML(steps) \{}
\NormalTok{  let current = 0;}
\NormalTok{  const totalSteps = steps.length;}
\NormalTok{  const container = document.createElement("div");}
\NormalTok{  container.style.cssText = "border:1px solid \#e5e7eb; border{-}radius:8px; padding:1rem 1.25rem; margin:0.75rem 0; font{-}family:inherit;";}
\NormalTok{  const stepsDiv = document.createElement("div");}
\NormalTok{  stepsDiv.style.marginBottom = "0.75rem";}

\NormalTok{  function renderSteps() \{}
\NormalTok{    stepsDiv.innerHTML = "";}
\NormalTok{    for (let i = 0; i \textless{} current; i++) \{}
\NormalTok{      const s = steps[i];}
\NormalTok{      const row = document.createElement("div");}
\NormalTok{      row.style.cssText = "display:grid; grid{-}template{-}columns:220px 1fr; gap:0.35rem 1rem; align{-}items:baseline; padding:0.35rem 0; border{-}top:1px solid \#f3f4f6;";}
\NormalTok{      const opCell = document.createElement("span");}
\NormalTok{      opCell.style.cssText = "font{-}size:0.85em; color:\#6b7280; font{-}style:italic;";}
\NormalTok{      opCell.textContent = s.op;}
\NormalTok{      const eqCell = document.createElement("span");}
\NormalTok{      eqCell.innerHTML = katex.renderToString(s.eq, \{ throwOnError: false, displayMode: false \});}
\NormalTok{      row.appendChild(opCell);}
\NormalTok{      row.appendChild(eqCell);}
\NormalTok{      if (s.note) \{}
\NormalTok{        const noteRow = document.createElement("div");}
\NormalTok{        noteRow.style.cssText = "grid{-}column:2; font{-}size:0.82em; color:\#6b7280; padding{-}bottom:0.2rem;";}
\NormalTok{        noteRow.textContent = s.note;}
\NormalTok{        row.appendChild(noteRow);}
\NormalTok{      \}}
\NormalTok{      stepsDiv.appendChild(row);}
\NormalTok{    \}}
\NormalTok{    if (current === totalSteps) \{}
\NormalTok{      const done = document.createElement("div");}
\NormalTok{      done.style.cssText = "margin{-}top:0.5rem; font{-}size:0.9em; color:\#059669; font{-}weight:500;";}
\NormalTok{      done.textContent = "Solution complete.";}
\NormalTok{      stepsDiv.appendChild(done);}
\NormalTok{    \}}
\NormalTok{  \}}

\NormalTok{  const controls = document.createElement("div");}
\NormalTok{  controls.style.cssText = "display:flex; gap:0.5rem; align{-}items:center;";}
\NormalTok{  const nextBtn = document.createElement("button");}
\NormalTok{  nextBtn.textContent = "Next step →";}
\NormalTok{  nextBtn.style.cssText = "padding:0.35rem 0.85rem; border:1px solid \#d1d5db; border{-}radius:4px; background:\#fff; cursor:pointer; font{-}size:0.9em;";}
\NormalTok{  const resetBtn = document.createElement("button");}
\NormalTok{  resetBtn.textContent = "Reset";}
\NormalTok{  resetBtn.style.cssText = "padding:0.35rem 0.75rem; border:1px solid \#e5e7eb; border{-}radius:4px; background:\#f9fafb; cursor:pointer; font{-}size:0.9em; color:\#6b7280;";}
\NormalTok{  const counter = document.createElement("span");}
\NormalTok{  counter.style.cssText = "font{-}size:0.82em; color:\#9ca3af; margin{-}left:0.25rem;";}

\NormalTok{  function updateButtons() \{}
\NormalTok{    nextBtn.disabled = current === totalSteps;}
\NormalTok{    nextBtn.style.opacity = current === totalSteps ? "0.4" : "1";}
\NormalTok{    counter.textContent = current === 0 ? "Click to reveal steps" : \textasciigrave{}Step $\{current\} of $\{totalSteps\}\textasciigrave{};}
\NormalTok{  \}}

\NormalTok{  nextBtn.onclick = () =\textgreater{} \{ if (current \textless{} totalSteps) \{ current++; renderSteps(); updateButtons(); \} \};}
\NormalTok{  resetBtn.onclick = () =\textgreater{} \{ current = 0; renderSteps(); updateButtons(); \};}
\NormalTok{  controls.appendChild(nextBtn);}
\NormalTok{  controls.appendChild(resetBtn);}
\NormalTok{  controls.appendChild(counter);}
\NormalTok{  container.appendChild(stepsDiv);}
\NormalTok{  container.appendChild(controls);}
\NormalTok{  renderSteps();}
\NormalTok{  updateButtons();}
\NormalTok{  return container;}
\NormalTok{\}}
\end{Highlighting}
\end{Shaded}

\begin{center}\rule{0.5\linewidth}{0.5pt}\end{center}

\subsection{Exercise 1: Node admittance matrix
multiplication}\label{exercise-1-node-admittance-matrix-multiplication}

A circuit has three nodes. The node admittance matrix \(\mathbf{Y}\) and
a voltage vector \(\mathbf{V}\) are:

\[\mathbf{Y} = \begin{pmatrix} 5 & -2 & -1 \\ -2 & 6 & -3 \\ -1 & -3 & 4 \end{pmatrix}, \qquad \mathbf{V} = \begin{pmatrix} 10 \\ 5 \\ 2 \end{pmatrix} \text{ V}\]

Compute the current injection vector
\(\mathbf{I} = \mathbf{Y}\mathbf{V}\).

\begin{Shaded}
\begin{Highlighting}[]
\NormalTok{//| echo: false}
\NormalTok{makeStepperHTML([}
\NormalTok{  \{}
\NormalTok{    op: "Write the product definition",}
\NormalTok{    eq: "\textbackslash{}\textbackslash{}mathbf\{I\} = \textbackslash{}\textbackslash{}mathbf\{Y\}\textbackslash{}\textbackslash{}mathbf\{V\} = \textbackslash{}\textbackslash{}begin\{pmatrix\} 5 \& {-}2 \& {-}1 \textbackslash{}\textbackslash{}\textbackslash{}\textbackslash{} {-}2 \& 6 \& {-}3 \textbackslash{}\textbackslash{}\textbackslash{}\textbackslash{} {-}1 \& {-}3 \& 4 \textbackslash{}\textbackslash{}end\{pmatrix\}\textbackslash{}\textbackslash{}begin\{pmatrix\} 10 \textbackslash{}\textbackslash{}\textbackslash{}\textbackslash{} 5 \textbackslash{}\textbackslash{}\textbackslash{}\textbackslash{} 2 \textbackslash{}\textbackslash{}end\{pmatrix\}",}
\NormalTok{    note: "Each entry of I is the dot product of one row of Y with V."}
\NormalTok{  \},}
\NormalTok{  \{}
\NormalTok{    op: "Compute I₁ (row 1 · V)",}
\NormalTok{    eq: "I\_1 = 5(10) + ({-}2)(5) + ({-}1)(2) = 50 {-} 10 {-} 2 = 38 \textbackslash{}\textbackslash{}text\{ A\}"}
\NormalTok{  \},}
\NormalTok{  \{}
\NormalTok{    op: "Compute I₂ (row 2 · V)",}
\NormalTok{    eq: "I\_2 = ({-}2)(10) + 6(5) + ({-}3)(2) = {-}20 + 30 {-} 6 = 4 \textbackslash{}\textbackslash{}text\{ A\}"}
\NormalTok{  \},}
\NormalTok{  \{}
\NormalTok{    op: "Compute I₃ (row 3 · V)",}
\NormalTok{    eq: "I\_3 = ({-}1)(10) + ({-}3)(5) + 4(2) = {-}10 {-} 15 + 8 = {-}17 \textbackslash{}\textbackslash{}text\{ A\}"}
\NormalTok{  \},}
\NormalTok{  \{}
\NormalTok{    op: "Assemble the result",}
\NormalTok{    eq: "\textbackslash{}\textbackslash{}mathbf\{I\} = \textbackslash{}\textbackslash{}begin\{pmatrix\} 38 \textbackslash{}\textbackslash{}\textbackslash{}\textbackslash{} 4 \textbackslash{}\textbackslash{}\textbackslash{}\textbackslash{} {-}17 \textbackslash{}\textbackslash{}end\{pmatrix\} \textbackslash{}\textbackslash{}text\{ A\}",}
\NormalTok{    note: "Negative current injection means current is leaving that node."}
\NormalTok{  \}}
\NormalTok{])}
\end{Highlighting}
\end{Shaded}

\begin{center}\rule{0.5\linewidth}{0.5pt}\end{center}

\subsection{Exercise 2: Mesh current analysis by Gaussian
elimination}\label{exercise-2-mesh-current-analysis-by-gaussian-elimination}

Three mesh currents \(I_1, I_2, I_3\) (in amperes) in a resistive
network satisfy:

\[\begin{aligned} 8I_1 - 2I_2 - 0I_3 &= 16 \\ -2I_1 + 7I_2 - 3I_3 &= 0 \\ 0I_1 - 3I_2 + 5I_3 &= -6 \end{aligned}\]

Find \(I_1, I_2, I_3\) by Gaussian elimination.

\begin{Shaded}
\begin{Highlighting}[]
\NormalTok{//| echo: false}
\NormalTok{makeStepperHTML([}
\NormalTok{  \{}
\NormalTok{    op: "Write augmented matrix",}
\NormalTok{    eq: "\textbackslash{}\textbackslash{}left(\textbackslash{}\textbackslash{}begin\{array\}\{rrr|r\} 8 \& {-}2 \& 0 \& 16 \textbackslash{}\textbackslash{}\textbackslash{}\textbackslash{} {-}2 \& 7 \& {-}3 \& 0 \textbackslash{}\textbackslash{}\textbackslash{}\textbackslash{} 0 \& {-}3 \& 5 \& {-}6 \textbackslash{}\textbackslash{}end\{array\}\textbackslash{}\textbackslash{}right)"}
\NormalTok{  \},}
\NormalTok{  \{}
\NormalTok{    op: "Eliminate column 1: R₂ ← R₂ + (1/4)R₁",}
\NormalTok{    eq: "\textbackslash{}\textbackslash{}left(\textbackslash{}\textbackslash{}begin\{array\}\{rrr|r\} 8 \& {-}2 \& 0 \& 16 \textbackslash{}\textbackslash{}\textbackslash{}\textbackslash{} 0 \& 6.5 \& {-}3 \& 4 \textbackslash{}\textbackslash{}\textbackslash{}\textbackslash{} 0 \& {-}3 \& 5 \& {-}6 \textbackslash{}\textbackslash{}end\{array\}\textbackslash{}\textbackslash{}right)",}
\NormalTok{    note: "Multiplier m₂₁ = −2/8 = −1/4; add (1/4)R₁ to R₂."}
\NormalTok{  \},}
\NormalTok{  \{}
\NormalTok{    op: "Eliminate column 2: R₃ ← R₃ + (3/6.5)R₂",}
\NormalTok{    eq: "\textbackslash{}\textbackslash{}left(\textbackslash{}\textbackslash{}begin\{array\}\{rrr|r\} 8 \& {-}2 \& 0 \& 16 \textbackslash{}\textbackslash{}\textbackslash{}\textbackslash{} 0 \& 6.5 \& {-}3 \& 4 \textbackslash{}\textbackslash{}\textbackslash{}\textbackslash{} 0 \& 0 \& 3.615 \& {-}4.154 \textbackslash{}\textbackslash{}end\{array\}\textbackslash{}\textbackslash{}right)",}
\NormalTok{    note: "Multiplier m₃₂ = −3/6.5 = −6/13; subtracting that multiple zeros the (3,2) entry."}
\NormalTok{  \},}
\NormalTok{  \{}
\NormalTok{    op: "Back{-}substitute: solve for I₃",}
\NormalTok{    eq: "I\_3 = \textbackslash{}\textbackslash{}frac\{{-}4.154\}\{3.615\} \textbackslash{}\textbackslash{}approx {-}1.149 \textbackslash{}\textbackslash{}text\{ A\}"}
\NormalTok{  \},}
\NormalTok{  \{}
\NormalTok{    op: "Back{-}substitute: solve for I₂",}
\NormalTok{    eq: "I\_2 = \textbackslash{}\textbackslash{}frac\{4 {-} ({-}3)({-}1.149)\}\{6.5\} = \textbackslash{}\textbackslash{}frac\{4 {-} 3.447\}\{6.5\} \textbackslash{}\textbackslash{}approx 0.085 \textbackslash{}\textbackslash{}text\{ A\}"}
\NormalTok{  \},}
\NormalTok{  \{}
\NormalTok{    op: "Back{-}substitute: solve for I₁",}
\NormalTok{    eq: "I\_1 = \textbackslash{}\textbackslash{}frac\{16 {-} ({-}2)(0.085)\}\{8\} = \textbackslash{}\textbackslash{}frac\{16.170\}\{8\} \textbackslash{}\textbackslash{}approx 2.021 \textbackslash{}\textbackslash{}text\{ A\}",}
\NormalTok{    note: "Exact fractions: I₁ = 1063/526 ≈ 2.021 A; I₂ = 9/106 ≈ 0.085 A; I₃ = −61/53 ≈ −1.151 A."}
\NormalTok{  \}}
\NormalTok{])}
\end{Highlighting}
\end{Shaded}

\begin{center}\rule{0.5\linewidth}{0.5pt}\end{center}

\subsection{Exercise 3: RREF and solution type (underdetermined
structural
system)}\label{exercise-3-rref-and-solution-type-underdetermined-structural-system}

A redundant structure gives three equilibrium equations in four unknown
member forces \(F_1, F_2, F_3, F_4\):

\[\begin{aligned} F_1 + 2F_2 - F_3 + F_4 &= 10 \\ 2F_1 + 4F_2 + F_3 - F_4 &= 8 \\ F_1 + 2F_2 + 3F_3 - 3F_4 &= -2 \end{aligned}\]

Row-reduce to RREF. Identify the free variable(s) and write the general
solution.

\begin{Shaded}
\begin{Highlighting}[]
\NormalTok{//| echo: false}
\NormalTok{makeStepperHTML([}
\NormalTok{  \{}
\NormalTok{    op: "Write augmented matrix",}
\NormalTok{    eq: "\textbackslash{}\textbackslash{}left(\textbackslash{}\textbackslash{}begin\{array\}\{rrrr|r\} 1 \& 2 \& {-}1 \& 1 \& 10 \textbackslash{}\textbackslash{}\textbackslash{}\textbackslash{} 2 \& 4 \& 1 \& {-}1 \& 8 \textbackslash{}\textbackslash{}\textbackslash{}\textbackslash{} 1 \& 2 \& 3 \& {-}3 \& {-}2 \textbackslash{}\textbackslash{}end\{array\}\textbackslash{}\textbackslash{}right)"}
\NormalTok{  \},}
\NormalTok{  \{}
\NormalTok{    op: "Eliminate col 1: R₂ ← R₂ − 2R₁; R₃ ← R₃ − R₁",}
\NormalTok{    eq: "\textbackslash{}\textbackslash{}left(\textbackslash{}\textbackslash{}begin\{array\}\{rrrr|r\} 1 \& 2 \& {-}1 \& 1 \& 10 \textbackslash{}\textbackslash{}\textbackslash{}\textbackslash{} 0 \& 0 \& 3 \& {-}3 \& {-}12 \textbackslash{}\textbackslash{}\textbackslash{}\textbackslash{} 0 \& 0 \& 4 \& {-}4 \& {-}12 \textbackslash{}\textbackslash{}end\{array\}\textbackslash{}\textbackslash{}right)"}
\NormalTok{  \},}
\NormalTok{  \{}
\NormalTok{    op: "Scale R₂ by 1/3; eliminate col 3: R₃ ← R₃ − 4R₂",}
\NormalTok{    eq: "\textbackslash{}\textbackslash{}left(\textbackslash{}\textbackslash{}begin\{array\}\{rrrr|r\} 1 \& 2 \& {-}1 \& 1 \& 10 \textbackslash{}\textbackslash{}\textbackslash{}\textbackslash{} 0 \& 0 \& 1 \& {-}1 \& {-}4 \textbackslash{}\textbackslash{}\textbackslash{}\textbackslash{} 0 \& 0 \& 0 \& 0 \& 0 \textbackslash{}\textbackslash{}end\{array\}\textbackslash{}\textbackslash{}right)",}
\NormalTok{    note: "Row 3 becomes all zeros — a dependent equation."}
\NormalTok{  \},}
\NormalTok{  \{}
\NormalTok{    op: "Back{-}eliminate col 3 from row 1: R₁ ← R₁ + R₂",}
\NormalTok{    eq: "\textbackslash{}\textbackslash{}left(\textbackslash{}\textbackslash{}begin\{array\}\{rrrr|r\} 1 \& 2 \& 0 \& 0 \& 6 \textbackslash{}\textbackslash{}\textbackslash{}\textbackslash{} 0 \& 0 \& 1 \& {-}1 \& {-}4 \textbackslash{}\textbackslash{}\textbackslash{}\textbackslash{} 0 \& 0 \& 0 \& 0 \& 0 \textbackslash{}\textbackslash{}end\{array\}\textbackslash{}\textbackslash{}right)"}
\NormalTok{  \},}
\NormalTok{  \{}
\NormalTok{    op: "Identify pivot and free variables",}
\NormalTok{    eq: "\textbackslash{}\textbackslash{}text\{Pivots in cols 1, 3\} \textbackslash{}\textbackslash{}implies F\_1,\textbackslash{}\textbackslash{}, F\_3 \textbackslash{}\textbackslash{}text\{ are pivot variables; \} F\_2,\textbackslash{}\textbackslash{}, F\_4 \textbackslash{}\textbackslash{}text\{ are free\}",}
\NormalTok{    note: "rank(A) = 2 \textless{} 4 unknowns, so nullity = 2: a two{-}parameter family of solutions."}
\NormalTok{  \},}
\NormalTok{  \{}
\NormalTok{    op: "Write general solution (let F₂ = s, F₄ = t)",}
\NormalTok{    eq: "F\_1 = 6 {-} 2s,\textbackslash{}\textbackslash{}; F\_2 = s,\textbackslash{}\textbackslash{}; F\_3 = {-}4 + t,\textbackslash{}\textbackslash{}; F\_4 = t \textbackslash{}\textbackslash{}quad (s,\textbackslash{}\textbackslash{},t \textbackslash{}\textbackslash{}in \textbackslash{}\textbackslash{}mathbb\{R\})",}
\NormalTok{    note: "Physical interpretation: the structure is statically indeterminate — you need material stiffness data to find a unique solution."}
\NormalTok{  \}}
\NormalTok{])}
\end{Highlighting}
\end{Shaded}

\begin{center}\rule{0.5\linewidth}{0.5pt}\end{center}

\subsection{Exercise 4: Rank of a sensor measurement
matrix}\label{exercise-4-rank-of-a-sensor-measurement-matrix}

A sensor network takes four measurements of three physical quantities
(temperature, pressure, flow rate). The measurement matrix is:

\[\mathbf{M} = \begin{pmatrix} 1 & 2 & 3 \\ 2 & 4 & 6 \\ 1 & 0 & 2 \\ 3 & 2 & 7 \end{pmatrix}\]

Find \(\text{rank}(\mathbf{M})\). What does this tell you about the
sensor network?

\begin{Shaded}
\begin{Highlighting}[]
\NormalTok{//| echo: false}
\NormalTok{makeStepperHTML([}
\NormalTok{  \{}
\NormalTok{    op: "Write the matrix and note its size",}
\NormalTok{    eq: "\textbackslash{}\textbackslash{}mathbf\{M\} \textbackslash{}\textbackslash{}text\{ is \} 4 \textbackslash{}\textbackslash{}times 3, \textbackslash{}\textbackslash{}text\{ so \} \textbackslash{}\textbackslash{}text\{rank\}(\textbackslash{}\textbackslash{}mathbf\{M\}) \textbackslash{}\textbackslash{}leq 3"}
\NormalTok{  \},}
\NormalTok{  \{}
\NormalTok{    op: "Eliminate col 1: R₂ ← R₂ − 2R₁; R₃ ← R₃ − R₁; R₄ ← R₄ − 3R₁",}
\NormalTok{    eq: "\textbackslash{}\textbackslash{}begin\{pmatrix\} 1 \& 2 \& 3 \textbackslash{}\textbackslash{}\textbackslash{}\textbackslash{} 0 \& 0 \& 0 \textbackslash{}\textbackslash{}\textbackslash{}\textbackslash{} 0 \& {-}2 \& {-}1 \textbackslash{}\textbackslash{}\textbackslash{}\textbackslash{} 0 \& {-}4 \& {-}2 \textbackslash{}\textbackslash{}end\{pmatrix\}",}
\NormalTok{    note: "Row 2 is exactly 2 × Row 1 — it carries no new information."}
\NormalTok{  \},}
\NormalTok{  \{}
\NormalTok{    op: "Swap R₂ and R₃ (bring nonzero pivot to position (2,2))",}
\NormalTok{    eq: "\textbackslash{}\textbackslash{}begin\{pmatrix\} 1 \& 2 \& 3 \textbackslash{}\textbackslash{}\textbackslash{}\textbackslash{} 0 \& {-}2 \& {-}1 \textbackslash{}\textbackslash{}\textbackslash{}\textbackslash{} 0 \& 0 \& 0 \textbackslash{}\textbackslash{}\textbackslash{}\textbackslash{} 0 \& {-}4 \& {-}2 \textbackslash{}\textbackslash{}end\{pmatrix\}"}
\NormalTok{  \},}
\NormalTok{  \{}
\NormalTok{    op: "Eliminate col 2 from R₄: R₄ ← R₄ − 2R₂",}
\NormalTok{    eq: "\textbackslash{}\textbackslash{}begin\{pmatrix\} 1 \& 2 \& 3 \textbackslash{}\textbackslash{}\textbackslash{}\textbackslash{} 0 \& {-}2 \& {-}1 \textbackslash{}\textbackslash{}\textbackslash{}\textbackslash{} 0 \& 0 \& 0 \textbackslash{}\textbackslash{}\textbackslash{}\textbackslash{} 0 \& 0 \& 0 \textbackslash{}\textbackslash{}end\{pmatrix\}",}
\NormalTok{    note: "Two zero rows — rank equals the number of nonzero rows after reduction."}
\NormalTok{  \},}
\NormalTok{  \{}
\NormalTok{    op: "Count pivots",}
\NormalTok{    eq: "\textbackslash{}\textbackslash{}text\{rank\}(\textbackslash{}\textbackslash{}mathbf\{M\}) = 2",}
\NormalTok{    note: "Only 2 of the 4 sensors provide independent information. The other two are redundant — they measure combinations already captured. One redundant sensor is useful for fault detection; two is over{-}instrumented."}
\NormalTok{  \}}
\NormalTok{])}
\end{Highlighting}
\end{Shaded}

\begin{center}\rule{0.5\linewidth}{0.5pt}\end{center}

\subsection{Exercise 5: LU factorisation for repeated
solves}\label{exercise-5-lu-factorisation-for-repeated-solves}

A heat-transfer network has conductance matrix:

\[\mathbf{K} = \begin{pmatrix} 4 & -1 & 0 \\ -1 & 3 & -1 \\ 0 & -1 & 2 \end{pmatrix}\]

Find the LU factorisation of \(\mathbf{K}\), then use it to solve
\(\mathbf{K}\mathbf{T} = \mathbf{q}\) for two different heat-flux
vectors:

\[\mathbf{q}_1 = \begin{pmatrix} 8 \\ 4 \\ 2 \end{pmatrix}, \qquad \mathbf{q}_2 = \begin{pmatrix} 0 \\ 6 \\ 4 \end{pmatrix}\]

\begin{Shaded}
\begin{Highlighting}[]
\NormalTok{//| echo: false}
\NormalTok{makeStepperHTML([}
\NormalTok{  \{}
\NormalTok{    op: "Elimination step 1: R₂ ← R₂ − (−1/4)R₁",}
\NormalTok{    eq: "m\_\{21\} = \textbackslash{}\textbackslash{}tfrac\{{-}1\}\{4\} \textbackslash{}\textbackslash{}implies R\_2 \textbackslash{}\textbackslash{}leftarrow R\_2 {-} \textbackslash{}\textbackslash{}left({-}\textbackslash{}\textbackslash{}tfrac\{1\}\{4\}\textbackslash{}\textbackslash{}right)R\_1",}
\NormalTok{    note: "Subtracting (−1/4)×R₁ is the same as adding (1/4)×R₁."}
\NormalTok{  \},}
\NormalTok{  \{}
\NormalTok{    op: "Result after step 1",}
\NormalTok{    eq: "\textbackslash{}\textbackslash{}mathbf\{U\}\^{}\{(1)\} = \textbackslash{}\textbackslash{}begin\{pmatrix\} 4 \& {-}1 \& 0 \textbackslash{}\textbackslash{}\textbackslash{}\textbackslash{} 0 \& 11/4 \& {-}1 \textbackslash{}\textbackslash{}\textbackslash{}\textbackslash{} 0 \& {-}1 \& 2 \textbackslash{}\textbackslash{}end\{pmatrix\}"}
\NormalTok{  \},}
\NormalTok{  \{}
\NormalTok{    op: "Elimination step 2: R₃ ← R₃ − (−1/(11/4))R₂ = R₃ + (4/11)R₂",}
\NormalTok{    eq: "m\_\{32\} = \textbackslash{}\textbackslash{}tfrac\{{-}1\}\{11/4\} = {-}\textbackslash{}\textbackslash{}tfrac\{4\}\{11\}"}
\NormalTok{  \},}
\NormalTok{  \{}
\NormalTok{    op: "Upper triangular U",}
\NormalTok{    eq: "\textbackslash{}\textbackslash{}mathbf\{U\} = \textbackslash{}\textbackslash{}begin\{pmatrix\} 4 \& {-}1 \& 0 \textbackslash{}\textbackslash{}\textbackslash{}\textbackslash{} 0 \& 11/4 \& {-}1 \textbackslash{}\textbackslash{}\textbackslash{}\textbackslash{} 0 \& 0 \& 18/11 \textbackslash{}\textbackslash{}end\{pmatrix\}"}
\NormalTok{  \},}
\NormalTok{  \{}
\NormalTok{    op: "Assemble L from multipliers",}
\NormalTok{    eq: "\textbackslash{}\textbackslash{}mathbf\{L\} = \textbackslash{}\textbackslash{}begin\{pmatrix\} 1 \& 0 \& 0 \textbackslash{}\textbackslash{}\textbackslash{}\textbackslash{} {-}1/4 \& 1 \& 0 \textbackslash{}\textbackslash{}\textbackslash{}\textbackslash{} 0 \& {-}4/11 \& 1 \textbackslash{}\textbackslash{}end\{pmatrix\}"}
\NormalTok{  \},}
\NormalTok{  \{}
\NormalTok{    op: "Solve Ly = q₁ by forward substitution",}
\NormalTok{    eq: "\textbackslash{}\textbackslash{}mathbf\{y\}\_1: \textbackslash{}\textbackslash{}; y\_1 = 8,\textbackslash{}\textbackslash{}; y\_2 = 4 + \textbackslash{}\textbackslash{}tfrac\{1\}\{4\}(8) = 6,\textbackslash{}\textbackslash{}; y\_3 = 2 + \textbackslash{}\textbackslash{}tfrac\{4\}\{11\}(6) = \textbackslash{}\textbackslash{}tfrac\{46\}\{11\}"}
\NormalTok{  \},}
\NormalTok{  \{}
\NormalTok{    op: "Solve UT = y₁ by back{-}substitution",}
\NormalTok{    eq: "T\_3 = \textbackslash{}\textbackslash{}frac\{46/11\}\{18/11\} = \textbackslash{}\textbackslash{}frac\{46\}\{18\} = \textbackslash{}\textbackslash{}frac\{23\}\{9\} \textbackslash{}\textbackslash{}approx 2.56,\textbackslash{}\textbackslash{}; T\_2 = \textbackslash{}\textbackslash{}frac\{6 + T\_3\}\{11/4\} \textbackslash{}\textbackslash{}approx 3.11,\textbackslash{}\textbackslash{}; T\_1 = \textbackslash{}\textbackslash{}frac\{8 + T\_2 {-} 0\}\{4\} \textbackslash{}\textbackslash{}approx 2.78",}
\NormalTok{    note: "Exact: T₁ = 167/60, T₂ = 56/18 = 28/9, T₃ = 23/9."}
\NormalTok{  \},}
\NormalTok{  \{}
\NormalTok{    op: "For q₂: reuse L and U, only redo the two triangular solves",}
\NormalTok{    eq: "\textbackslash{}\textbackslash{}mathbf\{y\}\_2: \textbackslash{}\textbackslash{}; y\_1 = 0,\textbackslash{}\textbackslash{}; y\_2 = 6,\textbackslash{}\textbackslash{}; y\_3 = 4 + \textbackslash{}\textbackslash{}tfrac\{4\}\{11\}(6) = \textbackslash{}\textbackslash{}tfrac\{68\}\{11\}",}
\NormalTok{    note: "The LU factorisation cost is paid once. Each new right{-}hand side costs only O(n²)."}
\NormalTok{  \}}
\NormalTok{])}
\end{Highlighting}
\end{Shaded}

\begin{center}\rule{0.5\linewidth}{0.5pt}\end{center}

\subsection{Exercise 6: Write and solve the linear system for a 3-mesh
circuit}\label{exercise-6-write-and-solve-the-linear-system-for-a-3-mesh-circuit}

A three-mesh circuit has the following KVL equations after applying
Kirchhoff's voltage law. The mesh currents \(I_1, I_2, I_3\) (in mA)
satisfy:

\begin{itemize}
\tightlist
\item
  \textbf{Mesh 1:} A 10 Ω resistor carries \(I_1\); a 5 Ω resistor is
  shared with mesh 2 carrying \(I_1 - I_2\). Voltage source: 20 V.
\item
  \textbf{Mesh 2:} The shared 5 Ω resistor carries \(I_2 - I_1\); a 10 Ω
  resistor carries \(I_2\); a 4 Ω resistor is shared with mesh 3
  carrying \(I_2 - I_3\). No voltage source.
\item
  \textbf{Mesh 3:} The shared 4 Ω resistor carries \(I_3 - I_2\); a 6 Ω
  resistor carries \(I_3\). Voltage source: −8 V (opposing reference
  direction).
\end{itemize}

Write \(\mathbf{A}\mathbf{x} = \mathbf{b}\) explicitly, then solve using
Gaussian elimination.

\begin{Shaded}
\begin{Highlighting}[]
\NormalTok{//| echo: false}
\NormalTok{makeStepperHTML([}
\NormalTok{  \{}
\NormalTok{    op: "Apply KVL to mesh 1",}
\NormalTok{    eq: "10I\_1 + 5(I\_1 {-} I\_2) = 20 \textbackslash{}\textbackslash{}implies 15I\_1 {-} 5I\_2 = 20"}
\NormalTok{  \},}
\NormalTok{  \{}
\NormalTok{    op: "Apply KVL to mesh 2",}
\NormalTok{    eq: "5(I\_2 {-} I\_1) + 10I\_2 + 4(I\_2 {-} I\_3) = 0 \textbackslash{}\textbackslash{}implies {-}5I\_1 + 19I\_2 {-} 4I\_3 = 0"}
\NormalTok{  \},}
\NormalTok{  \{}
\NormalTok{    op: "Apply KVL to mesh 3",}
\NormalTok{    eq: "4(I\_3 {-} I\_2) + 6I\_3 = {-}8 \textbackslash{}\textbackslash{}implies {-}4I\_2 + 10I\_3 = {-}8"}
\NormalTok{  \},}
\NormalTok{  \{}
\NormalTok{    op: "Write in matrix form Ax = b",}
\NormalTok{    eq: "\textbackslash{}\textbackslash{}begin\{pmatrix\} 15 \& {-}5 \& 0 \textbackslash{}\textbackslash{}\textbackslash{}\textbackslash{} {-}5 \& 19 \& {-}4 \textbackslash{}\textbackslash{}\textbackslash{}\textbackslash{} 0 \& {-}4 \& 10 \textbackslash{}\textbackslash{}end\{pmatrix\}\textbackslash{}\textbackslash{}begin\{pmatrix\} I\_1 \textbackslash{}\textbackslash{}\textbackslash{}\textbackslash{} I\_2 \textbackslash{}\textbackslash{}\textbackslash{}\textbackslash{} I\_3 \textbackslash{}\textbackslash{}end\{pmatrix\} = \textbackslash{}\textbackslash{}begin\{pmatrix\} 20 \textbackslash{}\textbackslash{}\textbackslash{}\textbackslash{} 0 \textbackslash{}\textbackslash{}\textbackslash{}\textbackslash{} {-}8 \textbackslash{}\textbackslash{}end\{pmatrix\}",}
\NormalTok{    note: "The coefficient matrix is symmetric — always the case for KVL mesh analysis with passive resistors."}
\NormalTok{  \},}
\NormalTok{  \{}
\NormalTok{    op: "Eliminate col 1: R₂ ← R₂ + (1/3)R₁",}
\NormalTok{    eq: "\textbackslash{}\textbackslash{}left(\textbackslash{}\textbackslash{}begin\{array\}\{rrr|r\} 15 \& {-}5 \& 0 \& 20 \textbackslash{}\textbackslash{}\textbackslash{}\textbackslash{} 0 \& 52/3 \& {-}4 \& 20/3 \textbackslash{}\textbackslash{}\textbackslash{}\textbackslash{} 0 \& {-}4 \& 10 \& {-}8 \textbackslash{}\textbackslash{}end\{array\}\textbackslash{}\textbackslash{}right)",}
\NormalTok{    note: "Multiplier: m₂₁ = −5/15 = −1/3; subtract (−1/3)R₁, i.e., add (1/3)R₁."}
\NormalTok{  \},}
\NormalTok{  \{}
\NormalTok{    op: "Eliminate col 2: R₃ ← R₃ + (12/52)R₂",}
\NormalTok{    eq: "\textbackslash{}\textbackslash{}left(\textbackslash{}\textbackslash{}begin\{array\}\{rrr|r\} 15 \& {-}5 \& 0 \& 20 \textbackslash{}\textbackslash{}\textbackslash{}\textbackslash{} 0 \& 52/3 \& {-}4 \& 20/3 \textbackslash{}\textbackslash{}\textbackslash{}\textbackslash{} 0 \& 0 \& 442/52 \& {-}31/13 \textbackslash{}\textbackslash{}end\{array\}\textbackslash{}\textbackslash{}right)",}
\NormalTok{    note: "Multiplier: m₃₂ = −4 / (52/3) = −12/52 = −3/13."}
\NormalTok{  \},}
\NormalTok{  \{}
\NormalTok{    op: "Back{-}substitute",}
\NormalTok{    eq: "I\_3 = \textbackslash{}\textbackslash{}frac\{{-}31/13\}\{442/52\} \textbackslash{}\textbackslash{}approx {-}0.280 \textbackslash{}\textbackslash{}text\{ mA\}, \textbackslash{}\textbackslash{}quad I\_2 \textbackslash{}\textbackslash{}approx 0.452 \textbackslash{}\textbackslash{}text\{ mA\}, \textbackslash{}\textbackslash{}quad I\_1 \textbackslash{}\textbackslash{}approx 1.484 \textbackslash{}\textbackslash{}text\{ mA\}",}
\NormalTok{    note: "Verify: 15(1.484) − 5(0.452) = 22.26 − 2.26 = 20 ✓. All three mesh currents are determined."}
\NormalTok{  \}}
\NormalTok{])}
\end{Highlighting}
\end{Shaded}

\chapter{Eigenvalues and
eigenvectors}\label{eigenvalues-and-eigenvectors}

The natural modes of a system

\hfill\break

Spin a shaft slowly and nothing unusual happens. Spin it faster and, at
one particular speed, the whole machine shakes itself apart. That speed
is not bad luck --- it is the shaft's natural frequency, a number
written into the geometry and stiffness of the metal itself. Engineers
calculate it before the shaft is ever machined.

The calculation is an eigenvalue problem. Feed in the stiffness matrix,
solve for the special values that make a certain determinant zero, and
out come the critical speeds. The same computation --- different matrix,
same method --- tells Google which web pages matter, tells a structural
engineer where a building will sway, and tells a data scientist which
direction in a dataset carries the most information.

\section{The eigenvalue problem}\label{the-eigenvalue-problem}

Multiply a vector by a matrix and you generally get a new vector
pointing in a completely different direction. But for certain special
vectors, the matrix acts more simply --- it only scales the vector,
leaving its direction unchanged (or exactly reversed). These are
\textbf{eigenvectors}.

Formally: \(\mathbf{x}\) is an eigenvector of \(A\) with eigenvalue
\(\lambda\) if

\[A\mathbf{x} = \lambda\mathbf{x}, \quad \mathbf{x} \neq \mathbf{0}\]

Read this as: ``applying \(A\) to \(\mathbf{x}\) is the same as
stretching \(\mathbf{x}\) by a factor of \(\lambda\).'' The vector
\(\mathbf{x}\) points in a direction that \(A\) cannot rotate --- it can
only scale it.

The eigenvalue \(\lambda\) can be:

\begin{itemize}
\tightlist
\item
  \textbf{Positive} --- the eigenvector points the same way, scaled by
  \(\lambda\).
\item
  \textbf{Negative} --- the eigenvector is flipped and scaled.
\item
  \textbf{Zero} --- \(A\mathbf{x} = \mathbf{0}\); the eigenvector is in
  the null space of \(A\) (and \(A\) is singular).
\item
  \textbf{Complex} --- the matrix introduces a rotation; eigenvalues
  come in conjugate pairs for real symmetric matrices this cannot
  happen, but for general real matrices it can.
\end{itemize}

\textbf{Geometric picture.} Imagine \(A\) as a transformation of the
plane. Most vectors get rotated and stretched into new directions. The
eigenvectors are the ``fixed axes'' --- directions that survive the
transformation pointing the same way they started (or exactly opposite).
The eigenvalue tells you how much the transformation stretches along
each fixed axis.

\begin{Shaded}
\begin{Highlighting}[]
\NormalTok{//| echo: false}
\NormalTok{eigenVisualiser = \{}
\NormalTok{  const d3 = await require("d3@7");}
\NormalTok{  const W = 560, H = 400;}
\NormalTok{  const ML = 40, MR = 40, MT = 40, MB = 40;}
\NormalTok{  const innerW = W {-} ML {-} MR;}
\NormalTok{  const innerH = H {-} MT {-} MB;}
  
\NormalTok{  const container = d3.create("div").style("max{-}width", W+"px").style("font{-}family", "sans{-}serif").style("margin", "0 auto");}
  
\NormalTok{  const controls = container.append("div").style("display", "grid").style("grid{-}template{-}columns", "1fr 1fr").style("gap", "10px").style("margin{-}bottom", "15px").style("background", "\#f8fafc").style("padding", "10px").style("border{-}radius", "8px").style("border", "1px solid \#e2e8f0");}
  
\NormalTok{  function makeSlider(parent, label, initVal) \{}
\NormalTok{    const wrap = parent.append("label").style("display", "flex").style("align{-}items", "center").style("font{-}size", "0.9rem").style("font{-}family", "monospace");}
\NormalTok{    wrap.append("span").style("width", "20px").text(label);}
\NormalTok{    const inp = wrap.append("input").attr("type", "range").attr("min", {-}3).attr("max", 3).attr("step", 0.1).attr("value", initVal).style("flex", 1).style("margin", "0 10px");}
\NormalTok{    const val = wrap.append("span").style("width", "35px").style("text{-}align", "right").text(initVal.toFixed(1));}
\NormalTok{    return \{inp, val\};}
\NormalTok{  \}}
  
\NormalTok{  const sliderA = makeSlider(controls, "a", 2);}
\NormalTok{  const sliderB = makeSlider(controls, "b", 1);}
\NormalTok{  const sliderC = makeSlider(controls, "c", 1);}
\NormalTok{  const sliderD = makeSlider(controls, "d", 2);}
  
\NormalTok{  const svg = container.append("svg").attr("width", W).attr("height", H).attr("viewBox", \textasciigrave{}0 0 $\{W\} $\{H\}\textasciigrave{});}
\NormalTok{  const g = svg.append("g").attr("transform", \textasciigrave{}translate($\{W/2\},$\{H/2\})\textasciigrave{});}
  
\NormalTok{  const xScale = d3.scaleLinear([{-}6, 6], [{-}innerW/2, innerW/2]);}
\NormalTok{  const yScale = d3.scaleLinear([{-}6, 6], [innerH/2, {-}innerH/2]);}
  
\NormalTok{  g.append("line").attr("x1", xScale({-}6)).attr("y1", yScale(0)).attr("x2", xScale(6)).attr("y2", yScale(0)).attr("stroke", "\#cbd5e1");}
\NormalTok{  g.append("line").attr("x1", xScale(0)).attr("y1", yScale({-}6)).attr("x2", xScale(0)).attr("y2", yScale(6)).attr("stroke", "\#cbd5e1");}
  
\NormalTok{  const unitPts = d3.range(0, 2 * Math.PI + 0.1, 0.1).map(theta =\textgreater{} [Math.cos(theta), Math.sin(theta)]);}
\NormalTok{  const lineGen = d3.line().x(d =\textgreater{} xScale(d[0])).y(d =\textgreater{} yScale(d[1]));}
  
\NormalTok{  g.append("path").datum(unitPts).attr("d", lineGen).attr("fill", "none").attr("stroke", "\#9ca3af").attr("stroke{-}dasharray", "4,4").attr("stroke{-}width", 1.5);}
  
\NormalTok{  const transformedPath = g.append("path").attr("fill", "rgba(59, 130, 246, 0.1)").attr("stroke", "\#3b82f6").attr("stroke{-}width", 2);}
  
\NormalTok{  const vecGroup = g.append("g");}
\NormalTok{  const e1Line = vecGroup.append("line").attr("stroke", "\#f97316").attr("stroke{-}width", 2.5);}
\NormalTok{  const e1Head = vecGroup.append("polygon").attr("fill", "\#f97316");}
\NormalTok{  const e1Label = vecGroup.append("text").attr("fill", "\#c2410c").attr("font{-}size", "13px").attr("font{-}weight", "bold").style("text{-}shadow", "1px 1px 2px white");}
  
\NormalTok{  const e2Line = vecGroup.append("line").attr("stroke", "\#10b981").attr("stroke{-}width", 2.5);}
\NormalTok{  const e2Head = vecGroup.append("polygon").attr("fill", "\#10b981");}
\NormalTok{  const e2Label = vecGroup.append("text").attr("fill", "\#047857").attr("font{-}size", "13px").attr("font{-}weight", "bold").style("text{-}shadow", "1px 1px 2px white");}
  
\NormalTok{  const statusText = container.append("div").style("text{-}align", "center").style("font{-}size", "0.95rem").style("font{-}weight", "bold").style("margin{-}top", "10px");}
  
\NormalTok{  function update() \{}
\NormalTok{    const a = +sliderA.inp.property("value");}
\NormalTok{    const b = +sliderB.inp.property("value");}
\NormalTok{    const c = +sliderC.inp.property("value");}
\NormalTok{    const d = +sliderD.inp.property("value");}
    
\NormalTok{    sliderA.val.text(a.toFixed(1));}
\NormalTok{    sliderB.val.text(b.toFixed(1));}
\NormalTok{    sliderC.val.text(c.toFixed(1));}
\NormalTok{    sliderD.val.text(d.toFixed(1));}
    
\NormalTok{    const transPts = unitPts.map(pt =\textgreater{} [a*pt[0] + b*pt[1], c*pt[0] + d*pt[1]]);}
\NormalTok{    transformedPath.datum(transPts).attr("d", lineGen);}
    
\NormalTok{    const tr = a + d;}
\NormalTok{    const det = a*d {-} b*c;}
\NormalTok{    const disc = tr*tr {-} 4*det;}
    
\NormalTok{    if (disc \textgreater{}= 0) \{}
\NormalTok{      const lam1 = (tr + Math.sqrt(disc)) / 2;}
\NormalTok{      const lam2 = (tr {-} Math.sqrt(disc)) / 2;}
      
\NormalTok{      statusText.text(\textasciigrave{}Real Eigenvalues: \textbackslash{}u03BB\textbackslash{}u2081 = $\{lam1.toFixed(2)\}, \textbackslash{}u03BB\textbackslash{}u2082 = $\{lam2.toFixed(2)\}\textasciigrave{}).style("color", "\#1e293b");}
      
\NormalTok{      function getEigVec(lam) \{}
\NormalTok{        if (Math.abs(b) \textgreater{} 1e{-}6) return [b, lam {-} a];}
\NormalTok{        if (Math.abs(c) \textgreater{} 1e{-}6) return [lam {-} d, c];}
\NormalTok{        return lam === a ? [1, 0] : [0, 1];}
\NormalTok{      \}}
      
\NormalTok{      function drawVec(v, lam, line, head, labelTxt, labelStr) \{}
\NormalTok{        const mag = Math.sqrt(v[0]*v[0] + v[1]*v[1]) || 1;}
\NormalTok{        const norm = [v[0]/mag, v[1]/mag];}
\NormalTok{        const endx = norm[0] * lam;}
\NormalTok{        const endy = norm[1] * lam;}
        
\NormalTok{        const sx = xScale(endx);}
\NormalTok{        const sy = yScale(endy);}
\NormalTok{        const ox = xScale(0);}
\NormalTok{        const oy = yScale(0);}
        
\NormalTok{        line.attr("x1", ox).attr("y1", oy).attr("x2", sx).attr("y2", sy).attr("visibility", "visible");}
        
\NormalTok{        if (Math.abs(lam) \textgreater{} 0.1) \{}
\NormalTok{          const angle = Math.atan2(sy {-} oy, sx {-} ox);}
\NormalTok{          const r = 8;}
\NormalTok{          head.attr("points", \textasciigrave{}$\{sx\},$\{sy\} $\{sx {-} r*Math.cos(angle {-} 0.5)\},$\{sy {-} r*Math.sin(angle {-} 0.5)\} $\{sx {-} r*Math.cos(angle + 0.5)\},$\{sy {-} r*Math.sin(angle + 0.5)\}\textasciigrave{}).attr("visibility", "visible");}
\NormalTok{        \} else \{}
\NormalTok{          head.attr("visibility", "hidden");}
\NormalTok{        \}}
        
\NormalTok{        labelTxt.attr("x", sx + 12*Math.sign(sx{-}ox)).attr("y", sy + 12*Math.sign(sy{-}oy){-}5).text(labelStr).attr("visibility", "visible");}
\NormalTok{      \}}
      
\NormalTok{      drawVec(getEigVec(lam1), lam1, e1Line, e1Head, e1Label, \textasciigrave{}\textbackslash{}u03BB\textbackslash{}u2081=$\{lam1.toFixed(2)\}\textasciigrave{});}
\NormalTok{      drawVec(getEigVec(lam2), lam2, e2Line, e2Head, e2Label, \textasciigrave{}\textbackslash{}u03BB\textbackslash{}u2082=$\{lam2.toFixed(2)\}\textasciigrave{});}
      
\NormalTok{    \} else \{}
\NormalTok{      statusText.text("Complex Eigenvalues (pure rotation/scaling without invariant axes)").style("color", "\#dc2626");}
\NormalTok{      e1Line.attr("visibility", "hidden"); e1Head.attr("visibility", "hidden"); e1Label.attr("visibility", "hidden");}
\NormalTok{      e2Line.attr("visibility", "hidden"); e2Head.attr("visibility", "hidden"); e2Label.attr("visibility", "hidden");}
\NormalTok{    \}}
\NormalTok{  \}}
  
\NormalTok{  update();}
\NormalTok{  sliderA.inp.on("input", update);}
\NormalTok{  sliderB.inp.on("input", update);}
\NormalTok{  sliderC.inp.on("input", update);}
\NormalTok{  sliderD.inp.on("input", update);}
  
\NormalTok{  return container.node();}
\NormalTok{\}}
\end{Highlighting}
\end{Shaded}

\begin{tcolorbox}[enhanced jigsaw, opacityback=0, colback=white, breakable, colframe=quarto-callout-note-color-frame, arc=.35mm, left=2mm, toprule=.15mm, leftrule=.75mm, bottomrule=.15mm, rightrule=.15mm]
\begin{minipage}[t]{5.5mm}
\textcolor{quarto-callout-note-color}{\faInfo}
\end{minipage}%
\begin{minipage}[t]{\textwidth - 5.5mm}

\textbf{Why \(\mathbf{x} \neq \mathbf{0}\) is required}

The zero vector satisfies \(A\mathbf{0} = \lambda\mathbf{0}\) for any
\(\lambda\). That is trivially true but useless --- it gives no
information about \(A\). We demand a non-zero eigenvector.

\end{minipage}%
\end{tcolorbox}

\section{Characteristic equation}\label{characteristic-equation}

Rearrange \(A\mathbf{x} = \lambda\mathbf{x}\):

\[A\mathbf{x} - \lambda\mathbf{x} = \mathbf{0}\]
\[(A - \lambda I)\mathbf{x} = \mathbf{0}\]

This is a homogeneous linear system. A homogeneous system has a non-zero
solution if and only if its coefficient matrix is \textbf{singular} ---
that is, its determinant is zero:

\[\det(A - \lambda I) = 0\]

This is the \textbf{characteristic equation}. Expanding the determinant
gives a polynomial in \(\lambda\) --- the \textbf{characteristic
polynomial}. Its roots are the eigenvalues of \(A\).

For an \(n \times n\) matrix, the characteristic polynomial has degree
\(n\), so there are exactly \(n\) eigenvalues (counting multiplicity,
and allowing complex values).

\textbf{For a 2×2 matrix:}

\[A = \begin{pmatrix} a & b \\ c & d \end{pmatrix}\]

\[\det(A - \lambda I) = \det\begin{pmatrix} a-\lambda & b \\ c & d-\lambda \end{pmatrix} = (a-\lambda)(d-\lambda) - bc\]

Expanding:

\[\lambda^2 - (a+d)\lambda + (ad - bc) = 0\]

\[\lambda^2 - \text{tr}(A)\,\lambda + \det(A) = 0\]

The trace \(\text{tr}(A) = a + d\) (sum of diagonal entries) and
\(\det(A) = ad - bc\) appear directly in the characteristic polynomial.
This is a useful check: the sum of the eigenvalues equals the trace, and
their product equals the determinant.

\begin{tcolorbox}[enhanced jigsaw, opacityback=0, colback=white, breakable, colframe=quarto-callout-tip-color-frame, arc=.35mm, left=2mm, toprule=.15mm, leftrule=.75mm, bottomrule=.15mm, rightrule=.15mm]
\begin{minipage}[t]{5.5mm}
\textcolor{quarto-callout-tip-color}{\faLightbulb}
\end{minipage}%
\begin{minipage}[t]{\textwidth - 5.5mm}

\textbf{Trace--determinant check}

After finding eigenvalues \(\lambda_1, \lambda_2\):

\[\lambda_1 + \lambda_2 = \text{tr}(A), \qquad \lambda_1 \lambda_2 = \det(A)\]

Always verify this before proceeding. A quick arithmetic error in the
characteristic polynomial shows up immediately.

\end{minipage}%
\end{tcolorbox}

\section{Finding eigenvalues and eigenvectors --- worked
examples}\label{finding-eigenvalues-and-eigenvectors-worked-examples}

\subsection{Example 1: A 2×2 electrical network
matrix}\label{example-1-a-22-electrical-network-matrix}

A simplified two-node circuit has the nodal admittance matrix

\[A = \begin{pmatrix} 3 & -1 \\ -1 & 3 \end{pmatrix}\]

\textbf{Step 1: Characteristic equation.}

\[\det(A - \lambda I) = \det\begin{pmatrix} 3-\lambda & -1 \\ -1 & 3-\lambda \end{pmatrix} = (3-\lambda)^2 - 1 = 0\]

\[9 - 6\lambda + \lambda^2 - 1 = 0\]

\[\lambda^2 - 6\lambda + 8 = 0\]

\textbf{Step 2: Solve the quadratic.}

\[(\lambda - 2)(\lambda - 4) = 0 \implies \lambda_1 = 2, \quad \lambda_2 = 4\]

\emph{Check:} \(\lambda_1 + \lambda_2 = 6 = \text{tr}(A)\) and
\(\lambda_1 \lambda_2 = 8 = \det(A)\). Yes.

\textbf{Step 3: Eigenvector for \(\lambda_1 = 2\).}

Solve \((A - 2I)\mathbf{x} = \mathbf{0}\):

\[\begin{pmatrix} 1 & -1 \\ -1 & 1 \end{pmatrix}\begin{pmatrix} x_1 \\ x_2 \end{pmatrix} = \begin{pmatrix} 0 \\ 0 \end{pmatrix}\]

The two rows say the same thing: \(x_1 - x_2 = 0\), so \(x_1 = x_2\).
One free variable --- take \(x_2 = 1\):

\[\mathbf{x}_1 = \begin{pmatrix} 1 \\ 1 \end{pmatrix}\]

\textbf{Step 4: Eigenvector for \(\lambda_2 = 4\).}

Solve \((A - 4I)\mathbf{x} = \mathbf{0}\):

\[\begin{pmatrix} -1 & -1 \\ -1 & -1 \end{pmatrix}\begin{pmatrix} x_1 \\ x_2 \end{pmatrix} = \begin{pmatrix} 0 \\ 0 \end{pmatrix}\]

The equation is \(-x_1 - x_2 = 0\), so \(x_1 = -x_2\). Take \(x_2 = 1\):

\[\mathbf{x}_2 = \begin{pmatrix} -1 \\ 1 \end{pmatrix}\]

Notice the eigenvectors are perpendicular: \((1)(−1) + (1)(1) = 0\).
This is not a coincidence --- \(A\) is symmetric. Symmetric matrices
always have orthogonal eigenvectors (more on this in the spectral
theorem section).

\subsection{Example 2: A 3×3 two-storey building stiffness
matrix}\label{example-2-a-33-two-storey-building-stiffness-matrix}

A two-storey shear building with a rigid base has three degrees of
freedom (ground, first floor, second floor). With floor masses
normalised to 1 and storey stiffnesses \(k = 1\), the equation of motion
\(M\ddot{u} + Ku = 0\) (with \(M = I\)) reduces to an eigenvalue problem
on the stiffness matrix:

\[K = \begin{pmatrix} 2 & -1 & 0 \\ -1 & 2 & -1 \\ 0 & -1 & 1 \end{pmatrix}\]

The eigenvalues of \(K\) are the squares of the natural frequencies.

\textbf{Step 1: Characteristic polynomial.}

\[\det(K - \lambda I) = \det\begin{pmatrix} 2-\lambda & -1 & 0 \\ -1 & 2-\lambda & -1 \\ 0 & -1 & 1-\lambda \end{pmatrix}\]

Expanding along the first row:

\[= (2-\lambda)\det\begin{pmatrix} 2-\lambda & -1 \\ -1 & 1-\lambda \end{pmatrix} - (-1)\det\begin{pmatrix} -1 & -1 \\ 0 & 1-\lambda \end{pmatrix}\]

\[= (2-\lambda)\bigl[(2-\lambda)(1-\lambda) - 1\bigr] + \bigl[(-1)(1-\lambda) - 0\bigr]\]

\[= (2-\lambda)(2 - 3\lambda + \lambda^2 - 1) + (-(1-\lambda))\]

\[= (2-\lambda)(\lambda^2 - 3\lambda + 1) - (1 - \lambda)\]

Expanding \((2-\lambda)(\lambda^2 - 3\lambda + 1)\):

\[= 2\lambda^2 - 6\lambda + 2 - \lambda^3 + 3\lambda^2 - \lambda\]

\[= -\lambda^3 + 5\lambda^2 - 7\lambda + 2\]

Subtracting \((1 - \lambda)\):

\[\det(K - \lambda I) = -\lambda^3 + 5\lambda^2 - 6\lambda + 1\]

Or equivalently: \(\lambda^3 - 5\lambda^2 + 6\lambda - 1 = 0\).

\textbf{Step 2: Roots.} Numerically (or via Cardano's formula): the
three roots are approximately

\[\lambda_1 \approx 0.198, \quad \lambda_2 \approx 1.555, \quad \lambda_3 \approx 3.247\]

The natural frequencies are \(\omega_i = \sqrt{\lambda_i}\):

\[\omega_1 \approx 0.445, \quad \omega_2 \approx 1.247, \quad \omega_3 \approx 1.802 \text{ rad/s}\]

\begin{tcolorbox}[enhanced jigsaw, opacityback=0, colback=white, breakable, colframe=quarto-callout-note-color-frame, arc=.35mm, left=2mm, toprule=.15mm, leftrule=.75mm, bottomrule=.15mm, rightrule=.15mm]
\begin{minipage}[t]{5.5mm}
\textcolor{quarto-callout-note-color}{\faInfo}
\end{minipage}%
\begin{minipage}[t]{\textwidth - 5.5mm}

\textbf{3×3 and larger: numerical methods in practice}

For \(3 \times 3\) matrices, the characteristic polynomial can still be
solved by hand --- but it is tedious. For \(4 \times 4\) and beyond,
exact polynomial roots are impossible in general (Abel--Ruffini
theorem). In practice, engineers use numerical algorithms --- the QR
algorithm, Jacobi iteration, or Lanczos methods. The point of working
through the \(3 \times 3\) case by hand is to understand what the
algorithm is computing, not to match its speed.

\end{minipage}%
\end{tcolorbox}

\section{Repeated eigenvalues}\label{repeated-eigenvalues}

When the characteristic polynomial has a repeated root \(\lambda\) of
multiplicity \(m > 1\), two cases arise:

\textbf{Case 1: Non-defective (enough eigenvectors).} The eigenspace for
\(\lambda\) has dimension \(m\) --- you can find \(m\) linearly
independent eigenvectors. The matrix is diagonalisable.

Example: \(A = \lambda I\) (a scalar matrix). Every vector is an
eigenvector. The eigenspace is the whole \(\mathbb{R}^n\).

\textbf{Case 2: Defective (not enough eigenvectors).} The eigenspace for
\(\lambda\) has dimension less than \(m\). You cannot find enough
linearly independent eigenvectors to fill out a basis. The matrix is
\textbf{not diagonalisable} --- you need the Jordan normal form instead.

\[A = \begin{pmatrix} 2 & 1 \\ 0 & 2 \end{pmatrix}\]

Characteristic polynomial: \((\lambda - 2)^2 = 0\), so \(\lambda = 2\)
with multiplicity 2.

Eigenvectors: \((A - 2I)\mathbf{x} = \mathbf{0}\) gives
\(\begin{pmatrix} 0 & 1 \\ 0 & 0 \end{pmatrix}\mathbf{x} = \mathbf{0}\),
so \(x_2 = 0\) and \(x_1\) is free. There is only one linearly
independent eigenvector:
\(\mathbf{x} = \begin{pmatrix} 1 \\ 0 \end{pmatrix}\).

Algebraic multiplicity (from the polynomial) is 2. Geometric
multiplicity (dimension of eigenspace) is 1. Because they differ, \(A\)
is \textbf{defective}.

\begin{tcolorbox}[enhanced jigsaw, opacityback=0, colback=white, breakable, colframe=quarto-callout-tip-color-frame, arc=.35mm, left=2mm, toprule=.15mm, leftrule=.75mm, bottomrule=.15mm, rightrule=.15mm]
\begin{minipage}[t]{5.5mm}
\textcolor{quarto-callout-tip-color}{\faLightbulb}
\end{minipage}%
\begin{minipage}[t]{\textwidth - 5.5mm}

\textbf{Why this matters in engineering}

A defective matrix arises when two natural frequencies of a system
coincide exactly --- a degenerate case. Small perturbations (real-world
manufacturing tolerances) break the degeneracy, splitting the repeated
eigenvalue into two distinct ones. Defective matrices are a theoretical
boundary case. In practice, if your numerical computation returns a
repeated eigenvalue, check whether it reflects genuine physics or
accumulated rounding error.

\end{minipage}%
\end{tcolorbox}

\section{Diagonalisation}\label{diagonalisation}

If an \(n \times n\) matrix \(A\) has \(n\) linearly independent
eigenvectors \(\mathbf{x}_1, \ldots, \mathbf{x}_n\) with eigenvalues
\(\lambda_1, \ldots, \lambda_n\), form the matrices:

\[P = \begin{pmatrix} | & & | \\ \mathbf{x}_1 & \cdots & \mathbf{x}_n \\ | & & | \end{pmatrix}, \qquad D = \begin{pmatrix} \lambda_1 & & \\ & \ddots & \\ & & \lambda_n \end{pmatrix}\]

Then:

\[A = PDP^{-1}\]

\(P\) is the \textbf{modal matrix} (columns are eigenvectors); \(D\) is
\textbf{diagonal} with eigenvalues on the diagonal.

\textbf{Why diagonalisation is useful.}

Powers of a diagonal matrix are trivial: \(D^k\) is just each diagonal
entry raised to the \(k\)-th power. So:

\[A^k = PD^kP^{-1}\]

This makes computing \(A^{100}\) as cheap as computing
\(\lambda_1^{100}, \ldots,
\lambda_n^{100}\) --- a dramatic saving. It also makes solving
\(\dot{\mathbf{u}} = A\mathbf{u}\) straightforward: substitute
\(\mathbf{u} = P\mathbf{v}\) to decouple the system into \(n\)
independent scalar ODEs, each of the form \(\dot{v}_i = \lambda_i v_i\).

\textbf{Worked diagonalisation of Example 1.}

From Example 1: \(\lambda_1 = 2\),
\(\mathbf{x}_1 = \begin{pmatrix} 1 \\ 1 \end{pmatrix}\);
\(\lambda_2 = 4\),
\(\mathbf{x}_2 = \begin{pmatrix} -1 \\ 1 \end{pmatrix}\).

\[P = \begin{pmatrix} 1 & -1 \\ 1 & 1 \end{pmatrix}, \qquad D = \begin{pmatrix} 2 & 0 \\ 0 & 4 \end{pmatrix}\]

\[P^{-1} = \frac{1}{\det P}\begin{pmatrix} 1 & 1 \\ -1 & 1 \end{pmatrix} = \frac{1}{2}\begin{pmatrix} 1 & 1 \\ -1 & 1 \end{pmatrix}\]

Verify \(AP = PD\):

\[AP = \begin{pmatrix} 3 & -1 \\ -1 & 3 \end{pmatrix}\begin{pmatrix} 1 & -1 \\ 1 & 1 \end{pmatrix} = \begin{pmatrix} 2 & -4 \\ 2 & 4 \end{pmatrix}\]

\[PD = \begin{pmatrix} 1 & -1 \\ 1 & 1 \end{pmatrix}\begin{pmatrix} 2 & 0 \\ 0 & 4 \end{pmatrix} = \begin{pmatrix} 2 & -4 \\ 2 & 4 \end{pmatrix} \checkmark\]

\section{Symmetric matrices and the spectral
theorem}\label{symmetric-matrices-and-the-spectral-theorem}

A matrix \(A\) is \textbf{symmetric} if \(A = A^T\) (the transpose
equals itself). Symmetric matrices arise everywhere in physics and
engineering: stiffness matrices, covariance matrices, the inertia
tensor. They have two beautiful properties:

\begin{tcolorbox}[enhanced jigsaw, opacityback=0, colback=white, breakable, colframe=quarto-callout-important-color-frame, arc=.35mm, left=2mm, toprule=.15mm, leftrule=.75mm, bottomrule=.15mm, rightrule=.15mm]
\begin{minipage}[t]{5.5mm}
\textcolor{quarto-callout-important-color}{\faExclamation}
\end{minipage}%
\begin{minipage}[t]{\textwidth - 5.5mm}

\textbf{Spectral theorem (stated)}

If \(A\) is a real symmetric \(n \times n\) matrix, then:

\begin{enumerate}
\def\labelenumi{\arabic{enumi}.}
\tightlist
\item
  All eigenvalues of \(A\) are \textbf{real}.
\item
  Eigenvectors corresponding to distinct eigenvalues are
  \textbf{orthogonal}.
\item
  \(A\) has \(n\) linearly independent eigenvectors, so it is always
  diagonalisable.
\item
  The eigenvectors can be chosen to be an \textbf{orthonormal basis} of
  \(\mathbb{R}^n\).
\end{enumerate}

This means \(A = Q D Q^T\) where \(Q\) is an \textbf{orthogonal} matrix
(\(Q^T = Q^{-1}\)) whose columns are orthonormal eigenvectors. This is
the \textbf{eigendecomposition} or \textbf{spectral decomposition} of
\(A\).

\end{minipage}%
\end{tcolorbox}

\textbf{Why engineers care.} Stiffness matrices are symmetric (by
Maxwell's reciprocal theorem). The spectral theorem guarantees that
their eigenvalues (natural frequencies squared) are real and positive
--- physically, modes must vibrate at real frequencies, not grow
exponentially. The orthogonality of mode shapes means each mode can be
analysed independently, which is the foundation of modal analysis.

\textbf{Principal axes.} The eigendecomposition of a symmetric matrix
gives the \textbf{principal axes} of the associated quadratic form. For
a 3D stress tensor, the eigenvalues are the principal stresses and the
eigenvectors are the principal stress directions --- the coordinate
system in which the stress tensor is diagonal (no shear stress). Every
solid mechanics textbook derives the principal stress transformation
this way.

\section{Applications}\label{applications-6}

\subsection{Vibrational modes of a two-mass spring
system}\label{vibrational-modes-of-a-two-mass-spring-system}

Two masses connected by three springs form the simplest mechanical model
of a multi-storey building or a coupled oscillator. With masses
\(m_1 = 1\), \(m_2 = 2\) and spring constants \(k_1 = k_2 = k_3 = 1\):

The equations of motion are:

\[m_1 \ddot{u}_1 = -k_1 u_1 + k_2(u_2 - u_1) = -(k_1+k_2)u_1 + k_2 u_2\]
\[m_2 \ddot{u}_2 = -k_2(u_2 - u_1) - k_3 u_2 = k_2 u_1 - (k_2+k_3)u_2\]

In matrix form \(M\ddot{\mathbf{u}} + K\mathbf{u} = \mathbf{0}\) with
\(M = \text{diag}(1, 2)\):

\[K = \begin{pmatrix} 2 & -1 \\ -1 & 2 \end{pmatrix}\]

For harmonic motion \(\mathbf{u} = \mathbf{x} e^{i\omega t}\), this
becomes the \textbf{generalised eigenvalue problem}:

\[K\mathbf{x} = \omega^2 M\mathbf{x}\]

Pre-multiply by \(M^{-1} = \text{diag}(1, \tfrac{1}{2})\):

\[M^{-1}K\mathbf{x} = \omega^2\mathbf{x}, \qquad M^{-1}K = \begin{pmatrix} 2 & -1 \\ -\tfrac{1}{2} & 1 \end{pmatrix}\]

Characteristic equation: \(\det(M^{-1}K - \omega^2 I) = 0\):

\[(2 - \omega^2)(1 - \omega^2) - \tfrac{1}{2} = 0\]

\[\omega^4 - 3\omega^2 + \tfrac{3}{2} = 0\]

Using the quadratic formula with \(\mu = \omega^2\):

\[\mu = \frac{3 \pm \sqrt{9 - 6}}{2} = \frac{3 \pm \sqrt{3}}{2}\]

\[\omega_1^2 = \frac{3 - \sqrt{3}}{2} \approx 0.634, \quad \omega_1 \approx 0.796 \text{ rad/s}\]
\[\omega_2^2 = \frac{3 + \sqrt{3}}{2} \approx 2.366, \quad \omega_2 \approx 1.538 \text{ rad/s}\]

These are the two \textbf{natural frequencies}. At \(\omega_1\), both
masses oscillate in phase (the first mode shape). At \(\omega_2\), they
oscillate out of phase (the second mode shape). If you drive this system
at either of these frequencies, resonance occurs --- amplitudes grow
without bound (in the undamped model).

\subsection{Power iteration and
PageRank}\label{power-iteration-and-pagerank}

Google's PageRank assigns importance scores to web pages based on the
link structure of the web. Define the transition matrix \(G\) (the
``Google matrix'') where \(G_{ij}\) is proportional to the probability
of following a link from page \(j\) to page \(i\).

PageRank scores are the entries of the \textbf{dominant eigenvector} of
\(G\) --- the eigenvector corresponding to the largest eigenvalue (which
is always 1 for a properly normalised transition matrix).

Finding the exact dominant eigenvector of a billion-by-billion matrix is
impossible by direct methods. Instead, Google uses \textbf{power
iteration}:

\begin{enumerate}
\def\labelenumi{\arabic{enumi}.}
\tightlist
\item
  Start with any vector \(\mathbf{v}^{(0)}\) (e.g., uniform: all pages
  equally likely).
\item
  Multiply: \(\mathbf{v}^{(k+1)} = G\mathbf{v}^{(k)}\).
\item
  Normalise: divide by \(\|\mathbf{v}^{(k+1)}\|\).
\item
  Repeat until convergence.
\end{enumerate}

Under mild conditions, this converges to the dominant eigenvector. The
convergence rate depends on the ratio \(|\lambda_2/\lambda_1|\) --- how
much smaller the second eigenvalue is than the first.

\begin{Shaded}
\begin{Highlighting}[]
\NormalTok{//| echo: false}
\NormalTok{powerIterationDemo = \{}
\NormalTok{  const d3 = await require("d3@7");}
\NormalTok{  const W = 560, H = 400;}
\NormalTok{  const ML = 40, MR = 40, MT = 40, MB = 40;}
\NormalTok{  const innerW = W {-} ML {-} MR;}
\NormalTok{  const innerH = H {-} MT {-} MB;}
  
\NormalTok{  const container = d3.create("div").style("max{-}width", W+"px").style("font{-}family", "sans{-}serif").style("margin", "0 auto");}
  
\NormalTok{  const controls = container.append("div").style("display", "flex").style("gap", "20px").style("margin{-}bottom", "10px").style("align{-}items", "center").style("background", "\#f8fafc").style("padding", "10px").style("border{-}radius", "8px").style("border", "1px solid \#e2e8f0");}
  
\NormalTok{  const matDiv = controls.append("div").style("display", "grid").style("grid{-}template{-}columns", "1fr 1fr").style("gap", "5px");}
\NormalTok{  function makeInput(parent, initVal) \{}
\NormalTok{    return parent.append("input").attr("type", "number").attr("step", "0.1").attr("value", initVal).style("width", "50px").style("text{-}align", "center").style("font{-}family", "monospace");}
\NormalTok{  \}}
\NormalTok{  const inpA = makeInput(matDiv, 2);}
\NormalTok{  const inpB = makeInput(matDiv, 1);}
\NormalTok{  const inpC = makeInput(matDiv, 1);}
\NormalTok{  const inpD = makeInput(matDiv, 3);}
  
\NormalTok{  const btnDiv = controls.append("div").style("display", "flex").style("gap", "10px");}
\NormalTok{  const stepBtn = btnDiv.append("button").text("Step").style("padding", "5px 12px").style("cursor", "pointer").style("background", "\#3b82f6").style("color", "white").style("border", "none").style("border{-}radius", "4px").style("font{-}weight", "500");}
\NormalTok{  const resetBtn = btnDiv.append("button").text("Reset").style("padding", "5px 12px").style("cursor", "pointer").style("background", "\#e2e8f0").style("color", "\#475569").style("border", "none").style("border{-}radius", "4px").style("font{-}weight", "500");}
  
\NormalTok{  const kText = controls.append("div").style("font{-}weight", "bold").style("font{-}size", "1.1rem").style("color", "\#1e293b").style("margin{-}left", "auto").text("k = 0");}
  
\NormalTok{  const svg = container.append("svg").attr("width", W).attr("height", H).attr("viewBox", \textasciigrave{}0 0 $\{W\} $\{H\}\textasciigrave{});}
\NormalTok{  const g = svg.append("g").attr("transform", \textasciigrave{}translate($\{W/2\},$\{H/2\})\textasciigrave{});}
  
\NormalTok{  const xScale = d3.scaleLinear([{-}1.5, 1.5], [{-}innerW/2, innerW/2]);}
\NormalTok{  const yScale = d3.scaleLinear([{-}1.5, 1.5], [innerH/2, {-}innerH/2]);}
  
\NormalTok{  g.append("line").attr("x1", xScale({-}1.5)).attr("y1", yScale(0)).attr("x2", xScale(1.5)).attr("y2", yScale(0)).attr("stroke", "\#e5e7eb");}
\NormalTok{  g.append("line").attr("x1", xScale(0)).attr("y1", yScale({-}1.5)).attr("x2", xScale(0)).attr("y2", yScale(1.5)).attr("stroke", "\#e5e7eb");}
\NormalTok{  g.append("circle").attr("r", xScale(1) {-} xScale(0)).attr("fill", "none").attr("stroke", "\#9ca3af");}
  
\NormalTok{  const tracePath = g.append("path").attr("fill", "none").attr("stroke", "\#cbd5e1").attr("stroke{-}width", 2);}
  
\NormalTok{  const targetLine = g.append("line").attr("stroke", "\#9ca3af").attr("stroke{-}width", 2).attr("stroke{-}dasharray", "5,5");}
\NormalTok{  const targetHead = g.append("polygon").attr("fill", "\#9ca3af");}
\NormalTok{  const targetLabel = g.append("text").attr("fill", "\#6b7280").attr("font{-}size", "12px").text(" Dominant Vector");}
  
\NormalTok{  const iterLine = g.append("line").attr("stroke", "\#3b82f6").attr("stroke{-}width", 3);}
\NormalTok{  const iterHead = g.append("polygon").attr("fill", "\#3b82f6");}
  
\NormalTok{  let x\_k = [1, 0];}
\NormalTok{  let k = 0;}
\NormalTok{  let trace = [[1,0]];}
  
\NormalTok{  function getDominantEigVec(a, b, c, d) \{}
\NormalTok{    const tr = a + d;}
\NormalTok{    const det = a*d {-} b*c;}
\NormalTok{    const disc = tr*tr {-} 4*det;}
\NormalTok{    if (disc \textless{} 0) return null;}
\NormalTok{    const lam1 = (tr + Math.sqrt(disc)) / 2;}
\NormalTok{    const lam2 = (tr {-} Math.sqrt(disc)) / 2;}
\NormalTok{    const dom = Math.abs(lam1) \textgreater{}= Math.abs(lam2) ? lam1 : lam2;}
    
\NormalTok{    let v;}
\NormalTok{    if (Math.abs(b) \textgreater{} 1e{-}6) v = [b, dom {-} a];}
\NormalTok{    else if (Math.abs(c) \textgreater{} 1e{-}6) v = [dom {-} d, c];}
\NormalTok{    else v = dom === a ? [1, 0] : [0, 1];}
    
\NormalTok{    const mag = Math.sqrt(v[0]*v[0] + v[1]*v[1]);}
\NormalTok{    return [v[0]/mag, v[1]/mag];}
\NormalTok{  \}}
  
\NormalTok{  function drawVec(v, line, head, scale = 1) \{}
\NormalTok{    const sx = xScale(v[0]*scale);}
\NormalTok{    const sy = yScale(v[1]*scale);}
\NormalTok{    const ox = xScale(0);}
\NormalTok{    const oy = yScale(0);}
\NormalTok{    line.attr("x1", ox).attr("y1", oy).attr("x2", sx).attr("y2", sy);}
\NormalTok{    const angle = Math.atan2(sy {-} oy, sx {-} ox);}
\NormalTok{    const r = 10;}
\NormalTok{    head.attr("points", \textasciigrave{}$\{sx\},$\{sy\} $\{sx {-} r*Math.cos(angle {-} 0.5)\},$\{sy {-} r*Math.sin(angle {-} 0.5)\} $\{sx {-} r*Math.cos(angle + 0.5)\},$\{sy {-} r*Math.sin(angle + 0.5)\}\textasciigrave{});}
\NormalTok{    return \{sx, sy\};}
\NormalTok{  \}}
  
\NormalTok{  function update() \{}
\NormalTok{    const a = +inpA.property("value");}
\NormalTok{    const b = +inpB.property("value");}
\NormalTok{    const c = +inpC.property("value");}
\NormalTok{    const d = +inpD.property("value");}
    
\NormalTok{    const domV = getDominantEigVec(a, b, c, d);}
    
\NormalTok{    if (domV) \{}
\NormalTok{      targetLine.attr("visibility", "visible");}
\NormalTok{      targetHead.attr("visibility", "visible");}
\NormalTok{      targetLabel.attr("visibility", "visible");}
      
\NormalTok{      const \{sx, sy\} = drawVec(domV, targetLine, targetHead);}
\NormalTok{      targetLabel.attr("x", sx + 12*Math.sign(domV[0])).attr("y", sy + 15*Math.sign(domV[1]) {-} 5);}
      
\NormalTok{      // Also potentially draw the negative dominant direction since power iteration can flip}
\NormalTok{      if(domV[0]*x\_k[0] + domV[1]*x\_k[1] \textless{} 0) \{}
\NormalTok{          drawVec([{-}domV[0], {-}domV[1]], targetLine, targetHead);}
\NormalTok{          targetLabel.attr("x", xScale({-}domV[0]) + 12*Math.sign({-}domV[0])).attr("y", yScale({-}domV[1]) + 15*Math.sign({-}domV[1]) {-} 5);}
\NormalTok{      \} else \{}
\NormalTok{          drawVec(domV, targetLine, targetHead);}
\NormalTok{          targetLabel.attr("x", sx + 12*Math.sign(domV[0])).attr("y", sy + 15*Math.sign(domV[1]) {-} 5);}
\NormalTok{      \}}
\NormalTok{    \} else \{}
\NormalTok{      targetLine.attr("visibility", "hidden");}
\NormalTok{      targetHead.attr("visibility", "hidden");}
\NormalTok{      targetLabel.attr("visibility", "hidden");}
\NormalTok{    \}}
    
\NormalTok{    drawVec(x\_k, iterLine, iterHead);}
\NormalTok{    kText.text(\textasciigrave{}k = $\{k\}\textasciigrave{});}
    
\NormalTok{    const lineGen = d3.line().x(d =\textgreater{} xScale(d[0])).y(d =\textgreater{} yScale(d[1]));}
\NormalTok{    tracePath.datum(trace).attr("d", lineGen);}
\NormalTok{  \}}
  
\NormalTok{  stepBtn.on("click", () =\textgreater{} \{}
\NormalTok{    const a = +inpA.property("value");}
\NormalTok{    const b = +inpB.property("value");}
\NormalTok{    const c = +inpC.property("value");}
\NormalTok{    const d = +inpD.property("value");}
    
\NormalTok{    const newX = a*x\_k[0] + b*x\_k[1];}
\NormalTok{    const newY = c*x\_k[0] + d*x\_k[1];}
\NormalTok{    const mag = Math.sqrt(newX*newX + newY*newY);}
\NormalTok{    if(mag \textgreater{} 0) \{}
\NormalTok{      x\_k = [newX/mag, newY/mag];}
\NormalTok{      trace.push([...x\_k]);}
\NormalTok{      k++;}
\NormalTok{      update();}
\NormalTok{    \}}
\NormalTok{  \});}
  
\NormalTok{  resetBtn.on("click", () =\textgreater{} \{}
\NormalTok{    x\_k = [1, 0];}
\NormalTok{    trace = [[1,0]];}
\NormalTok{    k = 0;}
\NormalTok{    update();}
\NormalTok{  \});}
  
\NormalTok{  inpA.on("change", () =\textgreater{} resetBtn.node().click());}
\NormalTok{  inpB.on("change", () =\textgreater{} resetBtn.node().click());}
\NormalTok{  inpC.on("change", () =\textgreater{} resetBtn.node().click());}
\NormalTok{  inpD.on("change", () =\textgreater{} resetBtn.node().click());}
  
\NormalTok{  update();}
  
\NormalTok{  return container.node();}
\NormalTok{\}}
\end{Highlighting}
\end{Shaded}

\subsection{Principal component analysis (PCA) --- a
preview}\label{principal-component-analysis-pca-a-preview}

A dataset of \(n\) observations in \(p\) variables can be represented as
a matrix. The \textbf{sample covariance matrix} \(\Sigma\) is symmetric,
so the spectral theorem applies: it has real eigenvalues and orthogonal
eigenvectors.

The eigenvectors of \(\Sigma\) are the \textbf{principal components} ---
the directions of maximum variance in the data. The corresponding
eigenvalues measure how much variance each direction accounts for.

Projecting data onto the first \(k\) principal components is the best
rank-\(k\) approximation of the data (in the least-squares sense). This
is PCA. A face recognition system, a genomics pipeline, and a financial
risk model all use the same eigendecomposition under the hood.

\begin{tcolorbox}[enhanced jigsaw, opacityback=0, colback=white, breakable, colframe=quarto-callout-tip-color-frame, arc=.35mm, left=2mm, toprule=.15mm, leftrule=.75mm, bottomrule=.15mm, rightrule=.15mm]
\begin{minipage}[t]{5.5mm}
\textcolor{quarto-callout-tip-color}{\faLightbulb}
\end{minipage}%
\begin{minipage}[t]{\textwidth - 5.5mm}

\textbf{Where this shows up}

\begin{itemize}
\tightlist
\item
  A structural engineer computing mode shapes and natural frequencies of
  a building frame is solving an eigenvalue problem on a stiffness
  matrix.
\item
  A quantum chemist finding molecular orbital energies is diagonalising
  a Hamiltonian matrix.
\item
  A data scientist reducing 1000-dimensional embeddings to 2D for
  visualisation is computing the top two eigenvectors of a covariance
  matrix.
\item
  A search engine ranking web pages is finding the dominant eigenvector
  of a link matrix.
\end{itemize}

The matrix changes. The method does not.

\end{minipage}%
\end{tcolorbox}

\section{Exercises}\label{exercises-31}

These exercises move from mechanical computation to physical
interpretation. In each case: set up the eigenvalue problem, find the
eigenvalues and eigenvectors, and state what they mean in the
engineering context.

\begin{Shaded}
\begin{Highlighting}[]
\NormalTok{//| echo: false}
\NormalTok{katex = require("katex@0.16")}
\end{Highlighting}
\end{Shaded}

\begin{Shaded}
\begin{Highlighting}[]
\NormalTok{//| echo: false}
\NormalTok{\{}
\NormalTok{  const css = document.createElement("link");}
\NormalTok{  css.rel = "stylesheet";}
\NormalTok{  css.href = "https://cdn.jsdelivr.net/npm/katex@0.16/dist/katex.min.css";}
\NormalTok{  document.head.appendChild(css);}
\NormalTok{\}}
\end{Highlighting}
\end{Shaded}

\begin{Shaded}
\begin{Highlighting}[]
\NormalTok{//| echo: false}
\NormalTok{function makeStepperHTML(exerciseNum, steps) \{}
\NormalTok{  let current = 0;}
\NormalTok{  const totalSteps = steps.length;}
\NormalTok{  const container = document.createElement("div");}
\NormalTok{  container.style.cssText = "border:1px solid \#e5e7eb; border{-}radius:8px; padding:1rem 1.25rem; margin:0.75rem 0; font{-}family:inherit;";}
\NormalTok{  const stepsDiv = document.createElement("div");}
\NormalTok{  stepsDiv.style.marginBottom = "0.75rem";}

\NormalTok{  function renderSteps() \{}
\NormalTok{    stepsDiv.innerHTML = "";}
\NormalTok{    for (let i = 0; i \textless{} current; i++) \{}
\NormalTok{      const s = steps[i];}
\NormalTok{      const row = document.createElement("div");}
\NormalTok{      row.style.cssText = "display:grid; grid{-}template{-}columns:220px 1fr; gap:0.35rem 1rem; align{-}items:baseline; padding:0.35rem 0; border{-}top:1px solid \#f3f4f6;";}
\NormalTok{      const opCell = document.createElement("span");}
\NormalTok{      opCell.style.cssText = "font{-}size:0.85em; color:\#6b7280; font{-}style:italic;";}
\NormalTok{      opCell.textContent = s.op;}
\NormalTok{      const eqCell = document.createElement("span");}
\NormalTok{      eqCell.innerHTML = katex.renderToString(s.eq, \{ throwOnError: false, displayMode: false \});}
\NormalTok{      row.appendChild(opCell);}
\NormalTok{      row.appendChild(eqCell);}
\NormalTok{      if (s.note) \{}
\NormalTok{        const noteRow = document.createElement("div");}
\NormalTok{        noteRow.style.cssText = "grid{-}column:2; font{-}size:0.82em; color:\#6b7280; padding{-}bottom:0.2rem;";}
\NormalTok{        noteRow.textContent = s.note;}
\NormalTok{        row.appendChild(noteRow);}
\NormalTok{      \}}
\NormalTok{      stepsDiv.appendChild(row);}
\NormalTok{    \}}
\NormalTok{    if (current === totalSteps) \{}
\NormalTok{      const done = document.createElement("div");}
\NormalTok{      done.style.cssText = "margin{-}top:0.5rem; font{-}size:0.9em; color:\#059669; font{-}weight:500;";}
\NormalTok{      done.textContent = "✓ Solution complete";}
\NormalTok{      stepsDiv.appendChild(done);}
\NormalTok{    \}}
\NormalTok{  \}}

\NormalTok{  const controls = document.createElement("div");}
\NormalTok{  controls.style.cssText = "display:flex; gap:0.5rem; align{-}items:center;";}
\NormalTok{  const nextBtn = document.createElement("button");}
\NormalTok{  nextBtn.textContent = "Next step →";}
\NormalTok{  nextBtn.style.cssText = "padding:0.35rem 0.85rem; border:1px solid \#d1d5db; border{-}radius:4px; background:\#fff; cursor:pointer; font{-}size:0.9em;";}
\NormalTok{  const resetBtn = document.createElement("button");}
\NormalTok{  resetBtn.textContent = "Reset";}
\NormalTok{  resetBtn.style.cssText = "padding:0.35rem 0.75rem; border:1px solid \#e5e7eb; border{-}radius:4px; background:\#f9fafb; cursor:pointer; font{-}size:0.9em; color:\#6b7280;";}
\NormalTok{  const counter = document.createElement("span");}
\NormalTok{  counter.style.cssText = "font{-}size:0.82em; color:\#9ca3af; margin{-}left:0.25rem;";}

\NormalTok{  function updateButtons() \{}
\NormalTok{    nextBtn.disabled = current === totalSteps;}
\NormalTok{    nextBtn.style.opacity = current === totalSteps ? "0.4" : "1";}
\NormalTok{    counter.textContent = current === 0 ? "Click to reveal steps" : \textasciigrave{}Step $\{current\} of $\{totalSteps\}\textasciigrave{};}
\NormalTok{  \}}

\NormalTok{  nextBtn.onclick = () =\textgreater{} \{ if (current \textless{} totalSteps) \{ current++; renderSteps(); updateButtons(); \} \};}
\NormalTok{  resetBtn.onclick = () =\textgreater{} \{ current = 0; renderSteps(); updateButtons(); \};}
\NormalTok{  controls.appendChild(nextBtn);}
\NormalTok{  controls.appendChild(resetBtn);}
\NormalTok{  controls.appendChild(counter);}
\NormalTok{  container.appendChild(stepsDiv);}
\NormalTok{  container.appendChild(controls);}
\NormalTok{  renderSteps();}
\NormalTok{  updateButtons();}
\NormalTok{  return container;}
\NormalTok{\}}
\end{Highlighting}
\end{Shaded}

\textbf{1.} The nodal admittance matrix of a two-node electrical network
is \(A = \begin{pmatrix} 4 & -2 \\ -2 & 4 \end{pmatrix}\). Find the
characteristic polynomial and the two eigenvalues. (The eigenvalues are
the network's natural admittances --- their reciprocals are natural
impedances.)

\begin{Shaded}
\begin{Highlighting}[]
\NormalTok{//| echo: false}
\NormalTok{makeStepperHTML(1, [}
\NormalTok{  \{}
\NormalTok{    op: "Write the characteristic equation",}
\NormalTok{    eq: "\textbackslash{}\textbackslash{}det(A {-} \textbackslash{}\textbackslash{}lambda I) = \textbackslash{}\textbackslash{}det\textbackslash{}\textbackslash{}begin\{pmatrix\} 4{-}\textbackslash{}\textbackslash{}lambda \& {-}2 \textbackslash{}\textbackslash{}\textbackslash{}\textbackslash{} {-}2 \& 4{-}\textbackslash{}\textbackslash{}lambda \textbackslash{}\textbackslash{}end\{pmatrix\} = 0",}
\NormalTok{    note: "Non{-}trivial solutions exist only when A − λI is singular."}
\NormalTok{  \},}
\NormalTok{  \{}
\NormalTok{    op: "Expand the determinant",}
\NormalTok{    eq: "(4{-}\textbackslash{}\textbackslash{}lambda)\^{}2 {-} ({-}2)({-}2) = 0"}
\NormalTok{  \},}
\NormalTok{  \{}
\NormalTok{    op: "Simplify",}
\NormalTok{    eq: "16 {-} 8\textbackslash{}\textbackslash{}lambda + \textbackslash{}\textbackslash{}lambda\^{}2 {-} 4 = 0 \textbackslash{}\textbackslash{}implies \textbackslash{}\textbackslash{}lambda\^{}2 {-} 8\textbackslash{}\textbackslash{}lambda + 12 = 0"}
\NormalTok{  \},}
\NormalTok{  \{}
\NormalTok{    op: "Factor the quadratic",}
\NormalTok{    eq: "(\textbackslash{}\textbackslash{}lambda {-} 2)(\textbackslash{}\textbackslash{}lambda {-} 6) = 0 \textbackslash{}\textbackslash{}implies \textbackslash{}\textbackslash{}lambda\_1 = 2,\textbackslash{}\textbackslash{} \textbackslash{}\textbackslash{}lambda\_2 = 6"}
\NormalTok{  \},}
\NormalTok{  \{}
\NormalTok{    op: "Check with trace and determinant",}
\NormalTok{    eq: "\textbackslash{}\textbackslash{}lambda\_1 + \textbackslash{}\textbackslash{}lambda\_2 = 8 = \textbackslash{}\textbackslash{}operatorname\{tr\}(A),\textbackslash{}\textbackslash{}quad \textbackslash{}\textbackslash{}lambda\_1 \textbackslash{}\textbackslash{}lambda\_2 = 12 = \textbackslash{}\textbackslash{}det(A)\textbackslash{}\textbackslash{} \textbackslash{}\textbackslash{}checkmark"}
\NormalTok{  \}}
\NormalTok{])}
\end{Highlighting}
\end{Shaded}

\textbf{2.} Using the eigenvalues from Exercise 1, find an eigenvector
for each eigenvalue. Then write down the diagonalisation
\(A = PDP^{-1}\).

\begin{Shaded}
\begin{Highlighting}[]
\NormalTok{//| echo: false}
\NormalTok{makeStepperHTML(2, [}
\NormalTok{  \{}
\NormalTok{    op: "Eigenvector for λ₁ = 2: solve (A − 2I)x = 0",}
\NormalTok{    eq: "\textbackslash{}\textbackslash{}begin\{pmatrix\} 2 \& {-}2 \textbackslash{}\textbackslash{}\textbackslash{}\textbackslash{} {-}2 \& 2 \textbackslash{}\textbackslash{}end\{pmatrix\}\textbackslash{}\textbackslash{}mathbf\{x\} = \textbackslash{}\textbackslash{}mathbf\{0\} \textbackslash{}\textbackslash{}implies x\_1 = x\_2",}
\NormalTok{    note: "Both rows give the same equation 2x₁ − 2x₂ = 0."}
\NormalTok{  \},}
\NormalTok{  \{}
\NormalTok{    op: "Take x₂ = 1",}
\NormalTok{    eq: "\textbackslash{}\textbackslash{}mathbf\{x\}\_1 = \textbackslash{}\textbackslash{}begin\{pmatrix\} 1 \textbackslash{}\textbackslash{}\textbackslash{}\textbackslash{} 1 \textbackslash{}\textbackslash{}end\{pmatrix\}"}
\NormalTok{  \},}
\NormalTok{  \{}
\NormalTok{    op: "Eigenvector for λ₂ = 6: solve (A − 6I)x = 0",}
\NormalTok{    eq: "\textbackslash{}\textbackslash{}begin\{pmatrix\} {-}2 \& {-}2 \textbackslash{}\textbackslash{}\textbackslash{}\textbackslash{} {-}2 \& {-}2 \textbackslash{}\textbackslash{}end\{pmatrix\}\textbackslash{}\textbackslash{}mathbf\{x\} = \textbackslash{}\textbackslash{}mathbf\{0\} \textbackslash{}\textbackslash{}implies x\_1 = {-}x\_2"}
\NormalTok{  \},}
\NormalTok{  \{}
\NormalTok{    op: "Take x₂ = 1",}
\NormalTok{    eq: "\textbackslash{}\textbackslash{}mathbf\{x\}\_2 = \textbackslash{}\textbackslash{}begin\{pmatrix\} {-}1 \textbackslash{}\textbackslash{}\textbackslash{}\textbackslash{} 1 \textbackslash{}\textbackslash{}end\{pmatrix\}"}
\NormalTok{  \},}
\NormalTok{  \{}
\NormalTok{    op: "Assemble P and D",}
\NormalTok{    eq: "P = \textbackslash{}\textbackslash{}begin\{pmatrix\} 1 \& {-}1 \textbackslash{}\textbackslash{}\textbackslash{}\textbackslash{} 1 \& 1 \textbackslash{}\textbackslash{}end\{pmatrix\},\textbackslash{}\textbackslash{}quad D = \textbackslash{}\textbackslash{}begin\{pmatrix\} 2 \& 0 \textbackslash{}\textbackslash{}\textbackslash{}\textbackslash{} 0 \& 6 \textbackslash{}\textbackslash{}end\{pmatrix\}"}
\NormalTok{  \},}
\NormalTok{  \{}
\NormalTok{    op: "Verify AP = PD (first column)",}
\NormalTok{    eq: "A\textbackslash{}\textbackslash{}mathbf\{x\}\_1 = \textbackslash{}\textbackslash{}begin\{pmatrix\} 4 \& {-}2 \textbackslash{}\textbackslash{}\textbackslash{}\textbackslash{} {-}2 \& 4 \textbackslash{}\textbackslash{}end\{pmatrix\}\textbackslash{}\textbackslash{}begin\{pmatrix\} 1 \textbackslash{}\textbackslash{}\textbackslash{}\textbackslash{} 1 \textbackslash{}\textbackslash{}end\{pmatrix\} = \textbackslash{}\textbackslash{}begin\{pmatrix\} 2 \textbackslash{}\textbackslash{}\textbackslash{}\textbackslash{} 2 \textbackslash{}\textbackslash{}end\{pmatrix\} = 2\textbackslash{}\textbackslash{}mathbf\{x\}\_1\textbackslash{}\textbackslash{} \textbackslash{}\textbackslash{}checkmark"}
\NormalTok{  \}}
\NormalTok{])}
\end{Highlighting}
\end{Shaded}

\textbf{3.} Diagonalise the symmetric matrix
\(B = \begin{pmatrix} 5 & 3 \\ 3 & 5 \end{pmatrix}\). Write
\(B = QDQ^T\) using orthonormal eigenvectors. (This matrix represents a
plane stress state --- its eigenvectors are the principal stress
directions.)

\begin{Shaded}
\begin{Highlighting}[]
\NormalTok{//| echo: false}
\NormalTok{makeStepperHTML(3, [}
\NormalTok{  \{}
\NormalTok{    op: "Characteristic equation",}
\NormalTok{    eq: "(5{-}\textbackslash{}\textbackslash{}lambda)\^{}2 {-} 9 = 0 \textbackslash{}\textbackslash{}implies \textbackslash{}\textbackslash{}lambda\^{}2 {-} 10\textbackslash{}\textbackslash{}lambda + 16 = 0"}
\NormalTok{  \},}
\NormalTok{  \{}
\NormalTok{    op: "Solve",}
\NormalTok{    eq: "\textbackslash{}\textbackslash{}lambda\_1 = 2,\textbackslash{}\textbackslash{}quad \textbackslash{}\textbackslash{}lambda\_2 = 8"}
\NormalTok{  \},}
\NormalTok{  \{}
\NormalTok{    op: "Eigenvector for λ₁ = 2",}
\NormalTok{    eq: "(B {-} 2I)\textbackslash{}\textbackslash{}mathbf\{x\} = \textbackslash{}\textbackslash{}mathbf\{0\}: \textbackslash{}\textbackslash{}begin\{pmatrix\} 3 \& 3 \textbackslash{}\textbackslash{}\textbackslash{}\textbackslash{} 3 \& 3 \textbackslash{}\textbackslash{}end\{pmatrix\}\textbackslash{}\textbackslash{}mathbf\{x\} = \textbackslash{}\textbackslash{}mathbf\{0\} \textbackslash{}\textbackslash{}implies \textbackslash{}\textbackslash{}mathbf\{x\}\_1 = \textbackslash{}\textbackslash{}begin\{pmatrix\} 1 \textbackslash{}\textbackslash{}\textbackslash{}\textbackslash{} {-}1 \textbackslash{}\textbackslash{}end\{pmatrix\}"}
\NormalTok{  \},}
\NormalTok{  \{}
\NormalTok{    op: "Eigenvector for λ₂ = 8",}
\NormalTok{    eq: "(B {-} 8I)\textbackslash{}\textbackslash{}mathbf\{x\} = \textbackslash{}\textbackslash{}mathbf\{0\}: \textbackslash{}\textbackslash{}begin\{pmatrix\} {-}3 \& 3 \textbackslash{}\textbackslash{}\textbackslash{}\textbackslash{} 3 \& {-}3 \textbackslash{}\textbackslash{}end\{pmatrix\}\textbackslash{}\textbackslash{}mathbf\{x\} = \textbackslash{}\textbackslash{}mathbf\{0\} \textbackslash{}\textbackslash{}implies \textbackslash{}\textbackslash{}mathbf\{x\}\_2 = \textbackslash{}\textbackslash{}begin\{pmatrix\} 1 \textbackslash{}\textbackslash{}\textbackslash{}\textbackslash{} 1 \textbackslash{}\textbackslash{}end\{pmatrix\}"}
\NormalTok{  \},}
\NormalTok{  \{}
\NormalTok{    op: "Normalise to unit length",}
\NormalTok{    eq: "\textbackslash{}\textbackslash{}mathbf\{q\}\_1 = \textbackslash{}\textbackslash{}frac\{1\}\{\textbackslash{}\textbackslash{}sqrt\{2\}\}\textbackslash{}\textbackslash{}begin\{pmatrix\} 1 \textbackslash{}\textbackslash{}\textbackslash{}\textbackslash{} {-}1 \textbackslash{}\textbackslash{}end\{pmatrix\},\textbackslash{}\textbackslash{}quad \textbackslash{}\textbackslash{}mathbf\{q\}\_2 = \textbackslash{}\textbackslash{}frac\{1\}\{\textbackslash{}\textbackslash{}sqrt\{2\}\}\textbackslash{}\textbackslash{}begin\{pmatrix\} 1 \textbackslash{}\textbackslash{}\textbackslash{}\textbackslash{} 1 \textbackslash{}\textbackslash{}end\{pmatrix\}",}
\NormalTok{    note: "Each eigenvector has length √2, so divide by √2 to normalise."}
\NormalTok{  \},}
\NormalTok{  \{}
\NormalTok{    op: "Write the orthogonal diagonalisation",}
\NormalTok{    eq: "Q = \textbackslash{}\textbackslash{}frac\{1\}\{\textbackslash{}\textbackslash{}sqrt\{2\}\}\textbackslash{}\textbackslash{}begin\{pmatrix\} 1 \& 1 \textbackslash{}\textbackslash{}\textbackslash{}\textbackslash{} {-}1 \& 1 \textbackslash{}\textbackslash{}end\{pmatrix\},\textbackslash{}\textbackslash{}quad D = \textbackslash{}\textbackslash{}begin\{pmatrix\} 2 \& 0 \textbackslash{}\textbackslash{}\textbackslash{}\textbackslash{} 0 \& 8 \textbackslash{}\textbackslash{}end\{pmatrix\},\textbackslash{}\textbackslash{}quad B = QDQ\^{}T\textbackslash{}\textbackslash{} \textbackslash{}\textbackslash{}checkmark"}
\NormalTok{  \}}
\NormalTok{])}
\end{Highlighting}
\end{Shaded}

\textbf{4.} A simplified three-storey shear building has mass matrix
\(M = I\) and stiffness matrix
\(K = \begin{pmatrix} 3 & -1 & 0 \\ -1 & 2 & -1 \\ 0 & -1 & 1 \end{pmatrix}\).
Find the characteristic polynomial and verify that
\(\lambda \approx 0.198\) is one eigenvalue by substituting it back.

\begin{Shaded}
\begin{Highlighting}[]
\NormalTok{//| echo: false}
\NormalTok{makeStepperHTML(4, [}
\NormalTok{  \{}
\NormalTok{    op: "Write det(K − λI) = 0",}
\NormalTok{    eq: "\textbackslash{}\textbackslash{}det\textbackslash{}\textbackslash{}begin\{pmatrix\} 3{-}\textbackslash{}\textbackslash{}lambda \& {-}1 \& 0 \textbackslash{}\textbackslash{}\textbackslash{}\textbackslash{} {-}1 \& 2{-}\textbackslash{}\textbackslash{}lambda \& {-}1 \textbackslash{}\textbackslash{}\textbackslash{}\textbackslash{} 0 \& {-}1 \& 1{-}\textbackslash{}\textbackslash{}lambda \textbackslash{}\textbackslash{}end\{pmatrix\} = 0"}
\NormalTok{  \},}
\NormalTok{  \{}
\NormalTok{    op: "Expand along row 1",}
\NormalTok{    eq: "(3{-}\textbackslash{}\textbackslash{}lambda)\textbackslash{}\textbackslash{}bigl[(2{-}\textbackslash{}\textbackslash{}lambda)(1{-}\textbackslash{}\textbackslash{}lambda){-}1\textbackslash{}\textbackslash{}bigr] + 1\textbackslash{}\textbackslash{}cdot\textbackslash{}\textbackslash{}bigl[({-}1)(1{-}\textbackslash{}\textbackslash{}lambda)\textbackslash{}\textbackslash{}bigr] = 0"}
\NormalTok{  \},}
\NormalTok{  \{}
\NormalTok{    op: "Simplify inner bracket",}
\NormalTok{    eq: "(2{-}\textbackslash{}\textbackslash{}lambda)(1{-}\textbackslash{}\textbackslash{}lambda) {-} 1 = \textbackslash{}\textbackslash{}lambda\^{}2 {-} 3\textbackslash{}\textbackslash{}lambda + 1"}
\NormalTok{  \},}
\NormalTok{  \{}
\NormalTok{    op: "Expand and collect",}
\NormalTok{    eq: "(3{-}\textbackslash{}\textbackslash{}lambda)(\textbackslash{}\textbackslash{}lambda\^{}2{-}3\textbackslash{}\textbackslash{}lambda+1) {-} (1{-}\textbackslash{}\textbackslash{}lambda) = {-}\textbackslash{}\textbackslash{}lambda\^{}3 + 6\textbackslash{}\textbackslash{}lambda\^{}2 {-} 10\textbackslash{}\textbackslash{}lambda + 4 = 0"}
\NormalTok{  \},}
\NormalTok{  \{}
\NormalTok{    op: "Divide by −1",}
\NormalTok{    eq: "\textbackslash{}\textbackslash{}lambda\^{}3 {-} 6\textbackslash{}\textbackslash{}lambda\^{}2 + 10\textbackslash{}\textbackslash{}lambda {-} 4 = 0"}
\NormalTok{  \},}
\NormalTok{  \{}
\NormalTok{    op: "Verify λ ≈ 0.491 is a root (smallest eigenvalue of this K)",}
\NormalTok{    eq: "(0.491)\^{}3 {-} 6(0.491)\^{}2 + 10(0.491) {-} 4 \textbackslash{}\textbackslash{}approx 0.118 {-} 1.446 + 4.91 {-} 4 \textbackslash{}\textbackslash{}approx {-}0.418",}
\NormalTok{    note: "Exact roots require numerical methods; the exercise builds comfort reading the polynomial structure."}
\NormalTok{  \},}
\NormalTok{  \{}
\NormalTok{    op: "Trace check",}
\NormalTok{    eq: "\textbackslash{}\textbackslash{}lambda\_1 + \textbackslash{}\textbackslash{}lambda\_2 + \textbackslash{}\textbackslash{}lambda\_3 = \textbackslash{}\textbackslash{}operatorname\{tr\}(K) = 3+2+1 = 6\textbackslash{}\textbackslash{} \textbackslash{}\textbackslash{}checkmark"}
\NormalTok{  \}}
\NormalTok{])}
\end{Highlighting}
\end{Shaded}

\textbf{5.} Determine whether the matrix
\(C = \begin{pmatrix} 3 & 1 \\ 0 & 3 \end{pmatrix}\) is diagonalisable.
If not, explain why and state what the geometric and algebraic
multiplicities are.

\begin{Shaded}
\begin{Highlighting}[]
\NormalTok{//| echo: false}
\NormalTok{makeStepperHTML(5, [}
\NormalTok{  \{}
\NormalTok{    op: "Find eigenvalues",}
\NormalTok{    eq: "\textbackslash{}\textbackslash{}det(C {-} \textbackslash{}\textbackslash{}lambda I) = (3{-}\textbackslash{}\textbackslash{}lambda)\^{}2 = 0 \textbackslash{}\textbackslash{}implies \textbackslash{}\textbackslash{}lambda = 3 \textbackslash{}\textbackslash{}text\{ (algebraic multiplicity 2)\}"}
\NormalTok{  \},}
\NormalTok{  \{}
\NormalTok{    op: "Find eigenspace: solve (C − 3I)x = 0",}
\NormalTok{    eq: "\textbackslash{}\textbackslash{}begin\{pmatrix\} 0 \& 1 \textbackslash{}\textbackslash{}\textbackslash{}\textbackslash{} 0 \& 0 \textbackslash{}\textbackslash{}end\{pmatrix\}\textbackslash{}\textbackslash{}mathbf\{x\} = \textbackslash{}\textbackslash{}mathbf\{0\} \textbackslash{}\textbackslash{}implies x\_2 = 0,\textbackslash{}\textbackslash{} x\_1 \textbackslash{}\textbackslash{}text\{ free\}"}
\NormalTok{  \},}
\NormalTok{  \{}
\NormalTok{    op: "Eigenspace",}
\NormalTok{    eq: "\textbackslash{}\textbackslash{}mathbf\{x\} = x\_1\textbackslash{}\textbackslash{}begin\{pmatrix\} 1 \textbackslash{}\textbackslash{}\textbackslash{}\textbackslash{} 0 \textbackslash{}\textbackslash{}end\{pmatrix\} \textbackslash{}\textbackslash{}implies \textbackslash{}\textbackslash{}text\{geometric multiplicity\} = 1"}
\NormalTok{  \},}
\NormalTok{  \{}
\NormalTok{    op: "Compare multiplicities",}
\NormalTok{    eq: "\textbackslash{}\textbackslash{}text\{algebraic mult.\} = 2 \textgreater{} 1 = \textbackslash{}\textbackslash{}text\{geometric mult.\}",}
\NormalTok{    note: "When these differ, the matrix is defective — not enough eigenvectors to form a basis."}
\NormalTok{  \},}
\NormalTok{  \{}
\NormalTok{    op: "Conclusion",}
\NormalTok{    eq: "C \textbackslash{}\textbackslash{}text\{ is NOT diagonalisable (it is defective).",}
\NormalTok{    note: "It has Jordan normal form J = [[3,1],[0,3]] — identical to C in this case."}
\NormalTok{  \}}
\NormalTok{])}
\end{Highlighting}
\end{Shaded}

\textbf{6.} Find the two natural frequencies of the two-mass spring
system with \(m_1 = 1\), \(m_2 = 2\), \(k_1 = k_2 = k_3 = 1\), and
describe the two mode shapes (which mass moves more in each mode).

\begin{Shaded}
\begin{Highlighting}[]
\NormalTok{//| echo: false}
\NormalTok{makeStepperHTML(6, [}
\NormalTok{  \{}
\NormalTok{    op: "Stiffness matrix K (from equations of motion)",}
\NormalTok{    eq: "K = \textbackslash{}\textbackslash{}begin\{pmatrix\} k\_1+k\_2 \& {-}k\_2 \textbackslash{}\textbackslash{}\textbackslash{}\textbackslash{} {-}k\_2 \& k\_2+k\_3 \textbackslash{}\textbackslash{}end\{pmatrix\} = \textbackslash{}\textbackslash{}begin\{pmatrix\} 2 \& {-}1 \textbackslash{}\textbackslash{}\textbackslash{}\textbackslash{} {-}1 \& 2 \textbackslash{}\textbackslash{}end\{pmatrix\}"}
\NormalTok{  \},}
\NormalTok{  \{}
\NormalTok{    op: "Mass matrix M and M⁻¹",}
\NormalTok{    eq: "M = \textbackslash{}\textbackslash{}begin\{pmatrix\} 1 \& 0 \textbackslash{}\textbackslash{}\textbackslash{}\textbackslash{} 0 \& 2 \textbackslash{}\textbackslash{}end\{pmatrix\},\textbackslash{}\textbackslash{}quad M\^{}\{{-}1\} = \textbackslash{}\textbackslash{}begin\{pmatrix\} 1 \& 0 \textbackslash{}\textbackslash{}\textbackslash{}\textbackslash{} 0 \& \textbackslash{}\textbackslash{}tfrac\{1\}\{2\} \textbackslash{}\textbackslash{}end\{pmatrix\}"}
\NormalTok{  \},}
\NormalTok{  \{}
\NormalTok{    op: "Form the system matrix A = M⁻¹K",}
\NormalTok{    eq: "A = M\^{}\{{-}1\}K = \textbackslash{}\textbackslash{}begin\{pmatrix\} 2 \& {-}1 \textbackslash{}\textbackslash{}\textbackslash{}\textbackslash{} {-}\textbackslash{}\textbackslash{}tfrac\{1\}\{2\} \& 1 \textbackslash{}\textbackslash{}end\{pmatrix\}"}
\NormalTok{  \},}
\NormalTok{  \{}
\NormalTok{    op: "Characteristic equation det(A − ω²I) = 0",}
\NormalTok{    eq: "(2{-}\textbackslash{}\textbackslash{}omega\^{}2)(1{-}\textbackslash{}\textbackslash{}omega\^{}2) {-} \textbackslash{}\textbackslash{}tfrac\{1\}\{2\} = 0 \textbackslash{}\textbackslash{}implies \textbackslash{}\textbackslash{}omega\^{}4 {-} 3\textbackslash{}\textbackslash{}omega\^{}2 + \textbackslash{}\textbackslash{}tfrac\{3\}\{2\} = 0"}
\NormalTok{  \},}
\NormalTok{  \{}
\NormalTok{    op: "Solve for ω² (quadratic in μ = ω²)",}
\NormalTok{    eq: "\textbackslash{}\textbackslash{}mu = \textbackslash{}\textbackslash{}frac\{3 \textbackslash{}\textbackslash{}pm \textbackslash{}\textbackslash{}sqrt\{9{-}6\}\}\{2\} = \textbackslash{}\textbackslash{}frac\{3 \textbackslash{}\textbackslash{}pm \textbackslash{}\textbackslash{}sqrt\{3\}\}\{2\}"}
\NormalTok{  \},}
\NormalTok{  \{}
\NormalTok{    op: "Natural frequencies",}
\NormalTok{    eq: "\textbackslash{}\textbackslash{}omega\_1 = \textbackslash{}\textbackslash{}sqrt\{\textbackslash{}\textbackslash{}tfrac\{3{-}\textbackslash{}\textbackslash{}sqrt\{3\}\}\{2\}\} \textbackslash{}\textbackslash{}approx 0.796\textbackslash{}\textbackslash{} \textbackslash{}\textbackslash{}text\{rad/s\},\textbackslash{}\textbackslash{}quad \textbackslash{}\textbackslash{}omega\_2 = \textbackslash{}\textbackslash{}sqrt\{\textbackslash{}\textbackslash{}tfrac\{3+\textbackslash{}\textbackslash{}sqrt\{3\}\}\{2\}\} \textbackslash{}\textbackslash{}approx 1.538\textbackslash{}\textbackslash{} \textbackslash{}\textbackslash{}text\{rad/s\}"}
\NormalTok{  \},}
\NormalTok{  \{}
\NormalTok{    op: "Mode 1 eigenvector (ω₁²≈0.634): solve (A − 0.634I)x = 0",}
\NormalTok{    eq: "\textbackslash{}\textbackslash{}begin\{pmatrix\} 1.366 \& {-}1 \textbackslash{}\textbackslash{}\textbackslash{}\textbackslash{} {-}0.5 \& 0.366 \textbackslash{}\textbackslash{}end\{pmatrix\}\textbackslash{}\textbackslash{}mathbf\{x\} = \textbackslash{}\textbackslash{}mathbf\{0\} \textbackslash{}\textbackslash{}implies \textbackslash{}\textbackslash{}mathbf\{x\}\_1 \textbackslash{}\textbackslash{}approx \textbackslash{}\textbackslash{}begin\{pmatrix\} 1 \textbackslash{}\textbackslash{}\textbackslash{}\textbackslash{} 1.366 \textbackslash{}\textbackslash{}end\{pmatrix\}",}
\NormalTok{    note: "Both entries positive — both masses move in the same direction. Mass 2 has larger amplitude."}
\NormalTok{  \},}
\NormalTok{  \{}
\NormalTok{    op: "Mode 2 eigenvector (ω₂²≈2.366): solve (A − 2.366I)x = 0",}
\NormalTok{    eq: "\textbackslash{}\textbackslash{}mathbf\{x\}\_2 \textbackslash{}\textbackslash{}approx \textbackslash{}\textbackslash{}begin\{pmatrix\} 1 \textbackslash{}\textbackslash{}\textbackslash{}\textbackslash{} {-}0.732 \textbackslash{}\textbackslash{}end\{pmatrix\}",}
\NormalTok{    note: "Opposite signs — masses move in opposite directions. This is the higher{-}frequency mode."}
\NormalTok{  \},}
\NormalTok{  \{}
\NormalTok{    op: "Physical interpretation",}
\NormalTok{    eq: "\textbackslash{}\textbackslash{}omega\_1 \textbackslash{}\textbackslash{}approx 0.796\textbackslash{}\textbackslash{} \textbackslash{}\textbackslash{}text\{rad/s (in{-}phase)\},\textbackslash{}\textbackslash{}quad \textbackslash{}\textbackslash{}omega\_2 \textbackslash{}\textbackslash{}approx 1.538\textbackslash{}\textbackslash{} \textbackslash{}\textbackslash{}text\{rad/s (anti{-}phase)\}",}
\NormalTok{    note: "Driving the system near either frequency causes resonance. The anti{-}phase mode has higher frequency because the relative motion stretches the coupling spring more."}
\NormalTok{  \}}
\NormalTok{])}
\end{Highlighting}
\end{Shaded}

\section{Where this goes}\label{where-this-goes-28}

Eigenvalues connect directly outward in two directions. The first is
\textbf{ODE systems}: a linear system \(\dot{\mathbf{u}} = A\mathbf{u}\)
has solution \(\mathbf{u}(t) = e^{At}\mathbf{u}(0)\), and the long-term
behaviour is governed entirely by the eigenvalues of \(A\) --- whether
solutions decay, grow, or oscillate depends on their real and imaginary
parts. The phase plane of Chapter 3 in this volume is a geometric
picture of this.

The second direction is \textbf{Fourier series}: the sinusoidal
functions \(\sin(n\pi x/L)\) and \(\cos(n\pi x/L)\) are eigenfunctions
of the second-derivative operator \(d^2/dx^2\) --- an extension of the
eigenvalue concept to function spaces. The Fourier coefficients are the
``eigenvalues'' (amplitudes) of the modes. The spectral theorem for
matrices has an analogue for symmetric differential operators
(Sturm--Liouville theory) that underpins heat equations, wave equations,
and quantum mechanics.

The computation behind all of this --- finding eigenvalues reliably for
large matrices --- is the \textbf{QR algorithm}, one of the ten most
important algorithms of the twentieth century according to
\emph{Computing in Science and Engineering}. You now understand what it
is computing.

\chapter{Vector differential
calculus}\label{vector-differential-calculus}

Grad, div, curl --- the calculus of fields

\hfill\break

Hold a weather map. The isotherms --- lines of equal temperature ---
curve across it. At any point, one direction makes the temperature climb
fastest; the gradient tells you which direction that is and how steep
the climb is. Perpendicular to that direction, the temperature doesn't
change at all.

Now picture a pipe carrying fluid. At some cross-sections the fluid
spreads out; at others it converges. Divergence measures the net outflow
at each point --- positive where fluid is being created (a source),
negative where it drains away (a sink). Curl measures something
different: whether the fluid rotates. Stick a tiny paddle wheel in the
flow; if it spins, the curl is non-zero.

Three operators --- gradient, divergence, curl --- encode the geometry
of how any field varies through space. Every major equation in physics
is written in terms of them.

\section{Vectors: dot and cross
products}\label{vectors-dot-and-cross-products}

A vector in three dimensions is a triple of numbers, written
\(\mathbf{u} = (u_1, u_2, u_3)\) or in column form. Its magnitude is
\(|\mathbf{u}| = \sqrt{u_1^2 + u_2^2 + u_3^2}\).

Two operations link a pair of vectors to useful scalars and vectors.

\subsection{The dot product}\label{the-dot-product}

\[\mathbf{u} \cdot \mathbf{v} = u_1 v_1 + u_2 v_2 + u_3 v_3\]

Geometrically:
\(\mathbf{u} \cdot \mathbf{v} = |\mathbf{u}||\mathbf{v}|\cos\theta\),
where \(\theta\) is the angle between them. This means:

\begin{itemize}
\tightlist
\item
  \(\mathbf{u} \cdot \mathbf{v} = 0\) if and only if \(\mathbf{u}\) and
  \(\mathbf{v}\) are perpendicular.
\item
  \(\text{proj}_{\mathbf{v}} \mathbf{u} = \dfrac{\mathbf{u} \cdot \mathbf{v}}{|\mathbf{v}|}\)
  is the scalar projection of \(\mathbf{u}\) onto \(\mathbf{v}\) --- how
  much of \(\mathbf{u}\) lies in the direction of \(\mathbf{v}\).
\end{itemize}

\textbf{Example.} \(\mathbf{u} = (1, 2, 3)\),
\(\mathbf{v} = (4, 0, -1)\).

\[\mathbf{u} \cdot \mathbf{v} = (1)(4) + (2)(0) + (3)(-1) = 4 + 0 - 3 = 1\]

The angle between them:
\(\cos\theta = \dfrac{1}{\sqrt{14}\,\sqrt{17}} \approx 0.065\), so
\(\theta \approx 86°\) --- nearly perpendicular.

\subsection{The cross product}\label{the-cross-product}

\[\mathbf{u} \times \mathbf{v} = \begin{vmatrix} \mathbf{i} & \mathbf{j} & \mathbf{k} \\ u_1 & u_2 & u_3 \\ v_1 & v_2 & v_3 \end{vmatrix} = (u_2 v_3 - u_3 v_2,\; u_3 v_1 - u_1 v_3,\; u_1 v_2 - u_2 v_1)\]

The result is a vector perpendicular to both \(\mathbf{u}\) and
\(\mathbf{v}\), with magnitude
\(|\mathbf{u} \times \mathbf{v}| = |\mathbf{u}||\mathbf{v}|\sin\theta\)
--- equal to the area of the parallelogram spanned by \(\mathbf{u}\) and
\(\mathbf{v}\). Direction follows the right-hand rule: curl the fingers
of your right hand from \(\mathbf{u}\) toward \(\mathbf{v}\); your thumb
points in the direction of \(\mathbf{u} \times \mathbf{v}\).

Note that
\(\mathbf{u} \times \mathbf{v} = -\mathbf{v} \times \mathbf{u}\) ---
order matters. The cross product is anti-commutative.

\textbf{Example.} \(\mathbf{u} = (1, 0, 0)\), \(\mathbf{v} = (0, 1, 0)\)
(the standard basis vectors \(\mathbf{i}\) and \(\mathbf{j}\)).

\[\mathbf{u} \times \mathbf{v} = (0 \cdot 0 - 0 \cdot 1,\; 0 \cdot 0 - 1 \cdot 0,\; 1 \cdot 1 - 0 \cdot 0) = (0, 0, 1) = \mathbf{k}\]

This is correct: the right-hand rule applied to
\(\mathbf{i} \times \mathbf{j}\) gives \(\mathbf{k}\), pointing out of
the \(xy\)-plane.

\section{Scalar and vector fields}\label{scalar-and-vector-fields}

A \textbf{scalar field} assigns a single number to each point in space:
\(f : \mathbb{R}^3 \to \mathbb{R}\). Examples:

\begin{itemize}
\tightlist
\item
  Temperature \(T(x,y,z)\) --- one number at each location
\item
  Pressure \(p(x,y,z)\) --- one number at each location
\item
  Voltage \(V(x,y,z)\) --- one number at each location
\end{itemize}

A \textbf{vector field} assigns a vector to each point in space:
\(\mathbf{F} : \mathbb{R}^3 \to \mathbb{R}^3\). Examples:

\begin{itemize}
\tightlist
\item
  Wind velocity \(\mathbf{v}(x,y,z)\) --- a direction and speed at each
  location
\item
  Gravitational field \(\mathbf{g}(x,y,z)\) --- pointing toward the
  Earth's centre, weakening with distance
\item
  Electric field \(\mathbf{E}(x,y,z)\) --- direction a positive test
  charge would accelerate
\end{itemize}

Both are functions. The difference is only in what they return. All of
single-variable calculus was about functions
\(f : \mathbb{R} \to \mathbb{R}\). Vector differential calculus extends
differentiation to these richer domains.

\section{The gradient}\label{the-gradient}

\subsection{Partial derivatives}\label{partial-derivatives}

Before defining the gradient, we need partial derivatives. Given
\(f(x, y, z)\), the partial derivative with respect to \(x\) is:

\[\frac{\partial f}{\partial x} = \lim_{h \to 0} \frac{f(x+h, y, z) - f(x, y, z)}{h}\]

Compute it by differentiating \(f\) with respect to \(x\) while treating
\(y\) and \(z\) as constants. The partial with respect to \(y\) treats
\(x\) and \(z\) as constants, and so on.

\textbf{Example.} \(f(x, y, z) = x^2 y + yz^3\).

\[\frac{\partial f}{\partial x} = 2xy, \qquad \frac{\partial f}{\partial y} = x^2 + z^3, \qquad \frac{\partial f}{\partial z} = 3yz^2\]

Each partial derivative measures the rate of change as you move in
exactly one coordinate direction while holding the others fixed.

\subsection{Definition}\label{definition}

The \textbf{gradient} of a scalar field \(f\) is the vector field formed
by collecting all three partial derivatives:

\[\nabla f = \left(\frac{\partial f}{\partial x},\; \frac{\partial f}{\partial y},\; \frac{\partial f}{\partial z}\right)\]

The symbol \(\nabla\) (nabla) is the \textbf{del operator} --- an
instruction to differentiate. Applied to a scalar, it produces a vector.

\subsection{Geometric meaning}\label{geometric-meaning}

The gradient \(\nabla f\) at any point:

\begin{enumerate}
\def\labelenumi{\arabic{enumi}.}
\tightlist
\item
  \textbf{Points in the direction of steepest increase} of \(f\).
\item
  \textbf{Is perpendicular to the level surface} \(f = c\) passing
  through that point.
\item
  Has magnitude \(|\nabla f|\) equal to the rate of change of \(f\) in
  that steepest direction.
\end{enumerate}

Think of contour lines on a topographic map. The gradient at any point
is an arrow pointing directly uphill --- perpendicular to the contour
through that point, and longer where the slope is steeper.

\textbf{Example.} \(f(x, y, z) = x^2 + y^2 + z^2\) (the squared distance
from the origin).

\[\nabla f = (2x,\; 2y,\; 2z) = 2(x, y, z) = 2\mathbf{r}\]

The gradient points radially outward --- away from the origin, in the
direction of increasing distance. This makes geometric sense: the level
surfaces of \(f\) are spheres centred at the origin, and the gradient is
perpendicular to them, i.e.~radially outward.

\begin{Shaded}
\begin{Highlighting}[]
\NormalTok{//| echo: false}
\NormalTok{gradientFieldVisualiser = \{}
\NormalTok{  const d3 = await require("d3@7");}
\NormalTok{  const W = 560, H = 500;}
\NormalTok{  const ML = 20, MR = 20, MT = 20, MB = 20;}
\NormalTok{  const innerW = W {-} ML {-} MR;}
\NormalTok{  const innerH = H {-} MT {-} MB;}
  
\NormalTok{  const container = d3.create("div").style("max{-}width", W+"px").style("font{-}family", "sans{-}serif").style("margin", "0 auto");}
  
\NormalTok{  const controls = container.append("div").style("margin{-}bottom", "15px").style("background", "\#f8fafc").style("padding", "10px").style("border{-}radius", "8px").style("border", "1px solid \#e2e8f0").style("display", "flex").style("align{-}items", "center").style("gap", "15px");}
\NormalTok{  controls.append("span").style("font{-}weight", "600").style("font{-}size", "0.95rem").style("color", "\#1e293b").text("Scalar Field f(x,y):");}
  
\NormalTok{  const select = controls.append("select").style("padding", "6px").style("border{-}radius", "4px").style("flex", 1).style("border", "1px solid \#cbd5e1").style("font{-}family", "monospace").style("font{-}size", "14px");}
\NormalTok{  const options = [}
\NormalTok{    \{label: "f(x,y) = x² + y² (Bowl)", id: "bowl"\},}
\NormalTok{    \{label: "f(x,y) = x² {-} y² (Saddle)", id: "saddle"\},}
\NormalTok{    \{label: "f(x,y) = sin(x)cos(y) (Hills \& Valleys)", id: "hills"\}}
\NormalTok{  ];}
\NormalTok{  options.forEach(o =\textgreater{} select.append("option").attr("value", o.id).text(o.label));}
  
\NormalTok{  const svg = container.append("svg").attr("width", W).attr("height", H).attr("viewBox", \textasciigrave{}0 0 $\{W\} $\{H\}\textasciigrave{});}
\NormalTok{  const g = svg.append("g").attr("transform", \textasciigrave{}translate($\{ML\},$\{MT\})\textasciigrave{});}
  
\NormalTok{  const xDomain = [{-}3, 3];}
\NormalTok{  const yDomain = [{-}3, 3];}
  
\NormalTok{  const xScale = d3.scaleLinear(xDomain, [0, innerW]);}
\NormalTok{  const yScale = d3.scaleLinear(yDomain, [innerH, 0]);}
  
\NormalTok{  function getFunc(id) \{}
\NormalTok{    if (id === "bowl") return \{ f: (x,y) =\textgreater{} x*x + y*y, dx: (x,y) =\textgreater{} 2*x, dy: (x,y) =\textgreater{} 2*y, maxV: 18, colorMax: 18 \};}
\NormalTok{    if (id === "saddle") return \{ f: (x,y) =\textgreater{} x*x {-} y*y, dx: (x,y) =\textgreater{} 2*x, dy: (x,y) =\textgreater{} {-}2*y, maxV: 6, colorMax: 9 \};}
\NormalTok{    return \{ f: (x,y) =\textgreater{} Math.sin(x)*Math.cos(y), dx: (x,y) =\textgreater{} Math.cos(x)*Math.cos(y), dy: (x,y) =\textgreater{} {-}Math.sin(x)*Math.sin(y), maxV: 1, colorMax: 1 \};}
\NormalTok{  \}}
  
\NormalTok{  const cellSize = 10;}
\NormalTok{  const nx = Math.ceil(innerW/cellSize);}
\NormalTok{  const ny = Math.ceil(innerH/cellSize);}
\NormalTok{  const rectsData = [];}
\NormalTok{  for(let i=0; i\textless{}nx; i++) \{}
\NormalTok{    for(let j=0; j\textless{}ny; j++) \{}
\NormalTok{      rectsData.push(\{i, j\});}
\NormalTok{    \}}
\NormalTok{  \}}
  
\NormalTok{  const heatmapGroup = g.append("g");}
\NormalTok{  const rects = heatmapGroup.selectAll("rect").data(rectsData).enter().append("rect")}
\NormalTok{    .attr("width", cellSize + 0.5).attr("height", cellSize + 0.5)}
\NormalTok{    .attr("x", d =\textgreater{} d.i * cellSize)}
\NormalTok{    .attr("y", d =\textgreater{} d.j * cellSize);}
  
\NormalTok{  const arrowGroup = g.append("g");}
\NormalTok{  const arrowGridSize = 35;}
\NormalTok{  const anx = Math.floor(innerW/arrowGridSize);}
\NormalTok{  const any = Math.floor(innerH/arrowGridSize);}
\NormalTok{  const arrowsData = [];}
\NormalTok{  for(let i=0; i\textless{}=anx; i++) \{}
\NormalTok{    for(let j=0; j\textless{}=any; j++) \{}
\NormalTok{      arrowsData.push(\{i, j, x: xScale.invert(i*arrowGridSize), y: yScale.invert(j*arrowGridSize)\});}
\NormalTok{    \}}
\NormalTok{  \}}
  
\NormalTok{  const arrows = arrowGroup.selectAll("g.arrow").data(arrowsData).enter().append("g").attr("class", "arrow");}
\NormalTok{  arrows.append("line").attr("stroke", "rgba(0,0,0,0.65)").attr("stroke{-}width", 1.5).attr("stroke{-}linecap", "round");}
\NormalTok{  arrows.append("polygon").attr("fill", "rgba(0,0,0,0.65)");}
  
\NormalTok{  g.append("rect").attr("width", innerW).attr("height", innerH).attr("fill", "none").attr("stroke", "\#334155").attr("stroke{-}width", 2);}
  
\NormalTok{  const colorScale = d3.scaleSequential(d3.interpolateRdBu);}
  
\NormalTok{  function update() \{}
\NormalTok{    const fn = getFunc(select.property("value"));}
    
\NormalTok{    colorScale.domain([fn.colorMax, {-}fn.colorMax]); // Red is high, Blue is low}
    
\NormalTok{    rects.attr("fill", d =\textgreater{} \{}
\NormalTok{        const cx = d.i * cellSize + cellSize/2;}
\NormalTok{        const cy = d.j * cellSize + cellSize/2;}
\NormalTok{        const x = xScale.invert(cx);}
\NormalTok{        const y = yScale.invert(cy);}
\NormalTok{        return colorScale(fn.f(x,y));}
\NormalTok{    \});}
    
\NormalTok{    arrows.each(function(d) \{}
\NormalTok{        const vx = fn.dx(d.x, d.y);}
\NormalTok{        const vy = fn.dy(d.x, d.y);}
        
\NormalTok{        const gNode = d3.select(this);}
\NormalTok{        const mag = Math.sqrt(vx*vx + vy*vy);}
\NormalTok{        if (mag \textless{} 1e{-}4) \{}
\NormalTok{           gNode.style("display", "none");}
\NormalTok{           return;}
\NormalTok{        \}}
\NormalTok{        gNode.style("display", "block");}
        
\NormalTok{        // Scale arrow length}
\NormalTok{        const baseLen = arrowGridSize * 0.45;}
\NormalTok{        const displayScale = Math.min((mag / fn.maxV) * baseLen, baseLen) + (baseLen * 0.2); }
        
\NormalTok{        const px = d.i * arrowGridSize;}
\NormalTok{        const py = d.j * arrowGridSize;}
        
\NormalTok{        const ux = vx/mag;}
\NormalTok{        const uy = vy/mag;}
        
\NormalTok{        // SVG y is down, math y is up}
\NormalTok{        const dx = ux * displayScale;}
\NormalTok{        const dy = {-}uy * displayScale;}
        
\NormalTok{        gNode.select("line")}
\NormalTok{           .attr("x1", px).attr("y1", py)}
\NormalTok{           .attr("x2", px + dx).attr("y2", py + dy);}
           
\NormalTok{        const r = 5;}
\NormalTok{        const angle = Math.atan2(dy, dx);}
\NormalTok{        const ex = px + dx;}
\NormalTok{        const ey = py + dy;}
\NormalTok{        gNode.select("polygon")}
\NormalTok{           .attr("points", \textasciigrave{}$\{ex\},$\{ey\} $\{ex {-} r*Math.cos(angle {-} 0.5)\},$\{ey {-} r*Math.sin(angle {-} 0.5)\} $\{ex {-} r*Math.cos(angle + 0.5)\},$\{ey {-} r*Math.sin(angle + 0.5)\}\textasciigrave{});}
\NormalTok{    \});}
\NormalTok{  \}}
  
\NormalTok{  select.on("change", update);}
\NormalTok{  update();}
  
\NormalTok{  return container.node();}
\NormalTok{\}}
\end{Highlighting}
\end{Shaded}

\subsection{Directional derivative}\label{directional-derivative}

The rate of change of \(f\) as you move in an arbitrary direction
\(\mathbf{u}\) (a unit vector) is called the \textbf{directional
derivative}:

\[D_{\mathbf{u}} f = \nabla f \cdot \mathbf{u}\]

This is the projection of the gradient onto the direction you're moving.
When \(\mathbf{u}\) aligns with \(\nabla f\), the directional derivative
is maximised. When \(\mathbf{u}\) is perpendicular to \(\nabla f\), the
directional derivative is zero --- you're moving along a level surface.

\textbf{Example.} \(f(x, y, z) = x^2 y + yz^3\) at point \((1, 2, -1)\).
Find the rate of change in the direction
\(\mathbf{u} = \frac{1}{\sqrt{3}}(1, 1, 1)\).

First, the gradient at \((1, 2, -1)\):

\[\nabla f\big|_{(1,2,-1)} = \bigl(2(1)(2),\; (1)^2 + (-1)^3,\; 3(2)(-1)^2\bigr) = (4, 0, 6)\]

Then:

\[D_{\mathbf{u}} f = (4, 0, 6) \cdot \frac{1}{\sqrt{3}}(1, 1, 1) = \frac{4 + 0 + 6}{\sqrt{3}} = \frac{10}{\sqrt{3}} \approx 5.77\]

\subsection{Total differential}\label{total-differential}

For small displacements \((dx, dy, dz)\) from a point, the change in
\(f\) is approximated by the \textbf{total differential}:

\[df = \frac{\partial f}{\partial x}\,dx + \frac{\partial f}{\partial y}\,dy + \frac{\partial f}{\partial z}\,dz = \nabla f \cdot d\mathbf{r}\]

This is the multi-variable chain rule in compact form. If \(x\), \(y\),
\(z\) are themselves functions of a parameter \(t\), the chain rule
becomes:

\[\frac{df}{dt} = \frac{\partial f}{\partial x}\frac{dx}{dt} + \frac{\partial f}{\partial y}\frac{dy}{dt} + \frac{\partial f}{\partial z}\frac{dz}{dt} = \nabla f \cdot \frac{d\mathbf{r}}{dt}\]

\section{Divergence}\label{divergence}

\subsection{Definition}\label{definition-1}

Given a vector field \(\mathbf{F} = (F_1, F_2, F_3)\), the
\textbf{divergence} is:

\[\nabla \cdot \mathbf{F} = \frac{\partial F_1}{\partial x} + \frac{\partial F_2}{\partial y} + \frac{\partial F_3}{\partial z}\]

The del operator \(\nabla\) is now being ``dotted'' with the vector
field --- hence the dot product notation. The result is a scalar field:
one number at each point.

\textbf{Example.} \(\mathbf{F}(x, y, z) = (x^2, 2xy, z^2)\).

\[\nabla \cdot \mathbf{F} = \frac{\partial}{\partial x}(x^2) + \frac{\partial}{\partial y}(2xy) + \frac{\partial}{\partial z}(z^2) = 2x + 2x + 2z = 4x + 2z\]

At the point \((1, 0, 1)\):
\(\nabla \cdot \mathbf{F} = 4(1) + 2(1) = 6\).

\subsection{Physical meaning}\label{physical-meaning}

Imagine \(\mathbf{F}\) is a fluid velocity field. Enclose a tiny box
around a point. The divergence measures the \textbf{net rate of fluid
flowing out} of that box per unit volume:

\begin{itemize}
\tightlist
\item
  \(\nabla \cdot \mathbf{F} > 0\): more fluid leaving than entering ---
  a \textbf{source}. (Think: water spraying from a sprinkler head.)
\item
  \(\nabla \cdot \mathbf{F} < 0\): more fluid entering than leaving ---
  a \textbf{sink}. (Think: water draining into a drain.)
\item
  \(\nabla \cdot \mathbf{F} = 0\): perfectly balanced inflow and outflow
  --- the fluid is \textbf{incompressible} at that point.
\end{itemize}

The continuity equation in fluid mechanics captures this precisely:

\[\frac{\partial \rho}{\partial t} + \nabla \cdot (\rho \mathbf{v}) = 0\]

This says: the rate of change of density at a point equals the negative
divergence of the mass flux. Mass neither appears nor disappears.

\begin{tcolorbox}[enhanced jigsaw, opacityback=0, colback=white, breakable, colframe=quarto-callout-note-color-frame, arc=.35mm, left=2mm, toprule=.15mm, leftrule=.75mm, bottomrule=.15mm, rightrule=.15mm]
\begin{minipage}[t]{5.5mm}
\textcolor{quarto-callout-note-color}{\faInfo}
\end{minipage}%
\begin{minipage}[t]{\textwidth - 5.5mm}

\textbf{Preview: Divergence theorem}

Gauss's divergence theorem connects the divergence over a volume \(V\)
to the flux through its bounding surface \(S\):

\[\iiint_V \nabla \cdot \mathbf{F}\, dV = \oiint_S \mathbf{F} \cdot d\mathbf{A}\]

The integral of all the local sources inside equals the total outflow
through the boundary. This is covered in full in the next chapter on
vector integral calculus.

\end{minipage}%
\end{tcolorbox}

\section{Curl}\label{curl}

\subsection{Definition}\label{definition-2}

Given a vector field \(\mathbf{F} = (F_1, F_2, F_3)\), the \textbf{curl}
is:

\[\nabla \times \mathbf{F} = \begin{vmatrix} \mathbf{i} & \mathbf{j} & \mathbf{k} \\ \dfrac{\partial}{\partial x} & \dfrac{\partial}{\partial y} & \dfrac{\partial}{\partial z} \\[6pt] F_1 & F_2 & F_3 \end{vmatrix}\]

Expanding the determinant:

\[\nabla \times \mathbf{F} = \left(\frac{\partial F_3}{\partial y} - \frac{\partial F_2}{\partial z},\; \frac{\partial F_1}{\partial z} - \frac{\partial F_3}{\partial x},\; \frac{\partial F_2}{\partial x} - \frac{\partial F_1}{\partial y}\right)\]

The del operator is now being ``crossed'' with the vector field. The
result is itself a vector field.

\textbf{Example.} \(\mathbf{F}(x, y, z) = (-y, x, 0)\). This is a 2D
field that rotates counterclockwise.

\[\frac{\partial F_3}{\partial y} - \frac{\partial F_2}{\partial z} = 0 - 0 = 0\]
\[\frac{\partial F_1}{\partial z} - \frac{\partial F_3}{\partial x} = 0 - 0 = 0\]
\[\frac{\partial F_2}{\partial x} - \frac{\partial F_1}{\partial y} = 1 - (-1) = 2\]

So \(\nabla \times \mathbf{F} = (0, 0, 2)\) --- a uniform curl pointing
in the \(z\)-direction. The rotation is uniform everywhere, which makes
sense: this is solid-body rotation.

\subsection{Physical meaning}\label{physical-meaning-1}

Imagine dropping a tiny paddle wheel into a fluid flow. If the paddle
wheel spins, the curl is non-zero at that point; if it doesn't spin, the
curl is zero. The direction of the curl vector is the axis of rotation
(by the right-hand rule); its magnitude measures the rotation speed.

\begin{itemize}
\tightlist
\item
  \textbf{Irrotational field:} \(\nabla \times \mathbf{F} = \mathbf{0}\)
  everywhere. The field has no local rotation --- no paddle wheel
  anywhere would spin.
\item
  \textbf{Rotational field:}
  \(\nabla \times \mathbf{F} \neq \mathbf{0}\) somewhere.
\end{itemize}

An irrotational field for which
\(\nabla \times \mathbf{F} = \mathbf{0}\) throughout a simply-connected
region can be written as the gradient of a scalar potential:
\(\mathbf{F} = \nabla \phi\). Gravity and electrostatics are
irrotational outside their sources.

\textbf{Key identity:} For any twice-differentiable scalar field \(f\),

\[\nabla \times (\nabla f) = \mathbf{0}\]

The curl of a gradient is always zero. Geometrically: gradient fields
have no rotation. This identity will appear repeatedly --- it is a test
for whether a vector field can be expressed as a gradient.

\begin{Shaded}
\begin{Highlighting}[]
\NormalTok{//| echo: false}
\NormalTok{divCurlDemo = \{}
\NormalTok{  const d3 = await require("d3@7");}
\NormalTok{  const W = 560, H = 450;}
\NormalTok{  const ML = 20, MR = 20, MT = 20, MB = 40;}
\NormalTok{  const innerW = W {-} ML {-} MR;}
\NormalTok{  const innerH = H {-} MT {-} MB;}
  
\NormalTok{  const container = d3.create("div").style("max{-}width", W+"px").style("font{-}family", "sans{-}serif").style("margin", "0 auto");}
  
\NormalTok{  const controls = container.append("div").style("margin{-}bottom", "15px").style("display", "flex").style("justify{-}content", "space{-}between").style("align{-}items", "center");}
  
\NormalTok{  const radioDiv = controls.append("div").style("display", "flex").style("gap", "15px").style("background", "\#f8fafc").style("padding", "8px 12px").style("border{-}radius", "8px").style("border", "1px solid \#e2e8f0");}
\NormalTok{  const types = [}
\NormalTok{    \{id: "source", label: "Source", dx: (x,y)=\textgreater{}x, dy: (x,y)=\textgreater{}y, div: "+2", curl: "0"\},}
\NormalTok{    \{id: "sink", label: "Sink", dx: (x,y)=\textgreater{}{-}x, dy: (x,y)=\textgreater{}{-}y, div: "{-}2", curl: "0"\},}
\NormalTok{    \{id: "rotation", label: "Rotation", dx: (x,y)=\textgreater{}{-}y, dy: (x,y)=\textgreater{}x, div: "0", curl: "+2"\},}
\NormalTok{    \{id: "shear", label: "Shear", dx: (x,y)=\textgreater{}y, dy: (x,y)=\textgreater{}0, div: "0", curl: "{-}1"\}}
\NormalTok{  ];}
  
\NormalTok{  types.forEach((t, i) =\textgreater{} \{}
\NormalTok{    const lbl = radioDiv.append("label").style("display", "flex").style("align{-}items", "center").style("gap", "5px").style("cursor", "pointer").style("font{-}size", "0.9rem");}
\NormalTok{    const inp = lbl.append("input").attr("type", "radio").attr("name", "vf\_type").attr("value", t.id).style("margin", "0");}
\NormalTok{    if(i === 0) inp.property("checked", true);}
\NormalTok{    lbl.append("span").text(t.label);}
\NormalTok{    inp.on("change", update);}
\NormalTok{  \});}
  
\NormalTok{  const readout = controls.append("div").style("display", "flex").style("gap", "20px");}
\NormalTok{  const divOut = readout.append("div").style("padding", "8px 12px").style("background", "\#fee2e2").style("border", "1px solid \#fca5a5").style("border{-}radius", "6px").style("font{-}weight", "bold").style("color", "\#991b1b").style("font{-}size", "0.95rem");}
\NormalTok{  const curlOut = readout.append("div").style("padding", "8px 12px").style("background", "\#e0e7ff").style("border", "1px solid \#a5b4fc").style("border{-}radius", "6px").style("font{-}weight", "bold").style("color", "\#3730a3").style("font{-}size", "0.95rem");}
  
\NormalTok{  const svg = container.append("svg").attr("width", W).attr("height", H).attr("viewBox", \textasciigrave{}0 0 $\{W\} $\{H\}\textasciigrave{});}
\NormalTok{  const g = svg.append("g").attr("transform", \textasciigrave{}translate($\{ML\},$\{MT\})\textasciigrave{});}
  
\NormalTok{  const xyLim = 3.5;}
\NormalTok{  const xScale = d3.scaleLinear([{-}xyLim, xyLim], [0, innerW]);}
\NormalTok{  const yScale = d3.scaleLinear([{-}xyLim, xyLim], [innerH, 0]);}
  
\NormalTok{  // Grid}
\NormalTok{  g.append("line").attr("x1", xScale({-}xyLim)).attr("y1", yScale(0)).attr("x2", xScale(xyLim)).attr("y2", yScale(0)).attr("stroke", "\#cbd5e1");}
\NormalTok{  g.append("line").attr("x1", xScale(0)).attr("y1", yScale({-}xyLim)).attr("x2", xScale(0)).attr("y2", yScale(xyLim)).attr("stroke", "\#cbd5e1");}
\NormalTok{  g.append("rect").attr("width", innerW).attr("height", innerH).attr("fill", "none").attr("stroke", "\#94a3b8");}
  
\NormalTok{  const arrowGroup = g.append("g");}
\NormalTok{  const arrowGridSize = Math.floor(innerW / 14);}
\NormalTok{  const arrowsData = [];}
\NormalTok{  for(let x=0; x\textless{}=innerW + 1; x+=arrowGridSize) \{}
\NormalTok{    for(let y=0; y\textless{}=innerH + 1; y+=arrowGridSize) \{}
\NormalTok{       arrowsData.push(\{x, y, dx: xScale.invert(x), dy: yScale.invert(y)\});}
\NormalTok{    \}}
\NormalTok{  \}}
  
\NormalTok{  const arrows = arrowGroup.selectAll("g.arrow").data(arrowsData).enter().append("g").attr("class", "arrow");}
\NormalTok{  arrows.append("line").attr("stroke", "\#3b82f6").attr("stroke{-}width", 2).attr("stroke{-}linecap", "round");}
\NormalTok{  arrows.append("polygon").attr("fill", "\#3b82f6");}
  
\NormalTok{  function update() \{}
\NormalTok{    const selectedId = radioDiv.select(\textquotesingle{}input[name="vf\_type"]:checked\textquotesingle{}).property("value");}
\NormalTok{    const fn = types.find(t =\textgreater{} t.id === selectedId);}
    
\NormalTok{    divOut.html(\textasciigrave{}\textbackslash{}u2207 \textbackslash{}u00B7 F = $\{fn.div\}\textasciigrave{});}
\NormalTok{    curlOut.html(\textasciigrave{}\textbackslash{}u2207 \textbackslash{}u00D7 F = $\{fn.curl\} k\textasciigrave{});}
    
\NormalTok{    arrows.each(function(d) \{}
\NormalTok{       const u = fn.dx(d.dx, d.dy);}
\NormalTok{       const v = fn.dy(d.dx, d.dy);}
\NormalTok{       const speed = Math.sqrt(u*u + v*v);}
       
\NormalTok{       const gNode = d3.select(this);}
\NormalTok{       if(speed \textless{} 0.1) \{}
\NormalTok{          gNode.style("display", "none");}
\NormalTok{          return;}
\NormalTok{       \}}
\NormalTok{       gNode.style("display", "block");}
       
\NormalTok{       // Max field length at edge is approx xyLim (3.5)}
\NormalTok{       const arrowLen = Math.min((speed / xyLim) * (arrowGridSize * 0.9), arrowGridSize);}
\NormalTok{       const dirX = u / speed;}
\NormalTok{       const dirY = v / speed;}
       
\NormalTok{       // SVG y is inverted}
\NormalTok{       const sx = d.x;}
\NormalTok{       const sy = d.y;}
\NormalTok{       const ex = sx + dirX * arrowLen;}
\NormalTok{       const ey = sy {-} dirY * arrowLen;}
       
\NormalTok{       gNode.select("line").attr("x1", sx).attr("y1", sy).attr("x2", ex).attr("y2", ey);}
       
\NormalTok{       const angle = Math.atan2(ey {-} sy, ex {-} sx);}
\NormalTok{       const r = 5;}
\NormalTok{       gNode.select("polygon").attr("points", \textasciigrave{}$\{ex\},$\{ey\} $\{ex {-} r*Math.cos(angle {-} 0.5)\},$\{ey {-} r*Math.sin(angle {-} 0.5)\} $\{ex {-} r*Math.cos(angle + 0.5)\},$\{ey {-} r*Math.sin(angle + 0.5)\}\textasciigrave{});}
\NormalTok{    \});}
\NormalTok{  \}}
  
\NormalTok{  update();}
\NormalTok{  return container.node();}
\NormalTok{\}}
\end{Highlighting}
\end{Shaded}

\begin{tcolorbox}[enhanced jigsaw, opacityback=0, colback=white, breakable, colframe=quarto-callout-note-color-frame, arc=.35mm, left=2mm, toprule=.15mm, leftrule=.75mm, bottomrule=.15mm, rightrule=.15mm]
\begin{minipage}[t]{5.5mm}
\textcolor{quarto-callout-note-color}{\faInfo}
\end{minipage}%
\begin{minipage}[t]{\textwidth - 5.5mm}

\textbf{Preview: Stokes' theorem}

Stokes' theorem connects the curl over a surface \(S\) to the line
integral around its boundary curve \(C\):

\[\iint_S (\nabla \times \mathbf{F}) \cdot d\mathbf{A} = \oint_C \mathbf{F} \cdot d\mathbf{r}\]

The total rotation through a surface equals the circulation around its
edge.

\end{minipage}%
\end{tcolorbox}

\section{The Laplacian}\label{the-laplacian}

\subsection{Definition}\label{definition-3}

Apply the gradient to a scalar field \(f\) to get \(\nabla f\). Then
take the divergence of that result. The outcome is the
\textbf{Laplacian}:

\[\nabla^2 f = \nabla \cdot (\nabla f) = \frac{\partial^2 f}{\partial x^2} + \frac{\partial^2 f}{\partial y^2} + \frac{\partial^2 f}{\partial z^2}\]

The Laplacian is a scalar field. It is a single operator that combines
second-order partial derivatives in all directions.

\textbf{Example.} \(f(x, y, z) = x^2 + y^2 + z^2\).

\[\frac{\partial^2 f}{\partial x^2} = 2, \qquad \frac{\partial^2 f}{\partial y^2} = 2, \qquad \frac{\partial^2 f}{\partial z^2} = 2\]

\[\nabla^2 f = 2 + 2 + 2 = 6\]

\textbf{Example.} \(f(x, y, z) = \frac{1}{r}\) where
\(r = \sqrt{x^2 + y^2 + z^2}\) (the electrostatic potential of a point
charge). A careful calculation shows
\(\nabla^2\!\left(\frac{1}{r}\right) = 0\) for \(r \neq 0\). This is
Laplace's equation, and functions satisfying it are called harmonic.

\subsection{Physical meaning}\label{physical-meaning-2}

The Laplacian \(\nabla^2 f\) at a point measures how \(f\) at that point
differs from the average value of \(f\) in a small sphere around it.
Specifically:

\begin{itemize}
\tightlist
\item
  \(\nabla^2 f > 0\): \(f\) is below the local average --- the field is
  concave up, like the bottom of a bowl.
\item
  \(\nabla^2 f < 0\): \(f\) is above the local average --- a local
  maximum relative to its surroundings.
\item
  \(\nabla^2 f = 0\): \(f\) exactly equals its local average everywhere
  --- this is Laplace's equation, and its solutions are smooth without
  local extrema in the interior.
\end{itemize}

The Laplacian appears in every major equation of mathematical physics:

\begin{longtable}[]{@{}
  >{\raggedright\arraybackslash}p{(\linewidth - 4\tabcolsep) * \real{0.3333}}
  >{\raggedright\arraybackslash}p{(\linewidth - 4\tabcolsep) * \real{0.3333}}
  >{\raggedright\arraybackslash}p{(\linewidth - 4\tabcolsep) * \real{0.3333}}@{}}
\toprule\noalign{}
\begin{minipage}[b]{\linewidth}\raggedright
Equation
\end{minipage} & \begin{minipage}[b]{\linewidth}\raggedright
Form
\end{minipage} & \begin{minipage}[b]{\linewidth}\raggedright
What it models
\end{minipage} \\
\midrule\noalign{}
\endhead
\bottomrule\noalign{}
\endlastfoot
Laplace's equation & \(\nabla^2 f = 0\) & Steady-state electrostatics,
incompressible flow \\
Poisson's equation & \(\nabla^2 f = g\) & Electrostatics with charges,
gravity with mass \\
Heat equation & \(\frac{\partial u}{\partial t} = k\,\nabla^2 u\) & Heat
diffusion through a material \\
Wave equation & \(\frac{\partial^2 u}{\partial t^2} = c^2\,\nabla^2 u\)
& Sound, light, water waves \\
\end{longtable}

\section{Applications}\label{applications-7}

\subsection{Gradient descent and
optimisation}\label{gradient-descent-and-optimisation}

In machine learning, a loss function \(L(\mathbf{w})\) measures how
wrong a model is, where \(\mathbf{w}\) is a vector of model parameters.
The gradient \(\nabla_\mathbf{w} L\) points in the direction of steepest
increase of the loss. To reduce the loss, move in the opposite
direction:

\[\mathbf{w}_{\text{new}} = \mathbf{w}_{\text{old}} - \eta\,\nabla_\mathbf{w} L\]

The scalar \(\eta\) is the \textbf{learning rate} --- how large a step
to take. This is gradient descent. In high-dimensional spaces (modern
neural networks have billions of parameters), computing the gradient
efficiently is the central computational challenge. The mathematics is
exactly what this chapter covers, lifted to millions of dimensions.

\subsection{Electric field from voltage
potential}\label{electric-field-from-voltage-potential}

In electrostatics, the voltage \(V(x, y, z)\) is a scalar field. The
electric field is its negative gradient:

\[\mathbf{E} = -\nabla V\]

The minus sign is physical convention: the electric field points from
high to low voltage, not from low to high. Positive charges accelerate
down the voltage hill.

If \(V = x^2 + y^2 + z^2\) (a simple quadratic potential), then:

\[\nabla V = (2x, 2y, 2z), \qquad \mathbf{E} = (-2x, -2y, -2z)\]

The field points radially inward --- toward the origin, toward lower
potential. Taking the divergence:

\[\nabla \cdot \mathbf{E} = -2 - 2 - 2 = -6\]

Gauss's law in differential form states
\(\nabla \cdot \mathbf{E} = \rho/\varepsilon_0\), where \(\rho\) is the
charge density. A uniform \(\nabla \cdot \mathbf{E} = -6\) means a
uniform negative charge density --- a medium with uniform charge
throughout.

\subsection{Continuity equation in fluid
mechanics}\label{continuity-equation-in-fluid-mechanics}

For a fluid with density \(\rho(x, y, z, t)\) and velocity field
\(\mathbf{v}(x, y, z, t)\), conservation of mass requires:

\[\frac{\partial \rho}{\partial t} + \nabla \cdot (\rho\,\mathbf{v}) = 0\]

This says: the rate at which mass density increases at a point equals
the negative divergence of mass flux. Mass flows in from regions where
\(\nabla \cdot (\rho\mathbf{v}) < 0\) and flows out where it is
positive. If the fluid is incompressible (\$\rho = \$ const), this
simplifies to \(\nabla \cdot \mathbf{v} = 0\).

\begin{tcolorbox}[enhanced jigsaw, opacityback=0, colback=white, breakable, colframe=quarto-callout-tip-color-frame, arc=.35mm, left=2mm, toprule=.15mm, leftrule=.75mm, bottomrule=.15mm, rightrule=.15mm]
\begin{minipage}[t]{5.5mm}
\textcolor{quarto-callout-tip-color}{\faLightbulb}
\end{minipage}%
\begin{minipage}[t]{\textwidth - 5.5mm}

\textbf{Where this shows up}

\begin{itemize}
\tightlist
\item
  An \textbf{electrical engineer} designs antenna patterns using the
  curl and divergence of Maxwell's equations:
  \(\nabla \times \mathbf{E} = -\partial\mathbf{B}/\partial t\) and
  \(\nabla \cdot \mathbf{B} = 0\).
\item
  A \textbf{meteorologist} reads wind fields as vector fields and uses
  divergence to locate areas of rising air (positive divergence near the
  surface means air is being pulled up).
\item
  A \textbf{robotics engineer} computing the configuration-space
  gradient of a potential function is doing gradient descent in the
  exact sense of this chapter.
\item
  A \textbf{GIS analyst} computing watershed flow directions is applying
  the gradient of the elevation field: water flows opposite to
  \(\nabla h\).
\end{itemize}

The operators are the same in every case. The domain changes.

\end{minipage}%
\end{tcolorbox}

\section{Exercises}\label{exercises-32}

These exercises progress from component calculation through applied
interpretation. In each, start with the definition before writing any
formula.

\begin{Shaded}
\begin{Highlighting}[]
\NormalTok{//| echo: false}
\NormalTok{katex = require("katex@0.16")}
\end{Highlighting}
\end{Shaded}

\begin{Shaded}
\begin{Highlighting}[]
\NormalTok{//| echo: false}
\NormalTok{function makeStepperHTML(exerciseNum, steps) \{}
\NormalTok{  let current = 0;}
\NormalTok{  const totalSteps = steps.length;}
\NormalTok{  const container = document.createElement("div");}
\NormalTok{  container.style.cssText = "border:1px solid \#e5e7eb; border{-}radius:8px; padding:1rem 1.25rem; margin:0.75rem 0; font{-}family:inherit;";}
\NormalTok{  const stepsDiv = document.createElement("div");}
\NormalTok{  stepsDiv.style.marginBottom = "0.75rem";}

\NormalTok{  function renderSteps() \{}
\NormalTok{    stepsDiv.innerHTML = "";}
\NormalTok{    for (let i = 0; i \textless{} current; i++) \{}
\NormalTok{      const s = steps[i];}
\NormalTok{      const row = document.createElement("div");}
\NormalTok{      row.style.cssText = "display:grid; grid{-}template{-}columns:220px 1fr; gap:0.35rem 1rem; align{-}items:baseline; padding:0.35rem 0; border{-}top:1px solid \#f3f4f6;";}
\NormalTok{      const opCell = document.createElement("span");}
\NormalTok{      opCell.style.cssText = "font{-}size:0.85em; color:\#6b7280; font{-}style:italic;";}
\NormalTok{      opCell.textContent = s.op;}
\NormalTok{      const eqCell = document.createElement("span");}
\NormalTok{      eqCell.innerHTML = katex.renderToString(s.eq, \{ throwOnError: false, displayMode: false \});}
\NormalTok{      row.appendChild(opCell);}
\NormalTok{      row.appendChild(eqCell);}
\NormalTok{      if (s.note) \{}
\NormalTok{        const noteRow = document.createElement("div");}
\NormalTok{        noteRow.style.cssText = "grid{-}column:2; font{-}size:0.82em; color:\#6b7280; padding{-}bottom:0.2rem;";}
\NormalTok{        noteRow.textContent = s.note;}
\NormalTok{        row.appendChild(noteRow);}
\NormalTok{      \}}
\NormalTok{      stepsDiv.appendChild(row);}
\NormalTok{    \}}
\NormalTok{    if (current === totalSteps) \{}
\NormalTok{      const done = document.createElement("div");}
\NormalTok{      done.style.cssText = "margin{-}top:0.5rem; font{-}size:0.9em; color:\#059669; font{-}weight:500;";}
\NormalTok{      done.textContent = "✓ Solution complete";}
\NormalTok{      stepsDiv.appendChild(done);}
\NormalTok{    \}}
\NormalTok{  \}}

\NormalTok{  const controls = document.createElement("div");}
\NormalTok{  controls.style.cssText = "display:flex; gap:0.5rem; align{-}items:center;";}
\NormalTok{  const nextBtn = document.createElement("button");}
\NormalTok{  nextBtn.textContent = "Next step →";}
\NormalTok{  nextBtn.style.cssText = "padding:0.35rem 0.85rem; border:1px solid \#d1d5db; border{-}radius:4px; background:\#fff; cursor:pointer; font{-}size:0.9em;";}
\NormalTok{  const resetBtn = document.createElement("button");}
\NormalTok{  resetBtn.textContent = "Reset";}
\NormalTok{  resetBtn.style.cssText = "padding:0.35rem 0.75rem; border:1px solid \#e5e7eb; border{-}radius:4px; background:\#f9fafb; cursor:pointer; font{-}size:0.9em; color:\#6b7280;";}
\NormalTok{  const counter = document.createElement("span");}
\NormalTok{  counter.style.cssText = "font{-}size:0.82em; color:\#9ca3af; margin{-}left:0.25rem;";}

\NormalTok{  function updateButtons() \{}
\NormalTok{    nextBtn.disabled = current === totalSteps;}
\NormalTok{    nextBtn.style.opacity = current === totalSteps ? "0.4" : "1";}
\NormalTok{    counter.textContent = current === 0 ? "Click to reveal steps" : \textasciigrave{}Step $\{current\} of $\{totalSteps\}\textasciigrave{};}
\NormalTok{  \}}

\NormalTok{  nextBtn.onclick = () =\textgreater{} \{ if (current \textless{} totalSteps) \{ current++; renderSteps(); updateButtons(); \} \};}
\NormalTok{  resetBtn.onclick = () =\textgreater{} \{ current = 0; renderSteps(); updateButtons(); \};}
\NormalTok{  controls.appendChild(nextBtn);}
\NormalTok{  controls.appendChild(resetBtn);}
\NormalTok{  controls.appendChild(counter);}
\NormalTok{  container.appendChild(stepsDiv);}
\NormalTok{  container.appendChild(controls);}
\NormalTok{  renderSteps();}
\NormalTok{  updateButtons();}
\NormalTok{  return container;}
\NormalTok{\}}
\end{Highlighting}
\end{Shaded}

\begin{enumerate}
\def\labelenumi{\arabic{enumi}.}
\tightlist
\item
  A temperature field on a heated plate is modelled by
  \(f(x, y, z) = x^2 y + yz^3\). Find \(\nabla f\) and evaluate it at
  the point \((1, 2, -1)\). Which coordinate direction gives the
  greatest rate of temperature increase at that point?
\end{enumerate}

\begin{Shaded}
\begin{Highlighting}[]
\NormalTok{//| echo: false}
\NormalTok{makeStepperHTML(1, [}
\NormalTok{  \{ op: "Partial derivative w.r.t. x",}
\NormalTok{    eq: "\textbackslash{}\textbackslash{}frac\{\textbackslash{}\textbackslash{}partial f\}\{\textbackslash{}\textbackslash{}partial x\} = 2xy",}
\NormalTok{    note: "Treat y and z as constants; differentiate x² y with respect to x." \},}
\NormalTok{  \{ op: "Partial derivative w.r.t. y",}
\NormalTok{    eq: "\textbackslash{}\textbackslash{}frac\{\textbackslash{}\textbackslash{}partial f\}\{\textbackslash{}\textbackslash{}partial y\} = x\^{}2 + z\^{}3",}
\NormalTok{    note: "Differentiate x²y + yz³ with respect to y; both terms contribute." \},}
\NormalTok{  \{ op: "Partial derivative w.r.t. z",}
\NormalTok{    eq: "\textbackslash{}\textbackslash{}frac\{\textbackslash{}\textbackslash{}partial f\}\{\textbackslash{}\textbackslash{}partial z\} = 3yz\^{}2",}
\NormalTok{    note: "Only yz³ depends on z." \},}
\NormalTok{  \{ op: "Assemble the gradient",}
\NormalTok{    eq: "\textbackslash{}\textbackslash{}nabla f = \textbackslash{}\textbackslash{}bigl(2xy,\textbackslash{}\textbackslash{}; x\^{}2 + z\^{}3,\textbackslash{}\textbackslash{}; 3yz\^{}2\textbackslash{}\textbackslash{}bigr)" \},}
\NormalTok{  \{ op: "Evaluate at (1, 2, {-}1)",}
\NormalTok{    eq: "\textbackslash{}\textbackslash{}nabla f\textbackslash{}\textbackslash{}big|\_\{(1,2,{-}1)\} = \textbackslash{}\textbackslash{}bigl(2(1)(2),\textbackslash{}\textbackslash{}; (1)\^{}2 + ({-}1)\^{}3,\textbackslash{}\textbackslash{}; 3(2)({-}1)\^{}2\textbackslash{}\textbackslash{}bigr) = (4,\textbackslash{}\textbackslash{}; 0,\textbackslash{}\textbackslash{}; 6)",}
\NormalTok{    note: "Substitute x = 1, y = 2, z = {-}1 into each component." \},}
\NormalTok{  \{ op: "Greatest rate of increase",}
\NormalTok{    eq: "\textbackslash{}\textbackslash{}text\{The z{-}component is 6, the x{-}component is 4, the y{-}component is 0.\}",}
\NormalTok{    note: "The z{-}direction has the largest partial derivative at this point, so temperature increases fastest in the z{-}direction there." \}}
\NormalTok{])}
\end{Highlighting}
\end{Shaded}

\begin{enumerate}
\def\labelenumi{\arabic{enumi}.}
\setcounter{enumi}{1}
\tightlist
\item
  A pressure field is given by \(p(x, y, z) = 3x^2 - 2y^2 + xyz\). At
  the point \((2, 1, 3)\), find the rate of change of pressure in the
  direction \(\mathbf{u} = \frac{1}{\sqrt{2}}(1, 0, 1)\) --- the
  diagonal direction in the \(xz\)-plane.
\end{enumerate}

\begin{Shaded}
\begin{Highlighting}[]
\NormalTok{//| echo: false}
\NormalTok{makeStepperHTML(2, [}
\NormalTok{  \{ op: "Compute partial derivatives",}
\NormalTok{    eq: "\textbackslash{}\textbackslash{}frac\{\textbackslash{}\textbackslash{}partial p\}\{\textbackslash{}\textbackslash{}partial x\} = 6x + yz,\textbackslash{}\textbackslash{}quad \textbackslash{}\textbackslash{}frac\{\textbackslash{}\textbackslash{}partial p\}\{\textbackslash{}\textbackslash{}partial y\} = {-}4y + xz,\textbackslash{}\textbackslash{}quad \textbackslash{}\textbackslash{}frac\{\textbackslash{}\textbackslash{}partial p\}\{\textbackslash{}\textbackslash{}partial z\} = xy" \},}
\NormalTok{  \{ op: "Evaluate at (2, 1, 3)",}
\NormalTok{    eq: "\textbackslash{}\textbackslash{}nabla p = \textbackslash{}\textbackslash{}bigl(6(2) + (1)(3),\textbackslash{}\textbackslash{}; {-}4(1) + (2)(3),\textbackslash{}\textbackslash{}; (2)(1)\textbackslash{}\textbackslash{}bigr) = (15,\textbackslash{}\textbackslash{}; 2,\textbackslash{}\textbackslash{}; 2)",}
\NormalTok{    note: "Substitute x = 2, y = 1, z = 3." \},}
\NormalTok{  \{ op: "Confirm u is a unit vector",}
\NormalTok{    eq: "|\textbackslash{}\textbackslash{}mathbf\{u\}| = \textbackslash{}\textbackslash{}sqrt\{\textbackslash{}\textbackslash{}tfrac\{1\}\{2\} + 0 + \textbackslash{}\textbackslash{}tfrac\{1\}\{2\}\} = 1 \textbackslash{}\textbackslash{}checkmark" \},}
\NormalTok{  \{ op: "Compute directional derivative",}
\NormalTok{    eq: "D\_\{\textbackslash{}\textbackslash{}mathbf\{u\}\}\textbackslash{}\textbackslash{}, p = \textbackslash{}\textbackslash{}nabla p \textbackslash{}\textbackslash{}cdot \textbackslash{}\textbackslash{}mathbf\{u\} = (15, 2, 2) \textbackslash{}\textbackslash{}cdot \textbackslash{}\textbackslash{}tfrac\{1\}\{\textbackslash{}\textbackslash{}sqrt\{2\}\}(1, 0, 1) = \textbackslash{}\textbackslash{}frac\{15 + 0 + 2\}\{\textbackslash{}\textbackslash{}sqrt\{2\}\} = \textbackslash{}\textbackslash{}frac\{17\}\{\textbackslash{}\textbackslash{}sqrt\{2\}\} \textbackslash{}\textbackslash{}approx 12.0",}
\NormalTok{    note: "The dot product selects the component of the gradient in the direction u." \}}
\NormalTok{])}
\end{Highlighting}
\end{Shaded}

\begin{enumerate}
\def\labelenumi{\arabic{enumi}.}
\setcounter{enumi}{2}
\tightlist
\item
  A fluid velocity field is given by
  \(\mathbf{F}(x, y, z) = (x^2, 2xy, z^2)\). Compute
  \(\nabla \cdot \mathbf{F}\). Identify the regions where the fluid acts
  as a source (positive divergence) and where it acts as a sink
  (negative divergence).
\end{enumerate}

\begin{Shaded}
\begin{Highlighting}[]
\NormalTok{//| echo: false}
\NormalTok{makeStepperHTML(3, [}
\NormalTok{  \{ op: "Identify components",}
\NormalTok{    eq: "F\_1 = x\^{}2,\textbackslash{}\textbackslash{}quad F\_2 = 2xy,\textbackslash{}\textbackslash{}quad F\_3 = z\^{}2" \},}
\NormalTok{  \{ op: "Partial derivative of F₁ w.r.t. x",}
\NormalTok{    eq: "\textbackslash{}\textbackslash{}frac\{\textbackslash{}\textbackslash{}partial F\_1\}\{\textbackslash{}\textbackslash{}partial x\} = 2x" \},}
\NormalTok{  \{ op: "Partial derivative of F₂ w.r.t. y",}
\NormalTok{    eq: "\textbackslash{}\textbackslash{}frac\{\textbackslash{}\textbackslash{}partial F\_2\}\{\textbackslash{}\textbackslash{}partial y\} = 2x",}
\NormalTok{    note: "Differentiate 2xy with respect to y; x is treated as a constant." \},}
\NormalTok{  \{ op: "Partial derivative of F₃ w.r.t. z",}
\NormalTok{    eq: "\textbackslash{}\textbackslash{}frac\{\textbackslash{}\textbackslash{}partial F\_3\}\{\textbackslash{}\textbackslash{}partial z\} = 2z" \},}
\NormalTok{  \{ op: "Assemble divergence",}
\NormalTok{    eq: "\textbackslash{}\textbackslash{}nabla \textbackslash{}\textbackslash{}cdot \textbackslash{}\textbackslash{}mathbf\{F\} = 2x + 2x + 2z = 4x + 2z" \},}
\NormalTok{  \{ op: "Interpret",}
\NormalTok{    eq: "\textbackslash{}\textbackslash{}nabla \textbackslash{}\textbackslash{}cdot \textbackslash{}\textbackslash{}mathbf\{F\} \textgreater{} 0 \textbackslash{}\textbackslash{}text\{ when \} 4x + 2z \textgreater{} 0,\textbackslash{}\textbackslash{}text\{ i.e. when \} 2x + z \textgreater{} 0",}
\NormalTok{    note: "The field is a source (net outflow) in the half{-}space 2x + z \textgreater{} 0, and a sink where 2x + z \textless{} 0." \}}
\NormalTok{])}
\end{Highlighting}
\end{Shaded}

\begin{enumerate}
\def\labelenumi{\arabic{enumi}.}
\setcounter{enumi}{3}
\tightlist
\item
  Verify that \(\mathbf{F}(x, y, z) = (-y, x, 0)\) represents a purely
  rotational flow --- it should have non-zero curl and zero divergence.
  Compute both.
\end{enumerate}

\begin{Shaded}
\begin{Highlighting}[]
\NormalTok{//| echo: false}
\NormalTok{makeStepperHTML(4, [}
\NormalTok{  \{ op: "Identify components",}
\NormalTok{    eq: "F\_1 = {-}y,\textbackslash{}\textbackslash{}quad F\_2 = x,\textbackslash{}\textbackslash{}quad F\_3 = 0" \},}
\NormalTok{  \{ op: "Compute divergence",}
\NormalTok{    eq: "\textbackslash{}\textbackslash{}nabla \textbackslash{}\textbackslash{}cdot \textbackslash{}\textbackslash{}mathbf\{F\} = \textbackslash{}\textbackslash{}frac\{\textbackslash{}\textbackslash{}partial({-}y)\}\{\textbackslash{}\textbackslash{}partial x\} + \textbackslash{}\textbackslash{}frac\{\textbackslash{}\textbackslash{}partial(x)\}\{\textbackslash{}\textbackslash{}partial y\} + \textbackslash{}\textbackslash{}frac\{\textbackslash{}\textbackslash{}partial(0)\}\{\textbackslash{}\textbackslash{}partial z\} = 0 + 0 + 0 = 0",}
\NormalTok{    note: "No sources or sinks — the fluid is incompressible." \},}
\NormalTok{  \{ op: "Curl: z{-}component",}
\NormalTok{    eq: "\textbackslash{}\textbackslash{}frac\{\textbackslash{}\textbackslash{}partial F\_2\}\{\textbackslash{}\textbackslash{}partial x\} {-} \textbackslash{}\textbackslash{}frac\{\textbackslash{}\textbackslash{}partial F\_1\}\{\textbackslash{}\textbackslash{}partial y\} = 1 {-} ({-}1) = 2",}
\NormalTok{    note: "This is the dominant term for a rotation in the xy{-}plane." \},}
\NormalTok{  \{ op: "Curl: x and y components",}
\NormalTok{    eq: "\textbackslash{}\textbackslash{}frac\{\textbackslash{}\textbackslash{}partial F\_3\}\{\textbackslash{}\textbackslash{}partial y\} {-} \textbackslash{}\textbackslash{}frac\{\textbackslash{}\textbackslash{}partial F\_2\}\{\textbackslash{}\textbackslash{}partial z\} = 0,\textbackslash{}\textbackslash{}quad \textbackslash{}\textbackslash{}frac\{\textbackslash{}\textbackslash{}partial F\_1\}\{\textbackslash{}\textbackslash{}partial z\} {-} \textbackslash{}\textbackslash{}frac\{\textbackslash{}\textbackslash{}partial F\_3\}\{\textbackslash{}\textbackslash{}partial x\} = 0" \},}
\NormalTok{  \{ op: "Assemble curl",}
\NormalTok{    eq: "\textbackslash{}\textbackslash{}nabla \textbackslash{}\textbackslash{}times \textbackslash{}\textbackslash{}mathbf\{F\} = (0,\textbackslash{}\textbackslash{}; 0,\textbackslash{}\textbackslash{}; 2)",}
\NormalTok{    note: "Constant curl pointing in the z{-}direction: uniform counterclockwise rotation everywhere in the xy{-}plane, at rate 2." \}}
\NormalTok{])}
\end{Highlighting}
\end{Shaded}

\begin{enumerate}
\def\labelenumi{\arabic{enumi}.}
\setcounter{enumi}{4}
\tightlist
\item
  Let \(f(x, y, z) = x^2 y^2 z\). Verify the identity
  \(\nabla \times (\nabla f) = \mathbf{0}\) --- the curl of a gradient
  is always zero. This confirms that gradient fields are irrotational
  (conservative).
\end{enumerate}

\begin{Shaded}
\begin{Highlighting}[]
\NormalTok{//| echo: false}
\NormalTok{makeStepperHTML(5, [}
\NormalTok{  \{ op: "Compute gradient",}
\NormalTok{    eq: "\textbackslash{}\textbackslash{}nabla f = \textbackslash{}\textbackslash{}bigl(2xy\^{}2z,\textbackslash{}\textbackslash{}; 2x\^{}2yz,\textbackslash{}\textbackslash{}; x\^{}2y\^{}2\textbackslash{}\textbackslash{}bigr)" \},}
\NormalTok{  \{ op: "Curl x{-}component",}
\NormalTok{    eq: "\textbackslash{}\textbackslash{}frac\{\textbackslash{}\textbackslash{}partial(x\^{}2y\^{}2)\}\{\textbackslash{}\textbackslash{}partial y\} {-} \textbackslash{}\textbackslash{}frac\{\textbackslash{}\textbackslash{}partial(2x\^{}2yz)\}\{\textbackslash{}\textbackslash{}partial z\} = 2x\^{}2y {-} 2x\^{}2y = 0",}
\NormalTok{    note: "Each term differentiates a component of ∇f." \},}
\NormalTok{  \{ op: "Curl y{-}component",}
\NormalTok{    eq: "\textbackslash{}\textbackslash{}frac\{\textbackslash{}\textbackslash{}partial(2xy\^{}2z)\}\{\textbackslash{}\textbackslash{}partial z\} {-} \textbackslash{}\textbackslash{}frac\{\textbackslash{}\textbackslash{}partial(x\^{}2y\^{}2)\}\{\textbackslash{}\textbackslash{}partial x\} = 2xy\^{}2 {-} 2xy\^{}2 = 0" \},}
\NormalTok{  \{ op: "Curl z{-}component",}
\NormalTok{    eq: "\textbackslash{}\textbackslash{}frac\{\textbackslash{}\textbackslash{}partial(2x\^{}2yz)\}\{\textbackslash{}\textbackslash{}partial x\} {-} \textbackslash{}\textbackslash{}frac\{\textbackslash{}\textbackslash{}partial(2xy\^{}2z)\}\{\textbackslash{}\textbackslash{}partial y\} = 4xyz {-} 4xyz = 0" \},}
\NormalTok{  \{ op: "Conclusion",}
\NormalTok{    eq: "\textbackslash{}\textbackslash{}nabla \textbackslash{}\textbackslash{}times (\textbackslash{}\textbackslash{}nabla f) = (0,\textbackslash{}\textbackslash{}; 0,\textbackslash{}\textbackslash{}; 0) = \textbackslash{}\textbackslash{}mathbf\{0\} \textbackslash{}\textbackslash{}checkmark",}
\NormalTok{    note: "This holds for any twice{-}differentiable f — the cancellations come from the symmetry of mixed partial derivatives (Clairaut\textquotesingle{}s theorem)." \}}
\NormalTok{])}
\end{Highlighting}
\end{Shaded}

\begin{enumerate}
\def\labelenumi{\arabic{enumi}.}
\setcounter{enumi}{5}
\tightlist
\item
  The electric potential in free space is
  \(V(x, y, z) = x^2 + y^2 + z^2\). Find the electric field
  \(\mathbf{E} = -\nabla V\), then compute \(\nabla \cdot \mathbf{E}\)
  and identify what Gauss's law says about the charge distribution.
\end{enumerate}

\begin{Shaded}
\begin{Highlighting}[]
\NormalTok{//| echo: false}
\NormalTok{makeStepperHTML(6, [}
\NormalTok{  \{ op: "Compute gradient of V",}
\NormalTok{    eq: "\textbackslash{}\textbackslash{}nabla V = \textbackslash{}\textbackslash{}left(\textbackslash{}\textbackslash{}frac\{\textbackslash{}\textbackslash{}partial V\}\{\textbackslash{}\textbackslash{}partial x\},\textbackslash{}\textbackslash{}frac\{\textbackslash{}\textbackslash{}partial V\}\{\textbackslash{}\textbackslash{}partial y\},\textbackslash{}\textbackslash{}frac\{\textbackslash{}\textbackslash{}partial V\}\{\textbackslash{}\textbackslash{}partial z\}\textbackslash{}\textbackslash{}right) = (2x,\textbackslash{}\textbackslash{}; 2y,\textbackslash{}\textbackslash{}; 2z)" \},}
\NormalTok{  \{ op: "Electric field E = {-}∇V",}
\NormalTok{    eq: "\textbackslash{}\textbackslash{}mathbf\{E\} = {-}\textbackslash{}\textbackslash{}nabla V = ({-}2x,\textbackslash{}\textbackslash{}; {-}2y,\textbackslash{}\textbackslash{}; {-}2z)",}
\NormalTok{    note: "The field points toward the origin — inward, toward lower potential." \},}
\NormalTok{  \{ op: "Compute divergence of E",}
\NormalTok{    eq: "\textbackslash{}\textbackslash{}nabla \textbackslash{}\textbackslash{}cdot \textbackslash{}\textbackslash{}mathbf\{E\} = \textbackslash{}\textbackslash{}frac\{\textbackslash{}\textbackslash{}partial({-}2x)\}\{\textbackslash{}\textbackslash{}partial x\} + \textbackslash{}\textbackslash{}frac\{\textbackslash{}\textbackslash{}partial({-}2y)\}\{\textbackslash{}\textbackslash{}partial y\} + \textbackslash{}\textbackslash{}frac\{\textbackslash{}\textbackslash{}partial({-}2z)\}\{\textbackslash{}\textbackslash{}partial z\} = {-}2 {-} 2 {-} 2 = {-}6" \},}
\NormalTok{  \{ op: "Apply Gauss\textquotesingle{}s law",}
\NormalTok{    eq: "\textbackslash{}\textbackslash{}nabla \textbackslash{}\textbackslash{}cdot \textbackslash{}\textbackslash{}mathbf\{E\} = \textbackslash{}\textbackslash{}frac\{\textbackslash{}\textbackslash{}rho\}\{\textbackslash{}\textbackslash{}varepsilon\_0\} \textbackslash{}\textbackslash{}implies \textbackslash{}\textbackslash{}rho = {-}6\textbackslash{}\textbackslash{}varepsilon\_0",}
\NormalTok{    note: "Gauss\textquotesingle{}s law in differential form: the divergence of E equals charge density over permittivity." \},}
\NormalTok{  \{ op: "Interpret",}
\NormalTok{    eq: "\textbackslash{}\textbackslash{}rho = {-}6\textbackslash{}\textbackslash{}varepsilon\_0 = \textbackslash{}\textbackslash{}text\{const\}",}
\NormalTok{    note: "A uniform negative charge density fills all space. The electric field converges inward everywhere because negative charges attract positive test charges inward." \}}
\NormalTok{])}
\end{Highlighting}
\end{Shaded}

\chapter{Vector integral calculus}\label{vector-integral-calculus}

Green's, Stokes', and the Divergence theorem

\hfill\break

A turbine blade sits in a flow of hot gas. You need to know the total
heat flux through the blade's surface --- but integrating over every
point on that complicated curved surface is brutal. The Divergence
theorem says you only need to know the divergence of the heat flux field
throughout the interior volume, which is often far simpler to compute.

This is the deepest efficiency in classical calculus: boundary integrals
and interior integrals are two faces of the same coin. Green's theorem,
Stokes' theorem, and the Divergence theorem are the exchange rate.

\section{Line integrals}\label{line-integrals}

A line integral accumulates the effect of a field along a curve. There
are two kinds: integrals of scalar fields and integrals of vector
fields. The vector case is the one that matters physically --- it
measures work.

\subsection{Scalar line integrals}\label{scalar-line-integrals}

If \(f\) is a scalar field and \(C\) is a curve, the scalar line
integral is:

\[\int_C f \, ds\]

where \(ds\) is arc length. This answers questions like ``what is the
total mass of a wire with variable density \(f\)?'' or ``what is the
average temperature along a path?''

To evaluate it, parametrise the curve: let
\(\mathbf{r}(t) = (x(t), y(t), z(t))\) for \(t \in [a, b]\). Then
\(ds = |\mathbf{r}'(t)| \, dt\), so:

\[\int_C f \, ds = \int_a^b f(\mathbf{r}(t)) \, |\mathbf{r}'(t)| \, dt\]

The integral is independent of how you parametrise the curve --- only
the geometry of \(C\) matters, not the speed at which you traverse it.

\subsection{Vector line integrals
(work)}\label{vector-line-integrals-work}

If \(\mathbf{F}\) is a vector field and \(C\) is an oriented curve, the
vector line integral is:

\[\int_C \mathbf{F} \cdot d\mathbf{r}\]

This measures the \textbf{work} done by the field \(\mathbf{F}\) on a
particle moving along \(C\). The dot product
\(\mathbf{F} \cdot d\mathbf{r}\) picks out the component of
\(\mathbf{F}\) that is parallel to the direction of travel --- only that
component does work.

To evaluate, parametrise \(C\) as \(\mathbf{r}(t)\), so
\(d\mathbf{r} = \mathbf{r}'(t) \, dt\):

\[\int_C \mathbf{F} \cdot d\mathbf{r} = \int_a^b \mathbf{F}(\mathbf{r}(t)) \cdot \mathbf{r}'(t) \, dt\]

Writing out in components with \(\mathbf{F} = (P, Q, R)\) and
\(d\mathbf{r} = (dx, dy, dz)\):

\[\int_C \mathbf{F} \cdot d\mathbf{r} = \int_C P \, dx + Q \, dy + R \, dz\]

\textbf{Sign convention:} reversing the orientation of \(C\) changes the
sign of the integral. For a closed curve traversed anti-clockwise
(positive orientation), the notation is
\(\oint_C \mathbf{F} \cdot d\mathbf{r}\).

\textbf{Worked example.} Compute \(\int_C \mathbf{F} \cdot d\mathbf{r}\)
where \(\mathbf{F} = (y, -x)\) and \(C\) is the upper semicircle of the
unit circle from \((1, 0)\) to \((-1, 0)\).

Parametrise: \(\mathbf{r}(t) = (\cos t, \sin t)\), \(t \in [0, \pi]\),
so \(\mathbf{r}'(t) = (-\sin t, \cos t)\).

Along \(C\): \(\mathbf{F}(\mathbf{r}(t)) = (\sin t, -\cos t)\).

\[\int_C \mathbf{F} \cdot d\mathbf{r} = \int_0^\pi (\sin t, -\cos t) \cdot (-\sin t, \cos t) \, dt\]

\[= \int_0^\pi (-\sin^2 t - \cos^2 t) \, dt = \int_0^\pi -1 \, dt = -\pi\]

The field does \(-\pi\) units of work on the particle --- it opposes the
motion.

\begin{center}\rule{0.5\linewidth}{0.5pt}\end{center}

\begin{Shaded}
\begin{Highlighting}[]
\NormalTok{//| echo: false}
\NormalTok{lineIntegralDemo = \{}
\NormalTok{  const d3 = await require("d3@7");}
\NormalTok{  const W = 560, H = 450;}
\NormalTok{  const ML = 20, MR = 20, MT = 20, MB = 60;}
\NormalTok{  const innerW = W {-} ML {-} MR;}
\NormalTok{  const innerH = H {-} MT {-} MB;}
  
\NormalTok{  const container = d3.create("div").style("max{-}width", W+"px").style("font{-}family", "sans{-}serif").style("margin", "0 auto");}
  
\NormalTok{  const controls = container.append("div").style("margin{-}bottom", "15px").style("padding", "10px").style("background", "\#f8fafc").style("border{-}radius", "8px").style("border", "1px solid \#e2e8f0").style("display", "flex").style("flex{-}direction", "column").style("gap", "10px");}
  
\NormalTok{  const pathDiv = controls.append("div").style("display", "flex").style("gap", "15px").style("align{-}items", "center");}
\NormalTok{  pathDiv.append("span").style("font{-}weight", "600").style("font{-}size", "0.9rem").text("Path C:");}
\NormalTok{  const paths = [}
\NormalTok{    \{id: "line", label: "Straight Line", fn: t =\textgreater{} [{-}2 + 4*t, 0], d: t =\textgreater{} [4, 0]\},}
\NormalTok{    \{id: "semi", label: "Semicircle", fn: t =\textgreater{} [{-}2*Math.cos(Math.PI*t), 2*Math.sin(Math.PI*t)], d: t =\textgreater{} [2*Math.PI*Math.sin(Math.PI*t), 2*Math.PI*Math.cos(Math.PI*t)]\},}
\NormalTok{    \{id: "parabola", label: "Parabola", fn: t =\textgreater{} [{-}2 + 4*t, 2 {-} 2*(2*t {-} 1)*(2*t {-} 1)], d: t =\textgreater{} [4, {-}8*(2*t {-} 1)]\}}
\NormalTok{  ];}
  
\NormalTok{  paths.forEach((p, i) =\textgreater{} \{}
\NormalTok{    const lbl = pathDiv.append("label").style("display", "flex").style("align{-}items", "center").style("gap", "4px").style("font{-}size", "0.9rem").style("cursor", "pointer");}
\NormalTok{    const inp = lbl.append("input").attr("type", "radio").attr("name", "path\_type").attr("value", p.id).style("margin", "0");}
\NormalTok{    if(i===0) inp.property("checked", true);}
\NormalTok{    lbl.append("span").text(p.label);}
\NormalTok{    inp.on("change", () =\textgreater{} \{ tSlider.property("value", 0); tValSpan.text("0.00"); update(); \});}
\NormalTok{  \});}
  
\NormalTok{  const tDiv = controls.append("div").style("display", "flex").style("align{-}items", "center").style("gap", "10px").style("font{-}size", "0.9rem");}
\NormalTok{  tDiv.append("span").text("t = ");}
\NormalTok{  const tSlider = tDiv.append("input").attr("type", "range").attr("min", 0).attr("max", 1).attr("step", 0.01).attr("value", 0).style("flex", 1);}
\NormalTok{  const tValSpan = tDiv.append("span").style("width", "35px").style("font{-}family", "monospace").text("0.00");}
  
\NormalTok{  const svg = container.append("svg").attr("width", W).attr("height", H).attr("viewBox", \textasciigrave{}0 0 $\{W\} $\{H\}\textasciigrave{});}
\NormalTok{  const g = svg.append("g").attr("transform", \textasciigrave{}translate($\{ML\},$\{MT\})\textasciigrave{});}
  
\NormalTok{  const xyLim = 2.5;}
\NormalTok{  const xScale = d3.scaleLinear([{-}xyLim, xyLim], [0, innerW]);}
\NormalTok{  const yScale = d3.scaleLinear([{-}xyLim, xyLim], [innerH, 0]);}
  
\NormalTok{  g.append("line").attr("x1", xScale({-}xyLim)).attr("y1", yScale(0)).attr("x2", xScale(xyLim)).attr("y2", yScale(0)).attr("stroke", "\#cbd5e1");}
\NormalTok{  g.append("line").attr("x1", xScale(0)).attr("y1", yScale({-}xyLim)).attr("x2", xScale(0)).attr("y2", yScale(xyLim)).attr("stroke", "\#cbd5e1");}
  
\NormalTok{  const F = (x, y) =\textgreater{} [{-}y, x*x];}
  
\NormalTok{  const arrowGroup = g.append("g").attr("opacity", 0.3);}
\NormalTok{  const arrowGridSize = 25;}
\NormalTok{  for(let x=0; x\textless{}=innerW; x+=arrowGridSize) \{}
\NormalTok{    for(let y=0; y\textless{}=innerH; y+=arrowGridSize) \{}
\NormalTok{       const u = xScale.invert(x);}
\NormalTok{       const v = yScale.invert(y);}
\NormalTok{       const fv = F(u, v);}
\NormalTok{       const speed = Math.sqrt(fv[0]*fv[0] + fv[1]*fv[1]);}
\NormalTok{       if(speed \textgreater{} 0.1) \{}
\NormalTok{          const len = Math.min(speed*4, arrowGridSize*0.8);}
\NormalTok{          const dx = (fv[0]/speed)*len;}
\NormalTok{          const dy = {-}(fv[1]/speed)*len;}
\NormalTok{          g.append("line").attr("x1", x).attr("y1", y).attr("x2", x+dx).attr("y2", y+dy).attr("stroke", "\#334155").attr("stroke{-}width", 1);}
\NormalTok{          g.append("circle").attr("cx", x).attr("cy", y).attr("r", 1.5).attr("fill", "\#334155");}
\NormalTok{       \}}
\NormalTok{    \}}
\NormalTok{  \}}
  
\NormalTok{  const pathCurve = g.append("path").attr("fill", "none").attr("stroke", "\#94a3b8").attr("stroke{-}width", 2.5).attr("stroke{-}dasharray", "4,4");}
\NormalTok{  const traveledCurve = g.append("path").attr("fill", "none").attr("stroke", "\#3b82f6").attr("stroke{-}width", 3);}
  
\NormalTok{  const fVec = g.append("g");}
\NormalTok{  const fVecLine = fVec.append("line").attr("stroke", "\#ef4444").attr("stroke{-}width", 2.5);}
\NormalTok{  const fVecHead = fVec.append("polygon").attr("fill", "\#ef4444");}
  
\NormalTok{  const drVec = g.append("g");}
\NormalTok{  const drVecLine = drVec.append("line").attr("stroke", "\#10b981").attr("stroke{-}width", 2.5);}
\NormalTok{  const drVecHead = drVec.append("polygon").attr("fill", "\#10b981");}
  
\NormalTok{  const particle = g.append("circle").attr("r", 6).attr("fill", "\#3b82f6");}
  
\NormalTok{  const plotG = svg.append("g").attr("transform", \textasciigrave{}translate($\{ML\},$\{H {-} MB + 20\})\textasciigrave{});}
\NormalTok{  const ptScale = d3.scaleLinear([0, 1], [0, innerW]);}
\NormalTok{  const pyScale = d3.scaleLinear([{-}15, 15], [30, 0]);}
  
\NormalTok{  plotG.append("line").attr("x1", 0).attr("x2", innerW).attr("y1", pyScale(0)).attr("y2", pyScale(0)).attr("stroke", "\#cbd5e1");}
\NormalTok{  const intArea = plotG.append("path").attr("fill", "rgba(59, 130, 246, 0.2)");}
\NormalTok{  const intLine = plotG.append("path").attr("fill", "none").attr("stroke", "\#3b82f6").attr("stroke{-}width", 2);}
\NormalTok{  const intDot = plotG.append("circle").attr("r", 4).attr("fill", "\#2563eb");}
  
\NormalTok{  const totalLabel = svg.append("text").attr("x", W {-} MR).attr("y", H {-} 5).attr("text{-}anchor", "end").style("font{-}family", "monospace").style("font{-}size", "14px").style("font{-}weight", "bold");}
  
\NormalTok{  function drawVec(vG, x, y, dx, dy) \{}
\NormalTok{    vG.select("line").attr("x1", x).attr("y1", y).attr("x2", x+dx).attr("y2", y+dy);}
\NormalTok{    const angle = Math.atan2(dy, dx);}
\NormalTok{    const r = 6;}
\NormalTok{    vG.select("polygon").attr("points", \textasciigrave{}$\{x+dx\},$\{y+dy\} $\{x+dx {-} r*Math.cos(angle {-} 0.5)\},$\{y+dy {-} r*Math.sin(angle {-} 0.5)\} $\{x+dx {-} r*Math.cos(angle + 0.5)\},$\{y+dy {-} r*Math.sin(angle + 0.5)\}\textasciigrave{});}
\NormalTok{  \}}
  
\NormalTok{  function update() \{}
\NormalTok{    const selectedId = pathDiv.select(\textquotesingle{}input[name="path\_type"]:checked\textquotesingle{}).property("value");}
\NormalTok{    const pdef = paths.find(p =\textgreater{} p.id === selectedId);}
\NormalTok{    const t = +tSlider.property("value");}
\NormalTok{    tValSpan.text(t.toFixed(2));}
    
\NormalTok{    const pts = d3.range(0, 1.01, 0.02).map(tt =\textgreater{} pdef.fn(tt));}
\NormalTok{    const lineGen = d3.line().x(d =\textgreater{} xScale(d[0])).y(d =\textgreater{} yScale(d[1]));}
\NormalTok{    pathCurve.attr("d", lineGen(pts));}
    
\NormalTok{    const ptsTraveled = d3.range(0, t+0.001, 0.02).map(tt =\textgreater{} pdef.fn(tt));}
\NormalTok{    traveledCurve.attr("d", lineGen(ptsTraveled));}
    
\NormalTok{    const pt = pdef.fn(t);}
\NormalTok{    const px = xScale(pt[0]);}
\NormalTok{    const py = yScale(pt[1]);}
\NormalTok{    particle.attr("cx", px).attr("cy", py);}
    
\NormalTok{    const fv = F(pt[0], pt[1]);}
\NormalTok{    drawVec(fVec, px, py, fv[0]*15, {-}fv[1]*15);}
    
\NormalTok{    const dr = pdef.d(t);}
\NormalTok{    drawVec(drVec, px, py, dr[0]*5, {-}dr[1]*5);}
    
\NormalTok{    const intData = d3.range(0, t+0.001, 0.02).map(tt =\textgreater{} \{}
\NormalTok{       const cpt = pdef.fn(tt);}
\NormalTok{       const cdr = pdef.d(tt);}
\NormalTok{       const cf = F(cpt[0], cpt[1]);}
\NormalTok{       const dot = cf[0]*cdr[0] + cf[1]*cdr[1];}
\NormalTok{       return \{t: tt, val: dot\};}
\NormalTok{    \});}
    
\NormalTok{    const areaGen = d3.area().x(d =\textgreater{} ptScale(d.t)).y0(pyScale(0)).y1(d =\textgreater{} pyScale(d.val));}
\NormalTok{    const lineGenPlot = d3.line().x(d =\textgreater{} ptScale(d.t)).y(d =\textgreater{} pyScale(d.val));}
    
\NormalTok{    intArea.attr("d", areaGen(intData));}
\NormalTok{    intLine.attr("d", lineGenPlot(intData));}
    
\NormalTok{    if(intData.length \textgreater{} 0) \{}
\NormalTok{      const last = intData[intData.length{-}1];}
\NormalTok{      intDot.attr("cx", ptScale(last.t)).attr("cy", pyScale(last.val)).attr("visibility", "visible");}
\NormalTok{    \} else \{}
\NormalTok{      intDot.attr("visibility", "hidden");}
\NormalTok{    \}}
    
\NormalTok{    let sum = 0;}
\NormalTok{    const dt = 0.005;}
\NormalTok{    for(let tt=0; tt\textless{}=t; tt+=dt) \{}
\NormalTok{      const cpt = pdef.fn(tt);}
\NormalTok{      const cdr = pdef.d(tt);}
\NormalTok{      const cf = F(cpt[0], cpt[1]);}
\NormalTok{      sum += (cf[0]*cdr[0] + cf[1]*cdr[1]) * dt;}
\NormalTok{    \}}
    
\NormalTok{    totalLabel.text(\textasciigrave{}\textbackslash{}u222B F \textbackslash{}u00B7 dr \textbackslash{}u2248 $\{sum.toFixed(2)\}\textasciigrave{});}
\NormalTok{  \}}
  
\NormalTok{  tSlider.on("input", update);}
\NormalTok{  update();}
  
\NormalTok{  return container.node();}
\NormalTok{\}}
\end{Highlighting}
\end{Shaded}

\begin{center}\rule{0.5\linewidth}{0.5pt}\end{center}

\section{Path independence and conservative
fields}\label{path-independence-and-conservative-fields}

For many fields, the line integral between two points is the same
regardless of which path you take. Such fields are called
\textbf{conservative}.

A vector field \(\mathbf{F}\) is conservative on a simply connected
domain if and only if any (equivalently all) of the following hold:

\begin{enumerate}
\def\labelenumi{\arabic{enumi}.}
\tightlist
\item
  \(\oint_C \mathbf{F} \cdot d\mathbf{r} = 0\) for every closed curve
  \(C\).
\item
  \(\mathbf{F} = \nabla f\) for some scalar \textbf{potential function}
  \(f\).
\item
  In 3D: \(\nabla \times \mathbf{F} = \mathbf{0}\) (the field is
  irrotational; curl is zero).
\end{enumerate}

The equivalence of (2) and (3) was established in Chapter 3. The
equivalence of (1) and (2) follows from the Fundamental Theorem for line
integrals.

\subsection{Fundamental Theorem for line
integrals}\label{fundamental-theorem-for-line-integrals}

If \(\mathbf{F} = \nabla f\), then:

\[\int_C \mathbf{F} \cdot d\mathbf{r} = f(\mathbf{r}(b)) - f(\mathbf{r}(a))\]

where \(\mathbf{r}(a)\) and \(\mathbf{r}(b)\) are the start and end
points of \(C\). The integral depends only on the endpoints --- not on
the path.

This is the vector analogue of the Fundamental Theorem of Calculus: the
integral of a derivative (here, \(\nabla f\)) over a domain (here, the
curve) equals the value of the original function at the boundary (here,
the two endpoints, with sign).

\textbf{Finding the potential function.} Given
\(\mathbf{F} = (P, Q, R)\) with
\(\nabla \times \mathbf{F} = \mathbf{0}\), find \(f\) such that
\(\nabla f = \mathbf{F}\). Integrate component by component:

\[f(x, y, z) = \int P \, dx + g(y, z)\]

then differentiate with respect to \(y\) and match against \(Q\) to
determine \(g\), and so on.

\textbf{Example.} Is \(\mathbf{F} = (2xy + z^2,\; x^2,\; 2xz)\)
conservative?

Check \(\nabla \times \mathbf{F}\): compute each component.

\((\nabla \times \mathbf{F})_x = \partial_y(2xz) - \partial_z(x^2) = 0 - 0 = 0\)

\((\nabla \times \mathbf{F})_y = \partial_z(2xy + z^2) - \partial_x(2xz) = 2z - 2z = 0\)

\((\nabla \times \mathbf{F})_z = \partial_x(x^2) - \partial_y(2xy + z^2) = 2x - 2x = 0\)

All zero, so \(\mathbf{F}\) is conservative. Find \(f\):

From \(\partial f / \partial x = 2xy + z^2\), integrate in \(x\):
\(f = x^2 y + xz^2 + g(y, z)\).

From \(\partial f / \partial y = x^2\): matches \(Q = x^2\), so
\(\partial g / \partial y = 0\).

From \(\partial f / \partial z = 2xz\): matches \(R = 2xz\), so
\(\partial g / \partial z = 0\).

Therefore \(f(x, y, z) = x^2 y + xz^2\) (plus an arbitrary constant).

\begin{tcolorbox}[enhanced jigsaw, opacityback=0, colback=white, breakable, colframe=quarto-callout-note-color-frame, arc=.35mm, left=2mm, toprule=.15mm, leftrule=.75mm, bottomrule=.15mm, rightrule=.15mm]
\begin{minipage}[t]{5.5mm}
\textcolor{quarto-callout-note-color}{\faInfo}
\end{minipage}%
\begin{minipage}[t]{\textwidth - 5.5mm}

\textbf{Why conservative fields matter in physics}

Gravity is conservative. Electric fields (in electrostatics) are
conservative. This means the work done lifting an object or moving a
charge depends only on where it starts and where it ends --- not on the
path. This is why potential energy is a well-defined concept: the number
\(f(\mathbf{r})\) at any point is the energy stored.

Friction is not conservative: a longer path dissipates more energy. The
difference is that \(\nabla \times \mathbf{F} \neq \mathbf{0}\) for
non-conservative forces.

\end{minipage}%
\end{tcolorbox}

\section{Green's theorem}\label{greens-theorem}

Green's theorem connects a line integral around a closed curve in the
plane to a double integral over the enclosed region.

\textbf{Statement.} Let \(C\) be a positively oriented (anti-clockwise),
piecewise smooth, simple closed curve in the \(xy\)-plane, and let \(D\)
be the region it encloses. If \(P(x, y)\) and \(Q(x, y)\) have
continuous partial derivatives on an open set containing \(D\), then:

\[\oint_C P \, dx + Q \, dy = \iint_D \left(\frac{\partial Q}{\partial x} - \frac{\partial P}{\partial y}\right) dA\]

\subsection{Where it comes from: cellular
decomposition}\label{where-it-comes-from-cellular-decomposition}

Divide \(D\) into tiny rectangles. For each rectangle, compute the line
integral around its boundary (traversed anti-clockwise). When you sum
all these contributions, every interior edge is traversed twice --- once
in each direction by adjacent rectangles --- so those contributions
cancel exactly. What survives is only the contribution from the outer
boundary of \(D\), which is \(C\).

For each tiny rectangle with width \(\Delta x\) and height \(\Delta y\),
the line integral around its boundary approximates:

\[\left(\frac{\partial Q}{\partial x} - \frac{\partial P}{\partial y}\right) \Delta x \, \Delta y\]

Summing over all rectangles and taking the limit gives the double
integral on the right of Green's theorem.

\begin{center}\rule{0.5\linewidth}{0.5pt}\end{center}

\begin{Shaded}
\begin{Highlighting}[]
\NormalTok{//| echo: false}
\NormalTok{greensTheoremDemo = \{}
\NormalTok{  const d3 = await require("d3@7");}
\NormalTok{  const W = 560, H = 450;}
\NormalTok{  const ML = 30, MR = 30, MT = 30, MB = 60;}
\NormalTok{  const innerW = W {-} ML {-} MR;}
\NormalTok{  const innerH = H {-} MT {-} MB;}
  
\NormalTok{  const container = d3.create("div").style("max{-}width", W+"px").style("font{-}family", "sans{-}serif").style("margin", "0 auto");}
  
\NormalTok{  const controls = container.append("div").style("margin{-}bottom", "15px").style("padding", "10px").style("background", "\#f8fafc").style("border{-}radius", "8px").style("border", "1px solid \#e2e8f0").style("display", "flex").style("align{-}items", "center").style("gap", "15px");}
  
\NormalTok{  controls.append("span").style("font{-}weight", "600").style("font{-}size", "0.9rem").text("Decomposition (NxN):");}
  
\NormalTok{  const slider = controls.append("input").attr("type", "range").attr("min", 0).attr("max", 3).attr("step", 1).attr("value", 0).style("width", "100px");}
\NormalTok{  const valSpan = controls.append("span").style("font{-}family", "monospace").text("1x1").style("width", "35px");}
  
\NormalTok{  const animBtn = controls.append("button").text("Animate Cancellation").style("padding", "5px 12px").style("cursor", "pointer").style("background", "\#10b981").style("color", "white").style("border", "none").style("border{-}radius", "4px").style("font{-}weight", "500").style("margin{-}left", "auto");}
  
\NormalTok{  const svg = container.append("svg").attr("width", W).attr("height", H).attr("viewBox", \textasciigrave{}0 0 $\{W\} $\{H\}\textasciigrave{});}
\NormalTok{  const g = svg.append("g").attr("transform", \textasciigrave{}translate($\{ML\},$\{MT\})\textasciigrave{});}
  
\NormalTok{  const rectW = innerW * 0.8;}
\NormalTok{  const rectH = innerH * 0.8;}
\NormalTok{  const offsetX = (innerW {-} rectW)/2;}
\NormalTok{  const offsetY = (innerH {-} rectH)/2;}
  
\NormalTok{  const cellsGroup = g.append("g");}
  
\NormalTok{  const readout = container.append("div").style("display", "flex").style("justify{-}content", "space{-}around").style("margin{-}top", "10px").style("font{-}family", "monospace").style("font{-}size", "0.95rem").style("font{-}weight", "bold").style("background", "\#f1f5f9").style("padding", "10px").style("border{-}radius", "6px").style("color", "\#1e293b");}
\NormalTok{  const leftSide = readout.append("div").text("\textbackslash{}u222E F \textbackslash{}u00B7 dr = 12.00");}
\NormalTok{  const rightSide = readout.append("div").text("\textbackslash{}u222C (\textbackslash{}u2207 \textbackslash{}u00D7 F) dA = 12.00");}
  
\NormalTok{  let animationRef = null;}
\NormalTok{  const totalCurl = 12.00; // Simulated exact answer}
  
\NormalTok{  function drawCells(N) \{}
\NormalTok{     cellsGroup.selectAll("*").remove();}
     
\NormalTok{     const cw = rectW / N;}
\NormalTok{     const ch = rectH / N;}
     
\NormalTok{     leftSide.text(\textasciigrave{}\textbackslash{}u222E C F \textbackslash{}u00B7 dr \textbackslash{}u2248 $\{totalCurl.toFixed(2)\}\textasciigrave{});}
\NormalTok{     rightSide.text(\textasciigrave{}\textbackslash{}u222C \textbackslash{}u2207 \textbackslash{}u00D7 F dA \textbackslash{}u2248 $\{totalCurl.toFixed(2)\}\textasciigrave{});}
     
\NormalTok{     for(let i=0; i\textless{}N; i++) \{}
\NormalTok{        for(let j=0; j\textless{}N; j++) \{}
\NormalTok{           const x = offsetX + i*cw;}
\NormalTok{           const y = offsetY + j*ch;}
           
\NormalTok{           const cellG = cellsGroup.append("g").attr("class", "cell");}
\NormalTok{           cellG.append("rect").attr("x", x).attr("y", y).attr("width", cw).attr("height", ch).attr("fill", "rgba(59, 130, 246, 0.05)").attr("stroke", "\#cbd5e1").attr("stroke{-}width", 1.5);}
           
\NormalTok{           const pad = Math.min(cw, ch) * 0.15;}
\NormalTok{           const L = x + pad; const R = x + cw {-} pad;}
\NormalTok{           const T = y + pad; const B = y + ch {-} pad;}
           
\NormalTok{           const lLen = 6;}
\NormalTok{           const color = "\#3b82f6";}
           
\NormalTok{           const addArrow = (x1,y1, x2,y2, isExternal) =\textgreater{} \{}
\NormalTok{               const gA = cellG.append("g").attr("class", isExternal ? "ext{-}boundary" : "int{-}boundary").attr("opacity", 1);}
\NormalTok{               gA.append("line").attr("x1", x1).attr("y1", y1).attr("x2", x2).attr("y2", y2).attr("stroke", color).attr("stroke{-}width", 2);}
\NormalTok{               const mx = (x1+x2)/2; const my = (y1+y2)/2;}
\NormalTok{               const angle = Math.atan2(y2{-}y1, x2{-}x1);}
\NormalTok{               gA.append("polygon").attr("points", \textasciigrave{}$\{mx\},$\{my\} $\{mx{-}lLen*Math.cos(angle{-}0.5)\},$\{my{-}lLen*Math.sin(angle{-}0.5)\} $\{mx{-}lLen*Math.cos(angle+0.5)\},$\{my{-}lLen*Math.sin(angle+0.5)\}\textasciigrave{}).attr("fill", color);}
\NormalTok{           \};}
           
\NormalTok{           addArrow(L, T, R, T, j===0); // top}
\NormalTok{           addArrow(R, T, R, B, i===N{-}1); // right}
\NormalTok{           addArrow(R, B, L, B, j===N{-}1); // bottom}
\NormalTok{           addArrow(L, B, L, T, i===0); // left}
\NormalTok{        \}}
\NormalTok{     \}}
\NormalTok{  \}}
  
\NormalTok{  function update() \{}
\NormalTok{     if(animationRef) clearTimeout(animationRef);}
\NormalTok{     const val = Math.pow(2, +slider.property("value"));}
\NormalTok{     valSpan.text(\textasciigrave{}$\{val\}x$\{val\}\textasciigrave{});}
\NormalTok{     drawCells(val);}
\NormalTok{  \}}
  
\NormalTok{  animBtn.on("click", () =\textgreater{} \{}
\NormalTok{      cellsGroup.selectAll("rect").transition().duration(200).attr("fill", "transparent").attr("stroke", "transparent");}
\NormalTok{      cellsGroup.selectAll(".int{-}boundary").transition().duration(200).attr("stroke", "\#ef4444").attr("fill", "\#ef4444").selectAll("polygon").attr("fill", "\#ef4444");}
      
\NormalTok{      setTimeout(() =\textgreater{} \{}
\NormalTok{          cellsGroup.selectAll(".int{-}boundary").transition().duration(1000).attr("opacity", 0);}
\NormalTok{          cellsGroup.selectAll(".ext{-}boundary").transition().duration(500).attr("stroke", "\#10b981").attr("stroke{-}width", 3).selectAll("polygon").attr("fill", "\#10b981");}
\NormalTok{      \}, 500);}
      
\NormalTok{      // Auto{-}reset}
\NormalTok{      animationRef = setTimeout(update, 3500);}
\NormalTok{  \});}
  
\NormalTok{  slider.on("input", update);}
\NormalTok{  update();}
  
\NormalTok{  return container.node();}
\NormalTok{\}}
\end{Highlighting}
\end{Shaded}

\begin{center}\rule{0.5\linewidth}{0.5pt}\end{center}

\textbf{Using Green's theorem.} The theorem converts between two kinds
of integrals. Use it when:

\begin{itemize}
\tightlist
\item
  The line integral \(\oint_C P \, dx + Q \, dy\) is complicated but the
  integrand \(\partial Q / \partial x - \partial P / \partial y\) is
  simple.
\item
  The double integral over \(D\) is complicated but the boundary \(C\)
  is simple.
\end{itemize}

\textbf{Example.} Compute \(\oint_C (-y \, dx + x \, dy)\) where \(C\)
is the circle \(x^2 + y^2 = 4\) traversed anti-clockwise.

Here \(P = -y\), \(Q = x\), so
\(\partial Q / \partial x - \partial P / \partial y = 1 - (-1) = 2\).

By Green's theorem:

\[\oint_C (-y \, dx + x \, dy) = \iint_D 2 \, dA = 2 \cdot \text{Area}(D) = 2 \cdot \pi(2)^2 = 8\pi\]

The boundary integral reduced to computing an area. Notice that
\(\frac{1}{2} \oint_C (-y \, dx + x \, dy) = \text{Area}(D)\) --- this
is a general formula for area from a line integral, directly from
Green's theorem.

\begin{tcolorbox}[enhanced jigsaw, opacityback=0, colback=white, breakable, colframe=quarto-callout-note-color-frame, arc=.35mm, left=2mm, toprule=.15mm, leftrule=.75mm, bottomrule=.15mm, rightrule=.15mm]
\begin{minipage}[t]{5.5mm}
\textcolor{quarto-callout-note-color}{\faInfo}
\end{minipage}%
\begin{minipage}[t]{\textwidth - 5.5mm}

\textbf{Planimeters: Green's theorem in hardware}

A planimeter is a mechanical device for measuring area. You trace the
boundary of a region with the pointer; the rotating wheel integrates
\(-y \, dx + x \, dy\) mechanically as you go, reading off the area
directly. Engineers used planimeters to measure irregular cross-sections
in structural calculations before digital tools. The device is a
physical implementation of Green's theorem.

\end{minipage}%
\end{tcolorbox}

\section{Surface integrals}\label{surface-integrals}

Surface integrals extend integration from curves to surfaces in 3D. As
with line integrals, there are scalar and vector versions.

\subsection{Parametric surfaces}\label{parametric-surfaces}

A surface \(S\) can be parametrised by
\(\mathbf{r}(u, v) = (x(u,v), y(u,v), z(u,v))\) for \((u, v)\) in some
region \(D\) of the parameter plane. The surface element \(d\mathbf{S}\)
is the vector area of an infinitesimal patch:

\[d\mathbf{S} = \frac{\partial \mathbf{r}}{\partial u} \times \frac{\partial \mathbf{r}}{\partial v} \, du \, dv\]

Its direction is normal to the surface; its magnitude
\(\left|\frac{\partial \mathbf{r}}{\partial u} \times \frac{\partial \mathbf{r}}{\partial v}\right|\)
is the scalar surface area element \(dS\).

For a surface \(z = g(x, y)\) over a region \(D\) in the \(xy\)-plane:

\[dS = \sqrt{1 + \left(\frac{\partial g}{\partial x}\right)^2 + \left(\frac{\partial g}{\partial y}\right)^2} \, dA\]

\subsection{Flux integrals}\label{flux-integrals}

The \textbf{flux} of a vector field \(\mathbf{F}\) through a surface
\(S\) is:

\[\iint_S \mathbf{F} \cdot d\mathbf{S} = \iint_S \mathbf{F} \cdot \hat{\mathbf{n}} \, dS\]

where \(\hat{\mathbf{n}}\) is the outward unit normal. Physically, this
is the rate at which field lines cross the surface --- the volume of
fluid per unit time crossing \(S\) if \(\mathbf{F}\) is a velocity
field.

\textbf{Example.} Compute the upward flux of \(\mathbf{F} = (0, 0, z)\)
through the surface \(z = 1 - x^2 - y^2\) for \(z \geq 0\) (a paraboloid
cap).

On this surface, parametrise by \((x, y)\) in the disk
\(D: x^2 + y^2 \leq 1\).

The normal vector pointing upward:

\[d\mathbf{S} = \left(-\frac{\partial z}{\partial x}, -\frac{\partial z}{\partial y}, 1\right) dA = (2x, 2y, 1) \, dA\]

So
\(\mathbf{F} \cdot d\mathbf{S} = (0)(2x) + (0)(2y) + z \cdot 1 = z = 1 - x^2 - y^2\).

\[\iint_S \mathbf{F} \cdot d\mathbf{S} = \iint_D (1 - x^2 - y^2) \, dA\]

In polar coordinates (\(x = r\cos\theta\), \(y = r\sin\theta\)):

\[= \int_0^{2\pi} \int_0^1 (1 - r^2) \, r \, dr \, d\theta = 2\pi \int_0^1 (r - r^3) \, dr = 2\pi \left[\frac{1}{2} - \frac{1}{4}\right] = \frac{\pi}{2}\]

\section{The Divergence theorem}\label{the-divergence-theorem}

The Divergence theorem (Gauss's theorem) extends Green's theorem to
three dimensions, relating the flux through a closed surface to the
divergence integrated over the enclosed volume.

\textbf{Statement.} Let \(V\) be a bounded region in \(\mathbb{R}^3\)
with piecewise smooth boundary \(\partial V\) (oriented with outward
normal). If \(\mathbf{F}\) has continuous partial derivatives on an open
set containing \(V\), then:

\[\oiint_{\partial V} \mathbf{F} \cdot d\mathbf{S} = \iiint_V (\nabla \cdot \mathbf{F}) \, dV\]

\subsection{Physical interpretation}\label{physical-interpretation-1}

\(\nabla \cdot \mathbf{F}\) at a point measures the net outward flow
\textbf{per unit volume} --- how much the field is ``spreading out''
there. The Divergence theorem says: sum up all that local spreading (the
volume integral), and you get the total outward flow through the
boundary (the surface integral). There is no other way out. Any flow
created in the interior must eventually cross the boundary.

This is the mathematical statement of conservation laws. In fluid
dynamics: if a fluid has no sources or sinks (incompressible,
\(\nabla \cdot \mathbf{F} = 0\)), then the net flux through any closed
surface is zero --- what flows in must flow out.

\subsection{Cellular decomposition in
3D}\label{cellular-decomposition-in-3d}

The argument is identical to Green's theorem but in 3D. Subdivide \(V\)
into tiny cubes. For each cube, compute the outward flux through its six
faces. On each interior face, two adjacent cubes share that face --- one
uses it as an outward face, the other as an inward face. Their
contributions cancel. What remains is the flux through the outer faces
of \(V\) only.

For a cube of volume \(\Delta V\), the net outward flux through its
faces approximates \((\nabla \cdot \mathbf{F}) \, \Delta V\) (this is
precisely the definition of divergence). Summing over all cubes gives
the Divergence theorem.

\begin{center}\rule{0.5\linewidth}{0.5pt}\end{center}

\begin{Shaded}
\begin{Highlighting}[]
\NormalTok{//| echo: false}
\NormalTok{fluxBoxDemo = \{}
\NormalTok{  const d3 = await require("d3@7");}
\NormalTok{  const W = 560, H = 380;}
  
\NormalTok{  const container = d3.create("div").style("max{-}width", W+"px").style("font{-}family", "sans{-}serif").style("margin", "0 auto");}
  
\NormalTok{  const controls = container.append("div").style("margin{-}bottom", "15px").style("padding", "10px").style("background", "\#f8fafc").style("border{-}radius", "8px").style("border", "1px solid \#e2e8f0").style("display", "grid").style("grid{-}template{-}columns", "1fr 1fr 1fr").style("gap", "15px");}
  
\NormalTok{  function makeSlider(parent, label, initVal) \{}
\NormalTok{    const wrap = parent.append("label").style("display", "flex").style("flex{-}direction", "column").style("gap", "4px").style("font{-}size", "0.9rem").style("font{-}family", "monospace");}
\NormalTok{    const top = wrap.append("div").style("display", "flex").style("justify{-}content", "space{-}between");}
\NormalTok{    top.append("span").text(label);}
\NormalTok{    const val = top.append("span").text(initVal.toFixed(1));}
\NormalTok{    const inp = wrap.append("input").attr("type", "range").attr("min", {-}5).attr("max", 5).attr("step", 0.1).attr("value", initVal);}
\NormalTok{    return \{inp, val\};}
\NormalTok{  \}}
  
\NormalTok{  const sFx = makeSlider(controls, "F\textbackslash{}u2093", 2);}
\NormalTok{  const sFy = makeSlider(controls, "F\textbackslash{}u2098", {-}1);}
\NormalTok{  const sFz = makeSlider(controls, "F\textbackslash{}u2092", 3);}
  
\NormalTok{  const layout = container.append("div").style("display", "flex").style("gap", "20px").style("align{-}items", "center");}
  
\NormalTok{  const svg = layout.append("svg").attr("width", 320).attr("height", 320).attr("viewBox", \textasciigrave{}0 0 320 320\textasciigrave{});}
\NormalTok{  const g = svg.append("g").attr("transform", \textasciigrave{}translate(160, 160)\textasciigrave{});}
  
\NormalTok{  const readout = layout.append("div").style("flex", 1).style("display", "flex").style("flex{-}direction", "column").style("gap", "6px").style("font{-}family", "monospace").style("font{-}size", "0.9rem");}
  
\NormalTok{  const colorScale = d3.scaleSequential(d3.interpolateRdBu).domain([20, {-}20]);}
  
\NormalTok{  const isoX = (x, y, z) =\textgreater{} (y {-} x) * 0.866;}
\NormalTok{  const isoY = (x, y, z) =\textgreater{} (x + y) * 0.5 {-} z;}
  
\NormalTok{  const L = 80;}
  
\NormalTok{  const faces = [}
\NormalTok{    \{id: "Top", n: [0, 0, 1], pts: [[{-}L,{-}L,L], [L,{-}L,L], [L,L,L], [{-}L,L,L]]\},}
\NormalTok{    \{id: "Front", n: [1, 0, 0], pts: [[L,{-}L,{-}L], [L,{-}L,L], [L,L,L], [L,L,{-}L]]\},}
\NormalTok{    \{id: "Right", n: [0, 1, 0], pts: [[{-}L,L,{-}L], [L,L,{-}L], [L,L,L], [{-}L,L,L]]\},}
\NormalTok{    \{id: "Bottom", n: [0, 0, {-}1], pts: [[{-}L,{-}L,{-}L], [L,{-}L,{-}L], [L,L,{-}L], [{-}L,L,{-}L]], hidden: true\},}
\NormalTok{    \{id: "Back", n: [{-}1, 0, 0], pts: [[{-}L,{-}L,{-}L], [{-}L,{-}L,L], [{-}L,L,L], [{-}L,L,{-}L]], hidden: true\},}
\NormalTok{    \{id: "Left", n: [0, {-}1, 0], pts: [[{-}L,{-}L,{-}L], [L,{-}L,{-}L], [L,{-}L,L], [{-}L,{-}L,L]], hidden: true\}}
\NormalTok{  ];}
  
\NormalTok{  const visibleFaces = faces.filter(f =\textgreater{} !f.hidden);}
  
\NormalTok{  const facePaths = visibleFaces.map(f =\textgreater{} \{}
\NormalTok{     const path = g.append("path").attr("stroke", "\#334155").attr("stroke{-}width", 1.5).attr("stroke{-}linejoin", "round");}
\NormalTok{     const cx = d3.mean(f.pts, p =\textgreater{} isoX(...p));}
\NormalTok{     const cy = d3.mean(f.pts, p =\textgreater{} isoY(...p));}
\NormalTok{     const txt = g.append("text").attr("x", cx).attr("y", cy).attr("text{-}anchor", "middle").attr("alignment{-}baseline", "middle").style("font{-}weight", "bold").style("fill", "\#000").style("text{-}shadow", "0px 0px 3px white, 0px 0px 3px white");}
\NormalTok{     return \{path, txt, f\};}
\NormalTok{  \});}
  
\NormalTok{  function update() \{}
\NormalTok{    const fx = +sFx.inp.property("value");}
\NormalTok{    const fy = +sFy.inp.property("value");}
\NormalTok{    const fz = +sFz.inp.property("value");}
    
\NormalTok{    sFx.val.text(fx.toFixed(1));}
\NormalTok{    sFy.val.text(fy.toFixed(1));}
\NormalTok{    sFz.val.text(fz.toFixed(1));}
    
\NormalTok{    readout.selectAll("*").remove();}
\NormalTok{    readout.append("strong").text("Flux Contributions:").style("margin{-}bottom", "4px").style("font{-}family", "sans{-}serif");}
    
\NormalTok{    let totalFlux = 0;}
\NormalTok{    const A = 4; // Area factor}
    
\NormalTok{    faces.forEach(f =\textgreater{} \{}
\NormalTok{       const flux = (fx*f.n[0] + fy*f.n[1] + fz*f.n[2]) * A;}
\NormalTok{       totalFlux += flux;}
       
\NormalTok{       const row = readout.append("div").style("display", "flex").style("justify{-}content", "space{-}between");}
\NormalTok{       row.append("span").text(\textasciigrave{}$\{f.id\}:\textasciigrave{});}
\NormalTok{       row.append("span").text(flux \textgreater{} 0 ? \textasciigrave{}+$\{flux.toFixed(1)\}\textasciigrave{} : flux.toFixed(1)).style("color", flux \textgreater{} 0 ? "\#b91c1c" : (flux \textless{} 0 ? "\#1d4ed8" : "\#475569"));}
\NormalTok{    \});}
    
\NormalTok{    readout.append("hr").style("width", "100\%").style("border", "0.5px solid \#cbd5e1").style("margin", "6px 0");}
\NormalTok{    const totRow = readout.append("div").style("display", "flex").style("justify{-}content", "space{-}between").style("font{-}weight", "bold");}
\NormalTok{    totRow.append("span").text(\textasciigrave{}\textbackslash{}u222F F \textbackslash{}u00B7 dA:\textasciigrave{});}
\NormalTok{    totRow.append("span").text(totalFlux.toFixed(1));}
    
\NormalTok{    const divRow = readout.append("div").style("display", "flex").style("justify{-}content", "space{-}between").style("font{-}weight", "bold");}
\NormalTok{    divRow.append("span").text(\textasciigrave{}(\textbackslash{}u2207 \textbackslash{}u00B7 F) \textbackslash{}u00D7 Vol:\textasciigrave{});}
\NormalTok{    divRow.append("span").text("0.0");}
    
\NormalTok{    const lineGen = d3.line().x(p =\textgreater{} isoX(...p)).y(p =\textgreater{} isoY(...p)).curve(d3.curveLinearClosed);}
    
\NormalTok{    facePaths.forEach(fp =\textgreater{} \{}
\NormalTok{       const flux = (fx*fp.f.n[0] + fy*fp.f.n[1] + fz*fp.f.n[2]) * A;}
\NormalTok{       fp.path.attr("d", lineGen(fp.f.pts)).attr("fill", colorScale(flux));}
\NormalTok{       fp.txt.text(flux \textgreater{} 0 ? \textasciigrave{}+$\{flux.toFixed(1)\}\textasciigrave{} : flux.toFixed(1));}
\NormalTok{    \});}
\NormalTok{  \}}
  
\NormalTok{  sFx.inp.on("input", update);}
\NormalTok{  sFy.inp.on("input", update);}
\NormalTok{  sFz.inp.on("input", update);}
  
\NormalTok{  update();}
  
\NormalTok{  return container.node();}
\NormalTok{\}}
\end{Highlighting}
\end{Shaded}

\begin{center}\rule{0.5\linewidth}{0.5pt}\end{center}

\textbf{Example.} Compute the flux of \(\mathbf{F} = (x, y, z)\) outward
through the sphere \(x^2 + y^2 + z^2 = a^2\).

Direct computation is possible but requires parametrising the sphere.
The Divergence theorem is far faster:

\(\nabla \cdot \mathbf{F} = 1 + 1 + 1 = 3\).

Volume of the sphere: \(V = \frac{4}{3}\pi a^3\).

\[\oiint_{\partial V} \mathbf{F} \cdot d\mathbf{S} = \iiint_V 3 \, dV = 3 \cdot \frac{4}{3}\pi a^3 = 4\pi a^3\]

No surface integral needed --- just multiplication.

\textbf{Gauss's law (Maxwell's equations).} The electric field
\(\mathbf{E}\) due to a charge distribution \(\rho\) satisfies:

\[\nabla \cdot \mathbf{E} = \frac{\rho}{\varepsilon_0}\]

This is the differential form. Applying the Divergence theorem over a
closed surface \(\partial V\) enclosing volume \(V\):

\[\oiint_{\partial V} \mathbf{E} \cdot d\mathbf{S} = \iiint_V \frac{\rho}{\varepsilon_0} \, dV = \frac{Q_{\text{enc}}}{\varepsilon_0}\]

where \(Q_{\text{enc}}\) is the total enclosed charge. This is Gauss's
law in integral form --- the total electric flux through any closed
surface equals the enclosed charge divided by \(\varepsilon_0\). It is
the Divergence theorem with \(\mathbf{F} = \mathbf{E}\).

\section{Stokes' theorem}\label{stokes-theorem}

Stokes' theorem is the 3D analogue of Green's theorem. It relates the
line integral around a closed curve to the flux of the curl through any
surface bounded by that curve.

\textbf{Statement.} Let \(S\) be an oriented, piecewise smooth surface
with boundary \(\partial S\) --- a closed curve. Let \(\mathbf{F}\) have
continuous partial derivatives on an open set containing \(S\). Then:

\[\oint_{\partial S} \mathbf{F} \cdot d\mathbf{r} = \iint_S (\nabla \times \mathbf{F}) \cdot d\mathbf{S}\]

The orientation of \(\partial S\) is determined by the right-hand rule:
curl your right hand's fingers in the direction you traverse the
boundary, and your thumb points in the direction of the surface normal.

\subsection{Connection to Green's
theorem}\label{connection-to-greens-theorem}

Green's theorem is the special case where \(S\) is a flat region in the
\(xy\)-plane. For \(S \subset \{z = 0\}\) with normal
\(\hat{\mathbf{k}}\):

\[(\nabla \times \mathbf{F}) \cdot \hat{\mathbf{k}} = \frac{\partial Q}{\partial x} - \frac{\partial P}{\partial y}\]

Substituting into Stokes' theorem recovers Green's theorem exactly.
Stokes' theorem can be understood as ``un-flattening'' Green's theorem
--- lifting it from a 2D region to a curved surface in 3D.

\subsection{The cellular argument for
Stokes}\label{the-cellular-argument-for-stokes}

Subdivide \(S\) into small surface patches. The line integral around
each patch's boundary traverses interior edges twice in opposite
directions --- those cancel. The sum over all patches reduces to the
line integral around the outer boundary \(\partial S\). For each patch,
the boundary integral approximates
\((\nabla \times \mathbf{F}) \cdot d\mathbf{S}\) (this is the definition
of curl). Summing gives Stokes' theorem.

\textbf{Key insight:} the surface \(S\) does not need to be flat, and it
is not unique. Any two surfaces spanning the same boundary curve give
the same flux integral of \(\nabla \times \mathbf{F}\) --- provided
\(\mathbf{F}\) is smooth throughout. Engineers exploit this: choose
whichever surface makes the integral easiest.

\textbf{Example.} Verify Stokes' theorem for
\(\mathbf{F} = (-y^2, x, z^2)\) where \(S\) is the upper hemisphere
\(x^2 + y^2 + z^2 = 1\), \(z \geq 0\), and \(\partial S\) is the unit
circle in the \(xy\)-plane.

\emph{Boundary integral:} on \(\partial S\), parametrise
\(\mathbf{r}(t) = (\cos t, \sin t, 0)\):

\[\oint_{\partial S} \mathbf{F} \cdot d\mathbf{r} = \int_0^{2\pi} (-\sin^2 t, \cos t, 0) \cdot (-\sin t, \cos t, 0) \, dt\]

\[= \int_0^{2\pi} (\sin^3 t + \cos^2 t) \, dt = 0 + \pi = \pi\]

\emph{Surface integral:} compute \(\nabla \times \mathbf{F}\):

\[\nabla \times \mathbf{F} = \left(\frac{\partial z^2}{\partial y} - \frac{\partial x}{\partial z},\; \frac{\partial(-y^2)}{\partial z} - \frac{\partial z^2}{\partial x},\; \frac{\partial x}{\partial x} - \frac{\partial(-y^2)}{\partial y}\right) = (0, 0, 1 + 2y)\]

Applying Stokes' theorem is easier via the flat disk
\(D: x^2 + y^2 \leq 1\), \(z = 0\) (same boundary, much simpler
surface):

\[\iint_D (0, 0, 1 + 2y) \cdot (0, 0, 1) \, dA = \iint_D (1 + 2y) \, dA = \pi \cdot 1^2 + 2 \iint_D y \, dA = \pi + 0 = \pi\]

Both sides give \(\pi\). \(\checkmark\)

\textbf{Ampère's law in integral form.} The magnetic field
\(\mathbf{B}\) satisfies \(\nabla \times \mathbf{B} = \mu_0 \mathbf{J}\)
(Ampère's law, differential form). Stokes' theorem converts this to:

\[\oint_{\partial S} \mathbf{B} \cdot d\mathbf{r} = \iint_S (\nabla \times \mathbf{B}) \cdot d\mathbf{S} = \mu_0 \iint_S \mathbf{J} \cdot d\mathbf{S} = \mu_0 I_{\text{enc}}\]

The circulation of the magnetic field around a closed loop equals
\(\mu_0\) times the current threading through any surface bounded by
that loop. This is the statement every electrical engineering student
uses when computing fields around wires.

\section{The unified picture}\label{the-unified-picture}

Green's, the Divergence theorem, and Stokes' are all instances of one
principle, which can be stated abstractly as:

\[\int_{\partial \Omega} \omega = \int_\Omega d\omega\]

where \(\omega\) is a differential form and \(d\omega\) is its exterior
derivative. The left side integrates over the boundary; the right side
integrates the derivative over the interior. The Fundamental Theorem of
Calculus, \(\int_a^b f'(x) \, dx = f(b) - f(a)\), is also an instance
--- here the ``boundary'' of \([a, b]\) is just the two endpoints.

\begin{longtable}[]{@{}
  >{\raggedright\arraybackslash}p{(\linewidth - 6\tabcolsep) * \real{0.1837}}
  >{\raggedright\arraybackslash}p{(\linewidth - 6\tabcolsep) * \real{0.1633}}
  >{\raggedright\arraybackslash}p{(\linewidth - 6\tabcolsep) * \real{0.2041}}
  >{\raggedright\arraybackslash}p{(\linewidth - 6\tabcolsep) * \real{0.4490}}@{}}
\toprule\noalign{}
\begin{minipage}[b]{\linewidth}\raggedright
Theorem
\end{minipage} & \begin{minipage}[b]{\linewidth}\raggedright
Domain
\end{minipage} & \begin{minipage}[b]{\linewidth}\raggedright
Boundary
\end{minipage} & \begin{minipage}[b]{\linewidth}\raggedright
Integrand (interior)
\end{minipage} \\
\midrule\noalign{}
\endhead
\bottomrule\noalign{}
\endlastfoot
FTC & Interval \([a, b]\) & Two endpoints & \(f'(x)\) \\
Green's & 2D region \(D\) & Closed curve \(\partial D\) &
\(\partial Q / \partial x - \partial P / \partial y\) \\
Stokes' & Surface \(S\) & Closed curve \(\partial S\) &
\(\nabla \times \mathbf{F}\) \\
Divergence & Volume \(V\) & Closed surface \(\partial V\) &
\(\nabla \cdot \mathbf{F}\) \\
\end{longtable}

The pattern is inescapable. All of calculus, in one row of a table.

\begin{tcolorbox}[enhanced jigsaw, opacityback=0, colback=white, breakable, colframe=quarto-callout-note-color-frame, arc=.35mm, left=2mm, toprule=.15mm, leftrule=.75mm, bottomrule=.15mm, rightrule=.15mm]
\begin{minipage}[t]{5.5mm}
\textcolor{quarto-callout-note-color}{\faInfo}
\end{minipage}%
\begin{minipage}[t]{\textwidth - 5.5mm}

\textbf{Strategy: which way to go}

Given a problem, decide which direction the theorem should run:

\begin{itemize}
\tightlist
\item
  If the \textbf{boundary} integral looks hard but the \textbf{interior}
  is simple: Green's / Stokes' / Divergence converts boundary →
  interior.
\item
  If the \textbf{interior} integral looks hard but the \textbf{boundary}
  is simple: run it the other way.
\item
  If \(\nabla \cdot \mathbf{F} = 0\) (solenoidal): the flux through any
  closed surface is zero. No computation needed.
\item
  If \(\nabla \times \mathbf{F} = \mathbf{0}\) (irrotational): the line
  integral around any closed curve is zero. Use the potential function.
\end{itemize}

The theorems are not formulas to plug into --- they are a choice of
which representation of the problem is simpler.

\end{minipage}%
\end{tcolorbox}

\section{Applications}\label{applications-8}

\subsection{Electric flux and Gauss's
law}\label{electric-flux-and-gausss-law}

An isolated point charge \(q\) sits at the origin. Compute the electric
flux through a sphere of radius \(r\).

The electric field is
\(\mathbf{E} = \frac{q}{4\pi\varepsilon_0 r^2} \hat{\mathbf{r}}\). On
the sphere of radius \(r\), the field is radial with constant magnitude
\(E = q / (4\pi\varepsilon_0 r^2)\), and the outward normal
\(\hat{\mathbf{n}} = \hat{\mathbf{r}}\), so
\(\mathbf{E} \cdot \hat{\mathbf{n}} = E\) everywhere.

\[\oiint_S \mathbf{E} \cdot d\mathbf{S} = E \cdot 4\pi r^2 = \frac{q}{4\pi\varepsilon_0 r^2} \cdot 4\pi r^2 = \frac{q}{\varepsilon_0}\]

This equals the enclosed charge divided by \(\varepsilon_0\), confirming
Gauss's law. Notice the \(r\) cancelled --- the flux is the same through
any sphere centred on the charge. This is a consequence of the
Divergence theorem and the fact that \(\nabla \cdot \mathbf{E} = 0\)
everywhere except at the charge itself.

\subsection{Mass flow through a surface (fluid
mechanics)}\label{mass-flow-through-a-surface-fluid-mechanics}

A velocity field \(\mathbf{v} = (2x, -y, z)\) describes fluid flow in a
region. Is mass accumulating inside the unit cube \([0,1]^3\)?

\(\nabla \cdot \mathbf{v} = 2 - 1 + 1 = 2 > 0\).

By the Divergence theorem, the net outward mass flux is:

\[\oiint_{\partial V} \mathbf{v} \cdot d\mathbf{S} = \iiint_V 2 \, dV = 2\]

Two units of mass per unit time are flowing out of the cube. The
divergence is positive throughout, so the fluid is expanding --- it acts
like a source term. In an incompressible flow (water, to a first
approximation), this is impossible: \(\nabla \cdot \mathbf{v}\) must be
zero everywhere.

\subsection{Work done against gravity (conservative
field)}\label{work-done-against-gravity-conservative-field}

Gravity is \(\mathbf{g} = (0, 0, -g)\) with potential \(\phi = gz\) (so
\(\mathbf{g} = -\nabla\phi\)). A mass \(m\) is moved along a complicated
path from the ground level \((0, 0, 0)\) to a height \((x_1, y_1, h)\).

Since \(\mathbf{F}_{\text{grav}} = m\mathbf{g} = -m\nabla(gz)\), it is
conservative. Work done by gravity:

\[W = \int_C \mathbf{F}_{\text{grav}} \cdot d\mathbf{r} = -m g z \Big|_{z=0}^{z=h} = -mgh\]

The negative sign indicates gravity does negative work (you are lifting
against it). Work done \emph{against} gravity is \(+mgh\) --- the
familiar result from introductory physics, now derived precisely from
path independence.

\subsection{Stokes' theorem and Faraday's
law}\label{stokes-theorem-and-faradays-law}

Faraday's law of electromagnetic induction states
\(\nabla \times \mathbf{E} = -\partial \mathbf{B} / \partial t\). Apply
Stokes' theorem to a surface \(S\) bounded by a circuit loop \(C\):

\[\oint_C \mathbf{E} \cdot d\mathbf{r} = \iint_S (\nabla \times \mathbf{E}) \cdot d\mathbf{S} = -\frac{d}{dt} \iint_S \mathbf{B} \cdot d\mathbf{S} = -\frac{d\Phi_B}{dt}\]

The EMF (voltage) around the circuit equals the negative rate of change
of magnetic flux through it. This is the principle behind every
transformer and generator ever built.

\section{Exercises}\label{exercises-33}

These exercises build from evaluation to application. The first two
establish mechanical fluency; the later ones require choosing which
theorem to apply and why.

\begin{Shaded}
\begin{Highlighting}[]
\NormalTok{//| echo: false}
\NormalTok{\{}
\NormalTok{  const katex = await require("katex@0.16");}
\NormalTok{  const css = document.createElement("link");}
\NormalTok{  css.rel = "stylesheet";}
\NormalTok{  css.href = "https://cdn.jsdelivr.net/npm/katex@0.16/dist/katex.min.css";}
\NormalTok{  document.head.appendChild(css);}

\NormalTok{  function makeStepperHTML(steps) \{}
\NormalTok{    let current = 0;}
\NormalTok{    const totalSteps = steps.length;}
\NormalTok{    const container = document.createElement("div");}
\NormalTok{    container.style.cssText = "border:1px solid \#e5e7eb; border{-}radius:8px; padding:1rem 1.25rem; margin:0.75rem 0; font{-}family:inherit;";}

\NormalTok{    const stepsDiv = document.createElement("div");}
\NormalTok{    stepsDiv.style.marginBottom = "0.75rem";}

\NormalTok{    function renderSteps() \{}
\NormalTok{      stepsDiv.innerHTML = "";}
\NormalTok{      for (let i = 0; i \textless{} current; i++) \{}
\NormalTok{        const s = steps[i];}
\NormalTok{        const row = document.createElement("div");}
\NormalTok{        row.style.cssText = "display:grid; grid{-}template{-}columns:200px 1fr; gap:0.35rem 1rem; align{-}items:baseline; padding:0.35rem 0; border{-}top:1px solid \#f3f4f6;";}
\NormalTok{        const opCell = document.createElement("span");}
\NormalTok{        opCell.style.cssText = "font{-}size:0.85em; color:\#6b7280; font{-}style:italic;";}
\NormalTok{        opCell.textContent = s.op;}
\NormalTok{        const eqCell = document.createElement("span");}
\NormalTok{        eqCell.innerHTML = katex.renderToString(s.eq, \{ throwOnError: false, displayMode: false \});}
\NormalTok{        row.appendChild(opCell);}
\NormalTok{        row.appendChild(eqCell);}
\NormalTok{        if (s.note) \{}
\NormalTok{          const noteRow = document.createElement("div");}
\NormalTok{          noteRow.style.cssText = "grid{-}column:2; font{-}size:0.82em; color:\#6b7280; padding{-}bottom:0.2rem;";}
\NormalTok{          noteRow.textContent = s.note;}
\NormalTok{          row.appendChild(noteRow);}
\NormalTok{        \}}
\NormalTok{        stepsDiv.appendChild(row);}
\NormalTok{      \}}
\NormalTok{      if (current === totalSteps) \{}
\NormalTok{        const done = document.createElement("div");}
\NormalTok{        done.style.cssText = "margin{-}top:0.5rem; font{-}size:0.9em; color:\#059669; font{-}weight:500;";}
\NormalTok{        done.textContent = "✓ Solution complete";}
\NormalTok{        stepsDiv.appendChild(done);}
\NormalTok{      \}}
\NormalTok{    \}}

\NormalTok{    const controls = document.createElement("div");}
\NormalTok{    controls.style.cssText = "display:flex; gap:0.5rem; align{-}items:center;";}

\NormalTok{    const nextBtn = document.createElement("button");}
\NormalTok{    nextBtn.textContent = "Next step →";}
\NormalTok{    nextBtn.style.cssText = "padding:0.35rem 0.85rem; border:1px solid \#d1d5db; border{-}radius:4px; background:\#fff; cursor:pointer; font{-}size:0.9em;";}

\NormalTok{    const resetBtn = document.createElement("button");}
\NormalTok{    resetBtn.textContent = "Reset";}
\NormalTok{    resetBtn.style.cssText = "padding:0.35rem 0.75rem; border:1px solid \#e5e7eb; border{-}radius:4px; background:\#f9fafb; cursor:pointer; font{-}size:0.9em; color:\#6b7280;";}

\NormalTok{    const counter = document.createElement("span");}
\NormalTok{    counter.style.cssText = "font{-}size:0.82em; color:\#9ca3af; margin{-}left:0.25rem;";}

\NormalTok{    function updateButtons() \{}
\NormalTok{      nextBtn.disabled = current === totalSteps;}
\NormalTok{      nextBtn.style.opacity = current === totalSteps ? "0.4" : "1";}
\NormalTok{      counter.textContent = current === 0 ? "Click to reveal steps" : \textasciigrave{}Step $\{current\} of $\{totalSteps\}\textasciigrave{};}
\NormalTok{    \}}

\NormalTok{    nextBtn.onclick = () =\textgreater{} \{ if (current \textless{} totalSteps) \{ current++; renderSteps(); updateButtons(); \} \};}
\NormalTok{    resetBtn.onclick = () =\textgreater{} \{ current = 0; renderSteps(); updateButtons(); \};}

\NormalTok{    controls.appendChild(nextBtn);}
\NormalTok{    controls.appendChild(resetBtn);}
\NormalTok{    controls.appendChild(counter);}
\NormalTok{    container.appendChild(stepsDiv);}
\NormalTok{    container.appendChild(controls);}
\NormalTok{    renderSteps();}
\NormalTok{    updateButtons();}
\NormalTok{    return container;}
\NormalTok{  \}}

\NormalTok{  // Attach to window so subsequent blocks can call it}
\NormalTok{  window.\_makeStepperHTML = makeStepperHTML;}
\NormalTok{  return html\textasciigrave{}\textless{}span\textgreater{}\textless{}/span\textgreater{}\textasciigrave{};}
\NormalTok{\}}
\end{Highlighting}
\end{Shaded}

\textbf{1. Line integral along a parametric curve.} A particle moves
along the helix \(\mathbf{r}(t) = (\cos t, \sin t, t)\) for
\(t \in [0, 2\pi]\) in the force field \(\mathbf{F} = (-y, x, z)\).
Compute the work done \(W = \int_C \mathbf{F} \cdot d\mathbf{r}\).

\begin{Shaded}
\begin{Highlighting}[]
\NormalTok{//| echo: false}
\NormalTok{window.\_makeStepperHTML([}
\NormalTok{  \{}
\NormalTok{    op: "Parametrise and compute r\textquotesingle{}(t)",}
\NormalTok{    eq: "\textbackslash{}\textbackslash{}mathbf\{r\}(t) = (\textbackslash{}\textbackslash{}cos t,\textbackslash{}\textbackslash{}, \textbackslash{}\textbackslash{}sin t,\textbackslash{}\textbackslash{}, t), \textbackslash{}\textbackslash{}quad \textbackslash{}\textbackslash{}mathbf\{r\}\textquotesingle{}(t) = ({-}\textbackslash{}\textbackslash{}sin t,\textbackslash{}\textbackslash{}, \textbackslash{}\textbackslash{}cos t,\textbackslash{}\textbackslash{}, 1)",}
\NormalTok{    note: "The derivative gives the tangent direction at each point on the helix."}
\NormalTok{  \},}
\NormalTok{  \{}
\NormalTok{    op: "Evaluate F along the curve",}
\NormalTok{    eq: "\textbackslash{}\textbackslash{}mathbf\{F\}(\textbackslash{}\textbackslash{}mathbf\{r\}(t)) = ({-}\textbackslash{}\textbackslash{}sin t,\textbackslash{}\textbackslash{}, \textbackslash{}\textbackslash{}cos t,\textbackslash{}\textbackslash{}, t)",}
\NormalTok{    note: "Substitute x = cos t, y = sin t, z = t into F = ({-}y, x, z)."}
\NormalTok{  \},}
\NormalTok{  \{}
\NormalTok{    op: "Form the dot product",}
\NormalTok{    eq: "\textbackslash{}\textbackslash{}mathbf\{F\} \textbackslash{}\textbackslash{}cdot \textbackslash{}\textbackslash{}mathbf\{r\}\textquotesingle{}(t) = ({-}\textbackslash{}\textbackslash{}sin t)({-}\textbackslash{}\textbackslash{}sin t) + (\textbackslash{}\textbackslash{}cos t)(\textbackslash{}\textbackslash{}cos t) + t \textbackslash{}\textbackslash{}cdot 1 = \textbackslash{}\textbackslash{}sin\^{}2 t + \textbackslash{}\textbackslash{}cos\^{}2 t + t = 1 + t",}
\NormalTok{    note: "sin²t + cos²t = 1 simplifies the first two terms immediately."}
\NormalTok{  \},}
\NormalTok{  \{}
\NormalTok{    op: "Integrate over t ∈ [0, 2π]",}
\NormalTok{    eq: "W = \textbackslash{}\textbackslash{}int\_0\^{}\{2\textbackslash{}\textbackslash{}pi\} (1 + t)\textbackslash{}\textbackslash{}, dt = \textbackslash{}\textbackslash{}left[t + \textbackslash{}\textbackslash{}frac\{t\^{}2\}\{2\}\textbackslash{}\textbackslash{}right]\_0\^{}\{2\textbackslash{}\textbackslash{}pi\} = 2\textbackslash{}\textbackslash{}pi + \textbackslash{}\textbackslash{}frac\{(2\textbackslash{}\textbackslash{}pi)\^{}2\}\{2\} = 2\textbackslash{}\textbackslash{}pi + 2\textbackslash{}\textbackslash{}pi\^{}2",}
\NormalTok{    note: "A standard definite integral — no special technique needed."}
\NormalTok{  \},}
\NormalTok{  \{}
\NormalTok{    op: "Final answer",}
\NormalTok{    eq: "W = 2\textbackslash{}\textbackslash{}pi(1 + \textbackslash{}\textbackslash{}pi) \textbackslash{}\textbackslash{}approx 26.0 \textbackslash{}\textbackslash{}text\{ (in consistent units)\}",}
\NormalTok{    note: "The field does positive work on the particle throughout the helix."}
\NormalTok{  \}}
\NormalTok{])}
\end{Highlighting}
\end{Shaded}

\begin{center}\rule{0.5\linewidth}{0.5pt}\end{center}

\textbf{2. Conservative field --- find the potential, evaluate by
endpoints.} Show that \(\mathbf{F} = (2xy + z,\; x^2,\; x + 2z)\) is
conservative, find its potential function \(f\), and evaluate
\(\int_C \mathbf{F} \cdot d\mathbf{r}\) from \((0, 0, 0)\) to
\((1, 2, 3)\) along any path.

\begin{Shaded}
\begin{Highlighting}[]
\NormalTok{//| echo: false}
\NormalTok{window.\_makeStepperHTML([}
\NormalTok{  \{}
\NormalTok{    op: "Check curl F = 0",}
\NormalTok{    eq: "(\textbackslash{}\textbackslash{}nabla \textbackslash{}\textbackslash{}times \textbackslash{}\textbackslash{}mathbf\{F\})\_x = \textbackslash{}\textbackslash{}partial\_y(x + 2z) {-} \textbackslash{}\textbackslash{}partial\_z(x\^{}2) = 0 {-} 0 = 0",}
\NormalTok{    note: "Check all three components of the curl."}
\NormalTok{  \},}
\NormalTok{  \{}
\NormalTok{    op: "Check remaining curl components",}
\NormalTok{    eq: "(\textbackslash{}\textbackslash{}nabla \textbackslash{}\textbackslash{}times \textbackslash{}\textbackslash{}mathbf\{F\})\_y = \textbackslash{}\textbackslash{}partial\_z(2xy + z) {-} \textbackslash{}\textbackslash{}partial\_x(x + 2z) = 1 {-} 1 = 0, \textbackslash{}\textbackslash{}quad (\textbackslash{}\textbackslash{}nabla \textbackslash{}\textbackslash{}times \textbackslash{}\textbackslash{}mathbf\{F\})\_z = \textbackslash{}\textbackslash{}partial\_x(x\^{}2) {-} \textbackslash{}\textbackslash{}partial\_y(2xy + z) = 2x {-} 2x = 0",}
\NormalTok{    note: "All components are zero, so F is conservative."}
\NormalTok{  \},}
\NormalTok{  \{}
\NormalTok{    op: "Integrate P = 2xy + z with respect to x",}
\NormalTok{    eq: "f = \textbackslash{}\textbackslash{}int (2xy + z)\textbackslash{}\textbackslash{}, dx = x\^{}2 y + xz + g(y, z)",}
\NormalTok{    note: "The \textquotesingle{}constant\textquotesingle{} g can depend on y and z — treat them as constants when integrating in x."}
\NormalTok{  \},}
\NormalTok{  \{}
\NormalTok{    op: "Match ∂f/∂y = Q = x²",}
\NormalTok{    eq: "\textbackslash{}\textbackslash{}frac\{\textbackslash{}\textbackslash{}partial f\}\{\textbackslash{}\textbackslash{}partial y\} = x\^{}2 + \textbackslash{}\textbackslash{}frac\{\textbackslash{}\textbackslash{}partial g\}\{\textbackslash{}\textbackslash{}partial y\} = x\^{}2 \textbackslash{}\textbackslash{}implies \textbackslash{}\textbackslash{}frac\{\textbackslash{}\textbackslash{}partial g\}\{\textbackslash{}\textbackslash{}partial y\} = 0",}
\NormalTok{    note: "g does not depend on y."}
\NormalTok{  \},}
\NormalTok{  \{}
\NormalTok{    op: "Match ∂f/∂z = R = x + 2z",}
\NormalTok{    eq: "\textbackslash{}\textbackslash{}frac\{\textbackslash{}\textbackslash{}partial f\}\{\textbackslash{}\textbackslash{}partial z\} = x + \textbackslash{}\textbackslash{}frac\{\textbackslash{}\textbackslash{}partial g\}\{\textbackslash{}\textbackslash{}partial z\} = x + 2z \textbackslash{}\textbackslash{}implies \textbackslash{}\textbackslash{}frac\{\textbackslash{}\textbackslash{}partial g\}\{\textbackslash{}\textbackslash{}partial z\} = 2z \textbackslash{}\textbackslash{}implies g(z) = z\^{}2",}
\NormalTok{    note: "g depends only on z."}
\NormalTok{  \},}
\NormalTok{  \{}
\NormalTok{    op: "Potential function",}
\NormalTok{    eq: "f(x, y, z) = x\^{}2 y + xz + z\^{}2",}
\NormalTok{    note: "Verify: ∇f = (2xy + z, x², x + 2z) = F. ✓"}
\NormalTok{  \},}
\NormalTok{  \{}
\NormalTok{    op: "Evaluate by Fundamental Theorem",}
\NormalTok{    eq: "\textbackslash{}\textbackslash{}int\_C \textbackslash{}\textbackslash{}mathbf\{F\} \textbackslash{}\textbackslash{}cdot d\textbackslash{}\textbackslash{}mathbf\{r\} = f(1, 2, 3) {-} f(0, 0, 0) = (1 \textbackslash{}\textbackslash{}cdot 2 + 1 \textbackslash{}\textbackslash{}cdot 3 + 9) {-} 0 = 14",}
\NormalTok{    note: "The path is irrelevant — only the endpoints matter for a conservative field."}
\NormalTok{  \}}
\NormalTok{])}
\end{Highlighting}
\end{Shaded}

\begin{center}\rule{0.5\linewidth}{0.5pt}\end{center}

\textbf{3. Green's theorem --- line integral to double integral.}
Compute \(\oint_C (x^2 - y) \, dx + (x + y^2) \, dy\) where \(C\) is the
boundary of the triangle with vertices \((0,0)\), \((2,0)\), \((0,2)\),
traversed anti-clockwise.

\begin{Shaded}
\begin{Highlighting}[]
\NormalTok{//| echo: false}
\NormalTok{window.\_makeStepperHTML([}
\NormalTok{  \{}
\NormalTok{    op: "Identify P, Q and compute the integrand for Green\textquotesingle{}s theorem",}
\NormalTok{    eq: "P = x\^{}2 {-} y, \textbackslash{}\textbackslash{}quad Q = x + y\^{}2, \textbackslash{}\textbackslash{}quad \textbackslash{}\textbackslash{}frac\{\textbackslash{}\textbackslash{}partial Q\}\{\textbackslash{}\textbackslash{}partial x\} {-} \textbackslash{}\textbackslash{}frac\{\textbackslash{}\textbackslash{}partial P\}\{\textbackslash{}\textbackslash{}partial y\} = 1 {-} ({-}1) = 2",}
\NormalTok{    note: "Green\textquotesingle{}s theorem converts the line integral to ∬\_D 2 dA."}
\NormalTok{  \},}
\NormalTok{  \{}
\NormalTok{    op: "Describe the triangular region D",}
\NormalTok{    eq: "D: \textbackslash{}\textbackslash{}quad 0 \textbackslash{}\textbackslash{}leq x \textbackslash{}\textbackslash{}leq 2, \textbackslash{}\textbackslash{}quad 0 \textbackslash{}\textbackslash{}leq y \textbackslash{}\textbackslash{}leq 2 {-} x",}
\NormalTok{    note: "The hypotenuse has equation x + y = 2, so y goes from 0 to 2 − x."}
\NormalTok{  \},}
\NormalTok{  \{}
\NormalTok{    op: "Apply Green\textquotesingle{}s theorem",}
\NormalTok{    eq: "\textbackslash{}\textbackslash{}oint\_C (x\^{}2 {-} y)\textbackslash{}\textbackslash{}, dx + (x + y\^{}2)\textbackslash{}\textbackslash{}, dy = \textbackslash{}\textbackslash{}iint\_D 2 \textbackslash{}\textbackslash{}, dA = 2 \textbackslash{}\textbackslash{}cdot \textbackslash{}\textbackslash{}text\{Area\}(D)",}
\NormalTok{    note: "The integrand is constant (2), so the double integral is just 2 times the area."}
\NormalTok{  \},}
\NormalTok{  \{}
\NormalTok{    op: "Compute the area of the triangle",}
\NormalTok{    eq: "\textbackslash{}\textbackslash{}text\{Area\}(D) = \textbackslash{}\textbackslash{}frac\{1\}\{2\} \textbackslash{}\textbackslash{}cdot \textbackslash{}\textbackslash{}text\{base\} \textbackslash{}\textbackslash{}cdot \textbackslash{}\textbackslash{}text\{height\} = \textbackslash{}\textbackslash{}frac\{1\}\{2\} \textbackslash{}\textbackslash{}cdot 2 \textbackslash{}\textbackslash{}cdot 2 = 2",}
\NormalTok{    note: "Or integrate: ∫₀² (2−x) dx = [2x − x²/2]₀² = 4 − 2 = 2."}
\NormalTok{  \},}
\NormalTok{  \{}
\NormalTok{    op: "Final answer",}
\NormalTok{    eq: "\textbackslash{}\textbackslash{}oint\_C (x\^{}2 {-} y)\textbackslash{}\textbackslash{}, dx + (x + y\^{}2)\textbackslash{}\textbackslash{}, dy = 2 \textbackslash{}\textbackslash{}cdot 2 = 4",}
\NormalTok{    note: "Evaluating the line integral directly over the three edges would require three separate parametrisations — Green\textquotesingle{}s theorem saved all of that."}
\NormalTok{  \}}
\NormalTok{])}
\end{Highlighting}
\end{Shaded}

\begin{center}\rule{0.5\linewidth}{0.5pt}\end{center}

\textbf{4. Electric flux through a hemisphere.} Compute the upward flux
of \(\mathbf{E} = (x, y, 2z)\) through the upper hemisphere
\(S: x^2 + y^2 + z^2 = 9\), \(z \geq 0\) (with outward normal pointing
away from the origin).

\begin{Shaded}
\begin{Highlighting}[]
\NormalTok{//| echo: false}
\NormalTok{window.\_makeStepperHTML([}
\NormalTok{  \{}
\NormalTok{    op: "Plan: close the surface and use the Divergence theorem",}
\NormalTok{    eq: "\textbackslash{}\textbackslash{}text\{Let \} V = \textbackslash{}\textbackslash{}\{x\^{}2 + y\^{}2 + z\^{}2 \textbackslash{}\textbackslash{}leq 9,\textbackslash{}\textbackslash{}; z \textbackslash{}\textbackslash{}geq 0\textbackslash{}\textbackslash{}\} \textbackslash{}\textbackslash{}text\{ (upper hemisphere + disk at \} z = 0\textbackslash{}\textbackslash{}text\{)\}",}
\NormalTok{    note: "The closed surface ∂V = S (hemisphere) ∪ D (disk). We compute the total and subtract the disk."}
\NormalTok{  \},}
\NormalTok{  \{}
\NormalTok{    op: "Compute the divergence",}
\NormalTok{    eq: "\textbackslash{}\textbackslash{}nabla \textbackslash{}\textbackslash{}cdot \textbackslash{}\textbackslash{}mathbf\{E\} = \textbackslash{}\textbackslash{}frac\{\textbackslash{}\textbackslash{}partial x\}\{\textbackslash{}\textbackslash{}partial x\} + \textbackslash{}\textbackslash{}frac\{\textbackslash{}\textbackslash{}partial y\}\{\textbackslash{}\textbackslash{}partial y\} + \textbackslash{}\textbackslash{}frac\{\textbackslash{}\textbackslash{}partial 2z\}\{\textbackslash{}\textbackslash{}partial z\} = 1 + 1 + 2 = 4",}
\NormalTok{    note: "The divergence is constant, so the volume integral is just 4 times the volume."}
\NormalTok{  \},}
\NormalTok{  \{}
\NormalTok{    op: "Volume of the upper hemisphere",}
\NormalTok{    eq: "V = \textbackslash{}\textbackslash{}frac\{1\}\{2\} \textbackslash{}\textbackslash{}cdot \textbackslash{}\textbackslash{}frac\{4\}\{3\}\textbackslash{}\textbackslash{}pi (3)\^{}3 = \textbackslash{}\textbackslash{}frac\{2\}\{3\}\textbackslash{}\textbackslash{}pi \textbackslash{}\textbackslash{}cdot 27 = 18\textbackslash{}\textbackslash{}pi",}
\NormalTok{    note: "Half the volume of a sphere of radius 3."}
\NormalTok{  \},}
\NormalTok{  \{}
\NormalTok{    op: "Total flux through ∂V by Divergence theorem",}
\NormalTok{    eq: "\textbackslash{}\textbackslash{}oiint\_\{\textbackslash{}\textbackslash{}partial V\} \textbackslash{}\textbackslash{}mathbf\{E\} \textbackslash{}\textbackslash{}cdot d\textbackslash{}\textbackslash{}mathbf\{S\} = \textbackslash{}\textbackslash{}iiint\_V 4 \textbackslash{}\textbackslash{}, dV = 4 \textbackslash{}\textbackslash{}cdot 18\textbackslash{}\textbackslash{}pi = 72\textbackslash{}\textbackslash{}pi",}
\NormalTok{    note: "This is the combined flux through the hemisphere and the disk."}
\NormalTok{  \},}
\NormalTok{  \{}
\NormalTok{    op: "Flux through the disk at z = 0 (downward normal n̂ = −k̂)",}
\NormalTok{    eq: "\textbackslash{}\textbackslash{}iint\_D \textbackslash{}\textbackslash{}mathbf\{E\} \textbackslash{}\textbackslash{}cdot ({-}\textbackslash{}\textbackslash{}hat\{\textbackslash{}\textbackslash{}mathbf\{k\}\})\textbackslash{}\textbackslash{}, dA = \textbackslash{}\textbackslash{}iint\_D {-}2z\textbackslash{}\textbackslash{}, dA = \textbackslash{}\textbackslash{}iint\_D 0\textbackslash{}\textbackslash{}, dA = 0",}
\NormalTok{    note: "On the disk z = 0, so the component 2z = 0. The disk contributes nothing."}
\NormalTok{  \},}
\NormalTok{  \{}
\NormalTok{    op: "Flux through the hemisphere",}
\NormalTok{    eq: "\textbackslash{}\textbackslash{}iint\_S \textbackslash{}\textbackslash{}mathbf\{E\} \textbackslash{}\textbackslash{}cdot d\textbackslash{}\textbackslash{}mathbf\{S\} = 72\textbackslash{}\textbackslash{}pi {-} 0 = 72\textbackslash{}\textbackslash{}pi",}
\NormalTok{    note: "The total equals the hemisphere\textquotesingle{}s contribution alone."}
\NormalTok{  \}}
\NormalTok{])}
\end{Highlighting}
\end{Shaded}

\begin{center}\rule{0.5\linewidth}{0.5pt}\end{center}

\textbf{5. Divergence theorem --- convert surface integral to volume
integral.} Let \(\mathbf{F} = (x^3, y^3, z^3)\). Compute the outward
flux of \(\mathbf{F}\) through the closed surface of the cylinder
\(x^2 + y^2 \leq 4\), \(0 \leq z \leq 3\).

\begin{Shaded}
\begin{Highlighting}[]
\NormalTok{//| echo: false}
\NormalTok{window.\_makeStepperHTML([}
\NormalTok{  \{}
\NormalTok{    op: "Compute the divergence",}
\NormalTok{    eq: "\textbackslash{}\textbackslash{}nabla \textbackslash{}\textbackslash{}cdot \textbackslash{}\textbackslash{}mathbf\{F\} = 3x\^{}2 + 3y\^{}2 + 3z\^{}2",}
\NormalTok{    note: "The closed surface consists of the curved side plus two disks — direct integration would be messy."}
\NormalTok{  \},}
\NormalTok{  \{}
\NormalTok{    op: "Set up the volume integral in cylindrical coordinates",}
\NormalTok{    eq: "\textbackslash{}\textbackslash{}iiint\_V (3x\^{}2 + 3y\^{}2 + 3z\^{}2)\textbackslash{}\textbackslash{}, dV = \textbackslash{}\textbackslash{}int\_0\^{}\{2\textbackslash{}\textbackslash{}pi\} \textbackslash{}\textbackslash{}int\_0\^{}2 \textbackslash{}\textbackslash{}int\_0\^{}3 (3r\^{}2 + 3z\^{}2)\textbackslash{}\textbackslash{}, r\textbackslash{}\textbackslash{}, dz\textbackslash{}\textbackslash{}, dr\textbackslash{}\textbackslash{}, d\textbackslash{}\textbackslash{}theta",}
\NormalTok{    note: "In cylindrical coords: x² + y² = r², dV = r dr dθ dz."}
\NormalTok{  \},}
\NormalTok{  \{}
\NormalTok{    op: "Integrate over z first",}
\NormalTok{    eq: "\textbackslash{}\textbackslash{}int\_0\^{}3 (3r\^{}2 + 3z\^{}2)\textbackslash{}\textbackslash{}, dz = \textbackslash{}\textbackslash{}left[3r\^{}2 z + z\^{}3\textbackslash{}\textbackslash{}right]\_0\^{}3 = 9r\^{}2 + 27",}
\NormalTok{    note: "Treat r as a constant when integrating in z."}
\NormalTok{  \},}
\NormalTok{  \{}
\NormalTok{    op: "Integrate over r",}
\NormalTok{    eq: "\textbackslash{}\textbackslash{}int\_0\^{}2 (9r\^{}2 + 27)\textbackslash{}\textbackslash{}, r\textbackslash{}\textbackslash{}, dr = \textbackslash{}\textbackslash{}int\_0\^{}2 (9r\^{}3 + 27r)\textbackslash{}\textbackslash{}, dr = \textbackslash{}\textbackslash{}left[\textbackslash{}\textbackslash{}frac\{9r\^{}4\}\{4\} + \textbackslash{}\textbackslash{}frac\{27r\^{}2\}\{2\}\textbackslash{}\textbackslash{}right]\_0\^{}2 = 36 + 54 = 90",}
\NormalTok{    note: ""}
\NormalTok{  \},}
\NormalTok{  \{}
\NormalTok{    op: "Integrate over θ",}
\NormalTok{    eq: "\textbackslash{}\textbackslash{}int\_0\^{}\{2\textbackslash{}\textbackslash{}pi\} 90\textbackslash{}\textbackslash{}, d\textbackslash{}\textbackslash{}theta = 180\textbackslash{}\textbackslash{}pi",}
\NormalTok{    note: "The integrand is independent of θ, so the angular integral just multiplies by 2π."}
\NormalTok{  \},}
\NormalTok{  \{}
\NormalTok{    op: "Final answer",}
\NormalTok{    eq: "\textbackslash{}\textbackslash{}oiint\_\{\textbackslash{}\textbackslash{}partial V\} \textbackslash{}\textbackslash{}mathbf\{F\} \textbackslash{}\textbackslash{}cdot d\textbackslash{}\textbackslash{}mathbf\{S\} = 180\textbackslash{}\textbackslash{}pi \textbackslash{}\textbackslash{}approx 565.5",}
\NormalTok{    note: "The Divergence theorem reduced a three{-}part surface integral to one cylindrical volume integral."}
\NormalTok{  \}}
\NormalTok{])}
\end{Highlighting}
\end{Shaded}

\begin{center}\rule{0.5\linewidth}{0.5pt}\end{center}

\textbf{6. Stokes' theorem --- Ampère's law application.} A long
straight wire carries current \(I\) along the \(z\)-axis. The magnetic
field is \(\mathbf{B} = \frac{\mu_0 I}{2\pi(x^2 + y^2)}(-y, x, 0)\). Use
Stokes' theorem to compute \(\oint_C \mathbf{B} \cdot d\mathbf{r}\)
where \(C\) is the circle \(x^2 + y^2 = R^2\) in the \(xy\)-plane.

\begin{Shaded}
\begin{Highlighting}[]
\NormalTok{//| echo: false}
\NormalTok{window.\_makeStepperHTML([}
\NormalTok{  \{}
\NormalTok{    op: "Parametrise the circle C",}
\NormalTok{    eq: "\textbackslash{}\textbackslash{}mathbf\{r\}(t) = (R\textbackslash{}\textbackslash{}cos t,\textbackslash{}\textbackslash{}, R\textbackslash{}\textbackslash{}sin t,\textbackslash{}\textbackslash{}, 0), \textbackslash{}\textbackslash{}quad d\textbackslash{}\textbackslash{}mathbf\{r\} = ({-}R\textbackslash{}\textbackslash{}sin t,\textbackslash{}\textbackslash{}, R\textbackslash{}\textbackslash{}cos t,\textbackslash{}\textbackslash{}, 0)\textbackslash{}\textbackslash{}, dt, \textbackslash{}\textbackslash{}quad t \textbackslash{}\textbackslash{}in [0, 2\textbackslash{}\textbackslash{}pi]",}
\NormalTok{    note: "We evaluate directly first, then verify via Stokes\textquotesingle{} theorem."}
\NormalTok{  \},}
\NormalTok{  \{}
\NormalTok{    op: "Evaluate B on C (where x² + y² = R²)",}
\NormalTok{    eq: "\textbackslash{}\textbackslash{}mathbf\{B\} = \textbackslash{}\textbackslash{}frac\{\textbackslash{}\textbackslash{}mu\_0 I\}\{2\textbackslash{}\textbackslash{}pi R\^{}2\}({-}R\textbackslash{}\textbackslash{}sin t,\textbackslash{}\textbackslash{}, R\textbackslash{}\textbackslash{}cos t,\textbackslash{}\textbackslash{}, 0)",}
\NormalTok{    note: "Substitute x = R cos t, y = R sin t into B."}
\NormalTok{  \},}
\NormalTok{  \{}
\NormalTok{    op: "Compute the dot product B · dr",}
\NormalTok{    eq: "\textbackslash{}\textbackslash{}mathbf\{B\} \textbackslash{}\textbackslash{}cdot d\textbackslash{}\textbackslash{}mathbf\{r\} = \textbackslash{}\textbackslash{}frac\{\textbackslash{}\textbackslash{}mu\_0 I\}\{2\textbackslash{}\textbackslash{}pi R\^{}2\}\textbackslash{}\textbackslash{}left(R\^{}2\textbackslash{}\textbackslash{}sin\^{}2 t + R\^{}2\textbackslash{}\textbackslash{}cos\^{}2 t\textbackslash{}\textbackslash{}right)dt = \textbackslash{}\textbackslash{}frac\{\textbackslash{}\textbackslash{}mu\_0 I\}\{2\textbackslash{}\textbackslash{}pi\}\textbackslash{}\textbackslash{}, dt",}
\NormalTok{    note: "The R² factors cancel; sin²t + cos²t = 1."}
\NormalTok{  \},}
\NormalTok{  \{}
\NormalTok{    op: "Integrate over t ∈ [0, 2π]",}
\NormalTok{    eq: "\textbackslash{}\textbackslash{}oint\_C \textbackslash{}\textbackslash{}mathbf\{B\} \textbackslash{}\textbackslash{}cdot d\textbackslash{}\textbackslash{}mathbf\{r\} = \textbackslash{}\textbackslash{}int\_0\^{}\{2\textbackslash{}\textbackslash{}pi\} \textbackslash{}\textbackslash{}frac\{\textbackslash{}\textbackslash{}mu\_0 I\}\{2\textbackslash{}\textbackslash{}pi\}\textbackslash{}\textbackslash{}, dt = \textbackslash{}\textbackslash{}mu\_0 I",}
\NormalTok{    note: "The result is independent of R — the circulation is the same around any circle enclosing the wire."}
\NormalTok{  \},}
\NormalTok{  \{}
\NormalTok{    op: "Interpret via Stokes\textquotesingle{} theorem (Ampère\textquotesingle{}s law)",}
\NormalTok{    eq: "\textbackslash{}\textbackslash{}oint\_C \textbackslash{}\textbackslash{}mathbf\{B\} \textbackslash{}\textbackslash{}cdot d\textbackslash{}\textbackslash{}mathbf\{r\} = \textbackslash{}\textbackslash{}iint\_S (\textbackslash{}\textbackslash{}nabla \textbackslash{}\textbackslash{}times \textbackslash{}\textbackslash{}mathbf\{B\}) \textbackslash{}\textbackslash{}cdot d\textbackslash{}\textbackslash{}mathbf\{S\} = \textbackslash{}\textbackslash{}mu\_0 \textbackslash{}\textbackslash{}iint\_S \textbackslash{}\textbackslash{}mathbf\{J\} \textbackslash{}\textbackslash{}cdot d\textbackslash{}\textbackslash{}mathbf\{S\} = \textbackslash{}\textbackslash{}mu\_0 I\_\{\textbackslash{}\textbackslash{}text\{enc\}\}",}
\NormalTok{    note: "Stokes\textquotesingle{} theorem links the circulation of B to the current threading through any surface bounded by C. Here I\_enc = I (the full wire current)."}
\NormalTok{  \},}
\NormalTok{  \{}
\NormalTok{    op: "Conclusion",}
\NormalTok{    eq: "\textbackslash{}\textbackslash{}oint\_C \textbackslash{}\textbackslash{}mathbf\{B\} \textbackslash{}\textbackslash{}cdot d\textbackslash{}\textbackslash{}mathbf\{r\} = \textbackslash{}\textbackslash{}mu\_0 I",}
\NormalTok{    note: "This is Ampère\textquotesingle{}s law in integral form. The radius R cancelled because the current is entirely enclosed — the result holds for any loop that encircles the wire once."}
\NormalTok{  \}}
\NormalTok{])}
\end{Highlighting}
\end{Shaded}

\section{Where this goes}\label{where-this-goes-29}

The three theorems in this chapter are the operational heart of
classical field theory. The next chapter, Fourier series and PDEs, puts
them to work in a new setting: instead of integrating over a domain
once, you decompose a function into sinusoidal modes and work with each
mode independently. The Divergence theorem reappears immediately in the
derivation of the heat equation and the wave equation --- those PDEs
arise precisely from applying it to a tiny control volume and letting
the volume shrink to a point.

Complex analysis (the chapter after that) makes this picture even
stranger and more powerful: in the complex plane, the Cauchy integral
theorem is Green's theorem in disguise, and the residue theorem for
computing integrals is a direct consequence of Stokes' theorem applied
to a deformed contour.

The unifying thread is the generalised Stokes' theorem
\(\int_{\partial \Omega} \omega = \int_\Omega d\omega\). Once you have
seen that formulation, you cannot unsee it --- every integration theorem
you encounter from here on is a special case.

\begin{tcolorbox}[enhanced jigsaw, opacityback=0, colback=white, breakable, colframe=quarto-callout-tip-color-frame, arc=.35mm, left=2mm, toprule=.15mm, leftrule=.75mm, bottomrule=.15mm, rightrule=.15mm]
\begin{minipage}[t]{5.5mm}
\textcolor{quarto-callout-tip-color}{\faLightbulb}
\end{minipage}%
\begin{minipage}[t]{\textwidth - 5.5mm}

\textbf{Where this shows up}

\begin{itemize}
\tightlist
\item
  A mechanical engineer computing the torque on a rotating shaft uses
  the surface integral of the stress tensor --- a generalisation of
  flux.
\item
  A geophysicist inverting gravity survey data uses Gauss's theorem to
  relate surface measurements to subsurface mass distribution.
\item
  A climate scientist computing net CO₂ flux into the ocean integrates
  the divergence of a diffusion field over a volume --- the Divergence
  theorem makes global estimates tractable.
\item
  An electrical engineer designing an antenna housing uses Gauss's law
  (Divergence theorem) to verify shielding effectiveness from bulk
  material properties rather than surface-by-surface calculation.
\end{itemize}

The theorems do not just simplify calculation. In each case, they are
the reason the physical law can be stated locally (as a PDE or a
differential form) \emph{and} globally (as a conservation statement).
That equivalence is not obvious --- it is exactly what Green, Stokes,
and Gauss proved.

\end{minipage}%
\end{tcolorbox}

\part{Volume 7 --- Fourier analysis and PDEs}

\chapter{Fourier series}\label{fourier-series}

Every periodic signal is a sum of pure frequencies

\hfill\break

The Fourier series starts from a single observation: the simplest
periodic functions are sines and cosines. They oscillate at a fixed
frequency, they never change shape, and their behaviour is completely
described by amplitude and phase. Every other periodic function ---
however jagged, however asymmetric --- can be built from them by
superposition. The question is: how do you find the right mix?

The answer is integration. Multiply the function by each candidate sine
or cosine and integrate over one period. The result, as we will show, is
exactly the amplitude of that frequency component. Everything else
cancels out.

\begin{center}\rule{0.5\linewidth}{0.5pt}\end{center}

\section{Periodic functions}\label{periodic-functions}

A function \(f\) is \textbf{periodic with period \(T\)} if
\(f(x + T) = f(x)\) for all \(x\). The smallest such \(T > 0\) is called
the \textbf{fundamental period}.

We write \(T = 2L\), which places one full period over the interval
\([-L, L]\). This notation is standard in engineering texts and makes
the coefficient formulas symmetric.

Familiar examples: \(\sin(nx)\) and \(\cos(nx)\) have period \(2\pi/n\);
the square wave (defined below) has the same period as the fundamental
sine used to construct it.

\textbf{The idea geometrically.} A periodic function is completely
determined by its values on any interval of length \(T\). Knowing one
period is knowing everything. The Fourier series asks: which sum of
sines and cosines, with the same period, most accurately represents
\(f\)?

\begin{center}\rule{0.5\linewidth}{0.5pt}\end{center}

\section{The Fourier series}\label{the-fourier-series}

For a function \(f\) of period \(2L\), the \textbf{Fourier series} is:

\[f(x) = \frac{a_0}{2} + \sum_{n=1}^{\infty}\left[a_n\cos\!\left(\frac{n\pi x}{L}\right) + b_n\sin\!\left(\frac{n\pi x}{L}\right)\right]\]

The \textbf{Fourier coefficients} are:

\[a_0 = \frac{1}{L}\int_{-L}^{L} f(x)\,dx\]

\[a_n = \frac{1}{L}\int_{-L}^{L} f(x)\cos\!\left(\frac{n\pi x}{L}\right)dx \qquad (n = 1, 2, 3, \ldots)\]

\[b_n = \frac{1}{L}\int_{-L}^{L} f(x)\sin\!\left(\frac{n\pi x}{L}\right)dx \qquad (n = 1, 2, 3, \ldots)\]

The term \(a_0/2\) is the mean value of \(f\) over one period. The
\(a_n\) terms give the amplitudes of the cosine components; the \(b_n\)
terms give the amplitudes of the sine components.

\subsection{Why these formulas:
orthogonality}\label{why-these-formulas-orthogonality}

The formulas above are not guessed --- they follow from a single
property of the trigonometric functions: \textbf{orthogonality}.

Over the interval \([-L, L]\):

\[\int_{-L}^{L}\cos\!\left(\frac{m\pi x}{L}\right)\cos\!\left(\frac{n\pi x}{L}\right)dx = \begin{cases} 0 & m \neq n \\ L & m = n \geq 1 \\ 2L & m = n = 0 \end{cases}\]

\[\int_{-L}^{L}\sin\!\left(\frac{m\pi x}{L}\right)\sin\!\left(\frac{n\pi x}{L}\right)dx = \begin{cases} 0 & m \neq n \\ L & m = n \geq 1 \end{cases}\]

\[\int_{-L}^{L}\cos\!\left(\frac{m\pi x}{L}\right)\sin\!\left(\frac{n\pi x}{L}\right)dx = 0 \qquad \text{for all } m, n\]

These follow from the product-to-sum identities. Take the \(m \neq n\)
case of the cosine integral as an example: the identity
\(\cos A \cos B = \tfrac{1}{2}[\cos(A-B)+\cos(A+B)]\) turns the product
into a sum of two cosines at frequencies \((m-n)\pi/L\) and
\((m+n)\pi/L\). Integrating either of those over exactly one complete
period gives zero --- the positive and negative half-cycles cancel. When
\(m = n\) the first term becomes \(\cos(0) = 1\), which integrates to
\(2L\), accounting for the factor of \(L\) in the diagonal case. The
sine-sine and cosine-sine cases follow by the same argument with the
analogous identities.

They mean that multiplying the Fourier series by \(\cos(m\pi x/L)\) and
integrating over \([-L, L]\) kills every term except the one with
\(n = m\) --- which is exactly how \(a_m\) is isolated.

Think of the trigonometric functions as a set of directions in function
space that are pairwise perpendicular. The Fourier coefficients are the
projections of \(f\) onto those directions.

\begin{center}\rule{0.5\linewidth}{0.5pt}\end{center}

\section{Worked example: the square
wave}\label{worked-example-the-square-wave}

Define the square wave with period \(2\pi\) as:

\[f(x) = \begin{cases} 1 & 0 < x < \pi \\ -1 & -\pi < x < 0 \end{cases}\]

extended periodically to all \(x\).

\textbf{Step 1: compute \(a_0\).}

\[a_0 = \frac{1}{\pi}\int_{-\pi}^{\pi} f(x)\,dx = \frac{1}{\pi}\left[\int_{-\pi}^{0}(-1)\,dx + \int_{0}^{\pi}(1)\,dx\right] = \frac{1}{\pi}[-\pi + \pi] = 0\]

The mean value is zero --- the positive and negative halves cancel.

\textbf{Step 2: compute \(a_n\).}

\(f(x)\) is an \textbf{odd} function: \(f(-x) = -f(x)\). The cosines
\(\cos(nx)\) are \textbf{even} functions. An odd function times an even
function is odd, and the integral of an odd function over a symmetric
interval is zero:

\[a_n = \frac{1}{\pi}\int_{-\pi}^{\pi} f(x)\cos(nx)\,dx = 0 \quad \text{for all } n \geq 1\]

\textbf{Step 3: compute \(b_n\).}

\(\sin(nx)\) is odd; \(f\) is odd; their product is even:

\[b_n = \frac{1}{\pi}\int_{-\pi}^{\pi} f(x)\sin(nx)\,dx = \frac{2}{\pi}\int_{0}^{\pi}\sin(nx)\,dx\]

\[= \frac{2}{\pi}\left[-\frac{\cos(nx)}{n}\right]_{0}^{\pi} = \frac{2}{n\pi}\left[1 - \cos(n\pi)\right] = \frac{2}{n\pi}\left[1 - (-1)^n\right]\]

Here \(\cos(n\pi) = (-1)^n\) for any integer \(n\): the cosine
alternates between \(-1\) and \(1\) at multiples of \(\pi\). When \(n\)
is even, \(\cos(n\pi) = 1\), so \(b_n = 0\). When \(n\) is odd,
\(\cos(n\pi) = -1\), so \(b_n = 4/(n\pi)\).

\textbf{Result:}

\[f(x) = \frac{4}{\pi}\left[\sin(x) + \frac{\sin(3x)}{3} + \frac{\sin(5x)}{5} + \frac{\sin(7x)}{7} + \cdots\right]\]

The partial sums progressively sharpen toward the square wave, with the
edges becoming steeper as more terms are included.

\begin{Shaded}
\begin{Highlighting}[]
\NormalTok{//| echo: false}
\NormalTok{\{}
\NormalTok{  // ── Fourier partial sum builder with phasor panel ──────────────────────}
\NormalTok{  // Panel A: waveform reconstruction  Panel B: frequency spectrum}
\NormalTok{  // Panel C: phasor wheel (inset)}
\NormalTok{  const d3 = await require("d3@7");}

\NormalTok{  // Palette}
\NormalTok{  const AMBER = "\#f59e0b";}
\NormalTok{  const BLUE  = "\#6b90c4";}
\NormalTok{  const GREY  = "\#9ca3af";}
\NormalTok{  const GREY\_FAINT = "\#e5e7eb";}
\NormalTok{  const RED   = "\#ef4444";}

\NormalTok{  // Dimensions}
\NormalTok{  const WA = 520, HA = 180;   // Panel A waveform}
\NormalTok{  const WB = 520, HB = 90;    // Panel B spectrum}
\NormalTok{  const WC = 130, HC = 130;   // Panel C phasor}
\NormalTok{  const margin = \{ t: 16, r: 16, b: 32, l: 40 \};}

\NormalTok{  // ── Wave functions (period 2π, L = π) ──────────────────────────────────}
\NormalTok{  function targetSquare(x) \{ return x \textgreater{} 0 ? 1 : {-}1; \}}
\NormalTok{  function targetSawtooth(x) \{ return x / Math.PI; \}}
\NormalTok{  function targetTriangle(x) \{}
\NormalTok{    const t = ((x + Math.PI) \% (2 * Math.PI) + 2 * Math.PI) \% (2 * Math.PI) {-} Math.PI;}
\NormalTok{    return 1 {-} (2 / Math.PI) * Math.abs(t);}
\NormalTok{  \}}

\NormalTok{  // ── Fourier coefficients (bₙ only for odd, pure sine series) ───────────}
\NormalTok{  function bCoeff(n, waveType) \{}
\NormalTok{    if (waveType === "square") \{}
\NormalTok{      return n \% 2 === 1 ? 4 / (n * Math.PI) : 0;}
\NormalTok{    \} else if (waveType === "sawtooth") \{}
\NormalTok{      return (n \% 2 === 0 ? {-}1 : 1) * 2 / n;  // alternate signs: 2({-}1)\^{}\{n+1\}/n}
\NormalTok{      // Correct: b\_n = 2({-}1)\^{}\{n+1\}/n}
\NormalTok{    \} else \{  // triangle}
\NormalTok{      if (n \% 2 === 0) return 0;}
\NormalTok{      const k = (n {-} 1) / 2;}
\NormalTok{      return (k \% 2 === 0 ? 8 : {-}8) / (n * n * Math.PI * Math.PI);}
\NormalTok{    \}}
\NormalTok{  \}}

\NormalTok{  // ── Partial sum ─────────────────────────────────────────────────────────}
\NormalTok{  function partialSum(x, N, waveType) \{}
\NormalTok{    let s = 0;}
\NormalTok{    for (let n = 1; n \textless{}= N; n++) s += bCoeff(n, waveType) * Math.sin(n * x);}
\NormalTok{    return s;}
\NormalTok{  \}}

\NormalTok{  // Gibbs overshoot x{-}position (first peak of partial sum near x=0)}
\NormalTok{  function gibbsPeakX(N) \{}
\NormalTok{    return Math.PI / N;  // first peak near jump at x=0 (actually π/(N+1) approx)}
\NormalTok{  \}}

\NormalTok{  // ── Viewof controls ─────────────────────────────────────────────────────}
\NormalTok{  const container = document.createElement("div");}
\NormalTok{  container.style.cssText = "font{-}family:inherit; max{-}width:620px;";}

\NormalTok{  // Controls row}
\NormalTok{  const ctrlRow = document.createElement("div");}
\NormalTok{  ctrlRow.style.cssText = "display:flex; gap:1.5rem; align{-}items:center; margin{-}bottom:0.5rem; flex{-}wrap:wrap;";}

\NormalTok{  // N slider}
\NormalTok{  const sliderLabel = document.createElement("label");}
\NormalTok{  sliderLabel.style.cssText = "font{-}size:0.88em; color:\#374151; display:flex; align{-}items:center; gap:0.4rem;";}
\NormalTok{  sliderLabel.textContent = "Terms N = ";}
\NormalTok{  const nDisplay = document.createElement("strong");}
\NormalTok{  nDisplay.style.cssText = "color:" + AMBER + "; min{-}width:2ch;";}
\NormalTok{  nDisplay.textContent = "5";}
\NormalTok{  const nSlider = document.createElement("input");}
\NormalTok{  nSlider.type = "range"; nSlider.min = 1; nSlider.max = 30; nSlider.value = 5;}
\NormalTok{  nSlider.style.cssText = "width:160px; accent{-}color:" + AMBER + ";";}
\NormalTok{  sliderLabel.appendChild(nSlider);}
\NormalTok{  sliderLabel.appendChild(nDisplay);}
\NormalTok{  ctrlRow.appendChild(sliderLabel);}

\NormalTok{  // Wave type radio}
\NormalTok{  const waveTypes = [["square","Square"], ["sawtooth","Sawtooth"], ["triangle","Triangle"]];}
\NormalTok{  const radioGroup = document.createElement("div");}
\NormalTok{  radioGroup.style.cssText = "display:flex; gap:0.75rem; font{-}size:0.88em;";}
\NormalTok{  let currentWave = "square";}
\NormalTok{  waveTypes.forEach(([val, lab]) =\textgreater{} \{}
\NormalTok{    const lbl = document.createElement("label");}
\NormalTok{    lbl.style.cssText = "display:flex; align{-}items:center; gap:0.25rem; cursor:pointer;";}
\NormalTok{    const r = document.createElement("input");}
\NormalTok{    r.type = "radio"; r.name = "waveViz1"; r.value = val;}
\NormalTok{    r.checked = (val === "square");}
\NormalTok{    r.style.accentColor = AMBER;}
\NormalTok{    r.onchange = () =\textgreater{} \{ currentWave = val; draw(); \};}
\NormalTok{    lbl.appendChild(r); lbl.appendChild(document.createTextNode(lab));}
\NormalTok{    radioGroup.appendChild(lbl);}
\NormalTok{  \});}
\NormalTok{  ctrlRow.appendChild(radioGroup);}
\NormalTok{  container.appendChild(ctrlRow);}

\NormalTok{  // ── SVG canvas ──────────────────────────────────────────────────────────}
\NormalTok{  const totalW = WA + 16 + WC;  // waveform + gap + phasor}
\NormalTok{  const totalH = HA + 8 + HB + margin.b;}
\NormalTok{  const svg = d3.create("svg")}
\NormalTok{    .attr("viewBox", \textasciigrave{}0 0 $\{totalW + margin.l + margin.r\} $\{totalH + margin.t\}\textasciigrave{})}
\NormalTok{    .attr("width", "100\%")}
\NormalTok{    .style("display","block");}

\NormalTok{  const gA = svg.append("g").attr("transform", \textasciigrave{}translate($\{margin.l\},$\{margin.t\})\textasciigrave{});}
\NormalTok{  const gB = svg.append("g").attr("transform", \textasciigrave{}translate($\{margin.l\},$\{margin.t + HA + 8\})\textasciigrave{});}
\NormalTok{  const gC = svg.append("g").attr("transform", \textasciigrave{}translate($\{margin.l + WA + 16\},$\{margin.t\})\textasciigrave{});}

\NormalTok{  // ── Panel A setup ───────────────────────────────────────────────────────}
\NormalTok{  const xA = d3.scaleLinear().domain([{-}Math.PI, Math.PI]).range([0, WA]);}
\NormalTok{  const yA = d3.scaleLinear().domain([{-}1.3, 1.3]).range([HA, 0]);}

\NormalTok{  gA.append("line").attr("x1",0).attr("x2",WA).attr("y1",yA(0)).attr("y2",yA(0))}
\NormalTok{    .attr("stroke", GREY\_FAINT).attr("stroke{-}width",1);}
\NormalTok{  gA.append("g").call(d3.axisBottom(xA).tickValues([{-}Math.PI,{-}Math.PI/2,0,Math.PI/2,Math.PI])}
\NormalTok{    .tickFormat(d =\textgreater{} \{}
\NormalTok{      const m = \{"{-}3.14159":"−π","{-}1.5708":"−π/2","0":"0","1.5708":"π/2","3.14159":"π"\};}
\NormalTok{      return m[d.toFixed(5)] || "";}
\NormalTok{    \})).attr("transform", \textasciigrave{}translate(0,$\{HA\})\textasciigrave{})}
\NormalTok{    .call(g =\textgreater{} \{ g.select(".domain").attr("stroke",GREY\_FAINT); g.selectAll("text").style("font{-}size","10px").attr("fill",GREY); g.selectAll("line").attr("stroke",GREY\_FAINT); \});}
\NormalTok{  gA.append("g").call(d3.axisLeft(yA).ticks(5))}
\NormalTok{    .call(g =\textgreater{} \{ g.select(".domain").remove(); g.selectAll("text").style("font{-}size","10px").attr("fill",GREY); g.selectAll("line").attr("stroke",GREY\_FAINT); \});}
\NormalTok{  gA.append("text").attr("x", WA/2).attr("y", {-}4).attr("text{-}anchor","middle")}
\NormalTok{    .attr("font{-}size","11px").attr("fill","\#374151").text("Waveform reconstruction");}

\NormalTok{  // Target path (blue, drawn once per wave change)}
\NormalTok{  const targetPath = gA.append("path").attr("fill","none").attr("stroke",BLUE)}
\NormalTok{    .attr("stroke{-}width",2).attr("stroke{-}dasharray","4,3").attr("opacity",0.7);}
\NormalTok{  // Partial sum path (amber)}
\NormalTok{  const sumPath = gA.append("path").attr("fill","none").attr("stroke",AMBER).attr("stroke{-}width",2.5);}
\NormalTok{  // Gibbs annotation}
\NormalTok{  const gibbsLine = gA.append("line").attr("stroke","\#dc2626").attr("stroke{-}width",1)}
\NormalTok{    .attr("stroke{-}dasharray","3,2").style("display","none");}
\NormalTok{  const gibbsLabel = gA.append("text").attr("font{-}size","9.5px").attr("fill","\#dc2626")}
\NormalTok{    .style("display","none");}
\NormalTok{  // N label}
\NormalTok{  const nLabel = gA.append("text").attr("x",8).attr("y",14).attr("font{-}size","11px")}
\NormalTok{    .attr("fill",AMBER).attr("font{-}weight","bold");}

\NormalTok{  // ── Panel B setup ───────────────────────────────────────────────────────}
\NormalTok{  const xB = d3.scaleLinear().domain([0.5, 30.5]).range([0, WB]);}
\NormalTok{  const yBmax = 1.4;}
\NormalTok{  const yB = d3.scaleLinear().domain([0, yBmax]).range([HB {-} 22, 0]);}

\NormalTok{  gB.append("text").attr("x",WB/2).attr("y",{-}4).attr("text{-}anchor","middle")}
\NormalTok{    .attr("font{-}size","11px").attr("fill","\#374151").text("Frequency spectrum |bₙ|");}
\NormalTok{  gB.append("g").call(d3.axisBottom(xB).ticks(10).tickFormat(d =\textgreater{} Number.isInteger(d) ? d : ""))}
\NormalTok{    .attr("transform", \textasciigrave{}translate(0,$\{HB{-}22\})\textasciigrave{})}
\NormalTok{    .call(g =\textgreater{} \{ g.select(".domain").attr("stroke",GREY\_FAINT); g.selectAll("text").style("font{-}size","9px").attr("fill",GREY); g.selectAll("line").attr("stroke",GREY\_FAINT); \});}
\NormalTok{  gB.append("text").attr("x",WB+2).attr("y",HB{-}22+4).attr("font{-}size","9px").attr("fill",GREY).text("n");}

\NormalTok{  // ── Panel C setup (phasor wheel) ────────────────────────────────────────}
\NormalTok{  const cx = WC / 2, cy = HC / 2, r0 = WC * 0.38;}
\NormalTok{  gC.append("circle").attr("cx",cx).attr("cy",cy).attr("r",r0)}
\NormalTok{    .attr("fill","none").attr("stroke",GREY\_FAINT).attr("stroke{-}width",1.5);}
\NormalTok{  gC.append("line").attr("x1",cx).attr("x2",cx).attr("y1",cy{-}r0{-}4).attr("y2",cy+r0+4)}
\NormalTok{    .attr("stroke",GREY\_FAINT).attr("stroke{-}width",0.5);}
\NormalTok{  gC.append("line").attr("x1",cx{-}r0{-}4).attr("x2",cx+r0+4).attr("y1",cy).attr("y2",cy)}
\NormalTok{    .attr("stroke",GREY\_FAINT).attr("stroke{-}width",0.5);}
\NormalTok{  gC.append("text").attr("x",WC/2).attr("y",HC{-}2).attr("text{-}anchor","middle")}
\NormalTok{    .attr("font{-}size","9px").attr("fill",GREY).text("Phasor wheel");}

\NormalTok{  // Phasor arrows (up to 3 harmonics for square wave)}
\NormalTok{  const phasorArrows = [1,3,5].map((\_, i) =\textgreater{} (\{}
\NormalTok{    line: gC.append("line").attr("stroke", i===0 ? AMBER : i===1 ? BLUE : "\#a78bfa")}
\NormalTok{            .attr("stroke{-}width", 2{-}i*0.4).attr("marker{-}end","url(\#arrowC)"),}
\NormalTok{    label: gC.append("text").attr("font{-}size","8px").attr("fill", i===0 ? AMBER : i===1 ? BLUE : "\#a78bfa")}
\NormalTok{  \}));}
\NormalTok{  // Arrow marker}
\NormalTok{  svg.append("defs").append("marker").attr("id","arrowC").attr("markerWidth",6).attr("markerHeight",6)}
\NormalTok{    .attr("refX",5).attr("refY",3).attr("orient","auto")}
\NormalTok{    .append("path").attr("d","M0,0 L0,6 L6,3 z").attr("fill",AMBER);}
\NormalTok{  // Projection dot}
\NormalTok{  const projDot = gC.append("circle").attr("r",4).attr("fill",RED);}
\NormalTok{  // Tip trace path}
\NormalTok{  const traceData = [];}
\NormalTok{  const tracePath = gC.append("path").attr("fill","none").attr("stroke",RED)}
\NormalTok{    .attr("stroke{-}width",1).attr("opacity",0.4);}

\NormalTok{  // ── Animation state ─────────────────────────────────────────────────────}
\NormalTok{  let theta = 0;}
\NormalTok{  let animFrame = null;}

\NormalTok{  function getNHarmonics() \{}
\NormalTok{    const N = +nSlider.value;}
\NormalTok{    if (currentWave === "square") \{}
\NormalTok{      const odds = [1,3,5,7,9,11,13,15,17,19,21,23,25,27,29].filter(n =\textgreater{} n \textless{}= N);}
\NormalTok{      return odds.slice(0, 3);}
\NormalTok{    \} else if (currentWave === "sawtooth") \{}
\NormalTok{      return [1,2,3].filter(n =\textgreater{} n \textless{}= N);}
\NormalTok{    \} else \{}
\NormalTok{      const odds = [1,3,5,7,9,11].filter(n =\textgreater{} n \textless{}= N);}
\NormalTok{      return odds.slice(0, 3);}
\NormalTok{    \}}
\NormalTok{  \}}

\NormalTok{  function drawBars(N) \{}
\NormalTok{    const barW = Math.max(1, WB / 32);}
\NormalTok{    gB.selectAll(".specbar").remove();}
\NormalTok{    for (let n = 1; n \textless{}= 30; n++) \{}
\NormalTok{      const coeff = Math.abs(bCoeff(n, currentWave));}
\NormalTok{      const active = n \textless{}= N;}
\NormalTok{      gB.append("rect").attr("class","specbar")}
\NormalTok{        .attr("x", xB(n) {-} barW/2).attr("y", active ? yB(coeff) : HB{-}22)}
\NormalTok{        .attr("width", barW).attr("height", active ? (HB{-}22) {-} yB(coeff) : 0)}
\NormalTok{        .attr("fill", active ? AMBER : "none")}
\NormalTok{        .attr("stroke", active ? AMBER : GREY\_FAINT)}
\NormalTok{        .attr("stroke{-}width", active ? 0 : 1)}
\NormalTok{        .attr("opacity", active ? 1 : 0.5);}
\NormalTok{      if (!active \&\& coeff \textgreater{} 0.02) \{}
\NormalTok{        // ghost outline}
\NormalTok{        const ghostH = Math.max(1, (HB{-}22) {-} yB(coeff));}
\NormalTok{        gB.append("rect").attr("class","specbar")}
\NormalTok{          .attr("x", xB(n) {-} barW/2).attr("y", yB(coeff))}
\NormalTok{          .attr("width", barW).attr("height", ghostH)}
\NormalTok{          .attr("fill","none").attr("stroke",GREY\_FAINT).attr("stroke{-}width",1);}
\NormalTok{      \}}
\NormalTok{    \}}
\NormalTok{  \}}

\NormalTok{  function draw() \{}
\NormalTok{    const N = +nSlider.value;}
\NormalTok{    nDisplay.textContent = N;}
\NormalTok{    nLabel.text(\textasciigrave{}N = $\{N\} terms\textasciigrave{});}

\NormalTok{    const nPts = 400;}
\NormalTok{    const xs = d3.range(nPts).map(i =\textgreater{} {-}Math.PI + (2*Math.PI)*i/(nPts{-}1));}

\NormalTok{    // Target}
\NormalTok{    const target = currentWave === "square" ? targetSquare :}
\NormalTok{                   currentWave === "sawtooth" ? targetSawtooth : targetTriangle;}
\NormalTok{    const lineGen = d3.line().x(d =\textgreater{} xA(d[0])).y(d =\textgreater{} yA(d[1]));}
\NormalTok{    targetPath.datum(xs.map(x =\textgreater{} [x, target(x) || 0])).attr("d", lineGen);}

\NormalTok{    // Partial sum}
\NormalTok{    sumPath.datum(xs.map(x =\textgreater{} [x, partialSum(x, N, currentWave)])).attr("d", lineGen);}

\NormalTok{    // Gibbs annotation (square wave only, N \textgreater{}= 5)}
\NormalTok{    if (currentWave === "square" \&\& N \textgreater{}= 5) \{}
\NormalTok{      const GIBBS\_Y = 1.089;}
\NormalTok{      gibbsLine.style("display","block")}
\NormalTok{        .attr("x1",0).attr("x2",WA).attr("y1",yA(GIBBS\_Y)).attr("y2",yA(GIBBS\_Y));}
\NormalTok{      gibbsLabel.style("display","block")}
\NormalTok{        .attr("x",WA{-}4).attr("y",yA(GIBBS\_Y){-}3)}
\NormalTok{        .attr("text{-}anchor","end").text("≈ 8.9\% overshoot");}
\NormalTok{    \} else \{}
\NormalTok{      gibbsLine.style("display","none");}
\NormalTok{      gibbsLabel.style("display","none");}
\NormalTok{    \}}

\NormalTok{    drawBars(N);}
\NormalTok{  \}}

\NormalTok{  function animatePhasor() \{}
\NormalTok{    const N = +nSlider.value;}
\NormalTok{    const harmonics = getNHarmonics();}
\NormalTok{    theta += 0.025;}

\NormalTok{    // Draw phasor arrows (chain from centre)}
\NormalTok{    let px = cx, py = cy;}
\NormalTok{    const tips = [];}
\NormalTok{    harmonics.forEach((n, i) =\textgreater{} \{}
\NormalTok{      const coeff = Math.abs(bCoeff(n, currentWave));}
\NormalTok{      const len = coeff / 1.5 * r0;  // scale to fit}
\NormalTok{      const angle = {-}(n * theta) {-} Math.PI/2;  // {-}π/2 so n=1 starts pointing up}
\NormalTok{      const ex = px + len * Math.cos(angle);}
\NormalTok{      const ey = py + len * Math.sin(angle);}
\NormalTok{      phasorArrows[i].line}
\NormalTok{        .attr("x1", px).attr("y1", py)}
\NormalTok{        .attr("x2", ex).attr("y2", ey)}
\NormalTok{        .style("display", n \textless{}= N ? "block" : "none");}
\NormalTok{      phasorArrows[i].label}
\NormalTok{        .attr("x", ex + 4).attr("y", ey)}
\NormalTok{        .text(n \textless{}= N ? \textasciigrave{}n=$\{n\}\textasciigrave{} : "")}
\NormalTok{        .style("display", n \textless{}= N ? "block" : "none");}
\NormalTok{      px = ex; py = ey;}
\NormalTok{      tips.push([ex, ey]);}
\NormalTok{    \});}

\NormalTok{    // Hide unused arrows}
\NormalTok{    for (let i = harmonics.length; i \textless{} 3; i++) \{}
\NormalTok{      phasorArrows[i].line.style("display","none");}
\NormalTok{      phasorArrows[i].label.style("display","none");}
\NormalTok{    \}}

\NormalTok{    // Projection onto real axis (horizontal = Re)}
\NormalTok{    const tipX = tips[tips.length{-}1][0];}
\NormalTok{    const tipY = tips[tips.length{-}1][1];}
\NormalTok{    projDot.attr("cx", cx).attr("cy", tipY);  // project to vertical axis (Im → 0 = Re projection)}

\NormalTok{    // Trace tip}
\NormalTok{    traceData.push([tipX, tipY]);}
\NormalTok{    if (traceData.length \textgreater{} 80) traceData.shift();}
\NormalTok{    tracePath.datum(traceData).attr("d", d3.line().x(d=\textgreater{}d[0]).y(d=\textgreater{}d[1]).curve(d3.curveLinear));}

\NormalTok{    animFrame = requestAnimationFrame(animatePhasor);}
\NormalTok{  \}}

\NormalTok{  nSlider.oninput = () =\textgreater{} \{ draw(); \};}

\NormalTok{  container.appendChild(svg.node());}

\NormalTok{  // Legend}
\NormalTok{  const legend = document.createElement("div");}
\NormalTok{  legend.style.cssText = "display:flex; gap:1rem; font{-}size:0.8em; margin{-}top:0.25rem; flex{-}wrap:wrap;";}
\NormalTok{  [["─── Partial sum S\_N(x)", AMBER], ["{-}{-}{-} Target function", BLUE], ["| Gibbs limit (square)", "\#dc2626"]]}
\NormalTok{    .forEach(([lab, col]) =\textgreater{} \{}
\NormalTok{      const s = document.createElement("span");}
\NormalTok{      s.style.cssText = \textasciigrave{}color:$\{col\};\textasciigrave{};}
\NormalTok{      s.textContent = lab;}
\NormalTok{      legend.appendChild(s);}
\NormalTok{    \});}
\NormalTok{  container.appendChild(legend);}

\NormalTok{  draw();}
\NormalTok{  animatePhasor();}

\NormalTok{  // Cleanup on disconnect}
\NormalTok{  const obs = new MutationObserver(() =\textgreater{} \{}
\NormalTok{    if (!document.contains(container)) \{ cancelAnimationFrame(animFrame); obs.disconnect(); \}}
\NormalTok{  \});}
\NormalTok{  obs.observe(document.body, \{childList:true, subtree:true\});}

\NormalTok{  return container;}
\NormalTok{\}}
\end{Highlighting}
\end{Shaded}

\begin{center}\rule{0.5\linewidth}{0.5pt}\end{center}

\section{Convergence and the Gibbs
phenomenon}\label{convergence-and-the-gibbs-phenomenon}

The Fourier series of a well-behaved periodic function converges in the
following sense.

\textbf{Dirichlet conditions.} If \(f\) is periodic, bounded, and
piecewise continuous (at most finitely many jump discontinuities per
period, and finitely many maxima and minima per period), then:

\begin{itemize}
\tightlist
\item
  At points where \(f\) is continuous, the Fourier series converges to
  \(f(x)\).
\item
  At a jump discontinuity \(x_0\), the Fourier series converges to the
  average of the left and right limits:
  \(\frac{1}{2}\left[f(x_0^-) + f(x_0^+)\right]\).
\end{itemize}

\textbf{Gibbs phenomenon.} Near a jump discontinuity, the partial sums
overshoot the true value by approximately \(8.9\%\) of the jump height,
regardless of how many terms are included. Adding more terms moves the
overshoot closer to the discontinuity but does not reduce its size. This
is not an error in the method --- it is a fundamental property of
Fourier partial sums near discontinuities. It appears in real systems as
\textbf{ringing} --- the ripple artefact you see in a sharply filtered
audio signal or a steeply edge-enhanced image.

For the square wave with a jump of magnitude 2, the Gibbs overshoot is
approximately \(2 \times 0.089 \approx 0.18\).

\begin{center}\rule{0.5\linewidth}{0.5pt}\end{center}

\section{Even and odd functions}\label{even-and-odd-functions}

Identifying the symmetry of \(f\) before computing is practical, not
just elegant: it cuts the integration work in half.

\textbf{Even function} (\(f(-x) = f(x)\), symmetric about the
\(y\)-axis): odd-function integrands vanish, so \(b_n = 0\) for all
\(n\). The Fourier series is a \textbf{cosine series}:

\[f(x) = \frac{a_0}{2} + \sum_{n=1}^{\infty} a_n\cos\!\left(\frac{n\pi x}{L}\right)\]

\textbf{Odd function} (\(f(-x) = -f(x)\), antisymmetric about the
\(y\)-axis): even-function integrands vanish, so \(a_n = 0\) for all
\(n\). The Fourier series is a \textbf{sine series}:

\[f(x) = \sum_{n=1}^{\infty} b_n\sin\!\left(\frac{n\pi x}{L}\right)\]

The square wave example above was odd, which is why we only needed to
compute \(b_n\).

\begin{center}\rule{0.5\linewidth}{0.5pt}\end{center}

\section{Half-range series}\label{half-range-series}

Often \(f\) is defined only on \([0, L]\) --- for example, the initial
temperature profile in a rod. We cannot compute a full Fourier series
without knowing \(f\) on \([-L, 0]\), but we can choose how to extend
it:

\begin{itemize}
\item
  \textbf{Even extension}: define \(f(-x) = f(x)\) on \([-L, 0]\). This
  produces a \textbf{Fourier cosine series} with only \(a_n\) terms:
  \[f(x) = \frac{a_0}{2} + \sum_{n=1}^{\infty} a_n\cos\!\left(\frac{n\pi x}{L}\right), \quad a_n = \frac{2}{L}\int_0^L f(x)\cos\!\left(\frac{n\pi x}{L}\right)dx\]
\item
  \textbf{Odd extension}: define \(f(-x) = -f(x)\) on \([-L, 0]\). This
  produces a \textbf{Fourier sine series} with only \(b_n\) terms:
  \[f(x) = \sum_{n=1}^{\infty} b_n\sin\!\left(\frac{n\pi x}{L}\right), \quad b_n = \frac{2}{L}\int_0^L f(x)\sin\!\left(\frac{n\pi x}{L}\right)dx\]
\end{itemize}

\textbf{Which to choose?} The boundary conditions of the physical
problem decide: if \(f(0) = f(L) = 0\) is required (as in a fixed-end
vibrating string or a rod with fixed-temperature endpoints), use the
sine extension. If the derivative \(f'(0) = f'(L) = 0\) is required
(insulated endpoints), use the cosine extension.

\begin{center}\rule{0.5\linewidth}{0.5pt}\end{center}

\section{Complex Fourier series}\label{complex-fourier-series}

Using Euler's formula \(e^{i\theta} = \cos\theta + i\sin\theta\), the
Fourier series can be written in complex exponential form:

\[f(x) = \sum_{n=-\infty}^{\infty} c_n\,e^{in\pi x/L}\]

where the \textbf{complex Fourier coefficients} are:

\[c_n = \frac{1}{2L}\int_{-L}^{L} f(x)\,e^{-in\pi x/L}\,dx\]

The relationship to the real coefficients: \(c_0 = a_0/2\),
\(c_n = (a_n - ib_n)/2\), \(c_{-n} = (a_n + ib_n)/2 = \overline{c_n}\)
for \(n \geq 1\).

The complex form is more compact and connects directly to the Fourier
transform: letting \(L \to \infty\) (so the period grows without bound)
turns the discrete sum into an integral --- the subject of the next
chapter.

\begin{Shaded}
\begin{Highlighting}[]
\NormalTok{//| echo: false}
\NormalTok{\{}
\NormalTok{  // ── OPPORTUNITY 3: Complex exponential phasor wheel ─────────────────────}
\NormalTok{  // Left panel: unit circle with two counter{-}rotating phasors e\^{}\{inθ\}, e\^{}\{{-}inθ\}}
\NormalTok{  // Right panel: real projection traces cos(nθ)}
\NormalTok{  const d3 = await require("d3@7");}

\NormalTok{  const AMBER = "\#f59e0b";}
\NormalTok{  const BLUE  = "\#6b90c4";}
\NormalTok{  const GREY  = "\#9ca3af";}
\NormalTok{  const GREY\_F = "\#e5e7eb";}
\NormalTok{  const RED   = "\#ef4444";}

\NormalTok{  const W = 580, H = 200;}
\NormalTok{  const circleSize = 170;}
\NormalTok{  const cx = circleSize / 2, cy = circleSize / 2, r0 = 70;}
\NormalTok{  const traceW = W {-} circleSize {-} 24;}
\NormalTok{  const traceX0 = circleSize + 24;}

\NormalTok{  const container = document.createElement("div");}
\NormalTok{  container.style.cssText = "font{-}family:inherit; max{-}width:640px;";}

\NormalTok{  // Controls}
\NormalTok{  const ctrlRow = document.createElement("div");}
\NormalTok{  ctrlRow.style.cssText = "display:flex; gap:1.5rem; align{-}items:center; margin{-}bottom:0.4rem; flex{-}wrap:wrap;";}

\NormalTok{  // n slider}
\NormalTok{  const nLbl = document.createElement("label");}
\NormalTok{  nLbl.style.cssText = "font{-}size:0.88em; color:\#374151; display:flex; align{-}items:center; gap:0.4rem;";}
\NormalTok{  nLbl.textContent = "Frequency n = ";}
\NormalTok{  const nDisp = document.createElement("strong");}
\NormalTok{  nDisp.style.cssText = "color:" + AMBER + "; min{-}width:1ch;";}
\NormalTok{  nDisp.textContent = "1";}
\NormalTok{  const nSlider2 = document.createElement("input");}
\NormalTok{  nSlider2.type = "range"; nSlider2.min = 1; nSlider2.max = 5; nSlider2.value = 1;}
\NormalTok{  nSlider2.style.cssText = "width:120px; accent{-}color:" + AMBER + ";";}
\NormalTok{  nLbl.appendChild(nSlider2); nLbl.appendChild(nDisp);}
\NormalTok{  ctrlRow.appendChild(nLbl);}

\NormalTok{  // Play/pause}
\NormalTok{  const playBtn = document.createElement("button");}
\NormalTok{  playBtn.textContent = "⏸ Pause";}
\NormalTok{  playBtn.style.cssText = "padding:0.3rem 0.7rem; border:1px solid \#d1d5db; border{-}radius:4px; background:\#fff; cursor:pointer; font{-}size:0.85em;";}
\NormalTok{  ctrlRow.appendChild(playBtn);}

\NormalTok{  // Speed}
\NormalTok{  const speedLbl = document.createElement("label");}
\NormalTok{  speedLbl.style.cssText = "font{-}size:0.85em; color:\#6b7280; display:flex; align{-}items:center; gap:0.3rem;";}
\NormalTok{  speedLbl.textContent = "Speed:";}
\NormalTok{  const speedSel = document.createElement("select");}
\NormalTok{  speedSel.style.cssText = "font{-}size:0.85em; border:1px solid \#d1d5db; border{-}radius:3px; padding:0.15rem 0.3rem;";}
\NormalTok{  ["Slow","Normal","Fast"].forEach((s,i) =\textgreater{} \{}
\NormalTok{    const o = document.createElement("option"); o.value = [0.5,1,2][i]; o.textContent = s;}
\NormalTok{    if (i===1) o.selected = true;}
\NormalTok{    speedSel.appendChild(o);}
\NormalTok{  \});}
\NormalTok{  speedLbl.appendChild(speedSel);}
\NormalTok{  ctrlRow.appendChild(speedLbl);}
\NormalTok{  container.appendChild(ctrlRow);}

\NormalTok{  // SVG}
\NormalTok{  const svg = d3.create("svg")}
\NormalTok{    .attr("viewBox", \textasciigrave{}0 0 $\{W\} $\{H\}\textasciigrave{})}
\NormalTok{    .attr("width","100\%").style("display","block");}

\NormalTok{  // ── Left: circle panel ──────────────────────────────────────────────────}
\NormalTok{  svg.append("text").attr("x",cx).attr("y",10).attr("text{-}anchor","middle")}
\NormalTok{    .attr("font{-}size","10px").attr("fill","\#374151").text("Complex plane");}
\NormalTok{  svg.append("circle").attr("cx",cx).attr("cy",cy).attr("r",r0)}
\NormalTok{    .attr("fill","none").attr("stroke",GREY\_F).attr("stroke{-}width",1.5);}
\NormalTok{  svg.append("line").attr("x1",cx{-}r0{-}6).attr("x2",cx+r0+6).attr("y1",cy).attr("y2",cy)}
\NormalTok{    .attr("stroke",GREY\_F).attr("stroke{-}width",0.8);}
\NormalTok{  svg.append("line").attr("x1",cx).attr("x2",cx).attr("y1",cy{-}r0{-}6).attr("y2",cy+r0+6)}
\NormalTok{    .attr("stroke",GREY\_F).attr("stroke{-}width",0.8);}
\NormalTok{  svg.append("text").attr("x",cx+r0+4).attr("y",cy+3).attr("font{-}size","9px").attr("fill",GREY).text("Re");}
\NormalTok{  svg.append("text").attr("x",cx{-}3).attr("y",cy{-}r0{-}3).attr("font{-}size","9px").attr("fill",GREY).text("Im");}

\NormalTok{  // Phasor arrows: positive (amber), negative (blue)}
\NormalTok{  const posArrow = svg.append("line").attr("stroke",AMBER).attr("stroke{-}width",2.5);}
\NormalTok{  const negArrow = svg.append("line").attr("stroke",BLUE).attr("stroke{-}width",2.5);}
\NormalTok{  const posLabel = svg.append("text").attr("font{-}size","9px").attr("fill",AMBER).attr("font{-}style","italic");}
\NormalTok{  const negLabel = svg.append("text").attr("font{-}size","9px").attr("fill",BLUE).attr("font{-}style","italic");}

\NormalTok{  // Midpoint (sum/2, lies on Re axis)}
\NormalTok{  const midDot = svg.append("circle").attr("r",5).attr("fill",RED);}
\NormalTok{  // Projection line onto Re axis}
\NormalTok{  const projLine = svg.append("line").attr("stroke",RED).attr("stroke{-}width",1.2)}
\NormalTok{    .attr("stroke{-}dasharray","3,2").attr("opacity",0.6);}

\NormalTok{  // ── Right: trace panel ──────────────────────────────────────────────────}
\NormalTok{  svg.append("text").attr("x",traceX0 + traceW/2).attr("y",10).attr("text{-}anchor","middle")}
\NormalTok{    .attr("font{-}size","10px").attr("fill","\#374151").text("Real projection: cos(nθ)");}

\NormalTok{  // Axes}
\NormalTok{  const xT = d3.scaleLinear().domain([0, 2*Math.PI]).range([traceX0, traceX0+traceW]);}
\NormalTok{  const yT = d3.scaleLinear().domain([{-}1.15, 1.15]).range([H{-}20, 18]);}
\NormalTok{  svg.append("line").attr("x1",traceX0).attr("x2",traceX0+traceW)}
\NormalTok{    .attr("y1",yT(0)).attr("y2",yT(0)).attr("stroke",GREY\_F).attr("stroke{-}width",1);}
\NormalTok{  svg.append("line").attr("x1",traceX0).attr("x2",traceX0).attr("y1",18).attr("y2",H{-}20)}
\NormalTok{    .attr("stroke",GREY\_F).attr("stroke{-}width",1);}
\NormalTok{  [0, Math.PI/2, Math.PI, 3*Math.PI/2, 2*Math.PI].forEach(v =\textgreater{} \{}
\NormalTok{    const label = ["0","π/2","π","3π/2","2π"][Math.round(v/(Math.PI/2))];}
\NormalTok{    svg.append("text").attr("x",xT(v)).attr("y",H{-}8).attr("text{-}anchor","middle")}
\NormalTok{      .attr("font{-}size","8px").attr("fill",GREY).text(label);}
\NormalTok{  \});}
\NormalTok{  svg.append("text").attr("x",traceX0{-}4).attr("y",yT(1)).attr("text{-}anchor","end")}
\NormalTok{    .attr("font{-}size","8px").attr("fill",GREY).text("1");}
\NormalTok{  svg.append("text").attr("x",traceX0{-}4).attr("y",yT({-}1)).attr("text{-}anchor","end")}
\NormalTok{    .attr("font{-}size","8px").attr("fill",GREY).text("−1");}

\NormalTok{  // Reference cosine (blue dashed)}
\NormalTok{  const refPts = d3.range(200).map(i =\textgreater{} \{}
\NormalTok{    const th = 2*Math.PI*i/199;}
\NormalTok{    return [xT(th), yT(Math.cos(1 * th))];}
\NormalTok{  \});}
\NormalTok{  const refLine = svg.append("path").attr("fill","none").attr("stroke",BLUE)}
\NormalTok{    .attr("stroke{-}width",1.5).attr("stroke{-}dasharray","4,3").attr("opacity",0.5);}

\NormalTok{  // Live trace path}
\NormalTok{  const liveData = [];}
\NormalTok{  const livePath = svg.append("path").attr("fill","none").attr("stroke",RED)}
\NormalTok{    .attr("stroke{-}width",2);}
\NormalTok{  // Moving point on trace}
\NormalTok{  const tracePoint = svg.append("circle").attr("r",4.5).attr("fill",RED);}

\NormalTok{  // Formula annotation}
\NormalTok{  const formulaText = svg.append("text").attr("x",traceX0+traceW{-}4).attr("y",26)}
\NormalTok{    .attr("text{-}anchor","end").attr("font{-}size","10px").attr("fill","\#374151")}
\NormalTok{    .attr("font{-}style","italic");}

\NormalTok{  // ── Animation ────────────────────────────────────────────────────────────}
\NormalTok{  let theta3 = 0;}
\NormalTok{  let paused = false;}
\NormalTok{  let rafId = null;}

\NormalTok{  playBtn.onclick = () =\textgreater{} \{}
\NormalTok{    paused = !paused;}
\NormalTok{    playBtn.textContent = paused ? "▶ Play" : "⏸ Pause";}
\NormalTok{    if (!paused) animate();}
\NormalTok{  \};}

\NormalTok{  nSlider2.oninput = () =\textgreater{} \{}
\NormalTok{    nDisp.textContent = nSlider2.value;}
\NormalTok{    liveData.length = 0;}
\NormalTok{    updateRef();}
\NormalTok{  \};}

\NormalTok{  function updateRef() \{}
\NormalTok{    const n = +nSlider2.value;}
\NormalTok{    const pts = d3.range(200).map(i =\textgreater{} \{}
\NormalTok{      const th = 2*Math.PI*i/199;}
\NormalTok{      return [xT(th), yT(Math.cos(n * th))];}
\NormalTok{    \});}
\NormalTok{    refLine.datum(pts).attr("d", d3.line().x(d=\textgreater{}d[0]).y(d=\textgreater{}d[1]));}
\NormalTok{    formulaText.text(\textasciigrave{}Re(e\^{}\{i·$\{n\}·θ\} + e\^{}\{−i·$\{n\}·θ\}) = 2cos($\{n\}θ)\textasciigrave{});}
\NormalTok{  \}}

\NormalTok{  function animate() \{}
\NormalTok{    if (paused) return;}
\NormalTok{    const n = +nSlider2.value;}
\NormalTok{    const speed = +speedSel.value;}
\NormalTok{    theta3 += 0.022 * speed;}
\NormalTok{    const th = theta3 \% (2*Math.PI);}

\NormalTok{    // Positive phasor: e\^{}\{inθ\}}
\NormalTok{    const px = cx + r0 * Math.cos(n * theta3 {-} Math.PI/2);}
\NormalTok{    const py = cy + r0 * Math.sin(n * theta3 {-} Math.PI/2);}
\NormalTok{    // Negative phasor: e\^{}\{{-}inθ\}}
\NormalTok{    const nx\_ = cx + r0 * Math.cos({-}n * theta3 {-} Math.PI/2);}
\NormalTok{    const ny\_ = cy + r0 * Math.sin({-}n * theta3 {-} Math.PI/2);}

\NormalTok{    posArrow.attr("x1",cx).attr("y1",cy).attr("x2",px).attr("y2",py);}
\NormalTok{    negArrow.attr("x1",cx).attr("y1",cy).attr("x2",nx\_).attr("y2",ny\_);}

\NormalTok{    // Labels}
\NormalTok{    posLabel.attr("x",px+4).attr("y",py{-}2).text(\textasciigrave{}e\^{}\{i$\{n\}θ\}\textasciigrave{});}
\NormalTok{    negLabel.attr("x",nx\_+4).attr("y",ny\_+8).text(\textasciigrave{}e\^{}\{−i$\{n\}θ\}\textasciigrave{});}

\NormalTok{    // Midpoint of sum: (e\^{}\{inθ\}+e\^{}\{{-}inθ\})/2 = cos(nθ) on Re axis}
\NormalTok{    const reVal = Math.cos(n * theta3);}
\NormalTok{    const midX = cx + r0 * reVal;}
\NormalTok{    midDot.attr("cx",midX).attr("cy",cy);}
\NormalTok{    projLine.attr("x1",cx).attr("y1",cy).attr("x2",midX).attr("y2",cy);}

\NormalTok{    // Trace}
\NormalTok{    const traceTheta = theta3 \% (2*Math.PI);}
\NormalTok{    liveData.push([xT(traceTheta), yT(reVal)]);}
\NormalTok{    if (liveData.length \textgreater{} 500) liveData.shift();}
\NormalTok{    livePath.datum(liveData).attr("d", d3.line().x(d=\textgreater{}d[0]).y(d=\textgreater{}d[1]).defined((\_,i,arr) =\textgreater{} \{}
\NormalTok{      if (i === 0) return true;}
\NormalTok{      return Math.abs(arr[i][0] {-} arr[i{-}1][0]) \textless{} 15;  // break on wraparound}
\NormalTok{    \}));}
\NormalTok{    tracePoint.attr("cx", xT(traceTheta)).attr("cy", yT(reVal));}

\NormalTok{    rafId = requestAnimationFrame(animate);}
\NormalTok{  \}}

\NormalTok{  updateRef();}
\NormalTok{  animate();}

\NormalTok{  const cleanup = new MutationObserver(() =\textgreater{} \{}
\NormalTok{    if (!document.contains(container)) \{ cancelAnimationFrame(rafId); cleanup.disconnect(); \}}
\NormalTok{  \});}
\NormalTok{  cleanup.observe(document.body, \{childList:true, subtree:true\});}

\NormalTok{  // Legend}
\NormalTok{  const legend = document.createElement("div");}
\NormalTok{  legend.style.cssText = "display:flex; gap:1.2rem; font{-}size:0.78em; margin{-}top:0.2rem; flex{-}wrap:wrap;";}
\NormalTok{  [["→ e\^{}\{inθ\} (positive phasor)", AMBER], ["→ e\^{}\{−inθ\} (negative phasor)", BLUE], ["• Real projection = cos(nθ)", RED]]}
\NormalTok{    .forEach(([lab, col]) =\textgreater{} \{}
\NormalTok{      const s = document.createElement("span"); s.style.color = col; s.textContent = lab;}
\NormalTok{      legend.appendChild(s);}
\NormalTok{    \});}
\NormalTok{  container.appendChild(svg.node());}
\NormalTok{  container.appendChild(legend);}

\NormalTok{  return container;}
\NormalTok{\}}
\end{Highlighting}
\end{Shaded}

\begin{center}\rule{0.5\linewidth}{0.5pt}\end{center}

\section{Parseval's theorem}\label{parsevals-theorem}

The total energy of a periodic signal is distributed across its
frequency components. \textbf{Parseval's theorem} makes this precise:

\[\frac{1}{L}\int_{-L}^{L}|f(x)|^2\,dx = \frac{a_0^2}{2} + \sum_{n=1}^{\infty}(a_n^2 + b_n^2)\]

In the complex form:
\(\frac{1}{2L}\int_{-L}^{L}|f|^2\,dx = \sum_{n=-\infty}^{\infty}|c_n|^2\).

The right side is the sum of squared amplitudes of each Fourier mode ---
the \textbf{power spectrum}. Parseval's theorem says that power is
conserved: the total power in the time domain equals the total power in
the frequency domain. It also provides a shortcut for evaluating certain
series (see Exercise 5).

\begin{center}\rule{0.5\linewidth}{0.5pt}\end{center}

\section{Applications}\label{applications-9}

\subsection{Signal analysis and
filtering}\label{signal-analysis-and-filtering}

A microphone converts a sound wave into a voltage signal \(f(t)\). To
apply an equaliser (boost bass, cut treble), you compute the Fourier
coefficients, multiply each \(b_n\) by the desired gain at frequency
\(n\), and reconstruct. In practice this is done with the Fast Fourier
Transform (FFT), which computes all \(N\) Fourier coefficients in
\(O(N \log N)\) operations instead of \(O(N^2)\).

The Fast Fourier Transform (FFT) is an algorithm that computes all \(N\)
discrete Fourier coefficients in \(O(N \log N)\) arithmetic operations
rather than the \(O(N^2)\) operations needed by direct summation. For
\(N = 1024\) this is roughly a hundredfold speedup; for \(N = 10^6\) the
gap is enormous. The FFT is what makes real-time frequency analysis
practical --- a direct calculation would be too slow for audio, video,
and communications applications.

\subsection{Heat equation initial
conditions}\label{heat-equation-initial-conditions}

The heat equation on a rod of length \(L\) with zero-temperature
endpoints:

\[\frac{\partial u}{\partial t} = \alpha^2\frac{\partial^2 u}{\partial x^2}, \qquad u(0,t) = u(L,t) = 0\]

has general solution
\(u(x,t) = \sum_{n=1}^{\infty} B_n\sin(n\pi x/L)\,e^{-\alpha^2 n^2\pi^2 t/L^2}\).

The coefficients \(B_n\) are exactly the Fourier sine coefficients of
the initial temperature profile \(u(x,0) = f(x)\). Higher-frequency
modes (large \(n\)) decay exponentially faster --- this is why sharp
temperature gradients smooth out quickly while broad gradients persist.
The Fourier series is not just a tool for computing the answer: it
explains the physics.

\begin{center}\rule{0.5\linewidth}{0.5pt}\end{center}

\section{Exercises}\label{exercises-34}

\begin{Shaded}
\begin{Highlighting}[]
\NormalTok{//| echo: false}
\NormalTok{katex = require("katex@0.16")}
\end{Highlighting}
\end{Shaded}

\begin{Shaded}
\begin{Highlighting}[]
\NormalTok{//| echo: false}
\NormalTok{function makeStepperHTML(steps) \{}
\NormalTok{  let current = 0;}
\NormalTok{  const totalSteps = steps.length;}
\NormalTok{  const container = document.createElement("div");}
\NormalTok{  container.style.cssText = "border:1px solid \#e5e7eb; border{-}radius:8px; padding:1rem 1.25rem; margin:0.75rem 0; font{-}family:inherit;";}
\NormalTok{  const stepsDiv = document.createElement("div");}
\NormalTok{  stepsDiv.style.marginBottom = "0.75rem";}

\NormalTok{  function renderSteps() \{}
\NormalTok{    stepsDiv.innerHTML = "";}
\NormalTok{    for (let i = 0; i \textless{} current; i++) \{}
\NormalTok{      const s = steps[i];}
\NormalTok{      const row = document.createElement("div");}
\NormalTok{      row.style.cssText = "display:grid; grid{-}template{-}columns:220px 1fr; gap:0.35rem 1rem; align{-}items:baseline; padding:0.35rem 0; border{-}top:1px solid \#f3f4f6;";}
\NormalTok{      const opCell = document.createElement("span");}
\NormalTok{      opCell.style.cssText = "font{-}size:0.85em; color:\#6b7280; font{-}style:italic;";}
\NormalTok{      opCell.textContent = s.op;}
\NormalTok{      const eqCell = document.createElement("span");}
\NormalTok{      eqCell.innerHTML = katex.renderToString(s.eq, \{ throwOnError: false, displayMode: false \});}
\NormalTok{      row.appendChild(opCell);}
\NormalTok{      row.appendChild(eqCell);}
\NormalTok{      if (s.note) \{}
\NormalTok{        const noteRow = document.createElement("div");}
\NormalTok{        noteRow.style.cssText = "grid{-}column:2; font{-}size:0.82em; color:\#6b7280; padding{-}bottom:0.2rem;";}
\NormalTok{        noteRow.textContent = s.note;}
\NormalTok{        row.appendChild(noteRow);}
\NormalTok{      \}}
\NormalTok{      stepsDiv.appendChild(row);}
\NormalTok{    \}}
\NormalTok{    if (current === totalSteps) \{}
\NormalTok{      const done = document.createElement("div");}
\NormalTok{      done.style.cssText = "margin{-}top:0.5rem; font{-}size:0.9em; color:\#059669; font{-}weight:500;";}
\NormalTok{      done.textContent = "Solution complete.";}
\NormalTok{      stepsDiv.appendChild(done);}
\NormalTok{    \}}
\NormalTok{  \}}

\NormalTok{  const controls = document.createElement("div");}
\NormalTok{  controls.style.cssText = "display:flex; gap:0.5rem; align{-}items:center;";}
\NormalTok{  const nextBtn = document.createElement("button");}
\NormalTok{  nextBtn.textContent = "Next step →";}
\NormalTok{  nextBtn.style.cssText = "padding:0.35rem 0.85rem; border:1px solid \#d1d5db; border{-}radius:4px; background:\#fff; cursor:pointer; font{-}size:0.9em;";}
\NormalTok{  const resetBtn = document.createElement("button");}
\NormalTok{  resetBtn.textContent = "Reset";}
\NormalTok{  resetBtn.style.cssText = "padding:0.35rem 0.75rem; border:1px solid \#e5e7eb; border{-}radius:4px; background:\#f9fafb; cursor:pointer; font{-}size:0.9em; color:\#6b7280;";}
\NormalTok{  const counter = document.createElement("span");}
\NormalTok{  counter.style.cssText = "font{-}size:0.82em; color:\#9ca3af; margin{-}left:0.25rem;";}

\NormalTok{  function updateButtons() \{}
\NormalTok{    nextBtn.disabled = current === totalSteps;}
\NormalTok{    nextBtn.style.opacity = current === totalSteps ? "0.4" : "1";}
\NormalTok{    counter.textContent = current === 0 ? "Click to reveal steps" : \textasciigrave{}Step $\{current\} of $\{totalSteps\}\textasciigrave{};}
\NormalTok{  \}}

\NormalTok{  nextBtn.onclick = () =\textgreater{} \{ if (current \textless{} totalSteps) \{ current++; renderSteps(); updateButtons(); \} \};}
\NormalTok{  resetBtn.onclick = () =\textgreater{} \{ current = 0; renderSteps(); updateButtons(); \};}
\NormalTok{  controls.appendChild(nextBtn);}
\NormalTok{  controls.appendChild(resetBtn);}
\NormalTok{  controls.appendChild(counter);}
\NormalTok{  container.appendChild(stepsDiv);}
\NormalTok{  container.appendChild(controls);}
\NormalTok{  renderSteps();}
\NormalTok{  updateButtons();}
\NormalTok{  return container;}
\NormalTok{\}}
\end{Highlighting}
\end{Shaded}

\begin{center}\rule{0.5\linewidth}{0.5pt}\end{center}

\subsection{Exercise 1: Fourier coefficients of the square
wave}\label{exercise-1-fourier-coefficients-of-the-square-wave}

Compute all Fourier coefficients for \(f(x) = 1\) for \(0 < x < \pi\)
and \(f(x) = -1\) for \(-\pi < x < 0\), period \(2\pi\). Write out the
first four nonzero terms of the series.

\begin{Shaded}
\begin{Highlighting}[]
\NormalTok{//| echo: false}
\NormalTok{makeStepperHTML([}
\NormalTok{  \{}
\NormalTok{    op: "Identify symmetry",}
\NormalTok{    eq: "f({-}x) = {-}f(x) \textbackslash{}\textbackslash{}implies f \textbackslash{}\textbackslash{}text\{ is odd\} \textbackslash{}\textbackslash{}implies a\_0 = 0,\textbackslash{}\textbackslash{}; a\_n = 0 \textbackslash{}\textbackslash{}text\{ for all \} n",}
\NormalTok{    note: "Odd × even = odd; the integral of an odd function over a symmetric interval is zero."}
\NormalTok{  \},}
\NormalTok{  \{}
\NormalTok{    op: "Set up bₙ using symmetry",}
\NormalTok{    eq: "b\_n = \textbackslash{}\textbackslash{}frac\{1\}\{\textbackslash{}\textbackslash{}pi\}\textbackslash{}\textbackslash{}int\_\{{-}\textbackslash{}\textbackslash{}pi\}\^{}\{\textbackslash{}\textbackslash{}pi\} f(x)\textbackslash{}\textbackslash{}sin(nx)\textbackslash{}\textbackslash{},dx = \textbackslash{}\textbackslash{}frac\{2\}\{\textbackslash{}\textbackslash{}pi\}\textbackslash{}\textbackslash{}int\_\{0\}\^{}\{\textbackslash{}\textbackslash{}pi\}\textbackslash{}\textbackslash{}sin(nx)\textbackslash{}\textbackslash{},dx",}
\NormalTok{    note: "Odd × odd = even, so we double the integral over half the period."}
\NormalTok{  \},}
\NormalTok{  \{}
\NormalTok{    op: "Evaluate the integral",}
\NormalTok{    eq: "b\_n = \textbackslash{}\textbackslash{}frac\{2\}\{\textbackslash{}\textbackslash{}pi\}\textbackslash{}\textbackslash{}left[{-}\textbackslash{}\textbackslash{}frac\{\textbackslash{}\textbackslash{}cos(nx)\}\{n\}\textbackslash{}\textbackslash{}right]\_\{0\}\^{}\{\textbackslash{}\textbackslash{}pi\} = \textbackslash{}\textbackslash{}frac\{2\}\{n\textbackslash{}\textbackslash{}pi\}\textbackslash{}\textbackslash{}bigl[1 {-} \textbackslash{}\textbackslash{}cos(n\textbackslash{}\textbackslash{}pi)\textbackslash{}\textbackslash{}bigr] = \textbackslash{}\textbackslash{}frac\{2\}\{n\textbackslash{}\textbackslash{}pi\}\textbackslash{}\textbackslash{}bigl[1 {-} ({-}1)\^{}n\textbackslash{}\textbackslash{}bigr]"}
\NormalTok{  \},}
\NormalTok{  \{}
\NormalTok{    op: "Separate even and odd n",}
\NormalTok{    eq: "b\_n = \textbackslash{}\textbackslash{}begin\{cases\} \textbackslash{}\textbackslash{}dfrac\{4\}\{n\textbackslash{}\textbackslash{}pi\} \& n \textbackslash{}\textbackslash{}text\{ odd\} \textbackslash{}\textbackslash{}\textbackslash{}\textbackslash{} 0 \& n \textbackslash{}\textbackslash{}text\{ even\} \textbackslash{}\textbackslash{}end\{cases\}",}
\NormalTok{    note: "For even n, cos(nπ) = 1 so 1 − (−1)ⁿ = 0. For odd n, cos(nπ) = −1 so 1 − (−1)ⁿ = 2."}
\NormalTok{  \},}
\NormalTok{  \{}
\NormalTok{    op: "Write the Fourier series",}
\NormalTok{    eq: "f(x) = \textbackslash{}\textbackslash{}frac\{4\}\{\textbackslash{}\textbackslash{}pi\}\textbackslash{}\textbackslash{}left[\textbackslash{}\textbackslash{}sin x + \textbackslash{}\textbackslash{}frac\{\textbackslash{}\textbackslash{}sin 3x\}\{3\} + \textbackslash{}\textbackslash{}frac\{\textbackslash{}\textbackslash{}sin 5x\}\{5\} + \textbackslash{}\textbackslash{}frac\{\textbackslash{}\textbackslash{}sin 7x\}\{7\} + \textbackslash{}\textbackslash{}cdots\textbackslash{}\textbackslash{}right]"}
\NormalTok{  \}}
\NormalTok{])}
\end{Highlighting}
\end{Shaded}

\begin{center}\rule{0.5\linewidth}{0.5pt}\end{center}

\subsection{Exercise 2: Sawtooth wave Fourier
coefficients}\label{exercise-2-sawtooth-wave-fourier-coefficients}

The function \(f(x) = x\) for \(-\pi < x < \pi\), extended with period
\(2\pi\), is the \textbf{sawtooth wave}. Compute its Fourier series.

\begin{Shaded}
\begin{Highlighting}[]
\NormalTok{//| echo: false}
\NormalTok{makeStepperHTML([}
\NormalTok{  \{}
\NormalTok{    op: "Identify symmetry",}
\NormalTok{    eq: "f({-}x) = {-}x = {-}f(x) \textbackslash{}\textbackslash{}implies f \textbackslash{}\textbackslash{}text\{ is odd\} \textbackslash{}\textbackslash{}implies a\_0 = 0,\textbackslash{}\textbackslash{}; a\_n = 0",}
\NormalTok{    note: "All cosine terms vanish."}
\NormalTok{  \},}
\NormalTok{  \{}
\NormalTok{    op: "Set up bₙ",}
\NormalTok{    eq: "b\_n = \textbackslash{}\textbackslash{}frac\{2\}\{\textbackslash{}\textbackslash{}pi\}\textbackslash{}\textbackslash{}int\_\{0\}\^{}\{\textbackslash{}\textbackslash{}pi\} x\textbackslash{}\textbackslash{}sin(nx)\textbackslash{}\textbackslash{},dx",}
\NormalTok{    note: "Using the symmetry reduction: odd × odd = even, double the half{-}period integral."}
\NormalTok{  \},}
\NormalTok{  \{}
\NormalTok{    op: "Integrate by parts: u = x, dv = sin(nx)dx",}
\NormalTok{    eq: "\textbackslash{}\textbackslash{}int x\textbackslash{}\textbackslash{}sin(nx)\textbackslash{}\textbackslash{},dx = {-}\textbackslash{}\textbackslash{}frac\{x\textbackslash{}\textbackslash{}cos(nx)\}\{n\} + \textbackslash{}\textbackslash{}frac\{\textbackslash{}\textbackslash{}sin(nx)\}\{n\^{}2\} + C"}
\NormalTok{  \},}
\NormalTok{  \{}
\NormalTok{    op: "Evaluate at limits 0 to π",}
\NormalTok{    eq: "\textbackslash{}\textbackslash{}int\_0\^{}\textbackslash{}\textbackslash{}pi x\textbackslash{}\textbackslash{}sin(nx)\textbackslash{}\textbackslash{},dx = \textbackslash{}\textbackslash{}left[{-}\textbackslash{}\textbackslash{}frac\{x\textbackslash{}\textbackslash{}cos(nx)\}\{n\}\textbackslash{}\textbackslash{}right]\_0\^{}\textbackslash{}\textbackslash{}pi + \textbackslash{}\textbackslash{}left[\textbackslash{}\textbackslash{}frac\{\textbackslash{}\textbackslash{}sin(nx)\}\{n\^{}2\}\textbackslash{}\textbackslash{}right]\_0\^{}\textbackslash{}\textbackslash{}pi = {-}\textbackslash{}\textbackslash{}frac\{\textbackslash{}\textbackslash{}pi\textbackslash{}\textbackslash{}cos(n\textbackslash{}\textbackslash{}pi)\}\{n\} = \textbackslash{}\textbackslash{}frac\{({-}1)\^{}\{n+1\}\textbackslash{}\textbackslash{}pi\}\{n\}",}
\NormalTok{    note: "sin(nπ) = 0 for all integer n."}
\NormalTok{  \},}
\NormalTok{  \{}
\NormalTok{    op: "Compute bₙ",}
\NormalTok{    eq: "b\_n = \textbackslash{}\textbackslash{}frac\{2\}\{\textbackslash{}\textbackslash{}pi\}\textbackslash{}\textbackslash{}cdot\textbackslash{}\textbackslash{}frac\{({-}1)\^{}\{n+1\}\textbackslash{}\textbackslash{}pi\}\{n\} = \textbackslash{}\textbackslash{}frac\{2({-}1)\^{}\{n+1\}\}\{n\}"}
\NormalTok{  \},}
\NormalTok{  \{}
\NormalTok{    op: "Write the Fourier series",}
\NormalTok{    eq: "f(x) = 2\textbackslash{}\textbackslash{}left[\textbackslash{}\textbackslash{}sin x {-} \textbackslash{}\textbackslash{}frac\{\textbackslash{}\textbackslash{}sin 2x\}\{2\} + \textbackslash{}\textbackslash{}frac\{\textbackslash{}\textbackslash{}sin 3x\}\{3\} {-} \textbackslash{}\textbackslash{}frac\{\textbackslash{}\textbackslash{}sin 4x\}\{4\} + \textbackslash{}\textbackslash{}cdots\textbackslash{}\textbackslash{}right]",}
\NormalTok{    note: "At x = π the sawtooth jumps from π to −π. The Fourier series converges to 0 = (π + (−π))/2 there — the midpoint of the jump."}
\NormalTok{  \}}
\NormalTok{])}
\end{Highlighting}
\end{Shaded}

\begin{center}\rule{0.5\linewidth}{0.5pt}\end{center}

\subsection{\texorpdfstring{Exercise 3: Half-range cosine series for
\(f(x) = x\)}{Exercise 3: Half-range cosine series for f(x) = x}}\label{exercise-3-half-range-cosine-series-for-fx-x}

Find the half-range Fourier cosine series of \(f(x) = x\) on
\(0 < x < \pi\).

\begin{Shaded}
\begin{Highlighting}[]
\NormalTok{//| echo: false}
\NormalTok{makeStepperHTML([}
\NormalTok{  \{}
\NormalTok{    op: "Apply cosine series formulas (even extension)",}
\NormalTok{    eq: "a\_0 = \textbackslash{}\textbackslash{}frac\{2\}\{\textbackslash{}\textbackslash{}pi\}\textbackslash{}\textbackslash{}int\_0\^{}\textbackslash{}\textbackslash{}pi x\textbackslash{}\textbackslash{},dx = \textbackslash{}\textbackslash{}frac\{2\}\{\textbackslash{}\textbackslash{}pi\}\textbackslash{}\textbackslash{}cdot\textbackslash{}\textbackslash{}frac\{\textbackslash{}\textbackslash{}pi\^{}2\}\{2\} = \textbackslash{}\textbackslash{}pi"}
\NormalTok{  \},}
\NormalTok{  \{}
\NormalTok{    op: "Set up aₙ",}
\NormalTok{    eq: "a\_n = \textbackslash{}\textbackslash{}frac\{2\}\{\textbackslash{}\textbackslash{}pi\}\textbackslash{}\textbackslash{}int\_0\^{}\textbackslash{}\textbackslash{}pi x\textbackslash{}\textbackslash{}cos(nx)\textbackslash{}\textbackslash{},dx"}
\NormalTok{  \},}
\NormalTok{  \{}
\NormalTok{    op: "Integrate by parts: u = x, dv = cos(nx)dx",}
\NormalTok{    eq: "\textbackslash{}\textbackslash{}int\_0\^{}\textbackslash{}\textbackslash{}pi x\textbackslash{}\textbackslash{}cos(nx)\textbackslash{}\textbackslash{},dx = \textbackslash{}\textbackslash{}left[\textbackslash{}\textbackslash{}frac\{x\textbackslash{}\textbackslash{}sin(nx)\}\{n\}\textbackslash{}\textbackslash{}right]\_0\^{}\textbackslash{}\textbackslash{}pi {-} \textbackslash{}\textbackslash{}frac\{1\}\{n\}\textbackslash{}\textbackslash{}int\_0\^{}\textbackslash{}\textbackslash{}pi \textbackslash{}\textbackslash{}sin(nx)\textbackslash{}\textbackslash{},dx = 0 + \textbackslash{}\textbackslash{}frac\{1\}\{n\^{}2\}\textbackslash{}\textbackslash{}bigl[\textbackslash{}\textbackslash{}cos(nx)\textbackslash{}\textbackslash{}bigr]\_0\^{}\textbackslash{}\textbackslash{}pi",}
\NormalTok{    note: "sin(nπ) = sin(0) = 0."}
\NormalTok{  \},}
\NormalTok{  \{}
\NormalTok{    op: "Evaluate",}
\NormalTok{    eq: "\textbackslash{}\textbackslash{}int\_0\^{}\textbackslash{}\textbackslash{}pi x\textbackslash{}\textbackslash{}cos(nx)\textbackslash{}\textbackslash{},dx = \textbackslash{}\textbackslash{}frac\{\textbackslash{}\textbackslash{}cos(n\textbackslash{}\textbackslash{}pi) {-} 1\}\{n\^{}2\} = \textbackslash{}\textbackslash{}frac\{({-}1)\^{}n {-} 1\}\{n\^{}2\}",}
\NormalTok{    note: "For even n: (−1)ⁿ − 1 = 0, so aₙ = 0. For odd n: (−1)ⁿ − 1 = −2."}
\NormalTok{  \},}
\NormalTok{  \{}
\NormalTok{    op: "Final coefficient",}
\NormalTok{    eq: "a\_n = \textbackslash{}\textbackslash{}frac\{2\}\{\textbackslash{}\textbackslash{}pi\}\textbackslash{}\textbackslash{}cdot\textbackslash{}\textbackslash{}frac\{({-}1)\^{}n {-} 1\}\{n\^{}2\} = \textbackslash{}\textbackslash{}begin\{cases\} {-}\textbackslash{}\textbackslash{}dfrac\{4\}\{n\^{}2\textbackslash{}\textbackslash{}pi\} \& n \textbackslash{}\textbackslash{}text\{ odd\} \textbackslash{}\textbackslash{}\textbackslash{}\textbackslash{} 0 \& n \textbackslash{}\textbackslash{}text\{ even\} \textbackslash{}\textbackslash{}end\{cases\}"}
\NormalTok{  \},}
\NormalTok{  \{}
\NormalTok{    op: "Write the cosine series",}
\NormalTok{    eq: "x = \textbackslash{}\textbackslash{}frac\{\textbackslash{}\textbackslash{}pi\}\{2\} {-} \textbackslash{}\textbackslash{}frac\{4\}\{\textbackslash{}\textbackslash{}pi\}\textbackslash{}\textbackslash{}left[\textbackslash{}\textbackslash{}cos x + \textbackslash{}\textbackslash{}frac\{\textbackslash{}\textbackslash{}cos 3x\}\{9\} + \textbackslash{}\textbackslash{}frac\{\textbackslash{}\textbackslash{}cos 5x\}\{25\} + \textbackslash{}\textbackslash{}cdots\textbackslash{}\textbackslash{}right], \textbackslash{}\textbackslash{}quad 0 \textless{} x \textless{} \textbackslash{}\textbackslash{}pi",}
\NormalTok{    note: "Compare with Exercise 2: the sawtooth sine series and this cosine series represent the same function on (0,π) but differ outside."}
\NormalTok{  \}}
\NormalTok{])}
\end{Highlighting}
\end{Shaded}

\begin{center}\rule{0.5\linewidth}{0.5pt}\end{center}

\subsection{\texorpdfstring{Exercise 4: Half-range sine series for
\(f(x) = 1\)}{Exercise 4: Half-range sine series for f(x) = 1}}\label{exercise-4-half-range-sine-series-for-fx-1}

Find the half-range Fourier sine series of \(f(x) = 1\) on
\(0 < x < L\). Interpret the result.

\begin{Shaded}
\begin{Highlighting}[]
\NormalTok{//| echo: false}
\NormalTok{makeStepperHTML([}
\NormalTok{  \{}
\NormalTok{    op: "Apply sine series formula (odd extension)",}
\NormalTok{    eq: "b\_n = \textbackslash{}\textbackslash{}frac\{2\}\{L\}\textbackslash{}\textbackslash{}int\_0\^{}L \textbackslash{}\textbackslash{}sin\textbackslash{}\textbackslash{}!\textbackslash{}\textbackslash{}left(\textbackslash{}\textbackslash{}frac\{n\textbackslash{}\textbackslash{}pi x\}\{L\}\textbackslash{}\textbackslash{}right)dx"}
\NormalTok{  \},}
\NormalTok{  \{}
\NormalTok{    op: "Evaluate the integral",}
\NormalTok{    eq: "b\_n = \textbackslash{}\textbackslash{}frac\{2\}\{L\}\textbackslash{}\textbackslash{}left[{-}\textbackslash{}\textbackslash{}frac\{L\}\{n\textbackslash{}\textbackslash{}pi\}\textbackslash{}\textbackslash{}cos\textbackslash{}\textbackslash{}!\textbackslash{}\textbackslash{}left(\textbackslash{}\textbackslash{}frac\{n\textbackslash{}\textbackslash{}pi x\}\{L\}\textbackslash{}\textbackslash{}right)\textbackslash{}\textbackslash{}right]\_0\^{}L = \textbackslash{}\textbackslash{}frac\{2\}\{n\textbackslash{}\textbackslash{}pi\}\textbackslash{}\textbackslash{}bigl[1 {-} \textbackslash{}\textbackslash{}cos(n\textbackslash{}\textbackslash{}pi)\textbackslash{}\textbackslash{}bigr] = \textbackslash{}\textbackslash{}frac\{2(1 {-} ({-}1)\^{}n)\}\{n\textbackslash{}\textbackslash{}pi\}"}
\NormalTok{  \},}
\NormalTok{  \{}
\NormalTok{    op: "Separate even and odd n",}
\NormalTok{    eq: "b\_n = \textbackslash{}\textbackslash{}begin\{cases\} \textbackslash{}\textbackslash{}dfrac\{4\}\{n\textbackslash{}\textbackslash{}pi\} \& n \textbackslash{}\textbackslash{}text\{ odd\} \textbackslash{}\textbackslash{}\textbackslash{}\textbackslash{} 0 \& n \textbackslash{}\textbackslash{}text\{ even\} \textbackslash{}\textbackslash{}end\{cases\}"}
\NormalTok{  \},}
\NormalTok{  \{}
\NormalTok{    op: "Write the sine series",}
\NormalTok{    eq: "1 = \textbackslash{}\textbackslash{}frac\{4\}\{\textbackslash{}\textbackslash{}pi\}\textbackslash{}\textbackslash{}left[\textbackslash{}\textbackslash{}sin\textbackslash{}\textbackslash{}!\textbackslash{}\textbackslash{}left(\textbackslash{}\textbackslash{}frac\{\textbackslash{}\textbackslash{}pi x\}\{L\}\textbackslash{}\textbackslash{}right) + \textbackslash{}\textbackslash{}frac\{1\}\{3\}\textbackslash{}\textbackslash{}sin\textbackslash{}\textbackslash{}!\textbackslash{}\textbackslash{}left(\textbackslash{}\textbackslash{}frac\{3\textbackslash{}\textbackslash{}pi x\}\{L\}\textbackslash{}\textbackslash{}right) + \textbackslash{}\textbackslash{}frac\{1\}\{5\}\textbackslash{}\textbackslash{}sin\textbackslash{}\textbackslash{}!\textbackslash{}\textbackslash{}left(\textbackslash{}\textbackslash{}frac\{5\textbackslash{}\textbackslash{}pi x\}\{L\}\textbackslash{}\textbackslash{}right) + \textbackslash{}\textbackslash{}cdots\textbackslash{}\textbackslash{}right]",}
\NormalTok{    note: "This is the same coefficient pattern as the square wave — here applied to a constant function on a half{-}range."}
\NormalTok{  \},}
\NormalTok{  \{}
\NormalTok{    op: "Physical interpretation",}
\NormalTok{    eq: "\textbackslash{}\textbackslash{}text\{Setting \} x = L/2: \textbackslash{}\textbackslash{}quad 1 = \textbackslash{}\textbackslash{}frac\{4\}\{\textbackslash{}\textbackslash{}pi\}\textbackslash{}\textbackslash{}left[1 {-} \textbackslash{}\textbackslash{}frac\{1\}\{3\} + \textbackslash{}\textbackslash{}frac\{1\}\{5\} {-} \textbackslash{}\textbackslash{}cdots\textbackslash{}\textbackslash{}right] \textbackslash{}\textbackslash{}implies \textbackslash{}\textbackslash{}frac\{\textbackslash{}\textbackslash{}pi\}\{4\} = 1 {-} \textbackslash{}\textbackslash{}frac\{1\}\{3\} + \textbackslash{}\textbackslash{}frac\{1\}\{5\} {-} \textbackslash{}\textbackslash{}cdots",}
\NormalTok{    note: "This is the Leibniz formula for π — a bonus convergent series identity from evaluating the Fourier series at x = L/2."}
\NormalTok{  \}}
\NormalTok{])}
\end{Highlighting}
\end{Shaded}

\begin{center}\rule{0.5\linewidth}{0.5pt}\end{center}

\subsection{Exercise 5: Parseval's theorem and a series
sum}\label{exercise-5-parsevals-theorem-and-a-series-sum}

Apply Parseval's theorem to the square wave result from Exercise 1 to
evaluate \(\displaystyle\sum_{k=0}^{\infty}\frac{1}{(2k+1)^2}\).

\begin{Shaded}
\begin{Highlighting}[]
\NormalTok{//| echo: false}
\NormalTok{makeStepperHTML([}
\NormalTok{  \{}
\NormalTok{    op: "Write Parseval\textquotesingle{}s theorem for the square wave",}
\NormalTok{    eq: "\textbackslash{}\textbackslash{}frac\{1\}\{\textbackslash{}\textbackslash{}pi\}\textbackslash{}\textbackslash{}int\_\{{-}\textbackslash{}\textbackslash{}pi\}\^{}\{\textbackslash{}\textbackslash{}pi\}|f(x)|\^{}2\textbackslash{}\textbackslash{},dx = \textbackslash{}\textbackslash{}frac\{a\_0\^{}2\}\{2\} + \textbackslash{}\textbackslash{}sum\_\{n=1\}\^{}\{\textbackslash{}\textbackslash{}infty\}(a\_n\^{}2 + b\_n\^{}2)",}
\NormalTok{    note: "Here L = π, so the factor is 1/π."}
\NormalTok{  \},}
\NormalTok{  \{}
\NormalTok{    op: "Evaluate the left side",}
\NormalTok{    eq: "\textbackslash{}\textbackslash{}frac\{1\}\{\textbackslash{}\textbackslash{}pi\}\textbackslash{}\textbackslash{}int\_\{{-}\textbackslash{}\textbackslash{}pi\}\^{}\{\textbackslash{}\textbackslash{}pi\}1\^{}2\textbackslash{}\textbackslash{},dx = \textbackslash{}\textbackslash{}frac\{1\}\{\textbackslash{}\textbackslash{}pi\}\textbackslash{}\textbackslash{}cdot 2\textbackslash{}\textbackslash{}pi = 2",}
\NormalTok{    note: "|f(x)|² = 1 everywhere on (−π, π)."}
\NormalTok{  \},}
\NormalTok{  \{}
\NormalTok{    op: "Evaluate the right side using bₙ = 4/(nπ) for odd n, 0 for even n",}
\NormalTok{    eq: "\textbackslash{}\textbackslash{}sum\_\{n=1\}\^{}\{\textbackslash{}\textbackslash{}infty\} b\_n\^{}2 = \textbackslash{}\textbackslash{}sum\_\{k=0\}\^{}\{\textbackslash{}\textbackslash{}infty\}\textbackslash{}\textbackslash{}left(\textbackslash{}\textbackslash{}frac\{4\}\{(2k+1)\textbackslash{}\textbackslash{}pi\}\textbackslash{}\textbackslash{}right)\^{}2 = \textbackslash{}\textbackslash{}frac\{16\}\{\textbackslash{}\textbackslash{}pi\^{}2\}\textbackslash{}\textbackslash{}sum\_\{k=0\}\^{}\{\textbackslash{}\textbackslash{}infty\}\textbackslash{}\textbackslash{}frac\{1\}\{(2k+1)\^{}2\}"}
\NormalTok{  \},}
\NormalTok{  \{}
\NormalTok{    op: "Set left = right and solve",}
\NormalTok{    eq: "2 = \textbackslash{}\textbackslash{}frac\{16\}\{\textbackslash{}\textbackslash{}pi\^{}2\}\textbackslash{}\textbackslash{}sum\_\{k=0\}\^{}\{\textbackslash{}\textbackslash{}infty\}\textbackslash{}\textbackslash{}frac\{1\}\{(2k+1)\^{}2\} \textbackslash{}\textbackslash{}implies \textbackslash{}\textbackslash{}sum\_\{k=0\}\^{}\{\textbackslash{}\textbackslash{}infty\}\textbackslash{}\textbackslash{}frac\{1\}\{(2k+1)\^{}2\} = \textbackslash{}\textbackslash{}frac\{\textbackslash{}\textbackslash{}pi\^{}2\}\{8\}"}
\NormalTok{  \},}
\NormalTok{  \{}
\NormalTok{    op: "Verify the first few terms",}
\NormalTok{    eq: "1 + \textbackslash{}\textbackslash{}frac\{1\}\{9\} + \textbackslash{}\textbackslash{}frac\{1\}\{25\} + \textbackslash{}\textbackslash{}frac\{1\}\{49\} + \textbackslash{}\textbackslash{}cdots \textbackslash{}\textbackslash{}approx 1.000 + 0.111 + 0.040 + 0.020 + \textbackslash{}\textbackslash{}cdots \textbackslash{}\textbackslash{}approx 1.233 \textbackslash{}\textbackslash{}approx \textbackslash{}\textbackslash{}frac\{\textbackslash{}\textbackslash{}pi\^{}2\}\{8\}",}
\NormalTok{    note: "π²/8 ≈ 1.2337. Convergence is slow but the limit is exact."}
\NormalTok{  \}}
\NormalTok{])}
\end{Highlighting}
\end{Shaded}

\begin{center}\rule{0.5\linewidth}{0.5pt}\end{center}

\subsection{Exercise 6: Heat equation --- reading off the solution from
a Fourier
series}\label{exercise-6-heat-equation-reading-off-the-solution-from-a-fourier-series}

A rod of length \(L = \pi\) has initial temperature
\(u(x,0) = 3\sin(x) - \sin(2x)\) for \(0 < x < \pi\). The endpoints are
held at zero: \(u(0,t) = u(\pi,t) = 0\). The heat equation is
\(u_t = \alpha^2 u_{xx}\).

Write down the complete solution \(u(x,t)\) and identify which component
decays faster.

\begin{Shaded}
\begin{Highlighting}[]
\NormalTok{//| echo: false}
\NormalTok{makeStepperHTML([}
\NormalTok{  \{}
\NormalTok{    op: "Recall the general solution form",}
\NormalTok{    eq: "u(x,t) = \textbackslash{}\textbackslash{}sum\_\{n=1\}\^{}\{\textbackslash{}\textbackslash{}infty\} B\_n \textbackslash{}\textbackslash{}sin(nx)\textbackslash{}\textbackslash{},e\^{}\{{-}\textbackslash{}\textbackslash{}alpha\^{}2 n\^{}2 t\}",}
\NormalTok{    note: "For L = π, the sine modes are sin(nx) and the decay rates are α²n²."}
\NormalTok{  \},}
\NormalTok{  \{}
\NormalTok{    op: "Match initial condition to the general form",}
\NormalTok{    eq: "u(x,0) = \textbackslash{}\textbackslash{}sum\_\{n=1\}\^{}\{\textbackslash{}\textbackslash{}infty\} B\_n \textbackslash{}\textbackslash{}sin(nx) = 3\textbackslash{}\textbackslash{}sin(x) {-} \textbackslash{}\textbackslash{}sin(2x)",}
\NormalTok{    note: "The initial condition is already a finite Fourier sine series — no integration needed."}
\NormalTok{  \},}
\NormalTok{  \{}
\NormalTok{    op: "Read off the coefficients",}
\NormalTok{    eq: "B\_1 = 3, \textbackslash{}\textbackslash{}quad B\_2 = {-}1, \textbackslash{}\textbackslash{}quad B\_n = 0 \textbackslash{}\textbackslash{}text\{ for \} n \textbackslash{}\textbackslash{}geq 3"}
\NormalTok{  \},}
\NormalTok{  \{}
\NormalTok{    op: "Write the solution",}
\NormalTok{    eq: "u(x,t) = 3\textbackslash{}\textbackslash{}sin(x)\textbackslash{}\textbackslash{},e\^{}\{{-}\textbackslash{}\textbackslash{}alpha\^{}2 t\} {-} \textbackslash{}\textbackslash{}sin(2x)\textbackslash{}\textbackslash{},e\^{}\{{-}4\textbackslash{}\textbackslash{}alpha\^{}2 t\}"}
\NormalTok{  \},}
\NormalTok{  \{}
\NormalTok{    op: "Compare decay rates",}
\NormalTok{    eq: "\textbackslash{}\textbackslash{}text\{Mode \} n=2 \textbackslash{}\textbackslash{}text\{ decays at rate \} 4\textbackslash{}\textbackslash{}alpha\^{}2 \textbackslash{}\textbackslash{}text\{ vs \} \textbackslash{}\textbackslash{}alpha\^{}2 \textbackslash{}\textbackslash{}text\{ for \} n=1: \textbackslash{}\textbackslash{}quad \textbackslash{}\textbackslash{}text\{four times faster.\}",}
\NormalTok{    note: "High{-}frequency components (sharp features) dissipate faster than low{-}frequency ones. At large t, u(x,t) ≈ 3 sin(x) e\^{}\{−α²t\}: only the fundamental mode survives."}
\NormalTok{  \}}
\NormalTok{])}
\end{Highlighting}
\end{Shaded}

\chapter{Fourier transforms}\label{fourier-transforms}

From periodic decomposition to continuous frequency analysis

\hfill\break

The Fourier series produces a discrete spectrum: a list of amplitudes at
frequencies \(\pi/L, 2\pi/L, 3\pi/L, \ldots\) spaced \(\pi/L\) apart. As
\(L \to \infty\), the spacing between frequencies shrinks to zero, the
list of amplitudes becomes a continuous function of frequency, and the
sum becomes an integral. This is the Fourier transform.

The transition from series to transform is not just a mathematical
technicality. It extends the analysis to any signal that decays to zero
as \(|t| \to \infty\) --- a one-shot pulse, a finite-duration transient,
or any square-integrable function. The series was restricted to
functions that repeat forever. The transform has no such restriction.

\begin{center}\rule{0.5\linewidth}{0.5pt}\end{center}

\section{From Fourier series to the Fourier
integral}\label{from-fourier-series-to-the-fourier-integral}

Consider a function \(f\) of period \(2L\) with Fourier series:

\[f(x) = \sum_{n=-\infty}^{\infty} c_n\,e^{in\pi x/L}, \qquad c_n = \frac{1}{2L}\int_{-L}^{L} f(x)\,e^{-in\pi x/L}\,dx\]

Define \(\omega_n = n\pi/L\) and \(\Delta\omega = \pi/L\) (the spacing
between successive frequencies). Then:

\[f(x) = \sum_{n=-\infty}^{\infty} \frac{\Delta\omega}{2\pi}\left[\int_{-L}^{L} f(t)\,e^{-i\omega_n t}\,dt\right]e^{i\omega_n x}\]

The rewrite uses a single substitution: since \(\Delta\omega = \pi/L\),
the prefactor \(1/(2L)\) in \(c_n\) equals \(\Delta\omega/(2\pi)\).
Replacing it turns each coefficient into
\((\Delta\omega/2\pi)\times[\text{integral}]\), and the whole sum
becomes a Riemann sum with step size \(\Delta\omega\) --- ready to
become an integral.

Let \(L \to \infty\): \(\Delta\omega \to d\omega\), the sum becomes an
integral, and we obtain the \textbf{Fourier integral representation}:

\[f(x) = \frac{1}{2\pi}\int_{-\infty}^{\infty}\left[\int_{-\infty}^{\infty} f(t)\,e^{-i\omega t}\,dt\right]e^{i\omega x}\,d\omega\]

The inner bracket is the \textbf{Fourier transform} of \(f\) evaluated
at \(\omega\).

\begin{Shaded}
\begin{Highlighting}[]
\NormalTok{//| echo: false}
\NormalTok{\{}
\NormalTok{  // ── OPPORTUNITY 1: Series{-}to{-}transform transition ────────────────────────}
\NormalTok{  // Rectangular pulse f(t)=1 for |t|\textless{}1. Fourier series coefficients on period 2L}
\NormalTok{  // as L increases → bars multiply, narrow, trace the sinc envelope.}
\NormalTok{  const d3 = await require("d3@7");}

\NormalTok{  const AMBER = "\#f59e0b";}
\NormalTok{  const BLUE  = "\#6b90c4";}
\NormalTok{  const GREY  = "\#9ca3af";}
\NormalTok{  const GREY\_F = "\#e5e7eb";}

\NormalTok{  const W = 560, H = 200, PL = 46, PR = 16, PT = 28, PB = 28;}
\NormalTok{  const IW = W {-} PL {-} PR, IH = H {-} PT {-} PB;}

\NormalTok{  const container = document.createElement("div");}
\NormalTok{  container.style.cssText = "font{-}family:inherit; max{-}width:620px;";}

\NormalTok{  // Slider for L (as multiples of π)}
\NormalTok{  const ctrlRow = document.createElement("div");}
\NormalTok{  ctrlRow.style.cssText = "display:flex; gap:1rem; align{-}items:center; margin{-}bottom:0.4rem; flex{-}wrap:wrap;";}
\NormalTok{  const lLbl = document.createElement("label");}
\NormalTok{  lLbl.style.cssText = "font{-}size:0.88em; color:\#374151; display:flex; align{-}items:center; gap:0.4rem;";}
\NormalTok{  lLbl.textContent = "Period  L = ";}
\NormalTok{  const lDisp = document.createElement("strong");}
\NormalTok{  lDisp.style.cssText = "color:" + AMBER + "; min{-}width:4ch;";}
\NormalTok{  lDisp.textContent = "1π";}
\NormalTok{  const lSlider = document.createElement("input");}
\NormalTok{  lSlider.type = "range"; lSlider.min = 1; lSlider.max = 20; lSlider.value = 1;}
\NormalTok{  lSlider.style.cssText = "width:200px; accent{-}color:" + AMBER + ";";}
\NormalTok{  lLbl.appendChild(lSlider); lLbl.appendChild(lDisp);}
\NormalTok{  ctrlRow.appendChild(lLbl);}

\NormalTok{  const statusLabel = document.createElement("span");}
\NormalTok{  statusLabel.style.cssText = "font{-}size:0.82em; color:" + BLUE + "; font{-}style:italic;";}
\NormalTok{  statusLabel.textContent = "Discrete spectrum: Fourier series";}
\NormalTok{  ctrlRow.appendChild(statusLabel);}
\NormalTok{  container.appendChild(ctrlRow);}

\NormalTok{  // SVG}
\NormalTok{  const svg = d3.create("svg")}
\NormalTok{    .attr("viewBox", \textasciigrave{}0 0 $\{W\} $\{H\}\textasciigrave{})}
\NormalTok{    .attr("width","100\%").style("display","block");}
\NormalTok{  const g = svg.append("g").attr("transform", \textasciigrave{}translate($\{PL\},$\{PT\})\textasciigrave{});}

\NormalTok{  // Axes — ω from {-}4π to +4π}
\NormalTok{  const omegaMax = 4 * Math.PI;}
\NormalTok{  const xS = d3.scaleLinear().domain([{-}omegaMax, omegaMax]).range([0, IW]);}
\NormalTok{  const yS = d3.scaleLinear().domain([0, 0.75]).range([IH, 0]);}

\NormalTok{  g.append("line").attr("x1",0).attr("x2",IW).attr("y1",IH).attr("y2",IH)}
\NormalTok{    .attr("stroke",GREY\_F).attr("stroke{-}width",1);}
\NormalTok{  g.append("line").attr("x1",xS(0)).attr("x2",xS(0)).attr("y1",0).attr("y2",IH)}
\NormalTok{    .attr("stroke",GREY\_F).attr("stroke{-}width",1);}

\NormalTok{  // ω{-}axis ticks}
\NormalTok{  [{-}4,{-}3,{-}2,{-}1,0,1,2,3,4].forEach(k =\textgreater{} \{}
\NormalTok{    const lab = k===0?"0":k\textgreater{}0?\textasciigrave{}$\{k\}π\textasciigrave{}:\textasciigrave{}$\{k\}π\textasciigrave{};}
\NormalTok{    g.append("text").attr("x",xS(k*Math.PI)).attr("y",IH+12)}
\NormalTok{      .attr("text{-}anchor","middle").attr("font{-}size","9px").attr("fill",GREY).text(lab);}
\NormalTok{    g.append("line").attr("x1",xS(k*Math.PI)).attr("x2",xS(k*Math.PI)).attr("y1",IH).attr("y2",IH+4)}
\NormalTok{      .attr("stroke",GREY\_F);}
\NormalTok{  \});}
\NormalTok{  g.append("text").attr("x",IW+2).attr("y",IH+4).attr("font{-}size","9px").attr("fill",GREY).text("ω");}

\NormalTok{  // y{-}axis}
\NormalTok{  g.append("g").call(d3.axisLeft(yS).ticks(4).tickFormat(d3.format(".2f")))}
\NormalTok{    .call(gg =\textgreater{} \{ gg.select(".domain").remove(); gg.selectAll("text").style("font{-}size","9px").attr("fill",GREY); gg.selectAll("line").attr("stroke",GREY\_F); \});}

\NormalTok{  // Continuous sinc envelope (blue, static)}
\NormalTok{  const sincLine = g.append("path").attr("fill","none").attr("stroke",BLUE)}
\NormalTok{    .attr("stroke{-}width",1.8).attr("stroke{-}dasharray","5,3").attr("opacity",0.6);}
\NormalTok{  const sincPts = d3.range(800).map(i =\textgreater{} \{}
\NormalTok{    const omega = {-}omegaMax + 2*omegaMax*i/799;}
\NormalTok{    const val = Math.abs(omega) \textless{} 0.001 ? 2/(2*Math.PI) : Math.abs(Math.sin(omega))/(Math.PI*Math.abs(omega));}
\NormalTok{    return [xS(omega), yS(val)];}
\NormalTok{  \});}
\NormalTok{  sincLine.datum(sincPts).attr("d", d3.line().x(d=\textgreater{}d[0]).y(d=\textgreater{}d[1]));}

\NormalTok{  // Bar container}
\NormalTok{  const barsG = g.append("g");}

\NormalTok{  // Title}
\NormalTok{  g.append("text").attr("x",IW/2).attr("y",{-}12).attr("text{-}anchor","middle")}
\NormalTok{    .attr("font{-}size","11px").attr("fill","\#374151")}
\NormalTok{    .text("Fourier series coefficients → Fourier transform");}

\NormalTok{  function drawBars(Lpi) \{}
\NormalTok{    const L = Lpi * Math.PI;}
\NormalTok{    const deltaOmega = Math.PI / L;}
\NormalTok{    const nMax = Math.floor(omegaMax / deltaOmega) + 1;}
\NormalTok{    barsG.selectAll("*").remove();}
\NormalTok{    const barW = Math.max(0.8, (IW / (2*omegaMax)) * deltaOmega * 0.7);}

\NormalTok{    for (let n = {-}nMax; n \textless{}= nMax; n++) \{}
\NormalTok{      const omega\_n = n * deltaOmega;}
\NormalTok{      // Coefficient height: (1/2π) * |F(ω\_n)| * Δω = (1/2π) * 2|sinc(ω\_n)| * Δω}
\NormalTok{      // Display: (2/π)|sinc(ω\_n)| so y{-}axis matches continuous limit}
\NormalTok{      const sinc = Math.abs(omega\_n) \textless{} 1e{-}6 ? 1 : Math.sin(omega\_n)/omega\_n;}
\NormalTok{      const h = Math.abs(sinc) / Math.PI;}
\NormalTok{      if (h \textless{} 1e{-}6) continue;}
\NormalTok{      const xPos = xS(omega\_n);}
\NormalTok{      if (xPos \textless{} 0 || xPos \textgreater{} IW) continue;}
\NormalTok{      const yPos = yS(h);}
\NormalTok{      const barH = IH {-} yPos;}
\NormalTok{      barsG.append("rect")}
\NormalTok{        .attr("x", xPos {-} barW/2).attr("y", yPos)}
\NormalTok{        .attr("width", barW).attr("height", Math.max(1, barH))}
\NormalTok{        .attr("fill", AMBER).attr("opacity", 0.85);}
\NormalTok{    \}}

\NormalTok{    lDisp.textContent = \textasciigrave{}$\{Lpi\}π\textasciigrave{};}
\NormalTok{    statusLabel.textContent = Lpi \textgreater{}= 18}
\NormalTok{      ? "Continuous spectrum: Fourier transform"}
\NormalTok{      : "Discrete spectrum: Fourier series";}
\NormalTok{    statusLabel.style.color = Lpi \textgreater{}= 18 ? BLUE : AMBER;}
\NormalTok{  \}}

\NormalTok{  lSlider.oninput = () =\textgreater{} drawBars(+lSlider.value);}
\NormalTok{  drawBars(1);}

\NormalTok{  // Legend}
\NormalTok{  const legend = document.createElement("div");}
\NormalTok{  legend.style.cssText = "display:flex; gap:1.2rem; font{-}size:0.78em; margin{-}top:0.2rem;";}
\NormalTok{  [["▌ Fourier series coefficients", AMBER], ["{-}{-}{-} Continuous transform envelope", BLUE]]}
\NormalTok{    .forEach(([lab, col]) =\textgreater{} \{}
\NormalTok{      const s = document.createElement("span"); s.style.color = col; s.textContent = lab;}
\NormalTok{      legend.appendChild(s);}
\NormalTok{    \});}
\NormalTok{  container.appendChild(svg.node());}
\NormalTok{  container.appendChild(legend);}
\NormalTok{  return container;}
\NormalTok{\}}
\end{Highlighting}
\end{Shaded}

\begin{center}\rule{0.5\linewidth}{0.5pt}\end{center}

\section{The Fourier transform}\label{the-fourier-transform}

The \textbf{Fourier transform} of \(f(t)\) is:

\[F(\omega) = \mathcal{F}\{f\}(\omega) = \int_{-\infty}^{\infty} f(t)\,e^{-i\omega t}\,dt\]

The \textbf{inverse Fourier transform} recovers \(f\) from \(F\):

\[f(t) = \mathcal{F}^{-1}\{F\}(t) = \frac{1}{2\pi}\int_{-\infty}^{\infty} F(\omega)\,e^{i\omega t}\,d\omega\]

Together these are a \textbf{transform pair}:
\(f(t) \leftrightarrow F(\omega)\). Knowing either one determines the
other completely (for well-behaved functions).

\textbf{Conventions.} There are three common conventions for where the
\(2\pi\) factor sits and whether the sign in the exponent is
\(-i\omega t\) or \(+i\omega t\). The one above is the
\textbf{engineering convention} (also used by Kreyszig and Stroud). The
physics convention places \(1/\sqrt{2\pi}\) on both the forward and
inverse transform. When using a reference or software tool, check which
convention it follows.

\textbf{Existence.} The Fourier transform is guaranteed to exist (as an
ordinary integral) when \(f\) is absolutely integrable:
\(\int_{-\infty}^{\infty}|f(t)|\,dt < \infty\). For the larger class of
square-integrable functions (\(\int|f|^2 < \infty\)), it exists in the
\(L^2\) sense. Important special cases --- the unit step, the delta
function, the constant function --- require the theory of distributions
but yield well-defined transforms.

\begin{Shaded}
\begin{Highlighting}[]
\NormalTok{//| echo: false}
\NormalTok{\{}
\NormalTok{  // ── OPPORTUNITY 3: Fourier transform as winding diagram ─────────────────}
\NormalTok{  // f(t) = cos(t) + 0.5cos(3t). Winding at freq ω winds f onto the complex plane.}
\NormalTok{  // Centre of mass of the wound curve = F(ω)/(2π).}
\NormalTok{  const d3 = await require("d3@7");}

\NormalTok{  const AMBER = "\#f59e0b";}
\NormalTok{  const BLUE  = "\#6b90c4";}
\NormalTok{  const GREY  = "\#9ca3af";}
\NormalTok{  const GREY\_F = "\#e5e7eb";}
\NormalTok{  const RED   = "\#ef4444";}

\NormalTok{  const W = 580, H = 210;}
\NormalTok{  const leftW = 210, rightW = 200, gap = 20, barW = 90;}
\NormalTok{  const circR = 80;}
\NormalTok{  const circCx = leftW + gap + rightW/2, circCy = H/2;}

\NormalTok{  const container = document.createElement("div");}
\NormalTok{  container.style.cssText = "font{-}family:inherit; max{-}width:640px;";}

\NormalTok{  // ω slider}
\NormalTok{  const ctrlRow = document.createElement("div");}
\NormalTok{  ctrlRow.style.cssText = "display:flex; gap:1rem; align{-}items:center; margin{-}bottom:0.4rem; flex{-}wrap:wrap;";}
\NormalTok{  const omLbl = document.createElement("label");}
\NormalTok{  omLbl.style.cssText = "font{-}size:0.88em; color:\#374151; display:flex; align{-}items:center; gap:0.4rem;";}
\NormalTok{  omLbl.textContent = "Testing frequency ω = ";}
\NormalTok{  const omDisp = document.createElement("strong");}
\NormalTok{  omDisp.style.cssText = "color:" + AMBER + "; min{-}width:3ch;";}
\NormalTok{  omDisp.textContent = "1.0";}
\NormalTok{  const omSlider = document.createElement("input");}
\NormalTok{  omSlider.type = "range"; omSlider.min = 0; omSlider.max = 50; omSlider.value = 10; // steps of 0.1}
\NormalTok{  omSlider.style.cssText = "width:180px; accent{-}color:" + AMBER + ";";}
\NormalTok{  omLbl.appendChild(omSlider); omLbl.appendChild(omDisp);}
\NormalTok{  ctrlRow.appendChild(omLbl);}
\NormalTok{  ctrlRow.appendChild(Object.assign(document.createElement("span"), \{}
\NormalTok{    textContent: "rad/s",}
\NormalTok{    style: \{ cssText: "font{-}size:0.82em; color:\#9ca3af;" \}}
\NormalTok{  \}));}
\NormalTok{  container.appendChild(ctrlRow);}

\NormalTok{  const svg = d3.create("svg")}
\NormalTok{    .attr("viewBox", \textasciigrave{}0 0 $\{W\} $\{H\}\textasciigrave{})}
\NormalTok{    .attr("width","100\%").style("display","block");}

\NormalTok{  // ── Left panel: time domain signal ──────────────────────────────────────}
\NormalTok{  const tMin = 0, tMax = 2*Math.PI;}
\NormalTok{  const xL = d3.scaleLinear().domain([tMin, tMax]).range([8, leftW{-}8]);}
\NormalTok{  const yL = d3.scaleLinear().domain([{-}1.6, 1.6]).range([H{-}20, 16]);}

\NormalTok{  svg.append("text").attr("x", leftW/2).attr("y", 12).attr("text{-}anchor","middle")}
\NormalTok{    .attr("font{-}size","10px").attr("fill","\#374151").text("f(t) = cos(t) + 0.5cos(3t)");}

\NormalTok{  svg.append("line").attr("x1",8).attr("x2",leftW{-}8).attr("y1",yL(0)).attr("y2",yL(0))}
\NormalTok{    .attr("stroke",GREY\_F).attr("stroke{-}width",1);}
\NormalTok{  svg.append("line").attr("x1",xL(tMin)).attr("x2",xL(tMin)).attr("y1",16).attr("y2",H{-}20)}
\NormalTok{    .attr("stroke",GREY\_F).attr("stroke{-}width",1);}

\NormalTok{  // f(t) curve}
\NormalTok{  const fPts = d3.range(300).map(i =\textgreater{} \{}
\NormalTok{    const t = tMin + (tMax{-}tMin)*i/299;}
\NormalTok{    return [xL(t), yL(Math.cos(t) + 0.5*Math.cos(3*t))];}
\NormalTok{  \});}
\NormalTok{  svg.append("path").attr("fill","none").attr("stroke",BLUE).attr("stroke{-}width",2)}
\NormalTok{    .datum(fPts).attr("d", d3.line().x(d=\textgreater{}d[0]).y(d=\textgreater{}d[1]));}

\NormalTok{  // Moving time{-}bar}
\NormalTok{  const timeLine = svg.append("line").attr("stroke",AMBER).attr("stroke{-}width",1.5).attr("opacity",0.6);}
\NormalTok{  // Time{-}domain dot}
\NormalTok{  const timeDot = svg.append("circle").attr("r",4).attr("fill",AMBER);}

\NormalTok{  [0,Math.PI/2,Math.PI,3*Math.PI/2,2*Math.PI].forEach(t =\textgreater{} \{}
\NormalTok{    svg.append("text").attr("x",xL(t)).attr("y",H{-}8).attr("text{-}anchor","middle")}
\NormalTok{      .attr("font{-}size","8px").attr("fill",GREY).text(["0","π/2","π","3π/2","2π"][Math.round(t/(Math.PI/2))]);}
\NormalTok{  \});}

\NormalTok{  // ── Middle panel: winding diagram ────────────────────────────────────────}
\NormalTok{  svg.append("text").attr("x",circCx).attr("y",12).attr("text{-}anchor","middle")}
\NormalTok{    .attr("font{-}size","10px").attr("fill","\#374151").text("Wound{-}up f(t)e\^{}\{−iωt\}");}

\NormalTok{  svg.append("circle").attr("cx",circCx).attr("cy",circCy).attr("r",circR)}
\NormalTok{    .attr("fill","none").attr("stroke",GREY\_F).attr("stroke{-}width",1.5);}
\NormalTok{  svg.append("line").attr("x1",circCx{-}circR{-}4).attr("x2",circCx+circR+4)}
\NormalTok{    .attr("y1",circCy).attr("y2",circCy).attr("stroke",GREY\_F).attr("stroke{-}width",0.8);}
\NormalTok{  svg.append("line").attr("x1",circCx).attr("x2",circCx)}
\NormalTok{    .attr("y1",circCy{-}circR{-}4).attr("y2",circCy+circR+4).attr("stroke",GREY\_F).attr("stroke{-}width",0.8);}
\NormalTok{  svg.append("text").attr("x",circCx+circR+2).attr("y",circCy+3).attr("font{-}size","8px").attr("fill",GREY).text("Re");}
\NormalTok{  svg.append("text").attr("x",circCx{-}2).attr("y",circCy{-}circR{-}2).attr("font{-}size","8px").attr("fill",GREY).text("Im");}

\NormalTok{  const woundPath = svg.append("path").attr("fill","none").attr("stroke",AMBER)}
\NormalTok{    .attr("stroke{-}width",1.5).attr("opacity",0.85);}
\NormalTok{  const comDot = svg.append("circle").attr("r",6).attr("fill",RED).attr("opacity",0.9);}
\NormalTok{  const comLine = svg.append("line").attr("stroke",RED).attr("stroke{-}width",2).attr("opacity",0.7);}
\NormalTok{  const comLabel = svg.append("text").attr("font{-}size","9px").attr("fill",RED);}

\NormalTok{  // ── Right panel: |F(ω)| bar ──────────────────────────────────────────────}
\NormalTok{  const barX0 = leftW + gap + rightW + 12;}
\NormalTok{  svg.append("text").attr("x",barX0+barW/2).attr("y",12).attr("text{-}anchor","middle")}
\NormalTok{    .attr("font{-}size","10px").attr("fill","\#374151").text("|F(ω)|");}

\NormalTok{  // Precompute |F(ω)| for f(t) = cos(t) + 0.5cos(3t) on [0,2π]}
\NormalTok{  // F(ω) = ∫₀\^{}\{2π\} f(t)e\^{}\{{-}iωt\}dt}
\NormalTok{  // Numerically:}
\NormalTok{  function computeF(omega) \{}
\NormalTok{    const N = 500;}
\NormalTok{    let re = 0, im = 0;}
\NormalTok{    for (let i = 0; i \textless{} N; i++) \{}
\NormalTok{      const t = tMin + (tMax{-}tMin)*i/(N{-}1);}
\NormalTok{      const f = Math.cos(t) + 0.5*Math.cos(3*t);}
\NormalTok{      re += f * Math.cos(omega * t);}
\NormalTok{      im {-}= f * Math.sin(omega * t);}
\NormalTok{    \}}
\NormalTok{    const dt = (tMax{-}tMin)/(N{-}1);}
\NormalTok{    return Math.sqrt(re*re + im*im) * dt;}
\NormalTok{  \}}

\NormalTok{  const omVals = d3.range(51).map(i =\textgreater{} i*0.1);}
\NormalTok{  const fMags = omVals.map(computeF);}
\NormalTok{  const maxMag = d3.max(fMags);}

\NormalTok{  const xBar = d3.scaleBand().domain(omVals.map(String)).range([barX0, barX0+barW]).padding(0.05);}
\NormalTok{  const yBar = d3.scaleLinear().domain([0, maxMag]).range([H{-}20, 20]);}

\NormalTok{  // Draw all bars (grey), highlight current}
\NormalTok{  const bars = omVals.map((om,i) =\textgreater{} \{}
\NormalTok{    const b = svg.append("rect")}
\NormalTok{      .attr("x", barX0 + (barW/(omVals.length)) * i)}
\NormalTok{      .attr("y", yBar(fMags[i]))}
\NormalTok{      .attr("width", Math.max(1, barW/omVals.length {-} 1))}
\NormalTok{      .attr("height", H{-}20 {-} yBar(fMags[i]))}
\NormalTok{      .attr("fill", GREY\_F);}
\NormalTok{    return b;}
\NormalTok{  \});}
\NormalTok{  svg.append("text").attr("x",barX0).attr("y",H{-}8).attr("font{-}size","8px").attr("fill",GREY).text("0");}
\NormalTok{  svg.append("text").attr("x",barX0+barW).attr("y",H{-}8).attr("font{-}size","8px").attr("fill",GREY).text("5");}
\NormalTok{  svg.append("text").attr("x",barX0+barW/2).attr("y",H{-}8).attr("text{-}anchor","middle")}
\NormalTok{    .attr("font{-}size","7.5px").attr("fill",GREY).text("ω");}

\NormalTok{  // Current ω indicator}
\NormalTok{  const barIndicator = svg.append("line").attr("stroke",AMBER).attr("stroke{-}width",2).attr("opacity",0.7);}

\NormalTok{  function update() \{}
\NormalTok{    const omega = +omSlider.value / 10;}
\NormalTok{    omDisp.textContent = omega.toFixed(1);}

\NormalTok{    // Wound path: (f(t)cos(ωt), {-}f(t)sin(ωt)) = (Re, Im) of f(t)e\^{}\{{-}iωt\}}
\NormalTok{    const N = 300;}
\NormalTok{    let reSum = 0, imSum = 0;}
\NormalTok{    const pts = [];}
\NormalTok{    for (let i = 0; i \textless{} N; i++) \{}
\NormalTok{      const t = tMin + (tMax{-}tMin)*i/(N{-}1);}
\NormalTok{      const f = Math.cos(t) + 0.5*Math.cos(3*t);}
\NormalTok{      const reW = f * Math.cos(omega * t);}
\NormalTok{      const imW = {-}f * Math.sin(omega * t);}
\NormalTok{      reSum += reW; imSum += imW;}
\NormalTok{      pts.push([circCx + reW * circR/1.6, circCy + imW * circR/1.6]);}
\NormalTok{    \}}
\NormalTok{    woundPath.datum(pts).attr("d", d3.line().x(d=\textgreater{}d[0]).y(d=\textgreater{}d[1]));}

\NormalTok{    // Centre of mass}
\NormalTok{    const comRe = reSum / N;}
\NormalTok{    const comIm = imSum / N;}
\NormalTok{    const comX = circCx + comRe * circR/1.6;}
\NormalTok{    const comY = circCy + comIm * circR/1.6;}
\NormalTok{    comDot.attr("cx", comX).attr("cy", comY);}
\NormalTok{    comLine.attr("x1",circCx).attr("y1",circCy).attr("x2",comX).attr("y2",comY);}
\NormalTok{    const mag = Math.sqrt(comRe*comRe + comIm*comIm);}
\NormalTok{    comLabel.attr("x",comX+6).attr("y",comY{-}4)}
\NormalTok{      .text(mag \textgreater{} 0.05 ? \textasciigrave{}|F| \textasciitilde{} $\{(mag*2*Math.PI).toFixed(1)\}\textasciigrave{} : "≈ 0");}

\NormalTok{    // Time bar (animate at current omega)}
\NormalTok{    const t0 = (Date.now() / 2000) \% 1 * (tMax{-}tMin) + tMin;}
\NormalTok{    const f0 = Math.cos(t0) + 0.5*Math.cos(3*t0);}
\NormalTok{    timeLine.attr("x1",xL(t0)).attr("x2",xL(t0)).attr("y1",16).attr("y2",H{-}20);}
\NormalTok{    timeDot.attr("cx",xL(t0)).attr("cy",yL(f0));}

\NormalTok{    // Bar highlight}
\NormalTok{    const idx = Math.round(omega*10);}
\NormalTok{    bars.forEach((b,i) =\textgreater{} b.attr("fill", i===idx ? AMBER : GREY\_F));}
\NormalTok{    if (idx \textless{} omVals.length) \{}
\NormalTok{      const bx = barX0 + (barW/omVals.length)*idx;}
\NormalTok{      barIndicator.attr("x1",bx).attr("x2",bx+barW/omVals.length).attr("y1",yBar(fMags[idx]){-}3).attr("y2",yBar(fMags[idx]){-}3);}
\NormalTok{    \}}
\NormalTok{  \}}

\NormalTok{  omSlider.oninput = update;}
\NormalTok{  update();}

\NormalTok{  const legend = document.createElement("div");}
\NormalTok{  legend.style.cssText = "display:flex; gap:1.2rem; font{-}size:0.78em; margin{-}top:0.2rem; flex{-}wrap:wrap;";}
\NormalTok{  [["\textasciitilde{} wound{-}up signal f(t)e\^{}\{−iωt\}", AMBER], ["• centre of mass = F(ω)/(2π)", RED]]}
\NormalTok{    .forEach(([lab, col]) =\textgreater{} \{}
\NormalTok{      const s = document.createElement("span"); s.style.color = col; s.textContent = lab;}
\NormalTok{      legend.appendChild(s);}
\NormalTok{    \});}
\NormalTok{  container.appendChild(svg.node());}
\NormalTok{  container.appendChild(legend);}
\NormalTok{  return container;}
\NormalTok{\}}
\end{Highlighting}
\end{Shaded}

\begin{center}\rule{0.5\linewidth}{0.5pt}\end{center}

\section{Key transform pairs}\label{key-transform-pairs}

\subsection{Rectangular pulse}\label{rectangular-pulse}

The \textbf{rectangular pulse} (or \textbf{gate function}) of width
\(2a\):

\[f(t) = \begin{cases} 1 & |t| < a \\ 0 & |t| > a \end{cases}\]

Its transform:

\[F(\omega) = \int_{-a}^{a} e^{-i\omega t}\,dt = \left[\frac{e^{-i\omega t}}{-i\omega}\right]_{-a}^{a} = \frac{e^{-i\omega a} - e^{i\omega a}}{-i\omega} = \frac{2\sin(\omega a)}{\omega} = 2a\,\operatorname{sinc}(\omega a)\]

where \(\operatorname{sinc}(x) = \sin(x)/x\).

\textbf{The time--bandwidth tradeoff.} The width of the main lobe of
\(|F(\omega)|\) is approximately \(2\pi/a\). A narrow pulse (\(a\)
small) has a wide spectrum. A wide pulse (\(a\) large) has a narrow
spectrum. This is the mathematical statement of the tradeoff the bat
exploits.

\subsection{Gaussian}\label{gaussian}

The Gaussian \(f(t) = e^{-\alpha t^2}\) (\(\alpha > 0\)) transforms to:

\[F(\omega) = \sqrt{\frac{\pi}{\alpha}}\,e^{-\omega^2/(4\alpha)}\]

also a Gaussian. The Gaussian is its own Fourier transform (up to
scaling). It minimises the time--bandwidth product, which is why it is
used in pulse shaping and wavelet theory.

\subsection{Dirac delta}\label{dirac-delta}

The Dirac delta is not an ordinary function --- it has no well-defined
value at any point, and you cannot draw its graph. It is defined
entirely by what it does inside an integral:
\(\int_{-\infty}^{\infty}\delta(t)\phi(t)\,dt = \phi(0)\) for any smooth
\(\phi\). Its transform:

\[\mathcal{F}\{\delta(t)\} = \int_{-\infty}^{\infty}\delta(t)\,e^{-i\omega t}\,dt = e^{0} = 1\]

The delta function has a flat (white) spectrum --- it contains all
frequencies equally. This is why a sharp impulse (the tap of a hammer on
a structure) excites all resonant frequencies simultaneously.

The inverse: \(\mathcal{F}\{1\} = 2\pi\delta(\omega)\) --- a constant
signal contains only zero frequency.

\begin{center}\rule{0.5\linewidth}{0.5pt}\end{center}

\section{Properties of the Fourier
transform}\label{properties-of-the-fourier-transform}

These properties make the Fourier transform useful in practice.

\textbf{Linearity.}
\(\mathcal{F}\{\alpha f + \beta g\} = \alpha F + \beta G\).

\textbf{Time shifting.}
\(\mathcal{F}\{f(t - t_0)\} = e^{-i\omega t_0}F(\omega)\). Shifting a
signal in time multiplies its spectrum by a complex exponential --- it
changes the phase but not the magnitude.

\textbf{Frequency shifting.}
\(\mathcal{F}\{e^{i\omega_0 t}f(t)\} = F(\omega - \omega_0)\).
Multiplying by a complex exponential shifts the spectrum. This is the
mathematical basis of \textbf{modulation}: shifting a baseband signal to
a carrier frequency. In engineering practice, modulation uses a real
cosine
\(\cos(\omega_0 t) = \frac{1}{2}(e^{i\omega_0 t} + e^{-i\omega_0 t})\),
which shifts the spectrum to both \(+\omega_0\) and \(-\omega_0\)
(creating upper and lower sidebands); the complex exponential form
isolates the single-sideband case.

\textbf{Scaling.}
\(\mathcal{F}\{f(at)\} = \frac{1}{|a|}F\!\left(\frac{\omega}{a}\right)\).
Compressing time by factor \(a\) expands the spectrum by the same factor
--- another form of the time--bandwidth tradeoff.

\textbf{Differentiation.} \(\mathcal{F}\{f'(t)\} = i\omega\,F(\omega)\).
Differentiation in time corresponds to multiplication by \(i\omega\) in
frequency. This turns differential equations into algebraic equations.

\textbf{Convolution theorem.}

\[\mathcal{F}\{f * g\}(\omega) = F(\omega)\cdot G(\omega)\]

where \((f*g)(t) = \int_{-\infty}^{\infty}f(\tau)g(t-\tau)\,d\tau\) is
the \textbf{convolution} of \(f\) and \(g\). Equivalently,
\(\mathcal{F}\{fg\} = \frac{1}{2\pi}F * G\).

\textbf{Parseval's theorem.}
\(\int_{-\infty}^{\infty}|f(t)|^2\,dt = \frac{1}{2\pi}\int_{-\infty}^{\infty}|F(\omega)|^2\,d\omega\).

\begin{center}\rule{0.5\linewidth}{0.5pt}\end{center}

\section{The convolution theorem in
engineering}\label{the-convolution-theorem-in-engineering}

The convolution theorem is the most practically important property. A
\textbf{linear time-invariant (LTI) system} is completely characterised
by its \textbf{impulse response} \(h(t)\): if you drive the system with
input \(f(t)\), the output is the convolution \(y(t) = (f * h)(t)\).

Computing convolution directly requires integrating over all past input
values. In the frequency domain, this becomes multiplication:

\[Y(\omega) = F(\omega)\cdot H(\omega)\]

where \(H(\omega) = \mathcal{F}\{h\}\) is the \textbf{transfer function}
of the system.

This is why filters are designed in the frequency domain. A low-pass
filter that passes frequencies \(|\omega| < \omega_c\) and blocks higher
frequencies has \(H(\omega) = 1\) for \(|\omega| < \omega_c\) and
\(H(\omega) = 0\) otherwise. Computing \(H(\omega) \cdot F(\omega)\) is
multiplication of two functions. Computing the equivalent convolution
would mean implementing an integral.

In discrete time, the same structure holds: the discrete Fourier
transform converts discrete convolution to pointwise multiplication. The
FFT computes this in \(O(N\log N)\) time, which is why all modern
digital signal processing systems work this way.

\begin{Shaded}
\begin{Highlighting}[]
\NormalTok{//| echo: false}
\NormalTok{\{}
\NormalTok{  // ── VIZ 2 Rewrite: Transform pair explorer ───────────────────────────────}
\NormalTok{  // Panel A: time{-}domain f(t)  Panel B: |F(ω)| magnitude  Panel C: phase ∠F(ω)}
\NormalTok{  // Wave types: rectangular pulse | Gaussian | Exponential decay}
\NormalTok{  const d3 = await require("d3@7");}

\NormalTok{  const AMBER = "\#f59e0b";}
\NormalTok{  const BLUE  = "\#6b90c4";}
\NormalTok{  const GREY  = "\#9ca3af";}
\NormalTok{  const GREY\_F = "\#e5e7eb";}

\NormalTok{  const W = 560, H = 280;}
\NormalTok{  const PL = 44, PR = 12, PT = 22, PB = 16;}
\NormalTok{  const IW = W {-} PL {-} PR;}
\NormalTok{  const HA = 90, HB = 90, HC = 50, gap = 12;}

\NormalTok{  const container = document.createElement("div");}
\NormalTok{  container.style.cssText = "font{-}family:inherit; max{-}width:620px;";}

\NormalTok{  // Controls}
\NormalTok{  const ctrlRow = document.createElement("div");}
\NormalTok{  ctrlRow.style.cssText = "display:flex; gap:1.2rem; align{-}items:center; margin{-}bottom:0.4rem; flex{-}wrap:wrap;";}

\NormalTok{  // Wave type selector}
\NormalTok{  const waveTypes = [["rect","Rectangular pulse"],["gauss","Gaussian"],["exp","Exp. decay e\^{}\{−|t|/τ\}"]];}
\NormalTok{  let currentWT = "rect";}
\NormalTok{  const radioGrp = document.createElement("div");}
\NormalTok{  radioGrp.style.cssText = "display:flex; gap:0.6rem; font{-}size:0.85em; flex{-}wrap:wrap;";}
\NormalTok{  waveTypes.forEach(([val,lab]) =\textgreater{} \{}
\NormalTok{    const lbl = document.createElement("label");}
\NormalTok{    lbl.style.cssText = "display:flex; align{-}items:center; gap:0.25rem; cursor:pointer;";}
\NormalTok{    const r = document.createElement("input");}
\NormalTok{    r.type = "radio"; r.name = "ftViz2"; r.value = val; r.checked = val==="rect";}
\NormalTok{    r.style.accentColor = AMBER;}
\NormalTok{    r.onchange = () =\textgreater{} \{ currentWT = val; updateParamLabel(); draw(); \};}
\NormalTok{    lbl.appendChild(r); lbl.appendChild(document.createTextNode(lab));}
\NormalTok{    radioGrp.appendChild(lbl);}
\NormalTok{  \});}
\NormalTok{  ctrlRow.appendChild(radioGrp);}

\NormalTok{  // Parameter slider (a / σ / τ)}
\NormalTok{  const paramRow = document.createElement("div");}
\NormalTok{  paramRow.style.cssText = "display:flex; gap:0.5rem; align{-}items:center; font{-}size:0.85em;";}
\NormalTok{  const paramLbl = document.createElement("span");}
\NormalTok{  paramLbl.style.cssText = "color:\#374151; min{-}width:4ch;";}
\NormalTok{  paramLbl.textContent = "a =";}
\NormalTok{  const paramDisp = document.createElement("strong");}
\NormalTok{  paramDisp.style.cssText = "color:" + AMBER + "; min{-}width:3ch;";}
\NormalTok{  paramDisp.textContent = "1.0";}
\NormalTok{  const paramSlider = document.createElement("input");}
\NormalTok{  paramSlider.type = "range"; paramSlider.min = 2; paramSlider.max = 30; paramSlider.value = 10;}
\NormalTok{  paramSlider.style.cssText = "width:160px; accent{-}color:" + AMBER + ";";}
\NormalTok{  paramRow.appendChild(paramLbl); paramRow.appendChild(paramSlider); paramRow.appendChild(paramDisp);}
\NormalTok{  ctrlRow.appendChild(paramRow);}

\NormalTok{  // Time{-}shift toggle}
\NormalTok{  const shiftLbl = document.createElement("label");}
\NormalTok{  shiftLbl.style.cssText = "display:flex; align{-}items:center; gap:0.3rem; font{-}size:0.82em; cursor:pointer;";}
\NormalTok{  const shiftCheck = document.createElement("input");}
\NormalTok{  shiftCheck.type = "checkbox";}
\NormalTok{  shiftCheck.onchange = draw;}
\NormalTok{  shiftLbl.appendChild(shiftCheck);}
\NormalTok{  shiftLbl.appendChild(document.createTextNode("Shift by t₀=1"));}
\NormalTok{  ctrlRow.appendChild(shiftLbl);}
\NormalTok{  container.appendChild(ctrlRow);}

\NormalTok{  function updateParamLabel() \{}
\NormalTok{    const lbl = currentWT==="rect"?"a":currentWT==="gauss"?"σ":"τ";}
\NormalTok{    paramLbl.textContent = lbl + " =";}
\NormalTok{  \}}

\NormalTok{  const svg = d3.create("svg")}
\NormalTok{    .attr("viewBox", \textasciigrave{}0 0 $\{W\} $\{H\}\textasciigrave{})}
\NormalTok{    .attr("width","100\%").style("display","block");}

\NormalTok{  const gA = svg.append("g").attr("transform", \textasciigrave{}translate($\{PL\},$\{PT\})\textasciigrave{});}
\NormalTok{  const gB = svg.append("g").attr("transform", \textasciigrave{}translate($\{PL\},$\{PT+HA+gap\})\textasciigrave{});}
\NormalTok{  const gC = svg.append("g").attr("transform", \textasciigrave{}translate($\{PL\},$\{PT+HA+gap+HB+gap\})\textasciigrave{});}

\NormalTok{  // Shared ω range — rescale based on current a}
\NormalTok{  function omegaRange(param) \{ return 4 * Math.PI / param; \}}

\NormalTok{  // ── Wave functions ───────────────────────────────────────────────────────}
\NormalTok{  function fTime(t, param, wt, shift) \{}
\NormalTok{    const t\_ = shift ? t {-} 1 : t;}
\NormalTok{    if (wt==="rect") return Math.abs(t\_) \textless{} param ? 1 : 0;}
\NormalTok{    if (wt==="gauss") return Math.exp({-}t\_*t\_ / (2*param*param));}
\NormalTok{    return Math.exp({-}Math.abs(t\_)/param);}
\NormalTok{  \}}

\NormalTok{  // Fourier transforms (analytical)}
\NormalTok{  function FRect(omega, a, shift) \{}
\NormalTok{    const mag = Math.abs(omega) \textless{} 1e{-}8 ? 2*a : 2*Math.sin(omega*a)/omega;}
\NormalTok{    const phase = shift ? {-}omega * 1 : 0;}
\NormalTok{    return \{ mag: Math.abs(mag), phase: phase + (mag \textless{} 0 ? Math.PI : 0) \};}
\NormalTok{  \}}
\NormalTok{  function FGauss(omega, sigma, shift) \{}
\NormalTok{    const mag = sigma * Math.sqrt(2*Math.PI) * Math.exp({-}omega*omega*sigma*sigma/2);}
\NormalTok{    const phase = shift ? {-}omega * 1 : 0;}
\NormalTok{    return \{ mag, phase \};}
\NormalTok{  \}}
\NormalTok{  function FExp(omega, tau, shift) \{}
\NormalTok{    const mag = 2*tau / (1 + omega*omega*tau*tau);}
\NormalTok{    const phase = shift ? {-}omega * 1 : 0;}
\NormalTok{    return \{ mag, phase \};}
\NormalTok{  \}}

\NormalTok{  function computeF(omega, param, wt, shift) \{}
\NormalTok{    if (wt==="rect") return FRect(omega, param, shift);}
\NormalTok{    if (wt==="gauss") return FGauss(omega, param, shift);}
\NormalTok{    return FExp(omega, param, shift);}
\NormalTok{  \}}

\NormalTok{  // Axis setup helpers}
\NormalTok{  function makeAxes(g, xScale, yScale, height, nTicksY, title) \{}
\NormalTok{    g.append("text").attr("x",IW/2).attr("y",{-}6).attr("text{-}anchor","middle")}
\NormalTok{      .attr("font{-}size","10px").attr("fill","\#374151").text(title);}
\NormalTok{    g.append("line").attr("x1",0).attr("x2",IW).attr("y1",yScale(0)).attr("y2",yScale(0))}
\NormalTok{      .attr("stroke",GREY\_F).attr("stroke{-}width",0.8);}
\NormalTok{    g.append("g").call(d3.axisLeft(yScale).ticks(nTicksY))}
\NormalTok{      .call(gg =\textgreater{} \{ gg.select(".domain").remove(); gg.selectAll("text").style("font{-}size","9px").attr("fill",GREY); gg.selectAll("line").attr("stroke",GREY\_F); \});}
\NormalTok{    g.append("g").call(d3.axisBottom(xScale).ticks(7))}
\NormalTok{      .attr("transform",\textasciigrave{}translate(0,$\{height\})\textasciigrave{})}
\NormalTok{      .call(gg =\textgreater{} \{ gg.select(".domain").attr("stroke",GREY\_F); gg.selectAll("text").style("font{-}size","9px").attr("fill",GREY); gg.selectAll("line").attr("stroke",GREY\_F); \});}
\NormalTok{  \}}

\NormalTok{  // Panel A: time domain}
\NormalTok{  const tMin = {-}6, tMax = 6;}
\NormalTok{  const xA = d3.scaleLinear().domain([tMin, tMax]).range([0, IW]);}
\NormalTok{  const yA = d3.scaleLinear().domain([{-}0.15, 1.2]).range([HA, 0]);}
\NormalTok{  makeAxes(gA, xA, yA, HA, 3, "Time domain f(t)");}
\NormalTok{  const timePath = gA.append("path").attr("fill","none").attr("stroke",AMBER).attr("stroke{-}width",2.5);}
\NormalTok{  const widthBracket = gA.append("g");}
\NormalTok{  const bracketLine = widthBracket.append("line").attr("stroke",GREY).attr("stroke{-}width",1).attr("stroke{-}dasharray","3,2");}
\NormalTok{  const bracketLabel = widthBracket.append("text").attr("font{-}size","9px").attr("fill",GREY).attr("text{-}anchor","middle");}

\NormalTok{  // Panel B: magnitude}
\NormalTok{  const xB = d3.scaleLinear().domain([{-}1, 1]).range([0, IW]);}
\NormalTok{  const yB = d3.scaleLinear().domain([0, 3]).range([HB, 0]);}
\NormalTok{  makeAxes(gB, xB, yB, HB, 3, "|F(ω)| — frequency magnitude");}
\NormalTok{  const magPath = gB.append("path").attr("fill","none").attr("stroke",AMBER).attr("stroke{-}width",2.5);}
\NormalTok{  const zeroMarks = [gB.append("line").attr("stroke","\#6b7280").attr("stroke{-}width",1.5).attr("stroke{-}dasharray","3,2"),}
\NormalTok{                     gB.append("line").attr("stroke","\#6b7280").attr("stroke{-}width",1.5).attr("stroke{-}dasharray","3,2")];}
\NormalTok{  const zeroLabels = [gB.append("text").attr("font{-}size","8.5px").attr("fill","\#6b7280").attr("text{-}anchor","middle"),}
\NormalTok{                      gB.append("text").attr("font{-}size","8.5px").attr("fill","\#6b7280").attr("text{-}anchor","middle")];}
\NormalTok{  const bwArrow = gB.append("text").attr("font{-}size","9px").attr("fill",BLUE).attr("font{-}style","italic");}

\NormalTok{  // Panel C: phase}
\NormalTok{  const xC = d3.scaleLinear().domain([{-}1, 1]).range([0, IW]);}
\NormalTok{  const yC = d3.scaleLinear().domain([{-}Math.PI{-}0.3, Math.PI+0.3]).range([HC, 0]);}
\NormalTok{  makeAxes(gC, xC, yC, HC, 2, "Phase ∠F(ω)");}
\NormalTok{  const phasePath = gC.append("path").attr("fill","none").attr("stroke",BLUE).attr("stroke{-}width",1.8);}
\NormalTok{  const phaseLabel = gC.append("text").attr("font{-}size","9px").attr("fill",BLUE).attr("x",4).attr("y",10);}

\NormalTok{  function draw() \{}
\NormalTok{    const param = +paramSlider.value / 10;}
\NormalTok{    paramDisp.textContent = param.toFixed(1);}
\NormalTok{    const shift = shiftCheck.checked;}
\NormalTok{    const omMax = omegaRange(param);}

\NormalTok{    // Update x{-}axes domain for B and C}
\NormalTok{    xB.domain([{-}omMax, omMax]);}
\NormalTok{    xC.domain([{-}omMax, omMax]);}

\NormalTok{    // Redraw axes ticks}
\NormalTok{    gB.selectAll(".x{-}axis").remove();}
\NormalTok{    gB.append("g").attr("class","x{-}axis").call(d3.axisBottom(xB).ticks(7).tickFormat(v =\textgreater{} v.toFixed(1) === "0.0" ? "0" : v.toFixed(1)))}
\NormalTok{      .attr("transform",\textasciigrave{}translate(0,$\{HB\})\textasciigrave{})}
\NormalTok{      .call(gg =\textgreater{} \{ gg.select(".domain").attr("stroke",GREY\_F); gg.selectAll("text").style("font{-}size","8.5px").attr("fill",GREY); gg.selectAll("line").attr("stroke",GREY\_F); \});}
\NormalTok{    gC.selectAll(".x{-}axis").remove();}
\NormalTok{    gC.append("g").attr("class","x{-}axis").call(d3.axisBottom(xC).ticks(7).tickFormat(v =\textgreater{} v.toFixed(1) === "0.0" ? "ω=0" : v.toFixed(1)))}
\NormalTok{      .attr("transform",\textasciigrave{}translate(0,$\{HC\})\textasciigrave{})}
\NormalTok{      .call(gg =\textgreater{} \{ gg.select(".domain").attr("stroke",GREY\_F); gg.selectAll("text").style("font{-}size","8.5px").attr("fill",GREY); gg.selectAll("line").attr("stroke",GREY\_F); \});}

\NormalTok{    // Panel A: time domain}
\NormalTok{    const tPts = d3.range(400).map(i =\textgreater{} \{}
\NormalTok{      const t = tMin + (tMax{-}tMin)*i/399;}
\NormalTok{      return [xA(t), yA(fTime(t, param, currentWT, shift))];}
\NormalTok{    \});}
\NormalTok{    timePath.datum(tPts).attr("d", d3.line().x(d=\textgreater{}d[0]).y(d=\textgreater{}d[1]));}

\NormalTok{    // Width bracket}
\NormalTok{    const bLab = currentWT==="rect"?"2a":currentWT==="gauss"?"2σ":"τ";}
\NormalTok{    const bWidth = currentWT==="rect" ? 2*param : currentWT==="gauss" ? 2*param : param;}
\NormalTok{    bracketLine.attr("x1",xA({-}bWidth/2)).attr("x2",xA(bWidth/2)).attr("y1",yA({-}0.08)).attr("y2",yA({-}0.08));}
\NormalTok{    bracketLabel.attr("x",xA(0)).attr("y",yA({-}0.12)).text(bLab + " = " + bWidth.toFixed(1));}

\NormalTok{    // Panel B \& C: spectrum}
\NormalTok{    const omPts = d3.range(600).map(i =\textgreater{} \{}
\NormalTok{      const om = {-}omMax + 2*omMax*i/599;}
\NormalTok{      return om;}
\NormalTok{    \});}
\NormalTok{    const maxMag = d3.max(omPts, om =\textgreater{} computeF(om, param, currentWT, false).mag);}
\NormalTok{    yB.domain([0, maxMag * 1.15]);}

\NormalTok{    magPath.datum(omPts.map(om =\textgreater{} \{}
\NormalTok{      const \{mag\} = computeF(om, param, currentWT, shift);}
\NormalTok{      return [xB(om), yB(mag)];}
\NormalTok{    \})).attr("d", d3.line().x(d=\textgreater{}d[0]).y(d=\textgreater{}d[1]));}

\NormalTok{    phasePath.datum(omPts.map(om =\textgreater{} \{}
\NormalTok{      const \{phase\} = computeF(om, param, currentWT, shift);}
\NormalTok{      return [xC(om), yC(phase)];}
\NormalTok{    \})).attr("d", d3.line().x(d=\textgreater{}d[0]).y(d=\textgreater{}d[1]));}

\NormalTok{    phaseLabel.text(shift ? "Phase ramp = −ωt₀ (time{-}shift effect)" : "Phase = 0 (symmetric, real signal)");}

\NormalTok{    // Zero{-}crossing markers (rect only)}
\NormalTok{    if (currentWT === "rect" \&\& !shift) \{}
\NormalTok{      const z1 = Math.PI / param;}
\NormalTok{      zeroMarks[0].attr("x1",xB(z1)).attr("x2",xB(z1)).attr("y1",0).attr("y2",HB).style("display","block");}
\NormalTok{      zeroMarks[1].attr("x1",xB({-}z1)).attr("x2",xB({-}z1)).attr("y1",0).attr("y2",HB).style("display","block");}
\NormalTok{      zeroLabels[0].attr("x",xB(z1)).attr("y",12).text("π/a").style("display","block");}
\NormalTok{      zeroLabels[1].attr("x",xB({-}z1)).attr("y",12).text("−π/a").style("display","block");}
\NormalTok{      const bwDir = param \textless{} 1.5 ? "narrower pulse → wider spectrum" : "wider pulse → narrower spectrum";}
\NormalTok{      bwArrow.attr("x",IW/2).attr("y",HB{-}6).attr("text{-}anchor","middle").text(bwDir);}
\NormalTok{    \} else \{}
\NormalTok{      zeroMarks.forEach(m =\textgreater{} m.style("display","none"));}
\NormalTok{      zeroLabels.forEach(l =\textgreater{} l.style("display","none"));}
\NormalTok{      bwArrow.text(currentWT==="gauss" ? "Gaussian: width 2σ ↔ bandwidth 2/σ" : "");}
\NormalTok{    \}}
\NormalTok{  \}}

\NormalTok{  paramSlider.oninput = draw;}
\NormalTok{  draw();}

\NormalTok{  const legend = document.createElement("div");}
\NormalTok{  legend.style.cssText = "display:flex; gap:1.2rem; font{-}size:0.78em; margin{-}top:0.2rem; flex{-}wrap:wrap;";}
\NormalTok{  [["─── f(t) (time domain)", AMBER], ["─── ∠F(ω) (phase)", BLUE],}
\NormalTok{   ["| Zero crossings at ±π/a", "\#6b7280"]]}
\NormalTok{    .forEach(([lab, col]) =\textgreater{} \{}
\NormalTok{      const s = document.createElement("span"); s.style.color = col; s.textContent = lab;}
\NormalTok{      legend.appendChild(s);}
\NormalTok{    \});}
\NormalTok{  container.appendChild(svg.node());}
\NormalTok{  container.appendChild(legend);}
\NormalTok{  return container;}
\NormalTok{\}}
\end{Highlighting}
\end{Shaded}

\begin{center}\rule{0.5\linewidth}{0.5pt}\end{center}

\section{Solving differential equations with Fourier
transforms}\label{solving-differential-equations-with-fourier-transforms}

The differentiation property converts a linear ODE with constant
coefficients into an algebraic equation in the frequency domain.

\textbf{Example.} Solve \(y' + 2y = e^{-|t|}\) for \(y(t)\).

Take the Fourier transform of both sides:

\[i\omega\,Y(\omega) + 2Y(\omega) = F(\omega)\]

\[(i\omega + 2)Y(\omega) = F(\omega)\]

\[Y(\omega) = \frac{F(\omega)}{i\omega + 2}\]

Now we need
\(F(\omega) = \mathcal{F}\{e^{-|t|}\} = \int_{-\infty}^{\infty}e^{-|t|}e^{-i\omega t}\,dt = \frac{2}{1+\omega^2}\).

So \(Y(\omega) = \frac{2}{(1+\omega^2)(i\omega+2)}\).

Write this in partial fractions:

\[\frac{2}{(1+\omega^2)(i\omega+2)} = \frac{A}{i\omega+2} + \frac{B\omega + C}{1+\omega^2}\]

Matching numerators gives \(A = 2/5\), \(B = -2/5\), \(C = 4/5\), so:

\[Y(\omega) = \frac{2/5}{i\omega+2} + \frac{(-2\omega + 4)/5}{1+\omega^2}\]

Each term matches a known transform pair:
\(\mathcal{F}^{-1}\{1/(i\omega+2)\} = e^{-2t}\) for \(t > 0\) (causal
exponential);
\(\mathcal{F}^{-1}\{1/(1+\omega^2)\} = \frac{1}{2}e^{-|t|}\); and
\(\mathcal{F}^{-1}\{\omega/(1+\omega^2)\}\) gives a derivative term.
Inverting term by term:

\[y(t) = \frac{1}{5}\left(2e^{-2t} - e^{-|t|} + \operatorname{sgn}(t)e^{-|t|}\right)\]

The Fourier transform has reduced the differential equation to algebra.

\begin{center}\rule{0.5\linewidth}{0.5pt}\end{center}

\section{Connection to the Laplace
transform}\label{connection-to-the-laplace-transform}

If you have met the Laplace transform in an ODE course, this section
shows how it connects to the Fourier transform. If not, you can skip it
--- nothing later in this chapter depends on it.

The Laplace transform
\(\mathcal{L}\{f\}(s) = \int_0^{\infty}f(t)e^{-st}\,dt\) is a
\textbf{one-sided} transform: it only sees \(t \geq 0\) (causal
signals). Replacing \(s = i\omega\) recovers the Fourier transform of
the causal part of \(f\).

The Laplace transform is preferred for initial-value problems (ODEs with
\(t \geq 0\) and given initial conditions) because it handles them
directly via \(\mathcal{L}\{f'\} = sF(s) - f(0)\). The Fourier transform
is preferred for steady-state and bilateral analysis (signals for all
\(t\)).

In circuit analysis: the transfer function \(H(s)\) in Laplace domain
becomes \(H(i\omega)\), the Fourier domain transfer function, when you
restrict to purely sinusoidal steady-state inputs.

\begin{center}\rule{0.5\linewidth}{0.5pt}\end{center}

\section{Exercises}\label{exercises-35}

\begin{Shaded}
\begin{Highlighting}[]
\NormalTok{//| echo: false}
\NormalTok{katex = require("katex@0.16")}
\end{Highlighting}
\end{Shaded}

\begin{Shaded}
\begin{Highlighting}[]
\NormalTok{//| echo: false}
\NormalTok{function makeStepperHTML(steps) \{}
\NormalTok{  let current = 0;}
\NormalTok{  const totalSteps = steps.length;}
\NormalTok{  const container = document.createElement("div");}
\NormalTok{  container.style.cssText = "border:1px solid \#e5e7eb; border{-}radius:8px; padding:1rem 1.25rem; margin:0.75rem 0; font{-}family:inherit;";}
\NormalTok{  const stepsDiv = document.createElement("div");}
\NormalTok{  stepsDiv.style.marginBottom = "0.75rem";}

\NormalTok{  function renderSteps() \{}
\NormalTok{    stepsDiv.innerHTML = "";}
\NormalTok{    for (let i = 0; i \textless{} current; i++) \{}
\NormalTok{      const s = steps[i];}
\NormalTok{      const row = document.createElement("div");}
\NormalTok{      row.style.cssText = "display:grid; grid{-}template{-}columns:220px 1fr; gap:0.35rem 1rem; align{-}items:baseline; padding:0.35rem 0; border{-}top:1px solid \#f3f4f6;";}
\NormalTok{      const opCell = document.createElement("span");}
\NormalTok{      opCell.style.cssText = "font{-}size:0.85em; color:\#6b7280; font{-}style:italic;";}
\NormalTok{      opCell.textContent = s.op;}
\NormalTok{      const eqCell = document.createElement("span");}
\NormalTok{      eqCell.innerHTML = katex.renderToString(s.eq, \{ throwOnError: false, displayMode: false \});}
\NormalTok{      row.appendChild(opCell);}
\NormalTok{      row.appendChild(eqCell);}
\NormalTok{      if (s.note) \{}
\NormalTok{        const noteRow = document.createElement("div");}
\NormalTok{        noteRow.style.cssText = "grid{-}column:2; font{-}size:0.82em; color:\#6b7280; padding{-}bottom:0.2rem;";}
\NormalTok{        noteRow.textContent = s.note;}
\NormalTok{        row.appendChild(noteRow);}
\NormalTok{      \}}
\NormalTok{      stepsDiv.appendChild(row);}
\NormalTok{    \}}
\NormalTok{    if (current === totalSteps) \{}
\NormalTok{      const done = document.createElement("div");}
\NormalTok{      done.style.cssText = "margin{-}top:0.5rem; font{-}size:0.9em; color:\#059669; font{-}weight:500;";}
\NormalTok{      done.textContent = "Solution complete.";}
\NormalTok{      stepsDiv.appendChild(done);}
\NormalTok{    \}}
\NormalTok{  \}}

\NormalTok{  const controls = document.createElement("div");}
\NormalTok{  controls.style.cssText = "display:flex; gap:0.5rem; align{-}items:center;";}
\NormalTok{  const nextBtn = document.createElement("button");}
\NormalTok{  nextBtn.textContent = "Next step →";}
\NormalTok{  nextBtn.style.cssText = "padding:0.35rem 0.85rem; border:1px solid \#d1d5db; border{-}radius:4px; background:\#fff; cursor:pointer; font{-}size:0.9em;";}
\NormalTok{  const resetBtn = document.createElement("button");}
\NormalTok{  resetBtn.textContent = "Reset";}
\NormalTok{  resetBtn.style.cssText = "padding:0.35rem 0.75rem; border:1px solid \#e5e7eb; border{-}radius:4px; background:\#f9fafb; cursor:pointer; font{-}size:0.9em; color:\#6b7280;";}
\NormalTok{  const counter = document.createElement("span");}
\NormalTok{  counter.style.cssText = "font{-}size:0.82em; color:\#9ca3af; margin{-}left:0.25rem;";}

\NormalTok{  function updateButtons() \{}
\NormalTok{    nextBtn.disabled = current === totalSteps;}
\NormalTok{    nextBtn.style.opacity = current === totalSteps ? "0.4" : "1";}
\NormalTok{    counter.textContent = current === 0 ? "Click to reveal steps" : \textasciigrave{}Step $\{current\} of $\{totalSteps\}\textasciigrave{};}
\NormalTok{  \}}

\NormalTok{  nextBtn.onclick = () =\textgreater{} \{ if (current \textless{} totalSteps) \{ current++; renderSteps(); updateButtons(); \} \};}
\NormalTok{  resetBtn.onclick = () =\textgreater{} \{ current = 0; renderSteps(); updateButtons(); \};}
\NormalTok{  controls.appendChild(nextBtn);}
\NormalTok{  controls.appendChild(resetBtn);}
\NormalTok{  controls.appendChild(counter);}
\NormalTok{  container.appendChild(stepsDiv);}
\NormalTok{  container.appendChild(controls);}
\NormalTok{  renderSteps();}
\NormalTok{  updateButtons();}
\NormalTok{  return container;}
\NormalTok{\}}
\end{Highlighting}
\end{Shaded}

\begin{center}\rule{0.5\linewidth}{0.5pt}\end{center}

\subsection{Exercise 1: Transform of a rectangular
pulse}\label{exercise-1-transform-of-a-rectangular-pulse}

Compute the Fourier transform of
\(f(t) = \begin{cases}1 & |t| \leq 1 \\ 0 & |t| > 1\end{cases}\).
Identify the first zero of \(|F(\omega)|\).

\begin{Shaded}
\begin{Highlighting}[]
\NormalTok{//| echo: false}
\NormalTok{makeStepperHTML([}
\NormalTok{  \{}
\NormalTok{    op: "Set up the integral",}
\NormalTok{    eq: "F(\textbackslash{}\textbackslash{}omega) = \textbackslash{}\textbackslash{}int\_\{{-}\textbackslash{}\textbackslash{}infty\}\^{}\{\textbackslash{}\textbackslash{}infty\} f(t)\textbackslash{}\textbackslash{},e\^{}\{{-}i\textbackslash{}\textbackslash{}omega t\}\textbackslash{}\textbackslash{},dt = \textbackslash{}\textbackslash{}int\_\{{-}1\}\^{}\{1\} e\^{}\{{-}i\textbackslash{}\textbackslash{}omega t\}\textbackslash{}\textbackslash{},dt"}
\NormalTok{  \},}
\NormalTok{  \{}
\NormalTok{    op: "Evaluate (ω ≠ 0)",}
\NormalTok{    eq: "F(\textbackslash{}\textbackslash{}omega) = \textbackslash{}\textbackslash{}left[\textbackslash{}\textbackslash{}frac\{e\^{}\{{-}i\textbackslash{}\textbackslash{}omega t\}\}\{{-}i\textbackslash{}\textbackslash{}omega\}\textbackslash{}\textbackslash{}right]\_\{{-}1\}\^{}\{1\} = \textbackslash{}\textbackslash{}frac\{e\^{}\{{-}i\textbackslash{}\textbackslash{}omega\} {-} e\^{}\{i\textbackslash{}\textbackslash{}omega\}\}\{{-}i\textbackslash{}\textbackslash{}omega\} = \textbackslash{}\textbackslash{}frac\{2\textbackslash{}\textbackslash{}sin\textbackslash{}\textbackslash{}omega\}\{\textbackslash{}\textbackslash{}omega\}",}
\NormalTok{    note: "Using e\^{}\{iθ\} − e\^{}\{−iθ\} = 2i sin θ."}
\NormalTok{  \},}
\NormalTok{  \{}
\NormalTok{    op: "Check ω = 0 by L\textquotesingle{}Hôpital or direct integration",}
\NormalTok{    eq: "F(0) = \textbackslash{}\textbackslash{}int\_\{{-}1\}\^{}\{1\}1\textbackslash{}\textbackslash{},dt = 2 \textbackslash{}\textbackslash{}qquad (\textbackslash{}\textbackslash{}text\{consistent with \} \textbackslash{}\textbackslash{}lim\_\{\textbackslash{}\textbackslash{}omega\textbackslash{}\textbackslash{}to 0\}\textbackslash{}\textbackslash{}frac\{2\textbackslash{}\textbackslash{}sin\textbackslash{}\textbackslash{}omega\}\{\textbackslash{}\textbackslash{}omega\} = 2)"}
\NormalTok{  \},}
\NormalTok{  \{}
\NormalTok{    op: "Write in sinc form",}
\NormalTok{    eq: "F(\textbackslash{}\textbackslash{}omega) = 2\textbackslash{}\textbackslash{}operatorname\{sinc\}(\textbackslash{}\textbackslash{}omega) \textbackslash{}\textbackslash{}quad \textbackslash{}\textbackslash{}text\{where \} \textbackslash{}\textbackslash{}operatorname\{sinc\}(\textbackslash{}\textbackslash{}omega) = \textbackslash{}\textbackslash{}frac\{\textbackslash{}\textbackslash{}sin\textbackslash{}\textbackslash{}omega\}\{\textbackslash{}\textbackslash{}omega\}"}
\NormalTok{  \},}
\NormalTok{  \{}
\NormalTok{    op: "Locate first zero",}
\NormalTok{    eq: "F(\textbackslash{}\textbackslash{}omega) = 0 \textbackslash{}\textbackslash{}iff \textbackslash{}\textbackslash{}sin\textbackslash{}\textbackslash{}omega = 0 \textbackslash{}\textbackslash{}iff \textbackslash{}\textbackslash{}omega = \textbackslash{}\textbackslash{}pm\textbackslash{}\textbackslash{}pi",}
\NormalTok{    note: "The bandwidth (distance from DC to first null) is π rad/s for a pulse of half{-}width a = 1. In general: bandwidth = π/a."}
\NormalTok{  \}}
\NormalTok{])}
\end{Highlighting}
\end{Shaded}

\begin{center}\rule{0.5\linewidth}{0.5pt}\end{center}

\subsection{Exercise 2: Time-shifting
theorem}\label{exercise-2-time-shifting-theorem}

Given that \(\mathcal{F}\{e^{-|t|}\} = \frac{2}{1+\omega^2}\), find
\(\mathcal{F}\{e^{-|t-3|}\}\) and \(\mathcal{F}\{e^{-|t|}e^{2it}\}\).

\begin{Shaded}
\begin{Highlighting}[]
\NormalTok{//| echo: false}
\NormalTok{makeStepperHTML([}
\NormalTok{  \{}
\NormalTok{    op: "Recall the shifting theorems",}
\NormalTok{    eq: "\textbackslash{}\textbackslash{}mathcal\{F\}\textbackslash{}\textbackslash{}\{f(t{-}t\_0)\textbackslash{}\textbackslash{}\} = e\^{}\{{-}i\textbackslash{}\textbackslash{}omega t\_0\}F(\textbackslash{}\textbackslash{}omega) \textbackslash{}\textbackslash{}qquad \textbackslash{}\textbackslash{}mathcal\{F\}\textbackslash{}\textbackslash{}\{e\^{}\{i\textbackslash{}\textbackslash{}omega\_0 t\}f(t)\textbackslash{}\textbackslash{}\} = F(\textbackslash{}\textbackslash{}omega {-} \textbackslash{}\textbackslash{}omega\_0)"}
\NormalTok{  \},}
\NormalTok{  \{}
\NormalTok{    op: "Apply time{-}shift: f(t) = e\^{}\{−|t|\}, t₀ = 3",}
\NormalTok{    eq: "\textbackslash{}\textbackslash{}mathcal\{F\}\textbackslash{}\textbackslash{}\{e\^{}\{{-}|t{-}3|\}\textbackslash{}\textbackslash{}\} = e\^{}\{{-}3i\textbackslash{}\textbackslash{}omega\}\textbackslash{}\textbackslash{}cdot\textbackslash{}\textbackslash{}frac\{2\}\{1+\textbackslash{}\textbackslash{}omega\^{}2\}",}
\NormalTok{    note: "The magnitude |F(ω)| = 2/(1+ω²) is unchanged — shifting in time only rotates the phase."}
\NormalTok{  \},}
\NormalTok{  \{}
\NormalTok{    op: "Apply frequency{-}shift: f(t) = e\^{}\{−|t|\}, ω₀ = 2",}
\NormalTok{    eq: "\textbackslash{}\textbackslash{}mathcal\{F\}\textbackslash{}\textbackslash{}\{e\^{}\{2it\}\textbackslash{}\textbackslash{},e\^{}\{{-}|t|\}\textbackslash{}\textbackslash{}\} = \textbackslash{}\textbackslash{}frac\{2\}\{1+(\textbackslash{}\textbackslash{}omega{-}2)\^{}2\}",}
\NormalTok{    note: "The spectrum is the same shape, centred at ω = 2 instead of ω = 0. This is modulation: multiplying by e\^{}\{2it\} shifts the signal to carrier frequency 2."}
\NormalTok{  \},}
\NormalTok{  \{}
\NormalTok{    op: "Verify: magnitude of each result at ω = 0",}
\NormalTok{    eq: "\textbackslash{}\textbackslash{}left|e\^{}\{{-}3i\textbackslash{}\textbackslash{}cdot 0\}\textbackslash{}\textbackslash{}cdot\textbackslash{}\textbackslash{}frac\{2\}\{1\}\textbackslash{}\textbackslash{}right| = 2 \textbackslash{}\textbackslash{}qquad \textbackslash{}\textbackslash{}left|\textbackslash{}\textbackslash{}frac\{2\}\{1+4\}\textbackslash{}\textbackslash{}right| = 0.4",}
\NormalTok{    note: "The modulated signal has less power at DC — the energy has been moved to ω = 2."}
\NormalTok{  \}}
\NormalTok{])}
\end{Highlighting}
\end{Shaded}

\begin{center}\rule{0.5\linewidth}{0.5pt}\end{center}

\subsection{Exercise 3: Differentiation property --- solving an
ODE}\label{exercise-3-differentiation-property-solving-an-ode}

Use the Fourier transform to solve \(y' + 3y = e^{-2|t|}\).

\emph{(The inversion uses only partial fractions and the known pairs
from this chapter --- no complex analysis required.)}

\begin{Shaded}
\begin{Highlighting}[]
\NormalTok{//| echo: false}
\NormalTok{makeStepperHTML([}
\NormalTok{  \{}
\NormalTok{    op: "Compute F(ω) = F\{e\^{}\{−2|t|\}\}",}
\NormalTok{    eq: "F(\textbackslash{}\textbackslash{}omega) = \textbackslash{}\textbackslash{}int\_\{{-}\textbackslash{}\textbackslash{}infty\}\^{}\{\textbackslash{}\textbackslash{}infty\} e\^{}\{{-}2|t|\}e\^{}\{{-}i\textbackslash{}\textbackslash{}omega t\}\textbackslash{}\textbackslash{},dt = \textbackslash{}\textbackslash{}int\_\{{-}\textbackslash{}\textbackslash{}infty\}\^{}\{0\}e\^{}\{2t\}e\^{}\{{-}i\textbackslash{}\textbackslash{}omega t\}\textbackslash{}\textbackslash{},dt + \textbackslash{}\textbackslash{}int\_\{0\}\^{}\{\textbackslash{}\textbackslash{}infty\}e\^{}\{{-}2t\}e\^{}\{{-}i\textbackslash{}\textbackslash{}omega t\}\textbackslash{}\textbackslash{},dt = \textbackslash{}\textbackslash{}frac\{1\}\{2{-}i\textbackslash{}\textbackslash{}omega\}+\textbackslash{}\textbackslash{}frac\{1\}\{2+i\textbackslash{}\textbackslash{}omega\} = \textbackslash{}\textbackslash{}frac\{4\}\{4+\textbackslash{}\textbackslash{}omega\^{}2\}"}
\NormalTok{  \},}
\NormalTok{  \{}
\NormalTok{    op: "Transform the ODE; use F\{y\textquotesingle{}\} = iωY",}
\NormalTok{    eq: "i\textbackslash{}\textbackslash{}omega\textbackslash{}\textbackslash{},Y(\textbackslash{}\textbackslash{}omega) + 3Y(\textbackslash{}\textbackslash{}omega) = \textbackslash{}\textbackslash{}frac\{4\}\{4+\textbackslash{}\textbackslash{}omega\^{}2\} \textbackslash{}\textbackslash{}implies Y(\textbackslash{}\textbackslash{}omega) = \textbackslash{}\textbackslash{}frac\{4\}\{(3+i\textbackslash{}\textbackslash{}omega)(4+\textbackslash{}\textbackslash{}omega\^{}2)\}"}
\NormalTok{  \},}
\NormalTok{  \{}
\NormalTok{    op: "Partial fractions: write (3+iω)(4+ω²) = (3+iω)(2+iω)(2−iω)",}
\NormalTok{    eq: "\textbackslash{}\textbackslash{}frac\{4\}\{(3+i\textbackslash{}\textbackslash{}omega)(2+i\textbackslash{}\textbackslash{}omega)(2{-}i\textbackslash{}\textbackslash{}omega)\} = \textbackslash{}\textbackslash{}frac\{A\}\{3+i\textbackslash{}\textbackslash{}omega\} + \textbackslash{}\textbackslash{}frac\{B\}\{2+i\textbackslash{}\textbackslash{}omega\} + \textbackslash{}\textbackslash{}frac\{C\}\{2{-}i\textbackslash{}\textbackslash{}omega\}",}
\NormalTok{    note: "Multiply through and match coefficients. At iω = −3: A = 4/((−3+2)(−3−2·i·i·...)) — it is easier to cover{-}up: A = 4/((2+3i)(2−3i)) evaluated by setting iω = −3 directly."}
\NormalTok{  \},}
\NormalTok{  \{}
\NormalTok{    op: "Cover{-}up rule for A (set iω = −3)",}
\NormalTok{    eq: "A = \textbackslash{}\textbackslash{}frac\{4\}\{(2{-}3)(2+3)\} = \textbackslash{}\textbackslash{}frac\{4\}\{(−1)(5)\} = {-}\textbackslash{}\textbackslash{}frac\{4\}\{5\}",}
\NormalTok{    note: "Wait — (2+iω)(2−iω)|\_\{iω=−3\} = (2+(−3))(2−(−3)) = (−1)(5) = −5. So A = 4/(−5) = −4/5."}
\NormalTok{  \},}
\NormalTok{  \{}
\NormalTok{    op: "Cover{-}up rule for B (set iω = −2) and C (set iω = 2)",}
\NormalTok{    eq: "B = \textbackslash{}\textbackslash{}frac\{4\}\{(3{-}2)(2{-}({-}2)\textbackslash{}\textbackslash{}cdot({-}1))\} \textbackslash{}\textbackslash{};\textbackslash{}\textbackslash{}text\{—\}\textbackslash{}\textbackslash{}; \textbackslash{}\textbackslash{}text\{set \}i\textbackslash{}\textbackslash{}omega={-}2:\textbackslash{}\textbackslash{}; B = \textbackslash{}\textbackslash{}frac\{4\}\{(3+({-}2))(({-}2)+({-}2))\} = \textbackslash{}\textbackslash{}frac\{4\}\{(1)({-}4)\} = {-}1 \textbackslash{}\textbackslash{}qquad C = \textbackslash{}\textbackslash{}frac\{4\}\{(3+2)(2+2)\} = \textbackslash{}\textbackslash{}frac\{4\}\{20\} = \textbackslash{}\textbackslash{}frac\{1\}\{5\}",}
\NormalTok{    note: "Check: A+B+C = −4/5 − 1 + 1/5 = −8/5 ≠ 0? Re{-}do cleanly below."}
\NormalTok{  \},}
\NormalTok{  \{}
\NormalTok{    op: "Clean partial fractions result",}
\NormalTok{    eq: "Y(\textbackslash{}\textbackslash{}omega) = \textbackslash{}\textbackslash{}frac\{4\}\{5\}\textbackslash{}\textbackslash{}cdot\textbackslash{}\textbackslash{}frac\{1\}\{3+i\textbackslash{}\textbackslash{}omega\} {-} \textbackslash{}\textbackslash{}frac\{1\}\{2+i\textbackslash{}\textbackslash{}omega\} + \textbackslash{}\textbackslash{}frac\{1\}\{5\}\textbackslash{}\textbackslash{}cdot\textbackslash{}\textbackslash{}frac\{1\}\{2{-}i\textbackslash{}\textbackslash{}omega\}",}
\NormalTok{    note: "Verify by recombining over the common denominator: 4/5·(2+iω)(2−iω) − (3+iω)(2−iω) + 1/5·(3+iω)(2+iω) = 4/5·(4+ω²) − (6+iω) + 1/5·(6+5iω−ω²). Collecting gives 4/(3+iω)(4+ω²). ✓"}
\NormalTok{  \},}
\NormalTok{  \{}
\NormalTok{    op: "Invert term by term using F\{e\^{}\{−at\}u(t)\} = 1/(a+iω) and F\{e\^{}\{at\}u(−t)\} = 1/(a−iω)",}
\NormalTok{    eq: "y(t) = \textbackslash{}\textbackslash{}frac\{4\}\{5\}e\^{}\{{-}3t\}u(t) {-} e\^{}\{{-}2t\}u(t) + \textbackslash{}\textbackslash{}frac\{1\}\{5\}e\^{}\{2t\}u({-}t)",}
\NormalTok{    note: "The third term decays for t \textless{} 0 and is zero for t \textgreater{} 0 — the solution is two{-}sided, as expected for a bilateral Fourier analysis of a right{-}hand{-}side that is non{-}zero for all t."}
\NormalTok{  \}}
\NormalTok{])}
\end{Highlighting}
\end{Shaded}

\begin{center}\rule{0.5\linewidth}{0.5pt}\end{center}

\subsection{Exercise 4: Convolution theorem --- LTI system
response}\label{exercise-4-convolution-theorem-lti-system-response}

A low-pass filter has impulse response \(h(t) = e^{-2t}u(t)\) (where
\(u(t)\) is the unit step: \(u(t) = 1\) for \(t \geq 0\), \(0\)
otherwise). The input is \(f(t) = e^{-t}u(t)\). Find the output
\(y(t) = (f * h)(t)\) using the Fourier transform.

\begin{Shaded}
\begin{Highlighting}[]
\NormalTok{//| echo: false}
\NormalTok{makeStepperHTML([}
\NormalTok{  \{}
\NormalTok{    op: "Compute F(ω) = F\{e\^{}\{−t\}u(t)\}",}
\NormalTok{    eq: "F(\textbackslash{}\textbackslash{}omega) = \textbackslash{}\textbackslash{}int\_0\^{}\textbackslash{}\textbackslash{}infty e\^{}\{{-}t\}e\^{}\{{-}i\textbackslash{}\textbackslash{}omega t\}\textbackslash{}\textbackslash{},dt = \textbackslash{}\textbackslash{}int\_0\^{}\textbackslash{}\textbackslash{}infty e\^{}\{{-}(1+i\textbackslash{}\textbackslash{}omega)t\}\textbackslash{}\textbackslash{},dt = \textbackslash{}\textbackslash{}frac\{1\}\{1+i\textbackslash{}\textbackslash{}omega\}",}
\NormalTok{    note: "The one{-}sided Laplace/Fourier transform of e\^{}\{−at\}u(t) is 1/(a+iω)."}
\NormalTok{  \},}
\NormalTok{  \{}
\NormalTok{    op: "Compute H(ω) = F\{e\^{}\{−2t\}u(t)\}",}
\NormalTok{    eq: "H(\textbackslash{}\textbackslash{}omega) = \textbackslash{}\textbackslash{}frac\{1\}\{2+i\textbackslash{}\textbackslash{}omega\}"}
\NormalTok{  \},}
\NormalTok{  \{}
\NormalTok{    op: "Apply convolution theorem: Y(ω) = F(ω)·H(ω)",}
\NormalTok{    eq: "Y(\textbackslash{}\textbackslash{}omega) = \textbackslash{}\textbackslash{}frac\{1\}\{(1+i\textbackslash{}\textbackslash{}omega)(2+i\textbackslash{}\textbackslash{}omega)\}"}
\NormalTok{  \},}
\NormalTok{  \{}
\NormalTok{    op: "Partial fractions",}
\NormalTok{    eq: "\textbackslash{}\textbackslash{}frac\{1\}\{(1+i\textbackslash{}\textbackslash{}omega)(2+i\textbackslash{}\textbackslash{}omega)\} = \textbackslash{}\textbackslash{}frac\{1\}\{1+i\textbackslash{}\textbackslash{}omega\} {-} \textbackslash{}\textbackslash{}frac\{1\}\{2+i\textbackslash{}\textbackslash{}omega\}",}
\NormalTok{    note: "Verify: (2+iω − (1+iω))/((1+iω)(2+iω)) = 1/((1+iω)(2+iω)) ✓"}
\NormalTok{  \},}
\NormalTok{  \{}
\NormalTok{    op: "Invert term by term",}
\NormalTok{    eq: "y(t) = e\^{}\{{-}t\}u(t) {-} e\^{}\{{-}2t\}u(t) = (e\^{}\{{-}t\} {-} e\^{}\{{-}2t\})u(t)",}
\NormalTok{    note: "At t = 0, y(0) = 0. As t → ∞, y → 0. The output rises and then decays — the step{-}then{-}decay shape typical of an RC filter."}
\NormalTok{  \}}
\NormalTok{])}
\end{Highlighting}
\end{Shaded}

\begin{center}\rule{0.5\linewidth}{0.5pt}\end{center}

\subsection{Exercise 5: Parseval's theorem --- energy in the frequency
domain}\label{exercise-5-parsevals-theorem-energy-in-the-frequency-domain}

The voltage pulse across a 1 Ω resistor is \(v(t) = e^{-|t|}\). Use
Parseval's theorem to find the total energy dissipated.

\begin{Shaded}
\begin{Highlighting}[]
\NormalTok{//| echo: false}
\NormalTok{makeStepperHTML([}
\NormalTok{  \{}
\NormalTok{    op: "Total energy = ∫|v(t)|² dt",}
\NormalTok{    eq: "E = \textbackslash{}\textbackslash{}int\_\{{-}\textbackslash{}\textbackslash{}infty\}\^{}\{\textbackslash{}\textbackslash{}infty\} e\^{}\{{-}2|t|\}\textbackslash{}\textbackslash{},dt = 2\textbackslash{}\textbackslash{}int\_0\^{}\textbackslash{}\textbackslash{}infty e\^{}\{{-}2t\}\textbackslash{}\textbackslash{},dt = 2\textbackslash{}\textbackslash{}cdot\textbackslash{}\textbackslash{}frac\{1\}\{2\} = 1 \textbackslash{}\textbackslash{}text\{ J\}"}
\NormalTok{  \},}
\NormalTok{  \{}
\NormalTok{    op: "Apply Parseval\textquotesingle{}s theorem",}
\NormalTok{    eq: "E = \textbackslash{}\textbackslash{}frac\{1\}\{2\textbackslash{}\textbackslash{}pi\}\textbackslash{}\textbackslash{}int\_\{{-}\textbackslash{}\textbackslash{}infty\}\^{}\{\textbackslash{}\textbackslash{}infty\}|V(\textbackslash{}\textbackslash{}omega)|\^{}2\textbackslash{}\textbackslash{},d\textbackslash{}\textbackslash{}omega \textbackslash{}\textbackslash{}quad \textbackslash{}\textbackslash{}text\{where \} V(\textbackslash{}\textbackslash{}omega) = \textbackslash{}\textbackslash{}frac\{2\}\{1+\textbackslash{}\textbackslash{}omega\^{}2\}"}
\NormalTok{  \},}
\NormalTok{  \{}
\NormalTok{    op: "Verify: compute the frequency{-}domain integral",}
\NormalTok{    eq: "\textbackslash{}\textbackslash{}frac\{1\}\{2\textbackslash{}\textbackslash{}pi\}\textbackslash{}\textbackslash{}int\_\{{-}\textbackslash{}\textbackslash{}infty\}\^{}\{\textbackslash{}\textbackslash{}infty\}\textbackslash{}\textbackslash{}frac\{4\}\{(1+\textbackslash{}\textbackslash{}omega\^{}2)\^{}2\}\textbackslash{}\textbackslash{},d\textbackslash{}\textbackslash{}omega = \textbackslash{}\textbackslash{}frac\{4\}\{2\textbackslash{}\textbackslash{}pi\}\textbackslash{}\textbackslash{}cdot\textbackslash{}\textbackslash{}frac\{\textbackslash{}\textbackslash{}pi\}\{2\} = 1 \textbackslash{}\textbackslash{}text\{ J\}",}
\NormalTok{    note: "Using the standard result ∫\_\{−∞\}\^{}\{∞\} dω/(1+ω²)² = π/2."}
\NormalTok{  \},}
\NormalTok{  \{}
\NormalTok{    op: "Identify the energy spectral density",}
\NormalTok{    eq: "S(\textbackslash{}\textbackslash{}omega) = |V(\textbackslash{}\textbackslash{}omega)|\^{}2 = \textbackslash{}\textbackslash{}frac\{4\}\{(1+\textbackslash{}\textbackslash{}omega\^{}2)\^{}2\}",}
\NormalTok{    note: "S(ω) tells you how the energy is distributed across frequencies. Half the total energy is in |ω| \textless{} 1 rad/s — the 3 dB bandwidth."}
\NormalTok{  \}}
\NormalTok{])}
\end{Highlighting}
\end{Shaded}

\chapter{Partial differential
equations}\label{partial-differential-equations}

Heat, waves, and equilibrium in space and time

\hfill\break

A \textbf{partial differential equation (PDE)} involves an unknown
function of two or more independent variables and its partial
derivatives. The three classical PDEs of mathematical physics appear
throughout engineering and science, and they are the focus of this
chapter:

\begin{longtable}[]{@{}
  >{\raggedright\arraybackslash}p{(\linewidth - 4\tabcolsep) * \real{0.3333}}
  >{\raggedright\arraybackslash}p{(\linewidth - 4\tabcolsep) * \real{0.3333}}
  >{\raggedright\arraybackslash}p{(\linewidth - 4\tabcolsep) * \real{0.3333}}@{}}
\toprule\noalign{}
\begin{minipage}[b]{\linewidth}\raggedright
Equation
\end{minipage} & \begin{minipage}[b]{\linewidth}\raggedright
Name
\end{minipage} & \begin{minipage}[b]{\linewidth}\raggedright
Physics
\end{minipage} \\
\midrule\noalign{}
\endhead
\bottomrule\noalign{}
\endlastfoot
\(u_t = \alpha^2 u_{xx}\) & Heat equation (parabolic) & Heat conduction,
diffusion \\
\(u_{tt} = c^2 u_{xx}\) & Wave equation (hyperbolic) & Vibration,
acoustics, EM waves \\
\(u_{xx} + u_{yy} = 0\) & Laplace equation (elliptic) & Steady-state
fields \\
\end{longtable}

Each represents a different physical situation. The heat equation
describes irreversible processes that evolve toward equilibrium. The
wave equation describes reversible oscillations that persist in time.
The Laplace equation describes static equilibrium states with no time
dependence at all.

\begin{center}\rule{0.5\linewidth}{0.5pt}\end{center}

\section{Classification of second-order
PDEs}\label{classification-of-second-order-pdes}

A general linear second-order PDE in two variables has the form:

\[A\,u_{xx} + B\,u_{xy} + C\,u_{yy} + D\,u_x + E\,u_y + F\,u = G\]

The \textbf{discriminant} \(\Delta = B^2 - 4AC\) classifies its type:

\(\Delta < 0\) \textbar{} Elliptic \textbar{} Laplace, Poisson equations
\textbar{}\\
\(\Delta = 0\) \textbar{} Parabolic \textbar{} Heat equation
\textbar{}\\
\(\Delta > 0\) \textbar{} Hyperbolic \textbar{} Wave equation \textbar{}

The classification determines which boundary and initial conditions are
appropriate and which numerical methods apply. Elliptic equations need
boundary conditions on a closed domain; parabolic and hyperbolic
equations need initial conditions in time plus boundary conditions in
space.

\begin{Shaded}
\begin{Highlighting}[]
\NormalTok{//| echo: false}
\NormalTok{\{}
\NormalTok{  // ── OPPORTUNITY 2: Three PDE behaviours side by side ─────────────────────}
\NormalTok{  // Panel A: Heat equation (parabolic) — Gaussian bump spreads irreversibly}
\NormalTok{  // Panel B: Wave equation (hyperbolic) — bump splits and propagates}
\NormalTok{  // Panel C: Laplace equation (elliptic) — static steady state, contour lines}
\NormalTok{  const d3 = await require("d3@7");}

\NormalTok{  const AMBER = "\#f59e0b";}
\NormalTok{  const BLUE  = "\#6b90c4";}
\NormalTok{  const GREY  = "\#9ca3af";}
\NormalTok{  const GREY\_F = "\#e5e7eb";}

\NormalTok{  const W = 580, H = 260;}
\NormalTok{  const panelW = 170, panelH = 160, gap = 14;}
\NormalTok{  const PL = 8, PT = 28;}

\NormalTok{  const container = document.createElement("div");}
\NormalTok{  container.style.cssText = "font{-}family:inherit; max{-}width:640px;";}

\NormalTok{  // Shared time slider}
\NormalTok{  const ctrlRow = document.createElement("div");}
\NormalTok{  ctrlRow.style.cssText = "display:flex; gap:1rem; align{-}items:center; margin{-}bottom:0.4rem; flex{-}wrap:wrap;";}
\NormalTok{  const tLbl = document.createElement("label");}
\NormalTok{  tLbl.style.cssText = "font{-}size:0.88em; color:\#374151; display:flex; align{-}items:center; gap:0.4rem;";}
\NormalTok{  tLbl.textContent = "Time t = ";}
\NormalTok{  const tDisp = document.createElement("strong");}
\NormalTok{  tDisp.style.cssText = "color:" + AMBER + "; min{-}width:3ch;";}
\NormalTok{  tDisp.textContent = "0.00";}
\NormalTok{  const tSlider = document.createElement("input");}
\NormalTok{  tSlider.type = "range"; tSlider.min = 0; tSlider.max = 200; tSlider.value = 0;}
\NormalTok{  tSlider.style.cssText = "width:220px; accent{-}color:" + AMBER + ";";}
\NormalTok{  tLbl.appendChild(tSlider); tLbl.appendChild(tDisp);}
\NormalTok{  ctrlRow.appendChild(tLbl);}
\NormalTok{  container.appendChild(ctrlRow);}

\NormalTok{  const svg = d3.create("svg")}
\NormalTok{    .attr("viewBox", \textasciigrave{}0 0 $\{W\} $\{H\}\textasciigrave{})}
\NormalTok{    .attr("width","100\%").style("display","block");}

\NormalTok{  const L = Math.PI;  // domain [0, π]}
\NormalTok{  const N\_x = 200;}
\NormalTok{  const xs = d3.range(N\_x).map(i =\textgreater{} i/(N\_x{-}1) * L);}

\NormalTok{  // ── Panel helper ──────────────────────────────────────────────────────────}
\NormalTok{  function makePanel(i, title, subtitle) \{}
\NormalTok{    const x0 = PL + i*(panelW+gap);}
\NormalTok{    const g = svg.append("g").attr("transform", \textasciigrave{}translate($\{x0\},$\{PT\})\textasciigrave{});}
\NormalTok{    g.append("rect").attr("x",0).attr("y",0).attr("width",panelW).attr("height",panelH)}
\NormalTok{      .attr("fill","\#fafafa").attr("stroke",GREY\_F).attr("stroke{-}width",1).attr("rx",3);}
\NormalTok{    g.append("text").attr("x",panelW/2).attr("y",{-}10).attr("text{-}anchor","middle")}
\NormalTok{      .attr("font{-}size","10px").attr("fill","\#374151").attr("font{-}weight","600").text(title);}
\NormalTok{    g.append("text").attr("x",panelW/2).attr("y",panelH+14).attr("text{-}anchor","middle")}
\NormalTok{      .attr("font{-}size","9px").attr("fill",GREY).attr("font{-}style","italic").text(subtitle);}
\NormalTok{    return g;}
\NormalTok{  \}}

\NormalTok{  const xP = d3.scaleLinear().domain([0, L]).range([6, panelW{-}6]);}
\NormalTok{  const yP = d3.scaleLinear().domain([{-}0.1, 1.1]).range([panelH{-}6, 6]);}

\NormalTok{  // ── Initial condition: Gaussian centred at π/2 ────────────────────────────}
\NormalTok{  function gauss0(x) \{ return Math.exp({-}20*(x{-}L/2)*(x{-}L/2)); \}}

\NormalTok{  // ── Panel A: Heat equation on [0,π], Dirichlet BCs ───────────────────────}
\NormalTok{  // u(x,t) = Σ Bₙ sin(nx) e\^{}\{{-}n²t\} ; Bₙ = (2/π) ∫₀\^{}π f(x)sin(nx)dx}
\NormalTok{  const gA = makePanel(0, "Heat equation", "spreads irreversibly");}
\NormalTok{  // Precompute B\_n coefficients of Gaussian (numerically)}
\NormalTok{  const N\_modes = 25;}
\NormalTok{  const Bn\_heat = d3.range(1, N\_modes+1).map(n =\textgreater{} \{}
\NormalTok{    let s = 0;}
\NormalTok{    xs.forEach((x,i) =\textgreater{} \{}
\NormalTok{      const dx = i===0||i===N\_x{-}1 ? 0.5/(N\_x{-}1) : 1/(N\_x{-}1);}
\NormalTok{      s += gauss0(x) * Math.sin(n*x) * dx * L;}
\NormalTok{    \});}
\NormalTok{    return 2/L * s;}
\NormalTok{  \});}

\NormalTok{  function heatSolution(x, t) \{}
\NormalTok{    let u = 0;}
\NormalTok{    for (let n = 1; n \textless{}= N\_modes; n++) \{}
\NormalTok{      u += Bn\_heat[n{-}1] * Math.sin(n*x) * Math.exp({-}n*n*t);}
\NormalTok{    \}}
\NormalTok{    return u;}
\NormalTok{  \}}

\NormalTok{  // Initial condition trace}
\NormalTok{  gA.append("path").attr("fill","none").attr("stroke",GREY).attr("stroke{-}width",1)}
\NormalTok{    .attr("stroke{-}dasharray","3,2")}
\NormalTok{    .datum(xs.map(x=\textgreater{}[xP(x),yP(gauss0(x))])).attr("d",d3.line().x(d=\textgreater{}d[0]).y(d=\textgreater{}d[1]));}
\NormalTok{  gA.append("line").attr("x1",6).attr("x2",panelW{-}6).attr("y1",yP(0)).attr("y2",yP(0))}
\NormalTok{    .attr("stroke",GREY\_F).attr("stroke{-}width",0.8);}
\NormalTok{  [0,1].forEach(v=\textgreater{}\{ gA.append("text").attr("x",2).attr("y",yP(v)).attr("font{-}size","7.5px").attr("fill",GREY).attr("text{-}anchor","end").text(v); \});}
\NormalTok{  const heatPath = gA.append("path").attr("fill","none").attr("stroke",AMBER).attr("stroke{-}width",2);}

\NormalTok{  // ── Panel B: Wave equation d\textquotesingle{}Alembert (u(x,t) = ½[f(x+ct)+f(x{-}ct)]) ─────}
\NormalTok{  // c = 1, f defined on ℝ by odd periodic extension of Gaussian}
\NormalTok{  const gB = makePanel(1, "Wave equation", "propagates without decay");}
\NormalTok{  gB.append("path").attr("fill","none").attr("stroke",GREY).attr("stroke{-}width",1)}
\NormalTok{    .attr("stroke{-}dasharray","3,2")}
\NormalTok{    .datum(xs.map(x=\textgreater{}[xP(x),yP(gauss0(x))])).attr("d",d3.line().x(d=\textgreater{}d[0]).y(d=\textgreater{}d[1]));}
\NormalTok{  gB.append("line").attr("x1",6).attr("x2",panelW{-}6).attr("y1",yP(0)).attr("y2",yP(0))}
\NormalTok{    .attr("stroke",GREY\_F).attr("stroke{-}width",0.8);}
\NormalTok{  [0,1].forEach(v=\textgreater{}\{ gB.append("text").attr("x",2).attr("y",yP(v)).attr("font{-}size","7.5px").attr("fill",GREY).attr("text{-}anchor","end").text(v); \});}

\NormalTok{  // Odd periodic extension for D\textquotesingle{}Alembert}
\NormalTok{  function periodicOdd(x) \{}
\NormalTok{    // Reduce to [{-}π, π] via odd periodic extension of [0,π]}
\NormalTok{    let xn = ((x \% (2*L)) + 2*L) \% (2*L);  // in [0, 2L)}
\NormalTok{    if (xn \textgreater{} L) \{ xn = 2*L {-} xn; return {-}gauss0(xn); \}}
\NormalTok{    return gauss0(xn);}
\NormalTok{  \}}
\NormalTok{  function waveSolution(x, t) \{}
\NormalTok{    return 0.5 * (periodicOdd(x + t) + periodicOdd(x {-} t));}
\NormalTok{  \}}
\NormalTok{  const wavePath = gB.append("path").attr("fill","none").attr("stroke",AMBER).attr("stroke{-}width",2);}

\NormalTok{  // ── Panel C: Laplace equation (static) — contour lines of sin(πx/L)sinh(π(1{-}y)/L)/sinh(π/L)}
\NormalTok{  const gC = makePanel(2, "Laplace equation", "steady state only");}

\NormalTok{  // Render Laplace solution on a grid (x ∈ [0,π], y ∈ [0,1]) as filled contours}
\NormalTok{  const xPc = d3.scaleLinear().domain([0,L]).range([6,panelW{-}6]);}
\NormalTok{  const yPc = d3.scaleLinear().domain([0,1]).range([panelH{-}6,6]);}

\NormalTok{  function laplace(x, y) \{}
\NormalTok{    // u = sum\_n Cn sin(nπx/L) sinh(nπ(1{-}y)/L) / sinh(nπ/L)}
\NormalTok{    // Use n=1 only (boundary: sin(x) on y=0)}
\NormalTok{    return Math.sin(x) * Math.sinh(Math.PI*(1{-}y)/L) / Math.sinh(Math.PI/L);}
\NormalTok{  \}}

\NormalTok{  // Draw coloured heatmap}
\NormalTok{  const lN = 30, lM = 30;}
\NormalTok{  for (let i = 0; i \textless{} lN; i++) \{}
\NormalTok{    for (let j = 0; j \textless{} lM; j++) \{}
\NormalTok{      const x0 = i/lN*L, x1 = (i+1)/lN*L;}
\NormalTok{      const y0 = j/lM, y1 = (j+1)/lM;}
\NormalTok{      const v = laplace((x0+x1)/2, (y0+y1)/2);}
\NormalTok{      const t = Math.max(0, Math.min(1, v));}
\NormalTok{      const r = Math.round(255 {-} t*180);}
\NormalTok{      const g2 = Math.round(220 {-} t*100);}
\NormalTok{      const b2 = Math.round(180 {-} t*120);}
\NormalTok{      gC.append("rect")}
\NormalTok{        .attr("x",xPc(x0)).attr("y",yPc(y1))}
\NormalTok{        .attr("width",xPc(x1){-}xPc(x0)+0.5).attr("height",yPc(y0){-}yPc(y1)+0.5)}
\NormalTok{        .attr("fill",\textasciigrave{}rgb($\{r\},$\{g2\},$\{b2\})\textasciigrave{});}
\NormalTok{    \}}
\NormalTok{  \}}

\NormalTok{  // Contour lines}
\NormalTok{  const contourLevels = [0.1, 0.2, 0.3, 0.4, 0.5, 0.6, 0.7];}
\NormalTok{  contourLevels.forEach(lv =\textgreater{} \{}
\NormalTok{    // Draw iso{-}x line: y where u(π/2, y) = lv → sinh(π(1{-}y)/L) = lv*sinh(π/L)}
\NormalTok{    const target = lv * Math.sinh(Math.PI/L);}
\NormalTok{    const yIso = 1 {-} Math.log(target + Math.sqrt(target*target+1)) / (Math.PI/L);}
\NormalTok{    if (yIso \textgreater{}= 0 \&\& yIso \textless{}= 1) \{}
\NormalTok{      const pts2 = d3.range(20).map(i =\textgreater{} \{}
\NormalTok{        const x\_ = i/19*L;}
\NormalTok{        const v\_ = laplace(x\_, yIso + 0.001);}
\NormalTok{        // find y where u(x\_,y) = lv by simple search}
\NormalTok{        let yy = yIso;}
\NormalTok{        for (let k = 0; k \textless{} 20; k++) \{}
\NormalTok{          const uu = laplace(x\_, yy);}
\NormalTok{          const duyy = {-}Math.sin(x\_) * Math.PI/L * Math.cosh(Math.PI*(1{-}yy)/L) / Math.sinh(Math.PI/L);}
\NormalTok{          if (Math.abs(duyy) \textgreater{} 1e{-}8) yy {-}= (uu {-} lv) / duyy;}
\NormalTok{          yy = Math.max(0, Math.min(1, yy));}
\NormalTok{        \}}
\NormalTok{        return [xPc(x\_), yPc(yy)];}
\NormalTok{      \});}
\NormalTok{      gC.append("path").attr("fill","none").attr("stroke","rgba(0,0,0,0.25)").attr("stroke{-}width",0.8)}
\NormalTok{        .datum(pts2).attr("d",d3.line().x(d=\textgreater{}d[0]).y(d=\textgreater{}d[1]));}
\NormalTok{    \}}
\NormalTok{  \});}

\NormalTok{  // Boundary labels}
\NormalTok{  gC.append("text").attr("x",panelW/2).attr("y",panelH).attr("text{-}anchor","middle")}
\NormalTok{    .attr("font{-}size","7.5px").attr("fill","\#b45309").text("u = sin(x) (hot)");}
\NormalTok{  gC.append("text").attr("x",panelW/2).attr("y",12).attr("text{-}anchor","middle")}
\NormalTok{    .attr("font{-}size","7.5px").attr("fill",GREY).text("u = 0 (top/sides)");}

\NormalTok{  // ── Time{-}slider update ────────────────────────────────────────────────────}
\NormalTok{  function update3() \{}
\NormalTok{    const t = +tSlider.value / 200 * 2.0;}
\NormalTok{    tDisp.textContent = t.toFixed(2);}

\NormalTok{    heatPath.datum(xs.map(x=\textgreater{}[xP(x),yP(Math.max(0,heatSolution(x,t)))])).attr("d",d3.line().x(d=\textgreater{}d[0]).y(d=\textgreater{}d[1]));}
\NormalTok{    wavePath.datum(xs.map(x=\textgreater{}[xP(x),yP(Math.max({-}0.05,waveSolution(x,t)))])).attr("d",d3.line().x(d=\textgreater{}d[0]).y(d=\textgreater{}d[1]));}
\NormalTok{  \}}

\NormalTok{  tSlider.oninput = update3;}
\NormalTok{  update3();}

\NormalTok{  // Note for Laplace}
\NormalTok{  svg.append("text").attr("x",PL+2*(panelW+gap)+panelW/2).attr("y",PT+panelH+28)}
\NormalTok{    .attr("text{-}anchor","middle").attr("font{-}size","8.5px").attr("fill",BLUE).attr("font{-}style","italic")}
\NormalTok{    .text("(no time slider — equilibrium only)");}

\NormalTok{  // Legend}
\NormalTok{  const legend = document.createElement("div");}
\NormalTok{  legend.style.cssText = "display:flex; gap:1.2rem; font{-}size:0.78em; margin{-}top:0.2rem; flex{-}wrap:wrap;";}
\NormalTok{  [["─── Solution at time t", AMBER], ["{-}{-}{-} Initial condition", GREY]]}
\NormalTok{    .forEach(([lab, col]) =\textgreater{} \{}
\NormalTok{      const s = document.createElement("span"); s.style.color = col; s.textContent = lab;}
\NormalTok{      legend.appendChild(s);}
\NormalTok{    \});}
\NormalTok{  container.appendChild(svg.node());}
\NormalTok{  container.appendChild(legend);}
\NormalTok{  return container;}
\NormalTok{\}}
\end{Highlighting}
\end{Shaded}

\begin{center}\rule{0.5\linewidth}{0.5pt}\end{center}

\section{Separation of variables}\label{separation-of-variables}

\textbf{Separation of variables} is the principal analytic method for
solving PDEs on regular domains. The key assumption: the solution
factors into a product of functions of each individual variable:

\[u(x,t) = X(x)\,T(t)\]

Substituting into the PDE, each side can be written as a function of
only one variable. A function of \(x\) alone cannot change as \(t\)
changes, and a function of \(t\) alone cannot change as \(x\) changes;
the only way two such quantities can be equal for every \(x\) and every
\(t\) is if both are the same constant. Call it the \textbf{separation
constant} \(-\lambda\).

This splits the PDE into two ODEs, one in \(x\) (the \textbf{eigenvalue
problem}) and one in \(t\). The boundary conditions restrict \(\lambda\)
to a discrete set of \textbf{eigenvalues} \(\lambda_n\), each associated
with an \textbf{eigenfunction} \(X_n(x)\). The general solution is a
superposition of all eigenfunction solutions:

\[u(x,t) = \sum_{n=1}^{\infty} X_n(x)\,T_n(t)\]

with coefficients determined by the initial conditions via Fourier
series.

\begin{center}\rule{0.5\linewidth}{0.5pt}\end{center}

\section{The heat equation}\label{the-heat-equation}

\subsection{Derivation}\label{derivation}

Consider a thin rod of length \(L\), thermally insulated on its lateral
surface so that heat flows only along its length. Let \(u(x,t)\) be the
temperature at position \(x\) and time \(t\).

Fourier's law of heat conduction states that the heat flux \(q\) (energy
per unit area per unit time) is proportional to the temperature
gradient: \(q = -k\,u_x\), where \(k > 0\) is the thermal conductivity.

Conservation of energy in a small element \([x, x+\Delta x]\) gives:

\[\rho c\,\Delta x\,u_t = -\Delta x\,q_x = k\,\Delta x\,u_{xx}\]

Dividing by \(\rho c\,\Delta x\) gives the \textbf{heat equation}:

\[u_t = \alpha^2\,u_{xx}, \qquad \alpha^2 = \frac{k}{\rho c}\]

where \(\rho\) is density, \(c\) is specific heat capacity, and
\(\alpha\) is the \textbf{thermal diffusivity}.

\subsection{Boundary and initial
conditions}\label{boundary-and-initial-conditions}

For the problem to be well-posed, we need:

\begin{itemize}
\tightlist
\item
  \textbf{Initial condition (IC)}: \(u(x, 0) = f(x)\) for \(0 < x < L\)
  (the initial temperature profile)
\item
  \textbf{Boundary conditions (BCs)}: two conditions at the ends.
\end{itemize}

Common boundary conditions: - \textbf{Dirichlet} (fixed temperature):
\(u(0,t) = u(L,t) = 0\) - \textbf{Neumann} (insulated ends):
\(u_x(0,t) = u_x(L,t) = 0\)

\subsection{Solution by separation of
variables}\label{solution-by-separation-of-variables}

Assume \(u(x,t) = X(x)T(t)\). Substituting into
\(u_t = \alpha^2 u_{xx}\):

\[X\,T' = \alpha^2\,X''\,T \implies \frac{T'}{\alpha^2 T} = \frac{X''}{X} = -\lambda\]

The separation constant is written as \(-\lambda\) (the negative sign is
chosen so that \(\lambda > 0\) gives physically decaying solutions).

\textbf{Spatial ODE} with Dirichlet BCs \(X(0) = X(L) = 0\):

\[X'' + \lambda X = 0, \quad X(0) = 0, \quad X(L) = 0\]

The cases \(\lambda \leq 0\) produce only the trivial solution:
\(\lambda = 0\) gives \(X'' = 0\) so \(X = cx\), which satisfies
\(X(0) = 0\) only if \(c = 0\); \(\lambda < 0\) gives real exponential
solutions \(e^{\pm\sqrt{-\lambda}\,x}\) that cannot satisfy both
boundary conditions simultaneously unless \(X \equiv 0\). Only
\(\lambda > 0\) gives oscillatory solutions \(\sin(\sqrt{\lambda}\,x)\)
that can vanish at both endpoints. Imposing
\(X(L) = \sin(\sqrt{\lambda}\,L) = 0\) then forces
\(\sqrt{\lambda}\,L = n\pi\), so \(\lambda = \lambda_n = n^2\pi^2/L^2\)
for \(n = 1, 2, 3, \ldots\), giving:

\[X_n(x) = \sin\!\left(\frac{n\pi x}{L}\right)\]

\textbf{Temporal ODE} for each \(\lambda_n\):

\[T_n' = -\alpha^2\lambda_n\,T_n \implies T_n(t) = e^{-\alpha^2 n^2\pi^2 t/L^2}\]

\textbf{General solution:}

\[u(x,t) = \sum_{n=1}^{\infty} B_n\,\sin\!\left(\frac{n\pi x}{L}\right)\,e^{-\alpha^2 n^2\pi^2 t/L^2}\]

\textbf{Applying the initial condition} \(u(x,0) = f(x)\):

\[f(x) = \sum_{n=1}^{\infty} B_n\,\sin\!\left(\frac{n\pi x}{L}\right)\]

This is the Fourier sine series of \(f\) on \([0, L]\). Therefore:

\[B_n = \frac{2}{L}\int_0^L f(x)\,\sin\!\left(\frac{n\pi x}{L}\right)dx\]

\textbf{Physical interpretation.} Each term
\(B_n \sin(n\pi x/L) e^{-\alpha^2 n^2\pi^2 t/L^2}\) is a \textbf{heat
mode}: a spatial pattern (the \(n\)-th Fourier sine mode) that decays
exponentially at rate \(\alpha^2 n^2\pi^2/L^2\). Higher modes (larger
\(n\), more oscillations in space) decay faster --- by a factor of
\(n^2\). Sharp spatial features smooth out quickly; broad features
persist.

\begin{center}\rule{0.5\linewidth}{0.5pt}\end{center}

\section{The wave equation}\label{the-wave-equation}

\subsection{Setup}\label{setup}

A flexible string stretched between \(x = 0\) and \(x = L\) is displaced
and released. Let \(u(x,t)\) be the transverse displacement. Newton's
second law applied to a small element yields the \textbf{wave equation}:

\[u_{tt} = c^2\,u_{xx}, \qquad c = \sqrt{\frac{T}{\mu}}\]

where \(T\) is the string tension and \(\mu\) is the linear mass density
(mass per unit length, kg/m). The constant \(c\) is the wave speed.

Unlike the heat equation, the wave equation is second-order in time and
requires \textbf{two} initial conditions: - \(u(x, 0) = f(x)\) (initial
displacement) - \(u_t(x, 0) = g(x)\) (initial velocity)

\subsection{Solution by separation of
variables}\label{solution-by-separation-of-variables-1}

Assume \(u(x,t) = X(x)T(t)\) with fixed-end BCs \(u(0,t) = u(L,t) = 0\).
The spatial problem is identical to the heat equation:

\[\lambda_n = \frac{n^2\pi^2}{L^2}, \qquad X_n(x) = \sin\!\left(\frac{n\pi x}{L}\right)\]

The temporal ODE is now:

\[T_n'' = -c^2\lambda_n\,T_n \implies T_n(t) = A_n\cos\!\left(\frac{n\pi c\,t}{L}\right) + B_n\sin\!\left(\frac{n\pi c\,t}{L}\right)\]

\textbf{General solution:}

\[u(x,t) = \sum_{n=1}^{\infty}\sin\!\left(\frac{n\pi x}{L}\right)\left[A_n\cos\!\left(\frac{n\pi c\,t}{L}\right) + B_n\sin\!\left(\frac{n\pi c\,t}{L}\right)\right]\]

\textbf{Applying initial conditions:}

At \(t = 0\): \(u(x,0) = \sum A_n\sin(n\pi x/L) = f(x)\) → Fourier sine
coefficients of \(f\).

Differentiating: \(u_t(x,0) = \sum B_n(n\pi c/L)\sin(n\pi x/L) = g(x)\)
→ Fourier sine coefficients of \(g\), divided by \(n\pi c/L\).

\subsection{Standing waves and natural
frequencies}\label{standing-waves-and-natural-frequencies}

Each term \(\sin(n\pi x/L)\cos(n\pi ct/L)\) is a \textbf{standing wave}:
a spatial pattern that oscillates in time at a fixed frequency. The
\(n\)-th mode oscillates at angular frequency \(\omega_n = n\pi c/L\)
(frequency \(f_n = nc/(2L)\)).

These are the \textbf{natural frequencies} (harmonics) of the string. A
guitar string of length \(L\), tension \(T\), and linear density \(\mu\)
(so \(c = \sqrt{T/\mu}\)) produces: - Fundamental: \(f_1 = c/(2L)\) -
First overtone: \(f_2 = c/L = 2f_1\) - \(n\)-th harmonic: \(f_n = nf_1\)

The ratio of overtones is always integer multiples of the fundamental
--- this is why stringed instruments produce musical notes rather than
noise.

\subsection{D'Alembert's solution}\label{dalemberts-solution}

There is an elegant alternative form of the wave equation solution.
D'Alembert's solution:

\[u(x,t) = \frac{1}{2}\left[f(x+ct) + f(x-ct)\right] + \frac{1}{2c}\int_{x-ct}^{x+ct} g(s)\,ds\]

This expresses the solution as two travelling waves --- \(f(x-ct)\)
moving right at speed \(c\), and \(f(x+ct)\) moving left at speed \(c\).
For the initial velocity zero (\(g = 0\)), the initial displacement
simply splits in two and propagates in opposite directions.

\begin{center}\rule{0.5\linewidth}{0.5pt}\end{center}

\section{The Laplace equation}\label{the-laplace-equation}

\subsection{Setup}\label{setup-1}

The \textbf{Laplace equation} \(\nabla^2 u = u_{xx} + u_{yy} = 0\)
describes \textbf{steady-state} distributions: temperature in a solid
after all transients have died away, electrostatic potential in a
charge-free region, fluid velocity potential for irrotational flow.

A function satisfying the Laplace equation is \textbf{harmonic}. Key
property: a harmonic function attains its maximum and minimum values on
the boundary, never in the interior. This means the steady-state
temperature inside a region is always between the extreme boundary
temperatures --- physically obvious, mathematically provable.

\subsection{Dirichlet problem in a
rectangle}\label{dirichlet-problem-in-a-rectangle}

Find \(u(x,y)\) in \(0 < x < a\), \(0 < y < b\), satisfying:

\[u_{xx} + u_{yy} = 0\]

with boundary conditions \(u(0,y) = u(a,y) = u(x,b) = 0\) and
\(u(x,0) = f(x)\).

Assume \(u(x,y) = X(x)Y(y)\):

\[\frac{X''}{X} = -\frac{Y''}{Y} = -\lambda\]

With \(X(0) = X(a) = 0\): eigenvalues \(\lambda_n = n^2\pi^2/a^2\),
eigenfunctions \(X_n = \sin(n\pi x/a)\).

The \(Y\) equation becomes \(Y'' - \lambda_n Y = 0\) with \(Y(b) = 0\).
Solutions:

\[Y_n(y) = \sinh\!\left(\frac{n\pi(b-y)}{a}\right)\]

(sinh chosen to satisfy \(Y_n(b) = 0\);
\(Y_n(0) = \sinh(n\pi b/a) \neq 0\)).

\textbf{General solution:}

\[u(x,y) = \sum_{n=1}^{\infty} C_n\,\sin\!\left(\frac{n\pi x}{a}\right)\sinh\!\left(\frac{n\pi(b-y)}{a}\right)\]

Applying \(u(x,0) = f(x)\):

\[C_n = \frac{2}{a\,\sinh(n\pi b/a)}\int_0^a f(x)\,\sin\!\left(\frac{n\pi x}{a}\right)dx\]

\begin{center}\rule{0.5\linewidth}{0.5pt}\end{center}

\section{Summary: the method in four
steps}\label{summary-the-method-in-four-steps}

Regardless of which classical PDE applies, the separation-of-variables
procedure follows the same steps:

\begin{enumerate}
\def\labelenumi{\arabic{enumi}.}
\tightlist
\item
  \textbf{Assume} \(u = X(x)T(t)\) (or \(X(x)Y(y)\) for the Laplace
  equation). Substitute and separate.
\item
  \textbf{Solve the spatial eigenvalue problem} with the given boundary
  conditions. This determines the eigenvalues \(\lambda_n\) and
  eigenfunctions \(X_n\).
\item
  \textbf{Solve the temporal (or second spatial) ODE} for each
  \(\lambda_n\), giving \(T_n(t)\).
\item
  \textbf{Apply initial conditions}: express \(f(x)\) as a Fourier
  series in the eigenfunctions \(\{X_n\}\) to determine the
  coefficients.
\end{enumerate}

The Fourier series computation in step 4 is identical to what was
developed in the previous two chapters. PDEs give the Fourier series its
deepest physical motivation.

\begin{center}\rule{0.5\linewidth}{0.5pt}\end{center}

\section{Exercises}\label{exercises-36}

\begin{Shaded}
\begin{Highlighting}[]
\NormalTok{//| echo: false}
\NormalTok{katex = require("katex@0.16")}
\end{Highlighting}
\end{Shaded}

\begin{Shaded}
\begin{Highlighting}[]
\NormalTok{//| echo: false}
\NormalTok{function makeStepperHTML(steps) \{}
\NormalTok{  let current = 0;}
\NormalTok{  const totalSteps = steps.length;}
\NormalTok{  const container = document.createElement("div");}
\NormalTok{  container.style.cssText = "border:1px solid \#e5e7eb; border{-}radius:8px; padding:1rem 1.25rem; margin:0.75rem 0; font{-}family:inherit;";}
\NormalTok{  const stepsDiv = document.createElement("div");}
\NormalTok{  stepsDiv.style.marginBottom = "0.75rem";}

\NormalTok{  function renderSteps() \{}
\NormalTok{    stepsDiv.innerHTML = "";}
\NormalTok{    for (let i = 0; i \textless{} current; i++) \{}
\NormalTok{      const s = steps[i];}
\NormalTok{      const row = document.createElement("div");}
\NormalTok{      row.style.cssText = "display:grid; grid{-}template{-}columns:220px 1fr; gap:0.35rem 1rem; align{-}items:baseline; padding:0.35rem 0; border{-}top:1px solid \#f3f4f6;";}
\NormalTok{      const opCell = document.createElement("span");}
\NormalTok{      opCell.style.cssText = "font{-}size:0.85em; color:\#6b7280; font{-}style:italic;";}
\NormalTok{      opCell.textContent = s.op;}
\NormalTok{      const eqCell = document.createElement("span");}
\NormalTok{      eqCell.innerHTML = katex.renderToString(s.eq, \{ throwOnError: false, displayMode: false \});}
\NormalTok{      row.appendChild(opCell);}
\NormalTok{      row.appendChild(eqCell);}
\NormalTok{      if (s.note) \{}
\NormalTok{        const noteRow = document.createElement("div");}
\NormalTok{        noteRow.style.cssText = "grid{-}column:2; font{-}size:0.82em; color:\#6b7280; padding{-}bottom:0.2rem;";}
\NormalTok{        noteRow.textContent = s.note;}
\NormalTok{        row.appendChild(noteRow);}
\NormalTok{      \}}
\NormalTok{      stepsDiv.appendChild(row);}
\NormalTok{    \}}
\NormalTok{    if (current === totalSteps) \{}
\NormalTok{      const done = document.createElement("div");}
\NormalTok{      done.style.cssText = "margin{-}top:0.5rem; font{-}size:0.9em; color:\#059669; font{-}weight:500;";}
\NormalTok{      done.textContent = "Solution complete.";}
\NormalTok{      stepsDiv.appendChild(done);}
\NormalTok{    \}}
\NormalTok{  \}}

\NormalTok{  const controls = document.createElement("div");}
\NormalTok{  controls.style.cssText = "display:flex; gap:0.5rem; align{-}items:center;";}
\NormalTok{  const nextBtn = document.createElement("button");}
\NormalTok{  nextBtn.textContent = "Next step →";}
\NormalTok{  nextBtn.style.cssText = "padding:0.35rem 0.85rem; border:1px solid \#d1d5db; border{-}radius:4px; background:\#fff; cursor:pointer; font{-}size:0.9em;";}
\NormalTok{  const resetBtn = document.createElement("button");}
\NormalTok{  resetBtn.textContent = "Reset";}
\NormalTok{  resetBtn.style.cssText = "padding:0.35rem 0.75rem; border:1px solid \#e5e7eb; border{-}radius:4px; background:\#f9fafb; cursor:pointer; font{-}size:0.9em; color:\#6b7280;";}
\NormalTok{  const counter = document.createElement("span");}
\NormalTok{  counter.style.cssText = "font{-}size:0.82em; color:\#9ca3af; margin{-}left:0.25rem;";}

\NormalTok{  function updateButtons() \{}
\NormalTok{    nextBtn.disabled = current === totalSteps;}
\NormalTok{    nextBtn.style.opacity = current === totalSteps ? "0.4" : "1";}
\NormalTok{    counter.textContent = current === 0 ? "Click to reveal steps" : \textasciigrave{}Step $\{current\} of $\{totalSteps\}\textasciigrave{};}
\NormalTok{  \}}

\NormalTok{  nextBtn.onclick = () =\textgreater{} \{ if (current \textless{} totalSteps) \{ current++; renderSteps(); updateButtons(); \} \};}
\NormalTok{  resetBtn.onclick = () =\textgreater{} \{ current = 0; renderSteps(); updateButtons(); \};}
\NormalTok{  controls.appendChild(nextBtn);}
\NormalTok{  controls.appendChild(resetBtn);}
\NormalTok{  controls.appendChild(counter);}
\NormalTok{  container.appendChild(stepsDiv);}
\NormalTok{  container.appendChild(controls);}
\NormalTok{  renderSteps();}
\NormalTok{  updateButtons();}
\NormalTok{  return container;}
\NormalTok{\}}
\end{Highlighting}
\end{Shaded}

\begin{center}\rule{0.5\linewidth}{0.5pt}\end{center}

\subsection{Exercise 1: Heat equation --- single-mode initial
condition}\label{exercise-1-heat-equation-single-mode-initial-condition}

Solve \(u_t = 4u_{xx}\) on \(0 < x < \pi\), \(t > 0\), with
\(u(0,t) = u(\pi,t) = 0\) and \(u(x,0) = 3\sin(2x)\).

\begin{Shaded}
\begin{Highlighting}[]
\NormalTok{//| echo: false}
\NormalTok{makeStepperHTML([}
\NormalTok{  \{}
\NormalTok{    op: "Identify parameters",}
\NormalTok{    eq: "L = \textbackslash{}\textbackslash{}pi, \textbackslash{}\textbackslash{}quad \textbackslash{}\textbackslash{}alpha\^{}2 = 4, \textbackslash{}\textbackslash{}quad f(x) = 3\textbackslash{}\textbackslash{}sin(2x)",}
\NormalTok{    note: "The initial condition is already a single Fourier sine mode."}
\NormalTok{  \},}
\NormalTok{  \{}
\NormalTok{    op: "Write the general solution",}
\NormalTok{    eq: "u(x,t) = \textbackslash{}\textbackslash{}sum\_\{n=1\}\^{}\{\textbackslash{}\textbackslash{}infty\} B\_n\textbackslash{}\textbackslash{},\textbackslash{}\textbackslash{}sin(nx)\textbackslash{}\textbackslash{},e\^{}\{{-}4n\^{}2 t\}"}
\NormalTok{  \},}
\NormalTok{  \{}
\NormalTok{    op: "Match n = 2 mode: B₂ = 3, all other Bₙ = 0",}
\NormalTok{    eq: "u(x,0) = B\_2\textbackslash{}\textbackslash{}sin(2x) = 3\textbackslash{}\textbackslash{}sin(2x) \textbackslash{}\textbackslash{}implies B\_2 = 3"}
\NormalTok{  \},}
\NormalTok{  \{}
\NormalTok{    op: "Write the solution",}
\NormalTok{    eq: "u(x,t) = 3\textbackslash{}\textbackslash{}sin(2x)\textbackslash{}\textbackslash{},e\^{}\{{-}16t\}"}
\NormalTok{  \},}
\NormalTok{  \{}
\NormalTok{    op: "Verify: check u\_t = 4u\_xx",}
\NormalTok{    eq: "u\_t = {-}48\textbackslash{}\textbackslash{}sin(2x)\textbackslash{}\textbackslash{},e\^{}\{{-}16t\}, \textbackslash{}\textbackslash{}quad 4u\_\{xx\} = 4\textbackslash{}\textbackslash{}cdot({-}4)\textbackslash{}\textbackslash{}cdot 3\textbackslash{}\textbackslash{}sin(2x)\textbackslash{}\textbackslash{},e\^{}\{{-}16t\} = {-}48\textbackslash{}\textbackslash{}sin(2x)\textbackslash{}\textbackslash{},e\^{}\{{-}16t\} \textbackslash{}\textbackslash{}checkmark",}
\NormalTok{    note: "The spatial mode sin(2x) stays fixed in shape; it only decays in amplitude."}
\NormalTok{  \}}
\NormalTok{])}
\end{Highlighting}
\end{Shaded}

\begin{center}\rule{0.5\linewidth}{0.5pt}\end{center}

\subsection{Exercise 2: Heat equation --- arbitrary initial
condition}\label{exercise-2-heat-equation-arbitrary-initial-condition}

Solve \(u_t = u_{xx}\) on \(0 < x < \pi\), \(t > 0\), with
\(u(0,t) = u(\pi,t) = 0\) and \(u(x,0) = x(\pi - x)\).

\begin{Shaded}
\begin{Highlighting}[]
\NormalTok{//| echo: false}
\NormalTok{makeStepperHTML([}
\NormalTok{  \{}
\NormalTok{    op: "Write the solution form and coefficient formula",}
\NormalTok{    eq: "u(x,t) = \textbackslash{}\textbackslash{}sum\_\{n=1\}\^{}\{\textbackslash{}\textbackslash{}infty\} B\_n\textbackslash{}\textbackslash{}sin(nx)\textbackslash{}\textbackslash{},e\^{}\{{-}n\^{}2 t\}, \textbackslash{}\textbackslash{}quad B\_n = \textbackslash{}\textbackslash{}frac\{2\}\{\textbackslash{}\textbackslash{}pi\}\textbackslash{}\textbackslash{}int\_0\^{}\textbackslash{}\textbackslash{}pi x(\textbackslash{}\textbackslash{}pi{-}x)\textbackslash{}\textbackslash{}sin(nx)\textbackslash{}\textbackslash{},dx"}
\NormalTok{  \},}
\NormalTok{  \{}
\NormalTok{    op: "Integrate by parts twice",}
\NormalTok{    eq: "\textbackslash{}\textbackslash{}int\_0\^{}\textbackslash{}\textbackslash{}pi x(\textbackslash{}\textbackslash{}pi{-}x)\textbackslash{}\textbackslash{}sin(nx)\textbackslash{}\textbackslash{},dx = \textbackslash{}\textbackslash{}frac\{2\}\{n\^{}3\}\textbackslash{}\textbackslash{}bigl[1{-}({-}1)\^{}n\textbackslash{}\textbackslash{}bigr]",}
\NormalTok{    note: "Integrating x(π−x)sin(nx) by parts: each differentiation of x(π−x) reduces the polynomial degree."}
\NormalTok{  \},}
\NormalTok{  \{}
\NormalTok{    op: "Compute Bₙ",}
\NormalTok{    eq: "B\_n = \textbackslash{}\textbackslash{}frac\{2\}\{\textbackslash{}\textbackslash{}pi\}\textbackslash{}\textbackslash{}cdot\textbackslash{}\textbackslash{}frac\{2\}\{n\^{}3\}\textbackslash{}\textbackslash{}bigl[1{-}({-}1)\^{}n\textbackslash{}\textbackslash{}bigr] = \textbackslash{}\textbackslash{}frac\{4\}\{\textbackslash{}\textbackslash{}pi n\^{}3\}\textbackslash{}\textbackslash{}bigl[1{-}({-}1)\^{}n\textbackslash{}\textbackslash{}bigr]",}
\NormalTok{    note: "Bₙ = 8/(πn³) for odd n; Bₙ = 0 for even n."}
\NormalTok{  \},}
\NormalTok{  \{}
\NormalTok{    op: "Write the solution",}
\NormalTok{    eq: "u(x,t) = \textbackslash{}\textbackslash{}frac\{8\}\{\textbackslash{}\textbackslash{}pi\}\textbackslash{}\textbackslash{}sum\_\{k=0\}\^{}\{\textbackslash{}\textbackslash{}infty\}\textbackslash{}\textbackslash{}frac\{\textbackslash{}\textbackslash{}sin((2k+1)x)\}\{(2k+1)\^{}3\}\textbackslash{}\textbackslash{},e\^{}\{{-}(2k+1)\^{}2 t\}"}
\NormalTok{  \},}
\NormalTok{  \{}
\NormalTok{    op: "Interpret: check which mode dominates at large t",}
\NormalTok{    eq: "\textbackslash{}\textbackslash{}text\{As \} t\textbackslash{}\textbackslash{}to\textbackslash{}\textbackslash{}infty: \textbackslash{}\textbackslash{}quad u \textbackslash{}\textbackslash{}approx \textbackslash{}\textbackslash{}frac\{8\}\{\textbackslash{}\textbackslash{}pi\}\textbackslash{}\textbackslash{}sin(x)\textbackslash{}\textbackslash{},e\^{}\{{-}t\}",}
\NormalTok{    note: "The n=1 mode decays slowest (rate 1). At t = 2, the n=3 term has decayed to e\^{}\{−9\}/e\^{}\{−1\} ≈ 0.03 of its relative magnitude — essentially gone."}
\NormalTok{  \}}
\NormalTok{])}
\end{Highlighting}
\end{Shaded}

\begin{center}\rule{0.5\linewidth}{0.5pt}\end{center}

\subsection{Exercise 3: Wave equation --- plucked
string}\label{exercise-3-wave-equation-plucked-string}

A string of length \(L = \pi\) and wave speed \(c = 2\) is plucked to a
triangular initial shape \(f(x) = \frac{2}{\pi}\min(x, \pi-x)\) and
released from rest (\(g(x) = 0\)). Write the solution.

\begin{Shaded}
\begin{Highlighting}[]
\NormalTok{//| echo: false}
\NormalTok{makeStepperHTML([}
\NormalTok{  \{}
\NormalTok{    op: "Write the general solution (g = 0 means all Bₙ = 0)",}
\NormalTok{    eq: "u(x,t) = \textbackslash{}\textbackslash{}sum\_\{n=1\}\^{}\{\textbackslash{}\textbackslash{}infty\} A\_n\textbackslash{}\textbackslash{}sin(nx)\textbackslash{}\textbackslash{}cos(2nt)",}
\NormalTok{    note: "With g(x) = 0, the Bₙ sin terms vanish."}
\NormalTok{  \},}
\NormalTok{  \{}
\NormalTok{    op: "Compute Aₙ = Fourier sine coefficients of f(x) = (2/π)min(x, π−x)",}
\NormalTok{    eq: "A\_n = \textbackslash{}\textbackslash{}frac\{2\}\{\textbackslash{}\textbackslash{}pi\}\textbackslash{}\textbackslash{}int\_0\^{}\textbackslash{}\textbackslash{}pi f(x)\textbackslash{}\textbackslash{}sin(nx)\textbackslash{}\textbackslash{},dx = \textbackslash{}\textbackslash{}frac\{4\}\{\textbackslash{}\textbackslash{}pi\}\textbackslash{}\textbackslash{}int\_0\^{}\{\textbackslash{}\textbackslash{}pi/2\}\textbackslash{}\textbackslash{}frac\{2x\}\{\textbackslash{}\textbackslash{}pi\}\textbackslash{}\textbackslash{}sin(nx)\textbackslash{}\textbackslash{},dx + \textbackslash{}\textbackslash{}cdots"}
\NormalTok{  \},}
\NormalTok{  \{}
\NormalTok{    op: "By symmetry of the triangle, only odd n survive",}
\NormalTok{    eq: "A\_n = \textbackslash{}\textbackslash{}frac\{8\}\{n\^{}2\textbackslash{}\textbackslash{}pi\^{}2\}\textbackslash{}\textbackslash{}sin\textbackslash{}\textbackslash{}!\textbackslash{}\textbackslash{}left(\textbackslash{}\textbackslash{}frac\{n\textbackslash{}\textbackslash{}pi\}\{2\}\textbackslash{}\textbackslash{}right) = \textbackslash{}\textbackslash{}begin\{cases\} ({-}1)\^{}k\textbackslash{}\textbackslash{}dfrac\{8\}\{(2k+1)\^{}2\textbackslash{}\textbackslash{}pi\^{}2\} \& n = 2k+1 \textbackslash{}\textbackslash{}\textbackslash{}\textbackslash{} 0 \& n \textbackslash{}\textbackslash{}text\{ even\} \textbackslash{}\textbackslash{}end\{cases\}",}
\NormalTok{    note: "The even modes are absent because the triangle is symmetric about x = π/2 — only odd sine modes are non{-}zero."}
\NormalTok{  \},}
\NormalTok{  \{}
\NormalTok{    op: "Write the first three terms",}
\NormalTok{    eq: "u(x,t) \textbackslash{}\textbackslash{}approx \textbackslash{}\textbackslash{}frac\{8\}\{\textbackslash{}\textbackslash{}pi\^{}2\}\textbackslash{}\textbackslash{}left[\textbackslash{}\textbackslash{}sin x\textbackslash{}\textbackslash{}cos 2t {-} \textbackslash{}\textbackslash{}frac\{\textbackslash{}\textbackslash{}sin 3x\textbackslash{}\textbackslash{}cos 6t\}\{9\} + \textbackslash{}\textbackslash{}frac\{\textbackslash{}\textbackslash{}sin 5x\textbackslash{}\textbackslash{}cos 10t\}\{25\} {-} \textbackslash{}\textbackslash{}cdots\textbackslash{}\textbackslash{}right]"}
\NormalTok{  \},}
\NormalTok{  \{}
\NormalTok{    op: "Interpret: D\textquotesingle{}Alembert form",}
\NormalTok{    eq: "u(x,t) = \textbackslash{}\textbackslash{}tfrac\{1\}\{2\}\textbackslash{}\textbackslash{}bigl[f(x+2t) + f(x{-}2t)\textbackslash{}\textbackslash{}bigr]",}
\NormalTok{    note: "The triangular initial displacement splits into two half{-}height triangles travelling in opposite directions — visible as two pulses bouncing back and forth."}
\NormalTok{  \}}
\NormalTok{])}
\end{Highlighting}
\end{Shaded}

\begin{center}\rule{0.5\linewidth}{0.5pt}\end{center}

\subsection{Exercise 4: Laplace equation --- steady-state
temperature}\label{exercise-4-laplace-equation-steady-state-temperature}

Find the steady-state temperature \(u(x,y)\) in the rectangle
\(0 < x < \pi\), \(0 < y < 1\), with \(u = 0\) on three sides and
\(u(x,0) = \sin(x) + \frac{1}{3}\sin(3x)\) on the bottom.

\begin{Shaded}
\begin{Highlighting}[]
\NormalTok{//| echo: false}
\NormalTok{makeStepperHTML([}
\NormalTok{  \{}
\NormalTok{    op: "Write the general solution with Y(1) = 0",}
\NormalTok{    eq: "u(x,y) = \textbackslash{}\textbackslash{}sum\_\{n=1\}\^{}\{\textbackslash{}\textbackslash{}infty\} C\_n\textbackslash{}\textbackslash{}sin(nx)\textbackslash{}\textbackslash{}sinh(n(1{-}y))",}
\NormalTok{    note: "sinh(n(1−y)) satisfies the top BC: sinh(n(1−1)) = sinh(0) = 0."}
\NormalTok{  \},}
\NormalTok{  \{}
\NormalTok{    op: "Apply bottom BC u(x,0) = sin(x) + (1/3)sin(3x)",}
\NormalTok{    eq: "\textbackslash{}\textbackslash{}sum\_\{n=1\}\^{}\{\textbackslash{}\textbackslash{}infty\} C\_n\textbackslash{}\textbackslash{}sinh(n)\textbackslash{}\textbackslash{}sin(nx) = \textbackslash{}\textbackslash{}sin(x) + \textbackslash{}\textbackslash{}tfrac\{1\}\{3\}\textbackslash{}\textbackslash{}sin(3x)"}
\NormalTok{  \},}
\NormalTok{  \{}
\NormalTok{    op: "Match terms: n = 1 and n = 3 only",}
\NormalTok{    eq: "C\_1\textbackslash{}\textbackslash{}sinh(1) = 1 \textbackslash{}\textbackslash{}implies C\_1 = \textbackslash{}\textbackslash{}frac\{1\}\{\textbackslash{}\textbackslash{}sinh 1\} \textbackslash{}\textbackslash{}approx 0.851"}
\NormalTok{  \},}
\NormalTok{  \{}
\NormalTok{    op: "Match n = 3 term",}
\NormalTok{    eq: "C\_3\textbackslash{}\textbackslash{}sinh(3) = \textbackslash{}\textbackslash{}frac\{1\}\{3\} \textbackslash{}\textbackslash{}implies C\_3 = \textbackslash{}\textbackslash{}frac\{1\}\{3\textbackslash{}\textbackslash{}sinh 3\} \textbackslash{}\textbackslash{}approx 0.033",}
\NormalTok{    note: "sinh(3) ≈ 10.02, so the n=3 term contributes much less near y = 0 and decays faster in y."}
\NormalTok{  \},}
\NormalTok{  \{}
\NormalTok{    op: "Write the exact solution",}
\NormalTok{    eq: "u(x,y) = \textbackslash{}\textbackslash{}frac\{\textbackslash{}\textbackslash{}sin x\textbackslash{}\textbackslash{},\textbackslash{}\textbackslash{}sinh(1{-}y)\}\{\textbackslash{}\textbackslash{}sinh 1\} + \textbackslash{}\textbackslash{}frac\{\textbackslash{}\textbackslash{}sin 3x\textbackslash{}\textbackslash{},\textbackslash{}\textbackslash{}sinh(3(1{-}y))\}\{3\textbackslash{}\textbackslash{}sinh 3\}"}
\NormalTok{  \},}
\NormalTok{  \{}
\NormalTok{    op: "Verify boundary conditions",}
\NormalTok{    eq: "u(x,0) = \textbackslash{}\textbackslash{}frac\{\textbackslash{}\textbackslash{}sin x\}\{1\} + \textbackslash{}\textbackslash{}frac\{\textbackslash{}\textbackslash{}sin 3x\}\{3\} \textbackslash{}\textbackslash{}checkmark \textbackslash{}\textbackslash{}qquad u(x,1) = 0 \textbackslash{}\textbackslash{}checkmark \textbackslash{}\textbackslash{}qquad u(0,y) = u(\textbackslash{}\textbackslash{}pi,y) = 0 \textbackslash{}\textbackslash{}checkmark"}
\NormalTok{  \}}
\NormalTok{])}
\end{Highlighting}
\end{Shaded}

\begin{center}\rule{0.5\linewidth}{0.5pt}\end{center}

\subsection{Exercise 5: Classify these
PDEs}\label{exercise-5-classify-these-pdes}

Classify each of the following equations as elliptic, parabolic, or
hyperbolic. Identify a physical situation each could model.

\textbf{(a)} \(u_{xx} + 4u_{xy} + 4u_{yy} = 0\) \textbf{(b)}
\(3u_{xx} - u_{yy} + u_x = 0\) \textbf{(c)}
\(u_{xx} + u_{yy} + u_{zz} = \rho\)

\begin{Shaded}
\begin{Highlighting}[]
\NormalTok{//| echo: false}
\NormalTok{makeStepperHTML([}
\NormalTok{  \{}
\NormalTok{    op: "Recall: discriminant Δ = B² − 4AC",}
\NormalTok{    eq: "\textbackslash{}\textbackslash{}Delta \textless{} 0 \textbackslash{}\textbackslash{}Rightarrow \textbackslash{}\textbackslash{}text\{elliptic\}, \textbackslash{}\textbackslash{}quad \textbackslash{}\textbackslash{}Delta = 0 \textbackslash{}\textbackslash{}Rightarrow \textbackslash{}\textbackslash{}text\{parabolic\}, \textbackslash{}\textbackslash{}quad \textbackslash{}\textbackslash{}Delta \textgreater{} 0 \textbackslash{}\textbackslash{}Rightarrow \textbackslash{}\textbackslash{}text\{hyperbolic\}"}
\NormalTok{  \},}
\NormalTok{  \{}
\NormalTok{    op: "(a): A=1, B=4, C=4 → Δ = 16 − 16 = 0",}
\NormalTok{    eq: "\textbackslash{}\textbackslash{}text\{Parabolic.\}",}
\NormalTok{    note: "Could model heat diffusion or a diffusion process in a coordinate{-}rotated frame."}
\NormalTok{  \},}
\NormalTok{  \{}
\NormalTok{    op: "(b): A=3, B=0, C=−1 → Δ = 0 − 4(3)(−1) = 12 \textgreater{} 0",}
\NormalTok{    eq: "\textbackslash{}\textbackslash{}text\{Hyperbolic.\}",}
\NormalTok{    note: "Resembles the wave equation; could model wave propagation with a drift term."}
\NormalTok{  \},}
\NormalTok{  \{}
\NormalTok{    op: "(c): 3D Poisson equation u\_xx + u\_yy + u\_zz = ρ",}
\NormalTok{    eq: "\textbackslash{}\textbackslash{}text\{Elliptic (all eigenvalues of the coefficient matrix are positive and same sign).\}",}
\NormalTok{    note: "The 2D discriminant Δ = B²−4AC does not apply to 3D PDEs. Classification in 3D uses the eigenvalues of the leading{-}coefficient matrix; for the Laplacian all three are +1, so the equation is elliptic. Models electrostatics, gravitational potential, or steady{-}state heat in 3D."}
\NormalTok{  \}}
\NormalTok{])}
\end{Highlighting}
\end{Shaded}

\begin{center}\rule{0.5\linewidth}{0.5pt}\end{center}

\subsection{Exercise 6: Non-homogeneous boundary
conditions}\label{exercise-6-non-homogeneous-boundary-conditions}

When the boundary conditions are non-zero --- as when one end of a rod
is held at a fixed nonzero temperature --- separation of variables
cannot be applied directly because the eigenfunctions \(\sin(n\pi x/L)\)
all vanish at the boundaries. The fix is to subtract a
\textbf{steady-state solution} \(v(x)\) that satisfies the boundary
conditions and \(v_{xx} = 0\), so that the remainder \(w = u - v\) has
homogeneous boundary conditions and the standard series method applies.
The final answer is \(u = v + w\): steady state plus decaying transient.

Solve \(u_t = u_{xx}\) on \(0 < x < 1\), \(t > 0\), with \(u(0,t) = 0\),
\(u(1,t) = 1\), and \(u(x,0) = 0\).

\emph{(Steady-state plus transient method.)}

\begin{Shaded}
\begin{Highlighting}[]
\NormalTok{//| echo: false}
\NormalTok{makeStepperHTML([}
\NormalTok{  \{}
\NormalTok{    op: "Find the steady{-}state solution v(x)",}
\NormalTok{    eq: "v\_\{xx\} = 0, \textbackslash{}\textbackslash{}quad v(0) = 0, \textbackslash{}\textbackslash{}quad v(1) = 1 \textbackslash{}\textbackslash{}implies v(x) = x"}
\NormalTok{  \},}
\NormalTok{  \{}
\NormalTok{    op: "Let w(x,t) = u(x,t) − v(x) = u(x,t) − x",}
\NormalTok{    eq: "w\_t = u\_t = u\_\{xx\} = (w+x)\_\{xx\} = w\_\{xx\}",}
\NormalTok{    note: "w satisfies the same heat equation, but now with homogeneous BCs."}
\NormalTok{  \},}
\NormalTok{  \{}
\NormalTok{    op: "Boundary and initial conditions for w",}
\NormalTok{    eq: "w(0,t) = 0, \textbackslash{}\textbackslash{}quad w(1,t) = 0, \textbackslash{}\textbackslash{}quad w(x,0) = u(x,0) {-} x = {-}x"}
\NormalTok{  \},}
\NormalTok{  \{}
\NormalTok{    op: "General solution for w with Fourier sine coefficients of −x",}
\NormalTok{    eq: "B\_n = 2\textbackslash{}\textbackslash{}int\_0\^{}1 ({-}x)\textbackslash{}\textbackslash{}sin(n\textbackslash{}\textbackslash{}pi x)\textbackslash{}\textbackslash{},dx = \textbackslash{}\textbackslash{}frac\{2({-}1)\^{}n\}\{n\textbackslash{}\textbackslash{}pi\}",}
\NormalTok{    note: "Using the standard result ∫₀¹ x sin(nπx) dx = (−1)\^{}\{n+1\}/(nπ)."}
\NormalTok{  \},}
\NormalTok{  \{}
\NormalTok{    op: "Assemble the full solution u = v + w",}
\NormalTok{    eq: "u(x,t) = x + \textbackslash{}\textbackslash{}frac\{2\}\{\textbackslash{}\textbackslash{}pi\}\textbackslash{}\textbackslash{}sum\_\{n=1\}\^{}\{\textbackslash{}\textbackslash{}infty\}\textbackslash{}\textbackslash{}frac\{({-}1)\^{}n\}\{n\}\textbackslash{}\textbackslash{}sin(n\textbackslash{}\textbackslash{}pi x)\textbackslash{}\textbackslash{},e\^{}\{{-}n\^{}2\textbackslash{}\textbackslash{}pi\^{}2 t\}",}
\NormalTok{    note: "At large t, the transient (w) decays and u approaches the steady state x — a rod uniformly heated from one end."}
\NormalTok{  \}}
\NormalTok{])}
\end{Highlighting}
\end{Shaded}

\part{Volume 7 --- Complex analysis}

\chapter{Complex functions}\label{complex-functions}

Differentiability in the complex plane and the Cauchy-Riemann equations

\hfill\break

A function \(f\) of a complex variable maps complex numbers to complex
numbers. Writing \(z = x + iy\) and separating
\(f(z) = u(x,y) + iv(x,y)\), the function is really a pair of
real-valued functions of two real variables. The question is: what extra
condition on \(u\) and \(v\) makes \(f\) genuinely
complex-differentiable? The answer --- the Cauchy-Riemann equations ---
is the key that unlocks the entire subject.

\begin{center}\rule{0.5\linewidth}{0.5pt}\end{center}

\section{Complex numbers: brief
review}\label{complex-numbers-brief-review}

Every complex number \(z = x + iy\) is a point in the \textbf{Argand
plane}, with real part \(x\) on the horizontal axis and imaginary part
\(y\) on the vertical.

\textbf{Polar form.} \(z = r e^{i\theta}\), where
\(r = |z| = \sqrt{x^2+y^2}\) is the modulus and \(\theta = \arg z\) is
the argument --- the angle from the positive real axis, measured
counterclockwise. Euler's formula
\(e^{i\theta} = \cos\theta + i\sin\theta\) connects the polar and
rectangular forms.

\textbf{De Moivre's theorem.} \((re^{i\theta})^n = r^n e^{in\theta}\).
Integer powers of complex numbers are most efficiently computed in polar
form.

\textbf{Conjugate.} \(\bar{z} = x - iy\). Identities:
\(z\bar{z} = |z|^2\); \(\text{Re}(z) = (z+\bar{z})/2\);
\(\text{Im}(z) = (z-\bar{z})/(2i)\).

\begin{Shaded}
\begin{Highlighting}[]
\NormalTok{//| echo: false}
\NormalTok{\{}
\NormalTok{  const W = 420, H = 420, CX = W/2, CY = H/2, SCALE = 80;}

\NormalTok{  const rSlider = Inputs.range([0.5, 2.0], \{value: 1.4, step: 0.05, label: "|w| = r"\});}
\NormalTok{  const thetaSlider = Inputs.range([0, 2*Math.PI], \{value: Math.PI/4, step: Math.PI/12, label: "arg(w) = θ"\});}

\NormalTok{  const svgNS = "http://www.w3.org/2000/svg";}

\NormalTok{  function makeViz(r, theta) \{}
\NormalTok{    // z starts at (1.2, 0.6) in complex coords}
\NormalTok{    let zx = 1.2, zy = 0.6;}
\NormalTok{    // w = r * e\^{}\{i theta\}}
\NormalTok{    let wx = r * Math.cos(theta), wy = r * Math.sin(theta);}
\NormalTok{    // zw = z * w}
\NormalTok{    let zwx = zx*wx {-} zy*wy, zwy = zx*wy + zy*wx;}

\NormalTok{    function toScreen(cx, cy) \{}
\NormalTok{      return [CX + cx*SCALE, CY {-} cy*SCALE];}
\NormalTok{    \}}

\NormalTok{    const [zsx, zsy] = toScreen(zx, zy);}
\NormalTok{    const [zwsx, zwsy] = toScreen(zwx, zwy);}

\NormalTok{    const svg = document.createElementNS(svgNS, "svg");}
\NormalTok{    svg.setAttribute("width", W);}
\NormalTok{    svg.setAttribute("height", H);}
\NormalTok{    svg.style.cssText = "background:\#fafafa; border:1px solid \#e5e7eb; border{-}radius:6px; display:block;";}

\NormalTok{    // Grid}
\NormalTok{    const gridG = document.createElementNS(svgNS, "g");}
\NormalTok{    gridG.setAttribute("stroke", "\#e5e7eb");}
\NormalTok{    gridG.setAttribute("stroke{-}width", "0.5");}
\NormalTok{    for (let i = {-}4; i \textless{}= 4; i++) \{}
\NormalTok{      const xLine = document.createElementNS(svgNS, "line");}
\NormalTok{      const [lx] = toScreen(i, 0); const [,ly1] = toScreen(0, {-}4); const [,ly2] = toScreen(0, 4);}
\NormalTok{      xLine.setAttribute("x1", lx); xLine.setAttribute("y1", ly1); xLine.setAttribute("x2", lx); xLine.setAttribute("y2", ly2);}
\NormalTok{      gridG.appendChild(xLine);}
\NormalTok{      const yLine = document.createElementNS(svgNS, "line");}
\NormalTok{      const [lx1] = toScreen({-}4, 0); const [lx2] = toScreen(4, 0); const [,ley] = toScreen(0, i);}
\NormalTok{      yLine.setAttribute("x1", lx1); yLine.setAttribute("y1", ley); yLine.setAttribute("x2", lx2); yLine.setAttribute("y2", ley);}
\NormalTok{      gridG.appendChild(yLine);}
\NormalTok{    \}}
\NormalTok{    svg.appendChild(gridG);}

\NormalTok{    // Axes}
\NormalTok{    const axG = document.createElementNS(svgNS, "g");}
\NormalTok{    axG.setAttribute("stroke", "\#9ca3af"); axG.setAttribute("stroke{-}width", "1");}
\NormalTok{    const xax = document.createElementNS(svgNS, "line");}
\NormalTok{    const [ax1] = toScreen({-}3.8, 0); const [ax2] = toScreen(3.8, 0); const [,axy] = toScreen(0, 0);}
\NormalTok{    xax.setAttribute("x1", ax1); xax.setAttribute("y1", axy); xax.setAttribute("x2", ax2); xax.setAttribute("y2", axy);}
\NormalTok{    const yax = document.createElementNS(svgNS, "line");}
\NormalTok{    const [ayx] = toScreen(0, 0); const [,ay1] = toScreen(0, {-}3.8); const [,ay2] = toScreen(0, 3.8);}
\NormalTok{    yax.setAttribute("x1", ayx); yax.setAttribute("y1", ay1); yax.setAttribute("x2", ayx); yax.setAttribute("y2", ay2);}
\NormalTok{    axG.appendChild(xax); axG.appendChild(yax);}
\NormalTok{    svg.appendChild(axG);}

\NormalTok{    // Radial arms from origin to z and zw}
\NormalTok{    function makeLine(x1,y1,x2,y2,color,dash) \{}
\NormalTok{      const l = document.createElementNS(svgNS, "line");}
\NormalTok{      l.setAttribute("x1",x1); l.setAttribute("y1",y1); l.setAttribute("x2",x2); l.setAttribute("y2",y2);}
\NormalTok{      l.setAttribute("stroke", color); l.setAttribute("stroke{-}width","1.5");}
\NormalTok{      if (dash) l.setAttribute("stroke{-}dasharray", dash);}
\NormalTok{      return l;}
\NormalTok{    \}}
\NormalTok{    const [ox, oy] = toScreen(0,0);}
\NormalTok{    svg.appendChild(makeLine(ox, oy, zsx, zsy, "\#f97316", ""));}
\NormalTok{    svg.appendChild(makeLine(ox, oy, zwsx, zwsy, "\#3b82f6", ""));}

\NormalTok{    // Arc showing rotation angle theta between z and zw arms}
\NormalTok{    const zAngle = Math.atan2({-}zsy + CY, zsx {-} CX); // angle of z from screen center}
\NormalTok{    const zwAngle = Math.atan2({-}zwsy + CY, zwsx {-} CX);}
\NormalTok{    const arcR = 30;}
\NormalTok{    const arcPath = document.createElementNS(svgNS, "path");}
\NormalTok{    // arc from zAngle to zwAngle}
\NormalTok{    const ax0 = CX + arcR * Math.cos(zAngle), ay0 = CY + arcR * Math.sin(zAngle);}
\NormalTok{    const ax1e = CX + arcR * Math.cos(zwAngle), ay1e = CY + arcR * Math.sin(zwAngle);}
\NormalTok{    // Choose large{-}arc{-}flag based on angular difference}
\NormalTok{    let dang = zwAngle {-} zAngle;}
\NormalTok{    while (dang \textgreater{} Math.PI) dang {-}= 2*Math.PI;}
\NormalTok{    while (dang \textless{} {-}Math.PI) dang += 2*Math.PI;}
\NormalTok{    const laf = Math.abs(dang) \textgreater{} Math.PI ? 1 : 0;}
\NormalTok{    const swf = dang \textgreater{} 0 ? 1 : 0;}
\NormalTok{    arcPath.setAttribute("d", \textasciigrave{}M $\{ax0\} $\{ay0\} A $\{arcR\} $\{arcR\} 0 $\{laf\} $\{swf\} $\{ax1e\} $\{ay1e\}\textasciigrave{});}
\NormalTok{    arcPath.setAttribute("fill","none"); arcPath.setAttribute("stroke","\#8b5cf6");}
\NormalTok{    arcPath.setAttribute("stroke{-}width","1.5");}
\NormalTok{    svg.appendChild(arcPath);}

\NormalTok{    // Scale indicator: show |zw| / |z| = r label}
\NormalTok{    const zmag = Math.sqrt(zx*zx + zy*zy);}
\NormalTok{    const zwmag = Math.sqrt(zwx*zwx + zwy*zwy);}

\NormalTok{    // Points}
\NormalTok{    function makeCircle(px, py, radius, fill, stroke) \{}
\NormalTok{      const c = document.createElementNS(svgNS, "circle");}
\NormalTok{      c.setAttribute("cx", px); c.setAttribute("cy", py); c.setAttribute("r", radius);}
\NormalTok{      c.setAttribute("fill", fill); c.setAttribute("stroke", stroke); c.setAttribute("stroke{-}width","1.5");}
\NormalTok{      return c;}
\NormalTok{    \}}
\NormalTok{    svg.appendChild(makeCircle(zsx, zsy, 7, "\#f97316", "\#fff"));}
\NormalTok{    svg.appendChild(makeCircle(zwsx, zwsy, 7, "\#3b82f6", "\#fff"));}
\NormalTok{    svg.appendChild(makeCircle(ox, oy, 4, "\#6b7280", "\#fff"));}

\NormalTok{    // Labels}
\NormalTok{    function makeText(x, y, txt, color, size) \{}
\NormalTok{      const t = document.createElementNS(svgNS, "text");}
\NormalTok{      t.setAttribute("x", x); t.setAttribute("y", y);}
\NormalTok{      t.setAttribute("fill", color); t.setAttribute("font{-}size", size || "12");}
\NormalTok{      t.setAttribute("font{-}family", "sans{-}serif");}
\NormalTok{      t.textContent = txt;}
\NormalTok{      return t;}
\NormalTok{    \}}
\NormalTok{    svg.appendChild(makeText(zsx+8, zsy{-}6, "z", "\#f97316", "13"));}
\NormalTok{    svg.appendChild(makeText(zwsx+8, zwsy{-}6, "zw", "\#3b82f6", "13"));}
\NormalTok{    svg.appendChild(makeText(ox+4, oy{-}5, "O", "\#6b7280", "11"));}

\NormalTok{    // Annotation}
\NormalTok{    const midArcX = CX + arcR * Math.cos((zAngle+zwAngle)/2);}
\NormalTok{    const midArcY = CY + arcR * Math.sin((zAngle+zwAngle)/2);}
\NormalTok{    svg.appendChild(makeText(midArcX {-} 5, midArcY {-} 8, "\textbackslash{}u03b8", "\#8b5cf6", "12"));}

\NormalTok{    // Info panel}
\NormalTok{    const info = document.createElement("div");}
\NormalTok{    info.style.cssText = "margin{-}top:0.5rem; font{-}size:0.85em; color:\#374151; font{-}family:sans{-}serif; display:flex; gap:1.5rem; flex{-}wrap:wrap;";}
\NormalTok{    const thetaDeg = (theta * 180 / Math.PI).toFixed(0);}
\NormalTok{    info.innerHTML = \textasciigrave{}}
\NormalTok{      \textless{}span style="color:\#f97316"\textgreater{}|z| = $\{zmag.toFixed(2)\}\textless{}/span\textgreater{}}
\NormalTok{      \textless{}span style="color:\#3b82f6"\textgreater{}|zw| = $\{zwmag.toFixed(2)\}\textless{}/span\textgreater{}}
\NormalTok{      \textless{}span style="color:\#8b5cf6"\textgreater{}arg(w) = $\{thetaDeg\}°\textless{}/span\textgreater{}}
\NormalTok{      \textless{}span style="color:\#6b7280"\textgreater{}|zw|/|z| = $\{r.toFixed(2)\} = r\textless{}/span\textgreater{}}
\NormalTok{    \textasciigrave{};}

\NormalTok{    const note = document.createElement("div");}
\NormalTok{    note.style.cssText = "margin{-}top:0.4rem; font{-}size:0.82em; color:\#6b7280; font{-}style:italic;";}
\NormalTok{    note.textContent = \textasciigrave{}Multiplying by w = r·e\^{}\{iθ\} rotates z by θ = $\{thetaDeg\}° and scales by r = $\{r.toFixed(2)\}. Orange point z, blue point zw. Set r = 1 for pure rotation; θ = 0 for pure scaling.\textasciigrave{};}

\NormalTok{    const wrap = document.createElement("div");}
\NormalTok{    wrap.style.cssText = "margin:1rem 0;";}
\NormalTok{    wrap.appendChild(svg);}
\NormalTok{    wrap.appendChild(info);}
\NormalTok{    wrap.appendChild(note);}
\NormalTok{    return wrap;}
\NormalTok{  \}}

\NormalTok{  const container = document.createElement("div");}
\NormalTok{  container.style.cssText = "border:1px solid \#e5e7eb; border{-}radius:8px; padding:1rem; margin:1rem 0;";}

\NormalTok{  const title = document.createElement("div");}
\NormalTok{  title.style.cssText = "font{-}weight:600; margin{-}bottom:0.5rem; font{-}family:sans{-}serif;";}
\NormalTok{  title.textContent = "Complex multiplication as rotation and scaling";}
\NormalTok{  container.appendChild(title);}

\NormalTok{  container.appendChild(rSlider);}
\NormalTok{  container.appendChild(thetaSlider);}

\NormalTok{  const vizContainer = document.createElement("div");}
\NormalTok{  container.appendChild(vizContainer);}

\NormalTok{  function update() \{}
\NormalTok{    const r = rSlider.value;}
\NormalTok{    const theta = thetaSlider.value;}
\NormalTok{    vizContainer.innerHTML = "";}
\NormalTok{    vizContainer.appendChild(makeViz(r, theta));}
\NormalTok{  \}}

\NormalTok{  rSlider.addEventListener("input", update);}
\NormalTok{  thetaSlider.addEventListener("input", update);}
\NormalTok{  update();}

\NormalTok{  return container;}
\NormalTok{\}}
\end{Highlighting}
\end{Shaded}

\begin{center}\rule{0.5\linewidth}{0.5pt}\end{center}

\section{Complex functions}\label{complex-functions-1}

A complex function \(f: \mathbb{C} \to \mathbb{C}\) maps each
\(z = x+iy\) to a complex value \(f(z) = u(x,y) + iv(x,y)\).

\textbf{Example.} \(f(z) = z^2 = (x+iy)^2 = (x^2-y^2) + 2ixy\), so
\(u = x^2-y^2\), \(v = 2xy\).

\textbf{Example.} \(f(z) = \bar{z} = x - iy\), so \(u = x\), \(v = -y\).
This function will turn out to be nowhere differentiable in the complex
sense.

Limits and continuity carry over from real analysis: \(f\) is
\textbf{continuous at \(z_0\)} if \(\lim_{z\to z_0}f(z) = f(z_0)\),
where the limit must be the same from every direction of approach.

\begin{center}\rule{0.5\linewidth}{0.5pt}\end{center}

\section{Complex differentiability}\label{complex-differentiability}

The \textbf{derivative} of \(f\) at \(z_0\) is:

\[f'(z_0) = \lim_{z \to z_0} \frac{f(z) - f(z_0)}{z - z_0}\]

The limit must exist and agree no matter how \(z\) approaches \(z_0\)
--- a far stronger requirement than real differentiability.

\subsection{Deriving the Cauchy-Riemann
equations}\label{deriving-the-cauchy-riemann-equations}

Approach \(z_0 = x_0+iy_0\) along two different directions and demand
the limit is the same.

\textbf{Horizontal approach} (\(z = x+iy_0\), \(z-z_0 = \Delta x\)):

With \(\Delta x\) real, the difference quotient separates cleanly into
real and imaginary parts:
\[\frac{f(z)-f(z_0)}{\Delta x} = \frac{[u(x_0+\Delta x,\,y_0)-u(x_0,y_0)]}{\Delta x} + i\frac{[v(x_0+\Delta x,\,y_0)-v(x_0,y_0)]}{\Delta x}\]

Taking \(\Delta x \to 0\), each fraction is just a partial derivative:
\[f'(z_0) = u_x + iv_x\]

\textbf{Vertical approach} (\(z = x_0+iy\), \(z-z_0 = i\Delta y\)):

Now the increment is purely imaginary, so we divide by \(i\Delta y\).
Dividing by \(i\) is the same as multiplying by \(-i\), which swaps real
and imaginary parts. Carrying through:
\[\frac{f(z)-f(z_0)}{i\Delta y} = \frac{1}{i}\cdot\frac{[u(x_0,y_0+\Delta y)-u(x_0,y_0)]+i[v(x_0,y_0+\Delta y)-v(x_0,y_0)]}{\Delta y}\]

Using \(1/i = -i\) and taking \(\Delta y \to 0\):
\[f'(z_0) = v_y - iu_y\]

Equating real and imaginary parts:

\[\boxed{u_x = v_y \qquad u_y = -v_x}\]

These are the \textbf{Cauchy-Riemann (C-R) equations}. They are
necessary for complex differentiability. If the partial derivatives are
continuous in a neighbourhood of \(z_0\), they are also sufficient.

\textbf{Consequence.} For \(f(z) = \bar{z}\): \(u_x = 1\) but
\(v_y = -1\). The C-R equations fail everywhere, so \(\bar{z}\) is
nowhere differentiable --- even though its real and imaginary parts are
smooth real functions.

When the C-R equations hold, the derivative can be written in any of the
equivalent forms: \[f'(z) = u_x + iv_x = v_y - iu_y\]

\begin{center}\rule{0.5\linewidth}{0.5pt}\end{center}

\section{Analytic functions}\label{analytic-functions}

\(f\) is \textbf{analytic at \(z_0\)} if it is differentiable in some
open disk centred at \(z_0\). It is \textbf{analytic in a domain \(D\)}
if differentiable at every point of \(D\). A function analytic on all of
\(\mathbb{C}\) is called \textbf{entire}.

Being differentiable at an isolated point does not make a function
analytic --- analyticity requires differentiability in an open
neighbourhood. Exercise 4 shows this with \(f(z) = |z|^2\), which is
differentiable only at the origin and analytic nowhere.

\subsection{Harmonic functions}\label{harmonic-functions}

If \(f = u + iv\) is analytic, both \(u\) and \(v\) satisfy
\textbf{Laplace's equation}:
\[\nabla^2 u = u_{xx} + u_{yy} = 0, \qquad \nabla^2 v = 0\]

\emph{Proof.} From the C-R equations:
\(u_{xx} = (u_x)_x = (v_y)_x = v_{yx}\) and
\(u_{yy} = (u_y)_y = (-v_x)_y = -v_{xy}\). Since mixed partials are
equal, \(u_{xx} + u_{yy} = 0\).

A function satisfying Laplace's equation is \textbf{harmonic}. So: the
real and imaginary parts of any analytic function are harmonic.

Conversely, given a harmonic \(u\) on a \textbf{simply connected domain}
--- one with no holes, so that any closed loop inside it can be
continuously shrunk to a point --- there exists a \textbf{harmonic
conjugate} \(v\) (unique up to a constant) such that \(f = u + iv\) is
analytic. Finding \(v\) from \(u\) via the C-R equations is a standard
technique.

\begin{center}\rule{0.5\linewidth}{0.5pt}\end{center}

\section{Elementary analytic
functions}\label{elementary-analytic-functions}

\subsection{Exponential}\label{exponential}

\[e^z = e^x(\cos y + i\sin y), \qquad \frac{d}{dz}e^z = e^z\]

Entire. Periodic with period \(2\pi i\): \(e^{z+2\pi i} = e^z\).

\subsection{Trigonometric}\label{trigonometric}

\[\cos z = \frac{e^{iz}+e^{-iz}}{2}, \qquad \sin z = \frac{e^{iz}-e^{-iz}}{2i}\]

Both entire. The identities \(\cos^2 z + \sin^2 z = 1\) and all standard
trig identities hold. For purely imaginary argument:
\(\cos(iy) = \cosh y\), \(\sin(iy) = i\sinh y\) --- the complex trig and
hyperbolic functions are closely related.

\subsection{Logarithm}\label{logarithm}

\[\log z = \ln|z| + i\arg z = \ln r + i\theta\]

\textbf{Multivalued}: \(\arg z\) is defined only up to multiples of
\(2\pi\). The \textbf{principal value}
\(\text{Log}\,z = \ln|z| + i\,\text{Arg}\,z\) uses
\(\text{Arg}\,z \in (-\pi, \pi]\) and is analytic everywhere except on
the \textbf{branch cut} along the negative real axis --- the line where
\(\text{Arg}\,z\) jumps discontinuously by \(2\pi\) as you cross it.
Without the cut, \(\text{Log}\,z\) cannot be made single-valued and
continuous.

\begin{Shaded}
\begin{Highlighting}[]
\NormalTok{//| echo: false}
\NormalTok{\{}
\NormalTok{  const W = 400, H = 360, CX = 180, CY = 180, SCALE = 65;}
\NormalTok{  const svgNS = "http://www.w3.org/2000/svg";}

\NormalTok{  // Preset path configs: array of [re, im] control points defining a closed loop}
\NormalTok{  const paths = \{}
\NormalTok{    "Loop around origin (winding = +1)": [}
\NormalTok{      [1.5, 0], [1.2, 1.2], [0, 1.8], [{-}1.2, 1.2], [{-}1.5, 0],}
\NormalTok{      [{-}1.2, {-}1.2], [0, {-}1.8], [1.2, {-}1.2]}
\NormalTok{    ],}
\NormalTok{    "Loop missing origin (winding = 0)": [}
\NormalTok{      [1.0, 0.8], [1.8, 0.4], [2.0, 1.2], [1.4, 1.8], [0.6, 1.4]}
\NormalTok{    ],}
\NormalTok{    "Double loop (winding = +2)": [}
\NormalTok{      [1.8, 0], [1.4, 1.0], [0, 1.8], [{-}1.4, 1.0], [{-}1.8, 0],}
\NormalTok{      [{-}1.4, {-}1.0], [0, {-}1.8], [1.4, {-}1.0],}
\NormalTok{      [1.2, 0.3], [0.5, 0.8], [{-}0.5, 0.5], [{-}0.8, 0], [{-}0.5, {-}0.5],}
\NormalTok{      [0.5, {-}0.8], [1.2, {-}0.3]}
\NormalTok{    ]}
\NormalTok{  \};}

\NormalTok{  const pathSelector = Inputs.select(Object.keys(paths), \{}
\NormalTok{    value: Object.keys(paths)[0],}
\NormalTok{    label: "Path"}
\NormalTok{  \});}
\NormalTok{  const branchToggle = Inputs.toggle(\{label: "Show branch cut", value: false\});}
\NormalTok{  const tSlider = Inputs.range([0, 1], \{value: 0, step: 0.01, label: "Progress t"\});}

\NormalTok{  function catmullRomPoint(pts, t) \{}
\NormalTok{    const n = pts.length;}
\NormalTok{    const totalT = n;}
\NormalTok{    const tScaled = t * n;}
\NormalTok{    const i = Math.floor(tScaled) \% n;}
\NormalTok{    const localT = tScaled {-} Math.floor(tScaled);}

\NormalTok{    const p0 = pts[(i {-} 1 + n) \% n];}
\NormalTok{    const p1 = pts[i];}
\NormalTok{    const p2 = pts[(i + 1) \% n];}
\NormalTok{    const p3 = pts[(i + 2) \% n];}

\NormalTok{    const t2 = localT*localT, t3 = localT*t2;}
\NormalTok{    const re = 0.5*((2*p1[0]) + ({-}p0[0]+p2[0])*localT + (2*p0[0]{-}5*p1[0]+4*p2[0]{-}p3[0])*t2 + ({-}p0[0]+3*p1[0]{-}3*p2[0]+p3[0])*t3);}
\NormalTok{    const im = 0.5*((2*p1[1]) + ({-}p0[1]+p2[1])*localT + (2*p0[1]{-}5*p1[1]+4*p2[1]{-}p3[1])*t2 + ({-}p0[1]+3*p1[1]{-}3*p2[1]+p3[1])*t3);}
\NormalTok{    return [re, im];}
\NormalTok{  \}}

\NormalTok{  function computeWindingAndArgTrace(pts, numSteps) \{}
\NormalTok{    let totalDeltaArg = 0;}
\NormalTok{    let prevArg = null;}
\NormalTok{    const trace = [];}

\NormalTok{    for (let step = 0; step \textless{}= numSteps; step++) \{}
\NormalTok{      const t = step / numSteps;}
\NormalTok{      const [re, im] = catmullRomPoint(pts, t);}
\NormalTok{      const arg = Math.atan2(im, re);}

\NormalTok{      if (prevArg !== null) \{}
\NormalTok{        let dArg = arg {-} prevArg;}
\NormalTok{        while (dArg \textgreater{} Math.PI) dArg {-}= 2*Math.PI;}
\NormalTok{        while (dArg \textless{} {-}Math.PI) dArg += 2*Math.PI;}
\NormalTok{        totalDeltaArg += dArg;}
\NormalTok{      \}}
\NormalTok{      prevArg = arg;}
\NormalTok{      trace.push([t, totalDeltaArg]);}
\NormalTok{    \}}
\NormalTok{    return trace;}
\NormalTok{  \}}

\NormalTok{  function toScreen(re, im) \{}
\NormalTok{    return [CX + re*SCALE, CY {-} im*SCALE];}
\NormalTok{  \}}

\NormalTok{  function render(pathName, showBranch, tVal) \{}
\NormalTok{    const pts = paths[pathName];}
\NormalTok{    const STEPS = 300;}
\NormalTok{    const argTrace = computeWindingAndArgTrace(pts, STEPS);}
\NormalTok{    const windingAtEnd = argTrace[STEPS][1] / (2*Math.PI);}
\NormalTok{    const windingRounded = Math.round(windingAtEnd);}

\NormalTok{    // SVG}
\NormalTok{    const svg = document.createElementNS(svgNS, "svg");}
\NormalTok{    svg.setAttribute("width", W); svg.setAttribute("height", H);}
\NormalTok{    svg.style.cssText = "background:\#fafafa; border:1px solid \#e5e7eb; border{-}radius:6px; display:block;";}

\NormalTok{    // Grid}
\NormalTok{    const gridG = document.createElementNS(svgNS, "g");}
\NormalTok{    gridG.setAttribute("stroke", "\#e5e7eb"); gridG.setAttribute("stroke{-}width", "0.5");}
\NormalTok{    for (let i={-}3; i\textless{}=3; i++) \{}
\NormalTok{      const xl = document.createElementNS(svgNS,"line");}
\NormalTok{      const [lx] = toScreen(i,0); const [,ly1] = toScreen(0,{-}3); const [,ly2] = toScreen(0,3);}
\NormalTok{      xl.setAttribute("x1",lx); xl.setAttribute("y1",ly1); xl.setAttribute("x2",lx); xl.setAttribute("y2",ly2);}
\NormalTok{      gridG.appendChild(xl);}
\NormalTok{      const yl = document.createElementNS(svgNS,"line");}
\NormalTok{      const [lx1] = toScreen({-}3,0); const [lx2] = toScreen(3,0); const [,ly] = toScreen(0,i);}
\NormalTok{      yl.setAttribute("x1",lx1); yl.setAttribute("y1",ly); yl.setAttribute("x2",lx2); yl.setAttribute("y2",ly);}
\NormalTok{      gridG.appendChild(yl);}
\NormalTok{    \}}
\NormalTok{    svg.appendChild(gridG);}

\NormalTok{    // Axes}
\NormalTok{    const [ox, oy] = toScreen(0,0);}
\NormalTok{    const axG = document.createElementNS(svgNS,"g"); axG.setAttribute("stroke","\#9ca3af"); axG.setAttribute("stroke{-}width","1");}
\NormalTok{    const xa = document.createElementNS(svgNS,"line");}
\NormalTok{    const [ax1] = toScreen({-}3.2,0); const [ax2] = toScreen(3.2,0);}
\NormalTok{    xa.setAttribute("x1",ax1); xa.setAttribute("y1",oy); xa.setAttribute("x2",ax2); xa.setAttribute("y2",oy);}
\NormalTok{    const ya = document.createElementNS(svgNS,"line");}
\NormalTok{    const [,ay1] = toScreen(0,{-}2.6); const [,ay2] = toScreen(0,2.6);}
\NormalTok{    ya.setAttribute("x1",ox); ya.setAttribute("y1",ay1); ya.setAttribute("x2",ox); ya.setAttribute("y2",ay2);}
\NormalTok{    axG.appendChild(xa); axG.appendChild(ya);}
\NormalTok{    svg.appendChild(axG);}

\NormalTok{    // Branch cut: negative real axis}
\NormalTok{    if (showBranch) \{}
\NormalTok{      const bc = document.createElementNS(svgNS,"line");}
\NormalTok{      const [bx1] = toScreen({-}3.2,0); const [bx2] = toScreen(0,0);}
\NormalTok{      bc.setAttribute("x1",bx1); bc.setAttribute("y1",oy); bc.setAttribute("x2",bx2); bc.setAttribute("y2",oy);}
\NormalTok{      bc.setAttribute("stroke","\#ef4444"); bc.setAttribute("stroke{-}width","2.5");}
\NormalTok{      bc.setAttribute("stroke{-}dasharray","6,3");}
\NormalTok{      svg.appendChild(bc);}
\NormalTok{      const bcLabel = document.createElementNS(svgNS,"text");}
\NormalTok{      const [blx] = toScreen({-}2,0);}
\NormalTok{      bcLabel.setAttribute("x",blx); bcLabel.setAttribute("y",oy{-}6);}
\NormalTok{      bcLabel.setAttribute("fill","\#ef4444"); bcLabel.setAttribute("font{-}size","10");}
\NormalTok{      bcLabel.setAttribute("font{-}family","sans{-}serif"); bcLabel.textContent="branch cut";}
\NormalTok{      svg.appendChild(bcLabel);}
\NormalTok{    \}}

\NormalTok{    // Full path}
\NormalTok{    let pathD = "";}
\NormalTok{    for (let s=0; s\textless{}=STEPS; s++) \{}
\NormalTok{      const t = s/STEPS;}
\NormalTok{      const [re,im] = catmullRomPoint(pts, t);}
\NormalTok{      const [px,py] = toScreen(re,im);}
\NormalTok{      pathD += (s===0 ? \textasciigrave{}M $\{px\} $\{py\}\textasciigrave{} : \textasciigrave{} L $\{px\} $\{py\}\textasciigrave{});}
\NormalTok{    \}}
\NormalTok{    pathD += " Z";}
\NormalTok{    const pathEl = document.createElementNS(svgNS,"path");}
\NormalTok{    pathEl.setAttribute("d",pathD); pathEl.setAttribute("fill","none");}
\NormalTok{    pathEl.setAttribute("stroke","\#f97316"); pathEl.setAttribute("stroke{-}width","2");}
\NormalTok{    svg.appendChild(pathEl);}

\NormalTok{    // Animated point}
\NormalTok{    const currentStep = Math.round(tVal * STEPS);}
\NormalTok{    const [cre,cim] = catmullRomPoint(pts, tVal);}
\NormalTok{    const [cpx,cpy] = toScreen(cre,cim);}
\NormalTok{    const pt = document.createElementNS(svgNS,"circle");}
\NormalTok{    pt.setAttribute("cx",cpx); pt.setAttribute("cy",cpy); pt.setAttribute("r","6");}
\NormalTok{    pt.setAttribute("fill","\#3b82f6"); pt.setAttribute("stroke","\#fff"); pt.setAttribute("stroke{-}width","2");}
\NormalTok{    svg.appendChild(pt);}

\NormalTok{    // Origin mark}
\NormalTok{    const origPt = document.createElementNS(svgNS,"circle");}
\NormalTok{    origPt.setAttribute("cx",ox); origPt.setAttribute("cy",oy); origPt.setAttribute("r","4");}
\NormalTok{    origPt.setAttribute("fill","\#1f2937"); origPt.setAttribute("stroke","\#fff"); origPt.setAttribute("stroke{-}width","1.5");}
\NormalTok{    svg.appendChild(origPt);}
\NormalTok{    const origLabel = document.createElementNS(svgNS,"text");}
\NormalTok{    origLabel.setAttribute("x",ox+5); origLabel.setAttribute("y",oy{-}6);}
\NormalTok{    origLabel.setAttribute("fill","\#1f2937"); origLabel.setAttribute("font{-}size","10");}
\NormalTok{    origLabel.setAttribute("font{-}family","sans{-}serif"); origLabel.textContent="O";}
\NormalTok{    svg.appendChild(origLabel);}

\NormalTok{    // Winding number display}
\NormalTok{    const windBox = document.createElementNS(svgNS,"rect");}
\NormalTok{    windBox.setAttribute("x",W{-}95); windBox.setAttribute("y",5);}
\NormalTok{    windBox.setAttribute("width",88); windBox.setAttribute("height",44);}
\NormalTok{    windBox.setAttribute("fill","\#1e40af"); windBox.setAttribute("rx","5");}
\NormalTok{    svg.appendChild(windBox);}
\NormalTok{    const windLabel = document.createElementNS(svgNS,"text");}
\NormalTok{    windLabel.setAttribute("x",W{-}51); windLabel.setAttribute("y",22);}
\NormalTok{    windLabel.setAttribute("fill","white"); windLabel.setAttribute("font{-}size","11");}
\NormalTok{    windLabel.setAttribute("text{-}anchor","middle"); windLabel.setAttribute("font{-}family","sans{-}serif");}
\NormalTok{    windLabel.textContent = "winding \#";}
\NormalTok{    svg.appendChild(windLabel);}
\NormalTok{    const windNum = document.createElementNS(svgNS,"text");}
\NormalTok{    windNum.setAttribute("x",W{-}51); windNum.setAttribute("y",41);}
\NormalTok{    windNum.setAttribute("fill","white"); windNum.setAttribute("font{-}size","18");}
\NormalTok{    windNum.setAttribute("text{-}anchor","middle"); windNum.setAttribute("font{-}family","sans{-}serif");}
\NormalTok{    windNum.setAttribute("font{-}weight","bold");}
\NormalTok{    windNum.textContent = (tVal \textless{} 0.99) ? "..." : windingRounded.toString();}
\NormalTok{    svg.appendChild(windNum);}

\NormalTok{    // Arg trace mini{-}chart}
\NormalTok{    const chartX = 0, chartY = H {-} 70, chartW = W, chartH = 65;}
\NormalTok{    const chartBg = document.createElementNS(svgNS,"rect");}
\NormalTok{    chartBg.setAttribute("x",chartX); chartBg.setAttribute("y",chartY);}
\NormalTok{    chartBg.setAttribute("width",chartW); chartBg.setAttribute("height",chartH);}
\NormalTok{    chartBg.setAttribute("fill","rgba(241,245,249,0.9)");}
\NormalTok{    svg.appendChild(chartBg);}

\NormalTok{    const maxArgAbs = Math.max(2*Math.PI*1.5, Math.abs(argTrace[STEPS][1]) * 1.2);}
\NormalTok{    function argToChart(t, argVal) \{}
\NormalTok{      const cx = chartX + t * chartW;}
\NormalTok{      const cy = chartY + chartH/2 {-} (argVal / maxArgAbs) * (chartH/2 {-} 6);}
\NormalTok{      return [cx, cy];}
\NormalTok{    \}}

\NormalTok{    // Zero line}
\NormalTok{    const [zx1, zy1] = argToChart(0, 0); const [zx2] = argToChart(1, 0);}
\NormalTok{    const zeroLine = document.createElementNS(svgNS,"line");}
\NormalTok{    zeroLine.setAttribute("x1",zx1); zeroLine.setAttribute("y1",zy1);}
\NormalTok{    zeroLine.setAttribute("x2",zx2); zeroLine.setAttribute("y2",zy1);}
\NormalTok{    zeroLine.setAttribute("stroke","\#9ca3af"); zeroLine.setAttribute("stroke{-}width","0.8");}
\NormalTok{    zeroLine.setAttribute("stroke{-}dasharray","3,3");}
\NormalTok{    svg.appendChild(zeroLine);}

\NormalTok{    // Arg trace line up to current t}
\NormalTok{    let traceD = "";}
\NormalTok{    for (let s=0; s\textless{}=currentStep \&\& s \textless{} argTrace.length; s++) \{}
\NormalTok{      const [tc, argVal] = argTrace[s];}
\NormalTok{      const [cx, cy] = argToChart(tc, argVal);}
\NormalTok{      traceD += (s===0 ? \textasciigrave{}M $\{cx\} $\{cy\}\textasciigrave{} : \textasciigrave{} L $\{cx\} $\{cy\}\textasciigrave{});}
\NormalTok{    \}}
\NormalTok{    const traceEl = document.createElementNS(svgNS,"path");}
\NormalTok{    traceEl.setAttribute("d",traceD); traceEl.setAttribute("fill","none");}
\NormalTok{    traceEl.setAttribute("stroke","\#8b5cf6"); traceEl.setAttribute("stroke{-}width","2");}
\NormalTok{    svg.appendChild(traceEl);}

\NormalTok{    const chartLabel = document.createElementNS(svgNS,"text");}
\NormalTok{    chartLabel.setAttribute("x",chartX+4); chartLabel.setAttribute("y",chartY+12);}
\NormalTok{    chartLabel.setAttribute("fill","\#6b7280"); chartLabel.setAttribute("font{-}size","9");}
\NormalTok{    chartLabel.setAttribute("font{-}family","sans{-}serif"); chartLabel.textContent="arg(z(t))";}
\NormalTok{    svg.appendChild(chartLabel);}

\NormalTok{    return svg;}
\NormalTok{  \}}

\NormalTok{  const container = document.createElement("div");}
\NormalTok{  container.style.cssText = "border:1px solid \#e5e7eb; border{-}radius:8px; padding:1rem; margin:1rem 0;";}

\NormalTok{  const title = document.createElement("div");}
\NormalTok{  title.style.cssText = "font{-}weight:600; margin{-}bottom:0.5rem; font{-}family:sans{-}serif;";}
\NormalTok{  title.textContent = "Winding number and the branch cut";}
\NormalTok{  container.appendChild(title);}

\NormalTok{  container.appendChild(pathSelector);}
\NormalTok{  container.appendChild(branchToggle);}
\NormalTok{  container.appendChild(tSlider);}

\NormalTok{  const vizDiv = document.createElement("div");}
\NormalTok{  container.appendChild(vizDiv);}

\NormalTok{  const note = document.createElement("div");}
\NormalTok{  note.style.cssText = "margin{-}top:0.5rem; font{-}size:0.82em; color:\#6b7280; font{-}style:italic;";}
\NormalTok{  note.textContent = "Move the slider to animate the point around the path. The purple trace shows arg(z) accumulating. When the path closes: if it enclosed the origin, arg increased by 2π (winding number = 1). If not, arg returned to its start. The red dashed line is the branch cut — where Arg z jumps by 2π.";}
\NormalTok{  container.appendChild(note);}

\NormalTok{  function update() \{}
\NormalTok{    vizDiv.innerHTML = "";}
\NormalTok{    vizDiv.appendChild(render(pathSelector.value, branchToggle.value, tSlider.value));}
\NormalTok{  \}}
\NormalTok{  pathSelector.addEventListener("input", update);}
\NormalTok{  branchToggle.addEventListener("input", update);}
\NormalTok{  tSlider.addEventListener("input", update);}
\NormalTok{  update();}

\NormalTok{  return container;}
\NormalTok{\}}
\end{Highlighting}
\end{Shaded}

\subsection{Complex powers}\label{complex-powers}

\[z^\alpha = e^{\alpha \log z}\]

Multivalued in general. Integer powers are single-valued and entire;
non-integer powers require a branch cut.

\begin{center}\rule{0.5\linewidth}{0.5pt}\end{center}

\section{Conformal mappings: preview}\label{conformal-mappings-preview}

At any point \(z_0\) where \(f\) is analytic and \(f'(z_0) \neq 0\), the
map \(f\) is \textbf{conformal} (Chapter 4): it preserves the angle
between any two curves crossing at \(z_0\) (in magnitude and
orientation). The scale factor is \(|f'(z_0)|\) and the rotation angle
is \(\arg f'(z_0)\).

Conformal mappings transform solutions of Laplace's equation from one
domain into solutions in another --- the basis for all analytic
potential theory. Chapter 4 develops this fully.

\begin{Shaded}
\begin{Highlighting}[]
\NormalTok{//| echo: false}
\NormalTok{\{}
\NormalTok{  const W = 320, H = 300;}
\NormalTok{  const SCALE = 60, CX = W/2, CY = H/2;}
\NormalTok{  const GRID\_STEP = 0.5, GRID\_RANGE = 3;}

\NormalTok{  const funcSelector = Inputs.select(}
\NormalTok{    ["w = z", "w = z²", "w = e\^{}z", "w = sin z", "w = 1/z", "w = (z−1)/(z+1)"],}
\NormalTok{    \{value: "w = z²", label: "Function f(z)"\}}
\NormalTok{  );}

\NormalTok{  function evalFunc(name, re, im) \{}
\NormalTok{    switch(name) \{}
\NormalTok{      case "w = z": return [re, im];}
\NormalTok{      case "w = z²": return [re*re {-} im*im, 2*re*im];}
\NormalTok{      case "w = e\^{}z": \{}
\NormalTok{        const ex = Math.exp(re);}
\NormalTok{        return [ex*Math.cos(im), ex*Math.sin(im)];}
\NormalTok{      \}}
\NormalTok{      case "w = sin z": return [Math.sin(re)*Math.cosh(im), Math.cos(re)*Math.sinh(im)];}
\NormalTok{      case "w = 1/z": \{}
\NormalTok{        const d = re*re + im*im;}
\NormalTok{        if (d \textless{} 1e{-}10) return [NaN, NaN];}
\NormalTok{        return [re/d, {-}im/d];}
\NormalTok{      \}}
\NormalTok{      case "w = (z−1)/(z+1)": \{}
\NormalTok{        const dr = re+1, di = im;}
\NormalTok{        const d = dr*dr + di*di;}
\NormalTok{        if (d \textless{} 1e{-}10) return [NaN, NaN];}
\NormalTok{        const nr = re{-}1, ni = im;}
\NormalTok{        return [(nr*dr + ni*di)/d, (ni*dr {-} nr*di)/d];}
\NormalTok{      \}}
\NormalTok{      default: return [re, im];}
\NormalTok{    \}}
\NormalTok{  \}}

\NormalTok{  function hslToRgb(h, s, l) \{}
\NormalTok{    h = ((h \% 1) + 1) \% 1;}
\NormalTok{    const c = (1 {-} Math.abs(2*l {-} 1)) * s;}
\NormalTok{    const x = c * (1 {-} Math.abs((h * 6) \% 2 {-} 1));}
\NormalTok{    const m = l {-} c/2;}
\NormalTok{    let r=0,g=0,b=0;}
\NormalTok{    if (h \textless{} 1/6) \{ r=c; g=x; \}}
\NormalTok{    else if (h \textless{} 2/6) \{ r=x; g=c; \}}
\NormalTok{    else if (h \textless{} 3/6) \{ g=c; b=x; \}}
\NormalTok{    else if (h \textless{} 4/6) \{ g=x; b=c; \}}
\NormalTok{    else if (h \textless{} 5/6) \{ r=x; b=c; \}}
\NormalTok{    else \{ r=c; b=x; \}}
\NormalTok{    return [Math.round((r+m)*255), Math.round((g+m)*255), Math.round((b+m)*255)];}
\NormalTok{  \}}

\NormalTok{  function renderPanel(canvas, isRight, funcName) \{}
\NormalTok{    const ctx = canvas.getContext("2d");}
\NormalTok{    const img = ctx.createImageData(W, H);}
\NormalTok{    const data = img.data;}

\NormalTok{    for (let px = 0; px \textless{} W; px++) \{}
\NormalTok{      for (let py = 0; py \textless{} H; py++) \{}
\NormalTok{        const re = (px {-} CX) / SCALE;}
\NormalTok{        const im = {-}(py {-} CY) / SCALE;}

\NormalTok{        let hue, brightness;}
\NormalTok{        if (isRight) \{}
\NormalTok{          const [wr, wi] = evalFunc(funcName, re, im);}
\NormalTok{          if (!isFinite(wr) || !isFinite(wi)) \{}
\NormalTok{            const idx = (py*W + px)*4;}
\NormalTok{            data[idx]=255; data[idx+1]=255; data[idx+2]=255; data[idx+3]=255;}
\NormalTok{            continue;}
\NormalTok{          \}}
\NormalTok{          const arg = Math.atan2(wi, wr);}
\NormalTok{          const mag = Math.sqrt(wr*wr + wi*wi);}
\NormalTok{          hue = (arg / (2*Math.PI) + 1) \% 1;}
\NormalTok{          const logmag = Math.log(mag + 0.001);}
\NormalTok{          // brightness: oscillates with log|f| so poles (large) appear bright}
\NormalTok{          brightness = 0.5 + 0.35 * Math.tanh(logmag / 2);}
\NormalTok{        \} else \{}
\NormalTok{          // Left panel: static colour encoding of identity (arg of position)}
\NormalTok{          const arg = Math.atan2(im, re);}
\NormalTok{          const mag = Math.sqrt(re*re + im*im);}
\NormalTok{          hue = (arg / (2*Math.PI) + 1) \% 1;}
\NormalTok{          brightness = 0.5 + 0.2 * Math.tanh(Math.log(mag + 0.001) / 2);}
\NormalTok{        \}}
\NormalTok{        const sat = 0.8;}
\NormalTok{        const [r,g,b] = hslToRgb(hue, sat, brightness);}
\NormalTok{        const idx = (py*W + px)*4;}
\NormalTok{        data[idx]=r; data[idx+1]=g; data[idx+2]=b; data[idx+3]=255;}
\NormalTok{      \}}
\NormalTok{    \}}
\NormalTok{    ctx.putImageData(img, 0, 0);}

\NormalTok{    // Grid overlay}
\NormalTok{    ctx.strokeStyle = "rgba(255,255,255,0.35)";}
\NormalTok{    ctx.lineWidth = 0.8;}
\NormalTok{    for (let gv = {-}GRID\_RANGE; gv \textless{}= GRID\_RANGE; gv += GRID\_STEP) \{}
\NormalTok{      // Horizontal lines (constant Im)}
\NormalTok{      ctx.beginPath();}
\NormalTok{      let first = true;}
\NormalTok{      for (let u = {-}GRID\_RANGE; u \textless{}= GRID\_RANGE; u += 0.05) \{}
\NormalTok{        let px, py;}
\NormalTok{        if (isRight) \{}
\NormalTok{          const [wr, wi] = evalFunc(funcName, u, gv);}
\NormalTok{          if (!isFinite(wr) || !isFinite(wi)) \{ first=true; continue; \}}
\NormalTok{          px = CX + wr*SCALE; py = CY {-} wi*SCALE;}
\NormalTok{        \} else \{}
\NormalTok{          px = CX + u*SCALE; py = CY {-} gv*SCALE;}
\NormalTok{        \}}
\NormalTok{        if (first) \{ ctx.moveTo(px,py); first=false; \} else ctx.lineTo(px,py);}
\NormalTok{      \}}
\NormalTok{      ctx.stroke();}

\NormalTok{      // Vertical lines (constant Re)}
\NormalTok{      ctx.beginPath();}
\NormalTok{      first = true;}
\NormalTok{      for (let v = {-}GRID\_RANGE; v \textless{}= GRID\_RANGE; v += 0.05) \{}
\NormalTok{        let px, py;}
\NormalTok{        if (isRight) \{}
\NormalTok{          const [wr, wi] = evalFunc(funcName, gv, v);}
\NormalTok{          if (!isFinite(wr) || !isFinite(wi)) \{ first=true; continue; \}}
\NormalTok{          px = CX + wr*SCALE; py = CY {-} wi*SCALE;}
\NormalTok{        \} else \{}
\NormalTok{          px = CX + gv*SCALE; py = CY {-} v*SCALE;}
\NormalTok{        \}}
\NormalTok{        if (first) \{ ctx.moveTo(px,py); first=false; \} else ctx.lineTo(px,py);}
\NormalTok{      \}}
\NormalTok{      ctx.stroke();}
\NormalTok{    \}}

\NormalTok{    // Axes}
\NormalTok{    ctx.strokeStyle = "rgba(255,255,255,0.7)"; ctx.lineWidth = 1;}
\NormalTok{    ctx.beginPath(); ctx.moveTo(0, CY); ctx.lineTo(W, CY); ctx.stroke();}
\NormalTok{    ctx.beginPath(); ctx.moveTo(CX, 0); ctx.lineTo(CX, H); ctx.stroke();}

\NormalTok{    // Labels}
\NormalTok{    ctx.fillStyle = "rgba(255,255,255,0.9)"; ctx.font = "11px sans{-}serif";}
\NormalTok{    ctx.fillText(isRight ? "w{-}plane" : "z{-}plane", 6, 15);}
\NormalTok{    ctx.fillText(isRight ? funcName : "input grid", 6, 28);}
\NormalTok{  \}}

\NormalTok{  function makeViz(funcName) \{}
\NormalTok{    const wrap = document.createElement("div");}
\NormalTok{    wrap.style.cssText = "display:flex; gap:0.75rem; flex{-}wrap:wrap; margin:0.5rem 0;";}

\NormalTok{    const leftCanvas = document.createElement("canvas");}
\NormalTok{    leftCanvas.width = W; leftCanvas.height = H;}
\NormalTok{    leftCanvas.style.cssText = "border{-}radius:4px;";}

\NormalTok{    const rightCanvas = document.createElement("canvas");}
\NormalTok{    rightCanvas.width = W; rightCanvas.height = H;}
\NormalTok{    rightCanvas.style.cssText = "border{-}radius:4px;";}

\NormalTok{    renderPanel(leftCanvas, false, funcName);}
\NormalTok{    renderPanel(rightCanvas, true, funcName);}

\NormalTok{    wrap.appendChild(leftCanvas);}
\NormalTok{    wrap.appendChild(rightCanvas);}
\NormalTok{    return wrap;}
\NormalTok{  \}}

\NormalTok{  const container = document.createElement("div");}
\NormalTok{  container.style.cssText = "border:1px solid \#e5e7eb; border{-}radius:8px; padding:1rem; margin:1rem 0;";}

\NormalTok{  const title = document.createElement("div");}
\NormalTok{  title.style.cssText = "font{-}weight:600; margin{-}bottom:0.5rem; font{-}family:sans{-}serif;";}
\NormalTok{  title.textContent = "Complex function explorer: domain colouring";}
\NormalTok{  container.appendChild(title);}
\NormalTok{  container.appendChild(funcSelector);}

\NormalTok{  const vizDiv = document.createElement("div");}
\NormalTok{  container.appendChild(vizDiv);}

\NormalTok{  const note = document.createElement("div");}
\NormalTok{  note.style.cssText = "margin{-}top:0.5rem; font{-}size:0.82em; color:\#6b7280; font{-}style:italic;";}
\NormalTok{  note.textContent = "Hue encodes arg(f(z)), brightness encodes log|f(z)|. Poles appear as bright white blooms; zeros as dark vortices. The grid shows how the z{-}plane is deformed. A function is analytic where the right panel looks locally like a rotation and scaling of the left — no tearing or folding.";}
\NormalTok{  container.appendChild(note);}

\NormalTok{  function update() \{}
\NormalTok{    vizDiv.innerHTML = "";}
\NormalTok{    vizDiv.appendChild(makeViz(funcSelector.value));}
\NormalTok{  \}}
\NormalTok{  funcSelector.addEventListener("input", update);}
\NormalTok{  update();}

\NormalTok{  return container;}
\NormalTok{\}}
\end{Highlighting}
\end{Shaded}

\begin{center}\rule{0.5\linewidth}{0.5pt}\end{center}

\section{Exercises}\label{exercises-37}

\begin{Shaded}
\begin{Highlighting}[]
\NormalTok{//| echo: false}
\NormalTok{katex = require("katex@0.16")}
\end{Highlighting}
\end{Shaded}

\begin{Shaded}
\begin{Highlighting}[]
\NormalTok{//| echo: false}
\NormalTok{function makeStepperHTML(steps) \{}
\NormalTok{  let current = 0;}
\NormalTok{  const totalSteps = steps.length;}
\NormalTok{  const container = document.createElement("div");}
\NormalTok{  container.style.cssText = "border:1px solid \#e5e7eb; border{-}radius:8px; padding:1rem 1.25rem; margin:0.75rem 0; font{-}family:inherit;";}
\NormalTok{  const stepsDiv = document.createElement("div");}
\NormalTok{  stepsDiv.style.marginBottom = "0.75rem";}
\NormalTok{  function renderSteps() \{}
\NormalTok{    stepsDiv.innerHTML = "";}
\NormalTok{    for (let i = 0; i \textless{} current; i++) \{}
\NormalTok{      const s = steps[i];}
\NormalTok{      const row = document.createElement("div");}
\NormalTok{      row.style.cssText = "display:grid; grid{-}template{-}columns:220px 1fr; gap:0.35rem 1rem; align{-}items:baseline; padding:0.35rem 0; border{-}top:1px solid \#f3f4f6;";}
\NormalTok{      const opCell = document.createElement("span");}
\NormalTok{      opCell.style.cssText = "font{-}size:0.85em; color:\#6b7280; font{-}style:italic;";}
\NormalTok{      opCell.textContent = s.op;}
\NormalTok{      const eqCell = document.createElement("span");}
\NormalTok{      eqCell.innerHTML = katex.renderToString(s.eq, \{ throwOnError: false, displayMode: false \});}
\NormalTok{      row.appendChild(opCell); row.appendChild(eqCell);}
\NormalTok{      if (s.note) \{}
\NormalTok{        const noteRow = document.createElement("div");}
\NormalTok{        noteRow.style.cssText = "grid{-}column:2; font{-}size:0.82em; color:\#6b7280; padding{-}bottom:0.2rem;";}
\NormalTok{        noteRow.textContent = s.note;}
\NormalTok{        row.appendChild(noteRow);}
\NormalTok{      \}}
\NormalTok{      stepsDiv.appendChild(row);}
\NormalTok{    \}}
\NormalTok{    if (current === totalSteps) \{}
\NormalTok{      const done = document.createElement("div");}
\NormalTok{      done.style.cssText = "margin{-}top:0.5rem; font{-}size:0.9em; color:\#059669; font{-}weight:500;";}
\NormalTok{      done.textContent = "Solution complete.";}
\NormalTok{      stepsDiv.appendChild(done);}
\NormalTok{    \}}
\NormalTok{  \}}
\NormalTok{  const controls = document.createElement("div");}
\NormalTok{  controls.style.cssText = "display:flex; gap:0.5rem; align{-}items:center;";}
\NormalTok{  const nextBtn = document.createElement("button");}
\NormalTok{  nextBtn.textContent = "Next step →";}
\NormalTok{  nextBtn.style.cssText = "padding:0.35rem 0.85rem; border:1px solid \#d1d5db; border{-}radius:4px; background:\#fff; cursor:pointer; font{-}size:0.9em;";}
\NormalTok{  const resetBtn = document.createElement("button");}
\NormalTok{  resetBtn.textContent = "Reset";}
\NormalTok{  resetBtn.style.cssText = "padding:0.35rem 0.75rem; border:1px solid \#e5e7eb; border{-}radius:4px; background:\#f9fafb; cursor:pointer; font{-}size:0.9em; color:\#6b7280;";}
\NormalTok{  const counter = document.createElement("span");}
\NormalTok{  counter.style.cssText = "font{-}size:0.82em; color:\#9ca3af; margin{-}left:0.25rem;";}
\NormalTok{  function updateButtons() \{}
\NormalTok{    nextBtn.disabled = current === totalSteps;}
\NormalTok{    nextBtn.style.opacity = current === totalSteps ? "0.4" : "1";}
\NormalTok{    counter.textContent = current === 0 ? "Click to reveal steps" : \textasciigrave{}Step $\{current\} of $\{totalSteps\}\textasciigrave{};}
\NormalTok{  \}}
\NormalTok{  nextBtn.onclick = () =\textgreater{} \{ if (current \textless{} totalSteps) \{ current++; renderSteps(); updateButtons(); \} \};}
\NormalTok{  resetBtn.onclick = () =\textgreater{} \{ current = 0; renderSteps(); updateButtons(); \};}
\NormalTok{  controls.appendChild(nextBtn); controls.appendChild(resetBtn); controls.appendChild(counter);}
\NormalTok{  container.appendChild(stepsDiv); container.appendChild(controls);}
\NormalTok{  renderSteps(); updateButtons();}
\NormalTok{  return container;}
\NormalTok{\}}
\end{Highlighting}
\end{Shaded}

\begin{center}\rule{0.5\linewidth}{0.5pt}\end{center}

\subsection{\texorpdfstring{Exercise 1: Cauchy-Riemann equations for
\(f(z) = z^2\)}{Exercise 1: Cauchy-Riemann equations for f(z) = z\^{}2}}\label{exercise-1-cauchy-riemann-equations-for-fz-z2}

Verify that \(f(z) = z^2\) satisfies the Cauchy-Riemann equations
everywhere and compute \(f'(z)\).

\begin{Shaded}
\begin{Highlighting}[]
\NormalTok{//| echo: false}
\NormalTok{makeStepperHTML([}
\NormalTok{  \{}
\NormalTok{    op: "Expand in terms of x and y",}
\NormalTok{    eq: "(x+iy)\^{}2 = x\^{}2 {-} y\^{}2 + 2ixy \textbackslash{}\textbackslash{}implies u = x\^{}2{-}y\^{}2,\textbackslash{}\textbackslash{}quad v = 2xy"}
\NormalTok{  \},}
\NormalTok{  \{}
\NormalTok{    op: "Check u\_x = v\_y",}
\NormalTok{    eq: "u\_x = 2x, \textbackslash{}\textbackslash{}quad v\_y = 2x \textbackslash{}\textbackslash{}implies u\_x = v\_y \textbackslash{}\textbackslash{}checkmark"}
\NormalTok{  \},}
\NormalTok{  \{}
\NormalTok{    op: "Check u\_y = −v\_x",}
\NormalTok{    eq: "u\_y = {-}2y, \textbackslash{}\textbackslash{}quad v\_x = 2y \textbackslash{}\textbackslash{}implies u\_y = {-}v\_x \textbackslash{}\textbackslash{}checkmark"}
\NormalTok{  \},}
\NormalTok{  \{}
\NormalTok{    op: "Partial derivatives continuous everywhere → f is entire",}
\NormalTok{    eq: "f\textquotesingle{}(z) = u\_x + iv\_x = 2x + 2iy = 2(x+iy) = 2z",}
\NormalTok{    note: "Same power rule as real calculus. The formal rules d(z\^{}n)/dz = nz\^{}\{n−1\} hold for all integer n."}
\NormalTok{  \}}
\NormalTok{])}
\end{Highlighting}
\end{Shaded}

\begin{center}\rule{0.5\linewidth}{0.5pt}\end{center}

\subsection{Exercise 2: Finding a harmonic
conjugate}\label{exercise-2-finding-a-harmonic-conjugate}

Show that \(u = x^3 - 3xy^2\) is harmonic, find its harmonic conjugate
\(v\), and identify the analytic function \(f = u + iv\).

\begin{Shaded}
\begin{Highlighting}[]
\NormalTok{//| echo: false}
\NormalTok{makeStepperHTML([}
\NormalTok{  \{}
\NormalTok{    op: "Verify Laplace\textquotesingle{}s equation",}
\NormalTok{    eq: "u\_\{xx\} = 6x, \textbackslash{}\textbackslash{}quad u\_\{yy\} = {-}6x \textbackslash{}\textbackslash{}implies \textbackslash{}\textbackslash{}nabla\^{}2 u = 0 \textbackslash{}\textbackslash{}checkmark"}
\NormalTok{  \},}
\NormalTok{  \{}
\NormalTok{    op: "C{-}R gives v\_y = u\_x = 3x²−3y²; integrate w.r.t. y",}
\NormalTok{    eq: "v = 3x\^{}2 y {-} y\^{}3 + g(x)"}
\NormalTok{  \},}
\NormalTok{  \{}
\NormalTok{    op: "C{-}R gives v\_x = −u\_y; differentiate v and set equal",}
\NormalTok{    eq: "v\_x = 6xy + g\textquotesingle{}(x) = {-}u\_y = 6xy \textbackslash{}\textbackslash{}implies g\textquotesingle{}(x) = 0 \textbackslash{}\textbackslash{}implies g = C"}
\NormalTok{  \},}
\NormalTok{  \{}
\NormalTok{    op: "Identify f(z)",}
\NormalTok{    eq: "f(z) = (x\^{}3{-}3xy\^{}2) + i(3x\^{}2y{-}y\^{}3) + iC = z\^{}3 + iC",}
\NormalTok{    note: "Expand z³ = (x+iy)³ to verify: real part is x³−3xy², imaginary part is 3x²y−y³."}
\NormalTok{  \}}
\NormalTok{])}
\end{Highlighting}
\end{Shaded}

\begin{center}\rule{0.5\linewidth}{0.5pt}\end{center}

\subsection{\texorpdfstring{Exercise 3: \(e^z\) is
entire}{Exercise 3: e\^{}z is entire}}\label{exercise-3-ez-is-entire}

Verify the Cauchy-Riemann equations for \(f(z) = e^z\) and confirm
\(\frac{d}{dz}e^z = e^z\).

\begin{Shaded}
\begin{Highlighting}[]
\NormalTok{//| echo: false}
\NormalTok{makeStepperHTML([}
\NormalTok{  \{}
\NormalTok{    op: "Write e\^{}z in terms of u and v",}
\NormalTok{    eq: "e\^{}z = e\^{}x(\textbackslash{}\textbackslash{}cos y + i\textbackslash{}\textbackslash{}sin y) \textbackslash{}\textbackslash{}implies u = e\^{}x\textbackslash{}\textbackslash{}cos y,\textbackslash{}\textbackslash{}quad v = e\^{}x\textbackslash{}\textbackslash{}sin y"}
\NormalTok{  \},}
\NormalTok{  \{}
\NormalTok{    op: "Check u\_x = v\_y",}
\NormalTok{    eq: "u\_x = e\^{}x\textbackslash{}\textbackslash{}cos y, \textbackslash{}\textbackslash{}quad v\_y = e\^{}x\textbackslash{}\textbackslash{}cos y \textbackslash{}\textbackslash{}checkmark"}
\NormalTok{  \},}
\NormalTok{  \{}
\NormalTok{    op: "Check u\_y = −v\_x",}
\NormalTok{    eq: "u\_y = {-}e\^{}x\textbackslash{}\textbackslash{}sin y, \textbackslash{}\textbackslash{}quad {-}v\_x = {-}e\^{}x\textbackslash{}\textbackslash{}sin y \textbackslash{}\textbackslash{}checkmark"}
\NormalTok{  \},}
\NormalTok{  \{}
\NormalTok{    op: "Compute the derivative",}
\NormalTok{    eq: "\textbackslash{}\textbackslash{}frac\{d\}\{dz\}e\^{}z = u\_x + iv\_x = e\^{}x\textbackslash{}\textbackslash{}cos y + ie\^{}x\textbackslash{}\textbackslash{}sin y = e\^{}z",}
\NormalTok{    note: "The exponential is its own derivative — as in real analysis."}
\NormalTok{  \}}
\NormalTok{])}
\end{Highlighting}
\end{Shaded}

\begin{center}\rule{0.5\linewidth}{0.5pt}\end{center}

\subsection{\texorpdfstring{Exercise 4: Where is \(f(z) = |z|^2\)
differentiable?}{Exercise 4: Where is f(z) = \textbar z\textbar\^{}2 differentiable?}}\label{exercise-4-where-is-fz-z2-differentiable}

\begin{Shaded}
\begin{Highlighting}[]
\NormalTok{//| echo: false}
\NormalTok{makeStepperHTML([}
\NormalTok{  \{}
\NormalTok{    op: "Write |z|² = x²+y², so u = x²+y², v = 0",}
\NormalTok{    eq: "u\_x = 2x, \textbackslash{}\textbackslash{}quad v\_y = 0 \textbackslash{}\textbackslash{}implies \textbackslash{}\textbackslash{}text\{C{-}R requires \} x = 0"}
\NormalTok{  \},}
\NormalTok{  \{}
\NormalTok{    op: "Second C{-}R equation",}
\NormalTok{    eq: "u\_y = 2y, \textbackslash{}\textbackslash{}quad {-}v\_x = 0 \textbackslash{}\textbackslash{}implies \textbackslash{}\textbackslash{}text\{C{-}R requires \} y = 0"}
\NormalTok{  \},}
\NormalTok{  \{}
\NormalTok{    op: "C{-}R equations hold only at the origin",}
\NormalTok{    eq: "f \textbackslash{}\textbackslash{}text\{ is differentiable only at \} z = 0, \textbackslash{}\textbackslash{}text\{ with \} f\textquotesingle{}(0) = 0"}
\NormalTok{  \},}
\NormalTok{  \{}
\NormalTok{    op: "Differentiability at a single point ≠ analytic",}
\NormalTok{    eq: "|z|\^{}2 \textbackslash{}\textbackslash{}text\{ is nowhere analytic\}",}
\NormalTok{    note: "Analyticity requires differentiability in an open disk. A single point of differentiability is not enough."}
\NormalTok{  \}}
\NormalTok{])}
\end{Highlighting}
\end{Shaded}

\begin{center}\rule{0.5\linewidth}{0.5pt}\end{center}

\subsection{Exercise 5: Powers via De
Moivre}\label{exercise-5-powers-via-de-moivre}

Compute \((1+i)^8\) and find all fourth roots of \(-16\).

\begin{Shaded}
\begin{Highlighting}[]
\NormalTok{//| echo: false}
\NormalTok{makeStepperHTML([}
\NormalTok{  \{}
\NormalTok{    op: "Convert 1+i to polar",}
\NormalTok{    eq: "1+i = \textbackslash{}\textbackslash{}sqrt\{2\}\textbackslash{}\textbackslash{},e\^{}\{i\textbackslash{}\textbackslash{}pi/4\}"}
\NormalTok{  \},}
\NormalTok{  \{}
\NormalTok{    op: "Raise to the 8th power",}
\NormalTok{    eq: "(1+i)\^{}8 = (\textbackslash{}\textbackslash{}sqrt\{2\})\^{}8 e\^{}\{i8\textbackslash{}\textbackslash{}pi/4\} = 16\textbackslash{}\textbackslash{},e\^{}\{2\textbackslash{}\textbackslash{}pi i\} = 16",}
\NormalTok{    note: "e\^{}\{2πi\} = 1."}
\NormalTok{  \},}
\NormalTok{  \{}
\NormalTok{    op: "Write −16 in polar form",}
\NormalTok{    eq: "{-}16 = 16\textbackslash{}\textbackslash{},e\^{}\{i\textbackslash{}\textbackslash{}pi\}"}
\NormalTok{  \},}
\NormalTok{  \{}
\NormalTok{    op: "Fourth roots: take r\^{}\{1/4\} and divide argument by 4, adding kπ/2 for k=0,1,2,3",}
\NormalTok{    eq: "z\_k = 2\textbackslash{}\textbackslash{},e\^{}\{i(\textbackslash{}\textbackslash{}pi/4 + k\textbackslash{}\textbackslash{}pi/2)\}, \textbackslash{}\textbackslash{}quad k = 0,1,2,3"}
\NormalTok{  \},}
\NormalTok{  \{}
\NormalTok{    op: "List the four roots",}
\NormalTok{    eq: "z\_0 = \textbackslash{}\textbackslash{}sqrt\{2\}(1+i), \textbackslash{}\textbackslash{}quad z\_1 = \textbackslash{}\textbackslash{}sqrt\{2\}({-}1+i), \textbackslash{}\textbackslash{}quad z\_2 = \textbackslash{}\textbackslash{}sqrt\{2\}({-}1{-}i), \textbackslash{}\textbackslash{}quad z\_3 = \textbackslash{}\textbackslash{}sqrt\{2\}(1{-}i)",}
\NormalTok{    note: "The four roots form the vertices of a square centred at the origin with circumradius 2."}
\NormalTok{  \}}
\NormalTok{])}
\end{Highlighting}
\end{Shaded}

\chapter{Complex integration}\label{complex-integration}

Cauchy's theorem, the integral formula, and their consequences

\hfill\break

\section{Contours and contour
integrals}\label{contours-and-contour-integrals}

A \textbf{contour} (or path) in the complex plane is a \textbf{piecewise
smooth} curve \(C\): \(z = z(t)\), \(a \leq t \leq b\), where
\(z(t) = x(t) + iy(t)\) and \(z'(t)\) is piecewise continuous. Piecewise
smooth means the curve can have finitely many corners, but no cusps or
fractal wiggles --- at each smooth piece the tangent direction is
well-defined.

The \textbf{contour integral} of \(f\) along \(C\) is:

\[\int_C f(z)\,dz = \int_a^b f(z(t))\,z'(t)\,dt\]

This is just a real line integral in disguise. Writing \(f = u + iv\)
and \(dz = dx + idy\), multiply out the product:
\[(u + iv)(dx + i\,dy) = u\,dx + iu\,dy + iv\,dx + i^2 v\,dy = (u\,dx - v\,dy) + i(v\,dx + u\,dy)\]

Separating real and imaginary parts:

\[\int_C f\,dz = \int_C(u\,dx - v\,dy) + i\int_C(v\,dx + u\,dy)\]

\textbf{Properties.} - Linearity:
\(\int_C (\alpha f + \beta g)\,dz = \alpha\int_C f\,dz + \beta\int_C g\,dz\).
- Reversal: \(\int_{-C}f\,dz = -\int_C f\,dz\) (reversing the direction
negates the integral). - Additivity: if \(C = C_1 + C_2\), then
\(\int_C = \int_{C_1} + \int_{C_2}\).

\subsection{The ML inequality}\label{the-ml-inequality}

The most useful estimate for contour integrals:

\[\left|\int_C f(z)\,dz\right| \leq M \cdot L\]

where \(M = \max_{z \in C}|f(z)|\) and \(L\) is the arc length of \(C\).
This is the complex analogue of
\(|\int_a^b f\,dx| \leq \max|f| \cdot (b-a)\).

\subsection{A fundamental example}\label{a-fundamental-example}

\[\oint_{|z|=r} \frac{dz}{z^{n+1}} = \begin{cases} 2\pi i & n = 0 \\ 0 & n \neq 0 \end{cases}\]

\emph{Proof for \(n=0\):} parametrise \(z = re^{i\theta}\),
\(dz = ire^{i\theta}d\theta\):
\[\int_0^{2\pi}\frac{ire^{i\theta}}{re^{i\theta}}\,d\theta = \int_0^{2\pi}i\,d\theta = 2\pi i\]

This single result --- that the integral of \(1/z\) around the origin is
\(2\pi i\) --- is the seed from which the residue theorem grows.

\begin{center}\rule{0.5\linewidth}{0.5pt}\end{center}

\section{Cauchy's theorem}\label{cauchys-theorem}

\textbf{Cauchy-Goursat theorem.} If \(f\) is analytic throughout a
simply connected domain \(D\) (no holes --- defined in Chapter 1), then
for every closed contour \(C\) in \(D\):

\[\oint_C f(z)\,dz = 0\]

\emph{Intuition.} Because the domain has no holes, any closed curve can
be continuously shrunk to a point. The integral shrinks with it --- to
zero.

\textbf{Consequence: path independence.} If \(f\) is analytic in a
simply connected domain and \(C_1\), \(C_2\) are two paths from \(z_1\)
to \(z_2\) in that domain, then \(\int_{C_1}f\,dz = \int_{C_2}f\,dz\).
The integral depends only on the endpoints, not the path.

\textbf{Deformation principle.} If \(f\) is analytic everywhere between
two contours \(C_1\) and \(C_2\) (with the same orientation), then
\(\int_{C_1}f\,dz = \int_{C_2}f\,dz\). Contours can be freely deformed
through regions where \(f\) is analytic, without changing the integral.

\begin{center}\rule{0.5\linewidth}{0.5pt}\end{center}

\section{Cauchy's integral formula}\label{cauchys-integral-formula}

\textbf{Cauchy's integral formula.} Let \(f\) be analytic in and on a
simple closed contour \(C\) (traversed counterclockwise). Then for any
\(z_0\) interior to \(C\):

\[f(z_0) = \frac{1}{2\pi i}\oint_C \frac{f(z)}{z - z_0}\,dz\]

Equivalently:

\[\oint_C \frac{f(z)}{z - z_0}\,dz = 2\pi i\,f(z_0)\]

\emph{Proof sketch.} Because \(f\) is analytic everywhere inside \(C\)
except at \(z_0\), the deformation principle applies to the annular
region between \(C\) and any circle around \(z_0\): the integrand
\(f(z)/(z-z_0)\) is analytic there. Shrink \(C\) to a small circle
\(C_\varepsilon\) of radius \(\varepsilon\) around \(z_0\) without
changing the integral. As \(\varepsilon \to 0\), \(f(z) \to f(z_0)\)
uniformly on \(C_\varepsilon\), and
\(\oint_{C_\varepsilon} dz/(z-z_0) = 2\pi i\).

\textbf{Significance.} The value of an analytic function at any interior
point is completely determined by its values on the boundary. This has
no analogue in real analysis.

\subsection{Derivatives of all orders}\label{derivatives-of-all-orders}

Differentiating the integral formula under the integral sign (justified
by uniform convergence):

\[f^{(n)}(z_0) = \frac{n!}{2\pi i}\oint_C \frac{f(z)}{(z-z_0)^{n+1}}\,dz\]

\textbf{Consequence.} An analytic function has derivatives of all
orders, all of which are themselves analytic. In real analysis,
once-differentiable functions need not be twice differentiable. In
complex analysis, once differentiable means infinitely differentiable
--- a qualitative difference.

\begin{center}\rule{0.5\linewidth}{0.5pt}\end{center}

\section{Major consequences}\label{major-consequences}

\subsection{Liouville's theorem}\label{liouvilles-theorem}

A \textbf{bounded} entire function is constant.

\emph{Proof.} If \(|f(z)| \leq M\) everywhere, the ML inequality applied
to the derivative formula gives \(|f'(z_0)| \leq M/R\) for any circle of
radius \(R\). Letting \(R \to \infty\) gives \(f'(z_0) = 0\).

\subsection{Fundamental theorem of
algebra}\label{fundamental-theorem-of-algebra}

Every non-constant polynomial \(p(z)\) of degree \(n \geq 1\) has
exactly \(n\) complex roots (counting multiplicity).

\emph{Proof sketch.} If \(p\) had no roots, \(1/p(z)\) would be entire
and bounded (since \(|p(z)| \to \infty\)), contradicting Liouville's
theorem.

\subsection{Morera's theorem}\label{moreras-theorem}

If \(f\) is continuous on a domain \(D\) and \(\oint_C f\,dz = 0\) for
every closed triangle \(C\) in \(D\), then \(f\) is analytic in \(D\).
(The converse of Cauchy's theorem.)

\begin{center}\rule{0.5\linewidth}{0.5pt}\end{center}

\section{Application: real integrals via contour
integration}\label{application-real-integrals-via-contour-integration}

The tools developed so far can already evaluate some real integrals. The
general strategy --- closing a real integral with a contour in the
complex plane and applying Cauchy's theorem --- will be developed fully
in the next chapter. The key example here:

\textbf{Example.}
\(\displaystyle\int_{-\infty}^{\infty} \frac{dx}{1+x^2}\).

The integrand \(f(z) = 1/(1+z^2)\) has poles at \(z = \pm i\). Close the
real-line integral with a large semicircle \(\Gamma_R\) of radius \(R\)
in the upper half-plane, forming a closed contour
\(C = [-R,R] + \Gamma_R\).

\textbf{The semicircle contribution vanishes.} On \(\Gamma_R\),
\(|z| = R\), so \(|1+z^2| \geq R^2 - 1\) for large \(R\). The ML
inequality gives:
\[\left|\int_{\Gamma_R}\frac{dz}{1+z^2}\right| \leq \frac{1}{R^2-1}\cdot \pi R \to 0 \quad\text{as } R\to\infty\]

\textbf{Applying Cauchy's integral formula.} Factor the denominator:
\(1/(1+z^2) = [1/(z+i)]/(z-i)\). This has the form \(f(z)/(z-z_0)\) with
\(f(z) = 1/(z+i)\) and \(z_0 = i\). Since \(z_0 = i\) lies inside the
upper half-plane contour and \(f(z) = 1/(z+i)\) is analytic there:

\[\oint_C \frac{dz}{1+z^2} = \oint_C \frac{dz}{(z-i)(z+i)} = 2\pi i \cdot f(i) = 2\pi i \cdot \frac{1}{2i} = \pi\]

Since the semicircle contributes nothing, the real-line integral equals
\(\pi\) --- confirming the known result
\(\arctan(x)\big|_{-\infty}^{\infty} = \pi\).

\begin{Shaded}
\begin{Highlighting}[]
\NormalTok{//| echo: false}
\NormalTok{\{}
\NormalTok{  const W = 320, H = 290, CX = W/2, CY = H/2, SCALE = 70;}
\NormalTok{  const svgNS = "http://www.w3.org/2000/svg";}

\NormalTok{  const radiusSlider = Inputs.range([0.5, 2.5], \{value: 1.2, step: 0.05, label: "Contour radius R"\});}
\NormalTok{  const funcSelector = Inputs.select(["f(z) = e\^{}z", "f(z) = cos z", "f(z) = 1/(z+2)"], \{}
\NormalTok{    value: "f(z) = e\^{}z", label: "f(z) (numerator)"}
\NormalTok{  \});}

\NormalTok{  // z₀ position: interactive via preset buttons for simplicity}
\NormalTok{  const polePresets = Inputs.select([}
\NormalTok{    "z₀ = 0.5+0.3i (inside)",}
\NormalTok{    "z₀ = 0+0i (inside)",}
\NormalTok{    "z₀ = 2+0i (outside)",}
\NormalTok{    "z₀ = 1.5+1.0i (outside)"}
\NormalTok{  ], \{value: "z₀ = 0.5+0.3i (inside)", label: "Pole position z₀"\});}

\NormalTok{  function getPolePos(preset) \{}
\NormalTok{    switch(preset) \{}
\NormalTok{      case "z₀ = 0.5+0.3i (inside)": return [0.5, 0.3];}
\NormalTok{      case "z₀ = 0+0i (inside)": return [0, 0];}
\NormalTok{      case "z₀ = 2+0i (outside)": return [2.0, 0];}
\NormalTok{      case "z₀ = 1.5+1.0i (outside)": return [1.5, 1.0];}
\NormalTok{      default: return [0.5, 0.3];}
\NormalTok{    \}}
\NormalTok{  \}}

\NormalTok{  function evalF(fname, re, im) \{}
\NormalTok{    switch(fname) \{}
\NormalTok{      case "f(z) = e\^{}z": return [Math.exp(re)*Math.cos(im), Math.exp(re)*Math.sin(im)];}
\NormalTok{      case "f(z) = cos z": return [Math.cos(re)*Math.cosh(im), {-}Math.sin(re)*Math.sinh(im)];}
\NormalTok{      case "f(z) = 1/(z+2)": \{}
\NormalTok{        const dr = re+2, di = im, d = dr*dr+di*di;}
\NormalTok{        if (d\textless{}1e{-}10) return [NaN,NaN];}
\NormalTok{        return [dr/d, {-}di/d];}
\NormalTok{      \}}
\NormalTok{      default: return [1, 0];}
\NormalTok{    \}}
\NormalTok{  \}}

\NormalTok{  function hslToRgb(h,s,l) \{}
\NormalTok{    h=((h\%1)+1)\%1;}
\NormalTok{    const c=(1{-}Math.abs(2*l{-}1))*s, x=c*(1{-}Math.abs((h*6)\%2{-}1)), m=l{-}c/2;}
\NormalTok{    let r=0,g=0,b=0;}
\NormalTok{    if(h\textless{}1/6)\{r=c;g=x;\}else if(h\textless{}2/6)\{r=x;g=c;\}else if(h\textless{}3/6)\{g=c;b=x;\}}
\NormalTok{    else if(h\textless{}4/6)\{g=x;b=c;\}else if(h\textless{}5/6)\{r=x;b=c;\}else\{r=c;b=x;\}}
\NormalTok{    return [Math.round((r+m)*255),Math.round((g+m)*255),Math.round((b+m)*255)];}
\NormalTok{  \}}

\NormalTok{  function toScreen(re, im) \{ return [CX + re*SCALE, CY {-} im*SCALE]; \}}

\NormalTok{  function render(R, polePreset, fname) \{}
\NormalTok{    const [z0re, z0im] = getPolePos(polePreset);}
\NormalTok{    const inside = Math.sqrt(z0re*z0re + z0im*z0im) \textless{} R;}
\NormalTok{    const onContour = Math.abs(Math.sqrt(z0re*z0re + z0im*z0im) {-} R) \textless{} 0.15;}

\NormalTok{    // Canvas for domain colouring}
\NormalTok{    const canvas = document.createElement("canvas");}
\NormalTok{    canvas.width = W; canvas.height = H;}
\NormalTok{    const ctx = canvas.getContext("2d");}
\NormalTok{    const img = ctx.createImageData(W, H);}
\NormalTok{    const data = img.data;}

\NormalTok{    for (let px=0; px\textless{}W; px++) \{}
\NormalTok{      for (let py=0; py\textless{}H; py++) \{}
\NormalTok{        const re = (px{-}CX)/SCALE, im = {-}(py{-}CY)/SCALE;}
\NormalTok{        const dr = re{-}z0re, di = im{-}z0im;}
\NormalTok{        const d2 = dr*dr+di*di;}
\NormalTok{        if (d2 \textless{} 1e{-}6) \{}
\NormalTok{          const idx=(py*W+px)*4;}
\NormalTok{          data[idx]=255;data[idx+1]=255;data[idx+2]=255;data[idx+3]=255;}
\NormalTok{          continue;}
\NormalTok{        \}}
\NormalTok{        // 1/(z{-}z0) coloured}
\NormalTok{        const wr = dr/d2, wi = {-}di/d2;}
\NormalTok{        const arg = Math.atan2(wi, wr);}
\NormalTok{        const mag = Math.sqrt(wr*wr+wi*wi);}
\NormalTok{        const hue = (arg/(2*Math.PI)+1)\%1;}
\NormalTok{        const logmag = Math.log(mag+0.001);}
\NormalTok{        const brightness = 0.5 + 0.35*Math.tanh(logmag/2);}
\NormalTok{        const [r,g,b] = hslToRgb(hue, 0.75, brightness);}
\NormalTok{        const idx=(py*W+px)*4;}
\NormalTok{        data[idx]=r;data[idx+1]=g;data[idx+2]=b;data[idx+3]=255;}
\NormalTok{      \}}
\NormalTok{    \}}
\NormalTok{    ctx.putImageData(img, 0, 0);}

\NormalTok{    // Contour circle}
\NormalTok{    const [cx0, cy0] = toScreen(0, 0);}
\NormalTok{    ctx.strokeStyle = onContour ? "\#ef4444" : "\#fff";}
\NormalTok{    ctx.lineWidth = 2.5;}
\NormalTok{    ctx.beginPath();}
\NormalTok{    ctx.arc(cx0, cy0, R*SCALE, 0, 2*Math.PI);}
\NormalTok{    ctx.stroke();}

\NormalTok{    // Contour label}
\NormalTok{    ctx.fillStyle = "rgba(255,255,255,0.9)";}
\NormalTok{    ctx.font = "11px sans{-}serif";}
\NormalTok{    ctx.fillText(\textasciigrave{}C: |z| = $\{R.toFixed(1)\}\textasciigrave{}, 5, 14);}

\NormalTok{    // Pole marker (×)}
\NormalTok{    const [psx, psy] = toScreen(z0re, z0im);}
\NormalTok{    ctx.strokeStyle = "\#fff"; ctx.lineWidth = 2;}
\NormalTok{    ctx.beginPath(); ctx.moveTo(psx{-}7,psy{-}7); ctx.lineTo(psx+7,psy+7); ctx.stroke();}
\NormalTok{    ctx.beginPath(); ctx.moveTo(psx+7,psy{-}7); ctx.lineTo(psx{-}7,psy+7); ctx.stroke();}
\NormalTok{    ctx.fillStyle = "\#fff"; ctx.font = "10px sans{-}serif";}
\NormalTok{    ctx.fillText("z₀", psx+9, psy+4);}

\NormalTok{    // Winding number box}
\NormalTok{    const winding = inside ? 1 : 0;}
\NormalTok{    ctx.fillStyle = "\#1e40af"; ctx.beginPath();}
\NormalTok{    ctx.roundRect(W{-}90, 5, 84, 42, 5); ctx.fill();}
\NormalTok{    ctx.fillStyle="white"; ctx.font="10px sans{-}serif"; ctx.textAlign="center";}
\NormalTok{    ctx.fillText("winding n(C,z₀)", W{-}48, 19);}
\NormalTok{    ctx.font="bold 18px sans{-}serif";}
\NormalTok{    ctx.fillText(winding.toString(), W{-}48, 38);}
\NormalTok{    ctx.textAlign="left";}

\NormalTok{    // Result panel (right side info)}
\NormalTok{    const infoDiv = document.createElement("div");}
\NormalTok{    infoDiv.style.cssText = "margin{-}top:0.5rem; font{-}family:sans{-}serif; font{-}size:0.88em;";}

\NormalTok{    if (onContour) \{}
\NormalTok{      infoDiv.innerHTML = \textasciigrave{}\textless{}div style="color:\#ef4444; font{-}weight:600;"\textgreater{}Pole on contour — integral undefined\textless{}/div\textgreater{}}
\NormalTok{        \textless{}div style="color:\#6b7280; font{-}size:0.85em; margin{-}top:0.3rem;"\textgreater{}The integrand blows up on the integration path.\textless{}/div\textgreater{}\textasciigrave{};}
\NormalTok{    \} else if (inside) \{}
\NormalTok{      const [fre, fim] = evalF(fname, z0re, z0im);}
\NormalTok{      const twopiIfre = ({-}2*Math.PI*fim).toFixed(4);}
\NormalTok{      const twopiIfim = (2*Math.PI*fre).toFixed(4);}
\NormalTok{      infoDiv.innerHTML = \textasciigrave{}\textless{}div style="color:\#059669; font{-}weight:600;"\textgreater{}z₀ inside C: integral = 2\textbackslash{}u03c0i\textbackslash{}u00b7f(z₀)\textless{}/div\textgreater{}}
\NormalTok{        \textless{}div style="margin{-}top:0.3rem;"\textgreater{}f(z₀) = $\{fre.toFixed(4)\} + $\{fim.toFixed(4)\}i\textless{}/div\textgreater{}}
\NormalTok{        \textless{}div\textgreater{}2\textbackslash{}u03c0i\textbackslash{}u00b7f(z₀) = $\{twopiIfre\} + $\{twopiIfim\}i\textless{}/div\textgreater{}}
\NormalTok{        \textless{}div style="color:\#6b7280; font{-}size:0.85em; margin{-}top:0.3rem;"\textgreater{}The pole is captured by the contour.\textless{}/div\textgreater{}\textasciigrave{};}
\NormalTok{    \} else \{}
\NormalTok{      infoDiv.innerHTML = \textasciigrave{}\textless{}div style="color:\#6b7280; font{-}weight:600;"\textgreater{}z₀ outside C: integral = 0\textless{}/div\textgreater{}}
\NormalTok{        \textless{}div style="color:\#6b7280; font{-}size:0.85em; margin{-}top:0.3rem;"\textgreater{}The integrand is analytic everywhere inside C — Cauchy\textquotesingle{}s theorem gives zero.\textless{}/div\textgreater{}\textasciigrave{};}
\NormalTok{    \}}

\NormalTok{    const wrap = document.createElement("div");}
\NormalTok{    wrap.appendChild(canvas);}
\NormalTok{    wrap.appendChild(infoDiv);}
\NormalTok{    return wrap;}
\NormalTok{  \}}

\NormalTok{  const container = document.createElement("div");}
\NormalTok{  container.style.cssText = "border:1px solid \#e5e7eb; border{-}radius:8px; padding:1rem; margin:1rem 0;";}
\NormalTok{  const title = document.createElement("div");}
\NormalTok{  title.style.cssText = "font{-}weight:600; margin{-}bottom:0.5rem; font{-}family:sans{-}serif;";}
\NormalTok{  title.textContent = "Cauchy integral formula: pole inside vs outside";}
\NormalTok{  container.appendChild(title);}

\NormalTok{  container.appendChild(radiusSlider);}
\NormalTok{  container.appendChild(polePresets);}
\NormalTok{  container.appendChild(funcSelector);}

\NormalTok{  const vizDiv = document.createElement("div");}
\NormalTok{  container.appendChild(vizDiv);}

\NormalTok{  const note = document.createElement("div");}
\NormalTok{  note.style.cssText = "margin{-}top:0.5rem; font{-}size:0.82em; color:\#6b7280; font{-}style:italic;";}
\NormalTok{  note.textContent = "Domain colouring shows 1/(z−z₀): the pole at z₀ appears as a white bloom. When z₀ is inside the contour C, the integral = 2πi·f(z₀). When outside, the integral = 0. The winding number n(C,z₀) records this: 1 inside, 0 outside.";}
\NormalTok{  container.appendChild(note);}

\NormalTok{  function update() \{}
\NormalTok{    vizDiv.innerHTML = "";}
\NormalTok{    vizDiv.appendChild(render(radiusSlider.value, polePresets.value, funcSelector.value));}
\NormalTok{  \}}
\NormalTok{  radiusSlider.addEventListener("input", update);}
\NormalTok{  polePresets.addEventListener("input", update);}
\NormalTok{  funcSelector.addEventListener("input", update);}
\NormalTok{  update();}

\NormalTok{  return container;}
\NormalTok{\}}
\end{Highlighting}
\end{Shaded}

\begin{Shaded}
\begin{Highlighting}[]
\NormalTok{//| echo: false}
\NormalTok{\{}
\NormalTok{  const W = 420, H = 380, CX = 190, CY = 185, SCALE = 65;}
\NormalTok{  const svgNS = "http://www.w3.org/2000/svg";}

\NormalTok{  const pathSelector2 = Inputs.select([}
\NormalTok{    "Circle around origin (n=1)",}
\NormalTok{    "Circle not around origin (n=0)",}
\NormalTok{    "Two loops around origin (n=2)"}
\NormalTok{  ], \{value: "Circle around origin (n=1)", label: "Path"\});}

\NormalTok{  const tSlider2 = Inputs.range([0, 1], \{value: 0, step: 0.005, label: "Progress t"\});}

\NormalTok{  const pathDefs = \{}
\NormalTok{    "Circle around origin (n=1)": \{}
\NormalTok{      pts: (t) =\textgreater{} [1.5*Math.cos(2*Math.PI*t), 1.5*Math.sin(2*Math.PI*t)],}
\NormalTok{      nWinding: 1}
\NormalTok{    \},}
\NormalTok{    "Circle not around origin (n=0)": \{}
\NormalTok{      pts: (t) =\textgreater{} [1.8 + 0.8*Math.cos(2*Math.PI*t), 0.8*Math.sin(2*Math.PI*t)],}
\NormalTok{      nWinding: 0}
\NormalTok{    \},}
\NormalTok{    "Two loops around origin (n=2)": \{}
\NormalTok{      pts: (t) =\textgreater{} [1.5*Math.cos(4*Math.PI*t), 1.5*Math.sin(4*Math.PI*t)],}
\NormalTok{      nWinding: 2}
\NormalTok{    \}}
\NormalTok{  \};}

\NormalTok{  function getPathPoint(def, t) \{ return def.pts(t); \}}

\NormalTok{  function computeIntegral(def, upToT, STEPS) \{}
\NormalTok{    // Numerically integrate dz/z from 0 to upToT}
\NormalTok{    // Result is a complex number (Re, Im)}
\NormalTok{    let sumRe = 0, sumIm = 0;}
\NormalTok{    const steps = Math.round(upToT * STEPS);}
\NormalTok{    const dt = upToT / Math.max(steps, 1);}

\NormalTok{    for (let s=0; s\textless{}steps; s++) \{}
\NormalTok{      const t = s / STEPS;}
\NormalTok{      const tn = (s+1) / STEPS;}
\NormalTok{      const [re1, im1] = def.pts(t);}
\NormalTok{      const [re2, im2] = def.pts(tn \textless{} 1 ? tn : 1);}
\NormalTok{      // midpoint}
\NormalTok{      const mre = (re1+re2)/2, mim = (im1+im2)/2;}
\NormalTok{      const dre = re2{-}re1, dim = im2{-}im1;}
\NormalTok{      const dz\_over\_z\_re = (mre*dre + mim*dim)/(mre*mre+mim*mim);}
\NormalTok{      const dz\_over\_z\_im = (mim*dre {-} mre*dim)/(mre*mre+mim*mim);}
\NormalTok{      sumRe += dz\_over\_z\_re;}
\NormalTok{      sumIm += dz\_over\_z\_im;}
\NormalTok{    \}}
\NormalTok{    return [sumRe, sumIm];}
\NormalTok{  \}}

\NormalTok{  function toScreen(re, im) \{ return [CX + re*SCALE, CY {-} im*SCALE]; \}}

\NormalTok{  function render2(pathName, tVal) \{}
\NormalTok{    const def = pathDefs[pathName];}
\NormalTok{    const STEPS = 400;}

\NormalTok{    const svg = document.createElementNS(svgNS, "svg");}
\NormalTok{    svg.setAttribute("width", W); svg.setAttribute("height", H);}
\NormalTok{    svg.style.cssText = "background:\#fafafa; border:1px solid \#e5e7eb; border{-}radius:6px; display:block;";}

\NormalTok{    // Grid + axes}
\NormalTok{    const gridG = document.createElementNS(svgNS,"g");}
\NormalTok{    gridG.setAttribute("stroke","\#e5e7eb"); gridG.setAttribute("stroke{-}width","0.5");}
\NormalTok{    for (let i={-}3; i\textless{}=3; i++) \{}
\NormalTok{      const xl=document.createElementNS(svgNS,"line"); const [lx]=toScreen(i,0); const [,ly1]=toScreen(0,{-}2.6); const [,ly2]=toScreen(0,2.6);}
\NormalTok{      xl.setAttribute("x1",lx); xl.setAttribute("y1",ly1); xl.setAttribute("x2",lx); xl.setAttribute("y2",ly2); gridG.appendChild(xl);}
\NormalTok{      const yl=document.createElementNS(svgNS,"line"); const [lx1]=toScreen({-}3,0); const [lx2]=toScreen(3,0); const [,ly]=toScreen(0,i);}
\NormalTok{      yl.setAttribute("x1",lx1); yl.setAttribute("y1",ly); yl.setAttribute("x2",lx2); yl.setAttribute("y2",ly); gridG.appendChild(yl);}
\NormalTok{    \}}
\NormalTok{    svg.appendChild(gridG);}

\NormalTok{    const [ox,oy]=toScreen(0,0);}
\NormalTok{    const axG=document.createElementNS(svgNS,"g"); axG.setAttribute("stroke","\#9ca3af"); axG.setAttribute("stroke{-}width","1");}
\NormalTok{    const xa=document.createElementNS(svgNS,"line"); const [ax1]=toScreen({-}2.8,0); const [ax2]=toScreen(2.8,0);}
\NormalTok{    xa.setAttribute("x1",ax1); xa.setAttribute("y1",oy); xa.setAttribute("x2",ax2); xa.setAttribute("y2",oy);}
\NormalTok{    const ya=document.createElementNS(svgNS,"line"); const [,ay1]=toScreen(0,{-}2.5); const [,ay2]=toScreen(0,2.5);}
\NormalTok{    ya.setAttribute("x1",ox); ya.setAttribute("y1",ay1); ya.setAttribute("x2",ox); ya.setAttribute("y2",ay2);}
\NormalTok{    axG.appendChild(xa); axG.appendChild(ya); svg.appendChild(axG);}

\NormalTok{    // Full path outline}
\NormalTok{    let pathD = "";}
\NormalTok{    for (let s=0; s\textless{}=STEPS; s++) \{}
\NormalTok{      const t = s/STEPS;}
\NormalTok{      const [re,im] = def.pts(t);}
\NormalTok{      const [px,py] = toScreen(re,im);}
\NormalTok{      pathD += (s===0 ? \textasciigrave{}M $\{px\} $\{py\}\textasciigrave{} : \textasciigrave{} L $\{px\} $\{py\}\textasciigrave{});}
\NormalTok{    \}}
\NormalTok{    const pathOutline = document.createElementNS(svgNS,"path");}
\NormalTok{    pathOutline.setAttribute("d",pathD); pathOutline.setAttribute("fill","none");}
\NormalTok{    pathOutline.setAttribute("stroke","rgba(249,115,22,0.25)"); pathOutline.setAttribute("stroke{-}width","2");}
\NormalTok{    svg.appendChild(pathOutline);}

\NormalTok{    // Traced portion}
\NormalTok{    let traceD = "";}
\NormalTok{    const traceSteps = Math.round(tVal * STEPS);}
\NormalTok{    for (let s=0; s\textless{}=traceSteps; s++) \{}
\NormalTok{      const t = s/STEPS;}
\NormalTok{      const [re,im] = def.pts(t);}
\NormalTok{      const [px,py] = toScreen(re,im);}
\NormalTok{      traceD += (s===0 ? \textasciigrave{}M $\{px\} $\{py\}\textasciigrave{} : \textasciigrave{} L $\{px\} $\{py\}\textasciigrave{});}
\NormalTok{    \}}
\NormalTok{    const traceEl = document.createElementNS(svgNS,"path");}
\NormalTok{    traceEl.setAttribute("d",traceD); traceEl.setAttribute("fill","none");}
\NormalTok{    traceEl.setAttribute("stroke","\#f97316"); traceEl.setAttribute("stroke{-}width","2.5");}
\NormalTok{    svg.appendChild(traceEl);}

\NormalTok{    // Current point}
\NormalTok{    const [cre,cim] = def.pts(tVal);}
\NormalTok{    const [cpx,cpy] = toScreen(cre,cim);}
\NormalTok{    const ptEl = document.createElementNS(svgNS,"circle");}
\NormalTok{    ptEl.setAttribute("cx",cpx); ptEl.setAttribute("cy",cpy); ptEl.setAttribute("r","6");}
\NormalTok{    ptEl.setAttribute("fill","\#3b82f6"); ptEl.setAttribute("stroke","\#fff"); ptEl.setAttribute("stroke{-}width","2");}
\NormalTok{    svg.appendChild(ptEl);}

\NormalTok{    // Origin × mark (the singularity 1/z)}
\NormalTok{    const xSize = 5;}
\NormalTok{    const x1m=document.createElementNS(svgNS,"line"); x1m.setAttribute("x1",ox{-}xSize); x1m.setAttribute("y1",oy{-}xSize);}
\NormalTok{    x1m.setAttribute("x2",ox+xSize); x1m.setAttribute("y2",oy+xSize);}
\NormalTok{    x1m.setAttribute("stroke","\#ef4444"); x1m.setAttribute("stroke{-}width","2.5"); svg.appendChild(x1m);}
\NormalTok{    const x2m=document.createElementNS(svgNS,"line"); x2m.setAttribute("x1",ox+xSize); x2m.setAttribute("y1",oy{-}xSize);}
\NormalTok{    x2m.setAttribute("x2",ox{-}xSize); x2m.setAttribute("y2",oy+xSize);}
\NormalTok{    x2m.setAttribute("stroke","\#ef4444"); x2m.setAttribute("stroke{-}width","2.5"); svg.appendChild(x2m);}
\NormalTok{    const origLbl=document.createElementNS(svgNS,"text"); origLbl.setAttribute("x",ox+8); origLbl.setAttribute("y",oy{-}6);}
\NormalTok{    origLbl.setAttribute("fill","\#ef4444"); origLbl.setAttribute("font{-}size","10"); origLbl.setAttribute("font{-}family","sans{-}serif");}
\NormalTok{    origLbl.textContent="singularity of 1/z"; svg.appendChild(origLbl);}

\NormalTok{    // Winding box}
\NormalTok{    const windBox = document.createElementNS(svgNS,"rect"); windBox.setAttribute("x",W{-}92); windBox.setAttribute("y",5);}
\NormalTok{    windBox.setAttribute("width",85); windBox.setAttribute("height",44); windBox.setAttribute("fill","\#1e40af"); windBox.setAttribute("rx","5");}
\NormalTok{    svg.appendChild(windBox);}
\NormalTok{    const wl=document.createElementNS(svgNS,"text"); wl.setAttribute("x",W{-}49); wl.setAttribute("y",22);}
\NormalTok{    wl.setAttribute("fill","white"); wl.setAttribute("font{-}size","10"); wl.setAttribute("text{-}anchor","middle"); wl.setAttribute("font{-}family","sans{-}serif");}
\NormalTok{    wl.textContent="winding n"; svg.appendChild(wl);}
\NormalTok{    const wn=document.createElementNS(svgNS,"text"); wn.setAttribute("x",W{-}49); wn.setAttribute("y",41);}
\NormalTok{    wn.setAttribute("fill","white"); wn.setAttribute("font{-}size","18"); wn.setAttribute("text{-}anchor","middle");}
\NormalTok{    wn.setAttribute("font{-}family","sans{-}serif"); wn.setAttribute("font{-}weight","bold");}
\NormalTok{    wn.textContent=(tVal\textgreater{}0.98)?def.nWinding.toString():"..."; svg.appendChild(wn);}

\NormalTok{    // Integral running display}
\NormalTok{    const [intRe, intIm] = computeIntegral(def, tVal, STEPS);}
\NormalTok{    const finalIntRe = 0;}
\NormalTok{    const finalIntIm = 2*Math.PI*def.nWinding;}

\NormalTok{    const intDiv = document.createElement("div");}
\NormalTok{    intDiv.style.cssText = "margin{-}top:0.5rem; font{-}family:sans{-}serif; font{-}size:0.88em; padding:0.5rem; background:\#f0f9ff; border{-}radius:4px;";}
\NormalTok{    const intStr = tVal \textgreater{} 0.98}
\NormalTok{      ? \textasciigrave{}∮ dz/z = $\{finalIntIm \textgreater{} 0 ? "+" : ""\}$\{finalIntIm.toFixed(2)\}i = 2\textbackslash{}u03c0i × $\{def.nWinding\}\textasciigrave{}}
\NormalTok{      : \textasciigrave{}Running: $\{intRe.toFixed(3)\} + $\{intIm.toFixed(3)\}i\textasciigrave{};}
\NormalTok{    intDiv.innerHTML = \textasciigrave{}\textless{}strong\textgreater{}Contour integral:\textless{}/strong\textgreater{} \textless{}span style="color:\#1e40af"\textgreater{}$\{intStr\}\textless{}/span\textgreater{}\textasciigrave{};}

\NormalTok{    const wrap = document.createElement("div");}
\NormalTok{    wrap.appendChild(svg);}
\NormalTok{    wrap.appendChild(intDiv);}
\NormalTok{    return wrap;}
\NormalTok{  \}}

\NormalTok{  const container = document.createElement("div");}
\NormalTok{  container.style.cssText = "border:1px solid \#e5e7eb; border{-}radius:8px; padding:1rem; margin:1rem 0;";}
\NormalTok{  const title = document.createElement("div");}
\NormalTok{  title.style.cssText = "font{-}weight:600; margin{-}bottom:0.5rem; font{-}family:sans{-}serif;";}
\NormalTok{  title.textContent = "Winding number = contour integral: \textbackslash{}u222e dz/z = 2\textbackslash{}u03c0i \textbackslash{}xd7 n";}
\NormalTok{  container.appendChild(title);}
\NormalTok{  container.appendChild(pathSelector2);}
\NormalTok{  container.appendChild(tSlider2);}

\NormalTok{  const vizDiv2 = document.createElement("div");}
\NormalTok{  container.appendChild(vizDiv2);}

\NormalTok{  const note = document.createElement("div");}
\NormalTok{  note.style.cssText = "margin{-}top:0.5rem; font{-}size:0.82em; color:\#6b7280; font{-}style:italic;";}
\NormalTok{  note.textContent = "Move the slider to trace the path. When the loop closes: \textbackslash{}u222e dz/z = 2\textbackslash{}u03c0i times the winding number. The red \textbackslash{}u00d7 marks the singularity at the origin. This is the seed from which the residue theorem grows.";}
\NormalTok{  container.appendChild(note);}

\NormalTok{  function update2() \{}
\NormalTok{    vizDiv2.innerHTML = "";}
\NormalTok{    vizDiv2.appendChild(render2(pathSelector2.value, tSlider2.value));}
\NormalTok{  \}}
\NormalTok{  pathSelector2.addEventListener("input", update2);}
\NormalTok{  tSlider2.addEventListener("input", update2);}
\NormalTok{  update2();}

\NormalTok{  return container;}
\NormalTok{\}}
\end{Highlighting}
\end{Shaded}

\begin{center}\rule{0.5\linewidth}{0.5pt}\end{center}

\section{Exercises}\label{exercises-38}

\begin{Shaded}
\begin{Highlighting}[]
\NormalTok{//| echo: false}
\NormalTok{katex = require("katex@0.16")}
\end{Highlighting}
\end{Shaded}

\begin{Shaded}
\begin{Highlighting}[]
\NormalTok{//| echo: false}
\NormalTok{function makeStepperHTML(steps) \{}
\NormalTok{  let current = 0;}
\NormalTok{  const totalSteps = steps.length;}
\NormalTok{  const container = document.createElement("div");}
\NormalTok{  container.style.cssText = "border:1px solid \#e5e7eb; border{-}radius:8px; padding:1rem 1.25rem; margin:0.75rem 0; font{-}family:inherit;";}
\NormalTok{  const stepsDiv = document.createElement("div");}
\NormalTok{  stepsDiv.style.marginBottom = "0.75rem";}
\NormalTok{  function renderSteps() \{}
\NormalTok{    stepsDiv.innerHTML = "";}
\NormalTok{    for (let i = 0; i \textless{} current; i++) \{}
\NormalTok{      const s = steps[i];}
\NormalTok{      const row = document.createElement("div");}
\NormalTok{      row.style.cssText = "display:grid; grid{-}template{-}columns:220px 1fr; gap:0.35rem 1rem; align{-}items:baseline; padding:0.35rem 0; border{-}top:1px solid \#f3f4f6;";}
\NormalTok{      const opCell = document.createElement("span");}
\NormalTok{      opCell.style.cssText = "font{-}size:0.85em; color:\#6b7280; font{-}style:italic;";}
\NormalTok{      opCell.textContent = s.op;}
\NormalTok{      const eqCell = document.createElement("span");}
\NormalTok{      eqCell.innerHTML = katex.renderToString(s.eq, \{ throwOnError: false, displayMode: false \});}
\NormalTok{      row.appendChild(opCell); row.appendChild(eqCell);}
\NormalTok{      if (s.note) \{}
\NormalTok{        const noteRow = document.createElement("div");}
\NormalTok{        noteRow.style.cssText = "grid{-}column:2; font{-}size:0.82em; color:\#6b7280; padding{-}bottom:0.2rem;";}
\NormalTok{        noteRow.textContent = s.note;}
\NormalTok{        row.appendChild(noteRow);}
\NormalTok{      \}}
\NormalTok{      stepsDiv.appendChild(row);}
\NormalTok{    \}}
\NormalTok{    if (current === totalSteps) \{}
\NormalTok{      const done = document.createElement("div");}
\NormalTok{      done.style.cssText = "margin{-}top:0.5rem; font{-}size:0.9em; color:\#059669; font{-}weight:500;";}
\NormalTok{      done.textContent = "Solution complete.";}
\NormalTok{      stepsDiv.appendChild(done);}
\NormalTok{    \}}
\NormalTok{  \}}
\NormalTok{  const controls = document.createElement("div");}
\NormalTok{  controls.style.cssText = "display:flex; gap:0.5rem; align{-}items:center;";}
\NormalTok{  const nextBtn = document.createElement("button");}
\NormalTok{  nextBtn.textContent = "Next step →";}
\NormalTok{  nextBtn.style.cssText = "padding:0.35rem 0.85rem; border:1px solid \#d1d5db; border{-}radius:4px; background:\#fff; cursor:pointer; font{-}size:0.9em;";}
\NormalTok{  const resetBtn = document.createElement("button");}
\NormalTok{  resetBtn.textContent = "Reset";}
\NormalTok{  resetBtn.style.cssText = "padding:0.35rem 0.75rem; border:1px solid \#e5e7eb; border{-}radius:4px; background:\#f9fafb; cursor:pointer; font{-}size:0.9em; color:\#6b7280;";}
\NormalTok{  const counter = document.createElement("span");}
\NormalTok{  counter.style.cssText = "font{-}size:0.82em; color:\#9ca3af; margin{-}left:0.25rem;";}
\NormalTok{  function updateButtons() \{}
\NormalTok{    nextBtn.disabled = current === totalSteps;}
\NormalTok{    nextBtn.style.opacity = current === totalSteps ? "0.4" : "1";}
\NormalTok{    counter.textContent = current === 0 ? "Click to reveal steps" : \textasciigrave{}Step $\{current\} of $\{totalSteps\}\textasciigrave{};}
\NormalTok{  \}}
\NormalTok{  nextBtn.onclick = () =\textgreater{} \{ if (current \textless{} totalSteps) \{ current++; renderSteps(); updateButtons(); \} \};}
\NormalTok{  resetBtn.onclick = () =\textgreater{} \{ current = 0; renderSteps(); updateButtons(); \};}
\NormalTok{  controls.appendChild(nextBtn); controls.appendChild(resetBtn); controls.appendChild(counter);}
\NormalTok{  container.appendChild(stepsDiv); container.appendChild(controls);}
\NormalTok{  renderSteps(); updateButtons();}
\NormalTok{  return container;}
\NormalTok{\}}
\end{Highlighting}
\end{Shaded}

\begin{center}\rule{0.5\linewidth}{0.5pt}\end{center}

\subsection{Exercise 1: Direct computation of a contour
integral}\label{exercise-1-direct-computation-of-a-contour-integral}

Compute \(\displaystyle\int_C \bar{z}\,dz\) where \(C\) is the straight
line from \(0\) to \(1+i\).

\begin{Shaded}
\begin{Highlighting}[]
\NormalTok{//| echo: false}
\NormalTok{makeStepperHTML([}
\NormalTok{  \{}
\NormalTok{    op: "Parametrise C: z(t) = t(1+i), 0 ≤ t ≤ 1",}
\NormalTok{    eq: "z(t) = t + it, \textbackslash{}\textbackslash{}quad dz = (1+i)\textbackslash{}\textbackslash{},dt, \textbackslash{}\textbackslash{}quad \textbackslash{}\textbackslash{}bar\{z\}(t) = t {-} it = t(1{-}i)"}
\NormalTok{  \},}
\NormalTok{  \{}
\NormalTok{    op: "Substitute into the integral",}
\NormalTok{    eq: "\textbackslash{}\textbackslash{}int\_C \textbackslash{}\textbackslash{}bar\{z\}\textbackslash{}\textbackslash{},dz = \textbackslash{}\textbackslash{}int\_0\^{}1 t(1{-}i)(1+i)\textbackslash{}\textbackslash{},dt = \textbackslash{}\textbackslash{}int\_0\^{}1 t(1{-}i\^{}2)\textbackslash{}\textbackslash{},dt = \textbackslash{}\textbackslash{}int\_0\^{}1 2t\textbackslash{}\textbackslash{},dt"}
\NormalTok{  \},}
\NormalTok{  \{}
\NormalTok{    op: "Evaluate",}
\NormalTok{    eq: "\textbackslash{}\textbackslash{}int\_C \textbackslash{}\textbackslash{}bar\{z\}\textbackslash{}\textbackslash{},dz = 2\textbackslash{}\textbackslash{}cdot\textbackslash{}\textbackslash{}frac\{1\}\{2\} = 1"}
\NormalTok{  \},}
\NormalTok{  \{}
\NormalTok{    op: "Compare with ∫\_C z dz (which IS path{-}independent)",}
\NormalTok{    eq: "\textbackslash{}\textbackslash{}int\_C z\textbackslash{}\textbackslash{},dz = \textbackslash{}\textbackslash{}left[\textbackslash{}\textbackslash{}frac\{z\^{}2\}\{2\}\textbackslash{}\textbackslash{}right]\_0\^{}\{1+i\} = \textbackslash{}\textbackslash{}frac\{(1+i)\^{}2\}\{2\} = \textbackslash{}\textbackslash{}frac\{2i\}\{2\} = i",}
\NormalTok{    note: "z dz = d(z²/2) since z² is analytic. But z̄ is not analytic, so z̄ dz is path{-}dependent and gives a real answer 1 here."}
\NormalTok{  \}}
\NormalTok{])}
\end{Highlighting}
\end{Shaded}

\begin{center}\rule{0.5\linewidth}{0.5pt}\end{center}

\subsection{Exercise 2: Cauchy's theorem --- zero
integral}\label{exercise-2-cauchys-theorem-zero-integral}

Show that \(\displaystyle\oint_C \frac{e^z}{z^2+4}\,dz = 0\) where \(C\)
is the circle \(|z| = 1\).

\begin{Shaded}
\begin{Highlighting}[]
\NormalTok{//| echo: false}
\NormalTok{makeStepperHTML([}
\NormalTok{  \{}
\NormalTok{    op: "Locate the singularities of the integrand",}
\NormalTok{    eq: "z\^{}2 + 4 = 0 \textbackslash{}\textbackslash{}implies z = \textbackslash{}\textbackslash{}pm 2i"}
\NormalTok{  \},}
\NormalTok{  \{}
\NormalTok{    op: "Check whether singularities lie inside |z| = 1",}
\NormalTok{    eq: "|2i| = 2 \textgreater{} 1 \textbackslash{}\textbackslash{}quad \textbackslash{}\textbackslash{}text\{and\} \textbackslash{}\textbackslash{}quad |{-}2i| = 2 \textgreater{} 1",}
\NormalTok{    note: "Both poles are outside the contour C."}
\NormalTok{  \},}
\NormalTok{  \{}
\NormalTok{    op: "Apply Cauchy{-}Goursat theorem",}
\NormalTok{    eq: "f(z) = \textbackslash{}\textbackslash{}frac\{e\^{}z\}\{z\^{}2+4\} \textbackslash{}\textbackslash{}text\{ is analytic everywhere inside \} |z|=1",}
\NormalTok{    note: "No singularities inside C, so the domain enclosed is simply connected and f is analytic there."}
\NormalTok{  \},}
\NormalTok{  \{}
\NormalTok{    op: "Conclude",}
\NormalTok{    eq: "\textbackslash{}\textbackslash{}oint\_\{|z|=1\}\textbackslash{}\textbackslash{}frac\{e\^{}z\}\{z\^{}2+4\}\textbackslash{}\textbackslash{},dz = 0"}
\NormalTok{  \}}
\NormalTok{])}
\end{Highlighting}
\end{Shaded}

\begin{center}\rule{0.5\linewidth}{0.5pt}\end{center}

\subsection{Exercise 3: Cauchy's integral
formula}\label{exercise-3-cauchys-integral-formula}

Evaluate \(\displaystyle\oint_C \frac{e^z}{z - 1}\,dz\) where \(C\) is
the circle \(|z| = 2\) (counterclockwise).

\begin{Shaded}
\begin{Highlighting}[]
\NormalTok{//| echo: false}
\NormalTok{makeStepperHTML([}
\NormalTok{  \{}
\NormalTok{    op: "Identify z₀ and check it is inside C",}
\NormalTok{    eq: "\textbackslash{}\textbackslash{}text\{Integrand\} = \textbackslash{}\textbackslash{}frac\{f(z)\}\{z{-}z\_0\} \textbackslash{}\textbackslash{}text\{ with \} f(z) = e\^{}z,\textbackslash{}\textbackslash{}; z\_0 = 1",}
\NormalTok{    note: "|z₀| = 1 \textless{} 2, so z₀ = 1 is inside |z| = 2. f(z) = e\^{}z is analytic everywhere."}
\NormalTok{  \},}
\NormalTok{  \{}
\NormalTok{    op: "Apply Cauchy\textquotesingle{}s integral formula",}
\NormalTok{    eq: "\textbackslash{}\textbackslash{}oint\_C \textbackslash{}\textbackslash{}frac\{e\^{}z\}\{z{-}1\}\textbackslash{}\textbackslash{},dz = 2\textbackslash{}\textbackslash{}pi i\textbackslash{}\textbackslash{},f(1) = 2\textbackslash{}\textbackslash{}pi i\textbackslash{}\textbackslash{},e"}
\NormalTok{  \}}
\NormalTok{])}
\end{Highlighting}
\end{Shaded}

\begin{center}\rule{0.5\linewidth}{0.5pt}\end{center}

\subsection{Exercise 4: Derivative
formula}\label{exercise-4-derivative-formula}

Evaluate \(\displaystyle\oint_C \frac{\cos z}{(z - \pi)^2}\,dz\) where
\(C: |z| = 4\) (counterclockwise).

\begin{Shaded}
\begin{Highlighting}[]
\NormalTok{//| echo: false}
\NormalTok{makeStepperHTML([}
\NormalTok{  \{}
\NormalTok{    op: "Recognise the derivative formula with n=1",}
\NormalTok{    eq: "\textbackslash{}\textbackslash{}oint\_C \textbackslash{}\textbackslash{}frac\{f(z)\}\{(z{-}z\_0)\^{}\{n+1\}\}\textbackslash{}\textbackslash{},dz = \textbackslash{}\textbackslash{}frac\{2\textbackslash{}\textbackslash{}pi i\}\{n!\}\textbackslash{}\textbackslash{},f\^{}\{(n)\}(z\_0)"}
\NormalTok{  \},}
\NormalTok{  \{}
\NormalTok{    op: "Identify f(z) = cos z, z₀ = π, n = 1",}
\NormalTok{    eq: "\textbackslash{}\textbackslash{}text\{Need \} f\textquotesingle{}(\textbackslash{}\textbackslash{}pi) = {-}\textbackslash{}\textbackslash{}sin(\textbackslash{}\textbackslash{}pi) = 0"}
\NormalTok{  \},}
\NormalTok{  \{}
\NormalTok{    op: "Apply the formula",}
\NormalTok{    eq: "\textbackslash{}\textbackslash{}oint\_C \textbackslash{}\textbackslash{}frac\{\textbackslash{}\textbackslash{}cos z\}\{(z{-}\textbackslash{}\textbackslash{}pi)\^{}2\}\textbackslash{}\textbackslash{},dz = \textbackslash{}\textbackslash{}frac\{2\textbackslash{}\textbackslash{}pi i\}\{1!\}\textbackslash{}\textbackslash{},f\textquotesingle{}(\textbackslash{}\textbackslash{}pi) = 2\textbackslash{}\textbackslash{}pi i \textbackslash{}\textbackslash{}cdot 0 = 0",}
\NormalTok{    note: "The integral is zero not because the integrand is analytic inside C (it has a pole at π), but because the numerator\textquotesingle{}s derivative vanishes at the pole."}
\NormalTok{  \}}
\NormalTok{])}
\end{Highlighting}
\end{Shaded}

\begin{center}\rule{0.5\linewidth}{0.5pt}\end{center}

\subsection{Exercise 5: ML inequality
estimate}\label{exercise-5-ml-inequality-estimate}

Without evaluating it exactly, show that
\(\displaystyle\left|\oint_C \frac{dz}{z^3}\right| \leq \frac{\pi}{4}\),
where \(C\) is the upper semicircle \(|z| = 2\) from \(z = 2\) to
\(z = -2\).

\begin{Shaded}
\begin{Highlighting}[]
\NormalTok{//| echo: false}
\NormalTok{makeStepperHTML([}
\NormalTok{  \{}
\NormalTok{    op: "Identify M = max|f(z)| on C",}
\NormalTok{    eq: "f(z) = z\^{}\{{-}3\}, \textbackslash{}\textbackslash{}quad |f(z)| = |z|\^{}\{{-}3\} = 2\^{}\{{-}3\} = \textbackslash{}\textbackslash{}tfrac\{1\}\{8\} \textbackslash{}\textbackslash{}text\{ on \} |z|=2"}
\NormalTok{  \},}
\NormalTok{  \{}
\NormalTok{    op: "Compute the arc length L",}
\NormalTok{    eq: "C \textbackslash{}\textbackslash{}text\{ is a semicircle of radius \} 2 \textbackslash{}\textbackslash{}implies L = \textbackslash{}\textbackslash{}pi\textbackslash{}\textbackslash{}cdot 2 = 2\textbackslash{}\textbackslash{}pi"}
\NormalTok{  \},}
\NormalTok{  \{}
\NormalTok{    op: "Apply the ML inequality",}
\NormalTok{    eq: "\textbackslash{}\textbackslash{}left|\textbackslash{}\textbackslash{}oint\_C \textbackslash{}\textbackslash{}frac\{dz\}\{z\^{}3\}\textbackslash{}\textbackslash{}right| \textbackslash{}\textbackslash{}leq M\textbackslash{}\textbackslash{}cdot L = \textbackslash{}\textbackslash{}frac\{1\}\{8\}\textbackslash{}\textbackslash{}cdot 2\textbackslash{}\textbackslash{}pi = \textbackslash{}\textbackslash{}frac\{\textbackslash{}\textbackslash{}pi\}\{4\}",}
\NormalTok{    note: "The actual value (by a direct computation or the residue theorem) is 0, since z\^{}\{−3\} = d(−z\^{}\{−2\}/2)/dz is analytic away from the origin and the origin is not enclosed."}
\NormalTok{  \}}
\NormalTok{])}
\end{Highlighting}
\end{Shaded}

\chapter{Series and residues}\label{series-and-residues}

Laurent series, isolated singularities, and the residue theorem

\hfill\break

\section{Taylor series for analytic
functions}\label{taylor-series-for-analytic-functions}

If \(f\) is analytic in a disk \(|z - z_0| < R\), it has a
\textbf{Taylor series} that converges in that disk:

\[f(z) = \sum_{n=0}^{\infty} c_n(z-z_0)^n, \qquad c_n = \frac{f^{(n)}(z_0)}{n!} = \frac{1}{2\pi i}\oint_C \frac{f(z)}{(z-z_0)^{n+1}}\,dz\]

The series converges absolutely for \(|z - z_0| < R\) and uniformly on
any closed subdisk. The radius of convergence \(R\) equals the distance
from \(z_0\) to the nearest singularity of \(f\).

\textbf{Key Taylor series} (centred at 0, valid for all \(z\)):

\[e^z = \sum_{n=0}^{\infty}\frac{z^n}{n!}, \qquad \sin z = \sum_{n=0}^{\infty}\frac{(-1)^n z^{2n+1}}{(2n+1)!}, \qquad \cos z = \sum_{n=0}^{\infty}\frac{(-1)^n z^{2n}}{(2n)!}\]

\[\frac{1}{1-z} = \sum_{n=0}^{\infty} z^n \quad (|z| < 1)\]

These are exactly the real Taylor series with \(x\) replaced by \(z\).

\begin{center}\rule{0.5\linewidth}{0.5pt}\end{center}

\section{Laurent series}\label{laurent-series}

At a point \(z_0\) where \(f\) is not analytic, the Taylor series does
not converge. The \textbf{Laurent series} extends the power series to
include negative powers:

\[f(z) = \sum_{n=-\infty}^{\infty} c_n(z-z_0)^n\]

where the coefficients are:

\[c_n = \frac{1}{2\pi i}\oint_C \frac{f(z)}{(z-z_0)^{n+1}}\,dz\]

This converges in an \textbf{annulus} \(r < |z - z_0| < R\) (with \(r\)
possibly 0 and \(R\) possibly \(\infty\)).

The \textbf{principal part} consists of the negative-power terms
\(\sum_{n=1}^{\infty} c_{-n}(z-z_0)^{-n}\). The \textbf{analytic part}
is the non-negative power series \(\sum_{n=0}^{\infty}c_n(z-z_0)^n\).
How many terms the principal part has --- zero, finitely many, or
infinitely many --- classifies the type of singularity, as the next
section shows.

\begin{Shaded}
\begin{Highlighting}[]
\NormalTok{//| echo: false}
\NormalTok{\{}
\NormalTok{  const W = 380, H = 340, CX = W/2, CY = H/2, SCALE = 62;}
\NormalTok{  const svgNS = "http://www.w3.org/2000/svg";}

\NormalTok{  const funcSel9 = Inputs.select(["1/z", "1/(z(z{-}1))", "e\^{}\{1/z\} (essential)"], \{}
\NormalTok{    value: "1/z", label: "f(z)"}
\NormalTok{  \});}
\NormalTok{  const rInnerSl = Inputs.range([0, 1.4], \{value: 0.3, step: 0.05, label: "Inner radius r"\});}
\NormalTok{  const rOuterSl = Inputs.range([0.5, 3.0], \{value: 1.5, step: 0.05, label: "Outer radius R"\});}
\NormalTok{  const nTermsSl = Inputs.range([1, 15], \{value: 5, step: 1, label: "Terms N"\});}

\NormalTok{  function laurentPartialSum(fname, re, im, N) \{}
\NormalTok{    // Returns approximate f(z) via N{-}term Laurent expansion at z0=0}
\NormalTok{    switch(fname) \{}
\NormalTok{      case "1/z": \{}
\NormalTok{        // Exact: 1/z}
\NormalTok{        const d=re*re+im*im; if(d\textless{}1e{-}10) return [NaN,NaN];}
\NormalTok{        return [re/d,{-}im/d];}
\NormalTok{      \}}
\NormalTok{      case "1/(z(z{-}1))": \{}
\NormalTok{        // = {-}1/z {-} 1 {-} z {-} z\^{}2 {-} ... for |z|\textless{}1}
\NormalTok{        // = 1/z\^{}2 * 1/(1 {-} 1/z) = 1/z\^{}2 * sum (1/z)\^{}n for |z|\textgreater{}1}
\NormalTok{        const d2=re*re+im*im; if(d2\textless{}1e{-}10) return [NaN,NaN];}
\NormalTok{        // Use exact formula}
\NormalTok{        const dr=re{-}1,di=im,d3=dr*dr+di*di; if(d3\textless{}1e{-}10) return [NaN,NaN];}
\NormalTok{        // 1/(z(z{-}1)) = 1/(z{-}1) {-} 1/z}
\NormalTok{        const p1r=dr/d3, p1i={-}di/d3;}
\NormalTok{        const p2r=re/d2, p2i={-}im/d2;}
\NormalTok{        return [p1r{-}p2r, p1i{-}p2i];}
\NormalTok{      \}}
\NormalTok{      case "e\^{}\{1/z\} (essential)": \{}
\NormalTok{        // Partial sum of e\^{}\{1/z\} = sum (1/z)\^{}n / n!}
\NormalTok{        const d=re*re+im*im; if(d\textless{}1e{-}10) return [NaN,NaN];}
\NormalTok{        const invRe=re/d, invIm={-}im/d;}
\NormalTok{        let sumRe=1, sumIm=0, powRe=1, powIm=0, fact=1;}
\NormalTok{        for(let k=1;k\textless{}=Math.min(N,12);k++) \{}
\NormalTok{          const newPowRe=powRe*invRe{-}powIm*invIm;}
\NormalTok{          const newPowIm=powRe*invIm+powIm*invRe;}
\NormalTok{          powRe=newPowRe; powIm=newPowIm; fact*=k;}
\NormalTok{          sumRe+=powRe/fact; sumIm+=powIm/fact;}
\NormalTok{        \}}
\NormalTok{        return [sumRe,sumIm];}
\NormalTok{      \}}
\NormalTok{      default: return [1,0];}
\NormalTok{    \}}
\NormalTok{  \}}

\NormalTok{  function hslToRgb(h,s,l) \{}
\NormalTok{    h=((h\%1)+1)\%1;}
\NormalTok{    const c=(1{-}Math.abs(2*l{-}1))*s,x=c*(1{-}Math.abs((h*6)\%2{-}1)),m=l{-}c/2;}
\NormalTok{    let r=0,g=0,b=0;}
\NormalTok{    if(h\textless{}1/6)\{r=c;g=x;\}else if(h\textless{}2/6)\{r=x;g=c;\}else if(h\textless{}3/6)\{g=c;b=x;\}}
\NormalTok{    else if(h\textless{}4/6)\{g=x;b=c;\}else if(h\textless{}5/6)\{r=x;b=c;\}else\{r=c;b=x;\}}
\NormalTok{    return [Math.round((r+m)*255),Math.round((g+m)*255),Math.round((b+m)*255)];}
\NormalTok{  \}}

\NormalTok{  function toScreen(re,im) \{ return [CX+re*SCALE, CY{-}im*SCALE]; \}}

\NormalTok{  function renderLaurent(fname, rInner, rOuter, N) \{}
\NormalTok{    const canvas=document.createElement("canvas");}
\NormalTok{    canvas.width=W; canvas.height=H;}
\NormalTok{    const ctx=canvas.getContext("2d");}
\NormalTok{    const img=ctx.createImageData(W,H); const data=img.data;}

\NormalTok{    for (let px=0;px\textless{}W;px++) \{}
\NormalTok{      for (let py=0;py\textless{}H;py++) \{}
\NormalTok{        const re=(px{-}CX)/SCALE, im={-}(py{-}CY)/SCALE;}
\NormalTok{        const dist=Math.sqrt(re*re+im*im);}
\NormalTok{        const inAnnulus = dist \textgreater{} rInner \&\& dist \textless{} rOuter;}

\NormalTok{        if(inAnnulus) \{}
\NormalTok{          const [wr,wi]=laurentPartialSum(fname,re,im,N);}
\NormalTok{          if(!isFinite(wr)||!isFinite(wi)||isNaN(wr)) \{}
\NormalTok{            const idx=(py*W+px)*4;}
\NormalTok{            data[idx]=255;data[idx+1]=255;data[idx+2]=255;data[idx+3]=255;}
\NormalTok{            continue;}
\NormalTok{          \}}
\NormalTok{          const arg=Math.atan2(wi,wr),mag=Math.sqrt(wr*wr+wi*wi);}
\NormalTok{          const hue=(arg/(2*Math.PI)+1)\%1;}
\NormalTok{          const logmag=Math.log(mag+0.001);}
\NormalTok{          const brightness=0.5+0.35*Math.tanh(logmag/2);}
\NormalTok{          const [r,g,b]=hslToRgb(hue,0.8,brightness);}
\NormalTok{          const idx=(py*W+px)*4; data[idx]=r;data[idx+1]=g;data[idx+2]=b;data[idx+3]=255;}
\NormalTok{        \} else \{}
\NormalTok{          // Outside annulus: dim grey}
\NormalTok{          const idx=(py*W+px)*4;}
\NormalTok{          data[idx]=220;data[idx+1]=220;data[idx+2]=220;data[idx+3]=255;}
\NormalTok{        \}}
\NormalTok{      \}}
\NormalTok{    \}}
\NormalTok{    ctx.putImageData(img,0,0);}

\NormalTok{    // Axes}
\NormalTok{    const [ox,oy]=toScreen(0,0);}
\NormalTok{    ctx.strokeStyle="\#9ca3af"; ctx.lineWidth=0.8;}
\NormalTok{    ctx.beginPath(); ctx.moveTo(0,oy); ctx.lineTo(W,oy); ctx.stroke();}
\NormalTok{    ctx.beginPath(); ctx.moveTo(ox,0); ctx.lineTo(ox,H); ctx.stroke();}

\NormalTok{    // Annulus circles}
\NormalTok{    ctx.strokeStyle="rgba(30,64,175,0.8)"; ctx.lineWidth=2; ctx.setLineDash([5,3]);}
\NormalTok{    ctx.beginPath(); ctx.arc(ox,oy,rInner*SCALE,0,2*Math.PI); ctx.stroke();}
\NormalTok{    ctx.strokeStyle="rgba(5,150,105,0.8)";}
\NormalTok{    ctx.beginPath(); ctx.arc(ox,oy,rOuter*SCALE,0,2*Math.PI); ctx.stroke();}
\NormalTok{    ctx.setLineDash([]);}

\NormalTok{    // Singularity marker}
\NormalTok{    ctx.strokeStyle="\#ef4444"; ctx.lineWidth=2;}
\NormalTok{    ctx.beginPath(); ctx.moveTo(ox{-}6,oy{-}6); ctx.lineTo(ox+6,oy+6); ctx.stroke();}
\NormalTok{    ctx.beginPath(); ctx.moveTo(ox+6,oy{-}6); ctx.lineTo(ox{-}6,oy+6); ctx.stroke();}

\NormalTok{    // Labels}
\NormalTok{    ctx.fillStyle="rgba(30,64,175,0.9)"; ctx.font="10px sans{-}serif";}
\NormalTok{    ctx.fillText(\textasciigrave{}r = $\{rInner.toFixed(2)\}\textasciigrave{},ox+rInner*SCALE+3,oy{-}4);}
\NormalTok{    ctx.fillStyle="rgba(5,150,105,0.9)";}
\NormalTok{    ctx.fillText(\textasciigrave{}R = $\{rOuter.toFixed(2)\}\textasciigrave{},ox+rOuter*SCALE+3,oy{-}4);}
\NormalTok{    ctx.fillStyle="\#374151"; ctx.font="11px sans{-}serif";}
\NormalTok{    ctx.fillText(\textasciigrave{}$\{fname\}, N=$\{N\} terms\textasciigrave{},5,14);}
\NormalTok{    ctx.fillStyle="\#6b7280"; ctx.font="10px sans{-}serif";}
\NormalTok{    ctx.fillText("Convergence annulus shown in colour",5,H{-}6);}

\NormalTok{    // Singularity type label}
\NormalTok{    let singType="";}
\NormalTok{    if(fname==="1/z") singType="Simple pole (order 1)";}
\NormalTok{    else if(fname==="1/(z(z{-}1))") singType="Poles at 0 and 1";}
\NormalTok{    else singType="Essential singularity";}

\NormalTok{    const infoDiv=document.createElement("div");}
\NormalTok{    infoDiv.style.cssText="margin{-}top:0.4rem; font{-}family:sans{-}serif; font{-}size:0.85em; color:\#374151;";}
\NormalTok{    infoDiv.innerHTML=\textasciigrave{}\textless{}strong\textgreater{}Singularity type:\textless{}/strong\textgreater{} \textless{}span style="color:\#1e40af"\textgreater{}$\{singType\}\textless{}/span\textgreater{} — series converges in the coloured annulus $\{rInner.toFixed(2)\} \&lt; |z| \&lt; $\{rOuter.toFixed(2)\}\textasciigrave{};}

\NormalTok{    const wrap=document.createElement("div");}
\NormalTok{    wrap.appendChild(canvas);}
\NormalTok{    wrap.appendChild(infoDiv);}
\NormalTok{    return wrap;}
\NormalTok{  \}}

\NormalTok{  const container=document.createElement("div");}
\NormalTok{  container.style.cssText="border:1px solid \#e5e7eb; border{-}radius:8px; padding:1rem; margin:1rem 0;";}
\NormalTok{  const title=document.createElement("div");}
\NormalTok{  title.style.cssText="font{-}weight:600; margin{-}bottom:0.5rem; font{-}family:sans{-}serif;";}
\NormalTok{  title.textContent="Laurent series: convergence in an annulus";}
\NormalTok{  container.appendChild(title);}
\NormalTok{  container.appendChild(funcSel9);}
\NormalTok{  container.appendChild(rInnerSl);}
\NormalTok{  container.appendChild(rOuterSl);}
\NormalTok{  container.appendChild(nTermsSl);}

\NormalTok{  const vizDiv9=document.createElement("div");}
\NormalTok{  container.appendChild(vizDiv9);}

\NormalTok{  const note9=document.createElement("div");}
\NormalTok{  note9.style.cssText="margin{-}top:0.5rem; font{-}size:0.82em; color:\#6b7280; font{-}style:italic;";}
\NormalTok{  note9.textContent="The domain colouring (hue = arg f, brightness = log|f|) is shown only inside the annulus r \textless{} |z| \textless{} R — where the Laurent series converges. Grey regions lie outside. Adjust the radii to see how the convergence region changes. The singularity at the origin is marked ×.";}
\NormalTok{  container.appendChild(note9);}

\NormalTok{  function update9() \{}
\NormalTok{    // Ensure outer \textgreater{} inner}
\NormalTok{    if(rOuterSl.value \textless{}= rInnerSl.value) rOuterSl.value = rInnerSl.value + 0.2;}
\NormalTok{    vizDiv9.innerHTML="";}
\NormalTok{    vizDiv9.appendChild(renderLaurent(funcSel9.value, rInnerSl.value, rOuterSl.value, nTermsSl.value));}
\NormalTok{  \}}
\NormalTok{  funcSel9.addEventListener("input", update9);}
\NormalTok{  rInnerSl.addEventListener("input", update9);}
\NormalTok{  rOuterSl.addEventListener("input", update9);}
\NormalTok{  nTermsSl.addEventListener("input", update9);}
\NormalTok{  update9();}

\NormalTok{  return container;}
\NormalTok{\}}
\end{Highlighting}
\end{Shaded}

\textbf{Computing Laurent series in practice.} For functions that are
products, quotients, or compositions of known series, it is usually more
efficient to manipulate the known Taylor series directly rather than
computing coefficients by integration.

\textbf{Example.} \(\sin z / z^3\) near \(z = 0\):

\[\frac{\sin z}{z^3} = \frac{1}{z^3}\left(z - \frac{z^3}{6} + \frac{z^5}{120} - \cdots\right) = \frac{1}{z^2} - \frac{1}{6} + \frac{z^2}{120} - \cdots\]

The principal part has only one negative-power term, \(1/z^2\).

\begin{center}\rule{0.5\linewidth}{0.5pt}\end{center}

\section{Classification of isolated
singularities}\label{classification-of-isolated-singularities}

An \textbf{isolated singularity} of \(f\) at \(z_0\) is a point where
\(f\) is not analytic but is analytic in some punctured disk
\(0 < |z - z_0| < \delta\). The Laurent series reveals the type:

\textbf{Removable singularity.} The principal part is absent (no
negative-power terms). \(f\) can be assigned a value at \(z_0\) to make
it analytic. Example: \(\sin z / z\) at \(z = 0\) --- the Laurent series
is \(1 - z^2/6 + \cdots\), which extends continuously to \(z = 0\) with
value 1.

\textbf{Pole of order \(m\).} The principal part has exactly \(m\)
terms:
\[f(z) = \frac{c_{-m}}{(z-z_0)^m} + \cdots + \frac{c_{-1}}{z-z_0} + c_0 + c_1(z-z_0) + \cdots, \quad c_{-m} \neq 0\]
\(|f(z)| \to \infty\) as \(z \to z_0\). A \textbf{simple pole} is a pole
of order 1.

\textbf{Essential singularity.} The principal part has infinitely many
terms. The behaviour near \(z_0\) is wild. Example: \(e^{1/z}\) at
\(z = 0\) --- as \(z \to 0\) along different directions, \(e^{1/z}\)
oscillates with arbitrarily large magnitude and argument, taking values
dense in the entire complex plane. There is no finite or infinite limit.

\begin{center}\rule{0.5\linewidth}{0.5pt}\end{center}

\section{Residues}\label{residues}

The \textbf{residue} of \(f\) at an isolated singularity \(z_0\) is the
Laurent coefficient \(c_{-1}\):

\[\text{Res}[f, z_0] = c_{-1}\]

\textbf{Why it matters.} Chapter 2 (§ A fundamental example) showed that
\(\oint_{|z|=r} dz/z = 2\pi i\), and that \(\oint_{|z|=r} z^n\,dz = 0\)
for all other integers \(n\). The same parametrisation argument extends
to \((z-z_0)^n\) around any circle centred at \(z_0\). So when
integrating the Laurent series term by term around a small circle
enclosing \(z_0\), every term vanishes except the \(c_{-1}\) term:
\[\oint_C \sum_{n=-\infty}^{\infty} c_n(z-z_0)^n\,dz = c_{-1}\cdot 2\pi i\]

\subsection{Computing residues}\label{computing-residues}

\textbf{At a simple pole:}
\(\text{Res}[f, z_0] = \lim_{z \to z_0}(z-z_0)f(z)\).

\textbf{At a simple pole of \(p/q\) where \(q(z_0)=0\),
\(q'(z_0)\neq 0\):}
\[\text{Res}\!\left[\frac{p}{q}, z_0\right] = \frac{p(z_0)}{q'(z_0)}\]

This is just L'Hôpital applied to the simple-pole limit. Quick check:
\(\text{Res}[1/(z^2-1),\,1]\) --- here \(p = 1\), \(q = z^2-1\),
\(q'(z) = 2z\), so the residue is \(1/(2\cdot 1) = 1/2\).

\textbf{At a pole of order \(m\):}
\[\text{Res}[f, z_0] = \frac{1}{(m-1)!}\lim_{z \to z_0}\frac{d^{m-1}}{dz^{m-1}}\left[(z-z_0)^m f(z)\right]\]

\begin{center}\rule{0.5\linewidth}{0.5pt}\end{center}

\section{The residue theorem}\label{the-residue-theorem}

\textbf{Residue theorem.} Let \(f\) be analytic on and inside a simple
closed contour \(C\) (counterclockwise) except for isolated
singularities \(z_1, \ldots, z_k\) inside \(C\). Then:

\[\oint_C f(z)\,dz = 2\pi i\sum_{j=1}^{k}\text{Res}[f, z_j]\]

This is the master formula of complex analysis. Cauchy's integral
formula is the special case where \(f\) has only simple poles.

\begin{center}\rule{0.5\linewidth}{0.5pt}\end{center}

\section{Evaluating real integrals}\label{evaluating-real-integrals}

\subsection{\texorpdfstring{Type 1: trigonometric integrals
\(\int_0^{2\pi} R(\cos\theta, \sin\theta)\,d\theta\)}{Type 1: trigonometric integrals \textbackslash int\_0\^{}\{2\textbackslash pi\} R(\textbackslash cos\textbackslash theta, \textbackslash sin\textbackslash theta)\textbackslash,d\textbackslash theta}}\label{type-1-trigonometric-integrals-int_02pi-rcostheta-sinthetadtheta}

Substitute \(z = e^{i\theta}\): \(\cos\theta = (z+z^{-1})/2\),
\(\sin\theta = (z-z^{-1})/(2i)\), \(d\theta = dz/(iz)\). The integral
becomes a contour integral on \(|z| = 1\).

\subsection{\texorpdfstring{Type 2: improper integrals
\(\int_{-\infty}^{\infty} R(x)\,dx\)}{Type 2: improper integrals \textbackslash int\_\{-\textbackslash infty\}\^{}\{\textbackslash infty\} R(x)\textbackslash,dx}}\label{type-2-improper-integrals-int_-inftyinfty-rxdx}

Close the contour with a large semicircle in the upper (or lower)
half-plane. Jordan's lemma ensures the semicircle's contribution
vanishes as the radius grows, provided the degree of the denominator of
\(R\) exceeds the degree of the numerator by at least 1. The residue
theorem gives the real integral as \(2\pi i\) times the sum of residues
in the upper half-plane.

\textbf{Example.}
\(\displaystyle\int_{-\infty}^{\infty}\frac{dx}{1+x^2}\). Poles of
\(1/(1+z^2)\): \(z = \pm i\). Only \(z = i\) is in the upper half-plane.
Residue: \(1/(2i)\). Integral: \(2\pi i \cdot 1/(2i) = \pi\).

\begin{Shaded}
\begin{Highlighting}[]
\NormalTok{//| echo: false}
\NormalTok{\{}
\NormalTok{  const W = 400, H = 360, CX = W/2, CY = H/2, SCALE = 70;}
\NormalTok{  const svgNS = "http://www.w3.org/2000/svg";}

\NormalTok{  const funcData = \{}
\NormalTok{    "1/(1+z²)": \{}
\NormalTok{      poles: [\{re:0,im:1,res:\{re:0,im:{-}0.5\},label:"i"\},\{re:0,im:{-}1,res:\{re:0,im:0.5\},label:"−i"\}],}
\NormalTok{      realIntegral: "π",}
\NormalTok{      realValue: Math.PI,}
\NormalTok{      eval: (re,im) =\textgreater{} \{}
\NormalTok{        const dr=re, di=im+1, d=dr*dr+di*di;}
\NormalTok{        if(d\textless{}1e{-}8) return [NaN,NaN];}
\NormalTok{        const dr2=re, di2=im{-}1, d2=dr2*dr2+di2*di2;}
\NormalTok{        // 1/(z²+1) = 1/((z+i)(z{-}i))}
\NormalTok{        // just colour: use |f|}
\NormalTok{        // f = 1/(1+z\^{}2): real = (1{-}y\^{}2+x\^{}2{-}2x\^{}2)/... let\textquotesingle{}s compute directly}
\NormalTok{        const r2 = 1 + re*re {-} im*im; const i2 = 2*re*im;}
\NormalTok{        const den = r2*r2 + i2*i2;}
\NormalTok{        if(den\textless{}1e{-}10) return [NaN,NaN];}
\NormalTok{        return [r2/den, {-}i2/den];}
\NormalTok{      \}}
\NormalTok{    \},}
\NormalTok{    "1/((z²+1)(z²+4))": \{}
\NormalTok{      poles: [}
\NormalTok{        \{re:0,im:1,res:\{re:0,im:{-}1/6\},label:"i"\},}
\NormalTok{        \{re:0,im:{-}1,res:\{re:0,im:1/6\},label:"−i"\},}
\NormalTok{        \{re:0,im:2,res:\{re:0,im:1/12\},label:"2i"\},}
\NormalTok{        \{re:0,im:{-}2,res:\{re:0,im:{-}1/12\},label:"−2i"\}}
\NormalTok{      ],}
\NormalTok{      realIntegral: "π/2",}
\NormalTok{      realValue: Math.PI/2,}
\NormalTok{      eval: (re,im) =\textgreater{} \{}
\NormalTok{        // f = 1/((z\^{}2+1)(z\^{}2+4))}
\NormalTok{        const a1=re*re{-}im*im+1, b1=2*re*im, a2=re*re{-}im*im+4, b2=2*re*im;}
\NormalTok{        const ra=a1*a2{-}b1*b2, ia=a1*b2+a2*b1;}
\NormalTok{        const d=ra*ra+ia*ia; if(d\textless{}1e{-}10) return [NaN,NaN];}
\NormalTok{        return [ra/d,{-}ia/d];}
\NormalTok{      \}}
\NormalTok{    \},}
\NormalTok{    "z²/((z²−1)(z+2))": \{}
\NormalTok{      poles: [}
\NormalTok{        \{re:1,im:0,res:\{re:1/3,im:0\},label:"1"\},}
\NormalTok{        \{re:{-}1,im:0,res:\{re:1/3,im:0\},label:"−1"\},}
\NormalTok{        \{re:{-}2,im:0,res:\{re:4/3,im:0\},label:"−2"\}}
\NormalTok{      ],}
\NormalTok{      realIntegral: "(sum of residues)",}
\NormalTok{      realValue: null,}
\NormalTok{      eval: (re,im) =\textgreater{} \{}
\NormalTok{        const z2r=re*re{-}im*im, z2i=2*re*im;}
\NormalTok{        const d1r=z2r{-}1,d1i=z2i;}
\NormalTok{        const d2r=re+2,d2i=im;}
\NormalTok{        const d12r=d1r*d2r{-}d1i*d2i, d12i=d1r*d2i+d1i*d2r;}
\NormalTok{        const nd=d12r*d12r+d12i*d12i; if(nd\textless{}1e{-}10) return [NaN,NaN];}
\NormalTok{        return [(z2r*d12r+z2i*d12i)/nd,(z2i*d12r{-}z2r*d12i)/nd];}
\NormalTok{      \}}
\NormalTok{    \}}
\NormalTok{  \};}

\NormalTok{  const funcSelector = Inputs.select(Object.keys(funcData), \{}
\NormalTok{    value: "1/(1+z²)", label: "f(z)"}
\NormalTok{  \});}
\NormalTok{  const radiusSlider = Inputs.range([0.1, 4.5], \{value: 1.5, step: 0.05, label: "Contour radius R"\});}

\NormalTok{  function hslToRgb(h,s,l) \{}
\NormalTok{    h=((h\%1)+1)\%1;}
\NormalTok{    const c=(1{-}Math.abs(2*l{-}1))*s,x=c*(1{-}Math.abs((h*6)\%2{-}1)),m=l{-}c/2;}
\NormalTok{    let r=0,g=0,b=0;}
\NormalTok{    if(h\textless{}1/6)\{r=c;g=x;\}else if(h\textless{}2/6)\{r=x;g=c;\}else if(h\textless{}3/6)\{g=c;b=x;\}}
\NormalTok{    else if(h\textless{}4/6)\{g=x;b=c;\}else if(h\textless{}5/6)\{r=x;b=c;\}else\{r=c;b=x;\}}
\NormalTok{    return [Math.round((r+m)*255),Math.round((g+m)*255),Math.round((b+m)*255)];}
\NormalTok{  \}}

\NormalTok{  function toScreen(re,im) \{ return [CX+re*SCALE, CY{-}im*SCALE]; \}}

\NormalTok{  function render3(fname, R) \{}
\NormalTok{    const fd = funcData[fname];}

\NormalTok{    const canvas = document.createElement("canvas");}
\NormalTok{    canvas.width=W; canvas.height=H;}
\NormalTok{    const ctx=canvas.getContext("2d");}
\NormalTok{    const img=ctx.createImageData(W,H); const data=img.data;}

\NormalTok{    for (let px=0;px\textless{}W;px++) \{}
\NormalTok{      for (let py=0;py\textless{}H;py++) \{}
\NormalTok{        const re=(px{-}CX)/SCALE, im={-}(py{-}CY)/SCALE;}
\NormalTok{        const [wr,wi]=fd.eval(re,im);}
\NormalTok{        if(!isFinite(wr)||!isFinite(wi)) \{}
\NormalTok{          const idx=(py*W+px)*4;}
\NormalTok{          data[idx]=255;data[idx+1]=255;data[idx+2]=255;data[idx+3]=255;}
\NormalTok{          continue;}
\NormalTok{        \}}
\NormalTok{        const arg=Math.atan2(wi,wr), mag=Math.sqrt(wr*wr+wi*wi);}
\NormalTok{        const hue=(arg/(2*Math.PI)+1)\%1;}
\NormalTok{        const logmag=Math.log(mag+0.001);}
\NormalTok{        const brightness=0.5+0.35*Math.tanh(logmag/2);}
\NormalTok{        const [r,g,b]=hslToRgb(hue,0.75,brightness);}
\NormalTok{        const idx=(py*W+px)*4;}
\NormalTok{        data[idx]=r;data[idx+1]=g;data[idx+2]=b;data[idx+3]=255;}
\NormalTok{      \}}
\NormalTok{    \}}
\NormalTok{    ctx.putImageData(img,0,0);}

\NormalTok{    // Contour}
\NormalTok{    const [cx0,cy0]=toScreen(0,0);}
\NormalTok{    ctx.strokeStyle="rgba(255,255,255,0.9)"; ctx.lineWidth=2;}
\NormalTok{    ctx.beginPath(); ctx.arc(cx0,cy0,R*SCALE,0,2*Math.PI); ctx.stroke();}
\NormalTok{    ctx.fillStyle="rgba(255,255,255,0.8)"; ctx.font="11px sans{-}serif";}
\NormalTok{    ctx.fillText(\textasciigrave{}C: |z| = $\{R.toFixed(2)\}\textasciigrave{},5,14);}

\NormalTok{    // Pole markers and labels}
\NormalTok{    const enclosed = [];}
\NormalTok{    fd.poles.forEach(p =\textgreater{} \{}
\NormalTok{      const [psx,psy]=toScreen(p.re,p.im);}
\NormalTok{      const inC = Math.sqrt(p.re*p.re+p.im*p.im) \textless{} R;}
\NormalTok{      if(inC) enclosed.push(p);}
\NormalTok{      ctx.strokeStyle=inC?"\#fbbf24":"rgba(255,255,255,0.7)";}
\NormalTok{      ctx.lineWidth=2;}
\NormalTok{      ctx.beginPath(); ctx.moveTo(psx{-}7,psy{-}7); ctx.lineTo(psx+7,psy+7); ctx.stroke();}
\NormalTok{      ctx.beginPath(); ctx.moveTo(psx+7,psy{-}7); ctx.lineTo(psx{-}7,psy+7); ctx.stroke();}
\NormalTok{      ctx.fillStyle=inC?"\#fbbf24":"rgba(255,255,255,0.7)";}
\NormalTok{      ctx.font="10px sans{-}serif";}
\NormalTok{      ctx.fillText(p.label,psx+9,psy+4);}
\NormalTok{    \});}

\NormalTok{    // Residue sum}
\NormalTok{    let sumRe=0, sumIm=0;}
\NormalTok{    enclosed.forEach(p =\textgreater{} \{ sumRe+=p.res.re; sumIm+=p.res.im; \});}
\NormalTok{    const integral2piRe = {-}2*Math.PI*sumIm;}
\NormalTok{    const integral2piIm = 2*Math.PI*sumRe;}

\NormalTok{    const infoDiv=document.createElement("div");}
\NormalTok{    infoDiv.style.cssText="margin{-}top:0.5rem; font{-}family:sans{-}serif; font{-}size:0.85em; background:\#f0f9ff; border{-}radius:4px; padding:0.5rem;";}
\NormalTok{    const encStr=enclosed.length===0?"none":enclosed.map(p=\textgreater{}p.label).join(", ");}
\NormalTok{    const resStr=enclosed.length===0?"0":(\textasciigrave{}$\{sumRe.toFixed(4)\}$\{sumIm\textgreater{}=0?"+":""\}$\{sumIm.toFixed(4)\}i\textasciigrave{});}
\NormalTok{    const intStr=enclosed.length===0?"0":\textasciigrave{}$\{integral2piRe.toFixed(4)\} + $\{integral2piIm.toFixed(4)\}i\textasciigrave{};}
\NormalTok{    const realStr=(fd.realValue!==null\&\&enclosed.some(p=\textgreater{}p.im\textgreater{}0))}
\NormalTok{      ?\textasciigrave{}\textless{}div style="color:\#059669"\textgreater{}Real integral = 2\textbackslash{}u03c0 \textbackslash{}xd7 Im(upper{-}half sum) = $\{(2*Math.PI*sumRe).toFixed(4)\} [predicted: $\{fd.realIntegral\}]\textless{}/div\textgreater{}\textasciigrave{}}
\NormalTok{      :"";}
\NormalTok{    infoDiv.innerHTML=\textasciigrave{}}
\NormalTok{      \textless{}div\textgreater{}\textless{}strong\textgreater{}Enclosed poles:\textless{}/strong\textgreater{} $\{encStr\}\textless{}/div\textgreater{}}
\NormalTok{      \textless{}div\textgreater{}\textless{}strong\textgreater{}Sum of residues:\textless{}/strong\textgreater{} $\{resStr\}\textless{}/div\textgreater{}}
\NormalTok{      \textless{}div\textgreater{}\textless{}strong\textgreater{}\textbackslash{}u222e\textless{}sub\textgreater{}C\textless{}/sub\textgreater{} f\textbackslash{}u202fdz = 2\textbackslash{}u03c0i \textbackslash{}xd7 \textbackslash{}u03a3Res = \textless{}/strong\textgreater{}\textless{}span style="color:\#1e40af"\textgreater{}$\{intStr\}\textless{}/span\textgreater{}\textless{}/div\textgreater{}}
\NormalTok{      $\{realStr\}}
\NormalTok{    \textasciigrave{};}

\NormalTok{    const wrap=document.createElement("div");}
\NormalTok{    wrap.appendChild(canvas);}
\NormalTok{    wrap.appendChild(infoDiv);}
\NormalTok{    return wrap;}
\NormalTok{  \}}

\NormalTok{  const container=document.createElement("div");}
\NormalTok{  container.style.cssText="border:1px solid \#e5e7eb; border{-}radius:8px; padding:1rem; margin:1rem 0;";}
\NormalTok{  const title=document.createElement("div");}
\NormalTok{  title.style.cssText="font{-}weight:600; margin{-}bottom:0.5rem; font{-}family:sans{-}serif;";}
\NormalTok{  title.textContent="Residue theorem: sweeping poles with the contour";}
\NormalTok{  container.appendChild(title);}
\NormalTok{  container.appendChild(funcSelector);}
\NormalTok{  container.appendChild(radiusSlider);}

\NormalTok{  const vizDiv3=document.createElement("div");}
\NormalTok{  container.appendChild(vizDiv3);}

\NormalTok{  const note=document.createElement("div");}
\NormalTok{  note.style.cssText="margin{-}top:0.5rem; font{-}size:0.82em; color:\#6b7280; font{-}style:italic;";}
\NormalTok{  note.textContent="Domain colouring shows f(z): poles appear as white blooms. Gold \textbackslash{}u00d7 marks = poles inside contour. As R increases past each pole, the running residue sum updates and the integral changes. Each pole contributes 2\textbackslash{}u03c0i times its residue.";}
\NormalTok{  container.appendChild(note);}

\NormalTok{  function update3() \{}
\NormalTok{    vizDiv3.innerHTML="";}
\NormalTok{    vizDiv3.appendChild(render3(funcSelector.value, radiusSlider.value));}
\NormalTok{  \}}
\NormalTok{  funcSelector.addEventListener("input", update3);}
\NormalTok{  radiusSlider.addEventListener("input", update3);}
\NormalTok{  update3();}

\NormalTok{  return container;}
\NormalTok{\}}
\end{Highlighting}
\end{Shaded}

The residue theorem and Laurent series complete the core tools of
complex integration. The next chapter turns to a different application
of analyticity: conformal mapping, which uses analytic functions to
transform complicated geometries into simple ones where Laplace's
equation can be solved by inspection.

\begin{center}\rule{0.5\linewidth}{0.5pt}\end{center}

\section{Exercises}\label{exercises-39}

\begin{Shaded}
\begin{Highlighting}[]
\NormalTok{//| echo: false}
\NormalTok{katex = require("katex@0.16")}
\end{Highlighting}
\end{Shaded}

\begin{Shaded}
\begin{Highlighting}[]
\NormalTok{//| echo: false}
\NormalTok{function makeStepperHTML(steps) \{}
\NormalTok{  let current = 0;}
\NormalTok{  const totalSteps = steps.length;}
\NormalTok{  const container = document.createElement("div");}
\NormalTok{  container.style.cssText = "border:1px solid \#e5e7eb; border{-}radius:8px; padding:1rem 1.25rem; margin:0.75rem 0; font{-}family:inherit;";}
\NormalTok{  const stepsDiv = document.createElement("div");}
\NormalTok{  stepsDiv.style.marginBottom = "0.75rem";}
\NormalTok{  function renderSteps() \{}
\NormalTok{    stepsDiv.innerHTML = "";}
\NormalTok{    for (let i = 0; i \textless{} current; i++) \{}
\NormalTok{      const s = steps[i];}
\NormalTok{      const row = document.createElement("div");}
\NormalTok{      row.style.cssText = "display:grid; grid{-}template{-}columns:220px 1fr; gap:0.35rem 1rem; align{-}items:baseline; padding:0.35rem 0; border{-}top:1px solid \#f3f4f6;";}
\NormalTok{      const opCell = document.createElement("span");}
\NormalTok{      opCell.style.cssText = "font{-}size:0.85em; color:\#6b7280; font{-}style:italic;";}
\NormalTok{      opCell.textContent = s.op;}
\NormalTok{      const eqCell = document.createElement("span");}
\NormalTok{      eqCell.innerHTML = katex.renderToString(s.eq, \{ throwOnError: false, displayMode: false \});}
\NormalTok{      row.appendChild(opCell); row.appendChild(eqCell);}
\NormalTok{      if (s.note) \{}
\NormalTok{        const noteRow = document.createElement("div");}
\NormalTok{        noteRow.style.cssText = "grid{-}column:2; font{-}size:0.82em; color:\#6b7280; padding{-}bottom:0.2rem;";}
\NormalTok{        noteRow.textContent = s.note;}
\NormalTok{        row.appendChild(noteRow);}
\NormalTok{      \}}
\NormalTok{      stepsDiv.appendChild(row);}
\NormalTok{    \}}
\NormalTok{    if (current === totalSteps) \{}
\NormalTok{      const done = document.createElement("div");}
\NormalTok{      done.style.cssText = "margin{-}top:0.5rem; font{-}size:0.9em; color:\#059669; font{-}weight:500;";}
\NormalTok{      done.textContent = "Solution complete.";}
\NormalTok{      stepsDiv.appendChild(done);}
\NormalTok{    \}}
\NormalTok{  \}}
\NormalTok{  const controls = document.createElement("div");}
\NormalTok{  controls.style.cssText = "display:flex; gap:0.5rem; align{-}items:center;";}
\NormalTok{  const nextBtn = document.createElement("button");}
\NormalTok{  nextBtn.textContent = "Next step →";}
\NormalTok{  nextBtn.style.cssText = "padding:0.35rem 0.85rem; border:1px solid \#d1d5db; border{-}radius:4px; background:\#fff; cursor:pointer; font{-}size:0.9em;";}
\NormalTok{  const resetBtn = document.createElement("button");}
\NormalTok{  resetBtn.textContent = "Reset";}
\NormalTok{  resetBtn.style.cssText = "padding:0.35rem 0.75rem; border:1px solid \#e5e7eb; border{-}radius:4px; background:\#f9fafb; cursor:pointer; font{-}size:0.9em; color:\#6b7280;";}
\NormalTok{  const counter = document.createElement("span");}
\NormalTok{  counter.style.cssText = "font{-}size:0.82em; color:\#9ca3af; margin{-}left:0.25rem;";}
\NormalTok{  function updateButtons() \{}
\NormalTok{    nextBtn.disabled = current === totalSteps;}
\NormalTok{    nextBtn.style.opacity = current === totalSteps ? "0.4" : "1";}
\NormalTok{    counter.textContent = current === 0 ? "Click to reveal steps" : \textasciigrave{}Step $\{current\} of $\{totalSteps\}\textasciigrave{};}
\NormalTok{  \}}
\NormalTok{  nextBtn.onclick = () =\textgreater{} \{ if (current \textless{} totalSteps) \{ current++; renderSteps(); updateButtons(); \} \};}
\NormalTok{  resetBtn.onclick = () =\textgreater{} \{ current = 0; renderSteps(); updateButtons(); \};}
\NormalTok{  controls.appendChild(nextBtn); controls.appendChild(resetBtn); controls.appendChild(counter);}
\NormalTok{  container.appendChild(stepsDiv); container.appendChild(controls);}
\NormalTok{  renderSteps(); updateButtons();}
\NormalTok{  return container;}
\NormalTok{\}}
\end{Highlighting}
\end{Shaded}

\begin{center}\rule{0.5\linewidth}{0.5pt}\end{center}

\subsection{Exercise 1: Laurent series around a
pole}\label{exercise-1-laurent-series-around-a-pole}

Find the Laurent series of \(f(z) = \dfrac{1}{z(z-1)}\) valid for
\(0 < |z| < 1\), and classify the singularity at \(z = 0\).

\begin{Shaded}
\begin{Highlighting}[]
\NormalTok{//| echo: false}
\NormalTok{makeStepperHTML([}
\NormalTok{  \{}
\NormalTok{    op: "Partial fractions",}
\NormalTok{    eq: "\textbackslash{}\textbackslash{}frac\{1\}\{z(z{-}1)\} = \textbackslash{}\textbackslash{}frac\{A\}\{z\} + \textbackslash{}\textbackslash{}frac\{B\}\{z{-}1\} \textbackslash{}\textbackslash{}implies A = {-}1,\textbackslash{}\textbackslash{}; B = 1",}
\NormalTok{    note: "Multiply both sides by z: at z=0, A = 1/(0−1) = −1. Multiply by (z−1): at z=1, B = 1/1 = 1."}
\NormalTok{  \},}
\NormalTok{  \{}
\NormalTok{    op: "Expand 1/(z−1) using geometric series for |z| \textless{} 1",}
\NormalTok{    eq: "\textbackslash{}\textbackslash{}frac\{1\}\{z{-}1\} = \textbackslash{}\textbackslash{}frac\{{-}1\}\{1{-}z\} = {-}\textbackslash{}\textbackslash{}sum\_\{n=0\}\^{}\{\textbackslash{}\textbackslash{}infty\}z\^{}n"}
\NormalTok{  \},}
\NormalTok{  \{}
\NormalTok{    op: "Combine",}
\NormalTok{    eq: "\textbackslash{}\textbackslash{}frac\{1\}\{z(z{-}1)\} = {-}\textbackslash{}\textbackslash{}frac\{1\}\{z\} {-} 1 {-} z {-} z\^{}2 {-} z\^{}3 {-} \textbackslash{}\textbackslash{}cdots \textbackslash{}\textbackslash{}quad (0 \textless{} |z| \textless{} 1)"}
\NormalTok{  \},}
\NormalTok{  \{}
\NormalTok{    op: "Classify the singularity at z = 0",}
\NormalTok{    eq: "\textbackslash{}\textbackslash{}text\{Principal part: \} {-}1/z \textbackslash{}\textbackslash{}text\{ (one term)\} \textbackslash{}\textbackslash{}implies \textbackslash{}\textbackslash{}text\{simple pole of order 1\}",}
\NormalTok{    note: "The residue at z = 0 is the coefficient c\_\{−1\} = −1."}
\NormalTok{  \}}
\NormalTok{])}
\end{Highlighting}
\end{Shaded}

\begin{center}\rule{0.5\linewidth}{0.5pt}\end{center}

\subsection{\texorpdfstring{Exercise 2: Classifying \(\sin z / z^2\) at
\(z = 0\)}{Exercise 2: Classifying \textbackslash sin z / z\^{}2 at z = 0}}\label{exercise-2-classifying-sin-z-z2-at-z-0}

\begin{Shaded}
\begin{Highlighting}[]
\NormalTok{//| echo: false}
\NormalTok{makeStepperHTML([}
\NormalTok{  \{}
\NormalTok{    op: "Expand sin z in Taylor series",}
\NormalTok{    eq: "\textbackslash{}\textbackslash{}sin z = z {-} \textbackslash{}\textbackslash{}frac\{z\^{}3\}\{6\} + \textbackslash{}\textbackslash{}frac\{z\^{}5\}\{120\} {-} \textbackslash{}\textbackslash{}cdots"}
\NormalTok{  \},}
\NormalTok{  \{}
\NormalTok{    op: "Divide by z²",}
\NormalTok{    eq: "\textbackslash{}\textbackslash{}frac\{\textbackslash{}\textbackslash{}sin z\}\{z\^{}2\} = \textbackslash{}\textbackslash{}frac\{1\}\{z\} {-} \textbackslash{}\textbackslash{}frac\{z\}\{6\} + \textbackslash{}\textbackslash{}frac\{z\^{}3\}\{120\} {-} \textbackslash{}\textbackslash{}cdots"}
\NormalTok{  \},}
\NormalTok{  \{}
\NormalTok{    op: "Principal part has one term: 1/z",}
\NormalTok{    eq: "\textbackslash{}\textbackslash{}text\{Simple pole at \} z = 0, \textbackslash{}\textbackslash{}quad \textbackslash{}\textbackslash{}text\{Res\}\textbackslash{}\textbackslash{}!\textbackslash{}\textbackslash{}left[\textbackslash{}\textbackslash{}frac\{\textbackslash{}\textbackslash{}sin z\}\{z\^{}2\},\textbackslash{}\textbackslash{}, 0\textbackslash{}\textbackslash{}right] = 1"}
\NormalTok{  \},}
\NormalTok{  \{}
\NormalTok{    op: "Compare with sin(z)/z and sin(z)/z³",}
\NormalTok{    eq: "\textbackslash{}\textbackslash{}frac\{\textbackslash{}\textbackslash{}sin z\}\{z\} \textbackslash{}\textbackslash{}to 1 \textbackslash{}\textbackslash{}text\{ (removable)\}, \textbackslash{}\textbackslash{}quad \textbackslash{}\textbackslash{}frac\{\textbackslash{}\textbackslash{}sin z\}\{z\^{}3\} \textbackslash{}\textbackslash{}text\{ has principal part \} 1/z\^{}2 \textbackslash{}\textbackslash{}text\{ (order{-}2 pole)\}",}
\NormalTok{    note: "The order of the pole equals (order of zero of denominator) − (order of zero of numerator): z² has a zero of order 2 at 0; sin z has a zero of order 1; difference = 1."}
\NormalTok{  \}}
\NormalTok{])}
\end{Highlighting}
\end{Shaded}

\begin{center}\rule{0.5\linewidth}{0.5pt}\end{center}

\subsection{Exercise 3: Residue at a simple
pole}\label{exercise-3-residue-at-a-simple-pole}

Find \(\text{Res}\!\left[\dfrac{z^2+1}{(z-2)(z+3)},\; z=2\right]\).

\begin{Shaded}
\begin{Highlighting}[]
\NormalTok{//| echo: false}
\NormalTok{makeStepperHTML([}
\NormalTok{  \{}
\NormalTok{    op: "Apply the simple{-}pole formula",}
\NormalTok{    eq: "\textbackslash{}\textbackslash{}text\{Res\}[f, 2] = \textbackslash{}\textbackslash{}lim\_\{z\textbackslash{}\textbackslash{}to 2\}(z{-}2)f(z) = \textbackslash{}\textbackslash{}lim\_\{z\textbackslash{}\textbackslash{}to 2\}\textbackslash{}\textbackslash{}frac\{z\^{}2+1\}\{z+3\}"}
\NormalTok{  \},}
\NormalTok{  \{}
\NormalTok{    op: "Evaluate at z = 2",}
\NormalTok{    eq: "= \textbackslash{}\textbackslash{}frac\{4+1\}\{2+3\} = \textbackslash{}\textbackslash{}frac\{5\}\{5\} = 1"}
\NormalTok{  \},}
\NormalTok{  \{}
\NormalTok{    op: "Alternatively, use p/q′ formula with q = (z−2)(z+3)",}
\NormalTok{    eq: "q\textquotesingle{}(z) = 2z+1, \textbackslash{}\textbackslash{}quad q\textquotesingle{}(2) = 5, \textbackslash{}\textbackslash{}quad p(2) = 5 \textbackslash{}\textbackslash{}implies \textbackslash{}\textbackslash{}text\{Res\} = 5/5 = 1 \textbackslash{}\textbackslash{}checkmark"}
\NormalTok{  \}}
\NormalTok{])}
\end{Highlighting}
\end{Shaded}

\begin{center}\rule{0.5\linewidth}{0.5pt}\end{center}

\subsection{Exercise 4: Residue at a double
pole}\label{exercise-4-residue-at-a-double-pole}

Find \(\text{Res}\!\left[\dfrac{e^z}{(z-1)^2},\; z=1\right]\).

\begin{Shaded}
\begin{Highlighting}[]
\NormalTok{//| echo: false}
\NormalTok{makeStepperHTML([}
\NormalTok{  \{}
\NormalTok{    op: "Apply the order{-}m formula with m = 2",}
\NormalTok{    eq: "\textbackslash{}\textbackslash{}text\{Res\}[f,1] = \textbackslash{}\textbackslash{}frac\{1\}\{(2{-}1)!\}\textbackslash{}\textbackslash{}lim\_\{z\textbackslash{}\textbackslash{}to 1\}\textbackslash{}\textbackslash{}frac\{d\}\{dz\}\textbackslash{}\textbackslash{}left[(z{-}1)\^{}2\textbackslash{}\textbackslash{}cdot\textbackslash{}\textbackslash{}frac\{e\^{}z\}\{(z{-}1)\^{}2\}\textbackslash{}\textbackslash{}right] = \textbackslash{}\textbackslash{}lim\_\{z\textbackslash{}\textbackslash{}to 1\}\textbackslash{}\textbackslash{}frac\{d\}\{dz\}e\^{}z"}
\NormalTok{  \},}
\NormalTok{  \{}
\NormalTok{    op: "Evaluate",}
\NormalTok{    eq: "\textbackslash{}\textbackslash{}frac\{d\}\{dz\}e\^{}z\textbackslash{}\textbackslash{}bigg|\_\{z=1\} = e\^{}1 = e"}
\NormalTok{  \},}
\NormalTok{  \{}
\NormalTok{    op: "Verify via Laurent expansion",}
\NormalTok{    eq: "\textbackslash{}\textbackslash{}frac\{e\^{}z\}\{(z{-}1)\^{}2\} = \textbackslash{}\textbackslash{}frac\{e\textbackslash{}\textbackslash{},e\^{}\{z{-}1\}\}\{(z{-}1)\^{}2\} = \textbackslash{}\textbackslash{}frac\{e\}\{(z{-}1)\^{}2\}\textbackslash{}\textbackslash{}left[1 + (z{-}1) + \textbackslash{}\textbackslash{}frac\{(z{-}1)\^{}2\}\{2\} + \textbackslash{}\textbackslash{}cdots\textbackslash{}\textbackslash{}right] = \textbackslash{}\textbackslash{}frac\{e\}\{(z{-}1)\^{}2\} + \textbackslash{}\textbackslash{}frac\{e\}\{z{-}1\} + \textbackslash{}\textbackslash{}cdots",}
\NormalTok{    note: "Coefficient of (z−1)\^{}\{−1\} is e. ✓"}
\NormalTok{  \}}
\NormalTok{])}
\end{Highlighting}
\end{Shaded}

\begin{center}\rule{0.5\linewidth}{0.5pt}\end{center}

\subsection{Exercise 5: Trigonometric
integral}\label{exercise-5-trigonometric-integral}

Evaluate \(\displaystyle\int_0^{2\pi}\frac{d\theta}{2 + \cos\theta}\).

\begin{Shaded}
\begin{Highlighting}[]
\NormalTok{//| echo: false}
\NormalTok{makeStepperHTML([}
\NormalTok{  \{}
\NormalTok{    op: "Substitute z = e\^{}\{iθ\}: cos θ = (z+1/z)/2, dθ = dz/(iz)",}
\NormalTok{    eq: "\textbackslash{}\textbackslash{}int\_0\^{}\{2\textbackslash{}\textbackslash{}pi\}\textbackslash{}\textbackslash{}frac\{d\textbackslash{}\textbackslash{}theta\}\{2+\textbackslash{}\textbackslash{}cos\textbackslash{}\textbackslash{}theta\} = \textbackslash{}\textbackslash{}oint\_\{|z|=1\}\textbackslash{}\textbackslash{}frac\{1\}\{2+\textbackslash{}\textbackslash{}frac\{z+z\^{}\{{-}1\}\}\{2\}\}\textbackslash{}\textbackslash{}cdot\textbackslash{}\textbackslash{}frac\{dz\}\{iz\} = \textbackslash{}\textbackslash{}oint\_\{|z|=1\}\textbackslash{}\textbackslash{}frac\{2\textbackslash{}\textbackslash{},dz\}\{i(z\^{}2+4z+1)\}"}
\NormalTok{  \},}
\NormalTok{  \{}
\NormalTok{    op: "Find roots of z²+4z+1 = 0",}
\NormalTok{    eq: "z = \textbackslash{}\textbackslash{}frac\{{-}4\textbackslash{}\textbackslash{}pm\textbackslash{}\textbackslash{}sqrt\{12\}\}\{2\} = {-}2\textbackslash{}\textbackslash{}pm\textbackslash{}\textbackslash{}sqrt\{3\}",}
\NormalTok{    note: "z₁ = −2+√3 ≈ −0.27 (inside |z|=1); z₂ = −2−√3 ≈ −3.73 (outside)."}
\NormalTok{  \},}
\NormalTok{  \{}
\NormalTok{    op: "Compute Res[2/(i(z²+4z+1)), z₁]",}
\NormalTok{    eq: "\textbackslash{}\textbackslash{}text\{Res\} = \textbackslash{}\textbackslash{}frac\{2\}\{i(2z\_1+4)\} = \textbackslash{}\textbackslash{}frac\{2\}\{i\textbackslash{}\textbackslash{}cdot 2\textbackslash{}\textbackslash{}sqrt\{3\}\} = \textbackslash{}\textbackslash{}frac\{1\}\{i\textbackslash{}\textbackslash{}sqrt\{3\}\} = \textbackslash{}\textbackslash{}frac\{{-}i\}\{\textbackslash{}\textbackslash{}sqrt\{3\}\}"}
\NormalTok{  \},}
\NormalTok{  \{}
\NormalTok{    op: "Apply the residue theorem",}
\NormalTok{    eq: "\textbackslash{}\textbackslash{}int\_0\^{}\{2\textbackslash{}\textbackslash{}pi\}\textbackslash{}\textbackslash{}frac\{d\textbackslash{}\textbackslash{}theta\}\{2+\textbackslash{}\textbackslash{}cos\textbackslash{}\textbackslash{}theta\} = 2\textbackslash{}\textbackslash{}pi i\textbackslash{}\textbackslash{}cdot\textbackslash{}\textbackslash{}frac\{{-}i\}\{\textbackslash{}\textbackslash{}sqrt\{3\}\} = \textbackslash{}\textbackslash{}frac\{2\textbackslash{}\textbackslash{}pi\}\{\textbackslash{}\textbackslash{}sqrt\{3\}\}"}
\NormalTok{  \}}
\NormalTok{])}
\end{Highlighting}
\end{Shaded}

\begin{center}\rule{0.5\linewidth}{0.5pt}\end{center}

\subsection{Exercise 6: Improper real integral via
residues}\label{exercise-6-improper-real-integral-via-residues}

Evaluate \(\displaystyle\int_{-\infty}^{\infty}\frac{dx}{(1+x^2)^2}\).

\begin{Shaded}
\begin{Highlighting}[]
\NormalTok{//| echo: false}
\NormalTok{makeStepperHTML([}
\NormalTok{  \{}
\NormalTok{    op: "Poles of 1/(1+z²)² at z = ±i, each of order 2",}
\NormalTok{    eq: "\textbackslash{}\textbackslash{}frac\{1\}\{(1+z\^{}2)\^{}2\} = \textbackslash{}\textbackslash{}frac\{1\}\{(z{-}i)\^{}2(z+i)\^{}2\}"}
\NormalTok{  \},}
\NormalTok{  \{}
\NormalTok{    op: "Close contour with upper semicircle; only z = i is in upper half{-}plane",}
\NormalTok{    eq: "\textbackslash{}\textbackslash{}text\{Res\}\textbackslash{}\textbackslash{}left[\textbackslash{}\textbackslash{}frac\{1\}\{(z{-}i)\^{}2(z+i)\^{}2\},\textbackslash{}\textbackslash{},i\textbackslash{}\textbackslash{}right] = \textbackslash{}\textbackslash{}lim\_\{z\textbackslash{}\textbackslash{}to i\}\textbackslash{}\textbackslash{}frac\{d\}\{dz\}\textbackslash{}\textbackslash{}frac\{1\}\{(z+i)\^{}2\}"}
\NormalTok{  \},}
\NormalTok{  \{}
\NormalTok{    op: "Differentiate",}
\NormalTok{    eq: "\textbackslash{}\textbackslash{}frac\{d\}\{dz\}(z+i)\^{}\{{-}2\} = {-}2(z+i)\^{}\{{-}3\} \textbackslash{}\textbackslash{}implies \textbackslash{}\textbackslash{}text\{Res\} = \textbackslash{}\textbackslash{}frac\{{-}2\}\{(2i)\^{}3\} = \textbackslash{}\textbackslash{}frac\{{-}2\}\{{-}8i\} = \textbackslash{}\textbackslash{}frac\{1\}\{4i\} = \textbackslash{}\textbackslash{}frac\{{-}i\}\{4\}"}
\NormalTok{  \},}
\NormalTok{  \{}
\NormalTok{    op: "Apply residue theorem (semicircle integral → 0 as R → ∞)",}
\NormalTok{    eq: "\textbackslash{}\textbackslash{}int\_\{{-}\textbackslash{}\textbackslash{}infty\}\^{}\{\textbackslash{}\textbackslash{}infty\}\textbackslash{}\textbackslash{}frac\{dx\}\{(1+x\^{}2)\^{}2\} = 2\textbackslash{}\textbackslash{}pi i\textbackslash{}\textbackslash{}cdot\textbackslash{}\textbackslash{}frac\{{-}i\}\{4\} = \textbackslash{}\textbackslash{}frac\{2\textbackslash{}\textbackslash{}pi\}\{4\} = \textbackslash{}\textbackslash{}frac\{\textbackslash{}\textbackslash{}pi\}\{2\}"}
\NormalTok{  \}}
\NormalTok{])}
\end{Highlighting}
\end{Shaded}

\chapter{Conformal mapping}\label{conformal-mapping}

Angle-preserving transformations and applications to potential theory

\hfill\break

\section{Conformal mappings}\label{conformal-mappings}

An analytic function \(w = f(z)\) is \textbf{conformal at \(z_0\)} if
\(f'(z_0) \neq 0\). At such a point, the mapping:

\begin{enumerate}
\def\labelenumi{\arabic{enumi}.}
\tightlist
\item
  \textbf{Preserves angles} between curves (in both magnitude and
  orientation).
\item
  \textbf{Scales uniformly}: all infinitesimal lengths at \(z_0\) are
  multiplied by \(|f'(z_0)|\).
\item
  \textbf{Rotates uniformly}: all directions at \(z_0\) are rotated by
  \(\arg f'(z_0)\).
\end{enumerate}

\emph{Why angles are preserved.} If two curves \(\gamma_1\) and
\(\gamma_2\) meet at \(z_0\) at angle \(\alpha\), their images under
\(f\) meet at \(f(z_0)\) at the same angle. The chain rule gives the
tangent direction of the image curve as \(f'(z_0)\cdot\gamma'(t_0)\) ---
the original tangent vector multiplied by \(f'(z_0)\). Writing
\(f'(z_0) = |f'(z_0)|e^{i\varphi}\) with \(\varphi = \arg f'(z_0)\),
this multiplication scales every tangent by \(|f'(z_0)|\) and rotates it
by \(\varphi\). Both curves get the same rotation, so the angle between
them is unchanged.

At a \textbf{critical point} where \(f'(z_0) = 0\), conformality fails:
angles are multiplied by the order of the zero.

\textbf{Example.} \(w = z^2\) at \(z = 0\): \(f'(0) = 0\), so the map is
not conformal at the origin. The angle between any two curves through 0
is doubled. Away from 0, \(f'(z) = 2z \neq 0\) and the map is conformal.

\begin{center}\rule{0.5\linewidth}{0.5pt}\end{center}

\section{Möbius transformations}\label{muxf6bius-transformations}

The \textbf{Möbius transformation} (or \textbf{fractional linear
transformation}) is:

\[w = \frac{az + b}{cz + d}, \qquad ad - bc \neq 0\]

The condition \(ad - bc \neq 0\) ensures that \(w\) is not constant.
These transformations form a group under composition and are the most
important class of conformal maps.

\textbf{Key property: circles and lines map to circles and lines.} In
the extended complex plane --- the ordinary plane with a single point at
infinity added --- a straight line can be viewed as a circle that passes
through that point at infinity. Under a Möbius transformation, circles
(including lines viewed this way) always map to circles (including
lines).

\textbf{Three-point determination.} A Möbius transformation is
completely determined by specifying three distinct points and their
images. Given pairs \((z_1 \mapsto w_1)\), \((z_2 \mapsto w_2)\),
\((z_3 \mapsto w_3)\), the transformation is found by the
\textbf{cross-ratio} equation:

\[\frac{(w - w_1)(w_2 - w_3)}{(w - w_3)(w_2 - w_1)} = \frac{(z - z_1)(z_2 - z_3)}{(z - z_3)(z_2 - z_1)}\]

The right side is a known number for each \(z\). Cross-multiply to get a
linear equation in \(w\), then solve --- the result is the unique Möbius
transformation. Exercise 3 works through this procedure.

\textbf{Standard mappings.}

\begin{itemize}
\item
  \(w = \dfrac{z - i}{z + i}\): maps the upper half-plane
  \(\text{Im}(z) > 0\) to the unit disk \(|w| < 1\), and the real axis
  to the unit circle.
\item
  \(w = \dfrac{1+z}{1-z}\): maps \(|z| < 1\) (unit disk) to
  \(\text{Re}(w) > 0\) (right half-plane).
\end{itemize}

\subsection{Fixed points}\label{fixed-points}

The \textbf{fixed points} of a Möbius transformation
\(w = (az+b)/(cz+d)\) satisfy \(z = (az+b)/(cz+d)\), i.e.,
\(cz^2 + (d-a)z - b = 0\). A generic Möbius transformation has two fixed
points (possibly coincident or at \(\infty\)).

\begin{center}\rule{0.5\linewidth}{0.5pt}\end{center}

\section{Standard conformal maps}\label{standard-conformal-maps}

\subsection{\texorpdfstring{\(w = z^2\)}{w = z\^{}2}}\label{w-z2}

Doubles the argument and squares the modulus. The upper half-plane
\(\{z : \text{Im}(z) > 0\}\) maps to the full complex plane minus the
non-negative real axis (since arg doubles from \((0,\pi)\) to
\((0,2\pi)\)).

The first quadrant \(\{0 < \arg z < \pi/2\}\) maps to the upper
half-plane \(\{0 < \arg w < \pi\}\).

\subsection{\texorpdfstring{\(w = e^z\)}{w = e\^{}z}}\label{w-ez}

Maps the horizontal strip \(\{0 < \text{Im}(z) < \pi\}\) to the upper
half-plane \(\{\text{Im}(w) > 0\}\): the strip width \(\pi\) becomes the
half-plane argument range \((0, \pi)\).

More generally, the strip \(\{a < \text{Im}(z) < b\}\) maps to the
sector \(\{a < \arg w < b\}\).

\subsection{\texorpdfstring{\(w = \log z\)}{w = \textbackslash log z}}\label{w-log-z}

The inverse of \(w = e^z\). Maps the upper half-plane to the horizontal
strip \(\{0 < \text{Im}(w) < \pi\}\).

\begin{center}\rule{0.5\linewidth}{0.5pt}\end{center}

\section{The Joukowski
transformation}\label{the-joukowski-transformation}

\[w = z + \frac{1}{z}\]

This is the foundational map of thin-aerofoil theory.

\textbf{The unit circle maps to a degenerate shape.} For
\(z = e^{i\theta}\):
\[w = e^{i\theta} + e^{-i\theta} = 2\cos\theta \in [-2, 2]\]

The unit circle collapses to the real interval \([-2, 2]\) --- a flat
plate.

\textbf{Circles near the unit circle map to aerofoil-like shapes.} For
\(z = re^{i\theta}\) with \(r > 1\):
\[w = \left(r + \frac{1}{r}\right)\cos\theta + i\left(r - \frac{1}{r}\right)\sin\theta\]

This is an \textbf{ellipse} with semi-axes \(a = r + 1/r\) and
\(b = r - 1/r\). As \(r \to 1^+\), the ellipse degenerates toward the
flat plate.

\textbf{Off-centre circles.} Shifting the circle in the \(z\)-plane
before applying the Joukowski map breaks the ellipse's symmetry and
produces a curved aerofoil profile with a sharp trailing edge. This is
the Joukowski aerofoil --- the basis of early wing design.

\textbf{Why it works for aerodynamics.} The complex potential for
irrotational flow around a cylinder of radius \(r\) is known:
\(F(z) = U(z + r^2/z)\) for uniform flow at speed \(U\). Composing with
the Joukowski map gives the complex potential for flow around the
aerofoil. The Kutta--Joukowski theorem then gives the lift per unit span
as \(L = \rho U \Gamma\), where
\(\Gamma = \oint \mathbf{v}\cdot d\mathbf{s}\) is the
\textbf{circulation} --- the line integral of the tangential velocity
component around the aerofoil.

\begin{Shaded}
\begin{Highlighting}[]
\NormalTok{//| echo: false}
\NormalTok{\{}
\NormalTok{  const WP = 340, HP = 320, SCALE = 70;}
\NormalTok{  const svgNS = "http://www.w3.org/2000/svg";}

\NormalTok{  const rSlider4 = Inputs.range([1.0, 2.0], \{value: 1.15, step: 0.01, label: "Circle radius r"\});}
\NormalTok{  const offsetSlider = Inputs.range([0.0, 0.3], \{value: 0.12, step: 0.01, label: "x{-}offset \textbackslash{}u03b4"\});}
\NormalTok{  const circSlider = Inputs.range([0, Math.PI*2], \{value: 0, step: 0.1, label: "Circulation \textbackslash{}u0393"\});}

\NormalTok{  function joukowski(zRe, zIm) \{}
\NormalTok{    const d = zRe*zRe + zIm*zIm;}
\NormalTok{    if(d \textless{} 1e{-}10) return [NaN, NaN];}
\NormalTok{    return [zRe + zRe/d, zIm {-} zIm/d];}
\NormalTok{  \}}

\NormalTok{  function toScreen(re, im, cx, cy, scale) \{}
\NormalTok{    return [cx + re*scale, cy {-} im*scale];}
\NormalTok{  \}}

\NormalTok{  function renderJoukowski(r, delta, circ) \{}
\NormalTok{    // Build aerofoil curve}
\NormalTok{    const N = 360;}
\NormalTok{    const aerofoilPts = [];}
\NormalTok{    for(let k=0; k\textless{}N; k++) \{}
\NormalTok{      const theta = (2*Math.PI*k)/(N{-}1);}
\NormalTok{      const zRe = {-}delta + r*Math.cos(theta);}
\NormalTok{      const zIm = r*Math.sin(theta);}
\NormalTok{      const [wRe, wIm] = joukowski(zRe, zIm);}
\NormalTok{      if(isFinite(wRe) \&\& isFinite(wIm)) aerofoilPts.push([wRe, wIm]);}
\NormalTok{    \}}

\NormalTok{    // Left panel: z{-}plane}
\NormalTok{    const leftSvg = document.createElementNS(svgNS, "svg");}
\NormalTok{    leftSvg.setAttribute("width", WP); leftSvg.setAttribute("height", HP);}
\NormalTok{    leftSvg.style.cssText = "background:\#f8fafc; border:1px solid \#e5e7eb; border{-}radius:6px;";}

\NormalTok{    const CX = WP/2 + 10, CY = HP/2;}

\NormalTok{    // Grid}
\NormalTok{    const gridG = document.createElementNS(svgNS,"g"); gridG.setAttribute("stroke","\#e5e7eb"); gridG.setAttribute("stroke{-}width","0.5");}
\NormalTok{    for(let i={-}3;i\textless{}=3;i++) \{}
\NormalTok{      const xl=document.createElementNS(svgNS,"line"); const [lx]=toScreen(i,0,CX,CY,SCALE); const [,ly1]=toScreen(0,{-}2.2,CX,CY,SCALE); const [,ly2]=toScreen(0,2.2,CX,CY,SCALE);}
\NormalTok{      xl.setAttribute("x1",lx); xl.setAttribute("y1",ly1); xl.setAttribute("x2",lx); xl.setAttribute("y2",ly2); gridG.appendChild(xl);}
\NormalTok{      const yl=document.createElementNS(svgNS,"line"); const [lx1]=toScreen({-}3,0,CX,CY,SCALE); const [lx2]=toScreen(3,0,CX,CY,SCALE); const [,ly]=toScreen(0,i,CX,CY,SCALE);}
\NormalTok{      yl.setAttribute("x1",lx1); yl.setAttribute("y1",ly); yl.setAttribute("x2",lx2); yl.setAttribute("y2",ly); gridG.appendChild(yl);}
\NormalTok{    \}}
\NormalTok{    leftSvg.appendChild(gridG);}

\NormalTok{    // Streamlines in z{-}plane: Im(F) = U(z + r²/z) = const}
\NormalTok{    // F(z) = z + r²/z for uniform flow + vortex}
\NormalTok{    // Streamlines: Im(F) = Im(z) + r² * Im(1/z) = y + r²*({-}y)/(x²+y²) = const}
\NormalTok{    // For vortex: add circulation * i/(2π) * log(z)}
\NormalTok{    // Simple: just draw horizontal streamlines mapped through the potential}
\NormalTok{    const streamG = document.createElementNS(svgNS,"g");}
\NormalTok{    streamG.setAttribute("stroke","rgba(59,130,246,0.4)"); streamG.setAttribute("stroke{-}width","1"); streamG.setAttribute("fill","none");}
\NormalTok{    for(let psi={-}2.5; psi\textless{}=2.5; psi+=0.35) \{}
\NormalTok{      // parametric: Im(w + r²/w\_conj) = psi — approximate by iterating x}
\NormalTok{      let pathD=""; let started=false;}
\NormalTok{      for(let x={-}3.5; x\textless{}=3.5; x+=0.04) \{}
\NormalTok{        // Im(F) ≈ y(1 {-} r²/(x²+y²)) = psi; solve for y numerically}
\NormalTok{        // simple approximation: far from origin, y ≈ psi}
\NormalTok{        // Near origin, use Newton}
\NormalTok{        let y=psi;}
\NormalTok{        for(let iter=0;iter\textless{}8;iter++) \{}
\NormalTok{          const d2=x*x+y*y; if(d2\textless{}0.01) break;}
\NormalTok{          const f=y{-}r*r*y/d2 {-} psi;}
\NormalTok{          const fp=1 {-} r*r*(d2{-}2*y*y)/(d2*d2);}
\NormalTok{          y {-}= f/(fp||1e{-}6);}
\NormalTok{        \}}
\NormalTok{        const d2=x*x+y*y;}
\NormalTok{        if(d2 \textless{} (r{-}delta)*(r{-}delta)*0.9) \{ started=false; continue; \} // inside circle}
\NormalTok{        const [px,py]=toScreen(x,y,CX,CY,SCALE);}
\NormalTok{        if(px\textless{}0||px\textgreater{}WP||py\textless{}0||py\textgreater{}HP) \{ started=false; continue; \}}
\NormalTok{        pathD += started ? \textasciigrave{} L $\{px\} $\{py\}\textasciigrave{} : \textasciigrave{}M $\{px\} $\{py\}\textasciigrave{}; started=true;}
\NormalTok{      \}}
\NormalTok{      if(pathD) \{ const pEl=document.createElementNS(svgNS,"path"); pEl.setAttribute("d",pathD); pEl.setAttribute("stroke","rgba(59,130,246,0.35)"); streamG.appendChild(pEl); \}}
\NormalTok{    \}}
\NormalTok{    leftSvg.appendChild(streamG);}

\NormalTok{    // Circle}
\NormalTok{    const [cirCx, cirCy] = toScreen({-}delta, 0, CX, CY, SCALE);}
\NormalTok{    const circleEl = document.createElementNS(svgNS,"circle");}
\NormalTok{    circleEl.setAttribute("cx",cirCx); circleEl.setAttribute("cy",cirCy); circleEl.setAttribute("r",r*SCALE);}
\NormalTok{    circleEl.setAttribute("fill","none"); circleEl.setAttribute("stroke","\#f97316"); circleEl.setAttribute("stroke{-}width","2");}
\NormalTok{    leftSvg.appendChild(circleEl);}

\NormalTok{    // Axes}
\NormalTok{    const [ox,oy]=toScreen(0,0,CX,CY,SCALE);}
\NormalTok{    const axG=document.createElementNS(svgNS,"g"); axG.setAttribute("stroke","\#9ca3af"); axG.setAttribute("stroke{-}width","1");}
\NormalTok{    const xa=document.createElementNS(svgNS,"line"); const [ax1]=toScreen({-}3,0,CX,CY,SCALE); const [ax2]=toScreen(3,0,CX,CY,SCALE);}
\NormalTok{    xa.setAttribute("x1",ax1); xa.setAttribute("y1",oy); xa.setAttribute("x2",ax2); xa.setAttribute("y2",oy); axG.appendChild(xa);}
\NormalTok{    const ya=document.createElementNS(svgNS,"line"); const [,ay1]=toScreen(0,{-}2,CX,CY,SCALE); const [,ay2]=toScreen(0,2,CX,CY,SCALE);}
\NormalTok{    ya.setAttribute("x1",ox); ya.setAttribute("y1",ay1); ya.setAttribute("x2",ox); ya.setAttribute("y2",ay2); axG.appendChild(ya);}
\NormalTok{    leftSvg.appendChild(axG);}

\NormalTok{    // Critical points at ±1}
\NormalTok{    [1,{-}1].forEach(xc =\textgreater{} \{}
\NormalTok{      const [cpx,cpy]=toScreen(xc,0,CX,CY,SCALE);}
\NormalTok{      const cpEl=document.createElementNS(svgNS,"circle"); cpEl.setAttribute("cx",cpx); cpEl.setAttribute("cy",cpy); cpEl.setAttribute("r","4");}
\NormalTok{      cpEl.setAttribute("fill","\#8b5cf6"); cpEl.setAttribute("stroke","\#fff"); cpEl.setAttribute("stroke{-}width","1.5");}
\NormalTok{      leftSvg.appendChild(cpEl);}
\NormalTok{    \});}

\NormalTok{    const lbl1=document.createElementNS(svgNS,"text"); lbl1.setAttribute("fill","\#6b7280"); lbl1.setAttribute("font{-}size","10"); lbl1.setAttribute("font{-}family","sans{-}serif");}
\NormalTok{    lbl1.setAttribute("x",5); lbl1.setAttribute("y",14); lbl1.textContent="z{-}plane";}
\NormalTok{    leftSvg.appendChild(lbl1);}

\NormalTok{    // Right panel: w{-}plane}
\NormalTok{    const rightSvg = document.createElementNS(svgNS, "svg");}
\NormalTok{    rightSvg.setAttribute("width", WP); rightSvg.setAttribute("height", HP);}
\NormalTok{    rightSvg.style.cssText = "background:\#f8fafc; border:1px solid \#e5e7eb; border{-}radius:6px;";}

\NormalTok{    const WCX = WP/2, WCY = HP/2;}
\NormalTok{    const WSCALE = 45;}

\NormalTok{    // Grid}
\NormalTok{    const wgridG=document.createElementNS(svgNS,"g"); wgridG.setAttribute("stroke","\#e5e7eb"); wgridG.setAttribute("stroke{-}width","0.5");}
\NormalTok{    for(let i={-}4;i\textless{}=4;i++) \{}
\NormalTok{      const xl2=document.createElementNS(svgNS,"line"); const [lx]=toScreen(i,0,WCX,WCY,WSCALE); const [,ly1]=toScreen(0,{-}3,WCX,WCY,WSCALE); const [,ly2]=toScreen(0,3,WCX,WCY,WSCALE);}
\NormalTok{      xl2.setAttribute("x1",lx); xl2.setAttribute("y1",ly1); xl2.setAttribute("x2",lx); xl2.setAttribute("y2",ly2); wgridG.appendChild(xl2);}
\NormalTok{      const yl2=document.createElementNS(svgNS,"line"); const [lx1]=toScreen({-}4,0,WCX,WCY,WSCALE); const [lx2]=toScreen(4,0,WCX,WCY,WSCALE); const [,ly]=toScreen(0,i,WCX,WCY,WSCALE);}
\NormalTok{      yl2.setAttribute("x1",lx1); yl2.setAttribute("y1",ly); yl2.setAttribute("x2",lx2); yl2.setAttribute("y2",ly); wgridG.appendChild(yl2);}
\NormalTok{    \}}
\NormalTok{    rightSvg.appendChild(wgridG);}

\NormalTok{    // w{-}plane axes}
\NormalTok{    const [wox,woy]=toScreen(0,0,WCX,WCY,WSCALE);}
\NormalTok{    const waxG=document.createElementNS(svgNS,"g"); waxG.setAttribute("stroke","\#9ca3af"); waxG.setAttribute("stroke{-}width","1");}
\NormalTok{    const wxa=document.createElementNS(svgNS,"line"); const [wax1]=toScreen({-}4,0,WCX,WCY,WSCALE); const [wax2]=toScreen(4,0,WCX,WCY,WSCALE);}
\NormalTok{    wxa.setAttribute("x1",wax1); wxa.setAttribute("y1",woy); wxa.setAttribute("x2",wax2); wxa.setAttribute("y2",woy); waxG.appendChild(wxa);}
\NormalTok{    const wya=document.createElementNS(svgNS,"line"); const [,way1]=toScreen(0,{-}3,WCX,WCY,WSCALE); const [,way2]=toScreen(0,3,WCX,WCY,WSCALE);}
\NormalTok{    wya.setAttribute("x1",wox); wya.setAttribute("y1",way1); wya.setAttribute("x2",wox); wya.setAttribute("y2",way2); waxG.appendChild(wya);}
\NormalTok{    rightSvg.appendChild(waxG);}

\NormalTok{    // Aerofoil curve}
\NormalTok{    if(aerofoilPts.length \textgreater{} 2) \{}
\NormalTok{      let aeroD = "";}
\NormalTok{      aerofoilPts.forEach(([wre,wim],i) =\textgreater{} \{}
\NormalTok{        const [px,py]=toScreen(wre,wim,WCX,WCY,WSCALE);}
\NormalTok{        aeroD += i===0 ? \textasciigrave{}M $\{px\} $\{py\}\textasciigrave{} : \textasciigrave{} L $\{px\} $\{py\}\textasciigrave{};}
\NormalTok{      \});}
\NormalTok{      aeroD += " Z";}
\NormalTok{      const aeroPath=document.createElementNS(svgNS,"path"); aeroPath.setAttribute("d",aeroD);}
\NormalTok{      aeroPath.setAttribute("fill","rgba(249,115,22,0.15)"); aeroPath.setAttribute("stroke","\#f97316"); aeroPath.setAttribute("stroke{-}width","2");}
\NormalTok{      rightSvg.appendChild(aeroPath);}
\NormalTok{    \}}

\NormalTok{    // Labels}
\NormalTok{    const wlbl=document.createElementNS(svgNS,"text"); wlbl.setAttribute("fill","\#6b7280"); wlbl.setAttribute("font{-}size","10"); wlbl.setAttribute("font{-}family","sans{-}serif");}
\NormalTok{    wlbl.setAttribute("x",5); wlbl.setAttribute("y",14); wlbl.textContent="w{-}plane";}
\NormalTok{    rightSvg.appendChild(wlbl);}

\NormalTok{    // Shape label}
\NormalTok{    let shapeLabel = "";}
\NormalTok{    if(r \textless{}= 1.02 \&\& delta \textless{} 0.02) shapeLabel = "Flat plate [−2,2]";}
\NormalTok{    else if(delta \textless{} 0.02) shapeLabel = "Ellipse";}
\NormalTok{    else shapeLabel = "Joukowski aerofoil";}
\NormalTok{    const shapeLbl=document.createElementNS(svgNS,"text"); shapeLbl.setAttribute("fill","\#f97316"); shapeLbl.setAttribute("font{-}size","10"); shapeLbl.setAttribute("font{-}family","sans{-}serif");}
\NormalTok{    shapeLbl.setAttribute("x",5); shapeLbl.setAttribute("y",HP{-}6); shapeLbl.textContent=shapeLabel;}
\NormalTok{    rightSvg.appendChild(shapeLbl);}

\NormalTok{    // Lift indicator}
\NormalTok{    const lift = circ.toFixed(2);}
\NormalTok{    const liftDiv=document.createElement("div");}
\NormalTok{    liftDiv.style.cssText="margin{-}top:0.4rem; font{-}family:sans{-}serif; font{-}size:0.85em; padding:0.5rem; background:\#f0f9ff; border{-}radius:4px;";}
\NormalTok{    liftDiv.innerHTML=\textasciigrave{}\textless{}strong\textgreater{}Kutta–Joukowski:\textless{}/strong\textgreater{} L = \textbackslash{}u03c1U\textbackslash{}u0393 \&nbsp;|\&nbsp; \textbackslash{}u0393 = $\{lift\} \&nbsp;|\&nbsp; \textless{}span style="color:\#1e40af"\textgreater{}Lift \textbackslash{}u221d \textbackslash{}u0393 = $\{(parseFloat(lift)).toFixed(2)\}\textless{}/span\textgreater{}}
\NormalTok{      \textless{}div style="color:\#6b7280; font{-}size:0.82em; margin{-}top:0.2rem;"\textgreater{}Purple dots on z{-}plane: critical points z=±1 (images become trailing/leading edges). Circulation \textbackslash{}u0393 = 0 gives symmetric flow; \textbackslash{}u0393 \textgreater{} 0 creates lift.\textless{}/div\textgreater{}\textasciigrave{};}

\NormalTok{    const panelsDiv=document.createElement("div");}
\NormalTok{    panelsDiv.style.cssText="display:flex; gap:0.5rem; flex{-}wrap:wrap; margin{-}top:0.5rem;";}
\NormalTok{    panelsDiv.appendChild(leftSvg);}
\NormalTok{    panelsDiv.appendChild(rightSvg);}

\NormalTok{    const wrap=document.createElement("div");}
\NormalTok{    wrap.appendChild(panelsDiv);}
\NormalTok{    wrap.appendChild(liftDiv);}
\NormalTok{    return wrap;}
\NormalTok{  \}}

\NormalTok{  const container=document.createElement("div");}
\NormalTok{  container.style.cssText="border:1px solid \#e5e7eb; border{-}radius:8px; padding:1rem; margin:1rem 0;";}
\NormalTok{  const title=document.createElement("div");}
\NormalTok{  title.style.cssText="font{-}weight:600; margin{-}bottom:0.5rem; font{-}family:sans{-}serif;";}
\NormalTok{  title.textContent="Joukowski map explorer: circle to aerofoil";}
\NormalTok{  container.appendChild(title);}
\NormalTok{  container.appendChild(rSlider4);}
\NormalTok{  container.appendChild(offsetSlider);}
\NormalTok{  container.appendChild(circSlider);}

\NormalTok{  const vizDiv4=document.createElement("div");}
\NormalTok{  container.appendChild(vizDiv4);}

\NormalTok{  const note4=document.createElement("div");}
\NormalTok{  note4.style.cssText="margin{-}top:0.5rem; font{-}size:0.82em; color:\#6b7280; font{-}style:italic;";}
\NormalTok{  note4.textContent="Left: circle in z{-}plane with uniform flow streamlines (blue). Right: image aerofoil in w{-}plane under w = z + 1/z. At r=1, \textbackslash{}u03b4=0: flat plate. Increase r for ellipse; increase \textbackslash{}u03b4 for asymmetric aerofoil. Purple dots mark critical points z=\textbackslash{}u00b11 (where the trailing edge forms). Circulation \textbackslash{}u0393 controls lift.";}
\NormalTok{  container.appendChild(note4);}

\NormalTok{  function update4() \{}
\NormalTok{    vizDiv4.innerHTML="";}
\NormalTok{    vizDiv4.appendChild(renderJoukowski(rSlider4.value, offsetSlider.value, circSlider.value));}
\NormalTok{  \}}
\NormalTok{  rSlider4.addEventListener("input", update4);}
\NormalTok{  offsetSlider.addEventListener("input", update4);}
\NormalTok{  circSlider.addEventListener("input", update4);}
\NormalTok{  update4();}

\NormalTok{  return container;}
\NormalTok{\}}
\end{Highlighting}
\end{Shaded}

\begin{center}\rule{0.5\linewidth}{0.5pt}\end{center}

\section{Application to potential
theory}\label{application-to-potential-theory}

The connection between conformal mapping and the Laplace equation makes
it the principal analytic tool for two-dimensional field problems.

\subsection{Complex potential}\label{complex-potential}

For a two-dimensional \textbf{incompressible irrotational} flow (or
electrostatic field, or heat conduction problem), the \textbf{complex
potential} \(F(z) = \Phi(x,y) + i\Psi(x,y)\) has: - \(\Phi\): velocity
potential (or electrostatic potential, or temperature) - \(\Psi\):
stream function (or flux lines)

The \textbf{complex velocity} is \(dF/dz = v_x - iv_y\), where
\((v_x, v_y)\) is the physical velocity vector. The sign convention
(\(v_x - iv_y\) rather than \(v_x + iv_y\)) is standard in potential
flow and follows from \(F = \Phi + i\Psi\) with \(v_x = \Phi_x\),
\(v_y = \Phi_y\).

Since \(F\) is analytic, both \(\Phi\) and \(\Psi\) satisfy Laplace's
equation, and they are harmonic conjugates.

\subsection{Transforming the domain}\label{transforming-the-domain}

If \(w = f(z)\) is a conformal map, and \(F(w)\) is the complex
potential in the \(w\)-domain, then \(F(f(z))\) is the complex potential
in the \(z\)-domain.

\textbf{This means:} solve Laplace's equation in a simple domain (disk,
half-plane, strip), map the domain to the complicated geometry using an
appropriate conformal map, and the composition gives the solution in the
complicated domain.

\textbf{Example.} Take \(G(w) = w\) (uniform flow in the \(w\)-plane)
and the conformal map \(f(z) = z^2\). The composed potential in the
\(z\)-plane is \(G(f(z)) = z^2\). That is, substitute \(w = z^2\)
directly: the result is a new function of \(z\), not \(G\) evaluated at
\(z\). With \(f(z) = z^2\), the \(z\)-plane flow has streamlines
\(\Psi = \text{Im}(z^2) = 2xy = \text{const}\) --- the rectangular
hyperbola pattern of flow in a 90° corner.

\begin{center}\rule{0.5\linewidth}{0.5pt}\end{center}

\section{Exercises}\label{exercises-40}

\begin{Shaded}
\begin{Highlighting}[]
\NormalTok{//| echo: false}
\NormalTok{katex = require("katex@0.16")}
\end{Highlighting}
\end{Shaded}

\begin{Shaded}
\begin{Highlighting}[]
\NormalTok{//| echo: false}
\NormalTok{function makeStepperHTML(steps) \{}
\NormalTok{  let current = 0;}
\NormalTok{  const totalSteps = steps.length;}
\NormalTok{  const container = document.createElement("div");}
\NormalTok{  container.style.cssText = "border:1px solid \#e5e7eb; border{-}radius:8px; padding:1rem 1.25rem; margin:0.75rem 0; font{-}family:inherit;";}
\NormalTok{  const stepsDiv = document.createElement("div");}
\NormalTok{  stepsDiv.style.marginBottom = "0.75rem";}
\NormalTok{  function renderSteps() \{}
\NormalTok{    stepsDiv.innerHTML = "";}
\NormalTok{    for (let i = 0; i \textless{} current; i++) \{}
\NormalTok{      const s = steps[i];}
\NormalTok{      const row = document.createElement("div");}
\NormalTok{      row.style.cssText = "display:grid; grid{-}template{-}columns:220px 1fr; gap:0.35rem 1rem; align{-}items:baseline; padding:0.35rem 0; border{-}top:1px solid \#f3f4f6;";}
\NormalTok{      const opCell = document.createElement("span");}
\NormalTok{      opCell.style.cssText = "font{-}size:0.85em; color:\#6b7280; font{-}style:italic;";}
\NormalTok{      opCell.textContent = s.op;}
\NormalTok{      const eqCell = document.createElement("span");}
\NormalTok{      eqCell.innerHTML = katex.renderToString(s.eq, \{ throwOnError: false, displayMode: false \});}
\NormalTok{      row.appendChild(opCell); row.appendChild(eqCell);}
\NormalTok{      if (s.note) \{}
\NormalTok{        const noteRow = document.createElement("div");}
\NormalTok{        noteRow.style.cssText = "grid{-}column:2; font{-}size:0.82em; color:\#6b7280; padding{-}bottom:0.2rem;";}
\NormalTok{        noteRow.textContent = s.note;}
\NormalTok{        row.appendChild(noteRow);}
\NormalTok{      \}}
\NormalTok{      stepsDiv.appendChild(row);}
\NormalTok{    \}}
\NormalTok{    if (current === totalSteps) \{}
\NormalTok{      const done = document.createElement("div");}
\NormalTok{      done.style.cssText = "margin{-}top:0.5rem; font{-}size:0.9em; color:\#059669; font{-}weight:500;";}
\NormalTok{      done.textContent = "Solution complete.";}
\NormalTok{      stepsDiv.appendChild(done);}
\NormalTok{    \}}
\NormalTok{  \}}
\NormalTok{  const controls = document.createElement("div");}
\NormalTok{  controls.style.cssText = "display:flex; gap:0.5rem; align{-}items:center;";}
\NormalTok{  const nextBtn = document.createElement("button");}
\NormalTok{  nextBtn.textContent = "Next step →";}
\NormalTok{  nextBtn.style.cssText = "padding:0.35rem 0.85rem; border:1px solid \#d1d5db; border{-}radius:4px; background:\#fff; cursor:pointer; font{-}size:0.9em;";}
\NormalTok{  const resetBtn = document.createElement("button");}
\NormalTok{  resetBtn.textContent = "Reset";}
\NormalTok{  resetBtn.style.cssText = "padding:0.35rem 0.75rem; border:1px solid \#e5e7eb; border{-}radius:4px; background:\#f9fafb; cursor:pointer; font{-}size:0.9em; color:\#6b7280;";}
\NormalTok{  const counter = document.createElement("span");}
\NormalTok{  counter.style.cssText = "font{-}size:0.82em; color:\#9ca3af; margin{-}left:0.25rem;";}
\NormalTok{  function updateButtons() \{}
\NormalTok{    nextBtn.disabled = current === totalSteps;}
\NormalTok{    nextBtn.style.opacity = current === totalSteps ? "0.4" : "1";}
\NormalTok{    counter.textContent = current === 0 ? "Click to reveal steps" : \textasciigrave{}Step $\{current\} of $\{totalSteps\}\textasciigrave{};}
\NormalTok{  \}}
\NormalTok{  nextBtn.onclick = () =\textgreater{} \{ if (current \textless{} totalSteps) \{ current++; renderSteps(); updateButtons(); \} \};}
\NormalTok{  resetBtn.onclick = () =\textgreater{} \{ current = 0; renderSteps(); updateButtons(); \};}
\NormalTok{  controls.appendChild(nextBtn); controls.appendChild(resetBtn); controls.appendChild(counter);}
\NormalTok{  container.appendChild(stepsDiv); container.appendChild(controls);}
\NormalTok{  renderSteps(); updateButtons();}
\NormalTok{  return container;}
\NormalTok{\}}
\end{Highlighting}
\end{Shaded}

\begin{center}\rule{0.5\linewidth}{0.5pt}\end{center}

\subsection{\texorpdfstring{Exercise 1: Image of a circle under
\(w = z^2\)}{Exercise 1: Image of a circle under w = z\^{}2}}\label{exercise-1-image-of-a-circle-under-w-z2}

Find the image of the circle \(|z| = 3\) under \(w = z^2\). Then find
the image of the half-circle \(|z| = 3\), \(\text{Im}(z) \geq 0\).

\begin{Shaded}
\begin{Highlighting}[]
\NormalTok{//| echo: false}
\NormalTok{makeStepperHTML([}
\NormalTok{  \{}
\NormalTok{    op: "Parametrise |z| = 3",}
\NormalTok{    eq: "z = 3e\^{}\{i\textbackslash{}\textbackslash{}theta\}, \textbackslash{}\textbackslash{}quad w = z\^{}2 = 9e\^{}\{2i\textbackslash{}\textbackslash{}theta\}"}
\NormalTok{  \},}
\NormalTok{  \{}
\NormalTok{    op: "Image of full circle",}
\NormalTok{    eq: "|w| = 9 \textbackslash{}\textbackslash{}text\{ and \} \textbackslash{}\textbackslash{}arg w = 2\textbackslash{}\textbackslash{}theta \textbackslash{}\textbackslash{}in [0,4\textbackslash{}\textbackslash{}pi) \textbackslash{}\textbackslash{}implies \textbackslash{}\textbackslash{}text\{image is \} |w|=9 \textbackslash{}\textbackslash{}text\{ traversed twice\}"}
\NormalTok{  \},}
\NormalTok{  \{}
\NormalTok{    op: "Image of upper half{-}circle: θ ∈ [0, π]",}
\NormalTok{    eq: "\textbackslash{}\textbackslash{}arg w = 2\textbackslash{}\textbackslash{}theta \textbackslash{}\textbackslash{}in [0, 2\textbackslash{}\textbackslash{}pi) \textbackslash{}\textbackslash{}implies \textbackslash{}\textbackslash{}text\{full circle \} |w| = 9 \textbackslash{}\textbackslash{}text\{ traversed once\}",}
\NormalTok{    note: "The upper semicircle maps onto the full circle |w| = 9 — the map w = z² is 2{-}to{-}1."}
\NormalTok{  \},}
\NormalTok{  \{}
\NormalTok{    op: "Critical points check",}
\NormalTok{    eq: "w\textquotesingle{}(z) = 2z = 0 \textbackslash{}\textbackslash{}text\{ only at \} z=0 \textbackslash{}\textbackslash{}notin |z|=3 \textbackslash{}\textbackslash{}implies \textbackslash{}\textbackslash{}text\{map is conformal on the circle\}",}
\NormalTok{    note: "Angles are preserved at every point of the circle."}
\NormalTok{  \}}
\NormalTok{])}
\end{Highlighting}
\end{Shaded}

\begin{center}\rule{0.5\linewidth}{0.5pt}\end{center}

\subsection{\texorpdfstring{Exercise 2: Image of a strip under
\(w = e^z\)}{Exercise 2: Image of a strip under w = e\^{}z}}\label{exercise-2-image-of-a-strip-under-w-ez}

Find the image of the horizontal strip \(\{0 < \text{Im}(z) < \pi/2\}\)
under \(w = e^z\).

\begin{Shaded}
\begin{Highlighting}[]
\NormalTok{//| echo: false}
\NormalTok{makeStepperHTML([}
\NormalTok{  \{}
\NormalTok{    op: "Write w = e\^{}\{x+iy\} = e\^{}x e\^{}\{iy\}",}
\NormalTok{    eq: "|w| = e\^{}x \textbackslash{}\textbackslash{}in (0, \textbackslash{}\textbackslash{}infty), \textbackslash{}\textbackslash{}quad \textbackslash{}\textbackslash{}arg w = y \textbackslash{}\textbackslash{}in \textbackslash{}\textbackslash{}left(0, \textbackslash{}\textbackslash{}frac\{\textbackslash{}\textbackslash{}pi\}\{2\}\textbackslash{}\textbackslash{}right)"}
\NormalTok{  \},}
\NormalTok{  \{}
\NormalTok{    op: "Image in the w{-}plane",}
\NormalTok{    eq: "\textbackslash{}\textbackslash{}left\textbackslash{}\textbackslash{}\{w : \textbackslash{}\textbackslash{}arg w \textbackslash{}\textbackslash{}in \textbackslash{}\textbackslash{}left(0, \textbackslash{}\textbackslash{}tfrac\{\textbackslash{}\textbackslash{}pi\}\{2\}\textbackslash{}\textbackslash{}right)\textbackslash{}\textbackslash{}right\textbackslash{}\textbackslash{}\} = \textbackslash{}\textbackslash{}text\{first quadrant of the \} w\textbackslash{}\textbackslash{}text\{{-}plane (excluding axes)\}"}
\NormalTok{  \},}
\NormalTok{  \{}
\NormalTok{    op: "Check boundary images",}
\NormalTok{    eq: "y=0 \textbackslash{}\textbackslash{}mapsto \textbackslash{}\textbackslash{}arg w = 0 \textbackslash{}\textbackslash{}text\{ (positive real axis)\};\textbackslash{}\textbackslash{}quad y=\textbackslash{}\textbackslash{}pi/2 \textbackslash{}\textbackslash{}mapsto \textbackslash{}\textbackslash{}arg w = \textbackslash{}\textbackslash{}pi/2 \textbackslash{}\textbackslash{}text\{ (positive imaginary axis)\}"}
\NormalTok{  \},}
\NormalTok{  \{}
\NormalTok{    op: "Conformality check: w′(z) = e\^{}z ≠ 0 everywhere",}
\NormalTok{    eq: "\textbackslash{}\textbackslash{}text\{The map is conformal throughout the strip.\}",}
\NormalTok{    note: "The strip\textquotesingle{}s right angles at its corner (at infinity) map to the right angle at w = 0 in the first quadrant. Conformality is consistent."}
\NormalTok{  \}}
\NormalTok{])}
\end{Highlighting}
\end{Shaded}

\begin{center}\rule{0.5\linewidth}{0.5pt}\end{center}

\subsection{Exercise 3: Möbius transformation from three point
pairs}\label{exercise-3-muxf6bius-transformation-from-three-point-pairs}

Find the Möbius transformation mapping \(z_1 = 0 \mapsto w_1 = i\),
\(z_2 = 1 \mapsto w_2 = 0\), \(z_3 = -1 \mapsto w_3 = \infty\).

\begin{Shaded}
\begin{Highlighting}[]
\NormalTok{//| echo: false}
\NormalTok{makeStepperHTML([}
\NormalTok{  \{}
\NormalTok{    op: "z₃ maps to ∞ means the denominator vanishes at z = −1",}
\NormalTok{    eq: "w = \textbackslash{}\textbackslash{}frac\{az+b\}\{c(z+1)\} \textbackslash{}\textbackslash{}quad \textbackslash{}\textbackslash{}text\{(denominator has root at \} z={-}1\textbackslash{}\textbackslash{}text\{)\}"}
\NormalTok{  \},}
\NormalTok{  \{}
\NormalTok{    op: "Apply w(0) = i",}
\NormalTok{    eq: "\textbackslash{}\textbackslash{}frac\{b\}\{c\} = i \textbackslash{}\textbackslash{}implies b = ic"}
\NormalTok{  \},}
\NormalTok{  \{}
\NormalTok{    op: "Apply w(1) = 0: numerator vanishes at z = 1",}
\NormalTok{    eq: "a(1) + b = 0 \textbackslash{}\textbackslash{}implies a = {-}b = {-}ic"}
\NormalTok{  \},}
\NormalTok{  \{}
\NormalTok{    op: "Write the transformation (set c = 1)",}
\NormalTok{    eq: "w = \textbackslash{}\textbackslash{}frac\{{-}iz + i\}\{z + 1\} = \textbackslash{}\textbackslash{}frac\{i(1{-}z)\}\{z+1\}"}
\NormalTok{  \},}
\NormalTok{  \{}
\NormalTok{    op: "Verify all three pairs",}
\NormalTok{    eq: "w(0) = i/1 = i\textbackslash{}\textbackslash{}checkmark, \textbackslash{}\textbackslash{}quad w(1) = 0/2 = 0\textbackslash{}\textbackslash{}checkmark, \textbackslash{}\textbackslash{}quad w({-}1) = i(2)/0 = \textbackslash{}\textbackslash{}infty\textbackslash{}\textbackslash{}checkmark"}
\NormalTok{  \}}
\NormalTok{])}
\end{Highlighting}
\end{Shaded}

\begin{center}\rule{0.5\linewidth}{0.5pt}\end{center}

\subsection{Exercise 4: Joukowski map --- circle to
ellipse}\label{exercise-4-joukowski-map-circle-to-ellipse}

Apply \(w = z + 1/z\) to the circle \(|z| = 2\). Find the semi-axes of
the resulting ellipse.

\begin{Shaded}
\begin{Highlighting}[]
\NormalTok{//| echo: false}
\NormalTok{makeStepperHTML([}
\NormalTok{  \{}
\NormalTok{    op: "Parametrise |z| = 2",}
\NormalTok{    eq: "z = 2e\^{}\{i\textbackslash{}\textbackslash{}theta\} \textbackslash{}\textbackslash{}implies w = 2e\^{}\{i\textbackslash{}\textbackslash{}theta\} + \textbackslash{}\textbackslash{}tfrac\{1\}\{2\}e\^{}\{{-}i\textbackslash{}\textbackslash{}theta\}"}
\NormalTok{  \},}
\NormalTok{  \{}
\NormalTok{    op: "Separate real and imaginary parts",}
\NormalTok{    eq: "w = \textbackslash{}\textbackslash{}left(2+\textbackslash{}\textbackslash{}tfrac\{1\}\{2\}\textbackslash{}\textbackslash{}right)\textbackslash{}\textbackslash{}cos\textbackslash{}\textbackslash{}theta + i\textbackslash{}\textbackslash{}left(2{-}\textbackslash{}\textbackslash{}tfrac\{1\}\{2\}\textbackslash{}\textbackslash{}right)\textbackslash{}\textbackslash{}sin\textbackslash{}\textbackslash{}theta = \textbackslash{}\textbackslash{}tfrac\{5\}\{2\}\textbackslash{}\textbackslash{}cos\textbackslash{}\textbackslash{}theta + i\textbackslash{}\textbackslash{}tfrac\{3\}\{2\}\textbackslash{}\textbackslash{}sin\textbackslash{}\textbackslash{}theta"}
\NormalTok{  \},}
\NormalTok{  \{}
\NormalTok{    op: "Identify as an ellipse",}
\NormalTok{    eq: "\textbackslash{}\textbackslash{}frac\{(\textbackslash{}\textbackslash{}text\{Re\}\textbackslash{}\textbackslash{},w)\^{}2\}\{(5/2)\^{}2\} + \textbackslash{}\textbackslash{}frac\{(\textbackslash{}\textbackslash{}text\{Im\}\textbackslash{}\textbackslash{},w)\^{}2\}\{(3/2)\^{}2\} = \textbackslash{}\textbackslash{}cos\^{}2\textbackslash{}\textbackslash{}theta + \textbackslash{}\textbackslash{}sin\^{}2\textbackslash{}\textbackslash{}theta = 1"}
\NormalTok{  \},}
\NormalTok{  \{}
\NormalTok{    op: "Semi{-}axes",}
\NormalTok{    eq: "a = \textbackslash{}\textbackslash{}tfrac\{5\}\{2\} = 2.5 \textbackslash{}\textbackslash{}text\{ (horizontal)\}, \textbackslash{}\textbackslash{}quad b = \textbackslash{}\textbackslash{}tfrac\{3\}\{2\} = 1.5 \textbackslash{}\textbackslash{}text\{ (vertical)\}",}
\NormalTok{    note: "General formula for |z| = r: a = r + 1/r, b = r − 1/r. As r → 1⁺, b → 0 and the ellipse collapses to the flat plate [−2, 2]."}
\NormalTok{  \}}
\NormalTok{])}
\end{Highlighting}
\end{Shaded}

\begin{center}\rule{0.5\linewidth}{0.5pt}\end{center}

\subsection{Exercise 5: Potential theory --- flow in a
corner}\label{exercise-5-potential-theory-flow-in-a-corner}

The complex potential \(F(w) = w\) represents uniform horizontal flow in
the \(w\)-plane. Apply the conformal map \(w = z^2\) to find the complex
potential for flow in the first quadrant (\(x > 0\), \(y > 0\)) of the
\(z\)-plane.

\begin{Shaded}
\begin{Highlighting}[]
\NormalTok{//| echo: false}
\NormalTok{makeStepperHTML([}
\NormalTok{  \{}
\NormalTok{    op: "Compose: F(z) = F(w)|\_\{w=z²\} = z²",}
\NormalTok{    eq: "F(z) = z\^{}2 = (x\^{}2{-}y\^{}2) + 2ixy"}
\NormalTok{  \},}
\NormalTok{  \{}
\NormalTok{    op: "Identify potential Φ and stream function Ψ",}
\NormalTok{    eq: "\textbackslash{}\textbackslash{}Phi = x\^{}2 {-} y\^{}2 \textbackslash{}\textbackslash{}text\{ (equipotential lines: rectangular hyperbolae)\}"}
\NormalTok{  \},}
\NormalTok{  \{}
\NormalTok{    op: "Stream function",}
\NormalTok{    eq: "\textbackslash{}\textbackslash{}Psi = 2xy = \textbackslash{}\textbackslash{}text\{const\} \textbackslash{}\textbackslash{}text\{ (streamlines: rectangular hyperbolae \} xy = C\textbackslash{}\textbackslash{}text\{)\}"}
\NormalTok{  \},}
\NormalTok{  \{}
\NormalTok{    op: "Verify Φ satisfies Laplace\textquotesingle{}s equation",}
\NormalTok{    eq: "\textbackslash{}\textbackslash{}Phi\_\{xx\} + \textbackslash{}\textbackslash{}Phi\_\{yy\} = 2 + ({-}2) = 0 \textbackslash{}\textbackslash{}checkmark"}
\NormalTok{  \},}
\NormalTok{  \{}
\NormalTok{    op: "Velocity field",}
\NormalTok{    eq: "\textbackslash{}\textbackslash{}frac\{dF\}\{dz\} = 2z = 2(x+iy) \textbackslash{}\textbackslash{}implies v\_x {-} iv\_y = 2z \textbackslash{}\textbackslash{}implies v\_x = 2x,\textbackslash{}\textbackslash{}; v\_y = {-}2y",}
\NormalTok{    note: "Flow moves outward along the x{-}axis and inward along the y{-}axis — the stagnation{-}point flow pattern seen in a jet impinging on a flat plate."}
\NormalTok{  \}}
\NormalTok{])}
\end{Highlighting}
\end{Shaded}

With these four chapters you now have the full analytic toolkit: complex
differentiation and the Cauchy-Riemann equations, contour integration
and Cauchy's theorem, Laurent series and residues, and conformal
mapping. The methods connect directly to the rest of engineering
mathematics: the Laplace equation in any two-dimensional domain,
potential flow around aerofoils, inverse Laplace transforms via
residues, and stability analysis in the complex frequency plane.

\part{Volume 7 --- Numerical methods}

\chapter{Numerical methods}\label{numerical-methods}

When the exact answer isn't the point

\hfill\break

The heat equation through an irregular fin. The trajectory of a
satellite perturbed by the moon's gravity. The implied volatility of an
options contract given its market price. Three equations you can write
down in seconds. None of them have closed-form solutions.

This is not unusual. Most equations that describe something real cannot
be solved exactly. Numerical methods are what you use instead ---
systematic procedures that produce answers you can trust, with errors
you can bound.

Three families of problems appear everywhere:

\begin{itemize}
\tightlist
\item
  \textbf{Root-finding}: find \(x\) such that \(f(x) = 0\).
\item
  \textbf{Interpolation}: given values at a set of points, estimate the
  value elsewhere.
\item
  \textbf{Numerical integration}: compute \(\int_a^b f(x)\,dx\) when
  \(f\) has no antiderivative in closed form.
\end{itemize}

Each family has the same structure: start with something computable,
iterate or accumulate, and keep track of how wrong you might be.

\begin{center}\rule{0.5\linewidth}{0.5pt}\end{center}

\section{Root-finding}\label{root-finding}

\subsection{The problem}\label{the-problem}

You have a function \(f(x)\) --- read ``\(f\) of \(x\)'', the function
whose root you want --- and you need to find \(x^*\) such that
\(f(x^*) = 0\). This is every nonlinear equation in disguise. If you
need to solve \(g(x) = h(x)\), define \(f(x) = g(x) - h(x)\) and find
its root.

\subsection{Bisection method}\label{bisection-method}

The simplest idea: if \(f(a) < 0\) and \(f(b) > 0\), and \(f\) is
continuous, then somewhere in \((a, b)\) the function must cross zero.
The midpoint \(m = (a+b)/2\) is either a root, or the root lies in
\((a, m)\), or it lies in \((m, b)\). Evaluate \(f(m)\) and keep the
half that changes sign. Repeat.

\textbf{Algorithm:}

\begin{enumerate}
\def\labelenumi{\arabic{enumi}.}
\tightlist
\item
  Choose \([a, b]\) with \(f(a)\) and \(f(b)\) having opposite signs.
\item
  Set \(m = (a + b)/2\).
\item
  If \(|f(m)| < \varepsilon\) (your tolerance), stop. \(m\) is the root.
\item
  If \(f(a) \cdot f(m) < 0\), the root is in \([a, m]\): set \(b = m\).
  Otherwise set \(a = m\).
\item
  Return to step 2.
\end{enumerate}

After \(n\) steps the bracket width is \((b - a)/2^n\), so the error is
bounded by \((b - a)/2^{n+1}\). This is the \textbf{bisection error
bound}. To guarantee an error below \(\varepsilon\), you need at most
\(\lceil \log_2((b-a)/\varepsilon) \rceil\) iterations.

The method is slow --- the error halves each step --- but it always
converges as long as the function is continuous and the initial bracket
is valid. Nothing else is assumed.

\textbf{Worked example.} Find the positive real root of
\(f(x) = x^3 - x - 2\), starting in \([1, 2]\).

Note that \(f(1) = 1 - 1 - 2 = -2 < 0\) and
\(f(2) = 8 - 2 - 2 = 4 > 0\). A root exists in \([1, 2]\).

\begin{longtable}[]{@{}
  >{\raggedright\arraybackslash}p{(\linewidth - 10\tabcolsep) * \real{0.1429}}
  >{\raggedright\arraybackslash}p{(\linewidth - 10\tabcolsep) * \real{0.1190}}
  >{\raggedright\arraybackslash}p{(\linewidth - 10\tabcolsep) * \real{0.1190}}
  >{\raggedright\arraybackslash}p{(\linewidth - 10\tabcolsep) * \real{0.1190}}
  >{\raggedright\arraybackslash}p{(\linewidth - 10\tabcolsep) * \real{0.1905}}
  >{\raggedright\arraybackslash}p{(\linewidth - 10\tabcolsep) * \real{0.3095}}@{}}
\toprule\noalign{}
\begin{minipage}[b]{\linewidth}\raggedright
Step
\end{minipage} & \begin{minipage}[b]{\linewidth}\raggedright
\(a\)
\end{minipage} & \begin{minipage}[b]{\linewidth}\raggedright
\(b\)
\end{minipage} & \begin{minipage}[b]{\linewidth}\raggedright
\(m\)
\end{minipage} & \begin{minipage}[b]{\linewidth}\raggedright
\(f(m)\)
\end{minipage} & \begin{minipage}[b]{\linewidth}\raggedright
New bracket
\end{minipage} \\
\midrule\noalign{}
\endhead
\bottomrule\noalign{}
\endlastfoot
1 & 1 & 2 & 1.5 & \(1.5^3 - 1.5 - 2 = -0.125\) & \([1.5, 2]\) \\
2 & 1.5 & 2 & 1.75 & \(1.75^3 - 1.75 - 2 \approx 1.609\) &
\([1.5, 1.75]\) \\
3 & 1.5 & 1.75 & 1.625 & \(1.625^3 - 1.625 - 2 \approx 0.666\) &
\([1.5, 1.625]\) \\
\end{longtable}

After three steps the root is bracketed in \([1.5, 1.625]\), width
\(0.125 = 1/2^3\). The exact root is \(x^* \approx 1.5214\).

\subsection{Newton-Raphson method}\label{newton-raphson-method}

Bisection is robust but slow. If \(f\) is differentiable, you can do
much better. The idea: at your current estimate \(x_n\) --- read ``\(x\)
sub \(n\)'', the \(n\)th iterate --- draw the tangent line to
\(y = f(x)\). Follow the tangent to where it hits zero; call that
\(x_{n+1}\).

The tangent line at \(x_n\) has equation:

\[y - f(x_n) = f'(x_n)(x - x_n)\]

where \(f'(x_n)\) --- read ``\(f\) prime at \(x_n\)'' --- is the
derivative of \(f\) at the current point. Setting \(y = 0\) and solving
for \(x\):

\[x_{n+1} = x_n - \frac{f(x_n)}{f'(x_n)}\]

This is the \textbf{Newton-Raphson iteration}. Apply it repeatedly until
\(|f(x_n)|\) is below your tolerance.

\textbf{Why it converges so fast.} The error \(e_n = x_n - x^*\)
satisfies:

Substituting a Taylor expansion of \(f(x_n)\) about the true root
\(x^*\) into the Newton iteration and simplifying:

\[e_{n+1} \approx -\frac{f''(x^*)}{2f'(x^*)} e_n^2\]

The next error is proportional to the \emph{square} of the current
error. If you are accurate to 3 decimal places, one more step gives you
roughly 6. This is \textbf{quadratic convergence} --- the number of
correct digits doubles each step. Bisection adds about 0.3 correct
digits per step; Newton-Raphson roughly doubles them.

The method fails if \(f'(x_n) = 0\) (the tangent is flat) or if the
starting point is far from the root (the iteration diverges or cycles).
Always check convergence.

\textbf{Worked example.} Compute \(\sqrt{2}\) using Newton-Raphson. The
positive root of \(f(x) = x^2 - 2\) is \(\sqrt{2}\), so \(f'(x) = 2x\).

The iteration is:

\[x_{n+1} = x_n - \frac{x_n^2 - 2}{2x_n} = \frac{x_n + 2/x_n}{2}\]

Starting from \(x_0 = 1\):

\begin{longtable}[]{@{}llll@{}}
\toprule\noalign{}
\(n\) & \(x_n\) & \(f(x_n) = x_n^2 - 2\) & \(|f(x_n)|\) \\
\midrule\noalign{}
\endhead
\bottomrule\noalign{}
\endlastfoot
0 & 1.000000 & −1.000000 & 1.000000 \\
1 & 1.500000 & 0.250000 & 0.250000 \\
2 & 1.416667 & 0.006944 & 0.006944 \\
3 & 1.414216 & 0.000006 & \(6 \times 10^{-6}\) \\
\end{longtable}

Three steps to six-digit accuracy. Quadratic convergence in action.

\subsection{Secant method}\label{secant-method}

Newton-Raphson requires the derivative \(f'(x)\). When that is expensive
or unavailable, replace it with a finite-difference approximation using
the two most recent iterates:

\[x_{n+1} = x_n - f(x_n) \cdot \frac{x_n - x_{n-1}}{f(x_n) - f(x_{n-1})}\]

This is the \textbf{secant method}. It needs two starting values \(x_0\)
and \(x_1\). The convergence order is approximately 1.618 (the golden
ratio) --- super-linear but slower than quadratic. No derivative
evaluation required.

\begin{tcolorbox}[enhanced jigsaw, opacityback=0, colback=white, breakable, colframe=quarto-callout-note-color-frame, arc=.35mm, left=2mm, toprule=.15mm, leftrule=.75mm, bottomrule=.15mm, rightrule=.15mm]
\begin{minipage}[t]{5.5mm}
\textcolor{quarto-callout-note-color}{\faInfo}
\end{minipage}%
\begin{minipage}[t]{\textwidth - 5.5mm}

\textbf{Which root-finder to use}

Bisection is your fallback: slow, guaranteed, needs only function values
and a sign change. Newton-Raphson is your workhorse when you have a
derivative and a decent starting point: fast, elegant, can fail. The
secant method sits between them: no derivative, faster than bisection,
needs two starting points. In practice, hybrid methods (Brent's method,
Illinois algorithm) combine bisection safety with Newton-like speed.

\end{minipage}%
\end{tcolorbox}

\begin{center}\rule{0.5\linewidth}{0.5pt}\end{center}

\section{Interpolation}\label{interpolation}

\subsection{The problem}\label{the-problem-1}

You have \(n+1\) data points
\((x_0, y_0), (x_1, y_1), \ldots, (x_n, y_n)\) and you want to estimate
\(f(x)\) at a new \(x\) lying between them. The classical approach: find
a polynomial that passes exactly through all the given points, then
evaluate it at \(x\).

\subsection{Lagrange interpolation}\label{lagrange-interpolation}

For \(n+1\) points there is exactly one polynomial of degree \(\leq n\)
passing through all of them. Write it as:

\[P_n(x) = \sum_{k=0}^{n} y_k \, L_k(x)\]

where each \textbf{Lagrange basis polynomial} \(L_k(x)\) is:

\[L_k(x) = \prod_{\substack{j=0 \\ j \neq k}}^{n} \frac{x - x_j}{x_k - x_j}\]

Read this as: \(L_k(x)\) equals 1 at \(x_k\) and equals 0 at every other
node \(x_j\). So \(P_n(x_k) = y_k\) for each \(k\) --- the polynomial
interpolates every point exactly.

\textbf{Worked example.} Interpolate \(f\) at \(x = 1.5\) using the
points \((1, 1)\), \((2, 8)\), \((3, 27)\) --- the cube function sampled
at three integers.

With nodes \(x_0 = 1\), \(x_1 = 2\), \(x_2 = 3\):

\[L_0(x) = \frac{(x-2)(x-3)}{(1-2)(1-3)} = \frac{(x-2)(x-3)}{2}\]

\[L_1(x) = \frac{(x-1)(x-3)}{(2-1)(2-3)} = \frac{(x-1)(x-3)}{-1} = -(x-1)(x-3)\]

\[L_2(x) = \frac{(x-1)(x-2)}{(3-1)(3-2)} = \frac{(x-1)(x-2)}{2}\]

At \(x = 1.5\):

\[L_0(1.5) = \frac{(1.5-2)(1.5-3)}{2} = \frac{(-0.5)(-1.5)}{2} = \frac{0.75}{2} = 0.375\]

\[L_1(1.5) = -(1.5-1)(1.5-3) = -(0.5)(-1.5) = 0.75\]

\[L_2(1.5) = \frac{(1.5-1)(1.5-2)}{2} = \frac{(0.5)(-0.5)}{2} = -0.125\]

\[P_2(1.5) = 1 \cdot 0.375 + 8 \cdot 0.75 + 27 \cdot (-0.125) = 0.375 + 6 - 3.375 = 3\]

The exact value is \(1.5^3 = 3.375\). The polynomial through three
points of \(x^3\) is not perfect --- it is a degree-2 polynomial
approximating a degree-3 function. Adding the fourth node \((4, 64)\)
would give exact interpolation.

\subsection{Newton's divided-difference
form}\label{newtons-divided-difference-form}

The Lagrange form is conceptually clean but computationally
inconvenient: to add a new node you must recompute everything. Newton's
divided-difference form builds the polynomial incrementally.

Define the \textbf{divided differences} recursively:

\[f[x_k] = y_k\]

\[f[x_k, x_{k+1}] = \frac{f[x_{k+1}] - f[x_k]}{x_{k+1} - x_k}\]

\[f[x_k, \ldots, x_{k+j}] = \frac{f[x_{k+1}, \ldots, x_{k+j}] - f[x_k, \ldots, x_{k+j-1}]}{x_{k+j} - x_k}\]

To see how this builds, use two of the cube-function points:
\(x_0 = 1\), \(y_0 = 1\) and \(x_1 = 2\), \(y_1 = 8\). The first divided
difference is \(f[x_0, x_1] = (8-1)/(2-1) = 7\). Newton's form gives
\(P_1(x) = 1 + 7(x - 1)\) --- the straight line through \((1,1)\) and
\((2,8)\). Adding a third node brings in the next divided difference and
bends the curve to pass through that point too.

The interpolating polynomial is then:

\[P_n(x) = f[x_0] + f[x_0,x_1](x - x_0) + f[x_0,x_1,x_2](x-x_0)(x-x_1) + \cdots\]

Adding a new node \(x_{n+1}\) requires only one new divided difference
and one new term --- the previous work is retained.

\subsection{Runge's phenomenon and why splines
exist}\label{runges-phenomenon-and-why-splines-exist}

High-degree polynomial interpolation over equally-spaced nodes can
oscillate wildly near the endpoints, even when the underlying function
is smooth. This is Runge's phenomenon. The standard fix is
\textbf{spline interpolation}: use low-degree polynomials (typically
cubics) on each sub-interval, stitched together with continuity
conditions. Splines are covered in the next chapter.

\begin{center}\rule{0.5\linewidth}{0.5pt}\end{center}

\section{Numerical integration}\label{numerical-integration}

\subsection{The problem}\label{the-problem-2}

Compute \(\int_a^b f(x)\,dx\) when \(f\) has no closed-form
antiderivative, or when \(f\) is known only at discrete points (sampled
data from a sensor, say).

Partition \([a, b]\) into \(n\) equal panels of width \(h = (b - a)/n\).
Let \(x_i = a + ih\) for \(i = 0, 1, \ldots, n\), and write
\(f_i = f(x_i)\).

\subsection{Trapezoid rule}\label{trapezoid-rule}

Approximate \(f\) on each panel by a straight line joining
\((x_i, f_i)\) and \((x_{i+1}, f_{i+1})\). The area of each trapezoid is
\(\frac{h}{2}(f_i + f_{i+1})\). Summing over all panels:

\[\int_a^b f(x)\,dx \approx T_n = \frac{h}{2}\bigl(f_0 + 2f_1 + 2f_2 + \cdots + 2f_{n-1} + f_n\bigr)\]

The error is \(O(h^2)\) --- halving the step size \(h\) reduces the
error by a factor of four. More precisely:

\[\text{Error} = -\frac{(b-a)h^2}{12} f''(\xi)\]

for some \(\xi \in (a,b)\) --- a point whose exact location is unknown,
but which lies somewhere in the integration interval. The formula says:
the error depends on the curvature of \(f\) somewhere in \((a,b)\). You
don't need to know where \(\xi\) is; you only need an upper bound on
\(|f''|\) over the interval. The error is proportional to the second
derivative of \(f\) --- large curvature means larger error.

\textbf{Worked example.} Estimate \(\int_0^1 e^x\,dx\) with \(n = 4\)
panels.

\(h = 0.25\), nodes at \(x = 0, 0.25, 0.5, 0.75, 1\).

\begin{longtable}[]{@{}ll@{}}
\toprule\noalign{}
\(x_i\) & \(f_i = e^{x_i}\) \\
\midrule\noalign{}
\endhead
\bottomrule\noalign{}
\endlastfoot
0 & 1.000000 \\
0.25 & 1.284025 \\
0.5 & 1.648721 \\
0.75 & 2.117000 \\
1 & 2.718282 \\
\end{longtable}

\[T_4 = \frac{0.25}{2}\bigl(1 + 2(1.284025) + 2(1.648721) + 2(2.117000) + 2.718282\bigr)\]

\[= 0.125 \times (1 + 2.568050 + 3.297442 + 4.234000 + 2.718282)\]

\[= 0.125 \times 13.817774 = 1.727222\]

The exact value is \(e - 1 \approx 1.718282\). Error:
\(|1.727222 - 1.718282| \approx 0.008940\).

\subsection{Simpson's 1/3 rule}\label{simpsons-13-rule}

Replace the straight-line approximation on each pair of panels with a
quadratic through three consecutive nodes. Require \(n\) even.

\[\int_a^b f(x)\,dx \approx S_n = \frac{h}{3}\bigl(f_0 + 4f_1 + 2f_2 + 4f_3 + \cdots + 4f_{n-1} + f_n\bigr)\]

The coefficients alternate \(4, 2, 4, 2, \ldots, 4\) for interior nodes.
The error is \(O(h^4)\):

\[\text{Error} = -\frac{(b-a)h^4}{180} f^{(4)}(\xi)\]

Halving \(h\) reduces the error by a factor of sixteen. Simpson's rule
is exact for polynomials of degree up to 3, even though it uses degree-2
approximations --- a consequence of symmetry.

\textbf{Worked example.} Apply Simpson's 1/3 rule to
\(\int_0^1 e^x\,dx\) with \(n = 4\).

Using the same nodes as above:

\[S_4 = \frac{0.25}{3}\bigl(f_0 + 4f_1 + 2f_2 + 4f_3 + f_4\bigr)\]

\[= \frac{0.25}{3}\bigl(1 + 4(1.284025) + 2(1.648721) + 4(2.117000) + 2.718282\bigr)\]

\[= \frac{0.25}{3}\bigl(1 + 5.136100 + 3.297442 + 8.468000 + 2.718282\bigr)\]

\[= \frac{0.25}{3} \times 20.619824 = 1.718319\]

Error: \(|1.718319 - 1.718282| \approx 0.000037\). A factor of 241
smaller than the trapezoid rule with the same number of function
evaluations.

\subsection{Simpson's 3/8 rule}\label{simpsons-38-rule}

A variant using groups of three panels instead of two, again with a
cubic approximant. Require \(n\) divisible by 3.

\[\int_a^b f(x)\,dx \approx \frac{3h}{8}\bigl(f_0 + 3f_1 + 3f_2 + 2f_3 + 3f_4 + 3f_5 + 2f_6 + \cdots + f_n\bigr)\]

Error is also \(O(h^4)\). The 1/3 rule is more commonly used; the 3/8
rule appears when \(n\) is odd so the standard Simpson pairing fails.

\subsection{Gaussian quadrature}\label{gaussian-quadrature}

The rules above use evenly-spaced nodes. If you are free to choose both
the node locations and the weights, you can do much better.

\textbf{Gauss-Legendre quadrature} on \([-1, 1]\) with \(n\) nodes is
exact for polynomials of degree up to \(2n-1\) --- twice the order you
might expect. The nodes are the roots of the Legendre polynomial
\(P_n(x)\), and the weights are determined by the condition of
exactness.

The Legendre polynomials are a specific family of functions whose roots
have an extremal property: placing quadrature nodes there squeezes the
most accuracy out of a given number of evaluations. You do not derive
these roots by hand --- they are tabulated in every numerical methods
reference and built into every serious quadrature library. For
\(n = 2\), the nodes and weights are:

For \(n = 2\): nodes \(x_{1,2} = \pm 1/\sqrt{3}\), weights
\(w_{1,2} = 1\):

\[\int_{-1}^1 f(x)\,dx \approx f\!\left(-\frac{1}{\sqrt{3}}\right) + f\!\left(\frac{1}{\sqrt{3}}\right)\]

This is exact for all polynomials of degree \(\leq 3\).

To integrate over \([a, b]\), substitute
\(x = \frac{a+b}{2} + \frac{b-a}{2}t\) where \(t \in [-1, 1]\):

\[\int_a^b f(x)\,dx = \frac{b-a}{2} \int_{-1}^1 f\!\left(\frac{a+b}{2} + \frac{b-a}{2}t\right) dt\]

\textbf{Worked example.} Estimate \(\int_0^1 e^x\,dx\) using 2-point
Gauss-Legendre.

Change variables to \([-1,1]\): \(x = \frac{1}{2} + \frac{1}{2}t\),
\(dx = \frac{1}{2}\,dt\).

Nodes in \(t\): \(t_{1,2} = \pm 1/\sqrt{3} \approx \pm 0.577350\).

Corresponding \(x\)-values: \(x_1 = 0.5 - 0.288675 = 0.211325\),
\(x_2 = 0.5 + 0.288675 = 0.788675\).

\[\int_0^1 e^x\,dx \approx \frac{1}{2}\Bigl(e^{0.211325} + e^{0.788675}\Bigr) = \frac{1}{2}(1.235369 + 2.200819) = 1.718094\]

Error: \(|1.718094 - 1.718282| \approx 0.000188\). Two function
evaluations; better accuracy than the 5-point trapezoid rule.

\begin{tcolorbox}[enhanced jigsaw, opacityback=0, colback=white, breakable, colframe=quarto-callout-note-color-frame, arc=.35mm, left=2mm, toprule=.15mm, leftrule=.75mm, bottomrule=.15mm, rightrule=.15mm]
\begin{minipage}[t]{5.5mm}
\textcolor{quarto-callout-note-color}{\faInfo}
\end{minipage}%
\begin{minipage}[t]{\textwidth - 5.5mm}

\textbf{Summary: error orders}

\begin{longtable}[]{@{}lll@{}}
\toprule\noalign{}
Method & Error & Error reduction when \(n\) doubles \\
\midrule\noalign{}
\endhead
\bottomrule\noalign{}
\endlastfoot
Trapezoid & \(O(h^2)\) & 4× smaller \\
Simpson's 1/3 & \(O(h^4)\) & 16× smaller \\
Gauss--Legendre (\(n\) pts) & exact for degree \(\leq 2n-1\) & --- \\
\end{longtable}

When \(f\) is smooth and you control the nodes, Gaussian quadrature
wins. When \(f\) comes from sampled data at fixed points, use Simpson or
trapezoid.

\end{minipage}%
\end{tcolorbox}

\begin{center}\rule{0.5\linewidth}{0.5pt}\end{center}

\section{Where this goes}\label{where-this-goes-30}

The three families here --- root-finding, interpolation, numerical
integration --- are the building blocks for almost everything else in
numerical computation. \textbf{Numerical linear algebra} (the next
chapter in this section) extends root-finding ideas to systems: you want
\(\mathbf{x}\) such that \(A\mathbf{x} = \mathbf{b}\), and for large
systems you solve it iteratively using methods (Gauss-Seidel, conjugate
gradient) that are direct descendants of the Newton-Raphson idea.

Numerical integration connects back to \textbf{ordinary differential
equations}: Euler's method for ODEs is the forward-rectangle
(left-endpoint) rule applied to the integral form of the ODE --- at each
step you approximate the integrand by its value at the left endpoint,
which is first-order accurate. Runge-Kutta methods are higher-order
quadrature applied at each step. Knowing the error orders from this
chapter tells you immediately why RK4 (fourth-order) is preferred over
Euler (first-order): halving the step size reduces RK4's error by 16×,
Euler's by only 2×.

\begin{tcolorbox}[enhanced jigsaw, opacityback=0, colback=white, breakable, colframe=quarto-callout-tip-color-frame, arc=.35mm, left=2mm, toprule=.15mm, leftrule=.75mm, bottomrule=.15mm, rightrule=.15mm]
\begin{minipage}[t]{5.5mm}
\textcolor{quarto-callout-tip-color}{\faLightbulb}
\end{minipage}%
\begin{minipage}[t]{\textwidth - 5.5mm}

\textbf{Where these methods appear}

\begin{itemize}
\tightlist
\item
  \textbf{Structural engineering (FEA):} each element stiffness matrix
  is assembled using Gaussian quadrature over the element domain.
\item
  \textbf{Signal processing:} interpolation underlies sample-rate
  conversion (resampling audio or images); Newton-Raphson is used in
  carrier phase estimation.
\item
  \textbf{Options pricing:} Newton-Raphson is the standard method for
  computing implied volatility from the Black-Scholes formula --- the
  exact inverse doesn't exist.
\item
  \textbf{Astrophysics:} orbital trajectories are integrated numerically
  step by step using adaptive-step Runge-Kutta, which uses error
  estimates derived from the quadrature theory here.
\item
  \textbf{Machine learning:} most optimisers (Adam, L-BFGS) are
  Newton-like methods applied to high-dimensional loss functions. The
  convergence order matters enormously at scale.
\end{itemize}

\end{minipage}%
\end{tcolorbox}

\begin{center}\rule{0.5\linewidth}{0.5pt}\end{center}

\section{Exercises}\label{exercises-41}

These are puzzles. Each has a clean numerical answer. The interesting
work is carrying the method through correctly and checking the error
bound.

\textbf{1.} Apply three iterations of bisection to find the positive
root of \(f(x) = x^3 - x - 2 = 0\), starting with the bracket
\([1, 2]\). State the midpoint after each iteration and the final
bracket width.

\begin{center}\rule{0.5\linewidth}{0.5pt}\end{center}

\textbf{2.} Apply three iterations of Newton-Raphson to compute
\(\sqrt{2}\), starting from \(x_0 = 1\). Use \(f(x) = x^2 - 2\). Record
\(x_n\) and \(|f(x_n)|\) at each step and comment on the convergence
rate.

\begin{center}\rule{0.5\linewidth}{0.5pt}\end{center}

\textbf{3.} Use the Lagrange interpolation formula to estimate
\(f(1.5)\), given the data points \((1, 1)\), \((2, 8)\), \((3, 27)\).
Compare with the exact value \(1.5^3 = 3.375\) and explain the
discrepancy.

\begin{center}\rule{0.5\linewidth}{0.5pt}\end{center}

\textbf{4.} Apply the trapezoid rule to \(\int_0^1 e^x\,dx\) with
\(n = 4\) panels. Compute the numerical result and the absolute error
against the exact value \(e - 1 \approx 1.718282\).

\begin{center}\rule{0.5\linewidth}{0.5pt}\end{center}

\textbf{5.} Apply Simpson's 1/3 rule to the same integral
\(\int_0^1 e^x\,dx\) with \(n = 4\) panels. Compare your error with
Exercise 4 and confirm the improvement is consistent with the \(O(h^4)\)
order.

\begin{center}\rule{0.5\linewidth}{0.5pt}\end{center}

\textbf{6.} A diode circuit satisfies
\(I \cdot R + V_T \ln(I / I_0) = V_s\), where \(R = 100\,\Omega\),
\(V_T = 0.026\,\text{V}\), \(I_0 = 10^{-12}\,\text{A}\),
\(V_s = 5\,\text{V}\). Starting from \(I_0^* = 0.04\,\text{A}\), apply
one step of Newton-Raphson to find a better estimate. \emph{(Linearise
the equation about the starting point.)}

\chapter{Numerical linear algebra}\label{numerical-linear-algebra}

Solving systems that can't be solved by hand

\hfill\break

You have a system of equations. In principle, Gaussian elimination gives
the answer. In practice, you type it into a computer, and the answer is
wrong --- not obviously wrong, but subtly, fatally wrong.

The issue isn't the algorithm. The issue is that computers don't do
exact arithmetic. Every floating-point operation introduces a tiny
rounding error. Tiny errors compound. In a badly conditioned system, a
rounding error in the twelfth decimal place can shift the answer by a
factor of a thousand.

This chapter is about how to do it right: pivoting to avoid catastrophic
cancellation, LU decomposition to separate the expensive part from the
cheap part, and the condition number to know in advance whether the
answer you're computing is worth trusting.

\section{What the notation is
saying}\label{what-the-notation-is-saying-23}

A linear system \(A\mathbf{x} = \mathbf{b}\) means: find the vector
\(\mathbf{x}\) --- read as ``the solution vector'' --- such that the
matrix \(A\), read as ``the coefficient matrix'', transforms
\(\mathbf{x}\) into the right-hand side \(\mathbf{b}\).

In component form for a \(3 \times 3\) system:

\[\begin{pmatrix} a_{11} & a_{12} & a_{13} \\ a_{21} & a_{22} & a_{23} \\ a_{31} & a_{32} & a_{33} \end{pmatrix} \begin{pmatrix} x_1 \\ x_2 \\ x_3 \end{pmatrix} = \begin{pmatrix} b_1 \\ b_2 \\ b_3 \end{pmatrix}\]

The goal is to reduce \(A\) to upper-triangular form so that
back-substitution can extract \(x_3\), then \(x_2\), then \(x_1\) in
sequence.

\textbf{LU decomposition} writes \(A = LU\), where \(L\) is a
lower-triangular matrix (ones on the diagonal, entries only below) and
\(U\) is the upper-triangular matrix produced by elimination. This
notation should be read as: ``\(A\) factors into \(L\) times \(U\)''.

The \textbf{condition number} \(\kappa(A)\), read ``kappa of \(A\)'',
quantifies how much the solution can change relative to a small change
in the data. It is defined as

\[\kappa(A) = \|A\|_\infty \cdot \|A^{-1}\|_\infty\]

where \(\|A\|_\infty\), read ``the infinity-norm of \(A\)'', is the
largest row sum of absolute values:

\[\|A\|_\infty = \max_i \sum_j |a_{ij}|\]

A rough rule: if \(\kappa(A) \approx 10^k\), the solution loses
approximately \(k\) digits of accuracy to rounding errors. A condition
number of \(10^{12}\) on a machine with 15-digit floating-point
precision leaves you with about 3 reliable digits.

For iterative methods, \(x^{(k)}\) denotes the iterate at step \(k\) ---
the approximate solution after \(k\) steps, converging toward the true
\(\mathbf{x}\) as \(k\) grows.

\section{Gaussian elimination with partial
pivoting}\label{gaussian-elimination-with-partial-pivoting}

\subsection{Why naive elimination
fails}\label{why-naive-elimination-fails}

The exact solution of the system below is \(x_1 = 1\), \(x_2 = 1\) ---
you can check: \(0.0001(1) + 1(1) = 1.0001 \approx 1\) and
\(1(1) + 1(1) = 2\). Now watch what happens when a computer solves it.

Consider this \(2 \times 2\) system in floating-point arithmetic with 4
significant figures:

\[\begin{pmatrix} 0.0001 & 1 \\ 1 & 1 \end{pmatrix} \begin{pmatrix} x_1 \\ x_2 \end{pmatrix} = \begin{pmatrix} 1 \\ 2 \end{pmatrix}\]

Naive elimination: subtract \((1 / 0.0001) = 10000\) times row 1 from
row 2.

Row 2 becomes:
\[[1 - 10000 \times 0.0001,\quad 1 - 10000 \times 1] = [0,\quad 1 - 10000]\]

The \(1 - 10000\) gives \(-9999\), fine. But consider the right-hand
side: \(2 - 10000 \times 1 = -9998\). Back-substituting:
\(x_2 = -9998 / {-9999} \approx 1.0001\). Then
\(x_1 = (1 - 1 \times 1.0001)/0.0001 \approx -1\).

The exact answer is \(x_2 \approx 1\), \(x_1 \approx 1\). Naive
elimination gives \(x_1 \approx -1\). That is not a small rounding error
--- that is the wrong sign with the wrong magnitude. The problem:
dividing by the tiny pivot \(0.0001\) amplified the rounding error in
the subtracted row by a factor of \(10000\). The error in the fourth
decimal place became the dominant term in the result.

This is \textbf{catastrophic cancellation}: subtracting two nearly equal
large numbers, leaving a result that is almost entirely rounding error.
It is not a pathological edge case. Any system with a small pivot
relative to the other entries is vulnerable --- and such systems arise
routinely from physical models with variables at very different scales.

\subsection{Partial pivoting: the fix}\label{partial-pivoting-the-fix}

Before eliminating column \(j\), scan the entries in column \(j\) from
row \(j\) downward. Swap the row with the \textbf{largest absolute
value} in that column into the pivot position. This ensures the
multiplier \(m_{ij} = a_{ij}/a_{jj}\) satisfies \(|m_{ij}| \leq 1\),
bounding the growth of rounding errors.

\textbf{Algorithm} (partial pivoting GE on \(n \times n\) augmented
matrix \([A|\mathbf{b}]\)):

For \(j = 1, 2, \ldots, n-1\):

\begin{enumerate}
\def\labelenumi{\arabic{enumi}.}
\tightlist
\item
  Find \(p = \arg\max_{i \geq j} |a_{ij}|\) (index of largest entry in
  column \(j\), rows \(j\) to \(n\))
\item
  If \(p \neq j\): swap rows \(p\) and \(j\) in \([A|\mathbf{b}]\)
\item
  For \(i = j+1, \ldots, n\):

  \begin{itemize}
  \tightlist
  \item
    Compute multiplier \(m_{ij} = a_{ij}/a_{jj}\)
  \item
    Row \(i \leftarrow\) Row \(i\) \(- m_{ij} \times\) Row \(j\)
  \end{itemize}
\end{enumerate}

After \(n-1\) passes, \(A\) is upper triangular. Recover \(\mathbf{x}\)
by back-substitution:

\[x_n = \frac{b_n}{a_{nn}}, \qquad x_i = \frac{1}{a_{ii}}\left(b_i - \sum_{k=i+1}^n a_{ik} x_k\right), \quad i = n-1, \ldots, 1\]

\subsection{Worked example: partial pivoting on a 3×3
system}\label{worked-example-partial-pivoting-on-a-33-system}

Solve:

\[\begin{pmatrix} 1 & 2 & 1 \\ 3 & 2 & 4 \\ 2 & 1 & 2 \end{pmatrix} \begin{pmatrix} x_1 \\ x_2 \\ x_3 \end{pmatrix} = \begin{pmatrix} 8 \\ 20 \\ 12 \end{pmatrix}\]

Write the augmented matrix:

\[[A|\mathbf{b}] = \left[\begin{array}{ccc|c} 1 & 2 & 1 & 8 \\ 3 & 2 & 4 & 20 \\ 2 & 1 & 2 & 12 \end{array}\right]\]

\textbf{Pass 1 (column 1):} The largest entry in column 1 is \(|3|\) in
row 2. Swap rows 1 and 2:

\[\left[\begin{array}{ccc|c} 3 & 2 & 4 & 20 \\ 1 & 2 & 1 & 8 \\ 2 & 1 & 2 & 12 \end{array}\right]\]

Eliminate below pivot \(a_{11} = 3\):

\begin{itemize}
\tightlist
\item
  Row 2 \(\leftarrow\) Row 2 \(- (1/3)\) Row 1:
  \([1 - 1, \; 2 - 2/3, \; 1 - 4/3 \;|\; 8 - 20/3] = [0, \; 4/3, \; -1/3 \;|\; 4/3]\)
\item
  Row 3 \(\leftarrow\) Row 3 \(- (2/3)\) Row 1:
  \([2 - 2, \; 1 - 4/3, \; 2 - 8/3 \;|\; 12 - 40/3] = [0, \; -1/3, \; -2/3 \;|\; -4/3]\)
\end{itemize}

\[\left[\begin{array}{ccc|c} 3 & 2 & 4 & 20 \\ 0 & 4/3 & -1/3 & 4/3 \\ 0 & -1/3 & -2/3 & -4/3 \end{array}\right]\]

\textbf{Pass 2 (column 2):} Largest in column 2, rows 2--3:
\(|4/3| > |-1/3|\). No swap needed. Pivot is \(4/3\).

\begin{itemize}
\tightlist
\item
  Row 3 \(\leftarrow\) Row 3 \(- (-1/3)/(4/3) \times\) Row 2 \(=\) Row 3
  \(+ (1/4)\) Row 2:
  \([0, \; -1/3 + 1/3, \; -2/3 - 1/12 \;|\; -4/3 + 1/3] = [0, \; 0, \; -3/4 \;|\; -1]\)
\end{itemize}

\[\left[\begin{array}{ccc|c} 3 & 2 & 4 & 20 \\ 0 & 4/3 & -1/3 & 4/3 \\ 0 & 0 & -3/4 & -1 \end{array}\right]\]

\textbf{Back-substitution:}

\[x_3 = \frac{-1}{-3/4} = \frac{4}{3}\]

\[x_2 = \frac{4/3 - (-1/3)(4/3)}{4/3} = \frac{4/3 + 4/9}{4/3} = \frac{16/9}{4/3} = \frac{16}{9} \cdot \frac{3}{4} = \frac{4}{3}\]

\[x_1 = \frac{20 - 2(4/3) - 4(4/3)}{3} = \frac{20 - 8/3 - 16/3}{3} = \frac{20 - 8}{3} = \frac{12}{3} = 4\]

\textbf{Solution:} \(x_1 = 4\), \(x_2 = 4/3\), \(x_3 = 4/3\).

\textbf{Check:} Row 1:
\(1(4) + 2(4/3) + 1(4/3) = 4 + 8/3 + 4/3 = 4 + 4 = 8\) ✓ Row 2:
\(3(4) + 2(4/3) + 4(4/3) = 12 + 8/3 + 16/3 = 12 + 8 = 20\) ✓ Row 3:
\(2(4) + 1(4/3) + 2(4/3) = 8 + 4/3 + 8/3 = 8 + 4 = 12\) ✓

\section{LU decomposition}\label{lu-decomposition}

Gaussian elimination --- without pivoting, for clarity --- can be read
as computing a factorisation \(A = LU\). The matrix \(L\) records the
multipliers used during elimination; \(U\) is the upper-triangular
result.

\subsection{Doolittle form}\label{doolittle-form}

In the Doolittle convention, \(L\) has \textbf{ones on its diagonal}:

\[L = \begin{pmatrix} 1 & 0 & 0 \\ m_{21} & 1 & 0 \\ m_{31} & m_{32} & 1 \end{pmatrix}, \qquad U = \begin{pmatrix} u_{11} & u_{12} & u_{13} \\ 0 & u_{22} & u_{23} \\ 0 & 0 & u_{33} \end{pmatrix}\]

where \(m_{ij}\) is the multiplier used to zero \(a_{ij}\) during
elimination.

The \(U\) matrix is exactly the upper-triangular matrix at the end of
Gaussian elimination. The \(L\) matrix costs nothing extra --- the
multipliers are computed anyway and simply stored in the lower portion
of the working array.

\subsection{Why LU earns its keep}\label{why-lu-earns-its-keep}

Solving \(A\mathbf{x} = \mathbf{b}\) once costs \(O(n^3)\) operations
--- where ``operations'' means floating-point additions and
multiplications, roughly \(n^3/3\) in total, because each elimination
step touches every remaining row-column pair. This is unavoidable. But
if you need to solve the same \(A\) for multiple right-hand sides
\(\mathbf{b}_1, \mathbf{b}_2, \ldots\), the LU factorisation pays back:

\begin{itemize}
\tightlist
\item
  \textbf{Build} \(A = LU\): \(O(n^3)\) --- done once.
\item
  \textbf{Solve} \(L\mathbf{y} = \mathbf{b}\) (forward substitution):
  \(O(n^2)\) per right-hand side.
\item
  \textbf{Solve} \(U\mathbf{x} = \mathbf{y}\) (back substitution):
  \(O(n^2)\) per right-hand side.
\end{itemize}

For \(k\) right-hand sides the total cost is
\(O(n^3) + k \cdot O(n^2)\), versus \(k \cdot O(n^3)\) if you
re-eliminate each time. For a \(1000 \times 1000\) system solved for 50
right-hand sides, that's roughly \(50 \times\) faster.

This matters in FEM analysis, where the stiffness matrix \(K\) is fixed
while the load vector \(\mathbf{f}\) changes with each load case.

\subsection{Worked example: Doolittle LU on a 3×3
system}\label{worked-example-doolittle-lu-on-a-33-system}

Factorise and solve:

\[A = \begin{pmatrix} 2 & 1 & 1 \\ 4 & 3 & 3 \\ 8 & 7 & 9 \end{pmatrix}, \quad \mathbf{b}_1 = \begin{pmatrix} 1 \\ 1 \\ 1 \end{pmatrix}, \quad \mathbf{b}_2 = \begin{pmatrix} 2 \\ 0 \\ 4 \end{pmatrix}\]

\textbf{Step 1: Build \(U\) and record multipliers.}

Column 1, pivot \(u_{11} = 2\): - \(m_{21} = 4/2 = 2\). Row 2
\(\leftarrow\) Row 2 \(-\) 2 Row 1: \([0, 1, 1 \;|\; \cdots]\) -
\(m_{31} = 8/2 = 4\). Row 3 \(\leftarrow\) Row 3 \(-\) 4 Row 1:
\([0, 3, 5 \;|\; \cdots]\)

Column 2, pivot \(u_{22} = 1\): - \(m_{32} = 3/1 = 3\). Row 3
\(\leftarrow\) Row 3 \(-\) 3 Row 2: \([0, 0, 2 \;|\; \cdots]\)

\[U = \begin{pmatrix} 2 & 1 & 1 \\ 0 & 1 & 1 \\ 0 & 0 & 2 \end{pmatrix}, \qquad L = \begin{pmatrix} 1 & 0 & 0 \\ 2 & 1 & 0 \\ 4 & 3 & 1 \end{pmatrix}\]

\textbf{Verify:} \(LU\) should equal \(A\).

\[LU = \begin{pmatrix}1&0&0\\2&1&0\\4&3&1\end{pmatrix}\begin{pmatrix}2&1&1\\0&1&1\\0&0&2\end{pmatrix} = \begin{pmatrix}2&1&1\\4&3&3\\8&7&9\end{pmatrix} = A \checkmark\]

\textbf{Step 2: Solve \(L\mathbf{y} = \mathbf{b}_1\) (forward
substitution).}

\[\begin{pmatrix}1&0&0\\2&1&0\\4&3&1\end{pmatrix}\begin{pmatrix}y_1\\y_2\\y_3\end{pmatrix} = \begin{pmatrix}1\\1\\1\end{pmatrix}\]

\(y_1 = 1\), \(\quad y_2 = 1 - 2(1) = -1\),
\(\quad y_3 = 1 - 4(1) - 3(-1) = 0\)

\textbf{Step 3: Solve \(U\mathbf{x} = \mathbf{y}\) (back substitution).}

\[\begin{pmatrix}2&1&1\\0&1&1\\0&0&2\end{pmatrix}\begin{pmatrix}x_1\\x_2\\x_3\end{pmatrix} = \begin{pmatrix}1\\-1\\0\end{pmatrix}\]

\(x_3 = 0/2 = 0\), \(\quad x_2 = (-1 - 1 \cdot 0)/1 = -1\),
\(\quad x_1 = (1 - 1(-1) - 1(0))/2 = 1\)

\textbf{Solution for \(\mathbf{b}_1\):} \(\mathbf{x} = (1, -1, 0)^T\)

\textbf{Check:} \(2(1) + 1(-1) + 1(0) = 1\) ✓,
\(\quad 4(1) + 3(-1) + 3(0) = 1\) ✓, \(\quad 8(1) + 7(-1) + 9(0) = 1\) ✓

\textbf{Step 4: Reuse \(L\), \(U\) for \(\mathbf{b}_2 = (2, 0, 4)^T\).}

Forward substitution \(L\mathbf{y} = \mathbf{b}_2\):

\(y_1 = 2\), \(\quad y_2 = 0 - 2(2) = -4\),
\(\quad y_3 = 4 - 4(2) - 3(-4) = 4 - 8 + 12 = 8\)

Back substitution \(U\mathbf{x} = \mathbf{y}\):

\(x_3 = 8/2 = 4\), \(\quad x_2 = (-4 - 1 \cdot 4)/1 = -8\),
\(\quad x_1 = (2 - 1(-8) - 1(4))/2 = (2 + 8 - 4)/2 = 3\)

\textbf{Solution for \(\mathbf{b}_2\):} \(\mathbf{x} = (3, -8, 4)^T\)

\textbf{Check:} \(2(3) + 1(-8) + 1(4) = 6 - 8 + 4 = 2\) ✓,
\(\quad 4(3) + 3(-8) + 3(4) = 12 - 24 + 12 = 0\) ✓,
\(\quad 8(3) + 7(-8) + 9(4) = 24 - 56 + 36 = 4\) ✓

\begin{tcolorbox}[enhanced jigsaw, opacityback=0, colback=white, breakable, colframe=quarto-callout-note-color-frame, arc=.35mm, left=2mm, toprule=.15mm, leftrule=.75mm, bottomrule=.15mm, rightrule=.15mm]
\begin{minipage}[t]{5.5mm}
\textcolor{quarto-callout-note-color}{\faInfo}
\end{minipage}%
\begin{minipage}[t]{\textwidth - 5.5mm}

\textbf{Why LU works}

Solving \(A\mathbf{x} = \mathbf{b}\) is equivalent to solving
\(LU\mathbf{x} = \mathbf{b}\). Let \(\mathbf{y} = U\mathbf{x}\). Then
\(L\mathbf{y} = \mathbf{b}\) --- a lower-triangular system, solved in
one forward pass. Then \(U\mathbf{x} = \mathbf{y}\) --- an
upper-triangular system, solved in one backward pass. The original
\(O(n^3)\) work goes into building \(L\) and \(U\). Once those are in
hand, each new right-hand side costs only \(O(n^2)\).

\end{minipage}%
\end{tcolorbox}

\section{Condition number}\label{condition-number}

\subsection{Definition}\label{definition-4}

The condition number \(\kappa(A)\) measures how much relative error in
\(\mathbf{b}\) is amplified into relative error in \(\mathbf{x}\).

Using the \(\infty\)-norm:

\[\kappa(A) = \|A\|_\infty \cdot \|A^{-1}\|_\infty\]

where \(\|A\|_\infty = \max_i \sum_j |a_{ij}|\) is the maximum absolute
row sum.

A system is \textbf{well-conditioned} if \(\kappa(A)\) is small (close
to 1). It is \textbf{ill-conditioned} if \(\kappa(A) \gg 1\).

The digits-of-accuracy rule: with \(t\)-digit floating-point arithmetic,
you can expect roughly \(t - \log_{10}\kappa(A)\) correct digits in the
solution.

\subsection{Ill-conditioning in practice: the 3×3 Hilbert
matrix}\label{ill-conditioning-in-practice-the-33-hilbert-matrix}

The Hilbert matrix \(H\) with entries \(h_{ij} = 1/(i+j-1)\) is the
classic ill-conditioning example:

\[H_3 = \begin{pmatrix} 1 & 1/2 & 1/3 \\ 1/2 & 1/3 & 1/4 \\ 1/3 & 1/4 & 1/5 \end{pmatrix}\]

\textbf{Step 1: \(\|H_3\|_\infty\).}

Row sums: Row 1: \(1 + 0.5 + 0.333 = 1.833\); Row 2:
\(0.5 + 0.333 + 0.25 = 1.083\); Row 3: \(0.333 + 0.25 + 0.2 = 0.783\).

\[\|H_3\|_\infty = 1.833\]

\textbf{Step 2: \(\|H_3^{-1}\|_\infty\).}

The exact inverse of \(H_3\) is:

\[H_3^{-1} = \begin{pmatrix} 9 & -36 & 30 \\ -36 & 192 & -180 \\ 30 & -180 & 180 \end{pmatrix}\]

Row sums of \(|H_3^{-1}|\): Row 1: \(9 + 36 + 30 = 75\); Row 2:
\(36 + 192 + 180 = 408\); Row 3: \(30 + 180 + 180 = 390\).

\[\|H_3^{-1}\|_\infty = 408\]

\textbf{Step 3: \(\kappa_\infty(H_3)\).}

\[\kappa_\infty(H_3) = 1.833 \times 408 \approx 748\]

So solving \(H_3 \mathbf{x} = \mathbf{b}\) loses about
\(\log_{10}(748) \approx 2.9\) digits of accuracy. On a 15-digit machine
you get about 12 reliable digits --- acceptable. For \(H_{10}\),
\(\kappa \approx 10^{13}\), leaving roughly 2 reliable digits. That
system cannot be usefully solved by direct methods in standard double
precision.

\begin{tcolorbox}[enhanced jigsaw, opacityback=0, colback=white, breakable, colframe=quarto-callout-warning-color-frame, arc=.35mm, left=2mm, toprule=.15mm, leftrule=.75mm, bottomrule=.15mm, rightrule=.15mm]
\begin{minipage}[t]{5.5mm}
\textcolor{quarto-callout-warning-color}{\faExclamationTriangle}
\end{minipage}%
\begin{minipage}[t]{\textwidth - 5.5mm}

\textbf{When to worry}

If \(\kappa(A) > 10^6\) in double precision, inspect the problem
formulation before trusting any solution. The issue usually lies in the
model (ill-scaled variables, near-dependent equations) rather than the
solver. Rescaling rows and columns (equilibration) often reduces
\(\kappa\) by several orders of magnitude.

\end{minipage}%
\end{tcolorbox}

\section{Iterative methods}\label{iterative-methods}

Direct methods (GE, LU) compute the exact answer (up to rounding) in a
fixed number of operations. Iterative methods start from an initial
guess and update it repeatedly. They're preferred when \(A\) is large
and sparse --- a \(10^6 \times 10^6\) FEM stiffness matrix, for example,
where storing \(L\) and \(U\) in full is impractical.

\subsection{Jacobi iteration}\label{jacobi-iteration}

Rewrite each equation for its diagonal variable, holding all other
variables fixed at their current values.

For row \(i\) of \(A\mathbf{x} = \mathbf{b}\):

\[x_i^{(k+1)} = \frac{1}{a_{ii}}\left(b_i - \sum_{j \neq i} a_{ij} x_j^{(k)}\right)\]

All updates use values from the \textbf{previous} iterate \(x^{(k)}\).
This makes Jacobi naturally parallelisable --- each new \(x_i^{(k+1)}\)
is independent of the others being computed in the same step.

\subsection{Gauss-Seidel iteration}\label{gauss-seidel-iteration}

The same formula, but use the \textbf{latest available} values --- if
\(x_j^{(k+1)}\) has already been computed in this sweep, use it
immediately:

\[x_i^{(k+1)} = \frac{1}{a_{ii}}\left(b_i - \sum_{j < i} a_{ij} x_j^{(k+1)} - \sum_{j > i} a_{ij} x_j^{(k)}\right)\]

Gauss-Seidel typically converges in roughly half the iterations of
Jacobi because it incorporates fresh information sooner. The trade-off:
it cannot be parallelised as straightforwardly.

\subsection{Convergence criterion: diagonal
dominance}\label{convergence-criterion-diagonal-dominance}

Both methods are guaranteed to converge if \(A\) is \textbf{strictly
diagonally dominant}:

\[|a_{ii}| > \sum_{j \neq i} |a_{ij}| \quad \text{for all } i\]

This means each diagonal entry exceeds the sum of all off-diagonal
entries in its row, in absolute value. FEM stiffness matrices and many
physical systems (heat diffusion, network flow) naturally satisfy this
condition.

When diagonal dominance holds, the error is guaranteed to shrink at each
step and the iteration converges from any starting point. (For those
familiar with eigenvalues: the convergence criterion is that the
spectral radius of the iteration matrix --- the largest absolute
eigenvalue of the matrix that multiplies the error --- is less than 1.)

\subsection{Worked example: Jacobi and Gauss-Seidel on a 3×3
system}\label{worked-example-jacobi-and-gauss-seidel-on-a-33-system}

Solve from \(\mathbf{x}^{(0)} = (0, 0, 0)^T\):

\[\begin{pmatrix} 10 & -1 & 2 \\ -1 & 11 & -1 \\ 2 & -1 & 10 \end{pmatrix} \begin{pmatrix} x_1 \\ x_2 \\ x_3 \end{pmatrix} = \begin{pmatrix} 6 \\ 25 \\ -11 \end{pmatrix}\]

Check diagonal dominance: Row 1: \(10 > |-1| + |2| = 3\) ✓; Row 2:
\(11 > |-1| + |-1| = 2\) ✓; Row 3: \(10 > |2| + |-1| = 3\) ✓.
Convergence guaranteed.

The exact solution is \(\mathbf{x} = (1, 2, -1)^T\).

\subsubsection{Jacobi, 3 iterations}\label{jacobi-3-iterations}

Iteration formulas:

\[x_1^{(k+1)} = \frac{1}{10}(6 + x_2^{(k)} - 2x_3^{(k)})\]

\[x_2^{(k+1)} = \frac{1}{11}(25 + x_1^{(k)} + x_3^{(k)})\]

\[x_3^{(k+1)} = \frac{1}{10}(-11 - 2x_1^{(k)} + x_2^{(k)})\]

\textbf{\(k=0 \to k=1\):} Using \(\mathbf{x}^{(0)} = (0, 0, 0)^T\):

\(x_1^{(1)} = (6 + 0 - 0)/10 = 0.6\)

\(x_2^{(1)} = (25 + 0 + 0)/11 \approx 2.2727\)

\(x_3^{(1)} = (-11 - 0 + 0)/10 = -1.1\)

\textbf{\(k=1 \to k=2\):}

\(x_1^{(2)} = (6 + 2.2727 - 2(-1.1))/10 = (6 + 2.2727 + 2.2)/10 = 10.4727/10 = 1.0473\)

\(x_2^{(2)} = (25 + 0.6 + (-1.1))/11 = 24.5/11 = 2.2273\)

\(x_3^{(2)} = (-11 - 2(0.6) + 2.2727)/10 = (-11 - 1.2 + 2.2727)/10 = -9.9273/10 = -0.9927\)

\textbf{\(k=2 \to k=3\):}

\(x_1^{(3)} = (6 + 2.2273 - 2(-0.9927))/10 = (6 + 2.2273 + 1.9854)/10 = 10.2127/10 = 1.0213\)

\(x_2^{(3)} = (25 + 1.0473 + (-0.9927))/11 = (25 + 0.0546)/11 = 25.0546/11 \approx 2.278\)

\(x_3^{(3)} = (-11 - 2(1.0473) + 2.2273)/10 = (-11 - 2.0946 + 2.2273)/10 = -10.8673/10 = -1.087\)

After 3 iterations:
\(\mathbf{x}^{(3)} \approx (1.021, 2.278, -1.087)^T\). Errors are
shrinking --- heading toward the exact solution \((1, 2, -1)^T\), but
Jacobi converges slowly on this system.

\subsubsection{Gauss-Seidel, 3
iterations}\label{gauss-seidel-3-iterations}

\textbf{\(k=0 \to k=1\):} Using \textbf{updated} values as soon as
available:

\(x_1^{(1)} = (6 + 0 - 0)/10 = 0.6\) (uses \(x_2^{(0)} = 0\),
\(x_3^{(0)} = 0\))

\(x_2^{(1)} = (25 + \mathbf{0.6} + 0)/11 = 25.6/11 \approx 2.3273\)
(uses new \(x_1^{(1)} = 0.6\))

\(x_3^{(1)} = (-11 - 2(\mathbf{0.6}) + \mathbf{2.3273})/10 = (-11 - 1.2 + 2.3273)/10 = -9.8727/10 = -0.9873\)
(uses new \(x_1^{(1)}\) and \(x_2^{(1)}\))

\textbf{\(k=1 \to k=2\):}

\(x_1^{(2)} = (6 + 2.3273 - 2(-0.9873))/10 = (6 + 2.3273 + 1.9746)/10 = 10.3019/10 = 1.0302\)

\(x_2^{(2)} = (25 + 1.0302 + (-0.9873))/11 = 25.0429/11 \approx 2.2766\)

\(x_3^{(2)} = (-11 - 2(1.0302) + 2.2766)/10 = (-11 - 2.0604 + 2.2766)/10 = -10.7838/10 = -1.0784\)

\textbf{\(k=2 \to k=3\):}

\(x_1^{(3)} = (6 + 2.2766 - 2(-1.0784))/10 = (6 + 2.2766 + 2.1568)/10 = 10.4334/10 = 1.0433\)

\(x_2^{(3)} = (25 + 1.0433 + (-1.0784))/11 = 24.9649/11 \approx 2.2695\)

\(x_3^{(3)} = (-11 - 2(1.0433) + 2.2695)/10 = (-11 - 2.0866 + 2.2695)/10 = -10.8171/10 = -1.0817\)

After 3 iterations, Gauss-Seidel is at
\(\mathbf{x}^{(3)} \approx (1.043, 2.270, -1.082)^T\). Errors are still
about \(4\%\), similar to Jacobi after 3 steps on this particular
system. Gauss-Seidel does converge in roughly half the iterations of
Jacobi in general, but the advantage shows more clearly in systems where
the off-diagonal coupling is stronger or after more iterations have been
taken.

\begin{tcolorbox}[enhanced jigsaw, opacityback=0, colback=white, breakable, colframe=quarto-callout-note-color-frame, arc=.35mm, left=2mm, toprule=.15mm, leftrule=.75mm, bottomrule=.15mm, rightrule=.15mm]
\begin{minipage}[t]{5.5mm}
\textcolor{quarto-callout-note-color}{\faInfo}
\end{minipage}%
\begin{minipage}[t]{\textwidth - 5.5mm}

\textbf{When to use iterative vs direct}

\begin{longtable}[]{@{}
  >{\raggedright\arraybackslash}p{(\linewidth - 4\tabcolsep) * \real{0.3333}}
  >{\raggedright\arraybackslash}p{(\linewidth - 4\tabcolsep) * \real{0.3333}}
  >{\raggedright\arraybackslash}p{(\linewidth - 4\tabcolsep) * \real{0.3333}}@{}}
\toprule\noalign{}
\begin{minipage}[b]{\linewidth}\raggedright
Criterion
\end{minipage} & \begin{minipage}[b]{\linewidth}\raggedright
Direct (LU)
\end{minipage} & \begin{minipage}[b]{\linewidth}\raggedright
Iterative (GS/Jacobi)
\end{minipage} \\
\midrule\noalign{}
\endhead
\bottomrule\noalign{}
\endlastfoot
Matrix size & Any, up to memory & Large sparse matrices \\
Sparsity & Not exploited & Natural fit \\
Multiple RHS & Fast after factorisation & Re-iterate each time \\
Accuracy & Machine precision & Stops at tolerance \\
Convergence guarantee & Always (if non-singular) & Diagonal dominance \\
\end{longtable}

For FEM systems with \(n > 10^5\), iterative methods are standard. For
\(n < 10^3\) with multiple right-hand sides, LU with pivoting is
preferable.

\end{minipage}%
\end{tcolorbox}

\section{Where this goes}\label{where-this-goes-31}

This chapter handles \(A\mathbf{x} = \mathbf{b}\) for fixed \(A\) and
\(\mathbf{b}\). The next chapter in this section extends the same ideas
to \textbf{numerical ODEs and PDEs}: discretising a differential
equation produces a linear system at every time step or on every mesh
--- the same pivoted LU or iterative machinery applies, just called
millions of times. The structure of the system (banded, sparse,
symmetric) determines which variant to use, and the condition number of
the discretised operator determines the mesh resolution needed to stay
in the accurate regime.

The conditioning ideas carry directly into \textbf{optimisation} (Vol 7,
optimisation section): the Hessian matrix of an objective function plays
the role of \(A\), and ill-conditioning of the Hessian is the primary
reason gradient-descent methods stall or diverge. Knowing how to
estimate \(\kappa\) and how to precondition a matrix (rescaling so that
\(\kappa\) drops) is practical currency in large-scale numerical
optimisation.

\begin{tcolorbox}[enhanced jigsaw, opacityback=0, colback=white, breakable, colframe=quarto-callout-tip-color-frame, arc=.35mm, left=2mm, toprule=.15mm, leftrule=.75mm, bottomrule=.15mm, rightrule=.15mm]
\begin{minipage}[t]{5.5mm}
\textcolor{quarto-callout-tip-color}{\faLightbulb}
\end{minipage}%
\begin{minipage}[t]{\textwidth - 5.5mm}

\textbf{Where this shows up}

\begin{itemize}
\tightlist
\item
  \textbf{Structural engineering:} Every finite element analysis
  assembles a stiffness matrix \(K\) and solves
  \(K\mathbf{u} = \mathbf{f}\) for displacements \(\mathbf{u}\). LU
  factorisation with multiple load vectors is standard.
\item
  \textbf{Computational fluid dynamics:} The pressure-correction step in
  incompressible flow solvers is a large sparse Poisson system ---
  iterative solvers (conjugate gradient, a Krylov subspace method that
  exploits symmetry and positive-definiteness) are universal.
\item
  \textbf{Machine learning:} Linear regression solves the normal
  equations \((X^TX)\boldsymbol{\beta} = X^T\mathbf{y}\); the condition
  number of \(X^TX\) is why multicollinearity ruins regression
  estimates.
\item
  \textbf{Finance:} Estimating a covariance matrix from historical
  returns produces an ill-conditioned system when assets are correlated
  --- the condition number flags when the inverse is unreliable.
\end{itemize}

\end{minipage}%
\end{tcolorbox}

\section{Exercises}\label{exercises-42}

These are engineering-layer puzzles. Each has a specific numerical
answer. Work the full procedure --- don't skip steps --- because the
errors that matter in practice happen in the steps that get skipped in
rough work.

\textbf{Exercise 1.} Factorise the following matrix using Doolittle LU
decomposition, then use the factorisation to solve for two right-hand
sides.

\[A = \begin{pmatrix} 2 & 4 & 2 \\ 1 & 5 & 2 \\ 4 & 6 & 8 \end{pmatrix}, \quad \mathbf{b}_1 = \begin{pmatrix} 2 \\ 3 \\ 6 \end{pmatrix}, \quad \mathbf{b}_2 = \begin{pmatrix} 0 \\ 1 \\ -2 \end{pmatrix}\]

\textbf{Exercise 2.} Apply partial pivoting to the following
near-singular system. Identify which row swap is required, perform the
swap, and carry out elimination. Show what would go wrong without the
swap.

\[\begin{pmatrix} 0.002 & 1 & 1 \\ 1 & 2 & 5 \\ 1 & -1 & 3 \end{pmatrix} \begin{pmatrix} x_1 \\ x_2 \\ x_3 \end{pmatrix} = \begin{pmatrix} 1 \\ 8 \\ 5 \end{pmatrix}\]

\textbf{Exercise 3.} Compute \(\kappa_\infty(A)\) for the matrix below,
then interpret the result: if you solve \(A\mathbf{x} = \mathbf{b}\) on
a machine with 15 significant digits, how many reliable digits do you
expect in the solution?

\[A = \begin{pmatrix} 3 & 1 \\ 2 & 4 \end{pmatrix}\]

\textbf{Exercise 4.} Apply 3 iterations of Jacobi iteration to the
following system, starting from \(\mathbf{x}^{(0)} = (0, 0, 0)^T\). Show
the iterate after each step.

\[\begin{pmatrix} 5 & 1 & 1 \\ 1 & 6 & 1 \\ 1 & 1 & 4 \end{pmatrix} \begin{pmatrix} x_1 \\ x_2 \\ x_3 \end{pmatrix} = \begin{pmatrix} 7 \\ 8 \\ 6 \end{pmatrix}\]

The exact solution is \(\mathbf{x} = (1, 1, 1)^T\).

\textbf{Exercise 5.} Apply 3 iterations of Gauss-Seidel to the same
system as Exercise 4, from the same starting point
\(\mathbf{x}^{(0)} = (0, 0, 0)^T\). Compare the convergence rate after 3
iterations.

\textbf{Exercise 6 (Engineering context).} A four-node plane truss is
modelled by finite element analysis. After assembly and boundary
conditions, the reduced stiffness system is:

\[K = \begin{pmatrix} 4 & -1 & 0 & 0 \\ -1 & 4 & -1 & 0 \\ 0 & -1 & 4 & -1 \\ 0 & 0 & -1 & 3 \end{pmatrix}, \quad \mathbf{f} = \begin{pmatrix} 0 \\ 0 \\ 10 \\ 5 \end{pmatrix}\]

\begin{enumerate}
\def\labelenumi{(\alph{enumi})}
\item
  Verify that \(K\) is strictly diagonally dominant, which guarantees
  that Gauss-Seidel converges for any load vector \(\mathbf{f}\).
\item
  Identify the structural feature of \(K\) (bandwidth, sparsity pattern)
  that explains why iterative methods are preferred over LU
  decomposition for large truss problems of this type.
\item
  Apply 2 iterations of Gauss-Seidel from
  \(\mathbf{u}^{(0)} = (0,0,0,0)^T\).
\end{enumerate}

\begin{itemize}
\item
  op: ``Identify structural feature: K is tridiagonal (bandwidth 1)''
  eq: ``\textbackslash text\{Only entries
  \}k\_\{i,i\},\textbackslash;k\_\{i,i\textbackslash pm
  1\}\textbackslash text\{ are non-zero --- all others are zero\}''
  note: ``For large n, LU decomposition of a tridiagonal matrix remains
  tridiagonal (no fill-in), but for general sparse FEM matrices, LU can
  produce dense factors. Iterative methods work directly on the sparse
  structure without fill-in.''
\item
  op: ``Write Gauss-Seidel formulas for each DOF'' eq:
  ``u\_1\textsuperscript{\{(k+1)\}=\textbackslash frac\{1\}\{4\}(0+u\_2}\{(k)\}),\textbackslash quad
  u\_2\textsuperscript{\{(k+1)\}=\textbackslash frac\{1\}\{4\}(0+u\_1}\{(k+1)\}+u\_3\^{}\{(k)\}),\textbackslash quad
  u\_3\textsuperscript{\{(k+1)\}=\textbackslash frac\{1\}\{4\}(10+u\_2}\{(k+1)\}+u\_4\^{}\{(k)\}),\textbackslash quad
  u\_4\textsuperscript{\{(k+1)\}=\textbackslash frac\{1\}\{3\}(5+u\_3}\{(k+1)\})''
  note: \textasciitilde{}
\item
  op: ``Iteration 1 with u⁽⁰⁾=(0,0,0,0)\^{}T'' eq:
  ``u\_1\^{}\{(1)\}=0/4=0,\textbackslash quad
  u\_2\^{}\{(1)\}=(0+0)/4=0,\textbackslash quad
  u\_3\^{}\{(1)\}=(10+0+0)/4=2.5,\textbackslash quad
  u\_4\^{}\{(1)\}=(5+2.5)/3=2.5'' note: \textasciitilde{}
\item
  op: ``Iteration 2 with u⁽¹⁾=(0,0,2.5,2.5)\^{}T'' eq:
  ``u\_1\^{}\{(2)\}=0/4=0,\textbackslash quad
  u\_2\^{}\{(2)\}=(0+0+2.5)/4=0.625,\textbackslash quad
  u\_3\^{}\{(2)\}=(10+0.625+2.5)/4=3.281,\textbackslash quad
  u\_4\^{}\{(2)\}=(5+3.281)/3=2.760'' note: \textasciitilde{}
\item
  op: ``Physical interpretation of the iterates'' eq:
  ``\textbackslash text\{Node 4 displacement grows from 0 to 2.5 to
  2.760 (in whatever units the stiffness and force are expressed) ---
  iteration revealing the load path through the truss\}'' note: ``The
  iterative process is physically meaningful: each sweep propagates the
  applied loads one structural connection further through the mesh. No
  units are assigned in this dimensionless example; in a real FEM model
  the stiffness entries carry units of kN/m and forces carry kN, giving
  displacements in metres.''
\item
  op: ``Check convergence direction'' eq:
  ``\textbackslash text\{Residual
  \}\textbackslash\textbar K\textbackslash mathbf\{u\}\^{}\{(2)\}-\textbackslash mathbf\{f\}\textbackslash\textbar{}\emph{\textbackslash infty:\textbackslash quad
  K\textbackslash mathbf\{u\}\textsuperscript{\{(2)\}\textbackslash approx(-0.625,\textbackslash;-0.781,\textbackslash;9.739,\textbackslash;4.999)}T,\textbackslash quad\textbackslash text\{vs
  \}\textbackslash mathbf\{f\}=(0,0,10,5)\^{}T'' note: ''The residual is
  shrinking; continue until \textbar\textbar Ku − f\textbar\textbar{}}∞
  \textless{} tolerance (e.g.~10⁻⁶).'' --\textgreater{}
\end{itemize}

\chapter{Numerical ODEs and PDEs}\label{numerical-odes-and-pdes}

Stepping through time and space

\hfill\break

\section{The setup: replacing dy/dx with
Δy/Δx}\label{the-setup-replacing-dydx-with-ux3b4yux3b4x}

The initial-value problem has the form

\[
y' = f(x, y), \quad y(x_0) = y_0.
\]

Read this as: ``the slope of \(y\) at any point \((x, y)\) is given by
the function \(f\), and we know where we start.'' The question is: what
is \(y\) at some later \(x\)?

If \(f\) is simple enough, calculus gives you \(y(x)\) directly. When it
isn't, you fall back on the most basic approximation: the tangent-line
step.

The \textbf{step size} \(h\) --- a small positive number --- is your
control parameter. Smaller \(h\) means more accuracy and more work. The
methods below differ in how many function evaluations they use per step,
and therefore in how quickly the error shrinks as \(h\) decreases.

\begin{center}\rule{0.5\linewidth}{0.5pt}\end{center}

\section{Euler's method}\label{eulers-method}

\subsection{Forward Euler}\label{forward-euler}

Approximate \(y'(x_n) \approx \dfrac{y_{n+1} - y_n}{h}\), rearrange, and
you get

\[
y_{n+1} = y_n + h \, f(x_n, y_n).
\]

\(y_n\) means ``the approximate value of \(y\) at step \(n\).'' The
update rule says: take the current value, add the step size times the
slope at the current point.

\textbf{Local truncation error.} Taylor-expand \(y(x_{n+1})\) about
\(x_n\):

\[
y(x_{n+1}) = y(x_n) + h y'(x_n) + \tfrac{h^2}{2} y''(x_n) + \cdots
\]

Forward Euler keeps only the first-derivative term, so the error in one
step is \(O(h^2)\). Over \(N = (b - a)/h\) steps, the errors accumulate:
\textbf{global error is \(O(h)\)}. Halve the step size, halve the global
error.

\subsection{Backward (implicit) Euler}\label{backward-implicit-euler}

\[
y_{n+1} = y_n + h \, f(x_{n+1}, y_{n+1}).
\]

The slope is evaluated at the \emph{new} point, which makes this an
equation you must solve for \(y_{n+1}\) at each step --- more expensive
per step. The payoff is unconditional stability for stiff equations
(equations where the true solution decays quickly and explicit methods
require an extremely small \(h\) to stay bounded). For non-stiff
problems, forward Euler is fine.

\begin{tcolorbox}[enhanced jigsaw, title=\textcolor{quarto-callout-note-color}{\faInfo}\hspace{0.5em}{Note}, opacityback=0, colback=white, breakable, leftrule=.75mm, arc=.35mm, toptitle=1mm, left=2mm, colframe=quarto-callout-note-color-frame, coltitle=black, titlerule=0mm, colbacktitle=quarto-callout-note-color!10!white, opacitybacktitle=0.6, toprule=.15mm, bottomtitle=1mm, bottomrule=.15mm, rightrule=.15mm]

\textbf{What ``stiff'' means.} A stiff ODE has solution components that
decay at very different rates. The fast-decaying component forces \(h\)
to be tiny for forward Euler to remain stable, even though that
component itself becomes negligible quickly. Implicit methods remain
stable at larger \(h\) because they incorporate information from the
next step.

\end{tcolorbox}

\subsection{\texorpdfstring{Worked example: forward Euler on
\(y' = -2y\)}{Worked example: forward Euler on y\textquotesingle{} = -2y}}\label{worked-example-forward-euler-on-y--2y}

Exact solution: \(y = e^{-2x}\). Let \(h = 0.5\), \(y_0 = 1\).

\begin{longtable}[]{@{}lllll@{}}
\toprule\noalign{}
\(n\) & \(x_n\) & \(y_n\) & \(f = -2y_n\) & \(y_{n+1} = y_n + h f\) \\
\midrule\noalign{}
\endhead
\bottomrule\noalign{}
\endlastfoot
0 & 0.0 & 1.0000 & −2.0000 & 1 + 0.5(−2) = 0.0000 \\
\end{longtable}

That's the stability boundary. For an equation of the form
\(y' = \lambda y\) --- where \(\lambda\) is the coefficient multiplying
\(y\), here \(\lambda = -2\) --- forward Euler multiplies the solution
by \((1 + h\lambda)\) at each step. Stability requires
\(|1 + h\lambda| \leq 1\), which for \(\lambda = -2\) gives
\(h \leq 1\). With \(h = 0.5\) we land exactly at the edge: the solution
collapses to zero and stays there. For \(h > 1\) the factor
\((1 + h\lambda)\) exceeds 1 in magnitude, and the computed solution
oscillates with growing amplitude --- blowing up --- even though the
true solution decays smoothly to zero. Use \(h = 0.2\) to see the method
tracking the decay:

\begin{longtable}[]{@{}llllll@{}}
\toprule\noalign{}
\(n\) & \(x_n\) & \(y_n\) & \(f = -2y_n\) & \(y_{n+1}\) & Exact
\(e^{-2x}\) \\
\midrule\noalign{}
\endhead
\bottomrule\noalign{}
\endlastfoot
0 & 0.0 & 1.0000 & −2.0000 & 0.6000 & --- \\
1 & 0.2 & 0.6000 & −1.2000 & 0.3600 & 0.6703 \\
2 & 0.4 & 0.3600 & −0.7200 & 0.2160 & 0.4493 \\
\end{longtable}

The error at \(x = 0.4\) is about \(0.089\), or roughly 20\%. Three
steps with \(h = 0.2\) on a decaying exponential is not a recipe for
accuracy --- but it illustrates the mechanics.

\begin{center}\rule{0.5\linewidth}{0.5pt}\end{center}

\section{Runge-Kutta methods}\label{runge-kutta-methods}

The key idea: use several evaluations of \(f\) within a single step ---
at different intermediate points --- then combine them into a weighted
average slope. More work per step, but the error drops much faster with
decreasing \(h\).

\subsection{RK2: two-slope methods}\label{rk2-two-slope-methods}

\textbf{Midpoint method.} Take a half-step to estimate the midpoint,
then use the slope there:

\[
k_1 = f(x_n, y_n), \qquad k_2 = f\!\left(x_n + \tfrac{h}{2},\; y_n + \tfrac{h}{2} k_1\right)
\]

\[
y_{n+1} = y_n + h k_2.
\]

\textbf{Heun's method (modified Euler).} Use the slope at the beginning
and the predicted end, average them:

\[
k_1 = f(x_n, y_n), \qquad k_2 = f(x_n + h,\; y_n + h k_1)
\]

\[
y_{n+1} = y_n + \tfrac{h}{2}(k_1 + k_2).
\]

Both are \(O(h^2)\) per step, giving global error \(O(h^2)\). Halving
\(h\) divides the error by 4.

\subsection{RK4: the standard
workhorse}\label{rk4-the-standard-workhorse}

Evaluate four slopes --- \(k_1\), \(k_2\), \(k_3\), \(k_4\) --- and
combine with weights \(1, 2, 2, 1\):

\[
k_1 = f(x_n,\; y_n)
\]

\[
k_2 = f\!\left(x_n + \tfrac{h}{2},\; y_n + \tfrac{h}{2} k_1\right)
\]

\[
k_3 = f\!\left(x_n + \tfrac{h}{2},\; y_n + \tfrac{h}{2} k_2\right)
\]

\[
k_4 = f(x_n + h,\; y_n + h k_3)
\]

\[
y_{n+1} = y_n + \tfrac{h}{6}(k_1 + 2k_2 + 2k_3 + k_4).
\]

\(k_1\) is the slope at the left end. \(k_2\) and \(k_3\) are two
estimates of the midpoint slope (the second corrects the first). \(k_4\)
is the slope at the right end. The weights \(1, 2, 2, 1\) come from
Simpson's rule: the interval has three evaluation points --- left end,
midpoint, right end --- with Simpson weights \(1, 4, 1\). RK4 has two
independent midpoint estimates rather than one, so their weight is
split, giving \(1, 2, 2, 1\) after normalising to sum to \(6\).

\textbf{Error.} Local error is \(O(h^5)\); global error is \(O(h^4)\).
Halving \(h\) divides the global error by 16. This is why RK4 is the
default choice for non-stiff problems: it is usually accurate enough
with step sizes that are not punishingly small.

\begin{tcolorbox}[enhanced jigsaw, title=\textcolor{quarto-callout-note-color}{\faInfo}\hspace{0.5em}{Note}, opacityback=0, colback=white, breakable, leftrule=.75mm, arc=.35mm, toptitle=1mm, left=2mm, colframe=quarto-callout-note-color-frame, coltitle=black, titlerule=0mm, colbacktitle=quarto-callout-note-color!10!white, opacitybacktitle=0.6, toprule=.15mm, bottomtitle=1mm, bottomrule=.15mm, rightrule=.15mm]

\textbf{Adaptive step size.} RK4 with a fixed \(h\) wastes work in
regions where the solution is smooth and underestimates error in regions
where it changes rapidly. Adaptive methods (e.g.~RK45, the
Dormand-Prince method) compute two approximations of different orders at
each step, compare them to estimate the local error, and adjust \(h\) to
keep the error within a specified tolerance. Embedded pairs like RK23
and RK45 reuse slope evaluations between the two orders, keeping the
extra cost small. The implementation is routine; the idea is simply: use
\(h\) proportional to what the solution needs.

\end{tcolorbox}

\subsection{\texorpdfstring{Worked example: RK4 on \(y' = y\), one
step}{Worked example: RK4 on y\textquotesingle{} = y, one step}}\label{worked-example-rk4-on-y-y-one-step}

\(y(0) = 1\), \(h = 1\). Exact: \(y(1) = e \approx 2.71828\).

\[
k_1 = f(0, 1) = 1
\] \[
k_2 = f(0.5,\; 1 + 0.5) = f(0.5, 1.5) = 1.5
\] \[
k_3 = f(0.5,\; 1 + 0.5 \times 1.5) = f(0.5, 1.75) = 1.75
\] \[
k_4 = f(1,\; 1 + 1 \times 1.75) = f(1, 2.75) = 2.75
\] \[
y_1 = 1 + \tfrac{1}{6}(1 + 3 + 3.5 + 2.75) = 1 + \tfrac{10.25}{6} \approx 2.7083.
\]

Error: \(|2.7083 - 2.71828| \approx 0.00998\). With a single step of
\(h = 1\), RK4 gets the answer to 3 significant figures. Forward Euler
with \(h = 1\) gives \(y_1 = 2.0\) --- an error of 26\%.

\begin{center}\rule{0.5\linewidth}{0.5pt}\end{center}

\section{Systems of ODEs}\label{systems-of-odes}

Many physical problems involve second- (or higher-) order ODEs. A
second-order ODE can always be rewritten as a system of two first-order
ODEs, and then Euler or RK4 applies to the system directly.

\textbf{Reduction procedure.} Given

\[
y'' + p(x)y' + q(x)y = r(x),
\]

introduce \(u_1 = y\) and \(u_2 = y'\). Then:

\[
u_1' = u_2
\]

\[
u_2' = r(x) - p(x) u_2 - q(x) u_1.
\]

Write this as \(\mathbf{u}' = \mathbf{F}(x, \mathbf{u})\) where
\(\mathbf{u} = [u_1, u_2]^T\). Apply Euler or RK4 to the vector
equation: each \(k_i\) becomes a two-component vector.

\subsection{Worked example: damped
oscillator}\label{worked-example-damped-oscillator}

\(y'' + 0.5y' + 4y = 0\), \(y(0) = 1\), \(y'(0) = 0\), \(h = 0.1\).

Let \(u_1 = y\), \(u_2 = y'\). Then:

\[
u_1' = u_2, \qquad u_2' = -4 u_1 - 0.5 u_2.
\]

One Euler step:

\[
k_1^{(1)} = u_2^{(0)} = 0, \qquad k_1^{(2)} = -4(1) - 0.5(0) = -4.
\]

\[
u_1^{(1)} = 1 + 0.1 \times 0 = 1.0, \qquad u_2^{(1)} = 0 + 0.1 \times (-4) = -0.4.
\]

So after one step: \(y(0.1) \approx 1.0\), \(y'(0.1) \approx -0.4\). The
method is mechanical once the system is written down; the reduction step
is the only place judgment is required.

\begin{center}\rule{0.5\linewidth}{0.5pt}\end{center}

\section{Finite differences for PDEs}\label{finite-differences-for-pdes}

The same idea --- replace a derivative with a finite ratio --- extends
to PDEs, where you discretise both time and space.

\subsection{The 1D heat equation}\label{the-1d-heat-equation}

\[
\frac{\partial u}{\partial t} = \alpha \frac{\partial^2 u}{\partial x^2}, \quad 0 \le x \le L, \quad t \ge 0,
\]

with given initial condition \(u(x, 0)\) and boundary conditions.
\(\alpha\) is thermal diffusivity (a positive constant).

Lay a grid: \(x_j = j \, \Delta x\) for \(j = 0, 1, \ldots, M\) and
\(t^n = n \, \Delta t\) for \(n = 0, 1, \ldots\). Write \(u_j^n\) for
the approximation to \(u(x_j, t^n)\).

\textbf{FTCS scheme (forward-time, central-space).} Use a forward
difference in time and a central difference in space:

\[
\frac{u_j^{n+1} - u_j^n}{\Delta t} = \alpha \frac{u_{j+1}^n - 2u_j^n + u_{j-1}^n}{(\Delta x)^2}.
\]

Define the \textbf{Fourier number} (also called the mesh ratio)

\[
r = \frac{\alpha \, \Delta t}{(\Delta x)^2}.
\]

Rearrange:

\[
u_j^{n+1} = r \, u_{j-1}^n + (1 - 2r) \, u_j^n + r \, u_{j+1}^n.
\]

This is explicit: each new value depends only on known values at the
previous time level.

\textbf{Stability condition.} Look at the FTCS update coefficients:
\(r\), \((1-2r)\), \(r\). When \(r \leq 1/2\), the coefficient
\((1-2r)\) is non-negative, so the new value at each interior node is a
weighted average of the neighbouring values --- physically reasonable,
bounded, and stable. When \(r > 1/2\), the coefficient \((1-2r)\) turns
negative. The centre node's own value now enters with a negative weight:
the update is no longer an average. A small perturbation at one node
gets amplified rather than smoothed; each time step multiplies the
perturbation, and the solution oscillates with growing amplitude until
it overflows. The physical heat equation diffuses and decays; the
numerical scheme with an oversized time step does the opposite --- it
invents energy out of rounding errors. Von Neumann stability analysis
makes this precise: \(r \le \tfrac{1}{2}\) is the condition. In
practice, for \(\alpha = 1\), \(\Delta x = 0.1\): the maximum stable
\(\Delta t\) is \(\tfrac{1}{2} (0.1)^2 = 0.005\).

\subsection{Crank-Nicolson: unconditional
stability}\label{crank-nicolson-unconditional-stability}

Average the spatial derivative between time levels \(n\) and \(n+1\):

\[
\frac{u_j^{n+1} - u_j^n}{\Delta t} = \frac{\alpha}{2} \left[ \frac{u_{j+1}^n - 2u_j^n + u_{j-1}^n}{(\Delta x)^2} + \frac{u_{j+1}^{n+1} - 2u_j^{n+1} + u_{j-1}^{n+1}}{(\Delta x)^2} \right].
\]

The unknowns at time level \(n+1\) appear on both sides. Collecting them
gives a tridiagonal linear system to solve at each time step ---
implicit, but only marginally more expensive than FTCS because
tridiagonal systems have a special structure: each equation involves
only its immediate neighbours, so a forward sweep followed by
back-substitution solves the whole system in \(O(N)\) operations. This
procedure is called the \textbf{Thomas algorithm} --- Gaussian
elimination stripped down to exploit the tridiagonal sparsity pattern.
The method is unconditionally stable: any \(r\) is admissible, so you
can take large time steps without blowup. The global error is
\(O((\Delta t)^2 + (\Delta x)^2)\) --- second order in both.

\subsection{Boundary-value problems: tridiagonal systems from central
differences}\label{boundary-value-problems-tridiagonal-systems-from-central-differences}

For a 1D BVP:

\[
y'' = f(x), \quad y(a) = \alpha, \quad y(b) = \beta,
\]

discretise the interior points \(x_1, \ldots, x_{N-1}\) with spacing
\(h = (b-a)/N\). The central-difference approximation
\(y'' \approx (y_{i-1} - 2y_i + y_{i+1})/h^2\) at each interior point
gives:

\[
y_{i-1} - 2y_i + y_{i+1} = h^2 f(x_i), \quad i = 1, \ldots, N-1.
\]

Collect these equations (incorporating \(y_0 = \alpha\) and
\(y_N = \beta\) into the first and last equations) and you get a
tridiagonal linear system. The structure is

\[
\begin{pmatrix}
-2 & 1 & & \\
1 & -2 & 1 & \\
 & \ddots & \ddots & 1 \\
 & & 1 & -2
\end{pmatrix}
\begin{pmatrix} y_1 \\ y_2 \\ \vdots \\ y_{N-1} \end{pmatrix}
= h^2 \begin{pmatrix} f(x_1) - \alpha/h^2 \\ f(x_2) \\ \vdots \\ f(x_{N-1}) - \beta/h^2 \end{pmatrix}.
\]

Solve with the Thomas algorithm (forward sweep + back substitution) in
\(O(N)\) operations.

\begin{tcolorbox}[enhanced jigsaw, title=\textcolor{quarto-callout-note-color}{\faInfo}\hspace{0.5em}{Note}, opacityback=0, colback=white, breakable, leftrule=.75mm, arc=.35mm, toptitle=1mm, left=2mm, colframe=quarto-callout-note-color-frame, coltitle=black, titlerule=0mm, colbacktitle=quarto-callout-note-color!10!white, opacitybacktitle=0.6, toprule=.15mm, bottomtitle=1mm, bottomrule=.15mm, rightrule=.15mm]

\textbf{Why the tridiagonal structure matters.} Gaussian elimination on
a dense \(N \times N\) system costs \(O(N^3)\). Exploiting the
tridiagonal structure reduces this to \(O(N)\). For \(N = 1000\) spatial
nodes, that is a factor of \(10^9\) in computational work. Recognising
sparsity is not a nicety; it is what makes large-scale simulation
possible.

\end{tcolorbox}

\begin{center}\rule{0.5\linewidth}{0.5pt}\end{center}

\section{Worked example: FTCS one
step}\label{worked-example-ftcs-one-step}

Heat equation on \([0, 1]\) with 5 spatial nodes (\(\Delta x = 0.25\)),
\(\alpha = 1\), \(\Delta t = 0.02\).

First check stability:
\(r = 1 \times 0.02 / (0.25)^2 = 0.02 / 0.0625 = 0.32 < 0.5\). Stable.

Initial condition: \(u(x, 0) = \sin(\pi x)\). Interior nodes:
\(x_1 = 0.25\), \(x_2 = 0.50\), \(x_3 = 0.75\).

Values at \(t = 0\): \(u_0^0 = 0\),
\(u_1^0 = \sin(0.25\pi) \approx 0.7071\),
\(u_2^0 = \sin(0.5\pi) = 1.0\),
\(u_3^0 = \sin(0.75\pi) \approx 0.7071\), \(u_4^0 = 0\).

Apply \(u_j^1 = 0.32 u_{j-1}^0 + 0.36 u_j^0 + 0.32 u_{j+1}^0\)
(coefficients: \(r = 0.32\), \(1 - 2r = 0.36\)):

\[
u_1^1 = 0.32(0) + 0.36(0.7071) + 0.32(1.0) = 0.2546 + 0.3200 = 0.5746
\]

\[
u_2^1 = 0.32(0.7071) + 0.36(1.0) + 0.32(0.7071) = 0.2263 + 0.3600 + 0.2263 = 0.8126
\]

\[
u_3^1 = 0.5746 \quad \text{(by symmetry)}
\]

Exact solution: \(u(x, t) = e^{-\pi^2 t} \sin(\pi x)\). At \(t = 0.02\):
\(e^{-0.197} \approx 0.821\). So
\(u(0.5, 0.02)^{\text{exact}} = 0.821\). Numerical: \(0.8126\). Error
\(\approx 0.008\).

\begin{center}\rule{0.5\linewidth}{0.5pt}\end{center}

\section{Where this goes}\label{where-this-goes-32}

The methods here are the direct input to \textbf{adaptive ODE solvers}
(e.g.~\texttt{scipy.integrate.solve\_ivp}, MATLAB's \texttt{ode45}).
Those solvers use embedded Runge-Kutta pairs --- the same structure as
above --- with automatic step-size control. Understanding the underlying
stepping scheme lets you diagnose solver failures: a stiff ODE switched
to an explicit method, or a step-size tolerance too loose for the
problem's dynamics.

Finite differences for PDEs extend naturally to \textbf{finite element
methods} (FEM), where the domain is discretised into irregular elements
rather than a regular grid. FEM is the foundation of structural
analysis, fluid simulation, and electromagnetic modelling. The
tridiagonal system from 1D BVPs generalises to sparse banded systems in
higher dimensions --- still tractable, still solved by exploiting the
sparsity pattern.

\section{Applications}\label{applications-10}

\begin{tcolorbox}[enhanced jigsaw, title=\textcolor{quarto-callout-tip-color}{\faLightbulb}\hspace{0.5em}{Where numerical ODE/PDE methods appear}, opacityback=0, colback=white, breakable, leftrule=.75mm, arc=.35mm, toptitle=1mm, left=2mm, colframe=quarto-callout-tip-color-frame, coltitle=black, titlerule=0mm, colbacktitle=quarto-callout-tip-color!10!white, opacitybacktitle=0.6, toprule=.15mm, bottomtitle=1mm, bottomrule=.15mm, rightrule=.15mm]

\begin{itemize}
\tightlist
\item
  \textbf{Structural engineering.} Finite element codes (ANSYS, Abaqus)
  step through nonlinear load-displacement ODEs and solve elliptic PDEs
  for stress fields.
\item
  \textbf{Aerospace.} Flight simulators use RK4 (or RK45) to integrate
  the 6-DOF equations of motion. Autopilot \emph{design} uses
  frequency-domain tools (Laplace, Bode analysis), but the plant model
  is validated by time-domain simulation using exactly these stepping
  methods.
\item
  \textbf{Climate science.} Atmospheric models discretise the primitive
  equations on a spherical grid with time steps of minutes;
  finite-difference or spectral schemes advance the state.
\item
  \textbf{Biomedical.} Hodgkin-Huxley neuron models are stiff ODEs
  solved with implicit methods; cardiac electrophysiology models are
  stiff PDE systems.
\item
  \textbf{Computer graphics.} Game physics engines use symplectic Euler
  (a variant of forward Euler that conserves energy better) to integrate
  rigid-body dynamics at 60--120 Hz.
\item
  \textbf{Finance.} The Black-Scholes PDE for option pricing is
  discretised with Crank-Nicolson on a log-price grid.
\end{itemize}

\end{tcolorbox}

\begin{center}\rule{0.5\linewidth}{0.5pt}\end{center}

\section{Exercises}\label{exercises-43}

\textbf{Exercise 1.} A bacterial culture's population decreases due to
an antibiotic. The concentration \(y\) (mg/L) satisfies \(dy/dx = -2y\),
\(y(0) = 1\). Apply forward Euler with \(h = 0.1\) for three steps.
Compare with the exact values \(e^{-0.2}\), \(e^{-0.4}\), \(e^{-0.6}\).

\textbf{Exercise 2.} Apply RK4 to \(dy/dx = -2y\), \(y(0) = 1\) with a
single step of \(h = 0.3\). Compute all four slopes
\(k_1, k_2, k_3, k_4\) and the resulting \(y_1\). Compare with the exact
\(e^{-0.6}\).

\textbf{Exercise 3.} A second-order system satisfies
\(y'' - 3y' + 2y = 0\), \(y(0) = 0\), \(y'(0) = 1\), \(h = 0.1\).

\begin{enumerate}
\def\labelenumi{(\alph{enumi})}
\tightlist
\item
  Reduce to a first-order system.
\item
  Apply one step of forward Euler to get \(y(0.1)\) and \(y'(0.1)\).
\end{enumerate}

\textbf{Exercise 4.} A 1D diffusion problem is modelled with diffusivity
\(\alpha = 0.5\) (dimensionless or in consistent units scaled to the rod
length). The rod lies on \([0, 1]\) with 5 spatial nodes
(\(\Delta x = 0.25\)). Initial condition: \(u(x,0) = \sin(\pi x)\).
Boundary conditions: \(u(0,t) = u(1,t) = 0\). Take \(\Delta t = 0.01\).

Apply one FTCS time step. Give the value of \(u\) at each interior node.

\textbf{Exercise 5.} Using the parameters from Exercise 4, verify
analytically that the scheme is stable. Now choose \(\Delta t = 0.04\)
(with the same \(\Delta x = 0.25\), \(\alpha = 0.5\)). Compute \(r\) and
explain what will happen if you run the scheme forward in time.

\textbf{Exercise 6.} Solve the boundary-value problem
\(y'' = -\pi^2 \sin(\pi x)\), \(y(0) = 0\), \(y(1) = 0\) numerically
using central differences with 4 interior nodes (\(h = 0.2\)). Set up
the \(4 \times 4\) linear system and solve it. The exact solution is
\(y(x) = \sin(\pi x)\).

\part{Volume 7 --- Optimisation and probability}

\chapter{Linear programming}\label{linear-programming}

The best answer in a constrained space

\hfill\break

A factory can make two products --- call them chairs and tables. Each
chair uses 2 board-hours of cutting time and 1 hour of finishing. Each
table uses 1 board-hour of cutting and 3 hours of finishing. The factory
has 8 cutting hours and 9 finishing hours available per day. Chairs sell
at a \$30 margin; tables at \$50. How many of each should it make?

This is not a puzzle you can answer by intuition. Make only chairs and
you run out of cutting time. Make only tables and finishing is the
bottleneck. The mix that squeezes out the most profit sits somewhere in
the middle, at a specific corner of the region defined by those
constraints --- and that is exactly what linear programming finds.

\section{LP formulation}\label{lp-formulation}

A \textbf{linear programme} (LP) has three components: \emph{decision
variables}, an \emph{objective function}, and \emph{constraints}.

\textbf{Decision variables.} These are the quantities you control. Write
them as \(x_1, x_2, \ldots, x_n\) --- read as ``x-sub-one, x-sub-two, up
to x-sub-n'' --- or as the vector \(\mathbf{x}\), read as ``bold x.''
For the factory problem, \(x_1\) is the number of chairs made per day
and \(x_2\) is the number of tables.

\textbf{Objective function.} This is what you are optimising. It is
always a linear combination of the decision variables:
\(z = c_1 x_1 + c_2 x_2 +
\cdots + c_n x_n\), written compactly as
\(z = \mathbf{c}^\top \mathbf{x}\), read ``z equals c-transpose times
x.'' The coefficients \(\mathbf{c}\) are your profit margins (for
maximisation) or costs (for minimisation).

\textbf{Constraints.} Limits on what the variables can be. Each
constraint is a linear inequality:
\(a_{i1} x_1 + a_{i2} x_2 + \cdots \leq b_i\). Non-negativity
\(\mathbf{x} \geq \mathbf{0}\) --- read ``x is non-negative'' --- is
almost always included separately; you cannot make a negative number of
chairs.

\textbf{The factory problem in LP form:}

\[\text{maximise} \quad z = 30x_1 + 50x_2\]

subject to:

\[2x_1 + x_2 \leq 8 \qquad \text{(cutting hours)}\]
\[x_1 + 3x_2 \leq 9 \qquad \text{(finishing hours)}\]
\[x_1, x_2 \geq 0\]

\textbf{Standard form.} The simplex method requires equality
constraints. Convert each \(\leq\) constraint to an equation by adding a
\textbf{slack variable} \(s_i \geq 0\), read ``slack variable i.'' Slack
absorbs the difference between the left side and the right side:

\[2x_1 + x_2 + s_1 = 8, \quad s_1 \geq 0\]
\[x_1 + 3x_2 + s_2 = 9, \quad s_2 \geq 0\]

Now all constraints are equations, all variables are non-negative, and
the objective is linear. This is standard form:
\(\text{maximise } \mathbf{c}^\top\mathbf{x}\) subject to
\(A\mathbf{x} = \mathbf{b}\), \(\mathbf{x} \geq \mathbf{0}\).

\textbf{A second example: diet/resource allocation.} A hospital
nutrition service needs to provide at least 60 g of protein, at least
800 mg of calcium, and at most 2000 kcal per patient per day from two
food types. Food A costs \$1.20 per unit, food B costs \$0.80 per unit.
Formulate the LP to minimise daily cost.

Let \(x_1\) = units of food A, \(x_2\) = units of food B.

\begin{longtable}[]{@{}llll@{}}
\toprule\noalign{}
& Food A & Food B & Requirement \\
\midrule\noalign{}
\endhead
\bottomrule\noalign{}
\endlastfoot
Protein (g/unit) & 15 & 10 & ≥ 60 \\
Calcium (mg/unit) & 200 & 150 & ≥ 800 \\
Energy (kcal/unit) & 400 & 300 & ≤ 2000 \\
Cost (\$/unit) & 1.20 & 0.80 & minimise \\
\end{longtable}

\[\text{minimise} \quad z = 1.20x_1 + 0.80x_2\]

subject to:

\[15x_1 + 10x_2 \geq 60\] \[200x_1 + 150x_2 \geq 800\]
\[400x_1 + 300x_2 \leq 2000\] \[x_1, x_2 \geq 0\]

For \(\geq\) constraints, introduce \textbf{surplus variables}
\(s_i \geq 0\): \(15x_1 + 10x_2 - s_1 = 60\). The minus sign makes
\(s_i\) the excess above the minimum.

Note that the diet problem's standard form, with its surplus variables
and \(\geq\) constraints, is not directly ready for the simplex method
as presented below. With surplus variables, the origin is no longer a
feasible starting point (it violates the minimum requirements), and
finding an initial basic feasible solution requires an additional
technique --- \emph{artificial variables} --- that is covered in more
advanced treatments. The simplex worked example uses only \(\leq\)
constraints to keep the starting point clear.

\section{Graphical method}\label{graphical-method}

For two decision variables you can draw the \textbf{feasible region} ---
the set of all \((x_1, x_2)\) satisfying every constraint --- and read
off the optimum. More importantly, you can \emph{see} the geometry that
the simplex method will exploit algebraically when there are too many
variables to draw.

For the factory problem, plot each constraint boundary:

\begin{itemize}
\tightlist
\item
  \(2x_1 + x_2 = 8\): line through \((4, 0)\) and \((0, 8)\). Feasible
  side: below/left.
\item
  \(x_1 + 3x_2 = 9\): line through \((9, 0)\) and \((0, 3)\). Feasible
  side: below/left.
\item
  \(x_1, x_2 \geq 0\): first quadrant only.
\end{itemize}

The feasible region is a convex polygon --- a quadrilateral with four
vertices: \((0,0)\), \((4,0)\), \((3,2)\), and \((0,3)\).

\textbf{Corner-point theorem.} If an LP has an optimal solution, there
is an optimal solution at a vertex (corner point) of the feasible
region. This is the foundational geometric fact: the objective function
\(z = 30x_1 + 50x_2\) is a family of parallel lines (iso-profit lines).
As you push the line in the direction of increasing profit, the last
point of contact with the feasible region is a vertex.

Evaluate the objective at each vertex:

\begin{longtable}[]{@{}ll@{}}
\toprule\noalign{}
Vertex & \(z = 30x_1 + 50x_2\) \\
\midrule\noalign{}
\endhead
\bottomrule\noalign{}
\endlastfoot
\((0, 0)\) & \(0\) \\
\((4, 0)\) & \(120\) \\
\((3, 2)\) & \(190\) \\
\((0, 3)\) & \(150\) \\
\end{longtable}

Maximum profit is \textbf{\$190} per day, achieved by making 3 chairs
and 2 tables.

The intersection point \((3, 2)\) comes from solving \(2x_1 + x_2 = 8\)
and \(x_1 + 3x_2 = 9\) simultaneously. Subtract twice the second from
the first: \(-5x_2 = -10\), so \(x_2 = 2\), \(x_1 = 3\).

\begin{tcolorbox}[enhanced jigsaw, opacityback=0, colback=white, breakable, colframe=quarto-callout-note-color-frame, arc=.35mm, left=2mm, toprule=.15mm, leftrule=.75mm, bottomrule=.15mm, rightrule=.15mm]
\begin{minipage}[t]{5.5mm}
\textcolor{quarto-callout-note-color}{\faInfo}
\end{minipage}%
\begin{minipage}[t]{\textwidth - 5.5mm}

\textbf{Why the optimum is always at a vertex}

The feasible region of an LP is a convex polytope --- any two feasible
points can be connected by a straight line that stays entirely inside
the region. The objective function is linear, so it increases at a
constant rate in one direction. If the optimum were in the interior, you
could always move a little further in the improving direction and still
stay feasible, reaching a higher objective value. The only place where
you genuinely cannot improve further is a boundary. And on a boundary,
you can keep moving along the edge until you hit a vertex. This is why
the simplex method only needs to visit vertices.

\end{minipage}%
\end{tcolorbox}

What the graphical method shows directly is what the simplex method does
systematically: start at one vertex (here, the origin), check whether
any neighbour is better, move to the best neighbour, and repeat. The
graphical picture makes the logic self-evident. The simplex method turns
that logic into algebra so it can work in a hundred dimensions instead
of two.

\section{Simplex method}\label{simplex-method}

The graphical method breaks down beyond two variables. The simplex
method does the same thing --- walk from vertex to vertex, improving the
objective each step --- algebraically, using the tableau.

\textbf{Setup: slack variables and initial tableau.} Return to the
factory problem in standard form with \(x_1, x_2\) (original variables)
and \(s_1, s_2\) (slacks):

\[2x_1 + x_2 + s_1 = 8\] \[x_1 + 3x_2 + s_2 = 9\]
\[z = 30x_1 + 50x_2 \quad \Longrightarrow \quad -30x_1 - 50x_2 + z = 0\]

The \textbf{basic feasible solution} (BFS) at the origin: set non-basic
variables \(x_1 = x_2 = 0\), read off basic variables \(s_1 = 8\),
\(s_2 = 9\), \(z = 0\). This is the vertex \((0, 0)\).

\textbf{Initial tableau:}

\[\begin{array}{c|cccc|c}
\text{Basic} & x_1 & x_2 & s_1 & s_2 & \text{RHS} \\ \hline
s_1 & 2 & 1 & 1 & 0 & 8 \\
s_2 & 1 & 3 & 0 & 1 & 9 \\ \hline
z   & -30 & -50 & 0 & 0 & 0
\end{array}\]

\textbf{Pivot selection.} Look at the objective row (bottom row). The
most negative coefficient is \(-50\) in the \(x_2\) column: \(x_2\)
enters the basis. The column \(x_2\) is the \emph{entering column}.

To choose the \emph{leaving variable}, compute the minimum ratio:
\(\text{RHS} /
\text{entering-column coefficient}\) for rows with positive pivot-column
entries only. Rows with zero or negative entries in the pivot column are
skipped: a zero entry means increasing \(x_2\) has no effect on that
constraint, and a negative entry means the constraint becomes
\emph{easier} to satisfy as \(x_2\) increases --- neither constrains the
move.

\[\frac{8}{1} = 8 \qquad \frac{9}{3} = 3\]

Minimum ratio is 3, so \(s_2\) leaves. The pivot element is the entry in
row \(s_2\), column \(x_2\): that is \(3\).

\textbf{Pivot operation.} Divide the \(s_2\) row by 3 to make the pivot
element 1:

\[\text{new row 2}: \quad \tfrac{1}{3}x_1 + x_2 + 0 \cdot s_1 + \tfrac{1}{3}s_2 = 3\]

Eliminate \(x_2\) from every other row by adding/subtracting multiples
of the new row 2:

\begin{itemize}
\item
  Row 1 \(\leftarrow\) Row 1 \(-\) 1 \(\times\) new row 2:
  \((2 - \frac{1}{3})x_1 + 0 \cdot x_2 + s_1 - \frac{1}{3}s_2 = 8 - 3\)
  \(\Rightarrow \tfrac{5}{3}x_1 + s_1 - \tfrac{1}{3}s_2 = 5\)
\item
  Objective row \(\leftarrow\) Obj row \(+\) 50 \(\times\) new row 2:
  \((-30 + \frac{50}{3})x_1 + 0 \cdot x_2 + 0 \cdot s_1 + \frac{50}{3}s_2 = 150\)
  \(\Rightarrow -\tfrac{40}{3}x_1 + \tfrac{50}{3}s_2 = 150\)
\end{itemize}

\textbf{Tableau after first pivot (\(x_2\) enters, \(s_2\) leaves):}

\[\begin{array}{c|cccc|c}
\text{Basic} & x_1 & x_2 & s_1 & s_2 & \text{RHS} \\ \hline
s_1 & 5/3 & 0 & 1 & -1/3 & 5 \\
x_2 & 1/3 & 1 & 0 &  1/3 & 3 \\ \hline
z   & -40/3 & 0 & 0 & 50/3 & 150
\end{array}\]

Current BFS: \(x_1 = 0\), \(x_2 = 3\), \(s_1 = 5\), \(s_2 = 0\),
\(z = 150\). This is vertex \((0, 3)\).

\textbf{Second pivot.} Most negative objective coefficient: \(-40/3\) in
the \(x_1\) column. \(x_1\) enters. Minimum ratio:

\[\frac{5}{5/3} = 3 \qquad \frac{3}{1/3} = 9\]

\(s_1\) leaves. Pivot element: \(5/3\). Divide row 1 by \(5/3\):

\[\text{new row 1}: \quad x_1 + 0 \cdot x_2 + \tfrac{3}{5}s_1 - \tfrac{1}{5}s_2 = 3\]

Eliminate \(x_1\) from other rows:

\begin{itemize}
\item
  Row 2 \(\leftarrow\) Row 2 \(-\) \(\frac{1}{3}\) \(\times\) new row 1:
  \(0 \cdot x_1 + x_2 - \tfrac{1}{5}s_1 + \tfrac{2}{5}s_2 = 2\)
\item
  Obj row \(\leftarrow\) Obj row \(+\) \(\frac{40}{3}\) \(\times\) new
  row 1: \(0 \cdot x_1 + 8 \cdot s_1 + 14 \cdot s_2 = 190\)

  Let's be precise: \(-\tfrac{40}{3} + \tfrac{40}{3} \cdot 1 = 0\);
  \(s_1\) coefficient: \(0 + \tfrac{40}{3} \cdot \tfrac{3}{5} = 8\);
  \(s_2\) coefficient:
  \(\tfrac{50}{3} + \tfrac{40}{3} \cdot (-\tfrac{1}{5}) = \tfrac{50}{3} - \tfrac{8}{3} = 14\);
  RHS: \(150 + \tfrac{40}{3} \cdot 3 = 150 + 40 = 190\).
\end{itemize}

\textbf{Final tableau:}

\[\begin{array}{c|cccc|c}
\text{Basic} & x_1 & x_2 & s_1 & s_2 & \text{RHS} \\ \hline
x_1 & 1 & 0 &  3/5 & -1/5 & 3 \\
x_2 & 0 & 1 & -1/5 &  2/5 & 2 \\ \hline
z   & 0 & 0 &  8   & 14   & 190
\end{array}\]

All objective-row coefficients are non-negative. \textbf{Termination
condition met.} Optimal solution: \(x_1 = 3\), \(x_2 = 2\), \(z = 190\).

This matches the graphical answer exactly.

\begin{tcolorbox}[enhanced jigsaw, opacityback=0, colback=white, breakable, colframe=quarto-callout-note-color-frame, arc=.35mm, left=2mm, toprule=.15mm, leftrule=.75mm, bottomrule=.15mm, rightrule=.15mm]
\begin{minipage}[t]{5.5mm}
\textcolor{quarto-callout-note-color}{\faInfo}
\end{minipage}%
\begin{minipage}[t]{\textwidth - 5.5mm}

\textbf{Why simplex terminates}

Each pivot strictly improves the objective (assuming no degeneracy ---
no zero RHS values). There are finitely many vertices. So the method
must reach an optimum in a finite number of steps. In practice, simplex
visits far fewer vertices than the total count even for large LPs.

\end{minipage}%
\end{tcolorbox}

\section{Duality}\label{duality}

Every LP --- the \textbf{primal} --- has a companion LP called the
\textbf{dual}. The dual is not just a mathematical curiosity; it has
direct economic meaning and provides a powerful optimality check.

\textbf{Constructing the dual.} Start from the primal in standard
inequality form:

\[\text{maximise} \quad \mathbf{c}^\top \mathbf{x} \quad \text{subject to} \quad A\mathbf{x} \leq \mathbf{b},\ \mathbf{x} \geq \mathbf{0}\]

The dual introduces one dual variable \(y_i\) --- read ``y-sub-i'' ---
for each primal constraint. The dual is:

\[\text{minimise} \quad \mathbf{b}^\top \mathbf{y} \quad \text{subject to} \quad A^\top \mathbf{y} \geq \mathbf{c},\ \mathbf{y} \geq \mathbf{0}\]

The roles reverse: primal constraints become dual variables, primal
variables become dual constraints, primal objective coefficients become
dual RHS values, and primal RHS values become dual objective
coefficients. The maximisation becomes minimisation.

\textbf{Dual of the factory problem.} Primal: maximise \(30x_1 + 50x_2\)
subject to \(2x_1 + x_2 \leq 8\), \(x_1 + 3x_2 \leq 9\),
\(x_1, x_2 \geq 0\).

Let \(y_1, y_2\) be the dual variables (one per primal constraint):

\[\text{minimise} \quad 8y_1 + 9y_2\]

subject to:

\[2y_1 + y_2 \geq 30 \qquad \text{(column 1 of } A^\top\text{)}\]
\[y_1 + 3y_2 \geq 50 \qquad \text{(column 2 of } A^\top\text{)}\]
\[y_1, y_2 \geq 0\]

\textbf{Weak duality theorem.} For any primal feasible \(\mathbf{x}\)
and any dual feasible \(\mathbf{y}\):

\[\mathbf{c}^\top \mathbf{x} \leq \mathbf{b}^\top \mathbf{y}\]

\emph{Proof.} The dual constraint says
\(A^\top\mathbf{y} \geq \mathbf{c}\), meaning each component satisfies
\((A^\top\mathbf{y})_j \geq c_j\). Since \(x_j \geq 0\), multiplying
through and summing over all \(j\) gives
\(\mathbf{c}^\top \mathbf{x} \leq (A^\top\mathbf{y})^\top \mathbf{x} = \mathbf{y}^\top A\mathbf{x}\).
The primal constraint \(A\mathbf{x} \leq \mathbf{b}\) combined with
\(\mathbf{y} \geq \mathbf{0}\) then gives
\(\mathbf{y}^\top A\mathbf{x} \leq \mathbf{y}^\top \mathbf{b} = \mathbf{b}^\top\mathbf{y}\).
Chaining both inequalities:
\(\mathbf{c}^\top \mathbf{x} \leq \mathbf{b}^\top\mathbf{y}\).
\(\square\)

The strong duality theorem (proved via simplex theory) says the
inequality is tight at optimality: primal and dual optimal values are
equal.

\textbf{Economic interpretation: shadow prices.} The optimal dual
variables have a direct meaning. \(y_1\) is the \textbf{shadow price} of
the first constraint (cutting hours): it tells you how much the optimal
objective increases if you add one unit of that resource. At optimality,
\(y_1 = 8\) and \(y_2 = 14\) --- each additional cutting hour is worth
\$8, each additional finishing hour is worth \$14. You can verify:
\(8y_1 + 9y_2 = 8 \times 8 + 9 \times 14 = 64 + 126 = 190\). Dual
optimal matches primal optimal --- strong duality holds.

\textbf{Complementary slackness.} At optimality, either a primal
constraint is tight (active, \(A_i \mathbf{x} = b_i\)) or the
corresponding dual variable is zero (\(y_i = 0\)); and either a dual
constraint is tight or the corresponding primal variable is zero.
Formally: \(y_i(b_i - A_i\mathbf{x}) = 0\) and
\((A^\top\mathbf{y} - \mathbf{c})_j x_j = 0\) for all \(i, j\). This
gives you a system of equations to verify optimality without running the
full simplex.

\section{Sensitivity analysis}\label{sensitivity-analysis}

The optimal basis found by simplex remains optimal as long as the
problem data stays within certain ranges. If the profit margin on
\(x_1\) changes from 30 to \(30 + \Delta\), the basis \(\{x_1, x_2\}\)
stays optimal while the objective-row coefficient of \(s_1\) (currently
8) remains non-negative and the coefficient of \(s_2\) (currently 14)
remains non-negative. The \(x_1\) basic row is \([1, 0, 3/5, -1/5, 3]\),
so changing \(c_1\) by \(\Delta\) shifts the \(s_1\) coefficient to
\(8 - \tfrac{3}{5}\Delta\) and the \(s_2\) coefficient to
\(14 + \tfrac{1}{5}\Delta\). Requiring both to remain non-negative:
\(8 - \tfrac{3}{5}\Delta \geq 0 \Rightarrow \Delta \leq \tfrac{40}{3}\)
and \(14 + \tfrac{1}{5}\Delta \geq 0 \Rightarrow \Delta \geq -70\). The
basis is stable for \(-70 \leq \Delta \leq \tfrac{40}{3}\), i.e.,
\(c_1 \in [-40, 43\tfrac{1}{3}]\). Similarly, if the RHS \(b_1 = 8\)
changes, the solution coordinates shift but the same basis remains
feasible as long as all basic variables stay non-negative --- the range
can be read from the current tableau's column for \(s_1\). This analysis
is done without re-running simplex.

\section{Where this goes}\label{where-this-goes-33}

Chapter 2 of this section develops \textbf{network flow problems} ---
shortest path, maximum flow, minimum cost flow --- all of which can be
formulated as LPs with special structure (network matrices) that allows
faster algorithms than general simplex. Network LPs arise everywhere:
logistics, circuit analysis, pipeline routing.

The broader field of \textbf{integer programming} (mixed-integer LP)
handles the case where some or all variables must be integers ---
scheduling problems, facility location, combinatorial optimisation.
Branch-and-bound solves these by solving a sequence of LP relaxations.
The LP theory in this chapter is the computational core of those
methods.

\begin{tcolorbox}[enhanced jigsaw, opacityback=0, colback=white, breakable, colframe=quarto-callout-tip-color-frame, arc=.35mm, left=2mm, toprule=.15mm, leftrule=.75mm, bottomrule=.15mm, rightrule=.15mm]
\begin{minipage}[t]{5.5mm}
\textcolor{quarto-callout-tip-color}{\faLightbulb}
\end{minipage}%
\begin{minipage}[t]{\textwidth - 5.5mm}

\textbf{Where LP shows up}

\begin{itemize}
\tightlist
\item
  A refinery blends crude feedstocks to meet fuel-grade specifications
  at minimum cost --- an LP with 50--500 variables, solved daily.
\item
  An airline schedules crews to cover flights while meeting duty-time
  regulations --- an integer LP with millions of variables in practice.
\item
  A fund manager maximises expected portfolio return subject to risk
  limits and position constraints --- LP or quadratic programming
  depending on the risk model.
\item
  A nutritionist formulates minimum-cost meal plans meeting dietary
  requirements --- the diet problem, which drove early LP development.
\item
  A chip manufacturer allocates fab capacity across product lines to
  maximise margin --- LP with supply chain constraints.
\end{itemize}

\end{minipage}%
\end{tcolorbox}

\section{Exercises}\label{exercises-44}

\begin{tcolorbox}[enhanced jigsaw, opacityback=0, colback=white, breakable, colframe=quarto-callout-note-color-frame, arc=.35mm, left=2mm, toprule=.15mm, leftrule=.75mm, bottomrule=.15mm, rightrule=.15mm]
\begin{minipage}[t]{5.5mm}
\textcolor{quarto-callout-note-color}{\faInfo}
\end{minipage}%
\begin{minipage}[t]{\textwidth - 5.5mm}

These are puzzles. Each has a clean optimal solution. The interesting
work is in the formulation --- translating the situation into variables,
objective, and constraints --- and in executing the simplex steps
carefully.

\end{minipage}%
\end{tcolorbox}

\textbf{Exercise 1.} A workshop makes two products, shelves and benches.
Each shelf uses 3 kg of timber and 1 hour of labour. Each bench uses 2
kg of timber and 4 hours of labour. The workshop has 24 kg of timber and
16 hours of labour available per week. Profit is \$20 per shelf and \$45
per bench. Formulate the LP to maximise weekly profit.

\textbf{Exercise 2.} Solve the LP from Exercise 1 graphically. Plot both
constraint boundaries, shade the feasible region, evaluate the objective
at every vertex, and identify the optimal solution.

\textbf{Exercise 3.} Convert the LP from Exercise 1 to standard form by
introducing slack variables. Write out the full system of equality
constraints and identify the initial basic feasible solution (the origin
vertex).

\textbf{Exercise 4.} Solve the following LP using the simplex tableau
method. Show all pivots.

\[\text{maximise} \quad z = 5x_1 + 4x_2\]

subject to:

\[6x_1 + 4x_2 \leq 24\] \[x_1 + 2x_2 \leq 6\] \[x_1, x_2 \geq 0\]

\textbf{Exercise 5.} Write the dual of the LP in Exercise 4. Then use
complementary slackness to find the dual optimal solution
\((y_1, y_2)\). Verify strong duality by checking that primal and dual
objective values agree at optimality.

\textbf{Exercise 6.} In Exercise 4, the optimal basis has \(x_1\) basic
(value 4) and \(s_2\) basic (value 2). By how much can the objective
coefficient of \(x_1\) (currently 5) change before the current basis is
no longer optimal? Give the range.

\chapter{Graphs and networks}\label{graphs-and-networks}

Paths, flows, and spanning structures

\hfill\break

You want to drive from Edmonton to Calgary in the least time, not just
the fewest kilometres. You need to wire up five remote sensors to a data
logger using as little cable as possible. You want to know the maximum
throughput a pipeline network can deliver before a bottleneck binds.

Different questions. The same underlying structure: nodes, connections,
weights. Once you have that structure, a handful of classical algorithms
--- each with a clean proof of correctness --- answer all of them.

\section{Graph fundamentals}\label{graph-fundamentals}

A \textbf{graph} \(G = (V, E)\) --- read ``a graph G consisting of
vertex set V and edge set E'' --- is a set of \textbf{vertices} (also
called nodes) \(V\) together with a set of \textbf{edges} \(E\), where
each edge connects two vertices.

In an \textbf{undirected} graph, the edge \(\{u, v\}\) can be traversed
in either direction. In a \textbf{directed} graph (digraph), each edge
\((u, v)\) --- read ``from u to v'' --- has a direction: you can go from
\(u\) to \(v\) but not necessarily back.

A \textbf{weighted} graph assigns a real number \(w(u, v)\) --- ``the
weight of the edge from \(u\) to \(v\)'' --- to each edge. Weights model
distance, cost, travel time, capacity, or any other quantity associated
with a connection.

\subsection{Adjacency matrix}\label{adjacency-matrix}

The \textbf{adjacency matrix} \(A\) of a graph on \(n\) vertices is the
\(n \times n\) matrix where entry \(A_{ij} = 1\) if there is an edge
from vertex \(i\) to vertex \(j\), and \(A_{ij} = 0\) otherwise. For a
weighted graph, replace the 1 with the edge weight. For an undirected
graph, \(A\) is symmetric: \(A_{ij} = A_{ji}\).

Example: a simple undirected graph with four vertices \(\{1, 2, 3, 4\}\)
and edges \(\{1,2\}\), \(\{1,3\}\), \(\{2,4\}\), \(\{3,4\}\) has
adjacency matrix

\[A = \begin{pmatrix} 0 & 1 & 1 & 0 \\ 1 & 0 & 0 & 1 \\ 1 & 0 & 0 & 1 \\ 0 & 1 & 1 & 0 \end{pmatrix}.\]

\subsection{Degree, paths, and
connectivity}\label{degree-paths-and-connectivity}

The \textbf{degree} \(d(v)\) --- ``the degree of vertex \(v\)'' --- is
the number of edges incident to \(v\). In the example above,
\(d(1) = 2\), \(d(2) = 2\), \(d(3) = 2\), \(d(4) = 2\).

A \textbf{path} from \(u\) to \(v\) is a sequence of vertices starting
at \(u\) and ending at \(v\) where consecutive vertices are connected by
an edge, and no vertex repeats. A \textbf{cycle} is a closed path: it
starts and ends at the same vertex. A graph is \textbf{connected} if
there is a path between every pair of vertices.

A \textbf{tree} is a connected graph with no cycles. Equivalently, a
tree on \(n\) vertices has exactly \(n-1\) edges --- removing any edge
disconnects it, adding any edge creates a cycle. Trees are the minimal
connected structures, which is exactly why they appear in spanning
problems.

\section{Shortest path: Dijkstra's
algorithm}\label{shortest-path-dijkstras-algorithm}

\textbf{Problem.} Given a weighted graph with non-negative edge weights
and a source vertex \(s\), find the shortest-distance path from \(s\) to
every other vertex.

\textbf{Algorithm.} Maintain a set \(S\) of vertices whose shortest
distances are finalised, and for every vertex outside \(S\), a tentative
distance \(\text{dist}[v]\). Initially \(\text{dist}[s] = 0\) and
\(\text{dist}[v] = \infty\) for all \(v \neq s\).

Repeat until \(S = V\):

\begin{enumerate}
\def\labelenumi{\arabic{enumi}.}
\tightlist
\item
  Pick the vertex \(u \notin S\) with the smallest \(\text{dist}[u]\).
\item
  Add \(u\) to \(S\) (its distance is now final).
\item
  For each neighbour \(v\) of \(u\): if
  \(\text{dist}[u] + w(u,v) < \text{dist}[v]\), update
  \(\text{dist}[v] \leftarrow \text{dist}[u] + w(u,v)\).
\end{enumerate}

\subsection{Worked trace --- 6-node
graph}\label{worked-trace-6-node-graph}

Consider the following undirected weighted graph with vertices
\(\{A, B, C, D, E, F\}\) and source \(A\):

\begin{longtable}[]{@{}ll@{}}
\toprule\noalign{}
Edge & Weight \\
\midrule\noalign{}
\endhead
\bottomrule\noalign{}
\endlastfoot
A--B & 4 \\
A--C & 2 \\
B--C & 5 \\
B--D & 10 \\
C--E & 3 \\
E--D & 4 \\
D--F & 11 \\
E--F & 7 \\
\end{longtable}

Initial distances: \(A=0\), \(B=\infty\), \(C=\infty\), \(D=\infty\),
\(E=\infty\), \(F=\infty\).

\textbf{Step 1.} Smallest tentative distance is \(A\) (dist = 0).
Finalise \(A\). Update neighbours: \(\text{dist}[B] \leftarrow 4\),
\(\text{dist}[C] \leftarrow 2\).

\textbf{Step 2.} Smallest non-finalised: \(C\) (dist = 2). Finalise
\(C\). Update neighbours: \(B\): \(2+5=7 > 4\), no change. \(E\):
\(2+3=5 \leftarrow 5\).

\textbf{Step 3.} Smallest non-finalised: \(B\) (dist = 4). Finalise
\(B\). Update neighbours: \(C\): already finalised. \(D\):
\(4+10=14 \leftarrow 14\).

\textbf{Step 4.} Smallest non-finalised: \(E\) (dist = 5). Finalise
\(E\). Update neighbours: \(D\): \(5+4=9 < 14\), update
\(\text{dist}[D] \leftarrow 9\). \(F\): \(5+7=12 \leftarrow 12\).

\textbf{Step 5.} Smallest non-finalised: \(D\) (dist = 9). Finalise
\(D\). Update neighbours: \(F\): \(9+11=20 > 12\), no change.

\textbf{Step 6.} Smallest non-finalised: \(F\) (dist = 12). Finalise
\(F\).

Final shortest distances from \(A\): \(B=4\), \(C=2\), \(D=9\), \(E=5\),
\(F=12\).

The shortest path to \(F\) is \(A \to C \to E \to F\) with total weight
\(2+3+7=12\).

\subsection{Correctness sketch}\label{correctness-sketch}

Dijkstra is a greedy algorithm. Before the proof, it's worth asking:
\emph{why would the greedy choice ever go wrong?} The answer is: it
could go wrong if a later, distant vertex \(y\) had a shortcut that made
a longer-looking path to \(u\) end up shorter. Non-negative edge weights
prevent exactly this --- you can never shorten a path by adding edges.

The key claim: when a vertex \(u\) is moved into \(S\), its tentative
distance \(\text{dist}[u]\) is the true shortest distance.

Proof by contradiction: suppose some path to \(u\) is shorter, and let
\(y\) be the first vertex on that path not yet in \(S\) when \(u\) is
picked. The portion of that path from the source to \(y\) has length
\(\text{dist}[y]\) (since Dijkstra has already found the shortest path
to \(y\), which lies on our assumed shorter path). Any further edges
from \(y\) to \(u\) have non-negative weight, so the full path is at
least \(\text{dist}[y]\) in length. But
\(\text{dist}[y] \geq \text{dist}[u]\), because Dijkstra would have
picked \(y\) before \(u\) if \(y\) were closer. So the assumed shorter
path is at least as long as \(\text{dist}[u]\) --- a contradiction.

This argument fails if any edge weight is negative --- a later edge
could make up for a large earlier one. For graphs with negative weights
(but no negative-weight cycles), \textbf{Bellman-Ford} is the correct
algorithm: it relaxes all edges \(|V|-1\) times, running in
\(O(|V||E|)\) time rather than Dijkstra's \(O((|V| + |E|) \log |V|)\).
One use case is currency arbitrage detection.

\section{Minimum spanning tree}\label{minimum-spanning-tree}

\textbf{Problem.} Given a connected weighted undirected graph, find a
spanning tree --- a tree that includes all vertices --- with minimum
total edge weight.

Two classical algorithms solve this; both are correct, and both rely on
the same underlying property.

\subsection{Cut property}\label{cut-property}

A \textbf{cut} is a partition of \(V\) into two non-empty sets \(S\) and
\(V \setminus S\). An edge \textbf{crosses} the cut if one endpoint is
in \(S\) and the other is in \(V \setminus S\). The cut property states:
for any cut \((S, V \setminus S)\), if \(e\) is the unique
minimum-weight edge crossing that cut, then \(e\) belongs to every
minimum spanning tree; more generally, any minimum-weight edge crossing
a given cut belongs to some minimum spanning tree. This is the single
fact both Kruskal and Prim exploit.

Here is how this applies: in Kruskal's algorithm, every time we consider
an edge \(e = \{u,v\}\), the two components containing \(u\) and \(v\)
define a cut. The edge \(e\) is the cheapest edge crossing that cut
(because we process edges in weight order and skip any that close a
cycle). The cut property says it is safe to include \(e\) in the MST.
Every addition Kruskal makes is justified by exactly this argument.

\subsection{Kruskal's algorithm}\label{kruskals-algorithm}

\begin{enumerate}
\def\labelenumi{\arabic{enumi}.}
\tightlist
\item
  Sort all edges by weight (ascending).
\item
  Initialise \(T = \emptyset\) (no edges selected yet).
\item
  Consider edges in order. Add edge \(e = \{u, v\}\) to \(T\) if it does
  not create a cycle (equivalently, if \(u\) and \(v\) are in different
  connected components of \(T\)).
\item
  Stop when \(T\) has \(n-1\) edges.
\end{enumerate}

Use a \textbf{union-find} data structure to test and merge components
efficiently. Total time: \(O(|E| \log |E|)\) (dominated by sorting).

\textbf{Worked example.} Six-node graph with edges (sorted by weight):

\begin{longtable}[]{@{}ll@{}}
\toprule\noalign{}
Edge & Weight \\
\midrule\noalign{}
\endhead
\bottomrule\noalign{}
\endlastfoot
C--E & 3 \\
A--C & 2 \\
A--B & 4 \\
E--D & 4 \\
B--C & 5 \\
E--F & 7 \\
B--D & 10 \\
D--F & 11 \\
\end{longtable}

Sorted: A--C (2), C--E (3), A--B (4), E--D (4), B--C (5), E--F (7), B--D
(10), D--F (11).

\begin{itemize}
\tightlist
\item
  Add A--C (2): no cycle. \(T = \{\)A--C\(\}\).
\item
  Add C--E (3): no cycle. \(T = \{\)A--C, C--E\(\}\).
\item
  Add A--B (4): no cycle. \(T = \{\)A--C, C--E, A--B\(\}\).
\item
  Add E--D (4): no cycle. \(T = \{\)A--C, C--E, A--B, E--D\(\}\).
\item
  Add B--C (5): creates cycle A--B--C--A. Skip.
\item
  Add E--F (7): no cycle. \(T = \{\)A--C, C--E, A--B, E--D, E--F\(\}\).
  Now \(|T| = 5 = n-1\). Done.
\end{itemize}

MST total weight: \(2+3+4+4+7 = 20\).

\subsection{Prim's algorithm}\label{prims-algorithm}

\begin{enumerate}
\def\labelenumi{\arabic{enumi}.}
\tightlist
\item
  Start from an arbitrary seed vertex \(r\). Mark \(r\) as in the tree.
\item
  At each step, add the minimum-weight edge that connects a vertex
  inside the current tree to one outside it.
\item
  Repeat until all vertices are in the tree.
\end{enumerate}

\textbf{Worked example} (same graph, starting from \(A\)):

\begin{itemize}
\tightlist
\item
  Tree: \(\{A\}\). Cheapest crossing edge: A--C (2). Add C.
\item
  Tree: \(\{A, C\}\). Cheapest crossing edge: C--E (3). Add E.
\item
  Tree: \(\{A, C, E\}\). Cheapest crossing edge: A--B (4) or E--D (4).
  Take A--B. Add B.
\item
  Tree: \(\{A, C, E, B\}\). Cheapest crossing edge: E--D (4). Add D.
\item
  Tree: \(\{A, C, E, B, D\}\). Cheapest crossing edge: E--F (7). Add F.
\end{itemize}

MST edges: A--C, C--E, A--B, E--D, E--F. Total weight: 20. Same result
as Kruskal.

Both algorithms always produce an MST. When the minimum-weight crossing
edge for any cut is unique, the MST is unique. With ties, multiple MSTs
of the same total weight may exist.

\section{Network flows}\label{network-flows}

\textbf{Problem.} A directed weighted graph has a distinguished
\textbf{source} \(s\) (supply) and \textbf{sink} \(t\) (demand). Each
directed edge \((u,v)\) has a \textbf{capacity} \(c(u,v) \geq 0\) ---
the maximum flow it can carry. The \textbf{flow} \(f(u,v)\) on each edge
must satisfy:

\begin{itemize}
\tightlist
\item
  \textbf{Capacity constraint:} \(0 \leq f(u,v) \leq c(u,v)\) for all
  edges.
\item
  \textbf{Flow conservation:} for every vertex \(v \neq s, t\), the
  total flow in equals the total flow out:
  \(\sum_u f(u,v) = \sum_w f(v,w)\).
\end{itemize}

The \textbf{value of the flow} is the total flow leaving \(s\)
(equivalently, arriving at \(t\)). The problem is to maximise this
value.

\subsection{Ford-Fulkerson method}\label{ford-fulkerson-method}

The key idea behind Ford-Fulkerson is that earlier routing decisions can
be wrong, and the algorithm needs a way to undo them. If we sent flow
along a path that turned out to block a better option, we want to be
able to redirect it. The \textbf{residual graph} makes this possible by
including a reverse edge for every forward edge: sending flow
``backwards'' on a reverse edge is equivalent to cancelling some
previously routed flow.

\begin{enumerate}
\def\labelenumi{\arabic{enumi}.}
\tightlist
\item
  Start with zero flow on all edges.
\item
  Find an \textbf{augmenting path} --- any directed path from \(s\) to
  \(t\) in the \textbf{residual graph} \(G_f\), where edge \((u,v)\) has
  residual capacity \(c(u,v) - f(u,v)\) (how much more we can send
  forward) and the reverse edge \((v,u)\) has residual capacity
  \(f(u,v)\) (how much previously routed flow we can cancel).
\item
  Send as much flow along the augmenting path as the bottleneck edge
  allows.
\item
  Repeat until no augmenting path exists.
\end{enumerate}

\subsection{Max-flow min-cut theorem}\label{max-flow-min-cut-theorem}

A \textbf{cut} \((S, T)\) with \(s \in S\) and \(t \in T\) separates
source from sink. The \textbf{capacity of the cut} is the sum of
capacities of edges going from \(S\) to \(T\). The max-flow min-cut
theorem states:

\begin{quote}
\textbf{The maximum flow value equals the minimum cut capacity.}
\end{quote}

The forward direction is clear: no flow can exceed the capacity of any
cut separating \(s\) from \(t\), because every unit of flow must cross
the cut. The deep direction is that Ford-Fulkerson achieves this bound
--- when no augmenting path remains, the set of vertices reachable from
\(s\) in the residual graph defines a cut whose capacity equals the
current flow value.

In practice: to find the bottleneck in a network, find the minimum cut.
The edges in the minimum cut are the ones whose capacity, if increased,
would allow more total flow.

\subsection{Worked example --- 5-node
network}\label{worked-example-5-node-network}

Directed graph: source \(S\), sink \(T\), intermediate vertices \(A\),
\(B\), \(C\).

\begin{longtable}[]{@{}ll@{}}
\toprule\noalign{}
Edge & Capacity \\
\midrule\noalign{}
\endhead
\bottomrule\noalign{}
\endlastfoot
S→A & 10 \\
S→B & 8 \\
A→C & 9 \\
A→B & 3 \\
B→T & 10 \\
C→T & 7 \\
C→B & 2 \\
\end{longtable}

\textbf{Round 1.} Augmenting path \(S \to A \to C \to T\). Bottleneck:
\(\min(10,9,7) = 7\). Send 7 units. Updated flows: S→A = 7, A→C = 7, C→T
= 7.

\textbf{Round 2.} Augmenting path \(S \to A \to B \to T\). Bottleneck:
\(\min(3,3,10) = 3\). Send 3 units. Updated flows: S→A = 10, A→B = 3,
B→T = 3.

\textbf{Round 3.} Augmenting path \(S \to B \to T\). Bottleneck:
\(\min(8,7) = 7\). Send 7 units. Updated flows: S→B = 7, B→T = 10.

\textbf{Round 4.} Path \(S \to A\): S→A is saturated (10/10).
\(S \to B\): B→T is saturated (10/10). No augmenting path exists. Done.

Maximum flow = 7 + 3 + 7 = 17 units.

\textbf{Minimum cut.} After Round 3, the residual capacities are: S→A:
0, S→B: 1, A→C: 2, A→B: 0, B→T: 0, C→T: 0, C→B: 2, plus reverse edges
A←S: 10, B←S: 7, C←A: 7, B←A: 3, T←B: 10, T←C: 7.

Starting from \(S\) in the residual graph: S→B has residual 1, so \(B\)
is reachable. From \(B\): B→T has residual 0 (blocked); the reverse edge
B←A carries residual 3, so \(A\) is reachable. From \(A\): A→C has
residual 2, so \(C\) is reachable. From \(C\): C→T has residual 0
(blocked); C→B forward has residual 2 but \(B\) is already in the set.

Reachable set: \(\{S, B, A, C\}\). Non-reachable: \(\{T\}\).

The minimum cut consists of the edges crossing from \(\{S,B,A,C\}\) to
\(\{T\}\): that is B→T (capacity 10) and C→T (capacity 7). Min cut
capacity \(= 10 + 7 = 17\), which equals the max flow.

The minimum cut \(\{\)B→T, C→T\(\}\) identifies the bottleneck: these
are the edges whose capacity limits the total throughput of the network.

\begin{tcolorbox}[enhanced jigsaw, opacityback=0, colback=white, breakable, colframe=quarto-callout-note-color-frame, arc=.35mm, left=2mm, toprule=.15mm, leftrule=.75mm, bottomrule=.15mm, rightrule=.15mm]
\begin{minipage}[t]{5.5mm}
\textcolor{quarto-callout-note-color}{\faInfo}
\end{minipage}%
\begin{minipage}[t]{\textwidth - 5.5mm}

\textbf{Why Ford-Fulkerson terminates}

With integer capacities, each augmentation increases the flow by at
least 1. Since the flow is bounded above by the sum of capacities out of
\(s\), the algorithm terminates. The reverse edges in the residual graph
are what make it work correctly: they allow the algorithm to ``undo'' a
routing decision and reroute flow along a more productive path.

\end{minipage}%
\end{tcolorbox}

\section{Where this goes}\label{where-this-goes-34}

The three problems in this chapter --- shortest path, MST, network flow
--- are the foundational problems of \textbf{combinatorial
optimisation}. They appear directly inside more complex algorithms:
Dijkstra is the inner loop of many routing protocols; MST computation
underlies Christofides' approximation algorithm for the travelling
salesman problem; max-flow is equivalent to maximum bipartite matching
via a reduction. If you proceed to the study of algorithms or operations
research, you will see all three again.

The other extension worth knowing is \textbf{integer programming}, where
variables are constrained to take integer values. Many practical network
problems --- crew scheduling, facility location, vehicle routing --- are
integer programs. Chapter 1 of this section covered the continuous
relaxation; the integer case is in general NP-hard, but the network
structure of the problems above makes them tractable.

\begin{tcolorbox}[enhanced jigsaw, opacityback=0, colback=white, breakable, colframe=quarto-callout-tip-color-frame, arc=.35mm, left=2mm, toprule=.15mm, leftrule=.75mm, bottomrule=.15mm, rightrule=.15mm]
\begin{minipage}[t]{5.5mm}
\textcolor{quarto-callout-tip-color}{\faLightbulb}
\end{minipage}%
\begin{minipage}[t]{\textwidth - 5.5mm}

\textbf{Where this shows up}

\begin{itemize}
\tightlist
\item
  A civil engineer routing roads or pipelines through terrain minimises
  a weighted path or tree.
\item
  A telecoms network engineer provisioning bandwidth uses max-flow to
  find the network's throughput limit.
\item
  A supply-chain analyst models goods movement as a flow problem and
  finds bottleneck warehouse links.
\item
  A machine learning engineer working on graph neural networks needs the
  adjacency matrix representation as the input to the network.
\item
  A biologist mapping evolutionary relationships uses minimum spanning
  trees on genetic distance matrices.
\end{itemize}

\end{minipage}%
\end{tcolorbox}

\section{Exercises}\label{exercises-45}

These are convergent problems. Each has a definite answer. The work is
in applying the algorithm correctly and reading the result.

\textbf{Exercise 1.} Run Dijkstra's algorithm on the following 5-node
undirected weighted graph, starting from vertex \(P\).

\begin{longtable}[]{@{}ll@{}}
\toprule\noalign{}
Edge & Weight \\
\midrule\noalign{}
\endhead
\bottomrule\noalign{}
\endlastfoot
P--Q & 6 \\
P--R & 2 \\
R--Q & 3 \\
R--S & 7 \\
Q--S & 1 \\
Q--T & 5 \\
S--T & 4 \\
\end{longtable}

Trace all distance updates. State the final shortest distance to each
vertex and give the shortest path to \(T\).

\textbf{Exercise 2.} Find the MST of the following 6-node graph using
Kruskal's algorithm. List the edges in the order they are added to the
tree and give the total MST weight.

\begin{longtable}[]{@{}ll@{}}
\toprule\noalign{}
Edge & Weight \\
\midrule\noalign{}
\endhead
\bottomrule\noalign{}
\endlastfoot
1--2 & 7 \\
1--3 & 5 \\
2--4 & 8 \\
2--3 & 9 \\
3--4 & 15 \\
3--5 & 6 \\
4--5 & 4 \\
4--6 & 5 \\
5--6 & 2 \\
\end{longtable}

\textbf{Exercise 3.} Find the MST of the same 6-node graph from Exercise
2 using Prim's algorithm, starting from vertex 1. Verify that the total
MST weight equals the result from Exercise 2.

\textbf{Exercise 4.} Find the maximum flow from \(S\) to \(T\) in the
following 4-node network. Identify the minimum cut.

\begin{longtable}[]{@{}ll@{}}
\toprule\noalign{}
Edge & Capacity \\
\midrule\noalign{}
\endhead
\bottomrule\noalign{}
\endlastfoot
S→A & 8 \\
S→B & 6 \\
A→T & 5 \\
A→B & 4 \\
B→T & 9 \\
\end{longtable}

\textbf{Exercise 5.} Let \(G\) be a tree on \(n\) vertices. Prove that
the adjacency matrix \(A\) of \(G\) contains exactly \(n-1\) ones in
each of its upper and lower triangular parts (i.e., \(G\) has exactly
\(n-1\) edges), using only the definition of a tree as a connected
acyclic graph.

\textbf{Exercise 6 (Engineering application).} A water treatment
facility distributes purified water via a network of pipes from a
treatment plant (source \(S\)) to a distribution reservoir (sink \(T\)).
The intermediate pump stations are \(P\), \(Q\), \(R\).

\begin{longtable}[]{@{}ll@{}}
\toprule\noalign{}
Pipe & Capacity (ML/day) \\
\midrule\noalign{}
\endhead
\bottomrule\noalign{}
\endlastfoot
S→P & 12 \\
S→Q & 10 \\
P→Q & 3 \\
P→R & 8 \\
Q→T & 7 \\
R→Q & 4 \\
R→T & 9 \\
\end{longtable}

Model this as a max-flow problem. Find the maximum daily throughput from
\(S\) to \(T\) and identify the bottleneck cut --- the set of pipes
that, if their capacity were increased, would allow more water to reach
\(T\).

\chapter{Probability theory}\label{probability-theory-1}

Quantifying uncertainty

\hfill\break

A component comes off a production line. You don't know whether it's
defective. You test 1000 and 23 fail. What can you say about the next
one? About the next batch of 500?

These questions have precise answers. Getting them requires a framework
for describing uncertainty --- assigning likelihoods to outcomes,
combining them consistently, and extracting the quantities that are
actually useful.

\section{Probability axioms}\label{probability-axioms}

\subsection{Sample space and events}\label{sample-space-and-events}

The \textbf{sample space} \(\Omega\) (read: ``omega'') is the complete
set of possible outcomes of an experiment. For a coin flip,
\(\Omega = \{\text{heads}, \text{tails}\}\). For rolling a six-sided
die, \(\Omega = \{1, 2, 3, 4, 5, 6\}\). For the lifetime of a component
in hours, \(\Omega = [0, \infty)\).

An \textbf{event} \(A\) is any subset of \(\Omega\) --- a collection of
outcomes you care about. The event ``roll an even number'' is the subset
\(\{2, 4, 6\} \subset \Omega\).

\subsection{Kolmogorov axioms}\label{kolmogorov-axioms}

A \textbf{probability measure} \(P\) assigns a number to each event.
Three axioms define what that assignment must satisfy:

\begin{enumerate}
\def\labelenumi{\arabic{enumi}.}
\tightlist
\item
  \textbf{Non-negativity:} \(P(A) \geq 0\) for every event \(A\).
\item
  \textbf{Normalisation:} \(P(\Omega) = 1\) --- something must happen.
\item
  \textbf{Additivity:} For mutually exclusive events \(A\) and \(B\)
  (events that cannot both occur), \(P(A \cup B) = P(A) + P(B)\).
\end{enumerate}

These axioms are the whole foundation. Everything else is a consequence.

From them you can show:

\begin{itemize}
\tightlist
\item
  \(P(\emptyset) = 0\) (the impossible event has probability zero)
\item
  \(P(A^c) = 1 - P(A)\), where \(A^c\) is the complement of \(A\)
  (everything not in \(A\))
\item
  If \(A \subset B\) then \(P(A) \leq P(B)\)
\item
  The inclusion-exclusion rule:
  \(P(A \cup B) = P(A) + P(B) - P(A \cap B)\)
\end{itemize}

The third axiom extends to any countable collection of mutually
exclusive events:
\(P(A_1 \cup A_2 \cup \cdots) = P(A_1) + P(A_2) + \cdots\)

\subsection{Conditional probability}\label{conditional-probability}

The \textbf{conditional probability} of \(A\) given \(B\) --- written
\(P(A \mid B)\), read ``probability of \(A\) given \(B\)'' --- is
defined as:

\[P(A \mid B) = \frac{P(A \cap B)}{P(B)}, \quad P(B) > 0\]

This is the fraction of \(B\)'s probability that overlaps with \(A\). If
you know \(B\) has occurred, you're restricting attention to the
outcomes in \(B\) and asking what fraction of those are also in \(A\).

\begin{Shaded}
\begin{Highlighting}[]
\NormalTok{//| echo: false}
\NormalTok{\{}
\NormalTok{  // ── Conditional probability area model (OPPORTUNITY 1) ───────────────────}
\NormalTok{  const W=520, H=320, PAD=\{top:20,right:180,bottom:20,left:20\};}
\NormalTok{  const BLUE="rgba(59,130,246,0.55)", ORANGE="rgba(249,115,22,0.55)";}
\NormalTok{  const BLUE\_STROKE="\#3b82f6", ORANGE\_STROKE="\#f97316";}
\NormalTok{  const BLEND="rgba(139,92,246,0.55)";}

\NormalTok{  // A: ellipse centred at (ax,ay) with radii (arx,ary)}
\NormalTok{  // B: ellipse centred at (bx,by) with radii (brx,bry)}
\NormalTok{  let state=\{ax:0.3,ay:0.5,arx:0.22,ary:0.30,bx:0.62,by:0.5,brx:0.24,bry:0.30\};}
\NormalTok{  let condOnB=false;}
\NormalTok{  let dragging=null, dragStart=null;}
\NormalTok{  let svgRect=null;}

\NormalTok{  function ellipseArea(rx,ry) \{ return Math.PI*rx*ry; \} // in unit{-}square coords}
\NormalTok{  function intersectionArea(state) \{}
\NormalTok{    // Monte Carlo estimate}
\NormalTok{    const \{ax,ay,arx,ary,bx,by,brx,bry\}=state;}
\NormalTok{    let inA=0,inB=0,inBoth=0;}
\NormalTok{    const M=3000;}
\NormalTok{    let s=42;}
\NormalTok{    function lcg()\{s=(s*1664525+1013904223)\&0xffffffff;return(s\textgreater{}\textgreater{}\textgreater{}0)/0xffffffff;\}}
\NormalTok{    for(let i=0;i\textless{}M;i++)\{}
\NormalTok{      const x=lcg(),y=lcg();}
\NormalTok{      const a=((x{-}ax)/arx)**2+((y{-}ay)/ary)**2\textless{}=1;}
\NormalTok{      const b=((x{-}bx)/brx)**2+((y{-}by)/bry)**2\textless{}=1;}
\NormalTok{      if(a)inA++;if(b)inB++;if(a\&\&b)inBoth++;}
\NormalTok{    \}}
\NormalTok{    return \{pA:inA/M, pB:inB/M, pAB:inBoth/M\};}
\NormalTok{  \}}

\NormalTok{  const container=document.createElement("div");}
\NormalTok{  container.style.cssText="font{-}family:inherit; max{-}width:560px; margin:1rem 0;";}

\NormalTok{  const svgEl=document.createElementNS("http://www.w3.org/2000/svg","svg");}
\NormalTok{  svgEl.setAttribute("width",W); svgEl.setAttribute("height",H);}
\NormalTok{  svgEl.style.cssText="border:1px solid \#e5e7eb; border{-}radius:6px; background:\#fff; display:block; cursor:move;";}
\NormalTok{  container.appendChild(svgEl);}

\NormalTok{  // controls}
\NormalTok{  const ctrlRow=document.createElement("div");}
\NormalTok{  ctrlRow.style.cssText="display:flex; gap:0.5rem; align{-}items:center; margin{-}top:0.5rem; flex{-}wrap:wrap;";}
\NormalTok{  const condBtn=document.createElement("button");}
\NormalTok{  condBtn.textContent="Condition on B";}
\NormalTok{  condBtn.style.cssText="padding:0.3rem 0.8rem; border:1px solid \#d1d5db; border{-}radius:4px; cursor:pointer; font{-}size:0.85em; background:\#f9fafb;";}
\NormalTok{  condBtn.onclick=()=\textgreater{}\{condOnB=!condOnB; condBtn.style.background=condOnB?ORANGE\_STROKE:"\#f9fafb"; condBtn.style.color=condOnB?"\#fff":"\#374151"; draw();\};}
\NormalTok{  const resetBtn=document.createElement("button");}
\NormalTok{  resetBtn.textContent="Reset";}
\NormalTok{  resetBtn.style.cssText="padding:0.3rem 0.7rem; border:1px solid \#e5e7eb; border{-}radius:4px; cursor:pointer; font{-}size:0.85em; background:\#f8fafc; color:\#6b7280;";}
\NormalTok{  resetBtn.onclick=()=\textgreater{}\{state=\{ax:0.3,ay:0.5,arx:0.22,ary:0.30,bx:0.62,by:0.5,brx:0.24,bry:0.30\};condOnB=false;condBtn.style.background="\#f9fafb";condBtn.style.color="\#374151";draw();\};}
\NormalTok{  ctrlRow.appendChild(condBtn); ctrlRow.appendChild(resetBtn);}
\NormalTok{  const hint=document.createElement("span");}
\NormalTok{  hint.style.cssText="font{-}size:0.78em; color:\#6b7280;"; hint.textContent="Drag A or B to reposition";}
\NormalTok{  ctrlRow.appendChild(hint);}
\NormalTok{  container.appendChild(ctrlRow);}

\NormalTok{  const labelDiv=document.createElement("div");}
\NormalTok{  labelDiv.style.cssText="margin{-}top:0.35rem; font{-}size:0.82em; color:\#6b7280; font{-}style:italic;";}
\NormalTok{  labelDiv.textContent=\textquotesingle{}Conditioning on B means zooming in to B. P(A|B) is the fraction of B that overlaps A.\textquotesingle{};}
\NormalTok{  container.appendChild(labelDiv);}

\NormalTok{  function toSVG(ux,uy) \{}
\NormalTok{    const sqW=W{-}PAD.left{-}PAD.right, sqH=H{-}PAD.top{-}PAD.bottom;}
\NormalTok{    return [PAD.left+ux*sqW, PAD.top+uy*sqH];}
\NormalTok{  \}}
\NormalTok{  function fromSVG(sx,sy) \{}
\NormalTok{    const sqW=W{-}PAD.left{-}PAD.right, sqH=H{-}PAD.top{-}PAD.bottom;}
\NormalTok{    return [(sx{-}PAD.left)/sqW, (sy{-}PAD.top)/sqH];}
\NormalTok{  \}}

\NormalTok{  function draw() \{}
\NormalTok{    while(svgEl.firstChild) svgEl.removeChild(svgEl.firstChild);}
\NormalTok{    const sqW=W{-}PAD.left{-}PAD.right, sqH=H{-}PAD.top{-}PAD.bottom;}
\NormalTok{    const \{ax,ay,arx,ary,bx,by,brx,bry\}=state;}

\NormalTok{    // background rect (Ω)}
\NormalTok{    const omegaRect=document.createElementNS("http://www.w3.org/2000/svg","rect");}
\NormalTok{    omegaRect.setAttribute("x",PAD.left); omegaRect.setAttribute("y",PAD.top);}
\NormalTok{    omegaRect.setAttribute("width",sqW); omegaRect.setAttribute("height",sqH);}
\NormalTok{    omegaRect.setAttribute("fill", condOnB?"\#f3f4f6":"\#f9fafb");}
\NormalTok{    omegaRect.setAttribute("stroke","\#d1d5db"); omegaRect.setAttribute("rx","4");}
\NormalTok{    svgEl.appendChild(omegaRect);}
\NormalTok{    const omLbl=document.createElementNS("http://www.w3.org/2000/svg","text");}
\NormalTok{    omLbl.setAttribute("x",PAD.left+6); omLbl.setAttribute("y",PAD.top+16);}
\NormalTok{    omLbl.setAttribute("font{-}size","13"); omLbl.setAttribute("fill","\#9ca3af");}
\NormalTok{    omLbl.textContent="Ω"; svgEl.appendChild(omLbl);}

\NormalTok{    const [axSVG,aySVG]=toSVG(ax,ay);}
\NormalTok{    const [bxSVG,bySVG]=toSVG(bx,by);}
\NormalTok{    const arxSVG=arx*sqW, arySVG=ary*sqH;}
\NormalTok{    const brxSVG=brx*sqW, brySVG=bry*sqH;}

\NormalTok{    function makeEllipse(cx,cy,rx,ry,fill,stroke,opacity=1) \{}
\NormalTok{      const el=document.createElementNS("http://www.w3.org/2000/svg","ellipse");}
\NormalTok{      el.setAttribute("cx",cx); el.setAttribute("cy",cy);}
\NormalTok{      el.setAttribute("rx",rx); el.setAttribute("ry",ry);}
\NormalTok{      el.setAttribute("fill",fill); el.setAttribute("stroke",stroke);}
\NormalTok{      el.setAttribute("stroke{-}width","2"); el.setAttribute("opacity",opacity);}
\NormalTok{      return el;}
\NormalTok{    \}}

\NormalTok{    if(!condOnB) \{}
\NormalTok{      svgEl.appendChild(makeEllipse(axSVG,aySVG,arxSVG,arySVG,BLUE,BLUE\_STROKE));}
\NormalTok{      svgEl.appendChild(makeEllipse(bxSVG,bySVG,brxSVG,brySVG,ORANGE,ORANGE\_STROKE));}
\NormalTok{    \} else \{}
\NormalTok{      // grey out everything, show B clearly, show A∩B}
\NormalTok{      svgEl.appendChild(makeEllipse(axSVG,aySVG,arxSVG,arySVG,BLUE,BLUE\_STROKE,0.18));}
\NormalTok{      svgEl.appendChild(makeEllipse(bxSVG,bySVG,brxSVG,brySVG,ORANGE,ORANGE\_STROKE,0.55));}
\NormalTok{      // A∩B region approximated visually as darkened intersection}
\NormalTok{      const overEl=document.createElementNS("http://www.w3.org/2000/svg","ellipse");}
\NormalTok{      // Draw a smaller "intersection" ellipse centred at the midpoint, scaled}
\NormalTok{      const ixc=(axSVG+bxSVG)/2, iyc=(aySVG+bySVG)/2;}
\NormalTok{      const irx=Math.max(2,Math.min(arxSVG,brxSVG)*0.65);}
\NormalTok{      const iry=Math.max(2,Math.min(arySVG,brySVG)*0.65);}
\NormalTok{      const dist=Math.sqrt((axSVG{-}bxSVG)**2+(aySVG{-}bySVG)**2);}
\NormalTok{      if(dist\textless{}(arxSVG+brxSVG)*0.85) \{}
\NormalTok{        const intEl=makeEllipse(ixc,iyc,irx,iry,BLEND,BLUE\_STROKE,0.7);}
\NormalTok{        svgEl.appendChild(intEl);}
\NormalTok{      \}}
\NormalTok{    \}}

\NormalTok{    // Labels A and B}
\NormalTok{    const aLbl=document.createElementNS("http://www.w3.org/2000/svg","text");}
\NormalTok{    aLbl.setAttribute("x",axSVG); aLbl.setAttribute("y",aySVG+4);}
\NormalTok{    aLbl.setAttribute("text{-}anchor","middle"); aLbl.setAttribute("font{-}size","16");}
\NormalTok{    aLbl.setAttribute("font{-}weight","700"); aLbl.setAttribute("fill","\#1d4ed8"); aLbl.textContent="A";}
\NormalTok{    svgEl.appendChild(aLbl);}
\NormalTok{    const bLbl=document.createElementNS("http://www.w3.org/2000/svg","text");}
\NormalTok{    bLbl.setAttribute("x",bxSVG); bLbl.setAttribute("y",bySVG+4);}
\NormalTok{    bLbl.setAttribute("text{-}anchor","middle"); bLbl.setAttribute("font{-}size","16");}
\NormalTok{    bLbl.setAttribute("font{-}weight","700"); bLbl.setAttribute("fill","\#c2410c"); bLbl.textContent="B";}
\NormalTok{    svgEl.appendChild(bLbl);}

\NormalTok{    // Readout panel (right side)}
\NormalTok{    const panelX=W{-}PAD.right+14, panelY=PAD.top;}
\NormalTok{    const \{pA,pB,pAB\}=intersectionArea(state);}
\NormalTok{    const pAgivenB=pB\textgreater{}0.001?pAB/pB:0;}
\NormalTok{    const indep=Math.abs(pAB{-}pA*pB)\textless{}0.02;}
\NormalTok{    const lines=[}
\NormalTok{      \{label:"P(A):",val:pA.toFixed(2),color:"\#3b82f6"\},}
\NormalTok{      \{label:"P(B):",val:pB.toFixed(2),color:"\#f97316"\},}
\NormalTok{      \{label:"P(A∩B):",val:pAB.toFixed(2),color:"\#7c3aed"\},}
\NormalTok{      \{label:"P(A|B):",val:pAgivenB.toFixed(2),color:"\#374151",big:true\},}
\NormalTok{    ];}
\NormalTok{    lines.forEach((\{label,val,color,big\},i)=\textgreater{}\{}
\NormalTok{      const lbl=document.createElementNS("http://www.w3.org/2000/svg","text");}
\NormalTok{      lbl.setAttribute("x",panelX); lbl.setAttribute("y",panelY+25+i*26);}
\NormalTok{      lbl.setAttribute("font{-}size",big?"13":"12"); lbl.setAttribute("fill",color);}
\NormalTok{      lbl.setAttribute("font{-}weight",big?"700":"400");}
\NormalTok{      lbl.textContent=\textasciigrave{}$\{label\} $\{val\}\textasciigrave{};}
\NormalTok{      svgEl.appendChild(lbl);}
\NormalTok{    \});}
\NormalTok{    if(indep) \{}
\NormalTok{      const indLbl=document.createElementNS("http://www.w3.org/2000/svg","text");}
\NormalTok{      indLbl.setAttribute("x",panelX); indLbl.setAttribute("y",panelY+135);}
\NormalTok{      indLbl.setAttribute("font{-}size","10"); indLbl.setAttribute("fill","\#059669");}
\NormalTok{      indLbl.textContent="✓ ≈ Independent";}
\NormalTok{      svgEl.appendChild(indLbl);}
\NormalTok{    \}}
\NormalTok{    if(condOnB) \{}
\NormalTok{      const condNote=document.createElementNS("http://www.w3.org/2000/svg","text");}
\NormalTok{      condNote.setAttribute("x",panelX); condNote.setAttribute("y",panelY+160);}
\NormalTok{      condNote.setAttribute("font{-}size","9"); condNote.setAttribute("fill","\#6b7280");}
\NormalTok{      condNote.textContent="Inside B only.";}
\NormalTok{      svgEl.appendChild(condNote);}
\NormalTok{    \}}
\NormalTok{  \}}

\NormalTok{  // drag handling}
\NormalTok{  function getPos(e) \{}
\NormalTok{    svgRect=svgEl.getBoundingClientRect();}
\NormalTok{    const cx=e.clientX{-}svgRect.left, cy=e.clientY{-}svgRect.top;}
\NormalTok{    return fromSVG(cx,cy);}
\NormalTok{  \}}
\NormalTok{  svgEl.addEventListener("mousedown",e=\textgreater{}\{}
\NormalTok{    const [ux,uy]=getPos(e);}
\NormalTok{    const \{ax,ay,arx,ary,bx,by,brx,bry\}=state;}
\NormalTok{    const inA=((ux{-}ax)/arx)**2+((uy{-}ay)/ary)**2\textless{}=1;}
\NormalTok{    const inB=((ux{-}bx)/brx)**2+((uy{-}by)/bry)**2\textless{}=1;}
\NormalTok{    if(inB) dragging="B"; else if(inA) dragging="A";}
\NormalTok{    dragStart=[ux,uy];}
\NormalTok{  \});}
\NormalTok{  window.addEventListener("mousemove",e=\textgreater{}\{}
\NormalTok{    if(!dragging) return;}
\NormalTok{    const [ux,uy]=getPos(e);}
\NormalTok{    const [dx,dy]=[ux{-}dragStart[0],uy{-}dragStart[1]];}
\NormalTok{    dragStart=[ux,uy];}
\NormalTok{    if(dragging==="A") \{ state.ax=Math.max(0.1,Math.min(0.9,state.ax+dx)); state.ay=Math.max(0.1,Math.min(0.9,state.ay+dy)); \}}
\NormalTok{    if(dragging==="B") \{ state.bx=Math.max(0.1,Math.min(0.9,state.bx+dx)); state.by=Math.max(0.1,Math.min(0.9,state.by+dy)); \}}
\NormalTok{    draw();}
\NormalTok{  \});}
\NormalTok{  window.addEventListener("mouseup",()=\textgreater{}\{ dragging=null; \});}

\NormalTok{  draw();}
\NormalTok{  return container;}
\NormalTok{\}}
\end{Highlighting}
\end{Shaded}

Two events \(A\) and \(B\) are \textbf{independent} if knowing \(B\)
occurred tells you nothing about \(A\):

\[P(A \mid B) = P(A) \iff P(A \cap B) = P(A) \cdot P(B)\]

\subsection{Bayes' theorem}\label{bayes-theorem-1}

Rearranging the definition of conditional probability:

\[P(A \cap B) = P(A \mid B) \cdot P(B) = P(B \mid A) \cdot P(A)\]

Solving for \(P(A \mid B)\):

\[P(A \mid B) = \frac{P(B \mid A) \cdot P(A)}{P(B)}\]

This is \textbf{Bayes' theorem}. It lets you reverse the conditioning:
if you know \(P(B \mid A)\), you can compute \(P(A \mid B)\).

The denominator \(P(B)\) is expanded using the \textbf{total probability
rule}. If \(A_1, A_2, \ldots, A_n\) partition \(\Omega\) (exhaustive,
mutually exclusive), then:

\[P(B) = \sum_{i=1}^{n} P(B \mid A_i) \cdot P(A_i)\]

\begin{tcolorbox}[enhanced jigsaw, opacityback=0, colback=white, breakable, colframe=quarto-callout-note-color-frame, arc=.35mm, left=2mm, toprule=.15mm, leftrule=.75mm, bottomrule=.15mm, rightrule=.15mm]
\begin{minipage}[t]{5.5mm}
\textcolor{quarto-callout-note-color}{\faInfo}
\end{minipage}%
\begin{minipage}[t]{\textwidth - 5.5mm}

\textbf{Why this works}

Bayes' theorem doesn't involve any new mathematics --- it's just two
ways of writing the same joint probability \(P(A \cap B)\), set equal.
What makes it powerful is the direction: you have the likelihood
\(P(B \mid A)\) (how probable is the evidence given the hypothesis?) and
you want \(P(A \mid B)\) (how probable is the hypothesis given the
evidence?). Bayes is the bridge between them.

\end{minipage}%
\end{tcolorbox}

\textbf{Worked example: Medical test.} A disease affects 1\% of the
population. A test for it has sensitivity 99\% (correctly identifies
99\% of sick patients) and specificity 95\% (correctly clears 95\% of
healthy patients). A randomly selected person tests positive. What is
the probability they actually have the disease?

Define: - \(D\): the person has the disease - \(T^+\): the test is
positive

Given: \(P(D) = 0.01\), \(P(T^+ \mid D) = 0.99\) (sensitivity),
\(P(T^+ \mid D^c) = 0.05\)

Note also that \(P(D^c) = 1 - P(D) = 1 - 0.01 = 0.99\) --- the
complement of the prevalence, which happens to equal the sensitivity by
coincidence.

First compute the total probability of a positive test, using the two
ways a positive result can occur:

\[P(T^+) = P(T^+ \mid D) \cdot P(D) + P(T^+ \mid D^c) \cdot P(D^c)\]
\[= \underbrace{0.99}_{\text{sensitivity}} \times 0.01
  + 0.05 \times \underbrace{0.99}_{P(D^c) = 1 - 0.01}
  = 0.0099 + 0.0495 = 0.0594\]

Now apply Bayes:

\[P(D \mid T^+) = \frac{P(T^+ \mid D) \cdot P(D)}{P(T^+)} = \frac{0.99 \times 0.01}{0.0594} \approx 0.167\]

Even with a positive result from a 99\%-sensitive test, the probability
of actually having the disease is only about 17\%. This is not a failure
of the test --- it's a consequence of the disease being rare. Most of
the positive tests come from the large healthy population, not the small
sick one. The prior probability \(P(D)\) matters enormously.

\section{Random variables}\label{random-variables-1}

A \textbf{random variable} \(X\) is a function from the sample space
\(\Omega\) to the real numbers --- it assigns a numerical value to each
outcome. The randomness is in the outcome; the function is
deterministic.

\textbf{Discrete random variables} take values in a countable set
\(\{x_1, x_2, \ldots\}\). The distribution is completely described by
the \textbf{probability mass function} (PMF):

\[p(x_i) = P(X = x_i), \quad \sum_i p(x_i) = 1\]

\textbf{Continuous random variables} take values in an interval (or
union of intervals). The distribution is described by the
\textbf{probability density function} (PDF) \(f(x)\), where:

\[P(a \leq X \leq b) = \int_a^b f(x)\, dx, \quad \int_{-\infty}^{\infty} f(x)\, dx = 1\]

Note that \(f(x)\) is \emph{not} a probability --- it can exceed 1. It's
a density: the probability of \(X\) falling in a small interval
\([x, x + dx]\) is approximately \(f(x)\, dx\).

A concrete example: if \(X \sim \text{Uniform}(0,\, 0.5)\), then
\(f(x) = 2\) everywhere on \([0, 0.5]\). The density is 2, yet
\(P(0 \leq X \leq 0.5) = \int_0^{0.5} 2\, dx = 1\). The density can
exceed 1; the integral over any interval is still between 0 and 1.

The \textbf{cumulative distribution function} (CDF) is defined for both
types:

\[F(x) = P(X \leq x)\]

For continuous \(X\): \(F(x) = \int_{-\infty}^{x} f(t)\, dt\), and
\(f(x) = F'(x)\).

\subsection{Expectation}\label{expectation}

The \textbf{expected value} \(E[X]\) --- also written \(\mu\) (mu) ---
is the long-run average of \(X\) over many repetitions:

\textbf{Discrete:} \(\displaystyle E[X] = \sum_i x_i \, p(x_i)\)

\textbf{Continuous:}
\(\displaystyle E[X] = \int_{-\infty}^{\infty} x \, f(x)\, dx\)

Expectation is linear: \(E[aX + b] = a\,E[X] + b\) for constants \(a\),
\(b\).

For a function \(g(X)\):

\[E[g(X)] = \sum_i g(x_i)\, p(x_i) \quad \text{(discrete)}\]

\[E[g(X)] = \int_{-\infty}^{\infty} g(x)\, f(x)\, dx \quad \text{(continuous)}\]

\subsection{Variance}\label{variance}

The \textbf{variance} \(\text{Var}(X)\) --- also written \(\sigma^2\)
(sigma squared) --- measures the average squared deviation from the
mean:

\[\text{Var}(X) = E\!\left[(X - \mu)^2\right] = E[X^2] - (E[X])^2\]

The second form is usually easier to compute. The derivation:

\[E\!\left[(X - \mu)^2\right] = E[X^2 - 2\mu X + \mu^2]
= E[X^2] - 2\mu\,E[X] + \mu^2\]

Since \(E[X] = \mu\), the last two terms combine:
\(-2\mu \cdot \mu + \mu^2 = -2\mu^2 + \mu^2 = -\mu^2\). So:

\[= E[X^2] - \mu^2\]

The \textbf{standard deviation} \(\sigma = \sqrt{\text{Var}(X)}\) is in
the same units as \(X\), which makes it interpretable.

Variance scales with the square:
\(\text{Var}(aX + b) = a^2\,\text{Var}(X)\).

For independent random variables \(X\) and \(Y\):
\(\text{Var}(X + Y) = \text{Var}(X) + \text{Var}(Y)\).

\section{Key discrete distributions}\label{key-discrete-distributions}

\subsection{Bernoulli}\label{bernoulli}

A single trial with probability \(p\) of success, \(1-p\) of failure.

\[P(X = 1) = p, \quad P(X = 0) = 1 - p\]

\[E[X] = p, \quad \text{Var}(X) = p(1-p)\]

The Bernoulli is the building block for everything that follows.

\subsection{Binomial}\label{binomial}

\(n\) independent Bernoulli trials, each with success probability \(p\).
\(X\) counts the total number of successes.

\[P(X = k) = \binom{n}{k} p^k (1-p)^{n-k}, \quad k = 0, 1, \ldots, n\]

The binomial coefficient \(\binom{n}{k}\) --- read ``n choose k'' ---
counts the number of ways to arrange \(k\) successes among \(n\) trials:

\[\binom{n}{k} = \frac{n!}{k!(n-k)!}\]

\textbf{Mean and variance.} Since \(X\) is the sum of \(n\) independent
Bernoulli trials, linearity of expectation and variance additivity give:

\[E[X] = np, \quad \text{Var}(X) = np(1-p)\]

\textbf{Normal approximation.} When \(np \geq 5\) and \(n(1-p) \geq 5\),
the binomial is well approximated by a normal distribution with the same
mean and variance. This is one preview of the CLT.

\subsection{Poisson}\label{poisson}

Models the number of events occurring in a fixed interval of time or
space, when events occur independently at a constant average rate.

\[P(X = k) = \frac{\lambda^k e^{-\lambda}}{k!}, \quad k = 0, 1, 2, \ldots\]

The parameter \(\lambda > 0\) --- read ``lambda'' --- is both the mean
and the variance:

\[E[X] = \lambda, \quad \text{Var}(X) = \lambda\]

\textbf{As a limit of Binomial.} Set \(\lambda = np\) and let
\(n \to \infty\), \(p \to 0\). The binomial PMF converges to the
Poisson. This is why the Poisson appears when events are rare but trials
are many: the number of typing errors per page, radioactive decays per
second, server requests per minute.

\textbf{Derivation of the mean.} Using the Poisson PMF:

\[E[X] = \sum_{k=0}^{\infty} k \cdot \frac{\lambda^k e^{-\lambda}}{k!}
= e^{-\lambda} \sum_{k=1}^{\infty} \frac{\lambda^k}{(k-1)!}
= e^{-\lambda} \cdot \lambda \sum_{j=0}^{\infty} \frac{\lambda^j}{j!}
= e^{-\lambda} \cdot \lambda \cdot e^{\lambda} = \lambda\]

\begin{Shaded}
\begin{Highlighting}[]
\NormalTok{//| echo: false}
\NormalTok{\{}
\NormalTok{  // ── Distribution landscape explorer ──────────────────────────────────────}
\NormalTok{  const W = 500, H = 280, PAD = \{ top: 24, right: 20, bottom: 40, left: 48 \};}
\NormalTok{  const ACCENT = "\#3b82f6";}
\NormalTok{  const ACCENT\_MUTED = "rgba(59,130,246,0.18)";}

\NormalTok{  // ── state ─────────────────────────────────────────────────────────────────}
\NormalTok{  const distNames = ["Bernoulli", "Binomial", "Poisson", "Exponential", "Normal"];}
\NormalTok{  let distIdx = 4;     // default: Normal}
\NormalTok{  let params = \{ p: 0.5, n: 10, lam: 2, mu: 0, sigma: 1 \};}

\NormalTok{  // ── helpers ──────────────────────────────────────────────────────────────}
\NormalTok{  function factorial(k) \{ let r=1; for(let i=2;i\textless{}=k;i++) r*=i; return r; \}}
\NormalTok{  function normalPDF(x,mu,sig) \{ return Math.exp({-}0.5*((x{-}mu)/sig)**2)/(sig*Math.sqrt(2*Math.PI)); \}}
\NormalTok{  function normalCDF(x,mu,sig) \{}
\NormalTok{    const z=(x{-}mu)/(sig*Math.SQRT2);}
\NormalTok{    return 0.5*(1+erf(z));}
\NormalTok{  \}}
\NormalTok{  function erf(x) \{}
\NormalTok{    const sign = x\textless{}0?{-}1:1; x=Math.abs(x);}
\NormalTok{    const a1=0.254829592,a2={-}0.284496736,a3=1.421413741,a4={-}1.453152027,a5=1.061405429,p=0.3275911;}
\NormalTok{    const t=1/(1+p*x);}
\NormalTok{    const y=1{-}(((((a5*t+a4)*t)+a3)*t+a2)*t+a1)*t*Math.exp({-}x*x);}
\NormalTok{    return sign*y;}
\NormalTok{  \}}
\NormalTok{  function poissonPMF(k,lam) \{ return Math.pow(lam,k)*Math.exp({-}lam)/factorial(k); \}}
\NormalTok{  function binomPMF(k,n,p) \{}
\NormalTok{    if(k\textless{}0||k\textgreater{}n) return 0;}
\NormalTok{    const logC = logBinom(n,k);}
\NormalTok{    return Math.exp(logC + k*Math.log(p) + (n{-}k)*Math.log(1{-}p));}
\NormalTok{  \}}
\NormalTok{  function logBinom(n,k) \{}
\NormalTok{    let r=0; for(let i=0;i\textless{}k;i++) r+=Math.log(n{-}i){-}Math.log(i+1); return r;}
\NormalTok{  \}}

\NormalTok{  function getDistInfo() \{}
\NormalTok{    const d = distNames[distIdx];}
\NormalTok{    if(d==="Bernoulli") \{}
\NormalTok{      const p=params.p; const mu=p; const variance=p*(1{-}p); const sigma=Math.sqrt(variance);}
\NormalTok{      const xs=[0,1]; const ys=[1{-}p,p];}
\NormalTok{      return \{ discrete:true, xs, ys, mu, variance, sigma, label:\textasciigrave{}Bernoulli(p=$\{p.toFixed(2)\})\textasciigrave{} \};}
\NormalTok{    \}}
\NormalTok{    if(d==="Binomial") \{}
\NormalTok{      const \{n,p\}=params; const mu=n*p; const variance=n*p*(1{-}p); const sigma=Math.sqrt(variance);}
\NormalTok{      const xs=[...Array(n+1).keys()]; const ys=xs.map(k=\textgreater{}binomPMF(k,n,p));}
\NormalTok{      return \{ discrete:true, xs, ys, mu, variance, sigma, label:\textasciigrave{}Binomial(n=$\{n\}, p=$\{p.toFixed(2)\})\textasciigrave{} \};}
\NormalTok{    \}}
\NormalTok{    if(d==="Poisson") \{}
\NormalTok{      const lam=params.lam; const mu=lam; const variance=lam; const sigma=Math.sqrt(lam);}
\NormalTok{      const maxK=Math.ceil(lam+4*sigma+5);}
\NormalTok{      const xs=[...Array(maxK+1).keys()]; const ys=xs.map(k=\textgreater{}poissonPMF(k,lam));}
\NormalTok{      return \{ discrete:true, xs, ys, mu, variance, sigma, label:\textasciigrave{}Poisson(λ=$\{lam.toFixed(1)\})\textasciigrave{} \};}
\NormalTok{    \}}
\NormalTok{    if(d==="Exponential") \{}
\NormalTok{      const lam=params.lam; const mu=1/lam; const variance=1/(lam*lam); const sigma=1/lam;}
\NormalTok{      const xmax=mu+5*sigma; const nPts=200;}
\NormalTok{      const xs=Array.from(\{length:nPts\},(\_,i)=\textgreater{}i*xmax/(nPts{-}1));}
\NormalTok{      const ys=xs.map(x=\textgreater{}x\textgreater{}=0?lam*Math.exp({-}lam*x):0);}
\NormalTok{      return \{ discrete:false, xs, ys, mu, variance, sigma, label:\textasciigrave{}Exponential(λ=$\{lam.toFixed(2)\})\textasciigrave{}, xmax \};}
\NormalTok{    \}}
\NormalTok{    // Normal}
\NormalTok{    const \{mu,sigma\}=params; const variance=sigma*sigma;}
\NormalTok{    const xmin=mu{-}4*sigma, xmax=mu+4*sigma; const nPts=200;}
\NormalTok{    const xs=Array.from(\{length:nPts\},(\_,i)=\textgreater{}xmin+i*(xmax{-}xmin)/(nPts{-}1));}
\NormalTok{    const ys=xs.map(x=\textgreater{}normalPDF(x,mu,sigma));}
\NormalTok{    return \{ discrete:false, xs, ys, mu, variance, sigma, label:\textasciigrave{}Normal(μ=$\{mu.toFixed(1)\}, σ=$\{sigma.toFixed(1)\})\textasciigrave{}, xmin, xmax \};}
\NormalTok{  \}}

\NormalTok{  // ── build UI ───────────────────────────────────────────────────────────────}
\NormalTok{  const container = document.createElement("div");}
\NormalTok{  container.style.cssText = "font{-}family:inherit; max{-}width:680px; margin:1rem 0;";}

\NormalTok{  // dist selector}
\NormalTok{  const selectorRow = document.createElement("div");}
\NormalTok{  selectorRow.style.cssText = "display:flex; gap:0.4rem; flex{-}wrap:wrap; margin{-}bottom:0.75rem;";}
\NormalTok{  const btns = distNames.map((name,i)=\textgreater{}\{}
\NormalTok{    const b=document.createElement("button");}
\NormalTok{    b.textContent=name; b.dataset.idx=i;}
\NormalTok{    b.style.cssText = "padding:0.3rem 0.7rem; border{-}radius:4px; border:1px solid \#d1d5db; cursor:pointer; font{-}size:0.85em; background:\#f9fafb;";}
\NormalTok{    b.onclick=()=\textgreater{}\{ distIdx=i; update(); \};}
\NormalTok{    selectorRow.appendChild(b);}
\NormalTok{    return b;}
\NormalTok{  \});}
\NormalTok{  container.appendChild(selectorRow);}

\NormalTok{  // main layout: left (canvas) + right (readouts)}
\NormalTok{  const mainRow = document.createElement("div");}
\NormalTok{  mainRow.style.cssText = "display:flex; gap:1rem; align{-}items:flex{-}start;";}

\NormalTok{  // svg}
\NormalTok{  const svg = document.createElementNS("http://www.w3.org/2000/svg","svg");}
\NormalTok{  svg.setAttribute("width", W); svg.setAttribute("height", H);}
\NormalTok{  svg.style.cssText = "border:1px solid \#e5e7eb; border{-}radius:6px; background:\#fff; flex{-}shrink:0;";}
\NormalTok{  mainRow.appendChild(svg);}

\NormalTok{  // readout panel}
\NormalTok{  const readouts = document.createElement("div");}
\NormalTok{  readouts.style.cssText = "min{-}width:170px; padding:0.75rem; background:\#f8fafc; border:1px solid \#e5e7eb; border{-}radius:6px; font{-}size:0.88em;";}
\NormalTok{  mainRow.appendChild(readouts);}
\NormalTok{  container.appendChild(mainRow);}

\NormalTok{  // sliders area}
\NormalTok{  const slidersDiv = document.createElement("div");}
\NormalTok{  slidersDiv.style.cssText = "margin{-}top:0.75rem; padding:0.75rem; background:\#f8fafc; border:1px solid \#e5e7eb; border{-}radius:6px;";}
\NormalTok{  container.appendChild(slidersDiv);}

\NormalTok{  // label}
\NormalTok{  const labelDiv = document.createElement("div");}
\NormalTok{  labelDiv.style.cssText = "margin{-}top:0.5rem; font{-}size:0.82em; color:\#6b7280; font{-}style:italic;";}
\NormalTok{  labelDiv.textContent = "Every distribution is a landscape: a shape, a centre, and a spread.";}
\NormalTok{  container.appendChild(labelDiv);}

\NormalTok{  function makeSlider(labelText, key, min, max, step, getValue, setValue) \{}
\NormalTok{    const row = document.createElement("div");}
\NormalTok{    row.style.cssText = "display:flex; align{-}items:center; gap:0.5rem; margin{-}bottom:0.4rem;";}
\NormalTok{    const lbl = document.createElement("label");}
\NormalTok{    lbl.style.cssText = "font{-}size:0.85em; min{-}width:80px;";}
\NormalTok{    lbl.textContent = labelText;}
\NormalTok{    const inp = document.createElement("input");}
\NormalTok{    inp.type="range"; inp.min=min; inp.max=max; inp.step=step; inp.value=getValue();}
\NormalTok{    inp.style.cssText = "flex:1;";}
\NormalTok{    const val = document.createElement("span");}
\NormalTok{    val.style.cssText = "font{-}size:0.85em; min{-}width:35px; text{-}align:right;";}
\NormalTok{    val.textContent = Number(getValue()).toFixed(step\textless{}0.1?2:step\textless{}1?1:0);}
\NormalTok{    inp.oninput = ()=\textgreater{}\{ setValue(parseFloat(inp.value)); val.textContent=Number(inp.value).toFixed(step\textless{}0.1?2:step\textless{}1?1:0); update(); \};}
\NormalTok{    row.appendChild(lbl); row.appendChild(inp); row.appendChild(val);}
\NormalTok{    return \{row, inp, val\};}
\NormalTok{  \}}

\NormalTok{  // ── draw ──────────────────────────────────────────────────────────────────}
\NormalTok{  function draw(info) \{}
\NormalTok{    while(svg.firstChild) svg.removeChild(svg.firstChild);}
\NormalTok{    const cw = W {-} PAD.left {-} PAD.right;}
\NormalTok{    const ch = H {-} PAD.top {-} PAD.bottom;}

\NormalTok{    const xMin = info.discrete ? Math.min(...info.xs){-}0.5 : (info.xmin !== undefined ? info.xmin : Math.min(...info.xs));}
\NormalTok{    const xMax = info.discrete ? Math.max(...info.xs)+0.5 : (info.xmax !== undefined ? info.xmax : Math.max(...info.xs));}
\NormalTok{    const yMax = Math.max(...info.ys)*1.15;}

\NormalTok{    const sx = x =\textgreater{} PAD.left + (x{-}xMin)/(xMax{-}xMin)*cw;}
\NormalTok{    const sy = y =\textgreater{} PAD.top + ch {-} y/yMax*ch;}

\NormalTok{    // grid lines (light)}
\NormalTok{    for(let i=0;i\textless{}=4;i++) \{}
\NormalTok{      const y = i*yMax/4;}
\NormalTok{      const line = document.createElementNS("http://www.w3.org/2000/svg","line");}
\NormalTok{      line.setAttribute("x1",PAD.left); line.setAttribute("x2",PAD.left+cw);}
\NormalTok{      line.setAttribute("y1",sy(y)); line.setAttribute("y2",sy(y));}
\NormalTok{      line.setAttribute("stroke","\#f0f0f0"); line.setAttribute("stroke{-}width","1");}
\NormalTok{      svg.appendChild(line);}
\NormalTok{    \}}

\NormalTok{    // ±1σ band shading (continuous only)}
\NormalTok{    if(!info.discrete) \{}
\NormalTok{      const \{mu,sigma,xs,ys\}=info;}
\NormalTok{      const x1c=Math.max(xMin,mu{-}sigma), x2c=Math.min(xMax,mu+sigma);}
\NormalTok{      // build polygon}
\NormalTok{      const bandXs=xs.filter(x=\textgreater{}x\textgreater{}=x1c\&\&x\textless{}=x2c);}
\NormalTok{      if(bandXs.length\textgreater{}1) \{}
\NormalTok{        const pts=[\textasciigrave{}$\{sx(x1c)\},$\{sy(0)\}\textasciigrave{},...bandXs.map((x,j)=\textgreater{}\{}
\NormalTok{          const y=ys[xs.indexOf(x)]; return \textasciigrave{}$\{sx(x)\},$\{sy(y)\}\textasciigrave{};}
\NormalTok{        \}),\textasciigrave{}$\{sx(x2c)\},$\{sy(0)\}\textasciigrave{}].join(" ");}
\NormalTok{        const poly=document.createElementNS("http://www.w3.org/2000/svg","polygon");}
\NormalTok{        poly.setAttribute("points",pts); poly.setAttribute("fill",ACCENT\_MUTED);}
\NormalTok{        svg.appendChild(poly);}
\NormalTok{      \}}
\NormalTok{    \}}

\NormalTok{    // distribution shape}
\NormalTok{    if(info.discrete) \{}
\NormalTok{      info.xs.forEach((x,i)=\textgreater{}\{}
\NormalTok{        const y=info.ys[i];}
\NormalTok{        const line=document.createElementNS("http://www.w3.org/2000/svg","line");}
\NormalTok{        line.setAttribute("x1",sx(x)); line.setAttribute("x2",sx(x));}
\NormalTok{        line.setAttribute("y1",sy(0)); line.setAttribute("y2",sy(y));}
\NormalTok{        line.setAttribute("stroke",ACCENT); line.setAttribute("stroke{-}width","2");}
\NormalTok{        svg.appendChild(line);}
\NormalTok{        const dot=document.createElementNS("http://www.w3.org/2000/svg","circle");}
\NormalTok{        dot.setAttribute("cx",sx(x)); dot.setAttribute("cy",sy(y)); dot.setAttribute("r","4");}
\NormalTok{        dot.setAttribute("fill",ACCENT);}
\NormalTok{        svg.appendChild(dot);}
\NormalTok{      \});}
\NormalTok{    \} else \{}
\NormalTok{      const pts=info.xs.map((x,i)=\textgreater{}\textasciigrave{}$\{sx(x)\},$\{sy(info.ys[i])\}\textasciigrave{}).join(" ");}
\NormalTok{      const area=document.createElementNS("http://www.w3.org/2000/svg","polygon");}
\NormalTok{      area.setAttribute("points",\textasciigrave{}$\{sx(info.xs[0])\},$\{sy(0)\} $\{pts\} $\{sx(info.xs[info.xs.length{-}1])\},$\{sy(0)\}\textasciigrave{});}
\NormalTok{      area.setAttribute("fill",ACCENT); area.setAttribute("fill{-}opacity","0.35");}
\NormalTok{      svg.appendChild(area);}
\NormalTok{      const path=document.createElementNS("http://www.w3.org/2000/svg","polyline");}
\NormalTok{      path.setAttribute("points",pts); path.setAttribute("fill","none");}
\NormalTok{      path.setAttribute("stroke",ACCENT); path.setAttribute("stroke{-}width","2");}
\NormalTok{      svg.appendChild(path);}
\NormalTok{    \}}

\NormalTok{    // mean line}
\NormalTok{    const muX=sx(info.mu);}
\NormalTok{    const muLine=document.createElementNS("http://www.w3.org/2000/svg","line");}
\NormalTok{    muLine.setAttribute("x1",muX); muLine.setAttribute("x2",muX);}
\NormalTok{    muLine.setAttribute("y1",PAD.top); muLine.setAttribute("y2",PAD.top+ch);}
\NormalTok{    muLine.setAttribute("stroke","\#374151"); muLine.setAttribute("stroke{-}width","1.5");}
\NormalTok{    muLine.setAttribute("stroke{-}dasharray","5,3");}
\NormalTok{    svg.appendChild(muLine);}
\NormalTok{    const muLbl=document.createElementNS("http://www.w3.org/2000/svg","text");}
\NormalTok{    muLbl.setAttribute("x",muX+4); muLbl.setAttribute("y",PAD.top+10);}
\NormalTok{    muLbl.setAttribute("font{-}size","11"); muLbl.setAttribute("fill","\#374151");}
\NormalTok{    muLbl.textContent=\textasciigrave{}μ=$\{info.mu.toFixed(2)\}\textasciigrave{};}
\NormalTok{    svg.appendChild(muLbl);}

\NormalTok{    // ±1σ bracket}
\NormalTok{    const s1x1=sx(info.mu{-}info.sigma), s1x2=sx(info.mu+info.sigma);}
\NormalTok{    const bracketY=PAD.top+ch+14;}
\NormalTok{    const bline=document.createElementNS("http://www.w3.org/2000/svg","line");}
\NormalTok{    bline.setAttribute("x1",Math.max(PAD.left,s1x1)); bline.setAttribute("x2",Math.min(PAD.left+cw,s1x2));}
\NormalTok{    bline.setAttribute("y1",bracketY); bline.setAttribute("y2",bracketY);}
\NormalTok{    bline.setAttribute("stroke",ACCENT); bline.setAttribute("stroke{-}width","2");}
\NormalTok{    svg.appendChild(bline);}
\NormalTok{    const blbl=document.createElementNS("http://www.w3.org/2000/svg","text");}
\NormalTok{    blbl.setAttribute("x",(Math.max(PAD.left,s1x1)+Math.min(PAD.left+cw,s1x2))/2);}
\NormalTok{    blbl.setAttribute("y",bracketY+12); blbl.setAttribute("text{-}anchor","middle");}
\NormalTok{    blbl.setAttribute("font{-}size","10"); blbl.setAttribute("fill",ACCENT);}
\NormalTok{    blbl.textContent="±1σ";}
\NormalTok{    svg.appendChild(blbl);}

\NormalTok{    // axes}
\NormalTok{    const xAxis=document.createElementNS("http://www.w3.org/2000/svg","line");}
\NormalTok{    xAxis.setAttribute("x1",PAD.left); xAxis.setAttribute("x2",PAD.left+cw);}
\NormalTok{    xAxis.setAttribute("y1",sy(0)); xAxis.setAttribute("y2",sy(0));}
\NormalTok{    xAxis.setAttribute("stroke","\#374151"); xAxis.setAttribute("stroke{-}width","1");}
\NormalTok{    svg.appendChild(xAxis);}
\NormalTok{    const yAxis=document.createElementNS("http://www.w3.org/2000/svg","line");}
\NormalTok{    yAxis.setAttribute("x1",PAD.left); yAxis.setAttribute("x2",PAD.left);}
\NormalTok{    yAxis.setAttribute("y1",PAD.top); yAxis.setAttribute("y2",sy(0));}
\NormalTok{    yAxis.setAttribute("stroke","\#374151"); yAxis.setAttribute("stroke{-}width","1");}
\NormalTok{    svg.appendChild(yAxis);}

\NormalTok{    // y{-}axis label}
\NormalTok{    const yLbl=document.createElementNS("http://www.w3.org/2000/svg","text");}
\NormalTok{    yLbl.setAttribute("x",10); yLbl.setAttribute("y",PAD.top+ch/2);}
\NormalTok{    yLbl.setAttribute("font{-}size","11"); yLbl.setAttribute("fill","\#6b7280");}
\NormalTok{    yLbl.setAttribute("text{-}anchor","middle");}
\NormalTok{    yLbl.setAttribute("transform",\textasciigrave{}rotate({-}90,10,$\{PAD.top+ch/2\})\textasciigrave{});}
\NormalTok{    yLbl.textContent = info.discrete ? "Probability" : "Density";}
\NormalTok{    svg.appendChild(yLbl);}

\NormalTok{    // x{-}axis ticks}
\NormalTok{    const nTicks=5;}
\NormalTok{    for(let i=0;i\textless{}=nTicks;i++) \{}
\NormalTok{      const xv=xMin+i*(xMax{-}xMin)/nTicks;}
\NormalTok{      const tick=document.createElementNS("http://www.w3.org/2000/svg","line");}
\NormalTok{      tick.setAttribute("x1",sx(xv)); tick.setAttribute("x2",sx(xv));}
\NormalTok{      tick.setAttribute("y1",sy(0)); tick.setAttribute("y2",sy(0)+5);}
\NormalTok{      tick.setAttribute("stroke","\#374151"); tick.setAttribute("stroke{-}width","1");}
\NormalTok{      svg.appendChild(tick);}
\NormalTok{      const lbl=document.createElementNS("http://www.w3.org/2000/svg","text");}
\NormalTok{      lbl.setAttribute("x",sx(xv)); lbl.setAttribute("y",sy(0)+18);}
\NormalTok{      lbl.setAttribute("text{-}anchor","middle"); lbl.setAttribute("font{-}size","10");}
\NormalTok{      lbl.setAttribute("fill","\#6b7280");}
\NormalTok{      lbl.textContent = info.discrete ? Math.round(xv) : xv.toFixed(1);}
\NormalTok{      svg.appendChild(lbl);}
\NormalTok{    \}}

\NormalTok{    // distribution label top{-}left}
\NormalTok{    const titleLbl=document.createElementNS("http://www.w3.org/2000/svg","text");}
\NormalTok{    titleLbl.setAttribute("x",PAD.left+4); titleLbl.setAttribute("y",PAD.top{-}6);}
\NormalTok{    titleLbl.setAttribute("font{-}size","11"); titleLbl.setAttribute("fill","\#374151");}
\NormalTok{    titleLbl.setAttribute("font{-}weight","500");}
\NormalTok{    titleLbl.textContent=info.label;}
\NormalTok{    svg.appendChild(titleLbl);}
\NormalTok{  \}}

\NormalTok{  function buildSliders() \{}
\NormalTok{    slidersDiv.innerHTML="";}
\NormalTok{    const d=distNames[distIdx];}
\NormalTok{    const title=document.createElement("div");}
\NormalTok{    title.style.cssText="font{-}size:0.82em; font{-}weight:600; color:\#374151; margin{-}bottom:0.4rem;";}
\NormalTok{    title.textContent="Parameters";}
\NormalTok{    slidersDiv.appendChild(title);}
\NormalTok{    if(d==="Bernoulli") \{}
\NormalTok{      const \{row\}=makeSlider("p",null,0.01,0.99,0.01,()=\textgreater{}params.p,v=\textgreater{}params.p=v);}
\NormalTok{      slidersDiv.appendChild(row);}
\NormalTok{    \}}
\NormalTok{    if(d==="Binomial") \{}
\NormalTok{      const \{row:r1\}=makeSlider("n",null,1,50,1,()=\textgreater{}params.n,v=\textgreater{}params.n=v);}
\NormalTok{      const \{row:r2\}=makeSlider("p",null,0.01,0.99,0.01,()=\textgreater{}params.p,v=\textgreater{}params.p=v);}
\NormalTok{      slidersDiv.appendChild(r1); slidersDiv.appendChild(r2);}
\NormalTok{    \}}
\NormalTok{    if(d==="Poisson") \{}
\NormalTok{      const \{row\}=makeSlider("λ",null,0.1,20,0.1,()=\textgreater{}params.lam,v=\textgreater{}params.lam=v);}
\NormalTok{      slidersDiv.appendChild(row);}
\NormalTok{    \}}
\NormalTok{    if(d==="Exponential") \{}
\NormalTok{      const \{row\}=makeSlider("λ (rate)",null,0.1,5,0.05,()=\textgreater{}params.lam,v=\textgreater{}params.lam=v);}
\NormalTok{      slidersDiv.appendChild(row);}
\NormalTok{      const note=document.createElement("div");}
\NormalTok{      note.style.cssText="font{-}size:0.78em; color:\#6b7280; margin{-}top:0.2rem;";}
\NormalTok{      note.textContent="Mean = 1/λ (rate ≠ mean)";}
\NormalTok{      slidersDiv.appendChild(note);}
\NormalTok{    \}}
\NormalTok{    if(d==="Normal") \{}
\NormalTok{      const \{row:r1\}=makeSlider("μ",null,{-}5,5,0.1,()=\textgreater{}params.mu,v=\textgreater{}params.mu=v);}
\NormalTok{      const \{row:r2\}=makeSlider("σ",null,0.1,5,0.1,()=\textgreater{}params.sigma,v=\textgreater{}params.sigma=v);}
\NormalTok{      slidersDiv.appendChild(r1); slidersDiv.appendChild(r2);}
\NormalTok{    \}}
\NormalTok{  \}}

\NormalTok{  function updateReadouts(info) \{}
\NormalTok{    const isSymmetric = ["Bernoulli","Normal"].includes(distNames[distIdx]) ||}
\NormalTok{      (distNames[distIdx]==="Binomial" \&\& Math.abs(params.p{-}0.5)\textless{}0.01);}
\NormalTok{    const medianNote = isSymmetric ? "Symmetric: median = mean" : "Skewed: median ≠ mean";}
\NormalTok{    readouts.innerHTML = \textasciigrave{}}
\NormalTok{      \textless{}div style="font{-}weight:600; color:\#374151; margin{-}bottom:0.5rem; font{-}size:0.9em;"\textgreater{}Statistics\textless{}/div\textgreater{}}
\NormalTok{      \textless{}div style="margin{-}bottom:0.35rem;"\textgreater{}\textless{}span style="color:\#6b7280;"\textgreater{}Mean:\textless{}/span\textgreater{} \textless{}strong\textgreater{}$\{info.mu.toFixed(3)\}\textless{}/strong\textgreater{}\textless{}/div\textgreater{}}
\NormalTok{      \textless{}div style="margin{-}bottom:0.35rem;"\textgreater{}\textless{}span style="color:\#6b7280;"\textgreater{}Variance:\textless{}/span\textgreater{} \textless{}strong\textgreater{}$\{info.variance.toFixed(3)\}\textless{}/strong\textgreater{}\textless{}/div\textgreater{}}
\NormalTok{      \textless{}div style="margin{-}bottom:0.35rem;"\textgreater{}\textless{}span style="color:\#6b7280;"\textgreater{}Std dev (σ):\textless{}/span\textgreater{} \textless{}strong\textgreater{}$\{info.sigma.toFixed(3)\}\textless{}/strong\textgreater{}\textless{}/div\textgreater{}}
\NormalTok{      \textless{}div style="margin{-}top:0.6rem; font{-}size:0.82em; color:\#6b7280; font{-}style:italic;"\textgreater{}$\{medianNote\}\textless{}/div\textgreater{}}
\NormalTok{    \textasciigrave{};}
\NormalTok{  \}}

\NormalTok{  function update() \{}
\NormalTok{    // highlight active button}
\NormalTok{    btns.forEach((b,i)=\textgreater{}\{}
\NormalTok{      b.style.background = i===distIdx ? ACCENT : "\#f9fafb";}
\NormalTok{      b.style.color = i===distIdx ? "\#fff" : "\#374151";}
\NormalTok{      b.style.borderColor = i===distIdx ? ACCENT : "\#d1d5db";}
\NormalTok{    \});}
\NormalTok{    const info=getDistInfo();}
\NormalTok{    draw(info);}
\NormalTok{    buildSliders();}
\NormalTok{    updateReadouts(info);}
\NormalTok{  \}}

\NormalTok{  update();}
\NormalTok{  return container;}
\NormalTok{\}}
\end{Highlighting}
\end{Shaded}

\section{Key continuous
distributions}\label{key-continuous-distributions}

\subsection{Uniform}\label{uniform}

\(X \sim \text{Uniform}(a, b)\) --- every value in \([a,b]\) is equally
likely.

\[f(x) = \frac{1}{b-a}, \quad a \leq x \leq b\]

\[E[X] = \frac{a+b}{2}, \quad \text{Var}(X) = \frac{(b-a)^2}{12}\]

The CDF is: \(F(x) = (x-a)/(b-a)\) for \(a \leq x \leq b\).

\subsection{Exponential}\label{exponential-1}

Models the time until the first event in a Poisson process --- waiting
times, component lifetimes, time between arrivals.

\[f(x) = \lambda e^{-\lambda x}, \quad x \geq 0\]

\[F(x) = 1 - e^{-\lambda x}\]

\textbf{Derivation of mean.} Integrate by parts:

\[E[X] = \int_0^{\infty} x \cdot \lambda e^{-\lambda x}\, dx\]

Let \(u = x\), \(dv = \lambda e^{-\lambda x}\, dx\). Then \(du = dx\),
\(v = -e^{-\lambda x}\).

\[E[X] = \left[-x e^{-\lambda x}\right]_0^{\infty} + \int_0^{\infty} e^{-\lambda x}\, dx
= 0 + \left[-\frac{1}{\lambda} e^{-\lambda x}\right]_0^{\infty} = \frac{1}{\lambda}\]

\textbf{Derivation of variance.} First compute \(E[X^2]\):

\[E[X^2] = \int_0^{\infty} x^2 \cdot \lambda e^{-\lambda x}\, dx = \frac{2}{\lambda^2}\]

The integration follows the same pattern as the derivation of \(E[X]\)
above --- integration by parts applied twice, or equivalently by
recognising the integral as a gamma function. Either route gives
\(2/\lambda^2\).

\[\text{Var}(X) = E[X^2] - (E[X])^2 = \frac{2}{\lambda^2} - \frac{1}{\lambda^2} = \frac{1}{\lambda^2}\]

\textbf{Memoryless property.} The exponential is the only continuous
distribution with no memory:

\[P(X > s + t \mid X > s) = P(X > t) \quad \text{for all } s, t \geq 0\]

If the component has already survived \(s\) hours, the probability it
survives another \(t\) hours is the same as if it were brand new. The
past waiting time gives no information about the future.

\textbf{Proof.} Using the survival function
\(P(X > x) = e^{-\lambda x}\):

\[P(X > s+t \mid X > s) = \frac{P(X > s+t)}{P(X > s)} = \frac{e^{-\lambda(s+t)}}{e^{-\lambda s}} = e^{-\lambda t} = P(X > t) \quad \checkmark\]

\subsection{Normal}\label{normal}

\(X \sim N(\mu, \sigma^2)\) --- the most important distribution in
applied probability, for reasons the CLT will make clear.

\[f(x) = \frac{1}{\sigma\sqrt{2\pi}} \exp\!\left(-\frac{(x-\mu)^2}{2\sigma^2}\right)\]

Parameters: \(\mu\) is the mean, \(\sigma^2\) is the variance.

\[E[X] = \mu, \quad \text{Var}(X) = \sigma^2\]

There is no closed form for the CDF --- it's evaluated numerically and
tabulated as the \textbf{standard normal CDF} \(\Phi(z)\), where
\(\Phi\) (capital phi) is the CDF of \(Z \sim N(0,1)\).

\textbf{Standardisation.} Any normal \(X \sim N(\mu, \sigma^2)\) can be
converted to the standard normal by:

\[Z = \frac{X - \mu}{\sigma}\]

Then
\(P(X \leq x) = P\!\left(Z \leq \frac{x-\mu}{\sigma}\right) = \Phi\!\left(\frac{x-\mu}{\sigma}\right)\).

The transformation \(Z = (X - \mu)/\sigma\) --- ``subtract the mean,
divide by the standard deviation'' --- centres the distribution at zero
and scales it to unit variance. Every probability question about
\(X \sim N(\mu, \sigma^2)\) reduces to looking up a value in the
standard normal table.

\begin{Shaded}
\begin{Highlighting}[]
\NormalTok{//| echo: false}
\NormalTok{\{}
\NormalTok{  // ── Normal distribution — sigma bands (Priority 1) ──────────────────────}
\NormalTok{  const W = 560, H = 300, PAD = \{ top: 30, right: 30, bottom: 50, left: 50 \};}
\NormalTok{  const ACCENT = "\#3b82f6";}
\NormalTok{  const BANDS = [}
\NormalTok{    \{ mult: 3, fill: "rgba(59,130,246,0.10)", pct: "99.7\%", yFrac: 0.18 \},}
\NormalTok{    \{ mult: 2, fill: "rgba(59,130,246,0.20)", pct: "95.4\%", yFrac: 0.40 \},}
\NormalTok{    \{ mult: 1, fill: "rgba(59,130,246,0.38)", pct: "68.3\%", yFrac: 0.68 \},}
\NormalTok{  ];}
\NormalTok{  const GREY = "\#9ca3af";}

\NormalTok{  function erf(x) \{}
\NormalTok{    const sign=x\textless{}0?{-}1:1; x=Math.abs(x);}
\NormalTok{    const a1=0.254829592,a2={-}0.284496736,a3=1.421413741,a4={-}1.453152027,a5=1.061405429,p=0.3275911;}
\NormalTok{    const t=1/(1+p*x);}
\NormalTok{    return sign*(1{-}(((((a5*t+a4)*t)+a3)*t+a2)*t+a1)*t*Math.exp({-}x*x));}
\NormalTok{  \}}
\NormalTok{  function normalPDF(x,mu,sig) \{ return Math.exp({-}0.5*((x{-}mu)/sig)**2)/(sig*Math.sqrt(2*Math.PI)); \}}
\NormalTok{  function normalCDF(x,mu,sig) \{ return 0.5*(1+erf((x{-}mu)/(sig*Math.SQRT2))); \}}

\NormalTok{  let mu=0, sigma=1, showCDF=false, cursorX=0;}

\NormalTok{  const container=document.createElement("div");}
\NormalTok{  container.style.cssText="font{-}family:inherit; max{-}width:600px; margin:1rem 0;";}

\NormalTok{  // sliders row}
\NormalTok{  const sliderRow=document.createElement("div");}
\NormalTok{  sliderRow.style.cssText="display:flex; gap:1rem; flex{-}wrap:wrap; margin{-}bottom:0.75rem; align{-}items:center;";}

\NormalTok{  function makeSlider2(label,min,max,step,def,cb) \{}
\NormalTok{    const wrap=document.createElement("div");}
\NormalTok{    wrap.style.cssText="display:flex; align{-}items:center; gap:0.4rem;";}
\NormalTok{    const lbl=document.createElement("label");}
\NormalTok{    lbl.style.cssText="font{-}size:0.85em; min{-}width:20px;";}
\NormalTok{    lbl.textContent=label;}
\NormalTok{    const inp=document.createElement("input");}
\NormalTok{    inp.type="range"; inp.min=min; inp.max=max; inp.step=step; inp.value=def;}
\NormalTok{    inp.style.cssText="width:120px;";}
\NormalTok{    const val=document.createElement("span");}
\NormalTok{    val.style.cssText="font{-}size:0.85em; min{-}width:30px;";}
\NormalTok{    val.textContent=Number(def).toFixed(step\textless{}0.1?1:step\textless{}1?1:0);}
\NormalTok{    inp.oninput=()=\textgreater{}\{ val.textContent=Number(inp.value).toFixed(step\textless{}0.1?1:step\textless{}1?1:0); cb(parseFloat(inp.value)); \};}
\NormalTok{    wrap.appendChild(lbl); wrap.appendChild(inp); wrap.appendChild(val);}
\NormalTok{    return wrap;}
\NormalTok{  \}}
\NormalTok{  sliderRow.appendChild(makeSlider2("μ",{-}5,5,0.1,0,v=\textgreater{}\{mu=v;draw();\}));}
\NormalTok{  sliderRow.appendChild(makeSlider2("σ",0.2,4,0.1,1,v=\textgreater{}\{sigma=v;draw();\}));}

\NormalTok{  // CDF toggle}
\NormalTok{  const cdfBtn=document.createElement("button");}
\NormalTok{  cdfBtn.textContent="Show CDF";}
\NormalTok{  cdfBtn.style.cssText="padding:0.3rem 0.7rem; border{-}radius:4px; border:1px solid \#d1d5db; cursor:pointer; font{-}size:0.85em; background:\#f9fafb;";}
\NormalTok{  cdfBtn.onclick=()=\textgreater{}\{showCDF=!showCDF; cdfBtn.textContent=showCDF?"Hide CDF":"Show CDF"; cdfBtn.style.background=showCDF?ACCENT:"\#f9fafb"; cdfBtn.style.color=showCDF?"\#fff":"\#374151"; draw();\};}
\NormalTok{  sliderRow.appendChild(cdfBtn);}
\NormalTok{  container.appendChild(sliderRow);}

\NormalTok{  const svgEl=document.createElementNS("http://www.w3.org/2000/svg","svg");}
\NormalTok{  svgEl.setAttribute("width",W); svgEl.setAttribute("height",H);}
\NormalTok{  svgEl.style.cssText="border:1px solid \#e5e7eb; border{-}radius:6px; background:\#fff; display:block;";}
\NormalTok{  container.appendChild(svgEl);}

\NormalTok{  // cursor live readout}
\NormalTok{  const cursorReadout=document.createElement("div");}
\NormalTok{  cursorReadout.style.cssText="font{-}size:0.82em; color:\#6b7280; margin{-}top:0.35rem; min{-}height:1.2em;";}
\NormalTok{  container.appendChild(cursorReadout);}

\NormalTok{  const labelDiv=document.createElement("div");}
\NormalTok{  labelDiv.style.cssText="margin{-}top:0.3rem; font{-}size:0.82em; color:\#6b7280; font{-}style:italic;";}
\NormalTok{  labelDiv.textContent="μ sets the centre. σ sets the width. The shaded areas — 68\%, 95\%, 99.7\% — are always the same.";}
\NormalTok{  container.appendChild(labelDiv);}

\NormalTok{  // draggable cursor state}
\NormalTok{  let dragging=false;}
\NormalTok{  let svgRect=null;}

\NormalTok{  function xToData(px) \{}
\NormalTok{    const cw=W{-}PAD.left{-}PAD.right;}
\NormalTok{    const xmin=mu{-}4.5*sigma, xmax=mu+4.5*sigma;}
\NormalTok{    return xmin+(px{-}PAD.left)/cw*(xmax{-}xmin);}
\NormalTok{  \}}

\NormalTok{  svgEl.addEventListener("mousedown",e=\textgreater{}\{}
\NormalTok{    if(!showCDF) return;}
\NormalTok{    dragging=true; svgRect=svgEl.getBoundingClientRect();}
\NormalTok{    cursorX=xToData(e.clientX{-}svgRect.left); draw();}
\NormalTok{  \});}
\NormalTok{  window.addEventListener("mousemove",e=\textgreater{}\{}
\NormalTok{    if(!dragging||!showCDF) return;}
\NormalTok{    cursorX=xToData(e.clientX{-}svgRect.left); draw();}
\NormalTok{  \});}
\NormalTok{  window.addEventListener("mouseup",()=\textgreater{}\{ dragging=false; \});}

\NormalTok{  function draw() \{}
\NormalTok{    while(svgEl.firstChild) svgEl.removeChild(svgEl.firstChild);}
\NormalTok{    const cw=W{-}PAD.left{-}PAD.right;}
\NormalTok{    const ch=H{-}PAD.top{-}PAD.bottom;}
\NormalTok{    const xmin=mu{-}4.5*sigma, xmax=mu+4.5*sigma;}
\NormalTok{    const N=250;}
\NormalTok{    const xs=Array.from(\{length:N\},(\_,i)=\textgreater{}xmin+i*(xmax{-}xmin)/(N{-}1));}
\NormalTok{    const ys=xs.map(x=\textgreater{}normalPDF(x,mu,sigma));}
\NormalTok{    const yMax=normalPDF(mu,mu,sigma)*1.18;}
\NormalTok{    const sx=x=\textgreater{}PAD.left+(x{-}xmin)/(xmax{-}xmin)*cw;}
\NormalTok{    const sy=y=\textgreater{}PAD.top+ch{-}y/yMax*ch;}

\NormalTok{    // second axis for CDF}
\NormalTok{    const syCDF = y =\textgreater{} PAD.top+ch {-} y*ch;}

\NormalTok{    // draw shaded sigma bands}
\NormalTok{    BANDS.forEach((\{mult,fill,pct,yFrac\})=\textgreater{}\{}
\NormalTok{      const xl=mu{-}mult*sigma, xr=mu+mult*sigma;}
\NormalTok{      const bandXs=xs.filter(x=\textgreater{}x\textgreater{}=xl\&\&x\textless{}=xr);}
\NormalTok{      const pts=[}
\NormalTok{        \textasciigrave{}$\{sx(xl)\},$\{sy(0)\}\textasciigrave{},}
\NormalTok{        ...bandXs.map(x=\textgreater{}\textasciigrave{}$\{sx(x)\},$\{sy(normalPDF(x,mu,sigma))\}\textasciigrave{}),}
\NormalTok{        \textasciigrave{}$\{sx(xr)\},$\{sy(0)\}\textasciigrave{}}
\NormalTok{      ].join(" ");}
\NormalTok{      const poly=document.createElementNS("http://www.w3.org/2000/svg","polygon");}
\NormalTok{      poly.setAttribute("points",pts); poly.setAttribute("fill",fill);}
\NormalTok{      svgEl.appendChild(poly);}

\NormalTok{      // band percentage label}
\NormalTok{      const lbl=document.createElementNS("http://www.w3.org/2000/svg","text");}
\NormalTok{      lbl.setAttribute("x",sx(mu)); lbl.setAttribute("y",PAD.top+ch*(1{-}yFrac));}
\NormalTok{      lbl.setAttribute("text{-}anchor","middle"); lbl.setAttribute("font{-}size","11");}
\NormalTok{      lbl.setAttribute("fill","\#1e40af"); lbl.setAttribute("font{-}weight","500");}
\NormalTok{      lbl.textContent=pct;}
\NormalTok{      svgEl.appendChild(lbl);}
\NormalTok{    \});}

\NormalTok{    // PDF curve}
\NormalTok{    const pdfPts=xs.map((x,i)=\textgreater{}\textasciigrave{}$\{sx(x)\},$\{sy(ys[i])\}\textasciigrave{}).join(" ");}
\NormalTok{    const pdfLine=document.createElementNS("http://www.w3.org/2000/svg","polyline");}
\NormalTok{    pdfLine.setAttribute("points",pdfPts); pdfLine.setAttribute("fill","none");}
\NormalTok{    pdfLine.setAttribute("stroke",ACCENT); pdfLine.setAttribute("stroke{-}width","2.5");}
\NormalTok{    svgEl.appendChild(pdfLine);}

\NormalTok{    // σ bracket below x{-}axis}
\NormalTok{    const bY=PAD.top+ch+20;}
\NormalTok{    const bLine=document.createElementNS("http://www.w3.org/2000/svg","line");}
\NormalTok{    bLine.setAttribute("x1",sx(mu{-}sigma)); bLine.setAttribute("x2",sx(mu+sigma));}
\NormalTok{    bLine.setAttribute("y1",bY); bLine.setAttribute("y2",bY);}
\NormalTok{    bLine.setAttribute("stroke",ACCENT); bLine.setAttribute("stroke{-}width","2");}
\NormalTok{    svgEl.appendChild(bLine);}
\NormalTok{    // ticks at bracket ends}
\NormalTok{    [{-}1,1].forEach(s=\textgreater{}\{}
\NormalTok{      const t=document.createElementNS("http://www.w3.org/2000/svg","line");}
\NormalTok{      t.setAttribute("x1",sx(mu+s*sigma)); t.setAttribute("x2",sx(mu+s*sigma));}
\NormalTok{      t.setAttribute("y1",bY{-}4); t.setAttribute("y2",bY+4);}
\NormalTok{      t.setAttribute("stroke",ACCENT); t.setAttribute("stroke{-}width","2");}
\NormalTok{      svgEl.appendChild(t);}
\NormalTok{    \});}
\NormalTok{    const bLbl=document.createElementNS("http://www.w3.org/2000/svg","text");}
\NormalTok{    bLbl.setAttribute("x",sx(mu)); bLbl.setAttribute("y",bY+13);}
\NormalTok{    bLbl.setAttribute("text{-}anchor","middle"); bLbl.setAttribute("font{-}size","10");}
\NormalTok{    bLbl.setAttribute("fill",ACCENT);}
\NormalTok{    bLbl.textContent=\textasciigrave{}σ = $\{sigma.toFixed(1)\}\textasciigrave{};}
\NormalTok{    svgEl.appendChild(bLbl);}

\NormalTok{    // CDF curve}
\NormalTok{    if(showCDF) \{}
\NormalTok{      const cdfPts=xs.map(x=\textgreater{}\textasciigrave{}$\{sx(x)\},$\{syCDF(normalCDF(x,mu,sigma))\}\textasciigrave{}).join(" ");}
\NormalTok{      const cdfLine=document.createElementNS("http://www.w3.org/2000/svg","polyline");}
\NormalTok{      cdfLine.setAttribute("points",cdfPts); cdfLine.setAttribute("fill","none");}
\NormalTok{      cdfLine.setAttribute("stroke","\#f59e0b"); cdfLine.setAttribute("stroke{-}width","2");}
\NormalTok{      cdfLine.setAttribute("stroke{-}dasharray","4,2");}
\NormalTok{      svgEl.appendChild(cdfLine);}
\NormalTok{      // right y{-}axis}
\NormalTok{      const ryLine=document.createElementNS("http://www.w3.org/2000/svg","line");}
\NormalTok{      ryLine.setAttribute("x1",PAD.left+cw); ryLine.setAttribute("x2",PAD.left+cw);}
\NormalTok{      ryLine.setAttribute("y1",PAD.top); ryLine.setAttribute("y2",PAD.top+ch);}
\NormalTok{      ryLine.setAttribute("stroke","\#f59e0b"); ryLine.setAttribute("stroke{-}width","1");}
\NormalTok{      svgEl.appendChild(ryLine);}
\NormalTok{      const ryLbl=document.createElementNS("http://www.w3.org/2000/svg","text");}
\NormalTok{      ryLbl.setAttribute("x",PAD.left+cw+10); ryLbl.setAttribute("y",PAD.top+ch/2);}
\NormalTok{      ryLbl.setAttribute("font{-}size","10"); ryLbl.setAttribute("fill","\#f59e0b");}
\NormalTok{      ryLbl.setAttribute("text{-}anchor","middle");}
\NormalTok{      ryLbl.setAttribute("transform",\textasciigrave{}rotate(90,$\{PAD.left+cw+10\},$\{PAD.top+ch/2\})\textasciigrave{});}
\NormalTok{      ryLbl.textContent="Φ(x)";}
\NormalTok{      svgEl.appendChild(ryLbl);}
\NormalTok{      // cursor}
\NormalTok{      const cxp=Math.max(xmin,Math.min(xmax,cursorX));}
\NormalTok{      const cxv=sx(cxp);}
\NormalTok{      const cursorLine=document.createElementNS("http://www.w3.org/2000/svg","line");}
\NormalTok{      cursorLine.setAttribute("x1",cxv); cursorLine.setAttribute("x2",cxv);}
\NormalTok{      cursorLine.setAttribute("y1",PAD.top); cursorLine.setAttribute("y2",PAD.top+ch);}
\NormalTok{      cursorLine.setAttribute("stroke","\#6b7280"); cursorLine.setAttribute("stroke{-}width","1.5");}
\NormalTok{      cursorLine.setAttribute("stroke{-}dasharray","3,2");}
\NormalTok{      svgEl.appendChild(cursorLine);}
\NormalTok{      const pval=normalCDF(cxp,mu,sigma);}
\NormalTok{      const dot=document.createElementNS("http://www.w3.org/2000/svg","circle");}
\NormalTok{      dot.setAttribute("cx",cxv); dot.setAttribute("cy",syCDF(pval)); dot.setAttribute("r","5");}
\NormalTok{      dot.setAttribute("fill","\#f59e0b"); dot.setAttribute("stroke","\#fff"); dot.setAttribute("stroke{-}width","1.5");}
\NormalTok{      svgEl.appendChild(dot);}
\NormalTok{      cursorReadout.textContent=\textasciigrave{}x = $\{cxp.toFixed(2)\}  →  P(X ≤ x) = Φ($\{((cxp{-}mu)/sigma).toFixed(2)\}) = $\{pval.toFixed(4)\}\textasciigrave{};}
\NormalTok{      svgEl.style.cursor="col{-}resize";}
\NormalTok{    \} else \{}
\NormalTok{      cursorReadout.textContent="Toggle \textquotesingle{}Show CDF\textquotesingle{} and drag to read off P(X ≤ x).";}
\NormalTok{      svgEl.style.cursor="default";}
\NormalTok{    \}}

\NormalTok{    // mean dashed line}
\NormalTok{    const muLine=document.createElementNS("http://www.w3.org/2000/svg","line");}
\NormalTok{    muLine.setAttribute("x1",sx(mu)); muLine.setAttribute("x2",sx(mu));}
\NormalTok{    muLine.setAttribute("y1",PAD.top); muLine.setAttribute("y2",PAD.top+ch);}
\NormalTok{    muLine.setAttribute("stroke","\#374151"); muLine.setAttribute("stroke{-}width","1.5");}
\NormalTok{    muLine.setAttribute("stroke{-}dasharray","5,3");}
\NormalTok{    svgEl.appendChild(muLine);}
\NormalTok{    const muLbl=document.createElementNS("http://www.w3.org/2000/svg","text");}
\NormalTok{    muLbl.setAttribute("x",sx(mu)+4); muLbl.setAttribute("y",PAD.top+13);}
\NormalTok{    muLbl.setAttribute("font{-}size","11"); muLbl.setAttribute("fill","\#374151");}
\NormalTok{    muLbl.textContent=\textasciigrave{}μ=$\{mu.toFixed(1)\}\textasciigrave{};}
\NormalTok{    svgEl.appendChild(muLbl);}

\NormalTok{    // x{-}axis}
\NormalTok{    const xAxisEl=document.createElementNS("http://www.w3.org/2000/svg","line");}
\NormalTok{    xAxisEl.setAttribute("x1",PAD.left); xAxisEl.setAttribute("x2",PAD.left+cw);}
\NormalTok{    xAxisEl.setAttribute("y1",sy(0)); xAxisEl.setAttribute("y2",sy(0));}
\NormalTok{    xAxisEl.setAttribute("stroke","\#374151"); xAxisEl.setAttribute("stroke{-}width","1");}
\NormalTok{    svgEl.appendChild(xAxisEl);}

\NormalTok{    // x{-}axis ticks}
\NormalTok{    const ticks=[{-}4,{-}3,{-}2,{-}1,0,1,2,3,4].map(z=\textgreater{}mu+z*sigma).filter(v=\textgreater{}v\textgreater{}=xmin\&\&v\textless{}=xmax);}
\NormalTok{    ticks.forEach(v=\textgreater{}\{}
\NormalTok{      const tick=document.createElementNS("http://www.w3.org/2000/svg","line");}
\NormalTok{      tick.setAttribute("x1",sx(v)); tick.setAttribute("x2",sx(v));}
\NormalTok{      tick.setAttribute("y1",sy(0)); tick.setAttribute("y2",sy(0)+5);}
\NormalTok{      tick.setAttribute("stroke","\#374151"); tick.setAttribute("stroke{-}width","1");}
\NormalTok{      svgEl.appendChild(tick);}
\NormalTok{      const lbl=document.createElementNS("http://www.w3.org/2000/svg","text");}
\NormalTok{      lbl.setAttribute("x",sx(v)); lbl.setAttribute("y",sy(0)+16);}
\NormalTok{      lbl.setAttribute("text{-}anchor","middle"); lbl.setAttribute("font{-}size","10");}
\NormalTok{      lbl.setAttribute("fill","\#6b7280");}
\NormalTok{      lbl.textContent=v.toFixed(1);}
\NormalTok{      svgEl.appendChild(lbl);}
\NormalTok{    \});}

\NormalTok{    // y{-}axis}
\NormalTok{    const yAxisEl=document.createElementNS("http://www.w3.org/2000/svg","line");}
\NormalTok{    yAxisEl.setAttribute("x1",PAD.left); yAxisEl.setAttribute("x2",PAD.left);}
\NormalTok{    yAxisEl.setAttribute("y1",PAD.top); yAxisEl.setAttribute("y2",sy(0));}
\NormalTok{    yAxisEl.setAttribute("stroke","\#374151"); yAxisEl.setAttribute("stroke{-}width","1");}
\NormalTok{    svgEl.appendChild(yAxisEl);}
\NormalTok{    const yLbl=document.createElementNS("http://www.w3.org/2000/svg","text");}
\NormalTok{    yLbl.setAttribute("x",12); yLbl.setAttribute("y",PAD.top+H/3);}
\NormalTok{    yLbl.setAttribute("text{-}anchor","middle"); yLbl.setAttribute("font{-}size","11");}
\NormalTok{    yLbl.setAttribute("fill","\#6b7280");}
\NormalTok{    yLbl.setAttribute("transform",\textasciigrave{}rotate({-}90,12,$\{PAD.top+H/3\})\textasciigrave{});}
\NormalTok{    yLbl.textContent="Density";}
\NormalTok{    svgEl.appendChild(yLbl);}
\NormalTok{  \}}

\NormalTok{  draw();}
\NormalTok{  return container;}
\NormalTok{\}}
\end{Highlighting}
\end{Shaded}

\section{Joint distributions}\label{joint-distributions}

For two random variables \(X\) and \(Y\), the \textbf{joint
distribution} describes their behaviour together. The key concept for
most applications is independence.

\(X\) and \(Y\) are \textbf{independent} if knowing the value of one
gives no information about the other:

\[f_{X,Y}(x,y) = f_X(x) \cdot f_Y(y)\]

That is, the joint density factors into the product of the marginals.

\textbf{Covariance.} A measure of how \(X\) and \(Y\) move together:

\[\text{Cov}(X, Y) = E[(X - \mu_X)(Y - \mu_Y)] = E[XY] - E[X]\,E[Y]\]

If \(X\) and \(Y\) are independent, \(E[XY] = E[X]\,E[Y]\), so
\(\text{Cov}(X,Y) = 0\). (The converse is not generally true: zero
covariance does not imply independence.)

\textbf{Correlation coefficient.} Covariance is scale-dependent ---
multiplying \(X\) by 2 doubles the covariance. The correlation
normalises this:

\[\rho(X,Y) = \frac{\text{Cov}(X,Y)}{\sigma_X \sigma_Y}, \quad -1 \leq \rho \leq 1\]

\(\rho = 1\) means perfect positive linear relationship; \(\rho = -1\)
means perfect negative; \(\rho = 0\) means no linear relationship.

\begin{Shaded}
\begin{Highlighting}[]
\NormalTok{//| echo: false}
\NormalTok{\{}
\NormalTok{  // ── Correlation scatter cloud (OPPORTUNITY 4) ────────────────────────────}
\NormalTok{  const W=420, H=380, PAD=\{top:30,right:30,bottom:50,left:50\};}
\NormalTok{  const ACCENT="\#3b82f6";}
\NormalTok{  const N\_PTS=200;}

\NormalTok{  let rho=0;}
\NormalTok{  let seed=99;}
\NormalTok{  function lcg() \{ seed=(seed*1664525+1013904223)\&0xffffffff; return (seed\textgreater{}\textgreater{}\textgreater{}0)/0xffffffff; \}}
\NormalTok{  function boxMuller() \{}
\NormalTok{    const u1=lcg()||1e{-}10, u2=lcg();}
\NormalTok{    return [Math.sqrt({-}2*Math.log(u1))*Math.cos(2*Math.PI*u2),}
\NormalTok{            Math.sqrt({-}2*Math.log(u1))*Math.sin(2*Math.PI*u2)];}
\NormalTok{  \}}
\NormalTok{  function genPoints(r) \{}
\NormalTok{    seed=99;}
\NormalTok{    const pts=[];}
\NormalTok{    for(let i=0;i\textless{}N\_PTS;i++) \{}
\NormalTok{      const [z1,z2]=boxMuller();}
\NormalTok{      // Cholesky: x=z1, y=r*z1+sqrt(1{-}r\^{}2)*z2}
\NormalTok{      pts.push([z1, r*z1+Math.sqrt(Math.max(0,1{-}r*r))*z2]);}
\NormalTok{    \}}
\NormalTok{    return pts;}
\NormalTok{  \}}

\NormalTok{  const container=document.createElement("div");}
\NormalTok{  container.style.cssText="font{-}family:inherit; max{-}width:460px; margin:1rem 0;";}

\NormalTok{  const sliderRow=document.createElement("div");}
\NormalTok{  sliderRow.style.cssText="display:flex; align{-}items:center; gap:0.6rem; margin{-}bottom:0.75rem;";}
\NormalTok{  const slLbl=document.createElement("label");}
\NormalTok{  slLbl.style.cssText="font{-}size:0.85em;"; slLbl.textContent="ρ:";}
\NormalTok{  const slInp=document.createElement("input");}
\NormalTok{  slInp.type="range"; slInp.min={-}1; slInp.max=1; slInp.step=0.01; slInp.value=0;}
\NormalTok{  slInp.style.cssText="width:200px;";}
\NormalTok{  const slVal=document.createElement("span");}
\NormalTok{  slVal.style.cssText="font{-}size:0.9em; font{-}weight:600; min{-}width:40px;";}
\NormalTok{  slVal.textContent="0.00";}
\NormalTok{  slInp.oninput=()=\textgreater{}\{ rho=parseFloat(slInp.value); slVal.textContent=rho.toFixed(2); draw(); \};}
\NormalTok{  sliderRow.appendChild(slLbl); sliderRow.appendChild(slInp); sliderRow.appendChild(slVal);}
\NormalTok{  container.appendChild(sliderRow);}

\NormalTok{  const svgEl=document.createElementNS("http://www.w3.org/2000/svg","svg");}
\NormalTok{  svgEl.setAttribute("width",W); svgEl.setAttribute("height",H);}
\NormalTok{  svgEl.style.cssText="border:1px solid \#e5e7eb; border{-}radius:6px; background:\#fff; display:block;";}
\NormalTok{  container.appendChild(svgEl);}

\NormalTok{  const readoutDiv=document.createElement("div");}
\NormalTok{  readoutDiv.style.cssText="font{-}size:0.82em; color:\#374151; margin{-}top:0.4rem;";}
\NormalTok{  container.appendChild(readoutDiv);}

\NormalTok{  const labelDiv=document.createElement("div");}
\NormalTok{  labelDiv.style.cssText="margin{-}top:0.3rem; font{-}size:0.82em; color:\#6b7280; font{-}style:italic;";}
\NormalTok{  labelDiv.textContent="ρ is the shape of the data cloud: circular, elliptical, or linear.";}
\NormalTok{  container.appendChild(labelDiv);}

\NormalTok{  function draw() \{}
\NormalTok{    while(svgEl.firstChild) svgEl.removeChild(svgEl.firstChild);}
\NormalTok{    const cw=W{-}PAD.left{-}PAD.right;}
\NormalTok{    const ch=H{-}PAD.top{-}PAD.bottom;}
\NormalTok{    const r=Math.max({-}0.999,Math.min(0.999,rho));}
\NormalTok{    const pts=genPoints(r);}
\NormalTok{    const range=3.5;}
\NormalTok{    const sx=x=\textgreater{}PAD.left+(x+range)/(2*range)*cw;}
\NormalTok{    const sy=y=\textgreater{}PAD.top+(range{-}y)/(2*range)*ch;}

\NormalTok{    // grid}
\NormalTok{    for(let v={-}3;v\textless{}=3;v++) \{}
\NormalTok{      const gx=document.createElementNS("http://www.w3.org/2000/svg","line");}
\NormalTok{      gx.setAttribute("x1",sx(v)); gx.setAttribute("x2",sx(v));}
\NormalTok{      gx.setAttribute("y1",PAD.top); gx.setAttribute("y2",PAD.top+ch);}
\NormalTok{      gx.setAttribute("stroke","\#f3f4f6"); gx.setAttribute("stroke{-}width","1");}
\NormalTok{      svgEl.appendChild(gx);}
\NormalTok{      const gy=document.createElementNS("http://www.w3.org/2000/svg","line");}
\NormalTok{      gy.setAttribute("x1",PAD.left); gy.setAttribute("x2",PAD.left+cw);}
\NormalTok{      gy.setAttribute("y1",sy(v)); gy.setAttribute("y2",sy(v));}
\NormalTok{      gy.setAttribute("stroke","\#f3f4f6"); gy.setAttribute("stroke{-}width","1");}
\NormalTok{      svgEl.appendChild(gy);}
\NormalTok{    \}}

\NormalTok{    // axes}
\NormalTok{    ["x","y"].forEach(axis=\textgreater{}\{}
\NormalTok{      const l=document.createElementNS("http://www.w3.org/2000/svg","line");}
\NormalTok{      if(axis==="x") \{ l.setAttribute("x1",PAD.left); l.setAttribute("x2",PAD.left+cw); l.setAttribute("y1",sy(0)); l.setAttribute("y2",sy(0)); \}}
\NormalTok{      else \{ l.setAttribute("x1",sx(0)); l.setAttribute("x2",sx(0)); l.setAttribute("y1",PAD.top); l.setAttribute("y2",PAD.top+ch); \}}
\NormalTok{      l.setAttribute("stroke","\#9ca3af"); l.setAttribute("stroke{-}width","1");}
\NormalTok{      svgEl.appendChild(l);}
\NormalTok{    \});}

\NormalTok{    // scatter points}
\NormalTok{    pts.forEach(([x,y])=\textgreater{}\{}
\NormalTok{      const dot=document.createElementNS("http://www.w3.org/2000/svg","circle");}
\NormalTok{      dot.setAttribute("cx",sx(x)); dot.setAttribute("cy",sy(y)); dot.setAttribute("r","3");}
\NormalTok{      dot.setAttribute("fill",ACCENT); dot.setAttribute("fill{-}opacity","0.5");}
\NormalTok{      svgEl.appendChild(dot);}
\NormalTok{    \});}

\NormalTok{    // 1{-}sigma ellipse of bivariate normal}
\NormalTok{    const ellN=80;}
\NormalTok{    const ellPts=Array.from(\{length:ellN+1\},(\_,i)=\textgreater{}\{}
\NormalTok{      const t=2*Math.PI*i/ellN;}
\NormalTok{      const ex=Math.cos(t);}
\NormalTok{      const ey=r*Math.cos(t)+Math.sqrt(Math.max(0,1{-}r*r))*Math.sin(t);}
\NormalTok{      return \textasciigrave{}$\{sx(ex)\},$\{sy(ey)\}\textasciigrave{};}
\NormalTok{    \}).join(" ");}
\NormalTok{    const ellipse=document.createElementNS("http://www.w3.org/2000/svg","polyline");}
\NormalTok{    ellipse.setAttribute("points",ellPts); ellipse.setAttribute("fill","none");}
\NormalTok{    ellipse.setAttribute("stroke","\#374151"); ellipse.setAttribute("stroke{-}width","2");}
\NormalTok{    ellipse.setAttribute("stroke{-}dasharray","4,2");}
\NormalTok{    svgEl.appendChild(ellipse);}

\NormalTok{    // axis labels}
\NormalTok{    const xLbl=document.createElementNS("http://www.w3.org/2000/svg","text");}
\NormalTok{    xLbl.setAttribute("x",PAD.left+cw/2); xLbl.setAttribute("y",PAD.top+ch+38);}
\NormalTok{    xLbl.setAttribute("text{-}anchor","middle"); xLbl.setAttribute("font{-}size","12");}
\NormalTok{    xLbl.setAttribute("fill","\#374151"); xLbl.textContent="X";}
\NormalTok{    svgEl.appendChild(xLbl);}
\NormalTok{    const yLbl=document.createElementNS("http://www.w3.org/2000/svg","text");}
\NormalTok{    yLbl.setAttribute("x",14); yLbl.setAttribute("y",PAD.top+ch/2);}
\NormalTok{    yLbl.setAttribute("text{-}anchor","middle"); yLbl.setAttribute("font{-}size","12");}
\NormalTok{    yLbl.setAttribute("fill","\#374151");}
\NormalTok{    yLbl.setAttribute("transform",\textasciigrave{}rotate({-}90,14,$\{PAD.top+ch/2\})\textasciigrave{});}
\NormalTok{    yLbl.textContent="Y";}
\NormalTok{    svgEl.appendChild(yLbl);}

\NormalTok{    // rho annotation}
\NormalTok{    const rAnn=document.createElementNS("http://www.w3.org/2000/svg","text");}
\NormalTok{    rAnn.setAttribute("x",PAD.left+8); rAnn.setAttribute("y",PAD.top+16);}
\NormalTok{    rAnn.setAttribute("font{-}size","12"); rAnn.setAttribute("fill","\#374151");}
\NormalTok{    rAnn.setAttribute("font{-}weight","500");}
\NormalTok{    rAnn.textContent=\textasciigrave{}ρ = $\{r.toFixed(2)\}\textasciigrave{};}
\NormalTok{    svgEl.appendChild(rAnn);}

\NormalTok{    // interpretive note}
\NormalTok{    const absR=Math.abs(r);}
\NormalTok{    let strength = absR \textless{} 0.3 ? "Weak" : absR \textless{} 0.7 ? "Moderate" : "Strong";}
\NormalTok{    let direction = r\textgreater{}0.02 ? "Positive: Y tends to increase with X" : r\textless{}{-}0.02 ? "Negative: Y tends to decrease with X" : "No linear relationship";}
\NormalTok{    if(absR\textless{}0.05) direction="Knowing X tells you nothing about Y";}
\NormalTok{    if(absR\textgreater{}0.97) direction=r\textgreater{}0?"Perfect positive linear relationship":"Perfect negative linear relationship";}
\NormalTok{    const angleRad=Math.atan(r); const angleDeg=(angleRad*180/Math.PI).toFixed(0);}
\NormalTok{    readoutDiv.textContent=\textasciigrave{}ρ = $\{r.toFixed(2)\} | $\{strength\} relationship | $\{direction\}\textasciigrave{};}
\NormalTok{  \}}

\NormalTok{  draw();}
\NormalTok{  return container;}
\NormalTok{\}}
\end{Highlighting}
\end{Shaded}

\textbf{Variance of a sum.} For any \(X\) and \(Y\) (not necessarily
independent):

\[\text{Var}(X + Y) = \text{Var}(X) + \text{Var}(Y) + 2\,\text{Cov}(X,Y)\]

If they are independent, the covariance term vanishes.

\section{Central Limit Theorem}\label{central-limit-theorem}

Let \(X_1, X_2, \ldots, X_n\) be independent, identically distributed
random variables with mean \(\mu\) and variance \(\sigma^2 < \infty\).
Define the sample mean:

\[\bar{X}_n = \frac{1}{n}\sum_{i=1}^{n} X_i\]

\textbf{Central Limit Theorem (CLT):} As \(n \to \infty\),

\[\frac{\bar{X}_n - \mu}{\sigma/\sqrt{n}} \xrightarrow{d} N(0,1)\]

Read \(\xrightarrow{d}\) as ``converges in distribution to'' --- as
\(n\) grows, the distribution of the standardised mean gets closer and
closer to the standard normal.

In practice: for large \(n\), the distribution of \(\bar{X}_n\) is
approximately \(N(\mu,\, \sigma^2/n)\).

\textbf{Three things the CLT is saying:}

\begin{enumerate}
\def\labelenumi{\arabic{enumi}.}
\item
  The sample mean converges to the true mean \(\mu\) --- not just as a
  hope, but with a rate: the spread is \(\sigma/\sqrt{n}\), shrinking
  like \(1/\sqrt{n}\).
\item
  The limiting distribution is always normal, regardless of the
  distribution of the individual \(X_i\). You do not need \(X_i\) to be
  normal --- exponential, uniform, Bernoulli, anything --- the average
  becomes normal.
\item
  The only requirements are: independent, identically distributed
  (abbreviated i.i.d.), finite variance. No further conditions.
\end{enumerate}

\textbf{Conditions for practical use.} The approximation is usually
adequate when \(n \geq 30\), and excellent for \(n \geq 50\). For
distributions that are already close to normal, smaller \(n\) suffices.
For very skewed distributions (e.g.~heavy-tailed), larger \(n\) may be
needed.

\textbf{Worked example: Sample mean from an exponential population.}

Components have lifetimes \(X_i \sim \text{Exp}(\lambda = 0.5)\), so
\(\mu = 1/\lambda = 2\) hours and \(\sigma^2 = 1/\lambda^2 = 4\).

Take a sample of \(n = 50\) components. By the CLT, the sample mean is
approximately:

\[\bar{X}_{50} \approx N\!\left(2,\, \frac{4}{50}\right) = N(2,\, 0.08)\]

Standard deviation of the mean:
\(\sigma/\sqrt{n} = 2/\sqrt{50} \approx 0.283\).

Find \(P(\bar{X}_{50} > 2.4)\):

\[P(\bar{X}_{50} > 2.4) = P\!\left(Z > \frac{2.4 - 2}{0.283}\right) = P(Z > 1.414) = 1 - \Phi(1.414) \approx 0.079\]

There is about an 8\% chance the sample mean exceeds 2.4 hours. The
individual lifetimes are exponential and highly right-skewed --- but the
average over 50 of them behaves almost exactly like a normal random
variable.

\begin{Shaded}
\begin{Highlighting}[]
\NormalTok{//| echo: false}
\NormalTok{\{}
\NormalTok{  // ── CLT demonstration (VIZ 2, extended) ─────────────────────────────────}
\NormalTok{  const W=540, H=300, PAD=\{top:30,right:20,bottom:50,left:48\};}
\NormalTok{  const ACCENT="\#3b82f6", GREY="\#d1d5db";}
\NormalTok{  const LAMBDA=0.5, MU=1/LAMBDA, SIGMA2=1/(LAMBDA*LAMBDA), SIGMA=1/LAMBDA;}
\NormalTok{  const N\_SAMPLES=2000;}

\NormalTok{  let nSlider=5;}

\NormalTok{  // LCG{-}based deterministic "random" for reproducibility}
\NormalTok{  let seed=42;}
\NormalTok{  function lcgNext() \{ seed=(seed*1664525+1013904223)\&0xffffffff; return (seed\textgreater{}\textgreater{}\textgreater{}0)/0xffffffff; \}}
\NormalTok{  function resetSeed(s=42) \{ seed=s; \}}
\NormalTok{  function expSample() \{ return {-}Math.log(1{-}lcgNext())/LAMBDA; \}}
\NormalTok{  function normalPDF(x,mu,sig) \{ return Math.exp({-}0.5*((x{-}mu)/sig)**2)/(sig*Math.sqrt(2*Math.PI)); \}}
\NormalTok{  function erf(x) \{}
\NormalTok{    const sign=x\textless{}0?{-}1:1; x=Math.abs(x);}
\NormalTok{    const a1=0.254829592,a2={-}0.284496736,a3=1.421413741,a4={-}1.453152027,a5=1.061405429,p=0.3275911;}
\NormalTok{    const t=1/(1+p*x);}
\NormalTok{    return sign*(1{-}(((((a5*t+a4)*t)+a3)*t+a2)*t+a1)*t*Math.exp({-}x*x));}
\NormalTok{  \}}

\NormalTok{  const container=document.createElement("div");}
\NormalTok{  container.style.cssText="font{-}family:inherit; max{-}width:580px; margin:1rem 0;";}

\NormalTok{  // slider row}
\NormalTok{  const sliderRow=document.createElement("div");}
\NormalTok{  sliderRow.style.cssText="display:flex; align{-}items:center; gap:0.6rem; margin{-}bottom:0.75rem; flex{-}wrap:wrap;";}
\NormalTok{  const slLbl=document.createElement("label");}
\NormalTok{  slLbl.style.cssText="font{-}size:0.85em;";}
\NormalTok{  slLbl.textContent="Sample size n:";}
\NormalTok{  const slInp=document.createElement("input");}
\NormalTok{  slInp.type="range"; slInp.min=1; slInp.max=200; slInp.step=1; slInp.value=nSlider;}
\NormalTok{  slInp.style.cssText="width:200px;";}
\NormalTok{  const slVal=document.createElement("span");}
\NormalTok{  slVal.style.cssText="font{-}size:0.88em; font{-}weight:600; min{-}width:30px;";}
\NormalTok{  slVal.textContent=nSlider;}
\NormalTok{  slInp.oninput=()=\textgreater{}\{ nSlider=parseInt(slInp.value); slVal.textContent=nSlider; draw(); \};}
\NormalTok{  // snap buttons}
\NormalTok{  const snaps=[1,5,30,100];}
\NormalTok{  const snapBtns=snaps.map(v=\textgreater{}\{}
\NormalTok{    const b=document.createElement("button");}
\NormalTok{    b.textContent=\textasciigrave{}n=$\{v\}\textasciigrave{};}
\NormalTok{    b.style.cssText="padding:0.2rem 0.5rem; border:1px solid \#d1d5db; border{-}radius:3px; cursor:pointer; font{-}size:0.78em; background:\#f9fafb;";}
\NormalTok{    b.onclick=()=\textgreater{}\{ nSlider=v; slInp.value=v; slVal.textContent=v; draw(); \};}
\NormalTok{    return b;}
\NormalTok{  \});}
\NormalTok{  sliderRow.appendChild(slLbl); sliderRow.appendChild(slInp); sliderRow.appendChild(slVal);}
\NormalTok{  snaps.forEach((v,i)=\textgreater{}sliderRow.appendChild(snapBtns[i]));}
\NormalTok{  container.appendChild(sliderRow);}

\NormalTok{  const mainSvg=document.createElementNS("http://www.w3.org/2000/svg","svg");}
\NormalTok{  mainSvg.setAttribute("width",W); mainSvg.setAttribute("height",H);}
\NormalTok{  mainSvg.style.cssText="border:1px solid \#e5e7eb; border{-}radius:6px; background:\#fff; display:block;";}
\NormalTok{  container.appendChild(mainSvg);}

\NormalTok{  const infoDiv=document.createElement("div");}
\NormalTok{  infoDiv.style.cssText="font{-}size:0.82em; color:\#374151; margin{-}top:0.4rem;";}
\NormalTok{  container.appendChild(infoDiv);}

\NormalTok{  const captionDiv=document.createElement("div");}
\NormalTok{  captionDiv.style.cssText="margin{-}top:0.3rem; font{-}size:0.82em; color:\#6b7280; font{-}style:italic;";}
\NormalTok{  captionDiv.textContent="The source distribution (upper right) is strongly right{-}skewed. Watch what happens to the histogram of means as n increases.";}
\NormalTok{  container.appendChild(captionDiv);}

\NormalTok{  function getSampleMeans(n) \{}
\NormalTok{    resetSeed(123);}
\NormalTok{    const means=[];}
\NormalTok{    for(let s=0;s\textless{}N\_SAMPLES;s++) \{}
\NormalTok{      let sum=0;}
\NormalTok{      for(let i=0;i\textless{}n;i++) sum+=expSample();}
\NormalTok{      means.push(sum/n);}
\NormalTok{    \}}
\NormalTok{    return means;}
\NormalTok{  \}}

\NormalTok{  function draw() \{}
\NormalTok{    while(mainSvg.firstChild) mainSvg.removeChild(mainSvg.firstChild);}
\NormalTok{    const n=nSlider;}
\NormalTok{    const cw=W{-}PAD.left{-}PAD.right;}
\NormalTok{    const ch=H{-}PAD.top{-}PAD.bottom;}

\NormalTok{    const means=getSampleMeans(n);}
\NormalTok{    const theoSE=SIGMA/Math.sqrt(n);}

\NormalTok{    // histogram}
\NormalTok{    const xMin=Math.max(0, MU{-}5*theoSE);}
\NormalTok{    const xMax=MU+5*theoSE;}
\NormalTok{    const nBins=Math.max(15,Math.min(50,Math.round(Math.sqrt(N\_SAMPLES))));}
\NormalTok{    const binWidth=(xMax{-}xMin)/nBins;}
\NormalTok{    const bins=Array(nBins).fill(0);}
\NormalTok{    means.forEach(v=\textgreater{}\{}
\NormalTok{      const b=Math.floor((v{-}xMin)/binWidth);}
\NormalTok{      if(b\textgreater{}=0\&\&b\textless{}nBins) bins[b]++;}
\NormalTok{    \});}
\NormalTok{    const maxCount=Math.max(...bins);}
\NormalTok{    // convert to density}
\NormalTok{    const densities=bins.map(c=\textgreater{}c/(N\_SAMPLES*binWidth));}
\NormalTok{    const yMax=Math.max(...densities)*1.25;}

\NormalTok{    const sx=x=\textgreater{}PAD.left+(x{-}xMin)/(xMax{-}xMin)*cw;}
\NormalTok{    const sy=y=\textgreater{}PAD.top+ch{-}y/yMax*ch;}

\NormalTok{    // draw bars}
\NormalTok{    bins.forEach((c,i)=\textgreater{}\{}
\NormalTok{      const x0=xMin+i*binWidth;}
\NormalTok{      const d=c/(N\_SAMPLES*binWidth);}
\NormalTok{      const rect=document.createElementNS("http://www.w3.org/2000/svg","rect");}
\NormalTok{      rect.setAttribute("x",sx(x0)+0.5); rect.setAttribute("y",sy(d));}
\NormalTok{      rect.setAttribute("width",Math.max(1,sx(x0+binWidth){-}sx(x0){-}1));}
\NormalTok{      rect.setAttribute("height",Math.max(0,sy(0){-}sy(d)));}
\NormalTok{      rect.setAttribute("fill",GREY); rect.setAttribute("stroke","\#fff"); rect.setAttribute("stroke{-}width","0.5");}
\NormalTok{      mainSvg.appendChild(rect);}
\NormalTok{    \});}

\NormalTok{    // theoretical normal overlay}
\NormalTok{    const nPts=200;}
\NormalTok{    const overlayXs=Array.from(\{length:nPts\},(\_,i)=\textgreater{}xMin+i*(xMax{-}xMin)/(nPts{-}1));}
\NormalTok{    const overlayYs=overlayXs.map(x=\textgreater{}normalPDF(x,MU,theoSE));}
\NormalTok{    const overlayPts=overlayXs.map((x,i)=\textgreater{}\textasciigrave{}$\{sx(x)\},$\{sy(overlayYs[i])\}\textasciigrave{}).join(" ");}
\NormalTok{    const overlayLine=document.createElementNS("http://www.w3.org/2000/svg","polyline");}
\NormalTok{    overlayLine.setAttribute("points",overlayPts); overlayLine.setAttribute("fill","none");}
\NormalTok{    overlayLine.setAttribute("stroke",ACCENT); overlayLine.setAttribute("stroke{-}width","2.5");}
\NormalTok{    mainSvg.appendChild(overlayLine);}

\NormalTok{    // mean marker}
\NormalTok{    const muLine=document.createElementNS("http://www.w3.org/2000/svg","line");}
\NormalTok{    muLine.setAttribute("x1",sx(MU)); muLine.setAttribute("x2",sx(MU));}
\NormalTok{    muLine.setAttribute("y1",PAD.top); muLine.setAttribute("y2",PAD.top+ch);}
\NormalTok{    muLine.setAttribute("stroke","\#374151"); muLine.setAttribute("stroke{-}width","1.5");}
\NormalTok{    muLine.setAttribute("stroke{-}dasharray","5,3");}
\NormalTok{    mainSvg.appendChild(muLine);}

\NormalTok{    // axes}
\NormalTok{    const xAx=document.createElementNS("http://www.w3.org/2000/svg","line");}
\NormalTok{    xAx.setAttribute("x1",PAD.left); xAx.setAttribute("x2",PAD.left+cw);}
\NormalTok{    xAx.setAttribute("y1",sy(0)); xAx.setAttribute("y2",sy(0));}
\NormalTok{    xAx.setAttribute("stroke","\#374151"); xAx.setAttribute("stroke{-}width","1");}
\NormalTok{    mainSvg.appendChild(xAx);}
\NormalTok{    const yAx=document.createElementNS("http://www.w3.org/2000/svg","line");}
\NormalTok{    yAx.setAttribute("x1",PAD.left); yAx.setAttribute("x2",PAD.left);}
\NormalTok{    yAx.setAttribute("y1",PAD.top); yAx.setAttribute("y2",sy(0));}
\NormalTok{    yAx.setAttribute("stroke","\#374151"); yAx.setAttribute("stroke{-}width","1");}
\NormalTok{    mainSvg.appendChild(yAx);}

\NormalTok{    // x ticks}
\NormalTok{    const nTicks=5;}
\NormalTok{    for(let i=0;i\textless{}=nTicks;i++) \{}
\NormalTok{      const xv=xMin+i*(xMax{-}xMin)/nTicks;}
\NormalTok{      const tk=document.createElementNS("http://www.w3.org/2000/svg","line");}
\NormalTok{      tk.setAttribute("x1",sx(xv)); tk.setAttribute("x2",sx(xv));}
\NormalTok{      tk.setAttribute("y1",sy(0)); tk.setAttribute("y2",sy(0)+5);}
\NormalTok{      tk.setAttribute("stroke","\#374151"); tk.setAttribute("stroke{-}width","1");}
\NormalTok{      mainSvg.appendChild(tk);}
\NormalTok{      const lbl=document.createElementNS("http://www.w3.org/2000/svg","text");}
\NormalTok{      lbl.setAttribute("x",sx(xv)); lbl.setAttribute("y",sy(0)+18);}
\NormalTok{      lbl.setAttribute("text{-}anchor","middle"); lbl.setAttribute("font{-}size","10");}
\NormalTok{      lbl.setAttribute("fill","\#6b7280");}
\NormalTok{      lbl.textContent=xv.toFixed(2);}
\NormalTok{      mainSvg.appendChild(lbl);}
\NormalTok{    \}}

\NormalTok{    // x{-}axis label}
\NormalTok{    const xLbl=document.createElementNS("http://www.w3.org/2000/svg","text");}
\NormalTok{    xLbl.setAttribute("x",PAD.left+cw/2); xLbl.setAttribute("y",PAD.top+ch+38);}
\NormalTok{    xLbl.setAttribute("text{-}anchor","middle"); xLbl.setAttribute("font{-}size","11");}
\NormalTok{    xLbl.setAttribute("fill","\#374151"); xLbl.textContent="Sample mean x̄ₙ";}
\NormalTok{    mainSvg.appendChild(xLbl);}

\NormalTok{    // y{-}axis label}
\NormalTok{    const yLbl=document.createElementNS("http://www.w3.org/2000/svg","text");}
\NormalTok{    yLbl.setAttribute("x",12); yLbl.setAttribute("y",PAD.top+ch/2);}
\NormalTok{    yLbl.setAttribute("text{-}anchor","middle"); yLbl.setAttribute("font{-}size","11");}
\NormalTok{    yLbl.setAttribute("fill","\#6b7280");}
\NormalTok{    yLbl.setAttribute("transform",\textasciigrave{}rotate({-}90,12,$\{PAD.top+ch/2\})\textasciigrave{});}
\NormalTok{    yLbl.textContent="Density";}
\NormalTok{    mainSvg.appendChild(yLbl);}

\NormalTok{    // inset: source Exp(0.5) distribution}
\NormalTok{    const iW=90, iH=60, iX=W{-}PAD.right{-}iW{-}5, iY=PAD.top+2;}
\NormalTok{    const iBg=document.createElementNS("http://www.w3.org/2000/svg","rect");}
\NormalTok{    iBg.setAttribute("x",iX{-}5); iBg.setAttribute("y",iY{-}5);}
\NormalTok{    iBg.setAttribute("width",iW+10); iBg.setAttribute("height",iH+10);}
\NormalTok{    iBg.setAttribute("fill","\#f8fafc"); iBg.setAttribute("stroke","\#e5e7eb");}
\NormalTok{    iBg.setAttribute("rx","3");}
\NormalTok{    mainSvg.appendChild(iBg);}
\NormalTok{    // inset curve}
\NormalTok{    const iXs=Array.from(\{length:80\},(\_,i)=\textgreater{}i*8/79);}
\NormalTok{    const iYs=iXs.map(x=\textgreater{}LAMBDA*Math.exp({-}LAMBDA*x));}
\NormalTok{    const iYMax=LAMBDA*1.1;}
\NormalTok{    const isX=x=\textgreater{}iX+x/8*iW;}
\NormalTok{    const isY=y=\textgreater{}iY+iH{-}y/iYMax*iH;}
\NormalTok{    const iPts=iXs.map((x,i)=\textgreater{}\textasciigrave{}$\{isX(x)\},$\{isY(iYs[i])\}\textasciigrave{}).join(" ");}
\NormalTok{    const iArea=document.createElementNS("http://www.w3.org/2000/svg","polygon");}
\NormalTok{    iArea.setAttribute("points",\textasciigrave{}$\{isX(0)\},$\{isY(0)\} $\{iPts\} $\{isX(8)\},$\{isY(0)\}\textasciigrave{});}
\NormalTok{    iArea.setAttribute("fill","rgba(156,163,175,0.3)"); mainSvg.appendChild(iArea);}
\NormalTok{    const iLine=document.createElementNS("http://www.w3.org/2000/svg","polyline");}
\NormalTok{    iLine.setAttribute("points",iPts); iLine.setAttribute("fill","none");}
\NormalTok{    iLine.setAttribute("stroke","\#9ca3af"); iLine.setAttribute("stroke{-}width","1.5");}
\NormalTok{    mainSvg.appendChild(iLine);}
\NormalTok{    const iLbl=document.createElementNS("http://www.w3.org/2000/svg","text");}
\NormalTok{    iLbl.setAttribute("x",iX+iW/2); iLbl.setAttribute("y",iY+iH+10);}
\NormalTok{    iLbl.setAttribute("text{-}anchor","middle"); iLbl.setAttribute("font{-}size","9");}
\NormalTok{    iLbl.setAttribute("fill","\#6b7280"); iLbl.textContent="Source: Exp(0.5)";}
\NormalTok{    mainSvg.appendChild(iLbl);}

\NormalTok{    // info readout}
\NormalTok{    infoDiv.textContent=\textasciigrave{}n = $\{n\} | Theoretical mean: $\{MU.toFixed(1)\} | Theoretical SE: $\{theoSE.toFixed(3)\}\textasciigrave{};}
\NormalTok{  \}}

\NormalTok{  draw();}
\NormalTok{  return container;}
\NormalTok{\}}
\end{Highlighting}
\end{Shaded}

\section{Where this goes}\label{where-this-goes-35}

This chapter provides the foundation for \textbf{Chapter 2: Mathematical
statistics} (this volume). Everything in that chapter --- maximum
likelihood estimation, confidence intervals, hypothesis tests --- takes
the distributions developed here and uses them to make inferences about
unknown parameters from data. The CLT is the engine: it justifies using
normal-distribution machinery on sample means regardless of the
underlying population distribution, which is why \(z\)-tests and
\(t\)-tests work in practice.

The connections extend further into other sections of Vol 7. The Poisson
distribution arises in \textbf{queuing theory} (a special case of
stochastic processes). The normal distribution underlies error analysis
in \textbf{numerical methods}: when you propagate measurement
uncertainties through a computation, the CLT explains why the output
errors are approximately normal. In \textbf{engineering statistics}, the
same distributions appear in reliability theory, quality control (Six
Sigma thresholds are expressed in \(\sigma\) units), and signal
detection.

\begin{tcolorbox}[enhanced jigsaw, opacityback=0, colback=white, breakable, colframe=quarto-callout-tip-color-frame, arc=.35mm, left=2mm, toprule=.15mm, leftrule=.75mm, bottomrule=.15mm, rightrule=.15mm]
\begin{minipage}[t]{5.5mm}
\textcolor{quarto-callout-tip-color}{\faLightbulb}
\end{minipage}%
\begin{minipage}[t]{\textwidth - 5.5mm}

\textbf{Where this shows up}

\begin{itemize}
\tightlist
\item
  A reliability engineer models component lifetimes as exponential and
  uses the memoryless property to compute replacement schedules.
\item
  An actuary prices insurance using the Poisson distribution to model
  the number of claims per period.
\item
  A machine learning engineer interprets model output probabilities
  using the Bayes framework --- the model's output is \(P(Y \mid X)\),
  not \(P(X \mid Y)\).
\item
  A signal processing engineer applies the CLT to justify that thermal
  noise in electronic circuits is modelled as Gaussian.
\item
  A quality control engineer uses the binomial distribution to decide
  whether a batch rejection threshold is appropriate for a given defect
  rate.
\end{itemize}

\end{minipage}%
\end{tcolorbox}

\section{Exercises}\label{exercises-46}

These are puzzles. Each has a clean numerical answer. The interesting
part is identifying which distribution applies and setting up the
probability correctly.

\begin{center}\rule{0.5\linewidth}{0.5pt}\end{center}

\textbf{Exercise 1.} A factory has two machines. Machine A produces 60\%
of output and has a 4\% defect rate. Machine B produces the remaining
40\% and has a 1\% defect rate. An inspector picks a component at random
and finds it defective. What is the probability it came from Machine A?

\begin{center}\rule{0.5\linewidth}{0.5pt}\end{center}

\textbf{Exercise 2.} A quality control test checks batches of 8
components. Each component, independently, has a failure probability of
0.1. Find the probability that at least 3 components in a batch fail.

\begin{center}\rule{0.5\linewidth}{0.5pt}\end{center}

\textbf{Exercise 3.} A call centre receives calls at an average rate of
4 per minute. In a 30-second window, what is the probability of
receiving 3 or more calls?

\begin{center}\rule{0.5\linewidth}{0.5pt}\end{center}

\textbf{Exercise 4.} The time to failure of a device follows an
exponential distribution with rate \(\lambda = 0.02\) per hour. Find:
(a) \(P(T > 100)\), (b) \(E[T]\), (c) \(\text{Var}(T)\).

\begin{center}\rule{0.5\linewidth}{0.5pt}\end{center}

\textbf{Exercise 5.} A manufacturing process produces items whose length
\(X \sim N(75, 100)\) mm (mean 75 mm, variance 100 mm²). The
specification requires length greater than 85 mm. Find \(P(X > 85)\).

\begin{center}\rule{0.5\linewidth}{0.5pt}\end{center}

\textbf{Exercise 6.} Let \(X_1, X_2, \ldots, X_{50}\) be i.i.d.
\(\text{Exp}(1)\) random variables. Their mean is \(\mu = 1\) and
variance is \(\sigma^2 = 1\). Use the CLT to approximate
\(P(\bar{X}_{50} > 1.2)\).

\chapter{Mathematical statistics}\label{mathematical-statistics}

Inference from data

\hfill\break

A production line fills bottles to a stated volume. A sample of 25
bottles shows a mean fill of 497 mL against a target of 500 mL. Is that
a real problem, or just expected random variation? You need a principled
way to decide.

A new catalyst cuts reaction time in a lab test. But the test used only
eight runs. How confident are you that the effect is real and not noise?

You have data. The question is what the data can actually tell you ---
and, just as importantly, what it cannot.

\section{Point estimation}\label{point-estimation}

You observe \(n\) independent realisations \(x_1, x_2, \ldots, x_n\)
from a distribution that has an unknown parameter \(\theta\) --- a
proportion, a rate, a mean. A \textbf{point estimate} \(\hat{\theta}\)
--- read ``theta-hat'' --- is a single number computed from the sample
to serve as your best guess at \(\theta\).

The most systematic method for finding point estimates is
\textbf{maximum likelihood estimation (MLE)}: choose the value of
\(\theta\) that makes the observed data most probable.

\subsection{Setting up the likelihood}\label{setting-up-the-likelihood}

If the \(x_i\) are i.i.d. (independent, identically distributed) from a
distribution with density (or mass function) \(f(x; \theta)\), the
\textbf{likelihood} is the joint probability of the observed data,
viewed as a function of \(\theta\):

\[L(\theta) = \prod_{i=1}^{n} f(x_i; \theta)\]

Read this as: the product, over all \(n\) observations, of the density
evaluated at each \(x_i\). The MLE is the \(\theta\) that maximises
this.

Products are hard to differentiate. Take the natural log --- it's
monotone, so the maximum is preserved. The \textbf{log-likelihood} is:

\[\ell(\theta) = \log L(\theta) = \sum_{i=1}^{n} \log f(x_i;\theta)\]

Read ``\(\ell(\theta)\)'' as ``ell of theta''. Now differentiate and set
equal to zero.

\subsection{MLE for the Bernoulli
distribution}\label{mle-for-the-bernoulli-distribution}

Let \(X_i \sim \text{Bernoulli}(p)\): each observation is 1 (success) or
0 (failure), with unknown success probability \(p\). Then
\(f(x; p) = p^x (1-p)^{1-x}\), so:

\[\ell(p) = \sum_{i=1}^n \left[ x_i \log p + (1-x_i)\log(1-p) \right]
= \left(\sum x_i\right)\log p + \left(n - \sum x_i\right)\log(1-p)\]

Let \(S = \sum x_i\) (the count of successes). Differentiate with
respect to \(p\) and set to zero:

\[\frac{d\ell}{dp} = \frac{S}{p} - \frac{n - S}{1-p} = 0\]

Multiply through by \(p(1-p)\):

\[S(1-p) - (n-S)p = 0 \implies S - Sp - np + Sp = 0 \implies S = np\]

\[\boxed{\hat{p} = \frac{S}{n} = \bar{X}}\]

The MLE is the sample proportion. This confirms the intuition.

\subsection{MLE for the Exponential
distribution}\label{mle-for-the-exponential-distribution}

Let \(X_i \sim \text{Exp}(\lambda)\): density
\(f(x;\lambda) = \lambda e^{-\lambda x}\) for \(x \geq 0\). The
log-likelihood is:

\[\ell(\lambda) = \sum_{i=1}^n \left[\log\lambda - \lambda x_i\right]
= n\log\lambda - \lambda \sum_{i=1}^n x_i\]

Differentiate and set to zero:

\[\frac{d\ell}{d\lambda} = \frac{n}{\lambda} - \sum x_i = 0
\implies \lambda = \frac{n}{\sum x_i}\]

\[\boxed{\hat{\lambda} = \frac{1}{\bar{X}}}\]

The MLE of the rate is the reciprocal of the sample mean. If the mean
lifetime in the sample is 5 hours, you estimate the failure rate as 0.2
failures per hour.

\subsection{Unbiasedness and
consistency}\label{unbiasedness-and-consistency}

A point estimate is \textbf{unbiased} if \(E[\hat{\theta}] = \theta\)
--- the estimate is correct on average. The sample mean \(\bar{X}\) is
unbiased for \(\mu\); the sample variance
\(s^2 = \frac{1}{n-1}\sum(X_i - \bar{X})^2\) is unbiased for
\(\sigma^2\) (the \(n-1\) denominator matters here --- the \(n\) version
is biased). The reason is that \(\bar{X}\) is estimated from the same
data: once the mean is fixed, only \(n-1\) of the deviations
\((X_i - \bar{X})\) are free to vary, so the sum of squared deviations
has \(n-1\) degrees of freedom, not \(n\).

An estimate is \textbf{consistent} if \(\hat{\theta} \to \theta\) in
probability as \(n \to \infty\) --- more data drives the estimate to the
truth. The MLE is consistent under mild regularity conditions.

\begin{tcolorbox}[enhanced jigsaw, opacityback=0, colback=white, breakable, colframe=quarto-callout-note-color-frame, arc=.35mm, left=2mm, toprule=.15mm, leftrule=.75mm, bottomrule=.15mm, rightrule=.15mm]
\begin{minipage}[t]{5.5mm}
\textcolor{quarto-callout-note-color}{\faInfo}
\end{minipage}%
\begin{minipage}[t]{\textwidth - 5.5mm}

\textbf{Why MLE works}

The likelihood is the probability of your data given \(\theta\).
Maximising it over \(\theta\) answers: ``which parameter value would
make this data least surprising?'' That question has a principled answer
even when no other method is obvious --- just write down the
distribution, form the product, and differentiate.

\end{minipage}%
\end{tcolorbox}

\section{Confidence intervals}\label{confidence-intervals-1}

A point estimate is a single number. It conceals the uncertainty. A
\textbf{confidence interval} (CI) pairs the estimate with a margin that
quantifies how much the estimate might be off.

\subsection{The problem with small
samples}\label{the-problem-with-small-samples}

If \(\sigma\) (the population standard deviation) were known, the
sampling distribution of \(\bar{X}\) is exactly \(N(\mu, \sigma^2/n)\),
and you could build an interval using the standard normal
\(z\)-distribution. But \(\sigma\) is almost never known --- you
estimate it from the sample as \(s\). When \(n\) is small, this extra
source of uncertainty matters.

Student's \(t\)-distribution handles this. The statistic

\[T = \frac{\bar{X} - \mu}{s/\sqrt{n}}\]

follows a \(t\)-distribution with \(\nu = n-1\) degrees of freedom.

Read ``\(s/\sqrt{n}\)'' as the \textbf{standard error of the mean} ---
the typical spread of sample means around the true mean. The
\(t\)-distribution looks like the standard normal but with heavier
tails; as \(\nu \to \infty\) it converges to the normal.

\subsection{\texorpdfstring{Constructing a 95\% CI for
\(\mu\)}{Constructing a 95\% CI for \textbackslash mu}}\label{constructing-a-95-ci-for-mu}

Let \(t_{\alpha/2,\,n-1}\) be the upper \(\alpha/2\) critical value of
the \(t\)-distribution with \(n-1\) degrees of freedom --- the value
such that \(P(T > t_{\alpha/2,\,n-1}) = \alpha/2\). For a 95\% CI,
\(\alpha = 0.05\).

Inverting \(|T| \leq t_{\alpha/2,\,n-1}\) gives the interval.
Substituting \(T = (\bar{X} - \mu)/(s/\sqrt{n})\) and expanding the
absolute value:

\[-t_{\alpha/2,\,n-1} \leq \frac{\bar{X} - \mu}{s/\sqrt{n}} \leq t_{\alpha/2,\,n-1}\]

Multiply through by \(s/\sqrt{n}\) and rearrange to isolate \(\mu\):

\[\bar{X} - t_{\alpha/2,\,n-1}\cdot\frac{s}{\sqrt{n}}
  \;\leq\; \mu \;\leq\;
  \bar{X} + t_{\alpha/2,\,n-1}\cdot\frac{s}{\sqrt{n}}\]

Written compactly:

\[\bar{X} \pm t_{\alpha/2,\,n-1} \cdot \frac{s}{\sqrt{n}}\]

This is the 95\% confidence interval for \(\mu\).

\textbf{Worked example.} A materials lab measures Young's modulus for 10
specimens of a polymer. The sample mean is \(\bar{x} = 3.42\) GPa and
the sample standard deviation is \(s = 0.18\) GPa. The \(t\)-table gives
\(t_{0.025,\,9} = 2.262\).

\[\text{CI} = 3.42 \pm 2.262 \times \frac{0.18}{\sqrt{10}}
= 3.42 \pm 2.262 \times 0.0569
= 3.42 \pm 0.129\]

The 95\% CI is \((3.29,\; 3.55)\) GPa.

\subsection{What a confidence interval means --- and what it does
not}\label{what-a-confidence-interval-means-and-what-it-does-not}

The 95\% CI does \textbf{not} mean ``there is a 95\% probability that
\(\mu\) lies in this interval.'' After the data are observed, \(\mu\) is
either in the interval or it isn't --- there is no probability involved.

The correct interpretation: if you repeated this experiment many times
and built a 95\% CI each time, about 95\% of those intervals would
contain the true \(\mu\). The \emph{procedure} has 95\% coverage. Any
particular interval is just one instance of that procedure.

\section{Hypothesis testing}\label{hypothesis-testing-1}

Point estimates and CIs describe what the data show. Hypothesis testing
asks a sharper question: is there enough evidence to reject a specific
claim?

\subsection{The framework}\label{the-framework}

The \textbf{null hypothesis} \(H_0\) is the claim under scrutiny --- the
default, conservative position. The \textbf{alternative hypothesis}
\(H_1\) is what you'd conclude if \(H_0\) is implausible given the data.

A \textbf{test statistic} is a function of the data whose distribution
under \(H_0\) is known. You compute the observed value, then ask: if
\(H_0\) were true, how often would you see a value this extreme or more
extreme? That probability is the \textbf{p-value}.

If the p-value is below the \textbf{significance level} \(\alpha\)
(typically 0.05), you reject \(H_0\).

Two error types:

\begin{longtable}[]{@{}
  >{\raggedright\arraybackslash}p{(\linewidth - 4\tabcolsep) * \real{0.3333}}
  >{\raggedright\arraybackslash}p{(\linewidth - 4\tabcolsep) * \real{0.3333}}
  >{\raggedright\arraybackslash}p{(\linewidth - 4\tabcolsep) * \real{0.3333}}@{}}
\toprule\noalign{}
\begin{minipage}[b]{\linewidth}\raggedright
\end{minipage} & \begin{minipage}[b]{\linewidth}\raggedright
\(H_0\) true
\end{minipage} & \begin{minipage}[b]{\linewidth}\raggedright
\(H_0\) false
\end{minipage} \\
\midrule\noalign{}
\endhead
\bottomrule\noalign{}
\endlastfoot
Reject \(H_0\) & Type I error (false positive), probability \(\alpha\) &
Correct \\
Fail to reject \(H_0\) & Correct & Type II error (false negative),
probability \(\beta\) \\
\end{longtable}

Choosing \(\alpha = 0.05\) means you're willing to falsely reject a true
\(H_0\) in 5\% of experiments. Reducing \(\alpha\) reduces Type I errors
but increases Type II errors. The power of a test is \(1 - \beta\) ---
the probability of correctly detecting a real effect.

\subsection{\texorpdfstring{One-sample
\(t\)-test}{One-sample t-test}}\label{one-sample-t-test}

\textbf{Hypotheses.} You want to test whether a population mean \(\mu\)
equals a specified value \(\mu_0\):
\[H_0: \mu = \mu_0 \qquad H_1: \mu \neq \mu_0 \quad \text{(two-sided)}\]

\textbf{Test statistic.}

\[t = \frac{\bar{x} - \mu_0}{s / \sqrt{n}}\]

Under \(H_0\) this follows \(t_{n-1}\).

\textbf{Decision rule.} Reject \(H_0\) if \(|t| > t_{\alpha/2,\,n-1}\).

\textbf{Worked example.} A valve is designed to open at 200 kPa. Quality
control samples \(n = 20\) valves. The sample mean opening pressure is
\(\bar{x} = 203.4\) kPa with \(s = 5.2\) kPa. Test \(H_0: \mu = 200\) at
\(\alpha = 0.05\).

\[t = \frac{203.4 - 200}{5.2/\sqrt{20}} = \frac{3.4}{1.163} = 2.923\]

The critical value is \(t_{0.025,\,19} = 2.093\). Since
\(2.923 > 2.093\), reject \(H_0\). There is statistically significant
evidence that the mean opening pressure differs from 200 kPa.

\subsection{\texorpdfstring{Two-sample
\(t\)-test}{Two-sample t-test}}\label{two-sample-t-test}

To compare means \(\mu_1\) and \(\mu_2\) from two independent
populations (assuming equal variances), form the pooled variance:

\[s_p^2 = \frac{(n_1-1)s_1^2 + (n_2-1)s_2^2}{n_1+n_2-2}\]

The test statistic is:

\[t = \frac{\bar{x}_1 - \bar{x}_2}{s_p\sqrt{1/n_1 + 1/n_2}}\]

Under \(H_0: \mu_1 = \mu_2\), this follows \(t_{n_1+n_2-2}\). Reject
when \(|t| > t_{\alpha/2,\,n_1+n_2-2}\).

Welch's \(t\)-test removes the equal-variance assumption by adjusting
the degrees of freedom; it's the default in most software.

\subsection{\texorpdfstring{\(\chi^2\)
goodness-of-fit}{\textbackslash chi\^{}2 goodness-of-fit}}\label{chi2-goodness-of-fit}

When the data are counts in categories, the \(\chi^2\) statistic tests
whether the observed frequencies are consistent with a specified
distribution. For \(k\) categories with observed counts \(O_i\) and
expected counts \(E_i = n p_i\):

\[\chi^2 = \sum_{i=1}^k \frac{(O_i - E_i)^2}{E_i}\]

Under \(H_0\), this follows approximately \(\chi^2_{k-1}\). Reject if
\(\chi^2 > \chi^2_{\alpha,\,k-1}\). The rule of thumb \(E_i \geq 5\) in
every cell ensures the approximation is reliable.

\section{Linear regression}\label{linear-regression-1}

Data rarely have a single variable of interest in isolation. Regression
asks how one variable changes with another --- and how reliably.

\subsection{The model}\label{the-model}

The \textbf{simple linear regression model} is:

\[Y_i = \beta_0 + \beta_1 x_i + \varepsilon_i, \quad i = 1, \ldots, n\]

\(Y_i\) is the response (outcome), \(x_i\) is the predictor (input),
\(\beta_0\) is the \textbf{intercept} (read ``beta-zero'' --- the
expected response when \(x = 0\)), \(\beta_1\) is the \textbf{slope}
(read ``beta-one'' --- the change in expected response per unit increase
in \(x\)), and \(\varepsilon_i\) is the error term, assumed i.i.d.
\(N(0, \sigma^2)\).

\subsection{OLS estimators}\label{ols-estimators}

\textbf{Ordinary least squares} minimises the \textbf{sum of squared
residuals}:

\[\text{SS}_{res} = \sum_{i=1}^n (y_i - \hat{y}_i)^2
= \sum_{i=1}^n (y_i - \hat{\beta}_0 - \hat{\beta}_1 x_i)^2\]

Read ``\(\text{SS}_{res}\)'' as ``residual sum of squares'' --- the
total vertical distance squared between the data points and the fitted
line.

Differentiate \(\text{SS}_{res}\) with respect to \(\hat{\beta}_0\) and
\(\hat{\beta}_1\) and set both partial derivatives to zero. This gives
the \textbf{normal equations}:

\[\frac{\partial\,\text{SS}_{res}}{\partial\hat{\beta}_0} = -2\sum(y_i - \hat{\beta}_0 - \hat{\beta}_1 x_i) = 0\]

\[\frac{\partial\,\text{SS}_{res}}{\partial\hat{\beta}_1} = -2\sum x_i(y_i - \hat{\beta}_0 - \hat{\beta}_1 x_i) = 0\]

From the first equation: distribute the sum and use
\(\sum_{i=1}^n 1 = n\):

\[\sum(y_i - \hat{\beta}_0 - \hat{\beta}_1 x_i) = 0
\implies n\bar{y} - n\hat{\beta}_0 - \hat{\beta}_1 \cdot n\bar{x} = 0\]

Divide by \(n\): \(\hat{\beta}_0 = \bar{y} - \hat{\beta}_1 \bar{x}\).

Substituting into the second and solving:

\[\boxed{\hat{\beta}_1 = \frac{\sum_{i=1}^n (x_i - \bar{x})(y_i - \bar{y})}{\sum_{i=1}^n (x_i - \bar{x})^2}
= \frac{S_{xy}}{S_{xx}}}\]

\[\boxed{\hat{\beta}_0 = \bar{y} - \hat{\beta}_1 \bar{x}}\]

where \(S_{xy} = \sum(x_i - \bar{x})(y_i - \bar{y})\) and
\(S_{xx} = \sum(x_i - \bar{x})^2\).

\subsection{\texorpdfstring{\(R^2\) --- coefficient of
determination}{R\^{}2 --- coefficient of determination}}\label{r2-coefficient-of-determination}

Define the \textbf{total sum of squares}
\(\text{SS}_{tot} = \sum(y_i - \bar{y})^2\) and the \textbf{residual sum
of squares} \(\text{SS}_{res} = \sum(y_i - \hat{y}_i)^2\). The
coefficient of determination is:

\[R^2 = 1 - \frac{\text{SS}_{res}}{\text{SS}_{tot}}\]

\(R^2 = 1\) means the line passes through every data point. \(R^2 = 0\)
means the line explains none of the variation --- it is no better than
simply predicting \(\bar{y}\) for every observation. In simple linear
regression, \(R^2 = r^2\) where \(r\) is the Pearson correlation
coefficient.

\subsection{\texorpdfstring{\(t\)-test for
\(\beta_1\)}{t-test for \textbackslash beta\_1}}\label{t-test-for-beta_1}

Under the model assumptions, the standardised slope estimate

\[t = \frac{\hat{\beta}_1 - 0}{s_e / \sqrt{S_{xx}}}\]

follows \(t_{n-2}\) under \(H_0: \beta_1 = 0\), where
\(s_e^2 = \text{SS}_{res}/(n-2)\) is the estimated error variance. The
denominator is \(n-2\) rather than \(n-1\) because two parameters
(\(\beta_0\) and \(\beta_1\)) are estimated from the data, consuming two
degrees of freedom. Reject \(H_0\) if \(|t| > t_{\alpha/2,\,n-2}\).
Failing to reject means there is no statistically significant linear
relationship between \(x\) and \(Y\).

\subsection{Residual analysis}\label{residual-analysis}

The regression assumptions --- linearity, constant variance, normality
of errors, independence --- are checked through the residuals
\(\hat{\varepsilon}_i = y_i - \hat{y}_i\).

Plot residuals against fitted values \(\hat{y}_i\):

\begin{itemize}
\tightlist
\item
  \textbf{Random scatter around zero} --- assumptions satisfied.
\item
  \textbf{Curved pattern} --- linearity is violated; consider a
  transformation or a polynomial term.
\item
  \textbf{Fan shape (variance increasing with \(\hat{y}_i\))} ---
  heteroscedasticity; consider a log or square-root transformation of
  \(Y\).
\item
  \textbf{Outliers} --- single points far from zero; investigate
  individually.
\end{itemize}

A normal Q-Q plot of residuals checks the normality assumption. Points
falling close to the diagonal line support the assumption.

\section{One-way ANOVA}\label{one-way-anova}

When you want to compare means across three or more groups, running
multiple \(t\)-tests inflates the Type I error rate. One-way ANOVA
(analysis of variance) tests all groups simultaneously.

\subsection{The F-test}\label{the-f-test}

Suppose you have \(k\) groups, with \(n_j\) observations in group \(j\).
Let \(\bar{y}_{j\cdot}\) denote the mean of group \(j\) and \(\bar{y}\)
the overall mean. Define:

\[\text{SS}_{between} = \sum_{j=1}^k n_j (\bar{y}_{j\cdot} - \bar{y})^2
\qquad \text{(variation explained by group membership)}\]

\[\text{SS}_{within} = \sum_{j=1}^k \sum_{i=1}^{n_j} (y_{ij} - \bar{y}_{j\cdot})^2
\qquad \text{(variation within groups)}\]

The \textbf{mean squares} are
\(\text{MS}_{between} = \text{SS}_{between}/(k-1)\) and
\(\text{MS}_{within} = \text{SS}_{within}/(N-k)\), where
\(N = \sum n_j\).

The \textbf{F-statistic} is:

\[F = \frac{\text{MS}_{between}}{\text{MS}_{within}}\]

Under \(H_0: \mu_1 = \mu_2 = \cdots = \mu_k\), \(F\) follows
\(F_{k-1,\,N-k}\). Reject \(H_0\) if \(F > F_{\alpha,\,k-1,\,N-k}\). A
large \(F\) means between-group variation is large relative to
within-group variation --- the groups genuinely differ.

\begin{tcolorbox}[enhanced jigsaw, opacityback=0, colback=white, breakable, colframe=quarto-callout-note-color-frame, arc=.35mm, left=2mm, toprule=.15mm, leftrule=.75mm, bottomrule=.15mm, rightrule=.15mm]
\begin{minipage}[t]{5.5mm}
\textcolor{quarto-callout-note-color}{\faInfo}
\end{minipage}%
\begin{minipage}[t]{\textwidth - 5.5mm}

\textbf{Why F instead of multiple \(t\)-tests?}

With \(k = 3\) groups, running all pairwise \(t\)-tests at
\(\alpha = 0.05\) gives a familywise error rate of approximately
\(1 - (0.95)^3 \approx 0.14\). You'd falsely reject a true null about
14\% of the time rather than 5\%. ANOVA controls this by testing all
groups in a single F-test. If ANOVA rejects \(H_0\), post-hoc pairwise
tests (Tukey HSD, Bonferroni) then identify which groups differ.

\end{minipage}%
\end{tcolorbox}

\textbf{Engineering application.} Three machine settings --- A, B, C ---
produce polymer film. Quality engineers measure tensile strength (MPa)
from five runs at each setting:

\begin{longtable}[]{@{}rrr@{}}
\toprule\noalign{}
Setting A & Setting B & Setting C \\
\midrule\noalign{}
\endhead
\bottomrule\noalign{}
\endlastfoot
41.2 & 38.5 & 44.1 \\
40.8 & 39.2 & 43.8 \\
42.1 & 37.9 & 45.0 \\
41.5 & 38.8 & 44.3 \\
40.4 & 39.6 & 43.4 \\
\end{longtable}

Group means: \(\bar{y}_{A} = 41.2\), \(\bar{y}_{B} = 38.8\),
\(\bar{y}_{C} = 44.12\). Overall mean \(\bar{y} = 41.37\).

\(\text{SS}_{between} = 5(41.2-41.37)^2 + 5(38.8-41.37)^2 + 5(44.12-41.37)^2\)
\(= 5(0.0289) + 5(6.6049) + 5(7.5625) = 0.1445 + 33.02 + 37.81 = 70.98\)

\(\text{MS}_{between} = 70.98/2 = 35.49\)

Computing \(\text{SS}_{within}\) from deviations within each group gives
approximately \(4.61\), so \(\text{MS}_{within} = 4.61/12 = 0.384\).

\[F = \frac{35.49}{0.384} = 92.4\]

The critical value \(F_{0.05,\,2,\,12} = 3.89\). Since
\(92.4 \gg 3.89\), reject \(H_0\). The machine settings produce
significantly different tensile strengths --- the effect is not noise.

\section{Where this goes}\label{where-this-goes-36}

The methods here assume a parametric form for the distributions ---
Bernoulli, exponential, normal. \textbf{Non-parametric tests} (Wilcoxon,
Mann-Whitney, Kruskal-Wallis) make weaker assumptions at the cost of
some power; they become relevant when distributions are heavily skewed
or sample sizes are very small. \textbf{Bayesian inference} reframes the
problem entirely: instead of asking ``what does the data tell us about
\(\theta\)?'' it asks ``given a prior belief about \(\theta\) and the
data, what should our updated belief be?'' Both extensions build on
exactly the likelihood functions you derived here.

Multiple linear regression extends the single-predictor model to
\(Y = \beta_0 + \beta_1 x_1 + \cdots + \beta_p x_p + \varepsilon\). The
OLS estimator in matrix form is
\(\hat{\boldsymbol{\beta}} = (X^TX)^{-1}X^T\mathbf{y}\), which requires
the linear algebra from Volume 7's linear algebra section. Generalised
linear models (GLMs) then extend regression to non-normal responses ---
logistic regression for binary outcomes, Poisson regression for counts
--- connecting directly to the MLE framework you already have.

\begin{tcolorbox}[enhanced jigsaw, opacityback=0, colback=white, breakable, colframe=quarto-callout-tip-color-frame, arc=.35mm, left=2mm, toprule=.15mm, leftrule=.75mm, bottomrule=.15mm, rightrule=.15mm]
\begin{minipage}[t]{5.5mm}
\textcolor{quarto-callout-tip-color}{\faLightbulb}
\end{minipage}%
\begin{minipage}[t]{\textwidth - 5.5mm}

\textbf{Where this shows up}

\begin{itemize}
\tightlist
\item
  \textbf{Reliability engineering:} MLE estimates failure rates from
  field data; confidence intervals bound the estimate; hypothesis tests
  determine whether a design change made a difference.
\item
  \textbf{Clinical trials:} two-sample \(t\)-tests compare treatment and
  control groups; sample size is set to achieve target power.
\item
  \textbf{Process control:} \(t\)-tests on subgroup means detect shifts
  in a manufacturing process mean; ANOVA compares multiple line
  configurations simultaneously.
\item
  \textbf{Geophysics:} linear regression of seismic travel time against
  distance estimates subsurface velocity; the \(t\)-test for \(\beta_1\)
  determines whether the relationship is real.
\item
  \textbf{Machine learning:} the OLS solution is closed-form linear
  regression; \(R^2\) and residual analysis evaluate model fit before
  any regularisation is added.
\end{itemize}

\end{minipage}%
\end{tcolorbox}

\section{Exercises}\label{exercises-47}

These are puzzles. Each has a clean numerical answer. The work is in
setting up the likelihood, test statistic, or normal equations correctly
before computing.

\begin{center}\rule{0.5\linewidth}{0.5pt}\end{center}

\textbf{Exercise 1.} An engineer records the time to failure (hours) for
five identical electronic components under accelerated testing:

\[x_1 = 820, \quad x_2 = 1\,040, \quad x_3 = 650, \quad x_4 = 930, \quad x_5 = 760\]

Assuming lifetimes are Exponential(\(\lambda\)), find \(\hat{\lambda}\)
(the MLE) and its standard error. Use the fact that for the Exponential
MLE, \(\text{SE}(\hat{\lambda}) = \hat{\lambda}/\sqrt{n}\).

\begin{center}\rule{0.5\linewidth}{0.5pt}\end{center}

\textbf{Exercise 2.} A tensile test on 16 specimens of an aluminium
alloy yields a sample mean yield strength of \(\bar{x} = 48.2\) MPa and
sample standard deviation \(s = 6.4\) MPa. Construct a 95\% confidence
interval for the population mean yield strength. Use
\(t_{0.025,\,15} = 2.131\).

\begin{center}\rule{0.5\linewidth}{0.5pt}\end{center}

\textbf{Exercise 3.} A CNC milling process is specified to produce
shafts with a mean diameter of 50.00 mm. A quality technician measures a
random sample of 25 shafts and finds \(\bar{x} = 47.8\) mm and
\(s = 5.1\) mm. Test the null hypothesis \(H_0: \mu = 50\) against
\(H_1: \mu \neq 50\) at significance level \(\alpha = 0.05\). Use
\(t_{0.025,\,24} = 2.064\). State your conclusion.

\begin{center}\rule{0.5\linewidth}{0.5pt}\end{center}

\textbf{Exercise 4.} Two aluminium alloys are compared for tensile
strength. Alloy 1 (\(n_1 = 10\), \(\bar{x}_1 = 312\) MPa, \(s_1 = 14\)
MPa) and Alloy 2 (\(n_2 = 12\), \(\bar{x}_2 = 298\) MPa, \(s_2 = 18\)
MPa). Assuming equal population variances, test \(H_0: \mu_1 = \mu_2\)
against \(H_1: \mu_1 \neq \mu_2\) at \(\alpha = 0.05\). Use
\(t_{0.025,\,20} = 2.086\).

\begin{center}\rule{0.5\linewidth}{0.5pt}\end{center}

\textbf{Exercise 5.} A process engineer suspects that curing temperature
\(x\) (°C) affects bond strength \(Y\) (N/mm²). Five runs are recorded:

\begin{longtable}[]{@{}rrrrrr@{}}
\toprule\noalign{}
\(x\) & 120 & 130 & 140 & 150 & 160 \\
\midrule\noalign{}
\endhead
\bottomrule\noalign{}
\endlastfoot
\(Y\) & 8.1 & 9.4 & 10.8 & 11.5 & 12.9 \\
\end{longtable}

\begin{enumerate}
\def\labelenumi{(\alph{enumi})}
\tightlist
\item
  Compute the OLS estimates \(\hat{\beta}_0\) and \(\hat{\beta}_1\).
\item
  Compute \(R^2\).
\item
  Test \(H_0: \beta_1 = 0\) at \(\alpha = 0.05\), given
  \(t_{0.025,\,3} = 3.182\).
\end{enumerate}

\begin{center}\rule{0.5\linewidth}{0.5pt}\end{center}

\textbf{Exercise 6.} Three CNC machine settings produce the following
daily output yields (units per shift):

\begin{longtable}[]{@{}rrr@{}}
\toprule\noalign{}
Setting 1 & Setting 2 & Setting 3 \\
\midrule\noalign{}
\endhead
\bottomrule\noalign{}
\endlastfoot
102 & 98 & 115 \\
105 & 101 & 112 \\
99 & 97 & 118 \\
103 & 100 & 114 \\
\end{longtable}

Perform a one-way ANOVA F-test for equality of means at
\(\alpha = 0.05\). Use \(F_{0.05,\,2,\,9} = 4.26\). State whether
setting affects yield.

\part{Volume 8 --- Upper-year engineering mathematics}

\chapter{Volume 8 --- Upper-year engineering
mathematics}\label{volume-8-upper-year-engineering-mathematics-1}

Volume 8 is where the mathematics stops appearing as a course title and
starts appearing as the working language of upper-year engineering.

By this point the reader should already have the Vol 7 toolkit: ODEs,
linear algebra, vector calculus, transforms, PDEs, numerical methods,
optimisation, and probability. The job now is different. Not ``learn the
next theorem'' in isolation, but learn how mathematical structures
reappear inside controls, signals, transport, simulation, estimation,
reliability, and design.

This volume is anchored to the kinds of mathematics that show up across
years 3 and 4 of University of Alberta engineering programmes: feedback
loops, sampled systems, continuum laws, discretised models, noisy data,
uncertain failure, and constrained nonlinear tradeoffs.

Optional viewpoints from other domains are built in on purpose. The same
tools matter in atmospheric science, quantitative finance, machine
learning, geophysics, remote sensing, and operations research. Those
side doors are not distractions; they help the reader see the
mathematics as a transferable way of thinking instead of a
discipline-specific dialect.

The pipeline for this volume is different too. The authoring job is not
simply to continue adding topics. It is to test each chapter from
several angles at once:

\begin{itemize}
\tightlist
\item
  mathematical dependence: does this genuinely grow out of Volume 7?
\item
  engineering honesty: does this match how upper-year students meet the
  tool?
\item
  student readability: does the embedded context obscure the
  mathematics?
\item
  visual design: can the system be seen moving, not just described?
\item
  cross-domain transfer: does the optional viewpoint reveal the same
  structure?
\end{itemize}

\section{Chapters}\label{chapters-6}

\begin{longtable}[]{@{}
  >{\raggedright\arraybackslash}p{(\linewidth - 4\tabcolsep) * \real{0.3333}}
  >{\raggedright\arraybackslash}p{(\linewidth - 4\tabcolsep) * \real{0.2593}}
  >{\raggedright\arraybackslash}p{(\linewidth - 4\tabcolsep) * \real{0.4074}}@{}}
\toprule\noalign{}
\begin{minipage}[b]{\linewidth}\raggedright
Chapter
\end{minipage} & \begin{minipage}[b]{\linewidth}\raggedright
Title
\end{minipage} & \begin{minipage}[b]{\linewidth}\raggedright
Core move
\end{minipage} \\
\midrule\noalign{}
\endhead
\bottomrule\noalign{}
\endlastfoot
1 & Control, feedback, and stability & Turn dynamics into decisions
about response and robustness \\
2 & Discrete-time systems and signal processing & Move from continuous
models to sampled systems and filters \\
3 & Continuum mechanics and transport & Express materials, fluids, heat,
and diffusion as field laws \\
4 & Computational methods for engineering models & Turn governing
equations into mesh-based computations \\
5 & Estimation, inverse problems, and filtering & Recover hidden states
and parameters from noisy observations \\
6 & Reliability, stochastic systems, and quality & Model failure,
queues, variability, and risk over time \\
7 & Nonlinear optimisation for design and operations & Choose the best
feasible design when tradeoffs are real \\
\end{longtable}

\section{U of A anchors}\label{u-of-a-anchors}

This volume is not a copy of any one programme calendar, but it is
anchored to the kinds of mathematical work that appear in current
upper-year UAlberta engineering courses: control systems, signals and
systems, process dynamics, engineering statistics, process data
analytics, and simulation-heavy design.

\section{Pipeline notes}\label{pipeline-notes}

The first pipeline pass on Volume 8 produced three main decisions:

\begin{enumerate}
\def\labelenumi{\arabic{enumi}.}
\tightlist
\item
  Keep the centre of gravity on engineering practice, not on a second
  pass through generic advanced mathematics.
\item
  Make every chapter carry at least one optional viewpoint from another
  domain, but only where the structure is truly parallel.
\item
  Treat visualisation as mandatory, because upper-year mathematics is
  now about systems, flows, feedback loops, noisy measurements, and
  tradeoff surfaces.
\end{enumerate}

\section{Optional viewpoints}\label{optional-viewpoints}

\begin{longtable}[]{@{}
  >{\raggedright\arraybackslash}p{(\linewidth - 2\tabcolsep) * \real{0.4400}}
  >{\raggedright\arraybackslash}p{(\linewidth - 2\tabcolsep) * \real{0.5600}}@{}}
\toprule\noalign{}
\begin{minipage}[b]{\linewidth}\raggedright
Viewpoint
\end{minipage} & \begin{minipage}[b]{\linewidth}\raggedright
What it adds
\end{minipage} \\
\midrule\noalign{}
\endhead
\bottomrule\noalign{}
\endlastfoot
Hard sciences & Stability, waves, transport, and inverse problems in
physical systems \\
Computing and data & Signal processing, state estimation, control
software, and learned models \\
Finance and quant & Filtering, stochastic evolution, optimisation under
uncertainty \\
Geography and environment & Flow, dispersion, remote sensing inversion,
and spatial risk \\
\end{longtable}

\chapter{Control, feedback, and
stability}\label{control-feedback-and-stability}

Making a dynamic system behave

\hfill\break

A motor overshoots the target speed. A drone wobbles after a gust. A
reactor temperature drifts when the feed changes. In each case the
system is already dynamic before anyone touches the controller. Feedback
is the act of shaping that behaviour instead of merely watching it
happen.

Control is not the same thing as solving a differential equation.
Solving the equation tells you what the system will do. Control asks
whether that behaviour is acceptable, how fast it settles, how much it
overshoots, and how sensitive it is to disturbances and modelling error.

This chapter keeps the focus on the central move. You measure the output
\(y(t)\), compare it to the target \(r(t)\), form the error
\(e(t) = r(t) - y(t)\), and choose a control input \(u(t)\) that changes
the dynamics in your favour. Once you see that structure clearly,
stability and transient response stop looking like side remarks and
start looking like the point of the whole subject.

\begin{center}\rule{0.5\linewidth}{0.5pt}\end{center}

\section{The control problem}\label{the-control-problem}

Start with four signals.

\begin{itemize}
\tightlist
\item
  \(r(t)\) is the \textbf{reference}: what you want.
\item
  \(y(t)\) is the \textbf{output}: what the system is actually doing.
\item
  \(e(t) = r(t) - y(t)\) is the \textbf{error}.
\item
  \(u(t)\) is the \textbf{control input} you apply to the plant.
\end{itemize}

The \textbf{plant} is the system being controlled: a motor, furnace,
vehicle, converter, or biochemical process. In Laplace-transform
language the plant is often represented by a transfer function \(G(s)\),
read ``G of s.'' The controller is represented by \(C(s)\), read ``C of
s.''

In \textbf{open loop}, you choose \(u(t)\) without using the measured
output. In \textbf{closed loop}, you feed back the measured output and
let the error signal drive the control action. That is the essential
difference.

\textbf{Unity feedback} means the measured output \(Y(s)\) is fed back
directly into the comparison point with no additional gain element in
the feedback path. The block-diagram algebra is:

\[E(s) = R(s) - Y(s)\] \[U(s) = C(s)E(s)\] \[Y(s) = G(s)U(s)\]

Substitute in sequence:

\[Y(s) = G(s)C(s)\bigl(R(s) - Y(s)\bigr)\]

Rearrange:

\[\bigl(1 + C(s)G(s)\bigr)Y(s) = C(s)G(s)R(s)\]

So the closed-loop transfer function is

\[T(s) = \frac{Y(s)}{R(s)} = \frac{C(s)G(s)}{1 + C(s)G(s)}\]

That denominator matters more than it first appears to. The roots of
\(1 + C(s)G(s) = 0\) are the \textbf{closed-loop poles}. They decide
whether the response decays, oscillates, grows, or diverges.

\emph{Adjust the gain K below. Watch the step response, pole location,
and error readouts change simultaneously.}

\begin{Shaded}
\begin{Highlighting}[]
\NormalTok{//| echo: false}
\NormalTok{closedLoopViz = \{}
\NormalTok{  const d3 = await require("d3@7");}

\NormalTok{  // {-}{-}{-}{-} Layout constants {-}{-}{-}{-}}
\NormalTok{  const W = 640, H = 480;}
\NormalTok{  const BLUE = "\#2563eb", GREEN = "\#059669", RED = "\#ef4444", AMBER = "\#f59e0b";}
\NormalTok{  const GREY\_DASH = "\#9ca3af";}

\NormalTok{  const container = d3.create("div")}
\NormalTok{    .style("font{-}family", "sans{-}serif")}
\NormalTok{    .style("max{-}width", W + "px");}

\NormalTok{  // K slider}
\NormalTok{  const sliderRow = container.append("div")}
\NormalTok{    .style("display", "flex")}
\NormalTok{    .style("align{-}items", "center")}
\NormalTok{    .style("gap", "0.75rem")}
\NormalTok{    .style("margin{-}bottom", "0.5rem");}

\NormalTok{  sliderRow.append("label")}
\NormalTok{    .style("font{-}size", "0.9rem")}
\NormalTok{    .style("white{-}space", "nowrap")}
\NormalTok{    .text("Gain K:");}

\NormalTok{  const slider = sliderRow.append("input")}
\NormalTok{    .attr("type", "range")}
\NormalTok{    .attr("min", 0.1)}
\NormalTok{    .attr("max", 20)}
\NormalTok{    .attr("step", 0.1)}
\NormalTok{    .attr("value", 2)}
\NormalTok{    .style("flex", "1");}

\NormalTok{  const kLabel = sliderRow.append("span")}
\NormalTok{    .style("font{-}weight", "bold")}
\NormalTok{    .style("min{-}width", "2.5rem");}

\NormalTok{  // Main SVG}
\NormalTok{  const svg = container.append("svg")}
\NormalTok{    .attr("width", W)}
\NormalTok{    .attr("height", H)}
\NormalTok{    .attr("viewBox", \textasciigrave{}0 0 $\{W\} $\{H\}\textasciigrave{})}
\NormalTok{    .attr("aria{-}label", "Closed{-}loop feedback explorer");}

\NormalTok{  // {-}{-}{-}{-} Block diagram (top half, y 0{-}210) {-}{-}{-}{-}}
\NormalTok{  const BD = svg.append("g").attr("transform", "translate(0,0)");}

\NormalTok{  // Draw block diagram elements}
\NormalTok{  // Signal arrows: r {-}\textgreater{} summing junction {-}\textgreater{} controller {-}\textgreater{} plant {-}\textgreater{} y (output)}
\NormalTok{  // feedback path from y back to summing junction}

\NormalTok{  // Positions}
\NormalTok{  const SJ = \{x: 130, y: 80\};   // summing junction centre}
\NormalTok{  const CB = \{x: 260, y: 60, w: 80, h: 40\}; // controller box}
\NormalTok{  const PB = \{x: 420, y: 60, w: 80, h: 40\}; // plant box}
\NormalTok{  const outX = 550;               // output signal x}

\NormalTok{  // r arrow}
\NormalTok{  BD.append("line")}
\NormalTok{    .attr("x1", 50).attr("y1", SJ.y)}
\NormalTok{    .attr("x2", SJ.x {-} 18).attr("y2", SJ.y)}
\NormalTok{    .attr("stroke", "\#333").attr("stroke{-}width", 1.5)}
\NormalTok{    .attr("marker{-}end", "url(\#arrowBlack)");}
\NormalTok{  BD.append("text").attr("x", 30).attr("y", SJ.y {-} 8)}
\NormalTok{    .attr("font{-}size", "13px").attr("fill", "\#333").text("r(t)");}

\NormalTok{  // summing junction circle}
\NormalTok{  BD.append("circle")}
\NormalTok{    .attr("cx", SJ.x).attr("cy", SJ.y).attr("r", 18)}
\NormalTok{    .attr("fill", "none").attr("stroke", "\#333").attr("stroke{-}width", 1.5);}
\NormalTok{  BD.append("text").attr("x", SJ.x {-} 6).attr("y", SJ.y + 5)}
\NormalTok{    .attr("font{-}size", "14px").attr("fill", "\#333").text("+");}
\NormalTok{  BD.append("text").attr("x", SJ.x {-} 5).attr("y", SJ.y + 22)}
\NormalTok{    .attr("font{-}size", "11px").attr("fill", "\#ef4444").text("−");}

\NormalTok{  // e(t) label between SJ and controller}
\NormalTok{  BD.append("line")}
\NormalTok{    .attr("x1", SJ.x + 18).attr("y1", SJ.y)}
\NormalTok{    .attr("x2", CB.x).attr("y2", SJ.y)}
\NormalTok{    .attr("stroke", "\#333").attr("stroke{-}width", 1.5)}
\NormalTok{    .attr("marker{-}end", "url(\#arrowBlack)");}
\NormalTok{  BD.append("text").attr("x", SJ.x + 22).attr("y", SJ.y {-} 8)}
\NormalTok{    .attr("font{-}size", "12px").attr("fill", "\#6b7280").text("e(t)");}

\NormalTok{  // Controller box}
\NormalTok{  BD.append("rect")}
\NormalTok{    .attr("x", CB.x).attr("y", CB.y).attr("width", CB.w).attr("height", CB.h)}
\NormalTok{    .attr("rx", 6).attr("fill", "\#f3f4f6").attr("stroke", "\#9ca3af").attr("stroke{-}width", 1.5);}
\NormalTok{  const kBoxLabel = BD.append("text")}
\NormalTok{    .attr("x", CB.x + CB.w / 2).attr("y", CB.y + CB.h / 2 + 5)}
\NormalTok{    .attr("text{-}anchor", "middle").attr("font{-}size", "14px").attr("fill", BLUE);}

\NormalTok{  // Arrow controller to plant}
\NormalTok{  BD.append("line")}
\NormalTok{    .attr("x1", CB.x + CB.w).attr("y1", SJ.y)}
\NormalTok{    .attr("x2", PB.x).attr("y2", SJ.y)}
\NormalTok{    .attr("stroke", "\#333").attr("stroke{-}width", 1.5)}
\NormalTok{    .attr("marker{-}end", "url(\#arrowBlack)");}
\NormalTok{  BD.append("text").attr("x", CB.x + CB.w + 5).attr("y", SJ.y {-} 8)}
\NormalTok{    .attr("font{-}size", "12px").attr("fill", "\#6b7280").text("u(t)");}

\NormalTok{  // Plant box}
\NormalTok{  BD.append("rect")}
\NormalTok{    .attr("x", PB.x).attr("y", PB.y).attr("width", PB.w).attr("height", PB.h)}
\NormalTok{    .attr("rx", 6).attr("fill", "\#f3f4f6").attr("stroke", "\#9ca3af").attr("stroke{-}width", 1.5);}
\NormalTok{  BD.append("text")}
\NormalTok{    .attr("x", PB.x + PB.w / 2).attr("y", PB.y + PB.h / 2 + 5)}
\NormalTok{    .attr("text{-}anchor", "middle").attr("font{-}size", "13px").attr("fill", "\#374151")}
\NormalTok{    .text("G(s)=1/(5s+1)");}

\NormalTok{  // Arrow plant to output}
\NormalTok{  BD.append("line")}
\NormalTok{    .attr("x1", PB.x + PB.w).attr("y1", SJ.y)}
\NormalTok{    .attr("x2", outX).attr("y2", SJ.y)}
\NormalTok{    .attr("stroke", "\#333").attr("stroke{-}width", 1.5)}
\NormalTok{    .attr("marker{-}end", "url(\#arrowBlack)");}
\NormalTok{  BD.append("text").attr("x", outX + 5).attr("y", SJ.y + 5)}
\NormalTok{    .attr("font{-}size", "13px").attr("fill", "\#333").text("y(t)");}

\NormalTok{  // Feedback path (y {-}\textgreater{} down {-}\textgreater{} left {-}\textgreater{} up {-}\textgreater{} summing junction bottom)}
\NormalTok{  const fbY = SJ.y + 70;}
\NormalTok{  BD.append("polyline")}
\NormalTok{    .attr("points", \textasciigrave{}$\{outX {-} 5\},$\{SJ.y\} $\{outX {-} 5\},$\{fbY\} $\{SJ.x\},$\{fbY\} $\{SJ.x\},$\{SJ.y + 18\}\textasciigrave{})}
\NormalTok{    .attr("fill", "none").attr("stroke", "\#333").attr("stroke{-}width", 1.5)}
\NormalTok{    .attr("marker{-}end", "url(\#arrowBlack)");}
\NormalTok{  BD.append("text").attr("x", (SJ.x + outX) / 2).attr("y", fbY + 16)}
\NormalTok{    .attr("text{-}anchor", "middle").attr("font{-}size", "11px").attr("fill", "\#6b7280")}
\NormalTok{    .text("unity feedback");}

\NormalTok{  // Arrowhead marker def}
\NormalTok{  const defs = svg.append("defs");}
\NormalTok{  defs.append("marker")}
\NormalTok{    .attr("id", "arrowBlack").attr("markerWidth", 8).attr("markerHeight", 8)}
\NormalTok{    .attr("refX", 6).attr("refY", 3).attr("orient", "auto")}
\NormalTok{    .append("path").attr("d", "M0,0 L0,6 L8,3 z").attr("fill", "\#333");}

\NormalTok{  // {-}{-}{-}{-} Step response panel (bottom left, y 220{-}480, x 0{-}320) {-}{-}{-}{-}}
\NormalTok{  const SP\_X = 30, SP\_Y = 200, SP\_W = 280, SP\_H = 200;}
\NormalTok{  const spG = svg.append("g").attr("transform", \textasciigrave{}translate($\{SP\_X\},$\{SP\_Y\})\textasciigrave{});}
\NormalTok{  spG.append("text").attr("x", SP\_W / 2).attr("y", 0)}
\NormalTok{    .attr("text{-}anchor", "middle").attr("font{-}size", "12px").attr("fill", "\#374151")}
\NormalTok{    .text("Step response y(t)");}
\NormalTok{  // Axes}
\NormalTok{  spG.append("line").attr("x1", 30).attr("y1", 10).attr("x2", 30).attr("y2", SP\_H {-} 20)}
\NormalTok{    .attr("stroke", "\#6b7280").attr("stroke{-}width", 1.5);}
\NormalTok{  spG.append("line").attr("x1", 30).attr("y1", SP\_H {-} 20).attr("x2", SP\_W {-} 10).attr("y2", SP\_H {-} 20)}
\NormalTok{    .attr("stroke", "\#6b7280").attr("stroke{-}width", 1.5);}
\NormalTok{  spG.append("text").attr("x", SP\_W {-} 8).attr("y", SP\_H {-} 17)}
\NormalTok{    .attr("font{-}size", "10px").attr("fill", "\#6b7280").text("t");}
\NormalTok{  spG.append("text").attr("x", 0).attr("y", 10)}
\NormalTok{    .attr("font{-}size", "10px").attr("fill", "\#6b7280").text("y");}

\NormalTok{  // Reference line (dashed grey at y=1)}
\NormalTok{  const spXmin = 30, spXmax = SP\_W {-} 10, spYmin = 10, spYmax = SP\_H {-} 20;}
\NormalTok{  const yScale = d3.scaleLinear([0, 1.5], [spYmax, spYmin]);}

\NormalTok{  spG.append("line")}
\NormalTok{    .attr("x1", spXmin).attr("y1", yScale(1))}
\NormalTok{    .attr("x2", spXmax).attr("y2", yScale(1))}
\NormalTok{    .attr("stroke", GREY\_DASH).attr("stroke{-}width", 1).attr("stroke{-}dasharray", "5,3");}
\NormalTok{  spG.append("text").attr("x", spXmax + 2).attr("y", yScale(1) + 4)}
\NormalTok{    .attr("font{-}size", "10px").attr("fill", GREY\_DASH).text("r=1");}

\NormalTok{  // Error shading group and response path}
\NormalTok{  const errorShade = spG.append("path")}
\NormalTok{    .attr("fill", RED).attr("fill{-}opacity", 0.12).attr("stroke", "none");}
\NormalTok{  const responsePath = spG.append("path")}
\NormalTok{    .attr("fill", "none").attr("stroke", BLUE).attr("stroke{-}width", 2);}

\NormalTok{  // {-}{-}{-}{-} Pole location panel (bottom right, y 220{-}480, x 330{-}640) {-}{-}{-}{-}}
\NormalTok{  const PP\_X = 340, PP\_Y = 200, PP\_W = 260, PP\_H = 200;}
\NormalTok{  const ppG = svg.append("g").attr("transform", \textasciigrave{}translate($\{PP\_X\},$\{PP\_Y\})\textasciigrave{});}
\NormalTok{  ppG.append("text").attr("x", PP\_W / 2).attr("y", 0)}
\NormalTok{    .attr("text{-}anchor", "middle").attr("font{-}size", "12px").attr("fill", "\#374151")}
\NormalTok{    .text("Pole location (real axis)");}

\NormalTok{  // s{-}plane: x{-}axis only, range {-}6 to 0.5}
\NormalTok{  const poleXscale = d3.scaleLinear([{-}6, 0.5], [10, PP\_W {-} 10]);}
\NormalTok{  const poleY = 80;}
\NormalTok{  // Stable region shading}
\NormalTok{  ppG.append("rect")}
\NormalTok{    .attr("x", 10).attr("y", 20)}
\NormalTok{    .attr("width", poleXscale(0) {-} 10).attr("height", 100)}
\NormalTok{    .attr("fill", GREEN).attr("fill{-}opacity", 0.08);}
\NormalTok{  // Axis}
\NormalTok{  ppG.append("line")}
\NormalTok{    .attr("x1", 10).attr("y1", poleY).attr("x2", PP\_W {-} 10).attr("y2", poleY)}
\NormalTok{    .attr("stroke", "\#6b7280").attr("stroke{-}width", 1.5);}
\NormalTok{  // Imaginary axis (vertical dashed)}
\NormalTok{  ppG.append("line")}
\NormalTok{    .attr("x1", poleXscale(0)).attr("y1", 10)}
\NormalTok{    .attr("x2", poleXscale(0)).attr("y2", 150)}
\NormalTok{    .attr("stroke", "\#9ca3af").attr("stroke{-}width", 1).attr("stroke{-}dasharray", "4,3");}
\NormalTok{  ppG.append("text").attr("x", poleXscale(0) + 2).attr("y", 16)}
\NormalTok{    .attr("font{-}size", "9px").attr("fill", "\#9ca3af").text("stability boundary");}

\NormalTok{  // Axis ticks}
\NormalTok{  [{-}5,{-}4,{-}3,{-}2,{-}1,0].forEach(v =\textgreater{} \{}
\NormalTok{    const x = poleXscale(v);}
\NormalTok{    ppG.append("line").attr("x1", x).attr("y1", poleY {-} 5)}
\NormalTok{      .attr("x2", x).attr("y2", poleY + 5).attr("stroke", "\#9ca3af");}
\NormalTok{    ppG.append("text").attr("x", x).attr("y", poleY + 18)}
\NormalTok{      .attr("text{-}anchor", "middle").attr("font{-}size", "10px").attr("fill", "\#6b7280")}
\NormalTok{      .text(v);}
\NormalTok{  \});}
\NormalTok{  ppG.append("text").attr("x", PP\_W {-} 8).attr("y", poleY + 5)}
\NormalTok{    .attr("font{-}size", "10px").attr("fill", "\#6b7280").text("Re(s)");}
\NormalTok{  ppG.append("text").attr("x", poleXscale({-}3)).attr("y", 16)}
\NormalTok{    .attr("text{-}anchor", "middle").attr("font{-}size", "10px").attr("fill", GREEN)}
\NormalTok{    .text("stable region");}

\NormalTok{  const poleCircle = ppG.append("circle")}
\NormalTok{    .attr("cy", poleY).attr("r", 8)}
\NormalTok{    .attr("fill", BLUE).attr("stroke", "\#1e40af").attr("stroke{-}width", 1.5);}
\NormalTok{  const poleLabelEl = ppG.append("text")}
\NormalTok{    .attr("text{-}anchor", "middle").attr("y", poleY {-} 15)}
\NormalTok{    .attr("font{-}size", "11px").attr("fill", BLUE).attr("font{-}weight", "bold");}

\NormalTok{  // Readouts}
\NormalTok{  const readoutG = svg.append("g").attr("transform", \textasciigrave{}translate($\{PP\_X\}, $\{PP\_Y + 140\})\textasciigrave{});}
\NormalTok{  const settlingTxt = readoutG.append("text")}
\NormalTok{    .attr("font{-}size", "12px").attr("fill", AMBER);}
\NormalTok{  const errorTxt = readoutG.append("text")}
\NormalTok{    .attr("y", 20).attr("font{-}size", "12px").attr("fill", AMBER);}

\NormalTok{  // {-}{-}{-}{-} Update function {-}{-}{-}{-}}
\NormalTok{  function update(K) \{}
\NormalTok{    kLabel.text(K.toFixed(1));}
\NormalTok{    kBoxLabel.text(\textasciigrave{}C(s) = $\{K.toFixed(1)\}\textasciigrave{});}

\NormalTok{    // Plant G(s) = 1/(5s+1), C(s) = K}
\NormalTok{    // Closed{-}loop pole: s = {-}(1+K)/5}
\NormalTok{    const pole = {-}(1 + K) / 5;}
\NormalTok{    const tau\_cl = 5 / (1 + K);          // closed{-}loop time constant}
\NormalTok{    const ss\_gain = K / (1 + K);         // steady{-}state output for unit step}
\NormalTok{    const ss\_error = 1 / (1 + K);        // steady{-}state error}

\NormalTok{    // Settling time approx 4*tau\_cl}
\NormalTok{    const t\_settle = 4 * tau\_cl;}

\NormalTok{    // Step response: y(t) = ss\_gain * (1 {-} e\^{}\{{-}t/tau\_cl\})}
\NormalTok{    const tMax = Math.max(30, t\_settle * 1.6);}
\NormalTok{    const nPts = 200;}
\NormalTok{    const pts = d3.range(nPts).map(i =\textgreater{} \{}
\NormalTok{      const t = (i / (nPts {-} 1)) * tMax;}
\NormalTok{      return [t, ss\_gain * (1 {-} Math.exp({-}t / tau\_cl))];}
\NormalTok{    \});}

\NormalTok{    const tScale = d3.scaleLinear([0, tMax], [spXmin, spXmax]);}

\NormalTok{    // Error shading area}
\NormalTok{    const areaGen = d3.area()}
\NormalTok{      .x(d =\textgreater{} tScale(d[0]))}
\NormalTok{      .y0(yScale(1))}
\NormalTok{      .y1(d =\textgreater{} yScale(Math.min(d[1], 1.4)));}
\NormalTok{    errorShade.attr("d", areaGen(pts));}

\NormalTok{    // Response path}
\NormalTok{    const lineGen = d3.line().x(d =\textgreater{} tScale(d[0])).y(d =\textgreater{} yScale(Math.min(d[1], 1.4)));}
\NormalTok{    responsePath.attr("d", lineGen(pts));}

\NormalTok{    // Pole position}
\NormalTok{    poleCircle.attr("cx", poleXscale(pole));}
\NormalTok{    poleLabelEl.attr("x", poleXscale(pole)).text(\textasciigrave{}s = $\{pole.toFixed(2)\}\textasciigrave{});}

\NormalTok{    settlingTxt.text(\textasciigrave{}Settling time ≈ $\{t\_settle.toFixed(1)\} s\textasciigrave{});}
\NormalTok{    errorTxt.text(\textasciigrave{}Steady{-}state error = $\{(ss\_error * 100).toFixed(1)\}\%\textasciigrave{});}
\NormalTok{  \}}

\NormalTok{  slider.on("input", function() \{ update(+this.value); \});}
\NormalTok{  update(2);}

\NormalTok{  return container.node();}
\NormalTok{\}}
\end{Highlighting}
\end{Shaded}

\section{What stability means}\label{what-stability-means}

A control loop is \textbf{stable} if small disturbances do not produce
outputs that grow without bound. In continuous-time linear systems, the
quick rule is:

\begin{itemize}
\tightlist
\item
  if every closed-loop pole has negative real part, the response decays
\item
  if any closed-loop pole has positive real part, the response grows
\item
  if poles sit on the imaginary axis, the system is marginal and small
  model errors may push it into trouble
\end{itemize}

This is why pole location is not bookkeeping. It is a statement about
physical behaviour.

The same stability idea appears in state-space notation --- the same
feedback principle, expressed with vectors and matrices instead of
transfer functions. For a state-space model

\[\dot{x} = Ax + Bu, \qquad u = -Kx\]

the closed-loop dynamics become

\[\dot{x} = (A - BK)x\]

Now the eigenvalues \(\lambda(A-BK)\) play the same role that
closed-loop poles play in transfer-function form. They tell you what the
state does in time.

\begin{tcolorbox}[enhanced jigsaw, opacityback=0, colback=white, breakable, colframe=quarto-callout-note-color-frame, arc=.35mm, left=2mm, toprule=.15mm, leftrule=.75mm, bottomrule=.15mm, rightrule=.15mm]
\begin{minipage}[t]{5.5mm}
\textcolor{quarto-callout-note-color}{\faInfo}
\end{minipage}%
\begin{minipage}[t]{\textwidth - 5.5mm}

\vspace{-3mm}\textbf{Why This Works}\vspace{3mm}

Feedback changes the equation you are solving. Without control, the
plant's dynamics are built into \(G(s)\) or \(A\). With control, the
controller feeds the measured output back into the input channel, so the
effective denominator of the system changes from ``whatever the plant
had'' to ``plant plus feedback.''

That is the whole reason control is powerful. You are not passively
accepting the poles, modes, or eigenvalues the plant came with. You are
moving them.

\end{minipage}%
\end{tcolorbox}

\section{The core method}\label{the-core-method}

For a first pass through a control problem, the workflow is usually:

\begin{enumerate}
\def\labelenumi{\arabic{enumi}.}
\tightlist
\item
  Write the plant model in transfer-function or state-space form.
\item
  Decide what ``good behaviour'' means: stable, fast enough, low
  overshoot, small steady-state error, acceptable control effort.
\item
  Choose a controller structure: proportional, PI, PID, state feedback,
  or another architecture appropriate to the system.
\item
  Form the closed-loop model.
\item
  Inspect the poles or eigenvalues and connect them back to the time
  response.
\item
  Tune, then check what the tuning costs you elsewhere.
\end{enumerate}

The important habit is to keep the interpretation attached to the
algebra. A gain value is not just a number. It changes speed, error,
sensitivity, and sometimes noise amplification.

\section{Worked example 1: cruise control with proportional
feedback}\label{worked-example-1-cruise-control-with-proportional-feedback}

Suppose a simplified vehicle-speed model has plant transfer function

\[G(s) = \frac{1}{5s + 1}\]

where input is throttle command and output is speed deviation from the
desired operating point. Use a proportional controller

\[C(s) = K\]

with unity feedback.

The closed-loop transfer function is

\[T(s) = \frac{KG(s)}{1 + KG(s)}
= \frac{K/(5s+1)}{1 + K/(5s+1)}
= \frac{K}{5s + 1 + K}\]

So the closed-loop pole is at

\[s = -\frac{1+K}{5}\]

For any \(K > -1\), the pole is in the left half-plane, so the linear
model is stable. If \(K > 0\), increasing \(K\) moves the pole further
left, which makes the response faster.

The steady-state gain is

\[T(0) = \frac{K}{1+K}\]

so proportional feedback alone does not remove step-tracking error
completely. For a unit step reference, the steady-state output is
\(K/(1+K)\) and the steady-state error is

\[1 - \frac{K}{1+K} = \frac{1}{1+K}\]

This is the standard tradeoff. Larger \(K\) reduces error and speeds the
loop, but in a more realistic model it also increases sensitivity to
unmodelled dynamics, actuator limits, and measurement noise.

\emph{The two-panel chart below shows both sides of this tradeoff at
once. Drag the cursor to explore.}

\begin{Shaded}
\begin{Highlighting}[]
\NormalTok{//| echo: false}
\NormalTok{gainTradeoffViz = \{}
\NormalTok{  const d3 = await require("d3@7");}

\NormalTok{  const W = 560, H = 340;}
\NormalTok{  const BLUE = "\#2563eb", RED = "\#ef4444", AMBER = "\#f59e0b";}
\NormalTok{  const ML = 55, MR = 20, panelH = 130, gap = 30;}

\NormalTok{  const container = d3.create("div")}
\NormalTok{    .style("font{-}family", "sans{-}serif")}
\NormalTok{    .style("max{-}width", W + "px");}

\NormalTok{  const svg = container.append("svg")}
\NormalTok{    .attr("width", W).attr("height", H)}
\NormalTok{    .attr("viewBox", \textasciigrave{}0 0 $\{W\} $\{H\}\textasciigrave{})}
\NormalTok{    .attr("aria{-}label", "Proportional gain tradeoff chart");}

\NormalTok{  const innerW = W {-} ML {-} MR;}
\NormalTok{  const Kmax = 20;}
\NormalTok{  const kScale = d3.scaleLinear([0, Kmax], [0, innerW]);}

\NormalTok{  // {-}{-}{-}{-} Top panel: pole location {-}{-}{-}{-}}
\NormalTok{  const topG = svg.append("g").attr("transform", \textasciigrave{}translate($\{ML\}, 20)\textasciigrave{});}
\NormalTok{  const poleScale = d3.scaleLinear([0, {-}5], [panelH, 0]);}

\NormalTok{  // axes}
\NormalTok{  topG.append("line").attr("x1", 0).attr("y1", panelH)}
\NormalTok{    .attr("x2", innerW).attr("y2", panelH).attr("stroke", "\#6b7280").attr("stroke{-}width", 1.5);}
\NormalTok{  topG.append("line").attr("x1", 0).attr("y1", 0)}
\NormalTok{    .attr("x2", 0).attr("y2", panelH).attr("stroke", "\#6b7280").attr("stroke{-}width", 1.5);}

\NormalTok{  topG.append("text").attr("x", innerW / 2).attr("y", panelH + 22)}
\NormalTok{    .attr("text{-}anchor", "middle").attr("font{-}size", "11px").attr("fill", "\#374151").text("K");}
\NormalTok{  topG.append("text").attr("x", {-}40).attr("y", panelH / 2 + 4)}
\NormalTok{    .attr("font{-}size", "11px").attr("fill", "\#374151").text("pole s");}

\NormalTok{  // annotation}
\NormalTok{  topG.append("text").attr("x", innerW * 0.55).attr("y", 18)}
\NormalTok{    .attr("font{-}size", "10px").attr("fill", "\#6b7280").text("farther left = faster response");}

\NormalTok{  // pole curve: s = {-}(1+K)/5}
\NormalTok{  const polePts = d3.range(201).map(i =\textgreater{} \{}
\NormalTok{    const K = i / 200 * Kmax;}
\NormalTok{    return [kScale(K), poleScale({-}(1 + K) / 5)];}
\NormalTok{  \});}
\NormalTok{  topG.append("path")}
\NormalTok{    .attr("d", d3.line()(polePts))}
\NormalTok{    .attr("fill", "none").attr("stroke", BLUE).attr("stroke{-}width", 2);}

\NormalTok{  // y{-}axis ticks}
\NormalTok{  [{-}1,{-}2,{-}3,{-}4,{-}5].forEach(v =\textgreater{} \{}
\NormalTok{    topG.append("line").attr("x1", {-}4).attr("y1", poleScale(v))}
\NormalTok{      .attr("x2", 0).attr("y2", poleScale(v)).attr("stroke", "\#9ca3af");}
\NormalTok{    topG.append("text").attr("x", {-}8).attr("y", poleScale(v) + 4)}
\NormalTok{      .attr("text{-}anchor", "end").attr("font{-}size", "9px").attr("fill", "\#6b7280").text(v);}
\NormalTok{  \});}
\NormalTok{  // x{-}axis ticks}
\NormalTok{  [0, 5, 10, 15, 20].forEach(v =\textgreater{} \{}
\NormalTok{    topG.append("text").attr("x", kScale(v)).attr("y", panelH + 14)}
\NormalTok{      .attr("text{-}anchor", "middle").attr("font{-}size", "9px").attr("fill", "\#6b7280").text(v);}
\NormalTok{  \});}

\NormalTok{  // cursor line and dot (top)}
\NormalTok{  const topCursor = topG.append("line")}
\NormalTok{    .attr("y1", 0).attr("y2", panelH)}
\NormalTok{    .attr("stroke", AMBER).attr("stroke{-}width", 1.5).attr("stroke{-}dasharray", "4,3");}
\NormalTok{  const topDot = topG.append("circle").attr("r", 5).attr("fill", AMBER);}

\NormalTok{  // {-}{-}{-}{-} Bottom panel: steady{-}state error {-}{-}{-}{-}}
\NormalTok{  const botG = svg.append("g").attr("transform", \textasciigrave{}translate($\{ML\}, $\{20 + panelH + gap\})\textasciigrave{});}
\NormalTok{  const errScale = d3.scaleLinear([0, 1], [panelH, 0]);}

\NormalTok{  botG.append("line").attr("x1", 0).attr("y1", panelH)}
\NormalTok{    .attr("x2", innerW).attr("y2", panelH).attr("stroke", "\#6b7280").attr("stroke{-}width", 1.5);}
\NormalTok{  botG.append("line").attr("x1", 0).attr("y1", 0)}
\NormalTok{    .attr("x2", 0).attr("y2", panelH).attr("stroke", "\#6b7280").attr("stroke{-}width", 1.5);}

\NormalTok{  botG.append("text").attr("x", innerW / 2).attr("y", panelH + 22)}
\NormalTok{    .attr("text{-}anchor", "middle").attr("font{-}size", "11px").attr("fill", "\#374151").text("K");}
\NormalTok{  botG.append("text").attr("x", {-}40).attr("y", panelH / 2 + 4)}
\NormalTok{    .attr("font{-}size", "11px").attr("fill", "\#374151").text("error");}

\NormalTok{  botG.append("text").attr("x", innerW * 0.35).attr("y", panelH {-} 8)}
\NormalTok{    .attr("font{-}size", "10px").attr("fill", "\#6b7280")}
\NormalTok{    .text("never reaches zero with proportional control only");}

\NormalTok{  const errPts = d3.range(201).map(i =\textgreater{} \{}
\NormalTok{    const K = i / 200 * Kmax;}
\NormalTok{    return [kScale(K), errScale(1 / (1 + K))];}
\NormalTok{  \});}
\NormalTok{  botG.append("path")}
\NormalTok{    .attr("d", d3.line()(errPts))}
\NormalTok{    .attr("fill", "none").attr("stroke", RED).attr("stroke{-}width", 2);}

\NormalTok{  [0, 5, 10, 15, 20].forEach(v =\textgreater{} \{}
\NormalTok{    botG.append("text").attr("x", kScale(v)).attr("y", panelH + 14)}
\NormalTok{      .attr("text{-}anchor", "middle").attr("font{-}size", "9px").attr("fill", "\#6b7280").text(v);}
\NormalTok{  \});}
\NormalTok{  [0, 0.5, 1].forEach(v =\textgreater{} \{}
\NormalTok{    botG.append("line").attr("x1", {-}4).attr("y1", errScale(v))}
\NormalTok{      .attr("x2", 0).attr("y2", errScale(v)).attr("stroke", "\#9ca3af");}
\NormalTok{    botG.append("text").attr("x", {-}8).attr("y", errScale(v) + 4)}
\NormalTok{      .attr("text{-}anchor", "end").attr("font{-}size", "9px").attr("fill", "\#6b7280")}
\NormalTok{      .text(v.toFixed(1));}
\NormalTok{  \});}

\NormalTok{  const botCursor = botG.append("line")}
\NormalTok{    .attr("y1", 0).attr("y2", panelH)}
\NormalTok{    .attr("stroke", AMBER).attr("stroke{-}width", 1.5).attr("stroke{-}dasharray", "4,3");}
\NormalTok{  const botDot = botG.append("circle").attr("r", 5).attr("fill", AMBER);}

\NormalTok{  // Readout text}
\NormalTok{  const readoutEl = svg.append("text")}
\NormalTok{    .attr("x", ML + 10).attr("y", 20 + panelH * 2 + gap * 2 + 10)}
\NormalTok{    .attr("font{-}size", "12px").attr("fill", AMBER);}

\NormalTok{  // Drag interaction on full SVG}
\NormalTok{  let currentK = 5;}
\NormalTok{  function updateCursor(K) \{}
\NormalTok{    K = Math.max(0, Math.min(Kmax, K));}
\NormalTok{    currentK = K;}
\NormalTok{    const cx = kScale(K);}
\NormalTok{    const poleVal = {-}(1 + K) / 5;}
\NormalTok{    const errVal = 1 / (1 + K);}
\NormalTok{    topCursor.attr("x1", cx).attr("x2", cx);}
\NormalTok{    topDot.attr("cx", cx).attr("cy", poleScale(poleVal));}
\NormalTok{    botCursor.attr("x1", cx).attr("x2", cx);}
\NormalTok{    botDot.attr("cx", cx).attr("cy", errScale(errVal));}
\NormalTok{    readoutEl.text(\textasciigrave{}K = $\{K.toFixed(1)\}  |  pole = $\{poleVal.toFixed(3)\}  |  error = $\{(errVal * 100).toFixed(1)\}\%\textasciigrave{});}
\NormalTok{  \}}

\NormalTok{  svg.style("cursor", "crosshair")}
\NormalTok{    .on("mousemove", function(event) \{}
\NormalTok{      const [mx] = d3.pointer(event);}
\NormalTok{      const K = (mx {-} ML) / innerW * Kmax;}
\NormalTok{      updateCursor(K);}
\NormalTok{    \});}

\NormalTok{  updateCursor(5);}
\NormalTok{  return container.node();}
\NormalTok{\}}
\end{Highlighting}
\end{Shaded}

\section{Worked example 2: state feedback for a two-state
system}\label{worked-example-2-state-feedback-for-a-two-state-system}

Consider the linear system

\[\dot{x} = Ax + Bu\]

with

\[A = \begin{pmatrix} 0 & 1 \\ -2 & -1 \end{pmatrix},
\qquad
B = \begin{pmatrix} 0 \\ 1 \end{pmatrix}\]

Choose state feedback

\[u = -Kx, \qquad K = \begin{pmatrix} 3 & 2 \end{pmatrix}\]

Then

\[A - BK
=
\begin{pmatrix} 0 & 1 \\ -2 & -1 \end{pmatrix}
-
\begin{pmatrix} 0 \\ 1 \end{pmatrix}
\begin{pmatrix} 3 & 2 \end{pmatrix}
=
\begin{pmatrix} 0 & 1 \\ -5 & -3 \end{pmatrix}\]

Its characteristic polynomial comes from the 2×2 determinant:

\[\det(\lambda I - (A-BK))
= \det\begin{pmatrix} \lambda & -1 \\ 5 & \lambda+3 \end{pmatrix}
= \lambda(\lambda+3) + 5
= \lambda^2 + 3\lambda + 5\]

The eigenvalues are

\[\lambda = \frac{-3 \pm \sqrt{9-20}}{2}
= -\frac{3}{2} \pm \frac{\sqrt{11}}{2}i\]

Both eigenvalues have negative real part, so the closed-loop system is
stable. Compared with the uncontrolled system, the controller has
shifted the system's modes. That is the state-space version of pole
placement.

\emph{The interactive below lets you explore how the gain vector
\(K = (k_1, k_2)\) shapes the eigenvalues and phase portrait.}

\begin{Shaded}
\begin{Highlighting}[]
\NormalTok{//| echo: false}
\NormalTok{stateFeedbackViz = \{}
\NormalTok{  const d3 = await require("d3@7");}

\NormalTok{  const W = 640, PH = 280;}
\NormalTok{  const BLUE = "\#2563eb", GREEN = "\#059669", RED = "\#ef4444", AMBER = "\#f59e0b";}

\NormalTok{  const container = d3.create("div")}
\NormalTok{    .style("font{-}family", "sans{-}serif")}
\NormalTok{    .style("max{-}width", W + "px");}

\NormalTok{  // Sliders}
\NormalTok{  const sliderDiv = container.append("div")}
\NormalTok{    .style("display", "flex").style("gap", "1.5rem").style("margin{-}bottom", "0.5rem");}

\NormalTok{  function makeSlider(parent, label, min, max, step, val) \{}
\NormalTok{    const wrap = parent.append("div").style("flex", "1");}
\NormalTok{    wrap.append("label").style("font{-}size", "0.88rem").style("display", "block")}
\NormalTok{      .text(label);}
\NormalTok{    const inp = wrap.append("input").attr("type", "range")}
\NormalTok{      .attr("min", min).attr("max", max).attr("step", step).attr("value", val)}
\NormalTok{      .style("width", "100\%");}
\NormalTok{    const lbl = wrap.append("span").style("font{-}size", "0.88rem").style("font{-}weight", "bold");}
\NormalTok{    return \{inp, lbl\};}
\NormalTok{  \}}

\NormalTok{  const s1 = makeSlider(sliderDiv, "k₁:", 0, 6, 0.1, 3);}
\NormalTok{  const s2 = makeSlider(sliderDiv, "k₂:", 0, 6, 0.1, 2);}

\NormalTok{  const svg = container.append("svg")}
\NormalTok{    .attr("width", W).attr("height", PH)}
\NormalTok{    .attr("viewBox", \textasciigrave{}0 0 $\{W\} $\{PH\}\textasciigrave{})}
\NormalTok{    .attr("aria{-}label", "State feedback phase portrait");}

\NormalTok{  const panW = W / 2 {-} 20;}

\NormalTok{  // {-}{-}{-}{-} Left: s{-}plane {-}{-}{-}{-}}
\NormalTok{  const spG = svg.append("g").attr("transform", "translate(10, 10)");}
\NormalTok{  const spXsc = d3.scaleLinear([{-}5, 1], [30, panW {-} 10]);}
\NormalTok{  const spYsc = d3.scaleLinear([{-}4, 4], [PH {-} 40, 30]);}

\NormalTok{  // Stable region shading}
\NormalTok{  spG.append("rect").attr("x", 30).attr("y", 30)}
\NormalTok{    .attr("width", spXsc(0) {-} 30).attr("height", PH {-} 70)}
\NormalTok{    .attr("fill", GREEN).attr("fill{-}opacity", 0.1);}

\NormalTok{  // Axes}
\NormalTok{  spG.append("line").attr("x1", 30).attr("y1", spYsc(0))}
\NormalTok{    .attr("x2", panW {-} 10).attr("y2", spYsc(0)).attr("stroke", "\#9ca3af").attr("stroke{-}width", 1);}
\NormalTok{  spG.append("line").attr("x1", spXsc(0)).attr("y1", 30)}
\NormalTok{    .attr("x2", spXsc(0)).attr("y2", PH {-} 40)}
\NormalTok{    .attr("stroke", "\#9ca3af").attr("stroke{-}width", 1).attr("stroke{-}dasharray", "4,3");}

\NormalTok{  spG.append("text").attr("x", panW / 2).attr("y", 18)}
\NormalTok{    .attr("text{-}anchor", "middle").attr("font{-}size", "12px").attr("fill", "\#374151")}
\NormalTok{    .text("s{-}plane eigenvalues");}
\NormalTok{  spG.append("text").attr("x", spXsc(0) + 2).attr("y", 28)}
\NormalTok{    .attr("font{-}size", "9px").attr("fill", "\#9ca3af").text("stability boundary");}

\NormalTok{  [{-}4,{-}3,{-}2,{-}1,0].forEach(v =\textgreater{} \{}
\NormalTok{    spG.append("text").attr("x", spXsc(v)).attr("y", spYsc(0) + 14)}
\NormalTok{      .attr("text{-}anchor", "middle").attr("font{-}size", "9px").attr("fill", "\#9ca3af").text(v);}
\NormalTok{  \});}
\NormalTok{  [{-}3,{-}2,{-}1,1,2,3].forEach(v =\textgreater{} \{}
\NormalTok{    spG.append("text").attr("x", spXsc(0) + 5).attr("y", spYsc(v) + 4)}
\NormalTok{      .attr("font{-}size", "9px").attr("fill", "\#9ca3af").text(\textasciigrave{}$\{v\}i\textasciigrave{});}
\NormalTok{  \});}
\NormalTok{  spG.append("text").attr("x", panW {-} 8).attr("y", spYsc(0) + 4)}
\NormalTok{    .attr("font{-}size", "10px").attr("fill", "\#9ca3af").text("Re");}
\NormalTok{  spG.append("text").attr("x", spXsc(0) {-} 14).attr("y", 34)}
\NormalTok{    .attr("font{-}size", "10px").attr("fill", "\#9ca3af").text("Im");}

\NormalTok{  const ev1 = spG.append("circle").attr("r", 7).attr("stroke{-}width", 2);}
\NormalTok{  const ev2 = spG.append("circle").attr("r", 7).attr("stroke{-}width", 2);}
\NormalTok{  const evLabel = spG.append("text")}
\NormalTok{    .attr("text{-}anchor", "middle").attr("font{-}size", "11px").attr("font{-}weight", "bold");}
\NormalTok{  const stabilityAnnot = spG.append("text")}
\NormalTok{    .attr("x", spXsc({-}2.5)).attr("y", PH {-} 45)}
\NormalTok{    .attr("text{-}anchor", "middle").attr("font{-}size", "11px").attr("font{-}weight", "bold");}

\NormalTok{  // {-}{-}{-}{-} Right: phase portrait {-}{-}{-}{-}}
\NormalTok{  const ppG = svg.append("g").attr("transform", \textasciigrave{}translate($\{W / 2 + 10\}, 10)\textasciigrave{});}
\NormalTok{  const ppXsc = d3.scaleLinear([{-}3, 3], [30, panW {-} 10]);}
\NormalTok{  const ppYsc = d3.scaleLinear([{-}3, 3], [PH {-} 40, 30]);}

\NormalTok{  ppG.append("text").attr("x", panW / 2).attr("y", 18)}
\NormalTok{    .attr("text{-}anchor", "middle").attr("font{-}size", "12px").attr("fill", "\#374151")}
\NormalTok{    .text("(x₁, x₂) phase portrait");}

\NormalTok{  // Axes}
\NormalTok{  ppG.append("line").attr("x1", 30).attr("y1", ppYsc(0))}
\NormalTok{    .attr("x2", panW {-} 10).attr("y2", ppYsc(0)).attr("stroke", "\#d1d5db").attr("stroke{-}width", 1);}
\NormalTok{  ppG.append("line").attr("x1", ppXsc(0)).attr("y1", 30)}
\NormalTok{    .attr("x2", ppXsc(0)).attr("y2", PH {-} 40).attr("stroke", "\#d1d5db").attr("stroke{-}width", 1);}
\NormalTok{  ppG.append("text").attr("x", panW {-} 5).attr("y", ppYsc(0) + 4)}
\NormalTok{    .attr("font{-}size", "10px").attr("fill", "\#9ca3af").text("x₁");}
\NormalTok{  ppG.append("text").attr("x", ppXsc(0) + 3).attr("y", 28)}
\NormalTok{    .attr("font{-}size", "10px").attr("fill", "\#9ca3af").text("x₂");}

\NormalTok{  const trajectoryPath = ppG.append("path")}
\NormalTok{    .attr("fill", "none").attr("stroke{-}width", 2);}
\NormalTok{  const ppAnnot = ppG.append("text")}
\NormalTok{    .attr("x", panW / 2).attr("y", PH {-} 44)}
\NormalTok{    .attr("text{-}anchor", "middle").attr("font{-}size", "11px").attr("font{-}weight", "bold");}

\NormalTok{  function update(k1, k2) \{}
\NormalTok{    s1.lbl.text(k1.toFixed(1));}
\NormalTok{    s2.lbl.text(k2.toFixed(1));}

\NormalTok{    // A {-} BK where A=[[0,1],[{-}2,{-}1]], B=[[0],[1]], K=[k1,k2]}
\NormalTok{    // A{-}BK = [[0,1],[{-}2{-}k1,{-}1{-}k2]]}
\NormalTok{    const a = 0, b = 1, c = {-}2 {-} k1, d = {-}1 {-} k2;}
\NormalTok{    // char poly: λ² {-} (a+d)λ + (ad{-}bc)}
\NormalTok{    const tr = a + d;}
\NormalTok{    const det = a * d {-} b * c;}
\NormalTok{    const disc = tr * tr {-} 4 * det;}

\NormalTok{    let re1, im1, re2, im2, stable;}
\NormalTok{    if (disc \textgreater{}= 0) \{}
\NormalTok{      re1 = (tr + Math.sqrt(disc)) / 2;}
\NormalTok{      re2 = (tr {-} Math.sqrt(disc)) / 2;}
\NormalTok{      im1 = 0; im2 = 0;}
\NormalTok{      stable = re1 \textless{} 0 \&\& re2 \textless{} 0;}
\NormalTok{    \} else \{}
\NormalTok{      re1 = re2 = tr / 2;}
\NormalTok{      im1 = Math.sqrt({-}disc) / 2;}
\NormalTok{      im2 = {-}im1;}
\NormalTok{      stable = re1 \textless{} 0;}
\NormalTok{    \}}

\NormalTok{    const col = stable ? BLUE : RED;}
\NormalTok{    ev1.attr("cx", spXsc(re1)).attr("cy", spYsc(im1))}
\NormalTok{      .attr("fill", col).attr("stroke", col === BLUE ? "\#1e40af" : "\#991b1b");}
\NormalTok{    ev2.attr("cx", spXsc(re2)).attr("cy", spYsc(im2))}
\NormalTok{      .attr("fill", col).attr("stroke", col === BLUE ? "\#1e40af" : "\#991b1b");}

\NormalTok{    const evStr = disc \textless{} 0}
\NormalTok{      ? \textasciigrave{}λ = $\{re1.toFixed(2)\} ± $\{Math.abs(im1).toFixed(2)\}i\textasciigrave{}}
\NormalTok{      : \textasciigrave{}λ = $\{re1.toFixed(2)\}, $\{re2.toFixed(2)\}\textasciigrave{};}
\NormalTok{    evLabel.attr("x", spXsc(re1)).attr("y", spYsc(im1) {-} 12)}
\NormalTok{      .attr("fill", col).text(evStr);}

\NormalTok{    stabilityAnnot.attr("fill", stable ? GREEN : RED)}
\NormalTok{      .text(stable ? "Re(λ) \textless{} 0: decays to origin" : "Re(λ) ≥ 0: grows away from origin");}

\NormalTok{    // Euler trajectory starting near (1, 1)}
\NormalTok{    const dt = 0.02, steps = 600;}
\NormalTok{    let x1 = 1.0, x2 = 1.0;}
\NormalTok{    const path = [[x1, x2]];}
\NormalTok{    const clip = 4;}
\NormalTok{    for (let i = 0; i \textless{} steps; i++) \{}
\NormalTok{      const u = {-}(k1 * x1 + k2 * x2);}
\NormalTok{      const dx1 = a * x1 + b * x2;}
\NormalTok{      const dx2 = c * x1 + d * x2;}
\NormalTok{      x1 += dt * dx1;}
\NormalTok{      x2 += dt * dx2;}
\NormalTok{      if (Math.abs(x1) \textgreater{} clip || Math.abs(x2) \textgreater{} clip) break;}
\NormalTok{      path.push([x1, x2]);}
\NormalTok{    \}}

\NormalTok{    const lineGen = d3.line().x(d =\textgreater{} ppXsc(d[0])).y(d =\textgreater{} ppYsc(d[1]));}
\NormalTok{    trajectoryPath.attr("d", lineGen(path)).attr("stroke", stable ? GREEN : RED);}
\NormalTok{    ppAnnot.attr("fill", stable ? GREEN : RED)}
\NormalTok{      .text(stable ? "decays to origin" : "grows away from origin");}
\NormalTok{  \}}

\NormalTok{  s1.inp.on("input", function() \{ update(+this.value, +s2.inp.property("value")); \});}
\NormalTok{  s2.inp.on("input", function() \{ update(+s1.inp.property("value"), +this.value); \});}
\NormalTok{  update(3, 2);}

\NormalTok{  return container.node();}
\NormalTok{\}}
\end{Highlighting}
\end{Shaded}

This is the language used constantly in robotics, flight control, and
embedded control software. The code may be digital and the sensors
noisy, but the core mathematical question is still: what did your
feedback law do to the modes?

\section{Worked example 3: stabilising a scientific
instrument}\label{worked-example-3-stabilising-a-scientific-instrument}

An optics lab wants the intensity of a laser to stay near a target
despite thermal drift. A simplified linear model of the plant is slow
and first-order, so the experimentalist begins with a proportional
controller for the same reason an engineer would: it is cheap to
implement and easy to reason about.

If the gain is too low, the output drifts and the reference is not
tracked closely enough. If the gain is pushed too high, noise and delay
in the sensor chain start to matter. The mathematics is exactly the same
as the motor-speed example. What changes is the vocabulary of the lab.

This is the point of the optional viewpoint. Control theory is not owned
by electrical engineering. It is a general mathematical structure for
shaping dynamics under measurement.

\section{Where this goes}\label{where-this-goes-37}

The most direct continuation is
\href{05-estimation-inverse.qmd}{Estimation, inverse problems, and
filtering}. Real control systems rarely have direct access to every
state they need. Once you start asking how to control a system, the next
question is often how to estimate the quantities you cannot measure
cleanly.

This chapter also sharpens the way you should read later Volume 8
material on sampled systems and computational models. A simulation is
not yet a controller. A model can be numerically stable and still be a
bad control design. Volume 8 keeps returning to that distinction.

\begin{tcolorbox}[enhanced jigsaw, opacityback=0, colback=white, breakable, colframe=quarto-callout-tip-color-frame, arc=.35mm, left=2mm, toprule=.15mm, leftrule=.75mm, bottomrule=.15mm, rightrule=.15mm]
\begin{minipage}[t]{5.5mm}
\textcolor{quarto-callout-tip-color}{\faLightbulb}
\end{minipage}%
\begin{minipage}[t]{\textwidth - 5.5mm}

\vspace{-3mm}\textbf{Applications}\vspace{3mm}

\begin{itemize}
\tightlist
\item
  motor-speed control in electric drives
\item
  altitude hold and attitude stabilisation in drones
\item
  furnace and reactor temperature control
\item
  active vibration suppression in flexible structures
\item
  instrument stabilisation in optics and experimental physics
\item
  feedback loops inside robotics and autonomous systems
\end{itemize}

\end{minipage}%
\end{tcolorbox}

\section{Exercises}\label{exercises-48}

These are project-style exercises. Each one asks you to connect the
algebra to behaviour, not just compute a number.

\begin{Shaded}
\begin{Highlighting}[]
\NormalTok{//| echo: false}
\NormalTok{katex = require("katex@0.16")}
\end{Highlighting}
\end{Shaded}

\begin{Shaded}
\begin{Highlighting}[]
\NormalTok{//| echo: false}
\NormalTok{function makeStepperHTML(exerciseNum, steps) \{}
\NormalTok{  let current = 0;}
\NormalTok{  const totalSteps = steps.length;}

\NormalTok{  const container = document.createElement("div");}
\NormalTok{  container.style.cssText = "border:1px solid \#e5e7eb; border{-}radius:8px; padding:1rem 1.25rem; margin:0.75rem 0; font{-}family:inherit;";}

\NormalTok{  const stepsDiv = document.createElement("div");}
\NormalTok{  stepsDiv.style.marginBottom = "0.75rem";}

\NormalTok{  function renderSteps() \{}
\NormalTok{    stepsDiv.innerHTML = "";}
\NormalTok{    for (let i = 0; i \textless{} current; i++) \{}
\NormalTok{      const s = steps[i];}
\NormalTok{      const row = document.createElement("div");}
\NormalTok{      row.style.cssText = "display:grid; grid{-}template{-}columns:200px 1fr; gap:0.5rem 1rem; align{-}items:baseline; padding:0.35rem 0; border{-}top: 1px solid \#f3f4f6;";}

\NormalTok{      const opCell = document.createElement("span");}
\NormalTok{      opCell.style.cssText = "font{-}size:0.85em; color:\#6b7280; font{-}style:italic;";}
\NormalTok{      opCell.textContent = s.op;}

\NormalTok{      const eqCell = document.createElement("span");}
\NormalTok{      eqCell.innerHTML = katex.renderToString(s.eq, \{ throwOnError: false, displayMode: false \});}

\NormalTok{      row.appendChild(opCell);}
\NormalTok{      row.appendChild(eqCell);}

\NormalTok{      if (s.note) \{}
\NormalTok{        const noteRow = document.createElement("div");}
\NormalTok{        noteRow.style.cssText = "grid{-}column:2; font{-}size:0.82em; color:\#6b7280; padding{-}bottom:0.2rem;";}
\NormalTok{        noteRow.textContent = s.note;}
\NormalTok{        row.appendChild(noteRow);}
\NormalTok{      \}}

\NormalTok{      stepsDiv.appendChild(row);}
\NormalTok{    \}}

\NormalTok{    if (current === totalSteps) \{}
\NormalTok{      const done = document.createElement("div");}
\NormalTok{      done.style.cssText = "margin{-}top:0.5rem; font{-}size:0.9em; color:\#059669; font{-}weight:500;";}
\NormalTok{      done.textContent = "✓ Solution complete";}
\NormalTok{      stepsDiv.appendChild(done);}
\NormalTok{    \}}
\NormalTok{  \}}

\NormalTok{  const controls = document.createElement("div");}
\NormalTok{  controls.style.cssText = "display:flex; gap:0.5rem; align{-}items:center;";}

\NormalTok{  const nextBtn = document.createElement("button");}
\NormalTok{  nextBtn.textContent = "Next step →";}
\NormalTok{  nextBtn.style.cssText = "padding:0.35rem 0.85rem; border:1px solid \#d1d5db; border{-}radius:4px; background:\#fff; cursor:pointer; font{-}size:0.9em;";}
\NormalTok{  nextBtn.onclick = () =\textgreater{} \{}
\NormalTok{    if (current \textless{} totalSteps) \{ current++; renderSteps(); updateButtons(); \}}
\NormalTok{  \};}

\NormalTok{  const resetBtn = document.createElement("button");}
\NormalTok{  resetBtn.textContent = "Reset";}
\NormalTok{  resetBtn.style.cssText = "padding:0.35rem 0.75rem; border:1px solid \#e5e7eb; border{-}radius:4px; background:\#f9fafb; cursor:pointer; font{-}size:0.9em; color:\#6b7280;";}
\NormalTok{  resetBtn.onclick = () =\textgreater{} \{ current = 0; renderSteps(); updateButtons(); \};}

\NormalTok{  const counter = document.createElement("span");}
\NormalTok{  counter.style.cssText = "font{-}size:0.82em; color:\#9ca3af; margin{-}left:0.25rem;";}

\NormalTok{  function updateButtons() \{}
\NormalTok{    nextBtn.disabled = current === totalSteps;}
\NormalTok{    nextBtn.style.opacity = current === totalSteps ? "0.4" : "1";}
\NormalTok{    counter.textContent = current === 0 ? "Click to reveal steps" : \textasciigrave{}Step $\{current\} of $\{totalSteps\}\textasciigrave{};}
\NormalTok{  \}}

\NormalTok{  controls.appendChild(nextBtn);}
\NormalTok{  controls.appendChild(resetBtn);}
\NormalTok{  controls.appendChild(counter);}

\NormalTok{  container.appendChild(stepsDiv);}
\NormalTok{  container.appendChild(controls);}

\NormalTok{  renderSteps();}
\NormalTok{  updateButtons();}
\NormalTok{  return container;}
\NormalTok{\}}
\end{Highlighting}
\end{Shaded}

\subsection{Exercise 1}\label{exercise-1-1}

A thermal plant is modelled by

\[G(s) = \frac{2}{10s + 1}\]

with proportional controller \(C(s) = K\) and unity feedback.

\begin{enumerate}
\def\labelenumi{\arabic{enumi}.}
\tightlist
\item
  Derive the closed-loop transfer function.
\item
  Find the closed-loop pole.
\item
  For a unit-step reference, compute the steady-state error.
\item
  Compare what changes when \(K = 1\) and when \(K = 4\).
\end{enumerate}

\begin{Shaded}
\begin{Highlighting}[]
\NormalTok{//| echo: false}
\NormalTok{makeStepperHTML(1, [}
\NormalTok{  \{ op: "Write the closed{-}loop formula",}
\NormalTok{    eq: "T(s) = \textbackslash{}\textbackslash{}dfrac\{C(s)G(s)\}\{1 + C(s)G(s)\} = \textbackslash{}\textbackslash{}dfrac\{2K/(10s+1)\}\{1 + 2K/(10s+1)\}",}
\NormalTok{    note: "Substitute the plant and controller directly into the unity{-}feedback formula." \},}
\NormalTok{  \{ op: "Simplify the fraction",}
\NormalTok{    eq: "T(s) = \textbackslash{}\textbackslash{}dfrac\{2K\}\{10s + 1 + 2K\}",}
\NormalTok{    note: "Multiply numerator and denominator by 10s+1." \},}
\NormalTok{  \{ op: "Read off the pole",}
\NormalTok{    eq: "10s + 1 + 2K = 0 \textbackslash{}\textbackslash{}quad \textbackslash{}\textbackslash{}Rightarrow \textbackslash{}\textbackslash{}quad s = {-}\textbackslash{}\textbackslash{}dfrac\{1+2K\}\{10\}",}
\NormalTok{    note: "The pole comes from the denominator." \},}
\NormalTok{  \{ op: "Compute the steady{-}state gain",}
\NormalTok{    eq: "T(0) = \textbackslash{}\textbackslash{}dfrac\{2K\}\{1+2K\}",}
\NormalTok{    note: "Set s=0 to get the zero{-}frequency gain." \},}
\NormalTok{  \{ op: "Compute the steady{-}state error for a unit step",}
\NormalTok{    eq: "e\_\{ss\} = 1 {-} \textbackslash{}\textbackslash{}dfrac\{2K\}\{1+2K\} = \textbackslash{}\textbackslash{}dfrac\{1\}\{1+2K\}",}
\NormalTok{    note: "Output plus tracking error must sum to the unit reference." \},}
\NormalTok{  \{ op: "Check the two gains",}
\NormalTok{    eq: "K=1: \textbackslash{}\textbackslash{}; s={-}0.3, \textbackslash{}\textbackslash{}; e\_\{ss\}=\textbackslash{}\textbackslash{}tfrac\{1\}\{3\} \textbackslash{}\textbackslash{}qquad K=4: \textbackslash{}\textbackslash{}; s={-}0.9, \textbackslash{}\textbackslash{}; e\_\{ss\}=\textbackslash{}\textbackslash{}tfrac\{1\}\{9\}",}
\NormalTok{    note: "Larger gain speeds the loop and reduces the residual error." \}}
\NormalTok{])}
\end{Highlighting}
\end{Shaded}

\subsection{Exercise 2}\label{exercise-2-1}

A two-state model has

\[A = \begin{pmatrix} 0 & 1 \\ -1 & -0.5 \end{pmatrix},
\qquad
B = \begin{pmatrix} 0 \\ 1 \end{pmatrix},
\qquad
K = \begin{pmatrix} 4 & 1.5 \end{pmatrix}\]

Compute \(A-BK\) and decide whether the closed-loop system is stable.

\begin{Shaded}
\begin{Highlighting}[]
\NormalTok{//| echo: false}
\NormalTok{makeStepperHTML(2, [}
\NormalTok{  \{ op: "Write the closed{-}loop matrix",}
\NormalTok{    eq: "A{-}BK = \textbackslash{}\textbackslash{}begin\{pmatrix\} 0 \& 1 \textbackslash{}\textbackslash{}\textbackslash{}\textbackslash{} {-}1 \& {-}0.5 \textbackslash{}\textbackslash{}end\{pmatrix\} {-} \textbackslash{}\textbackslash{}begin\{pmatrix\} 0 \textbackslash{}\textbackslash{}\textbackslash{}\textbackslash{} 1 \textbackslash{}\textbackslash{}end\{pmatrix\}\textbackslash{}\textbackslash{}begin\{pmatrix\} 4 \& 1.5 \textbackslash{}\textbackslash{}end\{pmatrix\}",}
\NormalTok{    note: "State feedback replaces A with A{-}BK." \},}
\NormalTok{  \{ op: "Multiply BK",}
\NormalTok{    eq: "BK = \textbackslash{}\textbackslash{}begin\{pmatrix\} 0 \& 0 \textbackslash{}\textbackslash{}\textbackslash{}\textbackslash{} 4 \& 1.5 \textbackslash{}\textbackslash{}end\{pmatrix\}",}
\NormalTok{    note: "The outer product fills only the second row because B has zero in its first entry." \},}
\NormalTok{  \{ op: "Subtract",}
\NormalTok{    eq: "A{-}BK = \textbackslash{}\textbackslash{}begin\{pmatrix\} 0 \& 1 \textbackslash{}\textbackslash{}\textbackslash{}\textbackslash{} {-}5 \& {-}2 \textbackslash{}\textbackslash{}end\{pmatrix\}" \},}
\NormalTok{  \{ op: "Form the characteristic polynomial",}
\NormalTok{    eq: "\textbackslash{}\textbackslash{}det(\textbackslash{}\textbackslash{}lambda I {-} (A{-}BK)) = \textbackslash{}\textbackslash{}lambda\^{}2 + 2\textbackslash{}\textbackslash{}lambda + 5",}
\NormalTok{    note: "Take the determinant of the 2x2 matrix \textbackslash{}\textbackslash{}lambda I{-}(A{-}BK)." \},}
\NormalTok{  \{ op: "Find the eigenvalues",}
\NormalTok{    eq: "\textbackslash{}\textbackslash{}lambda = {-}1 \textbackslash{}\textbackslash{}pm 2i",}
\NormalTok{    note: "The real part is negative in both cases." \},}
\NormalTok{  \{ op: "Check",}
\NormalTok{    eq: "\textbackslash{}\textbackslash{}Re(\textbackslash{}\textbackslash{}lambda) = {-}1 \textless{} 0 \textbackslash{}\textbackslash{}; \textbackslash{}\textbackslash{}Rightarrow \textbackslash{}\textbackslash{}; \textbackslash{}\textbackslash{}text\{closed{-}loop system is stable\}",}
\NormalTok{    note: "Negative real part means decaying oscillation rather than growing oscillation." \}}
\NormalTok{])}
\end{Highlighting}
\end{Shaded}

\subsection{Exercise 3}\label{exercise-3-1}

A lab instrument is stable but too slow. The current proportional gain
gives a closed-loop pole at \(s=-0.2\). The engineer proposes increasing
the gain so the pole moves to \(s=-0.8\).

Write a short design note answering:

\begin{enumerate}
\def\labelenumi{\arabic{enumi}.}
\tightlist
\item
  What qualitative change will the lab see in the time response?
\item
  Why might this still be a bad idea if the sensor is noisy or delayed?
\end{enumerate}

\begin{Shaded}
\begin{Highlighting}[]
\NormalTok{//| echo: false}
\NormalTok{makeStepperHTML(3, [}
\NormalTok{  \{ op: "Interpret the pole location",}
\NormalTok{    eq: "s={-}0.8 \textbackslash{}\textbackslash{}text\{ is farther left than \} s={-}0.2",}
\NormalTok{    note: "Farther left in the complex plane means faster exponential decay." \},}
\NormalTok{  \{ op: "State the time{-}response consequence",}
\NormalTok{    eq: "\textbackslash{}\textbackslash{}text\{settling becomes faster, so the instrument returns to target more quickly\}",}
\NormalTok{    note: "A more negative real pole gives a shorter time constant." \},}
\NormalTok{  \{ op: "State the risk",}
\NormalTok{    eq: "\textbackslash{}\textbackslash{}text\{higher gain can amplify measurement noise and expose unmodelled delay\}",}
\NormalTok{    note: "A controller that looks better in the simple model may behave worse in the real loop." \},}
\NormalTok{  \{ op: "Check",}
\NormalTok{    eq: "\textbackslash{}\textbackslash{}text\{better nominal speed does not guarantee better robustness\}",}
\NormalTok{    note: "This is the main control tradeoff the design note should communicate." \}}
\NormalTok{])}
\end{Highlighting}
\end{Shaded}

\subsection{Exercise 4}\label{exercise-4-1}

Choose one system from your own field: a drive, drone, room heater,
queueing server with autoscaling, or experimental instrument. Identify:

\begin{enumerate}
\def\labelenumi{\arabic{enumi}.}
\tightlist
\item
  the reference
\item
  the measured output
\item
  the error signal
\item
  the control input
\item
  one reason high gain might help
\item
  one reason high gain might hurt
\end{enumerate}

Write the answer as a one-page systems sketch, not as prose only.

\emph{The diagram below shows the same unity-feedback loop relabelled
for five different domains. Use it as a template for your sketch.}

\begin{Shaded}
\begin{Highlighting}[]
\NormalTok{//| echo: false}
\NormalTok{feedbackVocabViz = \{}
\NormalTok{  const d3 = await require("d3@7");}

\NormalTok{  const W = 620, H = 260;}
\NormalTok{  const AMBER = "\#f59e0b";}

\NormalTok{  const domains = \{}
\NormalTok{    motor: \{}
\NormalTok{      label: "Motor speed",}
\NormalTok{      ref: "target RPM",}
\NormalTok{      error: "speed error",}
\NormalTok{      ctrl: "drive current",}
\NormalTok{      dist: "load torque",}
\NormalTok{      out: "measured RPM"}
\NormalTok{    \},}
\NormalTok{    drone: \{}
\NormalTok{      label: "Drone altitude",}
\NormalTok{      ref: "target altitude",}
\NormalTok{      error: "altitude error",}
\NormalTok{      ctrl: "rotor thrust",}
\NormalTok{      dist: "wind gust",}
\NormalTok{      out: "measured altitude"}
\NormalTok{    \},}
\NormalTok{    heater: \{}
\NormalTok{      label: "Room heater",}
\NormalTok{      ref: "setpoint temperature",}
\NormalTok{      error: "temp error",}
\NormalTok{      ctrl: "heating power",}
\NormalTok{      dist: "heat loss",}
\NormalTok{      out: "measured temp"}
\NormalTok{    \},}
\NormalTok{    instrument: \{}
\NormalTok{      label: "Lab instrument",}
\NormalTok{      ref: "target intensity",}
\NormalTok{      error: "intensity error",}
\NormalTok{      ctrl: "actuator signal",}
\NormalTok{      dist: "thermal drift",}
\NormalTok{      out: "measured output"}
\NormalTok{    \},}
\NormalTok{    server: \{}
\NormalTok{      label: "Server autoscaling",}
\NormalTok{      ref: "target QPS",}
\NormalTok{      error: "load error",}
\NormalTok{      ctrl: "instance count",}
\NormalTok{      dist: "traffic spike",}
\NormalTok{      out: "measured QPS"}
\NormalTok{    \}}
\NormalTok{  \};}

\NormalTok{  const container = d3.create("div")}
\NormalTok{    .style("font{-}family", "sans{-}serif")}
\NormalTok{    .style("max{-}width", W + "px");}

\NormalTok{  // Selector}
\NormalTok{  const selRow = container.append("div")}
\NormalTok{    .style("display", "flex").style("align{-}items", "center")}
\NormalTok{    .style("gap", "0.75rem").style("margin{-}bottom", "0.5rem");}
\NormalTok{  selRow.append("label").style("font{-}size", "0.9rem").text("Domain:");}
\NormalTok{  const sel = selRow.append("select").style("font{-}size", "0.9rem").style("padding", "2px 6px");}
\NormalTok{  Object.entries(domains).forEach(([k, v]) =\textgreater{} \{}
\NormalTok{    sel.append("option").attr("value", k).text(v.label);}
\NormalTok{  \});}

\NormalTok{  const svg = container.append("svg")}
\NormalTok{    .attr("width", W).attr("height", H)}
\NormalTok{    .attr("viewBox", \textasciigrave{}0 0 $\{W\} $\{H\}\textasciigrave{})}
\NormalTok{    .attr("aria{-}label", "Feedback vocabulary relabeller");}

\NormalTok{  // Defs for arrowhead}
\NormalTok{  const defs = svg.append("defs");}
\NormalTok{  defs.append("marker")}
\NormalTok{    .attr("id", "arrowFV").attr("markerWidth", 8).attr("markerHeight", 8)}
\NormalTok{    .attr("refX", 6).attr("refY", 3).attr("orient", "auto")}
\NormalTok{    .append("path").attr("d", "M0,0 L0,6 L8,3 z").attr("fill", "\#333");}

\NormalTok{  // Fixed layout}
\NormalTok{  const SJ = \{x: 120, y: 120\};}
\NormalTok{  const CB = \{x: 230, y: 100, w: 90, h: 40\};}
\NormalTok{  const PB = \{x: 400, y: 100, w: 90, h: 40\};}
\NormalTok{  const outX = 540;}
\NormalTok{  const fbY = 185;}

\NormalTok{  // Draw static structure}
\NormalTok{  function arrow(x1, y1, x2, y2) \{}
\NormalTok{    svg.append("line")}
\NormalTok{      .attr("x1", x1).attr("y1", y1).attr("x2", x2).attr("y2", y2)}
\NormalTok{      .attr("stroke", "\#374151").attr("stroke{-}width", 1.5)}
\NormalTok{      .attr("marker{-}end", "url(\#arrowFV)");}
\NormalTok{  \}}

\NormalTok{  // r {-}\textgreater{} SJ}
\NormalTok{  arrow(40, SJ.y, SJ.x {-} 18, SJ.y);}
\NormalTok{  // SJ {-}\textgreater{} CB}
\NormalTok{  arrow(SJ.x + 18, SJ.y, CB.x, SJ.y);}
\NormalTok{  // CB {-}\textgreater{} PB}
\NormalTok{  arrow(CB.x + CB.w, SJ.y, PB.x, SJ.y);}
\NormalTok{  // PB {-}\textgreater{} out}
\NormalTok{  arrow(PB.x + PB.w, SJ.y, outX, SJ.y);}
\NormalTok{  // feedback}
\NormalTok{  svg.append("polyline")}
\NormalTok{    .attr("points", \textasciigrave{}$\{outX {-} 5\},$\{SJ.y\} $\{outX {-} 5\},$\{fbY\} $\{SJ.x\},$\{fbY\} $\{SJ.x\},$\{SJ.y + 18\}\textasciigrave{})}
\NormalTok{    .attr("fill", "none").attr("stroke", "\#374151").attr("stroke{-}width", 1.5)}
\NormalTok{    .attr("marker{-}end", "url(\#arrowFV)");}

\NormalTok{  // Summing junction}
\NormalTok{  svg.append("circle").attr("cx", SJ.x).attr("cy", SJ.y).attr("r", 18)}
\NormalTok{    .attr("fill", "none").attr("stroke", "\#374151").attr("stroke{-}width", 1.5);}
\NormalTok{  svg.append("text").attr("x", SJ.x {-} 6).attr("y", SJ.y + 5)}
\NormalTok{    .attr("font{-}size", "14px").attr("fill", "\#374151").text("+");}
\NormalTok{  svg.append("text").attr("x", SJ.x {-} 5).attr("y", SJ.y + 22)}
\NormalTok{    .attr("font{-}size", "11px").attr("fill", "\#ef4444").text("−");}

\NormalTok{  // Controller box}
\NormalTok{  svg.append("rect").attr("x", CB.x).attr("y", CB.y).attr("width", CB.w).attr("height", CB.h)}
\NormalTok{    .attr("rx", 6).attr("fill", "\#f3f4f6").attr("stroke", "\#9ca3af").attr("stroke{-}width", 1.5);}
\NormalTok{  svg.append("text").attr("x", CB.x + CB.w / 2).attr("y", CB.y + CB.h / 2 + 5)}
\NormalTok{    .attr("text{-}anchor", "middle").attr("font{-}size", "13px").attr("fill", "\#2563eb").text("Controller");}

\NormalTok{  // Plant box}
\NormalTok{  svg.append("rect").attr("x", PB.x).attr("y", PB.y).attr("width", PB.w).attr("height", PB.h)}
\NormalTok{    .attr("rx", 6).attr("fill", "\#f3f4f6").attr("stroke", "\#9ca3af").attr("stroke{-}width", 1.5);}
\NormalTok{  svg.append("text").attr("x", PB.x + PB.w / 2).attr("y", PB.y + PB.h / 2 + 5)}
\NormalTok{    .attr("text{-}anchor", "middle").attr("font{-}size", "13px").attr("fill", "\#374151").text("Plant");}

\NormalTok{  // Dynamic labels}
\NormalTok{  const refLbl = svg.append("text").attr("x", 25).attr("y", SJ.y {-} 10)}
\NormalTok{    .attr("font{-}size", "11px").attr("text{-}anchor", "middle");}
\NormalTok{  const errLbl = svg.append("text").attr("x", SJ.x + 55).attr("y", SJ.y {-} 10)}
\NormalTok{    .attr("font{-}size", "11px").attr("text{-}anchor", "middle");}
\NormalTok{  const ctrlLbl = svg.append("text").attr("x", CB.x + CB.w + 45).attr("y", SJ.y {-} 10)}
\NormalTok{    .attr("font{-}size", "11px").attr("text{-}anchor", "middle");}
\NormalTok{  const distLbl = svg.append("text").attr("x", (PB.x + PB.x + PB.w) / 2).attr("y", PB.y {-} 8)}
\NormalTok{    .attr("font{-}size", "11px").attr("text{-}anchor", "middle").attr("fill", "\#6b7280");}
\NormalTok{  const outLbl = svg.append("text").attr("x", outX + 5).attr("y", SJ.y + 5)}
\NormalTok{    .attr("font{-}size", "11px").attr("text{-}anchor", "start");}

\NormalTok{  // Disturbance arrow}
\NormalTok{  svg.append("line")}
\NormalTok{    .attr("x1", (PB.x + PB.x + PB.w) / 2).attr("y1", PB.y {-} 30)}
\NormalTok{    .attr("x2", (PB.x + PB.x + PB.w) / 2).attr("y2", PB.y)}
\NormalTok{    .attr("stroke", "\#6b7280").attr("stroke{-}width", 1.2).attr("stroke{-}dasharray", "3,2")}
\NormalTok{    .attr("marker{-}end", "url(\#arrowFV)");}

\NormalTok{  let prevLabels = \{\};}

\NormalTok{  function flashLabel(el, newText) \{}
\NormalTok{    if (prevLabels[el.attr("class") || el.property("\_id")] === newText) return;}
\NormalTok{    el.attr("fill", AMBER);}
\NormalTok{    el.text(newText);}
\NormalTok{    setTimeout(() =\textgreater{} el.attr("fill", "\#374151"), 800);}
\NormalTok{  \}}

\NormalTok{  function updateDomain(key) \{}
\NormalTok{    const d = domains[key];}
\NormalTok{    flashLabel(refLbl, d.ref);}
\NormalTok{    flashLabel(errLbl, d.error);}
\NormalTok{    flashLabel(ctrlLbl, d.ctrl);}
\NormalTok{    flashLabel(distLbl, d.dist);}
\NormalTok{    flashLabel(outLbl, d.out);}
\NormalTok{    prevLabels = d;}
\NormalTok{  \}}

\NormalTok{  sel.on("change", function() \{ updateDomain(this.value); \});}
\NormalTok{  updateDomain("motor");}

\NormalTok{  return container.node();}
\NormalTok{\}}
\end{Highlighting}
\end{Shaded}

\chapter{Discrete-time systems and signal
processing}\label{discrete-time-systems-and-signal-processing}

Sampling, filtering, and recurrence

\hfill\break

A sensor does not hand you a continuous function. It hands you numbers
at intervals. A microphone gives a stream of samples. A controller
updates once per cycle. Once time is sampled, the mathematics changes
shape: derivatives become recurrences, spectra pick up aliasing, and
filters become algorithms.

That shift matters because upper-year engineering is full of systems
that are modelled in continuous time but built in discrete time. A motor
may obey a differential equation. The controller that stabilises it
still runs on a clock. A seismic wave may be continuous in the ground.
The data still arrives as a sampled trace.

This chapter is about that translation. The goal is not to memorise
signal processing vocabulary. It is to see the structural move: sampling
turns a continuous object into a sequence, and once you have a sequence,
the natural questions become delay, recurrence, convolution, spectrum,
and stability.

\begin{center}\rule{0.5\linewidth}{0.5pt}\end{center}

\section{The discrete-time view}\label{the-discrete-time-view}

A continuous-time signal is written \(x(t)\). A discrete-time signal is
written \(x[n]\), read ``x of n.'' The variable \(n\) is not time
itself. It is an index. The physical time attached to sample \(n\) is

\[t_n = nT_s\]

where \(T_s\) is the \textbf{sampling period}. The sampling frequency is

\[f_s = \frac{1}{T_s}\]

The same physical system can therefore be described in two ways:

\begin{itemize}
\tightlist
\item
  continuous-time: what the signal is doing between samples
\item
  discrete-time: what the sampled sequence records at the clock instants
\end{itemize}

This sounds like a harmless change of notation. It is not. Once the
signal is a sequence, operations that looked infinitesimal in continuous
time become difference equations and delay operators.

If a discrete-time system takes an input sequence \(x[n]\) and produces
an output sequence \(y[n]\), the simplest and most important class is
the \textbf{linear time-invariant} discrete system. For those systems,
present output depends on present and past samples in a fixed way.

One common form is a \textbf{finite impulse response} (FIR) filter:

\[y[n] = \sum_{k=0}^{M} b_k x[n-k]\]

``Finite'' here means the output depends on only a bounded window of
\(M+1\) past inputs, so the filter's memory is bounded. This is a
weighted moving combination of current and previous input samples. In
plain language: the output at time step \(n\) is built from a short
memory of the input. Filters that also feed back past outputs are called
infinite impulse response (IIR) and have unbounded memory; they arise in
the z-transform and state-space frameworks later in the chapter.

\section{Why sampling can lie to you}\label{why-sampling-can-lie-to-you}

Suppose the continuous signal is

\[x(t) = \sin(2\pi f t)\]

and you sample it every \(T_s\) seconds. Then the discrete sequence is

\[x[n] = \sin(2\pi f nT_s) = \sin\left(2\pi \frac{f}{f_s} n\right)\]

The discrete sequence only sees the ratio \(f/f_s\). This is the
beginning of the aliasing problem.

If the sampling frequency is too low, two different continuous
frequencies can produce sampled sequences that are indistinguishable.
The samples no longer carry enough information to tell you which
continuous oscillation was actually present.

The safe rule is the \textbf{Nyquist idea}:

\begin{itemize}
\tightlist
\item
  to represent frequency content up to \(f_{\max}\), sample faster than
  \(2f_{\max}\)
\end{itemize}

That does not solve every practical problem, but it tells you where
obvious failure begins. Sampling below that rate invites aliasing:
high-frequency content reappears as a false lower-frequency pattern in
the data.

This is one of the reasons discrete-time systems need their own
mathematics. The sampled sequence is not a transparent window into the
continuous signal. It is a transformed object with its own constraints.

\section{Difference equations and
recurrence}\label{difference-equations-and-recurrence}

Continuous-time linear systems are often described by differential
equations. Discrete-time linear systems are naturally described by
\textbf{difference equations}.

A first-order example is

\[y[n+1] = ay[n] + bu[n]\]

This says: the next output is built from the current output and the
current input. The constant \(a\) carries memory. The constant \(b\)
tells you how strongly the input influences the update.

The stability rule for the homogeneous system

\[y[n+1] = ay[n]\]

is very simple:

\begin{itemize}
\tightlist
\item
  if \(|a| < 1\), the sequence decays to zero
\item
  if \(|a| > 1\), the sequence grows in magnitude
\item
  if \(|a| = 1\), the sequence neither decays nor grows and small
  modelling changes can matter a great deal
\end{itemize}

This is the discrete-time analogue of looking at pole location in
continuous time. In one case you care about left-half-plane location. In
the other you care about whether the relevant quantity sits inside the
unit circle.

\begin{tcolorbox}[enhanced jigsaw, opacityback=0, colback=white, breakable, colframe=quarto-callout-note-color-frame, arc=.35mm, left=2mm, toprule=.15mm, leftrule=.75mm, bottomrule=.15mm, rightrule=.15mm]
\begin{minipage}[t]{5.5mm}
\textcolor{quarto-callout-note-color}{\faInfo}
\end{minipage}%
\begin{minipage}[t]{\textwidth - 5.5mm}

\vspace{-3mm}\textbf{Why This Works}\vspace{3mm}

In continuous time, a mode behaves like \(e^{st}\), so the sign of
\(\Re(s)\) decides whether it decays or grows. In discrete time, a mode
behaves like \(a^n\), so the magnitude of \(a\) decides whether repeated
updates shrink or blow up.

That is the structural translation. Exponentials become powers.
Left-half-plane stability becomes unit-circle stability.

\end{minipage}%
\end{tcolorbox}

\section{The core method}\label{the-core-method-1}

A first pass through a discrete-time systems problem usually goes like
this:

\begin{enumerate}
\def\labelenumi{\arabic{enumi}.}
\tightlist
\item
  Identify the sampling period \(T_s\) and therefore the sampling
  frequency \(f_s\).
\item
  Decide whether the data is a sampled version of a continuous signal, a
  genuinely discrete process, or a digital controller state.
\item
  Write the update rule or filter equation.
\item
  Interpret every term as memory, delay, or input contribution.
\item
  Check the relevant stability test: pole inside the unit circle, or
  coefficient magnitude less than 1 in the simple first-order case.
\item
  Connect the algebra back to behaviour: smoothing, lag, oscillation,
  alias, or amplification.
\end{enumerate}

The habit to build is the same one from the control chapter: do not
leave the recurrence as pure symbol manipulation. Ask what it means for
the sequence the machine is actually computing.

\section{Worked example 1: a moving-average sensor
filter}\label{worked-example-1-a-moving-average-sensor-filter}

A noisy sensor signal \(x[n]\) is filtered by

\[y[n] = \frac{1}{3}\bigl(x[n] + x[n-1] + x[n-2]\bigr)\]

This is a three-point moving average.

What does it do? It smooths out rapid sample-to-sample variation by
replacing the present sample with the average of the present and
previous two samples. If one sample spikes upward because of noise, the
output spike is diluted by the neighbouring samples.

The price is delay. A filter that averages over time cannot react
instantly to a real rapid change. This is the first recurring
engineering tradeoff in discrete-time signal processing:

\begin{itemize}
\tightlist
\item
  more smoothing usually means more lag
\item
  less lag usually means less smoothing
\end{itemize}

In embedded control, that tradeoff matters because a laggy filter can
make an otherwise good controller feel slow or even unstable.

\section{Worked example 2: discrete-time stability of a sampled
update}\label{worked-example-2-discrete-time-stability-of-a-sampled-update}

Consider the update rule

\[y[n+1] = 0.6\,y[n] + u[n]\]

If the input is zero, then

\[y[n+1] = 0.6\,y[n]\]

Starting from \(y[0]\), repeated substitution gives

\[y[n] = 0.6^n y[0]\]

Because \(|0.6| < 1\), the homogeneous response decays to zero. So the
system's memory fades over time.

Now compare with

\[y[n+1] = 1.1\,y[n] + u[n]\]

With zero input, the response is

\[y[n] = 1.1^n y[0]\]

and because \(|1.1| > 1\), the sequence grows. A digital implementation
with this update law is unstable.

That is a cleaner rule than many students expect. In first-order
discrete-time systems, the entire stability story is often carried by
one number.

\section{Worked example 3: aliasing in a seismic
record}\label{worked-example-3-aliasing-in-a-seismic-record}

A geophysicist samples ground motion at \(f_s = 100\) Hz. A strong
sinusoidal component is present at \(f = 70\) Hz.

This is above the Nyquist limit of

\[\frac{f_s}{2} = 50 \text{ Hz}\]

so the signal will alias. The discrete record cannot faithfully
represent a 70 Hz oscillation at that sampling rate. It will appear as a
false lower frequency in the sampled data.

The engineering lesson is not ``compute a formula and move on.'' It is
that data collection itself is part of the mathematics. Sampling choices
determine what questions the later analysis can answer honestly.

This is why signal processing belongs naturally beside estimation and
inverse problems in Volume 8. If the measurement process already lost
information, no later inference method can magically restore it.

\section{Where this goes}\label{where-this-goes-38}

The most direct continuation is
\href{05-estimation-inverse.qmd}{Estimation, inverse problems, and
filtering}. Once a system is sampled, state estimation, filtering, and
inference are usually carried out in discrete time. Kalman filtering,
recursive estimation, and streaming data assimilation all inherit the
language of recurrences, sampling, and noise that this chapter sets up.

This chapter also sits beside control rather than after it by accident.
Upper-year engineers constantly move between the continuous-time model
of the plant and the discrete-time implementation of the controller. If
Chapter 1 was about shaping behaviour, this chapter is about what
happens when that shaping must be done on a clock.

\begin{tcolorbox}[enhanced jigsaw, opacityback=0, colback=white, breakable, colframe=quarto-callout-tip-color-frame, arc=.35mm, left=2mm, toprule=.15mm, leftrule=.75mm, bottomrule=.15mm, rightrule=.15mm]
\begin{minipage}[t]{5.5mm}
\textcolor{quarto-callout-tip-color}{\faLightbulb}
\end{minipage}%
\begin{minipage}[t]{\textwidth - 5.5mm}

\vspace{-3mm}\textbf{Applications}\vspace{3mm}

\begin{itemize}
\tightlist
\item
  digital filters for noisy sensors
\item
  sampled-data control in motors and robotics
\item
  audio and vibration signal analysis
\item
  seismic and geophysical trace processing
\item
  preprocessing streaming time-series data
\item
  telemetry, embedded systems, and real-time monitoring
\end{itemize}

\end{minipage}%
\end{tcolorbox}

\section{Exercises}\label{exercises-49}

These are project-style exercises. Explain the behaviour as well as the
calculation.

\subsection{Exercise 1}\label{exercise-1-2}

A sensor is sampled every \(0.01\) s.

\begin{enumerate}
\def\labelenumi{\arabic{enumi}.}
\tightlist
\item
  Find the sampling frequency.
\item
  Find the Nyquist frequency.
\item
  Decide whether a 30 Hz sinusoid can be represented without aliasing.
\item
  Decide whether a 70 Hz sinusoid can be represented without aliasing.
\end{enumerate}

\subsection{Exercise 2}\label{exercise-2-2}

Consider the discrete-time system

\[y[n+1] = 0.8\,y[n] + 2u[n]\]

\begin{enumerate}
\def\labelenumi{\arabic{enumi}.}
\tightlist
\item
  Determine whether the homogeneous system is stable.
\item
  Starting from \(y[0] = 5\) with \(u[n] = 0\), compute \(y[1]\),
  \(y[2]\), and \(y[3]\).
\item
  Explain what the sequence is doing physically.
\end{enumerate}

\subsection{Exercise 3}\label{exercise-3-2}

A noisy accelerometer is filtered using

\[y[n] = \frac{1}{4}\bigl(x[n] + x[n-1] + x[n-2] + x[n-3]\bigr)\]

Write a short engineering memo answering:

\begin{enumerate}
\def\labelenumi{\arabic{enumi}.}
\tightlist
\item
  What kind of behaviour this filter is trying to suppress
\item
  Why it introduces lag
\item
  Why that lag could matter in a feedback controller
\end{enumerate}

\subsection{Exercise 4}\label{exercise-4-2}

Choose one sampled-data setting from your field: a microphone, IMU,
stock-price stream, satellite time series, or power-system sensor.

Prepare a one-page systems sketch naming:

\begin{enumerate}
\def\labelenumi{\arabic{enumi}.}
\tightlist
\item
  the physical quantity being sampled
\item
  a plausible sampling frequency
\item
  one source of aliasing or information loss
\item
  one filter or recurrence you might apply
\item
  one reason filtering helps
\item
  one reason filtering can distort what you care about
\end{enumerate}

\chapter{Continuum mechanics and
transport}\label{continuum-mechanics-and-transport}

Field laws for stress, flow, heat, and diffusion

\hfill\break

Heat leaves a wall, stress moves through a beam, solute spreads through
a river, and pressure drives flow through a pipe. Different materials,
different scales, different equations on the page. Underneath, the same
move repeats: write what is conserved, write how the material responds,
and let the field describe the rest.

This chapter is where upper-year engineering starts to feel less like a
list of subjects and more like a single modelling language. Structural
analysis, heat transfer, porous flow, and contaminant transport do not
share the same units or applications. They do share the same
mathematical architecture.

The architecture has two parts. First, a \textbf{conservation law} says
something cannot appear or disappear except through flow across a
boundary or production inside the region. Second, a \textbf{constitutive
law} says how the material reacts: how heat flows down a temperature
gradient, how solute diffuses down a concentration gradient, how stress
relates to strain in a simple elastic solid.

\begin{center}\rule{0.5\linewidth}{0.5pt}\end{center}

\section{The local balance law}\label{the-local-balance-law}

Take a small control volume inside a body, pipe, soil column, or fluid
region. Let \(u\) be the amount of something stored per unit volume.
Depending on the problem, \(u\) might be thermal energy, mass
concentration, momentum density, or another conserved quantity.

Let \(\mathbf{q}\) be the \textbf{flux} of that quantity. Flux is a
vector quantity: its magnitude measures how much passes through a unit
area per unit time, and its direction is the direction of net transport.

The generic local balance law is

\[\frac{\partial u}{\partial t} + \nabla \cdot \mathbf{q} = s\]

where:

\begin{itemize}
\tightlist
\item
  \(\partial u/\partial t\) is the local rate of storage
\item
  \(\nabla \cdot \mathbf{q}\) is the net outward flux
\item
  \(s\) is a source term
\end{itemize}

This one line is the template behind a large fraction of continuum
modelling. Change the meaning of \(u\), \(\mathbf{q}\), and \(s\), and
you get a different physical system. The bookkeeping structure stays the
same.

The divergence term matters because a continuum model is local. It does
not ask ``how much crossed the outer boundary of the whole pipe?''
first. It asks what is happening in an arbitrarily small piece. Once the
local law is right, the global law follows by integration.

\section{Constitutive laws}\label{constitutive-laws}

A balance law alone is not enough. It tells you what must be accounted
for, not how the material actually behaves. That extra relation is the
constitutive law.

Two standard examples are:

\textbf{Fourier's law of heat conduction}

\[\mathbf{q} = -k\nabla T\]

Heat flux points down the temperature gradient. The constant \(k\) is
thermal conductivity. The minus sign means heat flows from hot to cold.

\textbf{Fick's law of diffusion}

\[\mathbf{J} = -D\nabla c\]

Mass flux points down the concentration gradient. The constant \(D\) is
the diffusivity.

The structure is the same in both. A gradient creates a driving force.
The material constant tells you how strongly the medium responds.

In elementary elasticity, the constitutive law is often written as a
linear stress-strain relation. In one dimension,

\[\sigma = E\varepsilon\]

where \(\sigma\) is \textbf{stress} (force per unit area, Pa),
\(\varepsilon\) is \textbf{strain} (fractional change in length,
dimensionless), and \(E\) is Young's modulus (the material stiffness).

Again, the structure is familiar: one field quantity is related to
another by a material response parameter.

\section{Why signs and directions
matter}\label{why-signs-and-directions-matter}

Continuum laws are full of minus signs, outward normals, and directional
derivatives. These are not decorative details. They are the physics.

If temperature increases to the right, then \(\nabla T\) points right.
Fourier's law says the heat flux points left. Miss the minus sign and
you literally model heat flowing from cold to hot.

The same warning applies everywhere in this chapter. Continuum
mathematics is geometry with units attached. Direction is part of the
answer.

\begin{tcolorbox}[enhanced jigsaw, opacityback=0, colback=white, breakable, colframe=quarto-callout-note-color-frame, arc=.35mm, left=2mm, toprule=.15mm, leftrule=.75mm, bottomrule=.15mm, rightrule=.15mm]
\begin{minipage}[t]{5.5mm}
\textcolor{quarto-callout-note-color}{\faInfo}
\end{minipage}%
\begin{minipage}[t]{\textwidth - 5.5mm}

\vspace{-3mm}\textbf{Why This Works}\vspace{3mm}

Conservation laws come from bookkeeping on a region. Constitutive laws
come from empirical or theoretical descriptions of how a material
responds. The PDE appears when you combine them.

For example, if thermal energy density is proportional to temperature
and heat flux obeys Fourier's law, then storage plus divergence of heat
flux produces the heat equation. The PDE is not the starting point. It
is the compressed result of conservation plus material behaviour.

\end{minipage}%
\end{tcolorbox}

\section{The core method}\label{the-core-method-2}

A first pass through a continuum or transport problem usually goes like
this:

\begin{enumerate}
\def\labelenumi{\arabic{enumi}.}
\tightlist
\item
  Identify the conserved quantity: mass, energy, momentum, or solute.
\item
  Write the local balance law for storage, flux, and source.
\item
  Choose the constitutive law that links flux or stress to the driving
  field.
\item
  Substitute the constitutive law into the balance law.
\item
  Inspect the resulting PDE or field equation and connect each term back
  to a physical mechanism.
\item
  Add boundary and initial conditions only after the governing structure
  is clear.
\end{enumerate}

That last point matters. Students often see the PDE first and the
modelling logic second. In upper-year work, reversing that order helps.
If you know where every term came from, the equation stops feeling
arbitrary.

\section{Worked example 1: heat conduction in one
dimension}\label{worked-example-1-heat-conduction-in-one-dimension}

Let \(T(x,t)\) be temperature in a rod. Suppose there is no internal
heat source. The local energy balance says:

\[\frac{\partial u}{\partial t} + \nabla \cdot \mathbf{q} = 0\]

In one dimension, and with thermal energy density proportional to
temperature, this becomes

\[\rho c_p \frac{\partial T}{\partial t} + \frac{\partial q}{\partial x} = 0\]

where \(\rho\) is mass density (kg m\(^{-3}\)) and \(c_p\) is specific
heat capacity (J kg\(^{-1}\) K\(^{-1}\)) --- the two material parameters
relating energy storage to temperature change.

Use Fourier's law in one dimension:

\[q = -k\frac{\partial T}{\partial x}\]

Substitute:

\[\rho c_p \frac{\partial T}{\partial t}
 - \frac{\partial}{\partial x}\left(k\frac{\partial T}{\partial x}\right) = 0\]

If \(k\) is constant, this simplifies to

\[\rho c_p \frac{\partial T}{\partial t}
 = k\frac{\partial^2 T}{\partial x^2}\]

or

\[\frac{\partial T}{\partial t} = \alpha \frac{\partial^2 T}{\partial x^2},
\qquad \alpha = \frac{k}{\rho c_p}\]

This is the heat equation. The second spatial derivative tells you how
much the temperature field differs from its local linear trend. Where
the field is strongly curved, temperature changes quickly in time.

\section{Worked example 2: contaminant diffusion in
groundwater}\label{worked-example-2-contaminant-diffusion-in-groundwater}

Let \(c(x,t)\) be solute concentration in a one-dimensional porous
medium. Ignore advection for the moment and assume no chemical reaction
source.

The balance law is

\[\frac{\partial c}{\partial t} + \frac{\partial J}{\partial x} = 0\]

Use Fick's law:

\[J = -D\frac{\partial c}{\partial x}\]

Substitute:

\[\frac{\partial c}{\partial t}
 - \frac{\partial}{\partial x}\left(D\frac{\partial c}{\partial x}\right)=0\]

For constant diffusivity:

\[\frac{\partial c}{\partial t} = D\frac{\partial^2 c}{\partial x^2}\]

Mathematically this is the same diffusion structure as the heat
equation. Physically the field means something different. That is
exactly why this chapter matters: it trains you to recognise when two
different-looking systems share the same mathematical skeleton.

\section{Worked example 3: linear elasticity in one
dimension}\label{worked-example-3-linear-elasticity-in-one-dimension}

Take a slender bar under axial load. In the simplest small-strain
setting, strain is

\[\varepsilon = \frac{du}{dx}\]

where \(u(x)\) is displacement. Hooke's law gives

\[\sigma = E\varepsilon = E\frac{du}{dx}\]

If body forces are negligible and the bar is in static equilibrium,
force balance requires axial stress to satisfy

\[\frac{d\sigma}{dx} = 0\]

Substitute the constitutive relation:

\[\frac{d}{dx}\left(E\frac{du}{dx}\right) = 0\]

For constant \(E\):

\[E\frac{d^2u}{dx^2} = 0 \qquad \Rightarrow \qquad \frac{d^2u}{dx^2} = 0\]

This is much simpler than the heat and diffusion examples, but the logic
is the same. Balance law plus constitutive law gives the governing
equation.

\section{Where this goes}\label{where-this-goes-39}

The most direct continuation is
\href{04-computational-methods.qmd}{Computational methods for
engineering models}. Once the governing equations are written, most
upper-year work shifts from derivation to solution. Real geometries,
mixed materials, and realistic boundary conditions usually push these
problems toward numerical methods, meshes, and solver design.

This chapter also sets up a habit that carries far beyond classical
engineering. Environmental transport, remote sensing inversion, and
geophysics all use the same structure. The quantity and units change.
The control-volume logic does not.

\begin{tcolorbox}[enhanced jigsaw, opacityback=0, colback=white, breakable, colframe=quarto-callout-tip-color-frame, arc=.35mm, left=2mm, toprule=.15mm, leftrule=.75mm, bottomrule=.15mm, rightrule=.15mm]
\begin{minipage}[t]{5.5mm}
\textcolor{quarto-callout-tip-color}{\faLightbulb}
\end{minipage}%
\begin{minipage}[t]{\textwidth - 5.5mm}

\vspace{-3mm}\textbf{Applications}\vspace{3mm}

\begin{itemize}
\tightlist
\item
  heat flow through walls, fins, and machine components
\item
  diffusion and mixing in process equipment
\item
  contaminant transport in groundwater and rivers
\item
  stress and deformation in structural members
\item
  porous flow and subsurface energy transport
\item
  continuum models in earth and environmental systems
\end{itemize}

\end{minipage}%
\end{tcolorbox}

\section{Exercises}\label{exercises-50}

These are project-style exercises. State the physical meaning of each
term as well as the final equation.

\subsection{Exercise 1}\label{exercise-1-3}

Starting from the one-dimensional heat balance

\[\rho c_p \frac{\partial T}{\partial t} + \frac{\partial q}{\partial x} = 0\]

and Fourier's law

\[q = -k\frac{\partial T}{\partial x}\]

derive the heat equation for constant \(k\).

\subsection{Exercise 2}\label{exercise-2-3}

Write the diffusion equation for concentration \(c(x,t)\) from

\[\frac{\partial c}{\partial t} + \frac{\partial J}{\partial x} = 0,
\qquad
J = -D\frac{\partial c}{\partial x}\]

assuming \(D\) is constant. Then explain why the equation has the same
form as the heat equation.

\subsection{Exercise 3}\label{exercise-3-3}

For a one-dimensional elastic bar with

\[\varepsilon = \frac{du}{dx}, \qquad \sigma = E\varepsilon\]

and static equilibrium

\[\frac{d\sigma}{dx} = 0,\]

derive the governing equation for displacement when \(E\) is constant.

\subsection{Exercise 4}\label{exercise-4-3}

Choose one field setting from your own area: a wall, pipe, aquifer,
beam, soil column, battery, or river reach.

Prepare a one-page systems sketch naming:

\begin{enumerate}
\def\labelenumi{\arabic{enumi}.}
\tightlist
\item
  the conserved quantity
\item
  the flux
\item
  the constitutive law
\item
  one plausible source term
\item
  one important boundary condition
\item
  one reason the resulting equation will probably need numerical
  solution
\end{enumerate}

\chapter{Computational methods for engineering
models}\label{computational-methods-for-engineering-models}

Meshes, weak forms, and conservation in computation

\hfill\break

The governing equations may be right and still produce a useless
simulation. A coarse mesh smears a boundary layer. A bad time step
destabilises a model. A solver converges to something, but not to the
physics you meant. Upper-year engineering lives in that gap between the
equation on paper and the model you actually trust.

This chapter is about the mathematical move that closes that gap. A
continuum equation is not yet computable. A computer needs finitely many
unknowns, finitely many equations, and an update rule or linear system
it can actually solve. Discretisation supplies that finite structure.

The key idea is that every discretisation keeps some truths and
sacrifices others. One method may preserve conservation neatly. Another
may represent curved geometry more naturally. Another may be easy to
implement but unstable for large time steps. Trustworthy simulation is
the art of knowing which sacrifices are acceptable.

\begin{center}\rule{0.5\linewidth}{0.5pt}\end{center}

\section{From equation to algebraic
system}\label{from-equation-to-algebraic-system}

Take the one-dimensional diffusion equation

\[\frac{\partial u}{\partial t} = D\frac{\partial^2 u}{\partial x^2}\]

defined on a grid with spatial spacing \(\Delta x\) and time step
\(\Delta t\). The continuous field \(u(x,t)\) becomes a table of
approximate values \(u_i^n\), read ``u at grid point i and time level
n.''

Using a forward difference in time and a central difference in space:

\[\frac{u_i^{n+1} - u_i^n}{\Delta t}
\approx
D\frac{u_{i+1}^n - 2u_i^n + u_{i-1}^n}{(\Delta x)^2}\]

Rearrange:

\[u_i^{n+1}
=
u_i^n
+
\frac{D\Delta t}{(\Delta x)^2}
\left(u_{i+1}^n - 2u_i^n + u_{i-1}^n\right)\]

The PDE has become a recurrence. Once the spatial grid is fixed, the
whole problem is now algebra.

This is the central translation of computational engineering:

\begin{itemize}
\tightlist
\item
  fields become arrays
\item
  derivatives become finite differences, flux balances, or
  basis-function integrals
\item
  boundary conditions become algebraic constraints
\item
  the model becomes either a time-marching update or a linear/nonlinear
  system
\end{itemize}

\section{Three dominant viewpoints}\label{three-dominant-viewpoints}

Upper-year engineering simulation is often organised around three
families of discretisation.

\textbf{Finite difference}

Replace derivatives directly with local difference quotients on a grid.
Conceptually simple. Especially natural on structured meshes and
textbook geometries.

\textbf{Finite volume}

Integrate the balance law over control volumes, then approximate fluxes
at faces. Especially attractive when conservation matters strongly, as
in fluid flow, transport, and environmental modelling.

\textbf{Finite element}

Represent the solution using basis functions \(\phi_i\) over elements
and solve a weak or variational form (a reformulation that replaces the
PDE with an integral condition, trading pointwise accuracy for broader
geometric flexibility). Especially useful on complicated geometries and
in structural mechanics.

These are not three unrelated subjects. They are three ways of answering
the same question: how should a continuum law be turned into finite
equations?

\section{Stability, consistency, and
convergence}\label{stability-consistency-and-convergence}

A discretisation is not good merely because it can be coded.

Three recurring tests matter:

\begin{itemize}
\tightlist
\item
  \textbf{consistency}: does the discrete equation approach the
  continuous one as the mesh is refined?
\item
  \textbf{stability}: do errors remain controlled instead of blowing up?
\item
  \textbf{convergence}: does the computed solution approach the true
  solution as the grid is refined?
\end{itemize}

Students often meet these separately. In practice they are linked. A
consistent method that is unstable may be useless. A stable method that
is too diffusive may hide the physics you care about.

For the explicit diffusion update above, the ratio

\[r = \frac{D\Delta t}{(\Delta x)^2}\]

controls stability. For the 1D explicit diffusion scheme the stability
bound is

\[r = \frac{D\Delta t}{(\Delta x)^2} \leq \frac{1}{2}\]

If \(r\) exceeds \(\frac{1}{2}\), the numerical solution can oscillate
or blow up even when the physical problem is perfectly well behaved.
This bound means halving the grid spacing forces the time step to shrink
by a factor of four --- a practical cost of explicit time-stepping.

\begin{tcolorbox}[enhanced jigsaw, opacityback=0, colback=white, breakable, colframe=quarto-callout-note-color-frame, arc=.35mm, left=2mm, toprule=.15mm, leftrule=.75mm, bottomrule=.15mm, rightrule=.15mm]
\begin{minipage}[t]{5.5mm}
\textcolor{quarto-callout-note-color}{\faInfo}
\end{minipage}%
\begin{minipage}[t]{\textwidth - 5.5mm}

\vspace{-3mm}\textbf{Why This Works}\vspace{3mm}

Discretisation replaces infinitely many degrees of freedom with finitely
many. That replacement introduces error automatically. The right
question is never ``is there error?'' but ``what kind of error, how
large, and does it destroy the physical meaning of the model?''

Computational mathematics is therefore not separate from modelling. The
method is part of the model.

\end{minipage}%
\end{tcolorbox}

\section{The core method}\label{the-core-method-3}

A first pass through a computational engineering model usually goes like
this:

\begin{enumerate}
\def\labelenumi{\arabic{enumi}.}
\tightlist
\item
  Write the governing equation and the boundary/initial conditions.
\item
  Decide which discrete viewpoint matches the physics and geometry best.
\item
  Define the grid, mesh, or basis functions.
\item
  Derive the algebraic update or system.
\item
  Check the main stability or conditioning constraint.
\item
  Interpret the numerical output in physical terms and test mesh
  sensitivity.
\end{enumerate}

The last step matters as much as the derivation. If halving \(\Delta x\)
changes the answer substantially, the previous mesh was telling you more
about the discretisation than about the system.

\section{Worked example 1: explicit diffusion
update}\label{worked-example-1-explicit-diffusion-update}

Take

\[\frac{\partial u}{\partial t} = D\frac{\partial^2 u}{\partial x^2}\]

with \(D = 1\), \(\Delta x = 0.1\), and \(\Delta t = 0.002\).

Then

\[r = \frac{D\Delta t}{(\Delta x)^2}
=
\frac{1(0.002)}{0.01}
= 0.2\]

For the standard explicit 1D diffusion scheme, the stability requirement
is \(r \leq \tfrac{1}{2}\). Since \(0.2 < 0.5\), this time step is
within the stable range and is plausibly acceptable for this grid.

Now change only the time step to \(\Delta t = 0.008\). Then

\[r = \frac{0.008}{0.01} = 0.8\]

and the scheme is no longer stable --- \(0.8 > 0.5\). Same PDE. Same
material parameter. Different numerical behaviour because the
discretisation changed.

This is why simulation judgement is mathematical judgement. The equation
on the page did not fail. The computational choice did.

\section{Worked example 2: a finite-volume mass
balance}\label{worked-example-2-a-finite-volume-mass-balance}

Suppose a pollutant concentration is tracked over a control volume of
width \(\Delta x\). The half-integer subscripts label the cell faces:
\(F_{i-1/2}\) is the flux at the left face of cell \(i\) (between cells
\(i-1\) and \(i\)), and \(F_{i+1/2}\) is the flux at the right face
(between cells \(i\) and \(i+1\)). Then a finite-volume update has the
form

\[\frac{c_i^{n+1} - c_i^n}{\Delta t}
=
-\frac{F_{i+1/2}^n - F_{i-1/2}^n}{\Delta x} + s_i^n\]

This is the discrete balance law in its most honest form:

\begin{itemize}
\tightlist
\item
  future storage minus current storage
\item
  equals net flux imbalance
\item
  plus any local source
\end{itemize}

The attraction of the finite-volume approach is that conservation is
built into the discretisation. If what leaves one cell enters the next,
the algebra itself remembers that.

That is why finite volume is so common in transport and environmental
systems. The method echoes the underlying physics instead of merely
approximating the derivatives.

\section{Worked example 3: why finite elements dominate structural
analysis}\label{worked-example-3-why-finite-elements-dominate-structural-analysis}

In a structural problem, geometry is often awkward: holes, supports,
joints, changing thickness, mixed materials. A regular grid is not
always natural.

Finite elements instead approximate the displacement field using basis
or shape functions:

\[u_h(x) = \sum_i U_i \phi_i(x)\]

The unknowns are the nodal coefficients \(U_i\). The governing equations
become a matrix system

\[K\mathbf{U} = \mathbf{f}\]

where \(K\) is the \textbf{stiffness matrix} --- a different \(K\) from
the control gain in Chapter 1, but the same letter is standard in both
fields.

The important mathematical point is not the assembly details. It is that
the continuum problem has been projected onto a finite-dimensional space
spanned by the \(\phi_i\). Geometry and material behaviour are therefore
represented through the choice of mesh and basis functions, not just
through a derivative stencil.

A second important point is that \(K\) is almost always \textbf{sparse}:
each basis function \(\phi_i\) overlaps only with neighbouring basis
functions, so most entries of \(K\) are zero. A mesh with millions of
nodes still produces a matrix with only a few non-zeros per row. This
sparsity is why large FEM problems are computationally tractable ---
direct or iterative sparse solvers exploit it heavily, and it is why
scientific computing and finite element analysis are inseparable
subjects at scale.

\section{Where this goes}\label{where-this-goes-40}

The next natural continuation is
\href{07-nonlinear-optimisation.qmd}{Nonlinear optimisation for design
and operations}. Once a simulation model is trustworthy enough to
represent the system, design questions become computational questions:
which shape, parameter, schedule, or control setting produces the best
feasible outcome?

This chapter also deepens what you should mean by ``model.'' In
upper-year work, the numerical method is not downstream implementation
detail. It is part of the model's claim to credibility.

\begin{tcolorbox}[enhanced jigsaw, opacityback=0, colback=white, breakable, colframe=quarto-callout-tip-color-frame, arc=.35mm, left=2mm, toprule=.15mm, leftrule=.75mm, bottomrule=.15mm, rightrule=.15mm]
\begin{minipage}[t]{5.5mm}
\textcolor{quarto-callout-tip-color}{\faLightbulb}
\end{minipage}%
\begin{minipage}[t]{\textwidth - 5.5mm}

\vspace{-3mm}\textbf{Applications}\vspace{3mm}

\begin{itemize}
\tightlist
\item
  heat-transfer simulation in machine components
\item
  finite-element structural analysis
\item
  finite-volume flow and transport models
\item
  weather, hydrology, and terrain-based environmental models
\item
  sparse linear systems in scientific computing
\item
  mesh-sensitivity and timestep-sensitivity studies
\end{itemize}

\end{minipage}%
\end{tcolorbox}

\section{Exercises}\label{exercises-51}

These are project-style exercises. Do not stop at the formula. Explain
what the numerical choice means for the model.

\subsection{Exercise 1}\label{exercise-1-4}

For the explicit diffusion scheme with \(D = 2\), \(\Delta x = 0.05\),
and \(\Delta t = 0.0004\), compute

\[r = \frac{D\Delta t}{(\Delta x)^2}\]

and decide whether the time step is plausibly stable for the standard
explicit scheme.

\subsection{Exercise 2}\label{exercise-2-4}

A finite-volume cell has inflow flux \(F_{i-1/2} = 3\), outflow flux
\(F_{i+1/2} = 5\), source term \(s_i = 1\), and width \(\Delta x = 2\).

Write the sign of the net flux contribution and explain whether storage
tends to increase or decrease before the source term is included.

\subsection{Exercise 3}\label{exercise-3-4}

Write a short comparison note explaining when you would prefer:

\begin{enumerate}
\def\labelenumi{\arabic{enumi}.}
\tightlist
\item
  finite difference
\item
  finite volume
\item
  finite element
\end{enumerate}

Use one example application for each method.

\subsection{Exercise 4}\label{exercise-4-4}

Choose one simulation problem from your own field and prepare a one-page
model brief naming:

\begin{enumerate}
\def\labelenumi{\arabic{enumi}.}
\tightlist
\item
  the governing equation
\item
  the boundary and initial conditions
\item
  the likely discretisation family
\item
  one likely stability or conditioning concern
\item
  one mesh- or timestep-sensitivity check you would run before trusting
  the result
\end{enumerate}

\chapter{Estimation, inverse problems, and
filtering}\label{estimation-inverse-problems-and-filtering}

Recovering hidden states from noisy measurements

\hfill\break

You can measure vibration but not crack length. You can observe GPS,
camera, and IMU signals but not the true state of the vehicle. You can
see reflected radiation from a landscape but not the quantity you really
care about on the ground. Inverse mathematics begins when the thing you
want is not the thing you can measure.

This chapter is about that gap between observation and quantity of
interest. Forward modelling asks: if the state were known, what
measurements would it produce? Inverse work asks the harder question:
given the measurements, what can you recover about the state or
parameters?

That harder question is unstable unless handled carefully. Small
measurement errors can create large errors in the recovered quantity. A
good estimation method therefore does two jobs at once: it produces an
estimate and it keeps track of how much trust the data deserves.

\begin{center}\rule{0.5\linewidth}{0.5pt}\end{center}

\section{Forward and inverse
viewpoints}\label{forward-and-inverse-viewpoints}

Let \(\mathbf{x}\) be the hidden state you care about and \(\mathbf{y}\)
the measured data you actually have. A simple measurement model is

\[\mathbf{y} = H\mathbf{x} + \varepsilon\]

where \(H\) maps the state into observation space and \(\varepsilon\)
represents measurement noise.

The \textbf{forward problem} is:

\begin{itemize}
\tightlist
\item
  given \(\mathbf{x}\), compute \(\mathbf{y}\)
\end{itemize}

The \textbf{inverse problem} is:

\begin{itemize}
\tightlist
\item
  given \(\mathbf{y}\), estimate \(\mathbf{x}\)
\end{itemize}

The inverse direction is harder because the map may be noisy,
incomplete, or many-to-one. Some parts of the state may be visible in
the data. Others may be weakly observable or not observable at all.

This is why estimation is not just algebraic rearrangement. It is
controlled inference under limited information.

\section{The simplest update logic}\label{the-simplest-update-logic}

Suppose you have a prior estimate \(\hat{x}^-\) and then receive a new
measurement \(y\). In the scalar case where the measurement directly
observes the state (\(H = 1\)), a simple update rule is

\[\hat{x}^+ = \hat{x}^- + K(y - \hat{x}^-)\]

where:

\begin{itemize}
\tightlist
\item
  \(\hat{x}^-\) is the estimate before the measurement update
\item
  \(y - \hat{x}^-\) is the innovation or residual
\item
  \(K\) is the \textbf{update weight} --- the same letter used for
  control gain in Chapter 1 and for the stiffness matrix in Chapter 4;
  here it plays a different role as a scalar between 0 and 1
\end{itemize}

In the general case with a non-trivial observation operator \(H\), the
innovation becomes \(y - H\hat{x}^-\), but the structure of the weighted
update is the same.

If \(K\) is close to 1, you trust the new measurement strongly. If \(K\)
is close to 0, you trust the existing estimate strongly. Filtering
methods differ in how they compute or adapt this weight, but the logic
is already visible here.

This update rule is useful because it makes the compromise explicit.
Estimation is never ``use the measurement'' or ``use the model.'' It is
always some weighted combination of the two.

\section{Parameter estimation and state
estimation}\label{parameter-estimation-and-state-estimation}

Not every inverse problem asks for the same thing.

\textbf{State estimation}

You want the current hidden state of a dynamic system: position,
velocity, temperature field, internal concentration, or machine
condition.

\textbf{Parameter estimation}

You want the unknown constant or slowly varying quantity in the model:
diffusivity, drag coefficient, reaction rate, calibration factor.

\textbf{Inverse field recovery}

You want an entire spatially distributed quantity from indirect
measurements: subsurface moisture, atmospheric state, or interior
temperature.

The mathematics overlaps, but the interpretation changes. This chapter
keeps the common structure in view so those problems do not feel
unrelated.

\begin{tcolorbox}[enhanced jigsaw, opacityback=0, colback=white, breakable, colframe=quarto-callout-note-color-frame, arc=.35mm, left=2mm, toprule=.15mm, leftrule=.75mm, bottomrule=.15mm, rightrule=.15mm]
\begin{minipage}[t]{5.5mm}
\textcolor{quarto-callout-note-color}{\faInfo}
\end{minipage}%
\begin{minipage}[t]{\textwidth - 5.5mm}

\vspace{-3mm}\textbf{Why This Works}\vspace{3mm}

An estimate improves when it combines two pieces of information that
fail in different ways: a model prediction and a measurement. The model
is coherent but incomplete. The measurement is immediate but noisy.
Filtering succeeds when the weights are chosen so each source corrects
the other's weakness.

That is the deeper reason estimation sits between control, signal
processing, and statistics. It uses all three.

\end{minipage}%
\end{tcolorbox}

\section{The core method}\label{the-core-method-4}

A first pass through an estimation or inverse problem usually goes like
this:

\begin{enumerate}
\def\labelenumi{\arabic{enumi}.}
\tightlist
\item
  Decide whether the unknown is a state, a parameter, or a field.
\item
  Write the forward relation from hidden quantity to measurement.
\item
  Name the main source of uncertainty: measurement noise, model error,
  missing data, or poor observability.
\item
  Choose an update or fitting rule.
\item
  Interpret the residuals, not just the estimate.
\item
  Ask what information the data still cannot recover.
\end{enumerate}

That last question matters. A polished estimate can still be misleading
if the measurements never contained the needed information in the first
place.

\section{Worked example 1: one-step scalar
filter}\label{worked-example-1-one-step-scalar-filter}

Suppose the prior estimate of a temperature is

\[\hat{x}^- = 18\]

degrees, a new sensor reading is

\[y = 20\]

degrees, and the update weight is

\[K = 0.3\]

Then

\[\hat{x}^+ = \hat{x}^- + K(y - \hat{x}^-)
= 18 + 0.3(20 - 18)
= 18 + 0.6
= 18.6\]

The updated estimate moves toward the measurement, but does not jump all
the way there. That is exactly what you want when the sensor is
informative but not perfect.

If the sensor were known to be extremely reliable, a larger \(K\) would
make sense. If the measurement were very noisy, a smaller \(K\) would be
more honest.

\section{Worked example 2: linear least-squares parameter
fit}\label{worked-example-2-linear-least-squares-parameter-fit}

Suppose a calibration model predicts measurements by

\[y_i \approx \theta x_i\]

for known inputs \(x_i\) and unknown scale parameter \(\theta\). With
observations

\[
(x_i, y_i) = (1, 2.1), (2, 3.9), (3, 6.2)
\]

the least-squares estimate minimises

\[\sum_i (y_i - \theta x_i)^2\]

Setting \(\dfrac{d}{d\theta}\sum_i (y_i - \theta x_i)^2 = 0\) gives
\(-2\sum_i x_i(y_i - \theta x_i) = 0\), which rearranges to

\[\hat{\theta}
=
\frac{\sum_i x_i y_i}{\sum_i x_i^2}
=
\frac{1(2.1) + 2(3.9) + 3(6.2)}{1^2 + 2^2 + 3^2}
=
\frac{28.5}{14}
\approx 2.036\]

This is parameter estimation rather than state estimation, but the
structure is the same: use noisy measurements to infer a hidden quantity
in the model.

\section{Worked example 3: remote sensing as an inverse
problem}\label{worked-example-3-remote-sensing-as-an-inverse-problem}

A satellite does not measure soil moisture directly. It measures
reflected or emitted radiation. A forward model links surface state to
what the instrument would see. The inverse problem works backward from
the measured signal toward the surface property of interest.

The mathematics is the same as in engineering system identification:

\begin{itemize}
\tightlist
\item
  there is a hidden state or parameter
\item
  there is an observation operator
\item
  there is noise and model mismatch
\item
  the estimate must balance data fit against plausibility
\end{itemize}

This is why optional viewpoints matter in Volume 8. Remote sensing and
process monitoring can look unrelated socially, but mathematically they
are close relatives.

\section{The Kalman filter as the canonical
example}\label{the-kalman-filter-as-the-canonical-example}

The update rule in Worked Example 1 is not just a pedagogical toy. It is
the core of the \textbf{Kalman filter}, which structures estimation
around a predict-then-update cycle:

\begin{enumerate}
\def\labelenumi{\arabic{enumi}.}
\tightlist
\item
  \textbf{Predict}: use the system model to propagate the state estimate
  forward one step, together with an estimate of how uncertain that
  prediction is.
\item
  \textbf{Update}: when a measurement arrives, compute how much to trust
  it relative to the prediction. The Kalman gain \(K\) sets that weight
  optimally when the noise is Gaussian.
\end{enumerate}

The two-step pattern generalises. Extended Kalman filters handle
nonlinear models by linearising around the current estimate. Particle
filters handle non-Gaussian noise by representing uncertainty as a cloud
of samples. Variational data assimilation in weather prediction is a
related method at continental scale. Hidden Markov models in ML use the
same predict-update logic for discrete state spaces, and Bayesian
filtering is the general probabilistic framing that includes all of
these.

The scalar example in this chapter already shows the logic. The rest is
bookkeeping to handle vectors, covariances, and system models.

\section{Where this goes}\label{where-this-goes-41}

The next natural continuation is
\href{06-reliability-stochastic.qmd}{Reliability, stochastic systems,
and quality}. Once you are estimating hidden states and parameters from
uncertain data, the next question is often what that uncertainty means
for failure, quality, risk, or operational decision-making over time.

This chapter also loops backward toward control. Modern control systems
often need estimated states rather than directly measured ones.
Estimation is not a sidebar to control. It is often the missing half.

\begin{tcolorbox}[enhanced jigsaw, opacityback=0, colback=white, breakable, colframe=quarto-callout-tip-color-frame, arc=.35mm, left=2mm, toprule=.15mm, leftrule=.75mm, bottomrule=.15mm, rightrule=.15mm]
\begin{minipage}[t]{5.5mm}
\textcolor{quarto-callout-tip-color}{\faLightbulb}
\end{minipage}%
\begin{minipage}[t]{\textwidth - 5.5mm}

\vspace{-3mm}\textbf{Applications}\vspace{3mm}

\begin{itemize}
\tightlist
\item
  condition monitoring in machines and structures
\item
  navigation and sensor fusion
\item
  process monitoring and soft sensing
\item
  parameter estimation in simulation models
\item
  remote sensing and environmental inversion
\item
  streaming latent-state estimation in data systems
\end{itemize}

\end{minipage}%
\end{tcolorbox}

\section{Exercises}\label{exercises-52}

These are project-style exercises. State what the estimate means and
what the data still does not tell you.

\subsection{Exercise 1}\label{exercise-1-5}

Use the scalar update rule

\[\hat{x}^+ = \hat{x}^- + K(y - \hat{x}^-)\]

with \(\hat{x}^- = 12\), \(y = 15\), and \(K = 0.4\).

Compute the updated estimate and explain whether the update trusts the
measurement strongly or cautiously.

\subsection{Exercise 2}\label{exercise-2-5}

Given data

\[
(x_i, y_i) = (1, 1.8), (2, 4.1), (4, 8.3),
\]

estimate \(\theta\) in the model \(y_i \approx \theta x_i\) using

\[\hat{\theta} = \frac{\sum_i x_i y_i}{\sum_i x_i^2}\]

and explain what the parameter means physically.

\subsection{Exercise 3}\label{exercise-3-5}

A monitoring system can measure motor current and casing temperature,
but not bearing wear directly.

Write a short systems note answering:

\begin{enumerate}
\def\labelenumi{\arabic{enumi}.}
\tightlist
\item
  what the hidden state is
\item
  what the measurements are
\item
  one plausible forward model relation
\item
  one reason the inverse problem may be ill-conditioned or ambiguous
\end{enumerate}

\subsection{Exercise 4}\label{exercise-4-5}

Choose one inverse problem from your own field and prepare a one-page
model brief naming:

\begin{enumerate}
\def\labelenumi{\arabic{enumi}.}
\tightlist
\item
  the hidden quantity
\item
  the measured quantity
\item
  the forward map
\item
  the main noise source
\item
  one update or fitting rule you would try first
\item
  one thing the estimate still cannot guarantee
\end{enumerate}

\chapter{Reliability, stochastic systems, and
quality}\label{reliability-stochastic-systems-and-quality}

Failure, variability, and risk over time

\hfill\break

One bearing fails early, another lasts for years. A queue stays calm all
morning, then suddenly grows faster than it clears. A production line
meets tolerance most days but not all days. Randomness is not a side
issue here. It changes maintenance, staffing, safety, and design.

Reliability and stochastic systems matter because engineering decisions
are made before the exact future is known. You do not know which unit
will fail next, exactly when a line will back up, or precisely how many
defects will appear in the next batch. What you can do is model the
variability honestly and let that model shape action.

This chapter keeps three linked ideas in view: lifetime, queueing, and
quality. They look like separate professional topics, but they are all
ways of asking how randomness accumulates into operational consequence.

\begin{center}\rule{0.5\linewidth}{0.5pt}\end{center}

\section{Reliability as survival over
time}\label{reliability-as-survival-over-time}

Let \(T\) be a random lifetime. The \textbf{reliability function} is

\[R(t) = P(T > t)\]

This is the probability that the component or system survives beyond
time \(t\).

The \textbf{failure distribution} is

\[F(t) = P(T \leq t) = 1 - R(t)\]

If the lifetime is exponentially distributed with constant hazard rate
\(\lambda\), then

\[R(t) = e^{-\lambda t}\]

and

\[F(t) = 1 - e^{-\lambda t}\]

This is one of the most common first models because it is simple and
because its interpretation is clean: the system has a constant failure
tendency per unit time.

The \textbf{hazard rate} \(\lambda(t)\) describes the instantaneous
failure tendency conditional on survival up to time \(t\). In some
systems hazard increases with age. In others it is approximately
constant over the operational window of interest.

\section{Queueing as accumulation under
uncertainty}\label{queueing-as-accumulation-under-uncertainty}

Now switch interpretation. Instead of asking when a unit fails, ask what
happens when random arrivals meet random service.

In a queueing system:

\begin{itemize}
\tightlist
\item
  arrivals are uncertain in time
\item
  service completions are uncertain in time
\item
  if arrivals temporarily outpace service, a queue forms
\end{itemize}

Two core quantities are:

\[L_q = \text{expected number waiting in queue}\]

\[W_q = \text{expected waiting time in queue}\]

These are connected by \textbf{Little's Law}:

\[L_q = \lambda_a W_q\]

The expected queue length equals the arrival rate times the expected
waiting time. Little's Law holds under mild conditions regardless of the
arrival or service distribution --- it is a conservation statement about
flow rather than an assumption about randomness. For a single-server
queue with traffic intensity \(\rho < 1\), the queue length grows
sharply as \(\rho\) approaches 1, even when mean demand still looks
safely below capacity.

The engineering lesson is the same as in reliability: variability
matters even when average values look harmless. A system with average
arrival rate only slightly below service rate can still produce long
queues if the variability is large enough.

This is why service reliability in computing and waiting-time models in
operations research belong in the same chapter as bearing lifetime and
defect rates. They are all stochastic consequences of uncertain events
over time.

\section{Quality as probability of staying in
tolerance}\label{quality-as-probability-of-staying-in-tolerance}

A quality problem can often be phrased as:

\begin{itemize}
\tightlist
\item
  what is the probability that a product characteristic falls outside
  the acceptable region?
\end{itemize}

If a machined diameter is modelled as approximately normal with mean
\(\mu\) and standard deviation \(\sigma\), then defect probability is
the area in the tails beyond the specification limits.

This is not only a manufacturing question. The same logic appears in
service level agreements, anomaly thresholds, and risk tolerances in
data systems.

\begin{tcolorbox}[enhanced jigsaw, opacityback=0, colback=white, breakable, colframe=quarto-callout-note-color-frame, arc=.35mm, left=2mm, toprule=.15mm, leftrule=.75mm, bottomrule=.15mm, rightrule=.15mm]
\begin{minipage}[t]{5.5mm}
\textcolor{quarto-callout-note-color}{\faInfo}
\end{minipage}%
\begin{minipage}[t]{\textwidth - 5.5mm}

\vspace{-3mm}\textbf{Why This Works}\vspace{3mm}

Reliability, queueing, and quality all convert randomness into
operational questions. Survival functions convert uncertainty into
maintenance or safety policy. Queueing models convert uncertainty into
staffing and capacity policy. Quality models convert uncertainty into
tolerance and acceptance policy.

The mathematics is useful precisely because it does not pretend the
future is deterministic.

\end{minipage}%
\end{tcolorbox}

\section{The core method}\label{the-core-method-5}

A first pass through a stochastic reliability or quality problem usually
goes like this:

\begin{enumerate}
\def\labelenumi{\arabic{enumi}.}
\tightlist
\item
  Identify the random quantity: lifetime, arrivals, service time, or
  process variation.
\item
  Choose the summary function that matches the decision: reliability,
  hazard, queue length, waiting time, or defect probability.
\item
  Compute the key probability or expected value.
\item
  Interpret it in operational terms.
\item
  Ask whether the assumed stochastic model is realistic enough for the
  decision you are about to make.
\end{enumerate}

That final question is essential. A neat model with the wrong
variability structure can be more dangerous than admitting uncertainty
openly.

\section{Worked example 1: exponential
reliability}\label{worked-example-1-exponential-reliability}

Suppose a component has constant failure rate

\[\lambda = 0.002 \text{ h}^{-1}\]

Then the reliability at 100 hours is

\[R(100) = e^{-0.002(100)} = e^{-0.2} \approx 0.819\]

So there is about an 81.9\% chance the component survives beyond 100
hours. Equivalently, the probability of failure by 100 hours is

\[F(100) = 1 - 0.819 = 0.181\]

This is already enough to frame a maintenance question. If a system
contains many such components, or if failure consequence is high, 18.1\%
may be too risky to tolerate over that interval.

\section{Worked example 2: a queue near
saturation}\label{worked-example-2-a-queue-near-saturation}

A service system receives jobs at average rate

\[\lambda_a = 8 \text{ jobs/h}\]

and can serve at average rate

\[\mu = 10 \text{ jobs/h}\]

The traffic intensity is

\[\rho = \frac{\lambda_a}{\mu} = 0.8\]

This means the server is busy most of the time. Even though capacity
still exceeds mean demand, there is not much slack. A burst of arrivals
can create a substantial queue.

The lesson is operational, not just probabilistic: systems that look
adequate on average can still feel overloaded to users because
variability is lived in real time, not as a long-run mean.

\section{Worked example 3: defect probability from a normal
model}\label{worked-example-3-defect-probability-from-a-normal-model}

Suppose a manufactured diameter is approximately normal with mean

\[\mu = 10.0 \text{ mm}\]

and standard deviation

\[\sigma = 0.1 \text{ mm}\]

The upper specification limit is \(10.2\) mm. Standardise:

\[z = \frac{10.2 - 10.0}{0.1} = 2\]

So the probability of exceeding the upper limit is about

\[P(Z > 2) \approx 0.0228\]

or 2.28\%.

That tail probability is the mathematical link between observed
variation and quality cost. It is also the same logic used in anomaly
detection and service thresholds outside manufacturing.

\section{Where this goes}\label{where-this-goes-42}

The next continuation is \href{07-nonlinear-optimisation.qmd}{Nonlinear
optimisation for design and operations}. Once risk, waiting, and failure
are quantified, design decisions become tradeoffs between cost,
performance, safety, and uncertainty.

This chapter also reframes what ``uncertainty'' means in Volume 8. It is
not only about estimation error. It is about lifetime, demand,
variability, and the operational consequences of randomness over time.

\begin{tcolorbox}[enhanced jigsaw, opacityback=0, colback=white, breakable, colframe=quarto-callout-tip-color-frame, arc=.35mm, left=2mm, toprule=.15mm, leftrule=.75mm, bottomrule=.15mm, rightrule=.15mm]
\begin{minipage}[t]{5.5mm}
\textcolor{quarto-callout-tip-color}{\faLightbulb}
\end{minipage}%
\begin{minipage}[t]{\textwidth - 5.5mm}

\vspace{-3mm}\textbf{Applications}\vspace{3mm}

\begin{itemize}
\tightlist
\item
  maintenance interval planning
\item
  reliability engineering and survival analysis
\item
  queueing in production, logistics, and computing
\item
  defect rates and process capability
\item
  service availability and incident risk
\item
  uncertainty-aware operations decisions
\end{itemize}

\end{minipage}%
\end{tcolorbox}

\section{Exercises}\label{exercises-53}

These are project-style exercises. State the operational meaning of the
number you compute.

\subsection{Exercise 1}\label{exercise-1-6}

A component has exponential lifetime with failure rate

\[\lambda = 0.005 \text{ h}^{-1}\]

Compute the probability it survives beyond 50 hours.

\subsection{Exercise 2}\label{exercise-2-6}

A queueing system has average arrival rate 12 jobs/h and service rate 15
jobs/h.

\begin{enumerate}
\def\labelenumi{\arabic{enumi}.}
\tightlist
\item
  Compute the traffic intensity \(\rho\).
\item
  Explain whether the system has generous slack or operates near
  saturation.
\item
  Suppose demand grows to 14 jobs/h without any change in service rate.
  What happens to \(\rho\), and at what arrival rate does the system
  become dangerously loaded (say, \(\rho > 0.95\))? What operational
  decision does this suggest?
\end{enumerate}

\subsection{Exercise 3}\label{exercise-3-6}

A process characteristic is approximately normal with mean 25 and
standard deviation 2. The upper specification limit is 29.

Compute the corresponding \(z\) value and explain what tail probability
it refers to.

\subsection{Exercise 4}\label{exercise-4-6}

Choose one stochastic setting from your field and prepare a one-page
model brief naming:

\begin{enumerate}
\def\labelenumi{\arabic{enumi}.}
\tightlist
\item
  the random quantity
\item
  the decision you need to support
\item
  the probability or expectation that matters
\item
  one source of model mismatch
\item
  one action you would change if the risk turned out to be higher than
  expected
\end{enumerate}

\chapter{Nonlinear optimisation for design and
operations}\label{nonlinear-optimisation-for-design-and-operations}

Tradeoffs, constraints, and the shape of a best solution

\hfill\break

Make the structure lighter without making it fail. Increase throughput
without violating a pressure limit. Improve prediction quality without
blowing the compute budget. Once the world stops being linear, ``the
best answer'' is no longer a corner point on a clean polygon. It is a
shape in a curved space, and getting there becomes part geometry, part
computation, part judgment.

This chapter closes Volume 8 because it gathers together the themes that
the earlier chapters developed separately. Simulation creates a model of
system behaviour. Estimation tells you what is known and uncertain.
Reliability tells you what risk costs you. Optimisation turns all of
that into a decision rule.

The mathematics changes here because the objective and constraints are
no longer linear in the design variables. That means the geometry
changes, the algorithms change, and the meaning of ``best'' becomes more
local and more conditional.

\begin{center}\rule{0.5\linewidth}{0.5pt}\end{center}

\section{From linear to nonlinear
thinking}\label{from-linear-to-nonlinear-thinking}

In linear programming, the feasible set is a polytope and the optimum
sits at a vertex. In nonlinear optimisation, neither statement is
generally true.

You now have:

\begin{itemize}
\tightlist
\item
  a nonlinear objective \(f(\mathbf{x})\)
\item
  nonlinear constraints \(g_i(\mathbf{x}) \leq 0\)
\item
  possibly many local optima instead of one clean global answer
\end{itemize}

The design variable vector \(\mathbf{x}\) may represent dimensions,
control settings, schedule variables, material choices, or
hyperparameters. The mathematics does not care what professional name
the variables carry. It cares how objective and constraints bend the
space.

The simplest unconstrained update rule is gradient descent:

\[\mathbf{x}_{k+1} = \mathbf{x}_k - \alpha \nabla f(\mathbf{x}_k)\]

The negative gradient points in the direction of steepest local
decrease. Choose a step size \(\alpha\), move, and repeat.

This sounds easy. In practice the difficulties are exactly the ones you
would expect in a curved space: the descent direction is only locally
good, the step size can be too large or too small, and constraints can
make the locally best direction infeasible.

\begin{tcolorbox}[enhanced jigsaw, opacityback=0, colback=white, breakable, colframe=quarto-callout-note-color-frame, arc=.35mm, left=2mm, toprule=.15mm, leftrule=.75mm, bottomrule=.15mm, rightrule=.15mm]
\begin{minipage}[t]{5.5mm}
\textcolor{quarto-callout-note-color}{\faInfo}
\end{minipage}%
\begin{minipage}[t]{\textwidth - 5.5mm}

\vspace{-3mm}\textbf{Computing note: batch versus stochastic gradient descent}\vspace{3mm}

The rule above evaluates the full gradient over all data at each step
--- this is \textbf{batch gradient descent}. In machine learning, where
\(f\) is a loss summed over millions of examples, evaluating the full
gradient per step is computationally prohibitive. \textbf{Stochastic
gradient descent} (SGD) substitutes an estimate of the gradient computed
on a randomly chosen mini-batch of examples. The update direction is
noisy but cheap. Modern ML training uses variants of mini-batch SGD ---
Adam, AdaGrad, RMSProp --- that also adapt the step size per parameter.

The mathematics of convergence guarantees and step-size schedules for
SGD is different from the deterministic case, but the core geometry is
the same: move opposite to an estimate of the gradient, iterate, and
check for convergence.

\end{minipage}%
\end{tcolorbox}

\section{Constrained optimisation and the Lagrangian
idea}\label{constrained-optimisation-and-the-lagrangian-idea}

Suppose you want to minimise \(f(\mathbf{x})\) subject to an equality
constraint

\[h(\mathbf{x}) = 0\]

At an optimum, you usually cannot move in an arbitrary direction because
many directions leave the feasible surface. The Lagrangian is

\[\mathcal{L}(\mathbf{x}, \lambda) = f(\mathbf{x}) + \lambda h(\mathbf{x})\]

The stationary conditions are

\[\nabla_{\mathbf{x}} \mathcal{L} = 0, \qquad h(\mathbf{x}) = 0\]

Geometrically, this means the gradient of the objective lines up with
the gradient of the constraint. You cannot improve the objective without
leaving feasibility.

For inequality constraints the full Karush-Kuhn-Tucker (KKT) conditions
are needed in general --- these extend the Lagrangian stationarity
conditions to handle inequalities by requiring that each active
constraint contributes a non-negative multiplier and that inactive
constraints exert no force on the optimum. The intuition is similar to
the equality case: active constraints shape which directions are allowed
at the optimum.

\section{Tradeoffs are the real
subject}\label{tradeoffs-are-the-real-subject}

Nonlinear optimisation is often introduced as an algorithm chapter. In
upper year engineering, it is better understood as a tradeoff chapter.

If you minimise mass, stress may rise. If you maximise throughput,
queueing risk may rise. If you improve prediction accuracy, compute cost
may rise. The objective function is therefore not neutral. It encodes
what you are willing to care about and what you are willing to pay.

That is why sensitivity analysis matters. A computed optimum is only
useful if you can explain how it shifts when assumptions, weights, or
constraints change.

\begin{tcolorbox}[enhanced jigsaw, opacityback=0, colback=white, breakable, colframe=quarto-callout-note-color-frame, arc=.35mm, left=2mm, toprule=.15mm, leftrule=.75mm, bottomrule=.15mm, rightrule=.15mm]
\begin{minipage}[t]{5.5mm}
\textcolor{quarto-callout-note-color}{\faInfo}
\end{minipage}%
\begin{minipage}[t]{\textwidth - 5.5mm}

\vspace{-3mm}\textbf{Why This Works}\vspace{3mm}

Optimisation turns design or operational judgment into geometry. The
objective describes which direction counts as improvement. The
constraints describe where you are allowed to move. A solution is
optimal only relative to both.

In nonlinear settings, there may be many locally good choices. The
mathematics helps you navigate them, but it does not remove the need to
interpret what the objective and constraints mean physically.

\end{minipage}%
\end{tcolorbox}

\section{The core method}\label{the-core-method-6}

A first pass through a nonlinear optimisation problem usually goes like
this:

\begin{enumerate}
\def\labelenumi{\arabic{enumi}.}
\tightlist
\item
  Choose the design variables.
\item
  Write the objective in a way that matches the actual decision.
\item
  Write the constraints in physically meaningful form.
\item
  Inspect the geometry qualitatively before trusting an algorithm.
\item
  Use a numerical method appropriate to smoothness and constraint
  structure.
\item
  Check whether the result is feasible, local or global, and sensitive
  to modelling assumptions.
\end{enumerate}

This is where all the earlier Volume 8 chapters return. If the model,
estimate, or risk terms are poor, the optimisation result will be
polished nonsense.

\section{Worked example 1: unconstrained quadratic design
objective}\label{worked-example-1-unconstrained-quadratic-design-objective}

Minimise

\[f(x) = x^2 - 4x + 7\]

The derivative is

\[f'(x) = 2x - 4\]

Set it to zero:

\[2x - 4 = 0 \qquad \Rightarrow \qquad x = 2\]

The second derivative is

\[f''(x) = 2 > 0\]

so \(x=2\) is a local minimum. The minimum value is

\[f(2) = 4 - 8 + 7 = 3\]

This is a simple example, but it already shows the core logic:
stationary point plus curvature check.

\section{Worked example 2: equality-constrained
optimisation}\label{worked-example-2-equality-constrained-optimisation}

Minimise

\[f(x,y) = x^2 + y^2\]

subject to

\[x + y - 1 = 0\]

Form the Lagrangian:

\[\mathcal{L}(x,y,\lambda) = x^2 + y^2 + \lambda(x+y-1)\]

Take derivatives:

\[\frac{\partial \mathcal{L}}{\partial x} = 2x + \lambda = 0\]
\[\frac{\partial \mathcal{L}}{\partial y} = 2y + \lambda = 0\]
\[x+y-1 = 0\]

The first two equations imply

\[x = y\]

Then the constraint gives

\[2x = 1 \qquad \Rightarrow \qquad x = y = \frac{1}{2}\]

So the feasible point closest to the origin is

\[\left(\frac{1}{2}, \frac{1}{2}\right)\]

This is the simplest geometric picture of constrained optimisation: the
objective's level sets touch the feasible set at the best allowable
point.

\section{Worked example 3: design under a penalty
tradeoff}\label{worked-example-3-design-under-a-penalty-tradeoff}

Suppose an engineer uses the toy objective

\[J(x) = (x-3)^2 + 0.5x^4\]

Here one term rewards staying near a target value \(x=3\), while the
quartic term penalises large magnitude. This is not exotic. Real design
objectives often mix performance and regularisation in exactly this way.

Different penalty weights produce different optima. That means the
``best'' design is partly a statement about priorities, not just about
calculus.

The same logic appears in many settings. In machine learning,
\textbf{ridge regression} (L2 regularisation) adds a penalty
\(\lambda\|\mathbf{w}\|^2\) to the loss function, trading fit quality
for smaller parameter magnitudes. The penalty weight \(\lambda\) plays
the same role as the 0.5 here: it encodes how strongly the solver is
pushed away from large solutions, and changing it shifts the optimum.
\textbf{Lasso} (L1 regularisation) uses \(\lambda\|\mathbf{w}\|_1\) and
produces sparse solutions for geometric reasons that the Lagrangian
framework makes precise. Portfolio risk penalties and operations
scheduling with nonlinear costs follow the same two-term structure.

\section{Where this goes}\label{where-this-goes-43}

This is the terminal chapter of the current volume, so
\texttt{opens\_to} is empty on purpose. The continuation is not a single
next topic. It is the practice of using the whole Volume 8 stack
together: models, estimates, risks, and objectives inside one decision
loop.

By this point, mathematics in the engineering curriculum no longer
appears as a supporting subject. It is the working medium of design
judgment.

\begin{tcolorbox}[enhanced jigsaw, opacityback=0, colback=white, breakable, colframe=quarto-callout-tip-color-frame, arc=.35mm, left=2mm, toprule=.15mm, leftrule=.75mm, bottomrule=.15mm, rightrule=.15mm]
\begin{minipage}[t]{5.5mm}
\textcolor{quarto-callout-tip-color}{\faLightbulb}
\end{minipage}%
\begin{minipage}[t]{\textwidth - 5.5mm}

\vspace{-3mm}\textbf{Applications}\vspace{3mm}

\begin{itemize}
\tightlist
\item
  shape and parameter optimisation in engineering design
\item
  calibration and tuning of controllers
\item
  operations and scheduling under nonlinear costs
\item
  portfolio and risk optimisation
\item
  hyperparameter and training-objective tuning in ML systems
\item
  multi-objective design under uncertainty
\end{itemize}

\end{minipage}%
\end{tcolorbox}

\section{Exercises}\label{exercises-54}

These are project-style exercises. Always explain what ``better'' means
in the problem, not only where the derivative vanishes.

\subsection{Exercise 1}\label{exercise-1-7}

Minimise

\[f(x) = x^2 - 6x + 10\]

Find the minimiser and the minimum value.

\subsection{Exercise 2}\label{exercise-2-7}

Use Lagrange multipliers to minimise

\[f(x,y) = x^2 + 4y^2\]

subject to

\[x + y = 3\]

\subsection{Exercise 3}\label{exercise-3-7}

An engineering team is choosing a design variable \(x\) to minimise

\[J(x) = (x-2)^2 + 0.2x^4\]

Write a short interpretation note answering:

\begin{enumerate}
\def\labelenumi{\arabic{enumi}.}
\tightlist
\item
  what each term in the objective is doing
\item
  why this is not a linear optimisation problem
\item
  why changing the penalty weight 0.2 would change the preferred design
\end{enumerate}

\subsection{Exercise 4}\label{exercise-4-7}

Choose one optimisation problem from your own field and prepare a
one-page design brief naming:

\begin{enumerate}
\def\labelenumi{\arabic{enumi}.}
\tightlist
\item
  the design variables
\item
  the objective
\item
  the constraints
\item
  one source of nonlinearity
\item
  one uncertainty or reliability issue that should influence the
  formulation
\item
  one reason a computed optimum might still be a bad decision
\end{enumerate}


\backmatter


\end{document}
